%% NOETHER EDITABLE LATEX COLLECTION 1
%% {"components": ["ja/reader.tex"]}
%% BEGIN NOETHER SOURCE "ja/reader.tex"
\documentclass[11pt]{article}
\usepackage[a4paper,margin=1.45cm]{geometry}
\usepackage{fontspec}
\usepackage{xeCJK}
\setmainfont{Noto Serif}
\setCJKmainfont[AutoFakeSlant=0.15]{Yu Gothic}
\setCJKsansfont{Yu Gothic}
\usepackage{amsmath,amssymb,mathtools,array,booktabs,pdflscape,adjustbox,diagbox,graphicx}
\setlength{\parindent}{1em}
\setlength{\parskip}{0pt}
\newcommand{\widebar}[1]{\overline{#1}}
\newcommand{\A}{\Delta}
\newcommand{\D}{\Delta}
\allowdisplaybreaks
\setlength{\emergencystretch}{12em}
\sloppy
\hbadness=10000
\vbadness=10000
\hfuzz=2pt
\vfuzz=10pt
\raggedbottom
\newcommand{\R}{\varrho}
\newcommand{\hata}{\widehat{a}}
\newcommand{\hatv}{\widehat{v}}
\newcommand{\modu}[1]{\;\operatorname{mod} #1}
\newcommand{\mcell}[1]{\ensuremath{\begin{gathered}#1\end{gathered}}}
\newcommand{\pair}[2]{(#1\,|\,#2)}
\newcommand{\eps}{\varepsilon}
\providecommand{\Hgrp}{\mathfrak{H}}
\providecommand{\GaloisRes}{\Phi}
\providecommand{\pol}[2]{\left(#1\frac{\partial}{\partial #2}\right)}
\providecommand{\Deltaop}{\Delta\!\left(\frac{\partial}{\partial x},\frac{\partial}{\partial y},\frac{\partial}{\partial z}\right)}
\providecommand{\Omegaop}{\Omega}
\makeatletter
\renewcommand\footnotesize{\@setfontsize\footnotesize{7.5}{8.7}\selectfont}
\makeatother
\newcommand{\bdelta}{\bar{\delta}}
\providecommand{\dd}{\mathrm d}
\providecommand{\Div}{\operatorname{Div}}
\providecommand{\G}{\mathfrak G}
\providecommand{\bd}{\bar{\delta}}
\providecommand{\p}{\partial}
\providecommand{\ideal}[1]{\mathfrak{#1}}
\providecommand{\eqzero}[1]{\equiv 0\,(#1)}
\providecommand{\gtr}{>}
% 組版用マクロ
\providecommand{\widebar}[1]{\overline{#1}}
\providecommand{\A}{\Delta}
\providecommand{\D}{\Delta}
\providecommand{\R}{\varrho}
\providecommand{\hata}{\widehat{a}}
\providecommand{\hatv}{\widehat{v}}
\providecommand{\modu}[1]{\;\operatorname{mod} #1}
\providecommand{\mcell}[1]{\ensuremath{\begin{gathered}#1\end{gathered}}}
\providecommand{\pair}[2]{(#1\,|\,#2)}
\providecommand{\eps}{\varepsilon}
\providecommand{\rhoidx}{\varrho}
\providecommand{\dd}{\mathrm d}
\providecommand{\Div}{\operatorname{Div}}
\providecommand{\G}{\mathfrak G}
\providecommand{\bd}{\bar{\delta}}
\providecommand{\p}{\partial}
\providecommand{\frS}{\mathfrak S}
\providecommand{\frp}{\mathfrak p}
\providecommand{\fra}{\mathfrak a}
\providecommand{\frb}{\mathfrak b}
\providecommand{\frc}{\mathfrak c}
\providecommand{\frf}{\mathfrak f}
\providecommand{\frr}{\mathfrak r}
\providecommand{\frm}{\mathfrak m}
\providecommand{\frD}{\mathfrak D}
\providecommand{\calA}{\mathcal A}
\providecommand{\calB}{\mathcal B}
\providecommand{\calN}{\mathcal N}
\providecommand{\mM}{\mathfrak{M}}
\providecommand{\mN}{\mathfrak{N}}
\providecommand{\mL}{\mathfrak{L}}
\providecommand{\mG}{\mathfrak{G}}
\providecommand{\mA}{\mathfrak{A}}
\providecommand{\mB}{\mathfrak{B}}
\providecommand{\mX}{\mathfrak{X}}
\providecommand{\mY}{\mathfrak{Y}}
\providecommand{\ma}{\mathfrak{a}}
\providecommand{\mb}{\mathfrak{b}}
\providecommand{\mx}{\mathfrak{x}}
\providecommand{\my}{\mathfrak{y}}
\providecommand{\mz}{\mathfrak{z}}
\providecommand{\me}{\mathfrak{e}}
\providecommand{\mbarA}{\overline{\mathfrak{A}}}
\providecommand{\mbarB}{\overline{\mathfrak{B}}}
\providecommand{\bara}{\overline{\mathfrak{a}}}
\providecommand{\mU}{\mathfrak{U}}
\providecommand{\mP}{\mathfrak{P}}
\providecommand{\mQ}{\mathfrak{Q}}
\providecommand{\mS}{\mathfrak{S}}
\providecommand{\mD}{\mathfrak{D}}
\providecommand{\mR}{\mathfrak{R}}
\providecommand{\mV}{\mathfrak{V}}
\providecommand{\ideal}[1]{\mathfrak{#1}}
\providecommand{\Amod}{\mathfrak A}
\providecommand{\Bmod}{\mathfrak B}
\providecommand{\Gmod}{\mathfrak G}
\providecommand{\Rmod}{\mathfrak R}
\providecommand{\Smod}{\mathfrak S}
\providecommand{\Mmod}{\mathfrak M}
\providecommand{\Nmod}{\mathfrak N}
\providecommand{\mideal}{\mathfrak m}
\providecommand{\nideal}{\mathfrak n}
\providecommand{\gideal}{\mathfrak g}
\providecommand{\hideal}{\mathfrak h}
\providecommand{\jideal}{\mathfrak j}
\providecommand{\tideal}{\mathfrak t}
\providecommand{\aideal}{\mathfrak a}
\providecommand{\bideal}{\mathfrak b}
\providecommand{\bb}{\mathfrak b}
\providecommand{\cc}{\mathfrak c}
\providecommand{\zero}[1]{\equiv 0\pmod{#1}}
\providecommand{\Norm}{N}
\providecommand{\frakS}{\mathfrak S}
\providecommand{\frakM}{\mathfrak M}
\providecommand{\frakG}{\mathfrak G}
\providecommand{\frakm}{\mathfrak m}
\providecommand{\frakn}{\mathfrak n}
\providecommand{\frakp}{\mathfrak p}
\providecommand{\frakq}{\mathfrak q}
\providecommand{\frakr}{\mathfrak r}
\providecommand{\fraka}{\mathfrak a}
\providecommand{\frakb}{\mathfrak b}
\providecommand{\frakc}{\mathfrak c}
\providecommand{\frakd}{\mathfrak d}
\providecommand{\frake}{\mathfrak e}
\providecommand{\frakg}{\mathfrak g}
\providecommand{\fraks}{\mathfrak s}
\providecommand{\frako}{\mathfrak o}
\providecommand{\frK}{\mathfrak K}
\providecommand{\barfrakm}{\overline{\mathfrak m}}
\providecommand{\frakh}{\mathfrak h}
\providecommand{\frakt}{\mathfrak t}
\providecommand{\barfraka}{\overline{\mathfrak a}}
\providecommand{\barfrakb}{\overline{\mathfrak b}}
\providecommand{\barfrakc}{\overline{\mathfrak c}}
\providecommand{\calP}{\mathcal P}
\providecommand{\calR}{\mathcal R}
\providecommand{\calG}{\mathcal G}
\providecommand{\calH}{\mathcal H}
\providecommand{\barr}{\overline}
\providecommand{\tagg}[1]{\tag{#1}}



% 組版用マクロ
\providecommand{\frakR}{\mathfrak R}
\providecommand{\frakT}{\mathfrak T}
\providecommand{\frakK}{\mathfrak K}
\providecommand{\frakL}{\mathfrak L}
\providecommand{\frakH}{\mathfrak H}
\providecommand{\frakP}{\mathfrak P}
\providecommand{\frakQ}{\mathfrak Q}
\providecommand{\frakZ}{\mathfrak Z}
\providecommand{\frakN}{\mathfrak N}
\providecommand{\frakB}{\mathfrak B}
\providecommand{\frakF}{\mathfrak F}
\providecommand{\frakW}{\mathfrak W}
\providecommand{\Trdeg}{\operatorname{trdeg}}
\providecommand{\Gal}{\operatorname{Gal}}
\providecommand{\Norm}{\operatorname{N}}


% 組版用マクロ
\providecommand{\widebar}[1]{\overline{#1}}
\providecommand{\mR}{\mathfrak{R}}
\providecommand{\mK}{\mathfrak{K}}
\providecommand{\mL}{\mathfrak{L}}
\providecommand{\mS}{\mathfrak{S}}
\providecommand{\mT}{\mathfrak{T}}
\providecommand{\mM}{\mathfrak{M}}
\providecommand{\mN}{\mathfrak{N}}
\providecommand{\mA}{\mathfrak{A}}
\providecommand{\mB}{\mathfrak{B}}
\providecommand{\mC}{\mathfrak{C}}
\providecommand{\mD}{\mathfrak{D}}
\providecommand{\mE}{\mathfrak{E}}
\providecommand{\mF}{\mathfrak{F}}
\providecommand{\mG}{\mathfrak{G}}
\providecommand{\mH}{\mathfrak{H}}
\providecommand{\mU}{\mathfrak{U}}
\providecommand{\mV}{\mathfrak{V}}
\providecommand{\mbar}[1]{\overline{#1}}
\providecommand{\ma}{\mathfrak{a}}
\providecommand{\mb}{\mathfrak{b}}
\providecommand{\mc}{\mathfrak{c}}
\providecommand{\md}{\mathfrak{d}}
\providecommand{\me}{\mathfrak{e}}
\providecommand{\mf}{\mathfrak{f}}
\providecommand{\mm}{\mathfrak{m}}
\providecommand{\mn}{\mathfrak{n}}
\providecommand{\mo}{\mathfrak{o}}
\providecommand{\mpideal}{\mathfrak{p}}
\providecommand{\mq}{\mathfrak{q}}
\providecommand{\mt}{\mathfrak{t}}
\providecommand{\mv}{\mathfrak{v}}
\providecommand{\frakp}{\mathfrak p}
\providecommand{\frakq}{\mathfrak q}
\providecommand{\modu}[1]{\;\operatorname{mod} #1}

% 組版用マクロ
\providecommand{\bR}{\overline{\mathfrak{R}}}
\providecommand{\bP}{\overline{P}}
\providecommand{\bq}{\overline{\mathfrak{q}}}
\providecommand{\bp}{\overline{\mathfrak{p}}}
\providecommand{\ba}{\overline{\mathfrak{a}}}
\providecommand{\mOmega}{\mathfrak{\Omega}}
\providecommand{\tr}{\operatorname{Tr}}
\providecommand{\Tr}{\operatorname{Tr}}
\providecommand{\class}[1]{(#1)}

\providecommand{\mZ}{\mathfrak{Z}}
\providecommand{\mO}{\mathfrak{O}}
\providecommand{\mC}{\mathfrak{C}}
\providecommand{\mE}{\mathfrak{E}}
\providecommand{\mf}{\mathfrak{f}}
\providecommand{\mOmega}{\mathfrak{\Omega}}
\providecommand{\mt}{\mathfrak{t}}

\providecommand{\mat}[1]{\begin{pmatrix}#1\end{pmatrix}}

% 組版用マクロ
\providecommand{\frakO}{\mathfrak O}
\providecommand{\frako}{\mathfrak o}
\providecommand{\Om}{\Omega}
\providecommand{\N}{\operatorname{N}}


% 組版用マクロ
\providecommand{\frR}{\mathfrak R}
\providecommand{\frS}{\mathfrak S}
\providecommand{\frT}{\mathfrak T}
\providecommand{\frM}{\mathfrak M}
\providecommand{\frN}{\mathfrak N}
\providecommand{\frA}{\mathfrak A}
\providecommand{\frB}{\mathfrak B}
\providecommand{\frX}{\mathfrak X}
\providecommand{\frY}{\mathfrak Y}
\providecommand{\frZ}{\mathfrak Z}
\providecommand{\frG}{\mathfrak G}
\providecommand{\frH}{\mathfrak H}
\providecommand{\frL}{\mathfrak L}
\providecommand{\frU}{\mathfrak U}
\providecommand{\frD}{\mathfrak D}
\providecommand{\frC}{\mathfrak C}
\providecommand{\frP}{\mathfrak P}
\providecommand{\Gcal}{\mathcal G}
\providecommand{\Sp}{\operatorname{Sp}}

\providecommand{\Gg}{\mathfrak G}
% 組版用マクロ
\providecommand{\mJ}{\mathfrak J}
\providecommand{\srcnumdisplay}[2]{\begin{equation*}\tag*{#1}\begin{aligned}#2\end{aligned}\end{equation*}}
\providecommand{\PXXIIIitem}[2]{\par\noindent\hangindent=2.2em\hangafter=1\makebox[2.2em][l]{#1.}#2\par}
\makeatletter
\providecommand{\srcfn}[2]{%
  \begingroup
  \renewcommand{\thefootnote}{#1}%
  \footnote{#2}%
  \addtocounter{footnote}{-1}%
  \endgroup}
\providecommand{\srcfnmark}[1]{%
  \begingroup
  \renewcommand{\thefootnote}{#1}%
  \footnotemark%
  \addtocounter{footnote}{-1}%
  \endgroup}
\providecommand{\srcfntext}[2]{%
  \begingroup
  \renewcommand{\thefootnote}{#1}%
  \footnotetext{#2}%
  \endgroup}
\makeatother
\RequirePackage{tikz}
\newcommand{\NoetherRosetteGlyph}{%
  \begin{tikzpicture}[x=0.1363584905660377pt,y=-0.1363584905660377pt,baseline=(current bounding box.south)]
    \path[use as bounding box] (0,0) rectangle (108,108);
    \fill[even odd rule]
    (33.5,84) -- (32.5,83) -- (30.5,83) -- (29,81.5) -- (29,76.5) -- (30.5,75) -- (31.5,75) -- (32.5,74) -- (34.5,74) -- (36,72.5) -- (36,71.5) -- (35.5,71) -- (33.5,71) -- (33,70.5) -- (33,69.5) -- (30.5,67) -- (28.5,69) -- (26.5,69) -- (24.5,71) -- (23.5,71) -- (22.5,72) -- (21.5,72) -- (20.5,73) -- (18.5,73) -- (17.5,72) -- (16.5,73) -- (16,72.5) -- (16,68.5) -- (18.5,66) -- (19.5,66) -- (20,65.5) -- (20,64.5) -- (26,58.5) -- (26,56.5) -- (24,54.5) -- (24,53.5) -- (17,46.5) -- (17,45.5) -- (16.5,45) -- (15.5,45) -- (15,44.5) -- (15,42.5) -- (15.5,42) -- (17.5,42) -- (19.5,44) -- (20.5,44) -- (21,44.5) -- (20,45.5) -- (21,46.5) -- (21,47.5) -- (22,48.5) -- (22,49.5) -- (29,56.5) -- (29,57.5) -- (27,59.5) -- (27,60.5) -- (26,61.5) -- (26,62.5) -- (24,64.5) -- (24,65.5) -- (22,67.5) -- (22.5,68) -- (24.5,68) -- (25.5,67) -- (28.5,67) -- (29.5,66) -- (31.5,66) -- (32.5,65) -- (33.5,65) -- (34.5,64) -- (36.5,64) -- (37.5,63) -- (38.5,63) -- (39.5,62) -- (41.5,62) -- (42.5,61) -- (44.5,61) -- (45.5,60) -- (46.5,60) -- (47.5,59) -- (48.5,59) -- (49,58.5) -- (49,57.5) -- (49.5,57) -- (51.5,57) -- (52,56.5) -- (50.5,55) -- (49.5,55) -- (48.5,54) -- (47.5,54) -- (46.5,53) -- (45.5,53) -- (43,50.5) -- (43,49.5) -- (42,48.5) -- (42,47.5) -- (41.5,47) -- (40.5,47) -- (39.5,48) -- (33.5,48) -- (32,46.5) -- (32,44.5) -- (34,42.5) -- (34,40.5) -- (33.5,40) -- (32.5,40) -- (31.5,39) -- (26.5,39) -- (26,38.5) -- (26,35.5) -- (27.5,34) -- (31.5,34) -- (33,35.5) -- (33,36.5) -- (34.5,38) -- (37.5,38) -- (38,37.5) -- (38,36.5) -- (41,33.5) -- (41,31.5) -- (40,30.5) -- (40,27.5) -- (39,26.5) -- (39,25.5) -- (38,24.5) -- (38,23.5) -- (37,22.5) -- (37,19.5) -- (37.5,19) -- (38.5,20) -- (39.5,20) -- (41,21.5) -- (41,24.5) -- (42,25.5) -- (42,27.5) -- (43,28.5) -- (43,30.5) -- (44,31.5) -- (44,35.5) -- (45,36.5) -- (45,37.5) -- (45.5,38) -- (46.5,38) -- (47,38.5) -- (47,40.5) -- (48,41.5) -- (48,42.5) -- (49,43.5) -- (49,44.5) -- (50,45.5) -- (50,48.5) -- (53,51.5) -- (52,52.5) -- (52,55.5) -- (53.5,57) -- (56.5,57) -- (58.5,59) -- (59.5,59) -- (60.5,60) -- (61.5,60) -- (62,60.5) -- (62,61.5) -- (63,62.5) -- (63,64.5) -- (63.5,65) -- (64.5,65) -- (65.5,64) -- (68.5,64) -- (69.5,65) -- (73.5,65) -- (74,65.5) -- (74,66.5) -- (75,67.5) -- (75,68.5) -- (75.5,69) -- (76,68.5) -- (76,66.5) -- (76.5,66) -- (78.5,68) -- (80.5,68) -- (81.5,69) -- (84.5,69) -- (85,68.5) -- (84,67.5) -- (84,65.5) -- (83.5,65) -- (82.5,65) -- (82,64.5) -- (82.5,64) -- (84.5,64) -- (86,65.5) -- (86,66.5) -- (88.5,69) -- (89.5,69) -- (90,69.5) -- (90,70.5) -- (91,71.5) -- (91,72.5) -- (93,74.5) -- (91.5,76) -- (90.5,76) -- (88.5,74) -- (87.5,74) -- (86.5,73) -- (85.5,73) -- (83.5,71) -- (82.5,71) -- (81.5,70) -- (80.5,70) -- (79.5,69) -- (77.5,69) -- (74.5,72) -- (73.5,72) -- (73,72.5) -- (73,74.5) -- (73.5,75) -- (78.5,75) -- (79.5,74) -- (80.5,74) -- (81,74.5) -- (81,78.5) -- (80,79.5) -- (80,80.5) -- (78.5,82) -- (73.5,82) -- (73,81.5) -- (73,79.5) -- (72,78.5) -- (72,77.5) -- (71.5,77) -- (70.5,77) -- (69.5,76) -- (68.5,76) -- (68,76.5) -- (68,77.5) -- (67,78.5) -- (67,79.5) -- (66.5,80) -- (65.5,80) -- (65,80.5) -- (66,81.5) -- (64.5,83) -- (63,81.5) -- (63,79.5) -- (62,78.5) -- (62,77.5) -- (61,76.5) -- (61,74.5) -- (60,73.5) -- (60,72.5) -- (59,71.5) -- (59,68.5) -- (58,67.5) -- (58,64.5) -- (56,62.5) -- (56,61.5) -- (55.5,61) -- (54.5,62) -- (53.5,62) -- (52,63.5) -- (52,64.5) -- (51,65.5) -- (51,66.5) -- (50,67.5) -- (50,68.5) -- (49,69.5) -- (49,70.5) -- (48,71.5) -- (48,75.5) -- (46,77.5) -- (46,79.5) -- (44.5,81) -- (42.5,81) -- (42,80.5) -- (44,78.5) -- (44,77.5) -- (43.5,77) -- (41.5,77) -- (39.5,75) -- (38.5,76) -- (37.5,76) -- (37,76.5) -- (37,78.5) -- (36,79.5) -- (38,81.5) -- (38,82.5) -- (37.5,83) -- (34.5,83) -- cycle
    (65.5,22) -- (63.5,22) -- (63,21.5) -- (63.5,21) -- (65.5,21) -- (66,21.5) -- cycle
    (55.5,28) -- (53.5,28) -- (53,27.5) -- (53.5,27) -- (54.5,27) -- (57.5,24) -- (59.5,24) -- (60,24.5) -- (57.5,27) -- (56.5,27) -- cycle
    (64.5,26) -- (64,25.5) -- (64.5,25) -- (65,25.5) -- cycle
    (63.5,29) -- (63,28.5) -- (63,27.5) -- (63.5,27) -- (64.5,27) -- (65,27.5) -- cycle
    (58.5,55) -- (57,53.5) -- (57,52.5) -- (56,51.5) -- (56,49.5) -- (57,48.5) -- (57,46.5) -- (60,43.5) -- (60,42.5) -- (59.5,42) -- (58.5,42) -- (58,41.5) -- (59,40.5) -- (59,39.5) -- (61.5,37) -- (62.5,37) -- (63.5,36) -- (65.5,36) -- (66.5,37) -- (68.5,37) -- (70,35.5) -- (69,34.5) -- (69,32.5) -- (68,31.5) -- (68.5,31) -- (74.5,31) -- (75,31.5) -- (75,32.5) -- (76,33.5) -- (76,35.5) -- (75,36.5) -- (75,37.5) -- (72,40.5) -- (72,41.5) -- (71,42.5) -- (72,43.5) -- (72,44.5) -- (72.5,45) -- (73.5,45) -- (74.5,46) -- (75.5,46) -- (76.5,45) -- (78.5,45) -- (79.5,44) -- (80,44.5) -- (80,47.5) -- (79.5,48) -- (77.5,48) -- (76.5,49) -- (74.5,49) -- (73.5,50) -- (71.5,50) -- (70.5,51) -- (68.5,51) -- (67.5,52) -- (65.5,52) -- (64.5,53) -- (63.5,53) -- (62.5,54) -- (59.5,54) -- cycle
    (86.5,50) -- (85.5,50) -- (85,49.5) -- (86,48.5) -- (86,46.5) -- (85.5,46) -- (82.5,46) -- (81.5,47) -- (81,46.5) -- (82,45.5) -- (82,44.5) -- (82.5,44) -- (84.5,44) -- (86.5,42) -- (88.5,42) -- (89.5,41) -- (92.5,41) -- (93,41.5) -- (93,42.5) -- (94,43.5) -- (94,44.5) -- (93.5,45) -- (92.5,45) -- (88.5,49) -- (87.5,49) -- cycle
    (44,42.5) -- (43.5,42) -- (43,42.5) -- (43.5,43) -- cycle
    (45,43.5) -- (44.5,43) -- (44,43.5) -- (44.5,44) -- cycle
    (80.5,60) -- (79.5,60) -- (79,59.5) -- (79,56.5) -- (81,54.5) -- (81,53.5) -- (82.5,52) -- (83,52.5) -- (83,53.5) -- (82,54.5) -- (82,55.5) -- (81.5,56) -- (80.5,56) -- (80,56.5) -- (80,57.5) -- (81,58.5) -- (81,59.5) -- cycle
    (55.5,54) -- (55,53.5) -- (55.5,53) -- (56,53.5) -- cycle
    (43,65.5) -- (42.5,65) -- (42,65.5) -- (42.5,66) -- cycle
    (61,67.5) -- (60.5,67) -- (60,67.5) -- (60.5,68) -- cycle
    (47,69.5) -- (47,68.5) -- (46.5,68) -- (46,68.5) -- (46,69.5) -- (46.5,70) -- cycle
    (62,68.5) -- (61.5,68) -- (61,68.5) -- (61.5,69) -- cycle
    (67.5,97) -- (65.5,97) -- (64,95.5) -- (64,94.5) -- (61.5,92) -- (60.5,92) -- (54.5,86) -- (53.5,86) -- (49.5,90) -- (48.5,90) -- (47.5,91) -- (46.5,91) -- (45.5,92) -- (45,91.5) -- (45,90.5) -- (45.5,90) -- (46.5,90) -- (48.5,88) -- (49.5,88) -- (54.5,83) -- (55.5,83) -- (60.5,88) -- (61.5,88) -- (62.5,89) -- (63.5,89) -- (64,89.5) -- (64,90.5) -- (64.5,91) -- (66.5,91) -- (67,90.5) -- (66,89.5) -- (67,88.5) -- (65,86.5) -- (65,85.5) -- (65.5,85) -- (66.5,85) -- (67,85.5) -- (67,86.5) -- (68,87.5) -- (68,90.5) -- (69,91.5) -- (69,94.5) -- (70,95.5) -- (69.5,96) -- (68.5,96) -- cycle
    (40.5,96) -- (36.5,96) -- (36,95.5) -- (37,94.5) -- (37,91.5) -- (40.5,88) -- (41.5,88) -- (42,88.5) -- (41,89.5) -- (41,92.5) -- (42,93.5) -- (42,94.5) -- cycle;
  \end{tikzpicture}%
}
\providecommand{\noethpIIrosette}{\makebox[9.44444pt][c]{\raisebox{-3.85pt}[0pt][0pt]{\NoetherRosetteGlyph}}}
\providecommand{\NoetherRosetteInline}{\makebox[17.72671698pt][c]{\raisebox{-3.85pt}[10.87671698pt][3.85pt]{\NoetherRosetteGlyph}}}
\usepackage{url}
\begin{document}
\begin{titlepage}
\centering
\vspace*{1.8cm}
{\LARGE\bfseries エミー・ネーター数学論文集\par}
\bigskip
{\large 日本語訳 — 基礎本文\par}
\bigskip
論文1–43、1929/30年冬学期講義、\par
KapfererとNoetherによる論文・補遺、文献目録\par
\vfill
\begin{minipage}{0.94\textwidth}\small
\textbf{著者・翻訳}\par
原著者はEmmy Noetherおよび各論文に記された共著者である。
1929/30年の講義録はM. Deuringの筆録による。
日本語訳はAIによる翻訳である。初版の制作主体はManuscript Typesetting Project / OpenAI Codexであり、当時の正確なモデルは現存する来歴資料からは確定していない。
本版のAIによる訂正・組版：OpenAI Codex — GPT-6 Astra、Ultra。AIによる訂正・組版：2026年10月23–24日；Claude Opus 5.5（Anthropic）、2026年10月26日。刊行日：2026年10月26日。
発行者による人間の翻訳または学術的査読を示すものではない。
\medskip
\textbf{底本と来歴}\par
日本語版2026-08-13 r10：\url{https://doi.org/10.5281/zenodo.21921485}。
ドイツ語版ED0014/v051：\url{https://doi.org/10.5281/zenodo.21908301}。
各作品の原刊情報、著者名、引用は本文に記載する。
\medskip
\textbf{権利について}\par
原著および保護される派生的寄与の権利は、それぞれの権利者に帰属する。
本版は、権利者が異議を唱えないであろうとの発行者の推定に基づき、非商用かつ純粋に教育的な目的のために無料で提供する。
これは推定であり、許諾が確認されたことを意味しない。無料で閲覧できること自体は、派生版の公開を許可するものではない。
権利者は\url{https://github.com/KokunoYumeto/emmy-noether-ja/issues/new}から異議を申し出ることができる。発行者は権利者の要請に応じ、該当資料を取り下げる。
本版による新たな包括的ライセンスの付与またはパブリックドメインへの提供は行わない。
旧版に付された権利表示は、その旧版と実際に付与された権利の範囲に限り存続する。
\end{minipage}
\end{titlepage}
\setcounter{page}{1}


\begin{center}
{\Large\bfseries 1. 三元四次形式の形式系の形成について}\par
\vspace{1.5em}
\emph{Sitzungsberichte der Physikalisch-medizinischen Sozietät in Erlangen} 39 (1907), pp. 176--179
\end{center}

\vspace{3em}
\begin{center}
（著者の学位論文からの抜粋。）
\end{center}

三元四次形式の形式系については、Gordan、Maisano、Pascal の研究\footnote{P. Gordan, ``Über das volle Formensystem der ternären biquadratischen Form'' \(f=x_1^3x_2+x_2^3x_3+x_3^3x_1\) (Math. Annalen, vol. XVII, pp. 217--233, 1880); G. Maisano, 1. ``Sistemi completi dei primi cinque gradi della forma ternaria biquadratica e degl' invarianti, covarianti e contravarianti di sesto grado'' (Giorn. di Battaglini XIX). 2. ``Sui covarianti indipendenti di \(6^\circ\) grado nei coefficienti della forma biquadratica ternaria'' (Rend. Circ. Mat. di Palermo I, 1887); E. Pascal, ``Contributo alla teoria della forma ternaria biquadratica e delle sue varie decomposizioni in fattori'' (Memoria premiata dalla R. Accademia delle scienze fisiche e matematiche di Napoli, 1905).}がある。Gordan は、特殊な自己同型形式
\[
f=x_1^3x_2+x_2^3x_3+x_3^3x_1
\]
に対して、54 個の形成から成る完全な形式系を、二元の場合の形式系について彼が与えたものと同種の原理に基づいて構成した。

Maisano の研究では、一般の四次形式について、五次の位数まで（それを含む）\footnote{ここで「位数」とは係数における次元を意味し、「次数」とは変数における次元を意味する。}の形式が、さらに高位数のいくつかの不変式・共変式・反変式とともに、Gordan が \emph{Mathematische Annalen} 第 I 巻で三元三次形式に用いた方法に従って構成されている。Pascal は Maisano の結果を用い、主として四次形式の因子への分解の問題を扱っている\footnote{Maisano は第二の短いノートで、位数 6 かつ次数 6 の三つの共変式の間の線形従属を証明しようとしている。この完全とはいえない証明に対して、Pascal は、位数 6 の三つの共変式の線形独立性、および位数 9 の三つの不変式の線形独立性を証明したと考えている。しかし、三つの共変式の間にも、また三つの不変式の間にも、実際には線形関係が存在することが私の計算の過程で明らかになった。Pascal の記法では、その明示公式は
\[
0=20\Omega_1+6\Omega_2-3\Omega_3+4fC_1-12fC_2-4A\cdot\Delta,
\]
\[
0=4A^3-15AB+30C-90D+9E.
\]
である。}。

\textbf{私の研究の目的は、一般三元四次形式の形式系を構成することである。さしあたり、主要な基礎だけを与え、いわゆる「相対的完全系」\footnote{用語については Gordan--Kerschensteiner, \emph{Vorlesungen über Invariantentheorie}, vol. II, p. 227 を参照。}を確立する。}

三元形式の系を作る基本的な考え方は、二元の場合と同じである。最初の相対的完全系、すなわち二元の場合から移された形式の系から出発し、一定の法則に従って、法が次第に高くなる系へと進む。その過程は、ある法を基本形式と見なしたときにその法の系が有限かつ既知になるか、またはある法が不変式を因子にもつ形式に還元できるようになるまで続けられる。第一の場合には、相対的完全系とその法の系とのトランスヴェクションによって絶対的完全系が得られ、第二の場合には相対的完全系と絶対的完全系とが一致する。Hilbert による形式系の有限性の一般的証明により、この手続きは必ず終わらなければならない。

本場合では、第一の相対的完全系の法 \((abc)\) は、法
\[
\Delta=(abc)^2a_x^2b_x^2c_x^2\quad\text{および}\quad \nu=(abu)^4
\]
に帰着できる。ここから、法 \(\nu\) に関する相対的完全系は容易に得られる。そこで法の列として次の形式を選ぶ：
\[
\nu;\qquad \nu^{(\nu)}=(\nu\nu_1x)^4=s_x^4,
\]
\[
\nu^{(s)}=(ss'u)^4\ldots,
\]
さらに、法 \(s\) に関する相対的完全系の形成に現れる二つの二次形式、すなわち \(u_\sigma^2\) と \(t_x^2\) を、追加の法として付け加える。すると、\(s\) に続く法、すなわち \((ss'u)^4\) は、法 \((\varrho,t)\) に還元可能であることを示すことができる。しかし二つの二次形式の同時系は有限かつ既知であるから、上に述べた終結点に到達したことになる。法 \((\varrho,t)\) に関する相対的完全系として、331 個の形成が得られる。

用いた方法について若干の詳細を加える。

形式を作る基本過程である折り畳み過程は、三元の場合には次のように定義できる。

記号的積
\[
s_x^m t_x^n u_\sigma^\mu u_\tau^\nu
\]
において、因子の対
\[
\begin{array}{c|c|c|c|c}
 & s_xt_x & u_\sigma u_\tau & s_xu_\sigma\ \text{または}\ s_xu_\tau & t_xu_\tau\ \text{または}\ t_xu_\sigma\\
\text{それぞれを} & (stu) & (\sigma\tau x) & s_\sigma\ \text{または}\ s_\tau & t_\tau\ \text{または}\ t_\sigma\\
\text{で置き換える折り畳み} & \mathrm{I} & \mathrm{II} & \mathrm{III} & \mathrm{IV},
\end{array}
\]
を行うなら、そのようにして生じる形式は、もとの形式から折り畳みによって得られたものである\footnote{Gordan, Math. Annalen, vol. XVII, p. 219.}。

ここではなお常に \(s_\sigma=0;\ t_\tau=0\) とおくことができる。個々の折り畳み相互の関係については、次の定理が成り立つ。

\textbf{折り畳み I と II は基本折り畳みであり、合成の順序に依存せず、折り畳み III と IV はこれらから合成できる。換言すれば、与えられた式から折り畳みによって生じるすべての形式を作るには、折り畳み I と II だけを適用すればよい}\footnote{この定理には、\(n\) と \(\mu\)、または \(m\) と \(\nu\) がそれぞれ同時に零になる場合に例外がある。このとき「形式列」は対角項に縮退する。}。

「形式列」とは、初期形式と、それを自己の中へ折り畳むことによって生じるすべての形式とを合わせたものをいう。上の定理により、形式列
\[
s_x^m t_x^n u_\sigma^\mu u_\tau^\nu
\]
は長方形の図式に配列することができる。この図式で進んでいくと、次が得られる：
\begin{enumerate}
\item 一列右へ進むことにより、隣接する形式から折り畳み I によって生じるすべての形式。
\item 一行下へ進むことにより、その上に立つ形式から折り畳み II によって生じるすべての形式、したがって折り畳みによって生じるすべての形式。
\end{enumerate}

これらの仮定のもとで、還元子に関する定理は、
\begin{center}
「還元子とは可約な形式列である」
\end{center}
という定義のもと、最も一般的な形で次のように述べられる。

\begin{center}
\textbf{形式列の初期形式が可約であるのは、その一つの項が還元子との折り畳みによって生じているためであり、かつその形式列の終形式がこの同じ項から折り畳みによって生じているならば、形式列全体は可約である。}
\end{center}

さらに挙げるべき還元方法は次の二つである：
\begin{enumerate}
\item いわゆる「二重還元」、すなわち高次の形式の間の関係を得るために、一つの形式を二通りの方法で還元すること。
\item 「分解した形式による折り畳み」。これは一部では還元子による還元の性格をもち、一部では「二重還元」の性格をもつ。
\end{enumerate}

\clearpage

\begin{center}
{\Large\bfseries 2. 三元四次形式の形式系の形成について}\par
\vspace{1.5em}
Journal f. d. reine u. angew. Math. 134 (1908), pp. 23--90 および二つの表
\end{center}

\vspace{3em}

\begin{center}
\textbf{序論。}
\end{center}

三元四次形式の形式系については、Gordan、Maisano、Pascal の研究\footnote{序論および第 I 章からの短い抜粋は、\emph{Sitzungsberichte der phys.-med. Sozietät Erlangen 1907}, pp. 176--179 に発表された。}\footnote{P. Gordan, 三元四次形式の完全な形式系について
\[
f=x_1^3x_2+x_2^3x_3+x_3^3x_1.
\]
(\emph{Math. Annalen}, vol. XVII (1880), pp. 217--233.)
G. Maisano, 1) \emph{Sistemi completi dei primi cinque gradi della forma ternaria biquadratica e degl' invarianti, covarianti e contravarianti di sesto grado}. (Giorn. di Battaglini XIX.)
2) \emph{Sui covarianti indipendenti di $6^\circ$ grado nei coefficienti della forma biquadratica ternaria}. (Rend. Circ. Mat. di Palermo I, 1887.)
E. Pascal, \emph{Contributo alla teoria della forma ternaria biquadratica e delle sue varie decomposizioni in fattori}. (Memoria premiata dalla R. Accademia delle scienze fisiche e matematiche di Napoli. 1905.)}がある。Gordan 氏は、特殊な自己同型形式
\[
f=x_1^3x_2+x_2^3x_3+x_3^3x_1
\]
について、54 個の形成から成る完全な形式系を、二元領域の形式系について彼が与えたものと同種の原理に基づいて構成している。

Maisano 氏においては、一般の四次形式について、五次の位数まで（それを含む）\footnote{ここで「位数」とは係数における次元を、「次数」とは変数における次元を意味する。}の形式が、さらに高位数のいくつかの不変式、共変式、反変式とともに、Gordan 氏が \emph{Math. Annalen} 第 I 巻で三元三次形式に用いた方法に従って構成されている。Pascal 氏は Maisano の結果を用い、主として四次形式の因子への分解の問題を扱っている\footnote{第二の短いノートで、Maisano 氏は位数 6 かつ次数 6 の三つの共変式の線形従属を証明しようとしている。この完全とはいえない証明に対して、Pascal 氏は、位数 6 の三つの共変式の線形独立性、さらに位数 9 の三つの不変式の線形独立性を証明したと考えている。しかし、三つの共変式の間にも、また三つの不変式の間にも、実際には線形関係が存在することは、本論文の式 (13), § 11 および § 17 の注の式 (22) が示しており、そこではこれらの関係が明示的に与えられている。Pascal 氏からの連絡によれば、この矛盾は、彼の反変式 $p$（ここでは $\varrho$ と呼ぶ）の表式、同論文 p. 46 における数値的誤りによって説明される。}。

\emph{以下の研究の目的は、一般三元四次形式の形式系を構成することである。すなわち本論文では主要な基礎だけを与え、いわゆる「相対的完全系」\footnote{この名称については § 3a を参照。}を構成する。}

本論文は Gordan の研究と密接に結びついている。ただし、そこではごく一般的にしか与えられていなかった原理を、まず個々に展開する必要があった。折り畳み相互の関係に関する一つの定理を用い、さらに「形式列」（§ 1）を導入することによって、還元定理を明確に定式化し補完する（§ 3）。一方、特別な級数展開を再帰的に構成すること（§ 2）が、還元を実際に遂行するための計算手段を与える。

三元形式の系を形成する基本思想は、二元領域の場合と同じである。最初の相対的完全系、すなわち二元の場合から移された形式の系から出発して、一定の法則に従い、法がしだいに高くなる系へ進む。この過程は、ある法の系が、その法を基本形式としてとったときに、有限かつ既知になるまで、またはある法が不変式を因子にもつ形式へ還元できるようになるまで続く。第一の場合には、相対的完全系とその法の系とのトランスヴェクションによって絶対的完全系が生じ、第二の場合には相対的完全系と絶対的完全系とが一致する。Hilbert 氏によって一般に証明された形式系の有限性により、この手続きは必ず終結しなければならない。

本場合では、第一の相対的完全系の法 $(abc)$（§ 4）は法
\[
\Delta=(abc)^2a_x^2b_x^2c_x^2\quad\text{および}\quad \nu=(abu)^4
\]
（§ 5）に帰着され、ついで法 $\nu$ に関する相対的完全系が求められる（§ 6）。これらの結果は、一般四次形式については Gordan の研究にすでに与えられている。しかし、ここでの導出方法は、より高い次数の形式へいっそう容易に移すことができる。

法の列として、いま次の形式を選ぶ（§ 7）：
\[
\nu,\qquad \nu(\nu)=(\nu\nu_1x)^4=s_x^4,\qquad \nu(s)=(ss'u)^4,\ldots .
\]

法 $s$ に関する相対的完全系を形成する際には、系の中に現れる二つの二次形式 $u_\varrho^2$ と $t_x^2$ を法として付け加え（§§ 9, 10）、それに応じて法 $(s,\varrho,t)$ に関する相対的完全系を作る。ついで、法 $(ss'u)^4$ が法 $\varrho$ と $t$ に還元可能であることを示す（§ 17）。これによって、次の高い系は法 $(\varrho,t)$ に関する相対的完全系へ移る。ところが二つの二次形式の同時系は有限かつ既知であり、一般の場合には 20 個の形成から成る\footnote{Clebsch--Lindemann, \emph{Vorlesungen über Geometrie}. (vol. I, p. 288 以下) 二つの共変的形式の系。Gordan, \emph{Über Büschel von Kegelschnitten}. (\emph{Math. Ann.} vol. XIX, p. 530.) 二つの反変的形式の系。}ので、法 $(\varrho,t)$ に関する相対的完全系の構成によって上に述べた終結点に到達する。この系を $(\varrho,t)$ の系とトランスヴェクションして絶対的完全系を作ることは、後に残しておく。

法 $(\varrho,t)$ に関する相対的完全系として、331 個の形成が得られ、それらは付された表において、変数 $x$ と $u$ における次数に従って配列されている。

最後に、導入した記号の概観を与える：
\begin{align*}
f&=a_x^4,\\
\theta&=\theta_x^4u_\vartheta^2=(abu)^2a_x^2b_x^2,\\
K&=K_x^6u_k^3=(a\theta u)u_\vartheta^2a_x^3\theta_x^3,\\
j&=u_j^6=(a\theta u)^4u_\vartheta^2,\\
\Delta&=\Delta_x^6=a_\vartheta^2a_x^2\theta_x^4=(abc)^2a_x^2b_x^2c_x^2,\\
N&=N_x^8u_n=(a\Delta u)a_x^3\Delta_x^5,\\
\nu&=u_\nu^4=(abu)^4,\\
H&=H_x^2u_\eta^4=(\nu\nu_1x)^2u_\nu^2u_{\nu_1}^2,\\
L&=L_x^3u_l^6=(\nu\eta x)H_x^2u_\nu^3u_\eta^3,\\
g&=g_x^6=(\nu\eta x)^4H_x^2,\\
\sigma&=u_\sigma^6=H_\nu^2u_\nu^2u_\eta^4,\\
s&=s_x^4=(\nu\nu_1x)^4,\\
Z&=Z_x^4u_\zeta^2=(ss'u)^2s_x^2s_x'^2,\\
\varrho&=u_\varrho^2=\theta_\nu^4u_\vartheta^2,\\
t&=t_x^2=a_\eta^4H_x^2,\\
i&=a_\nu^4,\\
J&=s_\nu^4.
\end{align*}

\begin{center}
{\Large\bfseries 第 I 章。\quad 三元形式に関する一般定理。}
\end{center}

\begin{center}
\textbf{§ 1. \quad 折り畳み過程。\quad 形式列。}
\end{center}

形式を作る基本過程は折り畳み過程であり、三元領域では次のように定義できる。記号的積
\[
s_x^m t_x^n u_\sigma^\mu u_\tau^\nu .
\]
が与えられているとする。因子の対
\[
s_xt_x\qquad u_\sigma u_\tau\qquad s_xu_\sigma\ \text{または}\ s_xu_\tau
\qquad t_xu_\tau\ \text{または}\ t_xu_\sigma
\]
をそれぞれ
\[
(stu)\qquad (\sigma\tau x)\qquad s_\sigma\ \text{または}\ s_\tau
\qquad t_\tau\ \text{または}\ t_\sigma,
\]
すなわち折り畳み
\[
\text{I}\qquad\text{II}\qquad\text{III}\qquad\text{IV},
\]
によって置き換えると、得られる形式は元の形式から折り畳みによって生じたものである\footnote{Gordan, \emph{Math. Annalen} vol. XVII, p. 219.}。

このとき式 $s_x^m t_x^n u_\sigma^\mu u_\tau^\nu$ は、二つの形式 $S\cdot T$ の実際の積として見ることもできるし、複数の記号列をもつ一つの形式
\[
s_x^m t_x^n u_\sigma^\mu u_\tau^\nu=A_x^{m+n}u_\lambda^{\mu+\nu}
\]
として見ることもできる。第一の場合には形式 $S$ と $T$ との折り畳みといい、第二の場合には「形式 $A$ の自己への折り畳み」という。簡潔のため、以下では常に（実際、後で考察されるすべての形式についてそうである）
\[
\tag{a}
s_\sigma^\lambda=0,\qquad t_\tau^\varkappa=0
\]
が、零でない任意の $\lambda$ と $\varkappa$ について、また任意の他の折り畳みと合わせて成り立つものと仮定する。これは、形式 $s_x^m t_y^n u_\sigma^\mu v_\tau^\nu$ がいわゆる「正規形式」である、すなわち $\Omega$ 過程を変数 $x,u$ についても、また変数 $y,v$ についても適用すると消える、ということを意味する。

したがって条件 (a) は何らの制限でもなく、常に達成できる\footnote{Gordan, \emph{Math. Annalen} vol. V, p. 104 を参照。}。

個々の折り畳みの関係については次の定理が成り立つ：

\emph{定理 I. 折り畳み I と II は基本折り畳みであり、合成の順序によらず、折り畳み III と IV はこれらから合成される。換言すれば、与えられた式から折り畳みによって生じるすべての形式を作るには、折り畳み I と II だけを適用すればよい。}

証明：同一性定理、または行列に対する積定理により、(a) を考慮すれば、
\[
(\widehat{st}\sigma x)=-t_\sigma,\qquad
(\widehat{\sigma\tau}su)=-s_\tau,
\]
\[
(\widehat{st}\tau x)=s_\tau,\qquad
(\widehat{\sigma\tau}tu)=t_\sigma,
\]
\[
(stu)(\sigma\tau x)=s_\tau+t_\sigma-u_x\cdot s_\tau t_\sigma .
\]
を得る。したがって折り畳み I と II を合成することで折り畳み III と IV が生じる。これは、折り畳み II を折り畳み I、すなわち因子 $(stu)u_\sigma u_\tau$ に適用する場合にも、折り畳み I を折り畳み II、すなわち因子 $(\sigma\tau x)s_xt_x$ に適用する場合にも成り立つ。

\emph{形式列}\footnote{Clebsch, \emph{Über eine Fundamentalaufgabe der Invariantentheorie} (Abh. der Gött. Ges. d. Wiss. vol. XVII): § 17 および § 18 の末尾を参照。「形式列」は、そこに導入された「本来的に還元された同値系」とは、より高い形式による順序づけの原理によって異なる。}とは、初期形式と、それを自己の中へ折り畳むことによって生じるすべての形式とを合わせたものをいう。また形式列をその初期形式によって表す。ここで「より高い形式」とは、自己への折り畳みをより多く含む任意の形式を意味する。ただし、同等の折り畳み $s_\tau$ と $t_\sigma$ の間では、一方をより高いものとして正規化しなければならない。形式列の形成法則から次が従う：

1) 形式列の各形式は初期形式の記号に関して線形である。

2) 形式列は、それぞれ二つの反変的記号列をもつ初期形式については一意に定まる。より多くの記号列をもつ初期形式については、同一性定理によって結びつけられる折り畳みのうちから特別なものを区別することによって、一意に正規化されなければならない。

定理 I により、形式列 $s_x^m t_x^n u_\sigma^\mu u_\tau^\nu$（$m>n;\ \mu>\nu$）は次の長方形の図式に従って配置できる：
\begingroup\small
\[
\begin{array}{c|c|c|c|c}
s_x^m t_x^n u_\sigma^\mu u_\tau^\nu
& (stu)
& (stu)^2
& \cdots
& \noethpIIrosette\\[0.3em]
(\sigma\tau x)
& t_\sigma,\ s_\tau
& t_\sigma(stu),\ s_\tau(stu)
& \cdots
& \vdots\\[0.3em]
(\sigma\tau x)^2
& t_\sigma(\sigma\tau x),\ s_\tau(\sigma\tau x)
& t_\sigma^2,\ t_\sigma s_\tau,\ s_\tau^2
& \cdots
& \vdots\\[0.3em]
\vdots&\vdots&\vdots&\ddots&\vdots\\[0.3em]
(\sigma\tau x)^\nu
& t_\sigma(\sigma\tau x)^{\nu-1},\ s_\tau(\sigma\tau x)^{\nu-1}
& t_\sigma^2(\sigma\tau x)^{\nu-2},\ t_\sigma s_\tau(\sigma\tau x)^{\nu-2},\ s_\tau^2(\sigma\tau x)^{\nu-2}
& \cdots
& \vdots\\[0.3em]
\vdots
& t_\sigma(\sigma\tau x)^\nu
& t_\sigma^2(\sigma\tau x)^{\nu-1},\ t_\sigma s_\tau(\sigma\tau x)^{\nu-1}
& \cdots
& \vdots\\[0.3em]
\vdots
& \vdots
& t_\sigma^2(\sigma\tau x)^\nu
& \cdots
& \vdots
\end{array}
\]
\endgroup
\begingroup\small
\[
\begin{array}{r@{\quad}c}
\noethpIIrosette\srcfnmark{*)}
&
\begin{array}{c|c|c}
(stu)^n & \cdots & \cdots\\[0.3em]
t_\sigma(stu)^{n-1},\ s_\tau(stu)^{n-1}
& s_\tau(stu)^n & \cdots\\[0.3em]
t_\sigma^2(stu)^{n-2},\ t_\sigma s_\tau(stu)^{n-2},\ s_\tau^2(stu)^{n-2}
& t_\sigma s_\tau(stu)^{n-1},\ s_\tau^2(stu)^{n-1}
& s_\tau^2(stu)^n .
\end{array}
\end{array}
\]
\endgroup
\srcfntext{*)}{ここで、および以下でも同様に、\NoetherRosetteInline{}で示された下部は、図式の右上に同じ印で示された場所に接続する。}

ここで、

1) 一列右へ進むと、隣接する形式から折り畳み I によって生じるすべての形式が得られ、

2) 一行下へ進むと、その上にある形式から折り畳み II によって生じるすべての形式が得られ、

したがって折り畳みによって生じるすべての形式が得られる。

全形式列の任意の形式 $t_\sigma^\kappa s_\tau^\lambda(stu)^\mu$ または $t_\sigma^\kappa s_\tau^\lambda(\sigma\tau x)^\nu$ は新しい形式列の始まりとなる。その形式列は、長方形の図式のうち $t_\sigma^\kappa s_\tau^\lambda(stu)^\mu$ または $t_\sigma^\kappa s_\tau^\lambda(\sigma\tau x)^\nu$ を左上隅とし、因子 $t_\sigma^\kappa s_\tau^\lambda$ をもつすべての形式から成る。

補足。定理 I には例外がある。数 $n$ と $\mu$（それぞれ $m$ と $\nu$）が零となり、折り畳み I, II, IV がもはや存在しない一方で、折り畳み III（それぞれ IV）はなお存在する場合である。この場合には、図式の対角項を形式列と呼ぶ：
\[
\begin{array}{c@{\qquad}c@{\qquad}c}
\begin{array}{ccccc}
s_x^m u_\tau^\nu &&&&\\
&s_\tau&&&\\
&&s_\tau^2&&\\
&&&\ddots&\\
&&&&s_\tau^\nu
\end{array}
&\text{それぞれ}&
\begin{array}{ccccc}
t_x^n u_\sigma^\mu &&&&\\
&t_\sigma&&&\\
&&t_\sigma^2&&\\
&&&\ddots&\\
&&&&t_\sigma^\mu
\end{array}
\end{array}
\]
折り畳み I と II が存在しないことに応じて、この場合には形式列の極形式による級数展開は常に初期項に縮退する。しかし、還元子に関する定理は有効に残る。

\begin{center}
\textbf{§ 2. \quad 形式列の極形式による級数展開。}
\end{center}

$n$ 個の変数をもつ形式については、二種類の級数展開が明示的に立てられている\footnote{Gordan, \emph{Über Kombinanten}, §§ 2, 5 (\emph{Math. Ann.} vol. V) を参照。}：

1) 反変的変数の列を一つずつもつ形式 $s_x^m u_\tau^\nu$ について、$u_x$ の冪によって進む級数で、その係数は「正規形式」である。

2) 共変的変数の二列をもつ形式 $s_x^m t_x^n$ について、基本共変式の極形式（形式列 $s_x^m t_x^n$ の極形式）による展開であり、三元の場合には次のように書かれる\footnote{Study, \emph{Methoden zur Theorie der ternären Formen} (Teubner 1889) II. § 7 を参照。そこに与えられた導出と異なり、われわれの導出は数値係数 $C_{pr\varrho}$ を迅速に計算する方法を与える。特殊化 a), a') のもとでの Clebsch--Study 展開は
\[
(stu)^\mu u_\tau^{\nu-\varkappa}v_\tau^\varkappa s_x^m t_x^{\,n-\lambda}(tuv)^\lambda
=\sum_{p,r,\varrho} C'_{pr\varrho}
\bigl[s_\tau^\varrho(stu)^{\mu+r-\varrho}\bigr]_{\widehat{uv}^{\lambda-r+\varrho}v^{\varkappa-\varrho},\,w^p,\,(w=x\widehat{uv})},
\]
となり、したがって、$u_x$ を因子にもつすべての項を省くべきことが分かっている場合に、われわれの展開で計算の途中ですでに生じる簡略化を許さない。}：
\[
\tag{I}
s_x^m t_x^{n-\lambda}t_y^\lambda
=\sum_{\varrho}
\frac{\binom{m}{\varrho}\binom{\lambda}{\varrho}}
{\binom{m+n-\varrho+1}{\varrho}}\cdot
\bigl[(st\widehat{xy})^\varrho s_x^{m-\varrho}t_x^{n-\varrho}\bigr]_{y^{\lambda-\varrho}}.
\]

この二つの展開を組み合わせ、反変的変数列を二つずつもつ形式
\[
s_x^m t_x^{n-\lambda}t_y^\lambda u_\sigma^\mu u_\tau^{\nu-\kappa}v_\tau^\kappa
\]
について、$u_x$ の冪に従う展開を作る。その係数は、形式列 $s_x^m t_x^n u_\sigma^\mu u_\tau^\nu$ の、変数 $x$ と $u$ に関して取られた合成極形式である。

記号（または変数）の二つの反変的列の場合に級数を立てれば十分である。というのも、記号（変数）の列を二つずつ一つの新しい列にまとめることによって、より多くの記号列（変数列）をもつ形式はこの場合に帰着されるからである。

形式列のすべての形式が記号列、または変数列を二つだけ含むようにするため、次の特殊化を行う：
\[
\begin{array}{ll}
\text{a) } \sigma=\widehat{st}, & \text{a') } y=\widehat{uv},\\[0.3em]
\text{b) } s=\widehat{\sigma\tau}, & \text{b') } v=\widehat{xy},
\end{array}
\]
そして特殊化 a), a') について展開を行う。

展開 I を直接転送すると、形式 $(stu)^\mu s_x^m t_x^{n-\lambda}t_y^\lambda$ については合成極形式が得られる。一方、形式 $(stu)^\mu u_\tau^{\nu-\kappa}v_\tau^\kappa s_x^m t_x^{n-\lambda}t_y^\lambda$ については、合成極形式の代わりに、評価のために次々と新しい補助変数を導入し、最後にだけ同一視しなければならない項が現れる。このため計算が不必要に複雑になる。

そこで間接的な道をとる。形式列のすべての形式の合成極形式、すなわち式
\[
\bigl[s_\tau^\varrho(stu)^{\mu-\varrho+r}\bigr]_{y^\lambda}^{v^\kappa}
\]
を、これらの極形式の初期項によって表し、その結果得られる再帰的な方程式系を反転することによって、求める極形式による展開を得る。

転送原理によれば、形式 $s_x^m t_x^n:[s_x^m t_x^n]_{y^\lambda}$ の $\lambda$ 番目の極形式（$\lambda\le n$）について次の展開が成り立つ（$\lambda>n$ の場合は $[s_y^m t_y^n]_{x^{m+n-\lambda}}$ の展開によって第一の場合に帰着される）\footnote{Gordan, \emph{Die Resultante binärer Formen} (chap. I, § 5). Rend. Circ. Mat. di Palermo XXII. 1906.}：
\[
\bigl[s_x^m t_x^n\bigr]_{y^\lambda}
=\sum_r (-1)^r\,
\frac{\binom{m}{r}\binom{\lambda}{r}}{\binom{m+n}{r}}\cdot
(st\widehat{xy})^r s_x^{m-r}t_x^{n-\lambda}t_y^{\lambda-r},
\]
同様に、$u_\sigma^\mu u_\tau^\nu:[u_\sigma^\mu u_\tau^\nu]^{v^\kappa}$ の $\kappa$ 番目の極形式について、
\[
\bigl[u_\sigma^\mu u_\tau^\nu\bigr]^{v^\kappa}
=\sum_\varrho (-1)^\varrho\,
\frac{\binom{\mu}{\varrho}\binom{\kappa}{\varrho}}{\binom{\mu+\nu}{\varrho}}\,
(\sigma\tau\widehat{uv})^\varrho u_\sigma^{\mu-\varrho}u_\tau^{\nu-\kappa}v_\tau^{\kappa-\varrho}.
\]

これらを掛け合わせ、特殊化 a) と a') を導入すると、
\[
\begin{aligned}
\bigl[s_x^m t_x^n(stu)^\mu u_\tau^\nu\bigr]_{\widehat{uv}^{\lambda}}^{v^\kappa}
={}&
\sum_{r,\varrho}c_{r\varrho}\cdot
(st\widehat{x\widehat{uv}})^r(\widehat{st}\tau\widehat{uv})^\varrho
s_x^{m-r}t_x^{n-\lambda}(tuv)^{\lambda-r}(stu)^{\mu-\varrho}
u_\tau^{\nu-\kappa}v_\tau^{\kappa-\varrho},
\end{aligned}
\]
したがって、$u_x$ の冪に従って展開し、第一項を区別し（$\sum'_{r,\varrho}$ は $r=0$, $\varrho=0$ の項を省くことを意味する）、(a) p. 26 を考慮すると、
\[
\tag{II}
\begin{aligned}
&s_x^m t_x^{n-\lambda}(tuv)^\lambda(stu)^\mu u_\tau^{\nu-\kappa}v_\tau^\kappa\\
&\quad =
\bigl[s_x^m t_x^n(stu)^\mu u_\tau^\nu\bigr]_{\widehat{uv}^{\lambda}}^{v^\kappa}\\
&\qquad
-\sum_{r,\varrho}' c_{r\varrho}\cdot s_\tau^\varrho(stu)^{\mu-\varrho+r}
u_\tau^{\nu-\kappa}v_\tau^{\kappa-\varrho}
s_x^{m-r}t_x^{n-\lambda}(tuv)^{\lambda-r+\varrho}\cdot v_x^r\\
&\qquad
+u_x\cdot\sum_{\varrho}\cdot\sum_{r=1}^{\lambda} c_{r\varrho}\cdot
s_\tau^\varrho(stu)^{\mu-\varrho+r-1}(stv)
u_\tau^{\nu-\kappa}v_\tau^{\kappa-\varrho}
s_x^{m-r}t_x^{n-\lambda}(tuv)^{\lambda-r+\varrho}\cdot v_x^{r-1}\\
&\qquad
\vdots\\
&\qquad
-(-1)^\lambda u_x^\lambda\cdot\sum_\varrho c_{\lambda\varrho}\cdot
s_\tau^\varrho(stu)^{\mu-\varrho}(stv)^\lambda
u_\tau^{\nu-\kappa}v_\tau^{\kappa-\varrho}
s_x^{m-\lambda}t_x^{n-\lambda}(tuv)^\varrho,
\end{aligned}
\]
ただし
\[
c_{r\varrho}=(-1)^{r+\varrho}\cdot
\frac{\binom{m}{r}\binom{\lambda}{r}}{\binom{m+n}{r}}\cdot
\frac{\binom{\mu}{\varrho}\binom{\kappa}{\varrho}}{\binom{\mu+\nu}{\varrho}}
\qquad
\begin{array}{l}
\text{ただし }\lambda\le n,\\
\kappa\le\nu.
\end{array}
\]
他の数の組合せ
\[\lambda,n;\varkappa,\nu.\]
についても同様である。式
\[s_x^m t_x^{n-\lambda}(tuv)^\lambda(stu)^\mu u_\tau^{\nu-\kappa}v_\tau^\kappa\]
は、したがって合成極形式へ帰着され、さらに恒等式
\[
(stv)u_\tau=(stu)v_\tau-s_\tau(tuv),
\]
を用いると、同様に作られた式の和へ帰着される。それらは形式列のより高い形式に対応する。

反復によって次の展開に至る：
\[
\tag{III}
s_x^m t_x^{n-\lambda}(tuv)^\lambda(stu)^\mu u_\tau^{\nu-\kappa}v_\tau^\kappa
=
\sum_{p,r,\varrho}u_x^p\cdot C_{pr\varrho}\cdot
\bigl[s_\tau^\varrho(stu)^{\mu-\varrho+r}\bigr]_{\widehat{uv}^{\lambda-r+\varrho}v^{\varkappa-\varrho+p}\cdot v_x^{r-p}}.
\]
数値係数 $C_{pr\varrho}$ は、(II.) を反復して得られる方程式系の係数 $c_{r\varrho}$, $c'_{r\varrho}$ から一意かつ再帰的に計算される（例 p. 42 を参照）。他の特殊化についても同様の展開が得られる。

\begin{center}
\textbf{§ 3. \quad 還元定理。}
\end{center}

形式および形式系の還元に関する既知の定理を、いま §§ 1, 2 と関連づける。そのために、二元形式の理論から取ったいくつかの定義が必要である。

a) あらかじめ与えられた形式列を法として、任意のこれらの形式の積を折り畳むことによって生じるすべての形成が、系の形式の整有理関数として表され、ただし系の形式を法の系と折り畳んで生じる式から成る加法項を除いてよい、という性質をもつ形式系を、\emph{相対的完全系}と呼ぶ\footnote{Gordan--Kerschensteiner, \emph{Vorlesungen über Invariantentheorie}. vol. II, p. 227.}。

b) 形式が、不変式を因子にもつ形式、または「より高い形式」によって表されるとき、その形式を\emph{可約}と呼ぶ。すなわち、その形式が属する全形式列のより高い形式、または法の記号をより高い位数で含む形式によって表される場合である。

還元子に関する定理\footnote{Gordan, \emph{Math. Annalen} vol. XVII, p. 222.}は、最も一般的な形で次のように述べられる：

定義：\emph{還元子とは、可約な形式列である。}

\emph{定理 II. 形式列の初期形式が可約であるのは、その一つの項が還元子との折り畳みによって生じている（還元子を因子にもつ）ためであり、さらに形式列の終形式がこの同じ項から折り畳みによって生じているならば、形式列全体は可約である}\footnote{定理 II による簡略化は次の例で見られる。特殊形式 $f=x_1^3x_2+x_2^3x_3+x_3^3x_1$ に対して、$a_\nu^2a_x^2u_\nu^2$ は還元子である。したがって定理 II により、形式列
\[
\theta_\nu(\vartheta\nu x)\theta_x^3u_\vartheta u_\nu^2
=
a_\nu^2(abu)-\underbrace{a_\nu b_\nu(abu)}_{0},
\]
の還元が従う。なぜなら $\theta_\nu^4=a_\nu^2b_\nu^2(abu)^2$ は項 $a_\nu^2(abu)$ から折り畳みによって生じているからである。Gordan, loc. cit. p. 231 では、これに対して形式
\[
\theta_\nu(\vartheta\nu x),\quad \theta_\nu(\vartheta\nu x)^2,\quad
\theta_\nu^2,\quad \theta_\nu^2(\vartheta\nu x),\quad
\theta_\nu^2(\vartheta\nu x)^2
\]
の還元が個別に行われている。}。

証明：形式列は初期項から自己への折り畳みによって生じる。したがって一方では還元子とのより高い折り畳みによって、他方では還元子の自己への折り畳みによって生じる。いずれの場合にも、§ 2 により、形式列のすべての形式は還元子の形式列上の極形式（トランスヴェクション）として表せる。

初期形式から全形式列へ結論することは、他の還元方法ではもはや成り立たない。それらは次のものである。

1) いわゆる「二重還元」。

これは、互いに独立な二つの還元子を因子にもつ式を二通りに還元することを意味し、それによって、還元先となったより高い形式の間に関係が生じる。

「二重還元」を体系的に適用することによって、一方ではすべての関係\footnote{例えば「二重還元」によって、形式 $a_\eta^3$ (§ 10) の還元が得られた。これは Maisano が不要に系に入れていた唯一の形式である。さらに、序論の注で述べた三つの共変式および三つの不変式の間の関係も得られた。}が得られなければならない。他方、さまざまな式に「二重還元」を適用しても新しい関係が得られない場合には、より高い形式の独立性の判定基準と計算の検算とを同時に得る。

2) \emph{分解可能形式による折り畳み。}

「分解可能形式」とは、通常の概念をわずかに一般化して、低い次数の形式の積および「より高い」形式によって表される形式を意味する。形式の積は、一回の折り畳みによって積に移る。同様に、非線形形式の積も、特殊化 a'), b') (§ 2) における二回の折り畳みによって、積定理
\[
a_yb_y a_xb_x=\frac12a_y^2\cdot b_x^2+\frac12b_y^2\cdot a_x^2-\frac12(ab\widehat{yx})^2
\]
に従って積に移る。
したがって、分解可能形式の第一極形式（トランスヴェクション）およびある第二極形式は、ふたたび分解可能形式である。よって次が従う。

分解可能形式との一回（または二回）の折り畳みによって生じた、分解可能形式上の一回（または二回）のトランスヴェクションの一項は、それ自身分解可能形式となる。

a) 分解可能形式の形式列における次に高い形式が分解するか可約になる場合、または一回の折り畳みで特殊化 $y=\widehat{uv}$ または $v=\widehat{xy}$ が現れる場合。

b) 他の一回の折り畳みがより高い形式へ導く場合（§ 12, B. $H_k$ および $H_k(k\eta x)$ 参照）。

そのようなトランスヴェクションの一項は、さらに次の場合にも分解し得る。

c) 他の一回の折り畳みが、分解可能形式との折り畳みとは独立に分解する低い形式へ導く場合、すなわち積および還元すべき形式によって表せる場合である。ただしその都度、該当する項の数値係数が恒等的に零でないことを計算によって確認しなければならない。これは低い形式の「二重還元」にほかならない（§ 23, C. $H_\vartheta(\theta H u)(\theta s u)$ 参照）。

分解可能形式は、関数行列式とのすべての一回の折り畳みによって生じる\footnote{Gordan, loc. cit. § 3 を参照。}ほか、ある高次の折り畳みによっても生じる。すなわち
\[
(st\widehat{yx})s_y=\frac12\,t\cdot s_y^2-\frac12\,s\cdot t_y^2+\frac12(st\widehat{yx})^2,
\]
\[
(st\widehat{yx})s_y^2=\frac13\,t\cdot s_y^3-\frac13\,s\cdot t_y^3+s_y(st\widehat{yx})^2-\frac13(st\widehat{yx})^3.
\]
双対的な形式
\[
(\sigma\tau\widehat{vu})v_\sigma
\quad\text{および}\quad
(\sigma\tau\widehat{vu})v_\sigma^2
\]
についても同様の公式が成り立つ。

この節の定理から、$S$ と $T$ の折り畳みから生じるすべての既約形式を構成するため、次の規則が得られる。

最も低い折り畳みから始め、あらかじめ定めた任意の道筋（例えば行ごと、または対角項に対して対称的）に従って、形式列 $ST$ の形成を進め、還元子を因子にもつ形式に到達するまで行う。その形式によって定められる形式列は、終形式が定理 II の条件を満たすならば、ただちに省く。すべての可約な形式列を除いた後、二重還元を用いて、残りのすべての可約形式およびすべての分解可能形式を取り除く。互いに独立に可約な形式列によって二重に覆われる全形式列の場所は、より高い形式の還元を生じさせる（§ 11, D 参照）。

\begin{center}
{\Large\bfseries 第 II 章。\quad 相対的完全基本系。}

\textbf{§ 4.}\quad\emph{最初の相対的完全系。}
\end{center}

以下では、最初の定理を除いて四次形式に限定する。しかし法を適切に選べば、§§ 5 と 6 の結果は高次形式へ容易に移せ、三次形式についても同様に成り立つ。\footnotemark
\footnotetext{三次形式について、対応する同様に導かれる公式は次のとおりである：
\[
(abc)a_x^2b_y^2c_z^2=\frac13(sxy)(sxz)(syz)+\frac23\Delta_x\Delta_y\Delta_z\cdot(xyz),\qquad (\mathrm{IV}.)
\]
\[
\begin{gathered}
s=(abc)(abu)(acu)(bcu),\qquad \Delta=(abc)^2a_xb_xc_x,\\
t=\Delta_\vartheta(\Delta\theta u)^2u_\vartheta=\theta_s^2u_su_\vartheta^2,
\qquad T=a_t^3,
\end{gathered}
\]
\[
\left.
\begin{aligned}
(a\theta u)^2&=\frac13 u_xs,\\
a_\vartheta&=\frac13u_x\cdot\Delta,
& a_\vartheta(a\theta u)&=0,
& a_\vartheta(a\theta u)^2&=s,\\
a_\vartheta^2&=\Delta,
& a_\vartheta^2(a\theta u)&=0.
\end{aligned}
\right\}\tag{1}
\]}

任意次数の形式に対する最初の相対的完全系は、次の周知の定理によって得られる。

\clearpage
\noindent\textbf{定理 III.}
\emph{二元領域から移された形式、すなわち、対応する二元形式の系の形式から Clebsch の転送原理による縁取りによって生じる形式は、法 $(abc)$ に関する相対的完全系を成す。}

\noindent\textbf{証明：}
二元の基本形式について、形式 $Q'_1,\ldots,Q'_n$ が完全系を成すことを表す二元関係
\[
Q'-F(Q')=0=\sum_\lambda\bigl\{(ab)c_x+(bc)a_x+(ca)b_x\bigr\}\varphi_\lambda\footnotemark
\]
\footnotetext{Gordan--Kerschensteiner, loc. cit. p. 134.}
には、縁取りによって三元関係
\[
Q-F(Q_i)=\sum_\lambda u_x\cdot(abc)\varphi_\lambda
\]
が対応する。これはまさに、法 $(abc)$ に関する相対的完全系の定義式である。四次形式の場合、移された形式の系は次の形式から成る：
\[
\begin{gathered}
f=a_x^4,\qquad
\theta=\theta_x^4u_\vartheta^2=(abu)^2a_x^2b_x^2,\qquad
K=K_x^6u_k^3=(a\theta u)u_\vartheta^2a_x^3\theta_x^3,\\
j=u_j^6=(a\theta u)^4u_\vartheta^2,\qquad
\nu=u_\nu^4=(abu)^4,
\end{gathered}
\]
これらのうち最初の四つは、法 $((abc),\nu)$ に関する相対的完全系を成す。

\begin{center}
\textbf{§ 5. \quad 法 $(abc)$ の法 $\A$ および $\nu$ への還元。}
\end{center}

単に括弧因子によって与えられる法 $(abc)$ は、定義 (§ 3, a) に従い、系の形式へ還元されなければならない。換言すれば、形式列 $(abc)$ は一意に正規化されなければならない（§ 1 参照）。

\[
a_s^2a_xu_s=\frac13a_s^3\cdot u_x=\frac13S\cdot u_x. \tag{2}
\]
\[
\left.
\begin{aligned}
(a\A u)^2&=a_s-\frac{S}{6}u_x^2,\\
(a\A u)^3&=0.
\end{aligned}
\right\} \tag{3}
\]
\[
\left.
\begin{aligned}
\theta(s)=(ss_1x)^2&=2a_t+\frac13S\cdot\theta-\frac13Tu_x^2,\\
\theta(t)=(tt_1x)^2&=-\frac13S\cdot a_t+\frac23T\cdot a_s+\frac1{18}S^2\cdot\theta-\frac1{18}u_x^2\cdot S\cdot T.
\end{aligned}
\right\} \tag{4}
\]

法の列：$(abc)$；$\A$；$s$；$t$（法 $t$ の系はただ一つの形式 $t$ から成る）。

次の形式
\[
\A=(abc)^2a_x^2b_x^2c_x^2,\qquad \nu=(abu)^4
\]
が法として十分であること、すなわち形式
\[
\begin{gathered}
\A=(abc)^2a_x^2b_x^2c_x^2,\qquad
a_\nu=(abc)(bcu)^3a_x^3,\\
a_\nu^2=(abc)^2(bcu)^2a_x^2,\qquad
i=a_\nu^4=(abc)^4
\end{gathered}
\]
が形式列 $(abc)$ を定めることを示す。形式列 $a_\nu$ では、恒等式
\[
a_\nu^3=(abc)^3a_x(bcu)=\frac13 u_x\cdot(abc)^4=\frac13 i\cdot u_x
\]
により、形式 $a_\nu^3$ が欠ける。

法をより高い法へ還元するためには、定義に従い、最も一般的な折り畳み、またはその法とともに考慮されるすべての特殊な折り畳みを、より高い法との折り畳みへ還元しなければならない。

本場合に最も一般的な折り畳みは、式
\[
(abc)a_x^3b_y^3c_z^3,\qquad
(abc)a_x^2b_y^2c_z^2a_xb_xc_x
\]
（三次形式から移されたもの）、
\[
(abc)a_xb_yc_za_x^2b_x^2c_x^2
\]
（二次形式から移されたもの）、およびこれらの式の極化から生じる。

これらの式を $\A$ の極形式および形式列 $a_\nu$ の極形式、あるいはそれらの初期項によって計算するため、§ 2 と同じく間接的な道をとる。極項
\[
(abc)^2a_x^2b_y^2c_z^2,\qquad
a_\nu(\nu xy)(\nu yz)(\nu zx)a_xa_ya_z
=(abc)(bc\widehat{xy})(bc\widehat{yz})(bc\widehat{zx})a_xa_ya_z\quad\text{等}
\]
を、記号因子 $(abc)$ をもつ式の線形結合として展開し、方程式系を反転することによって、求める関係、および変数の異なる組合せに従って取った極形式の間の関係を得る。評価には、行列式についての同一性定理および積定理を用いる。

$(abc)a_x^3b_y^3c^3$ に対する方程式系は次の通りである：
\begingroup\small
\[
\begin{aligned}
1)\quad
\sum_3 a_\nu(\nu yz)^3a_x^3
&=6(abc)a_x^3b_y^3c_z^3
-6\sum_3(abc)a_x^3b_y^2b_zc_z^2c_y\srcfnmark{*)},\\[0.4em]
2)\quad
3(abc)^2a_x^2b_y^2c_z^2\cdot(xyz)
-\frac12\sum_3 a_\nu^2(\nu yz)^2a_x^2\cdot(xyz)
&=2\sum_3(abc)a_x^3b_y^2b_zc_z^2c_y\\
&\quad-4\sum_3(abc)a_xa_ya_zb_y^2b_xc_z^2c_x,\\[0.4em]
3)\quad
a_\nu(\nu xy)(\nu yz)(\nu zx)a_xa_ya_z
&=2\sum_3(abc)a_xa_ya_zb_y^2b_xc_z^2c_x,\\[0.4em]
4)\quad
\sum_3 a_\nu(\nu yz)^3a_x^3
-\frac32\sum_3 a_\nu^2(\nu yz)^2a_x^2\cdot(xyz)
+\frac16 i\cdot(xyz)^3
&=6\sum_3(abc)a_xa_ya_zb_y^2b_xc_z^2c_x.
\end{aligned}
\]
\endgroup

したがって
\[
(abc)a_x^3b_y^3c_z^3
\]
について、また同様の計算により
\[
(abc)a_x^2b_y^2c_z^2a_xb_xc_x,\qquad
(abc)a_xb_yc_za_x^2b_x^2c_x^2:
\]
について、次を得る：
\[
\begin{aligned}
(abc)a_x^3b_y^3c_z^3
&=\frac32(abc)^2a_x^2b_y^2c_z^2\cdot(xyz)
+\frac32 a_\nu(\nu xy)(\nu yz)(\nu zx)a_xa_ya_z\\
&\quad-\frac1{12}i\cdot(xyz)^3,\\[0.4em]
(\text{IV.}.)\quad
(abc)a_x^2b_y^2c_z^2a_xb_xc_x
&=\frac23(abc)^2a_x^2b_yb_xc_zc_x\cdot(xyz)\\
&\quad+\frac13a_\nu(\nu xy)(\nu yz)(\nu zx)a_x^3
-\frac16 a_\nu^2(\nu xy)(\nu zx)a_x^2\cdot(xyz),\\[0.4em]
(abc)a_xb_yc_za_x^2b_x^2c_x^2
&=\frac16(abc)^2a_x^2b_x^2c_x^2\cdot(xyz).
\end{aligned}
\]

こうして、四つの形式 $f,\theta,K,j$ から成る、法 $(\A,\nu)$ に関する相対的完全系が得られた。したがって、二元領域から移されたものではない $f$ の $\theta$ 上のすべてのトランスヴェクション、および二元の場合に $(ab)^4$ へ還元可能なトランスヴェクション $(a\theta u)^2$ は、記号 $\A$ と $\nu$ によって表すことができる。\footnote{$\sum_3$ は、初期項から変数の循環置換によって得られる三項の和を指す。各方程式はその和の各項について別々にも成り立つ。}

以下で基本となる還元公式は、展開 IV の特殊化、またはより短く直接計算（記号の置換）によって得られる。すなわち、$a_\vartheta(a\theta u)^2$, $a_\vartheta^2(a\theta u)^2$, $(a\theta u)^2$ について、次の仮定から出発する：
\[
\begin{aligned}
a_\vartheta(a\theta u)^2
&=(abc)(bcu)(abu)^2-\frac13(abc)(bcu)\{(bcu)a_x-u_x\cdot(abc)\}^2,\\
a_\vartheta^2(a\theta u)^2
&=(abc)^2(abu)^2-\frac13(abc)^2\{(bcu)a_x-u_x\cdot(abc)\}^2,
\end{aligned}
\]
$(a\theta u)^2$ については\footnote{Gordan--Kerschensteiner, loc. cit. p. 92.}
\[
\begin{aligned}
(abu)a_y^3b_x^3&=\frac32u_\vartheta(\vartheta yx)\theta_y^2+\frac14(\nu yx)^3,\\
(abu)(acu)^3b_x^3&=\frac32u_\vartheta(\vartheta\widehat{au}x)(\theta au)^2+\frac14(\nu\widehat{au}x)^3.
\end{aligned}
\]
\[
\left.
\begin{aligned}
(a\theta u)^2&=\frac16\nu\cdot f-\frac23u_x\cdot a_\nu
+\frac23u_x^2\cdot a_\nu^2-\frac{i}{18}u_x^4,\\
a_\vartheta&=\frac13u_x\cdot\A,\\
a_\vartheta(a\theta u)&=0,\\
a_\vartheta^2&=\A,\\
(a\theta u)^3&=0,\\
a_\vartheta(a\theta u)^2&=-\frac56a_\nu+\frac76u_x\cdot a_\nu^2-\frac{i}{9}u_x^3,\\
a_\vartheta^2(a\theta u)&=0,\\
(a\theta u)^4&=j,\\
a_\vartheta(a\theta u)^3&=0,\\
a_\vartheta^2(a\theta u)^2&=\frac23a_\nu^2-\frac{i}{9}u_x^2.
\end{aligned}
\right\} \tag{1}
\]

これに、同じく法 $(abc)$ の還元から生じる公式
\[
a_\nu^3a_xu_\nu=\frac13 i\cdot u_x. \tag{2}
\]
を加える。

\emph{式 (1.) および (2.) と、基本形式 $\nu$ について同様に作られる公式から、以後のすべての還元公式は極化によって導かれる}。ただし、分解可能形式に関する関係を除く。式 (1.) から次が分かる。

形式列：
\[
\begin{array}{rcl}
a_\vartheta(a\theta u)^2 &\text{は}& \text{記号 }\nu\text{ のみへ導く},\\[0.2em]
a_\vartheta &\text{は}& \text{記号 }\A\text{ と }\nu\text{ へ導く},\\[0.2em]
(a\theta u)^2 &\text{は}& \text{記号 }j,\A\text{ と }\nu\text{ へ導く},\\[0.2em]
(a\theta u)^2\text{ の折り畳み I} &\text{は}& \text{記号 }j\text{ と }\nu\text{ へ導く},\\[0.2em]
(a\theta u)^2\text{ の折り畳み II} &\text{は}& \text{記号 }\A\text{ と }\nu\text{ へ導く}.
\end{array}
\]

したがって、次の置換が可能である：
\begin{enumerate}
\item 法 $(\A,\nu)$ に関して：$j$ との折り畳みを、$(a\theta u)^2$ または $(a\theta u)^3$ との対応する折り畳みに置き換える。
\item 法 $(\nu)$ に関して：$\A$ との折り畳みを、$a_\vartheta$ または $a_\vartheta(a\theta u)$ との対応する折り畳みに、または $(a\theta u)^2$ との特殊な折り畳みに置き換える（これは単なる一回の折り畳み I にしか導かない）。
\end{enumerate}
特に次が成り立つ：
\[
\begin{aligned}
(a\theta u)^2a_y^2\theta_x^2&=\frac13(jyx)^2\pmod{\nu},&
(a\theta u)^2a_y\theta_y&=-\frac16(jyx)^2\modu{\nu},\\[0.4em]
(a\theta u)^2v_\vartheta^2&=\frac13(\Delta uv)^2\modu{\nu},&
(a\theta u)(a\theta v)u_\vartheta v_\vartheta&=-\frac16(\Delta uv)^2\modu{\nu},\\[0.4em]
(a\theta u)^2a_yv_\vartheta^2&=\frac13(\Delta uv)^2\Delta_y\modu{\nu}.
\end{aligned}
\]

\begin{center}
\textbf{§ 6. \quad 法 $(\A,\nu)$ の法 $(\nu)$ への還元。}
\end{center}

前節の還元公式は、法 $\nu$ に関する相対的完全系を直接構成する手段を与える。移された形式の系には、形式 $\A$ および $(a\A u)=N$ だけを付け加えればよいことが分かる（これは三次形式の場合と同様であり、より一般にも成り立つ）。

法 $\A$ を還元するために、考慮すべきすべての特殊な折り畳み、すなわち法 $(\A,\nu)$ に関する相対的完全系と $\A$ の系との折り畳みを考え、係数における最も低い位数の形式、すなわち $f$ と $\A$ との折り畳みから出発する。

形式 $(a\A u)^2$, $(a\A u)^3$, $(a\A u)^4$ を還元するため、§ 5 に従い、記号 $\A$ を形式列の低い形式で置き換える。すなわち式
\[
a_\vartheta b_\vartheta(b\theta u)^2,\qquad
a_\vartheta(a\theta u)b_\vartheta(b\theta u)^2,\qquad
a_\vartheta(a\theta u)b_\vartheta(b\theta u)^3
\]
を、

1) 形式列 $a_\vartheta$（記号 $\A$ と $\nu$）の極形式に従って、

2) 形式列 $b_\vartheta(b\theta u)^2$（記号 $\nu$）の極形式に従って、

展開し、次の還元公式を得る（計算は下を参照）：
\begin{align}
\text{a)}\quad (a\A u)^2
&=\frac12(\vartheta\nu x)^2-\frac25 f\cdot a_\nu^2
-\frac45u_x\cdot\theta_\nu(\vartheta\nu x)^2 \notag\\
&\quad+u_x^2\left\{\frac16 i\cdot f-\frac1{40}s\right\},\notag\\
\text{b)}\quad (a\A u)^3
&=-\frac3{10}\theta_\nu(\vartheta\nu x),\tag{3}\\
\text{c)}\quad \underline{(a\Delta u)^4}
&=\frac{21}{10}\theta_\nu^2-\frac3{10}H
-\frac65u_x\cdot\theta_\nu^3+\frac1{10}u_x^2\cdot\varrho .\notag
\end{align}

（同じ結果は、$a_\vartheta^2(b\theta u)^2$, $a_\vartheta^2(b\theta u)^3$, $a_\vartheta^2(a\theta u)^2(b\theta u)^2$ および $a_\vartheta^2(a\theta u)(b\theta u)^3$ を二重還元し、$a_j^2$ を消去することによっても得られる。）

式 (3.) から、$(a\A u)^2$ は法 $\nu$ に関する還元子である。これより次が従う。

1) $\A$ の系の還元：形式列 $(\A\A' u)^2\equiv a_\vartheta^2(a\A u)^2$ は還元子 $(a\A u)^2$ を因子にもつ。

2) 法 $\A$ との折り畳みによって生じる系の還元：

形式列 $(\theta\A u)^2$ は還元子 $(a\A u)^2$ を因子にもつ。

形式 $(\theta\A u)$ と $\A_\vartheta$ については、
\[
(a\A u)^2(abu)=(\theta\A u)-u_x\cdot\A_\vartheta(\theta\A u),
\]
\[
(a\A u)(a\A b)=-\frac12\A_\vartheta+\frac12u_x\cdot\A_\vartheta^2.
\]
を得る。

しかし還元子 $(\theta\A u)$ から、法 $\A$ との折り畳みによって生じる系の還元が従う。

\bigskip

したがって、法 $\nu$ に関する相対的完全系は六つの形式
\[
f,\ \theta,\ K,\ j,\ \A,\ N
\]
から成る。

§ 2 の級数展開の例として、式 (3.)a の導出を詳しく与える。以後の計算では、最終公式 III を直接記すことにする。（零形式の極形式は計算中すでに省かれている。第一行は展開 II に従って記入した値を与え、第二行は方程式系の最後の式から始めて再帰的に求めた最終値を与える。）展開 II から次が従う：

\begin{align*}
\text{1)}\quad&m=3,\quad n=4,\quad \lambda=2,\quad \mu=0,\quad \nu=\varkappa=1,\\
a_\vartheta b_\vartheta(b\theta u)^2
&=[a_\vartheta]_{\widehat{ub^2}}b+\frac67\{a_\vartheta b_\vartheta(a\theta u)(b\theta u)-u_x\cdot a_\vartheta b_\vartheta(a\theta b)(b\theta u)\}\\
&\quad-\frac17\{a_\vartheta(a\theta u)^2b_\vartheta-2u_x\cdot a_\vartheta(a\theta u)(a\theta b)b_\vartheta+u_x^2\cdot a_\vartheta(a\theta b)^2b_\vartheta\}\\
&=\frac23(a\Delta u)^2+\frac15[a_\vartheta(a\theta u)^2]b+\frac1{30}f\cdot[a_\vartheta^2(a\theta u)^2]-\frac25u_x\cdot[a_\vartheta(a\theta u)^2]b^2\\
&\quad-\frac1{15}u_x\cdot[a_\vartheta^2(a\theta u)^2]b+\frac15u_x^2\cdot[a_\vartheta(a\theta u)^2]b^3+\frac1{30}u_x^2\cdot[a_\vartheta^2(a\theta u)^2]b^2;\\[.6em]
\text{2)}\quad&m=2,\quad n=3,\quad \lambda=1,\quad \mu=1,\quad \nu=\varkappa=1,\\
a_\vartheta(a\theta u)b_\vartheta(b\theta u)
&=\frac25\{a_\vartheta(a\theta u)^2b_\vartheta-u_x\cdot a_\vartheta(a\theta u)(a\theta b)b_\vartheta\}+\frac12a_\vartheta^2(b\theta u)^2\\
&\quad-\frac15\{a_\vartheta^2(b\theta u)(a\theta u)-u_x\cdot a_\vartheta^2(a\theta b)(b\theta u)\}\\
&=\frac12(a\Delta u)^2+\frac25[a_\vartheta(a\theta u)^2]b+\frac1{15}[a_\vartheta^2(a\theta u)^2]\cdot f\\
&\quad-\frac25u_x\cdot[a_\vartheta(a\theta u)^2]b^2-\frac1{10}u_x\cdot[a_\vartheta^2(a\theta u)^2]b+\frac1{30}u_x^2\cdot[a_\vartheta^2(a\theta u)^2]b^2;\\[.6em]
\text{3)}\quad&m=2,\quad n=3,\quad \lambda=1,\quad \mu=1,\quad \nu=1,\quad \varkappa=2,\\
a_\vartheta(a\theta b)b_\vartheta(b\theta u)
&=\frac25\{a_\vartheta(a\theta b)(a\theta u)b_\vartheta-u_x\cdot a_\vartheta(a\theta b)^2b_\vartheta\}\\
&=\frac25[a_\vartheta(a\theta u)^2]b^2+\frac1{30}[a_\vartheta^2(a\theta u)^2]b-\frac25u_x\cdot[a_\vartheta(a\theta u)^2]b^3-\frac1{30}u_x\cdot[a_\vartheta^2(a\theta u)^2]b^2;\\[.6em]
\text{4)}\quad&m=1,\quad n=2,\quad \lambda=0,\quad \mu=2,\quad \nu=\varkappa=1,\\
a_\vartheta(a\theta u)^2b_\vartheta
&=[a_\vartheta(a\theta u)^2]b+\frac23a_\vartheta^2(a\theta u)(b\theta u)\\
&=[a_\vartheta(a\theta u)^2]b+\frac16[a_\vartheta^2(a\theta u)^2]\cdot f-\frac16u_x\cdot[a_\vartheta^2(a\theta u)^2]b;\\[0.7em]
\text{5)}\quad&m=1,\quad n=2,\quad \lambda=0,\quad \mu=2,\quad \nu=1,\quad \varkappa=2\quad(\varkappa>\nu),\\
\underline{a_\vartheta(a\theta u)(a\theta b)b_\vartheta}
&=[a_\vartheta(a\theta u)^2]b^2+\frac13a_\vartheta^2(a\theta b)(b\theta u)\\
&=[a_\vartheta(a\theta u)^2]b^2+\frac1{12}[a_\vartheta^2(a\theta u)^2]b-\frac1{12}u_x\cdot[a_\vartheta^2(a\theta u)^2]b^2;\\[.6em]
\text{6)}\quad&m=1,\quad n=2,\quad \lambda=0,\quad \mu=2,\quad \nu=1,\quad \varkappa=3,\\
a_\vartheta(a\theta b)^2b_\vartheta&=[a_\vartheta(a\theta u)^2]b^3;\\[.6em]
\text{7)}\quad&m=2,\quad n=4,\quad \lambda=2,\quad \mu=\nu=\varkappa=0,\quad (\text{展開 I による}),\\
a_\vartheta^2(b\theta u)^2&=[a_\vartheta^2]_{\widehat{ub^2}}+\frac1{10}\{[a_\vartheta^2(a\theta u)^2]\cdot f-2u_x\cdot[a_\vartheta^2(a\theta u)^2]b+u_x^2\cdot[a_\vartheta^2(a\theta u)^2]b^2\};\\[.6em]
\text{8)}\quad&m=1,\quad n=3,\quad \lambda=1,\quad \mu=1,\quad \nu=\varkappa=0,\quad (\text{展開 I}),\\
a_\vartheta^2(a\theta u)(b\theta u)&=\frac14[a_\vartheta^2(a\theta u)^2]\cdot f-\frac14u_x\cdot[a_\vartheta^2(a\theta u)^2]b;\\[.6em]
\text{9)}\quad&m=1,\quad n=3,\quad \lambda=1,\quad \mu=\varkappa=1,\quad \nu=0,\quad (\text{展開 I}),\\
a_\vartheta^2(a\theta b)(b\theta u)&=\frac14[a_\vartheta^2(a\theta u)^2]b-\frac14u_x\cdot[a_\vartheta^2(a\theta u)^2]b^2.
\end{align*}

1) と 4) から次が従う：
\[
\begin{aligned}
(a\Delta u)^2
&=\frac65[a_\vartheta(a\theta u)^2]b+\frac15f\cdot[a_\vartheta^2(a\theta u)^2]+\frac35u_x\cdot[a_\vartheta(a\theta u)^2]b^2-\frac3{20}u_x\cdot[a_\vartheta^2(a\theta u)^2]b\\
&\quad-\frac3{10}u_x^2\cdot[a_\vartheta(a\theta u)^2]b^3-\frac1{20}u_x^2\cdot[a_\vartheta^2(a\theta u)^2]b^2,\\
(a\Delta u)^2
&=-a_\nu b_\nu+\frac35f\cdot a_\nu^2+\frac45u_x\cdot a_\nu^2b_\nu-\frac3{20}u_x^2\cdot a_\nu^2b_\nu^2-\frac1{12}u_x^2\cdot i\cdot f.
\end{aligned}
\]

（記号 $\theta$ への変換については § 8 および § 9 を参照。）式 (3.)a は、補助変数 $w=x\widehat{u}b$ を導入すれば、展開 I の直接転送によってさらに短く計算できる。しかし式 (3.)b および (3.)c 以降では、ここで採った道のほうが簡単になる。(3.)b については、§ 9 により、因子 $u_x$ を含むすべての項を省くべきことがあらかじめ分かっている。(3.)b と (3.)c を計算するために次を得る：
\[
\begin{aligned}
(3.)\mathrm{b)}\quad 1)&\quad a_\vartheta(a\theta u)b_\vartheta(b\theta u)^2=\frac12(a\Delta u)^3+\frac45[a_\vartheta(a\theta u)^2]_{\widehat{ub}}b+\frac{77}{360}[a_\vartheta^2(a\theta u)^2]_{\widehat{ub}},\\
2)&\quad a_\vartheta(a\theta u)b_\vartheta(b\theta u)^2=[a_\vartheta(a\theta u)^2]_{\widehat{ub}}b+\frac{13}{36}[a_\vartheta^2(a\theta u)^2]_{\widehat{ub}};\\[.5em]
(3.)\mathrm{c)}\quad 1)&\quad a_\vartheta(a\theta u)b_\vartheta(b\theta u)^3=\frac12(a\Delta u)^4+\frac65[a_\vartheta(a\theta u)^2]_{\widehat{ub^2}}b+\frac{53}{120}[a_\vartheta^2(a\theta u)^2]_{\widehat{ub^2}}\\
&\qquad-\frac65u_x\cdot[a_\vartheta(a\theta u)^2]_{\widehat{ub^2}}b^2-\frac35u_x\cdot[a_\vartheta^2(a\theta u)^2]_{\widehat{ub^2}}b+\frac{19}{120}u_x^2\cdot[a_\vartheta^2(a\theta u)^2]_{\widehat{ub^3}}b^2,\\
2)&\quad a_\vartheta(a\theta u)b_\vartheta(b\theta u)^3=\frac34[a_\vartheta^2(a\theta u)^2]_{\widehat{ub^2}}.
\end{aligned}
\]

最後の注意：§ 5 によって可約な $f$ の $j$ 上へのトランスヴェクションも同様に計算される。すなわち
\[
\tag{4}
\begin{aligned}
a_j&=-\theta_\nu+u_x\cdot\left\{\frac92\theta_\nu^2-\frac12H\right\}-3u_x^2\cdot\theta_\nu^3+\frac12u_x^3\cdot\varrho,\\
a_j^2&=\frac{21}{10}\theta_\nu^2-\frac3{10}H-2u_x\cdot\theta_\nu^3+\frac12u_x^2\cdot\varrho,\\
a_j^3&=-\frac7{10}\theta_\nu^3+\frac12u_x\cdot\varrho,\\
a_j^4&=\frac35\varrho.
\end{aligned}
\]

(3.) と (4.) から、二元領域からの転送により、
\[
(\theta\theta'u)^2=\frac13\,j\cdot f\pmod{\nu}.\footnotemark
\]
形式列 $(\theta\theta'u)^2$ の他の形式および形式 $\theta'_\vartheta$ は、記号 $\nu$ に還元可能である。したがって、次の置換が可能である：
\[
\bmod\ \nu:\quad \text{ } f\cdot j\text{ との折り畳みを}
\text{ }(\theta\theta'u)^2\text{ との対応する折り畳みで置き換える}.
\]
さらに
\[
(\vartheta\vartheta'x)^2=\frac43 f\cdot\A,
\qquad
\theta_{\vartheta'}(\vartheta\vartheta'x)=-N.
\]
\footnotetext{Gordan--Kerschensteiner, loc. cit. p. 181.}



\begin{center}
{\Large\bfseries 第 III 章。\quad 法 $(s,\varrho,t)$ に関する相対的完全系。}
\end{center}

\begin{center}
\textbf{\S\ 7.\quad 系の形成の概観。}
\end{center}

法 $\nu$ に関する相対的完全系から、次に高い法をもつ相対的完全系へ移るには、\S\ 3 の定義に従って、この高い法に関する $\nu$ の系を作り、それを法 $\nu$ に関する相対的完全系の形式と折り畳まなければならない。ところが $\nu=u_\nu^4$ は、形式 $f$ の四次反変式として $f$ と双対的に対応している。したがって $\nu$ を基本形式と見れば、法 $\nu(\nu)=(\nu\nu_1x)^4=s_x^4$ に関する六つの形式から成る相対的完全系が得られる。

したがって、法 $s$ に関する相対的完全系（系 III）を形成するためには、次をトランスヴェクションすればよい：
\[
\begin{array}{ll}
\text{系 I:} & f,\ \theta,\ K,\ j,\ \A,\ N\quad \text{を}\vphantom{\dfrac11}\\[0.3em]
\text{系 II:} & \nu,\ H=(\nu\nu_1x)^2u_\nu^2u_{\nu_1}^2,\ L=(\nu\eta x)H_x^2u_\nu^3u_\eta^3,\ g=(\nu\eta x)^4H_x^2,\\
& \sigma=H_\nu^2u_\nu^2u_\eta^4,\ (\nu\sigma x)
\end{array}
\]
とである。後で（式 (13)）示すように、形式 $g$ は還元可能であり、従って系 II は五つの形式だけから成る。

計算の途中で、二次形式
\[
\varrho=u_\varrho^2=\theta_\nu^4u_\vartheta^2,
\qquad
 t=t_x^2=a_\eta^4H_x^2
\]
を法として付け加える。そのため法 $(s,\varrho,t)$ に関する相対的完全系が得られる。ただし法 $(\varrho,t)$ は、法 $s$ より高い法とみなす。

個々の形式の配列は、係数における位数に従って行う。折り畳み
\[
\theta_\eta(K_\eta,\theta_\iota,\ldots)
\]
を折り畳み
\[
H_\vartheta(H_\kappa,L_\vartheta,\ldots)
\]
より高いものとして正規化する。これにより、\S\ 1 および \S\ 3b に従って、同じ位数をもつ個々の形式の順序は一意に定まる。

同じ位数の形式の形成は表によって実行する：

I. 係数における所定の位数を得るため、系 I と系 II のどの形式を互いに折り畳むべきかの指定。

II. 形式列の図式に従った既約形式の列挙。

III. 還元可能または分解可能な形式列あるいは形式の還元、すなわち全形式列が中断する箇所の還元。これにより、すべての既約な形成が、挙げられたものによって尽くされていることが示される。

IV. 帰結。すなわち 1) 新しく生じた還元子の指定、2) 同じ系の形式との積において、他方の系の形式との折り畳みに入らない形式の指定。

現れるものは次の場合だけであることが分かる：

1) 系 I の一つの形式と系 II の一つの形式との折り畳み。

2) $\theta$、または $H$ の冪の折り畳み、あるいは一方の系の二つの形式と第二の系の一つの形式との折り畳み。

折り畳みに関して次の図式を得る：
\[
\begin{array}{r@{\ }c@{\ }l@{\qquad}c@{\qquad}r@{\ }c@{\ }l}
\multicolumn{3}{c}{\text{系 I}} & \text{と折り畳む} & \multicolumn{3}{c}{\text{系 II}}\\[0.25em]
1. & \text{位数} & f & & 2. & \text{位数} & \nu\\
2. & ,, & \theta & & 4. & ,, & H\\
3. & ,, & K,j,\A & & 6. & ,, & L,\sigma\\
4. & ,, & N,\theta^2,f\cdot j & & 8. & ,, & (\nu\sigma x),H^2\\
5. & ,, & \theta\cdot K & & 10. & ,, & H\cdot L\\
6. & ,, & \theta^3,\A\cdot j & & 12. & ,, & H^3.
\end{array}
\]

計算の検査には、「二重還元」のほか、次の二つの特殊形式を用いる：

\[
\begin{array}{ll}
\text{I.}&\begin{array}{rlrlrl}
f&=\displaystyle\sum_3x_1^3x_2,& a_\nu^2&=\frac{i}{6}u_x^2,& s&=\frac{i}{3}f,\\
a_\eta^2&=-\frac19 i\theta,& \Delta_\nu^2&=\frac{i}{6}\theta,& \sigma&=-\frac{i}{6}j,\\
g&=-\frac23i\Delta,& \varrho&=0,& t&=0,
\end{array}\\[1.2em]
\text{II.}&\begin{array}{rlrlrl}
f&=\displaystyle\sum_3x_1^4,& s&=\frac43if,& a_\eta^2&=\frac29i\theta,\\
\Delta_\nu^2&=\frac1{15}i\theta,& \sigma&=\frac43ij,& g&=\frac43i\Delta,\\
\varrho&=0,& t&=0.
\end{array}
\end{array}
\]

後注：基本形式 $\nu$ について、式 (1.), (2.), (3.), (4.) に対応するものは次である：
\[
\tag{5.}
\begin{array}{c|c|c}
\cdot & \underline{H_\nu}=\frac13u_x\cdot\sigma & \cdot\\[0.4em]
\underline{(\nu\eta x)^2}=A & \underline{H_\nu(\nu\eta x)}=0 & \underline{H_\nu^2}=\sigma\\[0.4em]
\underline{(\nu\eta x)^3}=0 & \underline{H_\nu(\nu\eta x)^2}=B & \underline{H_\nu^2(\nu\eta x)}=0\\[0.4em]
\underline{(\nu\eta x)^4}=g & \underline{H_\nu(\nu\eta x)^3}=0 & \underline{H_\nu^2(\nu\eta x)^2}=C,
\end{array}
\]
\[
A=\frac16s\cdot\nu-\frac23u_x\cdot s_\nu+\frac23u_x^2\cdot s_\nu^2-\frac{J}{18}u_x^4,
\quad B=-\frac56s_\nu+\frac76u_x\cdot s_\nu^2-\frac{J}{9}u_x^3,
\quad C=\frac23s_\nu^2-\frac{J}{9}u_x^2.
\]
\[
\tag{6.}
s_\nu^3u_\nu s_x=\frac13u_x\cdot s_\nu^4=\frac13J\cdot u_x.
\]

\[
\tag{7}
\begin{aligned}
\text{a)}\quad (\nu\sigma x)^2
&=\frac12(Hsu)^2-\frac25\nu\cdot s_\nu^2-
\frac45u_x\cdot\,s_\eta(Hsu)^2
+u_x^2\left\{\frac16J\cdot\nu-\frac1{40}(ss'u)^4\right\},\\
\text{b)}\quad (\nu\sigma x)^3&=-\frac3{10}s_\eta(Hsu),\\
\text{c)}\quad (\nu\sigma x)^4&=\frac{21}{10}s_\eta^2-\frac3{10}Z-
\frac65u_x\cdot s_\eta^3+\frac1{10}u_x^2\cdot s_\eta^4.
\end{aligned}
\]
\[
\tag{8}
\begin{aligned}
\text{a)}\quad g_\nu&=-s_\eta+u_x\left\{\frac92s_\eta^2-\frac12Z\right\}-3u_x^2\cdot s_\eta^3+\frac12u_x^3\cdot s_\eta^4,\\
\text{b)}\quad g_\nu^2&=\frac{21}{10}s_\eta^2-\frac3{10}Z-2u_x\cdot s_\eta^3+\frac12u_x^2\cdot s_\eta^4,\\
\text{c)}\quad g_\nu^3&=-\frac7{10}s_\eta^3+\frac12u_x\cdot s_\eta^4,\\
\text{d)}\quad g_\nu^4&=\frac35s_\eta^4.
\end{aligned}
\]

\begin{center}
\textbf{\S\ 8.\quad 三次の形式（系 III）。}
\end{center}

\noindent\textbf{$f$ と $\nu$ の折り畳み。}

\noindent\emph{既約形式：}
\[
\begin{array}{cccc}
a_\nu&\cdot&\cdot&\cdot\\
\cdot&a_\nu^2&\cdot&\cdot\\
\cdot&\cdot&\cdot&\cdot\\
\cdot&\cdot&\cdot&a_\nu^4=i.
\end{array}
\]

\noindent\emph{還元：}
\[
a_\nu^3=\frac13 i u_x\qquad\text{(式 (2)).}
\]
\noindent\emph{帰結：} 1) $a_\nu^3$ は還元子である。2) $f$ は、反傾形式 $(j)$ との積においてのみ、$\nu$ との折り畳みに入る。$f$ に関する $\nu$ についても同じことが成り立つ。

\begin{center}
\textbf{\S\ 9.\quad 四次の形式（系 III）。}
\end{center}

\noindent\textbf{$\theta$ と $\nu$ の折り畳み。}

\noindent\emph{既約形式：}
\[
\begin{array}{cccc}
(\vartheta\nu x)&\theta_\nu&\cdot&\cdot\\
(\vartheta\nu x)^2&\theta_\nu(\vartheta\nu x)&\theta_\nu^2&\cdot\\
\cdot&\theta_\nu(\vartheta\nu x)^2&\cdot&\theta_\nu^3.
\end{array}
\]

\noindent\emph{還元：} 還元子 $a_\nu^3$ から、式 $a_\nu^3(abu)$, $a_\nu^3b_\nu$, $a_\nu^3b_\nu(abu)$ の二重還元を、次の設定に従って行う：

\[
\begin{aligned}
\theta_\nu^2(\vartheta\nu x)
&=a_\nu^2(\widehat{a}b\nu x)(abu)-\frac13(\nu\nu_1x)^3=a_\nu b_\nu(\widehat{a}b\nu x)(abu)+\frac16(\nu\nu_1x)^3\\
&=a_\nu^3(abu)-a_\nu^2b_\nu(abu)=2a_\nu^2b_\nu(abu)=\frac23a_\nu^3(abu),\\[.4em]
\theta_\nu^2(\vartheta\nu x)^2
&=a_\nu^2(\widehat{a}b\nu x)^2-\frac13(\nu\nu_1x)^4=a_\nu b_\nu(\widehat{a}b\nu x)^2+\frac16(\nu\nu_1x)^4\\
&=i\cdot f-2a_\nu^3b_\nu+a_\nu^2b_\nu^2-\frac13s=2a_\nu^3b_\nu-2a_\nu^2b_\nu^2+\frac16s=\frac23i\cdot f-\frac23a_\nu^3b_\nu-\frac16s,\\[.4em]
\theta_\nu^3(\vartheta\nu x)&=a_\nu^2b_\nu(\widehat{a}b\nu x)(abu)=a_\nu^3b_\nu(abu).
\end{aligned}
\]

得られる式は次である：
\[
\tag{9}
\begin{array}{c|c|c}
\underline{\theta_\nu^2(\vartheta\nu x)}=0
& \underline{\theta_\nu^3(\vartheta\nu x)}=0
&\underline{\theta_\nu^4u_\vartheta^2}=u_\varrho^2=\varrho\quad(\text{付加された法, \S\ 7}),\\[0.6em]
\underline{\theta_\nu^2(\vartheta\nu x)^2}=\dfrac49 i f-\dfrac16s
& &
\end{array}
\]

\noindent\emph{帰結：} 1) $\theta_\nu^2(\vartheta\nu x)^2$ は還元子である。2) $\nu$ は、反傾形式 $(g)$ との積においてのみ、$\theta$ との折り畳みに入る。

\begin{center}
\textbf{\S\ 10.\quad 五次の形式（系 III）。}
\end{center}

\[
\text{次の折り畳み： }\left\{\begin{array}{l}
K,j,\A \text{ と }\nu,\\
f\text{ と }H.
\end{array}\right.
\]

\noindent\emph{既約形式：}
\[
\begin{array}{c@{\qquad}c@{\qquad}c@{\qquad}c}
\text{A)} & \begin{array}{cccc}
\cdot&K_\nu&\cdot&\cdot\\
(k\nu x)^2&\cdot&K_\nu^2&\cdot\\
(k\nu x)^3&K_\nu(k\nu x)^2&\cdot&\cdot\\
\cdot&K_\nu(k\nu x)^3&\cdot&\cdot
\end{array}
&
\text{B)}\ \begin{array}{c}(j\nu x)\\(j\nu x)^2\\(j\nu x)^3\\ \cdot\end{array}
&
\text{C)}\ \begin{array}{c}\A_\nu\\\cdot\\\cdot\\\cdot\end{array}
\end{array}
\]
\[
\text{D)}\qquad
\begin{array}{cccc}
(aHu)&(aHu)^2&\cdot&\\
a_\eta&a_\eta(aHu)&a_\eta(aHu)^2&\\
\cdot&a_\eta^2&\cdot&\cdot
\end{array}
\]

\noindent\emph{還元：}

\noindent A) 形式列 $K_\nu^3$: 還元子 $a_\nu^3$。

形式 $K_\nu^2(k\nu x)^2$: 還元子 $\theta_\nu^2(\vartheta\nu x)^2$。
\[
(k\nu x),\quad K_\nu(k\nu x),\quad K_\nu^2(k\nu x)
\]
は \S\ 3 による分解可能形式である。

\noindent B) 形式
\[
\begin{aligned}
(j\nu x)^4&\equiv(\widehat{a\theta}\nu x)^4
=a_\nu^4\cdot\theta+\theta_\nu^4\cdot f-4a_\nu^3\theta_\nu
-4a_\nu\{a_\nu\cdot\theta_x-(\widehat{a\theta}\nu x)\}^3\\
&\quad+6a_\nu^2\{a_\nu\theta_x-(\widehat{a\theta}\nu x)\}^2
=\R\cdot f+\frac53 i\theta-6a_\nu^2(\widehat{a\theta}\nu x)^2
+4a_\nu(\widehat{a\theta}\nu x)^3,
\end{aligned}
\]
従って
\[
(j\nu x)^4=\R\cdot f+\frac53 i\theta\cdot\modu{(\A,\nu)},\qquad (\text{\S\ 5 の末尾参照}).
\]

\noindent C) 形式 $\A_\nu^4$: 還元子 $\theta_\nu^4$。

形式 $\A_\nu^2$ および $\A_\nu^3$: 形式 $\A$ を形式列の低い形式（\S\ 5 の末尾）で置き換えることにより、すなわち式
\[
a_\vartheta a_\nu^3,
\qquad a_\vartheta a_\nu^3\theta_\nu,
\qquad a_\vartheta\theta_\nu^4
\]
の二重還元により、還元子 $a_\nu^3$ または $\theta_\nu^4$ をもつ。$\A_\nu^3$ の二通りの計算は、高次形式 $a_\eta^3$ の還元を引き起こす。

\noindent\emph{還元公式：}
\[
\tag{10}
\begin{aligned}
\text{a)}\quad \A_\nu^2
&=\frac35a_\eta^2-\frac3{10}(asu)^2+\frac13 i\theta
+\frac15u_x^2\cdot(2a_\R^2-t),\\
\text{b)}\quad \A_\nu^3
&=-\frac3{10}a_\R+\frac35u_x\left\{\frac{13}{10}a_\R^2-\frac25t\right\},\\
\text{c)}\quad \A_\nu^4&=a_\R^2-\frac25t.
\end{aligned}
\]

\noindent\emph{公式の導出：}
\[
\begin{aligned}
\text{a)}\quad
 a_\vartheta a_\nu^3
&=\frac13 i\theta
=[a_\vartheta]_{\nu^3}+\frac67[a_\vartheta^2]_{\nu^2}
+\frac65[a_\vartheta(a\theta u)^2]_{\nu}\,\widehat{\nu x}^2
+\frac{11}{30}[a_\vartheta^2(a\theta u)^2]\,\widehat{\nu x}^2\\
&\quad-\frac67u_x\cdot[a_\vartheta^2]_\nu^3
-\frac16u_x\cdot[a_\vartheta^2(a\theta u)^2]_{\nu}\,\widehat{\nu x}^2,\\
\frac13 i\theta
&=\A_\nu^2-\frac23u_x\cdot\A_\nu^3-a_\nu a_{\nu_1}(\nu\nu_1x)^2
+\frac25a_\nu^2(\nu\nu_1x)^2
+\frac15u_x\cdot a_\nu^2a_{\nu_1}(\nu\nu_1x)^2;
\end{aligned}
\]
\[
\begin{aligned}
\text{b)}\quad
 a_\vartheta a_\nu^3\theta_\nu
&=0=[a_\vartheta]_{\nu^4}+\frac9{14}[a_\vartheta^2]_{\nu^3}
+\frac35[a_\vartheta(a\theta u)^2]_{\nu^2}\,\widehat{\nu}x^2
+\frac{17}{120}[a_\vartheta^2(a\theta u)^2]_{\nu}\,\widehat{\nu}x^2\\
&\quad-\frac9{14}u_x\cdot[a_\vartheta^2]_{\nu^4}
-\frac1{24}u_x\cdot[a_\vartheta^2(a\theta u)^2]_{\nu^2}\,\widehat{\nu x}^2,\\
0&=\frac56\A_\nu^3-\frac12u_x\cdot\A_\nu^4-\frac14a_\eta^3+\frac1{20}u_x\cdot a_\eta^4;
\end{aligned}
\]
\[
\begin{aligned}
\text{b')}\quad
 a_\vartheta\theta_\nu^4
&=a_\R=a_\vartheta\theta_\nu\{a_\nu-(a\theta\nu x)\}^3\\
&=-\frac32[a_\vartheta]_{\nu^3}+\frac35[a_\vartheta(a\theta u)^2]_{\nu^2}\,\widehat{\nu}x^2
-\frac{37}{120}[a_\vartheta^2(a\theta u)^2]_{\nu}\,\widehat{\nu}x^2\\
&\quad+\frac32u_x\cdot[a_\vartheta^2]_{\nu^4}
+\frac{49}{120}u_x\cdot[a_\vartheta^2(a\theta u)^2]_{\nu^2}\,\widehat{\nu x}^2,\\
a_\R&=-\frac32\A_\nu^3+\frac32u_x\cdot\A_\nu^4-\frac{11}{20}a_\eta^3+\frac7{20}u_x\cdot a_\eta^4;
\end{aligned}
\]
\[
\begin{aligned}
\text{c)}\quad
 a_\vartheta^2\theta_\nu^4
&=a_\R^2=[a_\vartheta^2]_{\nu^4}+\frac35[a_\vartheta^2(a\theta u)^2]_{\nu^2}\,\widehat{\nu}x^2,\\
a_\R^2&=\A_\nu^4+\frac25a_\eta^4.
\end{aligned}
\]

\noindent D) $\A_\nu^3$ の二重還元 (C) による $a_\eta^3$ の還元。

残りの還元は、\S\ 9 の計算と双対的に反対の計算によって得られる。

\noindent\emph{還元公式：}
\[
\tag{11}
\begin{array}{rl|rl}
a_\eta^2(aHu)&=-\dfrac16(asu)^3,&
 a_\eta^2(aHu)^2&=\dfrac49 i\nu-\dfrac16(asu)^4,\\[0.7em]
a_\eta^3&=-a_\R+\dfrac35u_x\cdot a_\R^2+\dfrac15u_x\cdot t,
& a_\eta^3(aHu)&=0,\\[0.7em]
&&a_\eta^4&=t\quad(\text{法として付け加えた}).
\end{array}
\]

\noindent\emph{帰結：}

1) $(j\nu x)^4$, $\A_\nu^2$, $a_\eta^2(aHu)$ は還元子である。

2) $\nu$ は、$K,j,\A$ との折り畳みにおいて、反傾形式との積にだけ入る。$H$ に関する $f$、および $\nu$ に関する $j$ についても同じことが成り立つ。なぜなら $\theta_\nu(\vartheta\cdot j,\nu,x^3)$ は \S\ 11 B により還元可能になるからである。

\begin{center}
\textbf{\S\ 11.\quad 六次の形式（系 III）。}
\end{center}
\[
\text{次の折り畳み： }\left\{\begin{array}{l}
\theta^2,
 f\cdot j,
 N \text{ と }\nu,\\
\theta\text{ と }H.
\end{array}\right.
\]

\noindent\emph{既約形式：}
\[
\begin{array}{c@{\quad}c@{\quad}c}
\text{A)}\ (\vartheta^2\nu x)^3,\quad \theta_\nu(\vartheta^2\nu x)^3
&
\text{B)}\ a_\nu(j\nu x),\quad a_\nu^2(j\nu x)
&
\text{C)}\ N_\nu
\end{array}
\]
\[
\text{D)}\quad
\begin{array}{c|c|c|c}
 & (\theta Hu)&(\theta Hu)^2&\cdot\\
(\vartheta\eta x)&\theta_\eta&H_\vartheta(\theta Hu),\ \theta_\eta(\theta Hu)&\theta_\eta(\theta Hu)^2\\
(\vartheta\eta x)^2&H_\vartheta(\vartheta\eta x),\ \theta_\eta(\vartheta\eta x)&\theta_\eta^2&\theta_\eta^2(\theta Hu)\\
\cdot&\theta_\eta(\vartheta\eta x)^2&\cdot&\cdot
\end{array}
\]

\noindent\emph{還元：}

\noindent A) $(\vartheta^2\nu x)^4$ は式 (3)a に従って分解する：
\[
(a\A\widehat{\vartheta}x)^2=\frac12(\vartheta^2\nu x)^4-\frac35 f\cdot a_\nu^2(\nu\vartheta x)^2
=\frac13\A^2+f\cdot\A_\vartheta^2.
\]

\noindent B) $a_\nu(j\nu x)^2$ は式 (9)a に従って分解する。ここでは \S\ 6 の末尾に従い $f\cdot j$ を $(\theta\theta'u)^2$ で置き換え、$\theta_\vartheta'$ の還元可能性に注意する：
\[
\frac13a_\nu(j\nu x)^2=(\widehat{\theta\theta'}\nu x)^2\theta_\nu
=\theta_\nu^3\cdot\theta-\theta_\nu^2\theta_\nu'
=\theta_\nu^3\cdot\theta+\theta_\nu^2(\widehat{\vartheta\nu}\theta'u)
=\theta_\nu^3\cdot\theta-\frac23\theta_\nu^3\cdot\theta.
\]
形式 $a_\nu(j\nu x)^3$ および $a_\nu^2(j\nu x)^2$：還元子 $a_j$ および $(j\nu x)^4$：
\[
a_\nu(j\nu x)^3=a_j(j\nu x)^3-(j\nu x)^3(j\nu\widehat{a}u),
\]
\[
a_\nu^2(j\nu x)^2=a_j^2(j\nu x)^2-2a_j(j\nu x)^2(j\nu\widehat{a}u)+(j\nu x)^2(j\nu\widehat{a}u)^2.
\]

\noindent C) 形式列 $N_\nu^2$；還元子 $\A_\nu^2$。$(n\nu x)$ および $N_\nu(n\nu x)$ は、\S\ 3 に従って関数行列式に関するトランスヴェクションとして分解する。

\noindent D) \emph{還元の概観：}

a) 形式列 $\theta_\eta^2H_\vartheta$: 還元子 $a_\eta^2(aHu)$。（高次記号をもつ形式へ導く。）

b) 形式列 $\theta_\eta H_\vartheta$: 二重還元による還元子 $a_\eta^2(aHu)$。（高次記号をもつ形式、および全形式列の高次形式、すなわち形式列 $\theta_\eta^2$ へ導く。）$\theta_\eta H_\vartheta$ の代わりに形式列
\[
\theta_\eta(\vartheta\eta x)(\theta Hu)=\theta_\eta H_\vartheta+\theta_\eta^2
\]
を考え、その中で記号 $\theta$ を低い形式 $(abu)$ で置き換え、こうして形式列
\[
a_\eta^2(aHu)(abu)=\theta_\eta H_\vartheta+\frac32\theta_\eta^2\modu{s}
\]
を、終端形式
\[
a_\eta^3b_\eta(bHu)(abu)=\theta_\eta^3H_\vartheta+\frac12\theta_\eta^4
\]
まで二重還元する。

c) 形式列 $\theta_\eta^3$: 形式列 $\theta_\eta H_\vartheta$ と $\theta_\eta^2H_\vartheta$ の両方に属する形式、すなわち形式列 $\theta_\eta^2H_\vartheta$ の形式を二重還元する。一方では高次記号をもつ形式へ、他方では高次記号をもつ形式と高次形式列 $\theta_\eta^3$ へ還元する。

d) 形式 $\theta_\eta^2(\vartheta\eta x)$, $\theta_\eta^2(\vartheta\eta x)^2$, $\theta_\eta^3(\vartheta\eta x)$: \S\ 9 の計算に類似して、還元子 $a$ による二重還元。

e) 形式 $\theta_\eta^2(\theta Hu)^2$ および $\theta_\eta^2(\vartheta\eta x)^2$: 還元子 $\A_\nu^2$ および $(a\A u)^4$ を用いる式 $\A_\nu^2(a\A u)^4$ および $(a\A\nu x)^4$ の二重還元。

f) 形式 $H_\vartheta$ および $H_\vartheta^2$: 分解可能形式。$H_\vartheta^2$ については、式 $\A_\nu^2(a\A u)^2$ の二重還元、または系の形式による積 $(a_\nu^2)^2$, $a_\nu^2\cdot b_{\nu_1}$ の直接計算から得られる。（$(a_\nu)^2$ の類似計算は分解可能形式間の関係を与える。）

g) 形式列 $\theta\cdot H$ の形式が二つの還元可能性を許す（すなわち二つの還元可能な形式列に属する）ことによる高次形式の還元。すなわち：
\begin{quote}\small
$g$: $\theta_\eta^2(\vartheta\eta x)^2$ の二重還元による。還元 d) と e) を許す。\par
$s_\vartheta^2(\theta su)$: $\theta_\eta^3(\vartheta\eta x)$ の二重還元による。還元 c) と d) を許す。\par
$s_\vartheta^2(\theta su)^2$: $\theta_\eta^2H_\vartheta^2$ の二重還元による。還元 a) を許し、還元子 $\theta_\nu^2(\vartheta\nu x)^2$ を因子にもつ。\par
$s_\nu^2$: $\theta_\eta^3H_\vartheta$ の二重還元による。還元 a) と c) を許す。
\end{quote}
残りの形式の独立性の証明は、形式列 $\A_\nu^2(a\A u)^2\equiv\theta_\eta^2$ の二重還元によって得られる。

\noindent\emph{還元の実行：}

還元 a), b), c), d) の存在はアプリオリに明らかである。なぜなら、示された形成法則に従えば、二重還元の際に生じる関係は互いに独立だからである。還元 e), f), g) については、恒等式が生じないことを計算により示さなければならない。

形式列 $\theta\cdot H$ の対角項に対して対称的に、形式 $\theta_\eta$ まで進め、計算の一部を示しながら、還元 a), b), c), d) から得られる公式も与える。これらは、一部は還元 e), f), g) に、一部は後に用いられるからである。

\noindent 1) $H_\vartheta$ および $H_\vartheta^2$、f) による。設定\footnote{等号の代わりに記号 $\sim$ を用いることは、$u_x$ と、折り畳み III および IV によって生じる形式より高次の形式との積を省略していることを意味する。}
\[
\begin{aligned}
a_\nu\cdot b_{\nu_1}^2
&=\nu\cdot a_\nu b_\nu^2-(aHu)(bHu)b_\eta
\sim \nu\cdot a_\nu b_\nu^2-f\cdot a_\eta(aHu)^2-(abu)(aHu)b_\eta+(abu)^2b_\eta,\\
a_\nu^2\cdot b_{\nu_1}^2
&=b_\nu^2\{a_\nu u_\nu-(a\nu\nu_1u)\}^2
\sim \nu\{a_\nu^2b_\nu^2-2a_\eta b_\eta(aHu)(bHu)
-\frac16(asu)^2(bsu)^2\}\\
&\sim \nu\cdot a_\nu^2b_\nu^2-2a^2(aHu)(abu)-\frac16(asu)^2(bsu)^2+(\vartheta\eta x)^2(\theta Hu)^2
\end{aligned}
\]
において、$\theta_\eta H_\vartheta$ の値を代入すると、次の公式が生じる：
\[
\tag{12}
\begin{aligned}
\text{a)}\quad H_\vartheta&\sim 2a_\nu\cdot a_\nu^2+2\nu b_\nu(\vartheta\nu x)^2+2f\cdot a_\eta(aHu)^2-2\theta_\eta,\\
\text{b)}\quad H_\vartheta^2&\sim (a^2)^2+2\theta_\eta^2+\frac23(\theta su)^2
+\frac1{12}s\cdot\nu-\frac19if\cdot\nu-\frac16f\cdot(asu)^4.
\end{aligned}
\]

\noindent 2) $\theta_\eta H_\vartheta$、b) による：
\[
\begin{aligned}
a_\eta^2(aHu)(abu)&\sim [a_\eta^2(aHu)]_{bu}+\frac12 f\cdot[a_\eta^2(aHu)^2]
=a_\eta b_\eta(aHu)(abu)\\
&\quad+a_\eta(a\widehat{b}\eta x)(abu)(aHu)
\sim \theta_\eta(\vartheta\eta x)(\theta Hu)+\frac12\theta_\eta^2+\frac14(\nu\eta x)^2.
\end{aligned}
\]
式 (5) および (11) から次が従う：
\[
\theta_\eta H_\vartheta\sim-\frac32\theta_\eta^2-\frac14(\theta su)^2-\frac1{12}s\cdot\nu+\frac29 i f\cdot\nu.
\]

\noindent 3) $\theta_\eta H_\vartheta(\theta Hu)$ は、$\theta_\eta H_\vartheta$ と同様に b) に従って計算される：
\[
\theta_\eta H_\vartheta(\theta Hu)\sim-\theta_\eta^2(\theta Hu)-\frac16(\theta su)^3.
\]

\noindent 4) $\theta_\eta H_\vartheta(\vartheta\eta x)$ は b) による。$\theta_\eta^2(\vartheta\eta x)$ は d) による（\S\ 9 参照）：
\[
\theta_\eta H_\vartheta(\vartheta\eta x)\sim-\theta_\eta^2(\vartheta\eta x)-\frac16s_\vartheta(\theta su)
\sim\frac13(\vartheta\R x)-\frac16s_\vartheta(\theta su),
\]
\[
\theta_\eta^2(\vartheta\eta x)\sim\frac23a_\eta^3(abu)\sim-\frac13(\vartheta\R x).
\]

\noindent 5)
\[
\begin{array}{ccc}
\theta_\eta H_\vartheta^2\text{ は b) による},&
\theta_\eta^2H_\vartheta\text{ は a) による},&
\theta_\eta^2H_\vartheta\text{ は b) による、従って }\theta_\eta^3\text{ は c) による},\\
\text{ }a_\eta^2(aHb)(bHu)\text{ から};&
\text{ }a_\eta^2(Hab)(abu)\text{ から};&
\text{ }a_\eta^3(bHu)(abu)\text{ から}.
\end{array}
\]
\[
\theta_\eta H_\vartheta^2\sim-\frac32\theta_\eta^2H_\vartheta-\frac14s_\vartheta(\theta su)^2+\frac16s_\nu-\frac49ia_\nu
\sim-\theta_\R-\frac16s_\nu-\frac19ia_\nu,
\]
\[
\theta_\eta^2H_\vartheta\sim\frac23\theta_\R-\frac16s_\vartheta(\theta su)^2+\frac29s_\nu-\frac29ia_\nu,
\]
\[
\theta_\eta^2H_\vartheta\sim\frac23\theta_\R-\frac23\theta_\eta^3+\frac5{36}s_\nu,
\qquad
\theta_\eta^3\sim\frac14s_\vartheta(\theta su)^2-\frac18s_\nu+\frac13i\cdot a_\nu.
\]

\noindent 6) $\theta_\eta^2(\theta Hu)^2$、e) による：
\[
\A_\nu^2(a\A u)^4=[(a\A u)^4]_{\nu^2}
=-\frac{12}{5}\theta_\eta^2(\theta Hu)^2-\frac3{10}\sigma-\frac1{20}(\theta su)^4+\nu\R
\quad(\text{式 (3)c による}),
\]
\[
\A_\nu^2(a\A u)^4=[\A_\nu^2]^4_{\widehat{u a}}=\frac35\theta_\eta^2(\theta Hu)^2+\frac15\sigma-\frac3{10}(\theta su)^4+\frac13ij
\quad(\text{式 (10)a による}),
\]
\[
\theta_\eta^2(\theta Hu)^2=-\frac16\sigma+\frac1{12}(\theta su)^4-\frac19ij+\frac13\nu\cdot\varrho.
\]

\noindent 7) $\theta_\eta^2(\vartheta\eta x)^2$、d) および e) による。これから高次形式 $g$ の還元が得られる。

\S\ 9 と同様の計算により、次が従う：
\[
\begin{aligned}
\theta_\eta^2(\vartheta\eta x)^2
&=\frac12 i f-\frac12a_\eta^2b_\eta-\frac1{12}g
=\frac23t\cdot f-\frac23a_\eta^3b_\eta-\frac16g\\
&=-\frac16g-\frac13(\vartheta\R x)^2+\frac4{15}f\cdot a_\R^2+\frac8{15}f\cdot t
\quad(\text{式 (11) による}).
\end{aligned}
\]
上と対応する計算により、次が従う：
\[
((a\A\nu x)^4)=[(a\A u)^4]_{\nu x^4}
=\frac{21}{10}\theta_\eta^2(\vartheta\eta x)^2+\frac7{10}s_\vartheta^2-\frac3{10}g
\quad(\text{式 (3)c による}),
\]
\[
\begin{aligned}
(a\A\widehat{\nu x})^4&=-\frac13i\A+6[\A_\nu^2]a^2-4[\A_\nu^3]a+\A_\nu^4\cdot f\\
&= -\frac{36}{5}\theta_\eta^2(\vartheta\eta x)^2-\frac95s_\vartheta^2-\frac35g-\frac35(\vartheta\R x)^2\\
&\quad+\frac53i\A+\frac{37}{25}f\cdot a_\varrho^2+\frac{74}{25}f\cdot t,\quad(\text{式 (10) による}).
\end{aligned}
\]
したがって
\[
\begin{aligned}
\frac{93}{10}\theta_\eta^2(\vartheta\eta x)^2
&=-\frac3{10}g-\frac52s_\vartheta^2-\frac35(\vartheta\R x)^2\\
&\quad+\frac53i\A+\frac{37}{25}f\cdot a_\varrho^2+\frac{74}{25}f\cdot t,
\end{aligned}
\]
となり、g) に従って $\theta_\eta^2(\vartheta\eta x)^2$ の二つの値を比較すると：
\[
\tag{13\srcfnmark{*)}}
0=g-2s_\vartheta^2+2(\vartheta\varrho x)^2+\frac43i\A-\frac45f\cdot a_\varrho^2-\frac85f\cdot t.
\]
\srcfntext{*)}{これは、序論の注で言及した、位数 6 かつ次数 6 の三つの共変式の間の関係である。}

\noindent 8) $\theta_\eta^2H_\vartheta(\theta Hu)$ は a) および b) による。これから $\theta_\eta^3(\theta Hu)$ は c) による（因子 $u_x$ をもつ形式は消える）：
\[
\theta_\eta^2H_\vartheta(\theta Hu)=-\frac16s_\vartheta(\theta su)^3-\frac13(\nu\R x),
\]
\[
\theta_\eta^2H_\vartheta(\theta Hu)=-\frac13\theta_\eta^3(\theta Hu)-\frac13(\nu\R x),
\]
\[
\theta_\eta^3(\theta Hu)=\frac12s_\vartheta(\theta su)^3.
\]

\noindent 9) $\theta_\eta^2H_\vartheta(\vartheta\eta x)$ は a) および b) による。これから $\theta_\eta^3(\vartheta\eta x)$ は c) による。また $\theta_\eta^3(\vartheta\eta x)$ は d) による。これから高次形式 $s_\vartheta^2(\theta su)$ の還元が g) に従って得られる（因子 $u_x$ をもつ形式は消える）：
\[
\theta_\eta^2H_\vartheta(\vartheta\eta x)=\frac13(atu),
\]
\[
\theta_\eta^2H_\vartheta(\vartheta\eta x)=\frac13(atu)+\frac13\theta_\eta^3(\vartheta\eta x)-\frac16s_\vartheta^2(\theta su),
\]
\[
\theta_\eta^3(\vartheta\eta x)=\frac12s_\vartheta^2(\theta su),
\]
\[
\theta_\eta^3(\vartheta\eta x)=\frac25\theta_\R(\vartheta\R x)+\frac15(atu),
\]
\[
\tag{14}
s_\vartheta^2(\theta su)=\frac45\theta_\R(\vartheta\R x)+\frac25(atu).
\]

\noindent 10) $\theta_\eta^2H_\vartheta^2$ は a) によるものと、還元子 $\theta_\nu^2(\vartheta\nu x)^2$ によるもの；$\theta_\eta^3H_\vartheta$ は a) と c) による；$\theta_\eta^4$ は c) による。ここから高次形式 $s_\vartheta^2(\theta su)^2$ および $s_\nu^2$ の還元が g) に従って得られる：
\[
\begin{aligned}
\text{a) によって}\quad
\theta_n^2H_\vartheta^2&=[a_\eta^2(aHu)^2]b^2+\frac13[a_\eta^4]_{\widehat{bu^2}}-\frac13H_\nu^2(\nu\eta x)^2\\
&=\frac49i\cdot a_\nu^2-\frac16s_\vartheta^2(\theta su)^2-\frac5{18}s_\nu^2+\frac13(atu)^2+\frac1{54}J\cdot u_x^2.
\end{aligned}
\]
\[
\begin{aligned}
\theta_\nu^2(\vartheta\nu x)^2\text{ によって}:
\quad\theta_\eta^2H_\vartheta^2
&=[\theta_\nu^2(\vartheta\nu x)^2]_{\nu_1^2}+\frac13[\theta_\nu^4]\widehat{\nu_1}x^2-\frac13s_\vartheta^2(\theta su)^2\\
&=\frac49i\cdot a_\nu^2-\frac16s_\nu^2-\frac13s_\vartheta^2(\theta su)^2+\frac13(\nu\R x)^2.
\end{aligned}
\]
\[
\begin{aligned}
\text{a) によって}\quad
\theta_\eta^3H_\vartheta&=-[a_\eta^2(aHu)]_{\widehat{bu}}b^2-\frac12[a_\eta^2(aHu)^2]b^2-\frac13[a_\eta^3]_{\widehat{bu^2}}b+\frac1{18}[a_\eta^4]_{\widehat{bu^2}}\\
&\quad+\frac12u_x\cdot[a_\eta^2(aHu)^2]b^3\\
&=-\frac29ia_\nu^2+\frac1{12}s_\nu^2+\frac4{15}\theta_\varrho^2-\frac7{90}(\nu\varrho x)^2+\frac1{30}(atu)^2+\frac2{27}i^2\cdot u_x^2-\frac1{36}J\cdot u_x^2.
\end{aligned}
\]
\[
\begin{aligned}
\text{c) によって}\quad
\theta_\eta^3H_\vartheta:
 a_\eta^3(Hab)(bHu)&=\frac12(atu)^2
=\theta_\eta^2H_\vartheta(\theta Hu)(\vartheta\eta x)+\frac14H_\nu^2(\nu\eta x)^2\\
&\quad+\frac12\theta_\eta^2H_\vartheta^2
=\frac32\theta_\eta^2H_\vartheta^2+\theta_\eta^3H_\vartheta-u_x\cdot[a_\eta^2(aHu)^2]b^3\\
&\quad+\frac14H_\nu^2(\nu\eta x)^2,
\end{aligned}
\]
従って
\[
\begin{aligned}
\theta_\eta^3H_\vartheta&=-\frac23i\cdot a_\nu^2+\frac5{12}s_\nu^2+\frac12(\nu\R x)^2-\frac12(atu)^2\\
&\quad+\frac4{27}i^2\cdot u_x^2-\frac1{12}J\cdot u_x^2.
\end{aligned}
\]
c) によって
\[
\begin{aligned}
\theta_\eta^4&=\frac49ia_\nu^2-\frac16s_\nu^2+\frac{16}{15}\theta_\R^2\\
&\quad-\frac{14}{45}(\nu\R x)^2+\frac2{15}(atu)^2.
\end{aligned}
\]
g) により、$\theta_\eta^2H_\vartheta^2$ の二つの値から：
\[
\tag{15}
\begin{aligned}
s_\vartheta^2(\theta su)^2
&=\frac23s_\nu^2-2(atu)^2+2(\nu\R x)^2-\frac19Ju_x^2\\
&=\frac89ia_\nu^2+\frac8{15}\theta_\R^2-\frac{14}{15}(atu)^2+\frac{38}{45}(\nu\R x)^2-\frac4{27}i^2u_x^2.
\end{aligned}
\]
g) により、$\theta_\eta^3H_\vartheta$ の二つの値から：
\[
\tag{16}
s_\nu^2=\frac43ia_\nu^2+\frac45\theta_\R^2+\frac85(atu)^2-\frac{26}{15}(\nu\R x)^2+\frac16Ju_x^2-\frac29i^2u_x^2.
\]
式 (16) は法 $(ss'u)^4$ の還元の基礎を与え、式 (13) は $s$ の系の還元の基礎を与える（\S\ 17）。

\noindent\emph{帰結：}

1) $a_\nu(j\nu x)^3$, $\theta_\eta H_\vartheta$, $\theta_\eta^2(\vartheta\eta x)$, $\theta_\eta^2(\theta Hu)^2$ は還元子であり、$a_\nu(j\nu x)^2$, $H_\vartheta$, $H_\vartheta^2$ は分解可能形式である。

2) 系 II の形式 $j(\nu)=g=(\nu\eta x)^4H_x^2$ は還元可能である。

3) 唯一の反傾形式 $g$ が還元可能になるので、$\nu$ は、系 I との折り畳みにおいて、冪にも、同じ系の形式との積にも、まったく入らない。

$\theta$ は、反変式と見たときにのみ、積または冪として折り畳みに入る。対応して $H$ は、共変式と見たときにのみ入る（\S\ 13 B 参照）。




\begin{center}
\textbf{\S\ 12.\quad 七次の形式。}
\end{center}
\[
\text{次の折り畳み：}\left\{\begin{array}{l}
\theta\cdot K\text{ と }\nu,\\
K,j,\A\text{ と }H,\\
f\text{ と }L,\sigma.
\end{array}\right.
\]

\noindent\emph{既約形式：}
\[
\text{B)}\quad
\begin{array}{c|c|c|c|c}
\cdot&\cdot&(KHu)^2&\cdot&\cdot\\
\cdot&K_\eta&\cdot&K_\eta(KHu)^2&\cdot\\
(k\eta x)^2&\cdot&K_\eta^2&\cdot&\cdot\\
(k\eta x)^3&H_k(k\eta x)^2,\ K_\eta(k\eta x)^2&\cdot&\cdot&\cdot\\
\cdot&K_\eta(k\eta x)^3&\cdot&\cdot&\cdot
\end{array}
\]
\[
\text{C)}\quad
\begin{array}{cc}
(j\eta x)&H_j\\
\cdot&H_j(j\eta x)\\
\cdot&H_j^2
\end{array}
\qquad
\text{D)}\quad
\begin{array}{c}
(\A Hu)\\\A_\eta
\end{array}
\]
\[
\text{E)}\quad
\begin{array}{cccc}
a_i&(aLu)^2&(aLu)^3&\cdot\\
\cdot&\cdot&a_l(aLu)^2&a_l(aLu)^3\\
\cdot&a_l^2&\cdot&\cdot
\end{array}
\]

\noindent\emph{還元：}

\noindent A) $(\vartheta\cdot k,\nu,x)^4$ は、分解可能形式 $(\vartheta^2\nu x)^4$ に関する一回のトランスヴェクションとして分解する。

\noindent B) 形式列 $H_kK_\eta$：還元子 $H_\vartheta\theta_\eta$。

形式列 $K_\eta^2(KHu)$：還元子 $a_\eta^2(aHu)$。

形式列 $K_\eta^2(k\eta x)$：還元子 $\theta_\eta^2(\vartheta\eta x)$。

ここから生じる形式列 $K_\eta^3$ の二重還元は、より高い形式列 $(j\eta x)^3$ の還元をもたらす。

$(KHu)$、$(k\eta x)$、$K_\eta(k\eta x)$ は \S\ 3 による分解可能形式である。

$H_k(KHu)$、$K_\eta(KHu)$ は、分解可能形式 $K_\nu(k\nu x)$ および関数行列式
\[
K_\nu^2\equiv\theta_\nu^2(a\theta u)\equiv a_\nu^2(a\theta u)\modu{\nu^2}
\]
との一回の折り畳みによって生じるものとして分解する。
$K_\nu^2$ の二重表示に対応して、分解可能形式の間の関係がある。すなわち
\[
H_k(KHu)=2K_\nu(\nu\nu_1k)=2\theta_{\nu_1}(\nu\nu_1a\theta)
=2\theta_\nu^2\cdot a_\nu-a_\nu^2\cdot\theta_\nu+f\cdot\theta_\eta(\theta Hu)^2-f\cdot \nu\cdot\theta_\nu^3
\modu{\nu^3},
\]
\[
\begin{aligned}
K_\eta(KHu)&=-K_\nu^2(\nu\nu_1x)
=-a_\nu^2(a\theta u)(\nu\nu_1x)\\
&=a_\eta(aHu)^2\cdot\theta+a_\nu^2\cdot\theta_\nu
=-\theta_\eta(\theta Hu)^2\cdot f-\theta_\nu^2\cdot a_\nu.
\end{aligned}
\]
従って
\[
\tag{17\srcfnmark{*)}}
0=a_\nu^2\theta_{\nu_1}+\theta_\nu^2a_{\nu_1}+f\theta_\eta(\theta Hu)^2+\theta a_\eta(aHu)^2\modu{\nu^3}.
\]
\srcfntext{*)}{反変式を含む形式の積、たとえばこの場合の $f\cdot\theta_\nu^3\cdot\nu$ は、分解可能形式の間のすべての関係から省かれている。それらは系 $s_{IV}$ に公式を適用する際に意味をもたないからである。}
形式 $H_k$、$H_k(k\eta x)$、$H_k^2$、$H_k^2(k\eta x)$ は、分解可能形式 $H_\vartheta$ および $H_\vartheta^2$ との一回の折り畳みによって生じるものとして分解する（\S\ 3. 2), a), b) 参照）。

実際、
\[
H_k=H_\vartheta(a\theta u)\sim [H_\vartheta]_{ua}-f\cdot H_\vartheta(\theta Hu)
\]
である（2)a の特殊化 $y=uv$）。
\[
H_k(k\eta x)=H_\vartheta a_\eta-H_\vartheta\theta_\eta\cdot f,
\qquad
H_\vartheta a_\eta=\frac54[H_\vartheta]a-\frac14H_\vartheta a_\vartheta
\quad(\S\ 3. 2)b),
\]
\[
H_\vartheta^2(a\theta u)=[H_\vartheta^2]_{\widehat{ua}},
\qquad
H_\vartheta^2a_\eta=[H_\vartheta^2]a=H_k^2(k\eta x)+H_\vartheta^2\theta_\eta\cdot f.
\]

\noindent C) 形式列 $(j\eta x)^3$：B) による形式列 $K_\eta^3$ の二重還元によって還元可能である。次の Ansatz による：
\[
0=a_\eta^3(a\theta u)=\theta_\eta^3(a\theta u)
=\theta_\eta+(a\theta\eta x)^3(a\theta u)-\theta_\eta^3(a\theta u)
=\frac12(j\eta x)^3\modu{(\nu^3,s,\R,t,i,J)}.
\]
同様の Ansatz は、形式列の終形式まで成り立つ：
\[
0=H_\vartheta^2(a\theta\eta x)a_\eta^3
=H_\vartheta^2(a\theta\eta x)\theta_\eta^3
=\frac12H_j^2(j\eta x)^4\modu{(\nu^3,s,\R,t,i,J)}.
\]
形式 $(j\eta x)^2$：1) 分解可能形式 $H_\vartheta^2$ に対する関係式の二重偏極によって分解する（(12.) b)）；

2) 積 $a_\nu\cdot\theta_\nu^3$ の直接計算によっても分解する。これにより高次形式 $(\A Hu)^2$ も分解する。すなわち
\[
\tag{18}
\left.\begin{array}{ll}
\text{a)}&\dfrac13(\A Hu)^2+\dfrac13(j\eta x)^2=\theta_\nu^2a_{\nu_1}^2,\\[0.6em]
\text{b)}&(j\eta x)^2=\dfrac43\theta_\nu^3a_{\nu_1},
\end{array}\right\}\modu{(\nu^3,s,\R,t,i,J)}.
\]
これは次の Ansatz による：
\[
H_\vartheta^2(a\theta u)^2=\frac13(\A Hu)^2
=2\theta_\eta^2(a\theta u)(aHu)+\theta_\nu^2a_{\nu_1}^2
=(a\theta u)(aHu)\{a_\eta-(a\theta\eta x)^2\}+\theta_\nu^2a_{\nu_1}^2,
\]
\[
\begin{aligned}
a_{\nu_1}\theta_\nu^3
&=-(a\widehat{\nu}\nu_1u)\theta_\nu^2(\theta\widehat{\nu}\nu_1u)
-\frac12(u\nu\nu_1u)\theta_\nu\theta_{\nu_1}(\vartheta\nu\nu_1u)\\
&=-\frac32\theta_\eta^2(\theta Hu)(aHu)
=-\frac32\theta_\eta^2(\theta Hu)(a\theta u)\\
&=-\frac32\{a_\eta-(a\theta\eta x)^2\}\{(aHu)-(a\theta u)\}(a\theta u).
\end{aligned}
\]

\noindent D) 形式列 $\A_\eta(\A Hu)$：還元子 $\A_\nu^2$。

$(\A Hu)^2$ は C) による分解可能形式である。

\noindent E) 形式列 $a_l^2(aLu)$：還元子 $a_\eta^2(aHu)$。形式 $(aLu)$ および $a_l(aLu)$ は、\S\ 3 により関数行列式に関するトランスヴェクションとして分解する。

\noindent F) 形式列 $a_\sigma^3$：還元子 $a_\R^3$。

形式 $a_\sigma^2$ および $a_\sigma^3$：対応する式
\[
H_\nu a_\nu^3\quad\text{および}\quad H_\nu a_\eta^3,
\qquad
H_\nu a_\nu^3a_\nu\quad\text{および}\quad H_\nu a_\eta^4
\]
の二重還元により、還元子は $a_\nu^3$ および $a_\eta^3$ である（\S\ 10. C), $\A_\nu^2$ および $\A_\nu^3$ の還元を参照）。

これから、(16.) を考慮して、高次形式 $a_\nu s_\nu(asu)^2$, $a_\nu^2s_\nu(asu)^2$ および記号 $\R,t(a_\nu,a_\varrho^2)$ における形式の還元が得られる：
\[
\tag{19}
\left.\begin{array}{rcl}
a_\nu s_\nu(asu)^2&=&\dfrac13 i\theta_\nu^2-\dfrac1{18}iH,\\[0.5em]
a_\nu^2s_\nu(asu)^2&=&0,
\end{array}\right\}\modu{(\R,t)}.
\]
形式 $a_\sigma$ は、(12.)b の偏極によって、またはより短く積 $a_\nu^2\cdot\theta_\nu^3$ を次の Ansatz に従って直接展開することによって分解する：
\[
a_\nu^2\cdot\theta_{\nu_1}^3=(a_\nu^2\theta_{\nu_1})a_{\nu_1}-(a\theta\nu_1x)^2
=-2a_\eta(aHu)(a\theta H)(a\theta\eta x)+(a\theta\nu_1x)^2a_\nu^2\theta_{\nu_1}
\]
さらに $Hj(j\eta x)^2$ の還元 (C) を考慮すれば、
\[
\tag{20}a_\sigma=2a_\nu^2\theta_{\nu_1}^3\modu{(s,\R,t,i)}.
\]

\noindent\emph{帰結：}

1) $(j\eta x)^3$, $\A_\eta(\A Hu)$, $a_\sigma^2$ は還元子であり、$(j\eta x)^2$, $(\A Hu)^2$, $a_\sigma$ は分解可能形式である。

2) $f$ は、反傾形式との積としてのみ、$L$ との折り畳みに入る。従って $f$、また $K$ と $N$ も、$\sigma$ したがって $(\nu\sigma x)$ との折り畳みにまったく入らない。$j$ は反傾形式との積としてのみ入る。$\A$ は $H$ および $L$ との折り畳みに積としてまったく入らない（これは $\A_\eta(\A Hu)$ と $(\A Hu)^2$ から従う）。同様に $\sigma$、従って $(\nu\sigma x)$ は、系 I のいかなる形式との積にも折り畳みとして入らない。\S\ 11 の帰結を考慮すると、系 II の積として残るのは $H$ の冪と $H$ と $L$ の積だけである。

\begin{center}
\textbf{\S\ 13.\quad 八次の形式（系 III）。}
\end{center}
\[
\text{次の折り畳み：}\left\{\begin{array}{l}
\A\cdot j\text{ と }\nu,\\
\theta^2,\\
 f\cdot j,\\
 N\text{ と }H,\\
\theta\text{ と }L,\sigma.
\end{array}\right.
\]

\noindent\emph{既約形式：}
\[
\begin{array}{c@{\quad}l@{\qquad}c@{\quad}l}
\text{A)}&\A_\nu(j\nu x)&
\text{B)}&\begin{array}{c}(\vartheta^2\eta x)^3\cdot\\(\vartheta^2\eta x)^4\end{array}
\quad H_\vartheta(\vartheta^2\eta x)^3,
\ \theta_\eta(\vartheta^2\eta x)^3,\\[1em]
\text{C)}&(aHu)H_j,
\quad a_\eta(aHu)H_j&
\text{D)}&H_n,
\ N_\eta,
\end{array}
\]
\[
\text{E)}\quad
\begin{array}{c|c|c|c}
\cdot&(\theta Lu)^2&(\theta Lu)^3&\cdot\\
\theta_l&\cdot&L_\vartheta(\theta Lu)^2,
\ \theta_l(\theta Lu)^2&\theta_l(\theta Lu)^3\\
(\vartheta lx)^2\cdot&\theta_l^2&\cdot&\cdot\\
\cdot&\theta_l(\vartheta lx)^2&\cdot&\cdot\cdot
\end{array}
\]

\noindent\emph{還元：}

\noindent A) $\A_\nu(j\nu x)^3$：還元子 $a_\nu(j\nu x)^3$。

$\A_\nu(j\nu x)^2$ は、\S\ 3. 2)b により分解可能形式 $a_\nu(j\nu x)^2$ によって還元可能であり、還元子 $\theta_{\vartheta'}$ を考慮する。実際、\S\ 11 B) により
\[
\frac3{10}\A_\nu(j\nu x)^2+\frac35a_\nu(j\nu\vartheta)a_\vartheta(j\nu x)
+\frac1{10}a_\nu(j\nu\vartheta)^2
=\frac35\theta_\nu^3\cdot\theta_{\vartheta''}^{\prime2}+\frac25\theta_\nu^3\theta_{\vartheta''}\theta_{\vartheta''}'.
\]
（還元子 $a_j$ および $\theta_{\vartheta'}$）。

\noindent B) 形式 $H_\vartheta(\vartheta'\eta x)^2$, $H_\vartheta^2(\vartheta'\eta x)$, $H_\vartheta^2(\vartheta'\eta x)^2$ は、\S\ 3. 2)b) および関係式
\[
H_\vartheta(\vartheta'\eta x)^2=2H_\vartheta a_\eta^2\cdot f-2H_\vartheta a_\eta b_\eta
\]
により、分解可能形式 $H_\vartheta$ と $H_\vartheta^2$、それぞれ $H_\vartheta a_\eta$ と $H_\vartheta^2a_\eta$ との逐次的な一回の折り畳みによって生じるものとして分解する。

\noindent C) 形式列 $a_\eta(j\eta x)^2$：還元子 $a_j$ および $(j\eta x)^3$。

形式 $a_\eta(j\eta x)$ は分解する。還元子 $a_j$ と、\S\ 3, 2)a（特殊化 $y=uv$）による分解可能形式 $(j\eta x)^2$ との一回の折り畳みである。

形式 $a_\eta H_j$ は分解する。1) 分解可能形式 $a_\nu(j\nu x)^2$ との一回の折り畳みによる。

2) 分解可能形式 $(j\eta x)^2$ との特殊な二回の折り畳みによる。これから分解可能形式間の関係が出る。

$a_\nu(j\nu x)^2$ から：
\[
0=\theta_\nu^3\cdot\theta_\nu+2a_\eta H_j\modu{(s,\R,t,i,J)}.
\]
$(j\eta x)^2$ から：
\[
0=\theta_\nu^3\cdot\theta_\nu-\frac32a_\eta H_j\modu{(s,\R,t,i,J)}
\]
（反変式との積は省略する）。これは次の Ansatz による：
\[
0=\frac12a_\nu(abu)^2\theta_\nu^3
+\frac12a_\nu(abu)\theta_{\nu_1}^3(\theta bu)
-\frac34(\widehat{j}\eta bu)^2+\frac34H_j(\widehat{j}\eta bu)
-\frac18f\cdot H_j^2
\]
および $a_\nu(abu)(\theta bu)\theta_{\nu_1}^3$ における記号の交換による。

\noindent D) 形式列 $N_\eta(NHu)$：還元子 $\A_\nu^2$。

$(NHu)$, $(n\eta x)$, $N_\eta(n\eta x)$, $H_n(NHu)$ は分解可能形式である（\S\ 12 B 参照）。$(NHu)^2$ は形式 $(\A Hu)^2$ との一回の折り畳みによって分解する。

\noindent E) 形式列 $\theta_lL_\vartheta$, $\theta_l^2(\theta Lu)^2$, $\theta_l^2(\vartheta lx)$：

還元子：$\theta_\eta H_\vartheta$, $\theta_\eta^2(\theta Hu)^2$, $\theta_\eta^2(\vartheta\eta x)$。

形式 $(\theta Lu)$, $(\vartheta lx)$, $L_\vartheta$, $L_\vartheta(\theta Lu)$, $\theta_l(\theta Lu)$, $L_\vartheta(\vartheta lx)$, $\theta_l(\vartheta lx)$, $L_\vartheta^2$, $L_\vartheta^2(\theta Lu)$, $\theta_l^2(\theta Lu)$ は分解する。（\S\ 12 B の分解可能形式に双対的に反対である。）

\noindent F) 形式列 $\theta_\sigma(\vartheta\sigma x)$：還元子 $a_\sigma^2$。

形式 $(\vartheta\sigma x)$ および $(\vartheta\sigma x)^2$ は、分解可能形式 $a_\sigma$ との一回の折り畳みによって生じるものとして分解する。二回の折り畳みによって生じる $\theta_\sigma$ も、次の形式列
\[
a_{\nu_2}^2(a\theta u)\theta_{\nu_1}^3\equiv H_j^2(j\eta x)\equiv0\modu{(s,\R,t,i,J)}:
\]
により分解する：
\[
\theta_\sigma=a_\sigma(abu)^2=a_\nu^2(abu)^2\cdot\theta_\nu^3.
\]

\noindent\emph{帰結：}

$\theta^3$ または $\theta^2K$ は $H$ との折り畳みに入らない。$\theta$ は、反変式と見たときにのみ、冪または積として $L$ との折り畳みに入るが、$\sigma$ および $(\nu\sigma x)$ との折り畳みにはまったく入らない。

$N$ は系 II のいかなる形式との積にも折り畳みに入らず、また系 II の形式の冪または積とも同様に入らない。直前の諸段落の帰結を考慮すると、系 I の積として残るのは、$f\cdot j$, $\A\cdot j$, $\theta$ の冪、および $\theta\cdot K$ の積だけである。

\begin{center}
\textbf{\S\ 14.\quad 九次の形式（系 III）。}
\end{center}
\[
\text{次の折り畳み：}\left\{\begin{array}{l}
\theta\cdot K\text{ と }H,\\
K,j,\A\text{ と }L,\\
j,\A\text{ と }\sigma,\\
f\text{ と }H^2.
\end{array}\right.
\]

\noindent\emph{既約形式：}
\[
\renewcommand{\arraystretch}{1.35}
\begin{array}{@{}r@{\quad}c@{\qquad}c@{\qquad\qquad}r@{\quad}c@{\qquad}c@{\qquad}c@{\qquad}c@{}}
&\cdot&\cdot&&\cdot&\cdot&(KLu)^3&\cdot\\
\text{A)}&(\vartheta\cdot k,\eta,x)^4\cdot&\cdot&\text{B)}&\cdot&\cdot&\cdot&K_l(KLu)^3\\
&\cdot&H_k(\vartheta\cdot k,\eta,x)^4&&\cdot&K_l^2&\cdot&\cdot\\
&&&& (klx)^3&\cdot&\cdot&\cdot\\
&&&&\cdot\cdot&K_l(Klx)^3&\cdot&\cdot
\end{array}
\]
\[
\renewcommand{\arraystretch}{1.35}
\begin{array}{@{}r@{\quad}c@{\qquad}c@{\qquad}r@{\quad}c@{\qquad}r@{\quad}c@{\qquad}r@{\quad}c@{\qquad}c@{}}
\text{C)}&L_j&\cdot&\text{D)}&\A_l,&\text{E)}&(j\sigma x),&\text{G)}&(aH^2u)^3&(aH^2u)^4\\
&\cdot&L_j^2,&&&&&&\cdot&a_\eta(aH^2u)^3
\end{array}
\]

\noindent\emph{還元：}

A) 形式 $H_\vartheta(k\eta x)^3$, $H_\vartheta^2(k\eta x)^2$, $H_\vartheta^2(k\eta x)^3$ は、分解可能形式との逐次的な一回の折り畳みによって生じるものとして分解する。

B) 形式列 $K_lL_k$, $K_l^2(KLu)$, $K_l^2(klx)$：

還元子：$\theta_\eta H_\vartheta$, $a_\eta^2(aHu)$, $\theta_\eta^2(\vartheta\eta x)$。

形式 $K_l$, $K_l(KLu)$, $K_l(klx)$, $(KLu)^2$, $(klx)^2$ は、関数行列式に関するトランスヴェクションとして分解する（\S\ 3）。

形式 $L_k$, $L_k(KLu)$, $L_k(KLu)^2$, $L_k(klx)$, $L_k(klx)^2$, $L_k^2$, $L_k^2(KLu)$, $L_k^2(klx)$, $L_k^3$, $K_l(KLu)^2$, $K_l(klx)^2$ は、次の分解可能形式との一回の折り畳みによって生じるものとして分解する：
\[
L_\vartheta,
\ H_k(KHu),
\ L_\vartheta(\vartheta lx),
\ H_k^2,
\ L_\vartheta^2,
\ L_\vartheta^2(\theta Lu),
\ K_\eta(KHu),
\ \theta_l(\vartheta lx)
\]
これは \S\ 3. 2), a), b) による。

C) 形式列 $(jlx)^3$：還元子 $(j\eta x)^3$。

形式 $(jlx)^2$ および $L_j(jlx)$：分解可能形式 $(j\eta x)^2$ との一回の折り畳みによって生じる。

D) 形式列 $\A_l(\A Lu)$：還元子 $\A_\nu^2$。

形式 $(\A Lu)^2$, $(\A Lu)^3$：分解可能形式 $(\A Hu)^2$ との一回の折り畳み。

E) 形式列 $(j\sigma x)^2$：$j$ を形式列の低次形式で置き換えることにより、還元子は $a_\sigma^2$ である（\S\ 5）：
\[
a_\sigma^2(a\theta u)^2=\frac13(j\sigma x)^2\ldots,
\qquad
a_\sigma^2(a\theta\sigma x)^2=\frac13(j\sigma x)^4.
\]
F) 形式列 $\A_\sigma^2$：還元子 $a_\sigma^2$。

形式 $\A_\sigma$：$\A$ を形式列の低次形式で置き換えることにより、還元子は $a_\sigma^2$ である：
\[
a_\vartheta a_\sigma^2=\frac23\A_\sigma\modu{\nu}.
\]

\noindent\emph{帰結：} $\sigma$ および $(\nu\sigma x)$ との折り畳みに入るのは形式 $j$ だけである。

$N$ は $L$ との折り畳みに入らない。なぜなら、$\A_l$ に対応する形式 $N_l$ は分解するからである。

\begin{center}
\textbf{\S\ 15.\quad 十次の形式（系 III）。}
\end{center}
\[
\text{次の折り畳み：}\left\{\begin{array}{l}
\theta^2,\\
\ f\cdot j\text{ と }L,\\
\theta\text{ と }H^2.
\end{array}\right.
\]

\noindent\emph{既約形式：}
\[
\begin{array}{c@{\quad}l@{\qquad}c@{\quad}l}
\text{A)}&(\vartheta^2lx)^4,
\ \theta_i(\vartheta^2lx)^4&
\text{B)}&(aLu)^2L_j,
\ a_l(aLu)^2L_j,\\[0.8em]
\text{C)}&(\theta H^2u)^3,
\ (\theta H^2u)^4,
\ H_\vartheta(\theta H^2u),
\ \theta_\eta(\theta H^2u)^3.
\end{array}
\]

\noindent\emph{還元：}

A) および B)：残りの形式は、分解可能形式との一回の折り畳みによって分解する。

C) 残りの形式は、\S\ 13 B) に従い、双対的に反対の考察によって分解する。

\noindent\emph{帰結：}

$\theta^2\cdot K$ は $L$ との折り畳みに入らない。対応して、$H^2L$ は $K$ との折り畳みに入らない。$H^3$ および $H^2L$ は $\theta$ との折り畳みに入らない。これにより \S\ 7 の折り畳み図式が証明された。

\begin{center}
\textbf{\S\ 16.\quad 11、12、13、15 次の形式（系 III）。}
\end{center}

\noindent\emph{11 次。}
\[
\text{次の折り畳み：}\left\{\begin{array}{l}
\theta\cdot K\text{ と }L,\\
K,j\text{ と }H^2,\\
j\text{ と }(\nu\sigma x),\\
f\text{ と }H\cdot L.
\end{array}\right.
\]
\noindent\emph{既約形式：}
\[
\begin{array}{c@{\quad}l@{\qquad}c@{\quad}l}
\text{A)}&(\vartheta\cdot k,l,x)^5,
\ \theta_i(\vartheta\cdot k,l,x)^5&
\text{B)}&(KH^2u)^4,
\ K_\eta(KH^2u)^4,\\[0.8em]
\text{C)}&(\nu\sigma j),&
\text{D)}&(a,H\cdot L,u)^4.
\end{array}
\]

\noindent\emph{12 次。}
\[
\text{次の折り畳み：}\left\{\begin{array}{l}
\theta^3\text{ と }L,\\
\theta\text{ と }H\cdot L.
\end{array}\right.
\]
\noindent\emph{既約形式：}
\[
\text{E)}\ (\vartheta^3lx)^6,
\qquad
\text{F)}\ (\theta,H\cdot L,u)^4,
\qquad L_\vartheta(\theta,H\cdot L,u^4).
\]

\noindent\emph{13 次。}
\[
K\text{ と }H\cdot L\text{ の折り畳み。}
\]
\noindent\emph{既約形式：}
\[
\text{G)}\ (K,H\cdot L,u)^5,
\qquad K_\eta(K,H\cdot L,u)^5.
\]

\noindent\emph{15 次。}
\[
\text{ と }K\text{ の折り畳み。}H_3
\]
\noindent\emph{既約形式：}
\[
\text{H)}\ (KH^3u)^6.
\]

\noindent\emph{還元：}

形式 $H_{j^3}^2$, $H_{j^4}^2$ は、次の関係式から還元子 $L_j^3$ をもつ：
\[
(\nu lx)L_x^3=-\frac12H^2\modu{(\sigma,s)}.
\]
残りの形式の分解または還元可能性は、以前の各節の考察から従う。

\noindent\emph{帰結：}

\S\ 7 の折り畳み図式および \S\S\ 8--15 の帰結により、系 III の既約形式（117 個）はこれで尽くされた。

\clearpage

\begin{center}
{\Large\bfseries 第 IV 章。\quad 法 $(\varrho,t)$ に関する相対的完全系。}
\end{center}

\begin{center}
\textbf{\S\ 17.\quad 法 $(ss'u)^4$ および $s$ の系の還元。}
\end{center}

次に高い相対的完全系を形成するには、\S\ 7 と同様に、次の法
\[
\nu(s)=(ss'u)^4
\]
に関する $s$ の系を作り、これを法 $s$ に関する相対的完全系の形式と折り畳む必要がある。高い法として法 $(\varrho,t)$ を導入すると、次が示される：
\[
\begin{array}{ll}
\text{A)}&\text{法 }(ss'u)^4\text{ は還元可能になる，}\vphantom{\dfrac11}\\
\text{B)}&s\text{ の系はただ一つの形式 }s\text{ から成る。}
\end{array}
\]

A) 法 $(ss'u)^4$ の還元は、\S\ 11 の式 (16.) で得られた還元子 $s_\nu^2$ から直ちに従う。すなわち
\[
 s_\nu^2=\frac43 i a_\nu^2+\frac45\theta_\varrho^2+\frac85(atu)^2-
 \frac{26}{15}(\nu\varrho x)^2+\frac16J\cdot u_x^2-
 \frac29i^2\cdot u_x^2
\]
から、$x$ を $\nu_1$ に偏極して、
\[
\tag{21}
\begin{aligned}
(ss'u)^4&=-\frac23J\cdot\nu+6s_\nu^2s_{\nu_1}^2\\
&=\frac{24}{5}\theta_\nu^2\theta_\varrho^2+
\frac{48}{5}a_\nu^2(atu)^2-
\frac{52}{5}H_\varrho^2+
\frac43i\cdot(asu)^4+\frac13J\cdot\nu-
\frac49i^2\cdot\nu,
\end{aligned}
\]
従って法 $(ss'u)^4$ の還元が得られる。\footnote{式 (21.) から、序論の注で述べられた 9 次の三つの不変式の間の関係が、式 (10.)c を用いて従う。}

B) 
\[
Z=(ss'u)^2=\theta(s)
\]
の還元、従って $s$ の系の還元は、\S\ 7 の式 (8.)b により、関係式
\[
 s_\eta^2=s_\nu^2(\nu\nu_1x)^2-\frac13Z
\]
を考慮して、形式 $Z$ をより低い形式 $g_\nu^2$ で置き換えることによって得られる。
形式 $g_\nu^2$ は、\S\ 11 の式 (13.) によって、還元子 $s_\nu^2$ を因子にもつ。式 (8.) および (16.) から
\[
 g_\nu^2=-Z+\frac{14}{5}ia_\eta^2+\frac{14}{15}i(asu)^2
 \modu{(\varrho,t)},
\]
また式 (13.), (14.), (15.), (16.) から
\[
\begin{aligned}
g_\nu^2&=2[s_\vartheta^2]_{\nu^2}-\frac43 i\A_\nu^2
      =2s_\vartheta^2s_\nu^2-\frac65[s_\vartheta^2(\theta su)^2]\widehat{\nu x}^{\,2}
        -\frac43i\A_\nu^2\\
&=\frac45i\cdot a_\eta^2-\frac25i\cdot(asu)^2+\frac13J\cdot\theta
\modu{(\varrho,t)}.
\end{aligned}
\]
従って
\[
\tag{23}
 Z=2i\cdot a_\eta^2+\frac43i\cdot(asu)^2-\frac13J\cdot\theta
 \modu{(\varrho,t)}.
\]

したがって、法 $((ss'u)^4,\varrho,t)$ に関する相対的完全系は、法 $(\varrho,t)$ に関する相対的完全系に移る。そしてここから、二つの二次形式の既知の系に関するトランスヴェクションによって、絶対的完全系が生じる。法 $(\varrho,t)$ に関する相対的完全系を形成するには、式 (23.) によって $s$ の系が形式 $s$ に還元されるので、法 $(s,\varrho,t)$ に関する相対的完全系を $s$ および $s$ の冪に関してトランスヴェクトしなければならない。$s$ の冪は系 I との折り畳みにのみ入ることが分かる。

予告した関係式は次の通りである：
\[
\tag{22}
\begin{aligned}
(ass')^4&=\frac{24}{5}\theta_\nu^2\theta_\varrho^2a_\nu^2a_\varrho^2+
\frac{48}{5}a_\nu^2b_\nu^2(tab)^2-
\frac{52}{5}H_\varrho^2a_\nu^4+
\frac53J\cdot i-\frac49i^3,\\
(ass')^4&=\frac{24}{5}a_\varrho^2a_{\varrho'}^2-
\frac{12}{5}t_\varrho^2+\frac53J\cdot i-\frac49i^3.
\end{aligned}
\]

同じ位数かつ同じ次数の形式のうち、「高次形式」とは、法 $(s,\varrho,t)$ に関する相対的完全系の高次形式から $s$ との折り畳みによって生じたものをいう。ただし、$s$ との折り畳みは元の形式の偏極にすぎないものと見る。これにより、個々の形式の配列は一意に定まる（\S\ 1、形式列 2 参照）。例えば
\[
(\theta_\eta(\theta Hu),s,u)^2\text{ は }s_\eta((\theta Hu)^2,s,u)\text{ より高い，}
\]
\[
(\theta_\eta(\theta Hu)^2,s,u)\text{ は }(\theta_\eta(\theta Hu),s,u)^2\text{ より高い。}
\]

得られる系を四つの部分系に分ける：
\[
\begin{array}{ll}
\text{系 }s_I:&s\text{ および }s\text{ の冪と系 I との折り畳みによって生じる；}\vphantom{\dfrac11}\\
\text{系 }s_{II}:&s\text{ と系 II との折り畳みによって生じる；}\vphantom{\dfrac11}\\
\text{系 }s_{III}:&s\text{ と系 III の個々の形式}\\
&\text{（積を除く）との折り畳み、および系 I と II から}\\
&\text{それぞれ一つずつ取った形式の積との折り畳みによって生じる；}\vphantom{\dfrac11}\\
\text{系 }s_{IV}:&s\text{ と、系 I と III、II と III、III と III からそれぞれ}\\
&\text{一つずつ取った形式の積との折り畳みによって生じる。}
\end{array}
\]
注意：以後、すべての公式は法 $(\varrho,t)$ で理解する。式 (27.) 以降は、これを毎回明記せずに、法 $(\varrho,t,i,J,\text{高次形式})$ で理解する。

還元のため、以前に得た公式をまとめる：
\[
\tag{24\srcfnmark{*)}}
\begin{aligned}
s_\nu^2&=\frac43ia_\nu^2+\frac16J\cdot u_x^2-\frac29i^2\cdot u_x^2,
& s_\nu^3&=\frac13J\cdot u_x,
& s_\nu^4&=J,\\[0.3em]
g&=2s_\vartheta^2-\frac43i\A,
& s_\vartheta^2(\theta su)&=0,\\[0.3em]
s_\vartheta^2(\theta su)^2&=\frac89i\cdot a_\nu^2-\frac4{27}i^2\cdot u_x^2,\\[0.3em]
Z&=2i\cdot a_\eta^2+\frac43i(asu)^2-\frac13J\cdot\theta,\\[0.3em]
g_\nu&=-s_\eta+u_x\left\{2ia_\eta^2-\frac23i(asu)^2+\frac23J\cdot\theta\right\},\\[0.3em]
g_\nu^2&=\frac45ia_\eta^2-\frac25i(asu)^2+\frac13J\cdot\theta,
&g_\nu^3&=0,
&g_\nu^4&=0.
\end{aligned}
\]
\srcfntext{*)}{記法を短くするため、完全な折り畳み $[(\theta Hu)^2]_{\widehat{su}}s$ の代わりに、その折り畳みの項を系の形式として採用する。公式を立てる際には、すべての折り畳みをその初項で置き換える。たとえば $(\theta_\eta(\theta Hu),s,u)^2$ を $\theta_\eta(\theta Hu)(\theta su)^2$ で置き換える。}

\noindent\emph{帰結：}
\[
 s_\nu^2,
\quad s_\vartheta^2(\theta su),
\quad s_\vartheta^2s_\nu,
\quad s_\vartheta^2\theta_\nu
\]
は還元子である。



\clearpage
\begin{center}
\textbf{\S\ 18.\quad 系 $s_I$.}
\end{center}

\[
\text{次の折り畳み：}\quad
\left\{\begin{array}{l}
s\text{ と } f;\ \theta;\ K,j,\A;\ f\cdot j,N;\ \A\cdot j,\\
s^2\text{ と }\theta,K.
\end{array}\right.
\]

\noindent\emph{既約形式。}

\noindent 5次：
\[
(asu);\qquad (asu)^2;\qquad (asu)^3;\qquad (asu)^4.
\]

\noindent 6次：
\[
\begin{array}{cccc}
(\theta su)&(\theta su)^2&(\theta su)^3&(\theta su)^4\\
s_\vartheta&s_\vartheta(\theta su)&s_\vartheta(\theta su)^2&s_\vartheta(\theta su)^3\\
\cdot&s_\vartheta^2&\cdot&\cdot
\end{array}
\]

\noindent 7次：
\[
\begin{array}{c@{\qquad}c@{\qquad}c@{\qquad}c}
\cdot&(Ksu)^2&(Ksu)^3&(Ksu)^4\\
\text{A)}\ s_k&s_k(Ksu)&s_k(Ksu)^2&s_k(Ksu)^3\\
\cdot&s_k^2&\cdot&\cdot
\end{array}
\qquad
\begin{array}{l}
\text{B)}\ s_j,\\[0.4em]
\text{C)}\ (\A su),\ (\A su)^2.
\end{array}
\]

\noindent 8次：
\[
\begin{array}{ll}
\text{A)}&s_j(asu),\ s_j(asu)^2,\ s_j(asu)^3,\\[0.4em]
\text{B)}&(Nsu)^2,\quad s_n,\quad s_n(Nsu).
\end{array}
\]

\noindent 10次：
\[
\text{A)}\ s_j(\A su),\qquad \text{B)}\ s_\vartheta(\theta s'u)^4.
\]

\noindent 11次：
\[
(Ks^2u)^6\cdot\qquad s_k(Ks^2u)^5,
\qquad s_k(Ks^2u)^6.
\]

\noindent\emph{還元：}

形式列
\[
 s_\vartheta^2(\theta su),\quad s_k^2(Ksu),\quad s_j^2,\quad (\A su)^3,
 \quad (Nsu)^3
\]
は、還元子 $s_\vartheta^2(\theta su)$ を因子にもつ。$s_j^2$ と $(\A su)^3$ は、$j$ と $\A$ を形式列の低次形式へ戻すことによって得られる：
\[
\begin{aligned}
s_\vartheta^2(\theta su)(a\theta u)^3&=-\frac12s_j^2,\\
s_\vartheta^2(\theta su)(a\theta u)^2&=\frac13(\A su)^3,\\
s_\vartheta^2(\theta su)^2(a\theta u)^2&=\frac13(\A su)^4+\frac13s_j^2.
\end{aligned}
\]
形式 $s_\vartheta^2s_{\vartheta'}$ と $s_\vartheta^2s_{\vartheta'}^2$：還元子 $s_\vartheta^2(\theta su)$ と $\theta_{\vartheta'}$：
\[
\begin{aligned}
s_\vartheta^2s_{\vartheta'}&=s_\vartheta^2\theta_{\vartheta'}-s_\vartheta^2(\theta s\widehat{\vartheta'}x),\\
s_\vartheta^2s_{\vartheta'}^2&=s_\vartheta^2s_{\vartheta'}\theta_{\vartheta'}-s_\vartheta^2s_{\vartheta'}(\theta s\widehat{\vartheta'}x).
\end{aligned}
\]
形式列 $s_j(\A su)^2$：還元子 $\A_j$ と $(\A su)^3$。

\noindent\emph{帰結：}

$s$ は $f,j,\A,N$ との折り畳みにおいて冪としては入らない。$f,j,\A,N$ は、反傾形式との積においてだけ折り畳みに入る。ただし、$f$ との積ではそれらの形式はなお $x$ に関して二次でありうるし、$\A$ との積ではなお $x$ に関して一次でありうる。従って、次の形式の積を考えればよい：
\[
 f\cdot(0,\lambda),\quad f\cdot(1,\lambda),\quad f\cdot(2,\lambda),\quad
 j\cdot(\mu,0),\quad N,\quad
 \Delta\cdot(0,\lambda)\quad \Delta\cdot(1,\lambda)\quad(\lambda>0),
\]
ここで $(\mu,\lambda)$ は、$x$ に関して次数 $\mu$、$u$ に関して次数 $\lambda$ の形式を表す。

$\theta$ も $K$ も、$s$ または $s^2$ との折り畳みにおいて、任意の形式との冪または積には入らない。実際、関数行列式 $K$ を含む積 $K\cdot\varphi_x$（$\varphi\ne f$ または $\theta$）については、分解可能形式間の関係式
\[
 K\cdot\varphi_x=(a\theta u)\varphi_x=f\cdot(\varphi\theta u)-\theta\cdot(\varphi au)
\]
が成り立つ。従って、上に掲げた系 $s_I$ の折り畳みの図式が証明された。

\begin{center}
\textbf{\S\ 19.\quad 系 $s_{II}$.}
\end{center}

$s$ と $\nu;\ H;\ L;\ H^2$ との折り畳み。

\noindent\emph{既約形式：}
\[
\begin{array}{c@{\qquad}c@{\qquad}c@{\qquad}c}
\text{6次} & \text{8次} & \text{10次} & \text{12次}\\[0.3em]
s_\nu & (Hsu),\ (Hsu)^2 & (Lsu)^2,\ (Lsu)^3 & (H^2su)^3\\
\cdot & s_\eta & s_l & \cdot\\
\cdot & s_\nu^4=J & \cdot & \cdot\cdot
\end{array}
\]

\noindent\emph{還元：}

形式列 $s_\eta(Hsu)$ と $s_l(Lsu)$：還元子 $s_\nu^2$。

形式列 $s_\sigma$：$H$ から来る還元子 $s_\nu^2$、すなわち $s_\nu^2=\dfrac23s_\sigma$。

$(H^2su)^4$ は式 (7.)a により分解する（\S\ 11 A 参照）。従って $(H\cdot L,s,u)^4$ も分解する。

\noindent\emph{帰結：}

$s^2$ は系 II との折り畳みに入らない。形式 $\nu$ は反傾形式との積においてだけ入る。形式 $H$ は、共変形式として見れば、任意の $\lambda$ について形式 $(1,\lambda)$、$(2,\lambda)$、$(3,\lambda)$ との積において折り畳みに入る。形式 $\sigma$ と $(\nu\sigma x)$ は全く入らず、$L$ は関数行列式としては折り畳みの積に入らない。系 III の共変形式に関する完全な超過は還元可能になる。系 I と II の積として残るのは
\[
 f\cdot\nu,\qquad \A\cdot\nu.
\]

\begin{center}
\textbf{\S\ 20.\quad 七次の形式（系 $s_{III}$）。}
\end{center}

$s$ と $f\cdot\nu$、$a_\nu$、$a_\nu^2$ との折り畳み。

\noindent\emph{既約形式：}
\[
\begin{gathered}
s_\nu(asu),\quad a_\nu(asu),\quad s_\nu(asu)^2,\quad a_\nu(asu)^2,
\quad s_\nu(asu)^3,\quad a_\nu(asu)^3,\\
a_\nu s_\nu,
\quad a_\nu s_\nu(asu),
\quad a_\nu^2(asu),
\quad a_\nu^2(asu)^2.
\end{gathered}
\]

\noindent\emph{還元：}

還元子 $s_\nu^2$ と $s_\eta(Hsu)$ との直接の折り畳みから生じる自明な還元は述べない。また、関数行列式を含むトランスヴェクションの分解も述べない。

形式 $a_\nu s_\nu(asu)^2$ と $a_\nu^2s_\nu(asu)^2$ については式 (19.)、\S\ 12 を見よ。さらに
\[
 a_\nu s_\nu(asu)^3=a_\nu(a\widehat{\nu}_1\nu_2u)^3(\nu_1\nu_2\nu)=0.
\]
形式列 $a_\nu^2s_\nu$：$a_\nu^2$ を低次形式へ戻すこと、すなわち式
\[
 s_\vartheta^2(a\theta s)(a\theta u),
 \qquad s_\vartheta^2(a\theta s)(a\theta u)^2
\]
の二重還元によって得られる還元子 $s_\vartheta^2(\theta su)$。Ansatz は
\[
\begin{aligned}
s_\vartheta^2(a\theta s)(a\theta u)
&=[(a\theta u)^2]s^3+\frac15[a_\vartheta(a\theta u)^2]s^2+\frac1{60}[a_\vartheta^2(a\theta u)^2]s\\
&=\frac12a_\nu^2s_\nu-\frac23a_\nu s_\nu^2,\\
s_\vartheta^2(a\theta s)(a\theta u)^2
&=\frac45[a_\vartheta(a\theta u)^2]_{\widehat{us}}s^2+\frac{23}{180}[a_\vartheta^2(a\theta u)^2]_{\widehat{us}}s\\
&=-\frac12a_\nu^2s_\nu(asu).
\end{aligned}
\]
ここで $(a_\nu^2b_\nu(abu)=0)$ である。

得られる公式は
\[
\tag{25}
\begin{array}{rcl@{\qquad}rcl}
a_\nu^2s_\nu&=&\dfrac49i\,a_\nu^2b_\nu,
& a_\nu s_\nu(asu)^2&=&\dfrac13i\theta_\nu^2-\dfrac1{18}iH,\\[0.7em]
a_\nu s_\nu(asu)^3&=&0,
& a_\nu^2s_\nu(asu)&=&0,\\[0.7em]
a_\nu^2s_\nu(asu)^2&=&0.
\end{array}
\]

\noindent\emph{帰結：}

1) $a_\nu s_\nu(asu)^2$ と $a_\nu^2s_\nu$ は還元子である。

2) \S\ 19 で既に還元可能と認められた形式 $(\A su)^3$、$(\A su)^4$ の値が得られる。同様に、対応する形式列 $(agu)^2$ についても、$\nu$ との折り畳みに際して還元子 $g_\nu$ により二重還元が生じる：
\[
\tag{26}
\begin{aligned}
\text{a)}\quad (\A su)^3&=-\frac65a_\nu^2(asu)+3a_\nu s_\nu(asu)+2i\theta_\nu(\theta\nu x),\\
\text{b)}\quad (\A su)^4&=-\frac25a_\nu^2(asu)^2+\frac43i\theta_\nu^2+\frac19iH.
\end{aligned}
\]
式 (24.) と (26.) から、法 $(i,J)$ で（p. 71 参照）、
\[
\tag{27}
\begin{aligned}
\text{a)}\quad (agu)^2&=\frac23(\A su)^2-\frac43a_\nu s_\nu+\frac35s\cdot a_\nu^2,\\
\text{b)}\quad (agu)^3&=2a_\nu s_\nu(asu)-a_\nu^2(asu)^2,\\
\text{c)}\quad (agu)^4&=-a_\nu^2(asu)^2.
\end{aligned}
\]

\begin{center}
\textbf{\S\ 21.\quad 八次の形式（系 $s_{III}$）。}
\end{center}

$s$ と四次の形式（系 III）との折り畳み（\S\ 9）。

\noindent\emph{既約形式：}
\[
\begin{array}{lll}
(\vartheta\nu x)s_\nu,& ((\vartheta\nu x),s,u)^2,\\[-0.1em]
((\vartheta\nu x)^2,s,u),& ((\vartheta\nu x)^2,s,u)^2,\ \theta s_\nu,\\[-0.1em]
(\vartheta\nu x)^2s_\nu,& ((\vartheta\nu x)^2,s,u)s_\nu,
\end{array}
\]
\[
\begin{array}{lll}
((\vartheta\nu x),s,u)^3,\ (\theta_\nu su)^2,&
((\vartheta\nu x),s,u)^4,\ (\theta_\nu su)^3,\\[-0.1em]
((\vartheta\nu x)^2,s,u)^3,\ (\theta_\nu(\vartheta\nu x),s,u)^2,\ (\theta_\nu^2su),&
((\vartheta\nu x),s,u)^3,\ (\theta_\nu^2su)^2,\ (\theta_\nu(\vartheta\nu x),s,u)^4.
\end{array}
\]

\noindent\emph{還元：}

$((\vartheta\nu x),s,u)s_\nu$ は、\S\ 3 により関数行列式に関するトランスヴェクションとして分解する。

形式列 $((\vartheta\nu x),s,u)^2s_\nu$：還元子 $s_\nu^2$ と $s_\vartheta^2\theta_\nu$：
\[
(\widehat\vartheta\nu su)s_\nu(\theta su)=s_\nu s_\vartheta(\theta su)=-\frac12(\widehat\vartheta\nu su)^2(\theta su),
\]
\[
\theta_\nu(\widehat\vartheta\nu su)s_\nu=\theta_\nu s_\nu s_\vartheta=-\frac12\theta_\nu(\widehat\vartheta\nu su)^2.
\]
形式列 $\theta_\nu s_\nu(\theta su)$：次の式の二重還元による還元子 $a_\nu s_\nu(asu)^2$：
\[
a_\nu s_\nu(asu)^2(abu)\quad\text{および}\quad a_\nu s_\nu(asu)^2(abu)(sbu);
\]
$\theta_\nu s_\nu(\theta su)^3$ は $(agu)^3(abu)(bgu)^3$ の二重還元による。

形式列 $\theta_\nu(\vartheta\nu x)s_\nu$：$\theta_\nu(\vartheta\nu x)s_\nu=a_\nu^2(abu)s_\nu$ からの還元子 $a_\nu^2s_\nu$。

形式列 $(\theta_\nu(\vartheta\nu x)^2,s,u)$：形式列 $\theta_\nu(\widehat\vartheta\nu su)s_\nu$ の二重還元。この列は同時に $((\vartheta\nu x)su)^2s_\nu$ と $\theta_\nu(\vartheta\nu x)s_\nu$ の形式列に属する。

形式 $(\theta_\nu(\vartheta\nu x),s,u)$ は、関数行列式 $\theta_\nu(\vartheta\nu x)=a_\nu^2(abu)$ に関する一回のトランスヴェクションとして分解する。

形式 $(\vartheta \nu x)^2,s,u)^4$：式 (27.) から、かつ $(j\eta x)^4$ の還元（\S\ 12 C）と同様に、式 $(a\A u)^2(\A su)^4$、または $(\theta gu)^4$ の二重還元。

二重還元により次の公式が得られる：
\[
\tag{28}
\begin{aligned}
0&=\theta_\nu s_\nu(\theta su)-\frac16(\widehat\vartheta\nu su)^2(\theta su)-\frac1{12}(Hsu),\\
0&=\theta_\nu s(\theta su)^2-\frac14(\widehat\vartheta\nu su)^2(\theta su)^2-\frac1{24}(Hsu)^2,\\
0&=(\widehat\vartheta\nu su)^2(\theta su)^2+2\theta_\nu(\widehat\vartheta\nu su)(\theta su)^2
  +3\theta_\nu^2(\theta su)^2-\frac13H(su)^2,\\
0&=\theta_\nu s_\nu(\theta su)^3+\frac12\theta_\nu(\widehat\vartheta\nu su)(\theta su)^3,\\
0&=\theta_\nu s_\nu s_\vartheta-\frac14s_\eta,
\quad 0=\theta_\nu s_\nu s_\vartheta(\theta su),
\quad 0=\theta_\nu s_\nu s_\vartheta(\theta su)^2.
\end{aligned}
\]
最後の群は $(\theta_\nu(\vartheta \nu x)^2,s,u)^2\ldots$ に対応する。

\noindent\emph{帰結：}

1) $\theta(\vartheta\nu x)s_\nu$、$(\widehat\vartheta\nu su)(\theta su)s_\nu$、$\theta_\nu(\widehat\vartheta\nu su)^2$、$(\widehat\vartheta\nu su)^2(\theta su)^2$ は還元子である。

2) 四次の形式 $(5,\lambda)$、$(6,\lambda)$ などは $s^2$ との折り畳みに入らない。

3) 形式列 $(\theta gu)^2$ について、式 (27.) と本節の還元を合わせると次の公式が得られる：
\[
\tag{29}
\begin{aligned}
\text{a)}\quad (\theta gu)^2&=\frac43\theta_\nu s_\nu+\frac23s_\nu s_\vartheta-\frac13s\theta_\nu^2-\frac19sH-\frac13f\,a_\nu^2(asu)^2+\frac13a_\nu^2(asu)^2,\\
\text{b)}\quad (\theta gu)^3&=-\frac13(\widehat\vartheta\nu su)^2(\theta su)-\theta_\nu(\widehat\vartheta\nu su)(\theta su)-\theta_\nu^2(\theta su),\\
\text{c)}\quad (\theta gu)^4&=2\theta_\nu^2(\theta su)^2-\frac13(Hsu)^2,\\
g\vartheta(\theta gu)&=(\widehat\vartheta\nu su)(\vartheta\nu x)s_\nu-\theta_\nu(\vartheta\nu x)(\widehat\vartheta\nu su),\\
g\vartheta(\theta gu)^2&=\frac23s\theta_\nu^3,
\quad g\vartheta(\theta gu)^3=\frac13\theta_\nu^3(\theta su),
\quad g\vartheta(\theta gu)^4=0.
\end{aligned}
\]

\begin{center}
\textbf{\S\ 22.\quad 九次の形式（系 $s_{III}$）。}
\end{center}

$s$ と五次の形式（系 III）および積 $\A\cdot\nu$ との折り畳み（\S\ 10）。

\noindent\emph{既約形式：}
\[
\begin{array}{ll}
\text{A)}& (k\nu x)^2s_\nu,
((k\nu x)^3,s,u),
(k\nu x)^3s_\nu,
((k\nu x)^2,s,u)^2,
K_\nu s_\nu,\\
& ((k\nu x)^3,s,u)^2,
((k\nu x)^2,s,u)s_\nu,
(K_\nu(k\nu x)^3,s,u),\\
& (K_\nu su)^2,\ (K_\nu su)^3,\ (K_\nu su)^4,
((k\nu x)^2,s,u)^3,
(K_\nu^2su)^4,\\[0.3em]
\text{B)}& (j\nu x)s_\nu,
((j\nu x)^2,s,u),
(j\nu x)^3,s,u,
((j\nu x)^2,s,u)^2,\\[0.3em]
\text{C)}& s_\nu(\A su),\quad (\A_\nu su),\\[0.3em]
\text{D)}& (aHu)s_\eta,
(a_\eta su),
((aHu),s,u)^2,
((aHu)^2,s,u),\\
& (aHu)^2s_\eta,
(a_\eta su)^2,
(a_\eta^2su),
((aHu),s,u)^3,
((aHu)^2,s,u)^2,\\
& ((aHu),s,u)^4,
(a_\eta su)^3,
(a_\eta(aHu),s,u)^2,
(a_\eta(aHu)^2,s,u),
(a_\eta(aHu),s,u)^3.
\end{array}
\]

\noindent\emph{還元：}

A) 還元は、\S\ 21 の還元から一回の折り畳みによって直接従う。形式列 $K_\nu s_\nu(Ksu)^2$ は還元子 $a_\nu s_\nu(asu)^2$ と $\theta_\nu s_\nu(\theta su)^2$ を因子にもつ。従って高次形式 $K_\nu^2(Ksu)^2$、$K_\nu^2(Ksu)^3$ は Ansatz
\[
0=a_\nu s_\nu(asu)^2(a\theta u)=K_\nu s_\nu(Ksu)^2=\theta_\nu s_\nu(\theta su)^2(a\theta u)
\]
に従って分解する。

B) 形式列 $(j\nu x)^2s_\nu$：
\[
(j\nu x)^2s_\nu=3a_\nu^2s_\nu(a\theta u)^2
\]
から得られる還元子 $a_\nu^2s_\nu$。
$((j\nu x)^3,su)^2$：
\[
((j\nu x)^3,s,u)^2=-6a_\nu^2s_\nu(a\theta u)^2(\theta su)
\]
から得られる還元子 $a_\nu^2s_\nu$。

C) 形式 $\A_\nu(\A su)^2$ と $s_\nu(\A su)^2$：式 (27.)a から得られる還元子 $g_\nu$。

形式 $\A_\nu s_\nu$：
1) $a_g a_\nu^2s_\nu=\frac23\A_\nu s_\nu$ から得られる還元子 $a_\nu^2s_\nu$；
2) 式 (27.)a により、式 $(\A s\widehat\nu x)^2$ の二重還元を通じて分解する。

従って高次形式 $a_\eta s_\eta$ が分解する。

D) 形式列 $(aHu)^2s_\eta(asu)$：形式列 $[a_\nu s_\nu(asu)^2]_{\nu_1x^2}$ の二重還元による還元子 $a_\nu s_\nu(asu)^2$；$(aHu)^2s_\eta(asu)^2$ は $a_\nu s_\nu(asu)^2$ と $a_\nu s_\nu(asu)^3$ から還元される。従って $(a_\eta,s,u)^4$ が還元される。

形式列 $a_\eta(aHu)s_\eta$：式 $(\A s\widehat\nu x)^3$ の二重還元による還元子 $a_\nu^2s_\nu$、$a_\eta^2s_\eta$。ただし $a_\eta^2s_\eta$ は、形式 $a_\eta(aHu)s_\eta$ の還元可能項 $a_\nu^2s_\nu(\nu\nu_1x)$ からの折り畳みでは生じない。

形式列 $(a_\eta^2su)^2$：
\[
a_\nu^2s_\nu s_{\nu_1}=-\frac12(a_\eta^2(Hsu))^2
\]
からの還元子 $a_\nu^2s_\nu$。
形式 $(a_\eta(aHu),s,u)$ は関数行列式に関するトランスヴェクションとして分解する。

得られる還元公式は
\[
\tag{30}
\begin{aligned}
\text{a)}\quad 0&=(aHu)^2(asu)s_\eta-a_\eta(asu)(Hsu)^2-a_\eta(aHu)(asu)(Hsu),\\
\text{b)}\quad 0&=(aHu)^2(asu)^2s_\eta-a_\eta(aHu)(asu)^2(Hsu),\\
\text{c)}\quad 0&=a_\eta(asu)^2(Hsu)^2,\\
0&=a_\eta s_\eta+\frac14a_\nu^2g-\frac14s\cdot a_\eta^2,\\
0&=a_\eta(aHu)s_\eta-\frac12a_\eta^2(Hsu),
\quad 0=a_\eta^2(Hsu)^2,\ldots,
\quad 0=a_\eta^2s_\eta.
\end{aligned}
\]

\noindent\emph{帰結：}

1) $(j\nu x)^2s_\nu$、$\A_\nu s_\nu$、$\A_\nu(\A su)^2$、従って $N_\nu(Nsu)^2$、$(aHu)^2(asu)s_\eta$、$a_\eta(aHu)s_\eta$、$a_\eta^2(Hsu)^2$ は還元子である。$a_\eta s_\eta$ は分解可能形式であり、すべての一回および二重の折り畳みで分解可能形式へ移る。

2) 五次の形式 $(5,\lambda)$、$(6,\lambda)$ などは $s^2$ との折り畳みに入らない。

\begin{center}\textbf{\S\ 23.\quad 十次の形式（系 $s_{III}$）。}\end{center}
$s$ と六次の形式（系 III）との折り畳み（\S\ 11）。

\noindent\emph{既約形式：}
\[
\begin{array}{rl}
\text{A)}& (\vartheta^2\nu x)^3s_\nu,
((\vartheta^2\nu x)^3,s,u),
((\vartheta^2\nu x)^3,s,u)^2,
(\vartheta^2\nu x)^3s_\nu s_\vartheta,
\theta_\nu(\vartheta^2\nu x)^3s_\vartheta,\\[0.3em]
\text{B)}& a_\nu(j\nu x)s_\nu,
(a_\nu(j\nu x),s,u)^2,
(a_\nu(j\nu x),s,u)^3,
(a_\nu(j\nu x),s,u)^4,
(a_\nu^2(j\nu x),s,u)^2,
(a_\nu^2(j\nu x),s,u)^3,\\[0.3em]
\end{array}
\]
\[
\begin{array}{r@{\quad}c}
\text{C)}&
\begin{array}{c|c|c|c}
\cdot&\cdot&((\theta Hu),s,u)^2,\ ((\theta Hu)^2,s,u)&\bigstar\\[0.35em]
\cdot&((\vartheta\eta x),s,u)^2,\ (\theta Hu)s_\eta,\ (\theta_\eta su)&((\vartheta\eta x),s,u)^3,\ (\theta_\eta su)^2&\\[0.35em]
((\vartheta\eta x)^2,s,u)&\cdot&(\theta_\eta^2su)&\\[0.35em]
\cdot&(\theta_\eta(\vartheta\eta x)^2,s,u)&\cdot&\cdot
\end{array}
\end{array}
\]
\[
\begin{array}{c@{\quad}|c|c}
\bigstar&((\theta Hu),s,u)^3,\ ((\theta Hu)^2,s,u)^2&((\theta Hu),s,u)^4\\[0.45em]
&(H_\vartheta(\theta Hu),s,u)^2,\ (\theta_\eta(\theta Hu),s,u)^2,\ (\theta_\eta(\theta Hu)^2,s,u)&(\theta_\eta su)^4,\ (\theta_\eta(\theta Hu),s,u)^3\\[0.45em]
&\cdot&(\theta_\eta^2(\theta Hu),s,u)^2
\end{array}
\]
\noindent\emph{還元：}
A) 形式 $((\vartheta^2\nu x)^3,s,u)^2$、$((\vartheta^2\nu x)^3,s,u)^3$ は分解する：
\[
(\vartheta\nu x)(\widehat{\vartheta\nu su})(\widehat{\vartheta\nu su})=(\vartheta\nu x)s_\vartheta s_\nu-(\vartheta\nu x)s_\nu s_\vartheta=-2\theta\cdot(\vartheta\nu x)s_\nu s_\vartheta+(\vartheta\nu x)\cdot s_\vartheta^2.
\]
残りの形式は、還元子を因子にもつか、あるいは分解された形式との一回の折り畳みである。

B) 還元子 $a_\nu^2s_\nu$ と $(\widehat{j\nu su})s_\nu$。

C) 1. 形式列 $(\theta Hu)^2s_\eta$：a) 還元子 $(aHu)^2(asu)s_\eta$；b) 形式列 $(\theta gu)^3g_\nu$ と $g_\nu(agu)^3(bgu)^2(abu)$ の二重還元。従って高次形式 $(\theta_\eta su)^3$、$(H_\vartheta(\theta Hu),s,u)^3$、$(a_\nu\cdot b_{\nu_1}^2,s,u)^4$［後者は $(H_\vartheta su)^4$ を置き換える形式］の還元が得られる。

2. $((\vartheta\eta x),s,u)^4$：還元子 $(a_\eta su)^4$。

3. $((\vartheta\eta x)^2,s,u)^3$：$(\widehat{\vartheta\eta su})^2(Hsu)$ からの還元子 $s_\vartheta^2s_\nu$。形式 $(\vartheta\eta x)s_\eta$、$(\vartheta\eta x)^2s_\eta$、$((\vartheta\eta x)^2,s,u)^2$、$\theta_\eta s_\eta$ は、分解された形式 $a_\eta s_\eta$ との折り畳みによって分解する。

4. 形式列 $H_\vartheta(\theta Hu)s_\eta$、$\theta_\eta(\theta Hu)s_\eta$、$H_\vartheta(\vartheta\eta x)s_\eta$、$\theta_\eta(\vartheta\eta x)s_\eta$：還元子 $a_\eta(aHu)s_\eta$ と $a_\eta^2s_\eta$。

形式列 $(\theta_\eta(\vartheta\eta x),s,u)^2$：還元子 $a_\eta^2(Hsu)^2$［ただし、折り畳みによって項 $a_\eta^2(\theta su)(Hsu)$ から生じる $(\theta_\eta^2(\theta Hu),s,u)^2$ は除く］。

5. $(H_\vartheta(\vartheta\eta x),s,u)$、$(H_\vartheta(\vartheta\eta x),s,u)^2$：$a_\eta(aHu)s_\eta(abu)$ と $g_\vartheta(\theta gu)g_\nu$ の二重還元。

形式列 $(H_\vartheta(\vartheta\eta x),s,u)^3$：還元子 $((\vartheta\nu x)^2,s,u)^2s_\nu$。

6. $(H_\vartheta(\theta Hu),s,u)$ は、\S\ 3. 2), c) に従って、分解された形式 $(\theta_\nu(\vartheta\nu x),s,u)$ との一回の折り畳みによって分解する。すなわち、低次の分解された形式 $(H_\vartheta su)^2$ の二重還元による：\footnote{記号 $\equiv$ は、より低次数の形式の積を省略したことを意味する。}
\begin{align*}
0&=\theta_\nu(\vartheta\nu_1)(\theta su)+\frac23\theta_\nu^2(\vartheta\nu_1x)(\theta su)+\theta_\nu(\vartheta\nu x)(\theta su)s_{\nu_1}\\
&\equiv -\frac12H_\vartheta(\theta Hu)(\theta su)+\theta_\eta(\theta su)(Hsu)+\frac12H_\vartheta(\theta su)(Hsu)\\
&\equiv -\frac12H_\vartheta(\theta Hu)(\theta su)+\theta_\eta(\theta su)(Hsu)+\frac12[H_\vartheta]_{sx}^2-\frac12\frac35H_\vartheta(\theta Hu)(\theta su)\\
&\equiv -\frac35H_\vartheta(\theta Hu)(\theta su).
\end{align*}
1. と 5. に従う二重還元により、次の公式が得られる：
\begin{equation}
\tag{31}
\begin{aligned}
0&=\theta_\eta(\theta su)(Hsu)^2+\frac14H_\vartheta(\theta Hu)(\theta su)^2\\
&\quad-\frac34\theta_\eta(\theta Hu)(\theta su)(Hsu)-\frac54\theta_\eta(\theta Hu)^2(\theta su)^3\\
&\quad-(asu)^3b_\eta(\theta Hu)^2-a_\nu(asu)^3b_\nu^2,\\
0&=H_\vartheta(\theta Hu)(\theta su)^3-7\theta_\eta(\theta Hu)(\theta su)^2(Hsu),\\
0&=a_\nu(asu)^2b_{\nu_1}^{2}(bsu)^2-\theta_\eta(\theta su)^2(Hsu)^2\\
&\quad+6\theta_\eta(\theta Hu)(\theta su)^2(Hsu),\\
0&\equiv (H_\vartheta(\vartheta\eta x),s,u),
&0&\equiv H_\vartheta(\widehat{\vartheta\eta su})(Hsu)-4\theta_\eta^2(Hsu).
\end{aligned}
\end{equation}
\noindent\emph{帰結：}
\begin{enumerate}
\item $(\theta Hu)^2s_\eta$、$H_\vartheta(\theta Hu)s_\eta$、$\theta_\eta(\theta Hu)s_\eta$、$H_\vartheta(\vartheta\eta x)s_\eta$、$\theta_\eta(\vartheta\eta x)s_\eta$、$(H_\vartheta(\vartheta\eta x),s,u)$、$(\theta_\eta(\vartheta\eta x),s,u)^2$ は還元子である。
\item $s^2$ は、6次の形式 $(5,\lambda)$ などとの折り畳みに入らない。
\end{enumerate}

\begin{center}\textbf{\S\ 24.\quad 十一次の形式（系 $s_{III}$）。}\end{center}
$s$ と七次の形式（系 III）との折り畳み（\S\ 12）。

\noindent\emph{既約形式：}
\[
\begin{array}{rl}
\end{array}
\]
\[
\text{A)}\quad
\begin{array}{c|c|c|c|c}
\cdot&\cdot&\cdot&((KHu)^2su)^2&\bigstar\\[0.35em]
\cdot&\cdot&(K_\eta su)^2&\cdot&\\[0.35em]
\cdot&\cdot&\cdot&\cdot&\cdot\\[0.35em]
((k\eta x)^3,s,u)&\cdot&\cdot&\cdot&\cdot\\[0.35em]
\cdot&(K_\eta(k\eta x)^3,s,u)&\cdot&\cdot&\cdot
\end{array}
\]
\[
\begin{array}{c@{\quad}|c|c}
\bigstar&((KHu)^2,s,u)^3&((KHu)^2,s,u)^4\\[0.45em]
&\cdot&(K_\eta(KHu)^2,s,u)^3
\end{array}
\]
\[
\begin{array}{rl}\text{B)}& ((j\eta x),s,u)^2,
\quad (H_jsu),
\quad ((j\eta x),s,u)^3,\\[0.3em]
\text{C)}& (\D_\eta su),\\[0.3em]
\text{D)}& (aLu)^2s_\eta,
\quad (a_\eta su)^2,
\quad ((aLu)^2,s,u)^2,
\quad ((aLu)^3,s,u)^2,
\quad ((aLu)^3,s,u)^3,
\quad (a_l(aLu)^2,s,u)^2.
\end{array}
\]
\noindent\emph{還元：}

A) 記号 $\theta-H$ における対応する形式との一回の折り畳みによる還元または分解。

$(K_\eta(KHu)^2,s,u)$ と $(K_\eta(KHu)^2,s,u)^2$：$K_\nu^2(Ksu)$、$K_\nu^2(Ksu)^2$ との折り畳みによって、\S\ 3. 2), a) に従って分解する。

$(K_\eta(KHu),s,u)^4\equiv(a_\nu^2\cdot\theta_{\nu_1},s,u)^4$：分解された形式 $K_\nu^2(Ksu)^3$ から、\S\ 3. 2), b) に従って分解する。

B) 形式列 $(j\eta x)s_\eta$：還元子 $(j\nu x)^2s_\nu$。

形式 $H_j(\widehat{j\eta su})(Hsu)$：還元子 $(j\nu x)(\widehat{j\nu su})^2$。

形式 $H_j(j\eta x)(Hsu)$：分解された形式 $((j\eta x)^2,s,u)^2$ を還元子 $(j\nu x)^2s_\nu$（式 (18.)a）で還元することにより分解する：
\[
0=-2(j\nu x)^2s_\nu s_{\nu_1}=[(j\eta x)^2]_{\widehat{su}^2}+H_j(j\eta x)(Hsu).
\]
C) 還元子 $\D_\nu s_\nu$ と $\D_\nu(\D su)^2$。

D) 記号 $f-H$ における対応する形式との一回の折り畳みによる還元および分解。

E) (30.)c から、分解された形式列 $(a_\sigma su)^2$ の還元が得られる：
\begin{align*}
0&=a_\eta(Hsu)^2(as\widehat{\nu x})^2
=a_\eta(Hsu)^2(a\widehat{\nu\eta}u)^2+2a_\eta(a\widehat{\nu\eta}u)(\widehat{\nu\eta}su)(Hsu)^2
=\frac13a_\sigma(asu)^2,\\
0&=a_\eta a_\nu(Hsu)^2(as\widehat{\nu x})(asu)
=a_\eta(a\widehat{\nu\eta}u)^2(Hsu)^2(asu)+a_\eta(a\widehat{\nu\eta}u)(\widehat{\nu\eta}su)(Hsu)^2(asu)
=\frac13a_\sigma(asu)^3.
\end{align*}
\noindent\emph{帰結：} $s^2$ は七次の形式との折り畳みに入らない。

\begin{center}\textbf{\S\ 25.\quad 十二次、十三次、十四次、十五次の形式（系 $s_{III}$）。}\end{center}
$s$ と八、九、十、十一次の形式（系 III）との折り畳み（\S\S\ 13--16）。

\noindent\emph{既約形式：}

\noindent 12次。
\[
\begin{array}{rl}
\text{A)}& ((\vartheta^2\eta x)^4,s,u),\quad \theta_\eta(\vartheta^2\eta x)^3s_\vartheta,\\
\text{B)}& ((aHu)H_j,s,u)^2,
\quad ((aH\cdot u)H_j,s,u)^3,\\
\text{C)}& (\theta_\nu su)^2,
\quad ((\theta Lu)^2,s,u)^2,
\quad ((\theta Lu)^3,s,u),
\quad ((\theta Lu)^2,s,u)^3,
\quad (\theta_l(\theta Lu)^2,s,u)^2.
\end{array}
\]
\noindent\emph{還元：} 還元子または分解された形式との直接の折り畳みから生じる還元は、もはや述べない。

B) $((aHu)H_j,s,u)$ と $(a_\eta(aHu)H_j,s,u)$ は、関数行列式に関するトランスヴェクションとして分解する。

$(a_\eta(aHu)H_j,s,u)^2$ の分解：分解された形式 $[\theta_\nu\cdot\theta_{\nu_1}^{\prime 3}]_{\widehat{su}^3}$（\S\ 13. C）の二重還元：
\begin{align*}
0&\equiv[\theta_{\nu_1}^{\prime 3}\cdot\theta_\nu]_{\widehat{su}^3}
\equiv-\frac34\theta_\nu(\theta su)^2(\theta\theta'u)\theta_{\nu_1}^{\prime 3}\\
&\equiv-\frac98(\widehat{\theta\theta'}\eta\widehat{su})(\theta\theta'u)(\theta Hu)(\theta'Hu)\theta_\eta'(\theta su)
\equiv-\frac38a_\eta(aHu)H_j(asu)^2.
\end{align*}

C) 公式 (31.) から、分解された形式
\[
 [L_\vartheta(\theta Lu)]_{su^3}\quad\hbox{および}\quad [\theta_l(\theta Lu)]_{s^2u^3}
\]
の還元、すなわち積 $[\theta_\nu^2\cdot H]_{su^3}$ と $[a_\nu^2\cdot(bHu)^2]_{su^3}$ の還元が従う：
\begin{equation}
\tag{32}
\begin{aligned}
0&\equiv\theta_\nu^2(\theta su)^2(Hsu),\\
0&\equiv [L_\vartheta(\theta Lu)]_{su^4}-7[\theta_l(\theta Lu)]_{su^4},\\
0&\equiv a_\nu^2(asu)^2(bHu)^2(bsu)\\
&\quad+6\theta_l(\theta Lu)^2(\theta su)(Lsu),\\
0&\equiv\theta_\nu^3(\theta su)(Hsu)^2\\
&\quad\hbox{から }0\equiv\theta_l^2(\theta Lu)(\theta su)(Lsu)^2\\
&\quad(\hbox{還元子 }a_\eta^2(Hsu)^2).
\end{aligned}
\end{equation}

\noindent 13次。
\[
\text{A)}\ ((KLu)^3,s,u)^2,\ ((KLu)^3,s,u)^3,\qquad
\text{B)}\ (L_jsu)^2,
\qquad
\text{C)}\ ((aH^2u)^3,s,u)^2.
\]
\noindent 14次。
\[
\text{A)}\ ((aLu)^2L_j,s,u)^2,
\qquad
\text{B)}\ ((\theta H^2u)^3,s,u)^2.
\]
\noindent 15次。
\[
 ((KH^2u)^4,s,u)^2.
\]
\noindent\emph{還元：} $(K_l(KLu)^3,s,u)^2$、$(H_\vartheta(\theta H^2u)^3,s,u)$、$((K,H\cdot L,u)^5,s,u)$ は、\S\ 3. 2), c) に従って分解する。すなわち、分解された形式 $(K_l(KLu)^2,s,u)^3$、$(H_\vartheta(\theta H'u)^2,s,u)^2$、$((K,H\cdot L,u)^4,s,u)^2$ を同時に考慮する。

\noindent\emph{帰結：} $s^2$ は系 III の個々の形式との折り畳みに入らない。

\begin{center}\textbf{\S\ 26.\quad 系 $s_{IV}$。}\end{center}
\noindent 1) \emph{$s$ と系 I--III との折り畳み。}

\noindent\emph{還元：} 最後の諸節の還元（$\cdot(a_\nu^2s_\nu(aHu)^2(asu)s_\eta$ など）により、形式 $(0,\lambda)$、$(1,\lambda)$、$(2,\lambda)$ は I と III の折り畳みを同時には許さない。従って \S\ 18 により、形式 $f,\D,N$ は系 $s_{IV}$ を形成するための折り畳みの積には入らない。

形式 $j$ は折り畳みに入らない。なぜなら、同じ還元によって、$j$ に反変な折り畳み
\[
(K_\nu(k\nu x)^3,s,u),\qquad ((\vartheta^2\eta x)^4,s,u)\quad\hbox{など}
\]
は、$j$ に共変な折り畳み
\[
K_\nu(k\nu x)^2s_k,
\qquad (\vartheta^2\eta x)^3s_\vartheta\quad\hbox{など}
\]
に対応するからである。

\noindent 2) \emph{$s$ と系 II--III との折り畳み}、すなわち $s$ と $H\cdot(1,\lambda)$、$H\cdot(2,\lambda)$、$H\cdot(3,\lambda)$ との折り畳み。

\noindent\emph{既約形式：}
\[
\begin{array}{rl}
\text{A)}&(a_\nu\cdot H,s,u)^4,
((aHu)^2\cdot H,s,u)^3,
((aHu)^2H,s,u)^4,\\
&((aLu)^2\cdot H,s,u)^4,
((aLu)^3\cdot H,s,u)^3,
((aH^2u)^3\cdot H,s,u)^4,\\[0.45em]
\text{B)}&(a_\nu^2\cdot H,s,u)^3,
(a_\eta(aHu)^2\cdot H,s,u)^3,
(a_l(aLu)^2\cdot H,s,u)^4,\\[0.45em]
\text{C)}&(\theta_\nu\cdot H,s,u)^4,
((\theta Hu)^2\cdot H,s,u)^3,
((\theta Hu)^2H,s,u)^4,\\
&((\theta Lu)^2\cdot H,s,u)^4,
((\theta Lu)^3\cdot H,s,u)^3,
((\theta H^2u)^3\cdot H,s,u)^4,\\[0.45em]
\text{D)}&(\theta_\eta(\theta Hu)^2\cdot H,s,u)^3,
(\theta_l(\theta Lu)^2\cdot H,s,u)^4,\\[0.45em]
\text{E)}&((KLu)^3\cdot H,s,u)^4,
((KH^2u)^4\cdot H,s,u)^4,\\[0.45em]
\text{F)}&((j\nu x)^2\cdot H,s,u)^3,
((j\nu x)^2\cdot H,s,u)^4,
((j\eta x)H,s,u)^4,\\
&(H_j\cdot H,s,u)^3,
(L_j\cdot H,s,u)^4.
\end{array}
\]
\noindent\emph{還元：} $\nu$ は $j$ と同様に折り畳みに入らない。

B) および D)。$(a_\nu^2\cdot H,s,u)^4$ と $(\theta_\nu^2\cdot H,s,u)^4$ は、$(gHu)^2=-\frac13s\cdot\sigma$ を考慮して、式 (27.)c) と (29.)c) により分解する：
\[
(a_\nu^2\cdot H,s,u)^4\equiv\frac13(asu)^4\cdot\sigma,
\qquad
(\theta_\nu^2\cdot H,s,u)^4\equiv-\frac16(\theta su)^4\cdot\sigma;
\]
従って $(a_\eta(aHu)\cdot H,s,u)^4$、$(\theta_\eta(\theta Hu)\cdot H,s,u)^4$ が分解する。

$(\theta_\nu^2\cdot H,s,u)^3$、$(\theta_\nu^3\cdot H,s,u)^3$、従って $(\theta_\eta^2(\theta Hu)\cdot H,s,u)^4$ は、\S\ 25、12次の形式 C) に従って分解する。

\noindent 3) \emph{$s$ と系 III--III との折り畳み}、すなわち $s$ と系 III の形式
\[
(3,\lambda)\cdot(3,\lambda),\quad (3,\lambda)\cdot(2,\lambda),\quad (3,\lambda)\cdot(1,\lambda),\quad
(2,\lambda)\cdot(2,\lambda),\quad (2,\lambda)\cdot(1,\lambda),\quad (1,\lambda)\cdot(1,\lambda)
\]
との折り畳み。
\noindent\emph{既約形式：}
\[
\begin{array}{rl}
\text{A)}&(a_\nu^2\cdot a_\nu^2,s,u)^3,
(a_\nu^2\cdot a_\nu^2,s,u)^4,
(a^2\cdot a_\eta(aHu)^2,s,u)^3,\\[0.3em]
\text{B)}&(a_\nu\cdot H_j,s,u)^4,
((aHu)^2\cdot H_j,s,u)^3,
((aLu)^2\cdot H_j,s,u)^4,
((aLu)^3\cdot H_j,s,u)^2,
((aH^2u)^3\cdot H_j,s,u)^3.\cdot
\end{array}
\]
\noindent\emph{還元：}

記号 $\nu$ と $f$ における同じ位数に対して、積 II--III は積 III--III より高いものと見なす。

還元は、積の間の関係（シジジー）から部分的に生じ、また積を対応する分解形式で置き換え、これらの形式と $s$ との折り畳みを還元することから部分的に生じる。反変形式を含む積は、積へ移るので常に省略される。

A) $f$ が二次、記号 $\nu$ は任意の位数。\S\ 11、式 (12.) および同様の計算により、
\begin{align*}
1)\quad 0&=a_\nu\cdot b_{\nu_1}+\frac12f\cdot(aHu)^2-\frac14\theta\cdot H+\frac12u_xH_\vartheta-\frac14u_x^2\cdot H_\vartheta^2,\\
2)\quad 0&=a_\nu\cdot b_{\nu_1}^2+f\cdot a_\eta(aHu)^2-\theta_\eta-\frac12H_\vartheta,\\
3)\quad 0&=a_\nu^2\cdot b_{\nu_1}^2-H_\vartheta^2+2\theta_\eta^2.
\end{align*}
従って
\begin{align*}
4)\quad 0&=a_\nu^2\cdot b_{\nu_1}^2(j\nu_1x),\qquad\text{（3) から）}\\
5)\quad L_\vartheta(\theta Lu)&=-a_\nu\cdot b_\eta(bHu)^2+\frac34a_\nu^2\cdot(bHu)^2+\frac34\theta_\nu^2\cdot H,\qquad\text{（2) から）}\\
6)\quad L_\vartheta(\theta Lu)&=-6a_\nu\cdot b_\eta(bHu)^2+2a_\nu^2\cdot(bHu)^2-\frac12\theta_\nu^2\cdot H,\qquad\text{（1) から）}\\
7)\quad 0&=a_\nu\cdot(bHu)^2+f\cdot(aLu)^3+\theta_\nu\cdot H,\qquad\text{（1) から）}\\
8)\quad 0&=a_\nu\cdot(bLu)^3+\frac34(\theta Hu)^2\cdot H,\\
9)\quad 0&=(aHu)^2\cdot(bHu)^2+2(\theta Hu)^2\cdot H,\qquad\text{（7) から）}\\
10)\quad 0&=a_\eta(aHu)^2\cdot b_\eta(bHu)^2,\qquad\text{（5) と 6) から）}\\
11)\quad 0&= [a_\nu^2\cdot b_\eta(bHu)]_{\widehat{su^4}}.\qquad\text{（式 (30.)b から）}
\end{align*}
\[
(H_\vartheta su)^4,
\qquad (L_\vartheta(\theta Lu),s,u)^3,
\qquad (L_\vartheta(\theta Lu),s,u)^4
\]
の還元（公式 (31.) と (32.)）は、公式 1) から 11) と合わせて、記号 $f$ に関して二次、$\nu$ に関して任意の積との $s$ のすべての折り畳みを還元する。

B) 記号 $f,\theta,j$ のうち二つずつの積、$\nu$ は任意の位数。式 (17.)、(18.)、(20.) および直接計算により、次の関係式が成り立つ：
\begin{align*}
1)\quad 0&=a_\nu\cdot\theta_{\nu_1}+\frac14\theta\cdot(aHu)^2+\frac14f\cdot(\theta Hu)^2,\\
2)\quad 0&=\theta_\nu\cdot\theta_{\nu_1}+\frac12\theta\cdot(\theta Hu)^2.\\
&\quad\text{（A1 と同様の直接計算）}
\end{align*}
これより
\begin{align*}
3)\quad 0&=3a_\nu\cdot(\theta Hu)^2+\theta\cdot(aLu)^3+2f\cdot(\theta Lu)^3,\\
4)\quad 0&=3\theta_\nu\cdot(aHu)^2+f\cdot(\theta Lu)^3+2\theta\cdot(aLu)^3,\\
&\quad\text{（1) と A6 から）}\\
5)\quad 0&=a_\nu\cdot(\theta Lu)^3,
&0&=\theta_\nu\cdot(aLu)^3,
&0&=(aHu)^2\cdot(\theta Hu)^2,\\
&\quad\text{（3), 4), A8 から）}\\
6)\quad 0&=a_\nu\cdot\theta_\nu^2+\theta_\nu\cdot a_\nu^2+f\cdot\theta_\eta(\theta Hu)^2+\theta\cdot a_\eta(aHu)^2,\\
&\quad\text{（式 17）}\\
7)\quad 0&=\theta_\nu\cdot\theta_\nu^2+\theta\cdot\theta_\eta(\theta Hu)^2+\frac16f\cdot H_j,\\
&\quad\text{［6 の導出と同様］}\\
8)\quad 0&=a_\nu\cdot(j\nu x)^2+f\cdot H_j,\\
&\quad\text{（6 から）}\\
9)\quad 0&=(aHu)^2\cdot(j\nu x)^2-2a_\nu\cdot H_j,\\
&\quad\text{（8 から）}\\
10)\quad 0&=\theta_\nu\cdot(j\nu x)^2,
&0&=\theta\cdot H_j,\\
&\quad\text{（7 と 8 から）}\\
11)\quad 0&=a_\nu\cdot\theta_\nu^3-\frac34(j\eta x)^2,\\
12)\quad 0&=a_\nu^2\cdot\theta_\nu^2-\frac13(j\eta x)^2-\frac13(\D Hu)^2,\\
&\quad\text{（式 18）}\\
13)\quad 0&=a_\nu^2\cdot\theta_\nu^3-\frac12a_\sigma,\\
&\quad\text{（式 20）}\\
14)\quad 0&=\theta_\nu\cdot\theta_\nu^3,
&0&=a_\eta H_j.\\
&\quad\text{（§13 C）}
\end{align*}
公式は、積が偏極可能になるところまで進められている。すなわち、折り畳みによって新しい積が一つだけ生じるのは、a) その積自身への折り畳みが還元可能になり、b) 折り畳みによって生じるその他の積が還元可能になるか、または反変形式に至るからである（\S\ 3. 2), a), b) 参照）。

公式 1) から 5) は、$a_\nu\theta_\nu$ から折り畳みによって生じる積との $s$ のすべての折り畳みの還元を与える。公式 6) から 10) は、$a_\nu\theta_\nu^2$ から折り畳みによって生じる積の還元を与える。\S\ 24 A) により、公式 6) から 10) を考慮すると、$a_\nu^2\theta_\nu$ から生じる積との $s$ の折り畳みは分解する。

すなわち
\begin{align*}
0&\equiv a_\nu^2(asu)^2\theta_{\nu_1}(\theta su)^2,
&0&\equiv a_\nu^2(asu)\theta_{\nu_1}(\theta su)^3-K_\eta(KHu)^2(Ksu)^3,\\
0&\equiv a_\nu^2(asu)^2\theta_{\nu_1}(\theta su),
&0&\equiv a_\nu^2(asu)\theta_{\nu_1}(\theta su)^2,\\
0&\equiv a_\nu^2(asu)\theta_{\nu_1}(\theta su).
\end{align*}
還元子
\[
 ((j\eta x)^2,s,u)^3\ [\hbox{由来 }(\widehat{\vartheta\eta su})^2(Hsu),\ \S\ 23.\ C3],
\]
\[
(H_j(j\eta x),s,u)^2\ [\S\ 24,\ B],
\quad ((\D Hu)^2,s,u)^2\ [\S\ 24,\ C],
\quad (a_\sigma su)^2\ [\S\ 24,\ E]
\]
は、公式 11) から 14) と合わせて、$a_\nu\theta_{\nu_1}^3$、$a_\nu^2\theta_{\nu_1}^2$、$a_\nu^2\theta_{\nu_1}^3$ から生じるすべての積の還元を与える。

以上の計算は次の\emph{結論}にまとめられる：

法 $(\R,t)$ に関する相対的完全系は、以下の 331 個の形式および部分系からなり、それらを続いて表としても整理して掲げる。
\clearpage
\thispagestyle{plain}
\begin{center}\bfseries 表 I（90頁参照）\end{center}
\vspace{0.25em}
\begingroup
\setlength{\tabcolsep}{2pt}
\renewcommand{\arraystretch}{1.35}
\tiny
\begin{center}
\begin{adjustbox}{max width=0.98\linewidth,max totalheight=0.72\textheight,center}
\begin{tabular}{|c|c|c|c|c|c|c|c|c|c|c|c|c|c|c|c|c|c|c|c|c|c|c|}
\hline
{} & 0 & 1 & 2 & 3 & 4 & 5 & 6 & 7 & 8 & 9 & 10 & 11 & 12 & 13 & 14 & 15 & 16 & 17 & 18 & 19 & 20 & 21 \\ \hline
0 & \mcell{i} &  &  &  & \mcell{f} &  & \mcell{\A} &  & \mcell{K_{\nu}(k\nu x)^3} &  & \mcell{K_{\eta}(k\eta x)^3} &  &  &  & \mcell{(\vartheta^2\eta x)^4} & \mcell{H_k(k\cdot\vartheta,\eta x)^4} &  & \mcell{\theta_l(\vartheta\cdot k,lx)^5} &  &  &  & \mcell{(\vartheta^3lx)^6} \\ \hline
1 &  & \mcell{u_x} &  &  &  & \mcell{\theta_\nu(\vartheta\nu x)^2} &  & \mcell{\theta_\eta(\vartheta\eta x)^2} & \mcell{N} & \mcell{(k\nu x)^3} & \mcell{\theta_\nu(\vartheta^2\nu x)^3} & \mcell{(k\eta x)^3} & \mcell{H_{\vartheta}(\vartheta^2\eta x)^3\\ \theta_\eta(\vartheta^2\eta x)^3} &  & \mcell{\theta_l(\vartheta^2lx)^4} &  & \mcell{(\vartheta k,\eta,x)^4} &  & \mcell{(\vartheta k,l,x)^5} &  &  &  \\ \hline
2 &  &  & \mcell{a_\nu^2} &  & \mcell{\theta\\ a_\eta^2} &  & \mcell{(\vartheta\nu x)^2} & \mcell{K_\nu(k\nu x)^2} & \mcell{(\vartheta\eta x)^2} & \mcell{H_k(k\eta x)^2\\K_\eta(k\eta x)^2} &  & \mcell{(\vartheta^2\nu x)^3\\ K_l(klx)^3} &  & \mcell{(\vartheta^2\eta x)^3} &  & \mcell{(\vartheta^2lx)^4} &  &  &  &  &  &  \\ \hline
3 &  & \mcell{\theta_\nu^3} &  & \mcell{a_\nu} & \mcell{\theta_\nu(\vartheta\nu x)} & \mcell{\A_\nu\\a_\eta} & \mcell{K\\H_{\vartheta}(\vartheta\eta x)\\\theta_\eta(\vartheta\eta x)} & \mcell{\A_\eta} & \mcell{(k\nu x)^2\\\theta_l(\vartheta lx)^2} &  & \mcell{(k\eta x)^2} &  & \mcell{(klx)^3} &  &  &  &  &  &  &  &  &  \\ \hline
4 & \mcell{\nu} &  & \mcell{H\\\theta_\nu^2} & \mcell{(j\nu x)^3\\a_\eta(aHu)} & \mcell{\theta_\eta^2} & \mcell{(\vartheta\nu x)\\a_l^2} &  & \mcell{N_\nu\\(\vartheta\eta x)} &  & \mcell{H_n\\N_\eta\\(\vartheta lx)^2} &  &  &  &  &  &  &  &  &  &  &  &  \\ \hline
5 &  & \mcell{a_\eta(aHu)^2} & \mcell{\theta_\eta^2(\theta Hu)} & \mcell{\theta_\nu} & \mcell{K_\nu^2\\(aHu)} & \mcell{\theta_\eta} & \mcell{K_\eta^2\\(\A Hu)\\a_l} &  & \mcell{\A_l} &  &  &  &  &  &  &  &  &  &  &  &  &  \\ \hline
6 & \mcell{j\\\sigma} &  & \mcell{(j\nu x)^2\\(aHu)^2} & \mcell{L\\a_\nu^2(j\nu x)\\H_{\vartheta}(\theta Hu)\\\theta_\eta(\theta Hu)} &  & \mcell{K_\nu\\\theta_l^2} &  & \mcell{K_\eta} &  &  &  &  &  &  &  &  &  &  &  &  &  &  \\ \hline
7 &  & \mcell{\theta_\eta(\theta Hu)^2} & \mcell{H_j(j\eta x)\\a_l(aLu)^2} &  & \mcell{a_\nu(j\nu x)\\(\theta Hu)} &  & \mcell{\A_\nu(j\nu x)\\\theta_l} & \mcell{K_l^2} &  &  &  &  &  &  &  &  &  &  &  &  &  &  \\ \hline
8 & \mcell{H_j^2\\a_l(aLu)^3} & \mcell{(\nu\sigma x)\\(j\nu x)} & \mcell{(\theta Hu)^2} & \mcell{K_\eta(KHu)^2\\(j\eta x)\\(aLu)^2} &  &  &  &  &  &  &  &  &  &  &  &  &  &  &  &  &  &  \\ \hline
9 &  & \mcell{H_j\\(aLu)^3} & \mcell{a_\eta(aHu)H_j\\L_{\vartheta}(\theta Lu)^2\\\theta_l(\theta Lu)^2} &  & \mcell{(KHu)} &  &  &  &  &  &  &  &  &  &  &  &  &  &  &  &  &  \\ \hline
10 & \mcell{\theta_l(\theta Lu)^3} & \mcell{L_j^2\\(j\sigma x)\\a_\eta(aH^2u)^3} &  & \mcell{H_j(aHu)\\(\theta Lu)^2} &  &  &  &  &  &  &  &  &  &  &  &  &  &  &  &  &  &  \\ \hline
11 &  & \mcell{(\theta Lu)^3} & \mcell{K_l(KLu)^3\\L_j\\(aH^2u)^3} &  &  &  &  &  &  &  &  &  &  &  &  &  &  &  &  &  &  &  \\ \hline
12 & \mcell{(aH^2u)^4} & \mcell{a_l(aLu)^2L_j\\H_{\vartheta}(\theta H^2u)^3\\\theta_\eta(\theta H^2u)^3} &  & \mcell{(KLu)^3} &  &  &  &  &  &  &  &  &  &  &  &  &  &  &  &  &  &  \\ \hline
13 & \mcell{(\nu\sigma j)} &  & \mcell{(aLu)^2L_j\\(\theta H^2u)^3} &  &  &  &  &  &  &  &  &  &  &  &  &  &  &  &  &  &  &  \\ \hline
14 & \mcell{(\theta H^2u)^4} & \mcell{K_\eta(KH^2u)^4\\(a,H\!\cdot L,u)^4} &  &  &  &  &  &  &  &  &  &  &  &  &  &  &  &  &  &  &  &  \\ \hline
15 & \mcell{L_\vartheta(\theta,H\!\cdot L,u)^4} &  & \mcell{(KH^2u)^4} &  &  &  &  &  &  &  &  &  &  &  &  &  &  &  &  &  &  &  \\ \hline
16 &  & \mcell{(\theta,H\!\cdot L,u)^4} &  &  &  &  &  &  &  &  &  &  &  &  &  &  &  &  &  &  &  &  \\ \hline
17 & \mcell{K_\eta(K,H\!\cdot L,u)^5} &  &  &  &  &  &  &  &  &  &  &  &  &  &  &  &  &  &  &  &  &  \\ \hline
18 &  & \mcell{(K,H\!\cdot L,u)^5} &  &  &  &  &  &  &  &  &  &  &  &  &  &  &  &  &  &  &  &  \\ \hline
19 &  &  &  &  &  &  &  &  &  &  &  &  &  &  &  &  &  &  &  &  &  &  \\ \hline
20 &  &  &  &  &  &  &  &  &  &  &  &  &  &  &  &  &  &  &  &  &  &  \\ \hline
21 & \mcell{(KH^3u)^6} &  &  &  &  &  &  &  &  &  &  &  &  &  &  &  &  &  &  &  &  &  \\ \hline
\end{tabular}
\end{adjustbox}
\end{center}
\vspace{0.2em}
\begin{center}\footnotesize （最上段の横列の数字は $x$ に関する次数を、最初の縦列の数字は $u$ に関する次数を表す。）\end{center}
\endgroup

\clearpage
\thispagestyle{plain}
\begin{center}\bfseries 表 II（90頁参照）\end{center}
\vspace{0.25em}
\begingroup
\setlength{\tabcolsep}{2pt}
\renewcommand{\arraystretch}{1.35}
\tiny
\begin{center}
\begin{adjustbox}{max width=0.98\linewidth,max totalheight=0.72\textheight,center}
\begin{tabular}{|c|c|c|c|c|c|c|c|c|c|c|c|c|c|c|c|c|c|}
\hline
{} & 0 & 1 & 2 & 3 & 4 & 5 & 6 & 7 & 8 & 9 & 10 & 11 & 12 & 13 & 14 & 15 & 16 \\ \hline
0 & \mcell{J} &  &  &  & \mcell{s} &  & \mcell{s_\vartheta^2} &  &  &  &  & \mcell{s_n} & \mcell{(k\nu x)^3s_\nu} & \mcell{(\vartheta^2\nu x)^3s_\nu s_\vartheta\\\theta_\nu(\vartheta^2\nu x)^3s_\vartheta} &  & \mcell{\theta_\eta(\vartheta^2\eta x)^3s_\vartheta} &  \\ \hline
1 &  &  &  &  &  &  & \mcell{(asu)} & \mcell{s_\vartheta} & \mcell{s_k^2\\(\Delta su)} & \mcell{s_n(Nsu)\\(\vartheta\nu x)^2s_\nu} & \mcell{((k\nu x)^3,s,u)s_\nu\\(K_\nu(k\nu x)^3,s,u)} &  & \mcell{(K_\eta(k\eta x)^3,s,u)} &  & \mcell{(\vartheta^2\nu x)^3s_\nu} &  & \mcell{((\vartheta^2\eta x)^4,s,u)} \\ \hline
2 &  &  &  &  & \mcell{(asu)^2} & \mcell{s_\vartheta(\theta su)} & \mcell{(\Delta su)^2\\a_\nu s_\nu} & \mcell{((\vartheta\nu x)^2,s,u)s_\nu\\(\theta_\nu(\vartheta\nu x)^2,s,u)} &  & \mcell{s_k\\(\theta_\eta(\vartheta\eta x)^2,s,u)} &  & \mcell{(k\nu x)^2s_\nu\\((k\nu x)^3,s,u)} &  & \mcell{((k\eta x)^3,s,u)} &  &  &  \\ \hline
3 &  &  & \mcell{(asu)^3} & \mcell{s_\vartheta(\theta su)^2\\s_\nu} & \mcell{a_\nu s_\nu(asu)\\a_\nu^2(asu)} & \mcell{s_\eta} & \mcell{(\theta su)\\(a_\eta^2su)} & \mcell{s_k(Ksu)} & \mcell{(Nsu)^2\\(\vartheta\nu x)s_\nu\\((\vartheta\nu x)^2,s,u)} & \mcell{((k\nu x)^3,s,u)^2} & \mcell{((\vartheta\eta x)^2,s,u)} &  &  &  &  &  &  \\ \hline
4 & \mcell{(asu)^4} & \mcell{s_\vartheta(\theta su)^3} & \mcell{a_\nu^2(asu)^2} & \mcell{(\theta_\nu^3su)} & \mcell{(\theta su)^2} & \mcell{s_k(Ksu)^2\\a_\nu(asu)\\s_\nu(asu)} & \mcell{((\vartheta\nu x)^2,s,u)^2\\\theta_\nu s_\nu} & \mcell{s_\nu(\Delta su)\\(\Delta_\nu su)\\(aHu)s_\eta\\(a_\eta su)} &  & \mcell{(\Delta_\eta su)} &  &  &  &  &  &  &  \\ \hline
5 &  &  & \mcell{(\theta su)^3} & \mcell{s_k(Ksu)^3,s_j,\\s_\vartheta(\theta s^\prime u)^4\\s_\nu(asu)^2\\a_\nu(asu)^2} & \mcell{(Hsu)\\((\vartheta\nu x)^2,s,u)^3\\(\theta_\nu(\vartheta\nu x),s,u)^2\\(\theta_\nu^2su)} & \mcell{((j\nu x)^3,s,u)\\(aHu)^2s_\eta\\(a_\eta su)^2} & \mcell{(Ksu)^2\\s_l\\(\theta_\eta^2su)} &  & \mcell{((k\nu x)^2,s,u)^2\\K_\nu s_\nu} &  &  &  &  &  &  &  &  \\ \hline
6 & \mcell{(\theta su)^4} & \mcell{s_\nu(asu)^3\\a_\nu(asu)^3} & \mcell{(Hsu)^2\\(\theta_\nu(\vartheta\nu x),s,u)^3\\(\theta_\nu^2su)^2} & \mcell{(a_\eta su)^3\\(a_\eta(aHu),s,u)^2\\(a_\eta(aHu)^2,s,u)} & \mcell{(Ksu)^3} & \mcell{s_j(asu)\\((\vartheta\nu x),s,u)^2\\(\theta_\nu su)} & \mcell{((k\nu x)^2,s,u)^3} & \mcell{s_j(\Delta su)\\a_\nu(j\nu x)s_\nu\\((\vartheta\eta x),s,u)^2\\(\theta Hu)s_\eta\\(\theta_\eta su)} &  &  &  &  &  &  &  &  &  \\ \hline
7 & \mcell{(\theta_\nu(\vartheta\nu x),s,u)^4} & \mcell{(a_\eta(aHu),s,u)^3} & \mcell{(Ksu)^4\\(\theta_\eta^2(\theta Hu),s,u)^2\\(a_\nu^2\cdot a_\nu^2,s,u)^3} & \mcell{s_j(asu)^2\\s_k(Ks^2u)^5\\((\vartheta\nu x),s,u)^3\\(\theta_\nu su)^2} & \mcell{(j\nu x)s_\nu\\((j\nu x)^2,s,u)\\((aHu),s,u)^2\\((aHu)^2,s,u)} & \mcell{((\vartheta\eta x),s,u)^3\\(\theta_\eta su)^2} & \mcell{(aLu)^2s_l\\(a_lsu)^2} &  &  &  &  &  &  &  &  &  &  \\ \hline

8 & \mcell{(a_\nu^2\cdot a_\nu^2,s,u)^4} & \mcell{s_j(asu)^3\\s_k(Ks^2u)^6\\((\vartheta\nu x),s,u)^4\\(\theta_\nu su)^3} & \mcell{((j\nu x)^2,s,u)^2\\((aHu),s,u)^3\\((aHu)^2,s,u)^2} & \mcell{(Lsu)^2\\(a_\nu^2(j\nu x),s,u)^2\\(H_\vartheta(\theta Hu),s,u)^2\\(\theta_\eta(\theta Hu),s,u)^2\\(\theta_\eta(\theta Hu)^2,s,u)} & \mcell{(a_lsu)^3} & \mcell{(K_\nu su)^2} &  & \mcell{(K_\eta su)^2} &  &  &  &  &  &  &  &  &  \\ \hline
9 & \mcell{(K_\nu^2su)^4\\((aHu),s,u)^4} & \mcell{(Lsu)^3\\(a_\nu^2(j\nu x),s,u)^3\\(\theta_\eta su)^4\\(\theta_\eta(\theta Hu),s,u)^3} & \mcell{(Ks^2u)^6\\(a_l(aLu)^2,s,u)^2\\(a_\nu^2\cdot H,s,u)^3} & \mcell{(K_\nu su)^3} & \mcell{(a_\nu(j\nu x),s,u)^2\\((\theta Hu),s,u)^2\\((\theta Hu)^2,s,u)} &  & \mcell{(\theta_lsu)^2} &  &  &  &  &  &  &  &  &  &  \\ \hline
10 &  & \mcell{(K_\nu su)^4\\(a_\nu^2\cdot a_\eta(aHu)^2,s,u)^3} & \mcell{(a_\nu(j\nu x),s,u)^3\\((\theta Hu),s,u)^3\\((\theta Hu),s,u)^2} & \mcell{((j\eta x),s,u)^2\\(H_jsu)\\((aLu)^2,s,u)^2\\((aLu)^3,s,u)} &  &  &  &  &  &  &  &  &  &  &  &  &  \\ \hline
11 & \mcell{(a_\nu(j\nu x),s,u)^4\\((\theta Hu),s,u)^4} & \mcell{(K_\eta(KHu)^2,s,u)^3\\((j\eta x),s,u)^3\\((aLu)^2,s,u)^3\\(a_\nu\cdot H,s,u)^4} & \mcell{(H^2su)^3\\(\theta_l(\theta Lu)^2,s,u)^2} &  & \mcell{((KHu)^2,s,u)^2} &  &  &  &  &  &  &  &  &  &  &  &  \\ \hline
12 &  & \mcell{(a_\eta(aHu)^2\cdot H,s,u)^3} & \mcell{((KHu)^2,s,u)^3} & \mcell{((aHu)H_j,s,u)^2\\((\theta Lu)^2,s,u)^2\\((\theta Lu)^3,s,u)} &  &  &  &  &  &  &  &  &  &  &  &  &  \\ \hline
13 & \mcell{((KHu)^2,s,u)^4} & \mcell{((aHu)H_j,s,u)^3\\((\theta Lu)^2,s,u)^3\\(\theta_\nu\cdot H,s,u)^4} & \mcell{(L_jsu)^2\\((aH^2u)^3,s,u)^2\\((j\nu x)^2\cdot H,s,u)^3\\((aHu)^2\cdot H,s,u)^3} &  &  &  &  &  &  &  &  &  &  &  &  &  &  \\ \hline
14 & \mcell{((j\nu x)^2\cdot H,s,u)^4\\((aHu)^2\cdot H,s,u)^4} & \mcell{(\theta_\eta(\theta Hu)^2\cdot H,s,u)^3} &  & \mcell{((KLu)^3,s,u)^2} &  &  &  &  &  &  &  &  &  &  &  &  &  \\ \hline
15 & \mcell{(a_l(aLu)^2\cdot H,s,u)^4} & \mcell{((KLu)^3,s,u)^3} & \mcell{((aLu)^2L_j,s,u)^2\\((\theta H^2u)^3,s,u)^2\\((\theta Hu)^2\cdot H,s,u)^3} &  &  &  &  &  &  &  &  &  &  &  &  &  &  \\ \hline
16 & \mcell{((\theta Hu)^2\cdot H,s,u)^4\\(a_\nu\cdot H_j,s,u)^4} & \mcell{((j\eta x)\cdot H,s,u)^4\\(H_j\cdot H,s,u)^3\\((aLu)^2\cdot H,s,u)^4\\((aLu)^3\cdot H,s,u)^3} &  &  &  &  &  &  &  &  &  &  &  &  &  &  &  \\ \hline
17 & \mcell{(\theta_l(\theta Lu)^2\cdot H,s,u)^4} &  & \mcell{((KH^2u)^4,s,u)^2} &  &  &  &  &  &  &  &  &  &  &  &  &  &  \\ \hline
18 &  & \mcell{((\theta Lu)^2\cdot H,s,u)^4\\((\theta Lu)^3\cdot H,s,u)^3\\((aHu)^2\cdot H_j,s,u)^3} &  &  &  &  &  &  &  &  &  &  &  &  &  &  &  \\ \hline
19 & \mcell{((aH^2u)^3\cdot H,s,u)^4\\(L_j\cdot H,s,u)^4} &  &  &  &  &  &  &  &  &  &  &  &  &  &  &  &  \\ \hline
20 &  & \mcell{((KLu)^3\cdot H,s,u)^4} & \mcell{((aLu)^3\cdot H_j,s,u)^2} &  &  &  &  &  &  &  &  &  &  &  &  &  &  \\ \hline
21 & \mcell{((\theta H^2u)^3\cdot H,s,u)^4\\((aLu)^2\cdot H_j,s,u)^4} &  &  &  &  &  &  &  &  &  &  &  &  &  &  &  &  \\ \hline
22 &  &  &  &  &  &  &  &  &  &  &  &  &  &  &  &  &  \\ \hline
23 & \mcell{((KH^2u)^4\cdot H,s,u)^4} & \mcell{((aH^2u)^3\cdot H,s,u)^3} &  &  &  &  &  &  &  &  &  &  &  &  &  &  &  \\ \hline
\end{tabular}
\end{adjustbox}
\end{center}
\vspace{0.2em}
\begin{center}\footnotesize （最上段の横列の数字は $x$ に関する次数を、最初の縦列の数字は $u$ に関する次数を表す。）\end{center}
\endgroup





\clearpage

\providecommand{\pair}[2]{(#1\,|\,#2)}
\providecommand{\eps}{\varepsilon}
\providecommand{\rhoidx}{\varrho}

\begin{center}
{\Large\bfseries 3. $n$ 変数形式の不変式論について\footnote{1909 年、ザルツブルク自然科学者集会で行われた講演。}}\par
\vspace{1.2em}
Jahresbericht der DMV 19 (1910), pp. 101--104
\end{center}

\vspace{1em}
$n$ 変数形式の射影的不変式論では、主問題、すなわち形式系の有限性は解決されている。しかし、さらに別の一連の問題、すなわち形式間の関連を調べる方法に関わる問題は扱われていない。私はそのような問題について得た結果を手短に述べたいが、明瞭さのため、まずそれらの問題設定を、それが知られている三元の領域で概説しておく。

まず念頭にあるのは、記号的方法の二つの基本定理である。第一は、すべての不変な形成が記号的に表示できるという定理であり、第二は、不変式の間に存在するすべての関係が、少数の記号的恒等式を逐次適用することによって得られるという定理である。したがって、記号的方法は、不変式の領域を離れることなくすべての関係を与える。\footnote{Study: \emph{Methoden zur Theorie der ternären Formen}. II. \S\ 6. (Leipzig, Teubner, 1889.)} その上に、形式を生成する過程、すなわち記号的な折り畳み過程と非記号的な微分過程、級数展開の理論などが置かれる。これに関連して、私の学位論文で証明した定理\footnote{\emph{Über die Bildung des Formensystems der ternären biquadratischen Form}. \S\ 1. (Crelle's Journal Bd. 134. 1908.)}も挙げておきたい。すなわち、すべての折り畳みは、二つの可能な「共変的」折り畳みを逐次適用することによって生成できる、という定理である。ここで、記号的積 $s_x t_x u_\sigma u_\tau$ に対して私がいう二つの折り畳みとは、$(stu)$ と $(\sigma\tau x)$ のことである。この定理の上に、二元領域の場合と同様に、還元子の理論および形式系の還元の理論を築くことができることは、私が学位論文で示したとおりである。

さて $n$ 元領域に移ると、すでに記号的方法の第一定理は、制限された範囲でしか成り立たない。よく知られているように、四元形式では、点座標および平面座標 $x$ と $u$ のほかに、二つの点座標列から作られる小行列式である直線座標 $p_{ik}$ が現れる。これと同様に、$n$ 元領域では変数列 $p_\varrho$ があり、その個々の要素は $\varrho$ 個の列 $x$ から作られる小行列式である。また記号列 $S^\varrho$ があり、その要素は、$u$ と共変的な $\rho$ 個の列 $s$ から作られる小行列式のように振る舞う。行列式による記号的表示は、知られているように、記号列 $S^\varrho$ を列 $s$ に分解し、それらを別々の行列式に分配した場合にしか成り立たない。このとき実際の意味をもつのは、$S^\varrho$ に依存する行列式の集合体だけである。しかし、どのような行列式の集合体がこれを実現するのかは見通せず、そのため三元領域で述べた他のすべての定理が失われてしまう。

そこで私は、列 $S^\varrho$ において明示的な記号的表示を得るための簡単な原理を示し、それによって残りの定理も回復したい。これは、対応行列に関するよく知られた定理の応用である。「各変数列 $p_\varrho$ には、一対一に別の列 $q_{n-\varrho}$ が対応する。その個々の要素は $n-\varrho$ 個の列 $u$ から作られる小行列式であり、これらの列は方程式 $(u\mid x)=0$ によって $\rho$ 個の列 $x$ と結ばれている。」これに対応して、各記号列 $S^\varrho$ には別の $S_{n-\varrho}$ が対応し、それは $x$ と共変的な $n-\varrho$ 個の列から構成されているかのように振る舞う。

この定理には、応用に適した第二の定式化もある。「各行列積 $(v^{(1)}\ldots v^{(\varrho)}\mid y^{(1)}\ldots y^{(\varrho)})$ には、\footnote{すなわち、$\varrho$ 個の列 $v$ の行列と $\varrho$ 個の列 $y$ の行列から作られる、対応する行列式の積の和である。}第二の行列の列が第一の行列の列に反変的であるとき、一つの行列式が対応する。また逆に、各行列式は、あらかじめ与えられた列数 $\varrho$ をもつ行列積に変換できる。」したがって、行列式と行列積とは完全に同等のものとして現れ、記号的方法の第一定理は次の形をとる。「すべての不変な形成は、行列積によって記号的に表示できる。」二つの記号列 $S^\sigma$, $T^\tau$ からは、変数列を用いて、これらの記号列を明示的に含む行列積、したがって形式 $(S^\sigma\mid p_\sigma)(T^\tau\mid p_\tau)$ の不変な形成を表す行列積を、非常に容易に作ることができる。すなわち次のものである。
\begin{equation}
 \pair{q_{n-\sigma-\lambda}S^\sigma}{T_{n-\tau}p_{\tau-\lambda}},
 \qquad (\lambda=1,2,\ldots,\tau\ \text{または}\ n-\sigma).
\tag{1}
\end{equation}

より重要なのはその逆である。すなわち、すべての不変な形成が行列積によって明示的に表示できるということであり、したがって上の過程は二元および三元の折り畳み過程の拡張を意味する。この逆を証明するには、記号的方法の第二定理を $n$ 元領域へ移す必要がある。すなわち、すべての列 $S^\varrho$, $p_\varrho$ を分解不能なものと見なす場合の、独立で分解不能な恒等式の全体を立てなければならない。これらの恒等式は、列 $x$ と $u$ だけを含む既知の恒等式\footnote{E. Pascal in \emph{Memorie d. R. Acc. d. Lincei}. (4) 5 (1888), p. 375.}から、行列式の積定理を行列に対する積定理の系で置き換え、さらにこれらの積定理そのものによって、いくつかの行列積を一つにまとめることによって得られる。このようにして二つの双対的な公式に到達するが、そのうち一つだけを記す。
\begin{equation}
\begin{aligned}
 \pair{S_1^{\rho_1}S_2^{\rho_2}\cdots S_k^{\rho_k}}{x p_{\rho-1}}
 &=\sum_{i=1}^{k}\eps\,\pair{S_1^{\rho_1}\cdots S_{i-1}^{\rho_{i-1}}S_{i+1}^{\rho_{i+1}}\cdots S_k^{\rho_k}q_{n-\rho+1}}{S_{n-\rho_i}x},\\
 \rho&=\sum_{i=1}^{k}\rho_i<n,\qquad \eps=\pm1.
\end{aligned}
\tag{2}
\end{equation}

これらの公式に含まれる唯一の特殊化、すなわち列 $x$ が入るということは避けられない。完全に一般の記号列および変数列の場合には、「分解恒等式」に到達する。これは、列 $p_\varrho$ が個々の列 $x$ に分解される恒等式である。分解恒等式の最も簡単な特殊例は、Clebsch がその「不変式論の基本問題」において、級数展開における基本共変式を導くためにすでに用いていた。非明示的な行列式表示から行列積への移行を媒介するこれらの恒等式によって、行列表現により望ましい明示的な記号的表示が実際に得られることを示すことができる。

しかし折り畳み過程については、三元領域の場合とまったく同様の定理が成り立ち、それに基づいて還元定理も同様に築くことができる。(1) において数 $\lambda$ を折り畳みの欠損数と呼び、$\sigma=\tau$ である折り畳みを共変的と呼ぶならば、すべての折り畳みは欠損数 1 の共変的折り畳みを逐次適用することによって生成できる。この定理の証明は、本質的に、最も一般の形式を「正規形」、すなわちすべての不変な微分演算の適用のもとで恒等的に消える形式で置き換えることができる、という事実に依拠する。四元領域ではこの後者の事実を Mertens\footnote{Über invariante Gebilde quaternärer Formen. Wiener Berichte. Bd. 98. 1889.} が証明した。最も一般の微分演算、すなわちよく知られているように記号列を微分記号で置き換えることによって折り畳み過程から生じる微分演算を我々が知っていることにより、分解恒等式を用いて Mertens の証明をほとんど字句どおりに $n$ 元領域へ移すことができる。

\clearpage
\begin{center}
{\Large\bfseries 4. $n$ 変数形式の不変式論について}\footnotemark\par
\footnotetext{ザルツブルク自然科学者集会への短報を参照。\emph{Jahresbericht der Deutschen Mathematiker-Vereinigung}, Bd. 19, S. 101.}
Erlangen の Frl. Emmy Noether による。
\vspace{1.2em}
Journal für die reine und angewandte Mathematik 139 (1911), pp. 118--154
\end{center}

\section*{序論.}
$n$ 変数形式の射影的不変式論では、主問題、すなわち形式系の有限性の証明に付随する問題は、二元および三元の領域を越えると、ほとんど扱われていない。それは、形式の相互依存性およびこの形式間の関連を把握するための方法に関する問題である。これらの方法は本質的に、二元および三元の領域で完全に証明されているように、Study が「記号的方法の基本定理」と呼んだ二つの定理、すなわち形式の記号的表示可能性に関する定理と記号的恒等式に関する定理に基づいている。\footnote{E. Study, \emph{Methoden zur Theorie der ternären Formen} (Leipzig, Teubner, 1889). II. \S\ 6. これらの基本定理が特殊な変換群の不変式論にも現れることについては、Study, \emph{Das Apollonische Problem} (Math. Ann. Bd. 49 [1897]) および \emph{Zur Differentialgeometrie der analytischen Kurven} (Transactions of the American Math. Society, Bd. 1 [1909]) を参照。} 後者の定理は、整な不変形成の間に存在するすべての関係が、少数の記号的恒等式を逐次適用することによって得られることを述べる。したがって、記号的方法、しかもそれだけが、不変式の領域を離れることなく形式間の関連についての知識を与えるのである。

Gordan はそのエルランゲン・プログラム\footnote{P. Gordan, \emph{Über das Formensystem binärer Formen}. S. 46. (Leipzig, Teubner, 1875).}において、すでに $n$ 変数への移行が困難になる理由を指摘していた。それは、記号的方法の第一定理が、もはやその一般的な形では有効でないという点にある。$n$ 元領域では、よく知られているように、高い段階の変数列 $p_\rho$\footnote{記法については \S\ 1 を参照。}が現れ、同様に高い段階の記号列 $S^\rho$ も現れる。もともと Clebsch\footnote{A. Clebsch, über eine Fundamentalaufgabe der Invariantentheorie. (Abh. der Ges. d. Wiss. zu Göttingen, Bd. XVII, 1872).} のように任意個の列 $p_\rho$ をもつ形式から出発する場合にも、F. Deruyts と A. Capelli\footnote{Cf. A. Capelli, \emph{Lezioni sulla teoria delle forme algebriche} (Napoli 1902), \S\ 22--24, およびそこに引用された文献。} に従って任意個の共変的な列 $x$ だけを含む形式から出発する場合にも、これらの列に到達する。列 $p_\rho$, $S^\rho$ を明示的に含む記号的表示\footnote{「明示的表示」のより詳しい定義は \S\ 2。}は存在しない。$p_\rho$, $S^\rho$ を個々の列 $x$ と、それに反変的な $s$ とに分解し、それらを別々の行列式に分配しなければならない。確かに、微分条件（公式 (19.) 参照）によって、与えられた行列式の集合体が列 $p_\rho$, $S^\rho$ だけに依存する性質をもつことを確かめることはできる。しかしこの方法では、不変な形成の全体およびそれらを生成する過程についての見通しは得られない。\footnote{Study による、計算上きわめて有用な記号法の展開（F. Engel, Ber. d. K. Sächs. Ges. d. Wiss., math.-phys. Kl., 1900, p. 126 によって伝えられ、四元領域については E. Study, \emph{Geometrie der Dynamen}, p. 135 による）もまた、ここでいう意味での明示的表示ではない。たとえば、それは形式を生成する非記号的な微分過程にはほとんど導かないであろう。}

そこで \S\S\ 2 および 6 で、「対応行列に関する定理」\footnote{Grassmann の 1862 年の \emph{Ausdehnungslehre}, no. 112 を別にすれば、最初は A. Brill, über zwei Berührungsprobleme, \S\ 2 (Math. Ann. Bd. 4. 1871)。}を用い、行列式による表示を「行列積」による表示で置き換えることによって、第一基本定理がその一般性において回復されることを示す。\S\ 2 では、列 $S^\rho,p_\rho$ を明示的に含む各行列積が基礎形式の不変な形成であることを示す。\S\ 6 で与えられる逆、すなわち各不変な形成が行列積によって明示的に表示できるということは、第二基本定理の移行（\S\S\ 3--5）を通して初めて成り立つ。

この定理は、変数列 $x$ だけを含む形式については知られている。\footnote{E. Pascal, Giorn. di mat. 26 (1888), pp. 33, 102; Rendic. d. Accad. d. Linc. (4) 4 (1888), p. 119; ib. Mem. (4) 5 (1888), p. 375.} すべての列 $S^\rho,p_\rho$ を分解不能と見なす場合の、分解不能な恒等式の全体を立てるという一般問題は、Study\footnote{``Methoden \dots'', p. 204, note 17).} によって定式化された。同時に彼は、現れる恒等式の数に、$n$ 元領域における方法の困難さの尺度を見ている。いまや恒等式の全体を一つの公式 (36.) にまとめることができるので、この困難は少なくとも最小限に抑えられる。しかし応用では、(36.) のある際立った特殊例、たとえば (20.), (21.) がしばしば現れる。それらは、新しい列を代入することなしに明示的表示を許すという点で特徴づけられる。また公式 (31.) は「分解恒等式」と呼ばれる。そこでは列 $p_\rho$ が列 $y$ に分解されているように現れるからである。

二つの基本定理の上に、さらなる理論を容易に築くことができる。第一定理は、形式を生成する過程、すなわち記号的折り畳み過程と非記号的微分過程\footnote{論文 ``Invariante Gebilde von Nullsystemen'' (Sitzungsber. d. Ak. d. Wiss. Wien, Bd. 97, Abt. IIa, 1888) において、Mertens は別の道筋により、直接には一般化できない四元の折り畳み過程に到達した。そこに現れる 6 個の列 $S^2$ の交代不変式は、ここで生じるものとは、置換 (11.) を一度適用することによって、偶不変式の集合体だけ異なる。特殊な場合については、四元の折り畳み過程は P. Gordan, \emph{Das simultane System von zwei quadratischen quaternären Formen} (Math. Ann. Bd. 56. 1903) にも見られる。}を直接に与える。一方、分解恒等式は、これらの過程のうち同値なものをすべて消去する（\S\ 6）。微分過程を知ることにより、さらに「正規形」の概念を $n$ 元領域へ移すことができ、また Mertens\footnote{F. Mertens, Über invariante Gebilde quaternärer Formen. (Sitzber. d. Ak. d. Wiss. Wien. Bd. 98. Abt. IIa. 1889.)} が四元領域で与えた証明手続きによって、不変な形成の系を同値な正規形の系で置き換えることができるという定理が得られる（\S\ 7）。この後者の仮定のもとで、私は学位論文\footnote{Über die Bildung des Formensystems der ternären biquadratischen Form. (Dieses Journal, Bd. 134. 1908.)}において、三元の折り畳み過程が二つの基本的な折り畳みを逐次適用することによって生成できることを示した。この定理と級数展開の理論に基づき、さらに「形式列」という概念を導入して、形式系の主要な還元定理を展開した。完全に同様に、$n$ 元領域では、折り畳み過程は $n-1$ 個の基本的な折り畳みから生成でき（\S\ 8）、還元定理を立てることができる（\S\ 9）。

\paragraph{後記。} ザルツブルクでの報告の際に、ウィーンの E. Müller 教授は、この論文の基礎にある行列積の計算が、原理的には Grassmann の拡張論の計算法と同一であることを私に指摘した。\footnote{Hermann Graßmanns gesammelte mathematische und physikalische Werke. Ersten Bandes zweiter Teil: Die Ausdehnungslehre von 1862. In Gemeinschaft mit H. Graßmann d. Jüngeren herausgegeben v. Fr. Engel (Leipzig, Teubner 1896).} 実際、Grassmann の「組合せ積」$[S^{\rho_1}S^{\rho_2}\cdots S^{\rho_i}]$ は、これらの列によって与えられた行列と見なすことができる（\S\ 2 参照）。具体的には、「進行積」はその列自身の行列、「退行積」は対応する列 $S_{n-\rho_1},S_{n-\rho_2},\ldots S_{n-\rho_i}$ の行列であり、「混合積」に対応する行列は、いくつかの列 $S$ が置換 (9.) によって新しい列 $P,Q$ にまとめられるときに生じる。我々の「対応する列」は Grassmann の「補元」に対応し、我々の「行列積」は「段階数 0 の混合積」に対応する。

この対応のもとで、1862 年の \emph{Ausdehnungslehre} no. 113 に与えられた展開は、我々の公式 (32.) に双対的な一公式と同一になる。特殊な場合 $\rho=1$ では、(32.) は (17.) へ移り、双対的公式は (16.) へ移る。Grassmann の展開のこれら特殊な場合から、Müller\footnote{E. Müller, Beiträge zur Graßmannschen Ausdehnungslehre. 1. Mitteilung: Einige allgemeine Sätze. (Sitzungsber. d. Ak. d. Wiss. Wien; Bd. 118. Abt. IIa. Juli 1909), Satz 16 und 18.} は最近、上に挙げた恒等式 (20.), (21.) を導いた。

Grassmann--Müller の導出と異なり、我々の導出は同時に恒等式の完全性の証明を与える。ここで考察している問題にとって本質的なのはまさにこの点であると思われる。さらに我々の導出は、(20.), (21.) から出発して、これまで注意されていなかったと思われる最も一般の分解不能な恒等式 (36.) へ進む。

\section*{\S\ 1. 記法と定義。既知の結果のまとめ。}
通常どおり、変数列 $x$ によって要素 $x_1,x_2,\ldots,x_n$ の列を表し、同様に変数列 $p_\rho$（$\rho=1,2,\ldots,n-1$）によって $\binom n\rho$ 個の要素 $p_{i_1i_2\ldots i_\rho}$ の列を表す。ここで $p_{i_1i_2\ldots i_\rho}$ は、$\rho$ 個の列 $x^{(1)},x^{(2)},\ldots,x^{(\rho)}$ の行列から作られる $\binom n\rho$ 個の行列式を表す。
\[
\left|\begin{array}{cccc}
 x_{i_1}^{(1)}&x_{i_2}^{(1)}&\cdots&x_{i_\rho}^{(1)}\\
 x_{i_1}^{(2)}&x_{i_2}^{(2)}&\cdots&x_{i_\rho}^{(2)}\\
 \vdots&\vdots&&\vdots\\
 x_{i_1}^{(\rho)}&x_{i_2}^{(\rho)}&\cdots&x_{i_\rho}^{(\rho)}
\end{array}\right|,
\qquad (i_1<i_2<\cdots<i_\rho=1,2,\ldots,n).
\]
対応行列に関する定理\footnote{たとえば Gordan--Kerschensteiner, \emph{Vorlesungen über Invariantentheorie}, Bd. I, p. 99 を参照。}によれば、各列 $p_\rho$ には、列の各要素について成り立つ次の関係によって、第二の列 $q_{n-\rho}$ が一対一に対応する。
\begin{equation}
 p_{i_1i_2\ldots i_\rho}=(-1)^{\frac{\rho(\rho+1)}{2}+i_1+i_2+\cdots+i_\rho}q_{j_1j_2\ldots j_{n-\rho}},
\tag{1}
\end{equation}
ここで $i$ と $j$ はそれぞれ整列しており、合わせて 1 から $n$ までの数の完全な置換をなす。列 $q_{n-\rho}$ は、$x$ に反変的な $n-\rho$ 個の列 $u^{(1)},u^{(2)},\ldots,u^{(n-\rho)}$ の行列から作られる、次数 $n-\rho$ の $\binom n\rho$ 個の行列式 $q_{j_1j_2\ldots j_{n-\rho}}$ から成る。これらの列は、$\rho\cdot(n-\rho)$ 個の条件方程式
\begin{equation}
 (u^{(k)}\mid x^{(l)})=\sum_{\alpha=1}^{n}u_\alpha^{(k)}x_\alpha^{(l)}=0;
 \qquad (k=1,2,\ldots,n-\rho;\ l=1,2,\ldots,\rho).
\tag{2}
\end{equation}
を満たす。以下では、Clebsch\footnote{a. a. O. \S\ 5.} に従って、(1.) に本来現れる比例係数が 1 になるように、これらを正規化しておく。

Clebsch\footnote{a. a. O. \S\ 12.} によれば（また F. Deruyts と A. Capelli の後の研究、序論参照、によっても）、最も一般の形式は、$n-1$ 個の列 $p_\rho$ の各々を高々一つずつ含むものに還元でき、したがって次のように記号的に表すことができる。
\begin{equation}
 (S^1\mid p_1)^{m_1}(S^2\mid p_2)^{m_2}\cdots(S^{n-1}\mid p_{n-1})^{m_{n-1}},
\tag{3}
\end{equation}
ここで (2.) と同様に
\[
 (S^\rho\mid p_\rho)=\sum S^{i_1i_2\ldots i_\rho}p_{i_1i_2\ldots i_\rho}
\]
と置いた。列 $p_\rho$ の要素の間に存在する非線形の恒等的関係のために、記号列 $S^\rho$ は、それらがまったく同じ恒等式を満たすように正規化できる。したがって、それらは、行 $s^{(1)},s^{(2)},\ldots s^{(\rho)}$ が $x$ に反変的である $\rho$ 行の行列の行列式のように振る舞う。\footnote{Clebsch, a. a. O. \S\ 4.} こうして、各列 $S^\rho$ には、次の関係を通じて別の $S_{n-\rho}$ が一対一に対応する。
\begin{equation}
 S_{i_1i_2\ldots i_{n-\rho}}=(-1)^{\frac{(n-\rho)(n-\rho+1)}2+i_1+i_2+\cdots+i_{n-\rho}}S^{j_1j_2\ldots j_\rho},
\tag{4}
\end{equation}
ここでも $i$ と $j$ はそれぞれ整列しており、合わせて 1 から $n$ までの数をちょうど尽くす。列 $S_{n-\rho}$ の要素 $S_{i_1i_2\ldots i_{n-\rho}}$ は、$x$ と共変的な $(n-\rho)$ 個の列 $\sigma^{(1)},\sigma^{(2)},\ldots\sigma^{(n-\rho)}$ から作られる行列式のように振る舞い、列 $s$ と $\sigma$ は $\rho\cdot(n-\rho)$ 個の条件
\begin{equation}
 (s^{(k)}\mid \sigma^{(l)})=0.\qquad(k=1,2,\ldots\rho;\ l=1,2,\ldots n-\rho).
\tag{5}
\end{equation}
によって結ばれる。関係 (1.) および (4.) の結果として、(3.) の各因子は、同値な四重の記号的表示をもつ。
\begin{equation}
 (S^\rho\mid p_\rho)=(S^\rho q_{n-\rho})=(S_{n-\rho}p_\rho)=(-1)^{\rho(n-\rho)}(q_{n-\rho}\mid S_{n-\rho}\footnotemark),
\tag{6}
\end{equation}
\footnotetext{Clebsch, a. a. O. \S\ 5 を参照。}
ここで、以下つねにそうであるように、$(S^\rho q_{n-\rho})$ および $(S_{n-\rho}p_\rho)$ は、それぞれ行列式 $(s^{(1)}\cdots s^{(\rho)}u^{(1)}\cdots u^{(n-\rho)})$ および $(\sigma^{(1)}\cdots\sigma^{(n-\rho)}x^{(1)}\cdots x^{(\rho)})$ の略記である。Laplace 展開によって、これらはそれぞれ列 $S^\rho,q_{n-\rho}$ および $S_{n-\rho},p_\rho$ のみに依存することが分かり、上の記法はそのことを示すためのものである。式 $(S^\rho\mid p_\rho)$ および $(q_{n-\rho}\mid S_{n-\rho})$ は、前述の議論によれば、列 $s$ と $x$ の行列、ならびに列 $u$ と $\sigma$ の行列の積を意味する。ここで行列積とは、対応する行列式の積の和、すなわち積行列の行列式の値を意味する。

記号列（それぞれ変数列）に属する特性数 $\rho\cdot(n-\rho)$ を、Clebsch に従ってその列の重みと呼ぶ。\footnote{a. a. O. \S\ 11.} 列 $s$ の数を与える数 $\rho$ を上重み添字と呼び、列 $\sigma$ の数を与える数 $n-\rho$ を下重み添字と呼ぶ。関係 (4.) から、一つの記号列は、同じ関係の中で一度は上重み添字をもって、また一度は下重み添字をもって現れ得ることが従う。対応する列 $p_\rho$ と $q_{n-\rho}$ の出現についても同じことが成り立つ。最後に、一般に、対応する行列が列 $x$ から成る高い段階の変数列を $p$ で表し、その行列が列 $u$ から成るものを $q$ で表すことを記しておく。$x$ と共変的な記号列は小ギリシア文字 $\sigma,\tau,\alpha,\beta$ で表し、$u$ と共変的な記号列は小ラテン文字 $s,t,a,b$ で表す。その他の記号列はすべて、重み添字を付した大ラテン文字 $S^\rho$, $T^\tau$, $A^\rho$ または $S_{n-\sigma}$, $T_{n-\tau}$, $A_{n-\rho}$ で表す。上または下の重み添字に対応して、たとえば (4.) のように、列の個々の要素も上付きまたは下付きの添字で書く。
\clearpage
\section*{\S\ 2. 行列積による表示.}
以下では常に、「行列積」とは、同じ数の行をもつ二つの行列について、対応する小行列式の積を総和したものを意味する。ただし第二の行列の列は第一の行列の列に反変的であるものとする。各行列の列は、1. 個別に与えられていてもよく、2. (6.) のように一つの列 $S^\rho$ にまとめられていてもよく、3. 異なる列 $S^{\rho_1},S^{\rho_2}\ldots$ に結合されていてもよい。行列積は (6.) で用いた記号 $(\ \mid\ )$ で表す。(6.) の表示は、対応行列に関する定理の第二の定式化の特別な応用である。「各 $\rho$ 行の行列積 $(v^{(1)}v^{(2)}\cdots v^{(\rho)}\mid y^{(1)}\cdots y^{(\rho)})$ には一つの行列式が対応し、逆に各行列式には、あらかじめ定めた行数 $\rho$ をもつ行列積を対応させることができる（$\rho=1,2,\ldots,n-1$）。」\footnote{$\rho>n$ では、行列積はよく知られているように消える。$\rho=n$ では二つの行列式の積になる。} 実際、与えられた行列積から行列式へ、またその逆へ移るには、列 $y$ を一つの列 $p_\rho$ にまとめ（または列 $v$ を一つの列 $q_\rho$ にまとめ）、(1.) で与えられた置換を実行すればよい。ただし、行列式の値はその個々の要素から直接与えられるのではなく、ラプラス展開によって初めて得られる。このように行列式と行列積が同値であることがわかるので、形式の記号的表示可能性に関する定理（第一基本定理）は、次のようにも表される。

\emph{定理 I.} 「すべての不変な形成は行列積によって記号的に表示でき、また逆にすべての行列積は不変な形成である。」\footnote{もちろん、指定された基礎形式の不変な形成とは、記号列 $S^\rho\ldots$ に依存する行列積の集合体だけである。}

行列積による表示は、行列式表示によって生じる非明示的表示の本質的困難を克服させる。\footnote{序論を参照。} ここで「明示的表示」とは、列 $S^\sigma,T^\tau,\ldots$ そのもの、またはそれらから一意に導かれる列 $R^\rho$ だけが入り、個々の成分列 $s,t,\ldots$ はもはや現れない表示をいう。\S\ 6 では、記号列 $S^\sigma,T^\tau,\ldots$ のみに依存するすべての不変な形成について、これらの記号列において明示的な表示を得る。そしてこれにより、生成過程、折り畳み過程、微分過程を $n$ 元領域へ移すことができる。明示的表示が可能であることは、行列式表示に現れる行列式に正確に対応しうるのが最も単純な行列積だけであり、したがって他のすべての行列積は、異なる行列式を一つの式にまとめたものになる、という事実から明らかである。実際、記号列および変数列を個々の列 $s,u,x$ に分解すれば、通常の形の第一基本定理によって、行列式の集合体が必然的に現れなければならない。逆に、指定された基礎形式の不変な形成を表す行列式の集合体がすべて明示的な行列積にまとめられることは、第二基本定理を用いて証明される（\S\ 6）。

二つの記号列 $S^\sigma,T^\tau$ から、変数列を用いて、基礎形式 $(S^\sigma\mid p_\sigma)(T^\tau\mid p_\tau)$ の不変な形成であってすべての列を明示的に含むもの（折り畳み過程）を得るには、この節の初めに与えた定義に従って第三の種類の行列積を作ればよい。
\begin{equation}
 \pair{q_{n-\sigma-\lambda}S^\sigma}{T_{n-\tau}p_{\tau-\lambda}},
 \qquad (\lambda=1,2,\ldots,\tau,n-\sigma).
\tag{7}
\end{equation}
展開すると (7.) は次を与える。
\begin{equation}
\begin{aligned}
&\pair{v^{(1)}\cdots v^{(n-\sigma-\lambda)}s^{(1)}\cdots s^{(\sigma)}}
       {\alpha^{(1)}\cdots\alpha^{(n-\tau)}y^{(1)}\cdots y^{(\tau-\lambda)}} \\
&=\sum
\left|\begin{array}{ccc}
 v_{i_1}^{(1)}&\cdots&v_{i_{n-\lambda}}^{(1)}\\
 \vdots&&\vdots\\
 v_{i_1}^{(n-\sigma-\lambda)}&\cdots&v_{i_{n-\lambda}}^{(n-\sigma-\lambda)}\\
 s_{i_1}^{(1)}&\cdots&s_{i_{n-\lambda}}^{(1)}\\
 \vdots&&\vdots\\
 s_{i_1}^{(\sigma)}&\cdots&s_{i_{n-\lambda}}^{(\sigma)}
\end{array}\right|
\left|\begin{array}{ccc}
 \alpha_{i_1}^{(1)}&\cdots&\alpha_{i_{n-\lambda}}^{(1)}\\
 \vdots&&\vdots\\
 \alpha_{i_1}^{(n-\tau)}&\cdots&\alpha_{i_{n-\lambda}}^{(n-\tau)}\\
 y_{i_1}^{(1)}&\cdots&y_{i_{n-\lambda}}^{(1)}\\
 \vdots&&\vdots\\
 y_{i_1}^{(\tau-\lambda)}&\cdots&y_{i_{n-\lambda}}^{(\tau-\lambda)}
\end{array}\right|\\
&=\sum_{i,k,l}\eps\,q_{k_1k_2\ldots k_{n-\sigma-\lambda}}
 S^{k_{n-\sigma-\lambda+1}\ldots k_{n-\lambda}}
 T_{l_1\ldots l_{n-\tau}}
 p_{l_{n-\tau+1}\ldots l_{n-\lambda}},
\end{aligned}
\tag{8}
\end{equation}
すなわち、列 $S,T,q,p$ のみに依存する式である。ここで $\eps=\pm1$\footnote{この記法は以下全体で保持する。}であり、和は、各列 $S\ldots$ の添字が正しい順序にあり、各組合せ $i_1i_2\ldots i_{n-\lambda}$ に対して、$k$ も $l$ もそれぞれ数 $i$ のすべての可能な順列を走るように理解する。(7.) に入る数 $\lambda$ は、$\sigma\ge\tau$ の場合に折り畳みの欠損数と呼ぶ。この最後の制限の理由は \S\ 6 で現れる。\footnote{数 $\lambda$ は、常に最後に示された二つの数の小さい方まで走る。}

(7.) に入る二つの行列の行列式は、新しい列 $Q^{n-\lambda},P_{n-\lambda}$ の要素と見なすことができ、それらに対応する列は $Q_\lambda,P^\lambda$ である。(8.) によれば、これらの列 $Q,P$ は、もとの $S$ と $q$、それぞれ $T$ と $p$ によって表せる。これを次の等式で示す。
\begin{equation}
\begin{aligned}
 q_{n-\sigma-\lambda}S^\sigma&=Q^{n-\lambda}\sim Q_\lambda, &\quad\text{または}\quad Q_\lambda&=[q_{n-\sigma-\lambda}S^\sigma]_\lambda,\\
 T_{n-\tau}p_{\tau-\lambda}&=P_{n-\lambda}\sim P^\lambda, &\quad\text{または}\quad P^\lambda&=[T_{n-\tau}p_{\tau-\lambda}]^\lambda.
\end{aligned}
\tag{9}
\end{equation}
(4.) を考慮すると、(6.) と同様に、行列積 (7.) について次の同値な四重の記号的表示を得る。
\begin{equation}
\pair{q_{n-\sigma-\lambda}S^\sigma}{T_{n-\tau}p_{\tau-\lambda}}
=(q_{n-\sigma-\lambda}S^\sigma P^\lambda)=(Q_\lambda T_{n-\tau}p_{\tau-\lambda})=(-1)^{\lambda(n-\lambda)}\pair{P^\lambda}{Q_\lambda}.
\tag{10}
\end{equation}
さらに、(7.) に対して新しい列を用いる別の置換を、次のように置くことによって行うことができる。
\begin{equation}
 \pair{q_{n-\sigma-\lambda}S^\sigma}{T_{n-\tau}p_{\tau-\lambda}}
 =\pair{R^{\sigma+\lambda}}{p_{\sigma+\lambda}}(R^{\tau-\lambda}p_{\tau-\lambda}).
\tag{11}
\end{equation}
ここでも (8.) は、列 $R$ の要素の積が、列 $S$ と $T$ によって一意に表されることを示す。置換 (9.) と (11.) は、特に (7.) によって他の列とのさらなる折り畳みを行うときに用いられる。

行列表現により、列 $p_\rho$ の要素間の二次関係（そこから、知られているように、他のすべての関係が従う）を不変的に述べることができる。したがって、\S\ 1 によれば、基礎形式の係数について線形な、列 $S^\rho$ が満たす関係も述べられる。それらは次の恒等式である。
\begin{equation}
 \pair{q_{n-\rho-\lambda}S^\rho}{S_{n-\rho}p_{\rho-\lambda}}=0,
 \qquad (\lambda=1,2,\ldots,\rho,n-\rho),
\tag{12}
\end{equation}
ここで $p$ と $q$ は任意量と見なす。\S\ 5 で、欠損数 $\lambda$ が奇数である恒等式は、より高い欠損の恒等式から従うことを見る。関係 (12.) は関係 (5.) の直接の帰結である。実際、
\[
 \pair{q_{n-\rho-\lambda}S^\rho}{S_{n-\rho}p_{\rho-\lambda}}
 =(v^{(1)}\cdots v^{(n-\rho-\lambda)}s^{(1)}\cdots s^{(\rho)}\mid \sigma^{(1)}\cdots\sigma^{(n-\rho)}y^{(1)}\cdots y^{(\rho-\lambda)}),
\]
であり、積定理を適用すると、(5.) によって消える $(n-\lambda)$ 行の行列式が得られる。恒等式 (12.) は、すべての折り畳みの型を見つけるためには、各変数列について線形な形式の積を折り畳めば足りることを述べる。あるいはまた、\S\ 6 により、正規化された形式 $(S^\rho\mid p_\rho)^m$ は係数について線形な不変形成をもたない、ということでもある。\footnote{Grassmann の著作、第一巻第二部 p. 510 にある Study の補注（量 $S^\rho$ が単純であるための条件）を参照。}

条件 (12.) は、列 $p_\rho$ を特徴づけるのに十分でもある。というのも、$p_\rho$ の個々の要素がある行列の行列式であるという性質は不変な性質であり、したがって不変的に表現されなければならないからである。しかし、定理 VI によれば、条件 (12.) は $p_\rho$ に対する可能な最大の不変的制限、すなわち列 $p_\rho$ を係数として作られる線形形式が不変形成をまったくもたないという制限を与える。

\section*{\S\ 3. 記号的恒等式.}
ここで、記号的方法の第二基本定理を移さなければならない。すなわち、すべての列 $S^\rho,p_\rho$ を分解不能と見なす場合の、既約で分解不能な恒等式の全体を設定するのである。さらに次の取り決めを置く。記号列と変数列の意味に対応して、変数列 $p_\rho$ を $p_{\rho-1}y$ に特殊化し、その結果の式に系の恒等式の一つを適用して生じる恒等式は合成的なものと見なす。一方、$k$ 個の記号列を含む恒等式は、上の過程によってより少ない記号列の恒等式から生じる場合でも、分解不能と数える。列 $S^\rho,p_\rho$ は必ずしも恒等式に明示的に入る必要はない。これらを分解不能な量として解釈するには、現れる行列積がこれらの列だけに依存することを示せば足りる。しかしこの場合にも、部分的には置換 (11.) を適用することで、常に明示的表示が得られることを示す。一般の恒等式は、行列表現の原理によって、Pascal が与えた特殊な恒等式から得られる。これらは列 $x$ と $u$ だけを含み、Study が三元領域で与えた恒等式の直接の移行である。\footnote{問題の定式化と文献については序論を参照。}
\begin{align}
\sum_{i=1}^n(-1)^{i+1}\pair{a^{(i)}}{x}(a^{(1)}\cdots a^{(i-1)}a^{(i+1)}\cdots a^{(n)}u)&=(-1)^{n+1}\pair{u}{x}(a^{(1)}\cdots a^{(n)}),\tag{13}\\
\sum\pm\pair{a^{(1)}}{x^{(1)}}\pair{a^{(2)}}{x^{(2)}}\cdots\pair{a^{(n)}}{x^{(n)}}&=(a^{(1)}\cdots a^{(n)})(x^{(1)}\cdots x^{(n)}),\tag{14}\\
\sum_{i=1}^n(-1)^{i+1}\pair{u}{\alpha^{(i)}}(\alpha^{(1)}\cdots\alpha^{(i-1)}\alpha^{(i+1)}\cdots \alpha^{n}x)&=(-1)^{n+1}\pair{u}{x}(\alpha^{(1)}\cdots\alpha^{(n)}).\tag{15}
\end{align}
(13.) と (15.) からは、置換 $(a^{(i)}\mid x)=(a^{(i)}q_{n-1})$ および $(u\mid \alpha^{(i)})=(p_{n-1}\alpha^{(i)})$ によってさらに二つの恒等式が生じるが、我々の解釈ではそれらは本質的に異なるものとは見なさない。(14.) は行列式の積定理を表す。我々は、(13.) と (14.)（また双対的に (15.) と (14.)）を、適切に形成された行列の積定理の系で置き換える方法を示す。$\rho$ 行の行列と $(\rho-1)$ 行の行列についてそれぞれ作られる積定理は、
\[
\pair{a^{(1)}a^{(2)}\cdots a^{(\rho)}}{xp_{\rho-1}}
=(a^{(1)}a^{(2)}\cdots a^{(\rho)}\mid xy^{(2)}\cdots y^{(\rho)})
=\sum\pm\pair{a^{(1)}}{x}\pair{a^{(2)}}{y^{(2)}}\cdots\pair{a^{(\rho)}}{y^{(\rho)}},
\]
\[
\sum\pm\pair{a^{(2)}}{y^{(2)}}\cdots\pair{a^{(\rho)}}{y^{(\rho)}}
=(a^{(2)}\cdots a^{(\rho)}\mid y^{(2)}\cdots y^{(\rho)})
=(a^{(2)}\cdots a^{(\rho)}\mid p_{\rho-1})=(a^{(2)}\cdots a^{(\rho)}q_{n-\rho-1});
\]
となる。ここから、$\rho=2,3,\ldots n$ について次の積定理の系が従う。
\begin{equation}
\begin{aligned}
\pair{a^{(1)}a^{(2)}\cdots a^{(\rho)}}{xp_{\rho-1}}
 &=\sum_{i=1}^{\rho}(-1)^{i+1}\pair{a^{(i)}}{x}(a^{(1)}\cdots a^{(i-1)}a^{(i+1)}\cdots a^{(\rho)}\mid p_{\rho-1})\\
 &=\sum_{i=1}^{\rho}(-1)^{i+1}\pair{a^{(i)}}{x}(a^{(1)}\cdots a^{(i-1)}a^{(i+1)}\cdots a^{(\rho)}q_{n-\rho+1}).\footnotemark
\end{aligned}
\tag{16}
\end{equation}
\footnotetext{恒等式 (16.) は、$\rho$ 個の変数の領域から移された恒等式 (13.) とみなすこともできる。}
$\rho=n$ では (16.) は (13.) に移る。さらに $p_{n-1}=y^{(2)}p_{n-2}$, $p_{n-2}=y^{(3)}p_{n-3}$ などと特殊化し、恒等式 (16.) を式 $(a^{(1)}\cdots a^{(i-1)}a^{(i+1)}\cdots a^{(n)}\mid y^{(2)}p_{n-2})$, $(a^{(1)}\cdots a^{(i-1)}a^{(i+1)}\cdots a^{(k-1)}a^{(k+1)}\cdots a^{(n)}\mid y^{(3)}p_{n-3})$ などに適用すると、(13.) から (14.) へ移る。したがって恒等式 (14.) は、我々の定義では、恒等式 (16.) から合成されたものである。同値に、系 (16.) は恒等式 (13.) と (14.) に同値である。系 (16.) は一般の恒等式に要求される条件の一つを満たす。すなわち列 $p_{\rho-1}$ を分解不能な量として含み、同時に明示的な形で含む。この条件だけを満たす系はこれだけである。というのも、列 $a,x,p$ からはこれ以上の行列式または行列積を作れず、またそれらの間に (16.) から独立な関係は存在しえないからである。もし存在すれば、$(a^{(1)}\cdots a^{(\rho)}\mid xp_{\rho-1})$ を消去することで、$\rho$ 個（$\rho\le n$）の式 $(a^{(i)}\mid x)(a^{(1)}\cdots a^{(i-1)}a^{(i+1)}a^{(\rho)}q_{n-\rho+1})$ の間に関係が生じるが、(13.) によれば少なくとも $n+1$ 個のそのような式が関係に入らなければならないからである。同様の考察から、(15.) と (14.) に同値な系を得る。
\begin{equation}
\begin{aligned}
\pair{u q_{\rho-1}}{\alpha^{(1)}\cdots \alpha^{(\rho)}}
 &=\sum_{i=1}^{\rho}(-1)^{i+1}\pair{u}{\alpha^{(i)}}(q_{\rho-1}\mid \alpha^{(1)}\cdots \alpha^{(i-1)}\alpha^{(i+1)}\cdots \alpha^{(\rho)})\\
 &=\sum_{i=1}^{\rho}(-1)^{i+1}\pair{u}{\alpha^{(i)}}(p_{n-\rho+1}\alpha^{(1)}\cdots \alpha^{(i-1)}\alpha^{(i+1)}\cdots \alpha^{(\rho)}).
\end{aligned}
\tag{17}
\end{equation}
(16.) から任意の列 $A^{\rho_1},B^{\rho_2},\ldots,x,p$ の間の恒等式へ移るため、次を観察する。これらの恒等式は、$A^{\rho_1}=a^{(1)}a^{(2)}\cdots a^{(\rho_1)}$, $B^{\rho_2}=b^{(1)}b^{(2)}\cdots b^{(\rho_2)}$ などと分解すれば直ちに (16.) と同一になるべきであり、また個々の式が (16.) に入る型である限り、最終的な式へまとめることも (16.) で定められた操作によってしか行えない。したがって、さらに
\begin{equation}
 \rho=\rho_1+\rho_2+\cdots\rho_k\leqq n,
\tag{18}
\end{equation}
と置けば、
\begin{align*}
&\pair{a^{(1)}\cdots a^{(\rho_1)}b^{(1)}\cdots b^{(\rho_2)}\cdots k^{(1)}\cdots k^{(\rho_k)}}{xp_{\rho-1}}
  =\pair{A^{\rho_1}B^{\rho_2}\cdots K^{\rho_k}}{xp_{\rho-1}} \\
&\quad=\sum_{i=1}^{\rho_1}(-1)^{i+1}\pair{a^{(i)}}{x}(a^{(1)}\cdots a^{(i-1)}a^{(i+1)}\cdots a^{(\rho_1)}B^{\rho_2}\cdots K^{\rho_k}q_{n-\rho+1})\\
&\qquad+(-1)^{\rho_1\rho_2}\sum_{i=1}^{\rho_2}(-1)^{i+1}\pair{b^{(i)}}{x}(b^{(1)}\cdots b^{(i-1)}b^{(i+1)}b^{(\rho_2)}A^{\rho_1}\cdots K^{\rho_k}q_{n-\rho+1})\\
&\qquad\vdots\\
&\qquad+(-1)^{(\rho_1+\cdots\rho_{k-1})\rho_k}\sum_{i=1}^{\rho_k}(-1)^{i+1}\pair{k^{(i)}}{x}(k^{(1)}\cdots k^{(i-1)}k^{(i+1)}\cdots k^{(\rho_k)}A^{\rho_1}\cdots (K-1)^{\rho_{k-1}}q_{n-\rho+1}).
\end{align*}
各個別の和（ただしまだ和の一部分ではない）は、実際に列 $A^{\rho_1},B^{\rho_2}\ldots$ を分解不能量として含む。なぜなら、それは必要十分条件\footnote{Capelli, a. a. O. \S\ 23 は十分条件 $\bigl(a^{(k)}\mid \frac{\partial}{\partial a^{(i)}}\bigr)=0$, $(k>i,\ i=1,2,\ldots,\rho_1-1)$ を与える。これらは、Wellstein, Von den Differentialgleichungen der projektiven Invarianten, Math. Ann. 67 (1909) \S\ 3 に従う Poisson 括弧過程を適用することで、最も簡単に必要十分条件を含意する。}
\begin{equation}
 \left(a^{(i+1)}\middle|\frac{\partial}{\partial a^{(i)}}\right)=0;\ \ldots\ ;
 \left(k^{(j+1)}\middle|\frac{\partial}{\partial k^{(j)}}\right)=0.
 \quad (i=1,2,\ldots,\rho_1-1;\ j=1,2,\ldots,\rho_k-1)
\tag{19}
\end{equation}
を満たすからである。さらに、それがすべての列を明示的に含むことを示す。実際、$B^{\rho_2}\cdots K^{\rho_k}q_{n-\rho+1}=q_{n-\rho_1+1}$ と置くと、(16.) と (10.) により、
\begin{align*}
\sum(-1)^{i+1}\pair{a^{(i)}}{x}(a^{(1)}\cdots a^{(i-1)}a^{(i+1)}\cdots a^{(\rho_1)}q_{n-\rho_1+1})
&=\pair{A^{\rho_1}}{xp_{\rho_1-1}}\\
&=(A_{n-\rho_1}xp_{\rho_1-1})=(-1)^{(\rho_1+1)(n+1)}\pair{q_{n-\rho_1+1}}{A_{n-\rho_1}x}\\
&=(-1)^{(\rho_1+1)(n+1)}\pair{B^{\rho_2}\cdots K^{\rho_k}q_{n-\rho+1}}{A_{n-\rho_1}x}
\end{align*}
を得る。残りの和を同様に計算し、ついで、列はすでに重み添字によって互いに区別されているので、$B^{\rho_2}=A^{\rho_2}$, $C^{\rho_3}=A^{\rho_3},\ldots$ と置くと、求める恒等式を得る。
\begin{equation}
\begin{aligned}
\pair{A^{\rho_1}A^{\rho_2}\cdots A^{\rho_k}}{xp_{\rho-1}}
&=\sum_{i=1}^{k}(-1)^{\delta_i}\pair{A^{\rho_1}\cdots A^{\rho_{i-1}}A^{\rho_{i+1}}\cdots A^{\rho_k}q_{n-\rho+1}}{A_{n-\rho_i}x},\\
\delta_i&=(\rho_1+\rho_2+\cdots+\rho_{i-1})\rho_i+(\rho_i+1)(n+1),
\end{aligned}
\tag{20}
\end{equation}
また (17.) から双対的な恒等式
\begin{equation}
\begin{aligned}
\pair{u q_{\sigma-1}}{A_{\sigma_1}A_{\sigma_2}\cdots A_{\sigma_k}}
&=\sum_{i=1}^{k}(-1)^{\eps_i}\pair{u A^{n-\sigma_i}}{p_{n-\sigma+1}A_{\sigma_1}\cdots A_{\sigma_{i-1}}A_{\sigma_{i+1}}\cdots A_{\sigma_k}},\\
\eps_i&=(\sigma_1+\sigma_2+\cdots+\sigma_{i-1})\sigma_i+(\sigma_i+1)(n+\sigma)
\end{aligned}
\tag{21}
\end{equation}
を得る。
\emph{定理 II.} (20.) と (21.) により、列 $A^{\rho_i},x,p$、それぞれ $A_{\sigma_i},u,q$ の間の分解不能な恒等式の全体が尽くされる。同時に、置換 (11.) を行わずに明示的表示を許す分解不能な恒等式の全体も尽くされる。

定理の第一部は (20.) と (21.) へ至る考察から従う。第二部は、二つの記号列の間の恒等式
\[
\pair{q_\rho A^\sigma}{B_{\rho+\sigma-\tau}p_\tau}
=\eps_1\pair{q_\rho B^{n-\rho-\sigma+\tau}}{A_{n-\sigma}p_\tau}
+\eps_2\pair{q_\rho q_{n-\tau}}{B_{\rho+\sigma-\tau}A_{n-\sigma}},
\]
が不可能であることから従う。ただし $\rho\leqq\tau\leqq\sigma$ で、どの列も重み $1\cdot(n-1)$ をもたないとする。（$\rho,\tau,\sigma$ に関する他の仮定は、列の名称を入れ替えることでこれに帰着できる。）この不可能性は、たとえば特殊化 $q_\rho=lM^{\rho-1}$, $q_{n-\tau}=lN^{n-\tau-1}$ のもとで (8.) を展開することにより従う。ここで $l_1=1$, $M^{2,3,\ldots\rho}=1$, $A^{\rho+1,\ldots\rho+\sigma}=1$ とし、列 $l,M,A$ の他のすべての要素を零にし、$N$ と $B$ は任意とする。任意の列の間の恒等式は、列 $p_\tau$ を $y^{(1)}y^{(2)}\cdots y^{(\tau)}$ に分解すると、(20.) から合成された恒等式へ移る。したがって、(20.) が二つの記号列だけをもつ恒等式から生成できるならば、第二部は証明される。実際、
\begin{equation}
\pair{A^{\rho_1}A^{\rho_2}}{xp_{\rho-1}}
=\eps\pair{A^{\rho_1}q_{n-\rho+1}}{A_{n-\rho_2}x}
+\pair{A^{\rho_1}}{xp_{\rho_1-1}},
\quad \text{ここで }p_{\rho_1-1}\sim A^{\rho_2}q_{n-\rho+1},
\tag{20a}
\end{equation}
から、$A^{\rho_1}=B^{\sigma_1}B^{\sigma_2}$ と特殊化すること、すなわち
\[
\pair{B^{\sigma_1}B^{\sigma_2}}{xp_{\rho_1-1}}
=\eps\pair{B^{\sigma_1}q_{n-\rho_1+1}}{B_{n-\sigma_2}x}
+\pair{B^{\sigma_1}}{xp_{\sigma_1-1}},
\quad \text{ここで } p_{\sigma_1-1}\sim B^{\sigma_2}q_{n-\rho_1+1},
\]
によって、三つの記号列をもつ恒等式 (20.) が得られる。さらに $B^{\sigma_1}=C^{\tau_1}C^{\tau_2}$ などと特殊化していけば、一般の恒等式 (20.) が得られる。したがって、二つの記号列と二つの任意の変数列の間の恒等式について明示的表示が存在しないことは、任意個の記号列の場合にもそのような表示が存在しないことを含意する。ただし、変数列の中に列 $x$ または $u$ が含まれない場合である。

\section*{\S\ 4. 分解恒等式.}
ここで、任意の記号列と変数列の間の恒等式の最も一般の場合を考える。その最終式では、置換 (11.) を行わずに、列 $p_\tau$ が個々の列 $y^{(1)}\cdots y^{(\tau)}$ に分解された形で現れる。このような恒等式を「分解恒等式」と呼ぶ。前節末の注意に従い、まず記号列が二つだけの場合についてそれらを決定する。すなわち式 $\pair{q_\rho A^\sigma}{B_{\rho+\sigma-\tau}p_\tau}$ を展開する。次を置く。
\begin{equation}
\begin{aligned}
q_{\rho-\alpha}^{(i_1\ldots i_\alpha)}&\sim p_{n-\rho}\,y^{(i_1)}\cdots y^{(i_\alpha)},\qquad
p_{\tau-\alpha}^{(i_1\ldots i_\alpha)}=y^{(i_{\alpha+1})}\cdots y^{(i_\tau)},\qquad
p_\tau=y^{(1)}\cdots y^{(\tau)}.
\end{aligned}
\tag{22}
\end{equation}
数 $\rho,\sigma,\tau$ に関しては四つの場合を区別しなければならない。
\[
\begin{array}{llll}
1)&\rho+\sigma\leqq n,& \tau\leqq\rho,& \tau\leqq\sigma;\\[2pt]
2)&\rho+\sigma\leqq n,& \tau\geqq\rho,& \tau\leqq\sigma;\\[2pt]
3)&\rho+\sigma\leqq n,& \tau\leqq\rho,& \tau\geqq\sigma;\\[2pt]
4)&\rho+\sigma\leqq n,& \tau\geqq\rho,& \tau\geqq\sigma.
\end{array}
\]
第一の場合を取り上げ、他の場合はこれから簡単に特徴づける。(20.), (10.), (9.) により、列 $y^{(1)}y^{(2)}\cdots y^{(\tau)}B_{\rho+\sigma-\tau}$ を $y^{(1)}$ と $y^{(2)}\cdots y^{(\tau)}B_{\rho+\sigma-\tau}$ に分解すると、
\begin{equation}
\begin{aligned}
\pair{q_\rho A^\sigma}{y^{(1)}y^{(2)}\cdots y^{(\tau)}B_{\rho+\sigma-\tau}}
&=\pair{q_{\rho-1}^{(1)}A^\sigma}{y^{(2)}\cdots y^{(\tau)}B_{\rho+\sigma-\tau}}\\
&\quad+(-1)^\rho\pair{q_\rho[A_{n-\sigma}y^{(1)}]^{\sigma-1}}{y^{(2)}\cdots y^{(\tau)}B_{\rho+\sigma-\tau}}
\end{aligned}
\tag{23}
\end{equation}
を得る。(23.) の最後の項は、もとの式とまったく同じ形をしている。ただし $\sigma$ は $\sigma-1$ に、$\tau$ は $\tau-1$ に置き換わっている。したがってそれは再び方程式 (23.) を許す。よって、$\tau\leqq\sigma$ であるから、操作 (23.) を $\tau$ 回適用でき、(10.) を考慮して次の関係を得る。
\begin{equation}
\begin{aligned}
\pair{q_\rho A^\sigma}{p_\tau B_{\rho+\sigma-\tau}}
&=(-1)^{\rho\sigma+(\sigma+\tau)(n+1)}
  \pair{q_\rho B^{n-\rho-\sigma+\tau}}{A_{n-\sigma}p_\tau}\\
&\quad+\sum_{i_1=1}^{\tau}(-1)^{\rho(i_1-1)}
 \pair{q_{\rho-1}^{(i_1)}[A_{n-\sigma}y^{(1)}y^{(2)}\cdots y^{(i_1-1)}]^{\sigma-i_1+1}}
 {y^{(i_1+1)}\cdots y^{(\tau)}B_{\rho+\sigma-\tau}}.
\end{aligned}
\tag{24}
\end{equation}
和記号の下の各項もまたもとの式の型をしている。$\rho$ は $\rho-1$ に、$\sigma$ は $\sigma-i_1+1$ に、$\tau$ は $\tau-i_1$ に置き換わるので、条件 1) はさらに強く満たされる。$\tau\leqq\rho$ であるから、操作 (24.) を $\tau$ 回適用でき、求める分解恒等式を得る。
\begin{equation}
\begin{aligned}
\pair{q_\rho A^\sigma}{p_\tau B_{\rho+\sigma-\tau}}
&=\eps_{i_0}\pair{q_\rho B^{n-\rho-\sigma+\tau}}{A_{n-\sigma}p_\tau}\\
&\quad+\sum_{i_1=1}^{\tau}\eps_{i_1}
 \pair{q_{\rho-1}^{(i_1)}B^{n-\rho-\sigma+\tau}}{A_{n-\sigma}p_{\tau-1}^{(i_1)}}\\
&\quad+\sum_{i_2=i_1+1}^{\tau}\eps_{i_2}
 \pair{q_{\rho-2}^{(i_1i_2)}B^{n-\rho-\sigma+\tau}}{A_{n-\sigma}p_{\tau-2}^{(i_1i_2)}}+\cdots\\
&\quad+\eps_{i_\tau}\pair{q_{\rho-\tau}^{(i_1i_2\ldots i_\tau)}B^{n-\rho-\sigma+\tau}}{A_{n-\sigma}}\\
&=\sum_{\alpha=0}^{\tau}\sum_{i_\alpha=i_{\alpha-1}+1}^{\tau}\eps_{i_\alpha}
 \pair{q_{\rho-\alpha}^{(i_1\ldots i_\alpha)}B^{n-\rho-\sigma+\tau}}{A_{n-\sigma}p_{\tau-\alpha}^{(i_1\ldots i_\alpha)}}.
\end{aligned}
\tag{25}
\end{equation}
符号因子 $\eps$ については (24.) から
\begin{equation}
\begin{aligned}
\eps_{i_\alpha}&=(-1)^{\delta_{i_\alpha}},\qquad
\delta_{i_\alpha}=\sum_{\beta=1}^{\alpha}(\rho-\beta+1)(i_{\beta-1}+i_\beta+1)\\
&\quad +(\rho-\alpha)(\sigma+i_\alpha+\alpha)+(\sigma+\tau+\alpha)(n+1)
\end{aligned}
\tag{26}
\end{equation}
を得る。ここで $i_0=0$ と置く。 (25.) の和記号は次のように理解する。$i_1<i_2<\cdots<i_{\alpha-1}$ である各組合せ $i_1i_2\ldots i_{\alpha-1}$ に対し、$i_{\alpha-1}<i_\alpha\le\tau$ を満たすすべての値 $i_\alpha$ を付け加える。(26.) を考慮すると、(25.) の個々の和が (19.) に類似する微分方程式
\begin{equation}
 \left(\frac{\partial}{\partial y^{(i)}}\middle|y^{(i+1)}\right)=0,
 \qquad (i=1,2,\ldots\tau-1)
\tag{27}
\end{equation}
を満たすことは容易に確かめられる。したがってそれらは列 $p_\tau$ のみに依存する。よって分解恒等式は列 $A,B,q,p$ の間の分解不能な恒等式である。その明示的表示には次節で戻る。

場合 2) では、操作 (23.) は同じく $\tau$ 回適用できるが、操作 (24.) は $\rho$ 回目の適用後に止まる。したがって (25.) の最終式では、第一の和記号を $\sum_{\alpha=0}^{\rho}$ に置き換えなければならない。場合 3) と 4) では、操作 (23.) は $\sigma$ 回だけ行える。そのとき (24.) の代わりに次の関係をもつ。
\begin{equation}
\begin{aligned}
\pair{q_\rho A^\sigma}{p_\tau B_{\rho+\sigma-\tau}}
&=(-1)^{\rho\sigma}
 (q_{\rho-(\tau-\sigma)}^{(\sigma+1\ldots\tau)}B^{n-\rho+\tau-\sigma})(A_{n-\sigma}p_{\tau-(\tau-\sigma)}^{(\sigma+1\ldots\tau)})\\
&\quad+\sum_{i_1=1}^{\sigma}(-1)^{\rho(i_1-1)}
 \pair{q_{\rho-1}^{(i_1)}[A_{n-\sigma}y^{(1)}\cdots y^{(i_1-1)}]^{\sigma-i_1+1}}
 {y^{(i_1+1)}\cdots y^{(\tau)}B_{\rho+\sigma-\tau}};
\end{aligned}
\tag{28}
\end{equation}
(28.) を $\tau-\sigma$ 回繰り返す。このとき和記号は、和記号下の項に (23.) を $(\sigma-i_1+1)$ 回適用することからわかるように、$\sum_{i_\beta=i_{\beta-1}+1}^{\sigma+\beta-1}$ になる。これにより、(25.) の個別和のうち添字 $\alpha=\tau-\sigma$ に対応するすべての項が、対応する符号とともに得られる。すると場合 3) と 4) は場合 1) と 2) に移る。なぜなら $\sigma,\tau$ に対応する数は $\sigma'=\sigma-i_{\tau-\sigma}+\tau-\sigma$, $\tau'=\tau-i_{\tau-\sigma}=\sigma'$ となり、手続きが続くにつれて $\tau'<\sigma'$ となるからである。したがって (25.) の最終式では、第一の和記号を $\sum_{\alpha=\tau-\sigma}^{\tau,\rho}$ に置き換えなければならない。

同様に (21.) からは、列 $q_\rho$ が個々の列 $v^{(1)},\ldots,v^{(\rho)}$ に分解される双対的な分解恒等式が得られる。もし
\begin{equation}
 \dot p_{\tau-\beta}^{(i_1\ldots i_\beta)}\sim v^{(i_1)}\cdots v^{(i_\beta)}q_{n-\tau};\qquad
 \dot q_{\rho-\beta}^{(i_1\ldots i_\beta)}=v^{(i_{\beta+1})}\cdots v^{(i_\rho)};\qquad
 q_\rho=v^{(1)}\cdots v^{(\rho)},
\tag{29}
\end{equation}
と置けば、
\begin{equation}
 \pair{q_\rho A^\sigma}{p_\tau B_{\rho+\sigma-\tau}}
 =\sum_{\beta=\max(0,\tau-\sigma)}^{\min(\tau,\rho)}
   \sum_{i_\beta=i_{\beta-1}+1}^{\rho}\eps
\pair{\dot q_{\rho-\beta}^{(i_1\ldots i_\beta)}B^{n-\rho-\sigma+\tau}}{A_{n-\sigma}\dot p_{\tau-\beta}^{(i_1\ldots i_\beta)}} ,
\tag{30}
\end{equation}
となる。ここで和の添字 $\beta$ は、下限の大きい方から上限の小さい方まで走る。(25.) と (30.) における個別和のすべての項は、一つの形式の極であり、それらは極化過程
\[
\left(\frac{\partial}{\partial p_{\tau-\alpha}}\middle|p_{\tau-\alpha}^{(i_1\ldots i_\alpha)}\right),\qquad
\left(q_{\rho-\alpha}^{(i_1\ldots i_\alpha)}\middle|\frac{\partial}{\partial q_{\rho-\alpha}}\right)
\]
によって作られる。これを表すため、こうした極化操作を定数倍して集めたものを $\Pi$ と書く。すると (25.) と (30.) の代わりに、
\begin{equation}
 \pair{q_\rho A^\sigma}{p_\tau B_{\rho+\sigma-\tau}}
 =\sum_{\alpha=\max(0,\tau-\sigma)}^{\min(\tau,\rho)}
 \Pi\pair{q_{\rho-\alpha}B^{n-\rho-\sigma+\tau}}{A_{n-\sigma}p_{\tau-\alpha}}
\tag{31}
\end{equation}
を得る。

\section*{\S\ 5. 明示的な形での分解恒等式.}
(25) の場合 4) で $\tau$ を $\sigma+\rho$ に特殊化すると、
\begin{equation}
 \pair{q_\rho A^\sigma}{y^{(1)}\cdots y^{(\rho+\sigma)}}
 =\sum_{1\le i_1<\cdots<i_\rho\le \rho+\sigma}
 \eps_{i_\rho}\pair{q_\rho}{y^{(i_1)}\cdots y^{(i_\rho)}}
 \pair{A^\sigma}{y^{(i_{\rho+1})}\cdots y^{(i_{\rho+\sigma})}},
\tag{32}
\end{equation}
を得る。ここで $\eps_{i_\rho}=(-1)^{\delta_{i_\rho}}$ かつ $\delta_{i_\rho}=\rho(\rho+1)/2+i_1+\cdots+i_\rho$ である。これは (16.) と (17.) より一般的な行列の積定理の形であり、積行列の行列式
\[
 \sum\pm(v^{(1)}y^{(1)})\cdots(v^{(\rho)}y^{(\rho)})(a^{(1)}y^{(\rho+1)})\cdots(a^{(\sigma)}y^{(\rho+\sigma)})
\]
をラプラス展開することで従う。したがって、より一般の積定理は恒等式 (20.) と (21.) からそれぞれ合成され、特殊な形 (16.), (17.) の選択もこれによって説明される。

関係 (32.) は、分解恒等式を明示的に表示する手段を与える。このため、(31.) に現れる形式を基礎形式と見なす。すなわち置換 (11.) を行う。この置換は、列 $A,B$ における明示的表示を何も変えない。なぜなら、(8.) により、新しく導入される列 $R$ は、成分列 $a^{(1)}a^{(2)}\ldots$, $b^{(1)}b^{(2)}\ldots$ ではなく、列 $A,B$ そのものによって表せるからである。そこで
\begin{equation}
 \pair{q_{\rho-\alpha}B^{n-\rho-\sigma+\tau}}{A_{n-\sigma}p_{\tau-\alpha}}
 =(R_{\rho-\alpha}p_{n-\rho+\alpha})\pair{R^{\tau-\alpha}}{p_{\tau-\alpha}}
\tag{33}
\end{equation}
と置く。すると (25.) の添字 $\alpha$ に対応する個別和について、
\begin{equation}
\begin{aligned}
&\sum \eps_{i_\alpha}\pair{R_{\rho-\alpha}p_{n-\rho}y^{(i_1)}\cdots y^{(i_\alpha)}}{R^{\tau-\alpha}y^{(i_{\alpha+1})}\cdots y^{(i_\tau)}}\\
&\quad=\sum \eps_{i_\alpha}\pair{[R_{\rho-\alpha}p_{n-\rho}]^\alpha y^{(i_1)}\cdots y^{(i_\alpha)}}{R^{\tau-\alpha}y^{(i_{\alpha+1})}\cdots y^{(i_\tau)}}\\
&\quad=\eps\pair{[R_{\rho-\alpha}p_{n-\rho}]^\alpha R^{\tau-\alpha}}{p_\tau}
 =\eps\pair{R^{\tau-\alpha}q_{n-\tau}}{R_{\rho-\alpha}p_{n-\rho}}
\end{aligned}
\tag{34}
\end{equation}
を得る。このように実際に明示的な最終式が得られ、また $q_\rho$ と $p_\tau$ が現れる対称的な形は、(30.) も同じ最終式へ導くことを示す。したがってこの最終式はもとの分解に依存しない。(25.) と (30.) の間に差があるのは、列 $y$ と $v$ が独立の意味をもっている間だけであり、そのとき分解恒等式は我々の定義により分解可能な恒等式へ移る。より多くの記号列をもつ恒等式へ進むには、\S\ 3（末尾）に従って、列 $q_\rho$ を $q_{\rho_1}C^{\sigma_1}$ に、または $p_\tau$ を $p_{\tau_1}C_{\tau-\tau_1}$ に分解すればよい。得られる式
\begin{equation}
 \pair{q_{\rho_1}C^{\sigma_1}}{[R^{\tau-\alpha}q_{n-\tau}]_\alpha R_{\rho_1+\sigma_1-\alpha}},
 \qquad
 \pair{[R_{\rho-\alpha}p_{n-\rho}]^\alpha R^{\tau-\alpha}}{p_{\tau_1}C_{\tau-\tau_1}}
\tag{35}
\end{equation}
は、二つの記号列の場合に現れる型とまったく同じであり、したがって (33.) に対応する置換の後に再び明示的表示を許す。これを繰り返すことにより、任意個の列について明示的表示を得る。また、(25.) および (30.) の最初と最後の和は、それぞれ (33.) を適用することなく、公式 (32.) だけで明示的表示を与えるか、あるいはすべての列を明示的に含む単一項に完全に還元されることも注意しておく。(32.) から、$q_\rho$ を $q_{\rho_1}C^{\sigma_1}$ などに分解して得られる関係は、最も一般的な形の積定理と見なすことができる。

要約すれば、次を示す。

\emph{定理 III:} 恒等式
\begin{equation}
 \pair{q_\rho A^\sigma}{p_\tau B_{\rho+\sigma-\tau}}
 =\sum_{\alpha=\max(0,\tau-\sigma)}^{\min(\tau,\rho)}
 \eps\pair{R^{\tau-\alpha}q_{n-\tau}}{R_{\rho-\alpha}p_{n-\rho}}
\tag{36}
\end{equation}
と、ここで列 $R$ は (33) によって決定されるものとし、さらに $q_\rho,p_\tau$ を $q_{\rho_1}C^{\sigma_1}$ などへ分解してそこから得られる恒等式を合わせると、分解不能な恒等式の全体が尽くされる。

実際、得られる恒等式は、その種類のうち最も一般的である。なぜなら、個々の列は、(20.) と (21.) の場合のように重みや数に関する制限を受けないからである。また、これらの列の間にはそれ以上の恒等式は存在しない。列 $p_\tau$ を $y^{(1)}\cdots y^{(\tau)}$ に分解すると、最も一般の恒等式は (20.) から、したがって \S\ 3（末尾）により (20a.) から合成される。そして (16.) の後に論じた合成は、まさに \S\ 4 で行われた合成である。任意個の列への移行は、(20a.) の議論が示すように、$q_\rho$ を $q_{\rho_1}C^{\sigma_1}$ などに分解して生じる式へ、常に同じ操作を適用することで与えられる。また、(20a.) ではなく (20.) から直接移行しても新しいものは得られない。これは計算によって容易に確かめられる。一方では (25.) において列 $q_\rho$ を一度特殊化し、他方では \S\ 4 の方法で (20.) から移行を行う。どちらの場合にも同じ和が生じ、それらは一般の積定理（上を見よ）の適用により、(36.) から得られる恒等式へ移る。恒等式 (20.), (21.) は特殊な場合 $\tau=1$, $\rho=1$ に対応する。そのとき個別和は二つだけ現れ、したがって前の注意に従って、置換 (33.) なしに、以前に得たものと一致する明示的表示が得られる。

最後に、分解恒等式の二つの特殊例を簡単に述べる。(31) で $\tau=\sigma$, $\rho=n-\sigma$ と置き、列 $p_\tau$ を $p'_\sigma\sim q'_{n-\sigma}$ と書く。すると
\begin{equation}
\begin{aligned}
\pair{A^\sigma}{p_\sigma}\pair{B^\sigma}{p'_\sigma}
-\pair{A^\sigma}{p'_\sigma}\pair{B^\sigma}{p_\sigma}
=\sum_{\alpha=1}^{\sigma,n-\sigma}
\Pi\pair{q_{n-\sigma-\alpha}B^\sigma}{A_{n-\sigma}p_{\sigma-\alpha}}.
\end{aligned}
\tag{37}
\end{equation}
これは二つの共変的列 $p_\sigma,p'_\sigma$ をもつ形式を基本共変式へ展開する Clebsch の公式である。\footnote{Clebsch, loc. cit., pp. 25--32. Clebsch は、右辺の個々の項が列 $A,B$ のみに依存することを証明している。明示的表示は、彼の式 $\Phi_h$ に (32.) を適用することで得られる。} 形式 $\pair{A^\sigma}{p_\sigma}^m\pair{B^\sigma}{p'_\sigma}^n$ に属する基本共変式は次のものである。
\begin{equation}
 \left\{\sum_{\alpha=1}^{\sigma,n-\sigma}
 \pair{q_{n-\sigma-\alpha}B^\sigma}{A_{n-\sigma}p_{\sigma-\alpha}}
 \right\}^{\rho}
 \pair{A^\sigma}{p_\sigma}^{m-\rho}\pair{B^\sigma}{p_\sigma}^{n-\rho},
 \qquad (\rho=0,1,\ldots m,n).
\tag{38}
\end{equation}
これらは、簡単に言えば、基礎形式から「欠損数 $\alpha$ の共変的折り畳み」によって生じる（\S\ 2 参照）。

第二に、$\tau=\sigma-\lambda$, $\rho=n-\sigma-\lambda$, $B^\sigma=A^\sigma$ と置く。すると
\begin{equation}
\begin{aligned}
\pair{q_{n-\sigma-\lambda}A^\sigma}{A_{n-\sigma}p_{\sigma-\lambda}}
&=(-1)^\lambda\pair{q_{n-\sigma-\lambda}A^\sigma}{A_{n-\sigma}p_{\sigma-\lambda}}\\
&\quad+(-1)^{(n-\sigma)(\sigma-\lambda)}
 \sum_{\alpha=1}^{\sigma-\lambda,n-\sigma-\lambda}
 \Pi\pair{q_{n-\sigma-\lambda-\alpha}A^\sigma}{A_{n-\sigma}p_{\sigma-\lambda}}.
\end{aligned}
\tag{39}
\end{equation}
したがって、\S\ 2 で述べたように、奇数の欠損 $\lambda$ に対して記号列を特徴づける条件 (12.) は、より高い欠損の条件から従う。線形形式 $\pair{A^\sigma}{p_\sigma}\equiv\pair{B^\sigma}{p_\sigma}$ については、(39.) は、係数について二次の既約な不変形成が偶数欠損のものに還元されることを含意する。これは、二元および三元領域で線形形式が不変形成をもたないという既知の事実と一致する。

一般の恒等式は、個々の列が条件 (12.) を満たすという仮定のもとで初めに導かれたことに注意すべきである。しかし、これらの個々の列は恒等式に線形にしか入らないので、恒等式は条件 (12.) を満たさないより一般の列についても成り立つ。

\section*{\S\ 6. 不変形成の明示的表示.\\折り畳み過程と微分過程.}
前節までの結果により、\S\ 2 で述べた記号的表示可能性の定理（第一基本定理）を、明示的表示可能性を加えて完成できる。すなわち次を示す。

\emph{定理 IV:} 「任意の基礎形式のすべての不変形成は行列積によって明示的に表示できる。すなわち、変数列 $p_\rho$ のほかには、もとの形式の列 $S^\sigma,\ldots T^\tau$、または置換 (9.) あるいは (11.) によってそれらから導かれる列 $P,Q,R$ だけが最終式に入る。」

証明のため、ある不変形成が、変数列のほかに列 $S^\sigma,\ldots T^\tau$ に依存するとする。また、これらの列が、行列式による表示を可能にする程度まで個々の列 $S^\sigma=S^{\sigma_1}\cdots S^{\sigma_i},\ldots T^\tau=T^{\tau_1}\cdots T^{\tau_k}$ に分解されているとする。列を、任意ではあるが一度限り固定された順序 $S^\sigma,\ldots T^\tau$ に並べる。最初の全体列 $S^\sigma=S^{\sigma_1}\cdots S^{\sigma_i}$ に依存し、個々の成分列だけには依存しない行列式の集合体を取り出す。この依存性は、\S\ 3--\S\ 5 により一般恒等式の帰結である。ここで変数列 $p_\rho$ を列 $y,v$ に、または後に現れる列 $T$ の一つを個々の列 $t$、すなわち $\tau$ に分解すると、最も一般の恒等式は (20.) と (21.) から合成される。しかしこれらの恒等式は常に、列 $S^\sigma=S^{\sigma_1}\cdots S^{\sigma_i}$ を明示的に含む式へ導く。すなわち (20.) と (21.) の左辺である。というのも、(19.) により、右辺の全和だけが、(20a.) の一項に対応する一部分ではなく、$S^\sigma$ のみに依存する性質をもつことが確認されるからである。手続きが続くにつれて、一度明示的に入ったすべての列は明示的なまま残る。恒等式の適用が要求しうるのは、せいぜい複数の列を一つにまとめること、すなわち置換 (9.) だけである。最後に、一つの列 $T^\tau$ だけを明示的な形にしなければならない場合には、二つの場合がある。まだ自由な変数列、すなわち置換 (9.) によって他の列と結びついていない列が残っているかもしれない。その列を成分列 $y$ または $v$ に分解すれば手続きは完了し、(31.) のように、すべての列を明示的に含む一つまたは複数の形式の極が得られる。自由な変数列が残っていなければ、個々の列 $t$ または $\tau$ を含むが列 $T^\tau$ に依存する式全体の中に、それらが生じた仕方から、分解恒等式の項を認めることができる。しかし \S\ 5 により、これらは置換 (11.) を適用した後、常にすべての列において明示的な表示をもつ。これで定理は完全な一般性において証明される。二つの記号列だけが現れる場合には、置換 (11.) は決して現れないことも注意しておく。なぜなら、$\sigma,n-\sigma,\tau,n-\tau<n$ なので、変数列がまったく入らないのは二つの記号列が一つの行列式に結合され、したがって明示的に現れる場合だけであり、しかし変数列が入るやいなや、(34.) と (33.) は置換 (11.) が不要になることを示すからである。

今述べたことから、すべての不変形成を生成するには、変数列を付け加えたのち、記号列を可能なすべての行列積に結合すれば足りることが従う。最も一般の行列積は今や
\begin{equation}
 \pair{P^{\mu_1}\cdots P^{\mu_i}}{Q_{\nu_1}\cdots Q_{\nu_k}},\qquad \sum\mu=\sum\nu,
\tag{40}
\end{equation}
の形をしている。ここで $P$ と $Q$ は、もとの形式の列でもよく、もとの列から置換 (9.) または (11.) の列を経て生じたものでもよく、最後には変数列でもよい。しかし式 (40.) は、変数列の助けを借りて、二つの記号列を一度に一つの行列積に結合していくことで生成できる。実際、
\[
 \pair{q_{n-\mu-\lambda}P^\mu}{Q_{\nu_1}p_{n-\nu_1-\lambda}}=(R_\lambda Q_{\nu_1}p_{n-\nu_1-\lambda})
\]
から、列 $Q_{\nu_2}$ と結合することにより、
\[
 \pair{[R_\lambda Q_{\nu_1}]^{n-\lambda-\nu_1}}{Q_{\nu_2}p_{n-\nu_1-\nu_2-\lambda}}
 =\pair{q_{n-\mu-\lambda}P^\mu}{Q_{\nu_1}Q_{\nu_2}p_{n-\nu_1-\nu_2-\lambda}}
\]
という式が得られる。したがってこの手続きを続ければ式 (40.) が得られる。$P$ と $Q$ がもとの形式の列でない場合にも、(9.) と (11.) は、今示したことと合わせて、それらも二つの記号列を一度に結合する列によって生じたことを意味する。したがって、

\emph{定理 V:} すべての不変形成を生成する過程、すなわち折り畳み過程は、変数列の助けを借りて、二つの記号列を一度に一つの行列積に結合していくことから成る。

二つの記号列 $S^\sigma,T^\tau$ について、$\sigma\geqq\tau$ とすると、次の行列積が作れる。
\begin{align}
&\pair{q_{n-\sigma-\lambda}S^\sigma}{T_{n-\tau}p_{\tau-\lambda}}; &&(\lambda=1,2,\ldots\tau,n-\sigma)\tag{41}\\
&\pair{q_{n-\tau-\mu}T^\tau}{S_{n-\sigma}p_{\sigma-\mu}}; &&(\mu=\sigma-\tau+\lambda)\tag{42}\\
&\pair{q_{n-\tau-\nu}T^\tau}{S_{n-\sigma}p_{\sigma-\nu}}. &&(\nu=1,2,\ldots\sigma-\tau).
\tag{43}
\end{align}
しかし (31.) から、(42.) と (43.) について次の値が得られる。
\begin{align}
&\pair{q_{n-\sigma-\lambda}S^\sigma}{T_{n-\tau}p_{\tau-\lambda}}
 +\sum_{\alpha}\Pi\pair{q_{n-\sigma-\lambda-\alpha}S^\sigma}{T_{n-\tau}p_{\tau-\lambda-\alpha}},\tag{42a}\\
&\Pi(q_{n-\sigma}S^\sigma)(T_{n-\tau}p_\tau)
 +\sum_{\beta}\Pi\pair{q_{n-\sigma-\beta}S^\sigma}{T_{n-\tau}p_{\tau-\beta}}.
\tag{43a}
\end{align}
したがって、形式 (42.) は (41.) とその極によって線形に表され、形式 (43.) は基礎形式と (41.) の形式の極によって表される。同様に (31.) は、形式 (41.) が (42.) とその極に線形に依存することを与える。したがって Clebsch の定義\footnote{Clebsch, loc. cit., \S\ 7.}によれば、形式 (42.) は形式 (41.) と同値であり、(43.) は基礎形式へ戻る。したがって生成過程は (41.) または (42.) によって同じ権利で与えられる。我々は、数 $\lambda$ に関してより簡単な過程 (41.) を折り畳み過程と定義する。数 $\lambda$、すなわち折り畳みの欠損数（\S\ 2 参照）は、行列式積から欠けている列の数を与える。より正確には、その折り畳みにおいて最大の列数を含む行列積において欠けている列数である。$\sigma\ge\tau$ のもとで、$\lambda$ は二つの値 $\tau,n-\sigma$ の小さい方までしか走らないので、
\begin{equation}
 \lambda\le\frac n2,\ n\text{ は偶数};\qquad
 \lambda\le\frac{n-1}{2},\ n\text{ は奇数}.
\tag{44}
\end{equation}
となる。(42.) と (43.) の (41.) への還元は、さらなる折り畳みによって列 $p,q$ が記号列に移った場合にも有効である。実際、(36.) により
\[
 \pair{A^{n-\tau-\varkappa}T^\tau}{S_{n-\sigma}B_{\sigma-\varkappa}}
 =\sum_{\alpha=0,\varkappa-\sigma+\tau}^{\tau,n-\sigma}
 \eps\pair{R^{\tau-\alpha}B^{n-\sigma+\varkappa}}{A_{\tau+\varkappa}R_{n-\sigma-\alpha}}
\]
をもつ。ここで列 $R$ は (33.) に従って $S$ と $T$ から生じたものである。したがって、記号列だけを含むこれらの形式は、$\sigma-\tau\gtreqless 2\varkappa$ に応じて、$\pair{q_\tau A^{n-\tau-\varkappa}}{B_{\sigma-\varkappa}p_{n-\sigma}}$ または $\pair{q_{\tau-\alpha}B^{n-\sigma+\varkappa}}{A_{\tau+\varkappa}p_{n-\sigma-\alpha}}$ を、基礎形式および形式 (41.) と折り畳むことによって得られる。

記号的折り畳み過程からは、通常どおり、記号列 $S^\sigma,T_{n-\tau}$ を微分記号の列 $\frac{\partial}{\partial p_\sigma}$, $\frac{\partial}{\partial q_{n-\tau}}$ で置き換えるだけで、不変な微分過程が得られる。よく知られているように、積 $\frac{\partial}{\partial p_{i_1\ldots i_\sigma}}\cdot\frac{\partial}{\partial q_{j_1\ldots j_{n-\tau}}}$ は、そのとき実際には $\frac{\partial^2}{\partial p_{i_1\ldots i_\sigma}\partial q_{j_1\ldots j_{n-\tau}}}$ を意味する。列 $\frac{\partial}{\partial p_\sigma}$ と $\frac{\partial}{\partial q_{n-\tau}}$ は列 $S^\sigma,T_{n-\tau}$ と共変的であるから、それらと同じ恒等式を満たし、とくに (42.) と (42a.) の間、および (43.) と (43a.) の間に存在する恒等式も満たす。要約すれば、次を得る。

\emph{定理 VI:} 任意の基礎形式のすべての不変形成は、それらから、記号的折り畳み過程 (41.) を反復適用することによって、または不変な微分過程
\begin{equation}
 \left(q_{n-\sigma-\lambda}\frac{\partial}{\partial p_\sigma}\middle|\frac{\partial}{\partial q_{n-\tau}}p_{\tau-\lambda}\right),
 \qquad(\sigma\geqq\tau,\ \lambda=1,2,\ldots\sigma,n-\tau)
\tag{45}
\end{equation}
によって得られる。
なお、各折り畳み (41.) において形式の総重量、すなわち個々の変数列の重みの和（\S\ 1 参照）は、正の数 $2(\lambda^2+\lambda(\sigma-\tau))$ だけ減少する。これは、一般の有限性証明とは独立に、形式列の有限性、すなわち、与えられた基礎形式の係数に線形な不変形成全体の有限性を含意する。\footnote{Clebsch, loc. cit., \S\ 11 と 12、基本共変式の系の有限性を参照。}

さらに、定理 VI は条件 (12.) を満たさない形式についてもなお成り立つことを注意しておく。(3) の各因子 $\pair{S^\rho}{p_\rho}^{m_\rho}$ は、$m_\rho$ 個の線形因子の積 $\pair{A^\rho}{p_\rho}\pair{B^\rho}{p_\rho}\cdots\pair{M^\rho}{p_\rho}$ で置き換えることができる。その場合、不変形成は線形形式の系の不変形成へ移り、後でそれらを互いに等しいものとすればよい。しかし各単独の線形形式については、条件 (12.) は、微分記号を導入すると二階の微分商だけを含むものであり、常に満たされる。

\section*{\S\ 7. 一般形式の正規形による展開.}
以下（\S\ 8）では、折り畳み過程 (41.) の主要な性質、すなわちそれが基本的折り畳みから合成できることを導く。そのために、四元領域で Mertens\footnote{F. Mertens, \emph{"Uber invariante Gebilde quatern"arer Formen}（序論参照）。以下「Mertens」と引用する。} が与えた重要な定理を $n$ 元領域へ移す必要がある。この定理は、任意の有限個の形式の系を、同値な正規形の系で置き換えることができるというものである。「正規形」とは、二元および三元の定義\footnote{Study, loc. cit., p. 54.}を一般化して、係数に線形な不変形成が、基礎形式自身を除いて恒等的に消える形式をいう。したがって定理 VI（\S\ 6）により、正規形 $\varphi$ は微分方程式系
\begin{equation}
 \left(q_{n-\sigma-\lambda}\frac{\partial}{\partial p_\sigma}\middle|\frac{\partial}{q_{n-\tau}}p_{\tau-\lambda}\right)=0,
 \qquad(\sigma\geqq\tau;\ \lambda=1,2,\ldots\tau,n-\sigma)
\tag{46}
\end{equation}
を満たす。ここで列 $q_{n-\sigma-\lambda},p_{\tau-\lambda}$ は任意量と見なす。四元領域では、$\sigma=\tau=2$ について (39.) を考慮すれば、これらの微分方程式は、Mertens が任意の列 $p,q$ を導入せずに与えた十個の微分方程式（loc. cit., 公式 28）と完全に一致する。それらを知ることと定理 VI とにより、彼の証明をほとんど字句どおり $n$ 変数へ移すことができる。付け加えるべきことはただ一つ、構成された形式が正規形であることの検証であり、これは分解恒等式 (30.) によって得られる。四元領域ではこの検証はさらなる変形なしに従っていた。したがって、ここでは証明の短い概説を与えれば十分である。簡単のため、以下で必要となる場合、すなわち一つの基礎形式 $f$ とその形式列（\S\ 6 末尾参照）の場合を扱う。任意個の基礎形式の係数に任意次数で現れる不変形成についても、証明は同じである。

証明の基本的な考えは、変換された係数を導入することにより、$n-1$ 個の列 $p_\rho$ をもつ形式を、$x$ と共変的な $n$ 個の列 $\xi$ をもつ形式に置き換えることである。これらの列は変換量の列である。求める正規形は、これらの形式から不変な微分過程によって直接得られる。一般形式の正規形による展開は、$\rho$ 個（$\rho\leqq n$）の共変的列 $\xi$ をもつ形式の級数展開によって媒介される。不変な微分過程は、元の列 $p_\rho$ における形式へ戻す。

変換量の列を $\xi^{(1)},\ldots\xi^{(n)}$ とし、
\begin{equation}
\begin{aligned}
x_i&=\xi_i^{(1)}x'_1+\xi_i^{(2)}x'_2+\cdots+\xi_i^{(n)}x'_n,\\
p_{i_1\ldots i_\rho}&=\sum_j(\xi_{i_1}^{(j_1)}\cdots \xi_{i_\rho}^{(j_\rho)})p'_{j_1\ldots j_\rho}
\end{aligned}
\tag{47}
\end{equation}
であるとする。形式 $f,\varphi,\ldots$ の変換された係数を $(f),(\varphi),\ldots$ と書き、
\[
 f\equiv\pair{S^1}{p_1}^{m_1}\pair{S^2}{p_2}^{m_2}\cdots\pair{S^{n-1}}{p_{n-1}}^{m_{n-1}}
\]
とする。すると
\begin{equation}
\begin{aligned}
(f)&=\pair{S^1}{\xi^{(1)}}^{m_{11}}\cdots\pair{S^1}{\xi^{(n)}}^{m_{1n}}
\pair{S^2}{\xi^{(1)}\xi^{(2)}}^{m_{21}}\cdots
\pair{S^{n-1}}{\xi^{(2)}\cdots\xi^{(n)}}^{m_{n-1,n}}\\
&=\pair{s^{(1)}}{\xi^{(1)}}^{k_1}\pair{s^{(2)}}{\xi^{(2)}}^{k_2}\cdots\pair{s^{(n)}}{\xi^{(n)}}^{k_n}
\end{aligned}
\tag{48}
\end{equation}
である。ここで $m_{11}+\cdots+m_{1n}=m_1$ などである。各係数 $(f)$ について、もし
\begin{equation}
 k_1\geqq k_2\geqq k_3\cdots\geqq k_n
\tag{49}
\end{equation}
ならば、列 $\xi$ をちょうど尽くす微分過程
\begin{equation}
\left(\frac{\partial}{\partial\xi^{(1)}}\middle|p_1\right)^{k_1-k_2}
\left(\frac{\partial}{\partial\xi^{(1)}}\frac{\partial}{\partial\xi^{(2)}}\middle|p_2\right)^{k_2-k_3}
\cdots
\left(\frac{\partial}{\partial\xi^{(1)}}\cdots\frac{\partial}{\partial\xi^{(n-1)}}\middle|p_{n-1}\right)^{k_{n-1}-k_n}
\left(\frac{\partial}{\partial\xi^{(1)}}\cdots\frac{\partial}{\partial\xi^{(n)}}\right)^{k_n}
\tag{50}
\end{equation}
によって、形式
\begin{equation}
 \varphi=\bigl(s^{(1)}\mid p_1\bigr)^{k_1-k_2}
 \bigl(s^{(1)}s^{(2)}\mid p_2\bigr)^{k_2-k_3}\cdots
 \bigl(s^{(1)}\cdots s^{(n-1)}\mid p_{n-1}\bigr)^{k_{n-1}-k_n}
 \bigl(s^{(1)}s^{(2)}\cdots s^{(n)}\bigr)^{k_n}
\tag{51}
\end{equation}
が得られる。これらの形式 $\varphi$ が求める正規形である。実際、それらは正規形を定義する次の性質をもつ。

1. これらは基礎形式 $f$ の係数に線形な不変形成である。これは、それらが $f$ から不変な微分過程 $\left(\frac{\partial}{\partial p_\rho}\middle|\xi^{(i_1)}\xi^{(i_2)}\cdots\xi^{(i_\rho)}\right)$ と (50.) によって生じることから従う。したがって（\S\ 6 末尾により）それらは $f$ の有限な形式列の極によって線形に表される。これらの極は、$\varphi$ に入る変数 $p_\rho$ を、変数列 $\mathfrak{p}_\sigma$ をもつ線形形式に分解された形式の係数と見なすことで得られる。
\begin{equation}
 H=(q_{n-1}\mid \mathfrak{p}_{n-1})^{k_1-k_2}(q_{n-2}\mid \mathfrak{p}_{n-2})^{k_2-k_3}\cdots(q_1\mid \mathfrak{p}_1)^{k_{n-1}-k_n}.
\tag{52}
\end{equation}
形式列 $f$ と $H$ の同時不変式は、必要な極をすべて与える。なぜなら、それらの各列 $p_\rho$ の変数に関する次数は、$\varphi$ 自身の同じ変数に関する次数と等しくなければならないからである。

2. それらは微分方程式 (46.) を満たす。微分過程 (45.) を $\varphi$ に適用すると、
\[
 \varphi'=\bigl(q_{n-\sigma-\lambda}s^{(1)}\cdots s^{(\sigma)}\mid [s^{(1)}\cdots s^{(\tau)}]_{n-\tau}p_{\tau-\lambda}\bigr).\qquad (\sigma\geqq\tau)
\]
となる。列 $s$ から作られる列 $q_\sigma=s^{(1)}\cdots s^{(\sigma)}$, $p_{n-\tau}=[s^{(1)}\cdots s^{(\tau)}]_{n-\tau}$ を、そこに現れる $q_\rho,p_\tau$ と同一視すると、分解恒等式 (30.) から
\[
 \varphi'=\sum_{\beta=\sigma-\tau+\lambda}^{n-\tau,\sigma}\sum_{i_\beta}\eps\,
 \pair{\dot{q}_{\sigma-\beta}^{(i_1\ldots i_\beta)}q_{n-\tau+\lambda}}{p_{\sigma+\lambda}\dot{p}_{n-\tau-\beta}^{(i_1\ldots i_\beta)}}
\]
が従う。ここで
\[
 \dot{p}_{n-\tau-\beta}^{(i_1\ldots i_\beta)}\sim s^{(1)}\cdots s^{(\tau)}s^{(i_1)}\cdots s^{(i_\beta)},
\]
であり、$i_1,\ldots i_\beta$ は 1 から $\sigma$ までの数から選んだ任意の相異なる $\beta$ 個の数を表す。$\beta>\sigma-\tau$ であるから、これらの数 $i_1,\ldots i_\beta$ の中には 1 から $\tau$ の間のものが必ずある。したがって $\dot{p}_{n-\tau-\beta}^{(i_1\ldots i_\beta)}$ は恒等的に消え、ゆえに各行列積も、したがって $\varphi'$ も消える。

正規形の系が形式列 $f$ と同値であることを示すには、形式列のすべての形式が、1. で定めた意味で、正規形の極に展開できることを示せばよい。$\Pi$ を極化操作
\begin{equation}
 \left(\frac{\partial}{\partial\xi^{(i)}}\middle|\xi^{(k)}\right),
 \qquad (i\ne k;\ i=1,2,\ldots,\rho;\ k=1,2,\ldots,\rho)
\tag{53}
\end{equation}
の集合体とし、$\Pi^\sigma$ を、$i=\sigma$ かつ $k<\sigma$ である特殊な操作 (53.) の集合体とする。列 $\xi^{(1)},\ldots,\xi^{(\rho)}$ の形式 $F$ が最後の列 $\xi^{(\rho)}$ に関して次数 $k_\rho$ をもつとき、展開\footnote{F. Mertens, \emph{"Uber eine Formel der Determinantentheorie}. Sitzungsberichte der Akademie der Wissenschaften, Wien, Math.-Naturw. Kl., vol. 91, 2. Abteilung (1885), formula 20).}
\begin{equation}
 F=\sum_{\alpha=0}^{k_\rho}\Pi F_\alpha,
 \qquad(\rho\le n)
\tag{54}
\end{equation}
が成り立つ。ここで $F_\alpha$ は最後の列 $\xi^{(\rho)}$ において次数 $\alpha$ をもち、$F$ から不変な微分操作
\begin{equation}
 \left(\frac{\partial}{\partial\xi^{(1)}}\cdots\frac{\partial}{\partial\xi^{(\rho)}}\middle|\xi^{(1)}\cdots\xi^{(\rho)}\right)^\alpha
\tag{55}
\end{equation}
および $\Pi^\rho$ によって生じる。$F_\alpha$ を構成する際に、過程 (55.) を
\begin{equation}
 \left(\frac{\partial}{\partial\xi^{(1)}}\cdots\frac{\partial}{\partial\xi^{(\rho)}}\middle|p_\rho\right)^\alpha
\tag{56}
\end{equation}
で置き換えると、$\rho-1$ 個の列 $\xi^{(1)},\ldots,\xi^{(\rho-1)}$ だけをもつ形式 $F_\alpha(p_\rho)$ が得られ、これに再び (54.) を適用できる。この手続きを続けると、形式 $G(p_\rho,p_{\rho-1},\ldots,p_1)$ に至る。これらから (54.) に従う $F$ の展開を得るには、個々の列 $p_\sigma$ を $\xi^{(1)},\ldots,\xi^{(\sigma)}$ で置き換え、得られた形式に極化過程 $\Pi$ (53.) を適用する。これにより（(48.) 参照）変換された係数 $(G)$ の和が得られる。$\rho=n$ の場合には、過程 (55.) は第一段階で残る。$c$ を数値定数とし、簡単のため
\begin{equation}
 (\xi^{(1)}\xi^{(2)}\cdots\xi^{(n)})=\Delta;
 \qquad
 \left(\frac{\partial}{\partial\xi^{(1)}}\frac{\partial}{\partial\xi^{(2)}}\cdots\frac{\partial}{\partial\xi^{(n)}}\right)=\nabla
\tag{57}
\end{equation}
と置くと、$\rho=n$ の場合には
\begin{equation}
 F=\sum c\Delta^\alpha(G),
\tag{58}
\end{equation}
となる。一方、$\rho<n$ では右辺には $\alpha=0$、したがって $\sum c(G)$ だけが現れる。

(58.) を変換された係数 $(f)$ (48.) に適用し、極化過程 $\Pi^\sigma$ が $\rho\ge\sigma$ であるすべての操作 (56.) と可換であること、また変換された係数の極は再び変換された係数を与えることを用いると、
\[
 G=\left(\frac{\partial}{\partial\xi^{(1)}}\middle|p_1\right)^{\beta_1}
 \left(\frac{\partial}{\partial\xi^{(1)}}\frac{\partial}{\partial\xi^{(2)}}\middle|p_2\right)^{\beta_2}
 \cdots
 \left(\frac{\partial}{\partial\xi^{(1)}}\cdots\frac{\partial}{\partial\xi^{(n-1)}}\middle|p_{n-1}\right)^{\beta_{n-1}}
 \nabla^{\beta_n}\Pi(f)=\sum c\varphi
\]
を得る。さらに (58.) から
\begin{equation}
 (f)=\sum c\Delta^\alpha(\varphi)\qquad\text{（Mertens, 公式 23）}
\tag{59}
\end{equation}
が従う。(59.) から $f$ 自身の展開へ移るため、再び変数 $p_\rho$ を、$f$ に対応する次数 $m_1,\ldots m_{n-1}$ をもつ形式 $H$ (52.) の係数と見なし、$m=m_1+\cdots+m_{n-1}$ と置く。すると不変式の定義により
\[
 f\cdot\Delta^m=\sum(f)(H)=\sum c\Delta^\alpha(\varphi)(H)
\]
であり、列 $\xi$ をちょうど尽くす過程 $\nabla^m$ を適用すると
\begin{equation}
 f=\sum c\Phi,
\tag{60}
\end{equation}
を得る。ここで形式 $\Phi$ は、変数に関する不変な微分過程によって $\varphi$ と $H$ から導かれ、したがって形式列 $\varphi$ と $H$ の同時不変式である。\footnote{直接の結論の代わりに、Mertens, \S\ 7 は、形式列 $H$ を形成するために本質的に用いられる一連のさらなる微分過程を導入する。これは我々の定理 VI によって行われる。} しかし $\varphi$ は正規形であるから、形式列 $\varphi$ は $\varphi$ 自身に還元される。形式列 $H$ は可能なすべての不変な変数結合を含む。したがって $\varphi$ と $H$ の同時不変式は、1. で示した意味での $\varphi$ の極である。言い換えれば、(60.) は $f$ の正規形の極による展開を与える。任意の形式列 $f$ の形式 $\theta$ についても同じ展開を確認する必要がある。$(\theta)$ から (50.) によって導かれる正規形を $\Gamma$ とし、
\begin{equation}
 (\theta)=\sum c\Delta^\beta(\Gamma)
\tag{61}
\end{equation}
とする。形式 $\theta$ は、各変換係数について有効な関係
\[
 (\theta)\Delta^\sigma=\sum c(f)\qquad\text{（Mertens, 公式 9）}
\]
によって特徴づけられる。しかし、因子 $\Delta^\sigma$ を含む $n$ 個の列 $\xi$ の任意の形式 $F$ は、第二の形式 $\nabla^\sigma\Pi F$ も含む（Mertens, 公式 20）。したがって
\[
 (\theta)=\sum c\nabla^\sigma(f),\qquad\text{したがって}\qquad \Gamma=\sum c\varphi,
\]
となる。(61.) から (59.) に対応する関係
\[
 (\theta)=\sum c\Delta^\beta(\varphi)
\]
を得る。これは $\theta$ について、正規形 $\varphi$ の極による展開 (60.) を含意する。したがって、

\emph{定理 VII.} 各形式列 $f$ は、同値な正規形の系で置き換えることができる。

\section*{\S\ 8. 折り畳み過程の基本的折り畳みへの還元.}
定理 VII に続いて、\S\ 7 の冒頭で述べた、折り畳み過程の基本的折り畳みへの還元を実行する。我々は（\S\ 5 参照）、二つの共変的列 $S^\rho,T^\rho$ の折り畳みをすべて「共変的折り畳み」と呼び、特に二つの列 $S^\rho,T^\rho$ の欠損数 1 の共変的折り畳みを「型 $\rho$ の基本的折り畳み」と呼ぶ。すると、定理 VI を補完して次を証明する。

\emph{定理 VIII:} 一般の折り畳み過程 (41.) は、型 $\rho$ の $(n-1)$ 個の基本的折り畳みから合成できる。ただし $\rho=1,2,\ldots n-1$ である。言い換えれば、すべての不変形成を生成するには、$n-1$ 個の基本的折り畳みを逐次適用すれば足りる。

二元領域では、折り畳み過程 (41.) は唯一可能な基本的折り畳みと直接一致する。三元領域では、私は学位論文の \S\ 1 で定理 VIII を証明した（序論参照）。そこでは (44.) のため、一般の共変的折り畳みはまだ基本的折り畳みと同一である。$n$ 変数についての定理 VIII が、単に共変的折り畳みへのはるかに弱い還元ではなく、基本的折り畳みへの還元を与えることは注目に値する。証明は原理的には三元の場合の証明に倣う。すなわち、最も一般の折り畳みを、記号的恒等式を用いて基本的折り畳みから構成する。生成するのは次の三つである。
\begin{enumerate}
\item 欠損数 1 の任意の折り畳みを、欠損数 1 の共変的折り畳み、すなわち基本的折り畳みから生成する（三元の場合の類似）。
\item 欠損数 $\lambda$ の任意の折り畳みを、欠損数 $\lambda$ の共変的折り畳みと欠損数 1 の任意の折り畳みから生成する。
\item 欠損数 $\lambda$ の共変的折り畳みを、欠損数 $\lambda-1$ の共変的折り畳みと欠損数 1 の任意の折り畳みから生成する。
\end{enumerate}
本質的なこととして、定理 VII により、折り畳まれる二つの形式 $S$ と $T$ は正規形であると仮定する。したがって (46.) により、
\begin{equation}
\begin{aligned}
\pair{q_{n-\sigma_1-\lambda}S^{\sigma_1}}{S_{n-\sigma_2}p_{\sigma_2-\lambda}}&=0,
&& (\sigma_1\ge\sigma_2,\ \lambda=1,2,\ldots,\sigma_2,n-\sigma_1),\\
\pair{q_{n-\tau_1-\lambda}T^{\tau_1}}{T_{n-\tau_2}p_{\tau_2-\lambda}}&=0,
&& (\tau_1\ge\tau_2,\ \lambda=1,2,\ldots,\tau_2,n-\tau_1)
\end{aligned}
\tag{62}
\end{equation}
が成り立つ。さらに \S\ 2（末尾）により、形式 $S$ と $T$ がすべての変数列に関して線形であると仮定すれば十分である。すなわち、構成される形式は形式列
\[
 f=\pair{S^1}{p_1}\pair{S^2}{p_2}\cdots\pair{S^{n-1}}{p_{n-1}}
   \pair{T^1}{p_1}\pair{T^2}{p_2}\cdots\pair{T^{n-1}}{p_{n-1}}
\]
に属する。
1) $\sigma>\tau$ である $\pair{q_{n-\sigma-1}S^\sigma}{T_{n-\tau}p_{\tau-1}}$ を、型 $\tau+\alpha$, $\alpha=0,1,\ldots\sigma-\tau$ の基本的折り畳みから生成する。型 $\tau$ の基本的折り畳み
\[
 \pair{q_{n-\tau-1}S^\tau}{T_{n-\tau}p_{\tau-1}}
\]
から、置換 $w\sim T_{n-\tau}p_{\tau-1}$ により、上の重み添字（\S\ 1 参照）$\tau+1$ をもつ三つの列をもつ形式列 $f$ の形式
\begin{equation}
 \pair{wS^\tau}{p_{\tau+1}}\pair{S^{\tau+1}}{p_{\tau+1}}\pair{T^{\tau+1}}{p_{\tau+1}}
\tag{63}
\end{equation}
を得る。(63.) に型 $\tau+1$ の基本的折り畳みを適用すると
\[
 \pair{q_{n-\tau-2}wS^\tau}{S_{n-\tau-1}p_\tau};
\]
が得られる。(31.) から、置換 $q_{n-\tau-2}w=q'_{n-\tau-1}$ により、これは値
\[
 \eps\bigl(q_{n-\tau-2}wS^{\tau+1}\bigr)\pair{S^\tau}{p_\tau}
 +\sum_{\alpha=1}^{\tau,n-\tau-1}\Pi\pair{q'_{n-\tau-1-\alpha}S^{\tau+1}}{S_{n-\tau}p_{\tau-\alpha}}
\]
をもつ。和記号下の式は (62.) によって恒等的に消える。したがって我々は
\[
 \pair{wS^{\tau+1}}{p_{\tau+2}}\pair{S^{\tau+2}}{p_{\tau+2}}\pair{T^{\tau+2}}{p_{\tau+2}}
\]
という形式に到達した。これは (63.) で $\tau$ を $\tau+1$ に変えたものである。この過程を $(\sigma-\tau)$ 回繰り返すと、求める折り畳み
\[
 \pair{wS^\sigma}{p_{\sigma+1}}=\eps\pair{q_{n-\sigma-1}S^\sigma}{T_{n-\tau}p_{\tau-1}}
\]
に到達する。今行った過程を公式
\begin{equation}
 T^\tau S^\tau,
\quad S^{\tau+1},\quad S^{\tau+2},\ldots,S^\sigma;
\tag{64}
\end{equation}
で示す。この過程では $S$ の正規形性だけが用いられ、$T$ のそれは用いられないことがわかる。(64.) に双対な過程も、$T$ の正規形性だけを用いて、逆順に行える。すなわち
\begin{equation}
 S^\sigma T^\sigma,
\quad T^{\sigma-1},\quad T^{\sigma-2},\ldots,T^\tau
\tag{65}
\end{equation}
であり、関係式
\begin{align*}
\pair{q_{n-\sigma-1}S^\sigma}{T_{n-\sigma}p_{\sigma-1}}&=\eps\pair{T_{n-\sigma}y}{p_{\sigma-1}},\quad y\sim q_{n-\sigma-1}S^\sigma,\\
\pair{q_{n-\sigma}T^{\sigma-1}}{T_{n-\sigma}yp_{\sigma-2}}&=\eps\pair{T_{n-\sigma+1}y}{p_{\sigma-2}},\\
&\vdots\\
\pair{q_{n-\tau-1}T^\tau}{T_{n-\tau-1}yp_{\tau-1}}&=\eps\pair{q_{n-\sigma-1}S^\sigma}{T_{n-\tau}p_{\tau-1}}
\end{align*}
に従う。\S\ 6（末尾）で与えた総重量の減少との関係も指摘しておく。欠損数 1 の折り畳みでは、この減少は $2(1+(\sigma-\tau))$ であり、基本的折り畳みでは $2\cdot1$ である。したがって、基本的折り畳みからの合成がそもそも可能であるならば、ちょうど $\sigma-\tau+1$ 個のそのような折り畳みが必要であり、上で実際にそれを行った。

2) 一般の折り畳みを共変的折り畳みと基本的折り畳みから生成する。一般の折り畳み (41) の重量減少は $2(\lambda^2+\lambda(\sigma-\tau))=2[\lambda^2+(\sigma-\tau)(1+(\lambda-1))]$ である。これは、それを欠損数 $\lambda$ の共変的折り畳み（重量減少 $2\lambda^2$）と、重み添字の差が $\lambda-1$ である $\sigma-\tau$ 個の欠損数 1 の折り畳み（重量減少 $2\{1+(\lambda-1)\}$）に分解できる可能性を示す。$\lambda=1$ の場合、この分解は 1) で与えたものと一致する。これが実際に実現できることを示す。

欠損数 $\lambda$ の共変的折り畳み
\[
 \pair{q_{n-\tau-\lambda}S^\tau}{T_{n-\tau}p_{\tau-\lambda}}
\]
から、置換 $R^\lambda\sim T_{n-\tau}p_{\tau-\lambda}$ により、上の重み添字 $\tau+\lambda$ と $\tau+1$ の二つの列をもつ形式
\begin{equation}
 \pair{R^\lambda S^\tau}{p_{\tau+\lambda}}\pair{S^{\tau+1}}{p_{\tau+1}}
\tag{66}
\end{equation}
を得る。より低い重み添字 $\tau+1$ をもつ列は正規形に属する。したがって (66.) は、(65.) の順序で $\lambda$ 個の基本的折り畳みから合成された欠損数 1 の折り畳みを許す。
\[
 (R^\lambda S^\tau)S^{\tau+\lambda},\quad S^{\tau+\lambda-1},\quad S^{\tau+\lambda-2},\ldots S^{\tau+1}.
\]
(31.) と (62.) を考慮すると、これは
\[
 \pair{q_{n-\tau-\lambda-1}R^\lambda S^\tau}{S_{n-\tau-1}p_\tau}=\eps\pair{R^\lambda S^{\tau+1}}{p_{\tau+\lambda+1}}
\]
を与える。すなわち、1) の (63.) と同様に、(66.) で $\tau$ を $\tau+1$ に置き換えた形式である。この過程を $(\sigma-\tau)$ 回繰り返すと、求める折り畳み
\[
 \pair{R^\lambda S^\sigma}{p_{\sigma+\lambda}}=\eps\pair{q_{n-\sigma-\lambda}S^\sigma}{T_{n-\tau}p_{\tau-\lambda}}
\]
に至る。全過程で用いたのは $S$ の正規形性だけである。$T$ を正規形と見ると、順序を完全に逆にした双対過程が、次の公式に従って得られる。
\begin{align*}
\pair{q_{n-\sigma-\lambda}S^\sigma}{T_{n-\sigma}p_{\sigma-\lambda}}&=\eps\pair{T_{n-\sigma}R_\lambda}{p_{\sigma-\lambda}},\quad R_\lambda\sim q_{n-\sigma-\lambda}S^\sigma,\\
\pair{q_{n-\sigma}T^{\sigma-1}}{T_{n-\sigma}R_\lambda p_{\sigma-\lambda-1}}&=\eps\pair{T_{n-\sigma+1}R_\lambda}{p_{\sigma-\lambda-1}},\\
&\vdots\\
\pair{q_{n-\tau-1}T^\tau}{T_{n-\tau-1}R_\lambda p_{\tau-\lambda}}&=\eps\pair{q_{n-\sigma-\lambda}S^\sigma}{T_{n-\tau}p_{\tau-\lambda}}.
\end{align*}
3) 欠損数 $\lambda$ の共変的折り畳みを、欠損数 $\lambda-1$ の共変的折り畳みと基本的折り畳みから生成する。欠損数 $\lambda$ の共変的折り畳みの重量減少は $2\lambda^2=2((\lambda-1)^2+2\lambda-1)$ である。したがってそれは、欠損数 $\lambda-1$ の共変的折り畳みと $2\lambda-1$ 個の基本的折り畳みに分解できる可能性を示す。これはまず $n=2\lambda$ と置くと最も容易に見いだされる。欠損数 $\lambda$ の可能な共変的折り畳みは $\pair{S^\lambda}{T_\lambda}=\pair{S^\lambda}{T_{n/2}}$ だけである。これは同じ列の欠損数 $\lambda-1$ の折り畳み $\pair{uS^\lambda}{T_\lambda x}$ よりも、列 $u$ と列 $x$ を一つずつ少なく含む。逆に、欠損数 $\lambda$ の折り畳みを、列 $u$ と列 $x$ の消滅を実現し（重量減少 $2(n-1)=2(2\lambda-1)$）、同時に新しい記号列 $R^\lambda$ を生成する基本的折り畳みと、欠損数 $\lambda-1$ の共変的折り畳みとから合成する。これを行う最も単純な折り畳みは、次の欠損数 1 の折り畳みである。
\[
 \pair{sT^\lambda}{\sigma p_\lambda}=\eps\pair{S^{2\lambda-1}}{R_{\lambda-1}p_\lambda}=\eps\pair{R^\lambda}{p_\lambda},
\]
ここで
\[
 R^{\lambda+1}=sT^\lambda,
 \quad s=S^1,
 \quad \sigma=S_1,
 \quad R^\lambda\sim\sigma R_{\lambda-1}
\]
である。(65.) と (64.) により、これは $2\lambda-1$ 個の基本的折り畳み
\begin{equation}
 T^\lambda S^\lambda,S^{\lambda-1},\ldots,S^1;
 \qquad R^{\lambda+1}S^{\lambda+1},S^{\lambda+2},\ldots,S^{2\lambda-1}
\tag{67}
\end{equation}
から合成される。(21.) と (62.) を考慮した欠損数 $\lambda-1$ の折り畳みは
\[
 \pair{uS^\lambda}{\sigma R_{\lambda-1}x}=\eps\pair{R^{\lambda+1}}{S_\lambda x}
 =\eps\pair{sT^\lambda}{S^\lambda x}=\eps\pair{S^\lambda}{T_\lambda}
\]
を与える。二つの列 $S^\rho,T^\rho$ と任意の $n$ についての一般の場合は、この特殊な場合から直接得られる。(67.) の代わりに、次の $2\lambda-1$ 個の基本的折り畳みをもつ。
\begin{equation}
 T^\rho S^\rho,S^{\rho-1},\ldots,S^{\rho-\lambda+1};
 \qquad R^{\rho+1}S^{\rho+1},S^{\rho+2},\ldots,S^{\rho+\lambda-1}.
\tag{68}
\end{equation}
折り畳みの前半は形式
\[
 \pair{q_{n-\rho-1}T^\rho}{S_{n-\rho+\lambda-1}p_{\rho-\lambda}}
\]
を与える。ここから置換
\[
 S_{n-\rho+\lambda-1}p_{\rho-\lambda}\sim w,
 \qquad T^\rho w=R^{\rho+1}
\]
と後半の折り畳みの適用により、
\[
 \pair{q_{n-\rho-\lambda}S^{\rho+\lambda-1}}{R_{n-\rho-1}p_\rho}
\]
を得る。さらに
\[
 q_{n-\rho-\lambda}S^{\rho+\lambda-1}\sim y
\]
と置換すると、欠損数 $\lambda-1$ の共変的折り畳みにより
\begin{equation}
 \pair{q_{n-\rho-\lambda+1}S^\rho}{yR_{n-\rho-1}p_{\rho-\lambda+1}}
\tag{69}
\end{equation}
を得る。これが求める欠損数 $\lambda$ の折り畳みであることを示す。
\[
 R_{n-\rho-1}p_{\rho-\lambda+1}=R_{n-\lambda}
\]
と置き、$y$ を分解すると
\[
 \pair{q_{n-\rho-\lambda+1}R^\lambda}{S_{n-\rho}y}
 =\eps\pair{q_{n-\rho-\lambda}S^{\rho+\lambda-1}}{S_{n-\rho}p'_{\rho-1}}=0
\]
となることを観察する。したがって (20.) により、(69.) の値は
\[
 \eps\pair{S^\rho R^\lambda}{p_{\rho+\lambda-1}y}
 =\eps\pair{q_{n-\rho-\lambda}S^{\rho+\lambda-1}}{[S^\rho R^\lambda]_{n-\rho-\lambda}p_{\rho+\lambda-1}}
\]
である。この折り畳みは型 (43.) であり、したがって
\[
 \pair{q_{n-\rho-\lambda}S^\rho}{R_{n-\rho-1}p_{\rho-\lambda+1}}
 =\eps\pair{wT^\rho}{R_\lambda p_{\rho-\lambda+1}},
 \qquad R_\lambda\sim q_{n-\rho-\lambda}S^\rho
\]
に同値である。ここで $R_\lambda$ はすでに望む最終形になっている。この式は特殊な場合の形式 $\pair{sT^\lambda}{S_\lambda x}$ に対応する。今、$w$ と $R_\lambda$ を分解すると
\[
 \pair{wR^{n-\lambda}}{T_{n-\rho}p_{\rho-\lambda+1}}
 =\eps\pair{q'_{\lambda-1}R^{n-\lambda}}{S_{n-\rho+\lambda-1}p_{\rho-\lambda}}
 =\eps\pair{q''_{n-\rho-1}S^\rho}{S_{n-\rho+\lambda-1}p_{\rho-\lambda}}=0
\]
となることから、(21.) により値
\[
\begin{aligned}
 \pair{wq_{n-\rho+\lambda-1}}{T_{n-\rho}R_\lambda}
 &=\eps\pair{q_{n-\rho+\lambda-1}[T_{n-\rho}R_\lambda]^{\rho-\lambda}}{S_{n-\rho+\lambda-1}p_{\rho-\lambda}}\\
 &\quad\text{これは}\quad
 \pair{q_{n-\rho-\lambda}S^\rho}{T_{n-\rho}p_{\rho-\lambda}}
 \text{に同値である}
\end{aligned}
\]
を得る。全過程で用いたのは $S$ の正規形性だけであり、$T$ のそれではない。したがって、(69.) を生成するために用いた欠損数 $\lambda-1$ の共変的折り畳みは、$2\lambda-3$ 個の基本的折り畳みと欠損数 $\lambda-2$ の共変的折り畳みに置き換えることができる。後者もさらに基本的折り畳みに分解でき、以下同様である。同様に、$T$ の正規形性を用い、双対的過程を適用することによっても生成が得られる。すなわち基本的折り畳み
\[
 S^\rho T^\rho,\quad T^{\rho+1},\ldots T^{\rho+\lambda-1};
 \qquad R^{\rho-1}T^{\rho-1},\quad T^{\rho-2},\ldots T^{\rho-\lambda+1},
\]
と、(69.) に双対な折り畳みである。そのとき前のものに同値な折り畳み
\[
\pair{q_{n-\rho-\lambda}T^\rho}{S_{n-\rho}p_{\rho-\lambda}}
\]
が得られる。

2) で得た分解を加えると、最も一般の折り畳みが確かに基本的折り畳みから生成されることがわかった。残るのは、この生成が順序の置換を除いて一意であること、すなわち各折り畳みに対して、型 $\rho$ の基本的折り畳みの数 $\alpha_\rho$ が一つだけ対応することを示すことである。ここで $\rho=1,2,\ldots n-1$ である。

初期形式 $f$ の変数列 $p_\rho$ における次数を $m_\rho$ とし、折り畳みによって作られた形式の次数を $l_\rho$ とする。型 $\rho$ の各基本的折り畳みは、次数 $m_{\rho-1},m_\rho,m_{\rho+1}$ をそれぞれ $m_{\rho-1}+1,m_\rho-2,m_{\rho+1}+1$ に変える。したがって $\alpha$ は方程式系
\begin{equation}
 m_\rho-l_\rho=-\alpha_{\rho-1}+2\alpha_\rho-\alpha_{\rho+1},
 \qquad (\rho=1,2,\ldots,n-1)
\tag{70}
\end{equation}
を満たす。ここで $\alpha_0=\alpha_n=0$ と置く。この方程式系の行列式の値は $n$ である。\footnote{$\Delta_n$ が $n$ 行の行列式を表すなら、漸化式 $\Delta_n=2\Delta_{n-1}-\Delta_{n-2}$ が成り立つ。すなわち $\Delta_n-\Delta_{n-1}=\Delta_{n-1}-\Delta_{n-2}=\cdots=\Delta_2-\Delta_1=1$ である。なぜなら $\Delta_2=3$, $\Delta_1=2$ だからである。} したがって $\alpha$ は一意に定まり、しかも定理 VIII により正の整数として定まる。これにより差 $m_\rho-l_\rho$ についての条件が得られる。

この節の結果は関係 (62.) に基づく。しかし、もし関係 (62.) が成り立たないなら、最終式は初期式と決して同値ではない。\S\ 6 で与えた同値性の定義は、次の定式化も許す。\footnote{次節末の考察を参照。}「二つの形式が同値であるとは、それらの差がより高い折り畳みの形式、すなわちより大きい重量減少を示す形式である場合、かつその場合に限る。」このより大きい重量減少は、折り畳みに入る二つの列 $q_{n-\sigma-\lambda},p_{\tau-\lambda}$ が実際に変数列であるときには常に起こる。これは (41.), (42.)、この節で同値として現れた折り畳み、また定理 VII の同値な形式において起こる。しかし、これらの列が置換によって記号列から導かれているときには、必ずしも起こるとは限らない。この節で、(62.) によって互いに表すことができた形式がその場合である。消える形式でさえ、部分的にはより高い総重量をもつことが容易にわかる。

\section*{\S\ 9. 形式列. 還元定理.\footnote{三元領域における対応する考察を参照。学位論文 §§1--3。}}
ここで、\S\ 6（末尾）で導入した形式列を、やや厳密に定義する。「与えられた初期形式の係数に線形な不変形成の全体、すなわち、その初期形式をそれ自身で折り畳むことによって生じる形式の全体を、より高い形式に従って配列したもの」である。

より高い形式を説明するため、定理 VII により初期形式が二つの正規形 $S\cdot T$ の積として与えられているとする。より高い形式とは、それ自身においてより高い折り畳みをもつ形式、そして同じ折り畳みをもつ形式の中では特別に正規化された形式をいう。より正確には、形式 $A$ が型 $\rho$ の基本的折り畳みを $\alpha_\rho$ 回行って生じ、形式 $B$ がそのような折り畳みを $\beta_\rho$ 回行って生じたとする。このとき $B$ が $A$ より高いとは、次の場合である。
\begin{enumerate}
\item 不等式系 $\beta_\rho\geqq\alpha_\rho$, $\rho=1,2,\ldots n-1$ において、少なくとも一つの $\rho$ について不等号が真に成り立つ場合。\footnote{不等式系の中で $>$ の符号と $<$ の符号が部分的に混在するなら、形式は比較できない。なぜなら、ある形式から折り畳みによって生じる形式は常に基本的折り畳みの正の値系をもち、したがってこの場合には正または零の値 $\beta_\rho-\alpha_\rho$ をもつからである。}
\item すべての $\rho$ について等号が成り立つが、より少ない記号列が一つの行列積に結合されている場合。この正規化は \S\ 8 の還元のために必要である。したがって、たとえば $\pair{S^{n-1}}{T_{n-1}}$ は $\pair{S^{n-1}}{[S^1T^\rho]_{n-\rho-1}p_\rho}$ より高い。
\item すべての $\rho$ について等号が成り立ち、一つの行列積に結合される記号列の数も同じである場合には、任意ではあるが一度限り固定された正規化によって $B$ を $A$ より高くする。この正規化は、値系 $\alpha$ に属する形式が有限個であるため常に可能である。たとえば一つの形式 $S$ を優先する（記号列 $S$ の数が多い、列 $S$ の上の重み添字の和が大きい、など）ことによって実現できる。
\end{enumerate}
形式列は、$(n-1)$ 次元の格子点の多様体として配列されていると考えると便利である。その $n-1$ 個の方向は $n-1$ 個の基本的折り畳みに対応する。各形式には一つ、かつただ一つの値系 $\alpha$、すなわち正の整数格子点が対応し、同じ値系 $\alpha$ に属する異なる形式は、上の正規化によってその順序が一意に定まる。形式列の任意の形式は、新しい形式列の初期形式を与え、その形式列は全形式列の一部として含まれる。実際、それは新しい初期形式から $n-1$ 個の正方向へ進むことによって得られる形式からなる部分として含まれる。

定理 VIII と級数展開の一般理論により、形式列については、二元および三元領域の還元定理（学位論文、\S\ 3 参照）がそのまま成り立つ。主定理を簡単に導く。まず定義を述べる。

「与えられた形式系において、ある有限個の形式が一つの加群にまとめられているとする。ある形式が、因子として不変式をもつ形式によって、または『より高い形式』によって表せるなら、その形式を可約と呼ぶ。すなわち、その形式が属する全形式列のより高い形式によって、またはその加群の記号をより高い次数で含む形式によって表せる場合である。可約な形式列を還元子と呼ぶ。」すると、

\emph{定理 IX:} ある形式列の初期形式が、還元子との折り畳みによって生じたために可約であるなら、この初期形式から生じる形式列全体は可約である。

証明のため、形式 $f$ と第二の形式 $g$ を折り畳んで生じる形式 $h$ は、(45.) により、いくつかの共変的な変数列 $p_\rho,p'_\rho$ を含む形式 $f$ と見なせることを観察する。このような形式は、級数展開の一般理論により、各対の共変的変数列 $p_\rho,p'_\rho$ に順に級数展開を適用することで、$f$ の基本共変式の極として表される。(38.) によれば、現れる基本共変式は共変的折り畳みによって生成される形式である。しかし、級数展開を適用する順序はなお任意であるから、とくに一般の折り畳みを基本的折り畳みから生成する順序に対応するものを選べる。得られる基本共変式の系は、形式列 $f$ と同一になる。欠損数の高い共変的折り畳みによって生じる形式は、基本的折り畳みによって生じる形式の中に配列される。しかし、級数展開の任意の順序に対しても形式列 $f$ の極が得られ、それに第二の基本共変式の系が対応することがある。形式列は不変形成の全体を尽くすので、第二の系の各形式は形式列の極によって線形に表され、しかも同じか、またはより大きい重量減少を示す形式の極によって表される。後者は、極が（\S\ 7 参照）形式 $h$ に対応する次数をもつ形式列 $H$ (52.) と形式列 $f$ との同時不変式と見なされることから従う。$H$ のそれ自身における各折り畳みは重量減少に対応し、置換 (11.) が行われた後、この減少は記号列へ、したがって $f$ の形式へ移される。\footnote{この考察は、同値性概念の第二の定式化（\S\ 8 末尾）の基礎にもなっている。三元領域における基本共変式の異なる系についてのさらなる説明は、Study, loc. cit., II, \S\ 7, 定理 7) および 8) にある。我々の定理 VII により、同じ結果は $n$ 変数についても成り立つ。}

形式列 $h$ の任意の形式は、$h$ をそれ自身で折り畳むことによって生じた。すなわち一方では $g$ と $f$ のより高い折り畳みによって、他方では $f$ または $g$ をそれ自身で折り畳むことによって生じた。どちらの場合にも、級数展開は、初期形式 $h$ について上で示したのと同じように、形式列 $f$ の極へ導く。特に $f$ が還元子であるなら、還元子の極による展開が得られる。したがって形式列 $h$ は可約となる。

完全な形式系の還元へのこの定理の応用は、二元および三元領域の場合と同様に従う。

\clearpage
%
\providecommand{\Kfield}{\mathfrak{K}}
\providecommand{\Ssys}{\mathfrak{S}}
\providecommand{\Mfield}{\mathfrak{M}}
\providecommand{\tq}{\tau}
\providecommand{\Lsys}{\mathfrak{L}}
\providecommand{\Jdom}{\mathfrak{J}}
\providecommand{\ds}{\displaystyle}

\begin{center}
{\Large\bfseries 5. 有理関数体\footnote{1913年ウィーン自然科学者集会で行われた講演。}}\par
\vspace{0.8em}
\textsc{Emmy Noether}（Erlangen）\\[0.5em]
\emph{Jahresbericht der DMV} 22 (1913), pp. 316--319
\end{center}

\vspace{1.5em}

以下の諸問題は、もともと E. Fischer との会話に由来する。なお、それらの問題のいくつかは、一つの不定元という特殊な場合について、しかも係数体に関するより一般的な仮定のもとで、すでに E. Steinitz によって提起され解決されている。\footnote{E. Steinitz: \emph{Algebraische Theorie der Körper}, とくに \S\ 24. Crelle's Journal 137, 1910.}

「有理関数体」とは、与えられた数体からの係数をもつ $n$ 個の不定元の有理関数を元とする体をいう。この係数体は、とくにすべての複素数を含んでもよい。このような体の例としては、$n$ 個の量の対称関数の体、またより一般に、ある群の置換を許す $n$ 個の不定元の有理関数全体（Lagrange の \emph{Gattungsbereiche}）、さらに不変式体がある。

最初に挙げた二つの体と最後の体との間には、特徴的な相違がある。前者は $n$ 個の代数的に独立な関数、すなわち基本対称関数を含むのに対し、代数的に独立な不変式の個数は常に不定元の個数より小さい。したがって、一般に、第二種のこのような体を、いくつかの不定元を数で置き換えることによって第一種の体と一対一に対応させられることが重要である。そこで以下では、第一種の体、すなわち $n$ 個の代数的に独立な関数をもつ体に限って考える。この対応により、結果は一般にも成り立つ。

問題は三つの基底概念、すなわち有理基底、最小基底、整基底のまわりに集まる。

1. 「有理基底」とは、その体の有限個の関数であって、体のすべての関数が、与えられた係数体からの係数を用いて、それら有限個の関数の有理的結合として表されるものをいう。簡単な考察によって、任意の有理関数体について有理基底の存在が示される。一つの不定元の場合には、この考察は E. Steinitz の前掲書にも見られる。例えば Lagrange の定理によれば、基本対称関数と群に属する一つの関数は、先に述べた Lagrange の \emph{Gattungsbereiche} の有理基底をなす。

2. 任意の有理基底は、第一種の体に対して、少なくとも $n$ 個の関数を含まなければならない。ちょうど $n$ 個の関数を含むならば、それらは代数的に独立でなければならず、そのような基底を「最小基底」と呼ぶ。最小基底の存在問題は、Lüroth, Castelnuovo, Enriques の定理によって部分的に答えられる。\footnote{J. Lüroth: Beweis eines Satzes über rationale Kurven. \emph{Math. Ann.} 9, 1875. --- G. Castelnuovo: Sulla razionalità delle involuzioni piane. \emph{Math. Ann.} 44, 1893. --- F. Enriques: Sopra una involuzione non razionale dello spazio. \emph{Rend. Acc. Linc.} vol. 21, 1912年1月21日.} それらによれば、一つの不定元の体には常に最小基底が存在する。すなわち、その体の Lüroth 関数であり、これは一次分数変換を除いて一意に定まる（Steinitz \S\ 24 参照）。二つの不定元の体については、係数体の代数拡大を許すならば、そして一般にはその場合に限って、常に最小基底が存在する。これに対して三つ以上の不定元については、一般には最小基底はまったく存在しない。最小基底をもつ特殊な体を特徴づける試みは、まだ行われていない。

私は Lagrange の \emph{Gattungsbereiche} における最小基底の問題をさらに追究した。この場合、それは次の意味をもつ。「群 $G$ に属する Lagrange の \emph{Gattungsbereich} が最小基底をもつならば、指定された群 $G$ をもつ方程式全体を、パラメータ表示によって有理的に構成できる。しかも、最小基底の係数体を含む任意の数体についてそうであり、特にこの係数体が有理数体であるならば、任意の数体についてそうである。」

最も簡単な例は対称群である。この場合、基本対称関数は有理数係数をもつ最小基底をなす。そこからパラメータ表示として、単に不定係数をもつ方程式が得られる。これから任意の数体上の任意に多くの affektlos な方程式へ移る方法は、Hilbert の既約性定理が示している。

さて、Lüroth と Castelnuovo の定理から、3 次および 4 次方程式に現れるすべての群について最小基底の存在が直接に従う。同時に、この最小基底は有理数係数をもつことも分かる。したがって、任意の数体について、指定された群をもつ 3 次および 4 次方程式全体を有理的に構成できる。巡回群および二面体群については、$n$ 乗根を係数体とする最小基底が存在する。さらにいくつかのメタ巡回群についても同様である。最小基底がそもそもある種の群を特徴づけるかどうかは、未決のままにしておかなければならない。

3. 最後の基底概念に到達するため、体の中に含まれる多項式全体を考える。「整基底」とは、これらの多項式の有限個であって、体の任意の多項式が、与えられた係数体からの係数を用いて、それら有限個の多項式の整有理結合として表されるものをいう。整基底の存在問題は、例えば不変式系の有限性を特殊な場合として含むが、Hilbert の「数学の問題」\footnote{D. Hilbert: Mathematische Probleme. パリ講演, 1900. 問題14. \emph{Göttinger Nachr.} 1900.}では、やや異なる形、すなわち相対的に整な関数の問題として提出されている。

当面、私は整基底が存在する体の一つの類だけを示すことができる。すなわち、その体が $n$ 個の多項式の系を含み、それらの最高次元項の同次終結式が零でないならば、整基底が存在する。それは、線形形式
\[
  u_0=u_1x_1+u_2x_2+\cdots+u_nx_n\footnotemark
\]
に対する、その体で既約な方程式に現れるすべての $u$ の係数によって与えられる。
\footnotetext{この既約方程式は、一つの不定元の場合には、Steinitz による Lüroth の定理の証明に見られる。}
とくに終結式の値が係数体の単元に等しいならば、表示の整性も保証される。この類の体の例は Lagrange の \emph{Gattungsbereiche} である。なぜなら基本対称関数の終結式は値 $\pm 1$ をもつからである。さらに別の例（整性なし）は、ただ一つの代数的に独立な多項式だけを含むすべての体によって与えられる。この場合、整基底は、すべての多項式を含む最小の部分体の Lüroth 関数と同一である。したがって、よく知られているように、$n$ 変数の二次形式のすべての整有理な射影不変式は、判別式の整関数である。

いま述べた条件は整基底の存在に十分であるにすぎず、決して必要ではない。任意の $n$ 個の多項式の終結式が消えるにもかかわらず、あの既約方程式の係数によって与えられる整基底をもつ体は、容易に構成できる。最も一般の場合にも、あの方程式の係数が整基底を与えるのではないか、という問題が生じる。

\clearpage

\begin{center}
{\Large\bfseries 6. 有理関数の体と系}\par
\vspace{0.6em}
Emmy Noether (Erlangen).\par
\vspace{1.2em}
\emph{Mathematische Annalen} 76 (1915), pp. 161--196
\end{center}

\vspace{1.5em}

本論文は、有理関数および整有理関数の任意の系に関する基底問題を扱う。ここで用いられる体論の方法により、有理関数からなる体、すなわち有理関数体\footnote{ウィーン自然科学者会議の機会における予備的報告「Rationale Funktionenkörper」, \emph{Jahresber. d. D. Math.-Ver.} 22 (1913) を参照。そこでは、問題設定と関数体に関する結果の概観が与えられている。}についてこれらの問題を解くことが本質的な部分として現れ、任意の系への結果の一般化はその帰結として得られる。

一般の系に対する基底問題のうち、これまで知られていたのは Hilbert の定理（\emph{Math. Ann.} 36）によって保証される加群基底の存在だけであった。以下では、任意の系に対する有理基底の存在（\S\ 7）によって、有理的表示可能性の問題が完全に解決される。体の有理基底はすでに \S\ 4 で得られる。この有理基底の存在により、常に抽象的に定義された体または系から出発することができ、したがって、特殊な有理基底の選択だけに由来し、系そのものには由来しない困難、例えば特殊な分母や基底関数の基本点の出現を避けることができる。

体については、最小基底、すなわち代数的に独立な関数からなる有理基底の問題も扱われる（\S\ 6）。さらに体論の方法は、特別な有理基底、すなわちインボリューション基底\footnote{この名称は、幾何学的解釈において、有理関数体が線形空間内の最も一般的なインボリューションを表すことによる。}（\S\ 5）へ導く。これは多項式からなる整域にとって特に重要となる（\S\ 8）。その場合、固定された分母（整域によって定まる関数の冪）をもつ表示を与える。これは、例えば不変式の典型表示という特殊な場合に知られているものである。

有理的表示可能性から、可能な限り、整有理的表示可能性、すなわち狭い意味での有限性の問題への帰結を引き出す。

これまで知られていた有限性定理はすべて、Hilbert の加群基底定理が、表示
\[
  F=A_1f_1+A_2f_2+\cdots+A_kf_k
\]
から第二の表示
\[
  F=B_1f_1+B_2f_2+\cdots+B_kf_k
\]
が従い、ここで $B_i$ がその系に属し、かつ $F$ より低い次数をもつようなすべての系に対して有限性を保証する、という事実に基づいている。

これに対して、有理基底に関する定理と体論の方法は、まったく異なる種類の仮定のもとで有限性を推論することを可能にする。

こうして Hilbert の問題（数学の問題 14）への部分的な答えとして、相対的に整な関数の一類（第一種の相対的に整な領域）が得られる。これらは有限な整域であり、その整基底は上に述べたインボリューション基底によって与えられる（\S\ 10）。さらに、Hilbert による Kronecker の代数的整数の基本系に関する定理の証明（\emph{Complete invariant systems}, \S\ 2; \emph{Math. Ann.} 42）と類似に、すべての正則多項式系、すなわち基本点をもたないものへ写すことのできるすべての系は整基底をもつことが示される（\S\ 12）。一方、整基底をもたない非正則系は容易に指定できる。この事実の簡単な例として、次の定理を挙げられる。「一つの不定元に関する任意の無限個の多項式の系において、その系の任意の多項式が有限個の多項式の整有理結合として表されるような有限個の多項式が存在する。」さらなる例は \S\ 13 にある。

最後に、末尾の二つの節（\S\S\ 14, 15）は、整性に関して基底定理を精密化する。

この研究の本質的な道具は \S\ 3 の転移原理である。それによれば、ここで提起されるすべての基底問題について、不定元の個数が系の代数的に独立な関数の個数と一致するような系だけを考えれば十分である。

本研究の発端は、E. Fischer 氏との会話、とくに Lagrange の \emph{Gattungsbereiche} の最小基底に関する同氏の問いであった。これに関わる研究は、指定された群をもつ方程式の構成へ導いたが（前掲の「有理関数体」第2部を参照）、それは後の公表に委ねる。

本論文が最初の報告から本質的に拡張されている点は、そこでは体または特殊な整域についてだけ述べられた基底定理が、ここでは任意の系へ移されていることである。この転移は、任意の系を含む最小の体という概念によって簡単に行われる。これは整域の商体の一般化である（\S\ 7）。私がこの概念を形成するに至ったのは、K. Hentzelt 氏が、体の有理基底を知った後、Hilbert の加群基底を経由する迂回によって任意の系の有理基底の存在を証明できた、という事実による。また、正則系の整基底に関する定理へ私を導いたのは、整域に対して私が用いた推論が、同じ仮定を満たす任意の系にも有効であるという K. Hentzelt の指摘であった。さらに、いくつかの個別の小さな指摘も Hentzelt 氏に負っている。

最後に、一つの不定元の場合、体の有理基底と最小基底は E. Steinitz の \emph{Algebraische Theorie der Körper}, \S\ 24（Crelle's Journal 137）に見られることを述べておく。そこでは係数体は最も一般的な意味での体として仮定されている。これらのより一般的な仮定のもとで、ここに与える定理がどこまで保たれるかという問題は、以下では扱わない。いずれにせよ証明は修正を必要とする。例えば、関数行列式の理論からの手段は標数 $p$ の体では機能しないからである。

\section*{\S\ 1. 体 $\Kfield_{n\rho}$ と系 $\Ssys_{n\rho}$.}

以下では、$n$ 個の不定元の有理関数の系、とくに有理関数体を調べる。

$f(x_1\cdots x_n)$, $g(x_1\cdots x_n)$, \ldots、または略して $f(x)$, $g(x)$, \ldots は、常に $x_1\cdots x_n$ の有理関数を意味するものとする。これらは既約表示、すなわち分子と分母が互いに素である表示にあると仮定する。整有理関数（多項式）は大文字 $F(x)$, $G(x)$, \ldots で表す。係数領域 $\Omega$ は何らかの数体と仮定され、したがってとくにすべての複素数を含んでもよい。$\Omega$ はさらに有限個のパラメータを含んでもよい。

$\Omega$ からの量、したがって零次のすべての関数は、常に系に属するものと数える。

これらの取り決めのもとで、有理関数からなる体、すなわち有理関数体は抽象的に定義できる。

\emph{定義 I: 有理関数の系は、次の条件を満たすとき体と呼ばれる。}
\begin{enumerate}
\item \emph{$f(x)$ とともに、$\Omega$ からの任意の量 $c$ に対して $c\cdot f(x)$ も含む。}
\item \emph{$f(x)$ と $g(x)$ とともに、常に $f(x)+g(x)$, $f(x)\cdot g(x)$ および、$g(x)\ne0$ のとき商 $f(x):g(x)$ も含む。}\footnote{ここで $f(x)$ と $g(x)$ は零次、すなわち $\Omega$ の量であってもよい。}
\end{enumerate}
特殊な体とは、有限個の有理関数
\[
  f_1(x),\ f_2(x),\ \cdots,\ f_k(x)
\]
を $\Omega$ に添加して得られるものであり、
\[
  \Omega\bigl(f_1(x)\cdots f_k(x)\bigr)
  \quad\text{または}\quad
  \Omega(f_1\cdots f_k)
\]
と表す。\footnote{これらの「特殊な体」がすでにすべての体を尽くすことは \S\ 4 で分かる。}
これら特殊な体の中には、$x_1\cdots x_n$ を $\Omega$ に添加して得られる体
\[
  \Omega(x_1,\ldots,x_n)
\]
も含めておく。これは $\Omega$ 係数の $x_1,\ldots,x_n$ の有理関数全体と同一である。

さらに、E. Steinitz に従って、中間体の概念を導入する。
\begin{quote}
「体 $\Mfield$ が $\mathfrak{L}_1$ と $\mathfrak{L}_2$ の間にあるとは、$\Mfield$ が $\mathfrak{L}_1$ を含み、かつ $\mathfrak{L}_2$ に含まれることをいう。」
\end{quote}
すると定義 I は次のようにも述べられる。

$n$ 個の不定元の有理関数体とは、$\Omega$ と $\Omega(x_1\cdots x_n)$ の間の中間体であり、少なくとも一つの関数において $n$ 個の不定元を実際に含むものである。

不定元の個数のほかに、有理関数の系にはさらに一つの特徴数、すなわち代数的に独立な関数の個数、あるいは代数的階数が、次の定義によって現れる。

\emph{定義 II: 有理関数の系が代数的階数 $\rho$ をもつとは、その系の中に代数的に独立な $\rho$ 個の関数を指定できるが、その系の任意の $(\rho+1)$ 個の関数は代数的に従属であることをいう。}\footnote{$\rho$ 個の関数 $f_1(x),\ldots,f_\rho(x)$ が代数的に独立であるとは、任意の関係
\[
F\bigl(f_1(x),\ldots,f_\rho(x)\bigr)=0\quad\text{$x_1,\ldots,x_n$ について恒等的に}
\]
から
\[
F(\lambda_1,\ldots,\lambda_\rho)=0\quad\text{$\lambda_1,\ldots,\lambda_\rho$ について恒等的に}
\]
が従うことをいう。反対の場合、それらは代数的に従属であると呼ばれる。}

$n$ 個の不定元の有理関数の（有限または無限の）系で、代数的階数 $\rho$ をもつものを、二重添字によって
\[
  \Ssys_{n\rho}
\]
と表す。特に
\[
  \Kfield_{n\rho}
\]
は、$n$ 個の不定元と代数的階数 $\rho$ をもつ体を意味する。次の不等式が成り立つ。
\[
  1\leq \rho\leq n.
\]

体 $\Kfield_{nn}$ および $\Kfield_{n\rho}$ の例をさらにいくつか挙げる。

1. 体 $\Kfield_{nn}$ として、Lagrange の \emph{Gattungsbereiche} がある。これは
\begin{quote}
「$x_1\cdots x_n$ の置換群 $G$ の置換を許す、$x_1\cdots x_n$ の有理関数（および整有理関数）全体」
\end{quote}
として定義できる。これらは常に対称関数の体を含み、したがって $n$ 個の代数的に独立な基本対称関数も含む。

2. 体 $\Kfield_{n\rho}$ として、射影不変式体がある。これは一つまたは複数の基本形式の有理不変式（および整有理不変式）全体から作られる。\footnote{行列式 $1$ の変換だけを考えるのが便利である。そのとき任意個の不変式の任意の有理結合も再び変換を許し、実際に体が得られる。斉次・等重の不変式だけを取り出したものは体をなさないが、系 $\Ssys_{n\rho}$ をなす。} ここでは $\rho<n$ である。不変式の代数的に独立な個数は、常に係数の個数より小さいからである。

射影群の部分群に関する不変式体についても、対応する主張が成り立つ。

\section*{\S\ 2. 代数的に従属な関数に関する補助定理.}

この節の補助定理により、\S\ 3 で、系の代数的階数がその本来的な特徴数であることが示される。すなわち補助定理は、$\rho$ 個の代数的に独立な関数が代数的に独立なままでいるように、いくつかの不定元を数値によって逐次特殊化しても、それら $\rho$ 個の関数に代数的に従属な任意の関数は、恒等的に消えることも不定になることもない、ということを示す。

したがって、系
\begin{equation}
  g_1(x),\ g_2(x),\ldots,g_\rho(x)
\tag{1}
\end{equation}
が代数的階数 $\rho$ をもち、ここで $\rho<n$ とし、$h(x)$ を関数 (1) に代数的に従属な有理関数とする。既知の定理により、(1) の関数行列の階数は $\rho$ に等しい。したがって不定元は、例えば次のように名づけられる。
\begin{equation}
 \frac{\partial(g_1,g_2,\ldots,g_\rho)}{\partial(x_1,x_2,\ldots,x_\rho)}
 =\frac{\Gamma(x)}{G(x)^{\rho+1}}\ne0,
\tag{2}
\end{equation}
ここで $G(x)$ は通常どおり、関数 (1) の共通分母を意味する。

\emph{いま数 $a$ を、積 $G(x)\cdot\Gamma(x)$ が $(x_i-a)$ で割り切れないように選ぶ}。ただし $i$ は $\rho+1,\rho+2,\ldots,n$ のいずれかとする。したがって $x_i=a$ としても、関数 (1) の共通分母は消えず、$\rho$ 個の関数
\begin{equation}
  [g_1(x)]_{x_i=a},\ldots,[g_\rho(x)]_{x_i=a}
\tag{3}
\end{equation}
もまた代数的に独立である。言い換えれば、系 (1) の代数的階数は、置換 $x_i=a$ のもとで保存される。

関数 (1) に関して $h(x)$ が満たす既約方程式を
\begin{equation}
\begin{aligned}
&A_0(g_1,\ldots,g_\rho)\cdot\,h(x)^\tau
 +A_1(g_1,\ldots,g_\rho)\cdot\,h(x)^{\tau-1}
 +\cdots \\
&\hspace{6em}+A_\tau(g_1,\ldots,g_\rho)=0
\end{aligned}
\tag{4}
\end{equation}
とする。ただし $A_0(g_1,\ldots,g_\rho)$ および $A_\tau(g_1,\ldots,g_\rho)$ は恒等的には消えない。これを多項式の間の恒等式に移すために、
\begin{equation}
  g_i(x)=G_i(x):G(x);\qquad h(x)=H_1(x):H_2(x)
\tag{5}
\end{equation}
と置く。そして、$h(x)$ が既約表示にあるという仮定により、$H_1(x)$ と $H_2(x)$ は互いに素であることに注意する。(5) によって (4) は
\begin{equation}
\begin{aligned}
&B_0(G,G_1,\ldots,G_\rho)\cdot\,G(x)^{\sigma_0\cdot}H_1(x)^\tau \\
&\quad+B_1(G,G_1,\ldots,G_\rho)\cdot\,G(x)^{\sigma_1\cdot}H_1(x)^{\tau-1}H_2(x)+\cdots \\
&\quad+B_\tau(G,G_1,\ldots,G_\rho)\cdot\,G(x)^{\sigma_\tau\cdot}H_2(x)^\tau=0
\end{aligned}
\tag{6}
\end{equation}
となる。ここで $B_k$ は $G,G_i$ の同次形式であり、非負指数 $\sigma_k$ の少なくとも一つは零でなければならない。

いま $H_1(x)$ が $(x_i-a)$ で割り切れると仮定する。すると (6) により
\[
  B_\tau(G\,G_1\cdots G_\rho)\cdot G(x)^{\sigma_\tau}\cdot H_2(x)^\tau
\]
も $(x_i-a)$ で割り切れる。これは $G(x)$ についての仮定により排除されている。したがって $B_\tau$ の割り切れ性から
\[
  A_\tau=B_\tau:G^\lambda
\]
の割り切れ性も従い、関数 (3) の間に関係 $A_\tau=0$ が存在することになり、仮定された代数的独立性に反する。最後に、$H_2(x)$ は $H_1(x)$ と互いに素であるから、これも $(x_i-a)$ で割り切れることはできない。\footnote{既約表示に基づくこの最後の結論は、残りの仮定が保たれるとしても、置換 $x_i=a$, $x_k=b$ に対してはもはや可能ではない。その場合 $H_1(x)$ と $H_2(x)$ は同時に消えることがあり、商はその置換のもとで $=\frac00$ に等しい不定形となりうる。例えば
\[
 h(x)=\frac{x_3(\alpha x_1+\beta x_2)+x_4(\gamma x_1+\delta x_2)}
 {x_3(\alpha' x_1+\beta' x_2)+x_4(\gamma' x_1+\delta' x_2)}
\]
を $x_3=0$, $x_4=0$ と置換する場合である。} よって $H_1(x)$ が $(x_i-a)$ で割り切れるという仮定は矛盾に至る。同様に、$H_2(x)$ も $(x_i-a)$ で割り切れないことが示される。したがって関数 $[h(x)]_{x_i=a}=k(x)$ は、恒等的に消えることも不定になることもない。

以上の結論を、再び既約表示に置いた系 (3) と関数 $k(x)$ に適用できる。この手続きを続けることで、最終的に次の補助定理に到達する。

\emph{補助定理: 二つの系}
\[
  g_1(x),\ g_2(x),\ldots,g_\rho(x),
\]
\[
  g_1(x),\ g_2(x),\ldots,g_\rho(x),\ h(x)
\]
\emph{はいずれも代数的階数 $\rho$ をもつものとし、ここで $\rho<n$ である。例えば}
\[
 \frac{\partial(g_1,\ldots,g_\rho)}{\partial(x_1,\ldots,x_\rho)}
 =\frac{\Gamma(x)}{G(x)^{\rho+1}}\ne0
\]
\emph{と仮定する。数}
\[
  a_n,\ a_{n-1},\ldots,a_{\rho+1}
\]
\emph{は、置換}
\[
  x_n=a_n,\quad x_{n-1}=a_{n-1},\ldots,\quad x_{\rho+1}=a_{\rho+1}
\]
\emph{のもとで積 $G(x)\cdot\Gamma(x)$ が消えないように、すなわち系 (1) の代数的階数が保存されるように選ばれるものとする。逐次置換}
\[
\begin{aligned}
 [h(x)]_{x_n=a_n}&=h^{(n-1)}(x),\qquad
 [h^{(n-1)}(x)]_{x_{n-1}=a_{n-1}}=h^{(n-2)}(x),\ldots,\\
 [h^{(\rho+1)}(x)]_{x_{\rho+1}=a_{\rho+1}}&=h^*(x_1,x_2,\ldots,x_\rho)
\end{aligned}
\]
\emph{において、関数 $h(x)$, $h^{(n-1)}(x)$, \ldots, $h^{(\rho+1)}(x)$ をそれぞれ既約表示に置く。これらの仮定のもとで、関数 $h^*(x_1,\ldots,x_\rho)$ では分子も分母も恒等的には消えず、また関数は $\frac{0}{0}$ になることもない。}\footnote{例えば前注の関数 $h(x)$ では、逐次置換により
\[
 [h(x)]_{x_4=0}=h^{(3)}(x)=\frac{(\alpha x_1+\beta x_2)}{(\alpha' x_1+\beta' x_2)}.\qquad
 h^*(x_1 x_2)=\frac{\alpha x_1+\beta x_2}{\alpha' x_1+\beta' x_2}
\]
が得られる。}

\section*{\S\ 3. 系 $\Ssys_{n\rho}$ から系 $\Ssys^*_{\rho\rho}$ への一対一対応.}

補助定理から、いくつかの不定元を数で置き換えることにより、系 $\Ssys_{n\rho}$ を系 $\Ssys^*_{\rho\rho}$ へ対応させることが直ちに得られる。このとき代数的階数 $\rho$ は本来的な特徴数として現れる。

$\Ssys_{n\rho}$ は代数的階数 $\rho$ をもつので、$\Ssys_{n\rho}$ から $\rho$ 個の関数
\begin{equation}
  g_1(x),\ldots,g_\rho(x)
\tag{1}
\end{equation}
が存在し、それらは代数的に独立である。さらに、$f(x)$ が $\Ssys_{n\rho}$ からの任意の関数を表すならば、系
\begin{equation}
  g_1(x),\ldots,g_\rho(x),\ f(x)
\tag{2}
\end{equation}
もまた代数的階数 $\rho$ をもち、系 (1) と (2) は補助定理の条件を満たす。\footnote{必要ならば、不定元を番号づけ直す。}

したがって、数
\[
  a_n,\ a_{n-1},\ldots,a_{\rho+1}
\]
を補助定理に従って、系 (1) の代数的階数 $\rho$ が置換によって保存されるように定める。このことは任意に多くの仕方で可能である。そのとき
\[
\begin{aligned}
 [f(x)]_{x_n=a_n}&=f^{(n-1)}(x),\qquad
 [f^{(n-1)}(x)]_{x_{n-1}=a_{n-1}}=f^{(n-2)}(x);\ldots,\\
 [f^{(\rho+1)}(x)]_{x_{\rho+1}=a_{\rho+1}}&=f^*(x_1,\ldots,x_\rho)
\end{aligned}
\]
で定義される関数、すなわち
\begin{equation}
  f^*(x_1,\ldots,x_\rho)
\tag{3}
\end{equation}
では、分子も分母も恒等的に消えることができない。

いま $f(x)$ に $\Ssys_{n\rho}$ のすべての（有限または無限個の）関数を走らせる。すべての関数 (3) 全体からなる系を考えると、それは次の性質をもつ。

1. 代数的階数 $\rho$ をもつ。実際、それは (1) から生じる関数
\[
  g_1^*(x_1\cdots x_\rho)\cdots g_\rho^*(x_1\cdots x_\rho)
\]
を含み、これらは $a_n,a_{n-1},\ldots,a_{\rho+1}$ の選び方により代数的に独立である。したがってこの系を
\[
  \Ssys^*_{\rho\rho}
\]
と表す。\footnote{対応して、$f^*(x)$, $g^*(x)$, \ldots は、以下いつでも対応関数、すなわち $x_1\cdots x_\rho$ の関数を意味するものとする。}

2. 関数 $f(x)$ と $f^*(x)$ の対応は例外なく一対一である。

実際、$f_1(x)$ と $f_2(x)$ が $\Ssys_{n\rho}$ に属すならば、その差
\[
  d(x)=f_1(x)-f_2(x)
\]
も系 (1) に代数的に従属であり、さらに
\[
  d^*(x)=f_1^*(x)-f_2^*(x)
\]
である。$d(x)$ は系 (1) とともに補助定理の条件を満たすので、
\[
  d(x)\ne0
\]
から必ず
\[
  d^*(x)\ne0
\]
が従う。すなわち、$\Ssys_{n\rho}$ の異なる関数には、$\Ssys^*_{\rho\rho}$ の異なる関数が対応する。

3. $\Ssys^*_{\rho\rho}$ の関数の間の任意の関係
\[
  F\bigl(f_1^*(x),f_2^*(x),\ldots,f_k^*(x)\bigr)=0
  \quad\text{$x_1,\ldots,x_\rho$ について恒等的に}
\]
から、一意に対応する $\Ssys_{n\rho}$ の関数について
\[
  F\bigl(f_1(x),f_2(x),\ldots,f_k(x)\bigr)=0
  \quad\text{$x_1,\ldots,x_n$ について恒等的に}
\]
が従う。なぜなら、$f_1(x),f_2(x),\ldots,f_k(x)$ とともに、これらの任意の有理結合も系 (1) に代数的に従属だからである。

さらにこれは
\begin{equation}
\begin{aligned}
 h(x)&=F\bigl(f_1(x),\ldots,f_k(x)\bigr):\\
 h^*(x)&=F\bigl(f_1^*(x),\ldots,f_k^*(x)\bigr)
\end{aligned}
\tag{4}
\end{equation}
からも従う。補助定理により、したがって
\[
  h(x)\ne0\quad\text{ならば}\quad h^*(x)\ne0
\]
である。証明終わり。

要約すれば次を得る。

\emph{定理 I: 有理関数の（有限または無限の）任意の系 $\Ssys_{n\rho}$ に対して、任意に多くの仕方で、一対一の系 $\Ssys^*_{\rho\rho}$ を割り当てることができる。その際、$\Ssys^*_{\rho\rho}$ の関数間に存在するすべての関係は $\Ssys_{n\rho}$ においても保存され、また逆も成り立つ。特に (4) により、体 $\Kfield_{n\rho}$ には再び体 $\Kfield^*_{\rho\rho}$ が対応する。}

定理 I から、以後のすべての証明を簡単にする帰結を引き出す。\emph{系の関数間の有限または無限個の関係によって特徴づけられる系 $\Ssys_{n\rho}$ の性質は、対応する系 $\Ssys^*_{\rho\rho}$ について成り立てば、もとの系 $\Ssys_{n\rho}$ についても成り立つ。}

この帰結を略して「転移原理」と呼ぶ。

\providecommand{\Lsys}{\mathfrak{L}}
\providecommand{\Jdom}{\mathfrak{J}}

\clearpage
この対応の最も簡単な例は、代数的不変式の理論によって与えられる。実際、線形変換によって、基本形式のいくつかの係数を数値的に固定し、特殊な形式に対して成り立つすべての関係が一般の場合にも成り立つようにすることが常にできる。

\section*{\S\ 4. 体 $\Kfield_{n\rho}$ の有理基底.}

以下の研究では、基底概念のまわりにまとめられた問題に対して、\S\ 3 の「転移原理」を一貫して用いる。まず体について基底の存在を証明するが、定義は一般に述べる。

\emph{定義 III: $\Ssys_{n\rho}$ の有限個の関数が有理基底と呼ばれるのは、$\Ssys_{n\rho}$ の任意の関数が、$\Omega$ からの係数を用いて、それら有限個の関数の有理結合として表されるときである。}

有理基底は $\rho$ 個の代数的に独立な関数を含まなければならない。なぜなら、有理基底から導かれる系の代数的階数は、有理基底の代数的階数に等しいからである。

まず、有理基底の存在証明が転移原理を許すことを示す。すなわち、定義 III は次のように定式化できる。

$\Ssys_{n\rho}$ の関数
\[
  g_1(x),\ g_2(x),\ldots, g_k(x)
\]
が有理基底をなすことと、無限個の関係
\begin{equation}
  f(x)=\gamma\bigl(g_1(x),\ldots,g_k(x)\bigr)
\tag{1}
\end{equation}
が満たされることは同値である。ここで $f(x)$ は $\Ssys_{n\rho}$ のすべての関数を走り、$\gamma$ は $f$ に応じて変わる $k$ 個の引数の有理関数を表す。

したがって、対応系 $\Ssys^*_{\rho\rho}$ において
\[
  g_1^*(x),\ g_2^*(x),\ldots,g_k^*(x)
\]
が有理基底を与え、それにより関係
\[
  f^*(x)=\gamma\bigl(g_1^*(x)\cdots g_k^*(x)\bigr)
\]
が生じるならば、定理 I によって、関係 (1) は $\Ssys_{n\rho}$ に対して成り立つ。したがって、\emph{有理基底に対する転移原理}は次のようにいう。

\emph{対応系 $\Ssys^*_{\rho\rho}$ に有理基底が存在することから、もとの系 $\Ssys_{n\rho}$ の有理基底が従う。}\footnote{転移原理は、ここで扱うすべての基底問題について同様に定式化される。}

すべての体 $\Kfield_{n\rho}$ の有理基底の存在を証明するためには、したがって体 $\Kfield^*_{\rho\rho}$ に限ればよい。ここでは E. Steinitz が用いた議論の直接の一般化が目的に至る。\footnote{E. Steinitz, \emph{Algebraische Theorie der Körper}, \S\ 24, Crelle's Journal 137 (1909).}

\begin{equation}
  g_1^*(x_1,\ldots,x_\rho),\ldots,
  g_\rho^*(x_1,\ldots,x_\rho)
\tag{2}
\end{equation}
を、$\Kfield^*_{\rho\rho}$ からの $\rho$ 個の代数的に独立な関数の系とする。$x_1,\ldots,x_\rho$ を消去すると、$\rho$ 個の関係
\begin{equation}
\begin{aligned}
  F_1\bigl(g_1^*(x),\ldots,g_\rho^*(x),x_1\bigr)&=0,\ldots,\\
  F_\rho\bigl(g_1^*(x),\ldots,g_\rho^*(x),x_\rho\bigr)&=0
\end{aligned}
\tag{3}
\end{equation}
が得られる。各関係 $F_i=0$ には不定元 $x_i$ が実際に現れる。そうでなければ、仮定に反して関数 (2) の間に代数的従属性が存在することになる。関係 (3) は
\[
  x_1,x_2,\ldots,x_\rho
\]
が
\[
  \Omega\bigl(g_1^*(x),\ldots,g_\rho^*(x)\bigr)
\]
に関して代数的元であることをいう。したがって、$x_1,\ldots,x_\rho$ を添加して得られる体
\[
  \Omega\bigl(g_1^*,\ldots,g_\rho^*,x_1,
  \ldots,x_\rho\bigr)=\Omega(x_1,
  \ldots,x_\rho)
\]
は、$\Omega(g_1^*,\ldots,g_\rho^*)$ の有限代数拡大である。既知の定理\footnote{例えば Weber, \emph{Lehrbuch d. Algebra} 1, \S\ 144（第1版）または \S\ 55（小版）を比較せよ。}により、$\Kfield^*_{\rho\rho}$ は、$\Omega(g_1^*\cdots g_\rho^*)$ と $\Omega(x_1\cdots x_\rho)$ の間の中間体として、それ自身 $\Omega(g_1^*\cdots g_\rho^*)$ の代数拡大であり、また $\Omega(x_1\cdots x_\rho)$ は $\Kfield^*_{\rho\rho}$ の代数拡大である。したがって $\Kfield^*_{\rho\rho}$ は、$\Omega(g_1^*\cdots g_\rho^*)$ に $\Kfield^*_{\rho\rho}$ からの有限個の元を添加することによって、または適当に選ばれた一つの関数を添加することによって得られる。すなわち
\[
\Kfield^*_{\rho\rho}
=\Omega(g_1^*\cdots g_\rho^*g_{\rho+1}^*\cdots g_k^*)
=\Omega(g_1^*\cdots g_\rho^*g_0^*)
\]
が成り立つ。転移原理によって次を得る。

\emph{定理 II: 任意の体 $\Kfield_{n\rho}$ は代数階数 $\rho$ の有理基底をもち、その有理基底は高々 $\rho+1$ 個の関数からなるように選べる。}

$\Kfield_{n\rho}$ から任意に選んだ $\rho$ 個の代数的に独立な関数 $h_1(x),\ldots,h_\rho(x)$ は、上の方法によって有理基底
\[
h_1(x),\ldots,h_\rho(x);\ h_{\rho+1}(x),\ldots,h_\sigma(x)
\]
へ拡張できる。定義により関係式
\[
\begin{array}{ll}
h_i(x)=\chi_i\bigl(g_1(x),\ldots,g_k(x)\bigr)&(i=1,2,\ldots,\sigma),\\
g_j(x)=\psi_j\bigl(h_1(x),\ldots,h_\sigma(x)\bigr)&(j=1,2,\ldots,k)
\end{array}
\]
が存在する。ここで $\chi_i,\psi_j$ はそれぞれの引数の有理関数である。したがって、任意の一つの有理基底から、他のすべての有理基底が双有理変換によって得られる。また、$h_1(x),\ldots,h_\rho(x)$ が $\Kfield_{n\rho}$ の任意の $\rho$ 個の代数的に独立な関数であるなら、$\Kfield_{n\rho}$ は
\[
\Omega\bigl(h_1(x),\ldots,h_\rho(x)\bigr)
\]
上の有限代数拡大である。

\section*{\S\ 5. 体 $\Kfield_{n\rho}$ のインボリューション形式とインボリューション基底.}

§4 で構成した各有理基底には、任意に選んだ出発関数に依存するという欠点があったが、ここでは\emph{体によって一意に定まる有理基底}の存在を示す。

§4 により、$\Omega(x_1\cdots x_\rho)$ は $\Kfield^*_{\rho\rho}$ 上の代数体、例えば次数 $i$ の代数体であり、
\[
\Kfield^*_{\rho\rho}(x_1\cdots x_\rho)=\Omega(x_1\cdots x_\rho)
\]
であるから、$\Kfield^*_{\rho\rho}$ に関して代数的な元 $x_1\cdots x_\rho$ を添加することによって生じる。そこで、根
\[
z=u_1x_1+u_2x_2+\cdots+u_\rho x_\rho=u(x)
\]
をもつ、$\Kfield^*_{\rho\rho}$ に関して既約な整関数 $\Phi^*(z,u)$ を考える。これは $z$ について $i$ 次である。さらに $\Phi^*(z,u)$ は、$[z-u(x)]$ と共役な量の積として、$z,u_1\cdots u_\rho$ に関する次数 $i$ の斉次形式であり、したがって $z^i$ の係数として $u$ に依存しない $\Kfield^*_{\rho\rho}$ の関数を含む。この最高係数を除して $1$ にすれば、既約関数 $\Phi^*(z,u)$ は一意に定まる。このように定義された斉次形式\begin{equation}
\begin{aligned}
  \Phi^*(z,u)
  &=z^i+
    \sum{}' g^*_{\alpha\alpha_1\ldots\alpha_\rho}(x_1,
    \ldots,x_\rho)
    z^\alpha u_1^{\alpha_1}\cdots u_\rho^{\alpha_\rho},\\
  &\hspace{8em}(\alpha+\alpha_1+\cdots+\alpha_\rho=i)
\end{aligned}
\tag{1}
\end{equation}
を $\Kfield^*_{\rho\rho}$ の\emph{インボリューション形式}\footnote{この名称の理由はこの段落の末尾で明らかになる。$\sum'$ は、和を $z^i$ を除くすべての項にわたってとることを意味する。}という。(1) から
\begin{equation}
\begin{aligned}
  \Phi(z,u)
  &=z^i+
    \sum{}' g_{\alpha\alpha_1\ldots\alpha_\rho}(x)\cdot
    z^\alpha u_1^{\alpha_1}\cdots u_\rho^{\alpha_\rho},\\
  &\hspace{8em}(\alpha+\alpha_1+\cdots+\alpha_\rho=i)
\end{aligned}
\tag{2}
\end{equation}
が $\Kfield_{n\rho}$ の\emph{インボリューション形式}として得られる。

インボリューション形式 $\Phi^*(z,u)$ の係数は、$\Kfield^*_{\rho\rho}$ の有理基底であることが分かる。すなわち関係
\begin{equation}
  \Omega\bigl(g^*_{\alpha\alpha_1\ldots\alpha_\rho}(x_1,
  \ldots,x_\rho)\bigr)=\Kfield^*_{\rho\rho}
  \quad(\alpha+\alpha_1+\cdots+\alpha_\rho=i)
\tag{3}
\end{equation}
が成り立つ。これを示すために、$\Phi^*(z,u)$ は $\Omega(g^*_{\alpha\alpha_1\ldots\alpha_\rho}(x))$ に属する係数をもち、しかも $\Kfield^*_{\rho\rho}$ に関して既約であるから、なおさらこの部分体に関して既約であることに注意する。したがって $\Phi^*(z,u)$ は、$\Omega(g^*_{\alpha\alpha_1\ldots\alpha_\rho}(x))$ に関する $u(x)$ の既約方程式を表す。これより、$\Omega(x_1,\ldots,x_\rho)$ は $\Omega(g^*_{\alpha\alpha_1\ldots\alpha_\rho}(x))$ 上の、次数 $i$ の代数体である。$\Kfield^*_{\rho\rho}$ は $\Omega(g^*_{\alpha\alpha_1\ldots\alpha_\rho}(x))$ と $\Omega(x_1\cdots x_\rho)$ の間にあるから、次数の一致から体の一致、すなわち関係 (3) が従うことが知られている。転移原理により、$\Phi(z,u)$ の係数は対応して $\Kfield_{n\rho}$ の有理基底を与える。すなわち次を得る。

\emph{定理 III: インボリューション形式の係数は、体によって定まる有理基底、すなわちインボリューション基底をなす。体 $\Kfield_{\rho\rho}$ については、この基底は一意に定まる。}\footnote{体 $\Kfield_{n\rho}$ については、インボリューション形式は対応の選択に依存しうる。不定係数をもつ対応を用いれば一意性を得ることができる。}

この事実に非有理的な解釈を付け加える。これはインボリューションの幾何学的理論とのつながりを与える。

$\Phi^*(z,u)$ の一次因子への分解は、拡大体において
\begin{equation}
\begin{aligned}
 \Phi^*(z,u)&=(z-u_1x_1-\cdots-u_\rho x_\rho)
 (z-u_1x_1^{(1)}-\cdots-u_\rho x_\rho^{(1)})\cdots\\
 &\qquad\cdot
 (z-u_1x_1^{(i-1)}-\cdots-u_\rho x_\rho^{(i-1)})
\end{aligned}
\tag{4}
\end{equation}
によって与えられる。(1) と比較すれば、関数
\[
  g^*_{\alpha\alpha_1\ldots\alpha_\rho}(x_1,
  \ldots,x_\rho)\quad(\alpha+\alpha_1+\cdots+\alpha_\rho=i)
\]
は、量の行
\begin{equation}
\begin{array}{cccc}
  x_1 & x_2 & \cdots & x_\rho\\
  x_1^{(1)} & x_2^{(1)} & \cdots & x_\rho^{(1)}\\
  \vdots & \vdots & & \vdots\\
  x_1^{(i-1)} & x_2^{(i-1)} & \cdots & x_\rho^{(i-1)}
\end{array}
\tag{5}
\end{equation}
の基本対称関数と同一であることが分かる。したがって $\Kfield^*_{\rho\rho}$ の任意の関数は、行 (5) に関して対称なこれらの基本関数の有理結合であり、逆に行 (5) の任意の対称関数は $\Kfield^*_{\rho\rho}$ に属する。

したがって
\[
  f^*(x)=\frac{F^*(x)}{H^*(x)}
\]
が $\Kfield^*_{\rho\rho}$ のすべての関数を走るとき、無限個の関係の系
\begin{equation}
  \frac{F^*(x)}{H^*(x)}=
  \frac{F^*(x^{(1)})}{H^*(x^{(1)})}=\cdots=
  \frac{F^*(x^{(i-1)})}{H^*(x^{(i-1)})}
\tag{6}
\end{equation}
が存在する。要約すれば、
\[
  F^*(x^{(k)})H^*(x^{(l)})-
  F^*(x^{(l)})H^*(x^{(k)})=0
  \quad(k,l=0,1,\cdots,i-1),
\]
である。ここで $F^*(x^{(k)})$ も $H^*(x^{(k)})$ も恒等的には消えない。これらは $F^*(x)$ および $H^*(x)$ の共役であり、形式的に対称だからである。

逆に、無限個の関係
\begin{equation}
  F^*(x)H^*(\xi)-H^*(x)F^*(\xi)=0,
  \qquad H^*(\xi)\ne0
\tag{7}
\end{equation}
を満たす任意の量の行 $\xi$ について、量の行 (5) への所属が得られる。

実際、$G^*(x)$ を $\Phi^*(z,u)$ の係数の共通分母とすると、仮定 (7) により関数
\[
  G^*(\xi)\cdot\Phi^*(z,u;\xi)
\]
は、$\xi$ に関しても整であり、恒等的には消えない。すなわち $\Phi^*(z,u;\xi)$ の係数は不定にならず、かつ零点
\[
  z=u_1\xi_1+u_2\xi_2+\cdots+u_\rho\xi_\rho
\]
をもつ。(7) により、したがって $\Phi^*(z,u;x)$ も零点 $z=u(\xi)$ をもつ。すなわち $\xi$ は量の行 (5) に属する。$\xi$ が共役な量の行に関する方程式系
\[
  F^*(x^{(k)})H^*(\xi)-H^*(x^{(k)})F^*(\xi)=0;
  \qquad H^*(\xi)\ne0
\]
を満たす場合にも、(6) により対応する主張が成り立つ。

したがって、量の行 (5) は、方程式系
\[
  F^*(x)H^*(\xi)-H^*(x)F^*(\xi)=0;
  \qquad H^*(\xi)\ne0
\]
の解 $\xi$ として定義される。ここで $F^*(x):H^*(x)$ は $\Kfield^*_{\rho\rho}$ のすべての関数を走る。また行 $x$ は任意の共役な行で置き換えてもよい。

幾何学的には、各行 $x$ を点と解釈すると、これは次のことを意味する。(7) の解は、$\rho$ 次元線形空間において、各々 $i$ 個の点からなる $\infty^\rho$ 個の群を表し、群の任意の点がその個々の群を一意に定める。このような点群の系を「\emph{線形空間におけるインボリューション}」と呼ぶ。\footnote{(7) の除外された解、すなわち $F^*(\xi)=0$, $H^*(\xi)=0$ となるもの、およびインボリューション基底に対応する方程式系の除外された解は、インボリューションの基本点と呼ばれる。これには同次化によって生じるもの、すなわち「無限遠にある」点も含まれる。}

対応して、体 $\Kfield_{n\rho}$ は、曲線、曲面などからなる線形空間内のインボリューションを特徴づける。

以上により、$\Omega(x_1,\ldots,x_\rho)$ の $\Kfield^*_{\rho\rho}$ に関する次数 $i$ は、側条件 $H^*(\xi)\ne0$ を満たす (7) の解系の個数に等しい。これを $\Kfield^*_{\rho\rho}$ の解系と呼ぶことができる。したがって、インボリューション基底の存在を証明するために用いた議論は、次の非有理的な定式化も許す。

「$\mathfrak{L}^*$ が $\Kfield^*_{\rho\rho}$ の約体であり、体 $\mathfrak{L}^*$ のすべての解系が $\Kfield^*_{\rho\rho}$ の解系でもあるならば、$\mathfrak{L}^*$ と $\Kfield^*_{\rho\rho}$ は同一である。」

対応する命題は、転移原理により、$\Kfield_{n\rho}$ の約体 $\mathfrak{L}$ についても成り立つ。

\section*{\S\ 6. 体 $\Kfield_{n\rho}$ の最小基底.}

この節では既知の定理を概観する。ただし、転移原理と有理基底の存在によって可能となった拡張された定式化で述べる。

\emph{定義 IV: $\Kfield_{n\rho}$ の有理基底が最小基底と呼ばれるのは、それが最小個数 $\rho$ の関数だけを含むときであり、その場合それらは代数的に独立でなければならない。}

Lüroth, Castelnuovo, Enriques の定理\footnote{J. Lüroth: Beweis eines Satzes über rationale Kurven, \emph{Math. Ann.} 9 (1875). --- G. Castelnuovo: Sulla razionalità delle involuzioni piane, \emph{Math. Ann.} 44 (1893). --- F. Enriques: Sopra una involuzione non razionale dello spazio, \emph{Rend. Acc. Linc.}, vol. 21, 1912年1月21日.}から次の事実が得られる。

I. 任意の体 $\Kfield_{n1}$ は最小基底をもつ。それはその体の Lüroth 関数であり、一次分数変換を除いて一意に定まる。

II. 体 $\Kfield_{n2}$ については、係数体 $\Omega$ の代数拡大を許すならば、そして一般にはその場合に限って、常に最小基底が存在する。

III. 体 $\Kfield_{n\rho}$（$\rho\ge3$）については、一般には最小基底は存在しない。最小基底をもつ特殊な体 $\Kfield_{n\rho}$ の特徴づけはまだ試みられていない。

E. Steinitz が与えた体 $\Kfield^*_{11}$ に対する Lüroth の定理の証明を簡単に示す。\footnote{前掲書, \S\ 24.}

$\Phi^*(z,u)$ を $\Kfield^*_{11}$ のインボリューション形式とし、$\psi^*(x)$ を、実際に不定元 $x$ を含む $\Phi^*(z,u)$ の任意の係数とする。すると、$\Omega(x)$ の $\Omega(\psi^*(x))$ に関する次数は、$\Omega(x)$ の $\Kfield^*_{11}$ に関する次数に等しいことが分かる。そこから体の同一性が従う。さらに、$\Omega(\chi^*(x))$ が $\Omega(\psi^*(x))$ に等しいならば、$\chi^*(x)$ と $\psi^*(x)$ の相互の有理従属性は、二つの関数の一次分数従属性を意味する。

したがって、転移原理により、Lüroth の定理は次の形でも述べられる。

\emph{体 $\Kfield_{n1}$ の最小基底は、インボリューション基底の任意の一つの関数からなる。また $\Kfield_{n1}$ のインボリューション基底の任意の二つの関数は、一次分数変換によって互いに依存する。}

G. Castelnuovo は、幾何学的なインボリューション理論を用いて、有理基底が三つの関数から成る体 $\Kfield^*_{22}$ について事実 II を証明している。転移原理は係数体 $\Omega$ の有限代数拡大のもとでも有効であるから、有理基底の存在によって、すべての体 $\Kfield_{n2}$ について証明が完結する。

III については、Enriques が三次元空間内の非有理的インボリューション、すなわち最小基底をもたない体 $\Kfield_{33}$ を構成していることを注意しておく。

%
\providecommand{\Kfield}{\mathfrak{K}}
\providecommand{\Ssys}{\mathfrak{S}}
\providecommand{\Jdom}{\mathfrak{J}}
\providecommand{\Gdom}{\mathfrak{G}}
\providecommand{\Lsys}{\mathfrak{L}}
\providecommand{\Omegaint}{[\Omega]}
\providecommand{\eps}{\varepsilon}

\section*{\S\ 7. 有理関数の任意の系 $\Ssys_{n\rho}$、線形族 $\Lsys_{n\rho}$、および整域 $\Jdom_{n\rho}$. 有理基底の存在.}

体 $\Kfield_{n\rho}$ の有理基底の存在から、いま、有理関数および整有理関数の任意の系に対する有理基底を導く。これらの系を、加法、乗法、除法という基本演算によって順次互いに生じる順に並べる。

I. いつものように、$\Ssys_{n\rho}$ は、$n$ 個の不定元の有理関数（または整有理関数）からなり、代数的階数 $\rho$ をもつ任意の系を表す。$\Omega$ の量も系に属するものとして数える。

II. 系 $\Lsys_{n\rho}$ が次の条件を満たすとき、これを\emph{線形族}と呼ぶ。

1. $f(x)$ とともに、$\Omega$ の任意の量 $c$ に対して $c\cdot f(x)$ も常に $\Lsys_{n\rho}$ に含まれる。

2. $f(x)$ と $g(x)$ とともに、$f(x)+g(x)$ も常に $\Lsys_{n\rho}$ に含まれる。

III. 線形族 $\Jdom_{n\rho}$ がさらに次の条件を満たすとき、これを\emph{整域}と呼ぶ。

3. $f(x)$ と $g(x)$ とともに、$f(x)\cdot g(x)$ も常に $\Jdom_{n\rho}$ に含まれる。

IV. 整域 $\Kfield_{n\rho}$ がさらに次の条件を満たすとき、これを\emph{体}と呼ぶ。

4. $f(x)$ と $g(x)$ とともに、ただし $g(x)\ne0$ のとき、商 $f(x):g(x)$ も常に $\Kfield_{n\rho}$ に含まれる。

条件 1 から 4 は、\S\ 1 で体について述べた条件である。\footnote{$f(x)$ と $g(x)$ は零次、すなわち $\Omega$ の量であってもよい。}

これらの領域を、以下で順次互いから導く。

a) 任意の系 $\Ssys_{n\rho}$ は、次の要請によって線形族 $\Lsys_{n\rho}$ に拡張される。

$\Lsys_{n\rho}$ は、$\Ssys_{n\rho}$ に加えて、$\Ssys_{n\rho}$ の有限個の関数によって $\Omega$ の係数で線形に表される関数 $h(x)$ のすべて、かつそのようなものだけを含む。
\[
  h(x)=c_1f_1(x)+c_2f_2(x)+\cdots+c_kf_k(x),
\]
ここで $f_1(x),\ldots,f_k(x)$ は $\Ssys_{n\rho}$ のすべての関数を走り、$c_1,\ldots,c_k$ は $\Omega$ のすべての量を走る。

実際、系の関数の有理的結合を付け加えても、系の代数的階数は保たれる。$\Lsys_{n\rho}$ は条件 1 および 2 を満たすので線形族である。さらに、$\Ssys_{n\rho}$ を含む任意の線形族は、1 と 2 によって $\Lsys_{n\rho}$ も含む。したがって $\Lsys_{n\rho}$ は $\Ssys_{n\rho}$ を含む最小の線形族である。

b) 同様に、線形族 $\Lsys_{n\rho}$ から、それを含む最小の整域 $\Jdom_{n\rho}$ が次の要請によって生じる。

$\Jdom_{n\rho}$ は、$\Lsys_{n\rho}$ に加えて、$\Lsys_{n\rho}$ の有限個の関数 $f_1(x),\ldots,f_k(x)$ の、$\Omega$ 係数の\emph{整有理結合}として表される関数 $h(x)$ のすべて、かつそのようなものだけを含む。ここで $f_1(x),\ldots,f_k(x)$ は $\Lsys_{n\rho}$ のすべての関数を走る。

c) 最後に、整域 $\Jdom_{n\rho}$ から、それを含む最小の体 $\Kfield_{n\rho}$ が次の要請によって生じる。すなわち、$\Kfield_{n\rho}$ は $\Jdom_{n\rho}$ に加えて、$\Jdom_{n\rho}$ の二つの関数の商として表される関数 $h(x)$ のすべて、かつそのようなものだけを含む。

三つの拡張を一つの段階にまとめると、次を得る。

\emph{任意の系 $\Ssys_{n\rho}$ は、次の要請によって、それを含む最小の体 $\Kfield_{n\rho}$ へ拡張できる。}

\emph{$\Kfield_{n\rho}$ は、$\Ssys_{n\rho}$ に加えて、$\Ssys_{n\rho}$ の有限個の関数 $f_1(x),\ldots,f_k(x)$ の、$\Omega$ 係数の有理的結合として表される関数 $h(x)$ のすべて、かつそのようなものだけを含む。ここで $f_1(x),\ldots,f_k(x)$ は $\Ssys_{n\rho}$ のすべての関数を走る。}

この最後の事実から、任意の領域に対する有理基底の存在が直ちに従う。

$\Kfield_{n\rho}$ を $\Ssys_{n\rho}$ を含む最小の体とし、$\Kfield_{n\rho}$ の有理基底が
\begin{equation}
  h_1(x),\ h_2(x),\ldots,h_\sigma(x)
\tag{1}
\end{equation}
で与えられているとする。$\Kfield_{n\rho}$ の定義により、関係式
\begin{equation}
\begin{aligned}
 h_1(x)&=\gamma_1\bigl(g_1(x),\ldots,g_\lambda(x)\bigr);\ldots;\\
 h_\sigma(x)&=\gamma_\sigma\bigl(g_\mu(x),\ldots,g_\nu(x)\bigr),
\end{aligned}
\tag{2}
\end{equation}
が存在する。ただし $\gamma_1,\ldots,r_\sigma$ はその引数の有理関数であり、
\begin{equation}
  g_1(x),\ldots,g_\lambda(x),\ldots,
  g_\mu(x),\ldots,g_\nu(x)
\tag{3}
\end{equation}
はすべて $\Ssys_{n\rho}$ に属する。関係式 (2) により、(3) も (1) と同じく $\Kfield_{n\rho}$ の有理基底を表す。また (3) の関数は $\Ssys_{n\rho}$ に属するから、これは同時に $\Ssys_{n\rho}$ の有理基底である。したがって次を得る。

\emph{定理 IV: 有理関数（または整有理関数）の任意の系 $\Ssys_{n\rho}$ は、その系の有限個の関数から成る、代数的階数 $\rho$ の有理基底をもつ。}

つぎに、特に $\Lsys_{n\rho}$ を線形族とし、
\begin{equation}
  g_1(x),\ldots,g_\rho(x),\ g_{\rho+1}(x),\ldots,g_\nu(x)
\tag{4}
\end{equation}
を $\Lsys_{n\rho}$ の有理基底とする。例えば
\[
  g_1(x)\cdots g_\rho(x)
\]
が代数的に独立であるとする。このとき
\[
  g(x)=c_1g_{\rho+1}(x)+c_2g_{\rho+2}(x)+\cdots+c_{\nu-\rho}g_\nu(x)
\]
と置けば、$c_i$ を $\Omega$ の数として適当に選ぶことにより
\begin{equation}
  g_1(x),\ldots,g_\rho(x),\ g(x)
\tag{5}
\end{equation}
が $\Lsys_{n\rho}$ を含む最小の体 $\Kfield_{n\rho}$ の有理基底となるようにできる。ところが $\Lsys_{n\rho}$ は線形族であるから、$g(x)$ は $\Lsys_{n\rho}$ に属し、(5) は $\Lsys_{n\rho}$ の有理基底を表す。すなわち、

\emph{定理 V: 有理関数（または整有理関数）の任意の線形族 $\Lsys_{n\rho}$ の有理基底、従って特に任意の整域 $\Jdom_{n\rho}$ および任意の体 $\Kfield_{n\rho}$ の有理基底は、高々 $\rho+1$ 個の関数（かつ少なくとも $\rho$ 個）を含むように選ぶことができる。}

したがって、\S\ 4 で体について得られた有理基底の関数個数の制限は、体が線形族であるという性質に基づく。一方、最小基底が存在する場合に $\rho$ 個の関数まで減らせることは、体のすべての仮定を用いる。

なお、\S\ 3 (4) によって、系およびそれを含む最小系のすべての特徴的性質は写像によって保存される。

\section*{\S\ 8. 多項式からなる整域 $\Jdom_{n\rho}$ のインボリューション基底.}

以下の節では、元が整有理関数（多項式）である系 $\Ssys_{n\rho}$ を扱う。まず、多項式からなる整域 $\Jdom_{n\rho}$ を考える。これらの領域 $\Jdom_{n\rho}$ については、最初に割り切りの概念を定めなければならない。

$F(x),G(x)$ を $\Jdom_{n\rho}$ の二つの多項式とし、$L(x)$ をそれらの共通因子、すなわち
\[
  F(x)=L(x)\cdot F_1(x);\qquad G(x)=L(x)\cdot G_1(x)
\]
が成り立つものとする。このとき $L(x)$ が\emph{$\Jdom_{n\rho}$ における共通因子}であるとは、三つの多項式 $L(x),F_1(x),G_1(x)$ がすべて $\Jdom_{n\rho}$ に属すること、かつそのときに限る。

$\Jdom_{n\rho}$ の二つの多項式は、$\Jdom_{n\rho}$ における共通因子をもたないとき、\emph{$\Jdom_{n\rho}$ において互いに素}であるという。

$\Kfield_{n\rho}$ を $\Jdom_{n\rho}$ を含む最小の体とする。すると定義により、$\Kfield_{n\rho}$ の任意の関数は
\[
  f(x)=F_1(x):F_2(x)
\]
の形に表される。ここで $F_1(x),F_2(x)$ は $\Jdom_{n\rho}$ に属し、$\Jdom_{n\rho}$ において互いに素である。同様に、$\Kfield_{n\rho}$ の複数の関数は $\Jdom_{n\rho}$ 内の共通分母にそろえることができる。
\[
  f_1(x)=\frac{F_1(x)}{F(x)},\ldots,
  f_k(x)=\frac{F_k(x)}{F(x)}.
\]
$F(x)$ が\emph{$\Jdom_{n\rho}$ における最小共通分母}であるとは、$\Jdom_{n\rho}$ の多項式
\[
  F(x),F_1(x),\ldots,F_k(x)
\]
が $\Jdom_{n\rho}$ において互いに素であること、かつそのときに限る。

このような表示は複数の仕方で可能でありうる。なぜなら $\Jdom_{n\rho}$ では、一般に既約関数への一意分解の法則が成り立たないからである。

つぎに、$\Kfield_{n\rho}$ のインボリューション基底が $\Jdom_{n\rho}$ に対してもつ意味を調べる。そのため、インボリューション形式とインボリューション基底の概念を整域へ移す。$\Kfield_{n\rho}$ のインボリューション形式を
\[
  \Phi(z,u)=z^i+\sum{}' g_{\alpha\alpha_1\ldots\alpha_\rho}(x)
  z^\alpha u_1^{\alpha_1}\cdots u_\rho^{\alpha_\rho},
  \qquad
  (\alpha+\alpha_1+\cdots+\alpha_\rho=i)
\]
とし、$G(x)$ を $\Phi(z,u)$ の係数の $\Jdom_{n\rho}$ における最小共通分母とする。$G(x)$ を掛けると、$\Phi(z,u)$ は、係数が $\Jdom_{n\rho}$ の多項式であり、かつ $\Jdom_{n\rho}$ において互いに素である形式 $\Psi(z,u)$ へ移る。
\[
  G(x)\cdot\Phi(z,u)=\Psi(z,u)
  =G(x)z^i+
  \sum{}' G_{\alpha\alpha_1\ldots\alpha_\rho}(x)
  z^\alpha u_1^{\alpha_1}\cdots u_\rho^{\alpha_\rho},
  \qquad
  (\alpha+\alpha_1+\cdots+\alpha_o=i).
\]

$\Kfield_{n\rho}$ のインボリューション形式 $\Phi(z,u)$ に $\Jdom_{n\rho}$ の多項式を掛けることによって得られ、係数が $\Jdom_{n\rho}$ の多項式で、かつ $\Jdom_{n\rho}$ において互いに素であるような形式 $\Psi(z,u)$ を、$\Jdom_{n\rho}$ の\emph{インボリューション形式}と呼ぶ。また $\Psi(z,u)$ の係数を対応する\emph{インボリューション基底}と呼ぶ。

$\Kfield_{n\rho}$ のインボリューション基底の意味から、$\Jdom_{n\rho}$ の任意のインボリューション基底が同時に有理基底であることはまず明らかである。さらに調べるため、対応する最小の体 $\Kfield^*_{\rho\rho}$ で
\[
  z=u_1x_1+\cdots+u_\rho x_\rho
\]
を根にもつ既約方程式が $\Phi^*(z,u)=0$ で与えられるような写像領域 $\Jdom^*_{\rho\rho}$ を考える。

\S\ 5 によれば、
\[
  \Phi^*(z,u)=z^i+\sum{}' g^*_{\alpha\alpha_1\ldots\alpha_\rho}(x)
  z^\alpha u_1^{\alpha_1}\cdots u_\rho^{\alpha_\rho},
  \qquad
  (\alpha+\alpha_1+\cdots+\alpha_\rho=i)
\]
の係数は、$\Kfield^*_{\rho\rho}$ に関して共役な量の行
\[
  x,\ x^{(1)}\cdots x^{(i-1)}
\]
の基本対称関数を表す。$F^*(x)$ を $\Kfield^*_{\rho\rho}$ の任意の多項式とする。すると
\[
  F^*(x)=\frac1i\Bigl(F^*(x)+F^*(x^{(1)})+
  \cdots+F^*(x^{(i-1)})\Bigr)
\]
であるから、$F^*(x)$ はこれらの量の行の整対称関数である。従って、それは基本対称関数、すなわち関数
\[
  g^*_{\alpha\alpha_1\ldots\alpha_\rho}(x)
\]
の整有理結合として表される。

いま
\[
  G^*(x)\cdot\Phi^*(z,u)=\Psi^*(z,u)
  =G^*(x)z^i+
  \sum{}'G^*_{\alpha\alpha_1\ldots\alpha_\rho}(x)
  z^\alpha u_1^{\alpha_1}\cdots u_\rho^{\alpha_\rho},
  \qquad
  (\alpha+\alpha_1+\cdots+\alpha_\rho=i)
\]
が $\Jdom^*_{\rho\rho}$ のインボリューション形式であるならば、$\Kfield^*_{\rho\rho}$ の任意の多項式、とくに $\Jdom^*_{\rho\rho}$ の任意の多項式は、商
\[
  G^*_{\alpha\alpha_1\ldots\alpha_\rho}(x):G^*(x)
\]
の整有理結合である。

転移原理を考慮すると、次を得る。

\emph{定理 VI: 多項式からなる任意の整域 $\Jdom_{n\rho}$ に対して、少なくとも一つの特別な有理基底、すなわちインボリューション基底が存在し、$\Jdom_{n\rho}$ のすべての多項式をこの基底の関数で有理表示するとき、分母にはただ一つの固定された関数（対応するインボリューション形式の最高係数）の冪だけが現れる。}

\section*{\S\ 9. 多項式からなる相対的に整な領域 $\Gdom_{n\rho}$.}

以下では、多項式からなる特殊な整域を考える。これは整域と体との中間段階をなす。

多項式の整域 $\Gdom_{n\rho}$ が次の条件を満たすとき、これを「相対的に整な領域」と呼ぶ。

\emph{4a) $\Gdom_{n\rho}$ は、$F(x)$ と $G(x)$ とともに、ただし $G(x)\ne0$ のとき、商 $F(x):G(x)$ を、その商が不定元について整、すなわち多項式である場合に、かつその場合に限り含む。}

まず、相対的に整な領域が、有理関数体 $\Kfield_{n\sigma}$（$\sigma\ge \rho$）の多項式全体と一致することを証明する。

$\Kfield_{n\rho}$ を $\Gdom_{n\rho}$ を含む最小の体とする。$\Kfield_{n\rho}$ の任意の関数 $f(x)$ には、$\Gdom_{n\rho}$ の二つの多項式の商としての表示
\[
  f(x)=F_1(x):F_2(x)
\]
がある。ここでは零次の多項式も許す。$f(x)$ が多項式になるならば、条件 4a) によって $\Gdom_{n\rho}$ に属する。また $\Gdom_{n\rho}$ はそれ以上の多項式を含むことはできないので、\emph{任意の相対的に整な領域 $\Gdom_{n\rho}$ は、それを含む最小の体の多項式全体と同一である。}

逆に、$\Kfield_{n\sigma}$ を任意の有理関数体（したがって、ある与えられた系を含む最小の体とは限らない）とし、$\Jdom$ を $\Kfield_{n\rho}$ に含まれる多項式全体とする。$\Jdom$ は整域の条件 1, 2, 3 および条件 4a) を満たすので、相対的に整な領域である。すなわち、\emph{体 $\Kfield_{n\sigma}$ に含まれる多項式全体は相対的に整な領域 $\Gdom_{n\rho}$（$\rho\le\sigma$）をなす。}\footnote{$\rho=0$ であってもよい。すなわち、$\Kfield_{n\sigma}$ の多項式全体が係数領域 $\Omega$ の元だけからなる場合である。}

さらに、相対的に整な領域が Hilbert の「相対的に整な関数の整域」と一致することを示す。\footnote{D. Hilbert: Mathematische Probleme. パリ講演 1900. 問題14 (Göttinger Nachrichten 1900).}

$\Gdom_{n\rho}$ の有理基底を
\begin{equation}
  G_1(x),\ldots,G_k(x)
\tag{1}
\end{equation}
とする。条件 1, 2, 3, 4a により、(1) の多項式の任意の有理結合で、不定元について整、すなわち多項式となるものは $\Gdom_{n\rho}$ に属する。逆に、$\Gdom_{n\rho}$ の任意の多項式は、有理基底の概念により、(1) の多項式で有理的に表される。従って $\Gdom_{n\rho}$ は、\emph{(1) の多項式の有理結合のうち、不定元について整、すなわち多項式であるもの全体と同一である。} こうして、特殊な有理基底から出発する Hilbert の定義が得られる。一方、われわれの定義は領域の抽象的性質だけを用いる。

相対的に整な領域については、\S\ 8 で一般の整域に対して与えた共通因子の定義が単純になる。

$F(x),G(x)$ を $\Gdom_{n\rho}$ の二つの多項式、$L(x)$ をそれらの共通因子とする。このとき $L(x)$ が\emph{$\Gdom_{n\rho}$ における共通因子}であるとは、$L(x)$ が $\Gdom_{n\rho}$ に属すること、かつそのときに限る。

というのは、商 $F(x):L(x)$ および $G(x):L(x)$ の属することは条件 4a) から従うからである。

相対的に整な領域のもう一つの本質的性質は、「$\Gdom_{n\rho}$ に関して代数的に整である」という概念を、代数的数体または体 $\Omega(x_1\cdots x_n)$ の場合とまったく同様に、任意の方程式によって定義できることである。

\emph{定義 V: 量 $\xi$ が $\Gdom_{n\rho}$ に関して代数的に整である、または相対的に代数的に整であるとは、最高係数が単位で、残りの係数が $\Gdom_{n\rho}$ の多項式であるような方程式を少なくとも一つ満たすことをいう。}

実際、相対的に代数的に整な量 $\xi$ が、例えば方程式
\[
  L(z)=z^\lambda+G_1(x)z^{\lambda-1}+\cdots+G_\lambda(x)=0
\]
を満たすとする。$L(z)$ は、最小の体 $\Kfield_{n\rho}$ に関して既約な整関数
\[
  M(z)=z^\mu+H_1(x)z^{\mu-1}+\cdots+H_\mu(x)
\]
で割り切れる。この割り切りから、Gauss の定理の既知の拡張\footnote{例えば J. König: \emph{Einleitung in die allgemeine Theorie der algebraischen Größen} (Leipzig, B.~G. Teubner, 1903), Kap. IX, \S\ 2, または Weber（小版）\S\ 20 を参照。}により、有理関数 $H_1(x)\cdots H_\mu(x)$ は $\Kfield_{n\rho}$ の多項式であり、従って $\Gdom_{n\rho}$ に属する。したがって、相対的に代数的に整な量 $\xi$ の定義は既約方程式から得られる。とくに、$\Gdom_{n\rho}$ の任意の多項式 $F(x)$ は、方程式 $z-F(x)=0$ を満たすから、$\Gdom_{n\rho}$ に関して代数的に整である。

さらに、多項式の整域 $\Jdom_{n\rho}$ に対して、最小の体の場合と同様に、それを含む最小の相対的に整な領域 $\Gdom_{n\rho}$ も次の要請で定義されることを述べておく。

$\Gdom_{n\rho}$ は、$\Jdom_{n\rho}$ に加えて、$\Jdom_{n\rho}$ の二つの多項式の商として表される多項式 $H(x)$ のすべて、かつそのようなものだけを含む。

この定義からさらに次が従う。$\Jdom_{n\rho}$ の任意のインボリューション形式およびインボリューション基底から、基底のすべての多項式の $\Gdom_{n\rho}$ における共通因子を、あっても一つ取り除くことにより、$\Gdom_{n\rho}$ の対応するものが得られる。

なぜなら $\Jdom_{n\rho}$ と $\Gdom_{n\rho}$ は同じ最小包含体 $\Kfield_{n\rho}$ をもつからである。また $\Jdom_{n\rho}$ の任意のインボリューション形式は、$\Kfield_{n\rho}$ のインボリューション形式に $\Jdom_{n\rho}$ の多項式を掛けることによって生じるので、$\Gdom_{n\rho}$ の係数をもつ。ただし、それらは $\Gdom_{n\rho}$ における共通因子をもつことがある。\footnote{例えば、$x_1^2$ と $x_1x_2$ のすべての整有理結合からなる整域 $\Jdom_{22}$ に対して、インボリューション形式として
\[
\Psi(z,u)=x_1^2\cdot z^2-x_1^4u_1^2-2x_1^3x_2u_1u_2-x_1^2x_2^2u_2^2
\]
が得られる。$x_1^2$ は $\Jdom_{22}$ に属し、従ってなおさら最小の相対的に整な包含領域 $\Gdom_{22}$ に属し、従って $\Gdom_{22}$ における共通因子であるから、$\Gdom_{22}$ のインボリューション形式として
\[
\Psi(z,u)=z^2-x_1^2u_1^2-2x_1x_2u_1u_2-x_2^2u_2^2
\]
が得られる。}

最後に、相対的に整な領域 $\Gdom_{n\rho}$ の\emph{写像領域}は、\S\ 7 により再び多項式の整域 $\Jdom^*_{\rho\rho}$ であるが、それ自体は必ずしも相対的に整な領域ではないことを注意する。実際、$F(x)$ と $G(x)$ が $\Gdom_{n\rho}$ の二つの多項式で、その商が多項式でなく従って $\Gdom_{n\rho}$ に属さないならば、商 $F^*(x):G^*(x)$ も $\Jdom^*_{\rho\rho}$ に属さない。しかし、この商は、いくつかの不定元を特殊化して生じるため、多項式でありうる。その場合、$\Jdom^*_{\rho\rho}$ では条件 4a) が満たされない。\footnote{例えば $\Gdom_{32}$ を、$F=x_1^2x_2$ と $G=x_1^2(1+x_3)$ の有理結合であるすべての多項式からなるものとする。許容される特殊化 $x_3=0$ は $F^*:G^*=x_2$ を与えるが、$F:G$ は多項式でない。従って $\Jdom^*_{22}$ は相対的に整な領域ではない。}

\section*{\S\ 10. 第一種の相対的に整な領域とその整基底.}

\S\ 9（定義 V）で導入した「相対的に代数的に整」という概念により、有限な整域である相対的に整な領域の一類を与えることができ、これによって D. Hilbert が提起した問題（数学の問題、問題 14）に、少なくともこの類について肯定的に答えることができる。次の定義を置く。

\emph{定義 VI: $\Gdom_{n\rho}$ の有限個の多項式であって、$\Gdom_{n\rho}$ の任意の多項式が、$\Omega$ の係数を用いてそれら有限個の整有理結合として表されるものを整基底という。整基底をもつ多項式の整域を有限整域という。}

ここで考える相対的に整な領域の類、すなわちインボリューション基底が整基底であることが示される類を、次の定義に従って第一種の相対的に整な領域と呼ぶ。

\emph{定義 VII: 相対的に整な領域 $\Gdom_{n\rho}$ が第一種であるとは、その写像領域として相対的に整な領域 $\Gdom^*_{\rho\rho}$ をもち、$x_1,\ldots,x_\rho$ の任意の多項式（係数は $\Omega$）が $\Gdom^*_{\rho\rho}$ 上代数的に整であることをいう。}

この定義により、不定元 $x_1,\ldots,x_\rho$、従って線形形式 $u(x)=u_1x_1+\cdots+u_\rho x_\rho$ も $\Gdom^*_{\rho\rho}$ 上代数的に整である。根 $z=u(x)$ をもつ既約整関数、すなわち最小の体 $\Kfield^*_{\rho\rho}$ のインボリューション形式は、従って最高係数として単位をもち、残りの係数として $\Gdom^*_{\rho\rho}$ の多項式をもつ。これは同時に $\Gdom^*_{\rho\rho}$ のインボリューション形式であり、$\Gdom^*_{\rho\rho}$ にはそれ以外のインボリューション形式は存在しない。従って転移原理により、$\Gdom_{n\rho}$ も最高係数が単位であるインボリューション形式をもつ。さらに定理 VI によれば、$\Gdom_{n\rho}$ のすべての多項式を対応するインボリューション基底で有理表示するとき、分母に現れるのはこの最高係数だけであるから、インボリューション基底は整基底になる。従って、

\emph{定理 VII. 第一種の任意の相対的に整な領域は有限整域である。その整基底は、それを特徴づける写像領域から得られるインボリューション基底で与えられる。特に、第一種の相対的に整な領域 $\Gdom_{\rho\rho}$ の整基底は、一意に定まる $\Gdom_{\rho\rho}$ のインボリューション基底によって与えられる。}

第一種の相対的に整な領域について判定条件を付け加える。$\Gdom_{n\rho}$ を第一種とし、$\Gdom_{n\rho}$ の整基底が
\begin{equation}
  F_1(x),F_2(x),\ldots,F_k(x)
\tag{1}
\end{equation}
で与えられているとする。すると Hilbert の補助定理\footnote{D. Hilbert: Über die vollen Invariantensysteme, Math. Ann. 42 (1893), \S\ 1. そこで同次多項式についてだけ述べられた補助定理は、非同次多項式についても同様に成り立つ。}により、$\Gdom_{n\rho}$ の $\rho$ 個の代数的に独立な多項式
\begin{equation}
  G_1(x),G_2(x),\ldots,G_\rho(x),
\tag{2}
\end{equation}
が存在し、(1) の多項式、従って $\Gdom_{n\rho}$ の任意の多項式は、(2) の多項式上代数的に整である。同じことは、写像領域 $\Gdom^*_{\rho\rho}$ の任意の多項式について、対応する多項式
\begin{equation}
  G_1^*(x),G_2^*(x),\ldots,G_\rho^*(x)
\tag{3}
\end{equation}
に関して成り立つ。

従って定義 VII により、$x_1,\ldots,x_\rho$ の任意の多項式も (3) の多項式上代数的に整である。

逆に、任意の相対的に整な領域 $\Gdom_{n\rho}$ の写像領域 $\Jdom^*_{\rho\rho}$ が、$x_1,\ldots,x_\rho$ の任意の多項式を代数的に整にするような $\rho$ 個の多項式 (3) を含むと仮定する。このとき $\Jdom^*_{\rho\rho}$ は相対的に整な領域 $\Gdom^*_{\rho\rho}$ になり、さらに $\Gdom^*_{\rho\rho}$、従って $\Gdom_{n\rho}$ が第一種であることを示す。

実際、$F^*(x)$ を $\Jdom^*_{\rho\rho}$ を含む最小の体 $\Kfield^*_{\rho\rho}$ の多項式とする。仮定により、それは (3) の多項式上代数的に整である。対応する $\Kfield_{n\rho}$ の関数は $\Gdom_{n\rho}$ の $\rho$ 個の多項式上代数的に整であるから、多項式 $F(x)$ であり、従って $\Gdom_{n\rho}$ に属する。したがって $F^*(x)$ は $\Jdom^*_{\rho\rho}$ に属する。よって写像領域 $\Jdom^*_{\rho\rho}$ は $\Kfield^*_{\rho\rho}$ のすべての多項式を含み、相対的に整な領域 $\Gdom^*_{\rho\rho}$ である。$\Gdom^*_{\rho\rho}$、従って $\Gdom_{n\rho}$ が第一種であることは、定義 V と VII および仮定から直ちに従う。すなわち、

\emph{第一種の相対的に整な領域 $\Gdom_{n\rho}$ は、写像領域 $\Jdom^*_{\rho\rho}$ の中に $\rho$ 個の多項式が存在し、それらの上で $x_1,\ldots,x_\rho$ の任意の多項式が代数的に整である、という事実によって特徴づけられる。この仮定から、写像領域そのものが相対的に整な領域であることが従う。}\footnote{例えば、\S\ 9 の末尾の注意で述べた領域 $\Gdom_{32}$ では、許容される特殊化 $x_1=1$ によって二つの多項式 $F^*=x_2$, $G^*=1+x_3$ が得られ、これらの上で $x_2,x_3$ の任意の多項式が代数的に整である。従って $\Gdom_{32}$ は第一種であり、実際 $F(x)$ と $G(x)$ を整基底にもつ。}

\section*{\S\ 11. 代数的整従属性に関する補助定理.}

いま、$x_1,\ldots,x_\rho$ の任意の多項式が、$x_1,\ldots,x_\rho$ の $\rho$ 個の多項式
\begin{equation}
  G_1^*(x),\ldots,G_\rho^*(x)
\tag{1}
\end{equation}
上代数的に整であるための判定条件を与える。

特に (1) の多項式が同次形式であるならば、(1) を消す任意の特殊な値系において、代数的に整従属なすべての形式も同時に消える。従って特に $x_1,\ldots,x_\rho$ も消える。

従って、形式 (1) は
\[
  x_1=0,
  \ldots,
  x_\rho=0
\]
以外の値系で同時に消えることはできない。それらの終結式は零であってはならない。

逆に、この条件も十分であり、次の形で非同次多項式へ拡張できる。

\emph{補助定理: $x_1,\ldots,x_\rho$ の任意の多項式が、$x_1,\ldots,x_\rho$ の $\rho$ 個の多項式}
\[
  \tag{1} G_1^*(x),\ldots,G_\rho^*(x)
\]
\emph{上代数的に整であるための十分条件は、これらの多項式 (1) の最高次元項}
\[
  H_1^*(x),\ldots,H_\rho^*(x)
\]
\emph{の、$x_1,\ldots,x_\rho$ に関する終結式が消えないことである。}
\begin{equation}
  R_0=\left[
  \begin{array}{ccc}
    H_1^*&\cdots&H_\rho^*\\
    x_1&\cdots&x_\rho
  \end{array}
  \right]\ne0.
\tag{2}
\end{equation}
\emph{同次形式 (1) に対しては、条件 (2) は必要である。}\footnote{終結式論については、F. Mertens, Zur Theorie der Elimination (I. and II. Teil), Sitzungsber. d. k. Ak. d. Wiss. Wien, Math.-naturw. Kl. Abt. IIa, Bd. 108 (1899) を参照。J. König, Theorie d. algebr. Größen, Kap. VI に部分的に再掲されている。}

証明のために、多項式
\begin{equation}
\begin{gathered}
  \bar G_1^*(x)=G_1^*(x)-\lambda_1,
  \ldots,
  \bar G_\rho^*(x)=G_\rho^*(x)-\lambda_\rho,\\
  \bar\Phi^*(x)=\Phi^*(x)-\mu
\end{gathered}
\tag{3}
\end{equation}
を作る。ここで $\lambda_1,\ldots,\lambda_\rho,\mu$ は不定元、$\Phi^*(x)$ は $x_1,\ldots,x_\rho$ の任意の多項式である。多項式 (3) の $x_1,\ldots,x_\rho,1$ に関する終結式\footnote{すなわち、多項式 (3) を追加の不定元 $x_0$ によって、$x$ に関する次数が保たれるように同次化したときの、$x_1,\ldots,x_\rho,x_0$ に関する同次終結式である。}を
\begin{equation}
  R(\lambda,\mu)=\left[
  \begin{array}{cccc}
   \bar G_1^*&\cdots&\bar G_\rho^*&\bar\Phi^*\\
   x_1&\cdots&x_\rho&1
  \end{array}
  \right]
\tag{4}
\end{equation}
と書く。これは $\lambda_1,\ldots,\lambda_\rho,\mu$ の整有理関数である。F. Mertens によれば、(4) を $\mu$ の冪で展開したとき、最高係数を明示的に指定できる。\footnote{同書 \S\ 3（定理 I の証明）および特に \S\ 6（終結式 $R$ の項のうち、与えられた冪積 $a_{\lambda\lambda}^{p_\lambda}a_{\mu\mu}^{p_\mu}\cdots a_{\sigma\sigma}^{p_\sigma}$ を因子として含むもの全体の決定）。本場合には $p_\lambda=0$, $p_\mu=0,\ldots,p_\sigma=\tau$ とし、$a_{\sigma\sigma}=(-\mu)$ と置く。König, Kap. VI, \S\ 9 に再掲。}すなわち
\begin{equation}
  (-1)^\tau\cdot R(\lambda,\mu)
  =R_0^\sigma\cdot\mu^\tau+A_1(\lambda)\mu^{\tau-1}+\cdots+A_\tau(\lambda),
\tag{5}
\end{equation}
を得る。ここで $A_1(\lambda),\ldots,A_\tau(\lambda)$ は、$\lambda_1,\ldots,\lambda_\rho$ および多項式 (3) の係数の、整数係数の整関数を表す。この展開はまず、仮定 (2)、すなわち $R_0\ne0$ のもとでは、$R(\lambda,\mu)$ も恒等的に消えないことを示す。$x_1,\ldots,x_\rho,\lambda_1,\ldots,\lambda_\rho,\mu$ に関する恒等式
\[
  R(\lambda,
\mu)\equiv 0\pmod{\bar G_1^*(x),\ldots,
\bar G_\rho^*(x),\bar\Phi^*(x)}
\]
は、代入
\[
  \lambda_i=G_i^*(x),\qquad \mu=\Phi^*(x)
\]
によって
\begin{equation}
\begin{aligned}
  R(G_i^*(x),\Phi^*(x))
  &=R_0^\sigma\cdot\Phi^{*\tau}(x)
    +A_1(G_i^*(x))\Phi^{*\tau-1}(x)+\cdots \\
  &\qquad +A_\tau(G_i(x))=0
\end{aligned}
\tag{6}
\end{equation}
となる。仮定により $R_0$ は $\Omega$ の数であるから、これで補助定理が証明される。

\section*{\S\ 12. 多項式からなる正則系 $\Ssys_{n\rho}$ の整基底.}

\S\ 11 の補助定理を \S\ 10 の結論と合わせると、直ちに次の事実が得られる。

\emph{$\Gdom_{n\rho}$ の写像領域 $\Jdom^*_{\rho\rho}$ に、最高次元項の同次終結式が恒等的に消えない、すなわち $R_0\ne0$ となる $\rho$ 個の多項式が存在するならば、$\Gdom_{n\rho}$ は第一種であり、従ってインボリューション基底を整基底とする有限整域である。}

条件 $R_0\ne0$ は、補助定理によって、整基底の存在の第二の証明を可能にする。この証明は本質的には D. Hilbert によるものである。\footnote{Über die vollen Invariantensysteme, \S\ 2. Hilbert の場合には、不変式体の整基底の存在が別の方法で得られていることを仮定して、その基底を Kronecker の体の代数的に整な量の基本系として解釈することだけが問題となる。} ただしこの証明では、インボリューション基底をそれとして認識することはできない。しかし、示した性質をもつすべての系 $\Ssys_{n\rho}$ に一様に適用される。

実際、$R_0\ne0$ となり、従って代数的に独立である $\Jdom^*_{\rho\rho}$ の $\rho$ 個の多項式を
\begin{equation}
  G_1^*(x),\ldots,G_\rho^*(x)
\tag{1}
\end{equation}
とする。補助定理により、$x_1,\ldots,x_\rho$ の任意の多項式、従って特に $\Ssys^*_{\rho\rho}$ を含む最小の体 $\Kfield^*_{\rho\rho}$ の任意の多項式は、(1) の多項式上代数的に整である。$\Ssys_{n\rho}$ における対応する多項式を
\begin{equation}
  G_1(x),\ldots,G_\rho(x)
\tag{2}
\end{equation}
とすれば、転移原理により、$\Ssys_{n\rho}$ を含む最小の体 $\Kfield_{n\rho}$ の任意の多項式、従って特に $\Ssys_{n\rho}$ の任意の多項式は、(2) の多項式上代数的に整である。逆に、(2) の多項式上代数的に整である $\Kfield_{n\rho}$ の任意の関数は多項式である。したがって、$\Kfield_{n\rho}$ を（\S\ 4 に従って）$\Omega(G_1,\ldots,G_\rho)$ 上の代数体と見ると、$\Kfield_{n\rho}$ の多項式全体、すなわち相対的に整な領域 $\Gdom_{n\rho}$ は、この体の代数的に整な量全体と同一である。$G(x)$ を (2) を有理基底へ補う $\Kfield_{n\rho}$ の多項式とし、$D(x)$ を $G(x)$ が $\Omega(G_1,\ldots,G_\rho)$ に関して満たす次数 $k$ の既約方程式の判別式とする。すると $\Kfield_{n\rho}$ の任意の多項式、特に $\Ssys_{n\rho}$ の任意の多項式 $F(x)$ は、代数的に整な量として表示
\begin{equation}
  F(x)=\frac{\Gamma_1(G_i)G(x)^{k-1}+\Gamma_2(G_i)G(x)^{k-2}+\cdots+\Gamma_k(G_i)}{D(x)}
\tag{3}
\end{equation}
をもつ。

いま $F(x)$ に $\Ssys_{n\rho}$ のすべての多項式を走らせる。すると、対応する関数
\[
  v_1\Gamma_1(G_i)+v_2\Gamma_2(G_i)+\cdots+v_k\Gamma_k(G_i)
\]
からなる系に Hilbert の加群基底に関する定理 I（Math. Ann. 36）を適用することで、$\Ssys_{n\rho}$ の有限個の多項式
\begin{equation}
  F_1(x),F_2(x),\ldots,F_\lambda(x)
\tag{4}
\end{equation}
が存在し、$\Ssys_{n\rho}$ の任意の多項式が
\begin{equation}
  F(x)=A_1(G_i)F_1(x)+A_2(G_i)F_2(x)+\cdots+A_\lambda(G_i)F_\lambda(x)
\tag{5}
\end{equation}
の形に表されることが分かる。ここで $A_1,\ldots,A_\lambda$ は $G_i(x)$ の多項式である。従って、多項式 (2) と (4) は $\Ssys_{n\rho}$ の整基底をなす。すなわち、

\emph{定理 VIII: 多項式系 $\Ssys_{n\rho}$ の写像系 $\Jdom^*_{\rho\rho}$ に、最高次元項の終結式が恒等的に消えない、すなわち $R_0\ne0$ となる $\rho$ 個の多項式が存在するならば、$\Ssys_{n\rho}$ は整基底をもつ。}

この定理は少し一般化できる。$R_0\ne0$ となる $\rho$ 個の多項式 (1) が、$\Ssys^*_{\rho\rho}$ を含む最小の整域 $\Jdom^*_{\rho\rho}$ に属していれば十分である。実際、定理 VIII に至る議論は、その場合にも表示 (5) が成り立つことを示す。ところで定義により、最小の包含整域 $\Jdom_{n\rho}$ の任意の多項式 $L(x)$ は、$L(x)$ に依存する有限個の $\Ssys_{n\rho}$ の多項式によって整有理的に表される。従って、$\Ssys_{n\rho}$ の $\sigma$ 個の多項式
\begin{equation}
  K_1(x),K_2(x),\ldots,K_\sigma(x)
\tag{6}
\end{equation}
が存在し、多項式 $G_i(x)$ は (6) の整有理結合となる。よって (6) と (4) が整基底を与える。次の用語を導入する。

\emph{定義 VIII: 多項式の系 $\Ssys_{n\rho}$ が正則であるとは、最小包含整域の写像領域 $\Ssys^*_{\rho\rho}$ に、その最高次元項の終結式が恒等的に消えない、すなわち $R_0\ne0$ となる $\rho$ 個の多項式が存在することをいう。}

従って次を得る。

\emph{定理 IX: 多項式からなる任意の正則系 $\Ssys_{n\rho}$ は整基底をもつ。}\footnote{整基底をもたない非正則系が存在することは、Hilbert が例によって示した（Gött. Nachr. 1891, p. 233）。次の例はさらに簡単である。
\[
  \Ssys_{22}=x_1,
  x_1x_2,
  x_1x_2^2,
  \ldots,
  x_1x_2^\nu,
  \ldots\quad\text{in inf.}
\]
ここで整基底が存在することは、正整数 $\nu$ が存在して恒等式
\[
  x_1x_2^{\nu+\sigma}=\sum c\cdot x_1^{i_0}(x_1x_2)^{i_1}(x_1x_2)^{i_2}\cdots(x_1x_2^\nu)^{i_\nu}
  \quad(\sigma=1,2,\ldots)
\]
が成り立つことと同値である。すでに $\sigma=1$ について $x_1,x_2$ の次元を比較すると、指数に関する二つの条件式
\[
  1=i_0+i_1+i_2+\cdots+i_\nu,
  \qquad
  \nu+1=i_1+2i_2+\cdots+\nu\cdot i_\nu
\]
が得られるが、これは明らかに非負整数解をもたない。従って $\Ssys_{22}$ は整基底をもたない。}

さらに、\S\ 5 のインボリューションの場合と同様に、正則系の幾何学的解釈を与える。そのため、そこでの (7) と同じく、方程式系
\begin{equation}
  L^*(x)-L^*(\xi)=0
\tag{7}
\end{equation}
を考える。ここで $L^*(x)$ は $\Jdom^*_{\rho\rho}$ のすべての多項式を走る。$x_0,\xi_0$ を用いて同次化すると、(7) は同次形式の間の方程式系
\begin{equation}
  \xi_0^{m\cdot}M^*(x)-x_0^{m\cdot}M^*(\xi)=0
\tag{8}
\end{equation}
へ移る。ここで $H^*(x)$ が $L^*(x)$ の最高次元項を表すならば、関係
\begin{equation}
  H^*(x)\equiv M^*(x)\pmod{x_0}
\tag{9}
\end{equation}
が成り立つ。

$\Jdom^*_{\rho\rho}$ の\emph{基本点}（\S\ 5 の注を参照）とは、(8) の解で $x$ に依存しないもの、すなわち $x$ に共役な行とは異なるものをいう。したがって基本点とは、$\xi_1=0,\ldots,\xi_\rho=0$ とは異なる値系で、同時に
\[
  \xi_0^m=0,
  \qquad M^*(\xi)=0
\]
を満たすもの、また (9) によれば
\[
  \xi_0=0,
  \qquad H^*(\xi)=0
\]
を満たすものである。

いま $R_0\ne0$ となる、従って共通解系をもたない $\rho$ 個の形式 $H^*(\xi)$ が存在するならば、$\Jdom^*_{\rho\rho}$ は基本点をもてない。特に、$\Ssys^*_{\rho\rho}$ が同次形式から成る場合には、$R_0\ne0$ ならば $\Ssys^*_{\rho\rho}$ も基本点をもてない。そのような点は同時に $\Jdom^*_{\rho\rho}$ の基本点だからである。しかし逆も成り立つ。そのため、形式 $H^*(x)$ 全体からなる系 $\Jdom(H^*)$ を考える。これは再び整域である。

Hilbert の加群基底に関する定理 I と Hilbert の補助定理（\S\ 10 の引用を参照）から、$\Jdom(H^*)$ の $\rho$ 個の形式が存在し、それらの消滅が $\Jdom(H^*)$ のすべての形式の消滅を帰結することが従う。もし $\Jdom^*_{\rho\rho}$ が基本点をもたなければ、これら $\rho$ 個の形式は同時に消えることができず、その終結式は $R_0\ne0$ となる。

従って、

\emph{多項式の正則系 $\Ssys_{n\rho}$ は、最小包含整域の写像領域 $\Jdom^*_{\rho\rho}$ が基本点をもたないことによって特徴づけられる。同次形式の正則系 $\Ssys^*_{\rho\rho}$ については、系そのものがすでに基本点をもたない。}\footnote{前の注で与えた非正則系 $\Ssys_{22}$ は基本点
\[
  \xi_1:\xi_2:\xi_0=0:1:0
\]
をもつ。}

\section*{\S\ 13. 第一種の相対的に整な領域および正則系の例.}

1. 第一種の相対的に整な領域 $\Gdom_{nn}$ の第一の例として、\emph{$x_1,\ldots,x_n$ の多項式のうち、与えられた置換群 $G$ を許すもの全体}、すなわち Lagrange の種領域の多項式全体を考える。恒等式
\begin{equation}
  x_k^n-\sigma_1(x)x_k^{n-1}+\sigma_2(x)x_k^{n-2}+\cdots\pm\sigma_n(x)=0
  \qquad(k=1,2,
\ldots,n)
\tag{1}
\end{equation}
により、$x_1,\ldots,x_n$ は $\Gdom_{nn}$ に属する基本対称関数上代数的に整であり、従って定義 V により $\Gdom_{nn}$ 上代数的に整である。したがって $\Gdom_{nn}$ は第一種である。さらに、$\Gdom_{nn}$ は第一種の領域であり、その元が同次形式で、かつ $n=\rho$ であるから、\S\ 10 の結論により正則系をなす。実際、(1) によって、基本対称関数の終結式 $R_0$ は恒等的に消えない。終結式論の計算法則により、$R_0$ は $\pm1$ と計算される。$\Gdom_{nn}$ の\emph{インボリューション形式}は、$G_{nn}$ が第一種で $n=\rho$ であるため、対応する Lagrange の種領域のインボリューション形式として一意に定まる。この領域に関して $x_1,\ldots,x_n$ に共役な量は群 $G$ の置換によって生じるから、このインボリューション形式は
\begin{equation}
\begin{aligned}
  \Psi(z,u)&=\prod_G
  (z-u_1x_{i_1}-u_2x_{i_2}-\cdots-u_nx_{i_n})\\
  &=z^i+\sum{}'G_{\alpha\alpha_1\ldots\alpha_n}(x)
  z^\alpha u_1^{\alpha_1}\cdots u_n^{\alpha_n}
\end{aligned}
\tag{2}
\end{equation}
で与えられる。積は $G$ のすべての置換にわたって取る。式 (2) は「群 $G$ のガロア形式」を表す。従って次の事実が成り立つ。

\emph{置換群 $G$ を許す $x_1,\ldots,x_n$ のすべての多項式は、群のガロア形式の係数 $G_{\alpha\alpha_1\ldots\alpha_n}(x)$ の整有理結合である。}

2. \emph{正則系}の例として、多項式の系 $\Ssys_{n1}$、すなわち代数的に独立な多項式を一つだけ含む系を考える。

実際、$\Ssys^*_{11}$ を $\Ssys_{n1}$ の写像系とし、
\[
  F^*(x)=a_0x^k+a_1x^{k-1}+\cdots+a_k\qquad(a_0\ne0)
\]
を $\Ssys^*_{11}$ の多項式とする。このとき
\begin{equation}
  R_0=a_0\ne0.\footnote{Mertens, 前掲書 §5。}
\tag{3}
\end{equation}
従って $\Ssys_{n1}$ は定義 VIII により正則であり、したがって、

\emph{多項式の任意の系 $\Ssys_{n1}$、特に一つの不定元の無限個の多項式の任意の系は、整基底をもつ。}\footnote{従って、整基底をもたない最も簡単な非正則系は $\Ssys_{22}$ である。そのような系が実際に存在することは、Hilbert の例および \S\ 12 の注で与えた例が示している。}

3. いま系 $\Ssys_{n1}$ を相対的に整な領域 $\Gdom_{n1}$ に特殊化する。これは \S\ 12 により正則系であるから、第一種の相対的に整な領域である。従ってその整基底はインボリューション基底で与えられ、この基底は同時に最小包含体 $\Kfield_{n1}$ のインボリューション基底でもある。しかし \S\ 6 により、$\Kfield_{n1}$ のインボリューション基底の任意の関数は同時に最小基底である。これを体の Lüroth 関数と呼んだ。そしてインボリューション基底の任意の関数は、この Lüroth 関数に一次分数変換によって依存する。現在の場合、Lüroth 関数は $\Gdom_{n1}$ に属するので多項式である。従って、インボリューション基底の他のすべての多項式は、この Lüroth 関数 $L(x)$ に一次分数的に依存し、かつ (3) によって代数的に整であるので、結局整かつ線形に依存する。したがって $L(x)$ が整基底となる。$\Gdom^*_{11}$ のインボリューション形式は、$(u_1=1)$ と置けば
\[
  \Psi^*(z)=L^*(z)-L^*(x)
\]
であり、ここからも $L(x)$ が整基底として現れる。よって、

\emph{相対的に整な領域 $\Gdom_{n1}$ は第一種である。その整基底は一つの多項式、すなわち最小包含体の Lüroth 関数からなる。}

4. 相対的に整な領域 $\Gdom_{n1}$ の例として、$s$ 変数の二次形式の整有理な射影不変式を特に挙げておく。前述と一致して、それらがその形式の行列式 $|a_{ik}|$ の整有理結合として表されることはよく知られている。そしてこの行列式は、対応する不変式体の Lüroth 関数でもある。

\section*{\S\ 14. 整数多項式からなる系.}

得られた基底定理を、整数上の整性に関して精密化する。

そのため以下では、係数領域 $\Omega$ を（有限）代数的数体と仮定し、$[\Omega]$ は $\Omega$ の代数的整数全体を表すものとする。これを簡単に整数と呼ぶ。整数多項式とは、$[\Omega]$ からの係数をもつ多項式をいう。以下では $F(x),G(x),\ldots$ は整数多項式だけを表す。

\S\ 7 と同様に、整数多項式から成るさまざまな系を考える。

整数多項式の任意の系を整数系 $\Ssys_{n\rho}$ と呼ぶ。

整数線形族 $\Lsys_{n\rho}$ とは、$F_1(x)$ と $F_2(x)$ とともに、$[\Omega]$ の任意の整数 $c_1,c_2$ に対して $c_1F_1(x)+c_2F_2(x)$ も常に含む整数系である。

整数整域 $\Jdom_{n\rho}$ とは、$F(x)$ と $G(x)$ とともに $F(x)\cdot G(x)$ も常に含む整数線形族である。

整数相対整領域 $\Gdom_{n\rho}$ とは、$F(x)$ と $G(x)$ とともに、ただし $G(x)\ne0$ のとき、商 $F(x):G(x)$ を、それが整数多項式である場合に、かつその場合に限って常に含む整域である。

\S\ 7 と同様に、整数系 $\Ssys_{n\rho}$ を含む最小の整数整域 $\Jdom_{n\rho}$ は、$\Ssys_{n\rho}$ の有限個の多項式の整数係数の整有理結合であるすべての多項式 $H(x)$ からなる。

最小の整数相対整領域 $\Gdom_{n\rho}$ は、$\Ssys_{n\rho}$ の有限個の多項式の有理結合である整数多項式 $H(x)$ 全体からなる。

いま基底定理の転移へ進む。有理基底については、体、最小包含体、有理基底という「有理性」だけに基づく概念は何ら変化しないことに注意する。さらに、整数系 $\Ssys_{n\rho}$ の写像系 $\Ssys^*_{\rho\rho}$ は、任意に多くの仕方で再び整数系として選ぶことができる。整数線形族についても、生成元数の制限に導く議論は有効であるから、次を得る。

\emph{定理 V$'$: 任意の整数系 $\Ssys_{n\rho}$ は有理基底をもつ。整数線形族 $\Lsys_{n\rho}$ の有理基底は、特に高々 $\rho+1$ 個の多項式を含むように選ぶことができる。}

つぎに、多項式の整域のインボリューション基底（\S\ 8）の類似を調べる。そのため、\emph{量の行の対称関数}に関する定理を、\emph{整数上の整性}に関して次の補助定理によって拡張しなければならない。

\emph{補助定理. $n$ 行の量 $(yz\cdots t)$}
\begin{equation}
\begin{array}{cccc}
  y_1&z_1&\cdots&t_1\\
  y_2&z_2&\cdots&t_2\\
  \vdots&\vdots&&\vdots\\
  y_n&z_n&\cdots&t_n
\end{array}
\tag{1}
\end{equation}
\emph{の整数係数の整対称関数は、それらのうち有限個、すなわち基本対称関数}
\begin{equation}
\begin{gathered}
  \sigma_1(y),\sigma_2(y),\ldots,\sigma_n(y);\quad
  \sigma_1(z),\sigma_2(z),\ldots,\sigma_n(z);\quad\ldots;\\
  \sigma_1(t),\ldots,\sigma_n(t);
\end{gathered}
\tag{2}
\end{equation}
\emph{および関数}
\begin{equation}
  \chi_1(y,z,\ldots,t),\ldots,\chi_k(y,z,\ldots,t)
\tag{3}
\end{equation}
\emph{の整数係数の整関数である。ただし (3) の関数は (2) の関数上で代数的に整かつ整数的である。}

実際、$n$ 個の量 $y_1,\ldots,y_n$ は $\sigma_1(y),\ldots,\sigma_n(y)$ 上で代数的に整かつ整数上整であり、同じ主張が $z_1,\ldots,z_n$ と $\sigma_1(z),\ldots,\sigma_n(z)$ について、また以後同様に成り立つ。従って、(1) の量の行の対称関数の体は、(2) の関数によって定まる体上の代数体をなし、この体の代数的に整かつ整数的な量は、(1) の量の行の整数係数の整対称関数と一致する。従って (3) を Kronecker の基本系として取れば、補助定理が証明される。

いま $\Jdom_{n\rho}$ を整数整域、$\Jdom^*_{\rho\rho}$ を整数写像領域とし、
\[
  \Psi^*(z,u)=G^*(x)\cdot z^i+
  \sum{}'G^*_{\alpha\alpha_1\ldots\alpha_\rho}(x)\cdot
  z^\alpha u_1^{\alpha_1}\cdots u_\rho^{\alpha_\rho}
\]
を $\Jdom^*_{\rho\rho}$ のインボリューション形式とする。ここで $G^*(x)$ は零次元、すなわち $[\Omega]$ の数であってもよい。関数 (2) に対応する基本対称関数は、このとき商
\begin{equation}
  G^*_{\alpha\alpha_1\ldots\alpha_\rho}(x):G^*(x)
\tag{4}
\end{equation}
のうちのあるものによって与えられる。

補助定理により、関数 (3) は (2) 上で代数的に整かつ整数上整に依存するから、(3) に対応する $\Kfield^*_{\rho\rho}$ の関数は、代数的整数係数をもつ多項式に対する Gauss の定理の拡張により、
\begin{equation}
  K_\lambda^*(x):(G^*(x))^\sigma
\tag{5}
\end{equation}
で与えられる。ここで $K_\lambda^*(x)$ は $\Kfield^*_{\rho\rho}$ の整数多項式である。したがって、相対整領域の場合には $K_\lambda^*(x)$ は領域に属し、一般には $\Jdom^*_{\rho\rho}$ の多項式の商である。

いま $F^*(x)$ を $\Jdom^*_{\rho\rho}$ の任意の整数多項式とすると、
\[
  F^*(x)=\frac1i\Bigl(F^*(x)+F^*(x^{(1)})+\cdots+F^*(x^{(i-1)})\Bigr)
\]
であるから、$i\cdot F^*(x)$ は量の行 $x,x^{(1)},\ldots,x^{(i-1)}$ の整数係数の整対称関数となる。

従って、補助定理と転移原理により、次を得る。

\emph{定理 VI$'$: 整数整域 $\Jdom_{n\rho}$ については、各インボリューション形式が特別な有理基底を定め、$\Jdom_{n\rho}$ の任意の多項式 $F(x)$ は表示}
\[
  F(x)=\frac{\Gamma(H(x),H_1(x),\ldots,H_\sigma(x))}{i\cdot H(x)^k}
\]
\emph{をもつ。ここで $\Gamma$ は整数係数の整関数を表す。整数相対整領域 $\Gdom_{n\rho}$ で、写像領域が整数相対整領域 $\Gdom^*_{\rho\rho}$ である場合には、分母 $H(x)$ は対応するインボリューション形式の最高係数によって与えられる。}

\section*{\S\ 15. 第一種の整数相対整領域と整数正則系.}

整数系については、整基底の代わりに整数整基底が用いられる。すなわち、系の有限個の整数多項式であって、系の任意の整数多項式がそれら有限個の整数係数の整結合として表されるものをいう。

整数整基底に関する定理は、本質的に先の定理と平行であり、証明だけがわずかに変更を要する。

まず、定義 V（\S\ 9）は直接に移される。

\emph{定義 V$'$: 量 $\xi$ が整数相対整領域 $\Gdom_{n\rho}$ に関して代数的に整かつ整数的であるとは、最高係数が $[\Omega]$ の単元で、残りの多項式が $\Gdom_{n\rho}$ に属するような方程式を少なくとも一つ満たすことをいう。}

代数的整数係数をもつ多項式に対する Gauss の定理の拡張により、\S\ 9 と同様、この定義は既約方程式から得られる。

定義 V$'$ から、さらに定義 VII の転移が得られる。

\emph{定義 VII$'$: 整数相対整領域 $\Gdom_{n\rho}$ が第一種であるとは、その写像領域として相対整領域 $\Gdom^*_{\rho\rho}$ をもち、$x_1,\ldots,x_\rho$ の任意の整数多項式が $\Gdom^*_{\rho\rho}$ 上で代数的に整かつ整数的であることをいう。}

これら第一種の領域の有限性を示すため、まず \S\ 10 と同じく、定義からインボリューション形式の最高係数が $[\Omega]$ の単元であることに注意する。

定理 VI$'$（\S\ 14）により、$\Gdom_{n\rho}$ の任意の多項式 $F(x)$ はしたがって表示
\[
  F(x)=\frac{\Gamma(H_1(x),\ldots,H_\sigma(x))}{i}
\]
をもつ。ここで、整数多項式 $\Gamma(H_1,\ldots,H_\sigma)$ からなる系に、Hilbert の整数加群基底に関する定理 II（Math. Ann. 36）を、もともとは整数有理数係数について述べられていたものを、代数的整数多項式への拡張として適用する。\footnote{$w_1,\ldots,w_k$ が $\Omega$ の代数的整数の基本系を表すとし、
\[
  \Gamma(H_1,\ldots,H_\sigma)=\Gamma_1(H)w_1+\cdots+\Gamma_k(H)w_k
\]
と置くならば、$\Gamma_1(H),\ldots,\Gamma_k(H)$ は整数有理数係数をもつ。形式 $v_1\Gamma_1(H)+\cdots+v_k\Gamma_k(H)$ からなる系に定理 II を適用すれば、この拡張の妥当性が直ちに示される。} これにより整数整基底の存在が従う。すなわち、

\emph{定理 VII$'$: 第一種の任意の整数相対整領域は整数整基底をもつ。}

さらに正則系の定義を移す。

\emph{定義 VIII$'$: 整数系 $\Ssys_{n\rho}$ が整数正則であるとは、最小包含整数整域の写像領域 $\Jdom^*_{\rho\rho}$ に、最高次元項の終結式が $[\Omega]$ の単元である、すなわち $R_0=\eps$ となる $\rho$ 個の多項式が存在することをいう。}

整基底の存在証明を移すため、\S\ 11 の補助定理が次のより鋭い形を許すことに注意する。この形は、\S\ 11 の式 (5), (6) と、$A_1(\lambda),\ldots,A_\tau(\lambda)$ が $\lambda_1,\ldots,\lambda_\rho$ の整数係数の整関数であることから従う。

\emph{$x_1,\ldots,x_\rho$ の任意の整数多項式が、$\rho$ 個の整数多項式上で代数的に整かつ整数的に依存するための十分条件は、それらの多項式の最高次元項の終結式が $[\Omega]$ の単元、すなわち $R_0=\eps$ であることである。}

この補助定理から、\S\ 12 とまったく同じように、整数正則系の任意の多項式 $F(x)$ について \S\ 12 の表示 (3) が従う。ただし今は $\Gamma_1(G_i),\ldots,\Gamma_k(G_i)$ は $G_i$ の整数多項式である。さらに \S\ 12 と同様に、Hilbert の定理 II（Math. Ann. 36）を代数的整数多項式への拡張として適用し、最小包含整数整域の定義を用いれば、整数整基底の存在が得られる。すなわち、

\emph{定理 IX$'$: 任意の整数正則系は整数整基底をもつ。}

第一種の整数相対整領域の例として、\S\ 13 の例 1 と同様に、Lagrange の種領域の整数多項式全体 $\Gdom_{nn}$ を挙げる。$\Gdom_{nn}$ が整数の第一種であることは恒等式
\[
  x_k^n-\sigma_1(x)x_k^{n-1}+\cdots\pm\sigma_n(x)=0
  \qquad(k=1,2,
\cdots,n);
\]
から従う。また基本対称関数の終結式は値 $\pm1$ をもつので、$\Gdom_{nn}$ は整数正則系でもある。

さらにもう一つの例は、\S\ 14 の補助定理によって、$n$ 行の量の整数係数の整対称関数全体によって与えられる。

Erlangen, 1914年5月.

\section*{7. 有限群の不変式に対する有限性定理}
\begin{center}
Math. Ann. 77 (1916), pp. 89--92

Emmy Noether（Erlangen）
\end{center}

以下では、有限群の不変式について、完全に初等的な有限性の証明を与える。この証明は対称関数の理論だけに基づくものであり、同時に完全系を実際に指定する。他方、Hilbert の加群基底定理（Ann. 36）に基づく通常の証明は、存在証明にすぎない。\footnote{例えば Weber, \emph{Lehrbuch der Algebra} (2nd ed.), vol. 2, \S\ 57 を参照。}

有限群 $\Hgrp$ は、行列式が 0 でない $h$ 個の線形変換 $A_1\cdots A_h$ から成るものとする。ここで $A_k$ は線形変換
\[
  x_1^{(k)}=\sum_{\nu=1}^{n}a_{1\nu}^{(k)}x_\nu,\cdots,
  x_n^{(k)}=\sum_{\nu=1}^{n}a_{n\nu}^{(k)}x_\nu,
\]
すなわち略記して $(x^{(k)})=A_k(x)$ を表す。したがって群 $\Hgrp$ は、成分 $x_1,
\ldots,x_n$ をもつ行 $(x)$ を、成分 $x_1^{(k)},
\ldots,x_n^{(k)}$ をもつ行 $(x^{(k)})$ へ移す。恒等変換は $A_1,
\ldots,A_h$ の中に含まれていなければならないから、行 $(x)$ 自身も行 $(x^{(k)})$ の中に含まれる。群の整有理（絶対）不変式とは、$x_1,
\ldots,x_n$ の整有理関数であって、$A_1,
\ldots,A_h$ の適用のもとで恒等的に変化しないものをいう。このような不変式 $f(x)$ に対しては、したがって
\begin{equation}
  f(x)=f(x^{(1)})=\cdots=f(x^{(h)})=
  \frac1h\sum_{k=1}^{h}f(x^{(k)}).
\tag{1}
\end{equation}
が成り立つ。

\subsection*{1.}
式 (1) は、$f(x)$ が数量行 $(x^{(1)}),
\ldots,(x^{(h)})$ の整有理対称関数であることをいう。しかも各項 $f(x^{(k)})$ はただ一つの行 $(x^{(k)})$ だけを含むから、これは通常\emph{一形式}の場合と呼ばれる最も単純な場合である。数量行の対称関数に関する既知の定理\footnote{2. への注を参照。}によって、$f(x)$ は、これらの行の基本対称関数、すなわち「Galois resolvent」
\[
\begin{aligned}
  \Phi(z,u)
  &=\prod_{k=1}^{h}\bigl(z+u_1x_1^{(k)}+\cdots+u_nx_n^{(k)}\bigr) \\
  &=z^h+
    \sum G_{\alpha\alpha_1\cdots\alpha_n}(x)
    z^{\alpha}u_1^{\alpha_1}\cdots u_n^{\alpha_n},
\end{aligned}
\qquad
\left(\begin{array}{c}
\alpha+\alpha_1+
\cdots+
\alpha_n=h\\
\alpha\ne h
\end{array}\right),
\]
の係数 $G_{\alpha\alpha_1\cdots\alpha_n}(x)$ によって整有理的に表される。ここで $G_{\alpha\alpha_1\ldots\alpha_n}(x)$ は、$x$ について次数 $\alpha_1+\cdots+\alpha_n$ の不変式である。したがって次が証明された。

\emph{Galois resolvent の係数 $G_{\alpha\alpha_1\cdots\alpha_n}(x)$ は群の不変式の完全系をなし、すべての不変式はこの有限個の不変式によって整有理的に表される。}

\subsection*{2.}
(1) から出発して、次の、なおいっそう初等的な考察を行うこともできる。\footnote{これは、上で述べた一形式の場合における数量行の対称関数の定理の証明に相当する。}同時にこれは第二の完全系へ導く。置く：
\[
  f(x)=a+b x_1^{\mu_1}\cdots x_n^{\mu_n}
       +\cdots+c x_1^{\nu_1}\cdots x_n^{\nu_n},
\]
ここで $a,b,\cdots,c$ は定数を表す。すると (1) から
\[
\begin{aligned}
 h\cdot f(x)
 &=h\cdot a+b\cdot\sum_{k=1}^{h}x_1^{(k)^{\mu_1}}\cdots x_n^{(k)^{\mu_n}}
     +\cdots
     +c\cdot\sum_{k=1}^{h}x_1^{(k)^{\nu_1}}\cdots x_n^{(k)^{\nu_n}} \\
 &=h\cdot a+b\cdot J_{\mu_1\cdots\mu_n}+\cdots+c\cdot J_{\nu_1\cdots\nu_n}
\end{aligned}
\]
を得る。したがって、すべての不変式は特殊な不変式
\[
  J_{\mu_1\cdots\mu_n}
  =\sum_{k=1}^{h}x_1^{(k)^{\mu_1}}\cdots x_n^{(k)^{\mu_n}}
\]
から整かつ線形に構成できるので、これら特殊な不変式について主張を証明すれば足りる。ところが、数値因子を別にすれば、$J_{\mu_1\cdots\mu_n}$ は式
\[
  S_\mu=\sum_{k=1}^{h}
  \bigl(u_1x_1^{(k)}+\cdots+u_nx_n^{(k)}\bigr)^{\mu}
\]
における $u_1^{\mu_1}
\cdots u_n^{\mu_n}$ の係数である。ただし $\mu=\mu_1+\cdots+\mu_n$ であり、この式は $h$ 個の線形形式
\[
  \xi_1=u_1x_1^{(1)}+\cdots+u_nx_n^{(1)},\cdots,
  \xi_h=u_1x_1^{(h)}+\cdots+u_nx_n^{(h)}
\]
の $\mu$ 番目の冪和である。無限に多くの冪和 $S_\mu$ は
\[
  S_1,S_2,\cdots,S_h
\]
の整有理結合であり、その係数は不変式
\[
  J_{\mu_1\cdots\mu_n}\qquad(\mu_1+\cdots+\mu_n\leqq h)
\]
によって与えられることが知られている。こうして第二の完全系が得られる。

\emph{群の不変式の完全系は、$\mu_1+\cdots+\mu_n\le h$ を満たすすべての不変式 $J_{\mu_1\cdots\mu_n}$ によって与えられる。ただし $h$ は群の位数を表す。}

1. との関係は、冪和 $S_1,
\ldots,S_h$ を $\xi_1,
\ldots,
\xi_h$ の基本対称関数で置き換えられることに注意すれば得られる。これにより、そこで与えた $G_{\alpha\alpha_1\ldots\alpha_n}(x)$ の完全系に至る。両方の結果から、すべての不変式は、$x$ に関する次数が群の位数 $h$ を超えないものによって整有理的に表されることが分かる。

\subsection*{3.}
いま証明したことから、有理表示に関する帰結が引き出される。任意の有理絶対不変式は、分子および分母に、群に関する分母の共役を掛ければただちに分かるように、必ずしも互いに素とは限らない二つの整有理絶対不変式の商として表せる。したがって、任意の有理絶対不変式は Galois resolvent $\Phi(z,u)$ の係数 $G_{\alpha\alpha_1\cdots\alpha_n}(x)$ によって有理的に表される。この定理は、有限性定理を経由しなくてもいくつかの方法で容易に証明できる。これはすでに Weber II, \S\ 58 に述べられている。\footnote{ただし、そこで与えられた証明は正しくない。式 (7) の関数 $\Psi(t)$ が不変式を係数にもつことは示されているが、それらの不変式が $\Phi(t)$ の係数によって有理的に表されることは示されていない。この欠陥は、$x_i^{(k)}$ を $\xi_k$ によって表す際に、Lagrange の補間公式ではなく既知の微分手続きを用いれば避けられる。これは、恒等式 $\Phi(-\xi_k,u)=0$ をすべての $u$ について微分して得られる関係である：
\[
 \left[\frac{\partial\Phi}{\partial u_i}
      -x_i^{(k)}\cdot\frac{\partial\Phi}{\partial z}\right]_{z=-\xi_k}=0.
\]
したがって、各有理関数 $\omega(x)$ に対しては、Weber の式 (7), (8) の代わりに次の式を置かなければならない：
\[
 \omega(x_1^{(k)}\cdots x_n^{(k)})=
 \left.
 \omega\!\left(
   \frac{\frac{\partial\Phi}{\partial u_1}}{\frac{\partial\Phi}{\partial z}}\cdots,
   \frac{\frac{\partial\Phi}{\partial u_n}}{\frac{\partial\Phi}{\partial z}}
 \right)\right|_{z=-\xi_k},
\]
これは $\Phi(z,u)$ の係数だけを含み、不変式 $\omega(x)$ について全ての $k$ にわたり和を取れば、求める表示を与える。}

\subsection*{4.}
最後に、ここで得られた結果は、私の論文「有理関数の体と系」\footnote{Math. Ann. 76, p. 161 (1915).}に暗に含まれていることを述べておく。より正確には、1. で与えた完全系は定理 VI および VII に含まれており、この有限性定理に依存しない有理表示の証明は定理 III に含まれている。\footnote{相対不変式については、定理 VIII および IX に、通常のものとは異なる基礎にもとづく有限性の証明が含まれているが、これもまた存在証明にすぎない。その理由は、相対不変式が体をなさないからである。}しかし現在の特殊な場合には、一般理論に比べて証明の道筋も結果もかなり簡単になる。有限群の不変式へ一般的研究を適用することへ私を導いたのは E. Fischer との会話であり、実質的には、2. で与えた証明（一般の場合にはそこに現れる完全系は有理的に定義できない）および 3. への注も彼に由来する。

Erlangen, 1915 年 5 月。

\clearpage
\section*{8. 任意に多くの基本形式からなる系の不変式の整有理表示について}
\begin{center}
Math. Ann. 77 (1916), pp. 93--102
\end{center}

Schwarz 記念論文集への寄稿「任意に多くの基本形式からなる系の不変式について」の末尾で、D. Hilbert は次の予想を述べている。これらの不変式は、有限個の不変式からなる系 $(J,PJ)$ によって整有理的に表される、というものである。ここで $J$ は線形独立な基本形式の不変式の完全系を表し、不変式 $PJ$ は $J$ から極化過程によって生じる。

以下 I で、この予想が実際に正しいことを示す。それは、各々 $n$ 個の変数からなる任意に多くの行の形式は、固定された $n$ 行だけを含み、しかももとの形式から極化過程によって導かれる形式の極化の和として表される、という還元定理の単純な帰結である。この還元定理は、F. Mertens が初めて定式化した特殊な場合であり、A. Capelli、F. Mertens、J. Deruyts が与えた、Clebsch--Gordan の級数展開を各々 $n$ 変数からなる任意に多くの行の形式へ一般化したものの特殊な場合である。この一般化は、微分過程の形式的変換によって証明される。\footnote{F. Mertens: ``Über eine Formel der Determinantentheorie.'' (Formula 5), Sitzb. d. Ak. d. Wiss. Wien, vol. 91, II. Abt. (1885), さらに詳しく（応用を含めて）: ``Über ganze Funktionen von $m$ Systemen von je $n$ Unbestimmten''. Monatsh. f. Math. u. Phys. 4th year (1893). Capelli のいくつかの証明（1882 年以降）については、例えば Math. Ann. 37、または自筆版 ``Lezioni sulla teoria delle forme algebriche'' (1902) を参照。J. Deruyts: \emph{Essai d'une théorie générale des formes algébriques}, Mém. d. 1. Soc. Roy. des Sciences de Liège (2) 17 (1892).}

II では、還元定理を、形式の線形族の同値性に関する問題として扱うことにより、より概念的な証明も与える。この見方、また用いられる考察の本質的な一部は、「代数的加群系について」\footnote{E. Fischer: Über algebraische Modulsysteme und lineare homogene partielle Differentialgleichungen mit konstanten Koeffizienten, \S\ 1--4, J. f. M. 140 (1911).}という論文、およびそれに続く E. Fischer の未公刊の定理から取ったものであり、彼はその使用を好意的に許してくれた。

最後に III では、対応する考察が最も一般的な級数展開にも導くことを示す。

\subsection*{I.}
$\theta$ を、$N>n$ 個の各行
\[
  A_1,A_2,\ldots,A_N
\]
に関して斉次な形式とする。一般に行 $A_k$ は $n$ 個の成分
\[
  A_k^{(1)},A_k^{(2)},\ldots,A_k^{(n)}
\]
から成る。$\theta$ の係数は不定元でもよいし、与えられた有理性領域の量でもよい。さらに $P$ は、極化過程
\[
  P_{hk}=\left(A_h\frac{\partial}{\partial A_k}\right)
  =A_h^{(1)}\frac{\partial}{\partial A_k^{(1)}}+
   A_h^{(2)}\frac{\partial}{\partial A_k^{(2)}}+
   \cdots+
   A_h^{(n)}\frac{\partial}{\partial A_k^{(n)}}
\]
の、整数係数をもつ整有理関数を表すものとする。積 $P_{hk}P_{ij}\theta$ は、$P_{ij}\theta$ に $P_{hk}$ を適用することとして理解される。

いま述べた還元定理は、恒等式
\begin{equation}
  \theta=\sum PZ,
\tag{1}
\end{equation}
によって与えられる。ここで形式 $Z$ は行
\[
  A_1,A_2,\ldots,A_n
\]
だけを含み、$\theta$ から極化過程 $P$ によって導かれる。

次に、同じ位数かつ同じ変数の数をもつ $N$ 個の形式
\[
  F_1,F_2,\ldots,F_N
\]
から成る基本形式の系 $(F')$ を考える。その係数をそれぞれ $N$ 個の行
\[
  A_1,A_2,\ldots,A_N
\]
として取る。$F_1,
\ldots,F_n$ の、有限個の不変式から成る完全系を $(J)$ と書き、$F_1,
\ldots,F_n$ の任意の不変式が $(J)$ の不変式によって整有理的に表されるものとする。証明すべき Hilbert の予想は、$F_1,
\ldots,F_N$ の同時不変式 $S$ に対し、完全系が有限個の不変式 $(J,PJ)$ によって与えられるというものである。ここで $PJ$ は、極化過程 $P$ によって $J$ から導かれるすべての不変式を表す。

実際、$(F')$ の有理整数係数をもつ任意の同時不変式 $S$ は、$N$ 個の各行 $A_1,
\ldots,A_N$ に関して斉次な形式であり、係数は有理整数であるから、恒等式 (1) をこれに適用できる。極化過程 $P$ を $S$ に適用しても再び不変式が得られることは知られているので、形式 $Z$ も不変式である。すなわち、それらは行 $A_1,
\ldots,A_n$ だけを含むから、$F_1,
\ldots,F_n$ の同時不変式である。したがって仮定により、それらは不変式 $(J)$ の整有理関数である。よって恒等式 (1) は
\[
  S=\sum PY
\]
となる。ここで $Y$ は不変式 $(J)$ の冪積を表す。ところが $P$ の定義により、$Y$ に対する過程 $P$ の実際の実行は、単純な極化操作 $P_{hk}$ を有限回反復して逐次実行することにほかならず、積の微分法則によって、不変式 $(J)$ と $(PJ)$ の積の和へ導く。同じことは、$(J)$ と $(PJ)$ からなる積に $P_{hk}$ を適用する場合にも成り立つ。したがって $PY$ も $(J,PJ)$ の整有理結合だけを与える。\emph{これによって、すべての不変式 $S$ が $(J,PJ)$ によって整有理的に表されることが証明された。}

\subsection*{II.}
還元定理を証明する前に、形式の線形族についていくつかの事実を思い出す。$K$ 上の\emph{形式の線形族}とは、同じ次元をもつ形式の系であり、$\theta_1$ と $\theta_2$ を含めば常に $c_1\theta_1+c_2\theta_2$ も含むという意味で完全なものをいう。ここで $c_1,c_2$ は所定の有理性領域 $K$ の量であり、$K$ は $\theta$ たちの係数領域を含むものとする。これらの形式の中には、常に有限個、たとえば $\rho$ 個、$K$ 上線形独立なものがあり、任意の $\rho+1$ 個の形式は $K$ からの係数をもつ線形関係を満たす。この族は $K$ に関して\emph{階数 $\rho$} をもつ。\footnote{より一般の形式の線形族、すなわち階数が有限である必要のないものは、同じ次元という制限を外すか、または $K$ が $\theta$ たちの係数領域を含むという制限を外すことによって得られる。}以下では、$\mathfrak L$ と $\mathfrak T$ が $K$ 上の同じ階数 $\rho$ の二つの線形族で、$\mathfrak T$ が $\mathfrak L$ の部分族であるなら、$\mathfrak L$ と $\mathfrak T$ は同一（同値）である、という容易に分かる事実を用いる。

これを準備として、還元定理へ進む。I の恒等式
\[
  \theta=\sum PZ
\]
は、同じ極化過程 $P$ を両辺に適用すると、$P\theta$ についての類似の恒等式へ移る。したがって還元定理は、同時に、$\theta$ のすべての極化を特殊な極化 $PZ$ の和で表示することを与える。逆に、これら特殊な極化はもちろんすべての極化の中に含まれるから、還元定理はこれら二つの線形族の同値性を主張している。とくに、各々の行において $\theta$ と同じ次元をもつ形式からなる部分族の同値性も主張している。この同値性を証明するために、形式 $Z$ によって定まる中間の族を挿入する。簡単のため、まず三つの二元行の場合 $(N=3,n=2)$ について証明を与え、その後、完全に平行な一般証明を短く示す。

そこで、$\theta(x,y,z)$ がそれぞれ行
\[
  x_1x_2;\quad y_1y_2;\quad z_1z_2;
\]
について次元 $\alpha,\beta,p$ をもつとする。まず $\theta$ の係数は不定元 $a$（または有理数）とする。次の三つの形式の線形族を考える。

1. 族 $\mathfrak L$ は、$\theta$ から極化過程 $P$ によって生じ、個々の行 $x,y,z$ において $\theta$ と同じ次元をもつすべての形式 $f(x,y,z)$ から成る。とくに $\mathfrak L$ は $\theta$ を含む。$f(x,y,z)$ を $x,y,z$ および不定元 $a$ の形式と見るなら、係数体は有理数体 $R$ である。$K$ としても $R$ を取らねばならない。なぜなら、$c_1,c_2$ が有理数である場合にのみ、$c_1P_1+c_2P_2$ も再び過程 $P$ になるからである。$\mathfrak L$ の $R$ に関する階数を $\rho$ とする。

2. 族 $\mathfrak S$ は、
\[
  g(\lambda;x,y)=f(x,y,\lambda_1x+\lambda_2y)
\]
によって定義されるすべての形式 $g(\lambda;x,y)$ から成る。あるいは極化過程で表せば、
\[
\begin{aligned}
  g(\lambda;xy)
  &=\frac1{p!}\left[(\lambda_1x+\lambda_2y)\frac{\partial}{\partial z}\right]^p f(x,y,z) \\
  &=g_0(xy)\lambda_1^p+g_1(x,y)\lambda_1^{p-1}\lambda_2+\cdots+g_p(xy)\lambda_2^p,
\end{aligned}
\]
ただし
\[
  i!(p-i)!\,g_i(xy)=
  \left(y\frac{\partial}{\partial z}\right)^i
  \left(x\frac{\partial}{\partial z}\right)^{p-i}f(xyz).
\]
\footnote{I と同じく、
\[
 \left(t\frac{\partial}{\partial z}\right)=t_1\frac{\partial}{\partial z_1}+t_2\frac{\partial}{\partial z_2},
\]
したがってとくに
\[
 \left[(\lambda_1x+\lambda_2y)\frac{\partial}{\partial z}\right]
 = (\lambda_1x_1+\lambda_2y_1)\frac{\partial}{\partial z_1}
 + (\lambda_1x_2+\lambda_2y_2)\frac{\partial}{\partial z_2},
\]
\[
 \left[z\left(\frac{\partial}{\partial\lambda_1}\frac{\partial}{\partial x}
 +\frac{\partial}{\partial\lambda_2}\frac{\partial}{\partial y}\right)\right]
 =z_1\left(\frac{\partial}{\partial\lambda_1}\frac{\partial}{\partial x_1}
 +\frac{\partial}{\partial\lambda_2}\frac{\partial}{\partial y_1}\right)
 +z_2\left(\frac{\partial}{\partial\lambda_1}\frac{\partial}{\partial x_2}
 +\frac{\partial}{\partial\lambda_2}\frac{\partial}{\partial y_2}\right).
\]}
したがって $g_i(xy)$ は恒等式 (1) の形式 $Z$ を与える。$g$ は $x,y,\lambda$ および不定元 $a$ に関する有理整数形式である。$\mathcal S$ の $R$ に関する階数を $\sigma$ とする。

3. 族 $\mathcal T$ は、$\mathcal S$ が $\mathcal L$ から得られるのと同じ仕方で、逆の極化過程によって $\mathcal S$ から得られる。$\mathcal T$ の形式 $h$ は
\[
\begin{aligned}
 h(x,y,z)
 &=\frac1{p!}\left[z\left(\frac{\partial}{\partial\lambda_1}\frac{\partial}{\partial x}
 +\frac{\partial}{\partial\lambda_2}\frac{\partial}{\partial y}\right)\right]^p g(\lambda;xy)\\
 &=h_0(xyz)+h_1(xyz)+\cdots+h_p(xyz)
\end{aligned}
\]
により定義される。ここで 2. により
\[
  h_i(xyz)=
  \left(z\frac{\partial}{\partial x}\right)^{p-i}
  \left(z\frac{\partial}{\partial y}\right)^i g_i(xy).
\]
個々の $h_i$、したがって $h$ も、形式 $PZ$ およびその和である。$\mathfrak T$ の $R$ に関する階数を $\tau$ とする。

$\mathcal T$ の形式 $h(x,y,z)$ の構成が示すように、これらの形式は $\theta$ から極化過程 $P$ によって生じ、行 $x,y,z$ において個々に $\theta$ と同じ次元をもつ。したがって族 $\mathfrak T$ は $\mathfrak L$ の部分族である。上で述べた事実により、残るのは階数の等しさを証明することだけである。この証明では、$f(xyz)$ の特殊な構造、すなわちそれらが同一の形式 $\theta$ の極化であることは、もはや用いない。

\emph{a) $\mathfrak L$ と $\mathfrak S$ の階数を比較するために、補題 a) を用いる：$f(x,y,\lambda_1x+\lambda_2y)=0$ からは、必然的に $f(x,y,z)=0$ が従う。}これは、三つの二元行のあいだの恒等関係
\[
  (xy)z+(yz)x+(zx)y=0,
  \qquad (xy)=(x_1y_2-x_2y_1),\ldots,
\]
からただちに従う。この関係は
\[
  f[x,y,(zy)x+(xz)y]=(xy)^p f(x,y,z)
\]
を含意する。したがって $f(x,y,\lambda_1x+\lambda_2y)=0$（$\lambda$ について恒等的）から、置換 $\lambda_1=(zy)$, $\lambda_2=(xz)$ により、かつ $(xy)\ne0$ であることから、必然的に $f(x,y,z)=0$ が従う。

この補題により、$\mathfrak S$ の形式と $\mathfrak L$ の形式との間には一対一対応がある。実際、
\[
  f_1(x,y,\lambda_1x+\lambda_2y)=g(\lambda;xy),\qquad
  f_2(x,y,\lambda_1x+\lambda_2y)=g(\lambda;xy)
\]
からは、必然的に
\[
  f_1(xyz)=f_2(xyz)
\]
が従う。さらに、$\mathfrak S$ における任意の関係は、$\mathfrak L$ における一意に対応する形式のあいだにも保たれる（また逆も成り立つ）。なぜなら、
\[
  c_1g^{(1)}(\lambda;xy)+c_2g^{(2)}(\lambda;xy)+\cdots+c_\varkappa g^{(\varkappa)}(\lambda;xy)=0
\]
から、補題により必然的に
\[
  c_1f^{(1)}(x,y,z)+c_2f^{(2)}(x,y,z)+\cdots+c_\varkappa f^{(\varkappa)}(x,y,z)=0
\]
が従うからである。\emph{これにより $\mathfrak L$ と $\mathfrak S$ の階数の等しさ、すなわち $\rho=\sigma$ が証明された。}

\emph{b) $\mathfrak S$ と $\mathfrak T$ の階数を比較するために、対応する補題 b) を用いる：$h(x,y,z)=0$ からは、必然的に $g(\lambda;xy)=0$ が従う。}証明のため、3. により $h(x,y,z)=0$ は、個々の冪積
\[
 \left(\frac{\partial}{\partial\lambda_1}\frac{\partial}{\partial x_1}
 +\frac{\partial}{\partial\lambda_2}\frac{\partial}{\partial y_1}\right)^{p_1}
 \left(\frac{\partial}{\partial\lambda_1}\frac{\partial}{\partial x_2}
 +\frac{\partial}{\partial\lambda_2}\frac{\partial}{\partial y_2}\right)^{p_2}g(\lambda;x,y),
 \qquad (p_1+p_2=p)
\]
の消滅と同値であることに注意する。さらに微分すれば、冪積
\[
\begin{aligned}
&\left(\frac{\partial}{\partial x_1}\right)^{\alpha_1}
 \left(\frac{\partial}{\partial x_2}\right)^{\alpha_2}
 \left(\frac{\partial}{\partial y_1}\right)^{\beta_1}
 \left(\frac{\partial}{\partial y_2}\right)^{\beta_2} \\
&\quad\cdot
 \left(\frac{\partial}{\partial\lambda_1}\frac{\partial}{\partial x_1}
 +\cdots\right)^{p_1}
 \left(\frac{\partial}{\partial\lambda_1}\frac{\partial}{\partial x_2}
 +\cdots\right)^{p_2}
 g(\lambda,xy)
\end{aligned}
\]
もまた消滅する。したがって、これらの冪積の任意の線形結合が消滅し、それゆえ任意の形式
\[
  f\left(\frac{\partial}{\partial x},\frac{\partial}{\partial y},
  \frac{\partial}{\partial\lambda_1}\frac{\partial}{\partial x}
  +\frac{\partial}{\partial\lambda_2}\frac{\partial}{\partial y}\right)g(\lambda,xy)
\]
も消滅する。とくに
\[
  f(x,y,\lambda_1x+\lambda_2y)=g(\lambda;xy)
\]
となるように $f$ を選べば、さらに
\[
  g\left(\frac{\partial}{\partial\lambda};
    \frac{\partial}{\partial x},
    \frac{\partial}{\partial y}\right)g(\lambda,x,y)=0
\]
を得る。あるいは、もし
\[
  g(\lambda,xy)=\sum G_i\lambda_1^{\varkappa_1}\lambda_2^{\varkappa_2}
  x_1^{\varkappa_3}x_2^{\varkappa_4}y_1^{\varkappa_5}y_2^{\varkappa_6}
\]
で、$G_i$ が不定元 $a$ に関する有理整数線形形式であるなら、
\[
  \sum \varkappa_1!\varkappa_2!\cdots\varkappa_6!\,G_i^2=0,
\]
したがって $g(\lambda;xy)=0$ となる。

補題 b) から、a) とまったく同様に、\emph{$\mathfrak S$ と $\mathfrak T$ の階数は等しい；$\sigma=\tau$}ことが従う。

したがって a) と b) により $\rho=\tau$ であり、従って、$\mathfrak T$ は $\mathfrak L$ の部分族であるから、これら二つの線形族の同値性が証明された。よって $\mathfrak L$ の任意の形式、とくに $\theta$ は、$\rho$ 個の線形独立な形式 $h(x,y,z)$ の有理整数線形結合として表せる：
\[
  \theta=\sum c_i h^{(i)}(x,y,z)=\sum PZ.
\]
$\theta$ の不定係数 $a$ に対して導かれたこの恒等式は、その不定元 $a$ を任意の有理性領域の量で置き換えても有効である。したがって、\emph{三つの二元行の場合に対して、還元定理は一般に証明された。}

各々 $n$ 変数からなる $N$ 行の場合の一般証明は、完全に平行である。$\theta(A)$ が行 $A_i$ に関して次元 $\alpha_i$ をもつとする。再び次の三つの形式の線形族を考える。

1. 族 $\mathfrak L$ は、$\theta(A)$ から極化過程 $P$ によって生じ、各行 $A_i$ において $\theta(A)$ と同じ次元をもつすべての形式 $f(A)$ から成る。

2. 族 $\mathfrak S$ は、置換
\[
\begin{aligned}
 A_{n+1}&=\lambda_{n+1}^{(1)}A_1+\lambda_{n+1}^{(2)}A_2+
        \cdots+\lambda_{n+1}^{(n)}A_n,\\
 &\vdotswithin{=}\\
 A_N&=\lambda_N^{(1)}A_1+\lambda_N^{(2)}A_2+
        \cdots+\lambda_N^{(n)}A_n;
\end{aligned}
\]
によって $f(A)$ から生じるすべての形式 $g(\lambda;A_1\cdots A_n)$ から成る。この過程で、$g(\lambda;A_1\cdots A_n)$ における $\lambda$ の個々の冪積の係数は、恒等式 (1) の形式 $Z$ を与える。

3. 族 $\mathfrak T$ は、次で定義されるすべての形式 $h(A)$ から成る：
\[
\begin{aligned}
 h(A)=&\left[A_{n+1}\left(
 \frac{\partial}{\partial\lambda_{n+1}^{(1)}}\frac{\partial}{\partial A_1}
 +\cdots+
 \frac{\partial}{\partial\lambda_{n+1}^{(n)}}\frac{\partial}{\partial A_n}\right)\right]^{\alpha_{n+1}} \\
&\cdots
\left[A_N\left(
 \frac{\partial}{\partial\lambda_N^{(1)}}\frac{\partial}{\partial A_1}
 +\cdots+
 \frac{\partial}{\partial\lambda_N^{(n)}}\frac{\partial}{\partial A_n}\right)\right]^{\alpha_N}
 g(\lambda;A_1\cdots A_n).
\end{aligned}
\]
ここで $h(A)$ は形式 $PZ$ の和である。

任意の $n+1$ 行のあいだの恒等関係により、補題 a) は有効なままである。補題 b) も同様に有効である。なぜなら、その証明には変数の数が入らなかったからである。したがって $\mathfrak L$ と $\mathfrak T$ は同じ階数をもち、かつ $\mathfrak T$ は $\mathfrak L$ の部分族であるから、両者は同一である。従って恒等式
\[
  \theta=\sum PZ
\]
は不定係数について成り立ち、ゆえに任意の有理性領域の量である係数についても成り立つ。\emph{これにより還元定理は一般に証明された。}

\subsection*{III.}
同様の考察が、\emph{一般級数展開}（$N\le n$ の場合）も与えることを簡単に示す。$Z$ を、各々 $n$ 個の量からなる $n$ 行 $A_1,
\ldots,A_n$ において別々に斉次な形式とし、$\Delta$ をこれらの行の行列式とする。まず (1) に対応する恒等式を証明する：
\begin{equation}
  Z\equiv\sum PH \mod{\Delta},
\tag{2}
\end{equation}
ここで形式 $H$ は行 $A_1,
\ldots,A_{n-1}$ だけを含み、$Z$ から過程 $P$ によって導かれる。

証明は、II で得られた関係を $\Delta$ を法とする合同式へ移すことから成る。簡単のため、三つの三元行 $x,y,z$ を取る。係数は不定元であるとする。

三つの線形族 $\mathfrak L,\mathfrak S,\mathfrak T$ は、二元行の場合の II と同じように定義する。その階数を $\rho,\sigma,\tau$ とし、$\rho^*$ および $\tau^*$ は $\Delta$ を法とする $\mathfrak L$ および $\mathfrak T$ の階数を表すものとする（すなわち、$\Delta$ を法として線形独立な $\mathcal L$ の形式は $\rho^*$ 個存在するが、$\rho^*+1$ 個は存在しない）。すると II と同様に、次が成り立つ。

\emph{補題 a)：$f(x,y,\lambda_1x+\lambda_2y)=0$ からは、必然的に $f(xyz)\equiv0\pmod{\Delta}$ が従う（また逆も成り立つ）。}実際、四つの三元行のあいだの関係
\[
  \Delta_1x-\Delta_2y+\Delta_3z=\Delta t,
  \qquad \Delta_1=(yzt),\ldots,
\]
により、三つの三元行のあいだには常に $\Delta$ を法とする線形関係
\[
  \Delta_1x-\Delta_2y+\Delta_3z\equiv0\mod{\Delta}
\]
がある。したがって II と同様に
\[
  f(x,y,-\Delta_1x+\Delta_2y)\equiv\Delta_3^p f(x,y,z)\pmod{\Delta}
\]
を得る。$f(x,y,\lambda_1x+\lambda_2y)=0$（$\lambda$ について恒等的）から、$\Delta$ は既約であり $\Delta_3$ は $\Delta$ で割り切れないので、必然的に $f(x,y,z)\equiv0\pmod{\Delta}$ が従う。逆は $\Delta(x,y,\lambda_1x+\lambda_2y)=0$ から明らかである。したがって II と同様に $\rho^*=\sigma$ である。

補題 b) はその証明とともに有効であり、したがって $\sigma=\tau$ である。

$\tau=\tau^*$ を証明するために次を用いる。

\emph{補題 c)：$h(x,y,z)\equiv0\pmod{\Delta}$ からは、必然的に $h(x,y,z)=0$ が従う。}$h(x,y,z)$ は極形式なので、既知の $\Omega$ 過程を適用すると消滅する。これは微分方程式
\[
  \Delta\!\left(\frac{\partial}{\partial x},\frac{\partial}{\partial y},\frac{\partial}{\partial z}\right)h(xyz)
  =\left(\frac{\partial}{\partial x}\frac{\partial}{\partial y}\frac{\partial}{\partial z}\right)h(x,y,z)=0.\footnotemark
\]
\footnotetext{3. により、
\[
 c\cdot\Omega(h)=\left(
 \frac{\partial}{\partial \xi}\frac{\partial}{\partial \eta}
 \left(\frac{\partial}{\partial\lambda_1}\frac{\partial}{\partial \xi}
 +\frac{\partial}{\partial\lambda_2}\frac{\partial}{\partial \eta}\right)
 \right)
 \left(x\frac{\partial}{\partial \xi}\right)^{\alpha-1}
 \left(y\frac{\partial}{\partial \eta}\right)^{\beta-1}
 \left[
 z\left(\frac{\partial}{\partial\lambda_1}\frac{\partial}{\partial \xi}
 +\frac{\partial}{\partial\lambda_2}\frac{\partial}{\partial \eta}\right)
 \right]^{p-1}\cdot g(\lambda;\xi,\eta)
\]
 である。従って \(\Omega(h)\) は行列式因子が消滅するために消滅する。}

いま \(h(x,y,z)=\Delta(x,y,z)\,k(x,y,z)\) とすれば、さらに微分することにより
\[
  k\!\left(\frac{\partial}{\partial x},\frac{\partial}{\partial y},\frac{\partial}{\partial z}\right)
  \Delta\!\left(\frac{\partial}{\partial x},\frac{\partial}{\partial y},\frac{\partial}{\partial z}\right)
  =
  h\!\left(\frac{\partial}{\partial x},\frac{\partial}{\partial y},\frac{\partial}{\partial z}\right)h(x,y,z)
\]
も消滅し、したがって \(h(x,y,z)\) も消滅する。ゆえに \(\rho^*=\sigma=\tau=\tau^*\) であり、\(\mathcal T\) は \(\mathcal L\) の部分族なので、\emph{恒等式 (2) が証明された}。\(n\) 変数の場合も同様の考察で証明される。

恒等式 (2) を $\Delta$ の乗数に適用すると、展開
\[
  Z=\varphi_0+\Delta\varphi_1+\Delta^2\varphi_2+\cdots+\Delta^\mu\varphi_\mu
\]
が得られる。ここで各 $\varphi_i$ は $\Omega$ 過程の適用で消滅する。この過程を $i$ 回繰り返すと、一般に $\Omega(\Delta\cdot\psi)=c_1\psi+c_2\Delta\cdot\Omega(\psi)$ であるから、
\[
  c\cdot \Omega^i(Z)\equiv\varphi_i\pmod{\Delta};
\]
他方、(2) により
\[
  c\cdot\Omega^i(Z)\equiv\sum PH_i\pmod{\Delta}
\]
である。ここで $H_i$ は、$H$ が $Z$ から導かれるのと同じ仕方で、形式 $\Omega^i(Z)$ から導かれる。したがって補題 c)（$n$ 変数へ拡張したもの）により $\varphi_i=\sum PH_i$ を得る。従って
\begin{equation}
  Z=\sum PH+\Delta\cdot\sum PH_1+\Delta^2\cdot\sum PH_2+
  \cdots\Delta^\mu\cdot\sum PH_\mu
\tag{3}
\end{equation}
\emph{これは、各々 $n$ 変数からなる $n$ 行の形式の、ただ $n-1$ 行の形式の極化による最も一般的な級数展開である。}

$N$ 行の $n$ 変数の形式 $(N<n)$ に対する級数展開は、よく知られているように、単純な置換によって得られる。$N=3$, $n>3$ に対し
\[
  u=\xi_1x+\xi_2y+\xi_3z,
  \quad v=\eta_1x+\eta_2y+\eta_3z,
  \quad w=\zeta_1x+\zeta_2y+\zeta_3z
\]
と置き、$H(u,v,w)=Z(\xi,\eta,\zeta)$ を (3) に従って展開する。すると $\xi,\eta,\zeta$ は組合せ $u,v,w$ においてだけ現れるので、
\[
  \left(\eta\frac{\partial}{\partial \xi}\right)=
  \left(v\frac{\partial}{\partial u}\right)\cdots,
  \qquad
  \Omega=\sum\left(\frac{\partial}{\partial u}
  \frac{\partial}{\partial v}
  \frac{\partial}{\partial w}\right)_{ikl}(xyz)_{ikl}
  =\nabla_{uvw}(xyz).
\]
特殊化 $\xi_1=\eta_2=\zeta_3=1$; $\xi_2=\xi_3=\cdots=\zeta_2=0$ のもとで、$H(u,v,w)$ は $H(x,y,z)$ へ移り、$\nabla_{uvw}(xyz)$ は、数値因子を除いて $\nabla_{xyz}(xyz)=\nabla_{xyz}$ へ移る。なぜなら、$(y\partial/\partial x)$ などを行列式 $(xyz)_{ikl}$ に適用するとそれらは消滅するからである。対応する主張は $N$ 行についても成り立つので、(3) は展開
\begin{equation}
  H=\sum P\Phi+\sum P\Phi_1+\cdots+\sum P\Phi_\mu
\tag{4}
\end{equation}
を与える。ここで $\Phi_i$ は、$\nabla^i_{A_1\ldots A_N}H$ から過程 $P$ によって生じ、行 $A_1,
\ldots,A_{N-1}$ だけを含む。ただし $\nabla^i$ に現れる行列式の組合せは別であり、それらはすでに述べたように、極化の際には考慮に入らない。

次に、$\nabla^i$ の中のこれらの行列式を不定元 $p$ で置き換えれば、得られる形式は再び (4) に従って展開できる。このようにして最終的に展開
\begin{equation}
  H=\left[\sum P\psi(p)\right]_{p=\text{$A_1\cdots A_N$ の行列式}A_1\cdots A_N},
\tag{5}
\end{equation}
が得られる。ここで $\psi$ は、もとの形式から過程 $P$ および $\nabla_{A_1\ldots A_N}(p)$, $\nabla_{A_1\ldots A_{N-1}}(p)$ などによって生じる。これらの展開およびその組合せ、例えば (1) に現れる形式 $Z$ に対して (3)、ついで (4) と (5) を適用することにより、各々 $n$ 個の共変な変数からなる任意に多くの行の形式について既知のすべての展開が尽くされる。

Erlangen, 1915 年 1 月 5 日。

\clearpage
\section*{9. 整超越数から成る最も一般の領域}
\begin{center}
Math. Ann. 77 (1916), pp. 103--128
\end{center}

E. Zermelo は整列定理を用いて「整超越数」の領域、すなわち次の抽象的性質をもつ数の領域 $\mathfrak G$ を構成した。
\begin{enumerate}
\item[I.] $\mathfrak G$ の二つの数の和、差、積は再び $\mathfrak G$ の数である。
\item[II.] すべての実数または複素数は、$\mathfrak G$ の二つの数の商である。
\item[III.] すべての有理整数（それぞれ代数的整数）は $\mathfrak G$ の元である。
\item[IV.] 整でない有理数（それぞれ整でない代数的数）は $\mathfrak G$ に属さない。
\footnote{E. Zermelo, Über ganze transzendente Zahlen, Math. Ann. 75, p. 434 (1914).}
\end{enumerate}

この構成は、整列定理が全ての数の\emph{代数的基底}の存在を含意するという事実に基づいている。すなわち、相互の間に代数関係が存在せず、しかもそれらによって他の全ての数を代数的に表せるような\emph{数} $\eta$ \emph{の系} $H$ が存在する。この基底数 $\eta$ は代数的に独立であるため、不定元とみなすことができる。したがって、求める領域は、不定元 $\eta$ の有理関数および代数関数から成る整域となり、\footnote{この理由から、本論文の考察は一部、著者の以前の論文 Körper und Systeme rationaler Funktionen, Math. Ann. 76, p. 161 (1915) の考察と平行している。}その係数は全ての代数的数から成る体 $\mathsf{K}$ に属する。Zermelo によって構成された特殊な領域 $\mathfrak G_\eta$ は、次の基底性質によって特徴づけられる。

$\mathfrak G_\eta$ は次を含む。
\begin{enumerate}
\item[1)] すべての基底数 $\eta$。
\item[2)] $\eta$ のすべての多項式で、係数が整であるもの。\footnote{本論文を通じて「整」とは、係数が全ての代数的整数から成る整域 $[\mathsf{K}]$ に属することを意味する。$\mathfrak G_\eta$ の定義については、1)--5) で有理整数係数だけを考えても十分である。しかし、より一般の領域については、その場合、基底性質がより単純でなくなる。}
\item[3)] $\eta$ のすべての代数的整関数で、係数が整であるもの。
\end{enumerate}
$\mathfrak G_\eta$ は次を排除する。
\begin{enumerate}
\item[4)] $\eta$ のすべての有理的に分数的な関数で、係数が整であるもの。
\item[5)] $\eta$ のすべての代数的に分数的な関数で、係数が整であるもの。\footnote{Zermelo, loc. cit., § 3.}
\end{enumerate}

いま、§§ 2, 3 の例と比較すれば、Zermelo 領域 $\mathfrak G_\eta$ と本質的に異なる領域 $\mathfrak G$ が存在することは容易に分かる。すなわち、どのような整列のもとでも基底性質 1)--5) がすべて成り立つことはなく、したがってどのような整列のもとでも Zermelo の構成によって生成することのできない領域 $\mathfrak G$ が存在する。言い換えると、$\mathfrak G_\eta$ の上へ一対一かつ同型的に写すことのできない領域 $\mathfrak G$ が存在する。そこで、整超越数から成る最も一般の領域という問題が生じる。この問題は以下で完全に答えられる。

そのためにはまず、基底性質のうちどれが抽象的な定義性質の帰結であり、どれが特殊な構成に由来するのかを決定しなければならない。任意に与えられた領域 $\mathfrak G$ に対して、基底性質 1), 2) が満たされるような整列が存在することを示す（§ 1）。したがって、すべての領域 $\mathfrak G$ は全ての数の代数的基底によって構成できる。これ以外の基底性質はどれも満たされる必要はない（§§ 2, 3）。特に、Zermelo 領域のさらに一つの抽象的性質、すなわち代数的整閉性はすべての領域 $\mathfrak G$ に属するわけではないことが従う。これは次のように定式化できる。
\begin{enumerate}
\item[V.] $\mathfrak G$ の有限個の数に代数的に整かつ整に依存するすべての数は $\mathfrak G$ に属する。
\end{enumerate}

I--IV に加えて V が成り立つ特殊な領域を、代数的に整な超越数から成る領域 $\mathfrak H$ と呼ぶ。最も一般の領域 $\mathfrak H$ については、次の単純な結果を得る（§ 4）。
\begin{enumerate}
\item[a)] 同じ基底 $\mathsf{H}$ によって構成できるすべての領域 $\mathfrak H$ の交わりは、基底のすべての代数的整関数、すなわち Zermelo 領域 $\mathfrak G_\eta$ によって与えられる。この $\mathfrak G_\eta$ 自身も領域 $\mathfrak H$ である。（これは Zermelo 領域の抽象的特徴づけも与える。）
\item[b)] すべての領域 $\mathfrak H$ は、任意の系 $\mathfrak S(\eta)$ を $\mathfrak G_\eta$ に代数的に整に添加することによって生じる。\footnote{この表現の説明は § 4 にある。} ただし $\mathfrak S(\eta)$ は、添加によって代数的に分数的な数が領域に入らないというただ一つの条件を満たせばよい。
\end{enumerate}

代数的基底を基礎にとるとき、最も一般の領域 $\mathfrak G$ に対する結果はこれほど単純ではない。この場合、すべての領域 $\mathfrak G$ の交わりはそれ自身が領域 $\mathfrak G$ ではなく、したがって最も一般の領域の構成は見通しにくい（§ 5）。

領域 $\mathfrak H$ の場合と類似の結果を得るために、代数的基底を\emph{有理基底} $\Theta$ で置き換える。これは、すべての数を有理的に表すことを可能にし、かつ少なくとも一つの整列のもとでどの基底数も有限個の先行するものによって有理的に表せないような\footnote{線形基底および代数的基底とは異なり、ここでは元の順序が役割を果たす。例えば順序 $\eta,\sqrt[\nu]{\eta}$ では両方の元を入れなければならないが、逆順では $\sqrt[\nu]{\eta}$ だけでよい。}\emph{数} $\vartheta$ \emph{の系} $\Theta$ である。この有理基底 $\Theta$ には、さらに、$\vartheta$ のすべての整多項式から成る整域 $[\Theta]$ が代数的に分数的な数を含まないという付帯条件を課す必要がある。（この付帯条件をもつ基底の構成は § 7 で示す。）すると、最も一般の領域 $\mathfrak G$ ---これは有理的に整な超越数から成る領域とも呼べる--- について、再び次の単純な結果を得る（§ 6）。
\begin{enumerate}
\item[a)] 同じ有理基底 $\Theta$ によって構成できるすべての領域 $\mathfrak G$ の交わりは、基底のすべての整有理関数、すなわち領域 $[\Theta]$ によって与えられ、この $[\Theta]$ 自身が領域 $\mathfrak G$ である。
\item[b)] すべての領域 $\mathfrak G$ は、任意の系 $\mathfrak S(\vartheta)$ を $[\Theta]$ に有理的に整に添加することによって生じる。ここで $\mathfrak S(\vartheta)$ は、添加によって代数的に分数的な数が領域に入らないというただ一つの条件を満たせばよい。
\end{enumerate}

領域 $\mathfrak H$ と完全な類似を得るには、領域 $\mathfrak G$ の定義性質 III と IV から代数的数を省かなければならない。III と IV をこのように修正すると、有理基底による整元の構成は任意の体へ移すことができる（§ 10）。これに対し Zermelo の構成は体の代数的閉性を前提としている。

$\mathsf{H}$ がすべての代数的基底を走り、また $\Theta$ が付帯条件を満たすすべての有理基底を走るとき、結果 a), b) はすべての領域 $\mathfrak H$ および $\mathfrak G$ の全体を与える。すべての同型な領域を一つの類にまとめるならば、この全体から、すべての類の全体を代表する部分集合を選ぶことができる。§ 9 により、少なくとも可算無限個の類がある（§ 8）。

\subsection*{§ 1. すべての領域 $\mathfrak G$ に対する基底性質 1), 2) の証明。}

$\mathfrak G$ を、抽象的性質 I--IV をもつ任意に与えられた整超越数の領域とする。E. Zermelo が全複素数体に対して用いた手続きに従い、$\mathfrak G$ の中から、\emph{$\mathfrak G$ のすべての数の代数的基底} $H$ を選ぶ。II により、すべての実数または複素数は $\mathfrak G$ の二つの数の商として表されるから、$\mathsf{H}$ は同時に\emph{すべての}数の代数的基底である。したがって、任意に与えられた領域 $\mathfrak G$ に対し、全ての数の代数的基底を $\mathfrak G$ に属するように選ぶことができる。ゆえに基底性質 1) は満たされる。さらに III により、$\mathfrak G$ は全ての代数的整数から成る整域 $[\mathsf{K}]$ を含み、したがって I により、基底元 $\eta$ の $[\mathsf{K}]$ 係数多項式、すなわち $\eta$ のすべての整多項式を含む。これらは整域 $[\mathsf{H}]$ をなす。ゆえに基底性質 2) も常に満たされる。言い換えると、すべての領域 $\mathfrak G$ は、全ての数の代数的基底 $\mathsf{H}$ によって、$\mathfrak G$ が整域 $[\mathsf{H}]$ を因子として含むように構成できる。

\textbf{注。} それでも、基底 $\mathsf{H}$ から出発して、$H$ 自身に関しては基底性質 1), 2) が成り立たないが、別の基底 $H'$ に関しては成り立つ領域 $\mathfrak G$ を構成することはもちろん可能である。例えば、Zermelo 領域 $\mathfrak G_\eta$ の全構成において、ある基底元 $\eta$ を多項式 $f(\eta)$ で置き換えて領域 $\mathfrak G'$ を得るとする。このとき $\mathfrak G'$ は $\eta$ も $[\mathsf{H}]$ も含まない。しかし基底 $\mathsf{H}$ において基底元 $\eta$ を $\zeta=f(\eta)$ で置き換えれば、$\mathsf{H}$ は再び代数的基底 $\mathsf{H}'$ となり、$\mathfrak G'$ は $\mathsf{H}'$ に関して基底性質 1), 2) を満たす。

\subsection*{§ 2. 基底性質 4), 5) の排除。}

ここで、領域 $\mathfrak G$ は残りの基底性質を何も満たす必要がないことを示す。まず 4), 5) は、Zermelo の構成において整域 $[\mathsf{H}]$ を整有理「関数式」の領域で置き換えることにより排除される。

Weber によれば、\footnote{H. Weber, Lehrbuch der Algebra. Kleine Ausgabe §§ 96, 97.} \emph{有理関数式}とは、次の一意に定まる表示における、有理数係数の任意の有理関数をいう。
\[
 A(\eta_1\cdots \eta_t)
   = a\cdot\frac{E_1(\eta_1\cdots \eta_t)}{E_2(\eta_1\cdots \eta_t)},
\]
ここで $E_1,E_2$ は、有理整数係数をもつ互いに素な原始多項式であり、$a$ は正の有理数（$A$ の絶対値）である。特に $a$ が\emph{有理整数}であるとき、$A$ は\emph{整有理関数式}と呼ばれる。整有理関数式の特殊な場合として有理整数および有理整数係数多項式が含まれ、一方、有理的に分数的な数および有理的に分数的な係数をもつ多項式は、分数的有理関数式に数えられる。

いま領域 $\mathfrak G$ を、一度に有限個の基底元 $\eta$ のすべての整有理関数式の全体 $\mathfrak J$ と、$\mathfrak J$ に代数的に整に依存するすべての関数、すなわち最高係数が 1 で残りの係数が整有理関数式であるようなある方程式を満たす関数、から成るものとする。

$\mathfrak G$ に対する性質 I--IV の証明は Zermelo の対応する証明とまったく平行であり、ここでは簡単に示すにとどめる。まず II と III は明らかである。なぜなら $\mathfrak G$ は Zermelo 領域 $\mathfrak G_\eta$ を因子として含むからである。有理整数係数多項式についての Gauss の定理により、整有理関数式の和、差、積は再び整有理関数式である。したがって $\mathfrak J$ は整域をなし、よく知られた議論\footnote{例えば Weber, § 97, 8 を参照。}により $\mathfrak G$ も整域となる。ゆえに I が成り立つ。さらに、拡張された Gauss の定理\footnote{Weber, § 20, 5.}から二つの補助命題が従う。すなわち「$z$ が\emph{ある}方程式によって $\mathfrak J$ に代数的に整に依存するならば、$z$ がすべての有理関数式の体に関して満たす\emph{既約}方程式によってもそうである」\footnote{Weber, § 97, 4 を参照。}、および「すべての有理数の体上既約で、ある代数的数が満たす方程式は、すべての有理関数式の体上でも既約である」である。したがって $\mathfrak G$ は分数的代数的数を含むことができない。よって IV、したがって抽象的性質 I--IV のすべてが検証された。

一方、どの整列のもとでも、$\mathfrak G$ が基底元のすべての有理的および代数的に分数的な関数を排除するように、$\mathfrak G$ から基底を選ぶことはできない。基底元として選ばれるべき領域の任意の関数を $z(\eta)$ とすれば、それは方程式
\[
 z^k+A_1(\eta)z^{k-1}+\cdots+A_k(\eta)=0
\]
によって定義される。ここで
\[
 A_i(\eta)=a_i\cdot\frac{E_1(\eta)}{E_2(\eta)}
\]
は整有理関数式である。すると関数 $\frac{a_k}{z}$ は $\mathfrak G$ に属し、同様に $\frac{a_k}{\sqrt z}$, $\frac{a_k}{\sqrt[3]z}$ なども属する。\footnote{まったく同様に、$\eta$ の任意の代数関数は、有理整数を掛けることによって関数式領域 $\mathfrak G$ の元へ変換できる。したがってこの領域は、代数的数の場合と同様に、任意の複素数を $\mathfrak G$ の元と有理整数との商として表せるという性質をもつ。} したがって基底性質 4), 5) は領域 $\mathfrak G$ の全体に対して排除される。

\textbf{注。} § 3 の例は、$\mathfrak G$ が、どの整列のもとでも有理的に分数的な関数式を含み得る一方、有理的または代数的に分数的な数を含まないこともあり得ることを示す。

\subsection*{§ 3. 基底性質 3) の排除。}

基底性質 3) を排除するため、一般化された\emph{加群領域}を用いる。この領域では、多項式の引数はもはや不定元ではなく、不定元のすべての代数的整関数であり、\footnote{「代数的整関数」とは常に\emph{整}係数をもつ関数を意味する。すなわち、最高係数が 1 で、残りの係数が代数的整数係数の $\eta$ の多項式、つまり $[\mathsf{H}]$ の多項式であるような方程式を満たす関数である。特に $[\mathsf{H}]$ のすべての多項式は代数的整関数に属する。}不定元自身が加群基底の役割を担う。

$\mathfrak G$ を\emph{すべての元}
\begin{equation}
 z=\alpha+g_1(\eta)\cdot\eta_1+g_2(\eta)\cdot\eta_2+\cdots+g_t(\eta)\cdot\eta_t
\tag{1}
\end{equation}
\emph{から成るものとする。ここで $\alpha$ はすべての代数的整数を走り、$\eta_1\cdots\eta_t$ はすべての基底元を走り、固定された各組合せ $\eta_1\cdots\eta_t$ に対して関数 $g_1(\eta)\cdots g_t(\eta)$ はそれぞれ独立に $\eta$ のすべての代数的整関数を走る。}\footnote{加群性は元 $z-\alpha$ に属する。}

この加群領域に対して性質 I--IV が成り立つことは容易に確かめられる。まず I が成り立つ。なぜなら二つの元 $z$ の和、差、積は再び (1) の形をとるからである。領域 $\mathfrak G$ はすべての基底元 $\eta$ も含み、さらに $\eta\cdot g(\eta)$ という形のすべての元も含む。ここで $g(\eta)$ は $\eta$ のすべての代数的整関数を走る。したがって $\eta$ のすべての代数的整関数、従って $\eta$ のすべての代数関数は、$\mathfrak G$ の二つの元の商として表せる。これにより II が証明される。表示 (1) には、$g_1=0\cdots g_t=0$ とすれば、すべての代数的整数も特に含まれている。しかし $z$ は分数的代数的数に等しくなり得ない。実際、$z$ が代数的数 $\beta$ を表すならば、
\[
 \beta=\alpha+g_1(\eta)\eta_1+\cdots+g_t(\eta)\eta_t
\]
が $\eta$ に関して恒等的に成り立つので、必ず $\beta=\alpha$ である。これは特殊化
\[
 \eta_1=0\cdots \eta_t=0
\]
によって示される。この特殊化のもとで、乗数 $g_1(\eta)\cdots g_t(\eta)$ は代数的整関数として有限に保たれる。なぜなら、それらは代数的整関数または代数的整数へ移るからである。\footnote{IV のこのより詳細な証明は、段落末の拡張例のために与える。ここでは、構成により $z$ が代数的整関数であるという注だけでも十分である。} したがって $\mathfrak G$ に対して条件 I--IV が満たされる。

一方、どの整列のもとでも、基底元のすべての代数的整関数が $\mathfrak G$ に属するように $\mathfrak G$ から基底を選ぶことはできない。領域の任意の関数 $z(\eta)$ を、基底元として選ぶべきものとする。したがってこれは実際に不定元 $\eta$ を含むものでなければならない。$z$ に依存するある有限の値以後の $\nu$ について、関数
\[
 \xi=\sqrt[\nu]{z-\alpha}
\]
はもはや $\mathfrak G$ に属せないことを示す。

実際、$\zeta$ が $\mathfrak G$ に属すると仮定する。すると (1) により、$\xi$ は
\[
 \zeta=\gamma+h_1(\eta)\eta_{i_1}+h_2(\eta)\eta_{i_2}
      +\cdots+h_\tau(\eta)\eta_{i_\tau}
\]
の形である。ここで再び $h_1(\eta)\cdots h_\tau(\eta)$ は $\eta$ の代数的整関数であり、$\eta_{i_1}\cdots\eta_{i_\tau}$ は任意の基底元で、特に $\eta_1\cdots\eta_t$ と一致してもよい。まず代数的整数 $\gamma$ が 0 であることを示さなければならない。$h_1(\eta)\cdots h_\tau(\eta)$ は代数的整関数であるから、残りの $\eta$ を任意の値にして $\eta_{i_1}=0\cdots\eta_{i_\tau}=0$ とおくと、$\zeta$ は $\gamma$ に等しくなる。特に、
\[
 \eta_{i_1}=0\cdots \eta_{i_\tau}=0; \qquad
 \eta_1=0\cdots \eta_t=0
\]
とすれば、$\xi$ は $\gamma$ に等しい。しかしこの特殊化のもとで、(1) により関数 $(z-\alpha)$ は消え、したがって $\zeta$ も消える。ゆえに $\gamma=0$ であり、
\[
 \zeta=h_1(\eta)\eta_{i_1}+h_2(\eta)\eta_{i_2}
      +\cdots+h_\tau(\eta)\eta_{i_\tau}
\]
が成り立つ。これを $\nu$ 乗すると
\[
 \zeta^\nu=z-\alpha=F_\nu(\eta)
\]
を得る。ここで $F_\nu(\eta)$ は、$\eta_{i_1}\cdots\eta_{i_\tau}$ に関する $\nu$ 次元の、恒等的には消えない同次形式で、その係数は $\eta$ の代数的整関数である。いま $z-\alpha$ は (1) により $\eta$ の代数的整関数であり、$[\mathsf{H}]$ に関して既約な方程式
\[
 (z-\alpha)^\varkappa+B(\eta)(z-\alpha)^{\varkappa-1}+\cdots+C(\eta)=0
\]
を満たす。ここで最後の係数 $C(\eta)$、すなわち $(z-\alpha)$ とその共役の積、は多項式である。その次数を $\lambda$ とする。一方、$z-\alpha=F_\nu(\eta)$ から、$F_\nu(\eta)$ に対応する共役を掛けることにより
\[
 C(\eta)=G_{\nu\varkappa}(\eta)
\]
を得る。ここで $G_{\nu\varkappa}(\eta)$ は $\eta_{i_1}\cdots\eta_{i_\tau}$ に関する $\nu\chi$ 次元の同次形式で、係数は $\eta$ の多項式である。したがって $G_{\nu\varkappa}$ は $\eta$ に関して少なくとも次数 $\nu\varkappa$ をもつので、$\lambda\geqq\nu\varkappa$ である。すなわち $\nu>\lambda/\chi$ となる関数 $\sqrt[\nu]{z-\alpha}$ は $\mathfrak G$ に属さない。したがって基底性質 3) は領域 $\mathfrak G$ の全体に対して排除される。

\textbf{注。} いま構成した領域をわずかに一般化すると、どの整列のもとでも分数的関数式をも含む領域 $\mathfrak G$ が得られる。$\mathfrak G$ をすべての元
\[
 z=\alpha+\gamma_1(\eta)\eta_1+\gamma_2(\eta)\eta_2
      +\cdots+\gamma_t(\eta)\eta_t
\]
から成るものとし、ここで $\gamma(\eta)$ は $\eta$ の代数的に整な、ただし必ずしも整ではない、すべての関数を走るものとする。性質 I--IV の証明は上と同じである。しかし $\mathfrak G$ は任意の関数 $z$ とともに $a\cdot(z-\alpha)$ も含む。ここで $a$ は任意の有理的または代数的に分数的な数である。したがってどの整列のもとでも分数的関数式が現れる。

\subsection*{§ 4. 代数的に整な超越数から成る最も一般の領域。}

以下を便利に述べるために、「元 $z$ が $\mathfrak T$ に代数的に整に依存する」という表現を導入する。これは、$\mathfrak T$ が任意に与えられた領域を表すとき、$z$ が、最高係数が 1 で、残りの係数が $\mathfrak T$ の元の整多項式であるようなある方程式を満たすことを意味する。\footnote{$\mathfrak T$ が整域の場合には、最高係数は 1 であり、残りの係数は $\mathfrak T$ の元である。}

領域 $\mathfrak T$ は、次の条件（序論を参照）を満たすとき、\emph{代数的整閉}であるという。
\begin{enumerate}
\item[V.] $\mathfrak T$ に代数的に整に依存するすべての元は $\mathfrak T$ に属する。
\end{enumerate}

まず、抽象的性質 V が Zermelo 領域 $\mathfrak G_\eta$ に属することを示す。元 $z$ が方程式 $f(z)=0$ によって $\mathfrak G_\eta$ に代数的に整に依存するとする。すると、$[H]$ に関する $f(z)$ の共役を掛け合わせることにより、$z$ は $[H]$ にも代数的に整に依存することがわかる。したがって $z$ は $\mathfrak G_\eta$ に属する。§ 2 で考えた関数式領域の代数的整閉性も同じように示される。

性質 V がすべての領域 $\mathfrak G$ に属するわけではないことは、§ 3 で考えた加群領域が示している。実際、それは基底性質 3) をもたない。より正確には、§ 3 は、加群領域が領域のどの超越元に関しても代数的整閉でないことを示した。一方、III と IV により、領域の任意の代数的元に関してはもちろんそうである。

これにより、まず領域 $\mathfrak G$ の全体から、抽象的性質 V ももつものを選び出す。これらを「代数的に整な超越数から成る領域 $\mathfrak H$」と呼ぶ。

$\mathfrak H$ をそのような領域とし、$\mathsf{H}$ を $\mathfrak H$ のすべての数の代数的基底とする。II により、$\mathsf{H}$ は同時にすべての数の代数的基底である。III, I, V により、$\mathfrak H$ は $\mathsf{H}$ のほかに、$\eta$ のすべての代数的整関数、すなわち Zermelo 領域 $\mathfrak G_\eta$ を含む。したがって § 1 と同様に、代数的に整な超越数から成るすべての領域 $\mathfrak H$ は、すべての数の代数的基底によって、$\mathfrak H$ が Zermelo 領域 $\mathfrak G_\eta$ を因子として含むように構成できる。

いま $\mathsf{H}$ をすべての数の任意の代数的基底とし、$\mathfrak M(H)$ を $H$ を含むすべての領域 $\mathfrak H$ の全体とする。$\mathfrak M(H)$ の任意の領域 $\mathfrak H$ は、いま証明したように、Zermelo 領域 $\mathfrak G_\eta$ を因子として含む。そして $\mathfrak G_\eta$ 自身が領域 $\mathfrak H$ であるから、$\mathfrak G_\eta$ は $\mathfrak M(H)$ のすべての領域 $\mathfrak H$ の最大公因子、すなわち交わりである。こうして $\mathfrak G_\eta$ は抽象的に定義された。また、$\mathfrak M(H)$ が交わりに関して閉じていることも直接に従う。すなわち、$\mathfrak M(H)$ の任意の部分系のすべての領域 $\mathfrak H$ の交わりは再び $\mathfrak M(H)$ に属する。実際、交わりとして I と V が成り立つ。II と III は $\mathfrak G_\eta$ が因子であることから従い、IV はその交わりが $\mathfrak H$ の因子であることから従う。

§ 2 の関数式領域の例が示すように、一般に領域 $\mathfrak H$ は $\mathfrak G_\eta$ 以外の関数を含む。そのような $\eta$ の有理関数または代数関数を $\psi(\eta)$ とする。すると I により、$\mathfrak H$ は $\psi$ と $\mathfrak G_\eta$ の有限個の関数とのすべての整有理結合も含む。このようにして生じる整域、すなわち $\mathfrak G_\eta$ 係数の $\psi$ のすべての多項式から成る整域を $[\mathfrak G_\eta,\psi]$ と書き、それへ至る過程を「$\psi$ の $\mathfrak G_\eta$ への有理的整添加」と呼ぶ。さらに V により、$\mathfrak H$ は $[\mathfrak G_\eta,\psi]$ に代数的に整に依存するすべての関数を含む。このようにして生じる領域を $\{\mathfrak G_\eta,\psi\}$ と書き、それへ至る過程を「$\psi$ の $\mathfrak G_\eta$ への代数的整添加」と呼ぶ。領域 $\{\mathfrak G_\eta,\psi\}$ は $\mathfrak G_\eta$ 係数の $\psi$ のすべての代数的整関数から成るので、よく知られた議論により整域である。\footnote{例えば Weber, § 97, 6, 7 を参照。}

$\{\mathfrak G_\eta,\psi\}$ も代数的整閉、すなわち条件 V を満たすことを示す。実際、$z$ が方程式
\[
 f(z;a\cdots c)=z^\nu+a z^{\nu-1}+\cdots+c=0
\]
によって $\{\mathfrak G_\eta,\psi\}$ に代数的に整に依存するとする。ここで $a\cdots c$ は $\{\mathfrak G_\eta,\psi\}$ の元として、$[\mathfrak G_\eta,\psi]$ の係数をもつ方程式
\[
\begin{gathered}
 a^\lambda+A_1a^{\lambda-1}+\cdots+A_\lambda=0,\\
 \vdots\\
 c^\mu+C_1c^{\mu-1}+\cdots+C_\mu=0
\end{gathered}
\]
を満たす。それらの根を
$a,a'\cdots a^{(\lambda-1)}\cdots c,c'\cdots c^{(\mu-1)}$ と書く。いま根のすべての組合せにわたって積
\[
 g(z)=f(z;a\cdots c)f(z;a'\cdots c')\cdots
      f(z;a^{(\lambda-1)}\cdots c^{(\mu-1)})
\]
を作れば、$g(z)$ は最高係数が 1 で、残りの係数が $[\mathfrak G_\eta,\psi]$ の元である。したがって $z$ は $\{\mathfrak G_\eta,\psi\}$ に属する。\footnote{このように、代数的整性を\emph{ある}方程式によって定義すれば、I と V の証明において基礎体に関する既約性の概念は不要になる。§ 2 も参照。そこでは IV の証明のためだけに、任意の方程式による代数的整性の定義から既約方程式へ移る補助命題が必要であった。この補助命題は現在の場合には一般に成り立たないので、IV は関数 $\psi$ に特別条件として課さなければならない。}

したがって領域 $\{\mathfrak G_\eta,\psi\}$ は性質 I と V をもつ。また $\mathfrak G_\eta$ がその領域の因子であるから、II と III ももつ。IV が成り立つかどうかは $\psi$ の選び方に依存する。例えば $\psi=\frac{1}{n\cdot\eta}$ のときは $\frac1n$ が領域に属するため、IV は成り立たない。しかし § 2 により、$\psi=\frac1\eta$ のとき、また $\psi=\frac{\eta}{n}$ のときは成り立つ。後者の場合、$\{\mathfrak G_\eta,\psi\}$ の関数が代数的数を表すなら、その値は $\eta=0$ としても変わらない。したがってそれは $\mathfrak G_\eta$ の関数へ移り、従って代数的整数へ移る。

ゆえに、$\psi$ に対し、代数的整添加によって代数的に分数的な数が領域に入らないという条件を課すと、$\{\mathfrak G_\eta,\psi\}$ は領域 $\mathfrak H$ となる。

まったく同様に、$\eta$ の有理関数または代数関数から成る任意の系 $\mathfrak S$ の $\mathfrak G_\eta$ への有理的整添加、また代数的整添加を定義する。有理的整添加は、一度に有限個の $\mathfrak G_\eta$ と $\mathfrak S$ の元の整多項式全体から成る整域 $[\mathfrak G_\eta,\mathfrak S]$ に導く。代数的整添加は、$[\mathfrak G_\eta,\mathfrak S]$ に代数的に整に依存するすべての元から成る領域 $\{\mathfrak G_\eta,\mathfrak S\}$ を与える。上とまったく同様に、この整域 $\{\mathfrak G_\eta,\mathfrak S\}$ は代数的整閉であり、また II と III も満たすことが示される。したがって、上と同様に次を得る。

$\mathfrak S$ に、代数的整添加によって代数的に分数的な数が領域に入らないという条件を課すと、すなわち $\mathfrak S$ が可容系となると、$\{\mathfrak G_\eta,\mathfrak S\}$ は領域 $\mathfrak H$ となる。\footnote{より一般に、同様に次が示される。$\mathfrak T$ を任意の領域とし、$\{\mathfrak T\}$ を $\mathfrak T$ に代数的に整に依存する元全体とすれば、$\{\mathfrak T\}$ は代数的整閉である。}

いま重要なのは逆も成り立つことである。すなわち、$\mathfrak S$ がすべての可容系を走るとき、領域 $\{\mathfrak G_\eta,\mathfrak S\}$ は実際に $\mathfrak M(H)$ のすべての領域 $\mathfrak H$ を尽くす。実際、$\mathfrak H$ を $\mathfrak M(\mathfrak H)$ の任意の領域とし、$\mathfrak L=\mathfrak H-\mathfrak G_\eta$ を $\mathfrak H$ のうち $\mathfrak G_\eta$ に属さないすべての関数から成る系とする。V により、$\mathfrak H$ は領域 $\{\mathfrak G_\eta,\mathfrak L\}$ を因子として含む。しかし $\{\mathfrak G_\eta,\mathfrak L\}$ も $\mathfrak H$ によって割り切られる。なぜなら、それは $\mathfrak G_\eta$ と $\mathfrak L$ を含み、これらは共通元をもたないので、$\mathfrak H$ も含むからである。したがって $\mathfrak H=\{\mathfrak G_\eta,\mathfrak L\}$ である。また IV が $\mathfrak H$ に対して成り立つので、$\mathfrak L$ はもちろん可容系である。\footnote{$\mathfrak L$ の部分系の添加ですでに十分なこともある。例えば § 2 の関数式領域は、多項式を除くすべての整有理関数式を $\mathfrak G_\eta$ に代数的整に添加することによって生じる。}

$\mathsf{H}$ が可能なすべての代数的基底を走るとき、$\mathfrak M(\mathfrak H)$ は冒頭で証明したようにすべての領域 $\mathfrak H$ の全体を走る。したがって、代数的に整な超越数から成る最も一般の領域 $\mathfrak H$ の問題は、この段落の結果によって完全に答えられる。
\begin{enumerate}
\item[a)] $\mathfrak M(H)$ のすべての領域 $\mathfrak H$ の交わりは Zermelo 領域 $\mathfrak G_\eta$ によって与えられ、これはそれ自身が領域 $\mathfrak H$ である。$\mathfrak M(H)$ は交わりの形成に関して閉じている。
\item[b)] $\mathfrak M(H)$ のすべての領域 $\mathfrak H$ は、可容系 $\mathfrak S$ を $\mathfrak G_\eta$ に代数的整添加することによって生じる。$\mathsf{H}$ が可能なすべての代数的基底を走るとき、$\mathfrak M(\mathfrak H)$ はすべての領域 $\mathfrak H$ の全体を走る。\footnote{可容系 $\mathfrak S$ は次のようにして見出せる。$\eta$ のすべての代数的に分数的な関数をある整列 $\Omega$ のもとに置き、$\mathfrak S$ には、$\mathfrak G_\eta$ およびすでに認められた関数に代数的整添加すると代数的に分数的な数を生成してしまうものを除き、すべての関数を入れる。この手続きは任意の場所で打ち切ることができる。$\Omega$ がすべての整列を走り、固定された各 $\Omega$ について任意の場所で手続きを打ち切るとき、すべての可容系 $\mathfrak S$ が尽くされる。}
\end{enumerate}

\section*{§ 5. 代数的基底に基づく、整超越数から作られる最も一般の領域。}

$\mathsf{H}$ をすべての数の任意の代数的基底とし、$\mathfrak M(\mathfrak G)$ を $H$ を含み、したがって $[H]$ を含むすべての領域 $\mathfrak G$ の全体とする。$H$ が可能なすべての基底を走るとき、$\mathfrak M(\mathfrak G)$ はすべての領域 $\mathfrak G$ の全体を走る。したがって § 4 と同様に、$\mathfrak M(\mathfrak G)$ の領域 $\mathfrak G$ だけを考えればよい。

$\mathfrak M(\mathfrak G)$ の領域 $\mathfrak G$ はすべて $[\mathsf{H}]$ を因子として含む。しかし $[\mathsf{H}]$ 自身は領域 $\mathfrak G$ ではないので、§ 4 のように $[\mathsf{H}]$ が最大公因子であるとは結論できない。それにもかかわらずそうであることは、§ 3 の加群領域をわずかに一般化して得られる $\mathfrak M(\mathfrak G)$ の可算個の領域 $\mathfrak G$ によって示される。それらは実際、$[H]$ 以外に共通元をもたない。

\emph{各固定された $\nu$ について、領域 $\mathfrak G_\nu$ $(\nu=1,2,\ldots\text{ in inf.})$ を、すべての量}
\begin{equation}
z=f(\eta)+g_1(\eta)\eta_1^\nu+g_2(\eta)\eta_2^\nu+\cdots+g_t(\eta)\eta_t^\nu
\tag{1}
\end{equation}
\emph{から成るものとする。ここで $f(\eta)$ は $\eta$ のすべての整多項式を走り、$\eta_1\cdots\eta_t$ は基底量を走り、固定された各組合せ $\eta_1^\nu,\ldots,\eta_t^\nu$ に対して、関数 $g_1(\eta)\cdots g_t(\eta)$ は互いに独立に $\eta$ のすべての代数的整関数を走る。}

§ 3 と同様に、領域 $\mathfrak G_\nu$ は抽象的性質 I--IV をもつ。$\mathfrak G_\nu$ は代数的整関数だけを含み、$\nu=1$ のとき § 3 の領域と一致する。いま、領域 $\mathfrak G_\nu$ が $[H]$ の多項式 $f(\eta)$ 以外の共通元をもたないことを示す。実際、$z$ を $\eta$ の代数的整数であるが有理ではない関数とし、$[\mathsf{H}]$ に関して既約な方程式
\begin{equation}
z^\varkappa+a_1(\eta)z^{\varkappa-1}+a_2(\eta)z^{\varkappa-2}+\cdots+a_\varkappa(\eta)=0
\tag{2}
\end{equation}
を満たすものとする。ここで $\varkappa>1$ であり、$\lambda$ を多項式 $a_1(\eta)\cdots a_\varkappa(\eta)$ の次数の最大値とする。もし $z$ が $\mathfrak G_\nu$ に属するなら、量 $\xi=z-f(\eta)$ は、同じく $[H]$ に関して既約な方程式
\begin{equation}
\xi^\chi+b_1(\eta)\xi^{\chi-1}+b_2(\eta)\xi^{\chi-2}+\cdots+b_\chi(\eta)=0
\tag{3}
\end{equation}
を満たす。ここで (1) により、$b_1(\eta),b_2(\eta)\cdots b_\varkappa(\eta)$ は $\xi$ の対称関数として、最低次元の項として少なくとも $\nu,2\nu,\cdots,\varkappa\nu$ 次元の項を含む。(2) と (3) を比較すると
\[
 b_1(\eta)=a_1(\eta)+\varkappa\cdot f(\eta)
\]
を得る。したがって、量 $\xi=z+\frac{a_1(\eta)}{\varkappa}$ が
\begin{equation}
\xi^\chi+c_2(\eta)\xi^{\chi-2}+\cdots+c_\chi(\eta)=0
\tag{4}
\end{equation}
を満たすならば、
\begin{equation}
 \xi-\zeta=\frac{b_1(\eta)}{\chi};
 \qquad
 c_i(\eta)\equiv b_i(\eta)\pmod{\frac{b_1(\eta)}{\chi}},
\tag{5}
\end{equation}
である。ここで $b_1(\eta)$ は少なくとも $\nu$ 次元の項から始まる。いま $c_i(\eta)$ は、(2) と (4) の比較が示すように、(2) の係数 $a_i(\eta)$ について次数 $\chi$ であり、したがって $\eta$ について高々次数 $\chi\lambda$ である。一方、すべての $b_i(\eta)$ は少なくとも $\nu$ 次元の項から始まる。ゆえに $\nu>\chi\lambda$ と選ぶと、(5) は
\[
 c_i(\eta)=0;\qquad
 \xi^\varkappa=0\quad\text{または}\quad
 \left(z+\frac{a_1(\eta)}{\varkappa}\right)^\varkappa=0
\]
を与える。したがって $z$ は仮定に反して $\eta$ の有理関数となる。ゆえに次が成り立つ。

\emph{$z$ が $\eta$ の代数的整数で非有理な関数であれば、$z$ は $\nu>\chi\lambda$ であるいかなる領域 $\mathfrak G_\nu$ にも属することができない。すなわち、}

\emph{$\mathfrak M(\mathfrak G)$ のすべての領域 $\mathfrak G$ の交わりは整域 $[H]$ によって与えられ、$[\mathsf{H}]$ 自身は領域 $\mathfrak G$ ではない。したがって $\mathfrak M(\mathfrak G)$ は交わりの形成に関して閉じていない。}

したがって、代数的基底による最も一般の領域 $\mathfrak G$ の構成も、領域 $\mathfrak H$ の構成より見通しにくい。領域 $[\mathsf{H}]$ は性質 I, III, IV をもつ。任意の系 $\mathfrak S$ を有理的に整に添加すると、得られる整域 $[\mathsf{H},\mathfrak S]$ については I と III が保たれ、IV は $\mathfrak S$ に課される特別条件となる。領域 $\mathfrak G$ が生じるためには、II も $\mathfrak S$ に対する条件として課されなければならない。別の言い方をすれば、$(\mathsf{H})$ 上有限な任意の体 $\mathfrak K$\footnote{$(\mathsf{H})$ は、代数的数係数の $\eta$ のすべての有理関数から成る体を意味する。}に対し、$\mathfrak K$ 上有限な $\mathfrak K$ 上の体 $\mathfrak L$ が存在し、$\mathfrak L$ のすべての関数が $[\mathsf{H},\mathfrak S]$ の関数の商として表されなければならない。\footnote{§ 7 を参照。}この条件は § 4 の可容系に対する唯一の条件よりもはるかに扱いにくい。そこで、いま有理基底へ移る。

\section*{§ 6. 有理基底に基づく、整超越数から作られる最も一般の領域。}

$\mathfrak G$ を整超越数から成る任意の領域とする。$\mathfrak G$ の元を整列し、超限帰納法によって、有限個の先行する元によって有理的に表すことのできない元を $\mathfrak G$ から選ぶ。こうして得られる系を $\Theta$ と書く。$\mathfrak G$ のすべての元は $\Theta$ によって有理的に表される。実際、$z_0$ がそうでない最初の元であったなら、$z_0$ は $\mathfrak G$ の有限個の先行する元によって有理的に表されるはずであり、それらを帰納法の仮定により\footnote{Zermelo, loc. cit., § 1 を参照。}$\vartheta$ の有理関数で置き換えれば、$z_0$ の $\vartheta$ による表示が得られて矛盾する。また構成により、$\Theta$ のどの元も有限個の先行する元によって有理的に表すことはできない。したがって $\Theta$ は $\mathfrak G$ の有理基底である。

II により、すべての数は $\mathfrak G$ の二つの元の商であるから、すべての数は $\Theta$ によって有理的に表される。すなわち $\Theta$ は同時にすべての数の有理基底である。\footnote{有理関数とは常に、係数が代数的数であるものを意味する。この代数的数係数は III と IV のため定義に含められている。本来なら有理数係数を許し、その後 III と IV によって代数的数係数を導入する方が自然である。§ 10 で扱う素標数の場合には、代数的数係数を素体で置き換えなければならない。}III と I により、$\mathfrak G$ は $\Theta$ のほかに、$\vartheta$ のすべての整多項式から成る整域 $[\Theta]$ を因子として含む。一方、IV により $[\Theta]$ は代数的に分数的な数を含まない。したがって、すべての数の有理基底 $\Theta$ は、$\Theta$ から得られる整域 $[\Theta]$ が代数的に分数的な数を含まないという付帯条件を満たす。\footnote{これが実際に条件であることは、例えば、二つの量 $\eta,\frac1{2\sqrt{\eta}}$ は基底元として許されないが、$\eta,\frac1{\sqrt{\eta}}$ は許されることから分かる。} $\Theta$ はすべての数の\emph{付帯条件付き有理基底}である。したがって § 2 と同様に、\emph{すべての領域 $\mathfrak G$ は、付帯条件付き有理基底によって、$\mathfrak G$ が整域 $[\Theta]$ を因子として含むように構成できる。}

$\Theta$ をすべての数の任意の付帯条件付き有理基底とする。その存在は領域 $\mathfrak G$ の存在により保証されている。$\mathfrak N(\mathfrak G)$ を、$\Theta$ を含み、したがって $[\Theta]$ を含むすべての領域 $\mathfrak G$ の全体とする。$\Theta$ が付帯条件付きの可能なすべての基底を走るとき、上により $\mathfrak N(\mathfrak G)$ はすべての領域 $\mathfrak G$ の全体を走る。したがって再び $\mathfrak N(\mathfrak G)$ の領域 $\mathfrak G$ だけを考えればよい。

いま $[\Theta]$ 自身は領域 $\mathfrak G$ である。実際、すべての数は $\Theta$ によって有理的に表されるので、$[\Theta]$ の二つの多項式の商として表される。またその構成により、$[\Theta]$ は残りの条件 I, III, IV も満たす。\emph{したがって $[\Theta]$ は $\mathfrak N(\mathfrak G)$ のすべての領域 $\mathfrak G$ の最大公因子、すなわち交わりである。さらに $\mathfrak N(\mathfrak G)$ は交わりに関して閉じている。すなわち、$\mathfrak N(\mathfrak G)$ の任意の部分系のすべての領域 $\mathfrak G$ の交わりも $\mathfrak N(\mathfrak G)$ に属する。} I は交わりに対して成り立つ。II と III は $[\Theta]$ が因子であることから従い、IV は交わりの領域がある領域 $\mathfrak G$ に因子として含まれることから従う。

$\vartheta$ の有理関数から成る任意の系 $\mathfrak S$ を $[\Theta]$ に有理的に整に添加すると、整域 $[\Theta,\mathfrak S]$ が得られる。ここで $\vartheta$ の代数関数は、有理基底の性質により、ある他の $\vartheta$ によって有理的に表すことができる。この整域は I, II, III をもち、したがって $\mathfrak S$ が「可容系」、すなわち有理的整添加によって代数的に分数的な数が領域に入らないという条件を満たすなら、領域 $\mathfrak G$ となる。逆に、$\mathfrak N(\mathfrak G)$ の任意の領域 $\mathfrak G$ は、$\mathfrak S$ を残余系 $\mathfrak G-[\Theta]$ とすれば、表示 $[\Theta,\mathfrak S]$ を許す。したがって、最も一般の領域 $\mathfrak G$ について、§ 4 と完全に類似して次が成り立つ。
\begin{enumerate}
\item[a)] \emph{$\mathfrak N(\mathfrak G)$ のすべての領域 $\mathfrak G$ の交わりは整域 $[\Theta]$ によって与えられ、$[\Theta]$ 自身が領域 $\mathfrak G$ である。$\mathfrak N(\mathfrak G)$ は交わりの形成に関して閉じている。}
\item[b)] \emph{$\mathfrak N(\mathfrak G)$ のすべての領域 $\mathfrak G$ は、可容系 $\mathfrak S$ を $[\Theta]$ に有理的整添加することによって生じる。$\Theta$ が付帯条件付きの可能なすべての有理基底を走るとき、$\mathfrak N(\mathfrak G)$ はすべての領域 $\mathfrak G$ の全体を走る。}\footnote{最も一般の可容系 $\mathfrak S$ については、§ 4 で述べたことが「代数的」を「有理的」に置き換えて成り立つ。}
\end{enumerate}

\emph{注。} $\Theta$ から $\Theta$ のすべての量の代数的基底 $H$ を抜き出すと、$\mathsf{H}$ は同時にすべての数の代数的基底である。逆に、任意の代数的基底は、整列において $\mathsf{H}$ を最初に置くことにより有理基底へ拡張できる。したがって、§ 5 の代数的基底による領域 $\mathfrak G$ の構成は次のように解釈できる。$[\mathsf{H}]$ に添加すべき系 $\mathfrak S$ を二つの部分系 $\mathfrak S_1$ と $\mathfrak S_2$ に分解する。$\mathfrak S_1$ の添加は、$\mathsf{H}$ を付帯条件付き有理基底へ拡張することに相当し、その基底に可容系 $\mathfrak S_2$ をさらに有理的に整に添加する。

\section*{§ 7. 付帯条件付きの最も一般の有理基底の構成。}

前節の結果にはなお補足が必要である。そこでは、すべての数に対する付帯条件付き有理基底の\emph{存在}は示されたが、与えられた整列からそれを構成する方法は示されていなかった。付帯条件のため、この構成は、領域 $\mathfrak G$ ではなく全複素数体が与えられると、より複雑な形をとる。\footnote{§ 6 b) では、$\Theta$ を、代数的基底とその基底量の代数的整関数から成る特殊な有理基底だけに走らせても十分である。その場合、付帯条件は自動的に満たされる。これを見るには、§ 6 により $\mathfrak G$ からの代数的基底を有理基底へ拡張し、得られた基底において、$\eta$ の各代数的に分数的な関数に、適当に選んだ $\eta$ の多項式を掛ければよい。すると基底性質は保たれる。この構成は全複素数体へそのまま移る。しかし、§ 6 から生じる付帯条件付きの最も一般の有理基底に関する問題は、それ自体としても興味をもつ。}

全複素数体を整列 $\Omega$ にかける。すると三種類の量が現れうる。
\begin{enumerate}
\item[1.] 量 $y$。これは、先行する基底量（3. で再帰的に定義される）への有理的整添加により、\footnote{帰納法の議論が示すように、このような再帰的定義は可能である。}得られる整域に代数的に分数的な数が入ることを強制するものである。基底量が一つも先行しない場合、この領域は $y$ のすべての整多項式から成る。特に、すべての代数的に分数的な数は量 $y$ である。なぜなら得られる整域は量 $y=\beta$ を含むからである。
\item[2.] 量 $z$。これは性質 1 をもたないが、$y$ とは異なる有限個の先行する数によって有理的に表せるものである。したがって $z=\varphi(\zeta_1\cdots\zeta_t)$ という関係を満たす。ここで $\varphi$ は $\mathsf{K}$ 係数の有理関数であり、$\zeta$ の中に $y$ はない。特に、すべての代数的整数 $\alpha$ はここに属する。なぜなら性質 1 はそれらに対して成り立たず、関係 $z=\alpha$ を満たすからである。
\item[3.] 量 $\vartheta$。これらは 1 も 2 も成り立たない量で、基底量と呼ばれる。特に、整列における最初の超越数 $\vartheta_0$ はそのような基底量である。なぜなら $\vartheta_0$ の整多項式は代数的に分数的な数を表せず、また $\vartheta_0$ は先行する代数的数に有理的に依存することもできないからである。
\end{enumerate}

\emph{いま、すべての基底量 $\vartheta$ の系 $\Theta$ が、すべての数に対する付帯条件付き有理基底であることを示す。}

量 $y$ は $\vartheta$ の中に現れないので、整域 $[\Theta]$ は代数的に分数的な数を含むことができない。したがって付帯条件は満たされる。さらに、量 $z$ も $\vartheta$ の中に現れないので、各 $\vartheta$ は先行する基底量に有理的に独立である。帰納法の議論は § 6 とまったく同様に、関係 $z=\varphi(\zeta)$（ただし $\zeta$ の中に $y$ はない）が常に関係 $z=\psi(\vartheta)$ を伴うことを示す。したがってすべての量 $z$ は $\vartheta$ によって有理的に表せる。

残るのは、より困難な主張、すなわち $y$ も $\vartheta$ によって有理的に表せることを証明することだけである。まず $y$ が $\vartheta$ に代数的に依存することを示す。仮定により、各 $y$ について、少なくとも一つの代数的に分数的な数 $\beta$ が存在し、表示
\begin{equation}
\beta=f_0(\vartheta)+f_1(\vartheta)y+\cdots+f_\varkappa(\vartheta)y^\chi
\tag{1}
\end{equation}
をもつ。ここで $f_0(\vartheta)\cdots f_\varkappa(\vartheta)$ は $[\Theta]$ の多項式であり、したがって代数的に分数的な数を表せない。もし $y$ が $\vartheta$ に代数的に独立であったなら、(1) は $y$ に関して恒等的に成り立ち、$\beta=f_0(\vartheta)$ となり、$[\Theta]$ の性質に反する。

代数的依存から有理的依存を導くため、$\Theta$ から $\Theta$ のすべての量の代数的基底 $H$ を抜き出す。すると各 $y$ は $\vartheta$ の代数関数として $\eta$ にも代数的に依存する。したがって $y=y(\eta)$ によって、$(H)$ 上の代数体 $(\mathsf{H},y(\eta))$ が定まる。すべての $\eta$ は同時に量 $\vartheta$ でもあるから、有理的依存の証明は、このような各体に対し、$\vartheta$ の有理関数である原始元が存在することを示すことに帰着する。より正確には、$y(\eta)$ に $\eta$ の多項式を掛けると性質 1 を失い、したがってそのような元になることを示す。

$(\mathsf{H})$ と $(\mathsf{H},y)$ の間には有限個の中間体しかないことが知られている。そのうち、$\vartheta$ によって有理的に表せる原始元をもつものを $\Omega_0=(H),\Omega_1,\ldots,\Omega_\sigma$ とする。つまり体の各元が $\vartheta$ の有理関数となるものを列挙する。この条件は少なくとも常に $\mathfrak L_0=(\mathsf{H})$ について満たされる。$\mathfrak L_i$ に属する原始関数を $\frac{g_i(\vartheta)}{h_i(\vartheta)}$ とし、$g_i,h_i$ は $[\Theta]$ の多項式とする。すると、$\alpha_i$ を適当に選ばれた整数として、$[\Theta]$ の多項式
\[
 \zeta_i=g_i(\vartheta)+\alpha_i h_i(\vartheta)
\]
は $\mathfrak L_i$ の原始関数、または $\mathfrak L_i$ 上の代数体の原始関数であり、\footnote{§ 5 の終わりを参照。}$\Omega_i$ の任意の関数 $\psi(\eta)$ は表示
\begin{equation}
\psi(\eta)=
\frac{a_1(\eta)\zeta_i^{\varkappa-1}+a_2(\eta)\zeta_i^{\varkappa-2}+\cdots+a_\varkappa(\eta)}
     {a(\eta)}
=\frac{A(\eta,\zeta_i)}{a(\eta)}
\tag{2}
\end{equation}
をもつ。ここで $a(\eta),a_1(\eta),\ldots,a_\kappa(\eta)$ は $\eta$ の整多項式であり、したがって $A(\eta,\zeta_i)$ と $a(\eta)$ は $[\Theta]$ の量である。

$y$ が $\mathfrak L_i$ に関して満たす既約方程式は、(2) により
\[
 f^{(i)}(\eta)y^{\lambda_i}+f_1^{(i)}(\eta,\zeta_i)y^{\lambda_i-1}
+\cdots+f_{\lambda_i}^{(i)}(\eta,\zeta_i)=0
\]
の形であり、この方程式のすべての係数は $[\Theta]$ に属する。すると関数 $y_i=f^{(i)}(\eta)\cdot y$ は、$\mathfrak L_i$ に関して同じく既約な方程式
\[
 y_i^{\lambda_i}+f_1^{(i)}(\eta,\zeta_i)y_i^{\lambda_i-1}
+\cdots+\bigl(f^{(i)}(\eta)\bigr)^{\lambda_i-1}
 f_{\lambda_i}(\eta,\zeta_i)=0
\]
を満たし、それによって $[\mathsf{H},\zeta_i]$ に代数的に整に依存する。同じことは $y_i$ に $\eta$ の任意の多項式を掛けたものについても成り立つ。特に、関数
\begin{equation}
Y=f^{(0)}(\eta)f^{(1)}(\eta)\cdots f^{(\sigma)}(\eta)\cdot y
\tag{3}
\end{equation}
は、すべての領域 $[\mathsf{H},\zeta_i]$ に代数的に整に依存し、それぞれ $\mathfrak L_i$ に関して既約な方程式
\begin{equation}
Y^{\lambda_i}+F_1^{(i)}(\eta,\zeta_i)Y^{\lambda_i-1}
+\cdots+F_{\lambda_i}^{(i)}(\eta,\zeta_i)=0
\qquad (i=0,1,\ldots,\sigma)
\tag{4}
\end{equation}
によってそうである。いま $Y$ が、求める、$\vartheta$ によって有理的に表せる原始元であることを示す。例えば、代数的数 $\gamma$ が $Y$ と $[\Theta]$ によって
\begin{equation}
\gamma=k_0(\vartheta)+k_1(\vartheta)Y+\cdots+k_\nu(\vartheta)Y^\nu
\tag{5}
\end{equation}
と表されると仮定する。ここで $k_0(\vartheta),\ldots,k_\nu(\vartheta)$ は $[\Theta]$ の多項式である。いま、(5) の係数領域から生じる体 $\mathfrak K=(H;k_0(\vartheta),\ldots,k_\nu(\vartheta))$ に関して $Y$ が満たす既約方程式を決定する。それは、$(\mathsf{H})$ と $(\mathsf{H},Y)$ の間にある $\mathfrak K$ と $(\mathsf{H},Y)$ の交体に関する $Y$ の既約方程式によって与えられることが知られている。ところが (3) により $(\mathsf{H},Y)$ は $(\mathsf{H},y)$ と同一である。また $\mathfrak K$ のすべての元、したがって交体のすべての元は有限個の $\vartheta$ によって有理的に表せるので、この交体は必ず $\mathfrak L_0,\mathfrak L_1,\cdots,\mathfrak L_\sigma$ の一つ、例えば $\mathfrak L_\tau$ である。したがって求める既約方程式は (4) により次数 $\lambda_\tau$ であり、
\begin{equation}
Y^{\lambda_\tau}+F_1^{(\tau)}(\eta,\xi_\tau)Y^{\lambda_\tau-1}
+\cdots+F_{\lambda_\tau}^{(\tau)}(\eta,\xi_\tau)=0
\tag{6}
\end{equation}
である。その係数は $[\Theta]$ の多項式である。(6) により、(5) は
\begin{equation}
\gamma=K_0(\vartheta)+K_1(\vartheta)Y+\cdots+
K_{\lambda_\tau-1}(\vartheta)Y^{\lambda_\tau-1}
\tag{7}
\end{equation}
の形に書き直せる。ここで $K_0(\vartheta),K_1(\vartheta),\cdots$ は $\mathfrak K$ に属し、他方で $[\Theta]$ の多項式である。しかし $\gamma$ は代数的数であるから、関係 (7) は $Y$ において次数 $\lambda_\tau$ に達しておらず、したがって $Y$ に関して\emph{恒等的に}満たされる。ゆえに
\[
 \gamma=K_0(\vartheta)
\]
であり、すなわち $\gamma$ は $[\Theta]$ に属し、したがって代数的整数である。

$\gamma$ が $Y$ と $[\Theta]$ によって表されるすべての代数的数を走るとき、$(\mathsf{H},Y)$ と対応する体 $\mathfrak K$ との交体は常に $\mathfrak L_0,\mathfrak L_1,\cdots,\mathfrak L_\sigma$ だけを走り、上の議論はすべて有効である。すなわち、\emph{$Y$ を $[\Theta]$ に有理的整添加しても、代数的に分数的な数は表されない。} したがって $Y$ はもはや性質 1 をもたない。2) または 3) により、$Y$ は量 $z$ または $\vartheta$ であり、したがって有限個の $\vartheta$ によって有理的に表される。(3) により同じことは $y$ にも成り立つ。\emph{$\Theta$ がすべての数の有理基底であることが証明された。} 逆に、付帯条件付きの任意の有理基底は、整列において $\Theta$ を先に置くことにより 1), 2), 3) に従って構成できる。したがって次が証明された。\emph{1), 2), 3) によって与えられた構成は、すべての数に対する付帯条件付きの最も一般の有理基底に導く。}

\section*{§ 8. 領域 $\mathfrak G$ および $\mathfrak H$ の類への分割。\srcfn{*)}{E. Steinitz, \emph{Algebraische Theorie der Körper}, Crelles Journal 137 (1910) を参照。とくに \S\S\,23, 24。（\S\,23 では代数的基底の存在がすでに証明されている。）}}

最も一般の領域 $\mathfrak G$ および $\mathfrak H$ の問題と並んで、さらに\emph{本質的に異なる}領域の問題が生じる。これは、以下でより正確にする意味で、同じ構成原理によって生成できない領域のことである。このことは、これらの領域を類へ分割することへ導く。

\emph{二つの領域の元を一対一に対応させることができ、しかも一方の領域の任意の二元の和および積が、対応する他方の元の和および積に常に対応するならば、その二つの領域は同じ類に属する、または同型であるという。} この定義から、任意の同型関係のもとで零元と単位元はそれ自身に対応し、有限個の元の間のすべての代数関係は対応する元について保たれ、逆も成り立つことが従う。特に単位元のほか、すべての有理数はそれ自身に対応する。一方、代数的数全体、同様に代数的整数全体はそれ自身へ移るが、個々の代数的数は共役に対応してもよい。

まず、固定された代数的基底 $H$ を含むすべての領域 $\mathfrak G$ の全体 $\mathfrak M(\mathfrak G)$ からの任意の領域は、別の代数的基底 $\mathsf{H}^*$ に属する全体 $\mathfrak M^*(\mathfrak G)$ からのある領域に同型であることに注意する。実際、全複素数体は任意の代数的基底から代数拡大によって生じ、したがって基底と同じ濃度をもつ。\footnote{Steinitz, §§ 23, 24.} したがって $\mathsf{H}$ と $\mathsf{H}^*$ は同じ濃度をもち、求める同型写像は基底量 $\eta_i$ を $\eta_i^*$ に対応させることで得られる。\footnote{もちろんこの結論は、二つの有理基底 $\Theta$ と $\Theta^*$ に属する集合 $\mathfrak N(\mathfrak G)$ と $\mathfrak N^*(\mathfrak G)$ には成り立たない。また、$\Theta$ と $\Theta^*$ の相互の有理的表現可能性から $\mathfrak N(\mathfrak G)$ と $\mathfrak N^*(\mathfrak G)$ の同型を推論することもできない。しかし、体 $(\Theta)$ と $(\Theta^*)$ の間の同型関係は得られる。} 他方、$\mathfrak M(\mathfrak G)$ は、§§ 2, 3, 5 の例から従うように、同型でない、すなわち本質的に異なる領域 $\mathfrak G$ を含む。実際、すべての代数関係は同型写像のもとで保たれるから、代数的基底は再び代数的基底に対応しなければならない。しかし §§ 2, 3 は、関数式領域および加群領域が、Zermelo 領域のすべての基底性質を満たすような代数的基底をもたないことを示した。§ 2 の関数式領域は代数的整閉であるから、部分集合 $\mathfrak M(\mathfrak H)$ も本質的に異なる領域 $\mathfrak H$ を含む。

いま、基底構成と同様に、$\mathfrak M(\mathfrak G)$ から部分集合 $\mathfrak L(\mathfrak G)$ を取り出す方法を示す。この部分集合の中のすべての領域 $\mathfrak G$ は本質的に異なり、一方で $\mathfrak M(\mathfrak G)$ の任意の領域、したがって任意の領域 $\mathfrak G$ は $\mathfrak L(\mathfrak G)$ のある領域に同型である。$\mathfrak M(\mathfrak G)$ を整列 $\Omega$ にかける。すると $\mathfrak M(\mathfrak G)$ の領域 $\mathfrak G$ は、先行するもののいずれかに同型である場合も、そうでない場合もある。先行するどの領域にも同型でない領域 $\mathfrak G$ は集合 $\mathfrak L(\mathfrak G)$ をなし、構成により、これは本質的に異なる領域 $\mathfrak G$ だけを含む。帰納法の議論はまた、$\mathfrak M(\mathfrak G)$ の任意の領域が $\mathfrak L(\mathfrak G)$ のある領域に同型であることを示す。もしこれが失敗する最初の領域が $\mathfrak G_1$ であったなら、$\mathfrak G_1$ は $\mathfrak L(\mathfrak G)$ に属するか、または先行する領域 $\mathfrak G_0$ に同型である。後者は仮定により $\mathfrak L(\mathfrak G)$ のある領域に同型である。したがって $\mathfrak G_1$ についても同じことが成り立つはずである。任意の領域 $\mathfrak G$ は何らかの集合 $\mathfrak M^*(\mathfrak G)$ に属し、したがって $\mathfrak M(\mathfrak G)$ のある領域に同型であるから、\emph{領域 $\mathfrak G$ の類全体は $\mathfrak L(\mathfrak G)$ によって代表される。}

$\mathfrak L(\mathfrak G)$ の領域 $\mathfrak G$ において、基底量 $\eta$ を別の基底の $\eta^*$ で置き換えると、上で示したように $\mathfrak G$ に同型な領域が生じる。逆に、$\mathfrak G^*$ を任意の領域とし、それが $\mathfrak L(\mathfrak G)$ の領域 $\mathfrak G$ に同型であると仮定する。基底量 $\eta$ が $\mathfrak G^*$ の量 $\eta^*$ に対応するなら、これらは再び代数的基底をなし、$\mathfrak G^*$ は、$\mathfrak G$ が $\eta$ の関数として含むのと同じ代数関数、またはその共役を、$\eta^*$ の関数として含む。\emph{したがって $\mathfrak G$ に同型な任意の領域は、代数的基底 $H$ を $\mathfrak G$ の中、および $(\mathsf{H})$ に関するその共役の中で可能なすべての代数的基底に走らせることによって生じる}\footnote{$\mathfrak G$ に共役なこれらの領域の元は、II により確かに $\mathfrak G$ の元によって有理的に表せるが、必ずしも整かつ有理的に表せるわけではない。したがって共役は $\mathfrak G$ と異なり得る。}\emph{---それらが $\mathfrak G$ と異なる場合には。} $\mathfrak L(\mathfrak G)$ の固定された各領域 $\mathfrak G_i$ から、このようにして類 $\mathfrak K_i$ が得られる。これは $\mathfrak G_i$ に同型なすべて、かつそれだけの領域 $\mathfrak G$ を含む。ただし、各領域が何度も、あるいは無限にしばしば現れる可能性がある。$\mathfrak R_i$ からは、上の手続きにより、$\mathfrak G_i$ に同型な各領域を一度だけ含む部分集合 $\mathfrak K_i^*$ を再び取り出すことができる。$\mathfrak K_i^*$ がすべての類を走ると、すべての領域 $\mathfrak G$ の全体が、各領域一度かぎりで得られる。類 $\mathfrak K_i$ を特徴づけるには、もちろん $\mathfrak G_i$ の代わりに、その類の任意の他の領域を用いることができ、その領域から上のように $\mathfrak K_i$ のすべての領域を導ける。要約すると次の通りである。

\emph{$\mathfrak M(\mathfrak G)$ から部分集合 $\mathfrak L(\mathfrak G)$ を取り出すことができ、その中の領域 $\mathfrak G$ は類全体と一対一に対応する（重複なし）。$\mathfrak G_i$ から、対応する類 $\mathfrak K_i$ は、$H$ を $\mathfrak G_i$ の中および $(H)$ に関するその共役の中で可能なすべての代数的基底に走らせることによって生じる。$\mathfrak K_i^*$ が $\mathfrak K_i$ の互いに異なるすべての領域を表し、$\mathfrak K_i^*$ がすべての類を走ると、すべての領域 $\mathfrak G$ の全体が各領域一度かぎりで生じる。}\footnote{ここで $\mathfrak M(\mathfrak G)$ は §§ 6, 7 に従って構成されるべきである。§ 7 により代数的基底 $\mathsf{H}$ を、それから生じるすべての付帯条件付き有理基底 $\Theta$ へ拡張し、各 $\Theta$ に対して § 6 b) により、$\Theta$ を含むすべての領域 $\mathfrak G$ の全体 $\mathfrak N(\mathfrak G)$ を形成する。}

代数的に整な超越数から成る領域 $\mathfrak H$ についても類似の主張が成り立つ。

\section*{§ 9. 領域 $\mathfrak G$ の可算無限個の類の例。}

§ 4 b) と § 6 b) で与えた構成の濃度は、それぞれの注から、すべての整列の濃度、したがって連続体の濃度を $c$ とすれば $2^c$ に等しいと読み取れる。しかしそこから、一度だけ数えたすべての領域 $\mathfrak G$ の全体の濃度について、ましてすべての類の濃度について、何らかの結論を引き出すことはできない。類の数が有限でないことは、§ 5 のものに類似する可算無限個の領域 $\mathfrak G$ の例によって示す。これらはすべて本質的に異なることが分かり、したがって可算無限個の類を与える。\footnote{§ 6 の領域 $\mathfrak G$ も同様に可算無限個の類を与えるが、証明はやや複雑である。}

§ 5, (1) と類似に、領域 $\mathfrak G_\sigma$ を量
\begin{equation}
 z\equiv a_0+a_1(\eta)+\cdots+a_{\sigma-1}(\eta)
 \pmod{\mathfrak M_\sigma(\eta)}
\tag{1}
\end{equation}
の全体から成るものとする。\emph{ここで $a_i(\eta)$ は $\eta$ に関する $i$ 次元の同次形式を表す。また加群 $\mathfrak M_\sigma$ は、$\eta$ に関する $\sigma$ 次元のすべての冪積を基底として含み\footnote{もちろん、個々の $z$ の中では、加群において有限個の $\eta$ の冪積だけが現れる。}、$\eta$ のすべての代数的整関数を乗数として含む。} 二つの領域 $\mathfrak G_\sigma$ と $\mathfrak G_\tau$ $(\sigma\ne\tau)$ の間の同型関係のもとで、加群 $\mathfrak M_\sigma$ と $\mathfrak M_\tau$ は必ず同型に対応しなければならないことを示す。すると不可能性は容易に従う。

$\mathfrak G_\tau$ の不定元を $\zeta$ と書く。同型関係のもとで、これら $\zeta$ に $\mathfrak G_\sigma$ の量 $f(\eta)$ が対応するとする。同型関係は商を取ることのもとで保たれるので、$\xi$ のすべての代数的整関数から成る領域 $\mathfrak G_\zeta$ は、$f(\eta)$ 上代数的に整なすべての関数から成る代数的整閉領域 $\mathfrak T$ に対応し、したがって $\eta$ においても代数的に整である。

(1) が $\mathfrak M_\tau$ の量に対応すると仮定する。すると $\mathfrak M_\tau$ の加群性により、(1) に $\mathfrak T$ の任意の量 $\psi(\eta)$ を掛けたものも $\mathfrak G_\sigma$ に属さなければならない。すなわち式で書けば
\begin{equation}
(a_0+a_1(\eta)+\cdots+a_{\sigma-1}(\eta))\cdot\psi(\eta)
\equiv b_0+b_1(\eta)+\cdots+b_{\sigma-1}(\eta)
\pmod{\mathfrak M_\sigma}
\tag{2}
\end{equation}
である。すべての $\eta$ を 0 に等しく置くと $a_0\cdot\psi(0)=b_0$ となり、(2) は
\begin{equation}
(a_0+a_1(\eta)+\cdots+a_{\sigma-1}(\eta))(\psi(\eta)-\psi(0))
\equiv c_1(\eta)+\cdots+c_{\sigma-1}(\eta)
\pmod{\mathfrak M_\sigma}
\tag{3}
\end{equation}
となる。$\psi(\eta)$ のほかに、関数
\[
 \chi_\nu(\eta)=\sqrt[\nu]{\psi(\eta)-\psi(0)}
\]
もすべての $\nu$ について $\mathfrak T$ に属する。また $\chi_\nu(0)=0$ であるから、(3) と類似の式は
\begin{equation}
(a_0+a_1(\eta)+\cdots+a_{\sigma-1}(\eta))\chi_\nu(\eta)
\equiv c_1^{(\nu)}(\eta)+\cdots+c_{\sigma-1}^{(\nu)}(\eta)
\pmod{\mathfrak M_\sigma}
\quad(\nu=1,2,\cdots)
\tag{4}
\end{equation}
である。$\chi_1=\psi(\eta)-\psi(0)$ とその $\kappa-1$ 個の共役との積を $A(\eta)$ とし、その次数を $\lambda$ とする。$\chi_1(0)=0$ なので、$A(\eta)$ は少なくとも一次元の項から始まらなければならない。(4) から
\begin{equation}
a_0\cdot\chi_\nu(\eta)\equiv0\pmod{\mathfrak M_1};\qquad
 a_0^\nu\cdot\chi_1(\eta)\equiv0\pmod{\mathfrak M_\nu};\qquad
 a_0^{\nu\kappa\cdot}A(\eta)\equiv0\pmod{\mathfrak M_{\nu\kappa}}
\tag{5}
\end{equation}
を得る。したがって $a_0\ne0$ なら、§ 3 と同様に $\lambda>\nu\kappa$ である。しかし $\nu<\frac{\lambda}{\kappa}$ は $\mathfrak T$ の代数的整閉性に矛盾するので、必ず $a_0=0$ である。いま (4) から
\begin{equation}
a_1(\eta)\cdot\chi_\nu(\eta)\equiv c_1^{(\nu)}(\eta)\pmod{\mathfrak M_2};\qquad
[a_1(\eta)]^{\nu\kappa\cdot}A(\eta)
\equiv[c_1^{(\nu)}(\eta)]^{\nu\kappa}\pmod{\mathfrak M_{\nu\kappa+1}}
\tag{6}
\end{equation}
を得る。ここから、$A(\eta)$ が少なくとも一次元の項から始まるので、$c_1^{(\nu)}(\eta)=0$ である。したがって (5) と完全に同様に、すべての $\nu$ について
\[
 a_1(\eta)\cdot\chi_\nu(\eta)\equiv0\pmod{\mathfrak M_2};\qquad
 [a_1(\eta)]^{\nu\kappa}A(\eta)\equiv0\pmod{\mathfrak M_{2\nu\kappa}}
\]
となり、従って $a_1(\eta)=0$ である。このように続けると、(5) と (6) を交互に適用して
\[
 a_0=0,\qquad a_1(\eta)=0,\quad \ldots,\quad a_{\sigma-1}(\eta)=0
\]
を得る。$\mathfrak G_\tau$ に関しても同様の考察が成り立つから、\emph{$\mathfrak G_\sigma$ と $\mathfrak G_\tau$ $(\sigma\ne\tau)$ の任意の同型関係のもとで、加群 $\mathfrak M_\sigma$ と $\mathfrak M_\tau$ は同型に対応する。}

例えば $\sigma>\tau$ と仮定する。不定元 $\zeta$ には $\mathfrak G_\sigma$ の量 $f(\eta)$ が対応していた。そして $\mathfrak G_\tau$ のすべての量は $\zeta$ の多項式 mod. $\mathfrak M_\tau$ によって表されるから、$\mathfrak M_\tau$ と $\mathfrak M_\sigma$ の対応は、$\mathfrak G_\sigma$ のすべての量が $f(\eta)$ の多項式 mod. $\mathfrak M_\sigma$ として得られることを意味する。$\sigma>\tau\geqq1$ であるから、不定元 $\eta$ は確かに $\mathfrak M_\sigma$ に属さない。したがって、$\eta$ が $f(\eta)$ によって mod. $\mathfrak M_\sigma$ で整かつ有理的に表せる以上、線形項が消えない $f(\eta)$ が存在しなければならない。そこで
\[
 \zeta:f(\eta)\equiv a_0+a_1(\eta)\pmod{\mathfrak M_2}\quad[a_1(\eta)\ne0],
\]
\[
 \xi^\tau:[f(\eta)]^\tau\equiv a_0^\tau\pmod{\mathfrak M_1}\quad\text{$a_0\ne0$ のとき},
\]
\[
 \hphantom{\zeta^\tau:[f(\eta)]^\tau}\equiv [a_1(\eta)]^\tau
 \pmod{\mathfrak M_{\tau+1}}\quad\text{$a_0=0$ のとき}
\]
とする。いま $\zeta^\tau$ は $\mathfrak M_\tau$ に属するが、$[f(\eta)]^\tau$ は 0 次または $\tau$ 次の消えない項から始まり、$\tau<\sigma$ であるため $\mathfrak M_\sigma$ に属し得ない。これは上で導いた $\mathfrak M_\sigma$ と $\mathfrak M_\tau$ の対応に反する。$\tau>\sigma$ の場合は $\eta$ と $\zeta$ を入れ替えれば同時に扱われる。したがって、領域 $\mathfrak G_\sigma$ と $\mathfrak G_\tau$ は\emph{同型でない、すなわち本質的に異なる}ことが示された。ゆえに\emph{領域 $\mathfrak G_i$ $(i=1,2,\cdots\text{ in inf.})$ は、領域 $\mathfrak G$ の可算無限個の類の例を与える。}

\emph{注。} 任意の二つの領域 $\mathfrak G_i$ は非同型であるが、二つの領域について、それぞれが他方の一部に同型であるということは十分起こり得る。\footnote{この振る舞いは体にも起こり得る。Steinitz, loc. cit., 結論を参照。} 例えば $\tau=1$ とし、$\sigma$ を任意とする。
\[
\mathfrak G_1:a_0+\mathfrak M_1(\zeta);\qquad
G_\sigma:a_0+a_1(\eta)+\cdots+a_{\sigma-1}(\eta)+\mathfrak M_\sigma(\eta).
\]
対応 $\zeta=\eta$ は $\mathfrak G_\sigma$ を $\mathfrak G_1$ の真の因子にする。逆に、$\mathfrak G_1$ と $\mathfrak G_\sigma$ の間の一対一対応 $\zeta=\eta^\sigma$, $\eta=\sqrt[\sigma]{\zeta}$ は、$\mathfrak G_1$ を $\mathfrak G_\sigma$ の真の因子にする。したがって各 $\mathfrak G_\sigma$ はそれ自身の真の因子にも同型である。交わりの概念とは対照的に、交わり類、すなわち、その各領域が他のすべての類のすべての領域の真の因子に同型であるような類、という概念は存在しない。少なくとも一意に定まる概念としては存在しない。

\section*{§ 10. 任意の体の整な量。}

前節までの展開は、領域 $\mathfrak G$ については、領域 $\mathfrak G$ の定義性質 III と IV に代数的数を含めないなら、全複素数の全体が体をなすという事実だけに依拠していた。\footnote{§ 7 でのみ、有理数の「素体」が無限個の元を含むという仮定をさらに用いた。しかし有限個の元だけをもつ素体の場合、この段落は、これから見るように、単に消える。} 領域 $\mathfrak H$ については、さらにこの体が代数的閉であるという事実に依拠していた。したがって得られた結果は、Weber および E. Steinitz\footnote{loc. cit., 序論。} に従って「二つの演算、加法と乗法をもつ元の系であり、それらが結合律と交換律に従い、分配律で結びつき、逆演算を無制限かつ一意に許すもの」\footnote{零による除法だけは除かなければならない。}として抽象的に定義される\emph{最も一般の体}へ直接一般化できる。

そのような体はすべて単位元を含み、したがって単位元から加法、乗法およびその逆演算によって導かれる「素体」を含む。この素体は、有理数の型（標数 $0$）であるか、または素数 $p$ の剰余類系の型（標数 $p$）である。\footnote{Steinitz, § 4.} そのとき\emph{素体の整な量}は、単位元から加法および減法によって生じる量としてよく定義される。標数 $p$ の素体では、整な量はすでに素体を尽くし、素体の分数的量は存在しない。任意の代数的閉体は「絶対代数体」を含む。これは素体に、素体上代数的なすべての元を添加して得られる。したがって標数 $0$ の体では、これは全代数的数の体の型である。すると\emph{絶対代数体の代数的に整な量}は、素体の整な量（または同じことだが単位元）上代数的に整なすべての量としてよく定義される。標数 $p$ の体では、したがって絶対代数体の代数的に分数的な量も存在しない。これらの準備的注意の後、「整」量および「代数的に整」量の定義を任意の体、あるいは代数的閉体に移すことができる。\emph{体 $\mathfrak R$ の整な量とは、$\mathfrak R$ の量から成る整域 $\mathfrak G$ で、その商体\footnote{すなわち、$\mathfrak G$ の量の対の商すべてから成る体。}が $\mathfrak R$ を尽くし、単位元を含むが素体の分数的量を含まないものとして定義される。}

\emph{代数的閉体 $\mathfrak R$ の代数的に整な量とは、$\mathfrak R$ の量から成る代数的整閉な整域 $\mathfrak H$ で、その商体が $\mathfrak R$ を尽くし、単位元を含むが絶対代数体の代数的に分数的な量を含まないものとして定義される。}

標数零の体については、前節までで得られた結果がそのまま適用される。ただし「代数的数」を、素体、あるいは絶対代数体の量で置き換えるだけでよい。標数 $p$ の体については、素体が分数的量を含まないため、結果はさらに単純化する。すなわち有理基底および添加される系に対する付帯条件が単に消える。\footnote{この段落の第一の注を参照。} したがって結果は次の通りである。\emph{抽象的定義は、標数 $p$ の体の整な量として、任意の有理基底から導かれるすべての整域、体自身、およびこれらの領域と体との間にある任意の整域を許す。対応して、素標数の代数的閉体の代数的に整な量としては、代数的基底から導かれる領域と体自身との間にある任意の代数的整閉な整域、二つの境界領域を含むものが得られる。} したがって素標数の最も一般の体では、対応する素体の場合と同様に、整と分数との区別は消える。

\medskip
\noindent Erlangen, 1915 年 3 月 30 日。

\clearpage
\section*{10. 同型写像の関数方程式}
\begin{center}
Math. Ann. 77 (1916), pp. 536--545
\end{center}

Dedekind によれば、\footnote{例えば Dirichlet-Dedekind, \emph{Vorlesungen über Zahlentheorie}（第4版）, Supplement XI, \S\ 161 を参照。群論にならった ``isomorphic mapping'' または ``isomorphism'' という名称は後の文献にはじめて現れる。Dedekind は「体の置換」と呼んでいる。}\emph{体} $\mathfrak A$ \emph{からある系} $\mathfrak B$ \emph{への同型写像}とは、次のような一価の対応をいう。すなわち、$\mathfrak A$ の各元に対して $\mathfrak B$ の元がただ一つ対応し、しかも $\mathfrak A$ の任意の二元の和、差、積、商には、それぞれ一意に対応する元の和、差、積、商が再び対応する、という対応である。このとき $\mathfrak B$ は後から同じく体であることが分かる。

そのような写像が、上の意味で一価な関数 $f(z)$ によって媒介されるものとする。$z$ が体 $\mathfrak A$ のすべての元 $x,y,\ldots$ を走るとき、$f(z)$ は次の関数方程式によって特徴づけられる。
\[
\begin{array}{ll}
(1)\quad f(x+y)=f(x)+f(y); & (2)\quad f(x-y)=f(x)-f(y);\\[0.45em]
(3)\quad f(x\cdot y)=f(x)\cdot f(y); & (4)\quad f\!\left(\dfrac{x}{y}\right)=\dfrac{f(x)}{f(y)}.
\end{array}
\]

\emph{以下では、} $x,y,\ldots$ \emph{がすべての実数および複素数を走るときの、これらの関数方程式の最も一般の一価解} $f(z)$ \emph{を与える。すなわち、全複素数体の同型写像を与える最も一般の関数} $f(z)$ \emph{を与える。}\footnote{この問題は Landau 氏から折にふれて私に提示された。後に気づいたことだが、解の一部はすでに H. Lebesgue, \emph{Sur les transformations ponctuelles, transformant les plans en plans}, Atti Acad. Torino, 1906/07 に現れており、また A. Ostrowski, \emph{Über einige Fragen der allgemeinen Körpertheorie}, Crelle's Journal 143 (1913), \S\ 2 にも暗黙に含まれている。} 任意の体に対する解もまったく同様に得られる。

G. Hamel\footnote{G. Hamel, \emph{Eine Basis aller Zahlen und die unstetigen Lösungen der Funktionalgleichung:} $f(x+y)=f(x)+f(y)$, Math. Ann. 60, p. 459 (1905).} は、すべての数の\emph{線形}基底によって関数方程式 (1) の最も一般の解を構成した。ここで (1) から (4) の解、したがって有理演算および代数演算が不変となる写像関数に対して与える構成は、すべての数の\emph{有理}基底および\emph{代数的}基底を用いる。1 で同型写像をより詳しく論じた後、2 で $f(z)$ に対する必要条件を与え、3 でそれらが十分条件でもあることを認め、それによって実際の構成に至る。\footnote{私の論文 \emph{Die allgemeinsten Bereiche aus ganzen transzendenten Zahlen}, Math. Ann. 77, p. 103 (1916), とくに \S\S\ 8, 10 の平行した考察を参照。} 既知の例外 $f(z)=z$ または $\overline z$ を除いて、(1) から (4) の解は不連続である。4 ではさらに、それらが「極端に不連続」であることを示す。これは G. Hamel が (1) の実解について証明した事実であるが、複素領域では必ずしも満たされない。\footnote{Hamel の用語「全く不連続」は、他の文献では別の意味で用いられている。}

\subsection*{1.}

まず、関数方程式 (1) から (4) からいくつかの帰結を引き出す。これは一般の体 $\mathfrak A$ に対して直接行う。(1) から、一価性により $f(0)=0$ が従う。$f(y)$ が消えなければ、元の $y$ も消えることはできない。したがって関数方程式 (1) から (4) は、\emph{すべての値} $f(z)$ \emph{から成る像の系} $\mathfrak B$ \emph{が再び体をなす}ことを示す。逆に、初期条件 $f(0)=0$ を課せば、(2) から関数 $f(z)$ の一価性が従う。したがって、\emph{一価性の要求と初期条件} $f(0)=0$ \emph{とは同値である}。(4) から $f(z)$ は恒等的には消えないことが従い、さらに (3) から $f(z)$ は $z=0$ に対してだけ消える。実際、$f(y)=0$ かつ $y\ne0$ ならば、$x/y$ も $\mathfrak A$ に属するので
\[
f(x)=f\!\left(\frac{x}{y}\cdot y\right)=f\!\left(\frac{x}{y}\right)\cdot f(y)=0
\]
となり、$f(z)$ が恒等的に消えることになるが、これは (4) に反する。したがって $f(x)=f(y)$ から $x=y$ が従う。すなわち $f(z)$ の各値にも $z$ のただ一つの値が対応する。\emph{$f(z)$ は} $z$ \emph{の一対一対応関数である}。関数方程式が示すように、逆関数によって媒介される写像もまた同型である。すべての有理演算は有限個の加法、乗法およびそれらの逆演算から構成されるから、次のようにも言える。

\emph{同型写像によって、$\mathfrak A$ の元と $\mathfrak B$ の元との間に一対一の関係が確立され、その関係において、$\mathfrak A$ の有限個の元の間に成り立つすべての有理関係
\[
F(x,y,\cdots,t)=0
\]
は $\mathfrak B$ において保存され、逆もまた成り立つ。}

さらに注意しておくべきことは、体に対しては、$f(z)$ はすでに関数方程式 (1) と (3)、$f(z)$ が恒等的には消えないという要求、および初期条件 $f(0)=0$ によって特徴づけられる、ということである。実際、(1) において $x$ を $\mathfrak A$ に属する元 $(x-y)$ で置き換えれば (2) が従う。そして (2) と $f(0)=0$ から $f(z)$ の一価性が得られる。上で示したように、さらに (3) において $x$ を $\mathfrak A$ に属する元 $x/y$ で置き換えれば、$y\ne0$ のとき $f(y)$ は消えないことが従う。したがって (4) の妥当性と同時に一対一性が得られる。体の代わりに整域 $\mathfrak J$ を基礎に置くならば、$x/y$ は $\mathfrak J$ に属するとは限らず、(3) から一対一性を結論することはできない。しかしここでは、次の最もよく知られた定式化で十分である。\emph{同型写像とは、$\mathfrak J$ と $\mathfrak J'$ の元の間の一対一対応であり、和と積には常に対応する元の和と積が対応するような対応である。}\footnote{第1項については Dedekind, loc. cit., \S\ 161 を参照。}

\subsection*{2.}

以下では簡単のため、$\mathfrak A$ の代わりに、すべての実数および複素数から成る体 $\mathfrak C$ を基礎に置く。まず問題となるのは、$f(z)$ がとる値の系の考察である。(3) または (4) から $f(1)=1$ が従い、さらに関数方程式によって、すべての\emph{有理数} $\alpha$ について $f(\alpha)=\alpha$ が従う。したがって、写像のもとで各有理数は自分自身に対応する。\emph{代数的数} $\beta$ は有理数係数の既約方程式 $F(\beta,\alpha)=0$ を満たすから、1 と先の注意により、$f(\beta)$ についても $F(f(\beta),\alpha)=0$ が成り立つ。したがって、\emph{各代数的数は自分自身またはその共役の一つへ移る}。また一対一性により、代数的数全体には再びすべての代数的数から成る体が対応する。

\emph{超越数}の挙動を認識し、同時に共役値を分離するために、\emph{すべての数の有理基底}\footnote{「整超越数」に関する私の引用論文 \S\ 6 を参照。}を基礎に置く。すなわち、すべての数を有理的に、つまり有理数係数で\footnote{「有理関数」とは、常に係数も有理数であるものを意味する。}表すことを可能にする数 $\vartheta$ の系 $\Theta$ であって、少なくとも一つの整列のもとで、どの基底数も有限個の先行する基底数によって有理的に表せないような系である。この基底の存在は、線形基底および代数的基底の場合とまったく同様に整列定理から従う。

$\Theta$ に整列 $\Omega$ を与える。すると各数は、有限個の先行する数によって有理的に表せるか、または先行するものから有理的に独立である。後者の数が基底 $\Theta$ をなし、帰納法により、実際にすべての数が有限個の $\vartheta$ によって有理的に表せることが分かる。

各基底数 $\vartheta$ に対して「切片体 $\mathfrak K(\vartheta)$」を定義する。これは整列において $\vartheta$ に先行する基底数のすべての有理結合から成る。最初の基底数の切片体としては、すべての有理数から成る体 $\mathfrak R$ を考える。$\vartheta$ は $\mathfrak K(\vartheta)$ に関して二種類の挙動を示し得る。すなわち、$\mathfrak K(\vartheta)$ に代数的に従属であるか、代数的に独立であるかである。1 により、$\mathfrak C$ で成り立つすべての関係は像の体でも成り立ち、逆もまた成り立つから、$\vartheta$ に対応する値 $f(\vartheta)$ の全体は像領域の基底をなし、$\Theta$ 自身の整列によって整列される。したがって、特に切片体 $\mathfrak K(\vartheta)$ には切片体 $\mathfrak K(f(\vartheta))$ が対応する。そして $\vartheta$ が $\mathfrak K(\vartheta)$ に代数的に従属または独立であるに従い、$f(\vartheta)$ も $\mathfrak K(f(\vartheta))$ に関して同じ性質をもつ。従属の場合には、$\vartheta$ と $f(\vartheta)$ の既約方程式も対応する。

それぞれ $\mathfrak K(\vartheta)$ から代数的に独立である値 $\vartheta$ の全体は、再び帰納法が示すように、\emph{すべての数の代数的基底}\footnote{E. Zermelo, \emph{Über ganze transzendente Zahlen}, \S\ 1, Math. Ann. 75 (1914) を参照。}をなす。すなわち、すべての数を代数的に表すことを可能にし、しかもそれら自身の間には代数関係が存在しないような数 $\eta$ の系 $H$ である。この代数的基底 $H$ には、像領域において再び代数的基底 $Z$ が対応する。一対一性により、この $Z$ は $H$ と同じ濃度をもつ。すなわち連続体の濃度である。代数拡大によって濃度は増加しないことが知られているからである。\footnote{E. Steinitz, \emph{Algebraische Theorie der Körper}, Crelle's Journal 137 (1910), \S\ 24 を参照。}

しかし $f(\vartheta)$ の値が分かれば、すべての値の系に対して $f(z)$ はただちに得られる。なぜなら、$z=\Psi(\vartheta)$ から常に $f(z)=\Psi(f(\vartheta))$ が従うからである。

\subsection*{3.}

いま、上で得られた必要条件が実際に $f(z)$ の構成も与えることを示す。代数的基底 $H$ に一対一に対応する基底 $Z$ を、次のように構成して $H$ に対応させる。すべての数の代数的基底 $Z'$ を与える第二の整列 $\Omega'$ を置く。$Z$ を $Z'$ の部分集合とし、その元を $\zeta$ とする。$Z$ は $Z'$ と同じ濃度、したがって $H$ と同じ濃度をもつものとする。$H$ の各元 $\eta_i$ に、同じ添字をもつ一つの $\zeta_i$ を一対一に割り当てる。その規則は、$\zeta_i$ が整列 $\Omega'$ において、$\eta_i$ に先行する $H$ のどの元にもまだ割り当てられていない $Z$ の元のうち最初のものである、というものである。

ここで $f(z)$ を次のように定義する。
\begin{enumerate}
\item[a)] すべての有理数 $\alpha$ に対して、$f(0)=0$, $f(\alpha)=\alpha$ とする。
\item[b)] 代数的基底 $H$ のすべての量に対して、$f(\eta_i)=\zeta_i$ とする。
\item[c)] 有理基底 $\Theta$ の残りのすべての量 $\vartheta$、すなわち切片体 $\mathfrak K(\vartheta)$ に代数的に従属し、したがって $\mathfrak K(\vartheta)$ で既約な方程式 $F(\vartheta;\vartheta_1,\ldots,\vartheta_r)=0$ を満たすものについては、$F(f(\vartheta);f(\vartheta_1),\ldots,f(\vartheta_r))=0$ の根として定める。
\item[d)] 残りのすべての値 $z$、すなわち $\Theta$ の定義により $z=\psi(\vartheta_{k_1},\ldots,\vartheta_{k_r})$ と表せるものについては、$f(z)=\psi(f(\vartheta_{k_1}),\ldots,f(\vartheta_{k_r}))$ とする。\footnote{(d) には (a) が特殊な場合として含まれる。}
\end{enumerate}

このように定義された $f(z)$ が実際に関数方程式 (1) から (4) を満たすことを示すには、第1項により、(1) と (3) について示せば十分である。なぜなら $\mathfrak C$ は体であり、$f(z)$ は (a) と (b) により恒等的には消えず、さらに初期条件 $f(0)=0$ を満たすからである。

有理基底 $\Theta$ によって $x$, $y$, $x+y$ を
\[
x=\psi_1(\vartheta_1,\ldots,\vartheta_t),\quad
 y=\psi_2(\vartheta_1,\ldots,\vartheta_t),\quad
 x+y=\psi_3(\vartheta_1,\ldots,\vartheta_t)
\]
と表す。このとき、$f(z)$ が関数方程式 (1)
\[
f(x)+f(y)-f(x+y)=0
\]
を満たすことの証明は、(d) により、
\[
\psi_1(\vartheta)+\psi_2(\vartheta)-\psi_3(\vartheta)=\frac{H_1(\vartheta)}{H_2(\vartheta)}=0
\]
から常に
\[
\psi_1(f(\vartheta))+\psi_2(f(\vartheta))-\psi_3(f(\vartheta))
=\frac{H_1(f(\vartheta))}{H_2(f(\vartheta))}=0
\]
が従うことと同一である。ここで $H_1,H_2$ はその引数の整有理関数を意味する。関数方程式 (3) もまったく同じ形の条件に導く。したがって、すべての関数方程式が満たされることは、次の条件が成り立つことと同一である。この条件は第1項および第2項により必要でもある。すなわち、すべての整有理関数 $H(\vartheta)$ について、$H(\vartheta)=0$ から常に $H(f(\vartheta))=0$ が従い、$H(\vartheta)\ne0$ から $H(f(\vartheta))\ne0$ が従う、という条件である。後者は分母が現れるために必要である。

これが実際に成り立つことは帰納法で示される。例えば $H(\vartheta_1,\ldots,\vartheta_n)$ において、基底量 $\vartheta_n$ が $\vartheta_1,\ldots,\vartheta_n$ のうち整列で最後のものであるとする。\footnote{$\vartheta$ について0次の式 $H(\vartheta)$ は写像によって自分自身に移るので、さらに調べる必要はない。} すると $H(\vartheta_1,\ldots,\vartheta_n)=0$ は、$\vartheta_n$ が切片体 $\mathfrak K(\vartheta_n)$ に関して満たす関係とみなせる。同様に $H(f(\vartheta_1),\ldots,f(\vartheta_n))=0$ は、$f(\vartheta_n)$ が $\mathfrak K(f(\vartheta_n))$ に関して満たす関係とみなせる。

もしすべての関係 $H(\vartheta)=0$ が $f(\vartheta)$ についても成り立ち、かつ逆も成り立つのでなければ、ある最初の基底量、例えば $\vartheta_n$ が存在して、$\mathfrak K(\vartheta_n)$ に関して成り立つ少なくとも一つの関係が像の体では保存されないか、またはその逆が起こるはずである。ただし、それ以前の切片体 $\mathfrak K(\vartheta_n)$ と $\mathfrak K(f(\vartheta_n))$ は同型である。

特に (a) により、$\Theta$ の最初の基底量の切片体、すなわちすべての有理数から成る体 $\mathfrak R$ は、その像の体と同型である。いま $H(\vartheta_n)$ が次の形であるとする。
\[
H(\vartheta_n)=\vartheta_n^{\mu}+a_1(\vartheta)\vartheta_n^{\mu-1}+\cdots+a_\mu(\vartheta),
\]
ここで $a_i(\vartheta)$ は $\mathfrak K(\vartheta_n)$ の量である。このとき二つの場合がある。

1) $\vartheta_n$ が $\mathfrak K(\vartheta_n)$ から代数的に\emph{独立}である場合。このとき $H(\vartheta_n)=0$ は $\vartheta_n$ について恒等的に満たされる。\footnote{$x,y,\ldots$ が $\vartheta$ によって表せるため、このような形の関係は十分起こり得る。} したがってすべての係数について $a_i(\vartheta)=0$ であり、$\mathfrak K(\vartheta_n)$ と $\mathfrak K(f(\vartheta_n))$ の同型性により、$a_i(f(\vartheta))=0$ も成り立つ。ゆえに $H(\vartheta_n)=0$ から常に $H(f(\vartheta_n))=0$ が従う。逆に $H(f(\vartheta_n))=0$ とする。$\vartheta_n$ は $\mathfrak K(\vartheta_n)$ から代数的に独立であるので代数的基底 $H$ に属し、したがって $f(\vartheta_n)$ は (b) により代数的基底 $Z$ に属し、$\mathfrak K(f(\vartheta_n))$ から代数的に独立である。したがってすべての $a_i(f(\vartheta))$ が消え、同型性によりすべての $a_i(\vartheta)$ が消える。ゆえに $H(f(\vartheta_n))=0$ から $H(\vartheta_n)=0$ が従い、また $H(\vartheta_n)\ne0$ から $H(f(\vartheta_n))\ne0$ が従う。

2) $\vartheta_n$ が $\mathfrak K(\vartheta_n)$ に代数的に\emph{従属}する場合。このとき $H(\vartheta_n)$ は、$\mathfrak K(\vartheta_n)$ に関して既約な方程式 $F(\vartheta_n)$ によって割り切れる。すなわち、$t$ について恒等的に
\[
H(t;a_i)=F(t;\vartheta_n)\cdot Q(t;a_i)
\]
が成り立つ。ところが現れるすべての係数は $\mathfrak K(\vartheta_n)$ に属するので、同型な体 $\mathfrak K(f(\vartheta_n))$ において対応する関係
\[
H(t;a_i(f))=F(t;f(\vartheta_n))\cdot Q(t;a_i(f))
\]
が成り立つ。しかも (c) により、$f(\vartheta_n)$ は $F(t;f(\vartheta_n))$ の根となるように選ばれていた。したがって $H(\vartheta_n)=0$ から $H(f(\vartheta_n))=0$ が従う。

逆に $H(f(\vartheta_n))=0$ であるとする。このとき、$H(t;a_i(f(\vartheta)))$ が $\mathfrak K(f(\vartheta_n))$ に関して既約な関数 $F(t;f(\vartheta_n))$ によって割り切れることから、同型な体 $\mathfrak K(\vartheta_n)$ に対して、$H(t;a_i(\vartheta))$ が $F(t;\vartheta_n)$ によって割り切れることが同様に従う。そして $F(t;f(\vartheta_n))$ は $\vartheta_n$ の既約方程式から同型関係によって生じたものであり、この関係は一対一であるから、$F(t;\vartheta_n)$ は必然的に $\vartheta_n$ を根にもつ既約関数を表す。したがって $H(f(\vartheta_n))=0$ から $H(\vartheta_n)=0$ が従い、また $H(\vartheta_n)\ne0$ から $H(f(\vartheta_n))\ne0$ が従う。

したがって、同型な切片体 $\mathfrak K(\vartheta_n)$ と $\mathfrak K(f(\vartheta_n))$ から、$\vartheta_n$ および $f(\vartheta_n)$ をそれぞれ添加することによって、やはり同型な体が生じる。これがある $\vartheta_n$ について成り立たないという仮定は矛盾に導く。

したがって、(a) から (d) によって定義された関数 $f(z)$ は、実際に関数方程式 (1) から (4) の解を与え、しかも 2 により最も一般の解を与えることが証明された。

$f(z)$ の構成に用いられたすべての考察は、すべての数の全体 $\mathfrak C$ が体をなすという事実だけを用いている。したがって、それらは任意の抽象的に定義された体 $\mathfrak A$ に対して成り立つ。ただし、$\mathfrak R$ という有理数体に対応するものは $\mathfrak A$ の「素体」である。これは、$\mathfrak A$ の単位元から加法と乗法およびそれらの逆演算によって導かれる体である。したがって (a) から (d) において有理数を素体の元で置き換えれば、そのようにして得られる関数 $f(z)$ は、\emph{任意の抽象的に定義された体の最も一般の同型写像}を与える。

\subsection*{4.}

最後に、すべての実数および複素数から成る体 $\mathfrak C$ に対して定義された関数 $f(z)$、すなわち $f(z)=z$ または $\overline z$ を除いて不連続であることが知られている関数が、\emph{極端に不連続}であることを示す。これは、任意の実数または複素数 $z_0$ の任意の近傍に、$f(z)$ が任意に指定された値 $Z_0$ にいくらでも近づくような値 $z$ が存在する、という意味である。

この極端な不連続性は、G. Hamel が前掲箇所で証明したように、すべての実数について定義された関数方程式 (1) の実解 $\varphi(x)$ に属する。しかし、下に挙げる例が示すように、複素数上で定義された解には必ずしも属さない。Hamel は次のように論じる。$\varphi(x)$ が $C\cdot x$ と異なるならば、線形基底の値\footnote{すべての（実）数の線形基底は、すべての（実）数を有理係数で線形に表せるような（実）数の系であって、有限個の基底数の間には有理係数の線形関係が存在しないような系として与えられる。}のうち少なくとも二つ、例えば $x_1,x_2$ が存在し、$\varphi(x_1):x_1\ne\varphi(x_2):x_2$ となる。このとき二つの方程式
\[
\alpha x_1+\beta x_2=a,
\qquad
\alpha\varphi(x_1)+\beta\varphi(x_2)=b
\]
は実数 $\alpha,\beta$ に解をもつ。したがって $a$ と $b$ は有理数 $\alpha,\beta$ によって任意に近似できる。一方、$\alpha,\beta$ が有理数であるため、$\alpha\varphi(x_1)+\beta\varphi(x_2)$ は実際に $\varphi(\alpha x_1+\beta x_2)$ を表す。この議論は複素値系では失敗する。実際、複素数上でも、(1) の解であって不連続ではあるが極端には不連続でないものを与えることができる。例えば
\[
\varphi(z)=\varphi(x+iy)=\varphi(x)+i\bigl(y+\varphi(x)\bigr),
\]
ここで $\varphi(x)$ は実数上で極端に不連続な実解を意味する。あるいは、さらに簡単に $\varphi(z)=\varphi(x)$ としてもよい。どちらの場合も $\varphi(z)$ の実部は極端に不連続であるが、虚部は実部によってともに決定される。

この挙動、そして同時に関数 $f(z)$ の挙動は、\emph{階数}の概念を導入すると完全に明らかになる。$z=x+iy$, $\varphi(z)=X+iY$ とおき、$\varphi(z)$ の階数を四、三、二、一と呼ぶのは、それぞれ、線形独立な関係
\[
c_1x+c_2y+c_3X+c_4Y=0
\]
が存在しない、または一つ、二つ、三つ存在する場合である。もちろん $c$ は実数である。

1) $\varphi(z)$ が階数四であるとする。このとき線形基底の値が少なくとも四つ、例えば $z_1,z_2,z_3,z_4$ 存在して、行列式
\[
\begin{vmatrix}
x_1&x_2&x_3&x_4\\
y_1&y_2&y_3&y_4\\
X_1&X_2&X_3&X_4\\
Y_1&Y_2&Y_3&Y_4
\end{vmatrix}
\]
が消えない。したがって方程式
\[
\begin{aligned}
\alpha x_1+\beta x_2+\gamma x_3+\delta x_4&=a,\\
\alpha y_1+{}\cdot{}\ \cdot{}\ \cdot{}+\delta y_4&=b,\\
\alpha X_1+{}\cdot{}\ \cdot{}\ \cdot{}+\delta X_4&=A,\\
\alpha Y_1+{}\cdot{}\ \cdot{}\ \cdot{}+\delta Y_4&=B
\end{aligned}
\]
は実数 $\alpha,\ldots,\delta$ に解をもち、$a,b,A,B$ は有理数 $\alpha,\beta,\gamma,\delta$ によって任意に近似できる。さらに
\[
\alpha(X_1+iY_1)+\beta(X_2+iY_2)+\gamma(X_3+iY_3)+\delta(X_4+iY_4)
=\varphi(\alpha z_1+\beta z_2+\gamma z_3+\delta z_4)
\]
である。「一般には」、したがって解 $\varphi(z)$ は複素数上でも極端に不連続である。$z_0=a+ib$ の近傍における $\varphi(z)$ の\emph{不連続値}は二次元線形多様体をなす。

2) $\varphi(z)$ が階数三であるとする。このとき線形基底の値系が少なくとも三つ、例えば $z_1,z_2,z_3$ 存在して、上の行列式の階数は三である。関係
\[
c_1a+c_2b+c_3A+c_4B=0
\]
を満たす値 $a,b,A,B$、そしてそれらだけが、任意に近似され得る。\emph{不連続値}は $z_0=a+ib$ によって定まる\emph{一次元}線形多様体をなす。挙げた例が示すように、この場合は実際に起こり得る。

3) $\varphi(z)$ が階数二であるとする。常に二つの複素数 $z_1,z_2$ を
\[
\begin{vmatrix}x_1&x_2\\ y_1&y_2\end{vmatrix}\ne0
\]
となるように選べるから、関係は
\[
X=\lambda_1x+\lambda_2y,
\qquad
Y=\mu_1x+\mu_2y
\]
の形をもつ。ここから生じる式
\[
\varphi(z)=(\lambda_1x+\lambda_2y)+i(\mu_1x+\mu_2y)
\]
は、複素数の領域における関数方程式 (1) の最も一般の\emph{連続}解を表す。

4) $\varphi(z)$ が階数一であるとする。この場合は複素数の領域では起こり得ない。実数の領域では\emph{連続}解 $\varphi(x)=C\cdot x$ を意味する。

いま、すべての実数および複素数に対して定義された関数方程式 (1) から (4) の\emph{解} $f(z)$ \emph{は、階数二または階数四しかもち得ない}ことを示す。階数二は、$f(1)=1$ および $f(i)=\pm i$ の特殊化を 3) と合わせると分かるように、二つの\emph{連続}関数 $f(z)=z$ または $\overline z$ だけを与える。階数一はすでに上で排除されたので、階数三を排除すればよい。そこで、少なくとも一つの関係
\[
c_1x+c_2y+c_3X+c_4Y=0
\]
が存在し、したがって $f(z)$ が階数三、または場合によってはより低い階数であると仮定し、この仮定のもとでは常に後者の場合が起こることを示す。再び $f(1)=1$, $f(i)=\pm i$ の特殊化から、関係は
\[
c_3(X-x)+c_4(Y\mp y)=0
\]
へ移る。ここで $c_3$ と $c_4$ は同時には消えない。まず例えば $c_4=0$ とする。このとき関係 $(X-x)=0$ は、純虚数 $z$ にも純虚数 $f(z)$ が対応することを意味する。したがって任意の実数 $y$ に対して $f(iy)=i\cdot\mathfrak Y$ と書ける。ここで $\mathfrak Y$ は再び実数である。ところが $f(iy)=\pm i f(y)$ であるから、$f(y)=\pm\mathfrak Y$ も従う。よって各実数 $z$ にも実数 $f(z)$ が対応し、$(X-x)$ によりすべての実値は自分自身へ移る。したがって $f(z)=z$ または $\overline z$ のみであり、実際には階数二である。$c_3=0$, $c_4\ne0$ の場合もまったく同様に結論される。

次に $c_3$ と $c_4$ がともに零でないとする。このとき関係は
\[
(Y\mp y)=c\cdot(X-x)
\]
の形に書ける。ここで $c$ は実数で、零でなく有限である。これは、$y=\pm cx$ に対して $Y=c\cdot X$ も成り立つことを意味する。したがって関数方程式を考慮すると、任意の実数 $r$ について
\[
f\{r\cdot(1\pm ic)\}=X(1+ic)=f(r)\cdot(1+i f(c))
\]
が成り立つ。ここで $X$ は再び実数である。しかし第1項により、$f(r)$ の因子が恒等的に消えることはできないので、割り算によって
\[
f(r)=X(\sigma+i\tau)
\]
を得る。ここで $\sigma$ と $\tau$ は実数で、$r$ と $X$ には依存しない。

特に $r=1$ に対しても実数 $X_0$ が属するので、条件 $1=X_0(\sigma+i\tau)$ は必然的に $\tau=0$ を示す。したがって $f(r)=\sigma\cdot X$ である。ゆえに各実数 $z$ にも実数 $f(z)$ が対応する。すなわち $y=0$ から必然的に $Y=0$ が従う。これにより、実数 $z$ に対する線形関係は $X-x=0$ に移る。\footnote{ここでは $c\ne0$、したがって $c_3$ も $c_4$ も消えないという事実を用いている。これに対し上では $c_4\ne0$ だけを用いた。そのため、二つのうち一つが消える場合を先に扱っておく必要があった。} 各実値は自分自身に対応する。したがって再び $f(z)=z$ または $\overline z$ のみであり、実際には階数二である。こうして階数三はすべての場合に排除された。一方、最初の各項の研究が示すように、不連続解は実際に存在し、それらはしたがって必然的に階数四でなければならない。ゆえに 1) により次の定理が成り立つ。\emph{「関数方程式 (1) から (4) のすべての解} $f(z)$ \emph{は、連続解} $f(z)=z$ \emph{および} $\overline z$ \emph{を除いて、実数上および複素数上で極端に不連続である。」}\footnote{Lebesgue は前掲箇所で、$f(z)$ の実部および虚部がともに $x,y$ の「非可測」関数であることを示している。しかし第4項の初めの例 $\varphi(z)=\varphi(x)+i(y+\varphi(x))$ が示すように、そこから複素領域での極端な不連続性はまだ従わない。}

\medskip
\noindent Erlangen, 1915年10月30日。

\newpage

\begin{center}
{\Large\bfseries 11. 指定された群をもつ方程式}\par
\vspace{1.2em}
\emph{Math. Ann.} 78 (1918), pp. 221--229

Emmy Noether（Göttingen）
\end{center}

\vspace{1.5em}

指定された群をもつ方程式を構成する問題は、二つの方向から攻めることができる。すなわち、根を特徴づける方向、簡単にいえば「非有理的」方向と、係数を特徴づける「有理的」方向である。

函数論的・算術的に働く「非有理的」方向には、有理数の領域におけるすべての Abel 体は円分体であるという Kronecker の定理、および二次数体に関する相対 Abel 体についての対応する定理が属する。これらの定理は、一方では、これらの特殊な有理性領域に関してすべての Abel 群が実現可能であることを与える。他方では、方程式の全体を与え、同時にそれらによって定義される数体の構造についてより深い洞察を与える。しかし各々の定理は特殊な有理性領域に限られており、新しい有理性領域ごとに全く新しい扱いを必要とする。

代数的に働く「有理的」方向は Hilbert の既約性定理に基づく。この定理によれば、係数の表示においてパラメータ表示に制限してよい。ここには、例えば任意の有理性領域に関して、交代群をもつ方程式が任意に多く存在することが含まれる。ここでは当然、上に述べたような、方程式で定義される体の構造についての洞察は欠ける。しかしこの「有理的」方向の利点は、群の実現可能性がすべての有理性領域に対して同時に証明される点にある。ただし、これまで知られているパラメータ表示は一般には方程式全体を与えない。\footnote{可解方程式については、係数のパラメータ表示を根のパラメータ表示で置き換えることができる。最近 F. Mertens は、次数 8 のある群について、最も一般的な根の表示を与えた: ``Gleichungen mit Quaternionengruppen'', Sitzb. d. Ak. d. Wiss., Wien, Abt. IIa, Bd. 125, S. 735。このノートによれば、それ以前に G. Bucht も、3 次および 4 次方程式、ならびに上記の群をもつ 8 次方程式について、常に作れる標準形に対する最も一般的な根の式を与えていた: ``Über einige algebraische Körper achten Grades'', Arkiv for Math., Astron. och Physik, Bd. 6, Nr. 30。[校正時の追記]}

以下はこの「有理的」方向への一つの寄与である。任意の有理性領域において、指定された群をもつ方程式の全体を与えるパラメータ表示が存在するための十分条件を与える。その条件とは、群に属する不変体、すなわち Lagrange の \emph{Gattungsbereich} が、その領域の独立な函数によって有理的に表示されること、つまり「最小基底」をもつことである。別の言い方をすれば、それが $n$ 個の不定元のすべての有理函数の体と同型であることである。

これは、\S 2 でなお示すように、例えば 3 次および 4 次方程式に現れるすべての群について成り立つ。これらの 3 次および 4 次方程式の実際の構成と詳しい議論は F. Seidelmann の学位論文\footnote{\emph{Die Gesamtheit der kubischen und biquadratischen Gleichungen mit Affekt bei beliebigem Rationalitätsbereich}. Erlangen, 1916.}にあり、その抜粋がこの報告の直後に続く。

Lagrange の \emph{Gattungsbereich} の最小基底の問題は、すでに論文 ``Körper und Systeme rationaler Funktionen'' (Math. Ann. 76) で述べたように、ここで説明した関連とは独立に、E. Fischer が私に対して提起したものであり、これが本研究の契機となった。\footnote{``Rationale Funktionenkörper'' に関する予備報告の第 2 部を参照。Jahresb. d. d. Math. Vereinig., Bd. 22, 1913.}

\subsection*{\S 1. パラメータ表示のための十分条件}

\textbf{1.} 指定された群をもつ方程式の構成をパラメータ表示に帰着できることは、Hilbert の既約性定理から従う。その内容は次の通りである。方程式
\begin{equation}\tag{1}
f(x)=x^n+F_1(\lambda_1,\ldots,\lambda_r)x^{n-1}+\cdots+F_n(\lambda_1,\ldots,\lambda_r)=0
\end{equation}
があり、ここで $F_1(\lambda_1,
\ldots,\lambda_r),\ldots,F_n(\lambda_1,
\ldots,\lambda_r)$ は、独立パラメータ $\lambda_1,
\ldots,\lambda_r$ の有理函数であり、係数は数体\footnote{数体の代わりに一般の有理性領域を置くこともできる。} $\Omega$ に属するとする。この方程式が、$\Omega$ 係数の $\lambda_1,
\ldots,\lambda_r$ のすべての有理函数から成る体 $\Omega(\lambda_1,
\ldots,\lambda_r)$ に関して群 $\Gamma$ をもつならば、パラメータ $\lambda$ を任意に多くの有理数で置き換えて、得られる数値方程式
\begin{equation}\tag{2}
g(x)=x^n+a_1x^{n-1}+\cdots+a_n=0
\end{equation}
が $\Omega$ に関してちょうど群 $\Gamma$ をもつようにできる。そのような $g(x)$ を、より正確には $g_\Gamma$ と記す。

そこで、次のようなパラメータ表示 (1) を与えられるかが問われる。すなわち、$\lambda_1,
\ldots,\lambda_r$ が $\Omega$ のすべての量を走るとき、ただし群の縮小を引き起こす $\Omega$ の量は除くとして、得られる数値方程式 (2) がすべての $g_\Gamma$ を走るような表示である。\footnote{ここには $g(x)=0$ の重根の出現も数え入れなければならない。} 言い換えると、非線型方程式系
\begin{equation}\tag{3}
F_1(\lambda_1,
\ldots,\lambda_r)=a_1;\quad \ldots;\quad F_n(\lambda_1,
\ldots,\lambda_r)=a_n
\end{equation}
が、$g_\Gamma$ に対応する各値系 $a_1,
\ldots,a_n$ に対して、しかも $\Omega$ に属する量 $\lambda_1,
\ldots,\lambda_r$ において解をもつべきである。方程式系 (3) は非線型であるから、$g_\Gamma$ の全体を一つのパラメータ表示で得るという要求には、初めから制限を加えなければならない。実際、この種の非線型系では、代数関係 $H(a_1,
\ldots,a_n)=0$ を満たす特異な値系 $a_1,
\ldots,a_n$ が常に現れ得て、そのような値系に対しては解が存在しない。\footnote{これは「非線型方程式系における交代」についての論文で、近く一般に実行される予定である。} その場合、パラメータ表示 (1) は失敗し、補足表示が必要になる。しかし本場合では、これらの特異値 $a_1,
\ldots,a_n$ は簡単に特徴づけられる。

\textbf{2.} このようなパラメータ表示を得るために、さらに、$\lambda_i$ が、根を不定元と見なしたとき、恒等的に群に属する自然な無理性であることを要求する。これは次の意味である。(1) においてパラメータ $\lambda_1,
\ldots,\lambda_r$ を、$\Omega$ 係数の、不定元 $x_1,
\ldots,x_n$ の適切に選んだ有理函数 $\varphi_1(x),\ldots,\varphi_r(x)$ で置き換えると、(1) は
\begin{equation}\tag{4}
h(x)=x^n-\sigma_1(x)x^{n-1}+\sigma_2(x)x^{n-2}\cdots \pm \sigma_n(x)
\end{equation}
に移るものとする。ここで $\sigma_1(x),\ldots,\sigma_n(x)$ は $x_1,
\ldots,x_n$ の基本対称函数を意味する。そして $h(x)$ は、$\Omega(\lambda_1,
\ldots,\lambda_r)$ からこのようにして生じる有理性領域 $\Omega(\varphi_1(x),\ldots,\varphi_r(x))$ に関して、群 $\Gamma$ をもつものとする。さらにパラメータの独立性により、$\varphi_1(x),\ldots,\varphi_r(x)$ も代数的に独立な函数でなければならない。

この要求を精密にするためには、$\Omega(\varphi_1(x),\ldots,\varphi_r(x))$ と群 $\Gamma$ との関係を明らかにする必要がある。そのため、まず $h(x)$ が $h_\Gamma$ となるような最も一般の有理性領域 $K$ を定める。$\mathfrak L_\Gamma$ で、$\Gamma$ を許す $x_1,
\ldots,x_n$ の有理整係数のすべての有理函数から成る体を表す。これは $\Gamma$ に属する「Lagrange の \emph{Gattungsbereich}」、または $\Gamma$ の不変体とも呼ばれる。また、$\Omega(\Omega_\Gamma)$ で、$\Omega_\Gamma$ を $\Omega$ に添加して得られる対応する体を表す。

したがって $\mathfrak L_\Gamma$ は、$x_1,
\ldots,x_n$ の有理整係数のすべての有理函数の体 $\mathfrak R(x_1,
\ldots,x_n)$ の一つの因子である。すると群の定義により、最も一般の領域 $K$ は、$\Omega_\Gamma$ を因子として含み、かつ $\mathfrak R(x_1,
\ldots,x_n)$ との交わりがちょうど $\Omega_\Gamma$ になるような任意の体によって与えられる。これに対応して、$\Omega$ を含む任意の領域 $K$ について、$\Omega(x_1,
\ldots,x_n)$ との交わりは $\Omega(\Omega_\Gamma)$ によって与えられる。とくに、われわれの要求により $\Omega(\varphi_1(x),\ldots,\varphi_r(x))$ は一つの領域 $K$ になるから、$\Omega(\varphi_1(x),\ldots,\varphi_r(x))$ と $\Omega(x_1,
\ldots,x_n)$ との交わりは $\Omega(\Omega_\Gamma)$ によって与えられる。他方、$\Omega(\varphi_1(x),\ldots,\varphi_r(x))$ は $\Omega(x_1,
\ldots,x_n)$ の因子であるから、ちょうど
\[
\Omega(\varphi_1(x),\ldots,\varphi_r(x))=\Omega(\Omega_\Gamma)
\]
となる。したがって要求は、Lagrange の \emph{Gattungsbereich} $\Omega(\mathfrak L_\Gamma)$ が、代数的に独立な函数 $\varphi_1(x),\ldots,\varphi_r(x)$ を $\Omega$ に添加することで生じるべきである、ということを意味する。このとき $r=n$ でなければならない。なぜなら $\mathfrak L_\Gamma$ は $n$ 個の代数的に独立な基本函数を含む一方、$n$ 個より多くの函数は代数的に独立ではあり得ないからである。$\varphi_1(x)\ldots\varphi_n(x)$ は $\Omega(\mathfrak L_\Gamma)$ の\emph{最小基底}を成す、と言うことにする。また、$\mathfrak L_\Gamma$ は $\Omega$ 係数の最小基底をもつ、とも言う。

\textbf{3.} いま重要なのは、この考察全体を反転でき、実際に求めるパラメータ表示の十分条件\footnote{要求は 2 で 1 より拡張されたので、条件は必要条件である必要はない。}へと導けることである。そこで、$\varphi_1(x)\ldots\varphi_n(x)$ が $\Omega(\mathfrak L_\Gamma)$ の最小基底であると仮定する。さらに $\Omega$ を、$\varphi_1(x)\ldots\varphi_n(x)$ の係数を含む最小の体とする。このとき $\Omega$ は必然的に代数的数体であり、とくにすべての有理数の体 $\mathfrak R$ であり得る。$\sigma_1(x),\ldots,\sigma_n(x)$ は $\Omega(\mathfrak L_\Gamma)$ に属するから、$x_1,
\ldots,x_n$ に関する恒等式が存在する。ここで $F_i$ はその引数の有理函数を意味する。
\begin{equation}\tag{5}
\sigma_1(x)=F_1(\varphi_1(x),\ldots,\varphi_n(x));\quad\ldots;\quad
\sigma_n(x)=F_n(\varphi_1(x),\ldots,\varphi_n(x)).
\end{equation}
逆に、$\varphi_1(x),\ldots,\varphi_n(x)$、したがってまた
$\varphi(x)=\mu_1\varphi_1(x)+\cdots+\mu_n\varphi_n(x)$ は、$\Omega(\sigma_1(x),\ldots,\sigma_n(x))$ に関して既約な方程式
\begin{equation}\tag{6}
G(\varphi(x);\sigma_1(x),\ldots,\sigma_n(x))=0
\end{equation}
によって、$\sigma_1(x),\ldots,\sigma_n(x)$ に代数的に依存する。この式もまた $x$ における恒等式である。

(5) により、この恒等式 (6) は
\begin{equation}\tag{7}
G\bigl(\varphi(x);F_1(\varphi_1,
\ldots,\varphi_n),\ldots,F_n(\varphi_1,
\ldots,\varphi_n)\bigr)
=H(\varphi_1(x),\ldots,\varphi_n(x))=0
\end{equation}
に移る。そして $\varphi_1(x),\ldots,\varphi_n(x)$ の代数的独立性により、(7) は $x$ においてだけでなく $\varphi_i$ においても恒等的に満たされる。すなわち、独立パラメータ $\lambda$ において恒等的に
\begin{equation}\tag{8}
H(\lambda_1,
\ldots,\lambda_n)=G\bigl(\mu_1\lambda_1+\cdots+\mu_n\lambda_n;
F_1(\lambda_1,
\ldots,\lambda_n),\ldots,F_n(\lambda_1,
\ldots,\lambda_n)\bigr)=0
\end{equation}
が成り立つ。恒等式 (8) から、方程式
\begin{equation}\tag{9}
f(x)=x^n-F_1(\lambda_1,
\ldots,\lambda_n)x^{n-1}+F_2(\lambda_1,
\ldots,\lambda_n)x^{n-2}\cdots \pm F_n(\lambda_1,
\ldots,\lambda_n)=0
\end{equation}
は $\Omega(\lambda_1,
\ldots,\lambda_n)$ に関して群 $\Gamma$ をもつことが従う。実際、この恒等式は、既知の有理的な量 $\lambda=\mu_1\lambda_1+
\cdots+
\mu_n\lambda_n$ が、$\Gamma$ に属する根の函数であることを示している。しかし、$\lambda_i=\varphi_i(x)$ という代入が示すように、$\Gamma$ の真の部分群に属する函数がさらに有理的に知られることはあり得ない。また同じ代入により、$\lambda$ はそのすべての共役とは異なる。さらに $\Omega^*$ が $\Omega$ を含む任意の数体であれば、$f(x)$ は $\Omega^*(\lambda_1,
\ldots,\lambda_n)$ に関して群 $\Gamma$ をもつ。なぜなら 2 により、$\lambda_i=\varphi_i(x)$ と代入しても、$\Omega^*$ を添加した場合に群は縮小されないからである。

つぎに、1 で修正された意味において、$f(x)$ が実際にすべての $g_\Gamma$ の全体を走ることを示さなければならない。そのため、(5) により $F_i(\lambda)$ は引数 $\lambda$ の代数的に独立な函数であることに注意する。したがって方程式系
\begin{equation}\tag{10}
F_1(\lambda_1,
\ldots,\lambda_n)=-a_1;\quad F_2(\lambda_1,
\ldots,\lambda_n)=a_2;\quad\ldots;\quad F_n(\lambda_1,
\ldots,\lambda_n)=\pm a_n
\end{equation}
は、任意の値系 $a_1,
\ldots,a_n$ に対して、まだ述べる特異な場合を除き、解をもつ。$\Omega^*$ に関してある $g_\Gamma$ に対応する任意の値系 $a_1,
\ldots,a_n$ に対して、対応する $\lambda_i$ は、群に属する自然な無理性\footnote{さらに下をも参照。}として、$\Omega^*$ の量になる。したがって $\lambda_1,
\ldots,\lambda_n$ がすべての値系を走れば、$g(x)$ はすべての方程式の全体を、修正された意味で走る。その中には、$\Omega^*$ の値が対応する $g_\Gamma$ の全体が含まれている。

逆に、Hilbert の既約性定理により、$\Omega^*$ の値 $\lambda$ には「一般に」$g_\Gamma$ も対応するから、1 で修正された意味において、パラメータ表示 (9) により $\Omega^*$ に関するすべての $g_\Gamma$ の全体が実際に得られる。

最後に、パラメータ表示が失敗する特異値系を求める必要がある。最小基底の函数を、分子と分母に分けて
\[
\varphi_i(x)=\frac{\psi_i(x)}{\chi_i(x)}
\]
と書く。これを恒等式 (5) に代入すると、これらが（約分なしに）
\begin{equation}\tag{11}
\sigma_k(x)=\frac{G_k(x)}{H_k(x)}\quad\text{または}\quad H_k(x)\cdot\sigma_k(x)=G_k(x)
\end{equation}
に移るものとする。したがってすべての $H_k(x)$ の積は、すべての分母 $\chi_i(x)$ の積で割り切れる。根系 $\alpha_1,
\ldots,\alpha_n$ について
\[
H_1(\alpha)\cdot H_2(\alpha)\cdots H_n(\alpha)\ne0
\]
が成り立ち、したがっていかなる分母 $\chi_i(\alpha)$ も消えないなら、(11) を $H_k(\alpha)$ で割ることができ、
\begin{equation}\tag{12}
\sigma_k(\alpha)=\frac{G_k(\alpha)}{H_k(\alpha)}=F_k(\varphi_1(\alpha),\ldots,\varphi_n(\alpha))
\end{equation}
を得る。ここで $\varphi_i(\alpha)$ は分母が消えないので有限値を取り、$g_\Gamma$ である方程式については $\Omega^*$ の量となる。$\alpha_1,
\ldots,\alpha_n$ を $a_k=(-1)^k\sigma_k(\alpha)$ となるように定めるなら、(12) は (10) の解系を表す。\footnote{方程式の解法はこうして、独立函数から成る系
$\varphi_1(\alpha_1,
\ldots,\alpha_n)=\lambda_1;\ldots;\varphi_n(\alpha_1,
\ldots,\alpha_n)=\lambda_n$
の解法に帰着される。}

したがって特異になり得るのは、
\[
H_1(\alpha)H_2(\alpha)\cdots H_n(\alpha)=0
\]
となるような係数系 $a_k=(-1)^k\sigma_k(\alpha)$ だけである。すなわち、(11) が表示 $0=0$ に移る場合である。その場合、このように定義される特異な $a_1,
\ldots,a_n$ の中に、方程式 $g_\Gamma$ に対応するものが含まれるか、またどの数体についてそうであるかを、別に調べる必要がある。\footnote{例えば F. Seidelmann が前掲書で示したように、3 次および 4 次方程式の交代群についてのみ、このような、特異値系に対応する方程式が現れる。しかもそれは、三乗根を含む数体 $\Omega^*$ についてのみである。この場合にはそれぞれ簡単な補足表示を与えることができる。} まとめると次の定理を得る。

\smallskip
\noindent\textbf{定理.} \emph{群 $\Gamma$ に属する Lagrange の \emph{Gattungsbereich} が、$\Omega$ 係数の最小基底\footnote{一つの最小基底の存在から、元のものから可逆有理変換で導かれる任意に多くの最小基底の存在がただちに従う。}をもつならば、$\Omega$ を含む任意の数体 $\Omega^*$ について、$\Omega^*$ に関して指定された群 $\Gamma$ をもつ方程式の全体は、パラメータ表示 (2) によって有理的に構成できる。ただし、パラメータ表示に関して特異なものは例外となることがあり、その場合は別に与えることができる。パラメータは、群の縮小を引き起こさない $\Omega^*$ のすべての値を走らなければならない。最小基底の係数領域 $\Omega$ がすべての有理数の体 $\mathfrak R$ に等しいならば、この構成は任意の数体に対して可能である。}

さらに Hilbert の既約性定理により、群の縮小を引き起こす $\Omega^*$ の値を除いても、パラメータ多様体の次元は低下しない。したがってなお次が成り立つ。\emph{最小基底の存在から、$\Omega^*$ に関して群 $\Gamma$ をもつ方程式の多様体は、Affekt をもたない方程式の多様体と等しいことが従う。多様体は Affekt によって低下しない。}

\subsection*{\S 2. 特殊な方程式への応用}

いま、3 次および 4 次方程式に現れるすべての群について、最小基底の存在を示す。そのためにまず一般に、$n$ 次方程式についてこの問題が $n-2$ 個の不定元の問題に還元されることを示す。第一の還元は、二番目に高い項を消去するための線型 Tschirnhaus 変換に対応する。第一の基底函数として $\varphi_1(x)=\sigma_1(x)$ を取る。さらに $\mathfrak L_\Gamma$ は、$x_1,
\ldots,x_n$ の、$\Gamma$ を許すすべての同次形式から有理的に導かれることに注意する。$p(x_1,
\ldots,x_n)$ をそのような形式とする。このとき
\[
q(x)=p\left(x_1-\frac{\sigma_1(x)}{n},\ldots,x_n-\frac{\sigma_1(x)}{n}\right)=p\left(x_i-\frac{\sigma_1(x)}{n}\right)
\]
も $\mathfrak L_\Gamma$ に属し、さらに
\[
p(x)\equiv p\left(x_i-\frac{\sigma_1(x)}{n}\right)\pmod{\sigma_1(x)},
\quad\text{または}\quad
p(x)=q(x)+\sigma_1(x)\cdot p_1(x)
\]
が成り立つ。ここで $p_1(x)$ は $\mathfrak L_\Gamma$ に属し、しかも $p(x)$ より低い次元をもつ。この手続きを続けると、$\mathfrak L_\Gamma$ のすべての同次形式、したがって一般に $\mathfrak L_\Gamma$ のすべての有理函数は、$\varphi_1(x)=\sigma_1(x)$ と同次形式
\[
q(x)=p\left(x_i-\frac{\sigma_1(x)}{n}\right)=p(y_i)
\]
によって有理的に表されることが分かる。これらの結合 $y_i$ の間には線型関係 $\sum y_i=0$ がある。したがって $q(x)$ は $(n-1)$ 個の引数にのみ依存する。

第二の還元は、$q(x)$ がその引数の同次形式であることに基づく。第二の基底函数として
\[
\varphi_2(x)=\frac{\tau_3(x)}{\tau_2(x)}
\]
を取る。ただし
\[
\tau_3(x)=\sigma_3\left(x_i-\frac{\sigma_1(x)}{n}\right)=\sigma_3(y_i),
\qquad
\tau_2(x)=\sigma_2\left(x_i-\frac{\sigma_1(x)}{n}\right)=\sigma_2(y_i)
\]
である。したがって $\varphi_2(x)$ は、$(n-1)$ 個の引数、すなわち $\sum y_i=0$ を満たす $y_i$ の、一次の同次分数函数である。$\varphi_2(x)$ によって、すべての $q(x)$ は、同じく $\mathfrak L_\Gamma$ に属する函数
\[
r(x)=\frac{q(x)}{\varphi_2(x)^\lambda}
=\frac{q(x)\cdot\tau_2(x)^\lambda}{\tau_3(x)^\lambda}
=\frac{p(y_i)\cdot\sigma_2(y_i)^\lambda}{\sigma_3(y_i)^\lambda}
\]
で有理的に表される。ここで $\lambda$ は $q(x)$ の次数を表す。したがって $r(x)$ は 0 次同次となり、$\sum y_i=0$ の下での比 $y_1:y_2:\cdots:y_n$ という $(n-2)$ 個の引数だけに依存する。こうして $\mathfrak L_\Gamma$ は
\[
\varphi_1(x)=\sigma_1(x),
\qquad
\varphi_2(x)=\frac{\tau_3(x)}{\tau_2(x)}
\]
と、これら $(n-2)$ 個の引数に依存する $\Omega_\Gamma$ のすべての函数 $r(x)$ から成る体によって有理的に表される。\footnote{第二項をもたない $n$ 次方程式
\[
f(x)=x^n+a_2x^{n-2}+a_3x^{n-3}+\cdots+a_n=0
\]
について、この第二の還元に対応するのは、$a_2\ne0$, $a_3\ne0$ なら常に可能な代入
\[
y=\frac{x a_2}{a_3}
\]
である。これにより $f(x)$ は、因子を除いて、
\[
g(y)=y^n+\frac{a_2^3}{a_3^2}y^{n-2}+\frac{a_2^3}{a_3^2}y^{n-3}
+\frac{a_4a_2^4}{a_3^4}y^{n-4}+\cdots+a_n\frac{a_2^n}{a_3^n}=0
\]
に移る。すなわち、ただ $n-2$ 個のパラメータだけを含む方程式
\[
g(y)=y^n+c\,y^{n-2}+c\,y^{n-3}+c_4y^{n-4}+\cdots+c_n=0
\]
になる。} これで望む還元が達成される。

3 次および 4 次方程式については、したがって、一つまたは二つの引数の函数体への還元が得られる。これらには Lüroth と Castelnuovo の定理により常に最小基底が存在する。\footnote{J. Lüroth: ``Beweis eines Satzes über rationale Kurven'', Math. Ann. 9 (1875); G. Castelnuovo: ``Sulla razionalità delle involuzioni piane'', Math. Ann. 44 (1893)。二つを超える不定元の体では最小基底が存在する必要がないことを、F. Enriques が反例によって示した。Rend. Acc. Linc., Vol. 21, 1912 年 1 月 21 日。} 一つの不定元の場合には、最小基底は有理演算によって導かれ、したがってとくに $\mathfrak L_\Gamma$ に対して有理数係数だけを含む。二つの不定元の場合、Castelnuovo によれば、なお代数的数体 $\Omega$ の係数をもつ基底である可能性がある。しかし実際の研究（F. Seidelmann 前掲書を参照）は、この場合にも各 $\Omega_\Gamma$ について有理数係数の基底が得られることを示す。これはほとんど計算なしに、すでに次のことから分かる。すなわち、四元群については有理数係数の基底
\[
\varphi_1(x)=\sigma_1(x);\quad
\varphi_2(x)=\frac{vw}{u};\quad
\varphi_3(x)=\frac{uw}{v};\quad
\varphi_4(x)=\frac{uv}{w},
\]
を選べる。ただし
\[
u=x_1+x_2-x_3-x_4;\quad v=x_1-x_2+x_3-x_4;\quad w=x_1-x_2-x_3+x_4
\]
である。さらに、交代群は三つの基底函数 $\varphi_2,\varphi_3,\varphi_4$ の巡回群へ移り、各八元群はこれらの函数のうち二つの対称群へ移る。これにより、3 次および 2 次方程式の群への還元がなされる。巡回群については、有理数係数の基底は Abel によってすでに知られていた。ただしそのようには定式化されていなかった。\footnote{Weber, \emph{Algebra}（小版）, \S 93, 式 (12) を参照。} したがって次が証明された。\emph{指定された群をもつ 3 次および 4 次方程式の全体は、パラメータ表示に関して特異なものを例外とする可能性を除き、任意の有理性領域についてパラメータ表示により有理的に構成できる。その多様体は Affekt をもたない方程式の多様体と等しい。}

実際の研究（F. Seidelmann 前掲書を参照）は、三乗根の体を含む有理性領域に関する交代群の場合にだけ補足表示が必要となり、それ以外のすべての場合にはパラメータ表示が端的に全体を与えることを示す。

なお、すべての Abel 群についても最小基底は知られている。\footnote{E. Fischer: ``Die Isomorphie der Invariantenkörper der endlichen Abelschen Gruppen linearer Transformationen'', Gött. Nachr. 1915, および ``Zur Theorie der endlichen Abelschen Gruppen'', Math. Ann. 77 (1915)。} しかしこの場合は、群の指標から導かれる円分体を係数にもつ基底である。したがって、指定された群をもつ方程式の構成について、この場合に新しい結果は得られない。

\medskip
\noindent Göttingen, 1916年7月。

\clearpage

\begin{center}
{\Large\bfseries 12. 任意の微分式の不変量}\par
\vspace{1.5em}
Nachr. v. d. Ges. d. Wiss. zu Göttingen 1918, pp. 37--44

Emmy Noether（Göttingen）
\end{center}

\vspace{3em}
\begin{center}
1918年1月25日の会合で F. Klein により提出。
\end{center}

微分式とは、\(n\) 個の変数 \(x_1,\ldots,x_n\) とその微分
\(dx_1,\ldots,dx_n\) との函数
\[
f(x,dx)=f(x_1,\ldots,x_n;\,dx_1,\ldots,dx_n)
\]
であり、二つの \(n\) 個ずつの引数の系に関して解析的であるが、必ずしも実であるとは仮定しないものをいう。いま基礎に置くのは、変数のすべての解析的変換から成る群であり、これに微分のすべての線形変換から成る群が対応する：
\begin{equation}
x_i=x_i(y_1,\ldots,y_n);\qquad
dx_i=\sum_k\frac{\partial x_i}{\partial y_k}\,dy_k;\qquad
\delta x_i=\sum_k\frac{\partial x_i}{\partial y_k}\,\delta y_k,
\tag{1}
\end{equation}
この変換によって \(f(x,dx)\) は \(g(y,dy)\) に移るものとする。\(f(x,dx)\) の不変量とは、この群に関する不変量であり、さらにこの群が高次微分 \(d^2x,d\delta x,\ldots\) および \(f(x,dx)\) の導函数に誘導する群に関する不変量である。すなわち、任意の \(f\) について、(1) により、変数および現れるすべての微分において恒等的に
\[
\begin{aligned}
&J\!\left(f,\frac{\partial f}{\partial dx}\cdots
\frac{\partial^{\varrho+\sigma}f}{\partial x^\varrho\partial dx^\sigma}
\cdots dx,\delta x,d^2x,\ldots\right)\\
&\qquad =
J\!\left(g,\frac{\partial g}{\partial dy}\cdots
\frac{\partial^{\varrho+\sigma}g}{\partial y^\varrho\partial dy^\sigma}
\cdots dy,\delta y,d^2y\ldots\right)\footnotemark
\end{aligned}
\]
を満たす（解析的）函数 \(J\) をいう。
\footnotetext{\(\displaystyle
\frac{\partial^{\rho+\sigma}f}{\partial x^\rho\partial dx^\sigma}\)
は、以下一貫して
\[
\frac{\partial^{\varrho+\sigma}f}
{\partial x_{i_1}\cdots\partial x_{i_\rho}
 \partial dx_{k_1}\cdots\partial dx_{k_\sigma}}
\]
の略記として用いる。ここで添字は \(1\) から \(n\) までの任意の数を表す。}
同時微分式系の不変量もまったく同様に定義される。とくに、そのような不変量が第一微分 \(dx,\delta x\) だけを含み、かつそれらの微分に関する導函数だけを含んで、変数に関する導函数を含まないならば、函数に変数が現れること、および微分の線形変換は考慮に入らない。この特殊な不変量は、不定係数をもつ微分 \(dx,\delta x\) のすべての線形変換の群に関する不変量であり、同時系の\emph{射影不変量}と呼ぶことにする。

二次の斉次微分形式については、Christoffel（Crelle 70）および Ricci（Math. Annalen 54）が、不変な構成から成る一つの系を与えている。その系により、すべての不変量はこの同時系の射影不変量になる。したがって、不変量全体の問題および同値性の問題は、線形不変式論の問題へ還元される。これを私は「還元定理」と呼びたい。完全系は可算無限個の形式から成るが、任意の有限位数 \(\rho\) の不変量、すなわち変数に関して \(\rho\) 次を超える導函数を含まない不変量については、有限個の形式だけから成る。

以下では、この「還元定理」が任意の微分式に対して成り立つことを示す。ここでも、問題となるのは可算無限個の不変な構成から成る完全系であるが、各有限位数については有限個しか現れない。ただし Christoffel と Ricci は、不変量の生成および還元定理の証明を、見通しにくい消去に依拠させているのに対し、私は不変な微分過程と変分過程によって不変量を直接生成する。後者は、単に別のパラメータに関する微分過程とみなすことができる。この不変な微分過程のアルゴリズムは、Lagrange（\emph{Mécanique analytique}, 第2部, IV節）に続いて、Riemann（全集 XXII）および Lipschitz（Crelle 70, 72）が二次微分形式について発展させたが、還元定理までは到達しなかった。還元定理の証明に用いる補助手段は、Riemann が二次形式について導入した「正規座標」（全集 XIII）である。これは、対応する変分問題において一点から出る極値曲線を直線へ変換するが、斉次微分式
\(f(x,\varkappa\cdot dx)=\Phi(\varkappa)\cdot f(x,dx)\footnote{容易に、\(\Phi(\varkappa)=\varkappa^c\) となることがわかる。ここで \(c\) は任意の実または複素量を表す。}\) に対してのみ存在する。しかし非斉次の場合はこの場合へ還元できる。

\begin{center}
\textbf{I. 不変量の生成。}
\end{center}

\(f(x,dx)\) の射影不変量の、一般には終わらない列は、極化形によって与えられる。微分の変換が線形であるため、\(f(x,dx)=g(y,dy)\) から
\(f(x,dx+\lambda\delta x)=g(y,dy+\lambda\delta y)\) も従う。これを \(\lambda\) で展開すると、不変性を特徴づける関係が得られる：
\begin{equation}
\begin{gathered}
f(dx)=g(dy),\\
f_\delta=\sum\frac{\partial f(dx)}{\partial dx}\,\delta x
       =\sum\frac{\partial g(dy)}{\partial dy}\,\delta y,\\
\vdots\\
f_{\delta^\sigma}=
\sum\frac{\partial^\sigma f(dx)}{\partial dx^\sigma}\,\delta x^\sigma
=
\sum\frac{\partial^\sigma g(dy)}{\partial dy^\sigma}\,\delta y^\sigma,\\
\vdots
\end{gathered}
\tag{2}
\end{equation}

各極化形は、変分過程 \(\bdelta\) を適用することにより、一般不変量の終わらない列を与える：
\begin{equation}
\begin{gathered}
\bdelta^\rho f(x,dx)=\bdelta^\rho g(y,dy),\\
\vdots\\
\bdelta^\rho f_{\delta^\sigma}
=\bdelta^\rho\sum\frac{\partial^\sigma f(dx)}{\partial dx^\sigma}\,
\delta x^\sigma
=\bdelta^\rho\sum\frac{\partial^\sigma g(dy)}{\partial dy^\sigma}\,
\delta y^\sigma,\\
\vdots\\[0.2em]
(\rho=0,1,2,\ldots).
\end{gathered}
\tag{3}
\end{equation}
関係 (2) と (3) から、(1) によって \(f(x,dx)\) の導函数の変換法則を読み取ることができるが、以下ではこれは用いない。不変量 (3) から、線形結合によって「\(\rho\) 次変分の正規形」を得ることができる。これは \(\rho=1\) では Lagrange に、\(\rho=2\) では Riemann に見られるもので、\((\rho+1)\) 次の混合微分、すなわち \(d^\rho\delta,d^{\rho-1}\delta^2,\ldots\) が消去されていることによって特徴づけられる。その形は次である：
\begin{equation}
\begin{gathered}
\Omega_1=\delta f(dx)-d f_\delta(dx),\\
\Omega_2=\delta^2 f-\delta d f_\delta+\frac12 d^2 f_{\delta^2},\\
\vdots\\
\Omega_\rho=\delta^\rho f-\delta^{\rho-1}d f_\delta
+\frac12\delta^{\rho-2}d^2 f_{\delta^2}
+\cdots+\frac{(-1)^\rho}{\rho!}\,d^\rho f_{\delta^\rho},\\
\vdots
\end{gathered}
\tag{4}
\end{equation}

不変量 (4) から、Riemann の一つの Ansatz を一般化することにより、第一微分だけを含む不変量に到達できる。このような不変量を「基本関数」と呼ぶことにする。実際、\(\Omega_1\) において微分を微分商で置き換えると、\(\Omega_1=0\)（\(\delta x\) に関して恒等的）は、ちょうど
\[
f\!\left(\frac{dx}{dt}\right)
\]
に属する変分問題の Lagrange 方程式を与える。そこで私は、方向 \((d+\lambda\delta)\) がこれらの Lagrange 方程式を満たす、という不変条件を課す：
\begin{equation}
\Omega_1(d+\lambda\delta)=
\bdelta f(dx+\lambda\delta x)
-(d+\lambda\delta)f_{\bdelta}(dx+\lambda\delta x)=0
\quad(\text{\(\delta x\) に関して恒等的}),
\tag{5}
\end{equation}
ここから、置換 \(\bdelta=d+\lambda\delta\) によって、斉次な \(f(dx)\) についてさらに
\begin{equation}
(d+\lambda\delta)f(dx+\lambda\delta x)=0
\tag{6}
\end{equation}
が従う。ただし一次斉次の場合は例外である。この場合には、方程式系 (5) が \(n-1\) 個の独立条件しか表さないため、(6) を新たな条件として加えることにする。(5)、または (5) と (6) において、\(\lambda\) の零次・一次・二次の係数を零とおくと、三つの不変な方程式系
\(\varphi=0\), \(\psi=0\), \(\chi=0\)（\(\delta x\) に関して恒等的）が得られ、これにより \(d^2x,d\delta x,\delta^2x\) を第一微分によって表すことができる。さらに次の不変な方程式系
\begin{equation}
\begin{gathered}
d^\sigma\varphi=0,\qquad d^\sigma\psi=0,\qquad
d^\sigma\chi=0,\qquad d^{\sigma-1}\delta\chi=0,\\
d^{\sigma-2}\delta^2\chi=0,\ldots,\delta^\sigma\chi=0
\qquad(\sigma=0,1,2,\ldots)
\end{gathered}
\tag{7}
\end{equation}
は、すべての高次微分を第一微分によって表すのにちょうど十分である。\footnote{線形形式 \(\sum A_i\,dx_i\) だけは特別な処理を要する。ここでは方程式系 (5) が第二微分を含まないからである。} したがって、函数 (4) と不変な方程式系 (7) とを組み合わせることにより、第一微分だけを含む基本関数 \([\Omega_\rho]\) が得られる。ただし \([\Omega_1]\) は恒等的に消える。

同様に、基本関数 \(h\) の微分 \(\delta h(dx,\delta x)\) から、(7) によって再び基本関数を得ることができ、これを \(h\) の「共変導関数」\([h^{(1)}]\) と呼ぶことにする。一般に \([h^{(\nu)}]\) は \([h^{(\nu-1)}]\) の共変導関数を表す。この二つの過程により、基本関数の二重無限列が得られる：
\begin{equation}
\begin{array}{cccccc}
[\Omega_2], & [\Omega_2^{(1)}], & [\Omega_2^{(2)}], & \ldots,
& [\Omega_2^{(\nu)}], & \ldots\\
\vdots &&&& \vdots&\\
{}[\Omega_\rho], & [\Omega_\rho^{(1)}], & \ldots, & \ldots,
& [\Omega_\rho^{(\nu)}], & \ldots\\
\vdots & \vdots &&& \vdots&
\end{array}
\tag{8}
\end{equation}

さらに、高次微分を第一微分によって表す方程式の左辺は、第一微分と共変的である、すなわち同じ線形変換を受けることがわかる。したがって、高次微分の代わりに、これらの左辺 \(p,q,r,\ldots\) を不変量の引数として導入すると、その不変量は、導函数を別にすれば、互いに共変的な量だけに依存する。上で述べた函数 (8) の形成は、単にこれらの新しい量を導入することに帰着する。任意の不変量はそれによって不変量の和へ移り、その和から、ちょうど \(dx,\delta x\) を含み、\(p,q,r,\ldots\) を含まない項を取り出すのである。

以下では、斉次性の仮定
\[
f(\varkappa\cdot dx)=\Phi(\varkappa)\cdot f(dx)
\]
の下で、基本関数 (8) が求める完全系を尽くすことを示す。

\begin{center}
\textbf{II. 正規座標、同値性および還元定理。}
\end{center}

私は、\(f\!\left(\frac{dx}{dt}\right)\) に属する変分問題を積分することによって正規座標に到達する：
\begin{equation}
\begin{gathered}
\Omega_1=\delta f\!\left(\frac{dx}{dt}\right)
-\frac{d}{dt}f_\delta\!\left(\frac{dx}{dt}\right)=0
\quad(\text{\(\delta x\) に関して恒等的}),\\
\frac{d}{dt}f\!\left(\frac{dx}{dt}\right)=0
\quad(\text{帰結または追加条件として}).
\end{gathered}
\tag{9}
\end{equation}
\(f(dx)\) の斉次性の仮定により、この方程式系はパラメータ置換 \(t=\varkappa\cdot\tau\) のもとで自分自身へ移る。したがって、(9) によって \(\frac{d^\rho x_i}{dt^\rho}\) を第一導函数の函数
\(\varphi_\rho^{(i)}\!\left(\frac{dx}{dt}\right)\) として表すなら、\(\varphi_\rho^{(i)}\) は \(\rho\) 次斉次である：
\begin{equation}
\varphi_\rho^{(i)}\!\left(\varkappa\cdot\frac{dx}{dt}\right)
=\varkappa^\rho\varphi_\rho^{(i)}\!\left(\frac{dx}{dt}\right).
\tag{10}
\end{equation}
いま (9) を積分するにあたり、\(t=t_0\) で初期値
\[
x=(x)_0,\qquad \frac{dx}{dt}=\left(\frac{dx}{dt}\right)_0
\]
を与え、
\begin{equation}
u_i=(t-t_0)\left(\frac{dx_i}{dt}\right)_0
\tag{11}
\end{equation}
とおく。ここで \(f\!\left(\frac{dx}{dt}\right)=\mathrm{const.}\) により、パラメータ \((t-t_0)\) は「極値曲線の長さ」に比例する。すると、\((t-t_0)\) の冪による展開は、(10) により
\begin{equation}
x_i-(x_i)_0=u_i+\frac12\varphi_2^{(i)}(u)+\cdots+
\frac{1}{\rho!}\varphi_\rho^{(i)}(u)+\cdots .
\tag{12}
\end{equation}
を与える。この展開は変数変換 \(x_i=x_i(u)\) と見ることができ、このとき量 \(u\) がまさに Riemann の正規座標である。したがって、(11) と (12) により、これらの座標は \((x)_0\) を通る極値曲線を直線に変換する。さらに任意の変数変換 \(x=x(y)\) に対して、(11) により、対応する正規座標の線形変換 \(u=u(v)\) が対応し、その係数は任意の値系 \((x)_0\) だけに依存する。したがって微分 \(du,\delta u\) は \(u\) 自身と同じ変換を受ける。

(12) によって \(f(x,dx)\) は \(F(u,du)\) に移るものとする。すなわち、\(u\) の冪で展開して
\begin{equation}
F(u,du)=F_0(u,du)+F_1(u,du)+\cdots+F_\rho(u,du)+\cdots,
\tag{13}
\end{equation}
ここで \(F_\rho(u,du)\) は \(u\) における \(\rho\) 次元の斉次形式を表す。変換 \(u=u(v)\) の線形性のため、\(F(u,du)=G(v,dv)\) によって、個々の斉次成分も対応する成分へ移る：
\begin{equation}
F_\rho(u,du)=G_\rho(v,dv).
\tag{14}
\end{equation}
したがって、\((x)_0\) に任意の固定された定数値系を取ると、(13) と (14) は、\(f(dx)\) と \(g(dy)\) の同値性が、線形変換に関する対応する函数系 \(F_\rho(u,du)\) の同値性と同じ意味であることを示す。さらに \(F_\rho(u,du)\)、または対応する \(F_\rho(\delta u,du)\) は、点 \(x=(x)_0\) で取った \(f(dx)\) の不変量を表し、実際、この点 \(x=(x)_0\) に対する基本関数の完全系を成す。すなわち
\[
F_\varrho(\delta u,du)=\frac1{\varrho!}\sum
\frac{\partial F_\varrho}{\partial \delta u^\varrho}\,\delta u^\varrho;
\qquad
\frac{\partial F_\varrho}{\partial \delta u^\varrho}
=\left(\frac{\partial F(du)}{\partial u^\varrho}\right)_0,
\]
であり、この最後の導函数は、(12) によって、\(x=(x)_0\) で取った \(f(dx)\) の導函数を用いて、\(G(dv)\) から対応して形成される導函数が \(g(dy)\) の導函数を用いて表されるのと全く同じ仕方で表すことができる。さらに
\[
\left(\frac{\partial^{\varrho+\sigma}F(du)}
{\partial u^\varrho\partial du^\sigma}\right)_0
=
\frac{\partial^{\varrho+\sigma}F_\varrho(\delta u,du)}
{\partial\delta u^\varrho\partial du^\sigma}
\]
であるから、点 \(x=(x)_0\) で取った \(f(dx)=F(du)\) の任意の不変量は、高次微分を \(dx,\delta x\) と共変的な \(p,q,r,\ldots\) で置き換えるだけで、\(F_\rho(\delta u,du)\) の射影不変量となる。ところが \((x)_0\) はパラメータと見ることができる。点 \(x=(x)_0\) における各不変量には、\((x)_0\) をふたたび \(x\) で置き換えれば、\(f(dx)\) の不変量が対応する。したがって、\(F_\rho(\delta u,du)\) が \(x=(x)_0\) で取った不変量 \(\Psi_\rho(\delta x,dx)\) になるならば、ここで \(x=(x)_0\) では \(du,\delta u\) が \(dx,\delta x\) へ移るので、\(\Psi_\rho\) 自身が完全系の基本関数を成す。残るのは、この基本関数系を \(f(dx)\) の明示的な不変量、すなわち函数 (8) によって表すことである。それが可能であることは、それらが微分過程によって生成されることから示される。最高位の項については、これらの過程は単純な極化過程へ移り、Clebsch--Gordan の級数展開に類似した変形を、無視した項に関する帰納法と組み合わせることで、主張が証明される。同時系の不変量を扱う場合には、\(f(dx)\) に属する正規座標をこれらの式に導入することから分かるように、基本関数 (8) にさらに残りの微分式の共変導関数を付け加えなければならない。

(14) に還元されていた \(f(dx)\) と \(g(dy)\) の同値性は、したがって、\(\Psi_\rho\) の同値性、または同じことだが函数 (8) の同値性、すなわち微分の線形変換に関する同値性とも同じ意味になる。このとき特別な可積分条件は不要であり、それは (14) への還元可能性から従う。これにより還元定理はすべての部分において証明された。とくに、函数列
\([\Omega_2],[\Omega_3],\ldots,[\Omega_\rho],\ldots\) が恒等的に消えることは、\(f(dx)\) が定数係数をもつ式へ変換できるための必要十分条件である。

最後に、すべての解析的変換から成る群の代わりに部分群を基礎に置くなら、対応する \(u\) の線形変換の群も射影群の部分群へ移る。\(f(dx)\) の不変量は、また函数 (8) の不変量となり、線形変換に関するものではあるが、今度はこの部分群に関するものとなる。したがって、斉次化によって、非斉次函数 \(f(dx)\) の場合はアフィン群へ還元できる。この場合、完全系は一つ多い変数について形成された函数 (8) から導き出される。

還元定理から、二次形式の場合、より一般に \(p\) 次元の形式の場合が、次のことにより得られる。ここでは函数系が \([\Omega_2]\)、より一般には \([\Omega_p]\) とその共変導関数で終わり、後続のすべての基本関数がこれらの射影不変量となるのである。これにより、二次形式における Riemann の曲率形式 \([\Omega_2]\) の特別な位置が説明される。二次形式、または \(p\) 次元の形式の同値性の問題は、まさに \([\Omega_2]\)、または \([\Omega_2]\ldots[\Omega_p]\)、およびそれらの共変導関数の、線形変換に関する同値性へ帰着する。そして \([\Omega_2]\)、または \([\Omega_2]\ldots[\Omega_p]\) が恒等的に消えることは、定数係数をもつ形式へ変換できるために必要十分である。これにより、二次形式について最もよく知られた結果の一つが述べられたことになる。

より詳しい叙述は、近く \emph{Mathematische Annalen} に掲載される予定である。

\clearpage

\begin{center}
{\Large\bfseries 13. 不変な変分問題}\par
\vspace{1.5em}
\emph{Nachrichten von der Gesellschaft der Wissenschaften zu Göttingen} 1918, pp. 235--257
\end{center}

\vspace{3em}
\begin{center}
1918年7月26日の会合で F. Klein により提出。\footnote{原稿の最終版が提出されたのは9月末になってからである。}
\end{center}

ここで扱うのは、（Lie の意味での）連続群を許す変分問題である。このことから対応する微分方程式について生じる帰結は、\S{}1 で定式化され、後続の節で証明される定理において最も一般的な表現を得る。変分問題から生じるこれらの微分方程式については、Lie の研究対象であった、群を許す任意の微分方程式についてよりも、はるかに精密な主張を行うことができる。したがって以下は、形式的変分計算の方法と Lie の群論の方法とを結合するものである。特殊な群と特殊な変分問題については、この方法の結合は新しいものではない。有限群の特殊な場合については Hamel と Herglotz を、無限群の特殊な場合については Lorentz とその弟子（例えば Fokker）、Weyl、Klein を挙げておく。\footnote{Hamel: \emph{Math. Ann.} vol. 59 および \emph{Zeitschrift f. Math. u. Phys.} vol. 50. Herglotz: \emph{Ann. d. Phys.} (4) vol. 36, とくに \S{}9, p. 511. Fokker, Verslag d. Amsterdamer Akad., 1917年1月27日。さらに文献については Klein の第二ノート、Göttinger Nachrichten, 1918年7月19日を参照。つい最近現れた Kneser の論文（\emph{Math. Zeitschrift} vol. 2）では、類似の方法による不変量の構成が問題になっている。} とくに、Klein の第二ノートと本稿の展開とは相互に影響し合っている。この点については Klein のノートの結語を参照してよい。

\section*{\S{}1. 予備的事項と定理の定式化}

以下に現れるすべての関数は、考察する領域で、解析的であるか、少なくとも連続で、有限回連続微分可能で、一価であると仮定する。

「変換群」とは、通常どおり、各変換に対してその逆変換が同じ系に含まれ、かつその系の任意の二つの変換の合成がふたたびその系に属するような変換の系をいう。群は、その変換が、\(\rho\) 個の本質的なパラメータ \(\eps\) に解析的に依存する最も一般の変換に含まれているとき、有限連続群 \(\G_\rho\) と呼ばれる（すなわち、これら \(\rho\) 個のパラメータは、より少数のパラメータの \(\rho\) 個の関数として表されてはならない）。これに対応して、無限連続群 \(\G_{\infty\rho}\) とは、その最も一般の変換が \(\rho\) 個の本質的な任意関数 \(p(x)\) とそれらの導関数に、解析的に、または少なくとも連続かつ有限回連続微分可能な仕方で依存する群をいう。両者の中間に、任意関数には依存しないが無限個のパラメータに依存する群がある。最後に、任意関数にもパラメータにも依存する群を混合群と呼ぶ。\footnote{Lie は \emph{Grundlagen für die Theorie der unendlichen kontinuierlichen Transformationsgruppen}（Berichte der Königl. Sächs. Gesellschaft der Wissenschaften, 1891）で、無限連続群を、変換が偏微分方程式系の最も一般の解で与えられ、かつそれらの解が有限個のパラメータだけに依存しないような変換群として定義している。これにより、上に挙げた型のうち、有限群とは異なる一つの型が得られる。しかし逆に、無限個のパラメータという極限の場合が、必ずしも微分方程式系を満たす必要はない。}

\(x_1,\ldots,x_n\) を独立変数、\(u_1(x),\ldots,u_\mu(x)\) をそれらに依存する関数とする。\(x\) と \(u\) をある群の変換に従わせるとき、変換が可逆であるという仮定により、変換された量の中には再びちょうど \(n\) 個の独立なもの \(y_1,\ldots,y_n\) が含まれていなければならない。これらに依存する残りの量を \(v_1(y),\ldots,v_\mu(y)\) と書く。変換には、\(x\) に関する \(u\) の導関数、すなわち \(\p u/\p x,\ \p^2u/\p x^2,\ldots\) も現れてよい。\footnote{添字は、和においても可能なかぎり省略する。したがって、例えば \(\p^2u/\p x^2\) は \(\p^2u_\nu/\p x_i\p x_k\) などを表す。} ある関数が群の不変量であるとは、関係式
\[
P\!\left(x,u,\frac{\p u}{\p x},\frac{\p^2u}{\p x^2},\cdots\right)
 =
P\!\left(y,v,\frac{\p v}{\p y},\frac{\p^2v}{\p y^2},\cdots\right)
\]
が成り立つことをいう。とくに、積分 \(I\) が群の不変量であるとは、関係式
\begin{equation}
\begin{aligned}
I&=\int\!\cdots\!\int f\!\left(x,u,\frac{\p u}{\p x},\frac{\p^2u}{\p x^2},\cdots\right)d x\\
 &=\int\!\cdots\!\int f\!\left(y,v,\frac{\p v}{\p y},\frac{\p^2v}{\p y^2},\cdots\right)d y
\end{aligned}
\tag{1}
\end{equation}
が成り立つことである。ここで積分は任意の実 \(x\)-領域と、対応する \(y\)-領域にわたる。\footnote{簡単のため \(\dd x,\dd y\) と書いて、それぞれ \(\dd x_1\cdots \dd x_n,\dd y_1\cdots \dd y_n\) を表す。変換に現れるすべての引数 \(x,u,\eps,p(x)\) は実とするが、係数は複素数であってよい。しかし最終結果は \(x,u\)、パラメータ、任意関数に関する恒等式であるから、現れる関数を解析的と仮定すれば複素値についても成立する。さらに、結果の大部分は積分を用いずに基礎づけることができるので、その証明においても実数への制限は必要ではない。これに対し、\S{}2 の末尾と \S{}5 の初めの考察は、積分なしには実行できないようである。}

他方、必ずしも不変とは限らない任意の積分 \(I\) に対して、私は第一変分 \(\delta I\) を作り、変分計算の規則に従って部分積分で変形する。よく知られているように、\(\delta u\) および境界に現れるそのすべての導関数を境界で零とし、その他では任意とすると、次を得る：
\begin{equation}
\delta I=\int\!\cdots\!\int \delta f\,d x
=\int\!\cdots\!\int\left(\sum\psi_i\!\left(x,u,\frac{\p u}{\p x},\cdots\right)\delta u_i\right)d x,
\tag{2}
\end{equation}
ここで \(\psi\) はラグランジュ式を表す。すなわち、対応する変分問題 \(\delta I=0\) のラグランジュ方程式の左辺である。この積分関係には、境界項も書き出すことによって生じる、\(\delta u\) とその導関数に関する積分を含まない恒等式が対応する。部分積分が示すように、これらの境界項は発散の積分、すなわち
\[
\Div A=\frac{\p A_1}{\p x_1}+\cdots+\frac{\p A_n}{\p x_n}
\]
という形の式の積分である。ただし \(A\) は \(\delta u\) とその導関数に関して線形である。したがって
\begin{equation}
\sum_i\psi_i\,\delta u_i=\delta f+\Div A.
\tag{3}
\end{equation}
を得る。とくに \(f\) が \(u\) の一階導関数だけを含むなら、単積分の場合には恒等式 (3) はいわゆる「ラグランジュの中央方程式」と同一である：
\begin{equation}
\sum_i \psi_i\,\delta u_i
=\delta f-\frac{\dd}{\dd x}\left(\sum_i\frac{\p f}{\p u_i'}\,\delta u_i\right),
\qquad \left(u_i'=\frac{d u_i}{d x}\right),
\tag{4}
\end{equation}
一方、\(n\)-重積分に対して (3) は
\begin{equation}
\sum_i\psi_i\,\delta u_i
=
\delta f
-\frac{\p}{\p x_1}\left(\sum_i\frac{\p f}{\p \left(\frac{\p u_i}{\p x_1}\right)}\,\delta u_i\right)
-\cdots
-\frac{\p}{\p x_n}\left(\sum_i\frac{\p f}{\p \left(\frac{\p u_i}{\p x_n}\right)}\,\delta u_i\right)
\tag{5}
\end{equation}
となる。単積分で、\(u\) の \(\kappa\) 階導関数まで含む場合、(3) は次で与えられる：
\begin{equation}
\begin{aligned}
\sum_i\psi_i\,\delta u_i
={}&\delta f
-\frac{\dd}{\dd x}\Bigg\{\sum_i\left[
\binom{1}{1}\frac{\p f}{\p u_i^{(1)}}\delta u_i
+\binom{2}{1}\frac{\p f}{\p u_i^{(2)}}\delta u_i^{(1)}
+\cdots+
\binom{\kappa}{1}\frac{\p f}{\p u_i^{(\kappa)}}\delta u_i^{(\kappa-1)}
\right]\Bigg\}\\
&+\frac{\dd^2}{\dd x^2}\Bigg\{\sum_i\left[
\binom{2}{2}\frac{\p f}{\p u_i^{(2)}}\delta u_i
+\binom{3}{2}\frac{\p f}{\p u_i^{(3)}}\delta u_i^{(1)}
+\cdots+
\binom{\kappa}{2}\frac{\p f}{\p u_i^{(\kappa)}}\delta u_i^{(\kappa-2)}
\right]\Bigg\}\\
&+\cdots
+(-1)^\varkappa\frac{d^\varkappa}{d x^\varkappa}
\left\{\sum\binom{\varkappa}{\varkappa}
\frac{\p f}{\p u_i^{(\kappa)}}\delta u_i\right\}.
\end{aligned}
\tag{6}
\end{equation}
対応する恒等式は \(n\)-重積分に対しても成り立つ。とくに \(A\) は \((\kappa-1)\) 階導関数までの \(\delta u\) を含む。(4), (5), (6) が実際にラグランジュ式 \(\psi_i\) を定義することは、右辺の組合せにおいて \(\delta u\) の高階導関数がすべて消去される一方で、部分積分により一意的に導かれる関係 (2) が満たされることから従う。

以下で問題となるのは、次の二つの定理である。

I. 積分 \(I\) が \(\G_\rho\) に関して不変であるなら、ラグランジュ式の \(\rho\) 個の線形独立な組合せは発散になる。逆に、このことから \(I\) の \(\G_\rho\) に関する不変性が従う。この定理は、無限個のパラメータという極限の場合にも成り立つ。

II. 積分 \(I\) が、任意関数が \(\sigma\) 階導関数まで現れる \(\G_{\infty\rho}\) に関して不変であるなら、ラグランジュ式とその \(\sigma\) 階までの導関数との間に \(\rho\) 個の恒等的な関係が存在する。ここでも逆が成り立つ。\footnote{ある自明な例外については \S{}2 の第二注を参照。}

混合群については、両定理の主張が適用される。したがって、互いに独立な従属関係と発散関係が同時に現れる。

これらの恒等式から対応する変分問題へ移り、したがって \(\psi=0\) とおくと、\footnote{少し一般化して \(\psi_i=T_i\) とおくこともできる。\S{}3 の第一注を参照。} 定理 I は一次元の場合、すなわち発散が全微分となる場合に、\(\rho\) 個の第一積分が存在することを述べる。ただしそれらの間に非線形の従属が生じることはありうる。\footnote{\S{}3 の末尾を参照。} 多次元の場合には、近年しばしば「保存則」と呼ばれている発散方程式が得られる。定理 II は、ラグランジュ方程式のうち \(\rho\) 個が残りの方程式の帰結であることを述べる。

定理 II の最も簡単な例（逆は除く）は Weierstrass のパラメータ表示である。ここでは一階の同次性の下で、独立変数 \(x\) を \(x\) の任意関数で置き換え、\(u\) を変えない \((y=p(x);\ v_i(y)=u_i(x))\) とき、積分が不変であることが知られている。したがって任意関数が一つ現れるが、導関数は現れない。これに対応して、ラグランジュ式そのものの間に有名な線形関係
\[
\sum_i \psi_i\,\frac{\dd u_i}{\dd x}=0
\]
がある。さらに別の例は、物理学者の「一般相対性理論」によって与えられる。ここではすべての \(x\) の変換 \(y_i=p_i(x)\) の群が問題となり、\(u\)（\(g_{\mu\nu}\) および \(q\) と記される）は、それによって二次および一次の微分形式の係数に誘導される変換を受ける。その変換には任意関数 \(p(x)\) の一階導関数が含まれる。これに対応して、ラグランジュ式とその一階導関数との間の既知の \(n\) 個の従属関係がある。\footnote{例えば Klein の叙述を参照。}

とくに、群を特殊化して、変換に \(u(x)\) の導関数が現れないようにし、さらに変換された独立量が \(u\) ではなく \(x\) だけに依存するようにすると、\S{}5 で示されるように、\(I\) の不変性から \(\sum_i\psi_i\delta u_i\) の相対不変性\footnote{すなわち、\(\sum_i\psi_i\delta u_i\) は変換の下である因子を獲得する。} が従う。また、パラメータを適切な変換に従わせるなら、定理 I に現れる発散についても同様である。さらに上に述べた第一積分も群を許すことが従う。定理 II についても同様に、任意関数によって結合された従属関係の左辺の相対不変性が得られ、その帰結として、発散が恒等的に消え、かつ群を許す関数も得られる。これは物理学者の相対性理論において、従属関係とエネルギー定理との連関を媒介する関数である。\footnote{Klein の第二ノートを参照。} 最後に、群論的な形で、定理 II は「一般相対性」において真性エネルギー定理が失敗することに関する Hilbert の関連する主張の証明を与える。これらの補足的な注を加えれば、定理 I は力学などで知られている第一積分に関するすべての定理を含み、定理 II は「一般相対性理論」の可能な限り最大の群論的一般化と述べることができる。

\section*{\S{}2. 発散関係と従属関係}

\(\G\) を、有限または無限の連続群とする。恒等変換が、パラメータ \(\eps\)、それぞれ任意関数 \(p(x)\) の値ゼロに対応するよう、常に整えることができる。\footnote{例えば Lie, \emph{Grundlagen}, p. 331 を参照。任意関数が関与する場合、パラメータの特殊値 \(a^\sigma\) は固定された関数 \(p^\sigma,\ \p p^\sigma/\p x,\ldots\) に置き換えられるべきであり、それに対応して値 \(a^\sigma+\eps\) は \(p^\sigma+p(x),\ \p p^\sigma/\p x+\p p/\p x\) などに置き換えられる。} したがって最も一般の変換は次の形をもつ：
\[
y_i=A_i\!\left(x,u,\frac{\p u}{\p x},\cdots\right)=x_i+\varDelta x_i+\cdots,\qquad
v_i(y)=B_i\!\left(x,u,\frac{\p u}{\p x},\cdots\right)=u_i+\varDelta u_i+\cdots,
\]
ここで \(\D x_i,\D u_i\) は、\(\eps\)、それぞれ \(p(x)\) とその導関数における最低次の項を表し、これらの量に関して線形であると仮定する。後で示すように、これは一般性を制限しない。

いま積分 \(I\) が \(\G\) に関して不変であり、したがって関係 (1) が満たされているとする。とくに \(I\) は、\(\G\) に含まれる infinitesimal 変換
\[
y_i=x_i+\varDelta x_i;\qquad v_i(y)=u_i+\varDelta u_i
\]
に関しても不変である。このとき関係 (1) は
\begin{equation}
\begin{aligned}
0=\D I
&=\int\!\cdots\!\int f\!\left(y,v(y),\frac{\p v}{\p y},\ldots\right)\dd y\\
&\qquad-\int\!\cdots\!\int f\!\left(x,u(x),\frac{\p u}{\p x},\ldots\right)\dd x
\end{aligned}
\tag{7}
\end{equation}
となる。第一の積分は、\(x\)-領域に対応する \(x+\D x\)-領域にわたって拡張される。しかしこの積分は、infinitesimal な \(\D x\) に対して有効な変換により、\(x\)-領域上の積分へ変換することもできる：
\begin{equation}
\begin{aligned}
\int\!\cdots\!\int f\!\left(y,v(y),\frac{\p v}{\p y},\ldots\right)\dd y
&=\int\!\cdots\!\int f\!\left(x,v(x),\frac{\p v}{\p x},\ldots\right)\dd x\\
&\quad+\int\!\cdots\!\int \Div(f\cdot \D x)\,\dd x.
\end{aligned}
\tag{8}
\end{equation}
したがって、infinitesimal 変換 \(\D u\) の代わりに変分
\begin{equation}
\bd u_i=v_i(x)-u_i(x)=\D u_i-\sum_\lambda\frac{\p u_i}{\p x_\lambda}\D x_\lambda
\tag{9}
\end{equation}
を導入すると、(7), (8) は
\begin{equation}
0=\int\!\cdots\!\int\{\bd f+\Div(f\cdot\D x)\}\dd x.
\tag{10}
\end{equation}
となる。関係 (10) は任意の領域上で積分したとき成立するので、被積分関数は恒等的に消えなければならない。したがって \(I\) の不変性に対する Lie の微分方程式は、関係
\begin{equation}
\bd f+\Div(f\cdot\D x)=0
\tag{11}
\end{equation}
へ移る。ここで (3) によって \(\bd f\) をラグランジュ式で表すと
\begin{equation}
\sum_i\psi_i\bd u_i=\Div B,\qquad (B=A-f\cdot\D x)
\tag{12}
\end{equation}
を得る。したがってこの関係は、任意の不変積分 \(I\) に対して、現れるすべての引数に関する恒等式を表す。これは \(I\) に対する Lie の微分方程式の求める形である。\footnote{\(\Div(f\cdot\D x)=0,\ \bd u=0\) となる自明な場合には、(12) は \(0=0\) になる。この場合は、\(\D x,\D u\) が \(u\) の導関数にも依存するときにだけ起こりうる。したがってこれらの infinitesimal 変換は常に群から切り離すべきであり、定理の定式化においては残りのパラメータまたは任意関数の数だけを数えるべきである。残りの infinitesimal 変換がなお常に群をなすかどうかは未決としておく必要がある。}

まず \(\G\) が有限連続群 \(\G_\rho\) であると仮定する。仮定により \(\D u\) と \(\D x\) はパラメータ \(\eps_1,\ldots,\eps_\rho\) に関して線形であるから、(9) により \(\bd u\) とその導関数にも同じことが成り立つ。したがって \(A\) と \(B\) は \(\eps\) に関して線形である。そこで
\[
B=B^{(1)}\eps_1+\cdots+B^{(\varrho)}\eps_\varrho;\qquad
\bd u=\bd u^{(1)}\eps_1+\cdots+\bd u^{(\varrho)}\eps_\varrho
\]
とおく。ここで \(\bd u^{(1)},\ldots\) は \(x,u,\p u/\p x,\ldots\) の関数である。(12) から求める発散関係
\begin{equation}
\sum_i\psi_i\bd u_i^{(1)}=\Div B^{(1)},\ldots,\qquad
\sum_i\psi_i\bd u_i^{(\rho)}=\Div B^{(\rho)}
\tag{13}
\end{equation}
が従う。したがって、ラグランジュ式の \(\rho\) 個の線形独立な組合せが発散になる。線形独立性は、(9) により \(\bd u=0,\D x=0\) が \(\D u=0,\D x=0\) を含意し、したがって infinitesimal 変換の間の従属関係を含意することから従う。しかし仮定により、そのような従属関係はいかなるパラメータ値に対しても存在しない。そうでなければ、infinitesimal 変換から積分により再び生じる \(\G_\rho\) は、\(\rho\) より少ない本質的パラメータに依存することになるからである。さらに \(\bd u=0,\ \Div(f\D x)=0\) という可能性は除外してある。これらの結論は、無限個のパラメータという極限の場合にも有効である。

次に \(\G\) が無限連続群 \(\G_{\infty\rho}\) であるとする。このときも \(\bd u\) とその導関数、したがって \(B\) も、任意関数 \(p(x)\) とその導関数に関して線形である。\footnote{\(p\) が \(u,\p u/\p x,\ldots\) に依存しないと仮定することが制限でないことは、逆定理によって示される。} (12) とはまだ独立に、\(\bd u\) の値を代入すると
\[
\sum\psi_i\bd u_i
=\sum_{\lambda,i}\psi_i\left\{
a_i^{(\lambda)}(x,u,\ldots)p^{(\lambda)}(x)
+b_i^{(\lambda)}(x,u,\ldots)\frac{\p p^{(\lambda)}}{\p x}
+\cdots+c_i^{(\lambda)}(x,u,\ldots)
\frac{\p^\sigma p^{(\lambda)}}{\p x^\sigma}\right\}
\]
となる。いま、部分積分の公式と同様に、恒等式
\[
\varphi(x,u,\ldots)\frac{\p^\tau p(x)}{\p x^\tau}
=(-1)^\tau\cdot\frac{\p^\tau\varphi}{\p x^\tau}\cdot p(x)
\quad \bmod \text{ 発散}
\]
によって、\(p\) の導関数は \(p\) 自身と、\(p\) およびその導関数に関して線形な発散とに置き換えられる。したがって
\begin{equation}
\begin{aligned}
\sum_i\psi_i\bd u_i
={}&\sum_\lambda\left\{
a_i^{(\lambda)}\psi_i-\frac{\p}{\p x}\bigl(b_i^{(\lambda)}\psi_i\bigr)
+\cdots+(-1)^\sigma
\frac{\p^\sigma}{\p x^\sigma}\bigl(c_i^{(\lambda)}\psi_i\bigr)
\right\}p^{(\lambda)}+\Div\Gamma
\end{aligned}
\tag{14}
\end{equation}
を得る。ただし括弧内では \(i\) について和をとる。これを (12) と組み合わせると
\begin{equation}
\sum_\lambda\left\{
a_i^{(\lambda)}\psi_i-\frac{\p}{\p x}\bigl(b_i^{(\lambda)}\psi_i\bigr)
+\cdots+(-1)^\sigma
\frac{\p^\sigma}{\p x^\sigma}\bigl(c_i^{(\lambda)}\psi_i\bigr)
\right\}p^{(\lambda)}
=\Div(B-\Gamma)
\tag{15}
\end{equation}
となる。いま (15) の \(n\)-重積分を任意の領域上で作り、\(p(x)\) を、\(B-\Gamma\) に現れるすべての導関数とともに境界で消えるように選ぶ。発散の積分は境界積分に帰着されるので、(15) の左辺の積分も、十分多くの導関数とともに境界で消える任意の \(p(x)\) について消える。通常の議論から、被積分関数はすべての \(p(x)\) について消えることが従い、\(\rho\) 個の関係
\begin{equation}
\sum_i\left\{
a_i^{(\lambda)}\psi_i-\frac{\p}{\p x}\bigl(b_i^{(\lambda)}\psi_i\bigr)
+\cdots+(-1)^\sigma
\frac{\p^\sigma}{\p x^\sigma}\bigl(c_i^{(\lambda)}\psi_i\bigr)
\right\}=0,\qquad(\lambda=1,2,\ldots,\rho)
\tag{16}
\end{equation}
が得られる。これらが、\(I\) が \(\G_{\infty\rho}\) に関して不変であるときの、ラグランジュ式とその導関数との間の求める従属関係である。線形独立性は前と同様に従う。というのも、逆定理は (12) へ戻すからであり、また infinitesimal 変換から有限変換へふたたび結論できるからである。このことは \S{}4 でより詳しく行う。したがって、\(\G_{\infty\rho}\) については、\(\rho\) 個の任意変換がすでに infinitesimal 変換に現れる。(15) と (16) からさらに \(\Div(B-\Gamma)=0\) が従う。

「混合群」に対応して \(\D x\) と \(\D u\) を \(\eps\) と \(p(x)\) に関して線形にとるなら、一度 \(p(x)\) を零に、もう一度 \(\eps\) を零におくことにより、発散関係 (13) と従属関係 (16) の双方が存在することがわかる。

\section*{\S{}3. 有限群の場合の逆}

逆を証明するには、まず本質的に前の考察を逆順にたどらなければならない。(13) が成り立つことから、\(\eps\) を掛けて加えると (12) が成り立つ。さらに恒等式 (3) により関係
\[
\bd f+\Div(A-B)=0
\]
を得る。したがって \(\D x=\frac1f(A-B)\) とおけば (11) に到達し、積分により最終的に (7), \(\D I=0\) を得る。すなわち \(I\) は \(\D x,\D u\) で決まる infinitesimal 変換に関して不変である。ここで \(\D u\) は (9) により \(\D x\) と \(\bd u\) から決まり、\(\D x,\D u\) はパラメータに関して線形になる。しかし \(\D I=0\) は、よく知られているように、同時系
\begin{equation}
\frac{\dd x_i}{\dd t}=\D x_i,\qquad
\frac{\dd u_i}{\dd t}=\D u_i,\qquad
\left\{\begin{array}{l}
x_i=y_i,\\
u_i=v_i
\end{array}\right\}
\quad(t=0\text{ のとき}).
\tag{17}
\end{equation}
の積分により生じる有限変換に関する \(I\) の不変性を含意する。これらの有限変換は、\(t\eps_1,\ldots,t\eps_\rho\) という組合せ、すなわち \(a_1,\ldots,a_\rho\) という \(\rho\) 個のパラメータを含む。\(\rho\) 個、しかも \(\rho\) 個だけの線形独立な発散関係 (13) があるという仮定から、さらに、有限変換は、それが導関数 \(\p u/\p x\) を含まないかぎり、常に群をなすことが従う。反対の場合、Lie の括弧過程によって生じる少なくとも一つの infinitesimal 変換が、残りの \(\rho\) 個の線形結合ではなくなる。そして \(I\) もこの変換を許すので、\(\rho\) 個を超える線形独立な発散関係が存在することになる。あるいは、この infinitesimal 変換は \(\bd u=0,\ \Div(f\D x)=0\) という特殊な形であろう。しかしその場合、\(\D x\) または \(\D u\) は導関数に依存し、仮定に反する。\(\D x\) または \(\D u\) に導関数が現れる場合にこの事態が起こりうるかどうかは未決としておかなければならない。その場合、群性をふたたび得るためには、上で決められた \(\D x\) に、\(\Div(f\D x)=0\) を満たすすべての関数 \(\D x\) をさらに加えなければならないが、こうして加えられるパラメータは約束により数えない。これで逆が証明された。

この逆からまた、\(\D x\) と \(\D u\) を実際にパラメータに関して線形と仮定してよいことも従う。もし \(\D u\) と \(\D x\) が \(\eps\) の高次形式であったなら、\(\eps\) の冪積の線形独立性により、(13) と完全に対応する関係が、ただしより多数に、従うであろう。逆定理により、それらは infinitesimal 変換がパラメータを線形に含むような群に関する \(I\) の不変性を含意する。この群がちょうど \(\rho\) 個のパラメータを含むべきなら、\(\eps\) の高次項から本来得られた発散関係の間に線形従属がなければならない。

さらに、\(\D x\) と \(\D u\) が \(u\) の導関数も含む場合、有限変換は \(u\) の無限に多くの導関数に依存しうることに注意すべきである。実際この場合、(17) を積分して
\[
\frac{\dd^2x_i}{\dd t^2},\quad \frac{\dd^2u_i}{\dd t^2}
\]
を決めるとき、
\[
\D\!\left(\frac{\p u}{\p x_\lambda}\right)
=\frac{\p\,\D u}{\p x_\lambda}
-\sum_\nu\frac{\p u}{\p x_\nu}
\frac{\p\,\D x_\nu}{\p x_\lambda}
\]
へ導かれるので、一般には各段階で \(u\) の導関数の数が増える。例は
\[
f=\frac12 u'^2,\qquad
\psi=-u'',\qquad
\psi\cdot x=\frac{d}{d x}(u-u'x);\qquad
\bd u=x\cdot\eps,
\]
\[
\varDelta x=\frac{-2u}{u'^2}\eps;\qquad
\varDelta u=\left(x-\frac{2u}{u'}\right)\cdot\eps.
\]
発散のラグランジュ式は恒等的に消えるから、逆は最後に次を示す。\(I\) が \(\G_\rho\) を許すなら、\(I\) と境界積分だけ、すなわち発散の積分だけ異なる任意の積分も、同じ \(\bd u\) をもつ \(\G_\rho\) を許す。ただしその infinitesimal 変換は一般に \(u\) の導関数を含む。したがって前の例に対応して
\[
f^*=\frac12\left\{u'^2-\frac{\dd}{\dd x}\left(\frac{u^2}{x}\right)\right\}
\]
は infinitesimal 変換 \(\D u=x\eps,\ \D x=0\) を許す。一方、\(f\) に対応する infinitesimal 変換には \(u\) の導関数が現れる。

変分問題へ移る、すなわち \(\psi_i=0\) とおくと、\footnote{\(\psi_i=0\)、またはやや一般に \(\psi_i=T_i\)（ここで \(T_i\) は新たに付加された関数）は、物理学では「場の方程式」と呼ばれる。\(\psi_i=T_i\) の場合、恒等式 (13) は \(\Div B^{(\lambda)}=\sum_i T_i\bd u_i^{(\lambda)}\) という方程式に移り、物理学ではこれもなお保存則と呼ばれる。} (13) は方程式
\[
\Div B^{(1)}=0,\ldots,\qquad \Div B^{(\rho)}=0
\]
になる。これはしばしば「保存則」と呼ばれる。一次元の場合にはこれは
\[
B^{(1)}=\mathrm{const.},\ldots,\qquad B^{(\rho)}=\mathrm{const.}
\]
を与える。ここで \(B\) は、\(\D u\) と \(\D x\) が \(f\) に現れる \(\kappa\) 階導関数より高い導関数を含まないかぎり、(6) により高々 \((2\kappa-1)\) 階導関数までの \(u\) を含む。\(\psi\) は一般に \(2\kappa\) 階導関数を含むので、\footnote{\(f\) が \(\kappa\) 階導関数に関して非線形である場合。} こうして \(\rho\) 個の第一積分の存在が得られる。それらの間に非線形従属がありうることは、再び上の \(f\) が示す。線形独立な \(\D u=\eps_1,\ \D x=\eps_2\) には、線形独立な関係
\[
u''=\frac{\dd}{\dd x}u',\qquad
u''u'=\frac12\frac{\dd}{\dd x}(u')^2
\]
が対応する。一方、第一積分 \(u'=\mathrm{const.},\ u'^2=\mathrm{const.}\) の間には非線形従属がある。これは \(\D u,\D x\) が \(u\) の導関数を含まない初等的な場合である。\footnote{そうでない場合、さらに各 \(\lambda\) について \(u''u'^{\lambda-1}=\mathrm{const.}\) があり、これは
\[
u''(u')^{\lambda-1}=\frac1\lambda\frac{\dd}{\dd x}(u')^\lambda
\]
に対応する。}

\section*{\S{}4. 無限群の場合の逆}

まず \(\D x\) と \(\D u\) の線形性の仮定が制限でないことを示す必要がある。ここではこれは、逆定理を使わなくても、\(\G_{\infty\rho}\) が形式的に \(\rho\) 個、しかも \(\rho\) 個だけの任意関数に依存するという事実から従う。すなわち非線形の場合には、変換を合成すると最低次の項が加わり、任意関数の数が増える。例えば
\[
\begin{aligned}
y&=A\!\left(x,u,\frac{\p u}{\p x},\cdots;p\right)\\
&=x+\sum a(x,u,\ldots)p^\nu
+b(x,u,\ldots)p^{\nu-1}\frac{\p p}{\p x}
+c\,p^{\nu-2}\left(\frac{\p p}{\p x}\right)^2+\cdots
+d\left(\frac{\p p}{\p x}\right)^\nu+\cdots,
\end{aligned}
\]
ただし
\[
p^\nu=(p^{(1)})^{\nu_1}\cdots(p^{(\rho)})^{\nu_\rho},
\]
であり、対応して
\[
v=B\!\left(x,u,\frac{\p u}{\p x},\ldots;p\right)
\]
であるとする。このとき
\[
z=A\!\left(y,v,\frac{\p v}{\p y},\ldots;q\right)
\]
との合成により、最低次の項は
\[
z=x+\sum a(p^\nu+q^\nu)+
b\left\{p^{\nu-1}\frac{\p p}{\p x}+q^{\nu-1}\frac{\p q}{\p x}\right\}
+c\left\{p^{\nu-2}\left(\frac{\p p}{\p x}\right)^2
+q^{\nu-2}\left(\frac{\p q}{\p x}\right)^2\right\}+\cdots
\]
を与える。ここで \(a\) と \(b\) 以外のある係数が零でなく、したがって
\[
p^{\nu-\sigma}\left(\frac{\p p}{\p x}\right)^\sigma
+q^{\nu-\sigma}\left(\frac{\p q}{\p x}\right)^\sigma
\]
という項が \(\sigma>1\) について実際に現れるなら、それは単一の関数の微分商としても、そのような商の冪の積としても書けない。したがって仮定に反して任意関数の数が増えている。\(a\) と \(b\) 以外の係数がすべて零なら、指数 \(\nu_1,\ldots,\nu_\rho\) の値に応じて、第二項は第一項の微分商である（例えば \(\G_{\infty1}\) では常にそうである）ので線形性が実際に現れるか、あるいはここでも任意関数の数が増える。したがって \(p(x)\) に関する線形性により、infinitesimal 変換は線形偏微分方程式系を満たす。そして群性が成り立つので、それらは Lie の意味での「infinitesimal 変換の無限群」をなす（\emph{Grundlagen}, \S{}10）。

逆は有限群の場合と同様に得られる。従属関係 (16) の存在から、\(p^{(\lambda)}(x)\) を掛けて加えると、恒等変換 (14) により
\[
\sum_i\psi_i\bd u_i=\Div\Gamma
\]
が導かれる。ここから \S{}3 と同様に \(\D x\) と \(\D u\) の決定、およびこれらの infinitesimal 変換に関する \(I\) の不変性が得られる。これらは実際に \(\rho\) 個の任意関数とその \(\sigma\) 階までの導関数に線形に依存する。これらの infinitesimal 変換が、\(\p u/\p x,\ldots\) を含まないとき、確かに群をなすことは \S{}3 と同じく従う。そうでなければ合成によりさらに多くの任意関数が現れるが、仮定では従属関係 (16) は \(\rho\) 個だけでなければならないからである。したがってそれらは「infinitesimal 変換の無限群」をなす。しかしそのような群は（\emph{Grundlagen}, 定理 VII, p. 391）、Lie の意味でそれによって定義されるある「有限変換の無限群 \(\G\)」の最も一般の infinitesimal 変換からなる。各有限変換は infinitesimal なものから生成される（\emph{Grundlagen}, \S{}7）。\footnote{ここから、とくに、\(\G_{\infty\rho}\) の infinitesimal 変換 \(\D x,\D u\) から生成される群 \(\G\) は、再び \(\G_{\infty\rho}\) に戻ることが従う。というのも \(\G_{\infty\rho}\) は、任意関数に依存し \(\D x,\D u\) と異なる infinitesimal 変換を含まず、またパラメータに依存する独立な変換も含みえないからである。もしそうであれば混合群になってしまう。しかし前の議論により、有限変換は infinitesimal 変換によって決定される。} したがってこれは同時系
\[
\frac{d x_i}{d t}=\varDelta x_i;\qquad
\frac{d u_i}{d t}=\varDelta u_i\quad
\left(\begin{array}{l}
x_i=y_i,\\
u_i=v
\end{array}\right\}
\quad(t=0\text{ のとき}).
\]
の積分によって生じる。ただし、任意の \(p(x)\) をさらに \(t\) に依存させる必要があるかもしれない。したがって \(\G\) は実際に \(\rho\) 個の任意関数に依存する。とくに \(p(x)\) を \(t\) に依存しないものとしてよいなら、この依存は任意関数 \(q(x)=t\,p(x)\) に解析的である。\footnote{この後者の場合が常に起こるのではないかという問題は、Lie によって別の定式化で提起された（\emph{Grundlagen}, \S{}7 および \S{}13 末尾）。} \(\p u/\p x,\ldots\) が現れるなら、同じ結論を引き出す前に infinitesimal 変換 \(\bd u=0,\ \Div(f\D x)=0\) を付加する必要があるかもしれない。

Lie の例（\emph{Grundlagen}, \S{}7）に従って、ここで、明示公式にまで到達できるかなり一般的な場合を示す。その公式は、任意関数の導関数が \(\sigma\) 階まで現れることも示している。したがってこの場合、逆は完全である。これは、すべての \(x\) の変換の群と、それによって「誘導される」\(u\) の変換の群が対応する infinitesimal 変換群である。すなわち \(\D u\)、したがって \(u\) が、\(\D x\) に現れる任意関数だけに依存するような \(u\) の変換である。なお \(\D u\) には導関数 \(\p u/\p x,\ldots\) が現れないと仮定する。したがって
\[
\varDelta x_i=p^{(i)}(x);\qquad
\varDelta u_i=
\sum_{\lambda=1}^{n}\left\{
a^{(\lambda)}(x,u)p^{(\lambda)}
+b^{(\lambda)}\frac{\p p^{(\lambda)}}{\p x}
+\cdots+
c_i^{(\lambda)}\frac{\p^\sigma p^{(\lambda)}}{\p x^\sigma}
\right\}.
\]
infinitesimal 変換 \(\D x=p(x)\) は任意の \(g(y)\) をもつすべての変換 \(x=y+g(y)\) を生成するので、特に \(p(x)\) を \(t\) の関数として、一次元パラメータ群
\begin{equation}
x_i=y_i+t\,g_i(y)
\tag{18}
\end{equation}
が生成されるよう決定できる。これは \(t=0\) で恒等変換となり、\(t=1\) で求める \(x=y+g(y)\) になる。(18) を微分すると
\begin{equation}
\frac{\dd x_i}{\dd t}=g_i(y)=p^{(i)}(x,t)
\tag{19}
\end{equation}
を得る。ここで \(p(x,t)\) は (18) の反転により \(g(y)\) から決まる。逆に (18) は、\(t=0\) で \(x_i=y_i\) という副条件により、(19) から生じる。この副条件により積分は一意に定まる。(18) により、\(\D u\) における \(x\) は積分定数 \(y\) と \(t\) で置き換えられる。そのとき \(g(y)\) はちょうど \(\sigma\) 階導関数まで現れる。なぜなら
\[
\frac{\p p}{\p x}=\sum\frac{\p g}{\p y_\varkappa}\frac{\p y_\varkappa}{\p x}
\]
において \(\p y/\p x\) は \(\p x/\p y\) によって表され、一般に \(\p^\sigma p/\p x^\sigma\) は
\[
\frac{\p g}{\p y},\ldots,\frac{\p x}{\p y},\ldots,\frac{\p^\sigma x}{\p y^\sigma}
\]
における値で置き換えられるからである。\(u\) の決定については、したがって方程式系
\[
\frac{d u_i}{d t}
=
F_i\!\left(g(y),\frac{\p g}{\p y},\cdots\frac{\p^\sigma g}{\p y^\sigma},u,t\right),
\qquad (u_i=v_i\quad t=0\text{ のとき}),
\]
を得る。この中では \(t\) と \(u\) だけが変数であり、\(g(y),\ldots\) は係数域に属する。したがって積分により
\[
u_i=v_i+
B_i\!\left(v,g(y),\frac{\p g}{\p y},\cdots\frac{\p^\sigma g}{\p y^\sigma},t\right)_{t=1}.
\]
が与えられる。こうして、任意関数の導関数 \(\sigma\) 階にちょうど依存する変換が得られる。(18) により \(g(y)=0\) の場合に恒等変換が含まれる。また群性は、いま述べた手続きがすべての変換 \(x=y+g(y)\) を与え、それによって \(u\) の誘導変換が一意に決定されるので、群 \(\G\) が尽くされるという事実から従う。

逆定理から、任意関数が \(u,\p u/\p x,\ldots\) ではなく \(x\) のみに依存すると仮定しても制限でないことも付随的に従う。後者の場合、恒等変換 (14)、したがって (15) において、\(p^{(\lambda)}\) のほかに \(\p p^{(\lambda)}/\p u,\ \p p^{(\lambda)}/\p(\p u/\p x),\ldots\) という導関数が現れるであろう。ここで \(p^{(\lambda)}\) を順に、\(u,\p u/\p x,\ldots\) に関する 0 次、1 次、さらに高次のものとし、その係数を \(x\) の任意関数とするなら、再び従属関係 (16) が生じる。ただし数は増える。しかし前の逆定理により、それらは \(x\) のみに依存する任意関数と組み合わせることによって先の場合に戻すことができる。同様に、混合群には従属関係と独立な発散関係の同時出現が対応することが示される。\footnote{\S{}3 と同じく、ここでも逆定理から、\(I\) のほか、それと発散の積分だけ異なる任意の積分 \(I^*\) も、同じ \(\bd u\) をもつ無限群を許すことが従う。ただしここでは \(\D x\) と \(\D u\) は一般に \(u\) の導関数を含む。Einstein は、エネルギー定理のより簡単な定式化を得るため、このような積分 \(I^*\) を一般相対性理論に導入した。Klein の第二ノートの記法をそのまま用いて、この \(I^*\) が許す infinitesimal 変換を与える。積分 \(I=\int\!\cdots\!\int K\,\dd\omega=\int\!\cdots\!\int \mathfrak{R}\,\dd S\) は、すべての \(w\) の変換の群と、それによって \(g_{\mu\nu}\) に誘導される群を許す。これに対応する従属関係（Klein の (30)）は
\[
\sum \mathfrak{R}_{\mu\nu}g_\tau^{\mu\nu}
+2\sum\frac{\p g^{\mu\nu}\mathfrak{R}_{\mu\tau}}{\p w^\sigma}=0.
\]
いま \(I^*=\int\!\cdots\!\int \mathfrak{R}^*\,\dd S\)、ただし \(\mathfrak{R}^*=\mathfrak{R}+\Div\) であり、したがって \(\mathfrak{R}_{\mu\nu}^*=\mathfrak{R}_{\mu\nu}\) である。ここで \(\mathfrak{R}_{\mu\nu}^*,\mathfrak{R}_{\mu\nu}\) はそれぞれラグランジュ式を表す。よって与えられた従属関係は \(\mathfrak{R}_{\mu\nu}^*\) に対する従属関係でもある。さらに \(p^\tau\) を掛けて加え、積の微分の変形を逆にたどると
\[
\sum \mathfrak{R}_{\mu\nu}p^{\mu\nu}
+2\Div\!\left(\sum g^{\mu\sigma}\mathfrak{R}_{\mu\tau}p^\tau\right)=0,
\]
\[
\delta\mathfrak{R}^*+\Div\!\left(\sum 2g^{\mu\sigma}\mathfrak{R}_{\mu\tau}p^\tau
-\frac{\p\mathfrak{R}^*}{\p g^{\mu\nu}_{\sigma}}p^{\mu\nu}\right)=0.
\]
Lie の微分方程式
\(\delta\mathfrak{R}^*+\Div(\mathfrak{R}^*\D w)=0\) と比較すると、\(I^*\) が許す infinitesimal 変換として
\[
\varDelta w^\sigma=\frac1{\mathfrak{R}^*}\cdot\left(\sum (2g^{\mu\sigma}\mathfrak{R}_{\mu\tau}p^\tau
-\frac{\p\mathfrak{R}^*}{\p g^{\mu\nu}_{\sigma}}p^{\mu\nu}\right),
\qquad
\varDelta g^{\mu\nu}=p^{\mu\nu}+\sum g^{\mu\nu}_{\sigma}\varDelta w^\sigma,
\]
を得る。したがってこれらの infinitesimal 変換は \(g^{\mu\nu}\) の一階および二階導関数に依存し、任意の \(p\) を一階導関数まで含む。}

\section*{\S{}5. 関係の各成分の不変性}

群 \(\G\) を最も単純で通常考察される場合、すなわち変換に \(u\) の導関数を含まず、変換された独立変数が \(u\) ではなく \(x\) だけに依存する場合に特殊化すると、公式中の個々の成分の不変性を推論できる。まずよく知られた議論により、
\[
\int\!\cdots\!\int\left(\sum\psi_i\delta u_i\right)\dd x
\]
の不変性が得られ、したがって \(\sum\psi_i\delta u_i\) の相対不変性が得られる。\footnote{すなわち、\(\sum\psi_i\delta u_i\) は変換のもとで因子を受け取り、代数的不変式論でつねにいう相対不変性を表す。}ここで \(\delta\) は任意の変分を表す。すなわち一方では
\[
\delta I
=\int\!\cdots\!\int \delta f\!\left(x,u,\frac{\p u}{\p x},\cdots\right)d x
=\int\!\cdots\!\int \delta f\!\left(y,v,\frac{\p v}{\p y},\cdots\right)d y.
\]
他方、\(\delta u,\delta(\p u/\p x),\ldots\) が境界で消えるとき、\(\delta u,\delta(\p u/\p x),\ldots\) の線形同次変換により、それに対応する \(\delta v,\delta(\p v/\p y),\ldots\) も境界で消えるので、
\[
\int\!\cdots\!\int \delta f\!\left(x,u,\frac{\p u}{\p x},\cdots\right)d x
=
\int\!\cdots\!\int\left(\sum\psi_i(u,\ldots)\delta u_i\right)d x,
\]
\[
\int\!\cdots\!\int \delta f\!\left(y,v,\frac{\p v}{\p x},\cdots\right)d y
=
\int\!\cdots\!\int\left(\sum\psi_i(v,\ldots)\delta v_i\right)d y.
\]
したがって境界で消える \(\delta u,\delta(\p u/\p x),\ldots\) について
\[
\int\!\cdots\!\int\left(\sum\psi_i(u,\ldots)\delta u_i\right)d x
=
\int\!\cdots\!\int\left(\sum\psi_i(v,\ldots)\delta v_i\right)d y
\]
\[
=
\int\!\cdots\!\int\left(\sum\psi_i(v,\ldots)\delta v_i\right)
\left|\frac{\p y_i}{\p x_\varkappa}\right|d x.
\]
第三の積分で \(y,v,\delta v\) を \(x,u,\delta u\) で表し、それを第一の積分に等置すれば、境界で消え、その他では任意の \(\delta u\) について関係
\[
\int\!\cdots\!\int\left(\sum\chi_i(u,\ldots)\delta u_i\right)d x=0
\]
を得る。ここから、よく知られているように、被積分関数は任意の \(\delta u\) に対して消える。したがって関係
\[
\sum\psi_i(u,\ldots)\delta u_i
=
\left|\frac{\p y_i}{\p x_\varkappa}\right|
\left(\sum\psi_i(v,\ldots)\delta v_i\right)
\]
は \(\delta u\) に関して恒等的に成り立つ。これは \(\sum\psi_i\delta u_i\) の相対不変性を表し、したがって
\[
\int\!\cdots\!\int\left(\sum\psi_i\delta u_i\right)\dd x
\]
の不変性を表す。\footnote{これらの議論は、\(y\) も \(u\) に依存する場合には失敗する。そのとき \(\delta f(y,v,\p v/\p y,\ldots)\) には \(\sum(\p f/\p y)\delta y\) という項も含まれ、発散変換はラグランジュ式へ導かないからである。\(u\) の導関数を許す場合にも同様に失敗する。このとき \(\delta v\) は \(\delta u,\delta(\p u/\p x),\ldots\) の線形結合となり、さらに発散変換を行った後に初めて \(\int\!\cdots\!\int(\sum\chi_i(u,\ldots)\delta u_i)\dd x=0\) という恒等式を導くので、右辺にはやはりラグランジュ式が現れない。\(\int\!\cdots\!\int(\sum\psi_i\delta u_i)\dd x\) の不変性からすでに発散関係の存在を結論できるかという問題は、逆定理により、同じ \(\D u,\D x\) に至る必要はないが同じ \(\bd u\) には至る群に関する \(I\) の不変性をそれが含意するかという問題と同値である。単積分で \(f\) に一階導関数だけが現れる特殊な場合には、有限群について、ラグランジュ式の不変性から第一積分の存在を推論できる（例えば Engel, Gött. Nachr. 1916, p. 270 を参照）。}

これを導かれた発散関係と従属関係に適用するには、まず \(\D u,\D x\) から導かれる \(\bd u\) が、実際に変分 \(\delta u\) の変換法則を満たすことを示さなければならない。ただし、\(\bd v\) におけるパラメータ、または任意関数は、\(y,v\) における同様の infinitesimal 変換群に対応するように決定されているものとする。\(T_q\) を \(x,u\) を \(y,v\) に移す変換とし、\(T_p\) を \(x,u\) における infinitesimal 変換とする。このとき \(y,v\) における同様のものは \(T_r=T_qT_pT_q^{-1}\) で与えられ、ここでパラメータまたは任意関数 \(r\) は \(p\) と \(q\) によって決まる。式で書けば次のようになる：
\[
T_p:\ \xi=x+\varDelta x(x,p),\qquad u^*=u+\varDelta u(x,u,p);
\]
\[
T_q:\ y=A(x,q),\qquad v=B(x,u,q);
\]
\[
T_qT_p:\ \eta=A(x+\D x(x,p),q),\quad
v^*=B(x+\varDelta x(p),\,u+\varDelta u(p),q).
\]
これから \(T_r=T_qT_pT_q^{-1}\) が生じるので、
\[
\eta=y+\varDelta y(r);\qquad v^*=v+\varDelta v(r),
\]
である。ただし \(T_q\) の逆により \(x\) を \(y\) の関数とみなし、infinitesimal な項だけを残す。したがって恒等式
\begin{equation}
\eta=y+\D y(r)=y+\sum\frac{\p A(x,q)}{\p x}\D x(p)
\tag{20a}
\end{equation}
および
\begin{equation}
v^*=v+\D v(r)=v+\sum\frac{\p B(x,u,q)}{\p x}\D x(p)
+\sum\frac{\p B(x,u,q)}{\p u}\D u(p)
\tag{20b}
\end{equation}
がある。ここで \(\xi=x+\D x\) を \(\xi-\D \xi\) に置き換え、\(\xi\) が \(x\) に戻り、したがって \(\D x\) が消えるようにすると、第一式 (20) により \(\eta\) も \(y=\eta-\D\eta\) に戻る。この置換の下で \(\D u(p)\) が \(\bd u(p)\) へ移るなら、\(\D v(r)\) も \(\bd v(r)\) へ移り、第二式 (20) は
\[
v+\bd v(y,v,\ldots r)=v+\sum\frac{\p B(x,u,q)}{\p u}\,\bd u(p),
\]
\[
\bd v(y,v,\ldots r)=\sum\frac{\p B}{\p u_\varkappa}\,\bd u_\varkappa(x,u,p)
\]
を与える。したがって \(\bd v\) がパラメータまたは任意関数 \(r\) に依存するものとして取られているかぎり、変分に関する変換公式は実際に満たされる。\footnote{結論が成り立つためには \(y\) が \(u\) などに依存しないと仮定しなければならないことが、ここでも見られる。例として、Klein が示した \(\delta g^{\mu\nu}\) と \(\delta q_\sigma\) を挙げられる。これらは \(p\) をベクトル変換に従わせると、変分に関する変換を満たす。}

とくに \(\sum\psi_i\bd u_i\) は相対不変であることが従う。したがって (12) により、発散関係は \(y,v\) でも満たされているので、\(\Div B\) の相対不変性も従う。さらに (14) と (13) により、\(\Div\Gamma\) と、\(p^{(\lambda)}\) と組み合わされた従属関係の左辺の相対不変性も従う。変換された公式では、任意の \(p(x)\)（またはパラメータ）は至る所で \(r\) に置き換えられる。ここから、\(\Div(B-\Gamma)\) の相対不変性も従う。すなわち、恒等的には消えないが、その発散は恒等的に消える関数系 \(B-\Gamma\) の発散の相対不変性である。

一次元の場合および有限群の場合には、\(\Div B\) の相対不変性から、第一積分の不変性についてさらに結論を引き出すことができる。infinitesimal 変換に対応するパラメータの変換は、(20) により線形同次となる。またすべての変換は可逆であるから、\(\eps\) も変換されたパラメータ \(\eps^*\) に関して線形同次である。\(\psi=0\) とおいてもこの可逆性は確かに保たれる。なぜなら (20) には \(u\) の導関数が現れないからである。
\[
\Div B(x,u,\ldots\eps)=\frac{d y}{d x}\cdot\Div B(y,v,\ldots\eps^*)
\]
で \(\eps^*\) の係数を比較すると、量 \(\frac{\dd}{\dd y}B^{(\lambda)}(y,v,\ldots)\) も \(\frac{\dd}{\dd x}B^{(\lambda)}(x,u,\ldots)\) の線形同次関数となる。したがって
\[
\frac{\dd}{\dd x}B^{(\lambda)}(x,u,\ldots)=0
\quad\text{または}\quad B^{(\lambda)}(x,u)=\mathrm{const.}
\]
から
\[
\frac{\dd}{\dd y}B^{(\lambda)}(y,v,\ldots)=0
\quad\text{または}\quad B^{(\lambda)}(y,v)=\mathrm{const.}
\]
が従う。したがって \(\G_\rho\) に対応する \(\rho\) 個の第一積分は、それぞれ群を許すので、さらに先の積分も単純化される。この最も単純な例は、\(f\) が \(x\) を含まない、またはある \(u\) を含まない場合であり、それぞれ infinitesimal 変換 \(\D x=\eps,\D u=0\)、\(\D x=0,\D u=\eps\) に対応する。そのとき \(\bd u=-\eps\,\dd u/\dd x\)、または \(\eps\) である。そして \(B\) は \(f\) と \(\bd u\) から微分と有理演算によって導かれるので、\(B\) も \(x\)、または \(u\) を含まず、対応する群を許す。\footnote{第一積分の存在がすでに \(\int(\sum\psi_i\delta u_i)\dd x\) の不変性から従う場合、それらは完全な群 \(\G_\rho\) を許さない。例えば \(\int(u''\delta u)\dd x\) は infinitesimal 変換 \(\D x=\eps_2,\D u=\eps_1+x\eps_3\) を許す。一方、\(\D x=0,\D u=x\eps_3\) に対応する第一積分 \(u-u'x=\mathrm{const.}\) は、\(u\) と \(x\) を明示的に含むので、他の二つの infinitesimal 変換を許さない。この第一積分に対応して、\(f\) に対する導関数を含む infinitesimal 変換がある。こうして、\(\int\!\cdots\!\int(\sum\psi_i\delta u_i)\dd x\) の不変性は、いずれにせよ \(I\) の不変性より少ないことしか達成しないことがわかる。これは前の注で提起された問題に関係する。}

\section*{\S{}6. Hilbert の主張}

以上から最後に、真性エネルギー定理の失敗と「一般相対性」との関連に関する Hilbert の主張（Klein の第一ノート、Göttinger Nachr. 1917、回答の第一段落）の証明が、一般化された群論的定式化で得られる。

積分 \(I\) が \(\G_{\infty\rho}\) を許すとし、\(\G_\sigma\) を、任意関数の特殊化によって生じる任意の有限群、したがって \(\G_{\infty\rho}\) の部分群とする。無限群 \(\G_{\infty\rho}\) には従属関係 (16) が対応し、有限群 \(\G_\sigma\) には発散関係 (13) が対応する。逆に、任意の発散関係の存在から、ある有限群に関する \(I\) の不変性が従う。この有限群が \(\G_\sigma\) と同一であるのは、\(\bd u\) が \(\G_\sigma\) から生じるものの線形結合である場合、かつその場合に限る。したがって \(\G_\sigma\) に関する不変性は、(13) と異なる発散関係へは導きえない。しかし (16) の存在から、任意の \(p(x)\) について、\(\G_{\infty\rho}\) の infinitesimal 変換 \(\D u,\D x\) に関する \(I\) の不変性が従うので、とくに特殊化によってこれらから生じる \(\G_\sigma\) の infinitesimal 変換に関して、したがって \(\G_\sigma\) に関して、\(I\) はすでに不変である。したがって発散関係
\[
\sum\psi_i\bd u_i^{(\lambda)}=\Div B^{(\lambda)}
\]
は、(16) の従属関係の帰結でなければならない。(16) はまた
\[
\sum\psi_i a_i^{(\lambda)}=\Div\chi^{(\lambda)}
\]
とも書ける。ここで \(\chi^{(\lambda)}\) はラグランジュ式とその導関数の線形結合である。\(\psi\) は (13) と (16) の双方に線形に現れるので、発散関係は特に従属関係 (16) の線形結合でなければならない。したがって
\[
\Div B^{(\lambda)}=\Div\left(\sum \alpha\cdot\chi^{(\varkappa)}\right);
\]
かつ \(B^{(\lambda)}\) 自身は、\(\chi\)、すなわちラグランジュ式とその導関数から、および発散が恒等的に消える関数から線形に構成されている。例として \S{}2 の末尾に現れる \(B-\Gamma\) があり、ここでは \(\Div(B-\Gamma)=0\) で、同時にその発散は不変量の性格をもつ。\(B^{(\lambda)}\) が示された仕方でラグランジュ式とその導関数から構成されるような発散関係を「非真性」と呼び、残りのすべてを「真性」と呼ぶことにする。

逆に、発散関係が従属関係 (16) の線形結合であり、したがって「非真性」であるなら、\(\G_\sigma\) に関する不変性は \(\G_{\infty\rho}\) に関する不変性から従う。\(\G_\sigma\) は \(\G_{\infty\rho}\) の部分群となる。したがって、有限群 \(\G_\sigma\) に対応する発散関係が非真性になるのは、\(\G_\sigma\) が、\(I\) が不変であるような無限群の部分群である場合、かつその場合に限る。

群を特殊化すると、Hilbert のもとの主張がこれから従う。「平行移動群」とは有限群
\[
y_i=x_i+\eps_i;\qquad v_i(y)=u_i(x)
\]
をいう。したがって
\[
\varDelta x_i=\eps_i,\qquad \varDelta u_i=0,\qquad
\bd u_i=-\sum_\lambda\frac{\p u_i}{\p x_\lambda}\eps_\lambda.
\]
平行移動群に関する不変性は、知られているように、
\[
I=\int\!\cdots\!\int f\!\left(x,u,\frac{\p u}{\p x},\ldots\right)\dd x
\]
において \(f\) に \(x\) が明示的に現れないことを意味する。対応する \(n\) 個の発散関係
\[
\sum\psi_i\frac{\p u_i}{\p x_\lambda}=\Div B^{(\lambda)}
\qquad(\lambda=1,2,\ldots n)
\]
を「エネルギー関係」と呼ぶことにする。なぜなら、変分問題に対応する「保存則」\(\Div B^{(\lambda)}=0\) が「エネルギー定理」に、\(B^{(\lambda)}\) が「エネルギー成分」に対応するからである。したがって次を得る。\(I\) が平行移動群を許すなら、エネルギー関係が非真性となるのは、\(I\) が、平行移動群を部分群として含む無限群に関して不変である場合、かつその場合に限る。\footnote{古典力学のエネルギー定理、および古い「相対性理論」のエネルギー定理（そこでは \(\sum \dd x^2\) がそれ自身へ移る）は「真性」である。ここでは無限群が現れないからである。}

このような無限群の例は、すべての \(x\) の変換と、\(u(x)\) に誘導される変換からなる群で与えられる。この群では任意関数 \(p(x)\) の導関数だけが現れる。平行移動群は特殊化 \(p^{(i)}(x)=\eps_i\) によって生じる。ただし、この手段と、\(I\) を境界積分で変えることによって生じる群とにより、すでにこのような群の最も一般のものが与えられているかどうかは、未決にしておかなければならない。示された種類の誘導変換は、例えば \(u\) を「全微分形式」の係数変換に従わせることによって生じる。すなわち、\(\dd x\) のほか高次微分も含む形式
\[
\sum a\,\dd^\lambda x_i+\sum b\,\dd^{\lambda-1}x_i\,\dd x_\varkappa+\cdots
\]
である。より特殊な誘導変換、すなわち \(p(x)\) が一階導関数にだけ現れるものは、通常の微分形式 \(\sum c\,\dd x_{i_1}\cdots \dd x_{i_\lambda}\) の係数変換によって与えられる。これらが通常考えられてきたものである。

さらに示された種類の群として、対数項が現れるため係数変換ではありえないものが、例えば次で与えられる：
\[
y=x+p(x),\qquad
v_i=u_i+\lg(1+p'(x))=u_i+\lg\frac{d y}{d x};
\]
\[
\varDelta x=p(x);\qquad
\varDelta u_i=p'(x)\footnote{これらの無限小変換から、有限変換は \S{}4 末尾に示した方法によって逆向きに計算される。};\qquad
\bd u_i=p'(x)-u_i'p(x).
\]
ここで従属関係 (16) は
\[
\sum_i\left(\psi_i u_i'+\frac{d\psi_i}{d x}\right)=0
\]
となり、非真性エネルギー関係は
\[
\sum\left(\psi_i u_i'+
\frac{d(\psi_i+\mathrm{const.})}{d x}\right)=0
\]
となる。この群の最も簡単な不変積分は
\[
I=\int \frac{e^{-2u_1}}{u_1'-u_2'}\,d x.
\]
最も一般の \(I\) は Lie の微分方程式 (11)
\[
\bd f+\frac{d}{d x}(f\cdot \varDelta x)=0
\]
を積分することによって決まる。\(\D x\) と \(\bd u\) の値を代入し、\(f\) が \(u\) の一階導関数だけに依存すると仮定すると、これは
\[
\frac{\p f}{\p x}p(x)
+\left\{\sum\frac{\p f}{\p u_i}
-\frac{\p f}{\p u_i'}u_i'+f\right\}p'(x)
+\left\{\sum\frac{\p f}{\p u_i''}\right\}p''(x)=0
\]
（\(p(x),p'(x),p''(x)\) に関する恒等式）となる。すでに二つの関数 \(u(x)\) に対して、この方程式系は導関数を実際に含む解をもつ。すなわち
\[
f=(n_1'-u_2')\,
\varPhi\!\left(u_1-u_2,\frac{e^{-u_1}}{u_1'-u_2'}\right),
\]
ここで \(\Phi\) は示された引数の任意関数である。

Hilbert は、真性エネルギー定理の失敗が「一般相対性理論」の特徴である、とその主張を述べている。この主張が文字どおり成り立つためには、「一般相対性」という名称を通常より広く理解し、上に述べた \(n\) 個の任意関数に依存する群にも拡張しなければならない。\footnote{これにより、物理学で普通に使われる「相対性」という語は「群に関する不変性」に置き換えられるべきだ、という Klein の注意がふたたび正しいことが確認される。（Über die geometrischen Grundlagen der Lorentzgruppe. \emph{Jhrber. d. d. Math. Vereinig.} vol. 19, p. 287, 1910; \emph{Physikalische Zeitschrift} に再録。）}

% 組版用マクロ
\providecommand{\p}{\partial}
\providecommand{\frS}{\mathfrak S}
\providecommand{\frp}{\mathfrak p}
\providecommand{\fra}{\mathfrak a}
\providecommand{\frb}{\mathfrak b}
\providecommand{\frc}{\mathfrak c}
\providecommand{\frf}{\mathfrak f}
\providecommand{\frr}{\mathfrak r}
\providecommand{\frm}{\mathfrak m}
\providecommand{\frD}{\mathfrak D}
\providecommand{\calA}{\mathcal A}
\providecommand{\calB}{\mathcal B}
\providecommand{\calN}{\mathcal N}
% 組版用マクロ

\section*{14. 一変数の代数関数の算術的理論と、他の諸理論および数体論との関係}
\begin{center}
\emph{Jahresbericht der Deutschen Mathematiker-Vereinigung} 38 (1919), pp. 182--203
\end{center}

以下の報告は、代数関数の個々の理論のあいだに連結環を得たいという願いから生まれた。同時にそれは、算術的理論の議論が本質的に欠けている Brill--Noether の報告\footnote{古い時代および新しい時代における代数関数論の発展。\emph{Jahresbericht der Deutschen Mathematiker-Vereinigung} vol. 3 (1894)［Brill--Noether として引用する］。}への補足と見ることもできる。\footnote{Kronecker の判別式に関する研究の議論を例外とする。（\emph{J. f. M.} vol. 91.）}

個々の理論の説明としては、Brill--Noether のほか、そこに掲げられた文献を含めて、次の \emph{Enzyklopädie} の論説が考慮される。

Riemann の理論については Wirtinger (II B 2), nos. 1--19；

Weierstrass の理論については Wirtinger (II B 2), nos. 23, 32, 33, 34；

Brill--Noether の代数幾何学的理論については Wirtinger (II B 2), nos. 21--31；Berzolari (III C 4), nos. 23--38；

Dedekind--Weber および Hensel--Landsberg の算術的理論については Hensel (II C 5)［Hensel として引用する］；\footnote{Dedekind--Weber に関する部分は A. Ostrowski による。}

より一般に、代数量の算術的理論については Landsberg (I C 5)、多変数代数関数の算術的理論については Jung (II C 6)。

数体論については Hilbert の「Zahlbericht」\footnote{代数的数体の理論。\emph{Jahresbericht der Deutschen Mathematiker-Vereinigung} vol. 4 (1894/95).}および Dirichlet--Dedekind の数論講義［Dedekind--数論として引用する］\footnote{Braunschweig, Vieweg, 4th ed. 1894.}を参照されたい。

\subsection*{目次}
\begin{enumerate}
\item さまざまな理論の基礎。
\item 代数的数と代数関数とのあいだの平行性（共役元、基底定理、加群概念）。
\item 代数的数と代数関数のイデアル論における平行性と相違。
\item 級数展開の算術的基礎づけとイデアル論との関係。
\item 代数関数の算術的理論と幾何学的理論との比較。異なる不変性概念。
\item 比較の続き。特異点。
\item 剰余定理とイデアル論の諸定理との比較。
\item 代数関数と代数的数における超越的問題。
\end{enumerate}

\subsection*{1. さまざまな理論の基礎}
Riemann の理論は、存在定理（Dirichlet の原理、Schwarz--Neumann の方法）に依拠する、唯一の純粋に超越的な理論である。ここでは代数関数の積分が第一のものであり、代数関数そのものは第二のものである。Riemann 面上の関数は、その零点と極によって特徴づけられる。これは算術的理論における多角形（または除数）に対応する。一方 Weierstrass は、関数そのものを、その零点および極とともに語る。存在定理を別にすれば、主要問題は、与えられた極をもつすべての関数を設定することであり、これは第一種および第二種の積分によって遂行される。さらに、無限遠点を指定したときになお入ってくる任意定数の数が問題になる。この後者の問いには Riemann--Roch の定理が答える。算術的理論ではこれに多角形類（または除数類）の次元の問いが対応し、代数幾何学的理論では完全線形系の次元の問いが対応するが、ここでも答えは Riemann--Roch の定理によって与えられる。関数の実際の構成は、算術的理論では基底定理によって、代数幾何学的理論では剰余定理によって行われる。とくに、第一種の被積分関数の零点には、算術的理論では微分類が対応し、代数幾何学的理論では特殊群の理論、すなわち \((n-3)\) 次の随伴曲線（\(\varphi\)-曲線）によって切り出される点群の束が対応する。

Riemann の理論が歴史的には他の理論に先行したが、存在定理の厳密な証明は後に初めて得られた。それと同様に、後にもその超越的な補助手段によって、今日に至るまで純粋に代数的または算術的な処理ではなお到達できない代数的問題が処理された。ここではとりわけ Hurwitz の特異対応の理論を挙げなければならない。\footnote{A. Hurwitz, Über algebraische Korrespondenzen und das erweiterte Korrespondenzprinzip. \emph{Math. Ann.} vol. 28.（Brill--Noether 第10節を参照。）}

代数的数体の理論における超越的な問題設定との類似は、最後の節で論じる。

残りの理論では、代数関数が第一のものであり、積分は第二のものである。この点で Weierstrass の理論は、冪級数展開と剰余定理という超越的基礎の上に立つ。しかしその後は、考慮すべきすべての関数を明示的に与えることによって、代数的に進む。

Hensel--Landsberg の形の算術的理論も、級数展開という超越的基礎を用いる。しかしこの基礎の上で個々の点に除数を対応させることにより、その後は純粋に算術的に進むので、Weierstrass と Dedekind--Weber との融合と見ることができる。後者と対照的に、この理論は、問題となるすべての基底関数を実際に構成する方法を同時に与える。細部では Dedekind--Weber に比べていくつかの簡略化があり、その本質は「分数的」除数の導入にある。最近 Hensel\footnote{\emph{Mathematische Zeitschrift} vol. 4。この基礎づけは彼の報告にも基礎として用いられている。} は、級数展開の理論を算術的に基礎づけた（収束証明のための優関数法を除く）。これによって理論はほとんど純粋に算術的なものになった。

Brill--Noether の代数幾何学的理論は、枝の「相続く点」という概念だけを前提にする。この概念は通常点における級数展開に帰着され\footnote{「特異」な相続く点については Brill--Noether 第6節、とくに no. 10, 11 を参照。通常の相続く点とは、例えば接線と曲線との交点、またはより高次の接触点（分岐点）である。}、Weierstrass の元概念に対応する。その後は代数的・有理的な方法、本質的には三元消去論の方法によって進み、これもまた明示的表示にまで到達する。

Dedekind--Weber の本来の純粋に算術的な理論\footnote{\emph{J. f. M.} 92（D.--W. として引用する）。}は、その諸定理がふたたび存在証明の性格をもつという点で Riemann の理論に近い。この理論は方法の完全な純粋性を保つ。\footnote{この理由から、例えば Hensel--Landsberg の理論よりも、他の理論との比較に適している。後者は以下の比較では時折現れるだけである。--- Abel 積分とその反転の理論は、連続性を前提とする限り、D.--W. にも代数的理論のさらなる展開（M. Noether, \emph{Ann.} 37）にも扱われていない。} 変数が不定元の役割だけを果たすこの理論の算術的基礎づけは、2 で詳しく説明するように、代数的数の理論と平行して進む。\footnote{Hensel--Landsberg の理論と数論、とくに \(p\)-進数の理論との比較は Hensel に見られる。} この基礎づけは、点概念を算術的に導入する手段を与え（Hensel no. 10a を参照）、したがって連続性に関するいかなる考察も用いずに絶対 Riemann 面の概念を設定する手段を与える。理論はそれによって形式的な性格を得る。例えば微分の概念も純粋に形式的に導入される。点概念によって得られた基礎の上で、代数幾何学的理論の方法と概念との近い親族性を示すことができる。というのも、後者も元概念を除けば超越的補助手段から自由だからである。これは次の節で行う。

\subsection*{2. 代数的数と代数関数との平行性。（共役元、基底定理、加群概念）}
代数的数体は、有理数体 \(R\) 上の \(n\) 次の体として定義される。一変数の代数関数体は、任意の複素係数をもつ一つの不定元 \(z\) のすべての有理関数からなる体 \(K(z)\) 上の \(n\) 次の体として定義される。したがって両方の体について、基礎体の特殊な性質を完全に捨象するすべての定理が共通に成り立つ。すなわち、基底、ノルム、トレース、判別式に関する定理である。\footnote{D.--W., § 1 および 2；Hensel no. 4（注 5 付き）。} Dedekind--Weber はこれらすべての概念を、共役元を用いずに導入する。しかし、Steinitz が Kronecker に従って示したように、共役体の概念は完全に形式的な仕方で定義できることを注意しておくべきである。\footnote{Steinitz, Algebraische Theorie der Körper.（\emph{J. f. M.} 137, § 6 および 8.）} すなわち、\(K\) を任意の体、\(\varphi(x)\) を \(K\) 上既約な \(x\) の多項式とすると、\(\varphi(x)\) を法とする剰余類の系は体をなす。\(x\) の剰余類に元 \(\xi\) を対応させると、同型対応のもとで、多項式 \(f(x)\) によって表される剰余類に元 \(f(\xi)\) が対応する。\(\varphi(x)\) で割り切れるすべての多項式は零類を表すので、\(\varphi(\xi)\) は零類に対応する。したがって元 \(\xi\) は記号的に \(\varphi(x)=0\) の零点と見なすことができ、代数体 \(K(\xi)\) を生成する。この手続きを反復すると、\(n\) 個の共役体 \(K(\xi),K(\xi'),\ldots,K(\xi^{(n-1)})\) に至る。したがって、とくに代数関数の場合にも、純粋に算術的な理論を離れることなく共役元について語ることができる。

二つの基礎体 \(R\) と \(K(z)\) は、その中ですべての整な量の系が定義されているという性質を共有する。それはそれぞれ、有理整数および一つの不定元の整有理関数である。これらの整な量は、任意の二つ \(a\) と \(b\) が最大公約元 \(d\) をもち、それが \(ma+nb\) の形に表されるという特徴的性質をもつ。ここで \(d,m,n\) もまた \(R\) または \(K(z)\) の整な量である。しかも \(d\) は、この形に表される絶対値最小の数、または次数最小の多項式である。有理整数および一変数関数の素因数への一意分解を導くこの性質だけに、代数的体または上体の整な量に関する諸定理は基づいている。最大公約元の存在を示す考察は、代数体の整な量の基本系または基底の存在も与える。整な量からなる体の基底だけを考えると、その判別式は有理整数または一つの不定元の整有理関数になる。絶対値最小の数、または次数最小の関数に対応する各基底は、基本系をなす。

基底の問題について、より一般の体も含めてより精密な見通しを与えるのが加群概念の導入である。この概念は、文献で元の性質に応じて異なる形で現れる定義を含むように定式化する。基本整域 \(\frS\) に関する加群とは、二つの元の差がふたたびその系に属し、また一つの元と \(\frS\) の任意の量との積もその系に属するような元の系である。（\(\frS\) が有理整数からなる場合、第二の条件は第一の条件から従う。）\footnote{加群概念は Dedekind に由来する。彼は数論の第二版で初めて数加群（\(\frS=\) 有理整数）を導入した。整係数線形形式の加群もそこに暗黙に含まれている。Dedekind--Weber には「関数加群」が現れる。その元は代数関数であり、\(\frS\) は一つの不定元のすべての整有理関数からなる。--- \(n\) 個の不定元の斉次形式、より一般に多項式の加群は、Hilbert（\emph{Annalen} 36）に初めて現れる。\(\frS\) はすべての多項式からなり、元も多項式であるから、この加群はわれわれのイデアル定義に入る。Kronecker の多項式の「加群系」は、抽象的定義からではなく基底から出発する。}

有限個の元からなる基底をもち、その基底によって各元が \(\frS\) の係数を用いて線形に表される加群を有限（有限生成）という。とくに一項加群は \(c\cdot\alpha\) の形である。ここで \(\alpha\) は基底元を表し、\(c\) は \(\frS\) のすべての量を走る。\emph{元そのものが \(\frS\) に属する加群を \(\frS\) におけるイデアルという。} ここから有限イデアルの概念、とくに一項イデアル（主イデアル）の概念が従う。

整な量とイデアルについてのすべての基底定理は、次の定理に基づいている。

\emph{\(M\) を \(\frS\) に関する加群とし、その元が \(\frS\) の係数をもつ \(n\) 個の不定元の線形形式からなるとする。もし \(\frS\) のすべてのイデアルが有限ならば、\(M\) も有限加群である。とくに \(\frS\) のすべてのイデアルが主イデアルならば、\(M\) は線形独立な形式からなる基底をもつ。}\footnote{この定理は、異なる係数領域 \(\frS\) について文献に別々に現れる。有理整数については、例えば Hilbert, Zahlbericht § 3 および 4 を参照。複数の不定元の多項式については、Hilbert, Über die vollen Invariantensysteme, § 2 の末尾（\emph{Math. Ann.} 42）を参照。}

実際、\(M\) の元が \(l(x)=a_1x_1+\cdots+a_nx_n\) で与えられているなら、係数 \(a_1\) 全体は \(\frS\) の一つのイデアルを走る。仮定によりそれは \(b_1a^{(1)}+\cdots+b_\rho a^{(\rho)}\) の形である。したがって \(M\) に属する線形形式
\[
m(x)=l(x)-b_1l^{(1)}(x)-\cdots-b_\rho l^{(\rho)}(x)
\]
は \((n-1)\) 個の不定元だけに依存する。ここから有限回反復することで主張が従う。各場合に \(\rho=1\) ならば、基底はそれぞれ \(n,n-1,\ldots,1\) 個の不定元の \(n\) 個の形式からなり、これによって零でない形式の線形独立性が得られる。

基本系についての定理は、これからさまざまな場合に次のように従う。

代数体が、整な代数元 \(\delta\) を基礎体に添加することによって得られ、\(\frS\) が基礎体のすべての整な量を表すとする。このとき上体の各整な量は次の表示をもつ（例えば Zahlbericht § 3 を参照）：
\[
\omega=\frac{a_1\delta^{\,n-1}+a_2\delta^{\,n-2}+\cdots+a_n}{\Delta(\delta)},
\]
ここで \(a_i\) は \(\frS\) の量であり、\(\Delta(\delta)\) は \(\delta\) の判別式を表す。置換
\[
x_i=\frac{\delta^{\,n-i}}{\Delta(\delta)}
\]
によって、すべての整な量の加群は線形形式の加群に移される。したがって、\(\frS\) のすべてのイデアルが有限であるすべての場合に、基本系の存在が従う。同じ考察は、より一般に、上体の整な量の領域におけるすべてのイデアルが有限であることを示す。さて、とくに有理整数および一変数関数については、すべてのイデアルが主イデアルである。したがって、代数的数体および一変数代数関数体における基本系は、\(n\) が次数を表すとき、ちょうど \(n\) 個の量からなることが従う。また、これらの体では上の議論によりすべてのイデアルが有限であるから、同時に相対体の基本系も得られる。それは一般に相対次数が示す数よりも多くの量からなる。\footnote{線形形式の加群に関する Steinitz の研究によれば、\(\frS\) が（有限）代数的数体の整数からなるとき、\(\nu\) を相対次数とすれば \(\nu+1\) 個の基底元で足りる（\emph{Math. Ann.} 71, § 4）。} さらに、\(\frS\) が複数の不定元のすべての多項式の領域を表すならば、\(\frS\) の各イデアルも有限である。\footnote{Hilbert の加群基底定理（\emph{Math. Ann.} 36）による。ここで一貫性のためイデアルと呼んだ多項式の加群は、通常「形式加群」または単に「加群」と呼ばれる。こうした加群 \(N\) が有限であるという Hilbert の定理も、線形形式に関する定理に帰着できる。すなわち、\(x_1,\ldots,x_r,x_{r+1}\) の \(r+1\) 個の不定元において、\(N\) に含まれ \(x_{r+1}^{\,n}\) の項を含む次数 \(n\) の多項式 \(f\) をとる（これは線形変換によって常に達成できる）。すると \(N\) のすべての多項式は \(f\) を法として \(a_1(x)x_{r+1}^{\,n-1}+\cdots+a_n(x)\) の形に表されるものに合同である。したがって \(x_{r+1}^{\,n-i}=X_i\) とおけば、それらは線形形式の加群をなす。この加群については、\(r\) 個の不定元の領域のすべてのイデアルが有限であると仮定してよい。これは一つの不定元について成り立つからである。} したがって、多変数の代数関数体についても基本系の存在が得られ、それはまた次数が示す数よりも多くの量からなる。

\subsection*{3. 代数的数と代数関数のイデアル論における平行性と相違}
すべてのイデアルが素イデアルの冪の積として一意的に表されるというイデアル論の基本定理の証明は、代数的数と代数関数について共通に実行できる。それは、有理整数および一変数関数についてすべてのイデアルが主イデアルであるという性質だけに基づく。\footnote{Hurwitz: Über die Theorie der Ideale（\emph{Gött. Nachr.} 1894）の序論の注を参照。Weber の代数学教科書では、Kronecker に従って、汎関数の導入により両方の体について平行に証明が行われている。［Kronecker の理論とイデアル論との関係は「拡張された Gauss の定理」によって媒介される；Hurwitz, loc. cit.］} 証明の主要点は、イデアル \(\frc\) がイデアル \(\fra\) で割り切れること（すなわち \(\frc\) の各元が \(\fra\) に含まれること）から、積表示 \(\frc=\fra\frb\) が従うことを示す点にある。\footnote{Dedekind（数論）でとくに強調されている。補遺 XI, 定理 VII, § 177；p. 554 の注を参照。この定理は代数体を扱っているという事実を用いる。多項式のイデアル（加群）についてはもはや成り立たず、そこでは積の代わりに最小公倍元が現れる。} Hurwitz ではこの証明は「拡張された Gauss の定理」に基づく。これは、二つの多項式 \(\varphi\) と \(\psi\) の積の各係数を割る整な量 \(\omega\) が、\(\varphi\) と \(\psi\) のすべての係数の積も割るというものである。Dedekind では、この定理と同値な加群定理に基づく。\footnote{数論、補遺 XI；定理 VI, § 173。Dedekind は二つの定理の同値性を指摘している（Über die Begründung der Idealtheorie, \emph{Göttinger Nachr.} 1895）。}

さらに進むと、補加群の概念（Hensel no. 12）、およびそこから分岐イデアル（基礎イデアル、差積）\(\frD\) を共通に定義できる（そのノルムが体判別式 \(D\) となる）。そして、分岐イデアル \(\frD\) は \emph{すべての主イデアル \(F'(\delta)\) の最大公約元} であるという基本定理が成り立つ。ここで \(\delta\) は体のすべての整な量を走り、\(F(\delta)=0\) はその都度対応する既約方程式を表す。これから、\emph{体判別式は、素イデアルの平方で割り切れる有理素数、または \(z\) の一次因子をすべて、かつそれだけを含む} ことが従う。\(F'(\delta)\) 自身については次の表示を得る：
\[
F'(\delta)=\frf\cdot\frD,
\]
さらにノルムをとると
\[
\Delta(\delta)=R^2\cdot D;
\]
ここで \(\frf\) は \(\delta\) から導かれる環（整環、整域）の導手を表す。すなわち \(\frf\) は、任意の整な量を掛けた後に加群 \([1,\delta,\ldots,\delta^{n-1}]\) で割り切れる（整な）量全体からなる。\(R^2=\operatorname{Norm}\frf\) は、基底 \(1,\delta,\ldots,\delta^{n-1}\) を整な量の基底に移す置換行列式の平方を表す。\footnote{Dedekind, Über die Diskriminante endlicher Körper（\emph{Abhandl. der Gött. Ges. d. Wissensch.} vol. 29, 1882）。そこで代数的数について与えられた定義は、代数関数に直接移すことができる。証明もほとんどそのまま移すことができる。D.--W. における証明の順序は異なる。下を見よ。引用した定理については Zahlbericht, § 31 および 32 も参照。（そこで \(F'(\delta)\) は代数的整数 \(\delta\) の「差積」と呼ばれる。）}

さらなる展開における\emph{相違}は、基礎体の特殊な性質によって条件づけられる。例えば代数的数の場合、\emph{基礎体の元が同時に指数でもある}という事実により、分岐イデアルが、\(p\) の \(e\) 乗に現れる素イデアル \(\frp\) の \((e-1)\) 乗より高い冪で割り切れる可能性が生じる（すなわち \(e\) が \(p\) で割り切れる場合である）。これはさらに慣性体と分岐体の概念\footnote{Zahlbericht, § 39--45 を参照。}へ導くが、これらは代数関数の理論には類似物をもたない。しかし何よりも、基礎体に\emph{無限に多くの定数が存在する}ため、代数関数ではイデアル類数の有限性が消え、それとともに類数の決定に関わるすべての問題も消える。それにもかかわらず、ある意味で、すべての（超越的な）素形式の領域\footnote{Klein, Zur Theorie der Abelschen Funktionen（\emph{Ann.} 36, 1890）。}を類体の超越的類似物と見なすことができる。ここでは各点が一つの形式、すなわち主イデアルによって生成されるからである。

他方、基礎体 \(K(z)\) に無限に多くの定数が存在することは、代数関数について、代数的数には類似物をもたない定理と問題を導く。第一に、このことから証明の配置にある簡略化が生じる。例えば D.--W. による、イデアルの素イデアルへの一意分解の証明では、任意の一次形式 \((z-c)\) が無限に多くの剰余類、すなわちすべての定数をもつという事実を用いることによって、「拡張された Gauss の定理」またはそれと同値な定理を避けることができる（これに対し、素数を法とする剰余類の数は有限である）。\footnote{D.--W. § 9。pp. 212, 213 の注を参照。}

さらに、\(z\) の整有理関数を係数にもつ任意の多項式が \((z-\alpha)\) を法として一次因子に分解するという事実（これは整係数多項式を \(p\) を法として考える場合には成り立たない）が、本当に新しい結果を導く。この事実は、すべての定数の代わりにすべての代数的数だけをとってもなお成り立つ。\footnote{D.--W. は序論で、この場合にも全理論がそのまま保たれることを述べている。} すなわち、ここからすべての素イデアルが一次のイデアルであること\footnote{D.--W., § 9, no. 7。}、さらに、体判別式 \(D\) が、\(\delta\) がすべての整な量を走るときのすべての判別式 \(\Delta(\delta)\) の交わりになること\footnote{D.--W., p. 224。}が従う。この二点はいずれも代数的数の場合と対照的である。続いてここから、分岐イデアルがすべての主イデアル \(F'(\delta)\) の交わりであることもふたたび従う。

さらに別の相違は、代数的数では有理数体 \(R\) が絶対的に特別なもの、すなわちその体に含まれる素体（単位から導かれる体）であるのに対し、基礎体 \(K(z)\) は相対的なものであるという事実によって条件づけられる。任意の非定数関数 \(\xi\) を独立変数として選ぶことができ、そのとき体は \(K(\xi)\) 上の代数体となる。各素イデアル \(\frp\) が一次のイデアルであり、したがって各関数が \(\frp\) を法として定数に合同であるという定理を、任意の関数を独立変数にとれるという可能性と合わせると、代数的数を超えて進む最も重要な概念、すなわち純粋に算術的な形での点概念に導かれる（Hensel 10a を参照）。「ここでは各点が、関数体に同型な定数体によって表される。」\footnote{よく知られているように、二つの体が「同型」であるとは、和、差、積、商も対応するように一対一対応に置けることをいう。} こう定義された点全体が「絶対 Riemann 面」をなす。この点の定義は体だけに依存し、特殊な変数には依存しない。しかし存在証明は、ある関数を独立変数 \(z\) として区別することによって行われる。点 \(P\) は、各関数に \(\frp\) を法とする定数剰余を割り当てることによって、素イデアル \(\frp\) により生成される。したがって、とくに \(\frp\) で割り切れる関数は \(P\) で消え、逆も成り立つ。\(\frp\) が \(z\) におけるすべての素イデアルを走ると、これにより \(z\) が有限であるすべての点が得られる。残りの点は独立変数 \(\xi=1/z\) を導入することによって生じ、主イデアル \(\xi\) を割る素イデアルに対応する。もちろん、点概念に付随する概念、例えば多角形、多角形類、多角形類の次元なども、代数的数には類似物をもちえない。

独立変数を自由に選べることは、分解 \(F'(\vartheta)=\frf\cdot\frD\) にさらに別の解釈を導くことも注意しておくべきである。すなわち、\(\vartheta\) を独立変数と見なすと、対応して
\[
\left(\frac{\p F}{\p z}\right)=\frf\cdot\fra
\]
を得る。なぜなら \(\vartheta\) から導かれる環は \(z\) と \(\vartheta\) に関して対称的に定義されていたからである。そして \(\frf\) は二つの主イデアル \(\frac{\p F}{\p z}\) と \(\frac{\p F}{\p\vartheta}\) の最大公約元となる。導手 \(\frf\) は同時に、\(\vartheta,z\) に関する「二重点イデアル」となる。\footnote{D.--W. p. 263。}

\subsection*{4. 級数展開の算術的基礎づけとイデアル論との関係}
Hensel--Landsberg では、級数展開という超越的補助手段と関数論から借りた点概念が、イデアル論の位置を占める。Hensel が最近与えた代数的数と代数関数のための級数展開の算術的基礎づけ\footnote{Eine neue Theorie der algebraischen Zahlen（\emph{Math. Zeitschr.} vol. 2, 1918）および Neue Begründung der arithmetischen Theorie der algebraischen Funktionen einer Variablen（\emph{Math. Zeitschr.} vol. 4, 1919）。より詳しい説明については Hensel no. 44ff. を参照。}は、したがってイデアル論の等価物と見なすことができる。実際、この算術的基礎づけは、各素イデアルについて、この素イデアルが主イデアルになるような（超越的）拡大体に移ることにほかならない。そうすれば通常の分解定理が成り立つ。そして、この考察は一つの素数 \(p\)、または一つの一次関数 \(z-a\) に現れるすべての素イデアルについて行えば十分である。

まず、\(z-a\)、または \(p\)（\(p\)-進数）の整冪級数展開全体からなる超越拡大体を導入する。ただし、それぞれ負冪は高々有限個だけ許す。この体において、代数的数体または関数体を定義する既約方程式は \(\rho\) 個の既約因子の積に分解する。これは \(p\)、または \(z-a\) が \(\rho\) 個の「一次的」イデアル \(\mathfrak q_1,\ldots,\mathfrak q_\rho\) の積に分解することに対応する。各一次的イデアルを素イデアルの冪としてさらに表示すること、
\[
\mathfrak q_i=\frp_i^{\,\ell_i},
\]
は、Hensel では超越拡大体に対して代数的なさらに一つの体を導入することに対応する。そして代数関数の場合、この体は純方程式によって定義できる。代数的数の場合には、より一般に、代数的に解ける方程式によって定義される。代数関数の場合の簡略化は、ここでは慣性体と分岐体（no. 3 参照）が現れないことによる。

これらの研究に現れる素元 \(\pi\) に関する「有限で終わる級数」\footnote{Neue Begründung \(\ldots\), 公式 (11)。}（これはまだ超越的な収束証明を必要としない）
\[
v=c_0+c_1\pi+\ldots+c_{\sigma-1}\pi^{\sigma-1}-v_\sigma\pi^\sigma,
\]
は、\(c_0,\ldots,c_{\sigma-1}\) が定数で \(v_\sigma\) が体の整元を表すとき、D.--W. にまったく対応する形で見いだされる。\footnote{D.--W., p. 231 および § 15。} そこではこれはイデアル論から導かれ、一つの点における関数の零の位数（または無限大になる位数）を定義するための基礎を与える。

\subsection*{5. 代数関数の算術的理論と幾何学的理論との比較。異なる不変性概念}
「不変」とは、体だけによって特徴づけられる任意の概念を意味する。算術的理論では、この不変性は一般にすでに定義の中に与えられている。その定義は、ある基底や 区別された変数に関してではなく、体のすべての元に関して述べられるからである。上で与えた点概念がその例である。これに対して残りの理論は、ある特殊な変数または基底から出発し、その後で概念がその特殊な選択に依存しないことを示す。ここでは不変性は双有理変換のもとでの不変性となる。

すなわち、体が一方では代数元 \(s\) を \(K(z)\) に添加することによって、他方では \(\sigma\) を \(K(\xi)\) に添加することによって生じるとする。ここで \(f(z,s)=0\)、また \(F(\xi,\sigma)=0\) は、それぞれ \(s\)、また \(\sigma\) の既約方程式を表す。あるいは、さらに一般に、\(K(\eta)\) に有限個の元を添加し、対応する既約方程式
\[
g(\eta,\tau_1)=0;\qquad g_2(\eta\tau_1\tau_2)=0\ldots .
\]
をもつとする。すると定義により、すべての元は \(z\) と \(s\) によって、\(\xi\) と \(\sigma\) によって、また \(\eta,\tau_1,\tau_2,\ldots\) によって、有理的に表される。とくに \(z\) と \(s\) は \(\xi\) と \(\sigma\) によって有理的に表され、その逆も成り立つ。同様に \(z\) と \(s\) は \(\eta,\tau_1,\tau_2,\ldots\) によって有理的に表され、その逆も成り立つ。したがって平面曲線 \(f(z,s)=0\) は「双有理変換」によって平面曲線 \(F(\xi,\sigma)=0\) に移り、あるいは \(g_1(\eta,\tau_1)=0; g_2(\eta,\tau_1,\tau_2)=0\ldots\) によって定義される変数 \(\eta,\tau_1,\tau_2,\ldots\) の空間内の曲線に移る。したがって不変性は、一つの特殊曲線について述べられた定義が、そこから双有理変換によって得られるすべての曲線（任意の次元数の空間内の曲線）について、すなわち Riemann の意味での「類」のすべての曲線について有効なままである、という事実の中に現れる。ここで「点」としては、\(f(a,b)=0\) となるような整合的な値系 \(z=a,s=b\) が定義される。双有理変換のもとで、この点は一般にふたたび \(F(\xi,\sigma)=0\) の整合的な値系 \(\alpha,\beta\) に移る。しかし特殊な場合、\(a,b\) が「特異」点であったときには、\(F(\xi,\sigma)=0\) 上の複数の「点」に対応することがある。そして \(f(z,s)=0\) の「特異」点が、分離して、または相続いて位置する有限個の単純点に対応するような曲線を常に指定できる（特異点の解消）。\footnote{例えば Brill--Noether 第6節を参照。「特異」点についてのより詳しい情報は次節に現れる。} 同様に、点群もまた点群に移る。そして適当な変換で対応する単純点の数だけ特異点を数えれば、その点数は不変である。任意の曲線は「通常多重点」をもつ曲線に双有理的に変換できるので、\footnotemark[\value{footnote}] 幾何学的理論はその概念をそのような特殊曲線についてだけ展開し、その後で、最も一般の特異点に至るものも含め、任意の双有理変換のもとでの不変性を証明する。

とくに、超楕円的な場合を除いて、その「類」の 特別な曲線として、\((p-1)\) 次元空間内のいわゆる「\(\varphi\) の正規曲線」を導入できる。これは特異点をもたない。ここで \(\varphi\) は、斉次化された \(f(z,s)=0\) に対する \(p\) 個の「\((n-3)\) 次の随伴曲線」である。\(f(z,s)=0\) から \(F(\xi,\sigma)=0\) への双有理変換には、\(\varphi\) の線形変換が対応する。したがってこれは射影的意味での不変性である。体のすべての関数を \(\varphi\) の有理関数として表す表示は、幾何学的理論では「不変表示」と呼ばれる。\footnote{例えば Brill--Noether 第8節を参照。多変数代数関数について、この不変性を線形変換のもとでのそのような不変性に還元できるかという問題は扱われていない。}

したがって \(\varphi\) の線形系、またそれらによって切り出される「特殊群」は、各双有理変換のもとで自分自身に移り、したがって体だけによって特徴づけられる。それは算術的理論における微分類に対応する。そこでは不変性が、体の各関数を独立変数とする性質によって与えられる定義の中に、ふたたび直接に示されている。

\subsection*{6. 比較の続き。特異点}
点概念の異なる把握と結びついて、代数幾何学的理論で特別な困難を生む曲線の「特異点」は、算術的理論の基礎づけにはまったく現れない。またその結果、幾何学的理論の場合とは違い、「随伴関数」は算術的理論を組み立てるためには必要ではなく、理論の特殊でより高次の部分に属する。一方、幾何学的理論の基礎づけには分岐点が現れない。実際、Aronhold の標準形による積分の表示（Hensel no. 27, 公式 80）でも、分岐点は消えてしまう。

5 で与えた特異点の定義は、次のように述べられる。\(f(z,s)=0\)（または高次元空間内の曲線）が \(z=a\), \(s=b\)（または \(z=a\), \(s_i=b_i\)）に特異点をもつとは、この値系が絶対 Riemann 面の複数の点（算術的意味での点）によってとられるか、または一つもしくは複数の点でより高い位数でとられることをいう。後者の場合、それは \(z\) と \(s\) の双方の分岐多角形の点である。これは幾何学的理論の「相続く点」に対応し、その点は \(f(z,s)=0\)（または高次元空間内の曲線）の（高次の）尖点となる。

体の任意の関数 \(\eta\) は、仮定により \(z\) と \(s\) によって有理的に表される。したがって、同型性により、\(z=a\), \(s=b\) となる絶対 Riemann 面の各点は、\(z\) と \(s\) によって表された各関数 \(\eta\) に、その \(z=a\), \(s=b\) での値を対応させることによって生じる。その点が特異であるためには、少なくとも一つの関数 \(\eta\) が、\(z=a\), \(s=b\) で複数の値をとるか、一つの値系を複数回とらなければならない。そして \(z\) と \(s\) によって有理的に表されるということから、これはその表示において \(z=a\), \(s=b\) で分子と分母が同時に消える場合にだけ起こる。したがって上の定義は、\(f\), \(\frac{\partial f}{\partial z}\), \(\frac{\partial f}{\partial s}\) が \(z=a\), \(s=b\) で消えるという通常の定義と同一である。実際、この場合 \(\eta=\frac{z-a}{s-b}\) は指定された種類の関数であり、\(z=a\), \(s=b\) で多価である。反対の場合には、たとえ \(\frac{0}{0}\) になっても、各関数は一つの値だけをとる。

算術的点概念では特異点は起こりえない。一つの関数にでも異なる値系が対応すれば、二つの点は異なるものと見なされるからである。しかし、算術的点概念の基礎づけに用いられるイデアル論においても、特異点の概念は現れない。これは、体のすべての整な量の全体を常に同時に考えるからである。\(z\) に関して代数的に整である \(\delta\) に対し、\(F(z,\delta)=0\) の有限部分にあるすべての特異点は、（第二の定義と no. 3 の末尾に従えば）「二重点イデアル」\(\frf\) によって生成される。さて、基本定理により分岐イデアル \(\frD\) はすべての主イデアル \(F'(\delta)\) の最大公約元であるから、任意の素イデアル \(\frp\) に対して、対応する \(\frf\) が \(\frp\) で割り切れないような \(\delta\) を指定できる。したがって、\(\frf\) は \(F'(\delta)\) と \(F'(z)\) の最大公約元であるので、少なくとも二つの主イデアル \(F'(\delta)\), \(F'(z)\) の一方は \(\frp\) で割り切れない。したがって、\emph{すべての \(F'(\delta)\) と \(F'(z)\) が同時に消える値系は存在しない}。その結果、値系 \(z=a,\delta_i=b_i\) は絶対 Riemann 面のただ一つの点でだけとられる。体の\emph{整代数関数の全体}（これは無限次元空間内の曲線と解釈してよい）は\emph{有限部分に特異点をもたない}。イデアル論で考慮されるのは有限値だけである。

しかし、すべての整な量の基底関数 \(\omega_1,\ldots,\omega_n\) も有限部分に特異点をもたない。上のことにより、算術的意味での各点は、すべての整関数に同型に対応させられた定数系によってすでに一意的に決定される。したがって、任意の整関数 \(s_i\) によって定義される「空間曲線」が \(s_i=\alpha_i\) において有限部分に特異点をもつならば、\(s_i=\alpha_i\) で複数の値系をとる整関数 \(\eta\) が少なくとも一つ存在しなければならない。いま \(\omega_1,\ldots,\omega_n\) がすべての整な量の基底を表すなら、整数的有理係数 \(a_i\) による \(\delta=a_1\omega_1+\cdots+a_n\omega_n\) のため、\(\omega\) の各値系には \(\delta\) の一つの値系だけが対応する。したがって \(\omega_1,\ldots,\omega_n\) によって定義される空間曲線は有限部分に特異点をもちえない。さらに、\emph{\(z\) が有限である絶対 Riemann 面の各点は、\(\omega\) に同型に対応させられた定数系によってすでに一意的に定義される。基底関数とはまさに、\(s_i=\alpha_i\) で多価になる関数 \(\eta\) である。} したがって、例えば基底 \(1,\delta,\ldots,\delta^{n-1}\) の代わりに整な量の基底を導入することによって、特異点の困難は回避される。

幾何学的理論では、絶対 Riemann 面のすべての点を一意的に生成する課題、より一般には与えられた点で無限大になるすべての関数を形成する課題は、随伴形式とその商によって遂行される。形式 \(\psi\) が随伴であるとは、基礎曲線 \(f(z,s)=0\) の各 \(i\) 重点で \(\psi=0\) が少なくとも \((i-1)\) 重点をもつことをいう。高次の特異点は、まず通常多重点に解消されなければならない。最も一般の関数 \(\eta=\frac{\psi}{\chi}\) について、分子と分母が \(f(z,s)=0\) のすべての特異点で消えなければならないことは上で示した。それらが随伴でなければならないことは、いたるところ有限な微分を考察すれば分かる。逆に、随伴性の条件が最も一般の関数を形成するために十分であることは、後で詳しく論じる「剰余定理」が示す。したがって随伴関数は、算術的理論の基底関数の位置を占める。

算術的には、特異点をすべて有限部分にあるものと仮定すれば（これは変換によって常に達成できる）、随伴形式は「二重点イデアル」\(\frf\) で割り切れる体の関数の分子に対応する。ここから、基底量 \(\omega_1,\ldots,\omega_n\) を随伴形式の商として表すこと、およびさまざまな表示方法の関係が、形式的に次のように示される。

\(R(z)\) を、基底 \(1,\delta,\ldots,\delta^{n-1}\) から基底 \(\omega_1,\ldots,\omega_n\) への置換行列式とする。すると逆に、各 \(\omega\) は \(z\) と \(\delta\) の整関数を \(R(z)\) で割った商として表される：
\[
\omega_i=\frac{G_i(z,\delta)}{R(z)}.
\]
\(R(z)\omega_i\) は \(\delta\) から導かれる環に属するから、導手の定義により \(R(z)\) は \(\frf\) で割り切れる［\((R(z))^2=\operatorname{Norm}\frf\) からは、この可除性は \((R(z))^2\) についてだけ従う］。いま \(\omega_i\) は整関数としてすべての有限な \(z\) で有限にとどまるので、分子関数 \(G_i(z,\delta)\) も \(\frf\) で割り切れなければならない。したがって、すべての \(\omega_i\) の分子と分母は随伴である。しかしこの \(\omega\) の表示から、体のすべての関数が随伴形式の商によって表示されることが従う。

以上により、幾何学的理論は、すべての整な量を同時に考える算術的理論と対比して、\emph{一つの量 \(\delta\) から導かれる環（整環）の理論}として特徴づけられる。したがって、その環の導手、すなわち特異点が決定的な役割を果たすことは明らかである。環（整環）の算術的理論と幾何学的理論とのあいだには、さらに別の類似も与えられる。例えば幾何学的理論は、随伴曲線との交わりにおいて特異点へ落ちる交点をどこにも数えず、それとは異なる点群だけを考える。これに対応して Dedekind\footnote{Über die Anzahl der Idealklassen in den verschiedenen Ordnungen eines endlichen Körpers.（Festschrift Braunschweig 1877）；§ 3--5。（これらの節で代数的数について述べられた定理は、一変数代数関数について文字通り成り立つ。）} は、導手イデアル \(\frf\) と互いに素な環イデアルだけを考え、このイデアル領域に対して、素イデアルへの一意分解を含むすべての可除性法則を立てている。

また、幾何学的理論の基礎にある「代数関数の基本定理」\footnote{M. Noether, Über einen Satz aus der Theorie der algebraischen Funktionen. \emph{Math. Ann.} vol. 6 (1873).}、あるいは少なくともその直接の帰結は、ある意味で環の算術的理論に類似物をもつ。この定理は、表示可能性
\[
\psi(z,s)=A(z,s)f(z,s)+B(z,s)\varphi(z,s)
\]
の条件を与える。ここで現れるすべての量は \(s,z\) に関して整有理である（より一般には \(x_1,x_2,x_3\) に関して整かつ斉次である）。これらの条件から、\footnote{Brill--Noether, no. 55.} \(f(z,s)=0\) によって、商 \(\frac{\psi}{\varphi}\) が \(f\) に随伴であるとき、常に整有理関数 \(B(z,s)\) に等しいことが従う。これに対応するのが、no. 3 の定理、すなわちこれらすべての比較の基礎にある定理である。一つの \(\delta\) から導かれる環の導手 \(\frf\) は、\(f(z,\delta)=0\) の二重点イデアルに等しい。なぜなら、導手の定義により、\(\frf\) で割り切れる各代数的整量はその環に属し、したがって \(z\) と \(\delta\) によって整かつ有理的に表示できるからである。一方、引用した定理は、そのような量が随伴関数であると言っている。\footnote{Dedekind に従って、ここでは一貫して \(\delta\) が代数的整数であると仮定している（これは変換によって常に達成できる）。幾何により正確に対応する場合、すなわち \(s\) が \(z\) に代数的に分数的に依存してよい \(f(z,s)=0\) は、Hensel--Landsberg が扱っている。Hensel no. 26 および 29 を参照。}

\subsection*{7. 剰余定理とイデアル論の諸定理との比較}
今述べた「代数関数の基本定理」の帰結は、幾何学的理論において直接に「剰余定理」へ導き、したがって幾何学的理論全体の基礎へ導く。しかし、ある点では、剰余定理は他の基礎よりも再び初等的な性格をもつ。というのも、それはイデアルの素イデアルへの一意分解、より正確にはこの分解定理の基礎にある事実からの直接の帰結として算術的に述べることができ、したがって高次の環の理論には依存しないからである。この算術的定式化へは次のように導かれる。

基礎曲線 \(f(z,\delta)=0\) が、変数の一次分数変換によって常に達成できるように、特異点および考慮している有限個の点群を有限部分にもつものと仮定する。さらに \(\delta\) は \(z\) に代数的に整に依存していると仮定する。これは変数に整な後続の線形変換によって常に達成できる。このとき、\(z=a,\delta=b\) となる絶対 Riemann 面の点 \(P\) を生成する素イデアル \(\frp\) に対して、\(f(z,\delta)=0\) の点 \(z=a,\delta=b\) が対応する。より一般に、イデアル \(\fra=\frp_1^{\,e_1}\cdots\frp_\rho^{\,e_\rho}\) は、対応する重複度をもつ点 \(z=a_i,\delta=b_i\) からなる点群を生成する。そして \(f(z,\delta)=0\) の有限部分にあるすべての点はこのように生成される。点 \(P\) では \(\frp\) で割り切れるすべての関数 \(g_1,g_2,\ldots\) が消え、逆にこの関数 \(g\) の全体がイデアル \(\frp\) を生成するから、\(f\) 上の対応する点は
\[
g_1(z,\delta)=0,\quad g_2(z,\delta)=0,\ldots
\]
と \(f=0\) との交わりである。イデアル \(\fra\) に対応する点群についても同じである。実際、各イデアルが二つの主イデアルの最大公約元であるという定理により、二つのそのような曲線で常に一つの点群を切り出すのに十分である。とくに主イデアルに対応する点群は曲線 \(g(z,\delta)=0\) によって切り出され、逆も成り立つ。主イデアルは完全交差に対応する。

二つの点群 \(A\) と \(B\) は、幾何学的理論において、同じ点群 \(R\) を加えることにより完全交差に完成できるとき、\emph{余剰対応} であると呼ばれる。これらの点群がそれぞれイデアル \(\fra,\frb,\frr\) によって生成されるなら、これらのイデアルのあいだには関係
\[
\fra\cdot\frr=(\alpha),\qquad \frb\cdot\frr=(\beta),
\]
があり、これはイデアル \(\fra\) と \(\frb\) が同値であることを示す。逆に、二つのイデアルの同値性 \(\fra=\eta\frb\) からは、各イデアルが別のイデアルを掛けることによって主イデアルに変換できるという定理に基づいて、常にこの関係が従う。\footnote{Dedekind, 数論 § 181 を参照。すなわち、\(\fra\frr\) が主イデアル \((\alpha)\) となるように \(\frr\) を選ぶと、\(\frb\frr=\eta^{-1}\cdot(\alpha)\) となる。そして \(\frb\frr\) は整な量だけを含むから、\(\eta^{-1}\cdot(\alpha)\) についても同じことが成り立ち、したがってこれも主イデアル \((\beta)\) である。} \emph{したがって、同値なイデアルと余剰対応する点群とは同一の概念である。}

剰余定理はいま次のように言う。\emph{\(A\) と \(B\) が余剰対応であるなら、これらは同じ剰余をもつ随伴曲線（特異点の外で）によって切り出される。} 関数の零点と極は常に余剰対応であるから、これは同時に、体の最も一般の関数が随伴形式の商として表されることを証明する。イデアル論の言葉では、剰余定理は単に、\(\fra\) と \(\frb\) のほかに \(\fra\frf\) と \(\frb\frf\) も同値であるということである。ここで \(\frf\) は二重点イデアルを表す。というのも、上によって、常にあるイデアル \(\frm\) が存在して
\[
\fra\frf\frm=(\alpha),\qquad \frb\frf\frm=(\beta)
\]
が主イデアルになるからである。\(\alpha=0\) と \(\beta=0\) は、点群 \(A\) と \(B\)、または特異点と異なるその部分群 \(A'\) と \(B'\) を切り出す随伴曲線である。\footnote{\(\fra=\frf\fra_1,\ \frb=\frf\frb_1,\ \frr=\frf\frr_1\) であり、\(\fra_1,\frb_1,\frr_1\) が互いに素であるなら、すでに \(\fra_1\frf\mathfrak n,\ \frb_1\frf\mathfrak n\) が主イデアルとなる。ここで \(\mathfrak n\) は \(\frf\) と互いに素である。任意のイデアルは、与えられたイデアルと互いに素なイデアルを掛けることによって主イデアルに変換できるからである（D.--W. p. 214；または数論 § 178, 定理 XI）。} したがって、剰余定理の証明は、算術的には、点を素イデアルに対応させられる可能性と、同値なイデアルを同一のイデアルを掛けることで主イデアルに変換できるという証明の中に含まれている。

余剰対応する二つの点群は、必ずしも同じ数の点からなる必要はないことに注意すべきである。これは同値なイデアルが同じ次数をもつ必要がないのと同様である。例えば、すべての主イデアルは同値であり、完全交差によって生成されるすべての点群は余剰対応である。同値な多角形（除数）\footnote{多角形 \(A\) は絶対 Riemann 面の有限個の点からなる。二つの多角形 \(A\) と \(B\) は、ある関数 \(\eta\) の零点と極であるとき同値と呼ばれる。このとき \(\eta\) は多角形商 \(\eta=\frac{A}{B}\) によって表される。（D.--W. § 17 および 18。）}との相違は、\(z\) が無限大になる点が \(z\) の素イデアルによって生成されないという事実によって説明される。したがって、\(\fra\) と \(\frb\) の同値性を表すイデアル商 \(\eta=\frac{\fra}{\frb}\) としての関数の表示には、\(z\) が無限大になる点での \(\eta\) の零点や極は現れない。例えば主イデアル \((\alpha)\) は、整関数 \(\alpha\) の零点に対応する素イデアルの積になる。\(\alpha\) の極は現れない。なぜなら \(\alpha\) が無限大になるのは、\(z\) が無限大になる点に限られるからである。この点で、イデアル論は幾何学的理論における形式論の正確な類似である。そこでも同様に零点だけがあり、極はない。

形式から関数への移行は、幾何学的には同じ次数の形式の商を作ることによって行われる。算術的には、これはイデアル類から、不変的に定義された多角形（除数）類への移行に対応する。というのも、ここではもはや特別な変数が区別されていないので、関数を多角形商で表示するとすべての零点と極が現れるからである。イデアル類（余剰対応点群の類）と多角形類との関係は次のとおりである。\(\calA\) と \(\calB\) をイデアル \(\fra\) と \(\frb\) によって生成される多角形とし、\(\calN\) を \(z\) が無限大になる点の多角形（\(z\) および \(z\) の整関数の分母多角形）とする。このとき、\(\fra\) と \(\frb\) の同値性から \(\calA\) と \(\calB\) の同値性が従うのは、\(\fra\) と \(\frb\) が同じ次数をもつ場合、かつその場合に限る。一般には、\(\calA\calN^\nu\) と \(\calB\calN^{\nu'}\)（\(\nu\ne\nu'\)）の同値性だけが従う。したがって、イデアル（または余剰対応点群）を類に分けることは、無限個の多角形類を一つの類にまとめることである。それはまた、多角形を \(\calN\) の冪を法として類に分けることとも見なせる。\footnote{Hensel--Landsberg, 第25講, p. 438。--- 二つの変数 \(z\) と \(s\) を多角形商で表すなら、これはこれら各変数の二元斉次化と理解できる。とくに \(z\) と \(s\) が同じ分母多角形をもつように選べば、幾何で通常用いられる三元斉次化が生じる。}

異なる数の点からなる余剰対応点群の場合の剰余定理は、純粋に幾何学的な問題でだけ現れる。代数関数への応用では、実際には多角形類に対応する場合、すなわち同じ数の点をもつ余剰対応点群の場合だけが現れる。一つの点群によって生成される「完全線形系」、すなわちその点群に余剰対応し同じ点数をもつ点群全体は、対応する多角形によって生成される多角形類に正確に対応する。「二つの完全系の和」と「二つの多角形類の積」という概念も一致する。そのような完全系が有限個の線形独立な点群だけを含むことは、剰余定理の直接の帰結である。剰余定理は、系の次元（線形独立な点群の数）を、任意だが固定された次数をもち、剰余群を通るすべての随伴曲線の次元に帰着する。この次元の実際の決定、すなわち Riemann--Roch の定理に到達するために、幾何学的理論はまず剰余定理から「還元定理」（固定点に関する定理）\footnote{Brill--Noether, no. 58。}を導き、そこから Riemann--Roch の定理を導く。

これとまったく対応して、算術的理論は、各多角形類が有限次元であることを示し、整な量の特殊な基底、すなわち補加群によって定義される正規基底によってその次元を決定する。こうして、補加群と分岐多角形との関係の帰結として Riemann--Roch の定理を得る。D.--W. の提示にも還元定理が現れる\footnote{D.--W., § 30, no. 2。}が、幾何学的理論におけるような基礎としてではない。no. 6 で述べた基底関数と随伴関数との対応の結果として、ここでも両理論の補助手段のあいだにある対応を立てることができる。要約すれば、二つの理論における概念形成は、本質的には一致していると言える。ただし、定式化を別にして、その配置は異なる。二つの理論は互いに独立に生じた。幾何学的理論のほうが十年早い。さらに注意すべきことは、Riemann--Roch の定理はもともと、Riemann と Roch においても Weierstrass においても、階数定理として現れるということである。問題となるのは行列 \((\varphi_i(x_j))\) の階数であり、ここで \(\varphi\) は \(n-3\) 次の \(p\) 個の随伴曲線を走り、\(x\) は剰余群の点を走る。

\subsection*{8. 代数関数と代数的数における超越的問題}
Hilbert は、とくに代数関数と代数的数における超越的問題の類似を指摘した。\footnote{Mathematische Probleme, problem 12.（\emph{Göttinger Nachrichten} 1900）} これは本質的には存在問題である。したがって、与えられた Riemann 面（したがって与えられた分岐点）に対して代数関数が存在するかという Riemann の問題は、指定された群と判別式をもつ代数的数体の存在問題に対応する。\footnote{これによって示唆された、群から出発する代数関数の扱いは、R. König によって、より一般的な問題の特殊な場合として遂行された。とくに Riemannsche Funktionen- und Differentialsysteme in der Ebene（\emph{J. f. M.} vol. 148）および \emph{Math. Ann.} 78, 79 を参照。代数体の代わりに、König には \(K(z)\) に関する特殊な有限加群がある。そこでは「類の性質」だけが存在する。} 代数関数の場合には存在証明を一般に実行できるのに対し、代数的数体の領域では、それは特殊な基礎体（有理数、二次数体）について、しかも Abel 群についてだけ達成されている。\footnote{三次および四次方程式に現れる群の単純な場合と、八次方程式のいくつかのきわめて特殊な群だけは、任意の基礎体について扱うことができる：G. Bucht, Über einige algebraische Körper 8. Grades（\emph{Arkiv f. Math., Astron. och Physik}, vol. 6, no. 30）。E. Noether, Gleichungen mit vorgeschriebener Gruppe およびそこで引用された Seidelmann の学位論文（\emph{Math. Ann.} 78）も参照。} これは、有理数上の各 Abel 数体が 1 の根の体から組み立てられるという Kronecker の定理、および相対 Abel 数体についての Fuëter\footnote{\emph{Math. Ann.} 75。} と Hecke\footnote{\emph{Math. Ann.} 76。} の対応する定理に関わる。したがって、求める数体は実際に、指数関数、モジュラー関数という超越的手段によって与えられる。これは代数関数についての純粋な存在証明と対照的である。

代数関数の存在証明には、いたるところ有限な積分の存在が先立つ。これに類似するのは、素数 \(\ell\) に関して定義される代数的数体の「特異主数」の存在である。その \(\ell\) 乗根は相対的に不分岐であり、類体の最初の構成要素を与える。\footnote{Furtwängler, Reziprozitätsgesetze für Potenzreste mit Primzahlexponenten in algebraischen Zahlkörpern. \emph{Math. Ann.} 67 および 72。} この存在証明は、数体において指定された剰余指標をもつ素イデアルが常に存在するという定理の本質的な助けを借りて行われる。したがってこの定理は、いたるところ有限な積分の存在が依拠する境界値問題の類似物と見なすことができる。

いたるところ有限な積分は、二つの多角形が同じ類に属するかどうかを決定する超越的判定条件を与える（同じ数の点からなる二つの点群が余剰対応であるかどうかを決定する）。すなわち Abel の定理による判定である。この定理は、第一種積分を「余剰対応する道」に沿って延長したものが消えると言う。あるいは、対応する微分定理とその逆によっても判定できる。\footnote{例えば M. Noether（\emph{Ann.} 37）の定式化を参照。剰余定理は歴史的には Abel の定理から生じたので、幾何学的理論では積分和に従って点群と線形系の「和」を語る。これは算術的理論における類の「積」と対照的である。}

代数的数体に対する Abel の定理の類似物としては、「\(\ell\) 乗冪剰余に対する相互法則」を取らなければならない。これは、体そのもの、または少なくとも適当な上体について、二つのイデアルが同じ類に属するための判定条件を与える。その内容は、この上体において、特異主数が同じ類のすべてのイデアルに関して同じ剰余指標をもつ、と述べることもできる。この後者の事実は体そのものでも成り立つが、そこでは相互法則とは一致しない。\footnote{Furtwängler, loc. cit.}

最後に、Hensel の \(p\)-進数の理論は、代数関数の冪級数展開の類似物を与える。これについては Hensel の報告を参照されたい。


% 組版用マクロ
\providecommand{\frH}{\mathfrak H}
\providecommand{\frK}{\mathfrak K}
\providecommand{\frM}{\mathfrak M}
\providecommand{\frN}{\mathfrak N}
\providecommand{\frG}{\mathfrak G}
\providecommand{\frS}{\mathfrak S}
\providecommand{\tmod}[1]{\;(\operatorname{mod} #1)}

\clearpage
\begin{center}
{\Large\bfseries 15. 二元形式の整係数不変式系の有限性}\par
\vspace{1em}
\emph{Nachrichten von der Gesellschaft der Wissenschaften zu Göttingen} 1919, pp. 138--156
\end{center}

\vspace{2em}
\begin{center}
1919年3月27日の会合で F. Klein により提出。
\end{center}

以下では、既知の有限性定理の拡張として、次の定理を証明する。
\emph{二元基本形式系のすべての整係数多項式不変式は、そのうち有限個を用いた、整数係数の多項式的有理関数である。}
ここで通常どおり、「不変式」とは常に、基本形式または基本形式系の多項式的有理不変式、すなわち、すべての線形変換の群によって誘導される係数変換に関する不変式を意味する。

証明は、Mertens--Hilbert の有限性証明を整性の面で強めたものに基づく。\footnote{F. Mertens, Wiener Sitzungsberichte Jan. 1889; D. Hilbert, \emph{Math. Ann.} vol. 33 (1889) [1888年3月の日付]. Mertens, Crelle vol. 100 に続いて互いに独立に生じた二つの証明は、内容上は一致している。}
この有限性証明は次のように特徴づけられる。基本形式をその線形因子に分解する、より正確には各基本形式に線形形式の積を対応させる。すると各不変式 \(I\) はこれらの線形形式の不変式、したがって同次化された根の関数に移る。この関数は三つの条件を満たさなければならない。1) 不変性の条件、2) 重み条件、3) 対称性条件である。条件 2) と 3) は、\(I\) が基本形式の係数の多項式的かつ有理的な関数となるために必要である。条件 1) は、\(I\) を線形形式の行列式の多項式的有理関数として表すことによって満たされる。条件 2) は、個々の行列式の代わりに、それらのある冪積を新しい変数として用いることによって満たされる。すると条件 3) は、\(I\) が量の行たちの多項式的有理対称関数になること、また逆に、そのような量の行たちの任意の対称関数が基本形式の不変式になることを述べるだけである。したがって、これらの量の行の基本対称関数が完全系を与える。

整係数不変式だけを考える場合も、線形形式の不変式について、より強い定理、すなわち\emph{二元線形形式のすべての整係数不変式は、それらの線形形式の行列式の整係数多項式関数である}ことが証明されれば、条件 1) は上とまったく同様に満たされる。私は § 3 で、これは実際に成り立つことを、これらの行列式の間の任意の素数を法とするすべての関係の全体が、すべての恒等関係の全体と一致する、という形で示す。すなわち、これらの関係に対応する加群は、任意の素数を法としても素加群であり続ける。この検証は、行列式の間の二次関係に対応する二次形式の加群に関して剰余類を正規化することによって行われる。その際、これらの二次形式は、任意の素数を法として加群の整基底をなすことが分かる。従来それらは、整性を無視したときにだけ加群基底として知られていたものである。

条件 2) を満たすためには変更は不要である。冪積の導入は整数係数に影響しないからである。条件 3) はいま、\(n_1!\cdots n_\mu! I\)（ここで \(n_1,\ldots,n_\mu\) は個々の基本形式の次数である）が、量の行たちの整係数対称関数になり、逆に、これらの行の任意の整係数対称関数が整係数不変式になることをいう。しかし容易に示されるように、これらの行の整係数対称関数は整基底をもち、その基底には基本関数だけではもはや十分でない。Hilbert の整加群基底に関する定理を適用すれば、したがって \(I\) 自身に対しても整基底が得られる。

§ 1 では、条件 1), 2), 3) に対応する特殊な有限性定理を集める。それらから § 2 で有限性の証明が従う。条件 3) に対応する特殊な有限性定理は、整数または代数的整数の置換からなる有限群の整係数不変式の有限性も与える。これを § 4 で簡潔に述べる。

\section*{§ 1. 特殊な有限性定理}

まず、導入で述べた不変性、重み、対称性の条件を満たすために用いる、次の特殊な有限性定理を述べる。

\textbf{I.} 二元線形形式の系の任意の整係数不変式は、これらの線形形式の行列式の、整数係数をもつ多項式的有理関数である。証明は § 3 で与える。

\textbf{II.} 非負指数をもつ \(n\) 個の変数の冪積からなる任意の系 \(\frS\) が、二つの冪積 \(f,g\) とともに、その商が変数に関する多項式である限りその商も含むとする。このとき \(\frS\) の関数からなる有限の乗法的基底 \(f_1,\ldots,f_k\) が存在し、\(\frS\) の任意の \(f\) に対して
\[
        f=f_1^{\lambda_1}\cdots f_k^{\lambda_k}
\]
という表示が、非負指数 \(\lambda_1,\ldots,\lambda_k\) によって成り立つ。

実際、Hilbert の加群基底定理により、\(\frS\) の関数の有限個 \(f_1,\ldots,f_k\) を、各々が実際に \(x\) を含むように選べて、\(\frS\) の任意の非定数 \(f\) はそれらで生成される加群に属する。
\[
        f\equiv0\pmod{(f_1,\ldots,f_k)}.
\]
したがって、任意の \(f=x_1^{i_1}\cdots x_n^{i_n}\) に対して、少なくとも一つの \(f_\nu=x_1^{\nu_1}\cdots x_n^{\nu_n}\) が存在し、\(\nu_1\le i_1,\ldots,\nu_n\le i_n\) を満たす。表示
\[
        f=x_1^{i_1-\nu_1}\cdots x_n^{i_n-\nu_n} f_\nu = A f_\nu
\]
において、因子 \(A\) は仮定により \(\frS\) に属するが、変数に関する次数は \(f\) より小さい。したがってこの議論を有限回繰り返すことで主張が従う。表示を \(f=1\) に対しても成り立たせるには、すべての指数 \(\lambda\) をゼロにすればよい。\footnote{この表示は直接にも証明でき、逆に Hilbert の定理はこの表示から従う。P. Gordan, Neuer Beweis des Hilbertschen Satzes über homogene Funktionen, Gött. Nachr. 1899 を参照。別の直接証明、およびそこからの定理 II の導出は、A. Ostrowski, Über die Existenz einer endlichen Basis bei gewissen Funktionensystemen, \emph{Math. Ann.} 78 (1916), § 2,1 および § 3,2 にある。§ 2 で用いる定理 II の特殊な場合、すなわち指数がディオファントス方程式系のすべての非負解を動く場合は、すでに Gordan-Kerschensteiner, \emph{Vorlesungen über Invariantentheorie}, vol. I, p. 199 に見られる。}

\textbf{III.} \(n\) 個の量の行に関するすべての多項式的整係数対称関数は、そのうち有限個の多項式的整係数関数である。

\(\tau_1(x),\ldots,\tau_n(x)\) を不定元 \(x_1,\ldots,x_n\) の基本対称関数とする。すると任意の指数 \(k\) について恒等式
\[
        x_i^k = G_0^{(k)}(\tau)+x_iG_1^{(k)}(\tau)+\cdots+x_i^{n-1}G_{n-1}^{(k)}(\tau),
        \qquad (i=1,2,\ldots,n),
\]
があり、ここで \(G_0^{(k)},\ldots,G_{n-1}^{(k)}\) は \(\tau(x)\) の整係数多項式関数である。\(x_i,y_i,\ldots,z_i\) が第 \(i\) 行の要素であるとき、これらの恒等式および \(y,\ldots,z\) についての対応する恒等式によって、すべての特殊な一行型対称関数
\[
        \sum_{i=1}^n x_i^{a}y_i^{b}\cdots z_i^{c}
\]
は、有限個の
\[
        \sum_{i=1}^n x_i^{a}y_i^{b}\cdots z_i^{c}
        \qquad
        (0\le a<n,\;0\le b<n,\;\ldots,\;0\le c<n)
\]
および基本対称関数 \(\tau(x),\tau(y),\ldots,\tau(z)\) の整係数多項式関数になる。\footnote{Hilbert, loc. cit. は、\(\tau(x),\ldots,\tau(z)\) の代わりに冪和を基底に入れる。そのため基底は一行型関数だけからなるが、表示の整性は失われる。}
\(n\) 行の最も一般の対称関数を考える場合（以下では生じない場合）も、同様にして、関数
\[
        \sum x_{i_1}^{a_1}y_{i_1}^{b_1}\cdots z_{i_1}^{c_1}
              x_{i_2}^{a_2}y_{i_2}^{b_2}\cdots z_{i_2}^{c_2}
              \cdots
              x_{i_n}^{a_n}y_{i_n}^{b_n}\cdots z_{i_n}^{c_n}
\]
\[
        (0\le a_\nu<n,\;0\le b_\nu<n,\;\ldots,\;0\le c_\nu<n),
\]
ここで \(i_1,i_2,\ldots,i_n\) はすべての置換を動くが、これらの関数は \(\tau(x),\tau(y),\ldots,\tau(z)\) とともに整基底をなすことが示される。

\textbf{IV.} \(F_1(x),\ldots,F_k(x)\) を \(x_1,\ldots,x_n\) の整係数多項式関数とし、\(a\) を固定された整数とする。このとき
\[
        \frac{G(F)}{a},
\]
ただし \(G\) は整係数多項式関数を表す、と書ける関数のうち、\(x\) に関して整になるものすべては、そのうち有限個の整係数多項式関数である。

実際、Hilbert の整加群基底に関する定理により、そのような各 \(G(F)/a\) について表示
\[
 \frac{G(F)}{a}
 =
 A_1(F_1\ldots F_k)\frac{G_1(F)}{a}
 +\cdots+
 A_\varrho(F_1\ldots F_k)\frac{G_\varrho(F)}{a},
\]
があり、\(A_1,\ldots,A_\varrho\) は \(F\) の整係数多項式関数であり、
\[
        \frac{G_1(F)}{a},\ldots,\frac{G_\varrho(F)}{a}
\]
はその系に属する。したがって関数
\[
        F_1=\frac{aF_1}{a},\ldots,F_k=\frac{aF_k}{a};
        \quad
        \frac{G_1(F)}{a},\ldots,\frac{G_\varrho(F)}{a}
\]
が求める基底をなす。

\setcounter{equation}{0}
\section*{§ 2. 二元整係数不変式の有限性}

\(x,y\) を不定係数をもつ変換
\begin{equation}
        x=s_{11}x'+s_{12}y',\qquad y=s_{21}x'+s_{22}y',
\end{equation}
を受ける二元変数とし、\(f\) を係数が不定元である二元基本形式とする。このとき (1) により
\begin{equation}
        f=a_0x^n+a_1x^{n-1}y+\cdots+a_ny^n
        =a_0'x'^n+a_1'x'^{\,n-1}y'+\cdots+a_n'y'^n.
\end{equation}
\(f\) の整係数不変式とは、周知のように、整係数多項式関数 \(I(a)\) であって
\begin{equation}
        I(a')=|s_{ik}|^\varrho I(a),\qquad \text{\(a,s_{ik}\) に関して恒等的に}
\end{equation}
を満たすものをいう。基本形式系の整係数不変式も対応して定義される。個々の基本形式の係数に関する同次性を定義に含めてもよい。なぜなら、残りの不変式は、そうした個別に同次な不変式から整線形に構成されるからである。

(2) に加えて、特別な二元形式 \(g\)、すなわち係数がやはり不定元である \(n\) 個の線形形式の積を考える。\footnote{不変式論でしばしば行われるように、\(f\) を二項係数付きで \(f=b_0x^n+\binom n1b_1x^{n-1}y+\cdots+b_ny^n\) と書くならば、すべての整係数不変式 \(I(b)\) の領域は \(I(a)\) の領域に比べて拡大し、ここで用いる有限性の議論は成り立たなくなる。\(I(b)\) はもはや根の整関数ではない。}
\begin{align}
        g&=(\alpha_1x+\beta_1y)(\alpha_2x+\beta_2y)\cdots(\alpha_nx+\beta_ny)\notag\\
         &=\sigma_0(\alpha,\beta)x^n+\sigma_1(\alpha,\beta)x^{n-1}y+\cdots+\sigma_n(\alpha,\beta)y^n\notag\\
         &=(\alpha_1'x'+\beta_1'y')(\alpha_2'x'+\beta_2'y')\cdots(\alpha_n'x'+\beta_n'y')\notag\\
         &=\sigma_0'x'^n+\sigma_1'x'^{\,n-1}y'+\cdots+\sigma_n'y'^n.
\end{align}
したがって \(\sigma_0(\alpha,\beta),\sigma_1(\alpha,\beta),\ldots,\sigma_n(\alpha,\beta)\) は同次化された基本対称関数を表し、変換後の係数 \(\sigma_i'\) は \(\sigma_i(\alpha',\beta')\) に等しくなる。

これらの基本関数の代数的独立性により、任意の関係
\begin{equation}
        G(\sigma_0,\sigma_1,\ldots,\sigma_n)=0
        \quad \text{\(\alpha,\beta\) に関して恒等的に}
\end{equation}
には、さらに関係
\begin{equation}
        G(\sigma_0,\sigma_1,\ldots,\sigma_n)=0
        \quad \text{\(\sigma_0,\ldots,\sigma_n\) に関して恒等的に}
\end{equation}
が対応し、したがって不定元 \(a_0,\ldots,a_n\) について
\begin{equation}
        G(a_0,a_1,\ldots,a_n)=0
        \quad \text{\(a_0,\ldots,a_n\) に関して恒等的に}
\end{equation}
が成り立つ。

いま \(f\) の任意の整係数不変式 \(I(a)\) は、置換 \(a_i=\sigma_i(\alpha,\beta)\) によって、\(g\) の整係数不変式 \(I(\sigma)=H(\alpha,\beta)\) を与える。このとき \(H\) は \(\alpha,\beta\) の整関数である。この不変式に対して関係 (3) は
\begin{equation}
        I(\sigma')=|s_{ik}|^\varrho I(\sigma)
        \quad \text{\(\sigma,s_{ik}\) に関して恒等的に}
\end{equation}
となり、\(\sigma'=\sigma(\alpha',\beta')\) であるから、(8) は
\begin{equation}
        H(\alpha',\beta')=|s_{ik}|^\varrho H(\alpha,\beta)
        \quad \text{\(\alpha,\beta;s_{ik}\) に関して恒等的に}
\end{equation}
を与える。したがって \(I(\sigma)=H(\alpha,\beta)\) は、\(n\) 個の線形形式 \((\alpha_i x+\beta_i y)\) の整係数不変式になる。逆に、\(H(\alpha,\beta)\) がこれら \(n\) 個の線形形式の整係数不変式、すなわち (9) を満たすものであり、同時に \(\sigma\) の整関数 \(I(\sigma)\) に等しいと仮定する。このとき
\[
        H(\alpha',\beta)=I(\sigma(\alpha',\beta'))=I(\sigma')
\]
であるから、式 (9) は
\begin{equation*}
        I(\sigma')=|s_{ik}|^\varrho I(\sigma)
        \quad \text{\(\alpha,\beta;s_{ik}\) に関して恒等的に}.\tag{9a}
\end{equation*}
とも書ける。(5) と (6) から (8) が従い、(7) から (3) が従う。したがって各整係数不変式 \(I(a)\) には \(n\) 個の線形形式の整係数不変式 \(H(\alpha,\beta)=I(\sigma)\) が対応し、逆も成り立つ。さらに、\(I_1(a),\ldots,I_k(a)\) がすべての \(I(a)\) の整基底ならば、\(H_1(\alpha,\beta),\ldots,H_k(\alpha,\beta)\) は、同時に整関数 \(I(\sigma)\) であるすべての \(H(\alpha,\beta)\) の整基底である。ここでも逆が (5), (6), (7) により成り立つ。したがって、すべての \(I(a)\) に対する整基底の存在を証明するには、同時に整関数 \(I(\sigma)\) であるすべての整係数不変式 \(H(\alpha,\beta)\) についてそれを証明すれば十分であり、基底
\[
        H_1(\alpha,\beta)=I_1(\sigma),\ldots,H_k(\alpha,\beta)=I_k(\sigma)
\]
が \(I(a)\) に対する基底 \(I_1(a),\ldots,I_k(a)\) を与える。

\(H(\alpha,\beta)=I(\sigma)\) に (4) のすべての整係数不変式を動かせる。(3) により各不変式は \(\sigma\) に関して同次であり、(4) により各 \(\sigma\) は各線形形式の係数に関して同次かつ一次であるから、\(H(\alpha,\beta)\) は各行 \(\alpha_i,\beta_i\) に関して同次で、各行において同じ次元をもつ。さらにもちろん、\(n\) 行の対称関数である。逆に、対称関数に関する基本定理により、\(n\) 行の任意の整係数対称関数 \(H(\alpha,\beta)\) が、各行に関して別々に同次で、各行において同じ次元をもつならば、それは \(\sigma\) の整係数多項式関数である。したがって前の議論により、\(H(\alpha,\beta)\) が不変であれば、ただちに不変式 \(I(\sigma)\) である。よって (4) の整係数不変式の全体は、次の性質をもつ整関数 \(H(\alpha,\beta)\) の全体として特徴づけられる。
\[
\begin{array}{ll}
1) & \text{\(n\) 個の線形形式 \(\alpha_i x+\beta_i y\) の整係数不変式である;}\\
2) & \text{\(n\) 個の各行 \(\alpha_i,\beta_i\) に関して別々に同次で、同じ次元をもつ;}\\
3) & \text{\(n\) 個の行 \(\alpha_i,\beta_i\) に関して対称である。}
\end{array}
\]
特殊有限性定理 I により、\(H(\alpha,\beta)\) は 1) のために、行列式
\[
        (ik)=\alpha_i\beta_k-\beta_i\alpha_k;
\]
の整係数多項式関数になる。したがって
\[
        H(\alpha,\beta)=G((ik)).
\]
この表示において、対称性は \(n\) 行のすべての置換を \(G((ik))\) に適用することで形式的にも表せる。
\begin{equation}
        n!\,H(\alpha,\beta)
        =
        G^{(1)}((ik))+G^{(2)}((ik))+\cdots+G^{(n!)}((ik))
        =G^*((ik)).
\end{equation}
したがって \(G^*\) は、\((ik)\) の各冪積 \(P\) とともに、\(P\) のすべての置換にわたる和 \(Q=P^{(1)}+\cdots+P^{(n!)}\) も含む。実際、\(G^*\) は有限個の \(Q\) の整線形結合 \(G^*=\sum cQ\) になる。\(Q\) は条件 1), 2), 3) を満たし、したがってそれ自身不変式になるから、\(G^*((ik))\) の代わりに \(Q\) を考えれば十分である。

いま \(P\) を、\(Q\) に現れるすべての \((ik)\) の冪積、すなわち条件 2) を満たす冪積の系 \(\frS\) の中で動かす。二つの \(P\) の商が \((ik)\) に関して多項式ならば、それも条件 2) を満たし、したがって \(\frS\) に属する。したがって定理 II により、このような \(P\) は有限個 \(P_1,\ldots,P_\varrho\) を用いて、任意の \(P\) が
\[
        P=P_1^{\lambda_1}\cdots P_\varrho^{\lambda_\varrho}
\]
という形に、非負指数 \(\lambda\) で表される。すべての置換を適用すると
\[
        Q=
        P_1^{(1)\lambda_1}\cdots P_\varrho^{(1)\lambda_\varrho}
        +\cdots+
        P_1^{(n!)\lambda_1}\cdots P_\varrho^{(n!)\lambda_\varrho}
\]
となる。こうして \(Q\) は、それぞれ \(\varrho\) 個の要素をもつ \(n!\) 個の量の行
\[
        P_1^{(1)}\ldots P_\varrho^{(1)};
        \quad
        P_1^{(2)}\ldots P_\varrho^{(2)};
        \quad \ldots;\quad
        P_1^{(n!)}\ldots P_\varrho^{(n!)};
\]
に関する整係数対称関数になる。よって定理 III により、それはこれらの行に関する有限個の対称関数の整係数多項式関数である。これらの基底関数は、それ自身 \(Q\) であるか、または \(n!\) 個の要素 \(P_i^{(1)},P_i^{(2)},\ldots,P_i^{(n!)}\) の基本対称関数であり、したがって条件 1), 2), 3) を満たし、整係数不変式
\begin{equation}
        H_1(\alpha,\beta)=I_1(\sigma);\ldots;\quad
        H_l(\alpha,\beta)=I_l(\sigma)
\end{equation}
になる。(10) と \(G^*=\sum cQ\) により、任意の \(H(\alpha,\beta)=I(\sigma)\) はしたがって
\[
        H(\alpha,\beta)=I(\sigma)=\frac1{n!}\Phi(I_1,\ldots,I_l)
\]
の形である。ここで \(\Phi\) は整係数多項式関数を表す。逆に、\(\frac1{n!}\Phi\) が \(\alpha,\beta\) に関して多項式であれば、それは整係数不変式である。したがって定理 IV により、さらに有限個の不変式 \(I_{l+1}(\sigma),\ldots,I_k(\sigma)\) を (11) に付け加えて、
\[
        H_1(\alpha,\beta)=I_1(\sigma);\ldots;H_l(\alpha,\beta)=I_l(\sigma);
\]
\[
        H_{l+1}(\alpha,\beta)=I_{l+1}(\sigma);\ldots;H_k(\alpha,\beta)=I_k(\sigma)
\]
がすべての不変式 \(H(\alpha,\beta)=I(\sigma)\) の整基底をなすようにできる。これで一つの基本形式の場合の定理が証明された。

基本形式の系を扱う場合には、先に述べたように、各基本形式の係数に関する不変式の同次性を仮定してよい。再び一般の基本形式を、それぞれ線形形式の積であるものに置き換える。その不変式は、1) これらの線形形式の不変式、2) 線形形式の係数からなる行に関して別々に同次で、各個別の基本形式に属する行において等しい次元をもつもの、3) 各個別の基本形式に属する行に関して対称なもの、になる。逆に、これらの条件を満たすすべての線形形式の整係数不変式は、再び基本形式系の整係数不変式に対応する。証明は今与えたものと完全に平行である。重み条件（定理 I と II）を満たす個々の行列式の冪積 \(P=P_1^{\lambda_1}\cdots P_\varrho^{\lambda_\varrho}\) ごとに、すべての
\[
        N=n_1!n_2!\cdots n_\mu!
\]
個の置換を適用する。ここで \(n_1,\ldots,n_\mu\) は基本形式の次数である。このようにして \(N\cdot I\) は、それぞれ \(\varrho\) 個の要素をもつ \(N\) 個の行の整係数対称関数に変換される。これからすべての \(N\cdot I\) に対する整基底が得られ、したがって定理 IV により、すべての \(I\) に対する整基底が得られる。

\setcounter{equation}{0}
\section*{§ 3. 二元線形形式の整係数不変式の基底}

特殊有限性定理 I の証明を残している。Clebsch によって初めて証明された既知の定理により、\(n\) 個の線形形式 \(\alpha_i x+\beta_i y\) のすべての多項式的有理不変式は、行列式 \((ik)\) の多項式的有理関数である。したがって、とくに、これらの線形形式のすべての整係数不変式は、これらの行列式の多項式的有理関数になるが、必ずしも整係数関数になるとは限らない。

以下で示すのは、\((ik)\) の間に存在する恒等関係を考慮すると、この表示は常に整係数で可能であるということである。

式で、かつ加群論の言葉で、これは次のように表される。\(\xi_{ik}\), \(i<k\), を不定元とする。このとき、\(F((ik))\) が \(\alpha,\beta\) に関して恒等的に消えるという性質をもつすべての整係数多項式 \(F(\xi_{ik})\) の全体は加群 \(\frM\) をなす。そのような二つの \(F\) の和、およびそのような \(F\) と \(\xi_{ik}\) の任意の整係数多項式との積は、再びこの全体に属する。式では
\[
        F((ik))=0\quad \text{\(\alpha,\beta\) に関して恒等的に}
\]
ならば
\[
        F(\xi_{ik})\equiv0\pmod{\frM}\quad \text{\(\xi_{ik}\) に関して恒等的に},
\]
かつその逆も成り立つ。

いま次の完全に類似した主張を示す。
\[
        F((ik))\equiv0\pmod p\quad \text{\(\alpha,\beta\) に関して恒等的に}
\]
ならば
\begin{equation}
        \text{次が従う:}\qquad
        F(\xi_{ik})\equiv0\pmod{(\frM,p)}
        \quad \text{\(\xi_{ik}\) に関して恒等的に},
\end{equation}
かつその逆も、任意の素数 \(p\) について成り立つ。ゆえに \(\frM\) は、\((ik)\) の間の \(p\) を法とするすべての関係の全体に対応する加群 \(\frM_p\) とも同一である。この加群は、\(F((ik))\) が \(\alpha,\beta\) に関して恒等的に \(p\) を法として消えるような、すべての整係数多項式 \(F(\xi_{ik})\) からなる。

(1) が整係数表示可能性を述べていることは次のように分かる。それはまた、\(F((ik))=p\cdot H(\alpha,\beta)\) が \(\alpha,\beta\) に関して恒等的に成り立つならば、
\[
        F(\xi_{ik})-p\cdot G(\xi_{ik})\equiv0\pmod{\frM}
        \quad \text{\(\xi_{ik}\) に関して恒等的に}
\]
が従う、という形にも書ける。したがって \(\frM\) の定義により、
\[
        F((ik))=p\cdot G((ik))\quad \text{\(\alpha,\beta\) に関して恒等的に}.\footnote{この最後の結論は同時に (1) の逆も証明している。}
\]
ゆえに \(p\cdot H(\alpha,\beta)=p\cdot G((ik))\)、すなわち
\[
        H(\alpha,\beta)=G((ik)).
\]
したがって、もし表示
\[
        H(\alpha,\beta)
        =
        \frac1{p_1^{\lambda_1}\cdots p_\varrho^{\lambda_\varrho}}
        F((ik)),
\]
ただし \(p_1,\ldots,p_\varrho\) は素数、が与えられているならば、同時に
\[
        H(\alpha,\beta)
        =
        \frac1{p_1^{\lambda_1-1}\cdots p_\varrho^{\lambda_\varrho}}
        G((ik))
\]
も従い、これを有限回繰り返せば表示の整性が得られる。

(1) の証明のため、
\[
        f_{ik,rs}=\xi_{ik}\xi_{rs}-\xi_{ir}\xi_{ks}+\xi_{is}\xi_{kr},
        \qquad i<k<r<s,
\]
で、\((ik)\) の間の二次関係
\[
        (ik)(rs)-(ir)(ks)+(is)(kr)=0
\]
に対応する \(\xi_{ik}\) の形式を表すものとし、これらから生成される整係数多項式の加群を \(\frN=(f_{ik,rs})\) とする。
\[
        F(\xi_{ik})\equiv0\pmod{\frN}
\]
からは
\[
        F((ik))=0\quad \text{\(\alpha,\beta\) に関して恒等的に}
\]
が従い、また
\begin{equation}
\begin{array}{rcl}
        \text{もし }F(\xi_{ik})&\equiv&0\pmod{(\frN,p)}\text{ ならば}\\
        F((ik))&\equiv&0\pmod p\quad\text{\(\alpha,\beta\) に関して恒等的に}
\end{array}
\end{equation}
が従う。以下では、\(\frN\) に関して剰余類を正規化することで、(2) の逆も成り立つことを示す。すると \(\frN=\frM=\frM_p\) となり、(1) が証明される。等式 \(\frN=\frM\) は、\(p=0\) の場合の \(\frN=\frM_p\) として得られる。以下の考察では、\(p=0\) の場合も常に含めておく。言い換えれば、\(f_{ik,rs}\) は \(\frM_p\) の整加群基底をなし、とくに \(\frM_0=\frM\) の基底でもある。

\subsection*{a) 二つまたは三つの行 \(\alpha_i,\beta_i\) の場合}

四つの異なる指数が現れないので、\(\frN\) は零加群である。したがって、
\begin{equation}
        F((ik))\equiv0\pmod p\quad \text{\(\alpha,\beta\) に関して恒等的に}
\end{equation}
から
\[
        F(\xi_{ik})\equiv0\pmod p\quad \text{\(\xi_{ik}\) に関して恒等的に}
\]
が従うことを示せばよい。これは \(\frM=\frM_p=0\) として (1) を証明する。

条件 \(F((ik))=p\cdot H(\alpha,\beta)\) は、個々の行に関して同次な成分ごとに別々に満たされなければならないので、以下では常に、\(F((ik))\) が各個別の行に関して同次、したがって \(F(\xi_{ik})\) が各指数に関して同次であると仮定して十分である。二つの行の場合には、行列式 \((12)\) と対応する不定元 \(\xi_{12}\) だけが存在する。仮定した同次性により
\[
        F=c\,\xi_{12}^x.
\]
\[
        c\,(12)^x\equiv0\pmod p\quad \text{\(\alpha,\beta\) に関して恒等的に}
\]
から、\((12)\not\equiv0\pmod p\) は \(\alpha,\beta\) に関して恒等的には成り立たないので、\(c\equiv0\pmod p\) が従う。ゆえに
\[
        F(\xi_{ik})\equiv0\pmod p\quad \text{\(\xi_{ik}\) に関して恒等的に}.
\]
これで (3)、したがって (1) が証明された。

三つの行の場合には、行列式 \((12),(13),(23)\) と対応する不定元 \(\xi_{12},\xi_{13},\xi_{23}\) がある。したがって
\[
        F=\sum c\,\xi_{12}^{\tau_{12}}\xi_{13}^{\tau_{13}}\xi_{23}^{\tau_{23}}.
\]
\(F\) の指数 \(1,2,3\) に関する次数をそれぞれ \(\sigma_1,\sigma_2,\sigma_3\) とする。このとき
\[
        \sigma_1=\tau_{12}+\tau_{13},\qquad
        \sigma_2=\tau_{12}+\tau_{23},\qquad
        \sigma_3=\tau_{13}+\tau_{23},
\]
であり、一意な逆写像
\[
        \tau_{12}=\frac12(\sigma_1+\sigma_2-\sigma_3),\quad
        \tau_{13}=\frac12(\sigma_1-\sigma_2+\sigma_3),\quad
        \tau_{23}=\frac12(-\sigma_1+\sigma_2+\sigma_3)
\]
をもつ。したがって各々の別々に同次な \(F\) は高々一項
\[
        F=c\,\xi_{12}^{\tau_{12}}\xi_{13}^{\tau_{13}}\xi_{23}^{\tau_{23}}
\]
を含むだけであり、そこから上と同様に (3)、したがって (1) が従う。

\subsection*{b) 四つの行 \(\alpha_i,\beta_i\) の場合}

六つの不定元 \(\xi_{12},\xi_{13},\xi_{14},\xi_{23},\xi_{24},\xi_{34}\) と対応する行列式が現れ、\(\frN=(f_{12,34})\) である。なぜなら
\begin{equation}
        \xi_{14}\xi_{23}\equiv \xi_{13}\xi_{24}-\xi_{12}\xi_{34}\pmod{\frN},
\end{equation}
であるから、任意の \(\xi_{ik}\) の冪積は、\(\xi_{14}\xi_{23}\) という積を含まない冪積の整線形結合に \(\frN\) を法として合同である。したがって任意の多項式 \(F(\xi_{ik})\) に対して
\begin{equation}
        F(\xi_{ik})\equiv F_0(\xi_{ik})\pmod{\frN},
\end{equation}
ただし \(F_0\) は \(\xi_{14}\xi_{23}\) という積を含まない。\(F_0\) を \(F\) の剰余類の正規化多項式と呼ぶ。(4) は別々に同次であるから、別々に同次な \(F\) には再び別々に同次な \(F_0\) が対応する。

いま
\begin{equation}
        F_0(\xi_{ik})=G(\xi_{ik})+\xi_{34}\cdot F_1(\xi_{ik}),
\end{equation}
とおく。ただし \(G\) は \(\xi_{34}\) で割れる冪積を含まない。さらに
\begin{equation}
        [G]_{\xi_{14}=\xi_{13}}=K(\xi_{12},\xi_{13},\xi_{23})
\end{equation}
とおく。
\[
        F((ik))\equiv0\pmod p\quad \text{\(\alpha,\beta\) に関して恒等的に}
\]
から、(2) を考慮した (5) により、
\begin{equation}
        F_0((ik))\equiv0\pmod p\quad \text{\(\alpha,\beta\) に関して恒等的に}
\end{equation}
が従う。\(\alpha_4=\alpha_3,\;\beta_4=\beta_3\) と特殊化すると、(6) と (7) は
\[
        K((ik))\equiv0\pmod p\quad \text{\(\alpha,\beta\) に関して恒等的に}
\]
を与える。したがって a) で証明した結果により、
\[
        K(\xi_{ik})\equiv0\pmod p\quad \text{\(\xi_{ik}\) に関して恒等的に}
\]
である。これから下で示すべき残りの主張
\[
        G(\xi_{ik})\equiv0\pmod p\quad \text{\(\xi_{ik}\) に関して恒等的に}
\]
が従う。したがって (6) は
\[
        F_0(\xi_{ik})\equiv \xi_{34}\cdot F_1(\xi_{ik})\pmod p
        \quad \text{\(\xi_{ik}\) に関して恒等的に}
\]
を与える。\((34)\not\equiv0\pmod p\) は \(\alpha,\beta\) に関して恒等的には成り立たないので、(8) から
\[
        F_1((ik))\equiv0\pmod p\quad \text{\(\alpha,\beta\) に関して恒等的に}
\]
が従う。ここで \(F_1(\xi_{ik})\) は再び正規化されており、\(\xi_{34}\) に関する次数は低くなっている。この手続きを有限回繰り返すと
\[
        F_0(\xi_{ik})\equiv\xi_{34}^{\lambda}F_\lambda(\xi_{ik})\pmod p,
\]
ただし \(F_\lambda\) は \(\xi_{34}\) に関して次数ゼロであり、したがって直前に示したことにより \(p\) を法として消える。よって
\begin{equation}
\begin{array}{rcl}
        F_0(\xi_{ik})&\equiv&0\pmod p\quad \text{\(\xi_{ik}\) に関して恒等的に},\quad\text{または}\\
        F(\xi_{ik})&\equiv&0\pmod{(\frN,p)}\quad \text{\(\xi_{ik}\) に関して恒等的に}
\end{array}
\end{equation}
を得る。これは (2) の逆を、したがって (1) を、同時に \(F_0(\xi_{ik})\) に関するより鋭い主張として証明して\mbox{いる。\footnotemark}\footnotetext{このより鋭い主張は、各剰余類がただ一つの正規化多項式を含むという意味である。}

あとは、\(K(\xi_{ik})\equiv0\pmod p\) から \(G(\xi_{ik})\equiv0\pmod p\) が従うことを付け加えるだけである。同値に、\(G(\xi_{ik})\not\equiv0\pmod p\) から \(K(\xi_{ik})\not\equiv0\pmod p\) が従う。正規化が必要なのは誘導のこの部分だけである。置換 \(\xi_{14}=\xi_{13}\) のもとで、\(G\) の個々の冪積は消えない。したがって \(K\) が消えるとすれば、それは \(G\) における異なる冪積が \(K\) では等しい冪積に対応するからでしかない。そしてこれは、それらの冪積が単に指数 3 と 4 の入れ替えによって異なる場合に限って起こり得る。たとえば、そのような正規化された冪積が
\[
        \xi_{12}^{\tau_{12}}\xi_{13}^{\tau_{13}}\xi_{14}^{\tau_{14}}\xi_{24}^{\tau_{24}}
\]
であるとしよう。各指数に関する別々の同次性の仮定により、一つの \(\xi_{ik}\) における 3 から 4 への置き換えは、別の \(\xi\) における 4 から 3 への置き換えを伴わなければならない。したがって関係する項は \(\xi_{13}\xi_{24}\) の代わりに \(\xi_{14}\xi_{23}\) という積を含むことになり、これは正規化によって排除されている。まったく同様に、
\[
        \xi_{12}^{\tau_{12}}\xi_{13}^{\tau_{13}}\xi_{14}^{\tau_{14}}\xi_{24}^{\tau_{24}}
\]
に対して相殺し得るのは、因子 \(\xi_{14}\xi_{23}\) を含む項だけであり、これも正規化により排除されている。ゆえに \(K(\xi_{ik})\equiv0\pmod p\) は \(G(\xi_{ik})\equiv0\pmod p\) を含意する。

\subsection*{c) \(n\) 個の行 \(\alpha_i,\beta_i\) の場合}

\(n\) 個の行の場合は、いま扱った四つの行の場合の厳密な一般化として、完全帰納法で解決される。\footnote{証明は \(n=3\) の場合にもなお適用される。}
まず再び、任意の多項式 \(F\) に対して合同
\begin{equation}
        F(\xi_{ik})\equiv F_0(\xi_{ik})\pmod{\frN}
        \quad \text{\(\xi_{ik}\) に関して恒等的に}
\end{equation}
が存在することを示す。ここで \(\frN=(f_{ik,rs})\) で、\(F_0(\xi_{ik})\) は正規化されている。正規化とは次のことを意味するものとする。\(i,k,r,s\) が四つの異なる指数で \(i<k<r<s\) のとき、\(F_0\) は積 \(\xi_{is}\xi_{kr}\) を含んではならない。すなわち、一つの \(\xi\) の二つの指数が、もう一つの \(\xi\) の二つの指数の間にともに入ってはならない。この正規化は、二つおよび三つの行の場合には何の制限も課さない。そこではすべての多項式がすでに正規化されているからである。四つの行の場合には、すでに示したように、各剰余類に正規化多項式が存在する。したがって \((n-1)\) 個の行については、各剰余類に正規化多項式が存在することを証明済みと仮定してよい。

\(F\) における任意の冪積を
\[
        \xi_{i_1n}\cdots\xi_{i_\sigma n}\,P
\]
の形に書く。ただし \(P\) はもはや指数 \(n\) を含まない。誘導仮定により、\(P\) は \(\frN\) を法として整係数の正規化多項式 \(P_0\) に合同である。
\[
        \xi_{i_1n}\cdots\xi_{i_\sigma n}P
        \equiv
        \xi_{i_1n}\cdots\xi_{i_\sigma n}P_0
        \pmod{\frN}.
\]
もし \(P_0\) が、少なくとも一つの指数 \(i_\lambda\) に対して二つの指数 \(r,s\) が \(i_\lambda\) と \(n\) の間にある、すなわち \(i_\lambda<r<s<n\) となるような \(\xi_{rs}\) を含まないならば、すでに正規化に到達している。なぜなら、\(P_0\) から取った任意の二つの \(\xi\) の積について仮定が満たされるからである。そうでない場合、そのような指数 \(i_\lambda\), \(i_\lambda<r<s<n\), を取り、\(P_0^*\) を \(\xi_{rs}\) を含む \(P_0\) の冪積とする。
\[
        \xi_{i_\lambda n}\xi_{rs}\equiv
        \xi_{i_\lambda s}\xi_{rn}-\xi_{i_\lambda r}\xi_{sn}
        \pmod{\frN}
\]
であるから、
\begin{align*}
 &\xi_{i_1n}\cdots\xi_{i_\lambda n}\cdots\xi_{i_\sigma n}P_0^*\\
 &\equiv
 \xi_{i_1n}\cdots\xi_{rn}\cdots\xi_{i_\sigma n}Q
 +
 \xi_{i_1n}\cdots\xi_{sn}\cdots\xi_{i_\sigma n}R
 \pmod{\frN}\\
 &\equiv
 \xi_{i_1n}\cdots\xi_{rn}\cdots\xi_{i_\sigma n}Q_0
 +
 \xi_{i_1n}\cdots\xi_{sn}\cdots\xi_{i_\sigma n}R_0
 \pmod{\frN},
\end{align*}
を得る。ここで \(Q_0\) と \(R_0\) は、\(Q\) と \(R\) の剰余類の正規化多項式を再び表す。これらは誘導仮定により存在する。\(Q\) と \(R\) は指数 \(n\) を含まないからであり、しかも \(Q\) と \(R\) は単なる冪積なので、それらは確かに整係数多項式である。\(i_\lambda<r<s<n\) であるから、さらに
\[
 ((n-i_1)+\cdots(n-r)+\cdots(n-i_\sigma))
 <
 ((n-i_1)+\cdots(n-i_\lambda)+\cdots(n-i_\sigma)),
\]
\[
 ((n-i_1)+\cdots(n-s)+\cdots(n-i_\sigma))
 <
 ((n-i_1)+\cdots(n-i_\lambda)+\cdots(n-i_\sigma))
\]
が成り立つ。したがって \(P_0^*\) から \(Q_0,R_0\) に移ると、この差の和は減少している。この和は必ず \(\ge\sigma\) であるから、この手続きは各冪積 \(P_0^*\) について有限回で停止しなければならない。ゆえに
\[
        F(\xi_{ik})=\sum \xi_{i_1n}\cdots\xi_{i_\sigma n}P^{(\nu)}
        \equiv\sum \xi_{j_1n}\cdots\xi_{j_\sigma n}S_0^{(\nu)}
        =F_0(\xi_{ik})\pmod{\frN},
\]
ここで整係数多項式 \(F_0(\xi_{ik})\) は必ず正規化されている。そうでなければ差の和をさらに下げられるからである。これで (10) が証明された。\((n-1)\) 個の指数についての誘導仮定により、各々の別々に同次な \(F\) は別々に同次な正規化 \(F_0\) をもつので、いまの手続きはそれが \(n\) 個の指数についても成り立つことを示す。したがって以下では、別々に同次な多項式だけを考えればよい。

b) と同様に、
\begin{equation}
        F_0(\xi_{ik})=G(\xi_{ik})+\xi_{n-1,n}F_1(\xi_{ik}),\footnote{(10) と (11) に対応する恒等式
        \[
          F((ik))=F_0((ik))=G((ik))+(n-1,n)F_1((ik)),
        \]
        は、二つの行 \((n-1)\) と \((n)\) を変数の行 \(x,y\) で置き換え、したがって \((n-1,n)\) を \((xy)\) で置き換えるならば、Clebsch--Gordan 級数展開の類似物である。しかし Clebsch--Gordan が極形式による展開を含むのに対し、\(F_0\) の正規化は極形式の「初項」に対応し、それによって Clebsch--Gordan では成り立たない整性を強制する。整性を無視したとき \(f_{ik,rs}\) が \(\frM\) の基底をなすことの既知の証明は、まさに Clebsch--Gordan 級数展開によって行われる (Gordan-Kerschensteiner, \emph{Vorlesungen über Invariantentheorie}, p. 132)。}
\end{equation}
とおく。ここで \(G\) は \(\xi_{n-1,n}\) で割れる冪積を含まない。さらに
\begin{equation}
        [G]_{\xi_{in}=\xi_{i,n-1}}=K(\xi_{ik})
\end{equation}
とおく。正規化の定義が示すように、\(K\) も \(G\) と同じく正規化されている。

再び \(G\) と \(K\) の間には一対一対応がある。\(G\) の異なる冪積が \(K\) で同じ積に対応することがあり得るのは、それらが単に指数 \(n\) と \((n-1)\) の入れ替えによって異なる場合だけである。たとえば、\(G\) が指数 \((n-1)\) に関して次数 \(\varrho\)、指数 \(n\) に関して次数 \(\sigma\) であり、
\[
        \chi_1\le\chi_2,\ldots,\le\chi_\varrho\le\chi_{\varrho+1},\ldots,\le\chi_{\varrho+\sigma}
\]
が \(G\) の冪積の中で \((n-1)\) と \(n\) に結びつく指数であるとする。(11) により不定元 \(\xi_{n-1,n}\) は \(G\) に現れないので、その個数は \(\varrho+\sigma\) でなければならない。正規化により、その冪積はしたがって
\[
        \xi_{\chi_1,n-1}\cdots\xi_{\chi_\varrho,n-1}
        \xi_{\chi_{\varrho+1},n}\cdots\xi_{\chi_{\varrho+\sigma},n}\,P(\xi_{ik}),
\]
の形である。ここで \(P(\xi_{ik})\) はもはや指数 \(n\) と \((n-1)\) を含まない。\(n\) と \((n-1)\) を入れ替えると、\(i\ne j\) に対して因子 \(\xi_{in}\xi_{j,n-1}\)（\(i<j\)）が生じる。すると \(j,n-1\) は \(i\) と \(n\) の間にあり、その項は正規化されていない。したがって、固定された \(\chi_1,\chi_2,\ldots,\chi_{\varrho+\sigma}\) と固定された \(P\) に対して、\(G\) は一つの冪積しか含まない。これにより、\(G\) の異なる冪積は \(K\) で異なる積に対応することが証明された。したがって
\begin{equation}
        K(\xi_{ik})\equiv0\pmod p\quad\Longrightarrow\quad
        G(\xi_{ik})\equiv0\pmod p.
\end{equation}

\[
        F((ik))\equiv0\pmod p\quad \text{\(\alpha,\beta\) に関して恒等的に}
\]
から、(2) を考慮した (10) により、
\begin{equation}
        F_0((ik))\equiv0\pmod p\quad \text{\(\alpha,\beta\) に関して恒等的に}
\end{equation}
が従う。\(\alpha_n=\alpha_{n-1},\;\beta_n=\beta_{n-1}\) と特殊化すると、(11) と (12) は
\begin{equation}
        K((ik))\equiv0\pmod p\quad \text{\(\alpha,\beta\) に関して恒等的に}
\end{equation}
を与える。\(K(\xi_{ik})\) は正規化されており、\((n-1)\) 個の指数だけを含むので、a) および b) の式 (9) で証明したことから
\[
        K(\xi_{ik})\equiv0\pmod p\quad \text{\(\xi_{ik}\) に関して恒等的に}
\]
が従う。すると (13) により
\[
        G(\xi_{ik})\equiv0\pmod p\quad \text{\(\xi_{ik}\) に関して恒等的に}
\]
である。したがって (11) は
\[
        F_0(\xi_{ik})\equiv \xi_{n-1,n}F_1(\xi_{ik})\pmod p
        \quad \text{\(\xi_{ik}\) に関して恒等的に}
\]
を与える。\((n,n-1)\not\equiv0\pmod p\) は \(\alpha,\beta\) に関して恒等的には成り立たないので、(14) から
\[
        F_1((ik))\equiv0\pmod p\quad \text{\(\alpha,\beta\) に関して恒等的に}
\]
が従う。ここで \(F_1(\xi_{ik})\) は再び正規化され、\(\xi_{n-1,n}\) に関する次数は \(F_0\) より低い。この過程を有限回繰り返すと
\[
        F_0(\xi_{ik})\equiv\xi_{n-1,n}^{\lambda}F_\lambda(\xi_{ik})\pmod p,
\]
ただし \(F_\lambda\) は \(\xi_{n-1,n}\) に関して次数ゼロであり、したがっていま証明したことによって \(p\) を法として消える。こうして
\begin{equation}
\begin{array}{rcl}
        F_0(\xi_{ik})&\equiv&0\pmod p\quad \text{\(\xi_{ik}\) に関して恒等的に},\quad\text{または}\\
        F(\xi_{ik})&\equiv&0\pmod{(\frN,p)}\quad \text{\(\xi_{ik}\) に関して恒等的に}
\end{array}
\end{equation}
を得る。これは (2) の逆を、したがって (1) を証明する。同時に、\(n\) 個の指数の場合、したがって一般に、誘導段階で用いた \(F_0(\xi_{ik})\) に対するより鋭い主張も証明している。

\setcounter{equation}{0}
\section*{§ 4. 有限群の整係数不変式の有限性}

私のノート「有限群の不変式に対する有限性定理」\footnote{\emph{Math. Ann.} 77, p. 89 (1915).}で与えた有限性の初等的証明は、§ 1 の特殊有限性定理 III と IV によって整性へと強めることができる。一方、Hilbert の加群基底定理に基づく通常の証明はこの整性を与えない。整加群基底を用いても、群の位数である整数 \(h\) の任意に高い冪が分母に現れる表示しか得られない。

以下では、いま引用したノートの記法を用いる。有限群 \(\frH\) は、非零行列式をもつ \(h\) 個の線形変換 \(A_1,\ldots,A_h\) からなるとする。\footnote{\(\frH\) が抽象的な有限群に単純に同型であると仮定するだけで十分であろう。すなわち、\(A\) のある変換行列式が消えるなら、\(\frH\) は群 \(\begin{pmatrix}\frS&0\\0&0\end{pmatrix}\) に相似であるべきであり、ここで \(\frS\) の置換の行列式は消えない。この場合、\(\frH\) の不変式は \(\frS\) の不変式と一致するので、行列式が消えないという条件は実際には一般性を制限しない。}ここで \(A_k\) は線形変換
\[
        x_1^{(k)}=\sum_{\nu=1}^{n}a_{1\nu}^{(k)}x_\nu,
        \quad\ldots,\quad
        x_n^{(k)}=\sum_{\nu=1}^{n}a_{n\nu}^{(k)}x_\nu,
\]
あるいは簡潔に \((x^{(k)})=A_k(x)\) を表す。したがって群 \(\frH\) は、要素 \(x_1,\ldots,x_n\) をもつ行 \((x)\) を、要素 \(x_1^{(k)},\ldots,x_n^{(k)}\) をもつ行 \((x^{(k)})\) へ送る。係数 \(a_{i\nu}^{(k)}\) は代数的整数であると仮定する。これは一般性を制限しない。J. Schur の定理\footnote{Über Gruppen linearer Substitutionen mit Koeffizienten aus einem algebraischen Zahlkörper. Satz VII. \emph{Math. Ann.} 71, p. 555 (1912).}により、任意の線形変換の有限群は、\(a_{i\nu}^{(k)}\) が代数的整数となるものに相似であるからである。

群の整絶対不変式とは、\(a_{i\nu}^{(k)}\) によって生成される数体 \(\frK\) の代数的整数を係数とする、\(x_1,\ldots,x_n\) の多項式的有理関数であって、\(A_1,\ldots,A_h\) の適用のもとで恒等的に変わらないものをいう。そのような不変式 \(f(x)\) に対しては
\begin{equation}
        f(x)=f(x^{(1)})=\cdots=f(x^{(h)})
        =\frac1h\sum_{k=1}^{h}f(x^{(k)}).
\end{equation}
逆に、量の \(h\) 個の行 \((x^{(k)})\) の任意の整係数対称関数も、群の整係数不変式になる。

\(\omega_1,\omega_2,\ldots,\omega_\varrho\) を \(\frK\) のすべての整数の基底とする。すると任意の不変式は
\[
        f(x)=\omega_1 f_1(x)+\cdots+\omega_\varrho f_\varrho(x),
\]
と表せる。ここで \(f_1(x),\ldots,f_\varrho(x)\) は有理整数係数をもつ。したがって (1) により
\[
        h.f(x)
        =
        \omega_1\sum_{k=1}^{h}f_1(x^{(k)})
        +\cdots+
        \omega_\varrho\sum_{k=1}^{h}f_\varrho(x^{(k)}).
\]
ここで個々の和は、量の \(h\) 個の行 \((x^{(k)})\) の対称関数であり、しかも有理整数係数をもつ。したがって定理 III により、それらは、これらの行の有限個の整有理対称関数を用いて、有理整数係数をもつ多項式的かつ有理的な形で表される。前の議論により、これらの基底関数 \(I_1(x),\ldots,I_\sigma(x)\) は、一般には代数的整数係数をもつ群の整係数不変式になる。というのも、\(\frH\) は \(A\) とともにその代数的共役変換のすべてを含むとは限らないからである。

したがって表示
\begin{equation}
        h\cdot f(x)=\omega_1G_1(I)+\cdots+\omega_\varrho G_\varrho(I)
\end{equation}
がある。ここで \(G_1,\ldots,G_\varrho\) は有理整数係数をもつ。

いま \(w_1,\ldots,w_\varrho\) を不定元とし、
\[
        L=w_1G_1(I)+\cdots+w_\varrho G_\varrho(I)
\]
に、置換 \(w_i=\omega_i/h\) によって表示 (2) の不変式 \(f(x)\) に対応するすべての \(w\) の線形形式を動かせる。Hilbert の整加群基底に関する定理は
\begin{equation}
        L=A_1(I)L_1+\cdots+A_\tau(I)L_\tau,
\end{equation}
を与える。ここで \(A_i\) は有理整数係数をもち、またすべての \(L\) は \(w\) に関して線形なので、\(A_i\) はもはや \(w\) を含まない。置換 \(\omega_i=w_i/h\) が \(L_1,\ldots,L_\tau\) を不変式 \(\mathfrak I_1(x),\ldots,\mathfrak I_\tau(x)\) に送るなら、(3) は
\begin{equation}
        f(x)=A_1(I)\mathfrak I_1(x)+\cdots+A_\tau(I)\mathfrak I_\tau(x)
\end{equation}
となる。この表示は
\[
        I_1(x),\ldots,I_\sigma(x);\mathfrak I_1(x),\ldots,\mathfrak I_\tau(x)
\]
が群 \(\frH\) のすべての整係数不変式の基底をなし、各不変式がこれら有限個の不変式の、有理整数係数をもつ多項式的有理関数として表されることを示している。

とくに、群 \(\frH\) が、各変換 \(A\) とともに代数的共役な数係数をもつすべての \(A',A'',\ldots\) も含むという性質をもつなら、\(h\) 個の行 \((x^{(k)})\) の整有理対称関数は、有理整数係数をもつ \(x\) の関数になる。これは、とくに対称関数の基底関数についても成り立つ。したがって、有理整数係数をもつ不変式 \(f(x)\) に制限するなら、基底関数
\[
        I_1(x),\ldots,I_\sigma(x);\mathfrak I_1(x),\ldots,\mathfrak I_\tau(x)
\]
も有理整数係数をもつ。この特殊な仮定は一般の群 \(\frH\) では満たされないが、その仮定の下では、整係数有限性は有理整数係数をもつ不変式の部分領域についても成り立つ。





%
\clearpage
\begin{center}
{\Large\bfseries 16. 形式論における級数展開について}\par
\emph{（\emph{Math. Ann.} 77, p. 93 掲載の第 III 項への補遺および訂正。）}\\[0.4em]
Emmy Noether（Göttingen）\\[0.5em]
\emph{Math. Ann.} 81 (1920), pp. 25--30
\end{center}

「形式論の級数展開」とは、周知のように、一方では Clebsch--Gordan の級数展開、すなわち二つの二元的な行を含む形式をそれらの行の行列式の冪によって展開し、その係数が一つの変数行だけに関する形式の極化形である、また同じことだが \(\Omega\) 過程を適用すると消える、という展開を、各々 \(n\) 個の変数からなる任意個の行の形式へ一般化することを意味する。他方では、「複素座標」\(t_{ik}\) で現れる正規形の一般化として、異なる段階の量 \(t_i,t_{ik},t_{ikl},\ldots\) を同時に含む形式に対し、正規形による展開を与えることを意味する。しかしこの第二の種類の展開は、Mertens と Deruyts がとくに遂行したように、不変な微分過程によって第一の種類から得ることができる。\footnote{F. Mertens, Über invariante Gebilde quaternärer Formen (Sitzber. d. Ak. d. Wiss. Wien 98, Abt. IIa (1889))、四元の場合。一般の正規形は J. Deruyts: Sur la réduction des fonctions invariantes dans le système des variables géométriques, Bulletins de l'Ac. R. de Belgique 24 (1892) にある。この論文を知らずに、Mertens に厳密に従って、私は論文 Zur Invariantentheorie der Formen von \(n\) Variabeln, J. f. M. 139 (1910) の § 7 で正規形による一般展開を与えた。}

第一種のさまざまな展開はすべて、たとえば \emph{Math. Ann.} 77, loc. cit. III で簡潔に素描されているように、各々 \(n\) 個の変数からなる \(n\) 個の行の形式を、それらの行の行列式の冪によって展開することから従う。このとき係数は、\((n-1)\) 個の行だけに関する形式の極化形であり、それ自身は出発形式から極化形成と \(\Omega\) 過程によって導かれる。\footnote{loc. cit. の文献指示を参照。その後、やはり形式的方法に基づくものとして、Schouten, Over reeksontwikkelingen van algebraische vormen met verschillende rijen van variabelen van verschillenden graad, Verslag Ak. Amsterdam 27 (1919), p. 1481 が現れた。} 元来の Clebsch--Gordan 展開 \((n=2)\) はさらに一歩進んでいた。すなわち、そこに現れる特殊な極化過程は、数係数を含めて明示的に与えられている。

そこで、級数展開が線形形式族の問題として、より概念的に捉えられていた上記のノートに私が加えたいのは、E. Fischer による一般加群系に関する研究を特殊化して得られる、数係数を除いて明示的な表示である。\footnote{E. Fischer, Über die Differentiationsprozesse der Algebra, J. f. M. 148 (1917).}

この位置づけへ導くのは、\(t_{ik}\) における正規形の解釈である。すなわち「複素座標」\(t_{ik}\) を「直線座標」
\[
        p_{ik}=(x_i y_k-x_k y_i)
\]
で置き換えると、\(p_{ik}\) の間に存在する恒等関係のため、形式 \(F(p_{ik})\) の表示はもはや一意でない。\(F\) の係数の間に、\(F\) に現れる \(p_{ik}\) の冪積の間と同じ線形関係が成り立つことを要求すれば、一意性を得ることができる。\footnote{Clebsch, Über eine Fundamentalaufgabe der Invariantentheorie, § 4, Abh. d. K. G. d. Wissensch. zu Göttingen 17 (1872) を参照。} したがってこの正規形 \(\Phi\) については
\begin{equation}
\begin{aligned}
        F(p_{ik})&=\Phi(p_{ik}) &&[\text{\(x,y\) に関して恒等的}],\quad\text{または}\notag\\
        F(t_{ik})&\equiv\Phi(t_{ik}) &&\modu{M},
\end{aligned}
\tag{1}
\end{equation}
が成り立つ。ここで \(M\) は \(p_{ik}\) の間のすべての恒等関係からなる加群、すなわち \(G(p_{ik})\) が \(x,y\) に関して恒等的に消えるすべての形式 \(G(t_{ik})\) の全体を意味する。\(\Phi(t_{ik})\) は複素座標における「正規形」であり、一方で \(M\) の基底関数は \(t_{ik}\) の二次形式になる。これらの基底関数の因子に関して手続きを反復すると、完全な級数展開が生じる。\(t_i,t_{ik},t_{ikl},\ldots\) が同時に現れる場合にも、対応する考察が成り立つ。

極化形による展開も、同様の考え方に帰着される。たとえば形式 \(F(x,y,z)\) において、三つの三元的な行 \(x=(x_1,x_2,x_3),y,z\) を、二つだけから線形に合成される行 \(\xi,\eta,\zeta\)、たとえば \(\xi=x,\eta=y;\zeta=\lambda_1x+\lambda_2y\) で置き換えると、\(F(\xi,\eta,\zeta)\) の表示は再び一意でなくなる。ここでも、係数に対して、\(F\) に現れる \(\xi,\eta,\zeta\) の冪積に対するのと同じ線形関係を指定することができる。この場合、すべての恒等関係からなる加群 \(M\) は、三つの行の行列式 \(\Delta=(xyz)\) を基底関数としてもつので（loc. cit. III, 補助補題 a)）、[1] の代わりに
\begin{equation}
\begin{aligned}
        F(\xi,\eta,\zeta)&=\Phi(\xi,\eta,\zeta)&&[\text{\(x,y\) に関して恒等的}],\quad\text{または}\notag\\
        F(x,y,z)&\equiv\Phi(x,y,z)&&\modu{\Delta}
\end{aligned}
\tag{2}
\end{equation}
が現れる。\(\Phi\) は極化形成によって生じることが分かる（loc. cit. の式 (2) を参照）。完全な級数展開も、やはり反復によって得られる。

したがって正規形 \(\Phi\) の形成は、どちらの場合にも次の問題に従属する。\(F(z)\) において不定元 \(z_i\) をパラメータ \(\xi\) の多項式 \(f_i(\xi)\) で置き換え、そのため表示 \(F(f_i(\xi))\) がもはや一意でないとき、一意性は、係数の間に \(F\) に現れる \(f_i(\xi)\) の冪積の間と同じ線形関係が成り立つようにして固定されるべきである。したがって
\begin{equation}
\begin{aligned}
        F(f_i(\xi))&=\Phi(f_i(\xi))&&[\text{\(\xi\) に関して恒等的}],\quad\text{または}\notag\\
        F(z)&\equiv\Phi(z)&&\modu{M},
\end{aligned}
\tag{3}
\end{equation}
となる。ここで \(M\) は \(f_i(\xi)\) の間のすべての恒等関係からなる加群を意味する。

この正規形 \(\Phi(z)\) について、E. Fischer は引用した論文の § 11, I において、数係数を除いて明示的な展開を、しかも線形作用によって与えた。同時に彼は、導入 III および § 6 II において、\(M\) に属する差 \(F-\Phi\) に対しても、加群基底を知っているという仮定の下で、同じ意味で明示的な式を与えた。そして § 11 II でこれらの展開を結合することにより、式 [3] を明示的表示で置き換えた。ここでなお未解決の数係数の問題は、同じ論文の § 12 によれば、線形作用の固有値と固有関数を探す問題（式 82）と同一である。

私は以下の 1. で、Fischer の公式を極化形による級数展開の場合に特殊化し、そこから反復によって級数展開を得る。\(\Phi\) を決定するための作用は、このとき特定の極化過程に移るので、上で述べた従来知られていた定式化よりかなり精密である。\(F-\Phi\) を決定するための作用には、行列式 \(\Delta\) による乗法と \(\Omega\) 過程が混合して現れる。通常の非明示的な形に至るには、さらに \(\Omega\Delta\) が極化過程によって表せることに注意しなければならない。

ここで一つ訂正も加えておく。私は上掲箇所（p. 101, 上から6行目）で、二元領域でだけ成り立つ恒等変形
\[
        \Omega(\Delta\psi)=c_1\psi+c_2\Delta\cdot\Omega(\psi)
\]
を誤って一般に有効なものと仮定し、それによって、[2] における正規形 \(\Phi\) の式を与える基本公式 (2) から級数展開 (3) へ移行した。したがってそこでの考察が与えるのは式 (2) だけであり、その特殊場合として還元定理 (1) だけである。しかし差 \(F-\Phi\) の明示的な式が欠けているため、級数展開の根拠づけには足りない。

2. でも同様に特殊化によって [1] の正規形 \(\Phi(t_{ik})\) を得る。より一般には正規形 \(\Phi(t_i,t_{ik},t_{ikl},\ldots)\) を得て、これは既知の明示的公式へ導く。しかし完全な級数展開を立てるには \(M\) の加群基底を知る必要があるので、ここでは不変式論的な道筋（引用した文献を参照）の方が望ましい。この道筋はその知識を要求せず、逆に級数展開によって加群基底を与えるからである。

1. Fischer の § 11（式 (19) および続く行）によれば、[3] の正規形 \(\Phi(z)\) は、そこでは \(N(\zeta)\) と書かれており、次で与えられる。
\begin{equation}
        \Phi(z)=(a_1S+a_2S^2+\cdots+a_rS^r)F(z),
\tag{4}
\end{equation}
ここで定数 \(a_1\ldots a_r\) は \(f(\xi)\) の係数領域に属し、\(F(z)\) の次数だけに依存する。作用素 \(S\) は
\begin{equation}
\begin{gathered}
        S=AB;\qquad BF(z)=F(f(\xi));\\
        AF(f(\xi))=\left\{\sum z_i f_i\!\left(\frac{\p}{\p\xi}\right)\right\}^{p}F(f(\xi))\footnotemark
\end{gathered}
\tag{5}
\end{equation}
で定義される。\footnotetext{式 (75) を参照。[3] の場合には、和 \(\alpha\) は \(p\) 次元のすべての冪積を動くので、(75) は [5] に特殊化される。}
[2] において \(F\) を \(z\) だけに関する \(p\) 次元の形式とみなすと、\(f_i(\xi)\) は \(\lambda_1x_i+\lambda_2y_i\) に特殊化される。したがって
\[
\begin{gathered}
        BF=F(x,y,\lambda_1x+\lambda_2y),\\
        SF=\left\{\sum z_i\left(\frac{\p}{\p\lambda_1}\frac{\p}{\p x_i}+
        \frac{\p}{\p\lambda_2}\frac{\p}{\p y_i}\right)\right\}^{p}
        F(x,y,\lambda_1x+\lambda_2y),
\end{gathered}
\]
また、極化過程で書けば（loc. cit. II, 2 および 3、また略記法についても参照）、
\begin{equation}
        SF(x,y,z)=\frac1{p!}\left[z\left(\frac{\p}{\p\lambda_1}\frac{\p}{\p x}+
        \frac{\p}{\p\lambda_2}\frac{\p}{\p y}\right)\right]^p
        \left[(\lambda_1x+\lambda_2y)\frac{\p}{\p z}\right]^p F(x,y,z)
\tag{6}
\end{equation}
となる。\emph{[4] と [6] によって、} \(\Phi(x,y,z)\) \emph{に対する求める明示的表示が与えられる。}同時に \(\Phi(x,y,z)\) は、冒頭で述べた一般形（loc. cit. の (2) における \(\sum PH\)）、すなわち、2 個の行だけに依存し \(F\) から極化過程によって得られる式の極化形、に従属する。まだ決まっていない係数 \(a_i\) は有理数である。ここでは \(f_i(\xi)\) が有理整数の数係数をもつからである。[2] を
\begin{equation}
        F=\Phi+\Delta F_1
\tag{2a}
\end{equation}
の形で書けば、Fischer が導入 III と § 6 II（p. 41）で与えた公式の特殊化により、
\begin{equation}
        F_1=(\alpha_1\Omega+\alpha_2\Omega\Delta\Omega+\cdots+\alpha_\lambda\Omega\Delta\Omega\cdots\Delta\Omega)F\footnotemark,
        \qquad \left[\Omega=\left(\frac{\p}{\p x}\frac{\p}{\p y}\frac{\p}{\p z}\right)\right],
\tag{7}
\end{equation}
が得られる。\footnotetext{この特殊化された公式を Kerim は、近く刊行される Erlangen の学位論文で、「慣性形式」への応用に用いている。}
ここで \(\alpha_i\) もまた、\(F\) の次数だけに依存する有理数である。

反復は次を導く。
\[
        F_1=\Phi_1+\Delta F_2;\quad F_2=\Phi_2+\Delta F_3;\quad\ldots;\quad
        F_i=\Phi_i+\Delta F_{i+1};\quad\ldots,
\]
ここで各 \(\Phi_i\) は \(F_i\) に対応する正規形である。\(F_i\) は \(z\) に関して次数 \((p-i)\) をもつので、[4] と [6] から
\begin{equation}
\begin{gathered}
        \Phi_i=(a_{i1}S_i+a_{i2}S_i^2+\cdots)F_i,\\
        S_iF_i=\frac1{(p-i)!}
        \left[z\left(\frac{\p}{\p\lambda_1}\frac{\p}{\p x}+
        \frac{\p}{\p\lambda_2}\frac{\p}{\p y}\right)\right]^{p-i}
        \left[(\lambda_1x+\lambda_2y)\frac{\p}{\p z}\right]^{p-i}F_i.
\end{gathered}
\tag{8}
\end{equation}
式 [7] には
\[
        F_i=(\alpha_{i1}\Omega+\alpha_{i2}\Omega\Delta\Omega+\cdots)F_{i-1},
        \quad\text{および反復により}
\]
\begin{equation}
        F_i=(\beta_1\Omega^i+\beta_2\Omega^i\Delta\Omega+\cdots
        +\beta_\rho\Omega^{k_1}\Delta\Omega^{k_2}\cdots\Delta\Omega^{k_\lambda}+\cdots)F,
        \quad(k_1+k_2+\cdots+k_\lambda-\lambda+1=i)
\tag{9}
\end{equation}
が対応する。\emph{[8] と [9] によって、級数展開}
\[
        F=\Phi+\Delta\Phi_1+\Delta^2\Phi_2+\cdots+\Delta^i\Phi_i+\cdots\Delta^p\Phi_p
\]
\emph{において、述べた意味で明示的な表示が達成される。}

[7] に恒等変形 \(\Omega\Delta F=PF\) を適用する。ただし \(P\) は極化過程を表す。\footnote{たとえば F. Mertens, Über ganze Funktionen von \(m\) Systemen von je \(n\) Unbestimmten, Monatsh. f. Math. u. Phys. 4, 式 (5) を参照。}すると、これらの過程が \(\Omega\) 過程と可換であるため、既知の公式
\[
        F_1=P\Omega F;\quad\ldots;\quad F_i=P\Omega^iF;
\]
が得られる。ここから [8] と結びつけて、または loc. cit. の考察と結びつけて、たとえば loc. cit. の (3) にある既知の級数展開の形が従う。\(n\) 個の変数からなる \(n\) 個の行 \(x^{(1)},x^{(2)},\ldots,x^{(n-1)},z\) を扱う場合には、それに応じて
\[
        f_i(\xi)=\lambda_1x_i^{(1)}+\lambda_2x_i^{(2)}+\cdots+\lambda_{n-1}x_i^{(n-1)};
\]
となり、同様に \(\Delta\) と \(\Omega\) も \(n\) 個の行に対して定義される。

2. \(F(t_i,t_{ik},t_{ikl},\ldots)\) の正規形については、\(f_i(\xi),f_{ik}(\xi),\ldots\) は、一行、二行、さらに \((n-1)\) 行までの行列式
\[
        \xi_i^{(1)},\quad (\xi^{(1)}\xi^{(2)})_{ik},\quad\ldots,\quad
        (\xi^{(1)}\xi^{(2)}\cdots\xi^{(n-1)})_{i_1i_2\cdots i_{n-1}}
\]
によって与えられる。したがって [5] の特殊化として
\[
\begin{gathered}
        S=AB;\qquad
        BF=F\bigl(\xi_i^{(1)},(\xi^{(1)}\xi^{(2)})_{ik},\ldots\bigr),\\
        ABF=\left(\sum t_i\frac{\p}{\p\xi_i^{(1)}}\right)^{\lambda_1}
        \left(\sum t_{ik}\left(\frac{\p}{\p\xi^{(1)}}\frac{\p}{\p\xi^{(2)}}\right)_{ik}\right)^{\lambda_2}
        \cdots F\bigl(\xi_i^{(1)},(\xi^{(1)}\xi^{(2)})_{ik},\ldots\bigr),
\end{gathered}
\]
を得る。ここで \(\lambda_1,\lambda_2,\ldots\) は \(t_i,t_{ik},\ldots\) に関する次数を意味する。

不変式論的な考察から、ここでは次の合同が従う。すなわち、\(F\) の不変式で係数に関して線形なものはすべて、\(\Phi\) と \(M\) の形式から線形に合成できるので、とくに \(M\) に属さない \(SF\) についてもそうである。\footnote{これは、たとえば私の引用論文 ``Zur Invariantentheorie der Formen von \(n\) Variabeln'' の § 6 および § 7, 2. から従う。}したがって \(\Phi\equiv a_1SF\modu{M}\) であり、正規形の性質により、[4] の代わりに単純な関係
\begin{equation}
        \Phi=a_1SF=a_1S\Phi
\tag{10}
\end{equation}
が得られる。第二の等式が成り立つのは、\(S\) を適用すると \(M\) の任意の形式が消えるからである。式 [10] は、正規形 \(\Phi\) 自身が同時に固有関数であることを示している。しかも、上の不変式論的考察により、同時に \(F\) の不変式でもある他の固有関数は存在しない（[10] については Deruyts の式 (10) および (10') を参照）。

\begin{center}
\textbf{訂正。}
\end{center}

1. 論文 ``Die allgemeinsten Bereiche aus ganzen transzendenten Zahlen''（\emph{Math. Ann.} 77）では、p. 121 の上から8行目、上から9行目、および下から11行目で、体 \((H,Y)\) を「体 \((H,Y,Y',\ldots,Y^{(\lambda_0-1)})\)、ここで \(Y',\ldots,Y^{(\lambda_0-1)}\) は \(H\) に関する \(Y\) の共役を意味する」で置き換えなければならない。これに応じて、p. 120 の上から5行目、および p. 121 の上から9行目の \((H,y)\) は \((H,y,y',\ldots,y^{(\lambda_0-1)})\) で置き換えられる。

2. 論文 ``Gleichungen mit vorgeschriebener Gruppe''（\emph{Math. Ann.} 78）では、Hilbert の既約性定理の定式化、p. 222 の下から3行目において、「数体 \(\Omega\)」ではなく、より正確に「有限代数的数体 \(\Omega\)」と言わなければならない。この点を Ostrowski 氏が私に注意してくださった。実際、たとえばすべての代数的数の体に関しては、代数的数係数をもつ任意の方程式の群は恒等群である。したがって導入部で「任意の有理性領域」と言うときには、なお有限個のパラメータを添加できる有限代数的数体、またはそのようにして生じる中間体の一つを意味すると理解しなければならない。

\begin{center}
（1919年9月受理。）
\end{center}


% 組版用マクロ
\providecommand{\p}{\partial}
\providecommand{\dd}{\mathrm d}
\providecommand{\modu}[1]{\;\operatorname{mod} #1}
\providecommand{\mM}{\mathfrak{M}}
\providecommand{\mN}{\mathfrak{N}}
\providecommand{\mL}{\mathfrak{L}}
\providecommand{\mG}{\mathfrak{G}}
\providecommand{\mA}{\mathfrak{A}}
\providecommand{\mB}{\mathfrak{B}}
\providecommand{\mX}{\mathfrak{X}}
\providecommand{\mY}{\mathfrak{Y}}
\providecommand{\mU}{\mathfrak{U}}
\providecommand{\mP}{\mathfrak{P}}
\providecommand{\mQ}{\mathfrak{Q}}
\providecommand{\mS}{\mathfrak{S}}
\providecommand{\mD}{\mathfrak{D}}
\providecommand{\mR}{\mathfrak{R}}
\providecommand{\mV}{\mathfrak{V}}
\providecommand{\ma}{\mathfrak{a}}
\providecommand{\mb}{\mathfrak{b}}
\providecommand{\mx}{\mathfrak{x}}
\providecommand{\my}{\mathfrak{y}}
\providecommand{\mz}{\mathfrak{z}}
\providecommand{\me}{\mathfrak{e}}
\providecommand{\mbarA}{\overline{\mathfrak{A}}}
\providecommand{\mbarB}{\overline{\mathfrak{B}}}
\providecommand{\bara}{\overline{\mathfrak{a}}}

\clearpage
\setcounter{footnote}{0}
\begin{center}
{\Large\bfseries 17. W. Schmeidler との共著：非可換領域における加群、特に微分式および差分式から生じるもの}\par
\vspace{1em}
\emph{Math. Zs.} 8 (1920), pp. 1--35
\end{center}

\begin{center}
\textbf{目次.}
\end{center}
\begin{center}
\begin{minipage}{0.88\linewidth}
序論.\\
\S~1. 微分式および差分式に関する形式的事項.\\
\S~2. 一般の多項式領域.\\
\S~3. 片側加群とその剰余群.\\
\S~4. 剰余群の二つの部分群への可約性と、加群との関係.\\
\S~5. 単位による計算と $k$ 個の部分群への可約性.\\
\S~6. 有限性定理.\\
\S~7. 分解が一意である場合.\\
\S~8. 完全可約群の素群への加法分解. 同型定理.\\
\S~9. 「同型」と「同種」の概念の関係.\\
\S~10. 群の無限に多くの分解の存在.\\
\S~11. 微分方程式系およびその積分との関係.\\
\S~12. 例.
\end{minipage}
\end{center}

\begin{center}\textbf{序論.}\end{center}

一変数の微分式の形式的理論は、ある意味で一意な既約因子への分解に導く。\footnote{E. Landau, Ein Satz über die Zerlegung homogener linearer Differentialausdrücke in irreduzible Faktoren, Journal für die reine und angewandte Mathematik 124 (1902), pp. 115--120、およびそれに続く A. Loewy, Über reduzible lineare homogene Differentialgleichungen, Mathematische Annalen 56 (1903), pp. 549--584 を参照。} しかし代数学における対応する定理は多変数多項式へ移すことができるのに対して、偏微分式についてはもはやそうはいかない。\footnote{H. Blumberg, Über algebraische Eigenschaften von linearen homogenen Differentialausdrücken, 学位論文, Göttingen 1912, 52 pp., pp. 51--52 に発表された Landau の反例を参照。} ただし一変数の微分式には、積表示のほかに最小公倍元としての表示もあり、同様の法則が成り立つ。二つの異なる分解において、個々の構成因子は互いに一対一に対応させることができる。すなわちそれらは「同種」である。\footnote{上掲の論文のほか、A. Loewy, Über vollständig reduzible lineare homogene Differentialgleichungen, Mathematische Annalen 62 (1906), pp. 89--117 を参照。以下 Loewy と引用する。}

この最小公倍元の概念は、いま加群概念として解釈でき、したがって $n$ 変数へ移すことができる。本論文ではこれが、いま述べた定理の自然な一般化を与える。さらに、微分式について用いられる事実は次のことだけである。すなわち、それらは非可換な乗法をもつ多項式として解釈でき、その場合、積の次数は因子の次数の和に等しい。この条件は、たとえば差分式についても満たされる。したがって、これらの定理は差分式にも成り立つが、そこでも従来はやはり一変数の場合だけが知られていた。\footnote{G. Wallenberg--A. Guldberg, Theorie der linearen Differenzengleichungen, Leipzig and Berlin 1911 を参照。}

そこで我々は一般に、非可換乗法をもつ多項式を元とする加群の理論を展開する。その際、加群の研究を剰余類の研究へ帰着させる。\footnote{可換領域の場合については W. Schmeidler, Über Moduln und Gruppen hyperkomplexer Größen, Mathematische Zeitschrift 3 (1919), pp. 29--42 を参照。} 可換の場合と違って、ここでは剰余類の和については語れるが、その積については語れない。語れるのは剰余類と多項式との積だけである。したがって「剰余群」の概念は可換の場合から修正されなければならない。一変数多項式という特殊な場合には、剰余群はさらに「有限」である。つまり線形独立な剰余類を有限個しかもたない。一般にはそのような有限な「位数」はなく、群は「無限」である。しかし我々の方法にとってこの相違は本質的ではない。

可換の場合と同じく、互いに素な加群の最小公倍元として加群を表示することに対応して、剰余群の加法分解というより単純な手続きがある。この関連で剰余群の同型性が本質的になる。一変数の微分式へ特殊化すると、この同型性は対応する加群に対する同種の概念になり、可換加群へ特殊化すると同一性になる。なお、同種の概念も $n$ 変数へ移すことができ、その概念が同型概念と一致することを証明できる。剰余群の加法分解において常に有限個の和因子だけが現れるという事実は、Gordan の証明法を用いて Hilbert の加群基底定理を移すことにより従う。Hilbert の元の証明法では、乗法法則に制限が必要になる（式 (6) で右側に $a_i$ ではなく $a$ が現れる場合であり、これは微分式では満たされるが差分式では満たされない）。

一意性定理へ到達するためには、Loewy と同様に、加群を「完全可約」なもの、すなわち素加群の最小公倍元として表示できるものに特殊化しなければならない。二つの異なる分解が与えられると、その構成因子は一対一に同一であるか、少なくとも一方の分解の各構成因子に対して他方の分解の同型な構成因子が対応し、また逆も成り立つ。同時に、そのとき各分解は少なくとも二つの同型な構成因子を含む。さらに、同じ分解の中に二つの同型な構成因子が現れることから、無限に多くの分解の存在が従う。これは素加群への制限なしにも成り立つ事実であり、従来は一変数の微分式という特殊な場合でも知られていなかった。

最後に、微分式から来る加群について、積分との関係を論じ、有限な剰余群については、その位数が線形独立な積分の最大個数に等しいことを示す。したがって一般積分に任意関数が現れるのは、ちょうど剰余群が有限でないときである。二つの加群が同種であるという概念は、ここでも一変数の場合と同様に、積分の間の「双有理的」関係を意味する。最後にいくつかの例を掲げる。

本論文への最初の刺激は、数年前に E. Landau が積定理の一般化について E. Noether に尋ねたことから来た。当時 E. Noether は、この問題が非可換多項式領域における加群論の問題であることだけを指摘したが、まだ取り組むことはできなかった。当時知られていた Lasker の加群定理は、この領域へ直接移すことができなかった。なぜなら Lasker の証明法は、多項式が既約因子へ一意に分解されるという定理に依存しているからである。\footnote{最近 E. Noether は、ここで用いる方法に近い方法により、この定理を使わずに Lasker の定理も確立できることに気づいた。} 移行の可能性を与えたのは Schmeidler の方法だけであり、それを我々は共同で実行した。細部について述べれば、剰余群の一般化、Hilbert の有限性定理の移行と応用、および無限に多くの分解の存在は E. Noether による。Loewy の完全可約性概念の一般化、同型概念、およびそれと同種概念との関係は W. Schmeidler による。

\subsection*{\S~1. 微分式および差分式に関する形式的事項.}

1. 一変数の微分式について、記号的な積形成が導入されている。いま
\[
 a_0(x)\frac{\dd^n y}{\dd x^n}+a_1(x)\frac{\dd^{n-1}y}{\dd x^{n-1}}+\cdots+a_n(x)y=A(y),
\]
\[
 b_0(x)\frac{\dd^m y}{\dd x^m}+b_1(x)\frac{\dd^{m-1}y}{\dd x^{m-1}}+\cdots+b_m(x)y=B(y)
\]
とする。このとき記号的積 $AB(y)$ を微分式
\[
 a_0(x)\frac{\dd^n B(y)}{\dd x^n}+a_1(x)\frac{\dd^{n-1}B(y)}{\dd x^{n-1}}+\cdots+a_n(x)B(y)
\]
として定義する。この積形成は実際には二つの作用素の乗法である。通常どおり $\frac{\dd^n}{\dd x^n}$ を $\left(\frac{\dd}{\dd x}\right)^n$ で置き換えると、最も簡単に書ける。すると $AB$ は
\[
 \left(a_0\left(\frac{\dd}{\dd x}\right)^n+\cdots+a_n\right)
 \left(b_0\left(\frac{\dd}{\dd x}\right)^m+\cdots+b_m\right)y
\]
となり、積を実行する際には和と積の微分に関する普通の規則を適用する。このような作用素はすべて、作用素が適用される関数を明示的に書かないことにし、変数 $\xi$ に積の微分に対応する乗法規則を定めることにより、一つの変数 $\frac{\dd}{\dd x}=\xi$ の多項式とみなせる。
\[
 \frac{\dd}{\dd x}\bigl(a(x)\cdot y\bigr)=a(x)\frac{\dd}{\dd x}(y)+a'(x)y,
\]
すなわち我々の記法では $(n=1,m=0)$、
\begin{equation}
 \xi\cdot a=a\cdot\xi+a'.
\tag{1}
\end{equation}

関数 $y$ の線形偏微分式
\[
 \left(\sum a_{\alpha_1\ldots\alpha_n}(x_1,\ldots,x_n)
 \frac{\p^{\alpha_1+\cdots+\alpha_n}}
 {\p x_1^{\alpha_1}\cdots\p x_n^{\alpha_n}}\right)y
\]
についてもまったく同じように進む。$n$ 個の作用素 $\frac{\p}{\p x_1},\ldots,\frac{\p}{\p x_n}$ に対して変数 $\xi_1,
\ldots,\xi_n$ を導入し、これにより多項式
\[
 \sum a_{\alpha_1\ldots\alpha_n}(x_1,\ldots,x_n)
 \xi_1^{\alpha_1}\cdots\xi_n^{\alpha_n}
\]
を得る。二つの多項式の積は次の法則に従って計算される。
\begin{equation}
 \xi_i a=a\xi_i+a_i \qquad (i=1,\ldots,n),
\tag{2}
\end{equation}
ここで $a_i$ は $x_i$ に関する $a$ の偏導関数を表す。$\xi_i$ 同士の積は可換である。この乗法規則により、任意の積は $\xi_i$ の多項式へ変形できる。

2. 差分式（Wallenberg, loc. cit. 参照）
\[
 a_x^{(0)}y_{x+n}+a_x^{(1)}y_{x+n-1}+\cdots+a_x^{(n)}y_x=A_x,
\]
\[
 b_x^{(0)}y_{x+m}+b_x^{(1)}y_{x+m-1}+\cdots+b_x^{(m)}y_x=B_x
\]
については、対応する積を
\[
 (AB)_x=a_x^{(0)}B_{x+n}+a_x^{(1)}B_{x+n-1}+\cdots+a_x^{(n)}B_x
\]
と定義し、乗法規則
\[
 (ay)_{x+1}=a_{x+1}y_{x+1}
\]
によって計算する。$x$ から $x+1$ への移行を表す作用素 $\xi$ を導入すると、$\xi\{ay\}=a^*\xi\{y\}$ を得る。ここで $a_x^*=a_{x+1}$ である。ここでも作用素が適用される関数 $y$ を省略でき、法則
\begin{equation}
 \xi a=a^*\xi
\tag{3}
\end{equation}
を得る。このとき $a$ が恒等的に消えなければ $a^*$ も恒等的には消えない。まったく同じように、$n$ 変数については
\begin{equation}
 \xi_i a=a_i\xi_i\qquad (i=1,\ldots,n),
\tag{4}
\end{equation}
となる。ここで $\xi_i$ は $x_i$ から $x_i+1$ への移行を表し、$a_i$ は $a$ において $x_i$ を $x_i+1$ で置き換えて得られる関数を表す。したがって $a$ が恒等的に消えなければ $a_i$ もそうではない。これらの取り決めにより、積の次数に関する規則を含む加法と乗法のすべての規則が、可換律を除いて再び成り立つ。

以下の節では、この二つの場合を、任意の係数をもつ一般に定義された多項式領域の下に置く。

\subsection*{\S~2. 一般の多項式領域.}

任意に抽象的に定義された体 $\mathsf{P}$ から出発する。その元を小ラテン文字で表す。この体では代数学のすべての規則が成り立ち、とくに乗法の可換律と一意的な逆元も成り立つ。この体に $n$ 個の不定元 $\xi_1,\ldots,\xi_n$ を添加する。すなわち、積 $a\xi_{\nu_1}b\xi_{\nu_2}\ldots c\xi_{\nu_\rho}$ とその有限和からなる整域 $J$ を考える。ここでは乗法の可換律を除いて、加法と乗法のすべての法則が成り立つものとする。とくに
\[
 F(\xi_1,\ldots,\xi_n)=\sum a_{\nu_1\ldots\nu_\rho}\xi_{\nu_1}\cdots\xi_{\nu_\rho}
\]
という形の有限和を「多項式」と呼ぶ（以下では常に大ラテン文字で表す）。可換律の代わりに、いま $J$ の元に、\emph{二つの多項式の積が再び多項式となる}ような積法則を課す。したがって $J$ における全多項式の領域がすでに $J$ を尽くす。これに必要かつ十分なのは明らかに、元 $\xi_1,
\ldots,\xi_n$ と $P$ の任意の元 $a$ に対して、
\begin{equation}
 \xi_i a=F_i(\xi_1,\ldots,\xi_n)\qquad (i=1,\ldots,n)
\tag{5}
\end{equation}
という形の積法則が成り立つことである。ここで $F_i(\xi_1,\ldots,\xi_n)$ は $\xi_i$ と $a$ に自然に依存するある多項式である。（とくに微分式と差分式に対する法則 (2), (4) はこの形である。）というのも、$\xi_i a$ はそれ自体としては多項式ではないが、仮定により領域に属するので、この式は必要である。逆に、式 (5) を有限回反復して適用することにより、領域 $J$ の任意の元は多項式へ変形できる。

体の単位を 1 と表し、さらに $\xi_i\cdot1=1\cdot\xi_i=\xi_i$ を要求する。また $\xi_1,\ldots,\xi_n$ 同士の乗法は可換であると仮定する。\footnote{これは \S~6 以降で有限性定理とその帰結にだけ必要である。} すると $J$ の各多項式は、$P$ の係数をもつ $\xi_1,\ldots,\xi_n$ の普通の多項式とまったく同じように書け、二つの多項式が一致するのはすべての係数が一致するときに限られる。

さらに以下では、\footnote{これも \S~6 以降でだけ用いられる。} (5) の代わりに規則
\begin{equation}
 \xi_i a=a_i\xi_i+b_i\qquad (a_i\ne0)\quad(i=1,\ldots,n)
\tag{6}
\end{equation}
が成り立つことを要求する。ここにも、とくに法則 (2), (4) が含まれる。すると積の次数は因子の次数の和に等しい。これは全次数についても各個の変数に関する次数についても成り立つので、恒等的に消えない二つの多項式の積もまた 0 とは異なる。

\subsection*{\S~3. 片側加群とその剰余群.}

我々の領域 $J$ では乗法が非可換であるため、次の種類の加群\footnote{加群は大ドイツ文字で表す。}を区別できる。\footnote{片側イデアルはすでに A. Hurwitz によって考察されている。Vorlesungen über die Zahlentheorie der Quaternionen (Berlin, Springer 1919), Lecture 6 を参照。ただしそこでは本質的には主イデアルが問題であり、そこで注 1 に引用された Du Pasquier の仕事についても同じである。}

1. \emph{右側加群}: $M$ とともに任意の多項式 $F$ に対する $FM$ を含み、また $M_1$ と $M_2$ とともに $M_1+M_2$ を含むもの。

2. \emph{左側加群}: $M$ とともに $MF$ を含み、また $M_1$ と $M_2$ とともに $M_1+M_2$ を含むもの。

これら二種類をまとめて\emph{片側加群}と呼ぶ。右側かつ左側である加群は\emph{両側}と呼ぶ。以下では右側加群に限り、左側加群の理論もまったく同じように構成できることを記しておく。

二つの多項式 $F$ と $G$ について、その差が $\mM$ で割り切れる、すなわち $\mM$ に属するとき、$F$ と $G$ は $\mM$ を法として合同であるといい、$F\equiv G(\mM)$ と書く。ある多項式と合同なすべての多項式の全体は、$\mM$ を法とする\emph{剰余類}をなす。（剰余類は小ドイツ文字で表す。）
\[
 F_1\equiv F_2(\mM)\qquad\text{かつ}\qquad G_1\equiv G_2(\mM)
\]
から $F_1+G_1\equiv F_2+G_2(\mM)$ が従うので、二つの剰余類 $\mx$ と $\my$ の和は定義され、再び剰余類 $\mx+\my$ となる。これに対して、加群は片側であるため、二つの剰余類の積については語れない。実際、
\[
 F_1=F_2+M_1,\qquad G_1=G_2+M_2
\]
からは
\[
 F_1G_1=F_2G_2+F_2M_2+M_1G_2+M_1M_2
\]
しか従わない。$M_1$ が $\mM$ に属しても $M_1G_2$ は $\mM$ に属するとは限らないので、一般には $F_1G_1$ と $F_2G_2$ が合同であるとは限らない。一方、$M_1=0$、すなわち $F_1=F_2=F$ ならば、$G_1\equiv G_2$ から $FG_1\equiv FG_2$ が従う。したがって任意の剰余類 $\my$ は前から多項式 $F$ で掛けることができ、よく定まった剰余類 $F\my=\mz$ を与える。

剰余類の計算では、加法の可換律と結合律、および加法の一意な逆操作の法則が成り立つ。また剰余類への多項式の乗法と加法については、分配律 $F(\my_1+\my_2)=F\my_1+F\my_2$ が成り立つ。これは $\mM$ を法とする多項式計算から直接従う。

単位が属する剰余類を単位類、または単位と呼び、$\me$ で表す。すると $F\me=\mx$ は多項式 $F$ が属する剰余類である。ただし一般の剰余類と異なり、この場合には積 $\mx\me$ もよく定義され、$\mx$ に等しい。なぜなら
\[
 F_1=F_2+M_1,\qquad G_1=1+M_2
\]
と $M_1\cdot1=M_1$ から
\[
 F_1G_1=F_2+M
\]
が従うからである。$\mM$ を法とする剰余類全体を $\mM$ の\emph{剰余群}という。したがって剰余群は次の性質をもつ。1. 任意の剰余類 $\mx$ と任意の多項式 $F$ とともに、剰余類 $F\mx$ を含み、また剰余類 $\mx_1,\mx_2$ とともに $\mx_1+\mx_2$ を含む。\footnote{この性質は、より一般の意味で、剰余群の「加群性」と呼べる。} 2. 単位 $\me$ をもち、各多項式 $F$ に対して $F\me$ が $F$ の剰余類 $\mx$ を表す。\emph{したがって} $F$ \emph{がすべての多項式を動くとき、} $F\me$ \emph{は剰余群全体を動く。}

\subsection*{\S~4. 剰余群の二つの部分群への可約性と、加群との関係.}

片側加群についても、次の定義が示すように、\emph{可除性}、\emph{最大公約数}、\emph{最小公倍元}の概念は両側の場合と同じままである。

$M\equiv0(\mM)$ から $M\equiv0(\mN)$ も従うとき、加群 $\mM$ は別の加群 $\mN$ で割り切れるという。

二つの加群 $\mM$ と $\mN$ について、\emph{最大公約数} $(\mM,\mN)$ は、$M\equiv0(\mM)$ かつ $N\equiv0(\mN)$ であるような形 $M+N$ で表されるすべての多項式の全体として定義される。これは再び加群である。複数の加群についても同じことが成り立つ。

$\mM$ と $\mN$ の両方で割り切れる全多項式は、加群 $[\mM,\mN]$ をなす。これは $\mM$ と $\mN$ の\emph{最小公倍元}である。この定義も、複数の加群に、さらには任意濃度の加群集合にも対応して適用される。

二つの加群は、その最大公約数が単位加群、すなわち全多項式の集合であるとき、互いに素と呼ばれる。

いま、互いに素な二つの加群 $\mN$ と $\mL$ の最小公倍元として表せる加群 $\mM$ から出発する。この二つはいずれも真の約因子であると仮定する。
\begin{equation}
 \mM=[\mN,\mL].
\tag{7}
\end{equation}
$\mM$ の剰余群 $\mG$ を調べる。仮定により、
\begin{equation}
 1\equiv N+L(\mM)
\tag{8}
\end{equation}
を満たす二つの多項式 $N\equiv0(\mN)$, $L\equiv0(\mL)$ がある。ここから
\begin{equation}
 \mL=(\mM,L),\qquad \mN=(\mM,N)
\tag{9}
\end{equation}
を示す。ここで、たとえば $(\mM,L)$ は $M+FL$ という形の全多項式からなる加群を意味する。実際、まず $(\mM,L)\equiv0(\mL)$ である。逆に (8) から、任意の多項式 $F$ に対して
\begin{equation}
 F\equiv FN(\mL)
\tag{10}
\end{equation}
が従う。いま $F\equiv0(\mL)$ ならば $FN\equiv0(\mL)$、したがって $\equiv0(\mM)$ であり、従って $F\equiv FL(\mM)$ である。同様に $\mN=(\mM,N)$ も証明される。なお、$\mN$ と $\mL$ は $\mM$ の真の約因子であるから、$L\equiv0$ も $N\equiv0(\mM)$ も成り立たないことが分かる。

$\ma$ を $\mM$ を法とする $N$ の剰余類、$\mb$ を $L$ の剰余類とする。このとき (8) により
\[
 \me=\ma+\mb\qquad(\ma\ne0,\ \mb\ne0)
\]
である。任意の剰余類は $F\me$ の形であるから、分配律により、任意の剰余類について
\begin{equation}
 F\me=F\ma+F\mb\qquad(\ma\ne0,\ \mb\ne0)
\tag{11}
\end{equation}
が成り立つ。任意の剰余類を $G\ma$ と $H\mb$ の形の二つの剰余類の和として表すこの表示は、しかし\emph{一意}である。なぜなら $0=G\ma+H\mb$ という関係があり、$K$ が剰余類 $G\ma=-H\mb$ に属する多項式なら、明らかに $K\equiv GN(\mM)$, $K\equiv-HL(\mM)$ であり、従って $K\equiv0(\mN)$, $K\equiv0(\mL)$、したがって $K\equiv0(\mM)$ である。すなわち $G\ma=0$, $H\mb=0$ である。

逆に、いま関係 (11) が任意の多項式 $F$ について、しかも上記の意味で一意に成り立つと仮定する。$N$ を $\ma$ に属する多項式、$L$ を $\mb$ に属する多項式とする。二つの加群 (9) を作り、(7) と (8) を主張する。後者は $F=1$ のときの (11) から従う。さらに $\mN$ と $\mL$ の定義により、$\mM$ はこれら二つの加群の公倍元であり、しかも最小のものである。実際、$F\equiv0(\mL)$ かつ $F\equiv0(\mN)$ から $F\equiv HL\equiv GN(\mM)$ が従う。従って $G\ma=H\mb$, $G\ma-H\mb=0$ であり、表示の一意性により $G\ma=0$, $H\mb=0$, $F\equiv0(\mM)$ である。したがって、\emph{加群} $\mM$ \emph{を互いに素な二つの真の約因子の最小公倍元として表すことは、} $\mM$ \emph{の剰余群} $\mG$ \emph{を式} (11) \emph{で与えられるように一意に加法分解することと同値である}ことが示された。

次に、$F\ma$ の形の剰余類からなる系 $\mA$ を考え、これを $\mL$ の剰余群 $\mbarA$ と関係づける。(11) により、任意の多項式 $F$ は、$\bara$ を $\mL$ を法とする $N$ の剰余類、したがって $\mbarA$ の単位類とするとき、$\mbarA$ において剰余類 $F\ma$ に属する。いま $\mA$ の剰余類 $F\ma$ を $\mbarA$ の剰余類 $F\bara$ に対応させると、$\mA$ の各剰余類に $\mbarA$ の剰余類が対応し、$\mA$ の二つの剰余類の和は $\mbarA$ の対応する剰余類の和に対応し、$\mA$ の剰余類を多項式で掛けたものは、$\mbarA$ の対応する剰余類を同じ多項式で掛けたものに対応する。この対応はさらに一対一である。$F\ma=0$ から (11) により $F\equiv0(\mL)$ が従い、従って $F\bara=0$ である。逆に $F\bara=0$、すなわち $F\equiv0(\mathfrak L)$ は、(10) により $FN\equiv0(\mL)$、従って $FN\equiv0(\mM)$、すなわち $F\ma=0$ も意味する。

これにより基本概念へ到達する。これを次のように定義する。

\emph{剰余類の二つの系 $\mA$ と $\mbarA$ の間の一対一対応は、$\mA$ の二つの剰余類の和に $\mbarA$ の対応する剰余類の和が割り当てられ、かつ $\mA$ の任意の剰余類と多項式との積に、$\mbarA$ の対応する剰余類と同じ多項式との積が割り当てられるとき、同型と呼ばれる。}

$\mG$ の中の剰余類の部分系 $\mA$ がある加群の剰余群に同型であるとき、これを $\mG$ の\emph{部分群}と呼ぶ。上の証明が示すように、ここでの対応は具体的には、$\bar{\mathfrak A}$ の各剰余類が $\mathfrak A$ の対応する剰余類に特定の多項式、すなわち $\mL$ に属するが $\mM$ には属さないすべての多項式を付け加えることで得られる、というものである。

同様に、$F\mb$ の形のすべての剰余類の全体は、$\mN$ の剰余群 $\mbarB$ に同型な $\mG$ の部分群 $\mB$ をなすことが示される。

剰余類が一意な加法分解 (11) を許す剰余群 $\mG$ は、二つの部分群 $\mA$ と $\mB$ へ\emph{可約}であるといい、$\mG=\mA+\mB$ と書く。可約でない群は\emph{既約}と呼ぶ。対応する加群も、ときに可約または既約と呼ぶ。したがって次が成り立つ。

\textbf{定理 I.}\footnote{Schmeidler, loc. cit., p. 39 を参照。} \emph{加群 $\mM$ が互いに素な二つの真の約因子 $\mL$ と $\mN$ の最小公倍元として表せるための必要十分条件は、その剰余群 $\mG$ が二つの部分群 $\mA$ と $\mB$ へ可約であることである。この場合、$\mA$ と $\mB$ はそれぞれ $\mL$ と $\mN$ の剰余群に同型である。}

\subsection*{\S~5. 単位による計算と $k$ 個の部分群への可約性.}

$\mG$ の剰余類 $\ma$ と $\mb$ は、単位 $\me$ と同じく、任意の剰余類 $\mx$ との積が定義されるという性質をもつ。したがってこれらも $\mG$ の単位と呼ぶ。実際、(11) で $F=M$ とすると
\[
 0=M\me=M\ma+M\mb
\]
が従い、従って一意性により $M\ma=0$, $M\mb=0$ であり、ゆえに
\[
 (F+M)\ma=F\ma=\mx\ma
\]
である。ここで $\mx$ は $F$ の剰余類を表す。従って関係 (11) の代わりに、同値な剰余類関係
\[
 \mx\me=\mx\ma+\mx\mb
\]
を書くことができる。とくに $\mx=\ma$ とおくと
\[
 \ma=\ma^2+\ma\mb
\]
を得るので、一意性により $\ma^2=\ma$, $\ma\mb=0$ である。同様に $\mb^2=\mb$, $\mb\ma=0$ を得る。

剰余群を二つの構成因子の和として表示することに対応して、$k$ 個の構成因子の和としての表示も定義できる。すなわち、一意な表示
\[
 F\me=F\ma_1+F\ma_2+\cdots+F\ma_k\qquad(\ma_1\ne0,\ldots,\ma_k\ne0)
\]
が与えられており、$0=G_1\ma_1+\cdots+G_k\ma_k$ からも $G_i\ma_i=0$ が従うと仮定する。このとき、とくに $F=M$ とすれば $M\ma_i=0$ が得られるので、表示はまた次の形にも書ける。\footnote{剰余群の加法分解については可換律と結合律が成り立つ。しかし可換の場合と違い、$\mU_1+\mU_2=\mU_1+\mU_3$ から $\mU_2=\mU_3$ は従わない。\S~12 の例 1 を参照。}
\begin{equation}
 \mx\me=\mx\ma_1+\cdots+\mx\ma_k\qquad(\ma_1\ne0,\ldots,\ma_k\ne0).
\tag{12}
\end{equation}
この表示は二つの構成因子への逐次分解から生じたものともみなせる。すべての $\ma_i\ne0$ を保ったままこれを無限に続けられるなら、無限個の構成因子の和という。

(12) で $\mx=\ma_i$ とおくと
\begin{equation}
 \ma_i^2=\ma_i,
 \qquad \ma_i\ma_j=0\quad(i\ne j),
 \qquad i,j=1,\ldots,k,
\tag{13}
\end{equation}
ただし $\ma_i\ne0$、を得る。逆に、$\mG$ の中に剰余類系 $\ma_1,\ldots,\ma_k$ があり、それらが直交条件 (13) を満たし、さらに
\begin{equation}
 \me=\ma_1+\cdots+\ma_k
\tag{14}
\end{equation}
を満たすと仮定する。これが群を $k$ 個の構成因子の和として表すことを示す。まず、仮定から、$\ma_i$ に属する任意の多項式 $A_i$ について $A_i\ma_i=(A_i+M)\ma_i$ であるから、$M\ma_i=0$ が従う。従って $\ma_i$ は任意の類で掛けることができる。よって (14) は表示 (12) を意味し、この表示は一意である。実際、$0=\mx_1\ma_1+\cdots+\mx_k\ma_k$ なら、右から $\ma_i$ を掛けることにより、(13) によって $0=\mx_i\ma_i$ が従う。これが一意性、したがって主張を証明する。単位類の分解に現れる剰余類 $\ma_1,
\ldots,\ma_k$ を、この分解に属する単位と呼ぶ。同値に、これらは直交関係 (13) と条件 (14) によって特徴づけられる。後の展開は本質的には単位による計算になる。

群を $k$ 個の構成因子へ加法分解することと、加群の対応する分解との関係を確立しておく。$k=2$ の場合と同様、帰納法により、加群を $k$ 個の互いに素な加群の最小公倍元として表すことが得られる。

分解 (12) が与えられているとし、これを
\begin{align}
 \mx\me&=\mx\mb+\mx\ma_k,\tag{15}\\
 \mx\mb&=\mx\ma_1+\cdots+\mx\ma_{k-1}.\tag{16}
\end{align}
と書く。(15) から定理 I により、剰余類 $\mx\mb$ は加群
\begin{equation}
 \mN=(\mM,A_k)
\tag{17}
\end{equation}
の剰余群に同型である。従ってその剰余類 $\mx\overline{\mb}$ については、(16) に対応する表示
\begin{equation}
 \mx\overline{\mb}=\mx\overline{\ma}_1+\cdots+\mx\overline{\ma}_{n-1}
\tag{18}
\end{equation}
も成り立つ。いま、すでに
\begin{equation}
 \mN=[\mM_1,\ldots,\mM_{k-1}]
\tag{19}
\end{equation}
が証明されていると仮定する。ここで
\begin{equation}
\left\{
\begin{aligned}
 \mM_1&=(\mN,A_2+\cdots+A_{k-1}),\\
 \mM_2&=(\mN,A_1+A_3+\cdots+A_{k-1}),\\
 &\hspace{1.8cm}\vdots\\
 \mM_{k-1}&=(\mN,A_1+\cdots+A_{k-2})
\end{aligned}\right.
\tag{20}
\end{equation}
である。二つの和因子の場合、この主張は先の結果と一致する。さらに、ここで多項式 $A_i$ は、$\overline{\ma}_i$ に属し同時に $\ma_i$ にも含まれる多項式として選べる。なぜなら \S~4 により後者の剰余類の多項式は前者に含まれるからである。このとき (15) により
\begin{equation}
 \mM=[\mN,\mM_k]
\tag{21}
\end{equation}
であり、ここで
\[
 \mM_k=(\mM,A_1+\cdots+A_{k-1}).
\]
(19) と (21) から、最小公倍元の定義により
\[
 \mM=[[\mM_1,
 \ldots,\mM_{k-1}],\mM_k]=[\mM_1,\ldots,\mM_{k-1},\mM_k]
\]
を得る。また (17) と (20) から
\[
 \mM_1=(\mM,A_k,A_2+\cdots+A_{k-1})
\]
を得る。従って
\[
 \mM_1=(\mM,A_2+\cdots+A_k)
\]
である。というのも、後者の加群は (13) により多項式
\[
 A_k\equiv A_k(A_2+\cdots+A_k)(\mM)
\]
も含むからである。$\mM_2,\ldots,\mM_{k-1}$ についても対応する式が得られる。さらに、$\mM_1,
\ldots,\mM_k$ は互いに素である。なぜなら (14) により
\[
 \frac{1}{k-1}\Bigl((A_2+\cdots+A_k)+(A_1+A_3+\cdots+A_k)+\cdots+(A_1+\cdots+A_{k-1})\Bigr)\equiv1\pmod{\mM}
\]
であるからである。さらにそれらは全体として互いに素な系をなす。すなわち、それらの任意の一部の最小公倍元は、残りの最小公倍元と互いに素である。\footnote{これは Loewy が導入した概念の一般化である。} これは式 (12) の和因子を二つの適当な群にまとめることから直ちに従う。

これで、剰余群を $k$ 個の和因子へ分解することに対し、加群を、全体として互いに素な系をなす $k$ 個の加群の最小公倍元として表すことが対応することが示された。逆定理も同様に帰納法で証明する。

従って
\[
 \mM=[\mM_1,\ldots,\mM_k]=[[\mM_1,\ldots,\mM_{k-1}],\mM_k]=[\mN,\mM_k]
\]
であり、$\mM_1,
\ldots,\mM_k$ が全体として互いに素な系をなすと仮定する。定理 I と組み合わせると、これは剰余群に対して表示 (15) を与える。$\mM_1,
\ldots,\mM_{k-1}$ も全体として互いに素な系をなす。もしこの系の加群を二つの群に分けたとき、一方の最小公倍元と他方の最小公倍元が共通の約因子をもつなら、加群 $\mM_k$ を二つの群の一方に加えてから最小公倍元を作っても、この約因子は残るからである。\footnote{この議論を続けると、$\mM_1,
\ldots,\mM_k$ の任意の二つが互いに素であることが分かる。その逆、すなわちこれが $\mM_1,
\ldots,\mM_k$ が全体として互いに素な系をなすことを含意する、ということは、一般には可換の場合に限って成り立つ。\S~12 の反例 1 を参照。} $k=2$ の場合、「全体として互いに素」は「互いに素」と同じ意味である。従って我々の主張は $k-1$ まで証明済みとしてよい。こうして剰余類 $\mx\overline{\mb}$ については一意表示 (18) が成り立ち、同型性により (16) も成り立つ。従って主張が従う。

\textbf{定理 II.} \emph{加群 $\mM$ が、全体として互いに素な系をなす $k$ 個の加群 $\mM_1,
\ldots,\mM_k$ の最小公倍元として表せるための必要十分条件は、その剰余群 $\mG$ が $k$ 個の部分群 $\mU_1,
\ldots,\mU_k$ へ可約であることである。適切な記法の下で、$\mU_i$ は $\mM_i$ の剰余群に同型である。}

\subsection*{\S~6. 有限性定理.}

この節では、任意の剰余群が有限個の既約構成因子の和として表せることを示す。そのためにまず Hilbert の加群基底定理を移す。

以下では \S~2 で定式化したより特殊な仮定を用いる。$\xi_1,
\ldots,\xi_n$ 同士の乗法は可換であり、多項式の乗法法則は
\begin{equation}
 \xi_i a=a_i\xi_i+b_i,
 \qquad (a_i\ne0)\quad(i=1,\ldots,n)
\tag{6}
\end{equation}
という形である。主張は加群についてだけ証明すれば十分である。実際、$\mS$ が任意の多項式系であれば、$\mS$ に属する有限個の多項式 $G_1,
\ldots,G_k$ から $A_1G_1+
\cdots+A_kG_k$ の形で作れるすべての多項式を $\mS$ に付け加えて得られる加群 $\mM$ を作る。もし $F_1,
\ldots,F_l$ が $\mM$ の加群基底で、
\[
 F_i=\sum_j A_{ij}G_j\qquad(i=1,\ldots,l)
\]
ならば、これらの表示に現れる有限個の多項式 $G_j$ が $\mS$ の基底をなす。

加群基底定理を Gordan (Göttingen Nachrichten 1899) に従って証明する。証明は二つの部分に分かれる。第一は無限に多くの冪積からなる系 $\mP$ に関するもので、$\xi_1,
\ldots,\xi_n$ の可換性により本質的に変わらない。冪積
\[
 P=\xi_1^{a_1}\cdots\xi_n^{a_n}
\]
は、それぞれ一度だけ系に入れ、次数の増加順に、同じ次数内では何らかの固定した方法で並べる。いま $\mP$ の部分系 $\mQ$ を、先行するどの $P$ によっても割り切れないすべての冪積 $P=Q$ からなるものとして考える。するとどの $Q$ も後に来る $Q$ によって割り切れない。後のものは次数が高いか、同じ次数なら前のものと異なるからである。他方、$\mP$ の任意の $P$ は少なくとも一つの $Q$ によって割り切れる。反対に、どの $Q$ によっても割り切れない最初の $P$ を $P_m$ とする。すると $P_m$ 自身は $Q$ ではありえないので、少なくとも一つの先行する $P$ によって割り切れる。仮定によりこの先行する $P$ はある $Q$ によって割り切れるので、$P_m$ 自身も $Q$ によって割り切れる。これは矛盾である。

いま $Q$ の個数が有限であると主張する。最初の $Q$ を
\[
 Q_0=\xi_1^{h_1}\cdots\xi_n^{h_n}
\]
とする。任意の $Q_\lambda=\xi_1^{h_{1\lambda}}\cdots\xi_n^{h_{n\lambda}}$ は $Q_0$ で割り切れないので、各 $Q$ について少なくとも一つの $h_{\rho\lambda}<h_\lambda$ が成り立つ。このことが成り立つ最小の添字を $\lambda$ とする。$h_{\rho\lambda}$ が同じ固定値 $c_\lambda$ をもつすべての $Q$ を一つの類 $K(c_\lambda)$ にまとめる。$c_\lambda<h_\lambda$ であるから、これらの類の数は有限である。しかも各類は有限個の元しか含まない。なぜなら、類 $K(c_\lambda)$ にある $Q$ を $\xi_\lambda^{c_\lambda}Q'$ と書けば、$Q'$ も互いに割り切らない冪積の系をなし、変数は $n-1$ 個だけである。$n=2$ なら $Q'$ は一変数だけを含むので、その個数は 1、従って有限である。よって $n-1$ 変数について有限性を仮定してよい。主張は証明された。

第二に、項を辞書式順序で並べた多項式の加群 $\mM$ を考える。最高項に現れる冪積を、各々一度だけ書いた系を $\mP$ とする。すると系 $\mQ=(Q_1,
\ldots,Q_\rho)$ に、有限個の多項式 $F_1,
\ldots,F_\rho$ が対応する。これは各 $Q_i$ に、最高項が $b_iQ_i$ に等しい $\mM$ のある多項式を割り当てて得られ、したがって $F_i=b_iQ_i+
\cdots$ で、以下は低い項である。

いま $F=bP+\cdots$ を $\mM$ の任意の多項式とし、$P=\xi_1^{\beta_1}\cdots\xi_n^{\beta_n}Q_i$ とする。積
\[
 \xi_1^{\beta_1}\cdots\xi_n^{\beta_n}F_i
 =\xi_1^{\beta_1}\cdots\xi_n^{\beta_n}(b_iQ_i+\cdots)
 =c\,\xi_1^{\beta_1}\cdots\xi_n^{\beta_n}Q_i+\cdots
 =cP+\cdots
\]
を作る。ここで $c\ne0$ は (6) による。また、以下のすべての項は辞書式順序で低い項である。なぜなら、乗法は低い項を除けば指数に関して可換だからである。\footnote{これは法則 (6) の代わりに、固定した $\xi$ の番号づけに対して、より弱い法則
\[
 \xi_i a=a_i\xi_i+a_{i,i+1}\xi_{i+1}+\cdots+a_{in}\xi_n+b_i\qquad(a_i\ne0)
\]
を用いることもできることを示している。} したがって差
\[
 F-\frac bc\,\xi_1^{\beta_1}\cdots\xi_n^{\beta_n}F_i
\]
は低い項だけを含み、再び $\mM$ に属する。従ってこの手続きは有限回で終わり、$F_1,
\ldots,F_\rho$ が加群基底であることが分かる。

\textbf{定理 III.} \emph{無限に多くの多項式からなる任意の系 $\mS$ は有限な加群基底 $G_1,
\ldots,G_k$ をもち、$\mS$ の任意の多項式 $G$ は $G=A_1G_1+
\cdots+A_kG_k$ と表せる。}

次は直ちに従う。

\textbf{系.} \emph{$\mM$ を法とする無限に多くの剰余類からなる任意の系 $\mS$ は基底 $\mx_1,
\ldots,\mx_k$ をもち、$\mS$ の任意の $\mx$ は $\mx=A_1\mx_1+
\cdots+A_k\mx_k$ の形で表せる。}

すなわち、各剰余類に代表元を一つ割り当て、それらの代表元の系と $\mM$ のすべての多項式を併せて考えればよい。

次の定理の証明に移る。

\textbf{定理 IV.} \emph{任意の剰余群は有限個の既約構成因子の和として表せる。同値に、定理 II により、任意の加群は全体として互いに素な系をなす有限個の既約加群の最小公倍元として表せる。}

反対に、$\mG$ が無限個の構成因子 $\mU_1,\mU_2,
\ldots$ の和であると仮定する（\S~5 参照）。対応する単位を $\ma_1,
\ma_2,
\ldots$ とし、これらは定義によりすべて非零である。これらすべての単位の系 $\mS$ を作る。定理 III の系により、これは有限基底 $\ma_{i_1},
\ldots,\ma_{i_v}$ をもち、$i_1<\cdots<i_v$ としてよい。$\mu>i_v$ とする。すると
\[
 \ma_\mu=\mx_1\ma_{i_1}+\cdots+\mx_v\ma_{i_v}
\]
である。これは単位の間の線形関係であり、定義によって排除されている。したがって剰余群の任意の加法分解では、有限回の後に既約構成因子が現れる。分解の可換性により、それを先頭に置くことができる。従って一般性を失わずに、$\mU_i$ 自身が既約であると仮定してよい。しかし前の議論により、現れ得る構成因子 $\mU_i$ は有限個だけである。主張は証明された。

\subsection*{\S~7. 分解が一意である場合.}

\[
 \mG=\mU_1+\cdots+\mU_k=\mB_1+\cdots+\mB_l
\]
を、構成因子 $\mU_1,
\ldots,\mU_k$, $\mB_1,
\ldots,\mB_l$ が既約である剰余群とする。このとき
\begin{equation}
 \mx\me=\mx\ma_1+\cdots+\mx\ma_k
 =\mx\mb_1+\cdots+\mx\mb_l
\tag{22}
\end{equation}
である。ここで $\ma_\chi$ は $\mU_\chi$ の単位類に同型的に割り当てられた $\mG$ の剰余類であり、$\mb_\lambda$ は $\mB_\lambda$ の単位類に割り当てられたものである。$\lambda$ を $1,
\ldots,l$ の範囲の固定した添字として、$\mx$ を $\mx\mb_\lambda$ で置き換えると
\begin{equation}
 \mx\mb_\lambda=\mx\mb_\lambda\ma_1+\cdots+\mx\mb_\lambda\ma_k
\tag{23}
\end{equation}
を得る。この $\mB_\lambda$ の表示は一意であるが、右辺の個々の剰余類が $\mB_\lambda$ に属するとは限らないので、$\mB_\lambda$ の加法分解ではない。$\mb_\lambda$ は 0 ではないから、すべての積 $\mb_\lambda\ma_\chi$ が 0 であることはない。一般性を失わずに
\[
 \mb_\lambda\ma_1\ne0,
 \ldots,
 \mb_\lambda\ma_\rho\ne0,
 \qquad
 \mb_\lambda\ma_{\rho+1}=0,
 \ldots,
 \mb_\lambda\ma_k=0
\]
としてよい。すると
\[
 \mx\mb_\lambda=\mx\mb_\lambda\ma_1+\cdots+\mx\mb_\lambda\ma_\rho
\]
である。

$1,
\ldots,\rho$ の範囲の固定した $\chi$ について、0 でなく、かつ $\mx\mb_\lambda\ma_\chi\ne0$ も満たす剰余類 $\mx\mb_\lambda$ 全体を考える。二つの系 $\mx\mb_\lambda$ と $\mx\mb_\lambda\ma_\chi$ は同型である。まず対応は一対一である。$\mx\mb_\lambda\ma_\chi=0$ から仮定により $\mx\mb_\lambda=0$ が従い、後者から、表示 (23) の一意性により $\mx\mb_\lambda\ma_\chi=0$ が従う。さらに、和は和に対応し、任意の多項式との積は対応する積に対応する。$\chi$ を固定したまま、同じ議論は $\mb_\lambda\ma_\chi\ne0$ であるすべての $\lambda$ について行える。各 $\ma_\chi$ に対して、そのような $\mb_\lambda$ が少なくとも一つ存在しなければならない。なぜなら $\ma_\chi=(\mb_1+
\cdots+\mb_l)\ma_\chi\ne0$ だからである。

まず、各 $\chi$ に対してただ一つの $\lambda$ が、また各 $\lambda$ に対してただ一つの $\chi$ が存在して、$\mb_\lambda\ma_\chi\ne0$ であると仮定する。

この対応は $k=l$ を意味する。さらに、$\mb_\chi\ma_\chi\ne0$ かつ $\lambda\ne\chi$ なら $\mb_\lambda\ma_\chi=0$ となるように記法を選べる。このとき積 $\mb_\lambda\ma_\chi$ は対角図式をなすという。この場合、(23) により単位は $\mb_\chi=\mb_\chi\ma_\chi$ であり、また
\begin{equation}
 \ma_\chi=\mb_1\ma_\chi+\cdots+\mb_l\ma_\chi=\mb_\chi\ma_\chi=\mb_\chi
\tag{24}
\end{equation}
であるから、一般に $\mU_\chi=\mB_\chi$ である。従って分解は一意であり、積 $\ma_\chi\mb_\lambda$ も対角図式をなす。

たとえば剰余類 $\ma_\chi\mb_\lambda$ について特に可換律が成り立つなら、上の仮定は常に満たされる。実際、(23) の各構成因子 $\mx\mb_\lambda\ma_\chi=\mx\ma_\chi\mb_\lambda$ はそのとき $\mB_\lambda$ に含まれるので、(23) は群 $\mB_\lambda$ の加法分解を与える。$\mB_\lambda$ は既約だから、$\mb_\lambda\ma_\chi$ のうち 0 でないものは一つだけである。対応して (24) により、固定した $\chi$ についても 0 でない $\mb_\lambda\ma_\chi$ は一つだけである。そうでなければ、$\ma_\chi$ と $\mb_\lambda$ の交換可能性により、(24) は既約と仮定した $\mU_\chi$ の加法分解を与えてしまう。\footnote{Schmeidler, loc. cit. で扱われた非可換両側の場合にも同じことが当てはまる。そこでは一方の部分群の任意の剰余類と他方の部分群の任意の剰余類との積が消える。すると $\mb_\lambda\ma_\chi=\mb_\lambda\ma_\chi\mb_1+
\cdots+\mb_\lambda\ma_\chi\mb_l$ かつ $\lambda\ne\mu$ なら $\mb_\lambda(\ma_\chi\mb_\mu)=0$ であるから、$\mb_\lambda\ma_\chi=\mb_\lambda\ma_\chi\mb_\lambda$ である。従ってこの類は $\mB_\lambda$ に含まれ、$\mB_\lambda$ の既約性により、ただ一つ、したがってちょうど一つの $\chi$ についてだけ $\mb_\lambda\ma_\chi\ne0$ となる。同じように、各 $\ma_\chi$ に対して $\ma_\chi\mb_\lambda\ne0$ となる $\mb_\lambda$ が一つだけ存在することが示され、従って $k=l$ である。よって $\mb_\lambda\ma_\chi$ は対角図式をなし、そこで確立された一意性が証明される。}

条件 $\mb_\lambda\ma_\lambda\ne0$, $\lambda\ne\chi$ なら $\mb_\lambda\ma_\chi=0$（$\chi=1,
\ldots,k$）、また各 $\mu=1,
\ldots,l$ について $\mb_\mu\ma_\chi=0$, $\mb_\mu\ma_\mu\ne0$（$\chi=1,
\ldots,\sigma$）が、すべてではなく一部の $\mb_\lambda$（$\lambda=1,
\ldots,\sigma$）についてだけ満たされる場合にも、同じ結論が成り立つ。$\chi=1,
\ldots,\sigma$ について $\mU_\chi=\mB_\chi$ を得る。また
\[
 \mU_{\sigma+1}+\cdots+\mU_k=\mU,
 \qquad
 \mB_{\sigma+1}+\cdots+\mB_l=\mB
\]
とおくと、$\mU$ と $\mB$ の単位 $\ma$ と $\mb$ についても、$\mb\ma_\chi=0$, $\mb_\chi\ma=0$（$\chi=1,
\ldots,\sigma$）および $\mb\ma\ne0$ が成り立ち、従って $\mU=\mB$ である。

\subsection*{\S~8. 完全可約群の素群への加法分解. 同型定理.}

次に、分解の一意性は成り立たないが、異なる分解の同型性が成り立つ場合を考える。これは Loewy が考察した場合の一般化である。

\textbf{定義.} \emph{加群が、自分自身と、全多項式を含む単位加群以外に約因子をもたないとき、その加群を素加群と呼ぶ。素加群の剰余群を素群と呼ぶ。}

素群はもちろん既約である。無限の素群も存在する（\S~11 の定義および \S~12, 2 の例を参照）。

与えられた剰余群 $\mG$ が、有限個の素群 $\mU_1,
\ldots,\mU_k$ および $\mB_1,
\ldots,\mB_l$ の和として二通りに表示できるとする。素群の和として少なくとも一つのこのような表示を許す剰余群、および対応する加群 $\mM$ を、Loewy の用語に従って完全可約と呼ぶ。

積 $\ma_\chi\mb_\lambda$（$\chi=1,
\ldots,k$, $\lambda=1,
\ldots,l$）を考え、これが $0$ であるか、または $\mx$ が $\mG$ のすべての剰余類を動くとき $\mx\ma_\chi\mb_\lambda$ が群 $\mB_\lambda$ 全体を動くかのどちらかであることを示す。実際、$\mB_\lambda$ に属する加群は
\[
 (\mM,\mb_1+\cdots+\mb_{\lambda-1}+\mb_{\lambda+1}+\cdots+\mb_l)=\mM_{\mb_\lambda};
\]
これは剰余類 $\mb_1+
\cdots+\mb_{\lambda-1}+\mb_{\lambda+1}+\cdots+\mb_l$ に属するすべての多項式を $\mM$ に付け加えて得られる。この加群は約因子
\[
 (\mM,\mb_1+\cdots+\mb_{\lambda-1}+\mb_{\lambda+1}+\cdots+\mb_l,
 \ma_\varkappa\mb_\lambda)
\]
をもつ。しかし $\mM_{\mb_\lambda}$ は仮定により素加群なので、この約因子は $\mM_{\mb_\lambda}$ に等しいか、単位加群に等しい。第一の場合、$\ma_\varkappa\mb_\lambda$ は $\mM_{\mb_\lambda}$ に属する。従って
\[
 \ma_\chi\mb_\lambda=\mx(\mb_1+
\cdots+\mb_{\lambda-1}+\mb_{\lambda+1}+\cdots+\mb_l)
\]
という関係があり、これは $\ma_\chi\mb_\lambda=0$ を与える。第二の場合には二つの類 $\mx_0$ と $\mx_1$ があり、
\[
 \mx_0\ma_\varkappa\mb_\lambda
 +\mx_1(\mb_1+\cdots+\mb_{\lambda-1}+\mb_{\lambda+1}+\cdots+\mb_l)
 =\me=\mb_1+\cdots+\mb_{\lambda-1}+\mb_{\lambda+1}+\cdots+\mb_l+\mb_\lambda
\]
となる。従って一意性により $\mx_0\ma_\varkappa\mb_\lambda=\mb_\lambda$ である。よって $F\mx_0\ma_\varkappa\mb_\lambda$、したがって $\mx\ma_\varkappa\mb_\lambda$ も、$F$ または $\mx$ がすべての多項式またはすべての剰余類を動くとき、群 $\mB_\lambda$ 全体を動く。これで主張が証明される。

もし $\ma_\chi\mb_\lambda\ne0$ なら、$\mU_\chi$ も $\mB_\lambda$ に同型である。\S~7 により、$\mx\ma_\chi\ne0$ かつ $\mx\ma_\chi\mb_\lambda\ne0$ であるすべての類 $\mx\ma_\chi$ の全体は、類 $\mx\ma_\chi\mb_\lambda$ の系に同型であり、後者は部分群 $\mB_\lambda$ を尽くす。残るのは、指示された類 $\mx\ma_\chi$ が $\mU_\chi$ を尽くすことを示すことだけである。そのために $\mx\ma_\chi\mb_\lambda=0$ であるすべての剰余類 $\mx\ma_\chi$ を考える。この系は、その二つの類の和および任意の多項式との積が再びその系に属するという性質をもつ。これらすべての剰余類を加群 $\mM_{\ma_\chi}$ に付け加えると、再び $\mM_{\ma_\chi}$ の約因子である加群が得られる。従ってそれは $\mM_{\ma_\chi}$ か単位加群である。しかし後者の場合には $\ma_\chi$ 自身がそこに含まれるので、$\ma_\chi\mb_\lambda=0$ となり、仮定に反する。従って $\mx\ma_\chi\ne0$ で $\mx\ma_\chi\mb_\lambda=0$ となる類は存在せず、$\mU_\chi$ と $\mB_\lambda$ の同型性が証明された。

いま固定した $\chi$ に対し、$\ma_\chi\mb_{\lambda_1}\ne0,
\ldots,\ma_\chi\mb_{\lambda_i}\ne0$ で $i>1$ であると仮定する。このとき $\mU_\chi$ は群 $\mB_{\lambda_1},\ldots,\mB_{\lambda_i}$ に同型であり、従ってこれらの群は互いに同型である。

次に、少なくとも二つの $\mU$ も同型でなければならないことを示す。もし積 $\mb_\lambda\ma_\chi$ において、各 $\mb_\lambda$ に対して $\mb_\lambda\ma_\chi\ne0$ となる $\ma_\chi$ が一つだけ対応し、そのため $\mb_\lambda=\mb_\lambda\ma_\chi$ となるなら、$\mB_\lambda=\mU_\chi$ である。なぜなら $\mU_\chi$ は素群として真の部分群を含めないからである。従って $k=l$ であり、番号づけを正しくすれば $\ma_\chi=\mb_\chi$ となる。しかしそのとき図式 $\ma_\chi\mb_\lambda$ は対角図式になるはずであり、これは今の場合ではない。従って次の定理が証明された。

\textbf{定理 V.} \emph{完全可約群 $\mG$ が、素群の和として二通りに
\[
 \mG=\mU_1+\cdots+\mU_k=\mB_1+\cdots+\mB_l
\]
と表されているなら、各群 $\mU$ は少なくとも一つの群 $\mB$ に同型であり、逆も成り立つ。各 $\mU$ がちょうど一つの $\mB$ に同型であるなら、分解は同一である。そうでない場合には、少なくとも二つの群 $\mU$ が互いに同型であり、同様に少なくとも二つの群 $\mB$ も互いに同型である。}

\subsection*{\S~9. 「同型」と「同種」の概念の関係.}

この節では、一変数の微分式について知られている「同種」の概念を自然に一般化すると、対応する剰余群の同型性へ導くことを示す。

\textbf{定義.}\footnote{一変数の微分式については、Blumberg の形式的定義 loc. cit., p. 25 を参照。} \emph{二つの加群 $\mM$ と $\mN$ は、$\mM$ に互いに素な多項式 $P$ と、$\mN$ に互いに素な多項式 $Q$ が存在し、任意の $M\equiv0(\mM)$ および $N\equiv0(\mN)$ に対して}
\begin{align}
 MQ&\equiv0(\mN),\tag{25}\\
 NP&\equiv0(\mM)\tag{26}
\end{align}
\emph{が成り立つとき、同種であると呼ばれる。ここで互いに素とは、二つの多項式 $X$, $Y$ が存在して}
\begin{align}
 YQ&\equiv1(\mN),\tag{27}\\
 XP&\equiv1(\mM)\tag{28}
\end{align}
\emph{が成り立つことを意味する。}

次の定理が成り立つ。

\textbf{定理 VI.} \emph{二つの加群が同種であるための必要十分条件は、それらの剰余群が同型であることである。}

1. $\mA$ を $\mM$ の剰余群、$\mB$ を $\mN$ の剰余群とし、$\mA$ が $\mB$ に同型であると仮定する。$\ma$ と $\mb$ を $\mA$ と $\mB$ の単位類とする。さらに、$\mA$ の単位 $\ma$ に対応する $\mB$ の類を $Q\mb$ とし、$\mB$ の単位 $\mb$ に対応する $\mA$ の類を $P\ma$ とする。この対応を
\begin{align}
 \ma&\sim Q\mb,\tag{29}\\
 P\ma&\sim\mb\tag{30}
\end{align}
と書く。すると
\[
 0=M\ma\sim MQ\mb,
 \qquad\text{すなわち}\qquad MQ\equiv0(\mN)
\]
が従い、(25) が証明される。同様に
\[
 NP\ma\sim N\mb=0,
 \qquad\text{すなわち}\qquad NP\equiv0(\mM)
\]
であり、(26) が証明される。さらに (29) により $P\ma\sim PQ\mb$ である。一方 (30) により $P\ma\sim\mb$ と仮定していたので、対応の一意性から $PQ\mb=\mb$、従って $PQ\equiv1(\mN)$ を得る。同様に $QP\equiv1(\mM)$ である。よって (27), (28) が $X=Q$, $Y=P$ によって証明された。

2. 逆に、$\mM$ と $\mN$ が同種であり、関係 (25) から (28) が成り立つと仮定する。$\mA$ の任意の類 $F\ma$ に、$\mB$ の類 $FQ\mb$ を割り当てる。これは $\mA$ の各類にただ一つの $\mB$ の類を割り当てる。実際、$F\ma=0$、すなわち $F=M$ なら、(25) により $FQ\mb=0$ が従う。さらに、この対応は群 $\mB$ 全体を尽くす。なぜなら特別な類 $Y\ma$ は $YQ\mb$ に対応し、これは (27) により $\mb$ に等しいからである。こうして $\mA$ から $\mB$ への全射的な写像 $\Phi$ が得られるが、まだ一対一であることは示されていない。同様に、式 (26) と (28) は $\mB$ から $\mA$ への全射的な写像 $\Psi$ を与えるが、これもまだ一対一とは示されていない。

この二つの写像のどちらかが一対一であれば、和と積に関する要請は満たされているので、同型性が証明される。しかしもし $\Phi$ と $\Psi$ がどちらも一対一でないなら、$\mA$ の中に、$\mB$ の零類 $GQ\mb=0$ に対応する類 $G\ma$ がある。すると他のすべての剰余類 $F\ma$ に対応する類 $FQ\mb$ だけで $\mB$ 全体をすでに尽くし、従って写像 $\Psi$ は類 $FQP\ma$ として $\mA$ 全体を返す。我々の知るように、$\Psi$ の下で $\mA$ の零類に対応する非零類も $\mB$ の中に存在する。$G_1\ma\ne0$ だが $G_1QP\ma=0$ であるような類をすべて $G_1\ma$ と表す。残りのすべての多項式 $F$ は、$FQP\ma$ が $\mA$ 全体を尽くすという性質をもつ。$FQP\ma=G_1\ma$ となる多項式 $F$ を $G_2$ と表す。このような多項式は、$G_1$ が存在すれば存在する。各 $G_2$ について $G_2QP\ma\ne0$ だが $G_2(QP)^2\ma=0$ である。対応して、任意の整数 $\nu$ について、$G_\nu(QP)^{\nu-1}\ma\ne0$ だが $G_\nu(QP)^\nu\ma=0$ となる類 $G_\nu\ma$ が存在する。

いま、すべての多項式 $G_1,G_2,
\ldots$ からなる系 $\mS$ を考える。これは有限な加群基底 $G_{\lambda_1},\ldots,G_{\lambda_e}$ をもつ。添字 $\lambda_1,
\ldots,\lambda_e$ の最大のものを $\lambda$ とし、$\nu>\lambda$ を選ぶ。すると
\[
 G_\nu=A_1G_{\lambda_1}+\cdots+A_eG_{\lambda_e}
\]
である。$(QP)^\lambda$ を掛けると
\[
 G_\nu(QP)^\lambda\equiv0(\mM)
\]
となり、これは仮定に反する。従って同型性が従い、さらに 1 により、より鋭い関係 (27), (28) が $X=Q$, $Y=P$ によって成り立つ。

定理 VI により、定理 V は次のようにも表現できる。

\textbf{定理 VII.}\footnote{Loewy, pp. 100--101 を参照。} \emph{完全可約加群 $\mM$ が、いずれも全体として互いに素な系をなす素加群 $\mP_1,
\ldots,\mP_k$ および $\mD_1,
\ldots,\mD_l$ の最小公倍元として二通りに表されるなら、各 $\mP$ は少なくとも一つの $\mD$ と同種であり、逆も成り立つ。各 $\mP$ がちょうど一つの $\mD$ と同種であるなら、分解は同一である。そうでない場合には、少なくとも二つの $\mP$ が互いに同種であり、同様に少なくとも二つの $\mD$ も互いに同種である。}

\setcounter{footnote}{19}
\subsection*{\S~10. 群の無限に多くの分解の存在.}

完全可約群の和表示が一意でないならば、定理 V により、どの表示にも少なくとも二つの互いに同型な部分群が存在する。そこで、そのような同型な部分群の系から作られる部分群
\[
  \mathfrak H=\mathfrak A_1+\cdots+\mathfrak A_\alpha
\]
を考える。ただし $\mathfrak A_1,\mathfrak A_2,\ldots,\mathfrak A_\alpha$ はすべて同型である。このとき、$\mathfrak H$ はこの種の相異なる分解を無限に多く許す、と主張する。ただし、有理性領域 $\mathsf{P}$ が無限に多くの「定数」、すなわち $\xi_i a=a\xi_i$ を満たす量 $a$ を含むことだけを仮定する。この定理は、群 $\mathfrak A_i$ について他に何も仮定しなくても、既約性さえ仮定しなくても成り立つ。すなわち、次が成り立つ。

\textbf{定理 VIII.} \emph{$\mathfrak G$ が有限個の互いに同型な部分群の和として表せるならば、有理性領域 $\mathsf{P}$ が無限に多くの定数を含む限り、その表示は無限に多くの異なる仕方で可能である。ここで定数とは、$\xi_i a=a\xi_i$ を満たす領域の量 $a$ をいう。}

\textbf{証明.} いま
\[
  \mathfrak G=\mathfrak A_1+\cdots+\mathfrak A_\alpha
\]
とし、$i,\varkappa=1,\ldots,\alpha$ に対して $\mathfrak A_i$ が $\mathfrak A_\varkappa$ と同型であるとする（$\alpha\ge2$）。明確な同型対応を指定するため、まず一つの部分群、たとえば $\mathfrak A_1$ を特に選ぶ。$\mathfrak A_1$ と $\mathfrak A_i$ の同型を
\begin{equation}
 a_1\sim t_{1i}a_i,
 \qquad
 a_i\sim t_{i1}a_1,
 \qquad
 t_{11}a_1=a_1
 \quad(i,\varkappa=1,
\ldots,\alpha)
\tag{31}
\end{equation}
によって定めるとする。つまり、群 $\mathfrak A_1$ の中では各剰余類が自分自身へ対応する。$M\mathfrak{a}_1=0$ および $M\mathfrak{a}_i=0$ であるから、$Mt_{1i}a_i=0$ および $M\mathfrak{t}_{i1}\mathfrak{a}_1=0$ が従う。したがって、クラス $t_{1i}a_i$ および $\mathfrak{t}_{i1}\mathfrak{a}_1$ は、単位と同様に、多項式だけでなくクラスによる乗法を許す。同型の定義から、このクラス乗法を許す二つの対応するクラス $\mathfrak y$ と $\bar{\mathfrak y}$ に対して、$\mathfrak x\mathfrak y$ と $\mathfrak x\bar{\mathfrak y}$ もまた対応するクラスであることが従う。したがって、関係 $\mathfrak{a}_r\mathfrak{a}_s=0$ $(r\ne s)$（$r=1,\ldots,\alpha$; $s=1,\ldots,\alpha$）は、$s=1$ のとき
\begin{equation}
\left\{
\begin{aligned}
 \mathfrak{a}_r \mathfrak{t}_{1\varkappa}\mathfrak{a}_\varkappa&=0 &&(\varkappa=1,
\ldots,
\alpha,
\ r\ne1),\\
 \text{また }s>1\text{ について}\qquad a_r t_{s1}a_1&=0 &&(r\ne s),\\
 \text{従ってまた}\qquad t_{1i}a_i&=a_1t_{1i}a_i,
 & \mathfrak{t}_{i1}\mathfrak{a}_1&=\mathfrak{a}_i\mathfrak{t}_{i1}\mathfrak{a}_1.
\end{aligned}\right.
\tag{32}
\end{equation}
を与える。さらに、(31) から、(27) と (28) の一般化として、対応の一意性により
\begin{equation}
 \mathfrak{t}_{1i}\mathfrak{t}_{i1}\mathfrak{a}_1=\mathfrak{a}_1,
 \qquad
 \mathfrak{t}_{\varkappa1}\mathfrak{t}_{1\varkappa}\mathfrak{a}_\varkappa=\mathfrak{a}_\varkappa.
\tag{33}
\end{equation}
が従う。いま $\mathfrak A_i$ と $\mathfrak A_\varkappa$ の同型として、(31) を合成して得られる関係を採用する。このような指定が必要なのは、個々の群が自分自身への同型を許し得るからである。したがって (31)、(32)、(33) から
\begin{equation}
 \mathfrak{a}_i\sim \mathfrak{t}_{i1}\mathfrak{t}_{1\varkappa}\mathfrak{a}_\varkappa=\mathfrak{t}_{i\varkappa}\mathfrak{a}_\varkappa,
 \qquad
 \mathfrak{t}_{ii}\mathfrak{a}_i=\mathfrak{a}_i,
 \qquad
 \mathfrak{a}_r\mathfrak{t}_{s\varkappa}\mathfrak{a}_\varkappa=0\quad (r\ne s),
\tag{34}
\end{equation}
を得る。ただし $\mathfrak{t}_{i1}\mathfrak{t}_{1\varkappa}=\mathfrak{t}_{i\varkappa}$ と置いた。(33) と (34) からさらに
\[
 \mathfrak{t}_{i\varkappa}\mathfrak{t}_{\varkappa l}\mathfrak{a}_l=\mathfrak{t}_{i1}\mathfrak{t}_{1\varkappa}\mathfrak{t}_{\varkappa1}\mathfrak{t}_{1l}\mathfrak{a}_l
 =\mathfrak{t}_{i1}\mathfrak{t}_{1l}\mathfrak{a}_l=\mathfrak{t}_{il}\mathfrak{a}_l
\]
が得られる。これにより、群 $\mathfrak A_1$ を特別に選んだことは再び消える。$\mathfrak A_i$ と $\mathfrak A_l$ の間の同型関係は、定義において $\mathfrak A_1$ の代わりにどの群 $\mathfrak A_\varkappa$ を用いても同じである。従って、$\mathfrak{t}_{i\varkappa}$ は、まとめると次の関係を満たすクラスである。
\begin{equation}
 M\mathfrak{t}_{i\varkappa}\mathfrak{a}_\varkappa=0;
 \quad \mathfrak{a}_r\mathfrak{t}_{i\varkappa}\mathfrak{a}_\varkappa=0\ (r\ne i);
 \quad \mathfrak{t}_{i\varkappa}\mathfrak{t}_{\varkappa l}\mathfrak{a}_l=\mathfrak{t}_{il}\mathfrak{a}_l;
 \quad \mathfrak{t}_{\varkappa\varkappa}\mathfrak{a}_\varkappa=\mathfrak{t}_{\varkappa i}\mathfrak{t}_{i\varkappa}\mathfrak{a}_\varkappa=\mathfrak{a}_\varkappa.
\tag{35}
\end{equation}
これらの関係は、$\mathfrak A_i$ および $\mathfrak A_\varkappa$ に属する加群 $\mathfrak M_i$ と $\mathfrak M_\varkappa$ が同種である、ということを正確に述べている（\S~9 の定義を参照）。そこに現れる $P$ と $Q$ の代わりに、ここではクラス $\mathfrak{t}_{i\varkappa}$ と $\mathfrak{t}_{\varkappa i}$ が現れ、(35) の最初の二つの関係を合わせたものが式 (25)、または (26) に対応する。逆に、これらの関係から定理 VI によって同型が従う。\footnote{しかも有限性定理を用いずにそう言える。なぜなら、式そのものがそれぞれ写像 $\varphi_{i\varkappa}$ とその逆 $\varphi_{\varkappa i}$ を示しているからである。実際、$\mathfrak{x}\mathfrak{a}_i=\mathfrak{x}\mathfrak{a}_i \mathfrak{t}_{i\varkappa}\mathfrak{a}_\varkappa \mathfrak{t}_{\varkappa i}\mathfrak{a}_i$ より、$\mathfrak{x}\mathfrak{a}_i\ne0$ ならば $\mathfrak{x}\mathfrak{a}_i \mathfrak{t}_{i\varkappa}\mathfrak{a}_\varkappa\ne0$ も従う。}

次に、クラス $\mathfrak{t}_{i\varkappa}\mathfrak{a}_\varkappa$ からクラス $\mathfrak{b}_\rho$ を作る。これらはふたたび群 $\mathfrak B_\rho$ の単位になることが分かり、従って表示
\[
 \mathfrak G=\mathfrak B_1+\cdots+\mathfrak B_\alpha
\]
が従う。$\lambda_{i\varkappa}$ を、行列式 $|\lambda_{i\varkappa}|$ が零でない定数とし、$\lambda^{i\varkappa}$ を対応する余因子を $|\lambda_{i\varkappa}|$ で割ったものとする。このとき
\begin{equation}
\left\{
\begin{aligned}
 \sum_{r=1}^{\alpha}\lambda_{ri}\lambda^{r\varkappa}&=\delta_{i\varkappa}=\begin{cases}0,&i\ne\varkappa,\\1,&i=\varkappa,
 \end{cases}\\
 \sum_{r=1}^{\alpha}\lambda_{ir}\lambda^{\varkappa r}&=\delta_{i\varkappa}
\end{aligned}\right.
\tag{36}
\end{equation}
が成り立つ。そこで
\begin{equation}
 \mathfrak{b}_\rho=\sum_{i,\varkappa}\lambda_{\rho i}\lambda^{\rho\varkappa}\mathfrak{t}_{i\varkappa}\mathfrak{a}_\varkappa
 \qquad(\rho=1,
\ldots,\alpha)
\tag{37}
\end{equation}
と定義する。(35)、(36)、(37) から $M\mathfrak{b}_\rho=0$ および
\begin{equation}
 \sum_{\rho=1}^{\alpha}\mathfrak{b}_\rho
 =\sum_{i,\varkappa,\rho}\lambda_{\rho i}\lambda^{\rho\varkappa}\mathfrak{t}_{i\varkappa}\mathfrak{a}_\varkappa
 =\sum_{i,\varkappa}\delta_{i\varkappa}\mathfrak{t}_{i\varkappa}\mathfrak{a}_\varkappa
 =\sum_i \mathfrak{t}_{ii}\mathfrak{a}_i=\sum_i \mathfrak{a}_i=\mathfrak{e},
\tag{38}
\end{equation}
さらに
\[
\begin{aligned}
 \mathfrak{b}_\rho \mathfrak{b}_\sigma
 &=\sum_{\substack{i,\varkappa\\ \mu,\nu}}
  \lambda_{\rho i}\lambda^{\rho\varkappa}\lambda_{\sigma\mu}\lambda^{\sigma\nu}
  \mathfrak{t}_{i\varkappa}\mathfrak{a}_\varkappa \mathfrak{t}_{\mu\nu}\mathfrak{a}_\nu \\
 &=\sum_{i,\mu,\nu}\lambda_{\rho i}\lambda^{\rho\mu}\lambda_{\sigma\mu}\lambda^{\sigma\nu}\mathfrak{t}_{i\nu}\mathfrak{a}_\nu
 =\delta_{\sigma\rho}\sum_{i,\nu}\lambda_{\rho i}\lambda^{\sigma\nu}\mathfrak{t}_{i\nu}\mathfrak{a}_\nu
 =\begin{cases}0&(\rho\ne\sigma),\\ \mathfrak{b}_\rho&(\rho=\sigma),\end{cases}
\end{aligned}
\]
も得る（$\rho,\sigma=1,\ldots,\alpha$）。したがって $\mathfrak r e=\mathfrak r b_1+\cdots+\mathfrak r b_\alpha$ は群 $\mathfrak G$ の加法分解である。$\lambda_{i\varkappa}$ が単に対角行列をなすのでない限り、群 $\mathfrak B_\rho$ は群 $\mathfrak A_\sigma$ と異なる。実際、$a_\sigma b_\rho$ を作ると、(35) により
\[
 a_\sigma b_\rho
 =\sum_{i,\varkappa}\lambda_{\rho i}\lambda^{\rho\varkappa}a_\sigma t_{i\varkappa}a_\varkappa
 =\lambda_{\rho\sigma}
 \sum_{\varkappa}\lambda^{\rho\varkappa}\mathfrak{t}_{\sigma\varkappa}\mathfrak{a}_\varkappa\ne0
 \quad\text{ のとき}\lambda_{\rho\sigma}\ne0,
\]
となる。もし最後の和が零なら各項が個別に消えなければならないが、それは起こらない。従って、固定した $\sigma$ に対して $\lambda_{\rho\sigma}\ne0$ である $\rho$ が少なくとも二つあれば、$\mathfrak A_\sigma$ はどの $\mathfrak B_\rho$ とも同一ではあり得ない。

しかし群 $\mathfrak B_\rho$ は相互にも、また群 $\mathfrak A_\sigma$ とも同型である。まず、$\mathfrak A_\sigma$ と $\mathfrak B_\rho$ の同型を媒介するクラス $\mathfrak{r}_{\rho\sigma}\mathfrak{a}_\sigma$ と $\mathfrak{r}^{\sigma\rho}\mathfrak{b}_\rho$ の存在を示す（$\rho,\sigma=1,\ldots,\alpha$）。すなわち、クラス
\[
 \mathfrak{r}_{\rho\sigma}\mathfrak{a}_\sigma=\sum_i\lambda_{\rho i}\mathfrak{t}_{i\sigma}\mathfrak{a}_\sigma,
 \qquad
 \mathfrak{r}^{\sigma\rho}=\sum_\varkappa\lambda^{\rho\varkappa}\mathfrak{t}_{\sigma\varkappa}\mathfrak{a}_\varkappa
 =\mathfrak{r}^{\sigma\rho}\mathfrak{b}_\rho\text{\footnotemark[21]}
\]
\footnotetext[21]{(37) により、$r^{\sigma\rho}b_\rho=\sum_{\varkappa,i,\mu}\lambda^{\rho\varkappa}t_{\sigma\varkappa}a_\varkappa\lambda_{\rho i}\lambda^{\rho\mu}t_{i\mu}a_\mu=\sum_{\varkappa,\mu}\lambda^{\rho\varkappa}\lambda_{\rho\varkappa}\lambda^{\rho\mu}t_{\sigma\mu}a_\mu=\sum_\mu\lambda^{\rho\mu}t_{\sigma\mu}a_\mu=r^{\sigma\rho}$ である。}
について、加群
\[
 (\mathfrak M,\mathfrak{b}_1+\cdots+\mathfrak{b}_{\rho-1}+\mathfrak{b}_{\rho+1}+\cdots+\mathfrak{b}_\alpha)
\]
と
\[
 (\mathfrak M,\mathfrak{a}_1+\cdots+\mathfrak{a}_{\sigma-1}+\mathfrak{a}_{\sigma+1}+\cdots+\alpha_\sigma)
\]
との間に「同種」の関係
\begin{equation}
\left\{
\begin{aligned}
 \mathfrak{b}_\varkappa \mathfrak{r}_{\rho\sigma}\mathfrak{a}_\sigma&=0 &&(\varkappa\ne\rho);&
 \mathfrak{r}^{\sigma\rho}\mathfrak{r}_{\rho\sigma}\mathfrak{a}_\sigma&=\mathfrak{a}_\sigma\text{\footnotemark[22]},\\
 \mathfrak{a}_\varkappa \mathfrak{r}^{\sigma\rho}\mathfrak{b}_\rho&=0 &&(\varkappa\ne\sigma);&
 \mathfrak{r}_{\rho\sigma}\mathfrak{r}^{\sigma\rho}\mathfrak{b}_\rho&=\mathfrak{b}_\rho,
\end{aligned}\right.
\tag{39}
\end{equation}
\footnotetext[22]{(35) と比べて公式の数が二倍になるのは、逆写像がそこにおけるような添字の交換によっては得られないからである。}
\setcounter{footnote}{22}
が満たされることを示す。これから (35) の場合と同様に、対応
\begin{equation}
 \mathfrak{b}_\rho\sim \mathfrak{r}_{\rho\sigma}\mathfrak{a}_\sigma;\qquad,
 \qquad
 \mathfrak{a}_\sigma\sim \mathfrak{r}^{\sigma\rho}\mathfrak{b}_\rho
\tag{40}
\end{equation}
により同型が従う。実際、
\[
\begin{aligned}
 \mathfrak{b}_\varkappa \mathfrak{r}_{\rho\sigma}\mathfrak{a}_\sigma
 &=\sum_{\mu,\nu,i}\lambda_{\varkappa\mu}\lambda^{\varkappa\nu}\mathfrak{t}_{\mu\nu}\mathfrak{a}_\nu\cdot\lambda_{\rho i}\mathfrak{t}_{i\sigma}\mathfrak{a}_\sigma\\
 &=\sum_{\mu,i}\lambda_{\varkappa\mu}\lambda^{\varkappa i}\lambda_{\rho i}\mathfrak{t}_{\mu\sigma}\mathfrak{a}_\sigma
 =\delta_{\rho\varkappa}\sum_\mu\lambda_{\varkappa\mu}\mathfrak{t}_{\mu\sigma}\mathfrak{a}_\sigma
 =\delta_{\rho\varkappa}\mathfrak{r}_{\varkappa\sigma}\mathfrak{a}_\sigma=0\quad(\varkappa\ne\rho),
\end{aligned}
\]
さらに
\[
 \mathfrak{r}^{\sigma\rho}\mathfrak{r}_{\rho\sigma}\mathfrak{a}_\sigma
 =\sum_{\varkappa,i}\lambda^{\rho\varkappa}\mathfrak{t}_{\sigma\varkappa}\mathfrak{a}_\varkappa\cdot\lambda_{\rho i}\mathfrak{t}_{i\sigma}\mathfrak{a}_\sigma
 =\sum_i\lambda^{\rho i}\lambda_{\rho i}\mathfrak{t}_{\sigma\sigma}\mathfrak{a}_\sigma=\mathfrak{a}_\sigma.
\]
逆向きにも同様に
\[
 \mathfrak{a}_\varkappa \mathfrak{r}^{\sigma\rho}\mathfrak{b}_\rho
 =\mathfrak{a}_\varkappa\sum_\mu\lambda^{\rho\mu}\mathfrak{t}_{\sigma\mu}\mathfrak{a}_\mu=0
 \quad(\varkappa\ne\sigma),
\]
および
\[
 \mathfrak{r}_{\rho\sigma}\mathfrak{r}^{\sigma\rho}\mathfrak{b}_\rho
 =\sum_{i,\mu}\lambda_{\rho i}\mathfrak{t}_{i\sigma}\mathfrak{a}_\sigma\lambda^{\rho\mu}\mathfrak{t}_{\sigma\mu}\mathfrak{a}_\mu
 =\sum_{i,\mu}\lambda_{\rho i}\lambda^{\rho\mu}\mathfrak{t}_{i\mu}\mathfrak{a}_\mu=\mathfrak{b}_\rho
\]
を得る。いま証明した同種の関係から、$\mathfrak A_\sigma$ と $\mathfrak B_\rho$ の同型が従う（ここでも有限性定理は用いない）。$\mathfrak B_\rho$ と $\mathfrak B_\tau$ の同型はこれを合成して得られる。公式を記しておく。
\[
 \ell_{\rho\tau}b_\tau=r_{\rho\sigma}r^{\sigma\tau}b_\tau
 =\sum_{i,j}\lambda_{\rho i}t_{i\sigma}a_\sigma\lambda^{\tau j}t_{\sigma j}a_j
 =\sum_{i,j}\lambda_{\rho i}\lambda^{\tau j}\mathfrak{t}_{ij}\mathfrak{a}_j
\]
と置く。このとき対応は $\mathfrak{b}_\rho\sim\mathfrak{s}_{\rho\tau}\mathfrak{b}_\tau$ である。ここでも同種の関係は形式的に検証できる。$t_{i\varkappa}$ の代わりに $\mathfrak{s}_{\rho\tau}$ を置けば、それらは (35) を満たす。実際、
\[
 b_\sigma\ell_{\rho\tau}b_\tau=b_\sigma r_{\rho\sigma}r^{\sigma\tau}b_\tau=0\quad(\sigma\ne\rho),
\]
\[
 \mathfrak{s}_{\rho\sigma}\mathfrak{s}_{\sigma\tau}\mathfrak{b}_\tau=\mathfrak{r}_{\rho\varkappa}\mathfrak{r}^{\varkappa\sigma}\mathfrak{r}_{\sigma\varkappa}\mathfrak{r}^{\varkappa\tau}\mathfrak{b}_\tau
 =\mathfrak{r}_{\rho\varkappa}\mathfrak{a}_\varkappa \mathfrak{r}^{\varkappa\tau}\mathfrak{b}_\tau=\mathfrak{s}_{\rho\tau}\mathfrak{b}_\tau.
\]

以上で定理 VIII は証明された。なぜなら、$\lambda_{i\varkappa}$ を無限に多くの仕方で選び、それらが対角行列をなさず、また対角行列によって互いに移されることもないようにすればよいからである。

定理 VIII から、定理 V または VII と合わせて、次が従う。

\textbf{定理 IX.}\footnote{Loewy, pp. 106--107 を参照。} \emph{加群 $\mathfrak M$ が完全可約で、完全に相対的素な系をなす素加群 $\mathfrak P_1,\ldots,\mathfrak P_k$ の最小公倍元であり、有理性領域 $P$ が無限に多くの定数を含むとする。このとき、$\mathfrak M$ が素加群の最小公倍元として無限に多くの異なる仕方で表せるための必要十分条件は、加群 $\mathfrak P_1,\ldots,\mathfrak P_k$ の少なくとも二つが同種であることである。}

ここで、加群 $\mathfrak M$ が無限に多くの素加群で割り切れるための一般判定法を導く。そのため、まず次を証明する。

\textbf{補助定理.}\footnote{Loewy, p. 96 を参照。} \emph{完全可約加群 $\mathfrak M=[\mathfrak P_1,\ldots,\mathfrak P_k]$ が、すべての $\mathfrak P_i$ と異なる素加群 $\mathfrak P'$ で割り切れるならば、$\mathfrak P_i$ の少なくとも二つは同種である。}

まず $\mathfrak P_1,\ldots,\mathfrak P_k$ から、まだ削除されていないものの最小公倍元に含まれる素加群を順次削除して、完全に相対的素な系を選び出す。従って、いま $\mathfrak P_1,\ldots,\mathfrak P_k$ 自身が完全に相対的素な系をなしていると仮定してよい。$a'$ を、$\mathfrak P'$ を法とする単位クラスに対応する $\mathfrak M$ のクラスとする。このとき任意の多項式 $F$ について
\[
 Fa'=Fa'a_1+
\cdots+Fa'a_k
\]
である。積 $\mathfrak{a}'\mathfrak{a}_i$ の少なくとも二つは零でない。たとえば $\mathfrak{a}'\mathfrak{a}_1$ と $\mathfrak{a}'\mathfrak{a}_2$ がそうである。実際、たとえば $\mathfrak{a}'\mathfrak{a}_1\ne0$ だけが成り立つならば、$F\mathfrak{a}'=F\mathfrak{a}'\mathfrak{a}_1$ となり、剰余類 $F\mathfrak{a}'$ は $\mathfrak A_1$ の部分群をなす。$\mathfrak A_1$ は素群であるから、それは $\mathfrak A_1$ と一致し、従って加群 $\mathfrak P'$ と $\mathfrak P_1$ も一致してしまう。\S~8 の議論により、$\mathfrak A'$ は $\mathfrak A_1$ および $\mathfrak A_2$ と同型である。

いま $\mathfrak M$ を任意の加群とする。$\mathfrak M$ を割り切るすべての素加群の最小公倍元 $\mathfrak N$ を考える。\footnote{それらが無限に多く存在するならば、定理 IV により、この表示に対応する剰余群の加法分解は存在し得ない。\S~12, 4 の例も参照。} 二つの場合があり得る。第一に、$\mathfrak N$ が有限個の素加群の最小公倍元として表せない場合である。この場合、$\mathfrak N$、したがって $\mathfrak M$ も、無限に多くの素加群を因子にもつ。第二に、$\mathfrak N$ が有限個の素加群の最小公倍元として表せる場合である。このとき $\mathfrak N$ を、Loewy の対応する構成に従って、$\mathfrak M$ の最大完全可約加群と呼ぶ。これら有限個の素加群の中に同種のものが二つあれば、定理 IX により、加群 $\mathfrak N$、従って $\mathfrak M$ は、無限に多くの素加群で割り切れる。逆に後者が成り立つならば、$\mathfrak N$ には、$\mathfrak N$ を表示する最小公倍元に用いられた有限個の加群のどれとも異なる素因子が少なくとも一つ存在する。補助定理により、$\mathfrak N$、従って $\mathfrak M$ は、同種である二つの素因子をもつ。こうして次を証明した。

\textbf{定理 X.} \emph{有理性領域 $\mathsf{P}$ が無限に多くの定数を含むならば、ある加群が無限に多くの相異なる素加群で割り切れるための必要十分条件は、最大完全可約加群が存在しないか、またはそのような加群が存在する場合には、それが同種である少なくとも二つの相異なる素加群で割り切れることである。}

\subsection*{\S~11. 微分方程式系およびその積分との関係.}

有限性定理により、微分式からなる各加群には、有限個の微分式の系を基底として割り当てることができる。その結果、加群の微分式を零に等置して得られるすべての微分方程式の同時積分全体は、基底微分式を零に等置して得られる有限個の微分方程式系の積分全体と一致する。これが線形偏微分方程式系の理論との関係を与える。ここでは、加群の剰余群が「有限」である場合に、この関係を少しさらに追う。

\textbf{定義.} \emph{剰余群は、ある数 $\mu$、すなわち群の位数が存在して、任意の $\mu+1$ 個の剰余類の間には $\mathsf{P}$ の係数をもつ線形関係が存在し、しかも少なくとも一つの $\mu$ 個の剰余類の系で、それらの間には関係が存在しないものがあるとき、有限と呼ばれる。このような系を剰余群の線形基底と呼ぶ。そうでない場合、剰余群は無限と呼ばれる。}

二つの有限な剰余群の和は有限位数をもち、その位数はそれぞれの位数の和である。二つの無限群の和、または有限群と無限群の和は無限群である。有限群の例は、一変数の加群のすべての剰余群によって与えられる。しかし多変数でも位数が有限になることがある。たとえば基底 $\xi_1^2$, $\xi_2^2$ をもつ加群では、すべての剰余類が $1,\xi_1,\xi_2,\xi_1\xi_2$ の剰余類から線形に構成できる。無限群の例については \S~12 を参照。

いま次が成り立つ。

\textbf{定理 XI.} \emph{ある加群が有限な剰余群をもつならば、その加群の基底を零に等置して得られる任意の微分方程式系は、その剰余群の位数が示すだけの個数の線形独立な積分を正確にもつ。逆に、有限個の線形偏微分方程式からなる与えられた系が有限個の線形独立な積分しかもたないならば、この個数は、それらの微分式が生成する加群の剰余群の位数に等しい。}

ここで有理性領域 $\mathsf{P}$ は、$n$ 個の変数 $x_1,\ldots,x_n$ の解析関数から成る。ただし、ある $n$ 次元領域内で現れる特異性は高々低次元の多様体をなすものとし、したがって正則な $n$ 次元近傍をもつ点が任意に多く存在するとする。

$\mathfrak M$ を有限な剰余群をもつ加群とする。次の性質をもつ冪積の線形基底が存在することを示す。すなわち、\emph{任意の剰余類をこの基底によって表すとき、分母に現れ得るのは、$\mathsf{P}$ の有限個の異なる量の冪だけである}。

この目的で、すべての冪積 $\xi_1^{a_1}\cdots\xi_n^{a_n}$ を辞書式に順序づけ、前にあるものから $\mathfrak M$ を法として線形従属しない剰余類の代表 $Q_1,\ldots,Q_m$ をすべて取る。仮定により、そのような積は有限個しかない。その他のものはすべて、これらによって線形に表される。

$\mathfrak S$ を、線形基底に採用されなかったすべての冪積の系とする。すると、定理 III の証明の第一部により、$\mathfrak S$ は冪積からなる有限部分系 $\mathfrak T$ をもち、$\mathfrak S$ の任意の冪積は $\mathfrak T$ の少なくとも一つで割り切れる。$P_1,\ldots,P_\rho$ を $\mathfrak T$ の冪積とし、$M_1=0,\ldots,M_\rho=0\;(\mathfrak M)$ を、$P_1,\ldots,P_\rho$ のそれぞれを $\mathfrak M$ を法として低い冪積の線形結合で表す関係とする。たとえば
\[
 M_i=b_iP_i+\text{低い項}.
\]
$P$ が $\mathfrak S$ の任意の冪積であるならば、それは
\[
 \xi_1^{h_1}\cdots\xi_n^{h_n}P_i
\]
の形をしているので、
\[
\begin{aligned}
 \xi_1^{h_1}\cdots\xi_n^{h_n}M_i
 &=b_i\xi_1^{h_1}\cdots\xi_n^{h_n}P_i+\text{低い項}\\
 &=b_iP+\text{低い項}
\end{aligned}
\]
となる。いま、$P$ より低いすべての冪積について、分母には $b_1,\ldots,b_\rho$ の冪だけが現れ、分子は $M_i$ の係数から加法、乗法、微分によって作られることがすでに証明されていると仮定する。すると同じことが $P$ 自身についても成り立つ。この仮定は許される。なぜなら、$\mathfrak S$ と $\mathfrak T$ の最初の元 $P_1$ について成り立っているからである。

さらに定理 III により、$M_1,\ldots,M_\rho$ は $\mathfrak M$ の加群基底である。いま、$x_1,\ldots,x_n$ の値の系 $\beta_1,\ldots,\beta_n$ で、$b_1\ne0,\ldots,b_\rho\ne0$ かつ $M_i$ の残りの係数が正則であるものを考える。すなわち、微分方程式系 $M_1=0,\ldots,M_\rho=0$ の正則点を考える。冪積 $Q_i=\xi_1^{h_{i1}}\cdots\xi_n^{h_{in}}$（$i=1,\ldots,m$）には偏導関数
\[
 \frac{\partial^{h_{i1}+\cdots+h_{in}}}{\partial x_1^{h_{i1}}\cdots\partial x_n^{h_{in}}}y
\]
が対応する。これらに対して、点 $(\beta_1,\ldots,\beta_n)$ における任意の定数値 $(\gamma_1,\ldots,\gamma_m)$ を指定する。すると微分方程式 $M_i=0$ は、その点における $y$ のすべての高階導関数の値を有限な値として一意に決定する。分母に現れるのは $b_i$ だけだからである。そしてこれは、$m$ 個の線形独立な解析的積分の存在を含意することが知られている。

他方、我々の仮定のもとでは、$m$ 個より多い線形独立な積分は存在し得ない。そうでなければ、正則点で適当に選んだ $m+1$ 個の導関数の値を任意に指定できることになり、任意の $(m+1)$ 個の冪積の間には $\mathfrak M$ を法として線形従属が存在しなければならないという事実に反する。

逆に、ある加群の剰余群が無限であれば、同じ方法により、無限に多くの線形独立な積分が存在することが分かる。したがって、線形独立な積分が $m$ 個しかない場合には、剰余群は有限であり、その位数は上の議論により $m$ に等しい。

この節の最後に、二つの微分式加群について、同型の概念、または同じことで同種であるという概念が、対応する微分方程式系の積分に対して、一変数の場合の Poincaré の種の概念と同じ意味をもつことを示す（Blumberg, loc. cit., Appendix I を参照）。

\textbf{定理 XII.} \emph{二つの微分式加群 $\mathfrak M$ と $\mathfrak N$ が同種であるならば、二つの微分式 $P$ と $Q$ が存在し、$y$ が $\mathfrak M$ のすべての積分を動くとき $P(y)$ は $\mathfrak N$ のすべての積分を動き、$z$ が $\mathfrak N$ のすべての積分を動くとき $Q(z)$ は $\mathfrak M$ のすべての積分を動く。}

仮定により、すなわち \S~9 により、適当に選ばれた二つの多項式 $P$ と $Q$ について
\[
 MQ\equiv0(\mathfrak N),
 \qquad
 PQ\equiv1(\mathfrak N),
\]
\[
 NP\equiv0(\mathfrak M),
 \qquad
 QP\equiv1(\mathfrak M)
\]
が成り立つ。したがって $NP(y)=0$ であるから、関数 $P(y)$ は $\mathfrak N$ の積分である。同様に $Q(z)$ は $\mathfrak M$ の積分である。$PQ(z)=z$ および $QP(y)=y$ であるから、$P(y)$ は $\mathfrak N$ のすべての積分を動き、$Q(z)$ は $\mathfrak M$ のすべての積分を動く。このようにして、零でない各 $y$ には零でない $z$ が対応し、逆も成り立つ。

\subsection*{\S~12. 例.}

1. 第一の例として、係数が積分の対称関数として表される二階の常微分方程式を考える。ここで有理性領域 $\mathsf{P}$ は、積分とその導関数のすべての有理関数から成る。微分方程式の正則点の十分小さい近傍に制限する。加群 $\mathfrak M$ を、$y_1$ と $y_2$ を積分にもつすべての微分式の全体として定義する。ここでも一つの不定元 $\xi$ に関する多項式として書く。すると加群基底は一元であり、$P$ の因子を除いて一意に定まる多項式
\[
 M=
 \begin{vmatrix} y_1&y_2\\ y_1'&y_2'\end{vmatrix}\xi^2
 -\begin{vmatrix} y_1&y_2\\ y_1''&y_2''\end{vmatrix}\xi
 +\begin{vmatrix} y_1'&y_2'\\ y_1''&y_2''\end{vmatrix}
\]
によって与えられる。いま
\begin{equation}
 \mathfrak M=[\mathfrak M_1,
\mathfrak M_2],
\tag{41}
\end{equation}
ただし $\mathfrak M_i$ は積分 $y_i$ をもつすべての微分式からなり、その基底は $\mathsf{P}$ の因子を除いて
\[
 M_i=y_i\xi-y_i'
 \qquad(i=1,2)
\]
で定まる。実際、$\mathfrak M$ は $y_1$ および $y_2$ で消えるから、$\mathfrak M_1$ と $\mathfrak M_2$ で割り切れ、したがって $[\mathfrak M_1,\mathfrak M_2]$ でも割り切れる。また、$[\mathfrak M_1,\mathfrak M_2]$ の基底多項式の位数は少なくとも 2 でなければならないから、$\mathfrak M=[\mathfrak M_1,\mathfrak M_2]$ である。さらに、$\mathfrak M_1$ と $\mathfrak M_2$ は相対的に素である。なぜなら
\begin{equation}
 \frac{y_2}{D}(y_1\xi-y_1')-\frac{y_1}{D}(y_2\xi-y_2')=1,
 \qquad
 D=D(y)=\begin{vmatrix}y_1&y_2\\ y_1'&y_2'\end{vmatrix}
\tag{42}
\end{equation}
だからである。これらは剰余群が位数 1 をもつので、素加群でもある。

しかし $\mathfrak M$ は、$y_1$ と $y_2$ のほかに、すべての積分
\[
 z_1=\alpha_{11}y_1+\alpha_{12}y_2,
 \qquad
 z_2=\alpha_{21}y_1+\alpha_{22}y_2,
 \qquad
 A=|\alpha_{ik}|\ne0
\]
ももつ。従って、それぞれ積分 $z_1$ と $z_2$ をもち、基底多項式
\[
 N_i=z_i\xi-z_i'=(\alpha_{i1}y_1+\alpha_{i2}y_2)\xi-(\alpha_{i1}y_1'+\alpha_{i2}y_2')
\]
をもつ二つの相対的に素な加群 $\mathfrak N_1$ と $\mathfrak N_2$ の最小公倍元と同一である。したがって $\mathfrak M$ は、相対的に素な素加群の最小公倍元として無限に多くの異なる表示を許す。定理 VII により、二つの加群 $\mathfrak M_1$ と $\mathfrak M_2$ は、互いに、またすべての $\mathfrak N_1$ および $\mathfrak N_2$ と同種である。実際、これらすべての加群の剰余群は位数 1 をもち、従って同型である。

次に、表示 (41) に対応する剰余群 $\mathfrak G=\mathfrak A_1+\mathfrak A_2$ の加法分解を書く。このため、\S~4 と同様に、$\mathfrak M_1$ と $\mathfrak M_2$ の相対的素性を表す関係 (42) から出発する。単位 $\mathfrak{a}_1$ と $\mathfrak{a}_2$ として、次の多項式によって生成される剰余類を得る。
\[
 -\frac{y_1}{D}(y_2\xi-y_2')
 \quad\text{および}\quad
 \frac{y_2}{D}(y_1\xi-y_1').
\]
したがって
\[
 \mathfrak M_1=\left(\mathfrak M,\frac{y_2}{D}(y_1\xi-y_1')\right),
 \qquad
 \mathfrak M_2=\left(\mathfrak M,-\frac{y_1}{D}(y_2\xi-y_2')\right).
\]
いま同じ仕方で、$y_i$ を $z_i$ に置き換え、分解 $\mathfrak M=[\mathfrak N_1,\mathfrak N_2]$ に対応する単位 $b_1$ と $b_2$ を作ると、
\[
 -\frac{z_1}{D(z)}(z_2\xi-z_2')
 =-\frac{\alpha_{11}y_1+\alpha_{12}y_2}{AD(y)}
 ((\alpha_{21}y_1+\alpha_{22}y_2)\xi-(\alpha_{21}y_1'+\alpha_{22}y_2')),
\]
\[
 \frac{z_2}{D(z)}(z_1\xi-z_1')
 =\frac{\alpha_{21}y_1+\alpha_{22}y_2}{AD(z)}
 ((\alpha_{11}y_1+\alpha_{12}y_2)\xi-(\alpha_{11}y_1'+\alpha_{12}y_2')),
\]
従って
\begin{equation}
\left\{
\begin{aligned}
 \mathfrak{b}_1&=\frac{\alpha_{11}y_1+\alpha_{12}y_2}{y_1}
 \left(\frac{\alpha_{22}}{A}\right)\mathfrak{a}_1
 +\frac{\alpha_{11}y_1+\alpha_{12}y_2}{y_2}
 \left(-\frac{\alpha_{21}}{A}\right)\mathfrak{a}_2
 =\frac{z_1}{y_1}\alpha^{11}\cdot \mathfrak{a}_1+\frac{z_1}{y_2}\alpha^{12}\cdot \mathfrak{a}_2,\\
 \mathfrak{b}_2&=\frac{\alpha_{21}y_1+\alpha_{22}y_2}{y_1}
 \left(-\frac{\alpha_{12}}{A}\right)\mathfrak{a}_1
 +\frac{\alpha_{21}y_1+\alpha_{22}y_2}{y_2}
 \left(\frac{\alpha_{11}}{A}\right)\mathfrak{a}_2
 =\frac{z_2}{y_1}\alpha^{21}\mathfrak{a}_1+\frac{z_2}{y_2}\alpha^{22}\mathfrak{a}_2.
\end{aligned}\right.
\tag{43}
\end{equation}
を得る。ここで $(\alpha^{ik})$ は $(\alpha_{ik})$ の逆行列である。

逆に、$z_i$ と $y_i$ を入れ替えると
\begin{equation}
\left\{
\begin{aligned}
 \mathfrak{a}_1&=\frac{y_1}{z_1}\alpha_{11}\mathfrak{b}_1+\frac{y_1}{z_2}\alpha_{12}\mathfrak{b}_2,\\
 \mathfrak{a}_2&=\frac{y_2}{z_1}\alpha_{21}\mathfrak{b}_1+\frac{y_2}{z_2}\alpha_{22}\mathfrak{b}_2.
\end{aligned}\right.
\tag{44}
\end{equation}
これらの公式は群 $\mathfrak A$ と $\mathfrak B$ の同型を示している。すなわち、一般論により、$\mathfrak{a}_1\mathfrak{b}_\sigma\ne0$ ならば対応する割当ては $\mathfrak{a}_1\sim \mathfrak{a}_1\mathfrak{b}_\sigma$ であり、従って (44) により
\[
 \mathfrak{a}_1\sim\frac{y_1}{z_\sigma}\alpha_{1\sigma}\mathfrak{b}_\sigma,
 \qquad(\alpha_{1\sigma}\ne0\text{ のとき}),
\]
また (43) により同様に
\[
 \mathfrak{b}_\sigma\sim\frac{z_\sigma}{y_2}\alpha^{\sigma2}\mathfrak{a}_2,
 \qquad(\alpha^{\sigma2}\ne0\text{ のとき}),
\]
であるから、
\[
 \frac{y_1}{z_\sigma}\alpha_{1\sigma}\mathfrak{b}_\sigma
 \sim
 \frac{y_1}{z_\sigma}\alpha_{1\sigma}\frac{z_\sigma}{y_2}\alpha^{\sigma2}\mathfrak{a}_2
 =\frac{y_1}{y_2}\alpha_{1\sigma}\alpha^{\sigma2}\mathfrak{a}_2.
\]
よって
\[
 \mathfrak{a}_1\sim\alpha_{1\sigma}\alpha^{\sigma2}\frac{y_1}{y_2}\mathfrak{a}_2
 \quad\text{ のとき}\alpha_{1\sigma}\alpha^{\sigma2}\ne0,
 \qquad
 \frac{1}{\alpha_{1\sigma}\alpha^{\sigma2}}\frac{y_2}{y_1}\mathfrak{a}_1\sim \mathfrak{a}_2,
\]
となり、これが逆対応を与える。$(\alpha_{i\varkappa})$ が対角行列でないならば、この条件 $\alpha_{1\sigma}\alpha^{\sigma2}\ne0$ は常にある $\sigma$ について満たされる。対角行列の場合は除かなければならない。なぜならそのとき加群は同一になるからである。たとえば $\alpha_{12}=0$ だけであれば、$\mathfrak A_1=\mathfrak B_1$ であるが、$\mathfrak A_2=\mathfrak B_2$ ではない。したがって $\mathfrak A_1+\mathfrak A_2=\mathfrak B_1+\mathfrak B_2$ および $\mathfrak A_1=\mathfrak B_1$ から $\mathfrak A_2=\mathfrak B_2$ は従わない（注 12 参照）。このとき三つの加群 $\mathfrak M_1$, $\mathfrak M_2$, $\mathfrak N_2$ は互いに二つずつ相対的に素であるが、完全に相対的素な系をなさない。なぜなら $[\mathfrak M_1,\mathfrak M_2]$ は $\mathfrak N_2$ で割り切れるからである（注 15 参照）。

さらに
\[
 \mathfrak{t}_{12}=\alpha_{1\sigma}\alpha^{\sigma2}\frac{y_1}{y_2}\mathfrak{a}_2,
 \qquad
 \mathfrak{t}_{21}=\frac{1}{\alpha_{1\sigma}\alpha^{\sigma2}}\frac{y_2}{y_1}\mathfrak{a}_1
\]
と置けば、関係 (35) が満たされることを直接計算で確かめられる。

2. 基底が一つの多項式から成る二つ以上の変数の任意の加群は、無限群の例を与える。次のやや複雑な例は、無限群が多元加群、すなわち基底が一つより多い多項式を含む加群にも現れることを示すためのものである。\footnote{Blumberg によって伝えられた Landau の例（loc. cit., pp. 51--52）を参照。ここおよび以下では独立変数を $x$ および $y$ で表す。}
\[
 \mathfrak M=[(\xi+x\eta),(\xi+1)]=[(P),(Q)]
 \qquad\left(\xi=\frac{\partial}{\partial x},\ \eta=\frac{\partial}{\partial y}\right).
\]
乗法規則は微分式のそれである。この剰余群は無限である。なぜなら $(\xi+x\eta)$ と $(\xi+1)$ の剰余群はいずれも無限であり、一方で加群 $(\xi+x\eta)$ と $(\xi+1)$ は相対的に素だからである。実際、
\[
 1=(-x\xi-x^2\eta+1)Q+(x\xi+x-1)P.
\]
いま $\mathfrak M$ が二次の多項式をもたないことは容易に示せる。すなわち
\[
 (u_1\xi+u_2\eta+u_3)(\xi+x\eta)=(v_1\xi+v_2\eta+v_3)(\xi+1)
\]
と置くと、$u_1=u_2=u_3=v_1=v_2=v_3=0$ が得られる。次に $\mathfrak M$ の中の三次多項式を探すと、線形独立なものが二つ見つかる。たとえば
\[
 (\xi^2+x\xi\eta+\xi+(2+x)\eta)Q=(\xi^2+2\xi+1)P
\]
および
\[
 (\xi^2+2x\xi\eta+x^2\eta^2+\eta)Q
 =(\xi^2+x\xi\eta+\xi+(x-1)\eta)P.
\]
もし $\mathfrak M$ が一元加群であったならば、両者は基底多項式で割り切れなければならず、しかも $\mathfrak M$ は二次多項式を含まないので、三次多項式で割り切れなければならない。従って両者は比例しなければならないが、実際には比例しない。$\mathfrak M$ が多元であることが、一変数の場合には有効な Landau の積定理が失敗する内的理由である。

3. さらに別の例は、無限群が素加群にも現れ得ることを示す。\footnote{我々の「素加群」の定義は、可換の場合に素加群と呼ばれるものとは異なる。可換の場合には、剰余群が有限であるとき、すなわち次元が $=0$ であるときを除いて、常により低次元の因子が存在するからである。} $P$ を二変数 $x,y$ のすべての有理関数の領域とする。基底
\[
 \xi-\frac{y}{x}=\xi-a
\]
をもつ加群 $\mathfrak M$ は無限群をもつ。なぜなら、この群は $y$ のすべての冪の剰余類を含むからである。他方、$\mathfrak M$ は素加群である。実際、任意の因子
\[
 \mathfrak M_1=(\xi-a,
F(\xi,\eta))
\]
が単位加群に等しいことを示す。任意の多項式 $F(\xi,\eta)$ は、$(\xi-a)$ を法として多項式 $G(\eta)$ によって表すことができる。
\[
 F(\xi,
\eta)\equiv c_0\eta^\nu+c_1\eta^{\nu-1}+\cdots+c_\nu\quad(\mathfrak M).
\]
一般性を失わず $c_0=1$ とする。乗法法則により $\xi\eta^\nu-\eta^\nu\xi=0$ である。従って
\[
 \xi\equiv a(\mathfrak M_1),
 \qquad
 \eta^\nu\equiv F_1(\mathfrak M_1)
\]
であり、ここで $F_1=-(c_1\eta^{\nu-1}+\cdots+c_\nu)$ とすると、
\[
 \xi F_1-\eta^\nu a\equiv0(\mathfrak M_1)
\]
である。
\[
 \eta^\nu a=a\eta^\nu+\nu\frac{\partial a}{\partial y}\eta^{\nu-1}+\cdots
\]
と展開し、再び $\eta^\nu$ を $F_1$ で置き換えると、
\[
 \xi(c_1\eta^{\nu-1}+\cdots+c_\nu)-a(c_1\eta^{\nu-1}+\cdots+c_\nu)
 +\nu\frac{\partial a}{\partial y}\eta^{\nu-1}+\cdots\equiv0(\mathfrak M_1),
\]
すなわち
\[
 c_1\eta^{\nu-1}a+\frac{\partial c_1}{\partial x}\eta^{\nu-1}+\cdots
 -ac_1\eta^{\nu-1}+\cdots+
\nu\frac{\partial a}{\partial y}\eta^{\nu-1}+\cdots
 \equiv0(\mathfrak M_1),
\]
また最後に
\[
 \left(\frac{\partial c_1}{\partial x}+\nu\frac{\partial a}{\partial y}\right)\eta^{\nu-1}+\cdots
 \equiv0(\mathfrak M_1)
\]
を得る。左辺には、いま $\eta$ に関する正確に次数 $\nu-1$ の多項式があり、これは恒等的には消えない。なぜなら最高係数
\[
 \frac{\partial c_1}{\partial x}+\nu\frac{\partial a}{\partial y}
 =\frac{\partial c_1}{\partial x}+\nu\frac{1}{x}
 \quad\text{は }\ne0,
\]
だからである。もしそうでなければ $c_1=-\nu\log x+\varphi(y)$ となり、有理関数ではなくなる。このように手続きを続けることができ、恒等的には零でない零次多項式へ到達する。これにより単位も $\mathfrak M_1$ に含まれることが証明される。\footnote{この手続きは、微分方程式系に必要な可積分条件が満たされていないことを証明することに相当する。}

4. 他方、有理性領域に関数 $\log x$ を付け加えると、$\mathfrak M$ は因子
\[
 \left(\xi-\frac{y}{x},\ \eta-\log x+\varphi(y)\right)
\]
をもつ。ただし $\varphi(y)$ は $y$ のすべての有理関数を動く。これらの因子はすべて異なり、それぞれについて可積分条件は満たされ、対応する積分は
\[
 y\log x-\int\varphi(y)\,dy+\alpha
\]
である。従って、単位はこの加群に含まれ得ない。そうでなければ、零関数だけが可能な積分となってしまうからである。他方、$\xi$ と $\eta$ は、この加群に関して、領域のある量に合同である。従って剰余群は有限で、位数 1 をもつ。さらに、無限に多くの因子の任意の二つは相対的に素であり、与えられた加群 $\mathfrak M$ はこれらすべての因子の最小公倍元 $\mathfrak N$ に等しい。実際、それはすべての因子で割り切れるので、$\mathfrak N$ でも割り切れる。もし $\mathfrak N$ が $\mathfrak M$ の真の因子であれば、$\mathfrak N$ はなお少なくとも一つの多項式 $F$、たとえば次数 $\rho$ で $\eta$ のみに依存するものを含まなければならない。$\varphi(y)$ に関数 $y,\ldots,y^{\rho+1}$ を動かさせると、対応する積分
\[
 y\log x-\frac{y^{i+1}}{i+1}+c_i
 \qquad(i=1,
\ldots,
\rho+1)
\]
の間には、$y$ に依存しない係数による線形関係は存在しない。しかし次数 $\rho$ の微分方程式 $F=0$ が $\rho+1$ 個の線形独立な積分をもつことはできない。

とはいえ、加群 $\mathfrak M=\mathfrak N$ は完全可約ではない。\footnote{従って、これは定理 X の第一の場合の例である。} もし完全可約ならば、それはこれらの素加群の有限個の最小公倍元であり、従ってその剰余群は有限位数をもつはずである。しかし $\xi-y/x$ の剰余群は $\eta$ のすべての冪の剰余類を含むので、無限である。

ゲッティンゲン、1919 年 8 月 1 日。

\begin{center}
（1919 年 8 月 4 日受理。）
\end{center}

\clearpage
\setcounter{footnote}{0}
\begin{center}
{\Large\bfseries 18. 戦没した K. Hentzelt の消去論に関する仕事について}\par
\vspace{1em}
\emph{Jahresbericht der DMV} 30 (1921), p. 101
\end{center}

\begin{center}
\textbf{第 11 回会合、1921 年 9 月 23 日金曜日、午前 11 時}\\
（議長：Loewy.）
\end{center}

\textbf{1. E. Noether: 戦没した K. Hentzelt の消去論に関する仕事について.}

Kurt Hentzelt は 1914 年 10 月以来 Dixmuiden の前線で行方不明であり、戦死者に数えなければならない。彼が Erlangen の学位論文として残したものは、内容上は欠けるところなく組み立てられているが、提示はほとんど理解不能である。私はその本質的部分を、近く \emph{Mathematische Annalen} に自由な編集で公刊するつもりである。

内容は、\textbf{消去論}の一般問題、すなわち多項式から成るイデアル $\mathfrak M$ のすべての多項式の共通零点を決定する問題に関わる。まず変数に不定係数をもつ線形変換を施すならば、主結果は次のように述べられる。

イデアル $\mathfrak M$ には、一意に終結式形式
\[
 R_{\mathfrak M}=R^{(1)}(x_1,
\ldots,x_n)\cdots R^{(i)}(x_i,
\ldots,x_n)\cdots R^{(n)}(x_n)\equiv0(\mathfrak M)
\]
を対応させることができる。この $R_{\mathfrak M}$ は $\mathfrak M$ のすべての零点で、かつそれらに限って零になる。$\mathfrak N$ が $\mathfrak M$ の因子であり、$R_{\mathfrak N}$ と $R_{\mathfrak M}$ が一致するならば、$\mathfrak N$ と $\mathfrak M$ 自身もまた一致する。

定理の第二部は、終結式形式が零点を特性的重複度つきで与えることを示している。それは、$\mathfrak M$ を $x_1,\ldots,x_{i-1}$ の冪積に関する線形形式の加群と見なすことができる、という事実に基づく。そのとき $R_{\mathfrak M}$ の個々の因子 $R^{(i)}(x_i,\ldots,x_n)$ は、$x_{i+1},\ldots,x_n$ を係数領域に付け加えたとき、この加群に関するそれぞれの「基礎加群」のノルムになる。抽象イデアル論を用いると、次の解釈が得られる。もし
\[
 \mathfrak M=[\mathfrak Q,\mathfrak Q_1,
\ldots,
\mathfrak Q_r]=[\mathfrak Q,\mathfrak L]
\]
が $\mathfrak M$ の準素イデアル $\mathfrak Q_i$ への分解であり、$\mathfrak L$ が $\mathfrak Q$ の補元であるならば、イデアル $\mathfrak Q$ には、$R^{(i)}(x_i,\ldots,x_n)$ の準素因子 $Q^{(i)}(x_i,\ldots,x_n)$ が対応する。この因子は、上に述べた意味で、$\mathfrak Q$ に関する $\mathfrak L$ のノルムを表す。



\clearpage

\section*{19. 環領域におけるイデアル論}
\begin{center}
Math. Ann. 83 (1921), pp. 24--66
\end{center}

\begin{center}\textbf{目次}\end{center}
\noindent 序論。\\
\S~1. 環領域、イデアル、有限性条件。\\
\S~2. 有限個の既約イデアルの最小公倍数としてのイデアルの表示。\\
\S~3. 既約イデアルへの二つの異なる分解における成分数の一致。\\
\S~4. 準素イデアル。既約イデアルへの二つの異なる分解における随伴素イデアルの一意性。\\
\S~5. 極大準素イデアルの最小公倍数としてのイデアルの表示。随伴素イデアルの一意性。\\
\S~6. 相対的に素な既約イデアルの最小公倍数としてのイデアルの一意表示。\\
\S~7. 孤立イデアルの一意性。\\
\S~8. 互いに素な既約イデアルの積としてのイデアルの一意表示。\\
\S~9. 研究の加群への拡張。既約加群への分解における成分数の一致。\\
\S~10. 多項式領域という特殊な場合。\\
\S~11. 数論および微分式論からの例。\\
\S~12. 初等因子論からの例。

\subsection*{序論。}
本論文の内容は、\emph{有理整数、あるいは代数的数体におけるイデアルについての分解定理を、任意の整域、さらに一般には環領域におけるイデアルへ移すこと}である。この移行を理解するために、まず有理整数の分解定理を、通常の定式化とは少し異なる形で述べておく。

\[
 a=p_1^{\varrho_1}p_2^{\varrho_2}\cdots p_\sigma^{\varrho_\sigma}=q_1q_2\cdots q_\sigma
\]
において、素数冪因子 $q_i$ を分解の成分と見なすと、これらの成分は次の特徴的性質をもつ。

1. それらは\emph{二つずつ互いに素}である。しかし、どの $q$ も二つずつ互いに素な数の積として表すことはできず、この意味で既約性が成り立つ。また二つずつ互いに素であることから、積 $q_1\cdots q_\sigma$ が最小公倍数 $[q_1\cdots q_\sigma]$ に等しいことが従う。

2. 任意の二つの成分 $q_i$ と $q_k$ は\emph{相対的に素}である。すなわち、$b\cdot q_i$ が $q_k$ で割り切れるならば、$b$ も $q_k$ で割り切れる。この意味でも既約性が成り立つ。

3. 各 $q$ は\emph{準素}である。すなわち、積 $b\cdot c$ が $q$ で割り切れ、しかも $b$ が割り切れないならば、$c$ のある冪\footnote{この冪が常に第一冪でよい場合には、よく知られているように素数を扱っていることになる。}が割り切れる。この表示はさらに\emph{極大準素成分}による表示である。なぜなら、異なる二つの $q$ の積はもはや準素でないからである。極大準素成分への分解という観点からも、$q$ は既約である。

4. 各 $q$ は、二つの真の約数の最小公倍数として表すことができない、という意味で\emph{既約}である。

これらの準素数 $q$ と素数 $p$ との関係は次の点にある。各 $q$ に対して、$q$ の約数であり、かつそのある冪が $q$ で割り切れるような $p$ が一つ存在し、符号を除けば一つしか存在しない。これが随伴素数である。$p^\varrho$ がこの種の最小冪であるとき、$\varrho$ を $q$ の指数というが、ここでは特に $p^\varrho=q$ である。したがって一意性定理は次のように述べられる。

\emph{有理整数を既約な極大準素成分 $q$ に分解する二つの異なる分解では、成分の個数、随伴素数（符号を除いて）、および指数が一致する。$p^\varrho=q$ であるから、これにより $q$ 自身も（符号を除いて）一致する。}

符号によって生じる不定性は、よく知られているように、数そのものの代わりにそれから得られるイデアル（$a$ で割り切れるすべての数）を考えれば取り除かれる。そのとき同じ定式化は、有限代数的数体のイデアルを素イデアル冪に一意分解する場合にそのまま成り立つ。

以下では（\S~1）、各イデアルが有限個のイデアル基底をもつ、という有限性条件だけを満たす一般の環領域を基礎に置く。このような有限性条件がなければ、既約イデアルも素イデアルも存在するとは限らない。実際、すべての代数的整数からなる領域では、素イデアルへの分解は存在しない。

明らかになるのは、成分 $q$ の四つの特徴的性質に対応して、一般には四つの別々の分解が存在し、それぞれが前の分解を細分することによって得られる、ということである。互いに素な既約イデアルへの分解は積表示であるが、他の三つの分解は最小公倍数としての既約（\S~2）表示である。準素イデアルと随伴素イデアルとの関係も保たれる。すなわち、既約イデアルも準素であり、各準素イデアル $\ideal Q$ には、それを割り、しかもそのある冪が $\ideal Q$ で割り切れるような随伴素イデアル $\ideal P$ が一意に定まる。$\ideal P^\varrho$ がこのような最小冪であり、$\varrho$ が $\ideal Q$ の指数であるとしても、ここでは $\ideal P^\varrho$ が $\ideal Q$ と一致するとは限らない。一意性定理は次のように表される。

分解 1 と 2 は一意である。分解 3 または 4 の二つの異なる分解では、成分の個数と随伴素イデアルが一致する。\footnote{おそらく、さらに指数も一致し、より一般には対応する成分が同型であると思われる。} 成分の中に現れる孤立イデアル（\S~7）は一意に定まる。

分解定理の証明では、まず有限性条件から、Dedekind が有限数加群に対して最初に述べた「有限鎖の定理」を導き、そこから各イデアルを有限個の既約イデアルの最小公倍数として表す第 4 の表示を得る。成分の可約性の概念を変形することにより、この表示から、既約イデアルへの第 4 の分解に関する基本的な一意性定理が得られる。ついで、有限個ずつの成分をまとめることにより、残りの分解へ進むことができ、その一意性定理は第 4 の一意性定理の帰結として得られる。

最後に（\S~9）、有限個の既約成分による表示は、より弱い仮定のもとでも成り立つことを示す。環領域の可換性は不要であり、イデアルの代わりにその領域に関する加群を考えれば十分である。このより一般の場合にも、二つの異なる分解における成分数の一致はなお成り立つが、素および準素の概念は可換性とイデアル概念に結びついている。これに対し、互いに素という概念は、非可換領域におけるイデアルについても保たれる。

四つの別々の分解が実際に現れる最も単純な環領域は、任意の複素係数をもつ $n$ 変数多項式全体の領域である。ここでは個々の分解を、代数的図形の挙動によって、非有理的に解釈することができ、随伴素イデアルの一意性定理は、代数的図形が既約図形へ一意に分解されるという消去論の基本定理に対応する。さらなる例は、多項式から作られるすべての有限整域によって与えられる（\S~10）。しかし、偶数全体からなる単純な領域、さらに一般には固定された数で割り切れるすべての数からなる領域でさえ、部分的に分離された分解の例を与える（\S~11）。非可換領域におけるイデアル論の例は初等因子論によって与えられる（\S~12）。そこでは既約イデアル、または既約類への一意分解が成り立つ。これらの既約類は初等因子の既約構成要素を完全に特徴づけるものであり、通常の初等因子論が失敗する領域におけるそれらの代替物と見なせるかもしれない。

既存文献については次のことを述べておく。任意の複素係数、または整数係数をもつ多項式領域における極大準素イデアルへの分解は Lasker によって与えられ、Macaulay によって個々の点でさらに発展させられた。\footnote{E. Lasker, Zur Theorie der Moduln und Ideale. Math. Ann. 60 (1905), p. 20, Theorems VII and XIII. -- F. S. Macaulay, On the Resolution of a given Modular System into Primary Systems including some Properties of Hilbert Numbers. Math. Ann. 74 (1913), p. 66.} 両者は消去論に依拠し、したがって多項式が既約多項式の積として一意に表されるという事実を用いている。実際には、代数的数体におけるイデアル論が示唆し、本論文が示すように、イデアルの分解定理はこの前提から独立である。準素イデアルもまた、Lasker と Macaulay によって消去論の概念を基礎として定義されている。

既約イデアルへの分解、および相対的に素な既約イデアルへの分解は、多項式領域の場合でさえ文献では注目されていなかったようである。Macaulay に、孤立準素イデアルの一意性に関する注意があるだけである。

互いに素な既約イデアルへの分解は、多項式領域について Schmeidler によって与えられた。\footnote{W. Schmeidler, Über Moduln und Gruppen hyperkomplexer Größen. Math. Zeitschr. 3 (1919), p. 29.} その有限性の証明には消去論が用いられている。しかしそこでは、一意性定理はイデアルの類についてだけ述べられており、イデアル自身については述べられていない。後者の一意性定理は共同論文\footnote{E. Noether--W. Schmeidler, Moduln in nichtkommutativen Bereichen, insbesondere aus Differential- und Differenzenausdrücken. Math. Zeitschr. 8 (1920), p. 1.}に現れる。そこでは非可換多項式領域におけるイデアルが関与している。そこでは有限なイデアル基底だけが用いられるので、定理と方法は一般の環領域にも有効であり、本論文では個数の一致（\S~11）についてさらに鋭くされる。本研究は、これら二つの論文の背後にある概念構成の強い一般化であり、さらなる発展である。二つの論文の本質的特徴は、最小公倍数としての表示から、剰余類の系の加法的分解へ移る点にある。本論文では、叙述を単純にするため、最小公倍数にとどまる。加法的分解に対応するものは、可約性の概念を補元の性質へ変換することである（\S~3）。とはいえ、共同論文の考察に本質的に対応する考察によって、ここに述べるすべての定理は、剰余類の系およびある部分系についての加法的分解定理としても理解できる。この剰余類の系は、最初に基礎とされた環と同じ一般性をもつ環をなす。実際、任意の環は、その環の元の間にある恒等関係全体に対応するイデアルの剰余類の系と見なすことができる。また、残りの関係がすでに領域内で満たされていると仮定すれば、それらの関係の部分系の剰余類の系としても見なせる。

この注意は Fraenkel の諸論文の位置づけも与える。\footnote{A. Fraenkel, Über die Teiler der Null und die Zerlegung von Ringen. J. f. M. 145 (1914), p. 139. Über gewisse Teilbereiche und Erweiterungen von Ringen. Habilitationsschrift, Leipzig, Teubner, 1916. Über einfache Erweiterungen zerlegbarer Ringe. J. f. M. 151 (1920), p. 121.} Fraenkel は、正則元の存在、それらによる除法、分解可能性条件といった制限条件を課された環の加法的分解を考える。対応するイデアルについては、そのような条件のもとで四つの分解が一致する。この一致のため、彼の有限性条件、すなわちイデアルが有限個の真の約数だけをもつという条件（彼は一部ではこれから離れる）も、われわれの条件より強い制限を意味しない。Fraenkel の出発点は、彼の諸論文の、異なる、主として代数的な目的によって規定されている。彼はその後、代数的拡大によって、より制限の少ない、より一般の環へ到達する。

\subsection*{\S~1. 環領域、イデアル、有限性条件。}

\textbf{1.} 基礎に置く領域 $\Sigma$ は、抽象的定義における（可換）環であるとする。\footnote{この定義は Fraenkel の教授資格論文から、彼の制限条件 6、I、II を省いて取ったものである。そのため加法の交換法則を含める必要があった。したがってこれは、乗法の可逆性を除いた体の定義法則である。} すなわち、$\Sigma$ は元 $a,b,c,\ldots,f,g,h,\ldots$ の系からなり、通常の条件を満たす関係が等号として定義されている。また二つの演算（結合法）、すなわち加法と乗法によって、任意の二つの環元 $a$ と $b$ から、それぞれ和 $a+b$ および積 $a\cdot b$ として第三の元が常に一意に得られる。環と、その他はまったく任意である演算は、次の法則を満たさなければならない。
\begin{enumerate}
\item 加法の結合法則：$(a+b)+c=a+(b+c)$。
\item 加法の交換法則：$a+b=b+a$。
\item 乗法の結合法則：$(a\cdot b)\cdot c=a\cdot(b\cdot c)$。
\item 乗法の交換法則：$a\cdot b=b\cdot a$。
\item 分配法則：$a(b+c)=ab+ac$。
\item 制限のない一意な減法の法則。方程式 $a+x=b$ を満たすただ一つの元 $x$ が $\Sigma$ に存在する。（$x=b-a$ と書く。）
\end{enumerate}
これらの性質から零元の存在が従う。しかし環は単位元をもつとは限らず、二つの元の積が消えても因子の一方が消えるとは限らない。積の消滅から常に一方の因子の消滅が従い、さらに単位元をもつ環を、真の整域と呼ぶ。有限和 $a+a\ldots+a$ に対しては、通常の省略記法 $na$ を導入する。ここで整数 $n$ は環元ではなく単なる省略記号とみなし、$a=1\cdot a$、$na+a=(n+1)a$ によって帰納的に定義する。

\textbf{2.} $\Sigma$ におけるイデアル $\ideal M$\footnote{イデアルは大文字のドイツ文字で表す。$\ideal M$ は、通常「加群」または形式加群と呼ばれる多項式のイデアルの例を想起させるものである。なお、\S\S~1--3 ではイデアル性ではなく加群性だけを用いる。\S~9 を参照。}とは、$\Sigma$ の元の系であって次の二条件を満たすものをいう。
\begin{enumerate}
\item $\ideal M$ は、$f$ とともに $a\cdot f$ も含む。ただし $a$ は $\Sigma$ の任意の元である。
\item $\ideal M$ は、$f$ と $g$ とともに差 $f-g$ も含む。したがって $f$ とともに任意の整数 $n$ に対する $nf$ も含む。
\end{enumerate}
$f$ が $\ideal M$ の元であることを、通常どおり $f\equiv0(\ideal M)$ と表し、$f$ は $\ideal M$ で割り切れるという。$\ideal N$ のすべての元が同時に $\ideal M$ の元であり、したがって $\ideal M$ で割り切れるとき、$\ideal N$ は $\ideal M$ で割り切れるという。記号では $\ideal N\equiv0(\ideal M)$ と書く。$\ideal M$ が $\ideal N$ とは異なる元を含み、したがって逆に $\ideal N$ で割り切れないとき、$\ideal M$ を $\ideal N$ の真の約数という。$\ideal N\equiv0(\ideal M)$ かつ $\ideal M\equiv0(\ideal N)$ からは $\ideal N=\ideal M$ が従う。

他の通常の概念もそのまま保たれる。二つのイデアル $\ideal A$ と $\ideal B$ の最大公約数、すなわち $\ideal D=(\ideal A,\ideal B)$ とは、$a$ が $\ideal A$ のすべての元を動き、$b$ が $\ideal B$ のすべての元を動くとき、$a+b$ の形で表される元全体をいう。$\ideal D$ は再びイデアルである。同様に、無限個のイデアルの最大公約数 $\ideal D=(\ideal A_1,\ideal A_2,\ldots,\ideal A_\nu,\ldots)$ は、毎回有限個のイデアルから取った元の和 $d=a_{i_1}+a_{i_2}+\cdots+a_{i_n}$ として表される元 $d$ 全体として定義される。この場合も $\ideal D$ は再びイデアルである。

特に、イデアル $\ideal M$ が有限個の元 $f_1,f_2,\ldots,f_\varrho$ を含み、
\[
 \ideal M=(f_1\ldots f_\varrho),\qquad
 f=a_1f_1+\cdots+a_\varrho f_\varrho+n_1f_1+\cdots+n_\varrho f_\varrho
\]
がすべての $f\equiv0(\ideal M)$ について成り立つとする。ここで $a_i$ は環領域の量、$n_i$ は整数である。このとき $\ideal M$ を有限イデアルといい、$f_1,\ldots,f_\varrho$ をイデアル基底という。

以下では、有限性条件、すなわち\emph{$\Sigma$ のすべてのイデアルが有限であり、したがってイデアル基底をもつ}という条件を満たす環 $\Sigma$ だけを基礎に置く。

\textbf{3.} 有限性条件から、以後のすべての考察の基礎となる結果がただちに従う。

\textbf{定理 I（有限鎖の定理）。}\footnote{数加群について Dedekind が最初に述べたもの：Zahlentheorie, Suppl. XI, \S~172, Theorem VIII（第 4 版）。われわれの証明と「鎖」という名称はそこから取った。多項式におけるイデアルについては Lasker, loc. cit., p. 56（補題）を参照。しかしいずれの場合にも、この定理は孤立した適用しか受けていない。われわれの適用は全体として選択原理に基づいている。} \emph{$\ideal M,\ideal M_1,\ideal M_2,\ldots,\ideal M_\nu,\ldots$ を $\Sigma$ における可算無限個のイデアルの系とし、各イデアルが次のものによって割り切れるとする。このとき、ある有限の添字 $n$ 以後、すべてのイデアルは同一であり、$\ideal M_n=\ideal M_{n+1}=\cdots$ となる。言い換えれば、$\ideal M,\ideal M_1,\ideal M_2,\ldots,\ideal M_\nu,\ldots$ が単純に順序づけられたイデアルの鎖をなし、各イデアルが直前のものの真の約数であるならば、この鎖は有限個の段階の後に停止する。}

実際、$\ideal D=(\ideal M,\ideal M_1,\ideal M_2,\ldots,\ideal M_\nu,\ldots)$ をこの系の最大公約数とし、有限性条件により存在する $\ideal D$ の基底を $f_1,\ldots,f_\varrho$ とする。割り切りの仮定から、$\ideal D$ の各元は同時に鎖の一つのイデアルの元であることが従う。なぜなら
\[
 f=g+h,\qquad g\equiv0(\ideal M_r),\qquad h\equiv0(\ideal M_s),\qquad (r\le s)
\]
ならば $g\equiv0(\ideal M_s)$ であり、したがって $f\equiv0(\ideal M_s)$ であるからである。$f$ がさらに多くの項の和である場合も同じである。したがって、有限の添字 $n$ が存在して
\[
 f_1\equiv0(\ideal M_n),\quad f_2\equiv0(\ideal M_n),\quad\ldots,\quad f_\varrho\equiv0(\ideal M_n)
\]
となる。逆に $\ideal M_n\equiv0(\ideal D)$ であるから $\ideal M_n=\ideal D$ であり、さらに $\ideal M_n\equiv0(\ideal D)$ および $\ideal D=\ideal M_n\equiv0(\ideal M_r)$ から、すべての $r>n$ について $\ideal M_r=\ideal D=\ideal M_n$ である。これで定理が証明された。

逆に、この定理からイデアル基底の存在も再び従うので、有限性条件は基底を用いない形で述べることもできたことを注意しておく。

\subsection*{\S~2. 有限個の既約イデアルの最小公倍数としてのイデアルの表示。}

イデアル $\ideal B_1,\ideal B_2,\ldots,\ideal B_k$ の最小公倍数 $[\ideal B_1,\ideal B_2,\ldots,\ideal B_k]$ は、通常どおり、$\ideal B_1$、$\ideal B_2$、$\ldots$、$\ideal B_k$ のすべてで割り切れる元全体として定義される。記号では
\[
 f\equiv0(\ideal B_i),\quad (i=1,2,\ldots,k),
 \quad\text{ならば}\quad
 f\equiv0([\ideal B_1,\ideal B_2,\ldots,\ideal B_k]),
\]
およびその逆が成り立つ、ということである。最小公倍数は再びイデアルであり、イデアル $\ideal B_i$ は分解の成分とも呼ばれる。

\textbf{定義 I.} 表示 $\ideal M=[\ideal B_1\cdots\ideal B_k]$ は、どの $\ideal B_i$ も他のイデアルの最小公倍数 $\ideal A_i$ に含まれず、かつどの $\ideal B_i$ も真の約数で置き換えることができないとき、既約表示と呼ばれる。\footnote{既約でない表示の例は、すべての指数 $\lambda\ge2$ について $(x,xy)=[(x),(x^2,xy,y^\lambda)]$ である。$\lambda=1$ に対応する表示 $[(x),(x^2,y)]$ はそれに随伴する既約表示である。} 条件がイデアル $\ideal B_i$ についてだけ満たされるとき、その表示は $\ideal B_i$ に関して既約であるという。最小公倍数 $\ideal A_i=[\ideal B_1,\ldots,\ideal B_{i-1},\ideal B_{i+1},\ldots,\ideal B_k]$ を $\ideal B_i$ の補元という。第一条件だけを満たす表示を最短表示という。

したがって、イデアルを最小公倍数として表すには、次の結果により、既約表示に限定すれば十分である。

\textbf{補題 I.} \emph{有限個のイデアルの最小公倍数としてのイデアルの任意の表示は、少なくとも一つの仕方で既約表示に置き換えられる。特に、そのような表示は逐次分解によって得ることができる。}

実際、$\ideal M=[\ideal B_1^*\cdots\ideal B_l^*]$ を $\ideal M$ の任意の表示とする。残されたイデアルの最小公倍数に含まれる $\ideal B_i$ を順に省く。残りのイデアルがなお $\ideal M$ を与えるので、得られた表示
\[
 \ideal M=[\ideal B_1\cdots\ideal B_k]=[\ideal A_i,\ideal B_i]
\]
では第一条件が満たされる。したがってこれは最短表示であり、この条件はどの $\ideal B_i$ を真の約数で置き換えても満たされ続ける。しかし第二条件は、有限鎖の定理（定理 I）によって常に達成できる。すなわち
\[
 \ideal B_i\supset \ideal B_i'\supset \ideal B_i''\supset\cdots\supset \ideal B_i^{(\nu)}\supset\cdots,
 \qquad
 \ideal M=[\ideal A_i,\ideal B_i]=[\ideal A_i,\ideal B_i']=\cdots=[\ideal A_i,\ideal B_i^{(\nu)}]=\cdots,
\]
であり、各 $\ideal B_i^{(\nu)}$ が直前のものの真の約数であるならば、その定理により鎖 $\ideal B_i,\ideal B_i',\ldots,\ideal B_i^{(\nu)},\ldots$ は有限段階で終わらなければならない。したがって表示 $\ideal M=[\ideal A_i,\ideal B_i^{(n)}]$ では、$\ideal B_i^{(\nu)}$ はもはや真の約数で置き換えられない。これは、$\ideal A_i$ を真の約数で置き換えた場合にはなおさら成り立つ。この手続きを各 $\ideal B_i$ に順に適用し、そのたびにすでに既約化された $\ideal B$ とともに補元を作れば、既約表示が得られる。\footnote{そのような表示が与えられた表示によって一意に定まるわけではないことは、前の例が示す。$(x^2,xy)=[(x),(x^2,xy,y^\lambda)]$、$\lambda\ge2$ に対して、$[(x),(x^2,y)]$ のほかに、任意の $u$ について $[(x),(x^2,\mu x+y)]$ も既約表示である。}

このような表示を逐次的に得るためには、個々の既約表示
\[
 \ideal M=[\ideal B_1,\ideal C_1],\qquad
 \ideal C_1=[\ideal B_2,\ideal C_2],\quad\ldots,
 \quad \ideal C_{l-1}=[\ideal B_l,\ideal C_l]
\]
から、結果として得られる表示 $\ideal M=[\ideal B_1\cdots\ideal B_k,\ideal C_l]$ も既約であることを示せばよい。そのためには、既約表示
\[
 \ideal M=[\ideal B,\ideal C],\qquad \ideal C=[\ideal C_1,\ideal C_2]
\]
から $\ideal M=[\ideal B,\ideal C_1,\ideal C_2]$ が既約表示になることを示せば十分である。実際、仮定により $\ideal B$ はその補元に入らない。もしある $\ideal C_i$ についてこれが起これば、第二の既約表示により $\ideal C_1$ と $\ideal C_2$ は $\ideal C$ の真の約数であるから、第一の表示の仮定に反して $\ideal C$ を真の約数で置き換えられることになる。したがって表示は最短である。さらに仮定により $\ideal B$ は真の約数で置き換えられない。もしある $\ideal C_i$ についてこれが可能であれば、$\ideal C$ の表示が既約であることから、仮定に反して $\ideal C$ を真の約数で置き換えることに対応する。これで補題が証明された。

\textbf{定義 II.} イデアル $\ideal M$ が二つの真の約数の最小公倍数として表せるとき、$\ideal M$ を可約という。そうでない場合、$\ideal M$ を既約という。

有限鎖についての定理 I を用い、既約表示を使って、次を証明する。

\textbf{定理 II.} \emph{すべてのイデアルは、有限個の既約イデアルの最小公倍数として表すことができる。}\footnote{このような表示が一般には一意でないことは前の例が示す。任意の $u$ について $(x^2,xy)=[(x),(x^2,\mu x+y)]$ である。二つの成分はいずれも、任意の $\mu$ について既約である。実際、$(x)$ のすべての約数は $(x,g(y))$ の形であり、ここで $g(y)$ は $y$ の多項式を表す。したがって二つの最小公倍数もこの形であり、ゆえに $(x)$ の真の約数である。イデアル $(x^2,\mu x+y)$ は約数 $(x,y)$ を一つだけもつので、やはり必然的に既約である。}

任意のイデアル $\ideal M$ は既約であるか、または可約である。既約ならば $\ideal M=[\ideal M]$ は定理 II が要求する型の表示である。可約ならば $\ideal M=[\ideal B_1,\ideal C_1]$ であり、$\ideal B_1$ と $\ideal C_1$ は $\ideal M$ の真の約数であって、その表示は補題 I によって既約と仮定できる。$\ideal C_1$ についても同じ選択が成り立つ。すなわち、$\ideal C_1$ は既約であるか、または既約表示 $\ideal C_1=[\ideal B_2,\ideal C_2]$ をもつ。

このように進めると、既約表示の列
\begin{equation}
 \ideal M=[\ideal B_1,\ideal C_1]=[\ideal B_1,\ideal B_2,\ideal C_2]=\cdots=[\ideal B_1,\ldots,\ideal B_n,\ideal C_n]=\cdots
\tag{1}
\end{equation}
が得られる。鎖 $\ideal C_1,\ideal C_2,\ldots,\ideal C_\nu,\ldots$ において、各 $\ideal C_i$ は直前のものの真の約数であり、したがって鎖は有限段階で終わる。すなわち、ある添字 $n$ があって $\ideal C_n$ は既約となる。さらに補題 I により表示 $\ideal M=[\ideal B_1\cdots\ideal B_n\ideal C_n]$ は既約である。したがって $\ideal C_n$ は補元 $\ideal A_n$ に入らず、表示 $\ideal M=[\ideal A_n,\ideal C_n]$ において $\ideal U_n$ を真の約数で置き換えることもできない。必要ならば $\ideal A_n$ を真の約数で置き換えて\footnote{実際には、この表示は $\ideal A_n$ に関しても既約であることが、\S~3（補題 IV）で補題 I の逆として示される。}表示を既約にすれば、すべての可約イデアルが、既約イデアルとそれに補完するイデアルとの最小公倍数として既約表示をもつことが示される。したがって列 (1) では、一般性を失わずにすべての $\ideal B_i$ を既約と仮定できる。上の議論を繰り返せば、既約な $\ideal C_n$ の存在が得られ、定理 II が証明される。

\subsection*{\S~3. 既約イデアルへの二つの異なる分解における成分数の一致。}

成分数の一致を証明するため、まずイデアルの可約性または既約性を、その補元の性質によって表す必要がある。

\textbf{定理 III。}\footnote{定理 III は、Schmeidler および Noether--Schmeidler の論文における加群から剰余群への移行に対応する（序論を参照）。剰余群に対応するのが $\ideal A$ であり、$\ideal N_1$ と $\ideal N_2$ は剰余群が分解される部分群に対応する。} \emph{最短表示 $\ideal M=[\ideal A,\ideal C]$ が $\ideal C$ に関して既約であるとする。このとき $\ideal C$ が可約であるための必要十分条件は、$\ideal M$ の真の約数である二つのイデアル $\ideal N_1$、$\ideal N_2$ が存在して}
\begin{equation}
 \ideal N_1\ideal N_2\equiv0(\ideal M),\qquad
 \ideal N_i\equiv0(\ideal U),\qquad
 [\ideal N_1,\ideal N_2]=\ideal M
\tag{2}
\end{equation}
\emph{を満たすことである。さらに次も従う。(2) が満たされ、$\ideal C$ が既約であるならば、少なくとも一つの $\ideal N_i$ は $\ideal M$ の真の約数ではなく、$\ideal N_i=\ideal M$ である。}

$\ideal C=[\ideal C_1,\ideal C_2]$ とし、$\ideal C_1,\ideal C_2$ が $\ideal C$ の真の約数であるとする。このとき
\[
 \ideal M=[\ideal A,\ideal C]=[\ideal A,\ideal C_1,\ideal C_2]
       =[[\ideal A,\ideal C_1],[\ideal A,\ideal C_2]].
\]
ここでイデアル $[\ideal A,\ideal C_i]$ は $\ideal M$ の真の約数である。そうでなければ $[\ideal A,\ideal C]$ は $\ideal C$ に関して既約でないからである。$\ideal A$ による割り切りも満たされているので、条件 (2) の必要性が示された。（表示 (2) は既約ではない。なぜなら、ある $[\ideal A,\ideal C_i]$ は $\ideal C_i$ に置き換えることができるからである。）

逆に (2) が満たされていると仮定する。イデアル
\[
 \ideal C_1=(\ideal C,\ideal N_1),\qquad
 \ideal C_2=(\ideal C,\ideal N_2),\qquad
 \ideal C^*=[\ideal C_1,\ideal C_2]
\]
を作る。すると $\ideal C$ は $\ideal C_1$ と $\ideal C_2$ のいずれによっても割り切れ、したがってそれらの最小公倍数 $\ideal C^*$ によっても割り切れる。$\ideal C^*$ が $\ideal C$ によって割り切れることを示すため、
\[
 f\equiv0(\ideal C^*);\quad f\equiv0(\ideal C_1);\quad f\equiv0(\ideal C_2),
 \quad\text{すなわち}\quad f=c+n_1;\quad f=\bar c+n_2
\]
とする。ここで $c,\bar c,n_1,n_2$ はそれぞれ $\ideal C,\ideal N_1,\ideal N_2$ の量であり、特に $n_i$ は $\ideal A$ で割り切れる。したがって差
\[
 g=c-\bar c=n_2-n_1
\]
は $\ideal C$ と $\ideal A$ の両方で割り切れ、ゆえに $\ideal M$ で割り切れる。さらに $n_1=n_2+m$ であるから、$n_1$（同様に $n_2$）は $\ideal N_1$ と $\ideal N_2$ の両方で割り切れ、したがって $\ideal M$ で割り切れる。よって
\[
 f=c+m,
 \qquad f\equiv0(\ideal C),
 \qquad \ideal C^*=\ideal C.
\]
ここで $\ideal C_1$ と $\ideal C_2$ は $\ideal C$ の真の約数である。もし $\ideal C_i=(\ideal C,\ideal N_i)=\ideal C$ であれば、$\ideal N_i$ は $\ideal C$ で割り切れ、さらに $\ideal U$ でも割り切れるから $\ideal M$ に等しくなり、仮定に反する。したがって $\ideal C=[\ideal C_1,\ideal C_2]$ は可約であることが分かった。定理 III が証明された。

ほとんど同じ議論により次も示される。

\textbf{補題 II。} \emph{最短表示 $\ideal M=[\ideal A,\ideal C]$ において、イデアル $\ideal C$ を真の約数で置き換えることができるならば、$\ideal C$ は可約である。}

すなわち
\[
 \ideal M=[\ideal A,\ideal C]=[\ideal A,\ideal C_1]
\]
とし、
\[
 \ideal C^*=[\ideal C_1,(\ideal A,\ideal C)]
\]
と置く。再び $\ideal C$ は $\ideal C^*$ で割り切れる。$f\equiv0(\ideal C^*)$ から
\[
 f=c_1=a+c
\]
が従う。差 $a=c_1-c$ は $\ideal A$ と $\ideal C_1$ の両方で割り切れ、したがって $\ideal M$ で割り切れる。ゆえに $c_1=c+m$、$f\equiv0(\ideal C)$、$\ideal C^*=\ideal C$ である。仮定により $\ideal C_1$ と $(\ideal A,\ideal C)$ はいずれも $\ideal C$ の真の約数なので、$\ideal C=\ideal C^*$ は可約である。

したがって既約な $\ideal C$ は真の約数で置き換えられない。

いま、有限個の既約イデアルの最小公倍数としての $\ideal M$ の二つの異なる最短表示
\[
 \ideal M=[\ideal B_1\cdots\ideal B_k]=[\ideal D_1\cdots\ideal D_l]
\]
が与えられているとする。補題 II の後の注意により、これらの表示は同時に既約である。まず次を証明する。

\textbf{補題 III。} \emph{各補元 $\ideal A_i=[\ideal B_1\cdots\ideal B_{i-1}\ideal B_{i+1}\cdots\ideal B_k]$ に対して、$\ideal M=[\ideal A_i,\ideal D_j]$ となるイデアル $\ideal D_j$ が存在する。}

実際、$\ideal M=[\ideal D_1,\ideal C_1]$、$\ideal C_1=[\ideal D_2,\ideal C_{12}]$ などと置けば、
\[
 \ideal M=[\ideal A_i,\ideal M]=[\ideal A_i,\ideal D_1,\ideal C_1]
 =[[\ideal A_i\ideal D_1],[\ideal A_i\ideal C_1]]
\]
を得る。ここで $\ideal N_1=[\ideal A_i,\ideal D_1]$、$\ideal N_2=[\ideal A_i,\ideal C_1]$ について定理 III の条件 (2) が満たされる。なぜなら $\ideal M=[\ideal A_i,\ideal B_i]$ は $\ideal B_i$ に関して既約であり、この表示は最短だからである。$\ideal B_i$ は既約と仮定されているので、$\ideal N_1$ と $\ideal N_2$ の一方は必ず $\ideal M$ に等しい。

$\ideal N_1=\ideal M$ なら補題は証明される。$\ideal N_2=\ideal M$ なら、同様に $\ideal M=[[\ideal A_i\ideal D_2],[\ideal A_i\ideal C_{12}]]$ となり、同じ議論により再び一つの成分が $\ideal M$ に等しくなる。このように続けると、ある $j<l$ について $\ideal M=[\ideal A_i,\ideal D_j]$ となるか、または $\ideal M=[\ideal A_i,\ideal C_{1\ldots l-1}]$ となる。$\ideal C_{1\ldots l-1}=\ideal D_l$ なので、いずれの場合にも補題が証明される。

これより次が従う。

\textbf{定理 IV。} \emph{既約イデアルの最小公倍数としてのイデアルの二つの異なる最短表示では、成分の個数は同じである。}

補題により、$i=1$ について
\[
 \ideal M=[\ideal A_1,\ideal B_1]=[\ideal A_1,\ideal D_{j_1}]
 = [\ideal D_{j_1},\ideal B_2,\ldots,\ideal B_k]
\]
である。いま二つの分解
\[
 \ideal M=[\ideal D_{j_1},\ideal B_2,\ldots,\ideal B_k]
          =[\ideal D_1,\ideal D_2,\ldots,\ideal D_l]
\]
を考え、$\ideal B_2$ の補元 $\bar{\ideal A}_2=[\ideal D_{j_1},\ideal B_3,\ldots,\ideal B_k]$ に関して前の議論を繰り返すと、
\[
 \ideal M=[\bar{\ideal A}_2,\ideal B_2]=[\bar{\ideal A}_2,\ideal D_{j_2}]
          =[\ideal D_{j_1},\ideal D_{j_2}\ideal B_3,\ldots,\ideal B_k]
\]
を得る。この手続きを続ければ
\[
 \ideal M=[\ideal D_{j_1},\ideal D_{j_2},\ldots,\ideal D_{j_k}]
\]
となる。仮定により $\ideal M=[\ideal D_1\cdots\ideal D_l]$ は最短表示であり、どの $\ideal D$ も省けないので、$\ideal D_{j_i}$ のうち相異なるものがすべての $\ideal D$ を尽くす。従って $k\ge l$ である。補題とその後の議論で $\ideal B$ と $\ideal D$ をすべて入れ替えると $l\ge k$ も得られ、ゆえに $k=l$ である。さらに $\ideal D_{j_i}$ はすべて互いに異なる。そうでなければ $\ideal D_{j_i}$ による最短表示の成分数が $k$ より少なくなるからである。従って $\ideal D_{j_i}=\ideal D_i$ となるように記号を選べる。同じ理由で、すべての中間表示 $\ideal M=[\ideal D_{j_1}\cdots\ideal D_{j_i},\ideal B_{i+1}\cdots\ideal B_k]$ も最短であり、補題 II の後の注意により既約である。

この個数の一致から、補題 I の逆として次が得られる。

\textbf{補題 IV。} \emph{既約表示の成分を群にまとめ、それぞれの最小公倍数を作ると、得られる表示も既約である。すなわち、既約表示}
\[
 \ideal M=[\ideal C_{11}\cdots\ideal C_{1\mu_1};\;\ldots;\;
 \ideal C_{\sigma1}\cdots\ideal C_{\sigma\mu_\sigma}]
\]
\emph{から、$\ideal N_i=[\ideal C_{i1}\cdots\ideal C_{i\mu_i}]$ と置けば、$\ideal M=[\ideal N_1\cdots\ideal N_\sigma]=[\ideal N_i,\ideal L_i]$ も既約であることが従う。}

まず $\ideal N_i$ はその補元 $\ideal L_i$ に含まれない。これはその約数 $\ideal C_{ij}$ のどれについても成り立たないからであり、従って表示は最短である。$\ideal N_i$ を真の約数で置き換えられないことを示すため、$\ideal C$ を既約イデアル $\ideal B$ に分解する。\footnote{ここでは常に $\ideal B$ による最短、従って既約な表示を意味する。} 補題 I により、既約表示
\[
 \ideal M=[\ideal B_{11}\cdots\ideal B_{1\lambda_1};\;\ldots;\;
 \ideal B_{\sigma1}\cdots\ideal B_{\sigma\lambda_\sigma}],
 \qquad
 \ideal N_i=[\ideal B_{i1}\cdots\ideal B_{i\lambda_i}]
\]
が得られる。いま $\ideal M=[\ideal N_i^*,\ideal L_i]$ が $\ideal N_i^*$ に関して既約であり、$\ideal N_i^*$ が $\ideal N_i$ の真の約数であるとする。補題 II により
\[
 \ideal N_i=[\ideal N_i^*,(\ideal N_i,\ideal L_i)]
\]
となる。この表示は必ず $\ideal N_i^*$ に関して既約である。そうでなければ $\ideal N_i^*$ を $\ideal M$ の中でも真の約数で置き換えられるからである。必要なら $(\ideal N_i,\ideal L_i)$ も真の約数で置き換えて、$\ideal N_i$ の既約表示を得る。次に $\ideal N_i$ の二成分を既約イデアルへ分解すると、$\ideal N_i$ の既約イデアル数 $\lambda_i$ は各成分の個数の和になる。従って真の約数 $\ideal N_i^*$ に対応する既約イデアル数は必ず $\lambda_i$ より小さい。しかしそのとき $\ideal M=[\ideal N_i^*,\ideal L]$ を既約イデアルへ分解すると $\sum_i\lambda_i$ より少ない個数になり、個数の一致に矛盾する。特別な場合 $\sigma=2$ には、表示 $\ideal M=[\ideal N_i,\ideal L_i]$ が補元 $\ideal L_i$ に関しても既約であることが従う。

\subsection*{\S~4. 準素イデアル。既約イデアルへの二つの異なる分解における随伴素イデアルの一意性。}

以下で問題となるのは、準素イデアルと既約イデアルとの関係である。

\textbf{定義 III。} $a\cdot b\equiv0(\ideal Q)$ かつ $a\not\equiv0(\ideal Q)$ から、指数 $\varkappa$ を有限の数として $b^\varkappa\equiv0(\ideal Q)$ が必ず従うとき、イデアル $\ideal Q$ を準素という。

この定義は次のようにも言える。積 $a\cdot b$ が $\ideal Q$ で割り切れるならば、一方の因子が割り切れるか、または各因子のある冪が割り切れる。特に $\varkappa$ を常に 1 と取れる場合、そのイデアルを素イデアルという。

準素イデアル（または素イデアル）の定義から、基底の存在により、イデアルの積だけを用いた次の定式化が従う。\footnote{二つのイデアル $\ideal A\cdot\ideal B$ の積は、通常の意味で、すべての量 $a\cdot b$ およびそれらの有限和から成るイデアルとして理解する。}

\textbf{定義 IIIa。} $\ideal A\cdot\ideal B\equiv0(\ideal Q)$ かつ $\ideal A\not\equiv0(\ideal Q)$ から、必ず $\ideal B^\lambda\equiv0(\ideal Q)$ が従うとき、イデアル $\ideal Q$ を準素という。$\lambda$ を常に 1 と取れる場合、そのイデアルを素イデアルという。したがって素イデアル $\ideal P$ については、$\ideal A\cdot\ideal B\equiv0(\ideal P)$ かつ $\ideal A\not\equiv0(\ideal P)$ から常に $\ideal B\equiv0(\ideal P)$ が従う。

IIIa は $\ideal A=(a)$、$\ideal B=(b)$ という特殊な場合として定義 III を含むので、IIIa に従って準素なイデアルは定義 III に従っても準素である。逆に、$\ideal Q$ が III に従って準素であり、IIIa の仮定 $\ideal A\cdot\ideal B\equiv0(\ideal Q)$ が満たされているとする。このとき、$\ideal A\equiv0(\ideal Q)$ が従うか、または $a\cdot\ideal B\equiv0(\ideal Q)$ かつ $a\not\equiv0(\ideal Q)$ となる少なくとも一つの量 $a$ が存在する。$b_1\ldots b_r$ を $\ideal B$ のイデアル基底とすると、定義 III により、$a\cdot b_i\equiv0(\ideal Q)$ であるから、
\[
 b_i^{\varkappa_i}\equiv0(\ideal Q)\qquad (i=1,\ldots,r)
\]
である。したがって
\[
 b=f_1b_1+\cdots+f_rb_r+n_1b_1+\cdots+n_rb_r
\]
であることから、$\lambda=\varkappa_1+\cdots+\varkappa_r$ とすれば、$\lambda$ 個の量 $b$ の任意の積は $\ideal Q$ で割り切れる。これにより、III に従って準素なイデアルは定義 IIIa を満たす。素イデアル $\ideal P$ の特別な場合には、$a\cdot\ideal B\equiv0(\ideal P)$、したがって任意の $b\equiv0(\ideal B)$ について $a\cdot b\equiv0(\ideal P)$ であり、かつ $a\not\equiv0(\ideal P)$ であることから、$b\equiv0(\ideal P)$、従って $\ideal B\equiv0(\ideal P)$ が従う。こうして二つの定義の同値性が証明された。

準素イデアルと素イデアルとの関係は、ある冪が $\ideal Q$ で割り切れるという性質をもつ元 $p$ 全体 $\ideal P$ が素イデアルをなす、という観察によって得られる。まず、$p_1$ と $p_2$ とともに $ap_1$ および $p_1-p_2$ もこの性質をもつので、$\ideal P$ がイデアルであることは明らかである。さらに IIIa の定義で用いたのと同じ基底に関する議論により、$\ideal P^\lambda\equiv0(\ideal Q)$ となる数 $\lambda$ の存在が従う。

いま
\[
 a\cdot b\equiv0(\ideal P);\qquad a\not\equiv0(\ideal P)
\]
とする。このとき $\ideal P$ の定義により
\[
 a^\lambda\cdot b^\lambda\equiv0(\ideal Q);\qquad a^\lambda\not\equiv0(\ideal Q)
\]
であり、したがって $\ideal Q$ の定義により
\[
 b^{\lambda\varkappa}\equiv0(\ideal Q)\qquad\text{従って}\qquad b\equiv0(\ideal P)
\]
となる。よって $\ideal P$ は素イデアルである。$\ideal P$ はまた、ある冪が $\ideal Q$ で割り切れるという性質をもつすべてのイデアル $\ideal B$ の最大公約数としても定義される。それらの各 $\ideal B$ は $\ideal P$ で割り切れるので、それらの最大公約数 $\ideal D$ も $\ideal P$ で割り切れる。逆に $\ideal P$ 自身がそのようなイデアル $\ideal B$ であり、したがって $\ideal D$ で割り切れる。よって $\ideal P=\ideal D$ である。したがって $\ideal P$ は、$\ideal Q$ の約数であり、かつそのある冪が $\ideal Q$ で割り切れる素イデアルである。この性質は $\ideal P$ を一意に定める。実際、
\[
 \ideal Q\equiv0(\ideal P),\qquad \ideal P^\varrho\equiv0(\ideal Q),
 \qquad
 \ideal Q\equiv0(\ideal B),\qquad \ideal B^\sigma\equiv0(\ideal Q)
\]
から
\[
 \ideal P^\varrho\equiv0(\ideal B),\qquad \ideal B^\sigma\equiv0(\ideal P)
\]
が従い、素イデアルの性質によって
\[
 \ideal P\equiv0(\ideal B),\qquad \ideal B\equiv0(\ideal P),\qquad \ideal B=\ideal P
\]
となる。まとめると次の通りである。

\textbf{定理 V。} \emph{各準素イデアル $\ideal Q$ に対して、$\ideal Q$ の約数であり、かつそのある冪が $\ideal Q$ で割り切れる素イデアル $\ideal P$ が一つ、かつただ一つ存在する。$\ideal P$ を「随伴素イデアル」と呼ぶ。}\footnote{逆が成り立たないことは、例 $\ideal M=(x^2,xy)$ によって示される。素イデアル $(x)$ はすべての条件を満たすが、$\ideal M$ は準素ではない。} \emph{$\ideal P$ は、ある冪が $\ideal Q$ で割り切れるという性質をもつすべてのイデアル $\ideal B$ の最大公約数として定義される。$\ideal P^\varrho\equiv0(\ideal Q)$ となる最小の数が $\varrho$ であるとき、$\varrho$ を $\ideal Q$ の指数という。}\footnote{有理整数、または代数的整数の領域の場合のように、一般には $\ideal P^\varrho=\ideal Q$ が成り立つとは限らない。例 $\ideal Q=(x^2,y)$、$\ideal P=(x,y)$ では、$\ideal P^2=(x^2,xy,y^2)\equiv0(\ideal Q)$ であるが、$\ideal Q\not\equiv0(\ideal P^2)$ である。したがって $\ideal Q$ は $\ideal P^2$ と異なる。}

次に、準素性と既約性の関係を証明する。

\textbf{定理 VI。} \emph{すべての非準素イデアルは可約である。言い換えれば、すべての既約イデアルは準素である。}\footnote{ここで逆が成り立たないことは、たとえば $\ideal Q=(x^2,xy,y^\lambda)=[(x^2,y),(x,y^\lambda)]$、$\lambda\ge2$ によって示される。このとき $\ideal Q$ は準素であるが可約である。$\ideal Q$ が準素であることは、$x,y$ における次数 $\lambda$ のすべての冪積を含むことから従う。したがって定数項をもたない任意の多項式について、そのある冪は $\ideal Q$ で割り切れる。一方、$a\cdot b\equiv0(\ideal Q)$ において多項式 $b$ が定数項を含むならば、すべての $\varkappa$ について $b^\varkappa\not\equiv0(\ideal Q)$ であり、$\ideal Q$ の基底多項式が斉次であるため $a\cdot b$ の各斉次成分が $\ideal Q$ で割り切れることから、$a$ は $\ideal Q$ で割り切れなければならない。}

$\ideal K$ を非準素イデアルとする。すなわち定義 III により、少なくとも一つの量の対 $a,b$ が存在して
\begin{equation}
 ab\equiv0(\ideal K),\qquad
 a\not\equiv0(\ideal K),\qquad
 b^\varkappa\not\equiv0(\ideal K)\quad\text{すべての }\varkappa\text{ について}
\tag{3}
\end{equation}
を満たす。いま二つのイデアル
\[
 \ideal L_0=(\ideal K,a),\qquad \ideal N_0=(\ideal K,b)
\]
を作る。これらは (3) により $\ideal K$ の真の約数であり、また
\begin{equation}
 \ideal L_0\ideal N_0\equiv0(\ideal K)
\tag{4}
\end{equation}
を満たす。最小公倍数 $\ideal K_0=[\ideal L_0,\ideal N_0]$ の元 $f$ について、次の二つの場合がある。

\emph{第一の場合}、
\[
 f\equiv0(\ideal L_0),\qquad f\equiv0(\ideal N_0),
 \qquad\text{すなわち}\qquad f\equiv a_1\cdot b\;(\ideal K)
\]
から常に
\[
 f\equiv l_0\cdot b\;(\ideal K),\qquad l_0\equiv0(\ideal L_0)
\]
という表示が従うなら、(4) により $f\equiv0(\ideal K)$ となる。従って $\ideal K_0\equiv0(\ideal K)$ であり、$\ideal K\equiv0(\ideal K_0)$ でもあるから $\ideal K=\ideal K_0$ となり、$\ideal K$ は可約である。

\emph{第二の場合}、そのような $l_0$ が存在しない $f\equiv0(\ideal K_0)$ が少なくとも一つある。その $f$ に属する $a_1$ を用いて
\[
 \ideal L_1=(\ideal L_0,a_1)=(\ideal K,a,a_1),\qquad
 \ideal N_1=(\ideal K,b^2)
\]
を作る。$a_1b\equiv0(\ideal L_0)$ なので、(4) により
\begin{equation}
 \ideal L_1\ideal N_1\equiv0(\ideal K)
\tag{4'}
\end{equation}
となり、$\ideal L_1$ は $\ideal L_0$ の真の約数である。

$\ideal K_1=[\ideal L_1,\ideal N_1]$ の元 $f$ についても同じ二つの場合がある。すなわち、
\[
 f\equiv0(\ideal L_1),\qquad f\equiv0(\ideal N_1)
\]
すなわち
\[
 f\equiv a_2b^2\;(\ideal K)
 \quad\text{から常に}\quad
 f\equiv l_1b^2\;(\ideal K),\qquad l_1\equiv0(\ideal L_1)
\]
が従うなら、(4') により $\ideal K=\ideal K_1$ となる。

そうでなければ、少なくとも一つの $f$ についてそのような $l_1$ が存在せず、$\ideal L_2=(\ideal L_1,a_2)$、$\ideal N_2=(\ideal K,b^{2^2})$ を作ることになる。このとき $\ideal L_2\cdot\ideal N_2\equiv0(\ideal K)$ であり、$\ideal L_2$ は $\ideal L_1$ の真の約数である。このように続けて、一般に
\[
 \ideal L_0=(\ideal K,a);\quad
 \ideal L_1=(\ideal L_0,a_1);\quad\ldots;\quad
 \ideal L_\nu=(\ideal L_{\nu-1},a_\nu);\quad\ldots,
\]
\[
 \ideal N_0=(\ideal K,b);\quad
 \ideal N_1=(\ideal K,b^2);\quad\ldots;\quad
 \ideal N_\nu=(\ideal K,b^{2^\nu});\quad\ldots
\]
と定義する。ここで $a_i$ は、ある $f$ が存在して
\[
 f\equiv0(\ideal L_{i-1});\qquad f\equiv0(\ideal N_{i-1}),
\]
\[
 \text{すなわち}\qquad f\equiv a_i\cdot b^{2^{i-1}}\;(\ideal K),
 \qquad\text{しかし}\qquad a_i\not\equiv0(\ideal L_{i-1})
\]
となるように定められる。従って一般に $\ideal L_i\cdot\ideal N_i\equiv0(\ideal K)$ であり、$\ideal N_i$ は (3) により $\ideal K$ の真の約数、$\ideal L_i$ は $\ideal L_{i-1}$ の真の約数である。有限鎖についての定理 I により、$\ideal L$ の鎖は有限段階、たとえば $\ideal L_n$ で停止しなければならない。従ってすべての $f\equiv0(\ideal L_n)$、$f\equiv0(\ideal N_n)$ に対して、$l_n\equiv0(\ideal L_n)$ を用いて $f\equiv l_n b^{2^n}\;(\ideal K)$ と表せる。前の議論により $\ideal K=[\ideal L_n,\ideal N_n]$ となり、$\ideal K$ は可約である。

\emph{随伴素イデアルの一意性}は、いま証明したことから従う。

\[
 \ideal M=[\ideal B_1\cdots\ideal B_k]=[\ideal D_1\cdots\ideal D_k]
\]
を、既約イデアルの最小公倍数としての $\ideal M$ の二つの最短、したがって既約な表示とし、その個数は定理 IV により一致しているとする。そのとき同じ定理により、そこで現れた中間表示（そこで述べたように添字 $j_i=i$ と置ける）
\[
 \ideal M=[\ideal D_1\cdots\ideal D_{i-1}\ideal B_i\ideal B_{i+1}\cdots\ideal B_k]
       =[\ideal D_1\cdots\ideal D_{i-1}\ideal D_i\ideal B_{i+1}\cdots\ideal B_k]
       =[\bar{\ideal A}_i,\ideal B_i]=[\bar{\ideal A}_i,\ideal D_i]
\]
も最短表示である。したがって
\[
 \bar{\ideal A}_i\cdot\ideal B_i\equiv0(\ideal D_i),\qquad \bar{\ideal A}_i\not\equiv0(\ideal D_i),
 \qquad
 \bar{\ideal A}_i\cdot\ideal D_i\equiv0(\ideal B_i),\qquad \bar{\ideal A}_i\not\equiv0(\ideal B_i)
\]
である。定理 VI により既約イデアル $\ideal B_i$ と $\ideal D_i$ は準素であるから、二つの数 $\lambda_i$ と $\mu_i$ が存在して
\begin{equation}
 \ideal B_i^{\lambda_i}\equiv0(\ideal D_i),
 \qquad
 \ideal D_i^{\mu_i}\equiv0(\ideal B_i)
\tag{5}
\end{equation}
となる。$\ideal B_i$ と $\ideal D_i$ の随伴素イデアルをそれぞれ $\ideal P_i$ と $\bar{\ideal P}_i$ とする。すなわち $\ideal P_i^{\varrho_i}\equiv0(\ideal B_i)$、$\bar{\ideal P}_i^{\sigma_i}\equiv0(\ideal D_i)$ である。(5) から
\[
 \ideal P_i^{\lambda_i\varrho_i}\equiv0(\bar{\ideal P}_i),
 \qquad
 \bar{\ideal P}_i^{\mu_i\sigma_i}\equiv0(\ideal P_i)
\]
が従い、素イデアルの性質により
\[
 \ideal P_i\equiv0(\bar{\ideal P}_i),
 \qquad
 \bar{\ideal P}_i\equiv0(\ideal P_i),
 \qquad
 \ideal P_i=\bar{\ideal P}_i
\]
となる。したがって次を得る。

\textbf{定理 VII。} \emph{既約イデアルの最小公倍数としてのイデアルの二つの異なる最短表示では、随伴素イデアルが一致する。その中には同じものも現れ得る\footnote{これは、たとえば注 19 の例 $(x^2,xy,y^\lambda)=[(x^2,y),(x,y^\lambda)]$、$\lambda\ge2$ によって示される。ここで二つの随伴素イデアルはいずれも $(x,y)$ である。}し、どの分解でも同じ回数だけ現れる。したがってイデアル自身は、少なくとも一つの仕方で、一方のイデアル $\ideal B_i$ のある冪が対応づけられた $\ideal D_i$ で割り切れ、かつ逆も成り立つように対応づけられる。その個数は定理 IV により一致する。}\footnote{既約イデアルの中に含まれる「孤立」イデアルの一意性については \S~7 を参照。}

\subsection*{\S~5. 極大準素イデアルの最小公倍数としてのイデアルの表示。随伴素イデアルの一意性。}

\textbf{定義 IV。} \emph{最短表示 $\ideal M=[\ideal Q_1\cdots\ideal Q_\alpha]$ は、すべての $\ideal Q$ が準素であり、しかも二つの $\ideal Q$ の最小公倍数がもはや準素でないとき、極大準素イデアルの最小公倍数と呼ばれる。}

そのような表示が少なくとも一つ常に存在することは、$\ideal M$ を既約イデアルの最小公倍数として表示できることから従う。これらのイデアルは準素である。すでに極大準素イデアルによる表示になっているか、または二つのイデアルの最小公倍数が再び準素である。後者の場合にはイデアルの個数が一つ減るので、この手続きを繰り返せば有限段階の後に求める表示に到達する。

この表示は補題 IV により既約である。逆に、極大準素イデアルによる任意の既約表示はこのようにして得られる。これは $\ideal Q$ を既約イデアルに分解することから分かる。

ここで定理 VII から対応する一意性定理を導くため、随伴素イデアルとの関係を調べなければならない。

\textbf{定理 VIII。} \emph{準素イデアル $\ideal N_1,\ideal N_2,\ldots,\ideal N_\lambda$ がすべて同じ随伴素イデアル $\ideal P$ をもつならば、それらの最小公倍数 $\ideal Q=[\ideal N_1,\ideal N_2,\ldots,\ideal N_\lambda]$ も準素であり、$\ideal P$ を随伴素イデアルにもつ。逆に、$\ideal Q=[\ideal N_1\cdots\ideal N_\lambda]$ が準素イデアル $\ideal Q$ の既約表示であるならば、すべての $\ideal N_i$ は準素であり、随伴素イデアルとして $\ideal Q$ の随伴素イデアル $\ideal P$ をもつ。}

第一部の証明では、まず各 $i$ について $\ideal P^{\varrho_i}\equiv0(\ideal N_i)$ から、$\tau$ を指数 $\varrho_i$ の最大値として $\ideal P^\tau\equiv0(\ideal Q)$ も従うことを注意する。$\ideal P$ はまた $\ideal Q$ の約数でもあるから、$\ideal Q$ が準素ならば $\ideal P$ は必然的にその随伴素イデアルである。
\[
 \ideal A\cdot\ideal B\equiv0(\ideal Q),
 \qquad
 \ideal B^k\not\equiv0(\ideal Q)\quad\text{（すべての }k\text{ について）}
\]
からは
\[
 \ideal B\not\equiv0(\ideal P),
 \quad\text{従って}\quad
 \ideal B^k\not\equiv0(\ideal N_i)\quad\text{（すべての }k\text{ について）}
\]
が従う。よって $\ideal A\equiv0(\ideal N_i)$ であり、したがって $\ideal A\equiv0(\ideal Q)$ である。これにより $\ideal Q$ が準素であり、$\ideal P$ が随伴素イデアルであることが証明された。

逆に、まず
\[
 \ideal Q=[\ideal N_1\cdots\ideal N_\lambda]=[\ideal N_i,\ideal L_i']
\]
を準素イデアル $\ideal N_i$ による $\ideal Q$ の最短表示とし、$\ideal P_i$ をそれらの随伴素イデアルとする。
\[
 \ideal N_i\ideal N_i'\equiv0(\ideal Q),
 \qquad
 \ideal N_i'\not\equiv0(\ideal N_i)
\]
（表示が最短であるため）から
\[
 \ideal N_i'\not\equiv0(\ideal P_i),
 \qquad\text{しかし}\qquad
 \ideal N_i'\equiv0(\ideal P)
\]
が従う。$\ideal P_i$ は同時に $\ideal Q$ の約数であるから、各 $i$ について $\ideal P_i$ は $\ideal Q$ の随伴素イデアル $\ideal P$ と一致する。

なお、任意の既約表示 $\ideal Q=[\ideal N_1\cdots\ideal N_\lambda]$ において $\ideal N_i$ が準素であることを示す必要がある。\footnote{ここで既約表示が本質的であることは、既約でない表示 $\ideal Q=[x^2,xy,y^\lambda]=[(x^2,xy,y^2,yz),(x,y^\lambda)]$、$\lambda\ge2$ の例が示す。このとき $(x^2,xy,y^2,yz)=[(x^2,y),(x,y^2,z)]$ は、いま証明したことにより準素ではない。なぜなら後者の表示は準素イデアルによる最短表示であるが、随伴素イデアル $(x,y)$ と $(x,y,z)$ は異なるからである。（$\ideal Q$ は注 19 により準素である。）} そのため、$\ideal N_i$ をその既約イデアル $\ideal B$ に分解する。こうして得られる最短表示 $\ideal Q=[\ideal B_1\cdots\ideal B_\mu]$ では、すべてのイデアルが準素であり、いま証明したことにより随伴素イデアルとして $\ideal P$ をもつ。すると上で証明した定理の第一部により、各 $\ideal N_i$ についても同じことが成り立つ。これで定理 VIII が完全に証明された。

\textbf{補足。} 素イデアルは必ず既約であることも従う。実際、既約表示 $\ideal P=[\ideal N,\ideal L]$ から、定理 VIII により $\ideal N\equiv0(\ideal P)$ が得られる。したがって $\ideal P\equiv0(\ideal N)$ でもあるので $\ideal P=\ideal N$ であり、同様に $\ideal P=\ideal L$ である。$\ideal P$ の既約性は直接にも従う。$\ideal P=[\ideal N,\ideal L]$ ならば、$\ideal N\cdot\ideal L\equiv0(\ideal P)$、$\ideal N\not\equiv(\ideal P)$、$\ideal L\not\equiv0(\ideal P)$ となり、素イデアルの定義性質に反するからである。

いま
\[
 \ideal M=[\ideal Q_1\cdots\ideal Q_\alpha]=[\bar{\ideal Q}_1\cdots\bar{\ideal Q}_\beta]
\]
を、極大準素イデアルの最小公倍数としての $\ideal M$ の二つの既約表示とする。$\ideal Q$ を既約イデアル $\ideal B$ に分解し、同様に $\bar{\ideal Q}$ を既約イデアル $\ideal D$ に分解すると、$\ideal M$ の既約イデアルの最小公倍数としての二つの既約表示が得られる。定理 VII により、成分の個数も随伴素イデアルも一致する。定理 VIII により、固定した $\ideal Q_i$ に現れるすべての既約イデアル $\ideal B$ は同じ随伴素イデアル $\ideal P_i$ をもつ。一方、$\ideal Q_k$ に随伴する $\ideal P_k$ は必ずこれと異なる。そうでなければ、定理 VIII により、極大準素イデアルによる表示ではなくなるからである。したがって $\ideal Q$ の個数 $\alpha$ は、$\ideal B$ の異なる随伴素イデアル $\ideal P$ の個数に等しい。これらの異なる $\ideal P$ が $\ideal Q$ の随伴素イデアルをなす。$\bar{\ideal Q}$ についても、その $\ideal D$ への分解に関して同じことが成り立つ。よって定理 VII は、$\ideal Q$ と $\bar{\ideal Q}$ の個数の一致と、それらの随伴素イデアルの一致を与える。同時に、節のはじめに述べた、既約イデアルを極大準素イデアルにまとめる操作は、同じ随伴素イデアルをもつものすべて、そしてそれらだけをまとめることからなることが分かる。さらに定理 VIII は、極大準素イデアルの既約性の性質を示す。すなわち、それらは極大準素イデアルの最小公倍数としての既約表示をもたない。

まとめると次の通りである。

\textbf{定理 IX。} \emph{極大準素イデアルの最小公倍数としてのイデアルの二つの既約表示では、成分の個数と、互いにすべて異なる随伴素イデアルが一致する。言い換えれば、各 $\ideal Q$ には、$\ideal Q$ のある冪が $\ideal Q'$ で割り切れ、かつ逆も成り立つような $\ideal Q'$ を一つ、かつただ一つ対応させることができる。}\footnote{異なる表示の例は、定理 II の注 12 に与えたもの、すなわち任意の $\mu$ について $(x^2,xy)=[(x),(x^2,\mu x+y)]$ である。随伴素イデアル $\ideal P_1=(x)$ と $\ideal P_2=(x,y)$ は異なるので、これらは極大準素イデアルである。--「孤立」極大準素イデアルの一意性については \S~7 を参照。} \emph{$\ideal Q$ と $\ideal Q'$ は、極大準素イデアルへの分解に関して既約性の性質をもつ。}

\textbf{補足。} 定理 IX は、既約表示ではなく最短表示だけを仮定した場合にも本質的に有効であることに注意する。たとえば
\[
 \ideal M=[\ideal Q_1\cdots\ideal Q_i^*\cdots\ideal Q_\alpha]
\]
が $\ideal Q_i^*$ に関して既約であり、$\ideal Q_i^*$ が $\ideal Q_i$ の真の約数であり、$\ideal L_i$ が $\ideal Q_i$ の補元であるとする。このとき補題 IV により
\[
 \ideal Q_i=[\ideal Q_i^*,(\ideal U_i,\ideal Q_i)]
\]
であり、この表示は $\ideal Q_i^*$ に関して既約である。定理 VIII により、定理を適用する際に必要なら $(\ideal L_i,\ideal Q_i)$ を真の約数で置き換えれば、$\ideal Q_i^*$ は準素であり、$\ideal Q_i$ と同じ随伴素イデアル $\ideal P_i$ をもつ。手続きを続けると、そのような各表示に対して、成分の個数と随伴素イデアルが一致するように、極大準素イデアルによる既約表示を対応させることができる。したがって、最短表示の場合にも、二つの異なる表示では成分の個数と随伴素イデアルが一致する。

このようにして一意に定義される、互いに異なる随伴素イデアルを、単に「$\ideal M$ の随伴素イデアル」と呼ぶことにする。




%
\subsection*{\S~6. 相対的に素な既約イデアルの最小公倍数としてのイデアルの一意表示。}

\textbf{定義 V.} \emph{イデアル $\ideal R$ が $\ideal S$ に対して相対的に素であるとは、
\[
 \ideal T\cdot\ideal R\equiv0(\ideal S)
\]
から必ず $\ideal T\equiv0(\ideal S)$ が従うことをいう。$\ideal R$ が $\ideal S$ に対して相対的に素であり、かつ $\ideal S$ も $\ideal R$ に対して相対的に素であるならば、$\ideal R$ と $\ideal S$ は互いに相対的に素であるという。}\footnote{相対的に素であるという条件は対称的でない。例えば、$\ideal R=(x^2,y)$ は $\ideal S=(x)$ に対して相対的に素である。しかし $\ideal S$ は $\ideal R$ に対して相対的に素でない。実際、$\ideal S^2\equiv0(\ideal R)$ であるが、$\ideal S\not\equiv0(\ideal R)$ である。}

\emph{イデアルが相対的に素な意味で既約であるとは、互いに相対的に素な真の約数の最小公倍数として表せないことをいう。}

特に、$\ideal T$ の代わりに、$\ideal T\ideal R\equiv0(\ideal S)$ となるすべての $\ideal T$ の最大公約数 $\ideal T_0$ を取るならば、やはり $\ideal T_0\ideal R\equiv0(\ideal S)$ かつ $\ideal S\equiv0(\ideal T_0)$ である。したがって、$\ideal R$ が $\ideal S$ に対して相対的に素であるとき $\ideal T_0=\ideal S$ であり、$\ideal R$ が $\ideal S$ に対して相対的に素でないとき $\ideal T_0$ は $\ideal S$ の真の約数である。\footnote{このように定義される $\ideal T_0$ は、Lasker の「剰余加群」と一致する。また、イデアルの代わりに数の加群へ拡張すると、Dedekind の二つの加群の「商」と一致する。Lasker, loc. cit., p. 49; Dedekind (Zahlentheorie), p. 504.}

一意性の証明は次に基づく。

\textbf{定理 X.} \emph{1. $\ideal R$ がイデアル $\ideal S_1\ldots\ideal S_\lambda$ のそれぞれに対して相対的に素であるならば、$\ideal R$ はそれらの最小公倍数 $\ideal S$ に対しても相対的に素である。}

\emph{2. イデアル $\ideal S_1,\ldots,\ideal S_\lambda$ が $\ideal R$ に対して相対的に素であるならば、それらの最小公倍数 $\ideal S$ も $\ideal R$ に対して相対的に素である。}

\emph{3. $\ideal R$ が $\ideal S$ に対して相対的に素であり、$\ideal S=[\ideal S_1\cdots\ideal S_\lambda]$ が $\ideal S$ の簡約表示であるならば、$\ideal R$ は各 $\ideal S_i$ に対しても相対的に素である。}

\emph{4. $\ideal S$ が $\ideal R$ に対して相対的に素であるならば、$\ideal S$ の任意の約数 $\ideal S_i$ も $\ideal R$ に対して相対的に素である。}

1. $\ideal T\cdot\ideal R\equiv0(\ideal S)$ は必ず $\ideal T\cdot\ideal R\equiv0(\ideal S_i)$ を与えるから、仮定により $\ideal T\equiv0(\ideal S_i)$ であり、したがって $\ideal T\equiv0(\ideal S)$ である。

2. 次のように置く。
\[
 \widehat{\ideal C}_1=[\ideal S_2\cdots\ideal S_\lambda],\quad
 \widehat{\ideal C}_{12}=[\ideal S_3\cdots\ideal S_\lambda],\quad \ldots,\quad
 \widehat{\ideal C}_{12\ldots\lambda-1}=\ideal S_\lambda.
\]
$\ideal T\cdot\ideal S\equiv0(\ideal R)$ から $\ideal T\cdot\ideal C_1\widehat{\ideal S}_1\equiv0(\ideal R)$ が従う。よって、仮定により $\ideal T\widehat{\ideal C}_1\equiv0(\ideal R)$ である。同じ議論を順に適用すると、$\ideal T\widehat{\ideal S}_{12}\equiv0(\ideal R)$ が得られ、最後に $\ideal T\widehat{\ideal C}_{12\ldots\lambda-1}=\ideal T\cdot\ideal S_\lambda\equiv0(\ideal R)$ が得られる。したがって $\ideal T\equiv0(\ideal R)$ である。

3. $\widehat{\ideal S}_i$ を $\ideal S_i$ の補元とする。$\ideal T_0\cdot\ideal R\equiv0(\ideal S_i)$ と $\ideal C_i\cdot\ideal R\equiv0(\ideal C_i)$ から、$[\ideal T_0,\ideal C_i]\cdot\ideal R\equiv0(\ideal S)$ が従う。もし $\ideal T_0$ が $\ideal S_i$ の真の約数であれば、$\ideal S=[\ideal S_i,\widehat{\ideal S}_i]$ は簡約されているから、$[\ideal T_0,\widehat{\ideal S}_i]$ は $\ideal S$ の真の約数となり、仮定に反する。

4. $\ideal T_0$ が $\ideal R$ の真の約数であるとして、$\ideal T\ideal S_i\equiv0(\ideal R)$ からは $\ideal T\ideal S\equiv0(\ideal R)$ も従う。したがって $\ideal T_0$ は $\ideal R$ の真の約数となり、仮定に反する。

定義 V は特に、任意のイデアルが $\Sigma$ の全要素から成る単位イデアル $\ideal O$ に対して相対的に素であることを示す。\footnote{$\ideal O$ が単位の役割を果たすのは、除法関係および最小公倍数に関してだけであり、積の形成に関してではない。例えば、定数項をもたない $x$ の整係数多項式全体の領域では $\ideal O=(x)$ であり、偶数全体の領域では $\ideal O=(2)$ である。} ただし、以下の定理は $\ideal O$ と異なるイデアルにだけ適用される。

いま証明した定理 X から、まず次が得られる。

\textbf{補題 V.} \emph{互いに相対的に素で、かつ $\ideal O$ と異なるイデアルたちの最小公倍数としてのイデアルの任意の表示は簡約されている。}

このような表示を
\[
 \ideal M=[\ideal R_1\ideal R_2\cdots\ideal R_\sigma]=[\ideal R_i,\ideal L_i]
\]
とする。定理 X, 2 により、$\ideal L_i$ も $\ideal R_i$ に対して相対的に素である。したがって $\ideal R_i$ は $\ideal L_i$ に含まれ得ず、この表示は最短である。実際、$\ideal L_i\equiv0(\ideal R_i)$ であり、$\ideal L_i$ が $\ideal R_i$ に対して相対的に素であれば、$\ideal O\cdot\ideal L_i\equiv0(\ideal R_i)$ から $\ideal O\equiv0(\ideal R_i)$ が従い、したがって $\ideal R_i=\ideal O$ となるが、これは仮定に反する。

もしここで $\ideal R_i$ を真の約数 $\ideal R_i^*$ で置き換えられるならば、補題 II により $\ideal R_i$ は可約であり、
\[
 \ideal R_i=[\ideal R_i^*,(\ideal R_i,\ideal L_i)]
\]
となる。必要ならば $(\ideal R_i,\ideal L_i)$ を真の約数 $(\ideal R_i,\ideal L_i)^*$ で置き換え、同様に $\ideal R_i^*$ も真の約数で置き換えて、$\ideal R_i$ の簡約表示が得られるようにする。このとき定理 X, 3 は、$\ideal L_i$ が $(\ideal R_i,\ideal L_i)^*$ に対しても相対的に素であることを示す。またこのイデアルは、もとの表示が最短であったから、$\ideal O$ とは異なる。しかし $(\ideal R_i,\ideal L_i)^*$ は $(\ideal R_i,\ideal L_i)$ に含まれ、$(\ideal R_i,\ideal L_i)$ は $\ideal L_i$ に含まれるので、矛盾が生じる。

定理 X はさらに、随伴素イデアルとの関係を仲介する次の定理を与える。

\textbf{定理 XI.} \emph{$\ideal R$ が $\ideal S$ に対して相対的に素であり、かつ $\ideal S$ が $\ideal O$ と異なるならば、$\ideal R$ の随伴素イデアル\footnote{$\ideal R$ の随伴素イデアルは \S~5 の終わりで定義された。}で $\ideal S$ の随伴素イデアルによって割り切れるものは存在しない。逆に、そのような割り切り関係が一つも存在しなければ、$\ideal R$ は $\ideal S$ に対して相対的に素であり、もちろん $\ideal S$ は $\ideal O$ と異なる。}

\[
 \ideal R=[\ideal Q_1\cdots\ideal Q_\alpha],
 \qquad
 \ideal S=[\ideal Q_1^*\cdots\ideal Q_\beta^*]
\]
を、$\ideal R$ と $\ideal S$ の極大準素イデアルによる簡約表示とし、$\ideal P_1,\ldots,\ideal P_\alpha$ および $\ideal P_1^*,\ldots,\ideal P_\beta^*$ をそれらの随伴素イデアルとする。主張を次の形で証明する。すなわち、ある $\ideal P$ がある $\ideal P^*$ で割り切れるならば、$\ideal R$ は $\ideal S$ に対して相対的に素ではあり得ず、また逆も成り立つ。

そこで $\ideal P_\mu\equiv0(\ideal P_\nu^*)$ と仮定する。このとき $\ideal Q_\mu\equiv0(\ideal P_\nu^*)$ でもあり、$\ideal P_\nu^*$ の定義により、これは $\ideal Q_\mu^{\sigma_\nu}\equiv0(\ideal Q_\nu^*)$ を与える。したがって $\ideal R^{\sigma_\nu}\equiv0(\ideal Q_\nu^*)$ である。$\ideal R^\tau$ を、$\ideal Q_\nu^*$ で割り切れる $\ideal R$ の最小冪とする。$\tau=1$ ならば $\ideal R$ は $\ideal Q_\nu^*$ で割り切れる。しかし $\ideal S\ne\ideal D$ であるから、$\ideal R$ は $\ideal Q_\nu^*$ に対して相対的に素でない。$\tau\ge2$ ならば、$\ideal R^{\tau-1}\cdot\ideal R\equiv0(\ideal Q_\nu^*)$ である一方、$\ideal R^{\tau-1}\not\equiv0(\ideal Q_\nu^*)$ であるから、やはり $\ideal R$ は $\ideal Q_\nu^*$ に対して相対的に素でない。\footnote{$\ideal R^0$ は定義されていない。なぜなら、$\Sigma$ が単位を含むとは限らないからである。したがって $\tau=1$ の場合を先に扱わなければならなかった。$\Sigma$ が単位をもつ場合にも、$\ideal S$ に関する仮定により $\tau=0$ の場合は排除される。} どちらの場合にも、定理 X, 3 によって $\ideal R$ は $\ideal S$ に対して相対的に素でないことが従う。

逆に、$\ideal R$ が $\ideal S$ に対して相対的に素でないならば、定理 X, 1 により、少なくとも一つの $\ideal Q_\nu^*$ に対して相対的に素でない。したがって
\[
 \ideal T_0\cdot\ideal R\equiv0(\ideal Q_\nu^*),
 \qquad
 \ideal T_0\not\equiv0(\ideal Q_\nu^*),
\]
であり、$\ideal Q_\nu^*$ は準素であるから、
\[
 \ideal R^\tau\equiv0(\ideal Q_\nu^*),
 \qquad
 \ideal Q_1^\tau\cdots\ideal Q_\alpha^\tau\equiv0(\ideal Q_\nu^*)
\]
となる。随伴素イデアルについては、これは
\[
 \ideal P_1^{\tau\varrho_1}\cdots\ideal P_\alpha^{\tau\varrho_\alpha}\equiv0(\ideal P_\nu^*)
\]
を意味し、したがって素イデアルの性質により、$\ideal P_\nu^*$ は少なくとも一つの $\ideal P$ に含まれる。これで定理 XI が証明された。

定理 X と定理 XI は、相対的に素な既約イデアルへの分解の存在\footnote{分解の存在は、\S~2 における有限個の既約イデアルへの分解の存在証明と完全に類似して、直接に証明することもできる。}と一意性を次のように与える。

$\ideal M=[\ideal Q_1\cdots\ideal Q_\alpha]$ を $\ideal M$ の極大準素イデアルによる簡約表示（または少なくとも最短表示）とし、$\ideal P_1,\ldots,\ideal P_\alpha$ を随伴素イデアルとする。これらの $\ideal P$ を、ある群のイデアルが別の群のイデアルによって割り切れないように、しかも一つの群を、この性質をともにもつ二つの部分群に分けられないように、群にまとめる。

このような群分けを構成するには、定義上、一つの群 $G$ は、各イデアル $\ideal P$ とともに、$\ideal P_1,\ldots,\ideal P_\alpha$ の中に現れる $\ideal P$ のすべての約数および倍数を含まなければならないことに注意する。例えば、$\ideal P_j^{(i)}$ を $\ideal P$ のすべての倍数とし、$\ideal P_{j_1}^{(i_1)}$ を $\ideal P_j^{(i)}$ のすべての約数とし、$\ideal P_{j_1}^{(i_1 i_2)}$ を $\ideal P_{j_1}^{(i_1)}$ のすべての倍数とする、というように進める。全体で有限個のイデアル $\ideal P$ しか現れないので、この手続きは有限回の後に停止しなければならない。こうして得られるイデアルの集合は、必要な性質をもつ群 $G$ である。すなわち、構成によって指定されたすべての約数と倍数を含み、その外のイデアルは中のイデアルの約数にも倍数にもなり得ない。

群 $G$ は既約性条件も満たす。もしそれが二つの部分群 $G^{(1)}$ と $G^{(2)}$ に分けられ、$G^{(1)}$ が交互の構成のある段階で現れるイデアルの一つを含むとすれば、それは先行するすべてのものを含む。なぜなら、それらは交互に倍数と約数（または約数と倍数）であるからである。したがって $\ideal P$ も含み、結局群 $G$ 全体を含む。$G$ に含まれないイデアルについてこの構成を繰り返せば、すべての $\ideal P$ の群 $G_1,\ldots,G_\sigma$ への分解が得られ、同じ議論によりこの群分けは一意である。

$\mu$ が $1$ から $\lambda_i$ まで動くとして、$\ideal Q_{i\mu}$ を、その随伴素イデアルが群 $G_i$ に集められる $\ideal P$ とする。$\ideal P$ はすべて相異なるので、最短表示に属する準素イデアル $\ideal Q_{i\mu}$ は一意に決まっている。そこで
\[
 \ideal R_i=[\ideal Q_{i1}\cdots\ideal Q_{i\lambda_i}],
 \qquad
 \ideal M=[\ideal R_1\cdots\ideal R_\sigma]
\]
と置く。$\ideal R_i$ が最短表示にすぎない場合にも、定理 XI はなお適用できる。なぜなら、定理 IX への付加によって、その随伴素イデアルはすでに明確に定義されているからである。$\ideal R_i$ の随伴素イデアルで $\ideal R_j$ の一つにより割り切れるものはなく、また逆もそうであるから、定理 XI は $\ideal R_i$ と $\ideal R_j$ が互いに相対的に素であることを示す。また各 $\ideal R_i$ は、群 $G_i$ の既約性により、相対的に素な意味で既約である。補題 V によりこの表示は簡約されている。逆に、相対的に素な既約イデアルへの任意の分解は、定理 XI により、いま述べた群分けへ導く。

いま
\[
 \ideal M=[\bar{\ideal R}_1\cdots\bar{\ideal R}_\tau]
\]
を、$\ideal M$ の相対的に素な既約イデアルによる第二の表示とする。補題 V によりこれは簡約されている。$\bar{\ideal R}$ を極大準素イデアルに分解すると、随伴素イデアルが一致すること、したがってこれらの素イデアルの群分けも一致することが分かる。ゆえに $\tau=\sigma$ であり、記法は $\ideal R_i$ と $\bar{\ideal R}_i$ が同じ群に属するように選べる。
\[
 \ideal M=[\ideal R_i,\ideal L_i]=[\bar{\ideal R}_i,\bar{\ideal L}_i]
\]
を、イデアルと補元による表示とする。このとき $\bar{\ideal R}_i$ は $\ideal R_i$ と同じ群に割り当てられているから、定理 XI により $\ideal L_i$ は $\bar{\ideal R}_i$ に対して相対的に素であり、$\bar{\ideal L}_i$ は $\ideal R_i$ に対して相対的に素である。したがって
\[
 \ideal R_i\cdot\ideal L_i\equiv0(\bar{\ideal R}_i),
 \qquad
 \bar{\ideal R}_i\cdot\bar{\ideal L}_i\equiv0(\ideal R_i)
\]
から
\[
 \ideal R_i\equiv0(\bar{\ideal R}_i),
 \qquad
 \bar{\ideal R}_i\equiv0(\ideal R_i),
 \qquad
 \ideal R_i=\bar{\ideal R}_i
\]
が従う。したがって次が成り立つ。

\textbf{定理 XII.} \emph{任意のイデアルは、有限個の互いに相対的に素で、かつ相対的に素な意味で既約なイデアルの最小公倍数として、一意に表すことができる。}

\subsection*{\S~7. 孤立イデアルの一意性。}

\textbf{定義 VI.} \emph{最短表示 $\ideal M=[\ideal R,\ideal L]$ が $\ideal L$ に関して簡約されているとき、$\ideal R$ の随伴素イデアルが $\ideal L$ の随伴素イデアルに含まれることが一つもなければ、$\ideal R$ を孤立イデアルという。これは同値に、$\ideal L$ が $\ideal R$ に対して相対的に素であることをいう。}

すると、表示 $\ideal M=[\ideal R,\ideal L]$ は補題 V における表示 $\ideal M=[\ideal R_i,\ideal L_i]$ の条件を満たし、補題 V は次を証明する。

\textbf{補題 VI.} \emph{$\ideal R$ が、$\ideal L$ に関して簡約された最短表示 $\ideal M=[\ideal R,\ideal L]$ の孤立イデアルであるならば、この表示は $\ideal R$ に関しても簡約されている。}

したがって孤立イデアルは簡約表示にしか現れない。$\ideal R$ および $\ideal L$ を既約イデアルへ分解したときに生じる随伴素イデアルは、対応する $\ideal M$ の分解に現れる、一意に定まる随伴素イデアルを互いに補い合う。したがって、$\ideal R$ の既約イデアルへの分解に属する随伴素イデアルで、$\ideal M$ の残りの随伴素イデアルに含まれるものはない。逆に、この条件が満たされ、かつ $\ideal R$ が少なくとも一つの表示 $\ideal M=[\ideal R,\ideal L]$ において $\ideal R$ に関して簡約された形で現れ、したがって簡約表示 $\ideal M=[\ideal R,\ideal L^*]$ にも現れるならば、定義 VI により $\ideal R$ は孤立している。このことは、特定の補元 $\ideal L$ に依存しない次の定義を与える。

\textbf{定義 VIa.} \emph{$\ideal R$ が孤立イデアルであるとは、その既約イデアルへの分解に属する素イデアルが、対応する $\ideal M$ の分解に属する他の随伴素イデアルに含まれず、かつ $\ideal R$ が少なくとも一つの表示 $\ideal M=[\ideal R,\ideal L]$ において $\ideal R$ に関して簡約された形で現れることをいう。}\footnote{通常の随伴素イデアル（\S~5 の終わり）を導入するならば、$\ideal L$ の随伴素イデアルがすべて $\ideal R$ のそれらと異なるという特別な条件を付け加えなければならない。定義 VIa によれば、もはや補元に関して表示が簡約されていると仮定する必要はない。}

\[
 \ideal M=[\ideal R,\ideal L]=[\bar{\ideal R},\bar{\ideal L}]
\]
を、孤立イデアル $\ideal R$ および $\bar{\ideal R}$ と補元 $\ideal L$ および $\bar{\ideal L}$ による $\ideal M$ の二つの表示とし、$\ideal R$ と $\bar{\ideal R}$ の随伴素イデアルが一致しているとする。$\ideal L$ と $\bar{\ideal L}$ を約数 $\ideal L^*$ と $\bar{\ideal L}^{*}$ で置き換えて、表示が簡約されるようにする。すると $\ideal L^*$ と $\bar{\ideal L}^{*}$ の随伴素イデアルも一致する。定理 XI により、$\ideal L^*$ は $\bar{\ideal R}$ に対して相対的に素であり、$\bar{\ideal L}^{*}$ は $\ideal R$ に対して相対的に素である。
\[
 \ideal R\cdot\ideal L^*\equiv0(\bar{\ideal R}),
 \qquad
 \bar{\ideal R}\bar{\ideal L}^{*}\equiv0(\ideal R)
\]
であるから、
\[
 \ideal R\equiv0(\bar{\ideal R}),
 \qquad
 \bar{\ideal R}\equiv0(\ideal R),
 \qquad
 \ideal R=\bar{\ideal R}
\]
を得る。したがって、孤立イデアルはその随伴素イデアルによって一意に決定される。これは既約イデアルへの分解、あるいは極大準素イデアルへの分解についての定理 VII と IX を精密化するものであり、定理 IX への付加により、最短表示を仮定するだけで足りる。まとめると次のようになる。

\textbf{定理 XIII.} \emph{イデアルを既約イデアル、あるいは極大準素イデアルの最小公倍数として表す任意の最短表示において、孤立した既約イデアル、あるいは孤立した極大準素イデアルは一意に定まる。曖昧さは、非孤立の既約イデアル、あるいは非孤立の極大準素イデアルだけに関わる。}\footnote{多項式領域のイデアルについては、極大準素イデアルへの分解の場合に、この定理はすでに Macaulay によって証明なしに述べられていた。彼の孤立準素イデアルおよび非孤立（埋め込まれた）準素イデアルの定義は、以下で与えたものの非合理的な形と見ることができる。} \emph{一般に、孤立イデアルはその随伴素イデアルによって一意に決定される。}

このような既約イデアル、あるいは極大準素イデアルによる最短表示において、イデアル $\ideal B_i$、あるいは $\ideal Q_j$ が非孤立であるならば、定義により、その補元 $\ideal A_i$、あるいは $\ideal L_j$ は $\ideal P_i$、あるいは $\ideal P_j$ で割り切れる。したがって
\[
 \ideal A_i^{\varrho_i}\equiv0(\ideal B_i)
 \qquad \text{それぞれ} \qquad
 \ideal L_j^{\sigma_j}\equiv0(\ideal Q_j).
\]
逆に、これらの関係が成り立つならば、$\ideal P_i$、あるいは $\ideal P_j$ は補元の少なくとも一つの随伴素イデアルに含まれ、したがって定理 IX への付加により、$\ideal L_j^*$ に関して簡約された表示を与える $\ideal L_j$ の約数 $\ideal L_j^*$ の随伴素イデアルに含まれる。したがって $\ideal B_i$、あるいは $\ideal Q_j$ は非孤立である。特に、$\ideal M$ の分解において随伴素イデアルが複数回現れる既約イデアル $\ideal B_i$ は、常に非孤立である。\emph{非孤立準素イデアルは、したがって、任意の補元のある冪がそれによって割り切れるという事実によっても特徴づけられる。孤立準素イデアルは、それが起こり得ないという事実によって特徴づけられる。}

\subsection*{\S~8. 互いに素な既約イデアルの積としてのイデアルの一意表示。}

基礎に置かれた環領域 $\Sigma$ が単位をもつ、すなわち $\Sigma$ のすべての元 $a$ に対して $\varepsilon a=a$ となる元 $\varepsilon$ をもつとき、\footnote{乗法は可換であるから、もちろん $\Sigma$ は高々一つの単位しかもち得ない。実際、二つの単位 $\varepsilon_1$ と $\varepsilon_2$ に対しては、$\varepsilon_1\varepsilon_2=\varepsilon_2=\varepsilon_1$ である。} 互いに素なイデアルを次のように定義できる。

\textbf{定義 VIII.} \emph{二つのイデアル $\ideal R$ と $\ideal S$ は、それらの最大公約数が $\Sigma$ の全要素から成る単位イデアル $\ideal O=(\varepsilon)$ であるとき、互いに素であるという。イデアルが互いに素な意味で既約であるとは、二つずつ互いに素なイデアルの最小公倍数として表せないことをいう。}

二つの互いに素なイデアルは、常に互いに相対的に素であることに注意する。定義により、$\varepsilon=r+s$ となる元 $r\equiv0(\ideal R)$ と $s\equiv0(\ideal S)$ が存在する。$\ideal T\cdot\ideal R\equiv0(\ideal S)$ から $\ideal T\cdot r\equiv0(\ideal S)$ が従い、したがって $\ideal T\cdot\varepsilon=\ideal T\equiv0(\ideal S)$ である。同様に、$\ideal T\cdot\ideal S\equiv0(\ideal R)$ から $\ideal T\equiv0(\ideal R)$ が従う。\footnote{逆は成り立たない。例えば、$\ideal R=(x)$ と $\ideal S=(y)$ は互いに相対的に素であるが、互いに素ではない。} したがって補題 V により、二つずつ互いに素なイデアルによる任意の表示も簡約されている。

一意性の証明は定理 X の類似に基づく。

\textbf{定理 XIV.} \emph{$\ideal R$ がイデアル $\ideal S_1,\ldots,\ideal S_\lambda$ のそれぞれと互いに素であるならば、$\ideal R$ は $\ideal S=[\ideal S_1\cdots\ideal S_\lambda]$ とも互いに素である。逆に、$\ideal R$ と $\ideal S$ が互いに素であるならば、$\ideal R$ は各 $\ideal S_j$ と互いに素である。$\ideal R=[\ideal R_1\cdots\ideal R_\mu]$ であり、すべての $\ideal R_i$ がすべての $\ideal S_j$ と互いに素であるならば、$\ideal R$ と $\ideal S$ は互いに素である。ここでも逆が成り立つ。}

実際、$\ideal R$ が各 $\ideal S_j$ と互いに素であるならば、次を満たす元 $s_j$ が存在する。
\[
 s_j\equiv0(\ideal S_j),
 \qquad
 s_j\equiv\varepsilon(\ideal R).
\]
したがって
\[
 s_1\cdot s_2\cdots s_\lambda\equiv0(\ideal S),
 \qquad
 s_1\cdot s_2\cdots s_\lambda\equiv\varepsilon(\ideal R),
 \qquad
 (\ideal R,\ideal S)=(\varepsilon).
\]
$(\ideal R,\ideal S)$ は各 $(\ideal R,\ideal S_j)$ で割り切れるから、逆も従う。同じ推論を繰り返し適用すれば第二部が得られる。すなわち、固定した $i$ について $\ideal R_i$ が $\ideal S_1,\ldots,\ideal S_\lambda$ と互いに素であるならば、$\ideal R_i$ は $\ideal S$ と互いに素である。これがすべての $i$ について成り立つならば、互いに素であるという概念は相互的であるから、$\ideal S$ は $\ideal R$ と互いに素である。逆に、$\ideal R$ と $\ideal S$ が互いに素であれば、$\ideal S$ は $\ideal R_i$ と互いに素であり、したがって $\ideal R_i$ は $\ideal S$ と互いに素である。

互いに素な意味で既約なイデアルへの分解の存在と一意性\footnote{分解の存在は、ふたたび \S~2 と同様に直接証明できる。一意性の証明も直接に実行できる（序論で Schmeidler および Noether--Schmeidler について述べたことを参照）。ここで与える証明は、互いに素な意味で既約なイデアルの構造についての洞察も与える。}の証明は、相対的に素なイデアルの場合の対応する証明と同じく、一意な群分けの問題に帰着する。しかし、$\ideal M$ の相対的に素な既約イデアル $\ideal R_1,\ldots,\ideal R_\sigma$ は定理 XII によって一意に定義されているので、ここでは随伴素イデアルへ戻る必要はない。

$\ideal M$ の一意に定義された相対的に素な既約イデアル $\ideal R_1,\ldots,\ideal R_\sigma$ を群にまとめる。ただし、\emph{一つの群の任意のイデアルは、他のどの群の任意のイデアルとも互いに素であり、一方で一つの群は、一方の部分群の任意のイデアルが他方の部分群の任意のイデアルと互いに素であるような二つの部分群には分けられない}ようにする。このような群分けは次のように得られる。定義により、群 $G$ は各イデアル $\ideal R$ とともに、$\ideal R$ と互いに素でないすべてのイデアルを含まなければならない。それらを $\ideal R_{i_1}$ とする。これらと互いに素でないイデアルを $\ideal R_{i_1i_2}$ とする。一般に、$\ideal R_{i_1\ldots i_\lambda}$ を $\ideal R_{i_1\ldots i_{\lambda-1}}$ と互いに素でないイデアルとする。有限個のイデアルしか現れないので、この過程は有限回の後に停止する。このように得られたイデアルの集合は望ましい性質をもつ群 $G$ をなす。すなわち、$G$ の外にあるすべてのイデアルは、構成により、$G$ のすべてのイデアルと互いに素である。もし $G$ が二つの部分群 $G^{(1)}$ と $G^{(2)}$ に分けられ、$\ideal R_{i_1\ldots i_\lambda}$ が $G^{(1)}$ に属するならば、互いに素であることは相互的であるから、$G^{(1)}$ は $\ideal R_{i_1\ldots i_{\lambda-1}},\ldots,\ideal R_{i_1}$ も含まなければならず、したがって $\ideal R$ も、ひいては群 $G$ 全体も含まなければならない。これにより既約性が証明される。残りのイデアルについてこの構成を繰り返すと、すべての $\ideal R$ の群分け $G_1,\ldots,G_\tau$ が得られる。

この群分けは一意である。$G'_1,\ldots,G'_\tau$ が第二の群分けであり、$\ideal R_{i_1\ldots i_\lambda}$ が $G'_i$ の元であるとする。このとき $G'_i$ は $\ideal R$ を含み、したがって $G_1$ を含む。そして既約性条件により、$G_1$ と異なる元を含むことはできない。ゆえに $G'_i=G_1$ である。

$\ideal T_i$ を、群 $G_i$ にまとめられたイデアル $\ideal R$ の最小公倍数とする。
\[
 \ideal M=[\ideal T_1\cdots\ideal T_\tau]
\]
が $\ideal M$ の互いに素な意味で既約なイデアルによる表示であることを示す。まず $\ideal T$ を $\ideal R$ に分解すれば、$[\ideal T_1\cdots\ideal T_\tau]$ が本当に $\ideal M$ を表すことが分かる。定理 XIV により $\ideal T$ は二つずつ互いに素であり、各 $\ideal T$ は互いに素な意味で既約である。これで、そのような表示の存在が証明された。

一意性について、$\ideal M=[\bar{\ideal T}_1\cdots\bar{\ideal T}_{\bar\tau}]$ をこの型の第二の表示とする。$\bar{\ideal T}$ を相対的に素な既約イデアル $\ideal R$ に分解すると、定理 XIV は、異なる $\bar{\ideal T}$ に現れるイデアル $\ideal R$ が互いに素であり、したがって互いに相対的に素でもあることを示す。ゆえにそれらは、$\ideal M$ の一意に定義された相対的に素な既約イデアル $\ideal R$ と一致する。またイデアル $\bar{\ideal T}_i$ は、定理 XIV により、述べた性質をもつ $\ideal R$ の群分け $G'_i$ を誘導する。この群分けは一意であり、各 $\bar{\ideal T}_i$ は群 $G'_i=G_i$ によって一意に定まるから、$\bar{\ideal T}_i=\ideal T_i$ である。一意性が証明された。

二つずつ互いに素なイデアルについては、最小公倍数はさらに積に等しい。実際、定理 XIV により、$\ideal M=[\ideal T_1\cdots\ideal T_\tau]$ のとき、補元 $\ideal L_i$ も $\ideal T_i$ と互いに素である。したがって元
\[
 t_i\equiv0(\ideal T_i),
 \qquad
 l_i\equiv0(\ideal L_i),
 \qquad
 \varepsilon=t_i+l_i
\]
が存在する。$f\equiv0(\ideal T_i)$ および $f\equiv0(\ideal L_i)$ から
\[
 f=f\varepsilon=ft_i+fl_i,
 \qquad
 f\equiv0(\ideal L_i\cdot\ideal T_i)
\]
が従う。逆に $\ideal L_i\cdot\ideal T_i$ は $[\ideal L_i,\ideal T_i]$ で割り切れるから、$[\ideal L_i,\ideal T_i]=\ideal L_i\cdot\ideal T_i$ である。この手続きを $\ideal L_i$ に対して続けると
\[
 \ideal M=[\ideal T_1\cdots\ideal T_\tau]
       =\ideal T_1\cdot\ideal T_2\cdots\ideal T_\tau
\]
を得る。したがって次が成り立つ。

\textbf{定理 XV.} \emph{任意のイデアルは、有限個の二つずつ互いに素で、かつ互いに素な意味で既約なイデアルの積として、一意に表すことができる。}

\subsection*{\S~9. 研究の加群への拡張。既約加群への分解における成分数の一致。}

ここで、準素イデアルや素イデアルではなく既約イデアルに関わる最初の三つの節の内容が、より弱い仮定のもとでも成り立つことを示す。これらの節は乗法の可換律を用いず、イデアルが加群であるという性質だけを用いる。したがって、それらは、これから定義する非可換領域に関する加群についても成り立つ。加群の定義は、次の性質をもつ二重領域 $(\Sigma,T)$ に基づける。

$\Sigma$ は抽象的に定義された非可換環である。すなわち、$\Sigma$ は元 $a,b,c,\ldots$ の系であって、そこには環の加法 $(+)$ と環の乗法 $(\times)$ という二つの結合が定義されており、\S~1 に述べた法則のうち環乗法の可換律 4 を除くものを満たす。

$T$ は元 $\alpha,\beta,\gamma,\ldots$ の系であり、$\Sigma$ と関連してやはり二つの結合が定義されている。一つは加法であり、任意の二元 $\alpha,\beta$ から一意な第三の元 $\alpha+\beta$ を作る。もう一つは、$T$ の元 $\alpha$ に $\Sigma$ の元 $c$ を掛ける乗法であり、$T$ の元 $c\alpha$ を一意に作る。\footnote{したがって、ここで扱っているのは「右側」乗法、「右側」領域 $T$、したがって「右側」加群およびイデアルである。もし $T$ を左側乗法 $\alpha c$ に基づけるならば、左側加群および左側イデアルの対応する理論が得られる。この場合、$M$ は $\alpha$ とともに $\alpha\cdot c$ を含む。結合律は $(\gamma\cdot b)\cdot a=\gamma\cdot(b\times a)$ の形をとる。}

これらの結合について、次の法則が成り立つ。
\begin{enumerate}
\item 加法の結合律：$(\alpha+\beta)+\gamma=\alpha+(\beta+\gamma)$。
\item 加法の可換律：$\alpha+\beta=\beta+\alpha$。
\item 制限のない一意な減法の法則：方程式 $\alpha+\xi=\beta$ を満たす元 $\xi$ が $T$ の中にただ一つ存在する（$\xi=\beta-\alpha$ と書く）。
\item 乗法の結合律：$a\cdot(b\cdot\gamma)=(a\times b)\cdot\gamma$。
\item 分配律：
\[
 (a+b)\cdot\gamma=a\cdot\gamma+b\cdot\gamma,
 \qquad
 c\cdot(\alpha+\beta)=c\cdot\alpha+c\cdot\beta.
\]
\end{enumerate}
これらの条件から、よく知られているように、零元の存在と、減法と乗法の組合せに対する分配律
\[
 (a-b)\cdot\gamma=a\cdot\gamma-b\cdot\gamma,
 \qquad
 c\cdot(\alpha-\beta)=c\cdot\alpha-c\cdot\beta
\]
も従う。ただし $(\dot{-})$ は $\Sigma$ における減法を表す。$\Sigma$ が単位 $\varepsilon$ を含むならば、$T$ の任意の元 $\alpha$ に対して $\varepsilon\alpha=\alpha$ が成り立つべきものとする。

$(\Sigma,T)$ における加群 $M$ とは、$T$ の元の系であって次の二条件を満たすものをいう。
\begin{enumerate}
\item $\alpha$ とともに、任意の $\Sigma$ の元 $c$ に対して $c\alpha$ を含む。
\item $\alpha$ と $\beta$ とともに差 $\alpha-\beta$ を含み、したがって任意の整数 $n$ に対して $n\alpha$ も含む。\footnote{整数はここでも環の元ではなく、略記記号として理解される。}
\end{enumerate}
この定義により、$T$ 自身が $(\Sigma,T)$ における加群をなす。特に、領域 $T$ とそこで定められた結合が、領域 $\Sigma$ とそこで成り立つ結合と一致するならば、加群 $M$ は $\Sigma$ における右イデアル $M$ になる。さらに $\Sigma$ が可換であると仮定すれば、通常のイデアル概念が得られ、それは加群概念の特殊な場合である。\footnote{加群の最も簡単な例は整数係数線形形式の加群である。この場合、$\Sigma$ は有理整数全体から成り、$T$ は整数係数線形形式全体から成る。もう少し一般の加群は、$\Sigma$ と $T$ において有理整数の代わりに代数的整数、あるいは例えば偶整数全体を取れば得られる。線形形式の代わりに各係数複合体を一つの元と見なせば、$\Sigma$ と $T$ における結合は実際には異なる。非可換多項式環領域におけるイデアルは Noether--Schmeidler の共同論文の対象である。他の特殊な非可換領域におけるイデアルは、四元数の数論に関する Hurwitz の講義（Berlin, Springer 1919）およびそこに引用された Du Pasquier の諸論文で扱われている。}

加群については、\S~1 のすべての定義がそのまま有効である。したがって、$\alpha\eqzero{M}$、または $N\eqzero{M}$ とは、$\alpha$、または $N$ のすべての元が $M$ の元であることを意味する。言い換えれば、$\alpha$、または $N$ が $M$ で割り切れるということである。$M$ が $N$ と異なる元を含むとき、$M$ は $N$ の真の約数である。$N\eqzero{M}$ かつ $M\eqzero{N}$ からは $M=N$ が従う。最大公約数と最小公倍数の定義も、文字どおり同じままである。特に、加群 $M$ が有限個の元 $\alpha_1,\ldots,\alpha_\rho$ を含み、
\[
 M=(\alpha_1,\ldots,\alpha_\rho),
\]
すなわち任意の $\alpha\eqzero{M}$ について
\[
 \alpha=c_1\alpha_1+\cdots+c_\rho\alpha_\rho+n_1\alpha_1+\cdots+n_\rho\alpha_\rho
\]
となるならば、ただし $c_i$ は $\Sigma$ の元、$n_i$ は整数であるが、このとき $M$ を有限加群、$\alpha_1\ldots\alpha_\varrho$ を加群基底という。

以下では、\S~1 と同じく、有限性条件を満たす領域 $(\Sigma,T)$ だけを仮定する。すなわち、$(\Sigma,T)$ におけるすべての加群は有限であり、したがって加群基底をもつ。

このような領域 $(\Sigma,T)$ に対しては、そこで与えられた証明が示すように、有限鎖についての定理 I は加群にも成り立つ。したがって、\S\S~2 と 3 のすべての仮定が満たされる。最短表示および簡約表示についての定義 I と補題 I は直接移される。同様に、任意の加群を有限個の既約加群の最小公倍数として表せるという定理 II も移り、補題 II はこの種の任意の最短表示が同時に簡約されていることを示す。加群の可約性をその補元の性質によって表す定理 III も有効なままであり、そこから補題 III が従い、最後に、加群を既約加群の最小公倍数として表す二つの異なる最短表示における成分数の一致を述べる定理 IV が従う。定理 IV はまた、簡約表示についての補題 I の逆として補題 IV を与える。

\emph{付加。} 同じ議論は、両側領域 $T$、すなわち右側でも左側でもある領域について、どこでも加群を両側加群、すなわち $\alpha$ とともに $c\cdot\alpha$ および $\alpha\cdot c$ を含み、$\alpha$ と $\beta$ とともに $\alpha-\beta$ も含むものと理解しても、これらすべての定理と定義が有効であることを示す。

これらの定理がすべて除法関係と最小公倍数の概念だけに基づくのに対し、さらに先の一意性定理は本質的に積の概念に基づくので、直接の移行を許さない。したがって、準素イデアルおよび素イデアルの定義は加群へ移せない。なぜなら、$T$ の二元の積は定義されていないからである。形式的には可能な非可換環への移行も意味を失う。ここでは随伴素イデアルの存在を証明できず\footnote{$p_1^\rho\eqzero{\ideal Q}$ および $p_2^\sigma\eqzero{\ideal Q}$ から、ここでは $(p_1-p_2)^{\rho+\sigma}\eqzero{\ideal Q}$ が従わない。同様に、$(ab)^\rho\eqzero{\ideal Q}$ から $a^\rho b^\rho\eqzero{\ideal Q}$ は従わない。したがって、$\ideal P$ がイデアルであることも、素イデアルの性質をもつことも証明できない。両側イデアルについては、$\ideal P$ がイデアルであることは証明できるが、素イデアルであることは証明できない。}、既約イデアルが準素であることの証明も失敗するからである。これら二つの事実が、後続の一意性定理の基礎をなす。これに対して、非可換環が単位をもつならば、互いに素なイデアルおよび互いに素な意味で既約なイデアルを定義でき、定理 II の証明で用いた議論により、任意のイデアルは有限個の二つずつ互いに素で互いに素な意味で既約なイデアルの最小公倍数として表せることが分かる。\footnote{特殊な「完全可約」イデアルについては、ここでも一意性定理を立てることができる。Noether--Schmeidler の共同論文を参照。}

最後に、$(\Sigma,T)$ において有限性条件が成り立つための十分条件を述べる。$\Sigma$ が単位を含み、有限性条件を満たし、かつ $T$ 自身が $(\Sigma,T)$ における有限加群であるならば、$(\Sigma,T)$ の任意の加群は有限である。

実際、$\Sigma$ に単位が存在することから、
\[
 \ideal M=(f_1,\ldots,f_e),
 \qquad
 f=\bar b_1f_1+\cdots+\bar b_\varrho f_\varrho+n_1f_1+\cdots+n_\varrho f_\varrho
\]
において、$\bar b_i$ が $\Sigma$ の元、$n_i$ が整数であるならば、
\[
 f=b_1f_1+\cdots+b_\varrho f_\varrho
\]
という表示も、$b_i$ を $\Sigma$ の元として成り立つ。なぜなら $f_i=\varepsilon f_i$ であるから $n_if_i=n_i\varepsilon f_i$ であり、$n_i\varepsilon=(\varepsilon+\cdots+\varepsilon)$ は $\Sigma$ に属するので、$b_i=\bar b_i+n_i\varepsilon$ は $\Sigma$ の元である。$T$ についての仮定は、$T$ の任意の元 $\alpha$ が
\[
 \alpha=\bar a_1\tau_1+\cdots+\bar a_k\tau_k+n_1\tau_1+\cdots+n_k\tau_k
\]
という表示をもつこと、したがって
\[
 \alpha=a_1\tau_1+\cdots+a_k\tau_k
\]
という表示をもつことをいう。第二の表示は、$\varepsilon\tau_i=\tau_i$ であるため、前と同様に第一の表示から従う。

$\alpha$ が $(\Sigma,T)$ の加群 $M$ の元を動くとき、$\tau_k$ の係数 $a_k$ は $\Sigma$ のイデアル $\ideal M_k$ を動く。上により、任意の $a_k\eqzero{\ideal M_k}$ に対して
\[
 a_k=b_1a_k^{(1)}+\cdots+b_\varrho a_k^{(\varrho)}
\]
である。$\alpha^{(i)}$ を、$\tau_k$ の係数が $a_k^{(i)}$ である $M$ の元とする。このとき $\alpha-b_1\alpha^{(1)}\ldots-b_\varrho\alpha^{(\varrho)}$ は $M$ に属し、$\tau_1,\ldots,\tau_{k-1}$ だけに依存する元である。この手続きを、もはや $\tau_k$ を含まず加群 $M'$ をなすこれらの元全体について繰り返すことができ、有限回の繰返しで主張が証明される。

\subsection*{\S~10. 多項式領域という特殊な場合。}

1. 基礎に置かれる環領域 $\Sigma$ を、任意の複素係数をもつ $x_1,\ldots,x_n$ の多項式全体とする。Hilbert の基底定理（Ann. 36）により、この領域は有限性条件を満たす。問題は、われわれの定理と、消去論および加群論の既知の定理との関係である。

この関係は、Hilbert の既知の定理の次の特殊な場合によって与えられる。\footnote{Über die vollen Invariantensysteme. Math. Ann. 42 (1893), \S~3, p. 313.}

$f$ が、素イデアル $\ideal P$ のすべての多項式の零点である $x_1,\ldots,x_n$ の任意の有限値系（$\ideal P$ の零点）で消えるならば、$f$ は $\ideal P$ で割り切れる。言い換えれば、素イデアル $\ideal P$ は、その零点で消える多項式全体から成る。\footnote{Lasker が示したように [Math. Ann. 60 (1905), p. 607]、この特殊な場合は斉次形式について直接にも証明できる。逆に、Hilbert の定理は斉次および非斉次の場合にそこから再び従う。その定理は次のように表せる。イデアル $\ideal R$ が $M$ のすべての零点で消えるならば、$\ideal R$ のある冪は $M$ で割り切れる。この定理も特殊な場合も、$x$ の値域が代数的閉体であるときにだけ成り立つ。したがって、これはわれわれの定理だけからは従わず、根の存在を用いなければならない。例えば、零点が $x_1=0,\ldots,x_n=0$ だけから成るイデアルが、ある次元以降の $x$ の冪積をすべて含むことを用いる箇所でそうである。残りの証明は、われわれの定理を用いれば Lasker に比べていくらか簡略化できる。Lasker は、イデアルを極大準素イデアルに分解することを証明するためにも Hilbert の定理に訴えなければならないからである。}

したがって、積 $fg$ が $\ideal P$ のすべての零点で消えるならば、少なくとも一方の因子が消える。零点は既約代数多様体をなす。逆に、既約多様体をこのように定義するならば、既約多様体上で消える多項式全体は素イデアルをなす。したがって、素イデアルと既約多様体は一対一に対応する。$\ideal P^\rho\eqzero{\ideal Q}$ および $\ideal Q\eqzero{\ideal P}$ であるから、準素イデアルの零点はその随伴素イデアルの零点と一致する。\footnote{Macaulay（序論を参照）は、準素イデアルのこの性質、すなわち既約多様体をもつことを定義としている。Lasker は定義に多様体の多様体という概念だけを含め、それ以外は抽象的に定義する。$x_1=0,\ldots,x_n=0$ でだけ消える準素イデアルは、Lasker において特殊な位置を占める。} したがって、イデアルを極大準素イデアルの最小公倍数として表すことは、そのイデアルのすべての零点を既約多様体に分解することを与える。そして Lasker が示したように、逆も成り立つ。したがって、随伴素イデアルの一意性の証明は、ここでは代数多様体の既約多様体への一意分解に関する消去論の基本定理に対応し、領域の既約多項式による多項式の積表示が一意でなく、したがって消去論が存在しない特殊な多項式領域においては、その消去論の定理の代替物として役立ち得る。

孤立準素イデアルに対応する既約多様体は、まさに「最小レゾルベント」に現れるものにほかならない。\footnote{例えば J. König, Einleitung in die allgemeine Theorie der algebraischen Größen (Leipzig, Teubner, 1903), p. 235 を参照。} なぜなら、非孤立極大準素イデアルの零点は、少なくとも一つの孤立なものの零点でもあり、その随伴素イデアルは非孤立イデアルのそれによって割り切れるからである。したがって孤立準素イデアルの一意性は、指数において新しい不変な重複度数を与える。準素イデアルを既約なものへ分解する一意性定理も、重複度の意味で消去論の補足と見なすことができる。

これらの注意の後では、種々の分解は代数的対象に対する振舞いによって解釈できる。二つずつ互いに素なイデアルには共通零点をもたない多様体が対応する。互いに相対的に素なイデアルでは、一方のイデアルの既約多様体が同時に共通零点になることはない。極大準素イデアルは、互いにすべて異なる既約多様体においてだけ消える。既約イデアルへの分解では、同じ既約多様体が重複して現れることもある。

なお、一般の多項式領域の代わりに、すべての斉次形式の領域を基礎に取ることもできる。この場合、加法は同じ次元の形式についてだけ定義されるが、そこで有効な結合についても一般定理が有効であることは容易に確かめられる。\footnote{曖昧な場合には斉次分解のほかに非斉次分解も存在し得ることを、例 $(x^3,xy,y^3)=[(x^3,y),(y^3,x)]=[(xy,x^3,y^3,x+y^2),(xy,x^3,y^3,y+x^2)]$ が示している。}

以下に公式を示す四つの異なる分解の最も簡単な例は、前述のことに従えば、一つの直線と、それに交わる二つの直線であってその一方が交点とは異なる一点を高い重複度で含むものによって与えられる。互いに素な意味で既約なイデアルへの分解は、その直線とそれに斜交する多様体への分解に対応する。相対的に素な既約イデアルへの分解では、この多様体が二つの直線に分かれる。極大準素イデアルへの分解は、高重複度の点を切り離すことに対応し、既約イデアルへの分解はその点の分解を強制する。

この点を原点とし、通る直線を $y$ 軸、それと交わる直線を $x$ 軸に平行、斜交する直線を $z$ 軸に平行とすれば、そのような配置は例えば次の既約イデアルによって表される。\footnote{これらのうち最初の三つは素イデアルとして既約である。$\ideal B_4$ は約数 $(x^2,y,z)$ と $(x,y,z)$ だけをもつからであり、$\ideal B_5$ は任意の約数が多項式 $xy$ を含むからである。}
\[
\begin{gathered}
 \ideal B_1=(x-1,y),\quad
 \ideal B_2=(y-1,z),\quad
 \ideal B_3=(x,z),\quad
 \ideal B_4=(x^3,y,z),\quad
 \ideal B_5=(x^2,y^2,z).
\end{gathered}
\]
随伴素イデアルは
\[
 \ideal P_1=\ideal B_1,
 \quad \ideal P_2=\ideal B_2,
 \quad \ideal P_3=\ideal B_3,
 \quad \ideal P_4=\ideal P_5=(x,y,z).
\]
極大準素イデアルは
\[
 \ideal Q_1=\ideal B_1,
 \quad \ideal Q_2=\ideal B_2,
 \quad \ideal Q_3=\ideal B_3,
 \quad \ideal Q_4=[\ideal B_4,\ideal B_5]=(x^3,y^2,x^2y,z).
\]
相対的に素な既約イデアルは
\[
 \ideal R_1=\ideal Q_1,
 \quad \ideal R_2=\ideal Q_2,
 \quad \ideal R_3=[\ideal Q_3,\ideal Q_4]=(x^3,x^2y,xy^2,z).
\]
互いに素な意味で既約なイデアルは
\[
 \ideal S_1=\ideal R_1,
 \quad
 \ideal S_2=[\ideal R_2,\ideal R_3]=((y-1)x^3,(y-1)x^2y,(y-1)xy^2,z).
\]
したがって全イデアルは
\[
\begin{aligned}
 \ideal M&=[\ideal S_1,\ideal S_2]=\ideal S_1\ideal S_2 \\
 &=((x-1)(y-1)x^3,(y-1)x^2y,(y-1)xy^2,(x-1)z,y(y-1)x^3,yz),
\end{aligned}
\]
であり、ここで
\[
 1=-(y-1)x^3+(y-1)(x^3-1)+y,
\]
\[
 (y-1)(x^3-1)+y\eqzero{\ideal S_1},
 \qquad -(y-1)x^3\eqzero{\ideal S_2}.
\]
ここでは、イデアル $\ideal B_1$, $\ideal B_2$, $\ideal B_3$ は孤立しており、したがって既約イデアルおよび極大準素イデアルへの分解においても一意に定まる。イデアル $\ideal B_4$ と $\ideal B_5$、あるいは $\ideal Q_4$ は非孤立であり、一意に定まらないが、例えば
\[
 \ideal D_4=(x^3,y+\lambda x^2,z),
 \qquad
 \ideal D_5=(x^2+\mu xy,y^2,z)
\]
によって置き換えることができ、また $\ideal Q_4$ も同様に
\[
 \overline{\ideal Q}_4=(x^3,x^2y,y^2+\lambda xy,z)
\]
で置き換えることができる。

2. 一般の多項式領域および整多項式領域が有限性条件を満たすのと同じく、任意の有限整多項式領域も、加群基底に関する Hilbert の定理によって有限性条件を満たす。\footnote{逆に、多項式領域が有限性条件を満たし、すべての多項式が、乗数が $x$ に関して低い次数をもつような表示を少なくとも一つもつならば、その領域は有限整領域である。} 係数は任意の体に属させてよい。もう一つの例として、偶多項式全体の領域を挙げる。これは、$x^2\cdot y^2=(xy)^2$ のために、領域の既約多項式による多項式の積表示が一意でない最も簡単な領域である。これは前の例と同じ配置であり、いま次の既約イデアルによって与えられる。
\[
\begin{gathered}
 \ideal B_1=(x^2-1,xy,y^2,yz),\quad
 \ideal B_2=(y^2-1,xz,yz,z^2),\\
 \ideal B_3=(x^2,xy,xz,yz,z^2),\quad
 \ideal B_4=(x^4,x^3y,y^2,xz,yz,z^2),\\
 \ideal B_5=(x^2,y^2,xz,yz,z^2).
\end{gathered}
\]
随伴素イデアルは
\[
 \ideal P_1=\ideal B_1,
 \quad \ideal P_2=\ideal B_2,
 \quad \ideal P_3=\ideal B_3,
 \quad \ideal P_4=\ideal P_5=(x^2,xy,y^2,xz,yz,z^2).
\]
$\ideal P_1$ が素イデアルであることは、領域の任意の多項式が
\[
 f=\varphi(z^2)+xz\,\psi(z^2)\pmod{\ideal P_1}
\]
の形をもつことから従う。したがって
\[
 f_1=\varphi_1(z^2)+xz\,\psi_1(z^2),
 \qquad
 f_2=\varphi_2(z^2)+xz\,\psi_2(z^2)
\]
ならば、$f_1f_2\eqzero{\ideal P_1}$ から方程式
\[
 \varphi_1\varphi_2+z^2\psi_1\psi_2=0,
 \qquad
 \varphi_2\psi_1+\varphi_1\psi_2=0
\]
が得られ、したがって $f_1\eqzero{\ideal P_1}$ または $f_2\eqzero{\ideal P_1}$ である。全く同様に $\ideal P_2$ が素イデアルであることが分かる。$\ideal P_3$ は、領域の任意の多項式が $\ideal P_3$ を法として $y^2$ の多項式に合同であるから素イデアルである。$\ideal P_4$ は領域のすべての多項式から成る。したがって $\ideal B_1$, $\ideal B_2$, $\ideal B_3$ は素イデアルとして既約であり、$\ideal B_4$ と $\ideal B_5$ はそれぞれ唯一の真の約数 $\ideal P_4$ だけをもつので、必然的に既約である。

既約イデアルから極大準素イデアル
\[
\begin{gathered}
 \ideal Q_1=\ideal B_1,
 \quad \ideal Q_2=\ideal B_2,
 \quad \ideal Q_3=\ideal B_3,
 \quad \ideal Q_4=[\ideal B_4,\ideal B_5]=(x^4,x^3y,y^2,xz,yz,z^2),
\end{gathered}
\]
相対的に素な既約イデアル
\[
 \ideal R_1=\ideal Q_1,
 \quad \ideal R_2=\ideal Q_2,
 \quad
 \ideal R_3=[\ideal Q_3,\ideal Q_4]=(x^4,x^3y,x^2y^2,xy^3,xz,yz,z^2),
\]
および互いに素な意味で既約なイデアル
\[
 \ideal S_1=\ideal R_1,
 \quad \ideal S_2=[\ideal R_2,\ideal R_3]
 =((y^2-1)x^4,(y^2-1)x^3y,(y^2-1)x^2y^2,
 (y^2-1)xy^3,xz,yz,z^2)
\]
が得られる。第一の例と同様に
\[
 [\ideal B_4,\ideal B_5]=[\ideal D_4,\ideal D_5],
 \qquad
 [\ideal Q_3,\ideal Q_4]=[\ideal Q_3,\overline{\ideal Q}_4],
\]
ただし
\[
 \ideal D_4=(x^4,xy+\lambda x^2,\ldots),
 \qquad
 \ideal D_5=(x^2+\mu xy,\ldots),
 \qquad \lambda\mu\ne1,
\]
であり、また
\[
 \overline{\ideal Q}_4=(x^4,x^3y,y^2+\lambda xy,\ldots)
\]
である。残る既約イデアル、あるいは極大準素イデアル $\ideal B_1$, $\ideal B_2$, $\ideal B_3$ は、孤立イデアルとして一意に決定される。

\subsection*{\S~11. 数論および微分式論からの例。}

1. 基礎領域 $\Sigma$ を偶有理整数全体とする。したがって、$\Sigma$ の各数 $2a$ に数 $a$ を対応させることにより、$\Sigma$ は有理整数全体の領域と一対一に対応させられる。このことから直ちに、$\Sigma$ の任意のイデアルは主イデアル $(2a)$ であることが従う。ただし、イデアルの任意の元 $2c$ について基底表示 $2c=n\cdot2a$ に現れる奇整数 $n$ は、有限和の略記にすぎない。

この領域の素イデアルは $\ideal P_0=\ideal O=(2)$ および $\ideal P=(2p)$ である。ただし $p$ は奇素数を表す。したがって任意の素イデアルは $\ideal P_0$ で割り切れるが、それ以外の素イデアルでは割り切れない。準素イデアルは
\[
 \ideal Q_0=(2\cdot2^{\varrho_0})
 \qquad\text{および}\qquad
 \ideal Q_\varrho=(2p^\varrho)
\]
である。それらは同時に既約イデアルであり、素イデアルについて述べたことにより、異なる奇素数に属する二つのイデアルは互いに相対的に素であるが、どの $\ideal Q$ もどの $\ideal Q_0$ に対しても相対的に素ではない。\footnote{実際、奇素数 $p_1\ne p_2$ に対して $2b\cdot2p_1^{\varrho_1}\eqzero{2p_2^{\varrho_2}}$ からは常に $2b\eqzero{2p_2^{\varrho_2}}$ が得られる。しかし $2b\cdot2p^\varrho\eqzero{2\cdot2^{\varrho_0}}$ からは $2b=2\cdot2^{e_0-1}\pmod{2\cdot2^{e_0}}$ だけが得られる。} したがって $\ideal Q_0$ は唯一の非孤立準素イデアルである。$a$ の素数冪への一意分解は、イデアル $(2a)$ の極大準素で同時に既約なイデアルによる一意表示
\[
 (2a)=[(2\cdot2^{e_0}),(2p_1^{e_1}),
\ldots,(2p_\mu^{e_\mu})]
\]
に対応する。したがってここでは、多項式領域からの例と対照的に、非孤立な極大準素イデアルまで一意に定まる。$\varrho_0=0$ の場合、これは同時に互いに相対的に素なイデアルによる表示であり、$\varrho_0>0$ の場合にはこのイデアルは相対的に素な意味で既約である。

したがって、整数全体の領域では四つの異なる分解が一致するのに対し、ここでは一方で極大準素イデアルへの分解と既約イデアルへの分解だけが一致し、他方で $\varrho_0>0$ のときに、互いに素な意味で既約なイデアル（領域に単位がないので任意のイデアルは互いに素な意味で既約である）と相対的に素な意味で既約なイデアルが一致する。$\varrho_0=0$ の場合には、互いに素な意味で既約な分解と相対的に素な意味で既約な分解は異なる。同時にここでは、一つの素イデアルが別の素イデアルで割り切れるが同一ではない例もすでに得られる。さらに一般に、積表示は除法関係からは従わないことが分かる。この後者の事実は、領域に単位がないことの帰結であり、また、すべてのイデアルが主イデアルであるにもかかわらず、領域の数を領域の既約数の積として一意に表示することが存在しない理由でもある。したがって最小公倍数の導入が必要になる。なお、すべての偶整数の代わりに、固定された素数または素数冪で割り切れるすべての整数を取っても、状況は全く同じである。

しかし、$\Sigma$ が合成数 $g=p_1^{\nu_1}\cdots p_s^{\nu_s}$ で割り切れるすべての数から成る場合には、既約イデアルと準素イデアルも異なる。ここでも任意のイデアルは主イデアル $(g\cdot a)$ であり、素イデアルは再び $\ideal P_0=\ideal O=(g)$ および $\ideal P=(g\cdot p)$ である。ただし $p$ は $g$ に現れる素数とは異なる素数である。しかし既約イデアルは $\ideal B_{\lambda_i}=(g\cdot p_i^{\lambda_i})$ および $\ideal B_e=(g\cdot p^e)$ であり、準素イデアルは $\ideal Q_e=\ideal B_e$ と、既約イデアルとは異なる
\[
 \ideal Q_{\lambda_1\ldots\lambda_s}=(g\cdot p_1^{\lambda_1}\cdots p_s^{\lambda_s})
\]
である。ここで $\ideal B_{\lambda_i}$ と $\ideal Q_{\lambda_1\ldots\lambda_\nu}$ はすべて同じ随伴素イデアル $\ideal P_0=(g)$ をもつ。ここでも既約イデアルへの分解は一意であり、したがって極大準素イデアルへの分解も一意である。よってここでも非孤立イデアルは一意に定まる。

2. 非可換環領域の例は、Noether--Schmeidler の論文で扱われた非可換多項式領域におけるイデアル論によって与えられる。特に「完全可約」イデアル、すなわち分解の成分が二つずつ互いに素で真の約数をもたないイデアルがある。したがってそれらの成分は \emph{a fortiori} 既約である。よって、そこにおいて \S~9 に従って証明された同型に加えて、二つの異なる分解における成分数の一致も得られる。その論文の特殊な場合として現れる、偏または常線形微分式の系の分解については、これは、一つの常線形微分式という既知の場合においてすら注意されていないと思われる結果を与える。

同時に、固定されたイデアル $M$ のすべての剰余類の系 $T$ は、非可換多項式領域 $\Sigma$ とともに二重領域 $(\Sigma,T)$ を与え、$T$ は $\Sigma$ に関して加群性をもつ。実際、二つの剰余類の差は再び剰余類であり、剰余類に任意の多項式を掛けたものも同様である。一方、二つの剰余類の積は存在しない（loc. cit., \S~3）。したがって、そこで「部分群」と呼ばれた剰余類の系は、基礎環領域 $\Sigma$ が非可換であるような二重領域 $(\Sigma,T)$ における加群の例を与える。

\subsection*{\S~12. 初等因子論からの例。}

これは一般的な展開によって要請される初等因子論の捉え方である。ただし、理論そのものは既知のものとして仮定する。

$\Sigma$ を、通常の行列の意味で加法と乗法を定義した、$n^2$ 個の成分をもつ整行列全体の領域とする。このとき $\Sigma$ は非可換環領域であり、したがってイデアルは一般には片側であり、両側イデアルは特殊な場合にすぎない。\footnote{これらの領域のイデアル論は Du Pasquier の諸研究の対象である。Zahlentheorie der Tettarionen, dissertation, Zürich, Vierteljahrsschr. d. Naturf. Ges. Zürich, 51 (1906); Zur Theorie der Tettarionenideale, ibid., 52 (1907)。第二論文の内容は、任意のイデアルが主イデアルであることの証明である。}

まず、任意のイデアルが主イデアルであることを示す。このため、右イデアルの場合には、各行列 $A=(a_{ik})$ に整数係数線形形式の加群
\[
 A=(a_{11}\xi_1+\cdots+a_{1n}\xi_n,
 \ldots,
 a_{n1}\xi_1+\cdots+a_{nn}\xi_n)
\]
を対応させる。逆に、この加群には $\mathsf{A}$ の基底を与える任意の行列が対応し、したがって $A$ のほかに、$U$ をユニモジュラーとして $UA$ も対応する。より一般に、積 $PA$ には $A$ の倍数である加群 $B$ が対応する。$\mathsf{A}$ からの一つの線形形式は、非零行を一つだけもつような行列 $P$ によって与えられる。

いま、$A_1,A_2,\ldots,A_\nu,\ldots$ をイデアル $\ideal M$ のすべての元とし、$\mathsf{A}_1,\mathsf{A}_2,\ldots,\mathsf{A}_\nu,\ldots$ をそれらに割り当てられた加群とする。それらの最大公約数を $\mathsf{A}$ とし、この加群 $A$ に割り当てられた最も一般の行列を $UA$ とする。すると、$\mathsf{A}$ の各一つの線形形式には、最大公約数の定義により、行列 $P_1A_{i_1}+\cdots+P_\sigma A_{i_\sigma}$ が対応する。ここで、前と同じく、$P$ は非零行を一つだけもつ。したがって、$\mathsf{A}$ の基底に対応する行列 $A$ も、この種の表示を、今度は一般の $P$ によってもつことになり、したがって $\ideal M$ の元となる。さらに任意の加群 $\mathsf{A}_i$ は $\mathsf{A}$ で割り切れるから、任意の行列 $A_i$ も $A$ で割り切れる。$A$、より一般には $UA$ が $\ideal M$ の基底をなす。左イデアルを扱う場合には、対応して各行列の列を加群の基底と見なすべきである。任意のイデアルは主イデアルであり、そこでは $A$ のほか、$V$ を任意のユニモジュラー行列として $AV$ も基底となる。

初等因子論との関係のため、以下では両側イデアルを基礎に置く。したがって特に、$A$ がイデアルに属すれば $PAQ$ もそのイデアルに属する。前述により、このようなイデアルの最も一般の基底は $UAV$ であり、$U$ と $V$ はユニモジュラーである。したがって基底元は同値な行列の一つの類を尽くす。\footnote{片側イデアルの基底元は右または左の類に対応する。} そしてイデアルと類の間、したがってイデアルとその類の初等因子系 $(a_1\mid a_2\mid\cdots\mid a_n)$ の間には一対一の関係がある。ここで $a_i$ は、よく知られているように、非負整数であり、それぞれ次のものを割り切る。したがって、初等因子によって定められる標準形に現れる類の行列は、イデアルの特殊な基底と見なせる。イデアル、あるいは類の除法関係は、初等因子の除法関係を意味し、逆も成り立つ。

しかし Du Pasquier は、loc. cit. \S~11 で、任意の両側イデアルについて階数は $n$ であり、すべての初等因子が一致することを示した。一般の初等因子の場合を考えるためには、したがってイデアルからではなく、直接に両側類（これも大文字のドイツ文字で表す）から出発しなければならない。類 $\mathfrak A=UAV$ が別の類 $\mathfrak B$ で割り切れるとは、$A=PBQ$ であることをいう。

一般に、\emph{二つの類の最小公倍数（最大公約数）は、対応する初等因子系の最小公倍数（最大公約数）を作ることによって得られる}。\footnote{論文 Zur Theorie der Moduln, Math. Ann. 52 (1899), p. 1 において、E. Steinitz は、類の最小公倍数（最大公約数）を初等系の最小公倍数および最大公約数によって定義している。独立に、類の最小公倍数は H. Brandt, Komposition der binären quadratischen Formen relativ einer Grundform, J. f. M. 150 (1919), p. 1 において「合同合成」として現れる。} 実際、
\[
 (a_1\mid a_2\mid\cdots\mid a_n),\quad
 (b_1\mid b_2\mid\cdots\mid b_n),\quad
 (c_1\mid c_2\mid\cdots\mid c_n),
\]
ただし $c_i=[a_i,b_i]$ とし、これらをそれぞれ $\mathfrak A$, $\mathfrak B$, $\mathfrak C^*$ の初等因子系とし、$\mathfrak C=[\mathfrak A,\mathfrak B]$ とする。このとき $\mathfrak C^*$ は $\mathfrak A$ と $\mathfrak B$ で割り切れ、したがって $\mathfrak C$ で割り切れる。逆に、$\mathfrak C$ の初等因子は $\mathfrak C^*$ のそれで割り切れるので、$\mathfrak C$ は $\mathfrak C^*$ で割り切れ、したがって $\mathfrak C=\mathfrak C^*$ である。最大公約数についての証明も同様に得られる。

したがって、初等因子 $a_i$ の素数冪の最小公倍数としての一意表示は、$\mathfrak A$ を、ただ一つの素数の冪によって与えられる初等因子をもつ類 $\mathfrak Q$ の最小公倍数として表すことに対応する。記号で書けば
\[
 \mathfrak Q\sim(p^{r_1}\mid p^{r_2}\mid\cdots\mid p^{r_\varrho}\mid0\mid\cdots\mid0),
 \qquad r_1\leqq r_2\leqq\cdots\leqq r_\varrho.
\]
とくに階数 $\varrho$ が $n$ に等しいなら、非可換領域であるにもかかわらず、互いに素で、かつ互いに素な意味で既約な類への一意分解が得られる。

類 $\mathfrak Q$ はさらに、次の初等因子系に対応する類へ分解できる。
\[
\begin{gathered}
 \mathfrak B_1\sim(p^{r_1}\mid\cdots\mid p^{r_1}),\quad
 \mathfrak B_2\sim(1\mid p^{r_2}\mid\cdots\mid p^{r_2});\quad \ldots;\\
 \mathfrak B_\varrho\sim(1\mid\cdots\mid1\mid p^{r_\varrho}\mid\cdots\mid p^{r_\varrho});\quad
 \mathfrak B_{\varrho+1}\sim(1\mid\cdots\mid1\mid0\mid\cdots\mid0).
\end{gathered}
\]
ここで $\mathfrak B_\nu$ では、そのつど数 $1$ が $(\nu-1)$ 回現れる。例えば $r_1=r_2=\cdots=r_\mu=0$ なら、$\mathfrak B_1,\ldots,\mathfrak B_\mu$ は単位類に等しいので、分解から省かなければならない。$\varrho=n$ のときの $\mathfrak B_{\varrho+1}$ についても同じである。さらに、例えば $r_\nu=r_{\nu+1}=\cdots=r_{\nu+\lambda}$ なら、$\mathfrak B_{\nu+1},\ldots \mathfrak B_{\nu+\lambda}$ は $\mathfrak B_\nu$ の真の約数となり、これらも除かなければならない。残るものを $\mathfrak B_{i_1},\ldots,\mathfrak B_{i_k}$ と記し、これらがいま最短表示を与えるものとすれば、
\[
 \mathfrak Q=[\mathfrak B_{i_1},\ldots,\mathfrak B_{i_k}]
\]
は $\mathfrak Q$ の既約類への一意分解である。実際、
\[
 \mathfrak B_\nu\sim(1\mid\cdots\mid1\mid p^{r_\nu}\mid\cdots\mid p^{r_\nu})
\]
が
\[
 \mathfrak C\sim(1\mid\cdots\mid1\mid p^{s_1}\mid\cdots\mid p^{s_\lambda})
 \quad\text{と}\quad
 \mathfrak D\sim(1\mid\cdots\mid1\mid p^{t_1}\mid\cdots\mid p^{t_\mu})
\]
の最小公倍数として表されるとする。このとき $\mathfrak C$ と $\mathfrak D$ の初等因子系では、最初の $(\nu-1)$ 箇所に数 $1$ がなければならず、$\nu$ 番目の箇所では指数 $s_1$ または $t_1$、例えば $s_1$ が $r_\nu$ に等しくなければならない。しかし
\[
 r_\nu=s_1\leqq s_2\leqq\cdots\leqq s_\lambda\leqq r_\nu
\]
となるので、$\mathfrak C=\mathfrak B_\nu$ であり、したがって $\mathfrak B_\nu$ は\emph{既約である}。\footnote{$\mathfrak B$ は片側類にはなお可約である。ここでは初等因子と類との間に一意な関係がもはや存在せず、したがって既約片側類への一意分解も存在しない。H. Brandt から私に与えられた次の例は、右側類への分解についてこれを示している（ここで類は基底行列、または対応する加群によって表される）：
\[
\begin{gathered}
 \mathfrak B=[\mathfrak C_1,\mathfrak C_2]=[\mathfrak D_1,\mathfrak D_2],\quad
 \mathfrak B\sim\begin{pmatrix}p&0\\0&p\end{pmatrix},\quad
 \mathfrak C_1\sim\begin{pmatrix}1&0\\0&p\end{pmatrix},\quad
 \mathfrak C_2\sim\begin{pmatrix}p&0\\0&1\end{pmatrix},\\
 \mathfrak D_1\sim\begin{pmatrix}1&0\\0&p\end{pmatrix},
 \quad
 \mathfrak D_2\sim\begin{pmatrix}p&0\\(p-1)&1\end{pmatrix}.
\end{gathered}
\]
実際、加群 $(\xi,p\eta)$、$(p\xi,\eta)$、および $(p\xi,(p-1)\xi+\eta)$ はそれぞれ唯一の真の約数として加群 $(\xi,\eta)$ だけをもち、したがって既約であり、互いに異なる。} $\mathfrak B_{\varrho+1}$ についても同じであり、その場合には $p$ を $0$ で置き換える。しかし、最小公倍数の形成が示すように、これらの既約類のそれぞれは、$\mathfrak Q$、あるいは $\mathfrak A$ の初等因子系において、ある定まった指数と、その指数が最初に現れる場所とを与え、$\mathfrak B_{\varrho+1}$ は階数を与える。これらの数は $\mathfrak Q$、あるいは $\mathfrak A$ の初等因子系によって一意に定まり、初等因子と類との関係は一対一であるから、$\mathfrak Q$、そして任意の類の、既約類への分解は一意である。

要約すれば、\emph{成分数が有界な整行列の任意の両側類 $\mathfrak A$ は、有限個の既約両側類の最小公倍数として一意に表すことができる。各既約類は、$\mathfrak A$ の初等因子系のある定まった素因子、随伴する指数、およびその指数が最初に現れる場所を表す。因子 $0$ に対応する既約類は $\mathfrak A$ の階数を与える。}

\begin{flushleft}
Erlangen, 1920年10月。
\end{flushleft}
\begin{center}
(1920年10月16日受理。)
\end{center}


\clearpage
\setcounter{footnote}{0}
\section*{20. 絶対既約性のための代数的判定法}
\begin{center}
Math. Ann. 85 (1922), pp. 26--33
\end{center}

任意の抽象的に定義された体から係数をとる多項式は、その係数領域を拡大して得られる代数的閉体においても既約であり続けるとき、\emph{絶対既約}であるという。\footnote{任意の体が、本質的に一意に、代数的閉体、すなわち一変数の任意の多項式が一次因子に分解する体へ拡大できることは、E. Steinitz が \emph{Algebraische Theorie der Körper} (J. f. M. 137 (1910), p.~167) で示した。これにより彼は代数学の基本定理の有理的同等物を与えた。} 以下で私は、与えられた体に関する既約性とは対照的に、絶対既約性が代数的に把握できることを示す。より正確には次の定理が成り立つ。

\emph{$n\geqq2$ 変数で不定係数をもつ任意の多項式に対して、それらの係数とさらに別の不定元についての、整数係数をもつ整有理関数、すなわち可約性形式を対応させることができる。その対応は、同じ次数の任意の特殊な多項式について、絶対既約性の必要十分条件が可約性形式の非消滅で与えられるようなものである。言い換えれば、可約性形式が不定元について恒等的に消滅し、かつ次数低下を生じさせないような係数の任意の特殊値系に対して、その特殊な多項式は可約であり、また逆も成り立つ。}

多項式の代わりに一つ多い変数に関する斉次形式を考えるなら、次数についての条件は、係数系が恒等的に消滅しないという、より単純な条件に置き換えられる。

この定理から直接の帰結として、A. Ostrowski が最初に述べた次の命題も得られる。\footnote{Zur arithmetischen Theorie der algebraischen Größen. Gött. Nachr. 1919, p.~279（補題 p.~296）。通常の数から成る体の場合については、私の知るところでは、Ostrowski も上の主定理に到達していた。彼はこの点に関する研究を近く公表する予定である。私の証明の配置において主定理が任意の抽象的に定義された体に対して成り立つのは、不定元について構成された可約性形式が、基礎体として特に選んだ体に依存しない整数係数をもつからである。}

\emph{係数が代数的数である任意の絶対既約多項式は、その係数領域を含む有限代数体 $\Omega$ の任意の素イデアルを法として既約であり続ける。ただし例外となる $\Omega$ の素イデアルは高々有限個である。特に、$\Omega$ は有理数体であってもよい。}

相対的整関数の理論へのさらなる応用\footnote{D. Hilbert, Mathematische Probleme. Gött. Nachr. 1900, p.~253, 問題 14。}については、別の機会に論じるつもりである。

1. まず、\emph{$n\geqq2$ 変数で不定係数をもつ任意の多項式が絶対既約である}ことを示す。すなわち
\begin{equation}
\begin{aligned}
F(x)&=\sum A_{i_1\ldots i_n}x_1^{i_1}\cdots x_n^{i_n} &&(i_1+\cdots+i_n\le l),\\
G(x)&=\sum B_{j_1\ldots j_n}x_1^{j_1}\cdots x_n^{j_n} &&(j_1+\cdots+j_n\le m),\\
H(x)&=F(x)G(x)=\sum C_{k_1\ldots k_n}(A,B)x_1^{k_1}\cdots x_n^{k_n} &&(k_1+\cdots+k_n\le l+m),
\end{aligned}
\tag{1}
\end{equation}
とおく。ここで $A$ および $B$ は不定元を表し、$C(A,B)$ は $A$ と $B$ の整数係数の双線形結合である。したがって商
\[
D=C(A,B)/C_{0\ldots0}(A,B)
\]
は、商 $A/A_{0\ldots0}$ および $B/B_{0\ldots0}$ の有理関数である。しかし、これらの関数の個数は引数の個数より多い。実際、それぞれの個数は
\[
\binom{l+m+n}{n}-1,
\qquad
\binom{l+n}{n}-1,
\qquad
\binom{m+n}{n}-1,
\]
であり、
\[
\binom{l+m+n}{n}+1>
\binom{l+n}{n}+\binom{m+n}{n}
\quad(n\ge2,\ l>0,\ m>0)
\]
が成り立つからである。ゆえに、$D(A,B)$ の間に少なくとも\emph{一つ}の代数的関係が存在し、したがって $C(A,B)$ の間にもそのような関係が存在する。
\begin{equation}
 \Phi(Z)\ne0\quad[\text{$Z_{k_1\ldots k_n}$ について恒等的}],
 \qquad
 \Phi(C(A,B))=0\quad[\text{$A,B$ について恒等的}],
\tag{2}
\end{equation}
ここで $\Phi$ は整数係数多項式を表す。

したがって、不定元 $Z$ を係数にもつ多項式
\begin{equation}
 E(x)=\sum Z_{k_1\ldots k_n}x_1^{k_1}\cdots x_n^{k_n}
 \qquad(k_1+\cdots+k_n\le l+m)
\tag{3}
\end{equation}
は、$n\geqq2$ では必然的に絶対既約である。\footnote{2 の用語でいえば、不定元 $Z$ は $\mathfrak P$ の零点ではないからである。}

2. 置換 $Z=C(A,B)$ の下で $A,B$ について恒等的に消滅するこれらの多項式 $\Phi(Z)$ 全体は、\emph{多項式のイデアル}\footnote{このような全体は通常、加群と呼ばれる。しかし実際には、その特徴的性質、すなわち $\varPhi_1$ と $\varPhi_2$ とともに $\varPhi_1\mp\varPhi_2$ も含み、$\varPhi$ とともに $\varPsi\cdot\varPhi$ も含むという性質（ここで $\varPsi$ は $Z$ に関する任意の整数係数多項式である）は、イデアルの性質である。}、より正確には素イデアル $\mathfrak P$ をなす。なぜなら、置換によって積 $\varPhi_1\varPhi_2$ が消滅するなら、少なくとも一方の因子が消滅するからである。(2) は $A,B$ について恒等的に満たされるので、$\mathsf{A},\mathsf{B}$ の任意の特殊値系に対しても成り立つ。このとき $\mathsf{A},\mathsf{B}$、したがって $\mathsf{\Gamma}=C(\mathsf{A},\mathsf{B})$ は何らかの体の元を表す。\emph{したがって、拡大体において可約である任意の多項式の係数 $\mathsf{\Gamma}$ は条件 $\varPhi(\mathsf{\Gamma})=0$ を満たし、$\mathfrak P$ の零点である}。同じことは $\mathfrak P$ の有限個の基底多項式についても言える。残るのは逆を証明することである。

このためには、新しい変数 $y$ を導入して (1) の多項式を次数 $l,m,l+m$ の斉次形式 $F(x,y)$、$G(x,y)$、$H(x,y)$ に変えても、$A,B,C$ の間のすべての関係が保たれることに注意する必要がある。したがって、たとえば $F(x,y)$\footnote{$F,G,\ldots,f,g,\ldots$ における不定元を特殊値系で置き換えて得られる多項式および形式は、以下を通じて上線を付して表す：$\bar F,\bar G,\ldots,\bar f,\bar g,\ldots$。}が $y$ の冪になる場合にも、その係数 $\Gamma$ は $\mathfrak P$ の零点となる。この場合、非斉次多項式への移行によって生じるのは $\bar H(x)$ の次数低下であって、因数分解ではない。以下では、この次数低下を除けば逆が常に成り立つことを示す。特に、\emph{斉次形式については、すべての $\mathsf{\Gamma}$ が消滅しないかぎり、逆は例外なく成り立つ。言い換えれば、$\mathfrak P$ の任意の零点は、ある拡大体の元 $A,B$ によるパラメータ表示 $\mathsf{\Gamma}=C(\mathsf{A},\mathsf{B})$ から得られる}。\footnote{「一般には」$\mathfrak P$ の零点がパラメータ表示によって与えられることは、素イデアル $\mathfrak P$ が既約代数多様体を定めるという性質から従う。しかし、よく知られているように、そこからは、任意の特殊値系に対してパラメータ表示が失敗しないということは従わない。これは私が当初見落としていた点であり、A. Ostrowski がしばらく前に注意を促してくれた点である。ここで与える証明は、より初等的な概念に基づく。}

この証明を、有限個の関数 $\varPhi(Z)$ を直接構成することによって行う。すなわち Kronecker の置換
\[
 x_\lambda=\xi^{d^{\,n-\lambda}}
\]
によって、(1) の多項式に\emph{一}変数の多項式を対応させる。ここで指数に空隙が生じることを示し、対応する係数のノルムを形成することによって、関数 $\varPhi(Z)$ と求める判定法に到達する。

3. 次のようにおく。\footnote{Festschrift, p.~11。}
\begin{equation}
 d=l+m+1,\qquad
 i=i_1d^{n-1}+i_2d^{n-2}+\cdots+i_n,
\quad
 j=j_1d^{n-1}+\cdots+j_n,
\quad
 k=k_1d^{n-1}+\cdots+k_n.
\tag{4}
\end{equation}
すると、Kronecker の置換によって、(1) の多項式 $F,G,H$ にはそれぞれ次の多項式が対応する。
\begin{equation}
\begin{aligned}
 f(\xi)&=\sum A_{i_1\ldots i_n}\xi^i &&(i_1+\cdots+i_n\le l),\\
 g(\xi)&=\sum B_{j_1\ldots j_n}\xi^j &&(j_1+\cdots+j_n\le m),\\
 h(\xi)&=\sum C_{k_1\ldots k_n}(A,B)\xi^k &&(k_1+\cdots+k_n\le l+m).
\end{aligned}
\tag{5}
\end{equation}
この対応は周知のように一対一である。実際、(5) に現れるすべての誘導された指数について、$k=0$ ならば $k_1=0,\ldots,k_n=0$ であり、したがって $i+j=i'+j'$ からも $i_1+j_1=i'_1+j'_1;\ldots;i_n+j_n=i'_n+j'_n$ が従う。ゆえに、異なる積 $\xi^i\xi^j$ が同じ指数 $\xi^k$ を与えるのは、対応する積 $x_1^{i_1}\ldots x_n^{i_n}x_1^{j_1}\ldots x_n^{j_n}$ も同じ指数に至る場合だけである。したがって
\begin{equation}
 h(\xi)=f(\xi)g(\xi)
\tag{6}
\end{equation}
が得られ、さらに、$f(\xi)$ と $g(\xi)$ に誘導された指数以外の指数が現れないような形で関係 (6) が成り立つなら、逆向きに
\begin{equation}
 H(x)=F(x)G(x)
\tag{7}
\end{equation}
が得られる。

ここで、\emph{$f(\xi)$ と $g(\xi)$ の誘導された指数には実際に空隙が生じる}。$f(\xi)$ の最高指数は $l\cdot d^{n-1}$ であり、第二の最高指数は $(l-1)d^{n-1}+d^{n-2}$ である。その差は $d^{n-2}\cdot(d-1)>1$ である。なぜなら $d=l+m+1\geqq3$ かつ $n\geqq2$ だからである。同じことが $g(\xi)$ についても成り立つ。

不定元 $A,B$ について成り立つ関係 (6) および (7) は、$\mathsf{A},\mathsf{B}$ の任意の特殊値系に対しても成り立つ。このとき次数を保つために、変数 $\eta$ と $y$ によって (6) および (7) を斉次化する。\emph{したがって、斉次形式 $\bar{H}(x,y)$ が代数的拡大体において分解するための必要十分条件は、その拡大体において対応する二元形式 $\bar{h}(\xi,\eta)$ が、各因子に誘導された指数以外の指数が現れないように二つの因子へ分解できることである。}

4. 2 の末尾で述べたノルムの形成を行うために、
\[
 a_1,b_1;\ldots; a_p,b_p;\quad a_{p+1},b_{p+1};\ldots; a_{p+q},b_{p+q}
\]
を新しい不定元とする。次のようにおく。
\begin{equation}
\begin{aligned}
 \varphi(\xi,\eta)&=(a_1\xi+b_1\eta)\cdots(a_p\xi+b_p\eta)
     =r_p\xi^p+r_{p-1}\xi^{p-1}\eta+\cdots+r_0\eta^p,\\
 \psi(\xi,\eta)&=(a_{p+1}\xi+b_{p+1}\eta)\cdots(a_{p+q}\xi+b_{p+q}\eta)
     =s_q\xi^q+s_{q-1}\xi^{q-1}\eta+\cdots+s_0\eta^q,\\
 \chi(\xi,\eta)&=\varphi(\xi,\eta)\psi(\xi,\eta)
     =t_{p+q}\xi^{p+q}+t_{p+q-1}\xi^{p+q-1}\eta+\cdots+t_0\eta^{p+q},
\end{aligned}
\tag{8}
\end{equation}
したがって $r,s,t$ は斉次化された基本対称関数を表す。すなわち、それらは各行に関して線形な、行 $a,b$ の対称関数であり、各個の行に関して斉次で同じ次元をもつすべての整対称関数は、これらから整数係数で整に、しかも一意に構成される。

さらに別の不定元 $u_{\mu\nu}$ を用いて、線形形式
\begin{equation}
 \zeta(u,a,b)=\sum u_{\mu\nu}r_\mu s_\nu
\tag{9}
\end{equation}
を作る。ここで和は所定の添字 $\mu$ および $\nu$ にわたる。$\zeta(u)$ と、行 $a,b$ の置換から生じる、$\zeta(u)$ とは異なり互いにも異なるすべての共役との積、すなわち $\zeta(u)$ を $r_\mu(a,b)$ および $s_\nu(a,b)$ の体の関数とみなすときの、$t(a,b)$ の体に関する $\zeta(u)$ のノルム $N(\zeta(u))$ は、すべての行 $a,b$ の整対称関数であり、各行に関して斉次で同じ次元をもつ。したがって上で述べたことにより、それは $t(a,b)$ および $u_{\mu\nu}$ の、整数係数をもつ一意に定まる整関数である。
\begin{equation}
 N(\zeta(u,a,b))=T(t(a,b),u)\qquad[\text{$a,b,u$ について恒等的}].
\tag{10}
\end{equation}
同じ考察から、
\[
N(z-\zeta(u))=z^\delta+T_{\delta-1}(t,u)z^{\delta-1}+\ldots+T(t,u)
\qquad[\text{$a,b,u,z$ について恒等的}]
\]
も得られる。ここで、すべての $T$ は $t$ と $u$ の整数係数をもつ整関数である。この恒等式の両辺は $z=\zeta(u)$ に対して消滅する。さらに、基本対称関数 $r(a,b)$ と $s(a,b)$ の代数的独立性により、(8) に従って $t(a,b)$ を $r$ と $s$ の整双線形結合 $w(r,s)$ に等しくおき、$\zeta(u)$ を $r$ と $s$ の関数とみなすと、両辺は $r,s$ について恒等的に消滅する。したがって
\begin{equation}
\zeta(u)^\delta+T_{\delta-1}(w(r,s),u)\zeta(u)^{\delta-1}+\cdots+T(w(r,s),u)=0
\quad[\text{$r,s,u$ について恒等的}].
\tag{11}
\end{equation}\footnote{(9) の和をすべての添字に拡張すると、(11) は Gauss の定理の Kronecker による拡張を表し、実質的には Kronecker の証明配置に従ったものである（Zur Theorie der Formen höherer Stufen, Berliner Ber. 1883, Werke, 2, p.~417）。}

恒等式 (10) および (11) は、$a,b$ のすべての特殊値系 $\alpha,\beta$、および $r,s$ のすべての特殊値系 $\varrho,\sigma$ に対して保たれる。特に $\varrho,\sigma$ が、(9) に現れる添字 $\mu,\nu$ に関するすべての積 $\varrho_\mu\sigma_\nu$ を消滅させるように選ばれ、したがって置換によって $\zeta(u)$ が消滅するとき、$w(\varrho,\sigma)=\tau$ とおけば、(11) から
\begin{equation}
 T(\tau,u)=0\qquad[\text{$u$ について恒等的}]
\tag{12}
\end{equation}
が得られる。

逆に、(12) を満たす、すべては零でない値系 $\tau_0,\ldots,\tau_{p+q}$ が与えられたとする。代数的拡大体においては分解
\begin{equation}
 \bar\chi(\xi,\eta)=\tau_{p+q}\xi^{p+q}+\cdots+\tau_0\eta^{p+q}
=(\alpha_1\xi+\beta_1\eta)\cdots(\alpha_{p+q}\xi+\beta_{p+q}\eta)
\tag{12'}
\end{equation}
が存在し、ここでどの因子においても $\alpha$ と $\beta$ は同時には消滅しない\footnote{この有理的な「代数学の基本定理」が、2 の末尾で述べた逆の例外なき妥当性を支える存在定理である。}が、いくつかの $\alpha$ または $\beta$ は消滅してもよいので、$\tau=t(\alpha,\beta)$ である。これらの値を (10) に代入し、(12) を用いると、$\zeta(u,\alpha,\beta)$ またはその共役因子の一つが $u$ について恒等的に消滅することが分かる。したがって、$\alpha,\beta$ を適当に番号付けした後で、$r(\alpha,\beta)=\varrho$、$s(\alpha,\beta)=\sigma$ とおくことは、まさに $\bar\chi(\xi,\eta)$ が少なくとも一つの分解
\begin{equation}
\bar\chi(\xi,\eta)=\bar\varphi(\xi,\eta)\bar\psi(\xi,\eta)
=\sum \varrho_\mu\xi^\mu\eta^{p-\mu}\cdot\sum \sigma_\nu\xi^\nu\eta^{q-\nu}
\tag{13}
\end{equation}
を許し、その際 $\zeta(u)$ に現れるすべての積 $\varrho_\mu\sigma_\nu$ が消滅し、しかし $\varrho$ の全部または $\sigma$ の全部が消滅するわけではない、ということを意味する。

5. (8) における次数 $p,q$ を、いま $ld^{n-1}$ および $md^{n-1}$、すなわち (5) における $f(\xi)$ と $g(\xi)$ の次数に等しいものとする。添字 $\mu,\nu$ を $\mu^{(1)},\nu^{(1)},\mu^{(2)},\nu^{(2)}$ に分ける。ここで $\mu^{(1)}$ は $f(\xi)$ に欠けているすべての指数を走り、$\nu^{(1)}$ は $\psi(\xi,\eta)$ における $\xi$ のすべての指数を走る。対応して、$\nu^{(2)}$ は $g(\xi)$ に欠けているすべての指数を走り、$\mu^{(2)}$ は $\varphi(\xi,\eta)$ における $\xi$ のすべての指数を走る。さらに
\[
 e(\xi)=\sum Z_{k_1\ldots k_n}\xi^k\qquad(k_1+\ldots+k_n\leqq l+m)
\]
を、(3) の多項式に対応する一変数多項式とし、
\[
 S(Z,u)
\]
を、$Z$ と $u$ の整数係数多項式とする。これは、(10) の一意に定まる右辺 $T(t,u)$（ただし (9) の和は指定された添字 $\mu^{(1)},\nu^{(1)},\mu^{(2)},\nu^{(2)}$ にわたる）から、$t$ を係数 $Z$ および $e(\xi)$ の対応する零点で置き換えて得られる。

いま $r,s$ を $f(\xi)$ と $g(\xi)$ の係数 $A,B$ および対応する零点で置き換えると、指定された和について $\zeta(u)$ はこの置換により消滅する。さらに $t=w(r,s)$ は $h(\xi)$ の係数 $C(A,B)$ および対応する零点へ移る。したがって $T(t,u)$ は $S(C(A,B),u)$ へ移る。恒等式 (11) から
\begin{equation}
 S(C(A,B),u)=0\qquad[\text{$A,B,u$ について恒等的}]
\tag{14}
\end{equation}
が得られる。(14) は任意の特殊値系に対して保たれるので、拡大体で二つの因子に分解する特殊な $H(x)$、またはより一般に $H(x,y)$ の係数 $\mathsf{\Gamma}=C(\mathsf{A},\mathsf{B})$ は、条件
\begin{equation}
 S(\Gamma,u)=0\qquad[\text{$u$ について恒等的}]
\tag{15}
\end{equation}
を満たす。

逆に、多項式 $\bar H(x)$、それぞれ形式 $\bar H(x,y)$ の、すべては零でない係数 $\mathsf{\Gamma}$ について条件 (15) が満たされると仮定する。上の議論により、これは $T(\tau,u)=0$ と同値である。ここで $\tau$ は $\bar h(\xi,\eta)$ のすべての係数を表す。したがって (13) により、$\mathsf{\Gamma}$ の拡大体において分解
\[
 \bar h(\xi,\eta)=\bar f(\xi,\eta)\cdot\bar g(\xi,\eta)
\]
が存在し、指定された和のために $\bar f$ と $\bar g$ には誘導された指数以外の指数が現れない。ゆえに
\[
 \bar H(x,y)=\bar F(x,y)\cdot\bar G(x,y)
\]
という関係が存在する。$y=1$ としてもどの因子も零次にならないなら、特に $\mathsf{\Gamma}$ が $H(x)$ の次数低下を生じさせないなら、さらに
\[
 \bar H(x)=\bar F(x)\cdot\bar G(x)
\]
が従う。

\emph{したがって条件 (15) は、$\bar H(x,y)$ が、また次数低下を除外したときには $\bar H(x)$ も、拡大体で二つの因子に分解するための必要十分条件である。}

さらに、$S(Z,u)$ は $Z$ と $u$ について恒等的に消滅しえない。なぜなら 1 で、$n\geqq2$ のとき係数を不定元とする多項式、したがって一つ多い変数に関する対応する斉次形式も、絶対既約であることが示されたからである。したがって、$U$ が $u$ の冪積を表すものとして、
\[
 S(Z,u)=\sum \varPhi_\lambda(Z)U_\lambda
\]
と書ける。ここで $\varPhi_\lambda(Z)$ は (14) により $\mathfrak P$ に属する多項式である。これにより、$\mathfrak P$ はパラメータ表示 $\mathsf{\Gamma}=C(\mathsf{A},\mathsf{B})$ によって与えられるもの以外の零点をもたないことが示された。ただし $\mathsf{A}$ と $\mathsf{B}$ は $\mathsf{\Gamma}$ の代数的拡大体に属する。よって、2 の末尾で述べた逆が証明された。

6. \emph{絶対既約性の条件}は、いま直接に定式化できる。(3) の $E(x)$ の次数を、可能なすべての方法で二つの正の和 $l+m$ に分解し、各分解について形式 $S(Z,u)$ を作ればよい。これらの形式の積
\begin{equation}
 R(Z,u)=\prod S(Z,u)
\tag{16}
\end{equation}
が $E(x)$ の可約性形式をなす。実際、上で見たように、$E(x)$、それぞれ $E(x,y)$ の任意の分解には (16) の一つの因子の消滅が対応し、また逆も成り立つ。これにより、序文で述べた主定理が証明され、同時に可約性形式を有限個の手順で実際に構成できることも示された。

7. 可約性形式 (16) の存在から、序文で述べた Ostrowski の定理が直ちに従う。
\[
 \bar E(x)=\sum \mathsf{\Gamma}_{k_1\ldots k_n}x_1^{k_1}\ldots x_n^{k_n}
 \qquad(k_1+\ldots+k_n\leqq v)
\]
を、代数的整数 $\mathsf{\Gamma}$ を係数にもつ絶対既約多項式とする。また $\Omega$ をすべての $\mathsf{\Gamma}$ を含む有限代数的数体とし、$\mathfrak p$ を $\Omega$ の素イデアルとする。

すると、$\bar E(x)$ の正確な次数に対して作られた可約性形式 $R(Z,u)$ は、置換 $Z=\mathsf{\Gamma}$ の下でも零と異なる。すなわち
\[
 R(\mathsf{\Gamma},u)=\sum \mathsf{P}_\lambda U_\lambda\ne0\qquad[\text{$u$ について恒等的}u]
\]
である。したがってイデアル $\mathfrak r=(\mathsf{P}_1,\mathsf{P}_2,\ldots)$ は、$\Omega$ の有限個の素イデアルによってしか割られない。$\mathfrak p$ がこれら有限個のものと異なれば、$R(\mathsf{\Gamma},u)$ は $\mathfrak p$ を法として恒等的には消滅しない。そして $\mathfrak p$ を法とする剰余類は再び体をなすので、$\bar E(x)$ は $\mathfrak p$ を法として既約、より正確には絶対既約であることが従う。同じことは $\bar E(x,y)$ についても成り立つので、$\mathfrak p$ を法として次数低下も生じない。

$\bar E(x)$ の係数が代数的分数である場合には、共通分母 $\Lambda$ を割る有限個の素イデアルからも $\mathfrak p$ を除外しなければならない。残りのすべての素イデアルを法として、$E(x)$ は $\Lambda\bar E(x)$ に置き換えられ、上の仮定が満たされる。これにより定理が証明された。

\begin{flushleft}
Erlangen, 1921 年 8 月。
\end{flushleft}
\begin{center}
（1921 年 8 月 28 日受理。）
\end{center}

\clearpage
%
\section*{21. 形式的変分法と微分不変式}
\begin{center}
Encyklopädie d. math. Wiss. III, 3 (1922), S. 68--71
（R. Weitzenböck, Differentialinvarianten 所収）
\end{center}

\begingroup
\renewcommand{\thefootnote}{\arabic{footnote})}
\setcounter{footnote}{148}
\subsection*{28. 形式的変分法と微分不変式。\footnote{この項目は E. Noether による。}}

すべての微分不変式の生成と主要定理の導出は、変分法の形式的方法によって最も統一的に行うことができる。この原理は Riemann\textsuperscript{78)} にさかのぼり、形式論の面では主として Lipschitz\textsuperscript{73)} によって展開された。

%
%
斉次形式 $f(dx)$ に属する変分問題から出発する。ただし
\[
\left|\frac{\partial^2 f}{\partial dx_i\,\partial dx_k}\right|\ne0
\]
とする。この変分問題は
\begin{equation}
\delta\int_{t_0}^{t_1} f(x')\,dt=0,
\qquad \left(x_i'=\frac{dx_i}{dt}\right)
\tag{140}
\end{equation}
である。部分積分により、不変な方程式系
\[
2\psi_i=\frac{d}{dt}\frac{\partial f}{\partial x_i'}-\frac{\partial f}{\partial x_i}=0,
\]
へ導かれる。ここでラグランジ\nolinebreak[4]ュ式 $\psi_i$（ここでの規格化では、対応する変分問題のラグランジ\nolinebreak[4]ュ方程式の左辺の二分の一）は、$dx$ に対して反傾の変換則に従う。これは、$\psi_i$ を、部分積分に対応する積分を含まない形式的恒等式（歴史的名称ではラグランジ\nolinebreak[4]ュの中心方程式\footnote{特殊な $f$ について、この名称は Heun による。彼の論文 IV, 1 II, p. 447 を参照。}）
\[
\delta f-df_\delta=-2\sum\psi_i(d,d)\,\delta x_i,
\qquad
\left(f_\delta=\sum\frac{\partial f}{\partial dx_i}\,\delta x_i\right),
\]
によって定義すれば、形式的に最も簡単に分かる。実際、$f$ と $df_\delta$ はともに不変式であり、したがって右辺も不変式である。特に二次形式の場合には
\[
\delta\sum g_{ik}dx_i dx_k-2d\sum g_{ik}dx_i\delta x_k
=-2\sum\psi_\mu(d,d)\,\delta x_\mu.
\]
を得る。

%
ここで $d$ を、それと同じ変換則に従う $(d+\lambda\delta)$ で置き換え、$\lambda$ の冪を比較すると、対応する恒等式系
\begin{equation}
\left\{
\begin{aligned}
D\sum g_{ik}dx_i dx_k-2d\sum g_{ik}dx_iDx_k
   &=-2\sum\psi_\mu(d,d)\,Dx_\mu,\\
D\sum g_{ik}dx_i\delta x_k-\delta\sum g_{ik}dx_iDx_k
   &\phantom{={}}{}-d\sum g_{ik}\delta x_iDx_k\\
   &=-2\sum\psi_\mu(d,\delta)\,Dx_\mu,\\
D\sum g_{ik}\delta x_i\delta x_k-2\delta\sum g_{ik}\delta x_iDx_k
   &=-2\sum\psi_\mu(\delta,\delta)\,Dx_\mu,
\end{aligned}
\right.
\tag{141}
\end{equation}
が得られる。これから $\psi_\mu(d,\delta)$、したがって特に $\psi_\mu(d,d)$ と $\psi_\mu(\delta,\delta)$ も、反傾ベクトルであることが分かる。(141) から
\[
\psi_\mu(d,\delta)=\sum g_{i\mu}\,d\delta x_i+
\sum\left[\begin{smallmatrix}ik\\ \mu\end{smallmatrix}\right]dx_i\delta x_k,
\]
と計算され、これより $dx$ と同じ変換則に従うベクトル
\begin{equation}
p^\sigma(d,\delta)=\sum g^{\sigma\mu}\psi_\mu
=d\delta x_\sigma+
\sum\left\{\begin{smallmatrix}ik\\ \sigma\end{smallmatrix}\right\}dx_i\delta x_k.
\tag{142}
\end{equation}
が得られる。$p^\sigma(d,d)=0$ とおけば、明らかに測地線の方程式を表す。

$dx$ と同じ変換則に従うベクトル $p^\sigma(d,\delta)$ が分かれば、共変微分（第19項参照）へ直接進むことができる。たとえば $h(d,\delta)$ が二つの系列 $dx$ と $\delta x$ の形式であるとき、式
\begin{equation}
\begin{aligned}
h^{(1)}(dx,\delta x)={}&\delta h-
\sum\frac{\partial h}{\partial dx_\sigma}p^\sigma(d,\delta)\\
&-\sum\frac{\partial h}{\partial \delta x_\sigma}p^\sigma(\delta,\delta)
\end{aligned}
\tag{143}
\end{equation}
は不変式の差としてそれ自体が不変式であり、しかも二階微分が相殺されるように作られている。したがって $h^{(1)}(dx,\delta x)$ は一階微分だけを含む形式、すなわち $h$ の共変微分を表す。$h$ がさらに多くの系列 $d^{(1)}x,\ldots,d^{(\lambda)}x$ を含む場合にも同様である。\footnote{Lipschitz, J. f. Math. 72 (1870), p. 17.}

%
Riemann と Lipschitz による曲率形式の構成も同じ原理に基づく。$\delta^2f$ から「第二変分の正規形」
\begin{equation}
\begin{aligned}
\Omega={}&\delta^2\sum g_{ik}dx_i dx_k
-2d\delta\sum g_{ik}dx_i\delta x_k\\
&+d^2\sum g_{ik}\delta x_i\delta x_k,
\end{aligned}
\tag{144}
\end{equation}
へ移る。この式は不変式の総和としてそれ自体ふたたび不変式であり、ここではラグランジ\nolinebreak[4]ュの中心方程式を一般化することによって三階微分が消去されている。二階微分は、共変微分との類推により、差
\begin{equation}
K=\Omega-2\sum\{p^\sigma(d,\delta)\psi_\sigma(d,\delta)
-p^\sigma(\delta,\delta)\psi_\sigma(d,d)\},
\tag{145}
\end{equation}
を作ることで消去される。これはまた、次のようにも書ける。\footnote{Lipschitz, J. f. Math. 82 (1876), § 1.}
\begin{equation}
\begin{aligned}
\frac12 K={}&d\sum\psi_\sigma(d,\delta)\delta x_\sigma
-\delta\sum\psi_\sigma(d,d)\delta x_\sigma\\
&-\sum\{p^\sigma(d,\delta)\psi_\sigma(d,\delta)
-p^\sigma(\delta,\delta)\psi_\sigma(d,d)\}.
\end{aligned}
\tag{146}
\end{equation}
$K$ の構成法則は、$\Omega$ において $p(d,\delta)$ と $p(\delta,\delta)$、すなわち (141) の右辺をそれぞれ零とおくことにより、二階微分を消去する、と述べることもできる。これは Riemann による曲率形式の定義である（第21項参照）。ただし、(146) が特に明瞭に示すとおり、この零置換は微分を行った後でなければならないことを明記しておく必要がある。Lipschitz の形式的構成法はこの困難を避ける。同様に、共変微分 $h^{(1)}$ (143) も、$\delta h$ の中で $p(d,\delta)\ldots$ を零とおくことによって二階微分を消去する、と定義できる。

%
曲率形式から共変微分を順次作ることにより、形式の系列 $[\Phi_4,\Phi_5,\ldots\text{（Christoffel）}]$ が得られる。Christoffel と Ricci によれば、$f$ を付け加えた後のこの系列の射影不変式は、二次形式の微分不変式のすべてを尽くす。また、線形変換に関するこれらの同値性は、追加の可積分条件なしに二つの二次微分形式の同値性を意味する（第22項参照）。この還元定理の内的な理由は、変分問題 (140) から生じる Riemann 正規座標（第20項参照）の存在にある。この座標は一点を通る測地線を直線へ変換し、したがって $x$ の任意の変換の下で線形に変換される。ゆえに、すべての不変式の構成と同値性の問題は、正規座標に関する $f$ の展開の各斉次成分についての対応する問題へ帰着する。そして、それらの成分が $f,\Phi_4,\Phi_5,\ldots$ から線形に構成されることが示される。二次形式については、Vermeil\textsuperscript{92)} が E. Noether\textsuperscript{93)} のノートに沿\nolinebreak[4]って、これを詳しく実行した。このノートによれば、任意の微分式を基礎に置く場合にも定理と方法は保たれる（第22項および第23項参照）。これにより、純粋な消去理論の方法に比して、形式的変分法の方法がどれほど広い射程をもつかが分かる。\footnote{以上の説明と、さらに詳しい幾何学的解釈については、「文献」に挙げた F. Klein のゼミナール講演を参照。}

Levi-Civita\textsuperscript{134)} は、$p(d,\delta)$ を零とおくことを「ベクトルの平行移動」として幾何学的に解釈した（Hessenberg\textsuperscript{133)} も参照）。ここでいう平行移動は接続によって定まる平行輸送であり、ユークリッド空間の単なる並進ではない。Weyl が公理的に捉えたこの概念は、本質的には二次形式に限られるが、別種の一般化を許す。群を拡張して $g_{ik}$ に任意関数を掛ける場合にも、二次形式に加えて一次形式を基礎に置けば、この概念は保たれる。そのとき $p(d,\delta)$ は一意に定まり、上と同様に不変式を構成し、測地座標を定義できる。しかし $p(d,d)=0$ は、いまや変分問題からは生じない！\footnote{H. Weyl, Reine Infinitesimalgeometrie, Math. Ztschr. 2 (1918), p. 384 bis 411; Raum, Zeit, Materie (3. u. 4. Aufl.).} 還元定理と同値性の問題は、ここではまだ一般には取り組まれていない。最低次の不変式についてだけ、R. Weitzenböck が還元定理を与えている。\footnote{Wien. Ber. 129 (1920), p. 683 u. 697.}

%
最後に、一般の「不変変分問題」、すなわち（単一または多重の）積分 $J$ がリーの意味での何らかの群に関して不変である問題に注意を促しておく。物理的に興味深い特殊な場合は、Hamel\footnote{Math. Ann. 59 (1904), p. 416--434; Ztschr. Math. Phys. 50 (1904), p. 1--54.}、Herglotz\footnote{Ann. d. Phys. (4) 36 (1911), p. 511, bes. § 9.}、Lorentz\footnote{Verslag der Amsterd. Akad., April, Sept. 1916, Okt. 1916, Mai 1917.}、Fokker\footnote{Ebenda, Jan. 1917.}、Weyl\footnote{Ann. d. Phys. (4) 54 (1917), p. 117, § 2.}、Klein\footnote{Gött. Nachr. 1917, p. 469--482; ebenda 1918, p. 171--189.} の研究で扱われている。E. Noether\footnote{Gött. Nachr. 1918, p. 235--257.} による原理的な定式化は、$J$ が $G_\varrho$（$\varrho$ 個の本質的パラメータをもつ有限連続群）に関して不変であるとき、これに $\varrho$ 個の線形独立な全発散式が対応することを示す。ここで全発散式とは、変分において境界項を与える発散式をいう。また、$\sigma$ 階導関数までを含む $\varrho$ 個の任意関数に依存する無限連続群に関する不変性には、ラグランジ\nolinebreak[4]ュ式とその $\sigma$ 階までの導関数との間の $\varrho$ 個の恒等的従属関係が対応する。いずれの場合にも逆が成り立つ。ラグランジ\nolinebreak[4]ュ式は群の（相対）不変式となるので、同時に不変式を生成する過程も得られる。

\begin{center}
（1921年3月完了。）
\end{center}
\endgroup

\clearpage
\setcounter{footnote}{0}

%
%
%
\begingroup
\renewcommand{\thefootnote}{\arabic{footnote})}
\setcounter{footnote}{0}
\section*{22. 多項式イデアルと終結式の理論について}
\begin{center}
Math. Ann. 88 (1923), S. 53--79\\[0.9em]
{\Large\bfseries 多項式イデアルと終結式の理論について。}\\[1.0em]
著者\\[0.45em]
Kurt Hentzelt $\dagger$。\\[0.8em]
Göttingen の Emmy Noether による改作\footnotemark。\\[0.6em]
\rule{4em}{0.4pt}
\end{center}

以下で扱うのは、\emph{消去理論の一般問題、すなわち多項式イデアルに属するすべての多項式の共通零点を決定する問題}である。これは、そのイデアルの一つの基底をなす有限個の多項式の零点を決定することと同じである。その際、現れる重複度の概念的な解釈も同時に得られる。各多項式の係数は任意の体に属してよく、零点はそこから得られる代数閉体に属する。さらに、一般性を失うことなく、イデアルは変換されたもの（§~3）と仮定できる。このとき本質的な結果は次のとおりである。

\emph{任意の多項式イデアル $\mideal$ には、次の終結式形式を対応させることができる：}
\[
 R_{\mideal}=R^{(1)}(x_1,\ldots,x_n)\cdots
 R^{(i)}(x_i,\ldots,x_n)\cdots R^{(n)}(x_n)\zero{\mideal}
\]
\footnotetext{Kurt Hentzelt は 1914 年 10 月以来 Dixmuiden 近郊の前線で行方不明であり、戦没者に数えなければならない。ここに掲げる論文は、彼が 1914 年夏に Erlangen で E. Fischer のもとで博士号を取得した博士論文 ``Zur Theorie der Formenmoduln und Resultanten'' の最も本質的な部分を、まったく自由に改作したものである。彼自身の着想だけに基づいて書かれたこの学位論文は、論理上の欠落なく構成されている。しかし補助定理が次々と連なり、すべての概念が四つ五つの添字をもつ公式によって記述され、文章による説明がほぼ全面的に欠けているため、理解には最大の困難が伴う。彼自身は予定していた改作を行うには至らなかった。私はこの仕事を純粋に概念的な形で提示する。\mbox{それによって}、基本思想が一貫して Hentzelt に由来する証明は大幅に簡明になり、また願わくは、この仕事の美しさが明らかになる。――学位論文のうち、\mbox{基底}が与えられたときに現れる関数を有限個の手順で構成する問題に関わる部分は、別稿での公刊に留保する。（E. N.）}
\emph{この $R_{\mideal}$ は $\mideal$ のすべての零点で消え、かつそれらでのみ消える。$\nideal$ が $\mideal$ を割り、終結式形式 $R_{\nideal}$ と $R_{\mideal}$ が一致するならば、イデアル $\nideal$ と $\mideal$ も一致する。}

主定理の第二部は、終結式形式が零点を、それぞれを特徴づける重複度とともに与えることを示す。この第二部は、代数的数体の二つのイデアルについて、一方が他方を割るとき、その二つが一致することと、それらの\emph{ノルム}が一致することとが同値である、という事実に類似する。両定理の内的根拠も同じである。

実際、各多項式を $x_1,\ldots,x_{i-1}$ の冪積に関する線型形式とみなし、$x_{i+1},\ldots,x_n$ を係数領域に添加すれば、$\mideal$ を線型形式の加群 $\Mmod_{i-1}$ とみなせる。そのとき終結式因子 $R^{(i)}$ は、対応する基本加群（§~1）$\Gmod_{i-1}$ の $\Mmod_{i-1}$ に関する\emph{ノルム}に等しくなる。これにより、\emph{$R^{(i)}$ の一因子の次数と指数との積として現れる\mbox{重複度}を、線型独立な剰余類の個数によって解釈すること}も直ちに得られる\footnote{これら後者の結果は、イデアルの準素イデアルへの分解を考慮するときに、その本来の意味を示す。準素イデアルに対応する重複度は、そのイデアルを法とする補元の線型独立な剰余類の個数として定義される。この点は別の場所で扱う予定である。（E. N.）}。この事実は、これまでは\mbox{不定係数}に対して終結式を定義できる特殊な場合にしか知られていなかった\footnote{Macaulay, \emph{The algebraic theory of modular systems}, Cambridge Tracts 19 (Cambridge University Press, 1916), Nr. 67 を参照。また Lasker, Zur Theorie der Moduln und Ideale, Math. Ann. 60 (1905), S. 20--116, S. 98 も参照。この場合に終結式と終結式形式が一致することは、注 1) で述べた Hentzelt の未公刊の諸定理で示されている。（E. N.）}。

これらの結果を導くため、§~1 および §~2 では、係数が有理整数または一つの不定元の多項式である\mbox{線型形式}の加群についての定理を与える。§~3 で導入する多項式イデアルとの結びつきは、§~4 の\mbox{同型定理}によって与えられる。§~5 は上に述べた主定理の第二部を、§~7 は零点に関する定理を与え、その前に §~6 で一般消去理論を展開する。そこでは同時に、新しい不定元を導入したとき、\mbox{リゾルベント}がそれらの不定元の線型形式へ分解されることが証明される。これにより、ここで行う消去について、Kronecker のある主張に対する厳密な証明が得られる\footnote{König, \emph{Algebraische Größen} (Leipzig 1903, Teubner), V, §~4 における Kronecker の消去理論のための証明の試みが成功していないことは、よく知られている。König が Kronecker の消去理論に導入した重複度についても Hentzelt の主定理の第二部が成り立たないことは、Macaulay, a. a. O. S. 23 に示されており、さらにそこでは Kronecker の消去理論が準素イデアルへの分解に対応しないことも示されている。（E. N.）}。

\endgroup
%
%
\begingroup
\renewcommand{\thefootnote}{\arabic{footnote})}
\setcounter{footnote}{4}
\section*{§ 1. 線型形式から成る加群。}

最初の二節において、整な量 $a,b,c,\ldots$ とは、有理整数、または任意の抽象的に定義された体 $\mathsf{P}$ を係数とする一つの不定元の多項式をいう。また線型形式 $a(\xi),b(\xi),\ldots$ とは、有限個の不定元 $\xi_1\ldots\xi_k$ に関する、整な量を係数とする線型形式をいう。

\emph{線型形式から成る加群 $\Amod$} は次の条件によって定義される。$\Amod$ は $a(\xi)$ と $b(\xi)$ とともに差 $a(\xi)-b(\xi)$ を含み、また $a(\xi)$ とともに $c\cdot a(\xi)$ を含む。ただし $c$ は任意の整な量である。

\emph{$\Amod$ の階数}とは、通常どおり、$\Amod$ に属する線型独立な線型形式の最大個数をいう。$\Amod$ の第 $\lambda$ \emph{行列式因子} $d_\lambda$ とは、$\Amod$ から得られるすべての $\lambda$ 次行列式、すなわち、$\Amod$ に属する線型形式を任意に $\lambda$ 個取ったとき、その係数行列から作れるすべての $\lambda\times\lambda$ 小行列式の最大公約元をいう。$\varrho$ を $\Amod$ の階数とし、したがって $d_1,\ldots,d_\varrho$ を零でない行列式因子とすると、$\Amod$ の\emph{単因子}は $e_1=d_1$, $e_2=d_2/d_1,\ldots,e_\varrho=d_\varrho/d_{\varrho-1}$, $e_{\varrho+i}=0$ によって定義される。周知のとおり、$\Amod$ は有限生成加群であり、ちょうど $\varrho$ 個の線型形式 $a_1(\xi),\ldots,a_\varrho(\xi)$ から成る加群基底をもつ。$\Amod$ に属する任意の線型形式は、この基底の、整な量を係数とする線型結合として表される：$\Amod=(a_1(\xi),\ldots,a_\varrho(\xi))$。これらの線型形式の係数行列を $A$ とすれば、$\Amod$ の行列式因子および単因子は、それぞれ $A$ の行列式因子および単因子と一致する。したがってまず、\emph{各単因子 $e_\lambda$ は、その次の単因子 $e_{\lambda+1}$ を割る}ことが従う。単因子論による行列等式
\[
 P A Q=\begin{pmatrix}
 e_1&&&0\\
 &e_2&&\\
 &&\ddots&\\
 0&&&e_\varrho
 \end{pmatrix}
\]
はさらに、単模変換 $\xi=Q(\eta)$ によって新しい不定元 $\eta_1\ldots\eta_k$ を導入すれば、次の加群等式が成り立つことを示す：
\begin{equation}
 \Amod=(e_1\eta_1,e_2\eta_2,\ldots,e_\varrho\eta_\varrho).
\tag{1}
\end{equation}

以下で最も本質的な概念は、次に定義する基本加群である\footnote{Steinitz, Rechteckige Systeme und Moduln in algebraischen Zahlkörpern II, Math. Ann. 72 (1912), S. 297--345, Nr. 36 を参照。Hentzelt は、ほかにも Steinitz との接点は見られるものの、少なくとも Steinitz とは独立に、彼のより単純な場合について、式 (4) の形でこの概念に到達したようである。（E. N.）}。

\endgroup
%
\begingroup
\renewcommand{\thefootnote}{\arabic{footnote})}
\setcounter{footnote}{5}

\medskip
\noindent\textbf{定義 I。} \emph{加群 $\Gmod$ を基本加群というのは、$\Gmod$ を割る同じ階数の真の約加群が存在しないときである}\footnote{ここで「割る」は加群の意味で用いる。$\Amod$ が $\Bmod$ で割り切れる、すなわち $\Amod\zero{\Bmod}$ とは、$\Amod$ の各元が $\Bmod$ に含まれることをいう。$\Bmod$ が $\Amod$ に属さない元を含むとき、$\Bmod$ を真の約加群という。}。

したがって、$\Gmod$ が基本加群であることは、
\begin{equation}
 c\cdot g(\xi)\zero{\Gmod};\quad c\ne0 \quad\text{ならば常に}\quad g(\xi)\zero{\Gmod}
\tag{2}
\end{equation}
が成り立つことと同値である。

任意の加群と基本加群との関係は、次の定理によって与えられる。

\medskip
\noindent\textbf{定理 I。} \emph{任意の加群 $\Amod$ は、それと同じ階数をもつただ一つの基本加群 $\Gmod$ で割り切れる。$\Gmod$ は、$\Amod$ を割る同じ階数の加群のうち最小のものとして定義され、}
\begin{equation}
 \Gmod=(\eta_1,\ldots,\eta_\varrho)
\tag{3}
\end{equation}
\emph{によって表される。この $\Gmod$ を $\Amod$ の基本加群と呼ぶ。}

$\Gmod$ が、$\Amod$ を割る同じ階数の加群のうち最小であるという条件は、次のように表される。すなわち $\Gmod$ は、ある関係
\begin{equation}
 b\cdot g(\xi)\zero{\Amod};\quad b\ne0
\tag{4}
\end{equation}
を満たす線型形式 $g(\xi)$ を、すべて、しかもそれだけ含む。これからまず $\Gmod$ が加群であることが従う。実際、
\[
 b_1g_1(\xi)\zero{\Amod};\quad b_1\ne0;
 \qquad
 b_2g_2(\xi)\zero{\Amod};\quad b_2\ne0
\]
とともに、
\[
 b_1b_2\bigl(g_1(\xi)-g_2(\xi)\bigr)\zero{\Amod};\quad b_1b_2\ne0
\]
も成り立ち、もちろん $b_1\cdot c g_1(\xi)\zero{\Amod}$ も成り立つ。さらに $\Gmod$ は基本加群でもある。実際、
\[
 c\cdot g(\xi)\zero{\Gmod};\quad c\ne0
\]
から (4) により
\[
 b c g(\xi)\zero{\Amod};\quad bc\ne0
\]
が従い、したがって再び (4) により $g(\xi)\zero{\Gmod}$ が従う。

さらに、以上の条件を満たす基本加群で、同じ階数の最小の約加群 $\Gmod$ と異なるもの $\Gmod_1$ は存在しない。実際、$\Gmod_1$ は $\Gmod$ で割り切れなければならないが、基本加群であるから同じ階数の真の約加群をもたず、したがって $\Gmod_1$ は $\Gmod$ と同一である。最後に、(3) の右辺で表される加群は基本加群である。なぜなら、その任意の真の約加群は、さらに一つの不定元 $\eta$ を含まなければならず、したがって、より高い階数をもつからである。ゆえに (3) は $\Amod$ の一意に定まる基本加群を与える。これで\emph{定理 I のすべての主張が証明された}。

\endgroup
%
\begingroup
\renewcommand{\thefootnote}{\arabic{footnote})}
\setcounter{footnote}{6}

以下でもう一つ本質的となる概念は、Dedekind による二つの加群の商の概念であり\footnote{Dedekind が扱うのは数の加群であり、そこでは乗法も定義されている。したがって Dedekind の商は、ここで定義する商とは一致しない。（E. N.）}、次のように\mbox{定義される}。

\medskip
\noindent\textbf{定義 II。} \emph{線型形式から成る二つの加群の商 $\cc=\Amod/\Bmod$ は、条件}
\[
 c\cdot\Bmod\zero{\Amod}
\]
\emph{を満たす整な量 $c$ 全体として\mbox{定義される}。$\cc$ は整な量から成るイデアル、したがって主イデアルをなす。これに対応して、線型形式から成る加群 $\Amod$ を整な量から成るイデアル $\bb$ で割った商 $\mathfrak C=\Amod/\bb$ は、条件}
\[
 c(\xi)\cdot\bb\zero{\Amod}
\]
\emph{を満たす線型形式 $c(\xi)$ 全体として\mbox{定義される}。$\mathfrak C$ は再び線型形式から成る加群をなし、$\bb$ が零イデアルでない限り、$\Amod$ と同じ階数をもつ。}

ここで注意すべきことは、定義から明らかなとおり、いずれの商も加群の性質をもつということだけである。そして、整な量から成り加群の性質をもつ系はイデアルである。

商の概念から、基本加群と最高単因子との間に次の関係が得られる。

\medskip
\noindent\textbf{定理 II。} \emph{$\Amod$ の基本加群と最高単因子との間には、相反関係}
\begin{equation}
 (e_\varrho)=\Amod/\Gmod;\qquad \Gmod=\Amod/(e_\varrho)
\tag{5}
\end{equation}
\emph{が成り立つ。ここで $(e_\varrho)$ は $e_\varrho$ が生成する主イデアルを表す。}

実際、各単因子は $e_\varrho$ を割るので、(1) と (3) から
\begin{equation}
 e_\varrho\Gmod\zero{\Amod}\quad \text{かつ、したがって}\quad (e_\varrho)\cdot\Gmod\zero{\Amod};
\tag{6}
\end{equation}
を得る。したがって $(e_\varrho)$ は商 $\Amod/\Gmod$ で\mbox{割り切れる}。逆も成り立つ。というのは、
\[
 b\Gmod\zero{\Amod}
 \quad\text{ならば}\quad b\eta_\varrho\zero{\Amod}
 \quad\text{であり、したがって}\quad b\equiv 0(e_\varrho)
\]
だからである。これで (5) の第一式が証明された\footnote{$\Amod_0=\Amod,\ldots,\Amod_\lambda=(\Amod,\eta_\varrho,\ldots,\eta_{\varrho-\lambda+1})$ とおく。ここで各 $\Amod_\lambda$ は最高単因子 $e_{\varrho-\lambda}$ をもつ。このとき、同じ推論によって
\[
 (e_\varrho)=\Amod_0/\Amod_1,\ldots,(e_{\varrho-\lambda})=\Amod_\lambda/\Amod_{\lambda+1},\ldots,(e_1)=\Amod_{\varrho-1}/\Amod_\varrho
\]
を得る。さらに一般に、$\zeta_\lambda=b_{\lambda1}\eta_1+\cdots+b_{\lambda\lambda}\eta_\lambda$ とおき、$(b_{\lambda\lambda},e_\lambda)=1$ とし、
\[
 \Bmod_0=\Amod,\ldots,\Bmod_\lambda=(\Amod,\zeta_\varrho,\ldots,\zeta_{\varrho-\lambda+1})
\]
と書く。このとき $\Bmod_\lambda$ も最高単因子 $e_{\varrho-\lambda}$ をもつので、
\[
 (e_{\varrho-\lambda})=\Bmod_\lambda/\Bmod_{\lambda+1}
\]
を得る。}。さらに (6) は、$\Gmod$ が商 $\Amod/(e_\varrho)$ で\mbox{割り切れる}ことを示す。この商は $\Amod$ を割る同じ階数の加群であるから、$\Amod$ の基本加群の定義により、逆にこの商も $\Gmod$ で\mbox{割り切れる}。これで (5) の第二式も証明された。

\endgroup
%
\begingroup
\renewcommand{\thefootnote}{\arabic{footnote})}
\setcounter{footnote}{8}

$d$ を、$\Amod$ から得られる恒等的に零でない任意の $\varrho$ 次行列式とする。このとき $d$ は第 $\varrho$ 行列式因子 $d_\varrho$ で\mbox{割り切れ}、したがって $e_\varrho$ でも\mbox{割り切れる}。それゆえ $d\Gmod\zero{\Amod}$ となり、上の結論から $\Gmod=\Amod/(d)$ が従う。しかし後の議論にとって本質的なのは、$\Gmod$ のこの最後の商表示が、単因子に立ち戻らずに\mbox{証明でき}、そのために、より一般的に成り立つということである。すなわち、

\medskip
\noindent\textbf{定理 III。} \emph{整な量を多変数多項式と解する場合にも}\footnote{実際に用いるのは、整な量が加法・減法・乗法によって再び整な量となり、通常の計算法則が成り立つこと、すなわち、それらが環をなすことだけである。さらに、この環では積が零となるのは一方の因子が零となる場合に限られるので、商を作ること（元の対を添加すること）によって、この環を体へ拡大できることだけを用いる。これに対して、加群の基底表示は用いない。したがって (7) は、たとえば、すべての\mbox{代数的整数}から成る領域を係数領域とする場合にも成り立つ。（E. N.）}、\emph{商表示}
\begin{equation}
 \Gmod=\Amod/(d)
\tag{7}
\end{equation}
\emph{はなお成り立つ。ここで $d$ は、$\Amod$ から得られる恒等的に零でない任意の $\varrho$ 次行列式である。}

まず、この拡張のもとでも、階数、基本加群、および加群の商という概念は保たれる。ただし、商は\mbox{この場合}には主イデアルとは限らない。また、基底表示を除けば、定理 I も保たれることは明らかである。

いま、
\[
 a_1(\xi)=a_{11}\xi_1+\cdots+a_{1k}\xi_k;\ \ldots;\
 a_\varrho(\xi)=a_{\varrho1}\xi_1+\cdots+a_{\varrho k}\xi_k
\]
を、$\Amod$ に属する $\varrho$ 個の線型独立な線型形式とする。これらは一般には加群基底をなさない。不定元 $\xi$ の番号を、
\[
 d=|a_{ij}|\ne0;\qquad i,j=1,2,\ldots,\varrho
\]
となるように定める。すると、次の特殊な形をもつ $\varrho$ 個の線型形式が $\Amod$ に属することが従う：
\begin{equation}
 d\xi_\lambda+b_{\lambda,\varrho+1}\xi_{\varrho+1}+\cdots+b_{\lambda,k}\xi_k\zero{\Amod}
 \quad(\lambda=1,2,\ldots,\varrho).
\tag{8}
\end{equation}
しかし $\Amod$ は、$\xi_{\varrho+1}\ldots\xi_k$ だけに依存する線型形式 $c_{\varrho+1}\xi_{\varrho+1}+\cdots+c_k\xi_k$ を含むことはできない。\mbox{そうでなければ}、少なくとも一つの $(\varrho+1)$ 次行列式 $d\cdot c_\sigma$ が零でないことになるからである。

さて、商 $\Amod/(d)$ は $\Amod$ を割る同じ階数の加群であるから、$\Gmod$ で\mbox{割り切れる}。しかし逆も成り立つ。実際、任意の $g(\xi)\zero{\Gmod}$ に対して、定義により
\[
 c\cdot g(\xi)\zero{\Amod};\quad c\ne0
 \quad\text{であり、したがって}\quad d\cdot c\cdot g(\xi)\zero{\Amod}
\]
が成り立つ。一方、(8) により
\[
 d\cdot g(\xi)\equiv g_{\varrho+1}\xi_{\varrho+1}+\cdots+g_k\xi_k\,(\Amod)
\]
となり、それゆえ
\[
 c g_{\varrho+1}\xi_{\varrho+1}+\cdots+c g_k\xi_k\zero{\Amod}
\]
を得る。したがって、いま述べたことにより $c\cdot g_{\varrho+1}=0,\ldots,c\cdot g_k=0$ であり、$c\ne0$ であるから $g_{\varrho+1}=0,\ldots,g_k=0$ でもある。ゆえに $d\cdot g(\xi)\zero{\Amod}$ となり、(7) が証明された。

なお、$d\cdot g(\xi)\zero{\Amod}$ は、二つの加群 $(a_1(\xi),\ldots,a_\varrho(\xi))$ と $(g(\xi),a_1(\xi),\ldots,a_\varrho(\xi))$ が同じ階数 $\varrho$ をもつことからも、直接に従う。実際、$g(\xi)=c_1\xi_1+\cdots+c_k\xi_k$ とおけば、$(\varrho+1)$ 次行列式
\[
\begin{vmatrix}
 a_1(\xi)&a_{11}&\cdots&a_{1\varrho}\\
 \vdots&\vdots&&\vdots\\
 a_\varrho(\xi)&a_{\varrho1}&\cdots&a_{\varrho\varrho}\\
 g(\xi)&c_1&\cdots&c_\varrho
\end{vmatrix}
\]
が零となるからである。しかし、上で与えた証明は、§~4 で式 (30) に続いて再び現れる。

\endgroup
%

\section*{§ 2. ノルムおよび\mbox{加群}の分解定理。}

ここでは再び、係数が有理整数、または $\mathsf{P}$ を係数とする一変数多項式である線型形式から成る\mbox{加群}を扱う。

\medskip
\noindent\textbf{定義 III。} \emph{通常どおり、$\Amod$ を法として互いに合同な $\xi$ の線型形式を一つの剰余類にまとめる。このとき\mbox{記号} $\Gmod\mid\Amod$ は、$\Amod$ を法とする $\Gmod$ の剰余類系、\mbox{すなわち}、$\Gmod$ の元から成る $\Amod$ を法とする剰余類すべての系を表す。$\Amod$ に関する $\Gmod$ のノルム---\mbox{記号}では $\Norm(\Gmod\mid\Amod)$---は、$\Gmod$ から $\Amod$ への移行変換の\mbox{行列式}として\mbox{定義される}。\mbox{すなわち}、$\Amod$ の任意の線型独立な\mbox{加群基底}を、$\Amod$ の\mbox{基本加群} $\Gmod$ の同様な基底によって表す変換の\mbox{行列式}である。}

(1) と (3)、または定理 II への注から、次を得る：
\begin{equation}
 \Norm(\Gmod\mid\Amod)=e_1e_2\cdots e_\varrho=d_\varrho;
 \qquad
 \Gmod=\Amod/\bigl(\Norm(\Gmod\mid\Amod)\bigr).
\tag{9}
\end{equation}

(1)、(3)、および (9) から、よく知られた仕方で次のことが従う。有理整数の場合、$\Norm(\Gmod\mid\Amod)$ は $\Amod$ を法とする $\Gmod$ の剰余類の個数を表す。一方、一変数多項式の場合には、剰余類の個数は一般に無限となるが、$\Norm(\Gmod\mid\Amod)$ の次数が、\emph{$\mathsf{P}$ に関して線型独立な剰余類の個数}に等しくなる。$\Bmod$ が $\Amod$ を割る同じ階数の\mbox{加群}\mbox{であれば}、$\Bmod$ は、そのような任意の剰余類の一つの代表元とともに、その剰余類のすべての元を含む。\mbox{したがって} $\Bmod$ は、$\Amod$ に有限個の剰余類、またはそれらの線型結合を添加して得られる。これから直ちに次が従う。

\medskip
\noindent\textbf{定理 IV。} \emph{$\Bmod$ が $\Amod$ を割る同じ階数の\mbox{加群}であり、\mbox{したがって}両者の\mbox{基本加群}が一致し、さらに \mbox{$\Norm(\Gmod\mid\Bmod)=\Norm(\Gmod\mid\Amod)$} であるならば、$\Bmod=\Amod$ である。}
%
\begingroup
\renewcommand{\thefootnote}{\arabic{footnote})}
\setcounter{footnote}{9}

ノルムの分解と加群の分解との関係について、次の定理が成り立つ。

\medskip
\noindent\textbf{定理 V。} \emph{いま}
\begin{equation}
 \Norm(\Gmod\mid\Amod)=r\cdot s,\qquad (r,s)=1
\tag{10}
\end{equation}
\emph{とする。このとき、ともに $\Amod$ を割る同じ階数の加群であって、}
\[
 r=\Norm(\Gmod\mid\Rmod),\qquad s=\Norm(\Gmod\mid\Smod)
\]
\emph{を満たす $\Rmod$ がただ一つ、同様に $\Smod$ がただ一つ存在する。$\Rmod$ と $\Smod$ は}
\[
 \Rmod=\Amod/(s);\qquad \Smod=\Amod/(r)
\]
\emph{によって定義され、$\Amod$ は $\Rmod$ と $\Smod$ の最小公倍数に等しい。}\footnote{最小公倍数 $[\Rmod,\Smod]$ は加群の意味で解する。すなわち、$\Rmod$ と $\Smod$ の双方で割り切れる線型形式全体である。これに対応して、加群の意味での最大公約数 $(\Rmod,\Smod)$ は、$\Rmod$ の一つの元と $\Smod$ の一つの元との和として表される線型形式全体と解する。}
\emph{$r_\varrho=(r,e_\varrho)$、$s_\varrho=(s,e_\varrho)$ とおけば、次の相反関係が成り立つ：}
\begin{equation}
 (s_\varrho)=\Amod/(\Rmod);\quad \Rmod=\Amod/(s_\varrho);
 \qquad
 (r_\varrho)=\Amod/\Smod;\quad \Smod=\Amod/(r_\varrho).
\tag{11}
\end{equation}

実際、一般に $r_i=(r,e_i)$、$s_i=(s,e_i)$ とおくと、(9) と (10) から
\[
 (r_i,s_i)=1;\qquad e_i=r_is_i;\qquad r=r_1\cdots r_\varrho;\qquad s=s_1\cdots s_\varrho
\]
を得る。また、各 $r_i$ はその次の $r_{i+1}$ を割り、各 $s_i$ はその次の $s_{i+1}$ を割る。そこで
\[
 \Rmod=(r_1\eta_1,\ldots,r_\varrho\eta_\varrho);\qquad
 \Smod=(s_1\eta_1,\ldots,s_\varrho\eta_\varrho)
\]
とおけば、$\Rmod$ と $\Smod$ は $\Amod$ を割る同じ階数の加群となる。$r_i$ と $s_i$ が互いに素であるため、$\Amod$ は $\Rmod$ と $\Smod$ の最小公倍数に等しく、さらに
\[
 \Norm(\Gmod\mid\Rmod)=r_1\cdots r_\varrho=r;\qquad
 \Norm(\Gmod\mid\Smod)=s_1\cdots s_\varrho=s
\]
を得る。さらに、
\[
 s_\varrho\Rmod\zero{\Amod};\qquad r_\varrho\Smod\zero{\Amod}
\]
が成り立つ。$c\Rmod\zero{\Amod}$ から $cr_\varrho\eta_\varrho\zero{\Amod}$ が従い、したがって $c\zero{(s_\varrho)}$ となる。これで $(s_\varrho)=\Amod/\Rmod$ が証明され、同様に $(r_\varrho)=\Amod/\Smod$ も証明される。また、
\[
 b(\eta)=b_1\eta_1+\cdots+b_\varrho\eta_\varrho\zero{\Amod/(s_\varrho)}
\]
から、すべての $r_i$ が $s_\varrho$ と互いに素であることにより、$b_i\zero{(r_i)}$ が従う。したがって $\Rmod=\Amod/(s_\varrho)$ であり、同様に $\Smod=\Amod/(r_\varrho)$ である。これで相反関係 (11) が証明された。特殊な場合 $r=1$ には、(11) は (5) となる。

なお、$\Rmod$ と $\Smod$ が\emph{一意に定まること}を示さなければならない。定理 V の条件を満たす加群を $\overline{\Rmod}$、$\overline{\Smod}$ とし、それらの単因子を $\bar r_i$、$\bar s_i$ とする。$\bar r_i$ と $\bar s_i$ は $e_i$ の因子であり、
\[
 r=\bar r_1\cdots\bar r_\varrho,
 \qquad
 s=\bar s_1\cdots\bar s_\varrho,
 \qquad
 (\bar r_i,\bar s_i)=1
\]
であるから、さらに
\[
 e_i=\bar r_i\bar s_i,
 \qquad
 \bar r_i=(r,e_i)=r_i,
 \qquad
 \bar s_i=(s,e_i)=s_i
\]
を得る。したがって $\overline{\Rmod}$ と $\overline{\Smod}$ は、単因子および基本加群に関して、それぞれ $\Rmod$ と $\Smod$ に一致する。そこで、単模変換 $\eta=Q(\zeta)$ によって新しい不定元 $\zeta$ を導入すれば、
\[
 \overline{\Rmod}=(r_1\zeta_1,\ldots,r_\varrho\zeta_\varrho);
\]
\[
 \Amod=\bigl(r_1s_1(q_{11}\zeta_1+\cdots+q_{1\varrho}\zeta_\varrho),\ldots,
 r_\varrho s_\varrho(q_{\varrho1}\zeta_1+\cdots+q_{\varrho\varrho}\zeta_\varrho)\bigr)
\]
となる。$\Amod\zero{\overline{\Rmod}}$ から、
\[
 r_i s_i(q_{i1}\zeta_1+\cdots+q_{i\varrho}\zeta_\varrho)\zero{\overline{\Rmod}}
\]
を得る。したがって $r_is_iq_{ij}\zero{(r_j)}$ である。$(r,s)=1$ であるから、これは $r_iq_{ij}\zero{(r_j)}$、すなわち $\Rmod\zero{\overline{\Rmod}}$ を与える。しかし $\Norm(\Gmod\mid\Rmod)=\Norm(\Gmod\mid\overline{\Rmod})$ であるから、定理 IV により $\Rmod=\overline{\Rmod}$ である。同様に $\Smod=\overline{\Smod}$ である。これで定理 V のすべての主張が証明された。

\endgroup
%
\begingroup
\renewcommand{\thefootnote}{\arabic{footnote})}
\setcounter{footnote}{10}

\section*{§ 3. 多項式から成るイデアル。}

$\overline{\mideal}$ を、任意の抽象的に定義された体 $P$ を係数とする、$n$ 個の不定元 $y_1\ldots y_n$ の\emph{多項式イデアル}とする。すなわち $\overline{\mideal}$ は、$\Phi_1(y)$ と $\Phi_2(y)$ とともに差 $\Phi_1(y)-\Phi_2(y)$ を含み、$\Phi(y)$ とともに $A(y)\cdot\Phi(y)$ を含む。\mbox{ここで} $A(y)$ は $\mathsf{P}$ を係数とする任意の多項式である\footnote{概念から自然に生じる「イデアル」という名称は、以前一般に用いられていた「加群」または「形式加群」に代わるものであり、ここですでに、線型形式から成る加群と区別するために必要となる。（E. N.）}。

$y$ に、不定元を係数とする変換
\begin{equation}
 y=U(x);\quad \text{たとえば}\quad
 \begin{cases}
 y_1=x_1,\\
 y_2=u_{21}x_1+x_2,\\
 \vdots\\
 y_n=u_{n1}x_1+\cdots+u_{n,n-1}x_{n-1}+x_n
 \end{cases}
\tag{12}
\end{equation}
を施し、不定元 $u_{\mu\nu}$ を多項式の係数領域に\mbox{添加}する。これによって係数領域は体 $\mathsf{P}(u)$ となる。この\mbox{添加}によって $\overline{\mideal}$ は $\overline{\mideal}_{(u)}$ となる。したがって $\overline{\mideal}_{(u)}$ は、$\overline{\mideal}$ に属する多項式 $\Phi_\lambda(y)$ とともに、有限個ずつ取った $\Phi_\lambda$ の線型結合 $\sum \alpha_\lambda(u)\Phi_\lambda(y)$ をすべて含む。\mbox{ここで} $\alpha_\lambda(u)$ は、$P$ を係数とする $u_{\mu\nu}$ の有理関数である。適当な $u$ の多項式を掛ければ、一般性を失わずに、$\alpha_\lambda(u)$ を $u$ の冪積と仮定できる。したがって、
\begin{equation}
 \sum \alpha_\lambda(u)\Phi_\lambda(y)\zero{\mideal_{(u)}};
 \qquad
 \Phi_\lambda(y)\zero{\overline{\mideal}}\quad\text{であり、逆も成り立つ。}
\tag{13}
\end{equation}
(12) によって $\overline{\mideal}_{(u)}$ は、$P(u)$ を係数とする $x_1,\ldots,x_n$ の多項式イデアル $\mideal$ となる：
\begin{equation}
 \overline{\mideal}_{(u)}=\mideal\quad \text{$y=U(x)$ によって。}
\tag{14}
\end{equation}
関係 (14) と (13) を組み合わせると、次のことが得られる：
\begin{equation}
 \begin{aligned}
 &\text{$F(x)\zero{\mideal}$ であり、}\\
 &F(x)=\Phi(y)=\sum \alpha_i(u)\Phi_i(y)=\sum \alpha_i(u)F_i(x)\\
 &\text{ならば、}\qquad F_i(x)\zero{\overline{\mideal}};
 \end{aligned}
\tag{15}
\end{equation}
一方、(14) と (15) から逆向きに (13) が再び得られる。

\medskip
\noindent\textbf{定義 IV。} \emph{$P(u)$ を係数とする $x_1,\ldots,x_n$ の多項式イデアル $\mideal$ で、関係 (15) を満たし、したがって $y=U(x)$ によって $\overline{\mideal}_{(u)}$ から生じるものを、変換イデアルと呼ぶ。}

変換イデアルは、\mbox{正則な多項式}が存在することによって特徴づけられる。\emph{次数 $r$ の多項式が $x_i$ に関して\mbox{正則である}とは、零でない係数をもつ項 $x_i^r$ を含むことをいう}。多項式 $F_\lambda(x)$ は $x_1$ に関して\mbox{正則である}ことが分かる。以下では、これらの変換イデアルを、$x$ のある冪積に関する線型形式から成る加群と\mbox{みなす}ことが本質となる。そのために、後で基本加群へ移行する基本イデアルの概念を\mbox{定めなければならない}。\mbox{まず}、以後一貫して用いる記法を定める。

\medskip
\noindent\textbf{記法。} \emph{$F^{(i)},f^{(i)},a^{(i)},\ldots$ は常に、$x_1,\ldots,x_{i-1}$ を含まない $x$ の多項式を表す。}

\medskip
\noindent\textbf{定義 V。} \emph{各イデアル $\mideal$ に対して、$n$ 個の基本イデアル $\gideal_0,\gideal_1,\ldots,\gideal_{n-1}$ を次の規定によって定義する。\mbox{第 $(i-1)$ 段}の基本イデアル $\gideal_{i-1}$ は、ある多項式 $b^{(i)}$---一般には $G(x)$ とともに変わる---が存在して、}
\begin{equation}
 b^{(i)}G(x)\zero{\mideal};\qquad b^{(i)}\ne0
\tag{16}
\end{equation}
\emph{を満たすような多項式 $G(x)$ を、すべて、しかもそれだけ含む。}

$\gideal$ が実際にイデアルであることは、式 (4) に続く $\Gmod$ の加群性の証明に正確に対応し、(16) はそこで (4) と同じ役割を果たす。$\gideal_i$ は $\gideal_{i-1}$ で割り切れる。なぜなら、各多項式 $b^{(i+1)}$ は同時に一つの $b^{(i)}$ でもあるからである。

基本イデアルについて、次が成り立つ。

\endgroup
%
\begingroup
\renewcommand{\thefootnote}{\arabic{footnote})}
\setcounter{footnote}{11}

\medskip
\noindent\textbf{定理 VI。} \emph{変換イデアルの基本イデアルは変換イデアルである}\footnote{定理 VI は §~7 ではじめて用いられる。}。

証明は、よく知られた Dedekind--Mertens の定理に基づく\footnote{Dedekind: Über einen arithmetischen Satz von Gauß, Mitt. d. deutsch. math. Ges. zu Prag 1892. F. Mertens: Über einen algebraischen Satz, Ber. d. Ak. d. Wissensch. Wien 101 (1892), S. 1560--1566。---不定係数をもつ多項式への Gauss の定理の Kronecker による拡張は、この定理の直接の帰結である。}。$a$ と $b$ を不定元とし、
\[
 \sum a_{i_1\ldots i_s}z_1^{i_1}\cdots z_s^{i_s}\cdot
 \sum b_{j_1\ldots j_s}z_1^{j_1}\cdots z_s^{j_s}
 =
 \sum c(a,b)_{k_1\ldots k_s}z_1^{k_1}\ldots z_s^{k_s}
\]
および
\[
 \sum i\le \nu,
 \qquad
 \sum j\le \mu,
 \qquad
 \sum k\le \nu+\mu
\]
とおく。さらに $\Amod,\Bmod,\mathfrak C$ を、それぞれ $a,b,c$ から有理整数によって生成される加群とする。このとき、\mbox{恒等式}
\begin{equation}
 \Amod^q\Bmod=\Amod^{q-1}\mathfrak C
\tag{17}
\end{equation}
が成り立つような指数 $q$ が存在する。ここで $\Amod^q$ は、$a$ の次数 $q$ の冪積と、それらの整係数線型結合から成る。

定理 VI を証明するため、いま
\begin{equation}
\begin{cases}
 G(x)=\Gamma(y)=\sum \alpha_\lambda(u)\Gamma_\lambda(y)=\sum \alpha_\lambda(u)G_\lambda(x),\\
 b^{(i)}(x)=\varphi(y)=\sum \beta_\mu(u)\varphi_\mu(y)
\end{cases}
\tag{18}
\end{equation}
とおく。ここで $\alpha_\lambda(u)$ と $\beta_\mu(u)$ は $u$ の冪積である。すると (16)、(14)、および (18) により、
\begin{equation}
 \sum\beta_\mu(u)\varphi_\mu(y)\cdot
 \sum\alpha_\lambda(u)\Gamma_\lambda(y)\zero{\overline{\mideal}_{(u)}};
 \qquad
 \sum\beta_\mu(u)\varphi_\mu(y)\ne0
\tag{19}
\end{equation}
が成り立つ。左辺は $u$ に関する二つの多項式の積であり、それらの係数は多項式 $\varphi_\mu(y)$ および $\Gamma_\lambda(y)$ である。(19) の右側の関係により、$u$ に関する積多項式の係数は $\overline{\mideal}$ に属する多項式である。そこで (17) の $a,b,c$ を、これらの $y$ の多項式で置き換えると、
\[
 \left\{\sum\beta_\mu(u)\varphi_\lambda(y)\right\}^q\cdot \Gamma_\lambda(u)\zero{\overline{\mideal}_{(u)}}
\]
を得る。これは (18) により、
\begin{equation}
 b^{(i)}(x)^q\cdot G_\lambda(x)\zero{\mideal};
 \qquad
 G_\lambda(x)\zero{\gideal_{i-1}}
\tag{20}
\end{equation}
と同値である。関係 (16)、(18)、および (20) は、\emph{基本イデアル $\gideal_{i-1}$ が実際に変換イデアルである}ことを示す。

\endgroup
%

\section*{§ 4. 多項式から成るイデアルと線型形式から成る加群との関係。}

常に変換イデアルであると仮定する多項式イデアル $\mideal$ を、線型形式から成る加群 $\Mmod_{i-1}$（$(i=1\ldots n)$）とみなすため、各多項式 $F(x)$ を、$x_1\ldots x_{i-1}$ の冪積 $\xi_\lambda$ に関する線型形式と考える：
\[
 F(x)=\sum a^{(i)}_\lambda\xi_\lambda.
\]
ここで §~3 の記法により、$a^{(i)}_\lambda$ は $x_i\ldots x_n$ の多項式である。イデアルの定義から、$\Mmod_{i-1}$ が、これらの整な量 $a^{(i)},b^{(i)},\ldots$ に関して加群の性質をもつことが直ちに従う。さらに、$x_1\ldots x_{i-1}$ の無限個の冪積 $\xi$ は\emph{線型に}しか現れないので、その線型独立性により、$\Mmod_{i-1}$ に対して不定元の役割を果たす。各加群 $\Mmod_{i-1}$ は無限個の不定元 $\xi$ を含むが、各線型形式に現れる不定元は有限個だけである。同様に、各基本イデアル $\gideal_{i-1}$ も線型形式から成る加群 $\Gmod_{i-1}$ とみなす。この場合にも、不定元は $x_1\ldots x_{i-1}$ の冪積である。$i=1$ の場合には、単位に対応する\emph{一つ}の不定元 $\xi_0$ だけがある。

以下のすべての議論の基礎となるのは、このように不定元 $\xi$ が無限個あっても、有限個の不定元に関する線型形式から成る加群へ帰着できるということである。すなわち、

\medskip
\noindent\textbf{定理 VII。} \emph{剰余類系 $\Gmod_{i-1}\mid \Mmod_{i-1}$ は、剰余類系 $\Gmod^*_{i-1}\mid \Mmod^*_{i-1}$ と同型である。ここで $\Mmod^*_{i-1}$ は有限個の不定元に関する加群であり、$\Gmod^*_{i-1}$ はその基本加群である。$x_{i+1}\ldots x_n$ を $P(u)$ に添加すると、加群 $\Mmod_{i-1}$ がもつ $1$ と異なる単因子は有限個だけであり、それらは $\Mmod^*_{i-1}$ の $1$ と異なる単因子である。}

ここで $\Mmod_{i-1}$ の単因子は、注 8) と同じように定義する。定理 VII に現れる同型は、次の定義に基づく。

\medskip
\noindent\textbf{定義 V。} \emph{二つの剰余類系を一対一に対応させ、二つの類の差には対応する二つの類の差を対応させ、また一つの類と一つの整な量との積には、対応する類と同じ整な量との積を対応させることができるとき、この二つの剰余類系を同型という。}

定理 VII は完全帰納法によって証明される。$i=1$ の場合、定理 VII は明らかである。実際、$\Mmod_0$ は一つの不定元 $\xi_0$ に関する線型形式から成る加群であり、$\xi_0$ は $x$ の零次の冪積、すなわち単位に対応する。整な量は $x_1\ldots x_n$ のすべての多項式である。基本イデアル $\gideal_0$ は $x_1\ldots x_n$ の全多項式から成る単位イデアルに等しく、したがって $\Gmod_0$ は加群 $(\xi_0)$、すなわち $\Mmod_0$ の基本加群に等しい。それゆえ、ここでは $\Gmod^*_0\mid\Mmod^*_0$ は $\Gmod_0\mid\Mmod_0$ と一致する。最後に、$\Mmod_0$ は階数 $1$ であるから、$1$ と異なる単因子を高々一つしかもたない。
%

まず $i=1$ から $i=2$ へ進み、ここで得られる $\Gmod_1\mid\Mmod_1$ の構造についての知見を用いて、一般の\mbox{帰納段階}を行う。そのため、はじめに $\Gmod_0\mid\Mmod_0$ について $x_1$ に関して有界な代表系を選べることを示す。$\gideal_1$ が $\gideal_0$ で割り切れることにより、これは $\Gmod^*_1$ の有限個の不定元へ導く。

$\mideal$ は変換イデアルであるから、§~3 により、$x_1$ に関して正則な多項式 $C^{(1)}(x)$ を少なくとも一つ含む。その次数を $k$ とする。したがって $\Mmod_0$ は線型形式 $C^{(1)}\xi_0$ を含む。$C^{(1)}(x)$ の正則性により、任意の多項式 $H(x)$ は次の表示をもつ：
\begin{equation}
 H(x)\equiv b^{(2)}_0+b^{(2)}_1x_1+\cdots+b^{(2)}_{k-1}x_1^{k-1}\quad (C^{(1)}(x)).
\tag{21}
\end{equation}
それゆえ、$C^{(1)}\xi_0$ が $\Mmod_0$ に属することから、$\Gmod_0\mid\Mmod_0$ に対して、実際に $x_1$ に関する次数が $k$ に達しない代表系を選べる。

特に、$\gideal_1$ に属する多項式 $G(x)$ と $\mideal$ に属する多項式 $F(x)$ について
\begin{equation}
 \left\{\begin{aligned}
  G(x)&\equiv g^{(2)}_0+g^{(2)}_1x_1+\cdots+g^{(2)}_{k-1}x_1^{k-1} &&(C^{(1)}(x)),\\
  F(x)&\equiv a^{(2)}_0+a^{(2)}_1x_1+\cdots+a^{(2)}_{k-1}x_1^{k-1} &&(C^{(1)}(x))
 \end{aligned}\right.
\tag{22}
\end{equation}
とおく。さらに $\Gmod^*_1$ を、すべての線型形式 $g(\xi)=g^{(2)}_0\xi_0+\cdots+g^{(2)}_{k-1}\xi_{k-1}$ から成る系とし、これに\mbox{対応して} $\Mmod^*_1$ を、すべての線型形式 $a(\xi)=a^{(2)}_0\xi_0+\cdots+a^{(2)}_{k-1}\xi_{k-1}$ から成る系とする。すると、$\gideal_1$ と $\mideal$ の\mbox{イデアル性から}、$\Gmod^*_1$ と $\Mmod^*_1$ は、整な量 $b^{(2)},c^{(2)},\ldots$ に関する、$\xi_0\ldots\xi_{k-1}$ の線型形式から成る加群であることが従う。\mbox{同様に}、$C^{(1)}(x)$ が生成する主イデアルも、これらの整な量に関する加群を与える：
\[
 \mathfrak C_1=(\zeta_0,\zeta_1,\ldots,\zeta_\nu,\ldots),
\]
ここで $\zeta_\nu=x_1^\nu C^{(1)}(x)$ とおく。$\zeta_\nu$ は、整な量 $b^{(2)},c^{(2)},\ldots$ に関して、互いに線型独立であり、また有限個の $\xi$ からも線型独立である。$\mathfrak C_1$ が $\Mmod_1$ で、したがって $\Gmod_1$ で割り切れることから、$\Gmod^*_1$ も $\Gmod_1$ で割り切れ、$\Mmod^*_1$ も $\Mmod_1$ で割り切れる。(22) を考慮すると、したがって
\begin{equation}
 \Gmod_1=(\Gmod^*_1,\mathfrak C_1),\qquad
 \Mmod_1=(\Mmod^*_1,\mathfrak C_1)
\tag{23}
\end{equation}
を得る。

(23) と $\xi$ および $\zeta$ の線型独立性から、$i=2$ に対する定理 VII は次のように従う。まず、$\mathfrak C_1$ が $\Mmod_1$ で\mbox{割り切れる}ことにより、$\Gmod_1\mid\Mmod_1=\Gmod^*_1\mid\Mmod_1$ となる。$\Gmod^*_1\mid\Mmod^*_1$ との同型を示すため、$\Gmod^*_1$ に\mbox{属するある線型形式族} $g^*(\xi)$ が $\Gmod^*_1\mid\Mmod^*_1$ の代表系をなすとする。$\xi$ と $\zeta$ の線型独立性により、\mbox{その各元について} $g^*(\xi)\equiv 0(\Mmod_1)$ から $g^*(\xi)\equiv 0(\Mmod^*_1)$ が従う。それゆえ、これらの $g^*(\xi)$ は $\Mmod_1$ を法としても互いに合同ではなく、$\Gmod^*_1\mid\Mmod_1$ の代表系をなし、この代表系によって与えられる $\Gmod^*_1\mid\Mmod_1$ から $\Gmod^*_1\mid\Mmod^*_1$ への対応は、定義 V の意味で同型である。

まったく同様に、$b^{(2)}g(\xi)\equiv0(\Mmod_1)$、$b^{(2)}\ne0$ から $b^{(2)}g(\xi)\equiv0(\Mmod^*_1)$、$b^{(2)}\ne0$ が従う。\mbox{このこと}は $\Gmod^*_1$ に属するすべての $g(\xi)$ について、しかもそれらについてだけ成り立つ。なぜなら (23) \mbox{により}、$\Gmod^*_1$ は、基本イデアル $\gideal_1$ に属し、同時に $\xi$ の線型形式である多項式すべてから成るからである。\mbox{したがって} $\Gmod^*_1$ は $\Mmod^*_1$ の基本加群として特徴づけられる。

最後に、$x_3\ldots x_n$ を $\mathsf{P}(u)$ に添加すると、
\begin{equation}
\left\{\begin{aligned}
 \Gmod^*_1&=(\eta_1\ldots\eta_\rho);&
 \Mmod^*_1&=(e_1\eta_1,\ldots,e_\rho\eta_\rho),\\
 \Gmod_1&=(\eta_\rho\ldots\eta_1,\zeta_0\ldots\zeta_\nu\ldots);&
 \Mmod_1&=(e_\rho\eta_\rho,\ldots,e_1\eta_1,\zeta_0\ldots\zeta_\nu\ldots)
\end{aligned}\right.
\tag{24}
\end{equation}
を得る。したがって、
\[
 (e_\rho)=\Mmod_1/\Gmod_1=\Mmod_1/(\Mmod_1,\eta_\rho),\qquad
 (e_{\rho-1})=(\Mmod_1,\eta_\rho)/\Gmod_1\quad\hbox{など}
\]
である。注 8) を考慮すると、これにより $e_\rho,e_{\rho-1},\ldots,e_1,1,\ldots,1,\ldots$ が $\Mmod_1$ の単因子であることが分かる。これで $i=2$ に対する定理 VII のすべての主張が証明された。
%

(22) と (23) では、$C^{(1)}(x)$ が $x_1$ に関して正則であることだけを用いたことに注意しよう。したがって、(12) の代わりに特殊な変換
\begin{equation}
 y_1=x_1;\quad y_2=u_{21}x_1+z_2;\quad\ldots;\quad y_n=u_{n1}x_1+z_n
\tag{25}
\end{equation}
を基礎とする場合にも、これらの式と、それらに結びついて定理 VII へ至る議論はそのまま成り立つ。(25) を
\begin{equation}
\left\{\begin{array}{l}
 z=U_1(x)\quad\hbox{すなわち}\\[2mm]
 z_2=x_2;\quad z_3=u_{32}x_2+x_3;\quad\ldots;\quad z_n=u_{n2}x_2+\cdots+u_{n,n-1}x_{n-1}+x_n
\end{array}\right.
\tag{26}
\end{equation}
と合成すれば、再び (12) が得られる。(25) によって $\overline{\mideal}_{(u)}$ が $\widetilde{\mideal}$ となり、$\widetilde{\gideal}_1,\widetilde{\Gmod}_1\ldots$ が $\widetilde{\mideal}$ に対して、$\gideal_1,\Gmod_1\ldots$ が $\mideal$ に対してもつのと同じ意味をもつとする。このとき、
\begin{equation}
 \widetilde{\Gmod}_1=(\widetilde{\Gmod}^{*}_1,\widetilde{\mathfrak C}_1),\qquad
 \widetilde{\Mmod}_1=(\widetilde{\Mmod}^{*}_1,\widetilde{\mathfrak C}_1)
\tag{27}
\end{equation}
が成り立つ。

実際、(27) は (26) の逆変換によって (23) からそのまま得られる。このことは $\mideal$ と $C^{(1)}(x)$ について明らかに成り立ち、したがって $\widetilde{\mathfrak C}_1$ と $\widetilde\Mmod_1$ についても成り立つ。また $\gideal_1$ についても成り立つ。なぜなら、
\[
 b^{(2)}(x)G(x)\equiv0\pmod{\mideal},\quad b^{(2)}\ne0,
 \qquad
 \widetilde b^{(2)}(z)\widetilde G(x,z)\equiv0\pmod{\widetilde\mideal},\quad \widetilde b^{(2)}\ne0
\]
は、$z=U_1(x)$ によって互いに同値だからである。したがって (26) によって $\widetilde\gideal_1=\gideal_1$ となり、それゆえ $\widetilde\Gmod_1=\Gmod_1$ となる。さらに、(22) によって与えられる $\Gmod_1^*$ と $\Mmod_1^*$ の定義から、$\widetilde\Gmod_1^*$ と $\widetilde\Mmod_1^*$ についても同じことが成り立ち、したがって表示 (27) 全体について成り立つ。
%

一般の帰納段階のため、次を仮定する：
\begin{equation}
\left\{\begin{aligned}
 \Gmod_{i-1}&=(\Gmod_{i-1}^*,\mathfrak C_{i-1}),&
 \Mmod_{i-1}&=(\Mmod_{i-1}^*,\mathfrak C_{i-1}),\\
 \mathfrak C_{i-1}&=(\zeta_0,\zeta_1,\ldots,\zeta_\nu,\ldots)
\end{aligned}\right.
\tag{28}
\end{equation}
ここで $\Gmod^*_{i-1}$ と $\Mmod^*_{i-1}$ は、整な量 $b^{(i)},c^{(i)},\ldots$ に関する、有限個の不定元 $\xi$ の線型形式から成る加群であり、これらの $\xi$ は $x_1,\ldots,x_{i-1}$ のある冪積である。また $\zeta$ は、これらの整な量に関して、互いに線型独立であり、かつ $\xi$ からも線型独立であるとする。(12) の代わりに特殊な変換
\begin{equation}
\begin{aligned}
 y_1&=x_1;\quad \ldots;\quad y_{i-1}=u_{i-1,1}x_1+\cdots+x_{i-1};\\
 y_i&=u_{i,1}x_1+\cdots+u_{i,i-1}x_{i-1}+z_i;\quad\ldots;\quad
 y_n=u_{n,1}x_1+\cdots+u_{n,i-1}x_{i-1}+z_n
\end{aligned}
\tag{29}
\end{equation}
を基礎とするときにも、(28) に対応する表示と、$\xi$ および $\zeta$ の線型独立性が成り立つものとする。この変換を、
\[
 z=U_{i-1}(x)\quad\hbox{すなわち}\quad
 z_i=x_i;\quad z_{i+1}=u_{i+1,i}x_i+x_{i+1};\quad\ldots;\quad
 z_n=u_{n,i}x_i+\cdots+u_{n,n-1}x_{n-1}+x_n
\]
と合成することによって (12) が得られ、さらに、この特殊な表示は $z=U_{i-1}(x)$ の逆変換によって (28) からそのまま得られるものとする。

仮定 (28) と $\xi$ および $\zeta$ の線型独立性から、同じ代表系を対応させることにより、$\Gmod_{i-1}\mid\Mmod_{i-1}$ と $\Gmod^*_{i-1}\mid\Mmod^*_{i-1}$ の同型が従う。これは (23) に続いて用いた議論とまったく同じである。同様に、$\Gmod^*_{i-1}$ は $\Mmod^*_{i-1}$ の基本加群となり、$\Mmod_{i-1}$ は $1$ と異なる単因子を有限個しかもたず、それらが $\Mmod^*_{i-1}$ の $1$ と異なる単因子であることも従う。

帰納段階を実行するため、まず再び、$\Gmod^*_{i-1}\mid\Mmod^*_{i-1}$ に対して、したがって $\Gmod_{i-1}\mid\Mmod_{i-1}$ に対して、$x_i$ に関して有界な代表系を選べることを示さなければならない。これは定理 III の (7) で証明した商表示
\begin{equation}
 \Gmod^*_{i-1}=\Mmod^*_{i-1}/(C^{(i)})
\tag{30}
\end{equation}
から従う。ここで $C^{(i)}$ は、$\Mmod^*_{i-1}$ から得られる恒等的に零でない任意の $\rho$ 次行列式であり、$\rho$ は $\Mmod^*_{i-1}$ の階数である。特殊変換 (29) に関する仮定から、$C^{(i)}$ は常に $x_i$ に関して正則であると仮定できる。実際、$\Mmod^*_{i-1}$ に対応する加群 $\widetilde\Mmod^*_{i-1}$ は同じ階数をもつ。したがって、$\widetilde\Mmod^*_{i-1}$ から得られる恒等的に零でない任意の $\rho$ 次行列式 $\widetilde C^{(i)}$ が、$z=U_{i-1}(x)$ によって正則な $C^{(i)}$ へ移るならば、仮定により $C^{(i)}$ は $\Mmod^*_{i-1}$ から得られる $\rho$ 次行列式となる。
%

$C^{(i)}$ が存在することから、(8) と同様に、$\xi$ に適切な番号を付ければ、特殊な形をもつ $\rho$ 個の線型形式
\[
 L_\lambda(\xi)=C^{(i)}\xi_\lambda+c^{(i)}_{\lambda,1}\xi_{\rho+1}+\cdots+c^{(i)}_{\lambda,s-\rho}\xi_s
 \qquad(\lambda=1\ldots\rho)
\]
が $\Mmod^*_{i-1}$ に属することが従う。さて、$\xi$ の任意の線型形式は次の形に直すことができる：
\begin{equation}
 H(\xi)\equiv a^{(i)}_1\xi_1+\cdots+a^{(i)}_\rho\xi_\rho+b^{(i)}_1\xi_{\rho+1}+\cdots+b^{(i)}_{s-\rho}\xi_s
 \quad (L_1(\xi),\ldots,L_\rho(\xi)).
\tag{31}
\end{equation}
ここで $a^{(i)}_1\ldots a^{(i)}_\rho$ は $x_i$ に関して有界であり、$C^{(i)}$ の次数 $k$ に達しない。この表示は一意である。実際、
\[
 a^{(i)}_1\xi_1+\cdots+a^{(i)}_\rho\xi_\rho+b^{(i)}_1\xi_{\rho+1}+\cdots+b^{(i)}_{s-\rho}\xi_s
 \equiv0\quad (L_1(\xi),\ldots,L_\rho(\xi))
\]
ならば、$x_i$ に関する次数を考慮しながら $\xi_1\ldots\xi_\rho$ の係数を比較することにより、すべての $a^{(i)}$ と $b^{(i)}$ が\mbox{恒等的に零である}ことが従う。

特に、いま
\[
 H(\xi)\equiv0\pmod{\Gmod^*_{i-1}};\qquad
 \hbox{したがって}\qquad
 C^{(i)}H(\xi)\equiv0\pmod{\Mmod^*_{i-1}}\quad\hbox{(30) により}
\]
とする。すると、定理 III の\mbox{証明}と同じように、
\[
 C^{(i)}H(\xi)\equiv\sum_{\tau=1}^{s-\rho}
 \left\{C^{(i)}b^{(i)}_\tau-\sum_\lambda a^{(i)}_\lambda c^{(i)}_{\lambda,\tau}\right\}\xi_{\rho+\tau}
 \equiv0\pmod{\Mmod^*_{i-1}}
\]
を得る。したがって、定理 III の場合と同様に $\xi_{\rho+1}\ldots\xi_s$ の係数は零でなければならないから、
\[
 C^{(i)}b^{(i)}_\tau=\sum_\lambda a^{(i)}_\lambda c^{(i)}_{\lambda,\tau}
 \qquad(\tau=1\ldots s-\rho)
\]
となる。これから $b^{(i)}_\tau$ も有界であり、$c^{(i)}_{\lambda,\tau}$ の最大次数に達しないことが従う。これで、$\Gmod^*_{i-1}\mid\Mmod^*_{i-1}$、\mbox{すなわち} $\Gmod_{i-1}\mid\Mmod_{i-1}$ に対して、$x_i$ に関して有界で、ある固定次数 $r$ に達しない代表系が存在することが\mbox{証明された}。
%

いま $\vartheta_1\ldots\vartheta_t$ によって、有限個の冪積
\[
 x_i^\alpha\xi_\lambda\quad(\alpha=0,1,\ldots,k-1;\ \lambda=1,2,\ldots,\rho),\qquad
 x_i^\beta\xi_{\rho+\tau}\quad(\beta=0,1,\ldots,r-1;\ \tau=1,2,\ldots,s-\rho)
\]
を表す。また、可算無限個の式
\[
 x_i^\gamma L_\lambda\quad(\gamma=0,1\ldots\ldots;\ \lambda=1,2,\ldots,\rho),
 \qquad
 x_i^\gamma\zeta_\nu\quad(\gamma,\nu=0,1\ldots\ldots)
\]
を一つの無限列 $\omega_0,\omega_1,\ldots,\omega_\mu\ldots$ に並べる。このとき $\vartheta$ と $\omega$ は、整な量 $b^{(i+1)},c^{(i+1)},\ldots$ に関して、それぞれ相互に線型独立であり、また両者の間でも線型独立である。実際、
\[
 \sum c^{(i+1)}_{\alpha\lambda} x_i^\alpha\xi_\lambda+
 \sum c^{(i+1)}_{\beta\tau} x_i^\beta\xi_{\rho+\tau}+
 \sum d^{(i+1)}_{\gamma\lambda} x_i^\gamma L_\lambda(\xi)+
 \sum e^{(i+1)}_{\gamma\nu}x_i^\gamma\zeta_\nu=0
\]
とする。各和は有限個の項だけにわたる。仮定した $\xi$ と $\zeta$ の線型独立性、および整な量 $b^{(i)}$ に関する $\xi$ 相互の線型独立性により、すべての $e^{(i+1)}_{\gamma\nu}$ が零となる。さらに表示 (31) の一意性から $c^{(i+1)}_{\alpha\lambda}$ と $c^{(i+1)}_{\beta\tau}$ が零となり、最後に $L_\lambda(\xi)$ の線型独立性から $d^{(i+1)}_{\gamma\lambda}$ が零となる。

いま加群 $(\mathfrak C_{i-1},L_1(\xi),\ldots,L_\rho(\xi))$ を、整な量 $b^{(i+1)}$ に関する加群 $\mathfrak C_i$ とみなす。$\Gmod_{i-1},L_1(\xi),\ldots,L_\rho(\xi)$ が $\Mmod_{i-1}$ で割り切れることにより、$\mathfrak C_i$ は $\Mmod_i$ で割り切れ、
\[
 \mathfrak C_i=(\omega_0,\omega_1,\ldots,\omega_\mu\ldots)
\]
となる。各 $H(x)\equiv0(\gideal_{i-1})$ に対して、(28)、(31)、および上で証明したことから、
\[
 H(x)\equiv b^{(i+1)}\vartheta_1+\cdots+b^{(i+1)}_t\vartheta_t\quad(\mathfrak C_i)
\]
が、整な量 $b^{(i+1)},c^{(i+1)}$ に関する加群等式として成り立つ。一方、$\gideal_i$ は $\gideal_{i-1}$ で割り切れ、$\mideal$ は $\gideal_i$ で割り切れる。そこで特に、$\mathfrak C_i$ に属する各 $G(x)$ と $\Mmod_i$ に属する各 $F(x)$ について
\[
 G(x)\equiv g^{(i+1)}_1\vartheta_1+\cdots+g^{(i+1)}_t\vartheta_t\quad(\mathfrak C_i),\qquad
 F(x)\equiv a^{(i+1)}\vartheta_1+\cdots+a^{(i+1)}\vartheta_t\quad(\mathfrak C_i)
\]
とおき、すべての $g(\vartheta)$ から成る加群を $\Gmod_i^*$、すべての $a(\vartheta)$ から成る加群を $\Mmod_i^*$ と表す。すると、(23) へ導いた議論とまったく同じ議論により、
\begin{equation}
 \Gmod_i=(\Gmod_i^*,\mathfrak C_i),\qquad
 \Mmod_i=(\Mmod_i^*,\mathfrak C_i),\qquad
 \mathfrak C_i=(\omega_0,\omega_1,\ldots,\omega_\mu\ldots)
\tag{32}
\end{equation}
を得る。

(32) の導出でも、$C^{(i)}$ が $x_i$ に関して正則であることだけを用いた。したがって、特殊変換 (29) の添字を一つ増やして基礎とする場合にも、議論全体がそのまま成り立つ。また、(27) に続く議論とまったく同じ議論により、この特殊な表示は $z=U_i(x)$ の逆変換によって (32) からそのまま得られる。

これによって、一般の帰納段階の仮定 (28) などが、添字を一つ増やした場合についてすべて証明された。そして、すでに示したように、これらの仮定から定理 VII が直ちに従うので、定理 VII は一般に証明された。同時に、$C^{(i)}$ の任意性から生じる $\Gmod^*_{i-1},\Mmod^*_{i-1}$ の任意性が、剰余類系 $\Gmod^*_{i-1}\mid\Mmod^*_{i-1}$ の同型によって解消されることも示された。
%
\begingroup
\renewcommand{\thefootnote}{\arabic{footnote})}
\setcounter{footnote}{13}

特に、§~3 で係数領域として基礎に置いた体 $\mathsf{P}$ の標数が零であるとする。すなわち、$\mathsf{P}$ に\mbox{含まれ}、\mbox{単位から生成される}素体が有理数型であるとする。このとき不定元 $u_{\mu\nu}$ を $P$ の元 $\bar u_{\mu\nu}$ に\mbox{特殊化して}、$i=1\ldots n$ のそれぞれについて $C^{(i)}$ が $x_i$ に関して正則であり続けるようにできる。上の議論は、\mbox{この特殊化}のもとでも定理 VII が保たれることを示す\footnote{Hentzelt は、任意の $\mathsf{P}$ の代わりに全複素数体だけを取り、主として後者の場合を扱っている。不定元を基礎に置くのは、本来の\mbox{消去理論において}だけである。さらに Hentzelt は、ここで与えた推論と本質的に同じ推論によって、終結式形式を作るには、各 $\Mmod_{i-1}$ において、\mbox{固定された限界}を超える何らかの有限次数まで $x$ について進めば十分であることを示している。ただし、定理 VII で与えた $\Gmod_{i-1}\mid\Mmod_{i-1}$ と $\Gmod^*_{i-1}\mid\Mmod^*_{i-1}$ との同型が欠けており、この同型こそが次数の値からの独立性の理由を明らかにする。（E. N.）}。しかし、この場合、基本加群と\mbox{単因子}が、\mbox{一般の場合}から\mbox{特殊化} $u_{\mu\nu}=\bar u_{\mu\nu}$ によって必ず生じるとは限らない\footnote{$\mideal=(x^2,ux+y)$ に対して $\gideal_0=(1)$、$\gideal_1=(1)$ となる。$C^{(1)}=x^2$ を選べば、$\Gmod_1^*=(1,x)$、$\Mmod_1^*=(ux+y,xy)$ である。ここで $\Mmod_1^*$ は\mbox{単因子} $y^2$ と $1$ をもち、$C^{(2)}=y^2$ を選べる。$u=0$ とすると、$C^{(1)}$ と $C^{(2)}$ はそれぞれ $x$ と $y$ に関して正則であり続けるが、$\widetilde{\Mmod}_1^*=(y,xy)$ は\mbox{単因子} $y,y$ をもつ。したがって $\widetilde{\Gmod}_1\mid\widetilde{\Mmod}_1$ はなお、有限個の不定元に関する加群の剰余類系と同型であるが、$\Gmod_1\mid\Mmod_1$ とは同型でない。\mbox{これに対して} $C^{(1)}=ux+y$ を選び、$C^{(1)}$ が正則であり続けるように\mbox{特殊化}したならば、同型も保たれたであろう。（E. N.）}。

\endgroup
%

\section*{§ 5. 多項式から成るイデアルの終結式形式と単因子形式。}

$x_{i+1}\ldots x_n$ を $\mathsf{P}(u)$ に添加する。すると $\Gmod_{i-1}$ と $\Mmod_{i-1}$、したがって $\Gmod^*_{i-1}$ と $\Mmod^*_{i-1}$ も、整な量を\emph{一つ}の\mbox{不定元} $x_i$ の多項式とみなせる線型形式から成る加群となる。定理 VII により、この場合、$\Mmod_{i-1}$ の $1$ と異なる\mbox{単因子}と、高々有限個の $1$ に等しい単因子は、$\Mmod^*_{i-1}$ の単因子と一致する。\mbox{これら}の単因子の積、\mbox{すなわち} $\Mmod^*_{i-1}$ の第 $\rho$ 行列式因子は、$\Gmod^*_{i-1}$ から $\Mmod^*_{i-1}$ への移行変換の行列式に等しく、\mbox{したがって}\mbox{定義 III} により $N(\Gmod^*_{i-1}\mid\Mmod^*_{i-1})$ に等しかった。(24)、または一般の $i$ について (28) から得られる同様の式により、\mbox{これら}の単因子の積は、$\Gmod_{i-1}$ から $\Mmod_{i-1}$ への移行変換の行列式にも等しいことが分かる。\mbox{したがって}、\mbox{これは} $N(\Gmod_{i-1}\cdot\Mmod_{i-1})$ と表すべきである。ゆえに、
\[
 N(\Gmod_{i-1}\mid\Mmod_{i-1})=N(\Gmod^*_{i-1}\mid\Mmod^*_{i-1})=R^{(i)}(x)
\]
となる。ここで多項式 $R^{(i)}(x)$、\mbox{すなわち} $\Mmod^*_{i-1}$ から得られるすべての $\rho$ 次行列式の、多項式の意味での最大公約元は、$x_{i+1}\ldots x_n$ に関して整かつ\mbox{原始的}であると仮定できる。したがって、$C^{(i)}$ の因子として $x_i$ に関して正則となる。同様に、$\Mmod_{i-1}$ または $\Mmod^*_{i-1}$ の最高単因子 $E^{(i)}(x)$ も、$x_{i+1}\ldots x_n$ に関して\mbox{整かつ}\mbox{原始的}であり、したがって $x_i$ に関して正則であると仮定できる。これにより、次の定義を得る。

\medskip
\noindent\textbf{定義 VI。} \emph{多項式 $R^{(i)}(x)$ を $\mideal$ の\mbox{第 $i$ 段終結式}と呼び、$E^{(i)}(x)$ を\mbox{第 $i$ 段単因子}と呼ぶ。したがって、\mbox{第 $i$ 段終結式}は基本イデアル $\gideal_{i-1}$ の $\mideal$ に関するノルムであり、$x_{i+1}\ldots x_n$ を $P(u)$ に添加したもとで、加群 $\Gmod_{i-1}$ の $\Mmod_{i-1}$ に関するノルムとみなされる。同様に、同じ添加のもとで $(E^{(i)}(x))$ は商 $\Mmod_{i-1}/\Gmod_{i-1}$ に\mbox{等しくなる}。$\mideal$ の終結式形式および単因子形式とは、それぞれ積}
\[
 R_\mideal=R^{(1)}\cdot R^{(2)}\cdot\ldots\cdot R^{(n)};\qquad
 E_\mideal=E^{(1)}\cdot E^{(2)}\cdot\ldots\cdot E^{(n)}
\]
\emph{のことである。}

終結式形式と単因子形式について、次が成り立つ。
%

\medskip
\noindent\textbf{定理 VIII。} \emph{終結式形式は単因子形式で割り切れる。逆に、$E_\mideal$ のある冪は $R_\mideal$ で割り切れる。$E_\mideal$ と $R_\mideal$ は $\mideal$ で割り切れる。より一般に、$E^{(i)}\cdots E^{(n)}\gideal_{i-1}\equiv0\pmod{\mideal}$ であり、したがって $R^{(i)}\cdots R^{(n)}\gideal_{i-1}\equiv0\pmod{\mideal}$ である。}

実際、ノルムは最高単因子で割り切れ、逆に、この最高単因子のある冪はノルムで割り切れる。したがって $x_{i+1}\ldots x_n$ を添加すれば、この割り切り関係は $R^{(i)}(x)$ と $E^{(i)}(x)$ について成り立つ。しかし $R^{(i)}(x)$ と $E^{(i)}(x)$ は $x_{i+1}\ldots x_n$ に関して原始多項式であると仮定したので、この割り切り関係はすべての不定元 $x_i,x_{i+1},\ldots,x_n$ に関して成り立つ。それゆえ、すべての不定元 $x$ に関するこれらの式の定義により、$R_\mideal$ は $E_\mideal$ で割り切れ、$E_\mideal$ のある冪は $R_\mideal$ で割り切れる。

さらに、定理 VII と定理 II により、$x_{i+1}\ldots x_n$ を添加したもとで、最高単因子 $E^{(i)}(x)$ は商 $\Mmod_{i-1}/\Gmod_{i-1}$ として定義される。したがって、すべての不定元 $x$ の多項式へ戻ると、各
\begin{equation}
 G(x)\equiv0(\gideal_{i-1})
 \quad\hbox{に対して}\quad b^{(i+1)}\ne0
 \quad\hbox{が存在し、}\quad b^{(i+1)}E^{(i)}(x)G(x)\equiv0(\mideal)
\tag{33}
\end{equation}
となる。特に $\gideal_0$ は単位イデアルに等しいから、
\[
 b^{(2)}E^{(1)}(x)\equiv0(\mideal);\quad b^{(2)}\ne0,
 \qquad\hbox{したがって}\qquad E^{(1)}(x)\equiv0(\gideal_1)
\]
を得る。一般に、
\[
 E^{(1)}(x)\cdots E^{(i-1)}(x)\equiv0(\gideal_{i-1})
\]
に対して、(33) により $b^{(i+1)}\ne0$ が存在し、$b^{(i+1)}E^{(1)}\cdots E^{(i)}\equiv0(\mideal)$ となる。したがって $E^{(1)}\cdots E^{(i)}\equiv0(\gideal_i)$ となる。$\gideal_n=\mideal$ であるから、
\begin{equation}
 E_\mideal=E^{(1)}\cdots E^{(n)}\equiv0(\mideal)
 \quad\hbox{であり、したがって}\quad
 R_\mideal=R^{(1)}\cdots R^{(n)}\equiv0(\mideal).
\tag{34}
\end{equation}

(33) において $G(x)$ が $\gideal_{i-1}$ のイデアル基底をなす有限個の多項式を遍歴するものとし、それらに対応する $b^{(i+1)}(x)$ の積を $c^{(i+1)}(x)$ と書く。このとき、
\[
 c^{(i+1)}E^{(i)}(x)\gideal_{i-1}\equiv0\pmod{\mideal};
 \qquad c^{(i+1)}\ne0
 \quad\hbox{であり、したがって}\quad E^{(i)}(x)\gideal_{i-1}\equiv0(\gideal_i)
\]
となる。これを上と同じように有限回繰り返すと、
\[
 E^{(n)}(x)\cdots E^{(i)}(x)\gideal_{i-1}\equiv0(\mideal)
\]
となり、したがってまた
\[
 R^{(n)}(x)\cdots R^{(i)}(x)\gideal_{i-1}\equiv0(\mideal)
\]
となる。これで定理 VIII のすべての主張が証明された。
%

定理 VIII は本質的に単因子形式の性質に基づく。一方、終結式形式が $\mideal$ に対して特徴的な意味をもつことは、次の定理が示す。

\medskip
\noindent\textbf{定理 IX。} \emph{$\nideal$ が $\mideal$ を割る---ここで「割る」はイデアルの意味で解する---ものとする。終結式形式 $R_\nideal$ と $R_\mideal$ が一致するならば、イデアル $\nideal$ と $\mideal$ は一致する。}

$\gideal_0\ldots\gideal_{n-1},\gideal_n=\mideal$ と $\hideal_0\ldots\hideal_{n-1},\hideal_n=\nideal$ を、それぞれ $\mideal$ と $\nideal$ の基本イデアルとする。まず $\gideal_0$ と $\hideal_0$ はともに単位イデアルに等しいから、$\gideal_0=\hideal_0$ である。いま $\gideal_{i-1}=\hideal_{i-1}$ と仮定する。このとき $\Mmod_{i-1}$ と $\Nmod_{i-1}$ は同じ基本加群 $\Gmod_{i-1}$ をもつ。仮定 $R_\nideal=R_\mideal$ は、$x_{i+1}\ldots x_n$ を添加すれば、
\[
 N(\Gmod_{i-1}\mid\Nmod_{i-1})=N(\Gmod_{i-1}\mid\Mmod_{i-1})
 \quad\hbox{または}\quad
 N(\Gmod^*_{i-1}\mid\Nmod^*_{i-1})=N(\Gmod^*_{i-1}\mid\Mmod^*_{i-1})
\]
と書ける。ここで (32) に対応して $\Nmod_{i-1}=(\Nmod^*_{i-1},\mathfrak C_{i-1})$ とおいた。したがって $\Nmod^*_{i-1}$ も $\xi$ の線型形式から成る加群であり、$\Gmod^*_{i-1}$ がその基本加群である。$\Mmod_{i-1}$ が $\Nmod_{i-1}$ で割り切れることから、$\Mmod^*_{i-1}$ も $\Nmod^*_{i-1}$ で割り切れる。それゆえ定理 IV により、この添加のもとで $\Nmod^*_{i-1}$ と $\Mmod^*_{i-1}$ は一致し、したがって $\Nmod_{i-1}$ と $\Mmod_{i-1}$ も一致する。

ところで、$x_{i+1}\ldots x_n$ の添加とは、$\Mmod_{i-1}$ または $\mideal$ に、
\[
 b^{(i+1)}(x)G(x)\equiv0\pmod{\mideal};
 \qquad b^{(i+1)}\ne0
\]
を満たすすべての多項式 $G(x)$ を加えることを意味する。すなわち、この添加によって $\Mmod_{i-1}$ または $\mideal$ は $\gideal_i$ となり、これに対応して $\nideal$ は $\hideal_i$ となる。したがって $\gideal_i=\hideal_i$ である。ゆえに $\gideal_n=\hideal_n$、すなわち $\mideal=\nideal$ が証明された。
%

最後に、同じ移行原理を定理 V に適用すると、さらに次を得る。

\medskip
\noindent\textbf{定理 X。} \emph{$R^{(i)}=S^{(i)}T^{(i)}$ を、$R^{(i)}$ の、多項式の意味で互いに素な多項式 $S^{(i)}$ と $T^{(i)}$ への\mbox{分解とする}。\mbox{このとき}、ともに $\mideal$ を割り、第 $(i-1)$ 段の\mbox{基本イデアル} $\gideal_{i-1}$ を共有する\mbox{イデアル} $\fraks$ がただ一つ、\mbox{同様に}\mbox{イデアル} $\tideal$ がただ一つ存在し、$S^{(i)}=N(\Gmod_{i-1}\mid\Smod_{i-1})$、$T^{(i)}=N(\Gmod_{i-1}\mid\mathfrak T_{i-1})$ となる。$\fraks$ と $\tideal$ は、\mbox{それぞれ次を満たす多項式} $S(x)$ と $T(x)$ 全体として定義される：}
\[
 s^{(i+1)}(x)T^{(i)}(x)S(x)\equiv0(\mideal);\quad s^{(i+1)}\ne0,
\]
\[
 t^{(i+1)}(x)S^{(i)}(x)T(x)\equiv0(\mideal);\quad t^{(i+1)}\ne0.
\]
\emph{また、$\gideal_i$ は $\fraks$ と $\tideal$ の最小公倍数に等しい：}
\[
 \gideal_i=[\fraks,\tideal].
\]

ここで注意すべきことは、$\fraks$ と $\tideal$ に対する $s^{(i+1)}$ と $t^{(i+1)}$ は、それぞれ $\fraks$ と $\tideal$ の基底元に対応する $s^{(i+1)}$ と $t^{(i+1)}$ の積として、固定して選べるということだけである。また、上で述べたように、$x_{i+1}\ldots x_n$ を\mbox{添加する}と $\mideal$ は $\gideal_i$ となる。特に、$R^{(i)}$ の分解は、既約多項式の冪---準素因子---に至るまで続けることが\mbox{でき}、それに対応して $\gideal_i$ は最小公倍数として表される。
%
\begingroup
\renewcommand{\thefootnote}{\arabic{footnote})}
\setcounter{footnote}{15}

\section*{§ 6. 本来の消去理論。}

定理 VIII または式 (34) により、終結式形式と単因子形式は $\mideal$ で割り切れるので、$\mideal$ のすべての零点で零となる。ここで「$\mideal$ の零点」とは、$\mathsf{P}(u;x_{i+1}\ldots x_n)$ から得られる代数的閉体に属する $x_1\ldots x_i$ の値の組\footnote{任意の体が、本質的に一意な仕方で代数的閉体へ拡大できることを、Steinitz は Algebraische Theorie der Körper（J. f. M. 137 (1910), S. 167--309）で示し、それによって代数学の基本定理の有理的な対応物を与えた。実際には、各 $i=1,\ldots,n$ について、$\mathsf{P}(u,x_{i+1},\ldots,x_n)$ の\mbox{有限拡大体}を扱う。（E. N.）}であって、その値の組において $\mideal$ のすべての多項式が零となるものをいう。ここでは $x_{i+1}\ldots x_n$ を係数領域に添加したものと考える（$i=1,\ldots,n$）。後で示すように、これらの添加された $x_{i+1}\ldots x_n$ は $\mathsf{P}(u)$ の元によって置き換えることもできる。消去理論の内容は、逆に、終結式形式の零点から $\mideal$ の零点を導くことにある。この逆向きの推論は、本節で展開する逐次消去に基づく。次節では、そこで現れる形式 $D^{(i)}(x)$ が $E^{(i)}(x)$ で割り切れることを示す。したがって、$E^{(i)}(x)$ のある冪が $R^{(i)}(x)$ で割り切れることから、求める逆向きの推論が従う。

ここで述べる逐次消去の理論は、次の定理に基づく。

\medskip
\noindent\textbf{定理 XI。} \emph{$\bideal$ を、$P$ を係数とする一つの不定元 $t$ の多項式から成るイデアル、\mbox{したがって}\mbox{主イデアルとする}。$h(t)$ を $\bideal$ に属する次数 $k$ の多項式とし、$\Bmod$ を、$h(t)$ を法として $\bideal$ が移る、$1,t,\ldots,t^{k-1}$ の\mbox{線型形式}\allowbreak\mbox{から成る}\allowbreak\mbox{加群}とする。このとき $\Bmod$ の階数は $k-p$ である。ここで $p$ は $\bideal$ の\mbox{基底多項式} $f(t)$ の次数である。}

仮定により $h(t)=h_1(t)\cdot f(t)$ であり、$h_1(t)$ の次数は $k-p$ である。したがって、多項式 $f(t),tf(t),\ldots,t^{k-p-1}f(t)$ が表す、$h(t)$ を法とする $\bideal$ の剰余類は線型独立である。また、$h(t)$ を法とする $\bideal$ の任意の剰余類は、この $k-p$ 個の特別な剰余類によって、$P$ の係数を用いて線型に表される。一方、$h(t)$ を法として $\bideal$ は $1,t,\ldots,t^{k-1}$ の\mbox{線型形式}\allowbreak\mbox{から成る}\allowbreak\mbox{加群}となるので、その階数はちょうど $k-p$ である。

\endgroup
%
\begingroup
\renewcommand{\thefootnote}{\arabic{footnote})}
\setcounter{footnote}{16}

いま $\aideal=\aideal_1$ を変換多項式イデアルとし、$A^{(1)}(x)$ を $\aideal_1$ に属する、$x_1$ に関して正則な次数 $k_1$ の\mbox{多項式}\allowbreak\mbox{とする}。また $\mathfrak A_1$ を、$A^{(1)}(x)$ を法として $\aideal_1$ が移る、$1,x_1,\ldots,x_1^{k_1-1}$ の線型形式から成る加群とする。\mbox{その}\mbox{係数}は\mbox{多項式} $a^{(2)}(x)$ である。さらに $\aideal_2$ を、$\mathfrak A_1$ から得られるすべての $k_1$ 次行列式が生成する $x_2\ldots x_n$ の\mbox{イデアル}とする。定理 XI により、$\aideal_2$ が零イデアルであるための必要十分条件は、$\aideal_1$ に\mbox{属する}\mbox{多項式}が、\mbox{多項式}の意味で $x_1$ に関する次数 $p_1>0$ の共通因子をもつことである。$\aideal_2$ が\mbox{零イデアル}\allowbreak\mbox{でなければ}、$\aideal$ は\mbox{変換イデアル}であると仮定したので、$\aideal_2$ は $x_2$ に関して正則な次数 $k_2$ の多項式 $A^{(2)}(x)$ を\mbox{含む}。このことは、\mbox{特殊変換} (25) と (26) から変換 (12) を合成すれば直接に分かる。再び $\mathfrak A_2$ を、$A^{(2)}(x)$ を\mbox{法として} $\aideal_2$ が\mbox{移る}、$1,x_2,\ldots,x_2^{k_2-1}$ の線型形式から成る加群とし、$\aideal_3$ を、$\mathfrak A_2$ から得られるすべての $k_2$ 次\mbox{行列式}が生成する $x_3\ldots x_n$ のイデアルとする。このイデアルは $x_3$ に関して正則な多項式 $A^{(3)}(x)$ を\mbox{含まなければならない}。

この仕方で、零イデアルに達するか、または $\aideal_{n+1}$ が単位イデアルとなるまで手続きを続ける。$\aideal_{n+1}$ は $x$ を含まないので、零イデアルまたは単位イデアルでしかありえない。イデアル $\aideal_1,\aideal_2,\ldots$ は、基礎に\mbox{選んだ}多項式 $A^{(\lambda)}(x)$ には依存せず、$\aideal$ によって一意に定まることが分かる。これは \mbox{$\aideal_1=\aideal$} \mbox{について}\mbox{明らかであるから}、$\aideal_1,\ldots,\aideal_\lambda$ について成り立つと仮定する。$A^{(i)}(x)$ を法として $\aideal_\lambda$ が\mbox{線型形式}\allowbreak\mbox{から成る}加群 $\mathfrak A_\lambda$ へ移るとき、$\mathfrak A_\lambda$ は、$x_i$ に関する次数が $k_i$ に達しない $\aideal_i$ の多項式全体\mbox{としても}\allowbreak\mbox{特徴づけられる}。\mbox{したがって} $\mathfrak A_i$ は $A^{(i)}(x)$ の次数 $k_i$ だけに依存する。いま $\bar A^{(i)}(x)$ を $\aideal_\lambda$ に\mbox{属する}、$x_i$ に関して正則な次数 $\bar k_\lambda>k_\lambda$ の多項式とし、$\overline{\mathfrak A}_\lambda$ を対応する線型形式\mbox{から成る}加群、$\bar\aideal_{\lambda+1}$ を $\overline{\mathfrak A}_i$ \mbox{から}\allowbreak\mbox{得られる} $\bar k_\lambda$ 次行列式が\mbox{生成するイデアル}とする。このとき加群等式
\[
 \overline{\mathfrak A}_\lambda=(\mathfrak A_\lambda,A^{(\lambda)},x_\lambda A^{(\lambda)},\ldots,x^{\bar k_\lambda-k_\lambda-1}A^{(\lambda)})
\]
が成り立つ。$A^{(\lambda)}(x)$ は $x_\lambda$ に関して正則であるから、多項式 $A^{(\lambda)},x_\lambda A^{(\lambda)},\ldots$ を線型形式の\mbox{新しい}\allowbreak\mbox{不定元として}導入できる。したがって $\mathfrak A_i$ から得られる $k_i$ 次行列式は $\overline{\mathfrak A}_i$ から得られる $\bar k_\lambda$ 次行列式と\mbox{一致し}、\mbox{それゆえ} $\bar\aideal_{\lambda+1}=\aideal_{\lambda+1}$ となる\footnote{原理的には、これは $\Gmod^*_{i-1}\mid\Mmod^*_{i-1}$ と $\Gmod_{i-1}\mid\Mmod_{i-1}$ との同型の基礎となった推論と同じである。}。

さらに $\aideal_{\lambda+1}$ は常に $\aideal_\lambda$ で割り切れる。$\aideal_{\lambda+1}$ が零イデアルならば明らかである。そうでなければ $\mathfrak A_\lambda$ の\mbox{階数}は $k_\lambda$ である。したがって $\aideal_{\lambda+1}$ は、定義により、$\mathfrak A_i$ の $k_i$ 次行列式から成るイデアルであり、$\mathfrak A_\lambda$ で、\mbox{したがって} $\aideal_\lambda$ で割り切れる。それゆえ各 $\aideal_{\lambda+1}$ は $\aideal$ で割り切れる。したがって $\aideal_{n+1}$ が\mbox{単位イデアル}\mbox{となるならば}、$\aideal$ も単位イデアルである。

いま $\aideal$ が単位イデアルではなく、$\aideal_1\ne0,\ldots,\aideal_i\ne0;\ \aideal_{i+1}=0$ とする。$D^{(i)}(x)$ を、定理 XI \mbox{により}\mbox{存在する}、$\aideal_i$ の全多項式の最大公約元---多項式の意味での因子---とする。$D^{(i)}(x)$ は $A^{(i)}(x)$ の因子\mbox{であるから}、\mbox{それ自身} $x_i$ に関して正則であり、次数 $p_i>0$ をもつ。$x_{i+1}\ldots x_n$ を $\mathsf{P}(u)$ に添加する。\mbox{このとき} $D^{(i)}(x)$ は、$\mathsf{P}(u;x_{i+1}\ldots x_n)$ の代数拡大体において $p_i$ 個の一次因子に分解する。その一つを $x_i-\bar x_i$ \mbox{とする}。\mbox{すると} $x_i=\bar x_i$ において $\aideal_i$ のすべての多項式が零となり、したがって定理 XI により、この特殊化のもとで $\aideal_{i-1}$ の多項式は最大公約元 $D^{(i-1)}(x)$ をもつ。$D^{(i-1)}(x)$ は $A^{(i-1)}(x)$ の因子であるから、\mbox{それ自身}も $x_{i-1}$ \mbox{に関して}正則であり、次数 $p_{i-1}>0$ をもつ。特に恒等的に零ではなく、$\mathsf{P}(u,\bar x_i,x_{i+1},\ldots,x_n)$ の拡大体で $p_{i-1}$ 個の一次因子に分解する。その一つを $x_{i-1}-\bar x_{i-1}$ とすれば、これに $\aideal_{i-2}$ から得られる\mbox{最大公約元}が\mbox{対応する}。このように続けると、$P(u;x_{i+1}\ldots x_n)$ の代数拡大体に属する、$\aideal$ の有限個の共通零点に\mbox{達する}。

逆に、$\aideal_\lambda$ は $\aideal$ で割り切れるから、$\aideal$ の任意の零点は $\aideal_\lambda$ の零点でもある。特に、$x_1=\bar x_1,\ldots,\allowbreak\mbox{$x_i=\bar x_i$},\allowbreak\mbox{$x_{i+1}=x_{i+1}$},\allowbreak\ldots,x_n=x_n$ が $\aideal$ の零点であるならば、$\aideal_{i+1}=0$ が従い、$\aideal_i$ の多項式は、一次因子 $x_i-\bar x_i$ を\mbox{もつ最大公約元}をもつ。以上をまとめると、次を得る。

\medskip
\noindent\textbf{定理 XII。} \emph{$\aideal$ は単位イデアルではなく、$\aideal_1\ne0\ldots;\ \aideal_i\ne0;\ \aideal_{i+1}=0$ とする。このとき $x_{i+1}\ldots x_n$ を $\mathsf{P}(u)$ に添加すれば、$P(u,x_{i+1},\ldots,x_n)$ の代数拡大体に属する、有限個の互いに対応する値系 $\bar x_1\ldots\bar x_i$ が\mbox{存在し}、$x_1=\bar x_1,\ldots,x_i=\bar x_i,x_{i+1}=x_{i+1},\ldots,x_n=x_n$ とおくと、$\aideal$ のすべての多項式が零となる。\mbox{逆に}、\mbox{このような} $\aideal$ の零点が存在するならば、$\aideal_{i+1}=0$ が従い、$\aideal_i$ のすべての多項式の共通因子 $D^{(i)}(x)$ で、\mbox{一次因子} $x_i-\bar x_i$ をもつものが現れる。}

定理 XII は特に、零点をもたない任意のイデアル $\aideal$ は単位イデアルであることを示す。

\endgroup
%
\begingroup
\renewcommand{\thefootnote}{\arabic{footnote})}
\setcounter{footnote}{17}

値系どうしの対応も明示する分解を得るため、新しい不定元 $v$ を導入し、これらを $\mathsf{P}(u,x_{i+1},\ldots,x_n)$ に添加する。次のようにおく：
\begin{equation}
 x_i=-v_1x_1-\cdots-v_{i-1}x_{i-1}+z_i.
\tag{35}
\end{equation}
すると、(12) と (35) を合成したものは再び (12) 型の置換となる。ただし、$u_{i\kappa}$ は $w_{i\kappa}=u_{i\kappa}-v_\kappa$ に\mbox{置き換えられ}、より一般に $u_{i+\lambda,\kappa}$ は $w_{i+\lambda,\kappa}=u_{i+\lambda,\kappa}-u_{i+\lambda,i}v_\kappa$ に\mbox{置き換えられる}。そこで、\mbox{置換} (35) が、\mbox{仮定により}変換イデアルである $\aideal$ を $\bideal$ へ移すとする。このとき $\bideal$ は、$\aideal$ の $u_{\mu\nu}$ を $w_{\mu\nu}$ で\mbox{置き換え}、$v$ を\mbox{添加する}だけで得られる。同じ過程により、$\bideal$ から定義されるイデアル $\bideal_1,\bideal_2,\ldots$ は $\aideal_1,\aideal_2,\ldots$ \mbox{から得られる}。逆に $v=0$ とすれば、$\bideal_i$ から $\aideal_\lambda$ が\mbox{得られる}。したがって $\aideal_\lambda$ と $\bideal_\lambda$ は\mbox{同時に}\mbox{零イデアルとなるか}、\mbox{同時に}\mbox{零イデアル}でなくなる。

再び $\aideal_1\ne0;\ldots;\aideal_i\ne0;\aideal_{i+1}=0$ とする。したがって $\bideal_1\ne0;\ldots;\bideal_i\ne0;\bideal_{i+1}=0$ でもある。$\bar x_1\ldots\bar x_i,x_{i+1}\ldots x_n$ を $\aideal$ の互いに対応する一組の零点系とする。$\bar z_i=\bar x_i+v_1\bar x_1+\cdots+v_{i-1}\bar x_{i-1}$ と\mbox{定義すれば}、(35) により、$\bar x_1\ldots\bar x_{i-1},\bar z_i,x_{i+1}\ldots x_n$ は $\bideal$ の零点となる。そこで $H^{(i)}(z,x)$ を $\bideal_i$ の全多項式の最大公約元とし、$v$ に関して整かつ原始的であると仮定する。定理 XII の最後の部分により、$H^{(i)}(z,x)$ は因子
\[
 z_i-\bar z_i=z_i-(\bar x_i+v_1\bar x_1+\cdots+v_{i-1}\bar x_{i-1})
\]
をもつ。$x_{i+1}\ldots x_n$ を添加した後に現れる $\aideal$ の各零点には、このような因子が対応する。

これで $H^{(i)}(z,x)$ の分解が尽くされること、すなわち、$z$ \emph{および} $v$ の一次因子に\mbox{分解しない}\mbox{因子} $H^{(i)}_1$ を $H^{(i)}$ から分離できないことを示さなければならない。反対に、そのような因子を\mbox{分離できる}と\mbox{仮定する}。$v$ に関する原始性の仮定により、$H^{(i)}$ は少なくとも一つの一次因子 $z_i-\bar{\bar z}_i$ を\mbox{含み}、\mbox{ここで} $\bar{\bar z}_i$ は $\mathsf{P}(u,v,x_{i+1}\ldots x_n)$ の代数拡大体に属する。$\bar{\bar x}_1\ldots\bar{\bar x}_{i-1},\bar{\bar x}_i,x_{i+1}\ldots x_n$ を、\mbox{これに対応する} $\aideal$ の\mbox{零点とする}。この零点は、$x_{i+1}\ldots x_n$ を添加した後に現れる $\aideal$ の有限個の零点の一つと\mbox{一致しなければならず}、\mbox{したがって} $v$ に依存しない。それゆえ $\bar{\bar z}_i$ は必然的に $v$ を線型に\mbox{含み}、\mbox{代数的}または\mbox{有理的}な\mbox{非線型}\allowbreak{}の\mbox{仕方では}\allowbreak\mbox{含まない}。\par\noindent\mbox{したがって}、仮定に反して、\allowbreak{}$H^{(i)}_1$ から $z$ と $v$ の\allowbreak\mbox{もう一つ}\allowbreak{}の一次因子 $z_i-(\bar x_i+v_1\bar x_1+\cdots+v_{i-1}\bar x_{i-1})$ を\mbox{分離できる}。よって、明示的な分解は
\[
 H^{(i)}(z,x)=\prod \left\{z_i-(\bar x_{i\lambda}+v_1\bar x_{1\lambda}+\cdots+v_{i-1}\bar x_{i-1,\lambda})\right\}^{\alpha_\lambda}
\]
となる。

$H^{(i)}$ の次数と各指数 $\alpha_\lambda$ は、$D^{(i)}(x)$ の対応する次数および指数と一致しなければならない。\mbox{なぜなら}、$H^{(i)}$ は特殊化 $u_{\mu\nu}=w_{\mu\nu}$ によって $D^{(i)}$ から生じ、$D^{(i)}$ は特殊化 $v=0$ によって $H^{(i)}$ から\mbox{生じるからである}。この分解はさらに、$u_{\mu\nu}$ を $\mathsf{P}$ に添加したため、$\aideal$ から始める消去に現れる $D^{(i)}$ の各零点 $\bar x_i$ には、$\aideal_{i-1}\ldots\aideal_1$ から得られる線型の最大公約元、またはその冪が対応することを示す。したがって、それぞれの場合に得られる $\aideal$ の互いに対応する零点系 $\bar x_1\ldots\bar x_i,x_{i+1}\ldots x_n$ はただ一つである\footnote{ここで与えた逐次消去は、Kronecker の方法とは異なり、与えられたイデアル $\aideal$ だけに依存し、どのイデアル基底にも依存しないことを指摘しておく。特に、指数 $\alpha_\lambda$ についても同じである。Hentzelt によれば、この方法で消去を完全に実行するには、商 $\aideal_i/D^{(i)}$ に対して手続きを続ける必要があるが、次節を考慮すれば、これは省いてよい。（E. N.）}。

\endgroup
%

\section*{§ 7. 終結式形式と消去理論。}

終結式形式と消去理論との関係を確立するため、§~6 で述べた逐次消去を、特にイデアルを\mbox{第 $(i-1)$ 段基本イデアル}で割った商 $\aideal=\mideal/\gideal_{i-1}$ に適用する。\par\noindent まず、この商が\mbox{変換イデアル}\allowbreak\mbox{である}ことを\mbox{示さなければならない}。\par\noindent $\mideal$ は変換イデアルであり、したがって §~3 の定理 VI により $\gideal_{i-1}$ も\mbox{変換イデアル}である。
\[
 H(x)\equiv0(\mideal/\gideal_{i-1});\qquad
 H(x)=\overline H(y)=\sum\alpha_\lambda(u)\overline H_\lambda(y)=\sum\alpha_\lambda(u)H_\lambda(x)
\]
から、
\[
 \overline H(y)\cdot\overline{\gideal}_{i-1,(u)}\equiv0\pmod{\overline{\mideal}_{(u)}};
 \quad\hbox{したがってまた}\quad
 \overline H(y)\overline{\gideal}_{i-1}\equiv0\pmod{\overline{\mideal}_{(u)}}
\]
が従う。すなわち、
\[
 \overline H_\lambda(y)\overline{\gideal}_{i-1}\equiv0\pmod{\overline\mideal},
 \quad\hbox{したがって}\quad H_\lambda(x)\gideal_{i-1}\equiv0\pmod{\mideal}
\]
である。これで $H_\lambda(x)$ が $\mideal/\gideal_{i-1}$ で割り切れることが証明され、したがって、このイデアルが\mbox{変換イデアルである}ことも証明された。ゆえに §~6 の消去法を適用できる。

一方、(28)---§~4---を考慮すれば、(30) は $C^{(i)}(x)$ が $\mideal/\gideal_{i-1}$ で割り切れることを示す。ここで $C^{(i)}(x)$ は、$\Mmod^*_{i-1}$ から得られる恒等的に零でない任意の $\rho$ 次行列式である。\par\noindent $i>1$ のとき $C^{(i)}$ は \mbox{$x_1$ に関して零次である}。\par\noindent それゆえ、イデアル $\aideal=\aideal_1$ に対応する加群 $\mathfrak A_1$ は、多項式 $C^{(i)},x_1C^{(i)},\ldots,x_1^{k_1-1}C^{(i)}$ を\mbox{含み}、\mbox{したがって} $(C^{(i)})^{k_1}$ は $\aideal_2$ で\mbox{割り切れる}。同様に $(C^{(i)})^{k_1k_2}$ は $\aideal_3$ で\mbox{割り切れ}、最後に $(C^{(i)})^{k_1k_2\cdots k_{i-1}}$ は $\aideal_i$ で\mbox{割り切れる}。したがって $\aideal_1,\ldots,\aideal_i$ は\mbox{零イデアル}ではない。このことは $i=1$ についても成り立つ。

いま $D^{(i)}(x)$ を $\aideal_i$ の全多項式の最大公約元とする。$D^{(i)}(x)$ が次数 $p_i>0$ をもつのは、$\aideal_{i+1}$ が\mbox{零イデアル}となる場合に限る。$D^{(i)}(x)$ の定義と、$\aideal_i$ が $\aideal$ で割り切れることから、
\[
 d^{(i+1)}D^{(i)}(x)\equiv0(\mideal/\gideal_{i-1});\qquad d^{(i+1)}\ne0
\]
を得る。したがって $d^{(i+1)}D^{(i)}$ は、$x_1\ldots x_{i-1}$ を含まない、$\mideal/\gideal_{i-1}$ に属する多項式である。一方、$E^{(i)}(x)$ は、$\Mmod_{i-1}$ または $\Mmod^*_{i-1}$ の最高単因子として、定理 II と §~4 により、$\mideal/\gideal_{i-1}$ に属し $x_1\ldots x_{i-1}$ を\mbox{含まない}\allowbreak\mbox{全多項式}の、多項式の意味での最大公約元と定義されていた。したがって $d^{(i+1)}D^{(i)}$ は $E^{(i)}(x)$ で\mbox{割り切れ}、さらに $E^{(i)}(x)$ が $x_i$ に関して正則であるため、$D^{(i)}(x)$ も $E^{(i)}(x)$ で割り切れる。

同じ議論は、変換 (12) が最初から (35) と合成されていた場合、すなわち不定元 $u_{\mu\nu}$ が $w_{\mu\nu}$ に\allowbreak\mbox{置き換えられて}\allowbreak{}いた場合にも行える。\allowbreak{}これにより $H^{(i)}(z,x)$ は、対応する $\bar E^{(i)}(z,x)$ で割り切れる。$D^{(i)}(x)$ と $H^{(i)}(z,x)$ が特殊化によって相互に得られるのと同様に、$E^{(i)}(x)$ と $\bar E^{(i)}(z,x)$、また $R^{(i)}(x)$ と $\bar R^{(i)}(z,x)$ も特殊化によって相互に得られる。したがって、一次因子への分解に現れる重複度も相互に\mbox{一致しなければならない}。さらに $E^{(i)}(x)$ のある冪が $R^{(i)}(x)$ で割り切れることを考慮すると、以上を次のように\mbox{まとめられる}。
%

\medskip
\noindent\textbf{定理 XIII。} \emph{$R^{(i)}(x)$ を $\mathsf{P}(u,x_{i+1}\ldots x_n)$ の代数拡大体上で一次因子 $x_i-\bar x_i$ に分解するとき、$\bar x_i$ は\mbox{ただ一通りに}、互いに対応する一組の値系 $\bar x_1\ldots\bar x_i,x_{i+1}\ldots x_n$ へ補うことができ、この値系は $\mideal/\gideal_{i-1}$ の零点、したがって $\mideal$ の零点となる。このように $R^{(i)}$ の一次因子から生じる $\mideal$ の有限個の零点---$x_{i+1}\ldots x_n$ を\mbox{添加した後}には有限個となる---は、次の明示的な分解によって一つにまとめられる：}
\begin{equation}
 \bar R^{(i)}(z,x)=\prod\left\{z_i-(\bar x_{i\nu}+v_1\bar x_{1\nu}+\cdots+v_{i-1}\bar x_{i-1,\nu})\right\}^{\lambda_\nu}.
\tag{36}
\end{equation}
\emph{$\bar R^{(i)}(z,x)$ は $z_i$ に関して正則であるから、この分解は、$\mathsf{P}(u)$ に添加された $x_{i+1}\ldots x_n$ を、$P(u)$ または\mbox{そこから}導かれる代数的閉体に属する任意の特殊な値系 $\bar x_{i+1}\ldots\bar x_n$ で置き換えても保たれる。}

これにより同時に、終結式形式の個々の因子に従って $\mideal$ の零点を分離することも達成された。$\mathsf{P}(u)$ に添加された不定元 $x$ の個数は、対応する代数的図形の次元とも呼ばれる。
%
\begingroup
\renewcommand{\thefootnote}{\arabic{footnote})}
\setcounter{footnote}{18}

$\bar R^{(i)}(z,x)$ の分解、またはそれに対応する $R^{(i)}(x)$ の分解において、$\mathsf{P}(u,x_{i+1}\ldots x_n)$ に関して共役な因子を一つにまとめると、一般の $\mathsf{P}$ ではそれらの全部または一部が同一となりうるが、$R^{(i)}(x)$ は $\mathsf{P}(u,x_{i+1}\ldots x_n)$ に関して既約な多項式の冪へと分解される：
\[
 R^{(i)}(x)=S^{(i)}_1(x)^{\beta_1}\cdots S^{(i)}_\varkappa(x)^{\beta_\varkappa}.
\]
定理 X により、この分解には $\gideal_i=[\fraks_1\ldots\fraks_\varkappa]$ という表示が対応し、$S^{(i)}_\mu(x)^{\beta_\mu}$ は $\gideal_{i-1}$ のイデアル $\fraks_\mu$ \mbox{に関するノルム}、または $\Gmod_{i-1}$ の $\fraks_\mu$ に対応する加群に関するノルムに等しくなる。これにより指数 $\beta_\mu$ の意味も明らかとなる。とりわけ、$\gamma_\mu$ を $S^{(i)}_\mu$ の次数とすれば、重複度 $\beta_\mu\gamma_\mu$ は、$\mathsf{P}(u,x_{i+1}\ldots x_n)$ \mbox{に関して}\mbox{線型独立な}、$\gideal_{i-1}$ の $\fraks_\mu$ を法とする剰余類の個数に等しい。

最後に、$\mathsf{P}$ が標数零である特殊な場合を簡単に考える。$u_{\mu\nu}$ を $P$ に属する量 $\bar u_{\mu\nu}$ へ特殊化し、\mbox{$n$ 個}の\mbox{行列式} $C^{(i)}(x)$（$(i=1,2,\ldots,n)$）がそれぞれ $x_i$ に関して正則であり続けるようにすれば\footnote{Hentzelt は初めから $u_{\mu\nu}$ にこの特殊化を施している。そのため、$H^{(i)}(z,x)$ および $\bar R^{(i)}(z.x)$ が $z$ と $v$ の一次因子へ分解できることを\mbox{証明する}にも、はるかに複雑な考察を用いなければならない。その際に現れる指数は、上で得た指数と必ずしも一致しない。（E. N.）}、$R^{(i)}$ の\mbox{正則性}も\mbox{保たれる}。この特殊化により、$R^{(i)}$ と $R^{(i)}(z,x)$ はいずれも $\Mmod^*_{i-1}$ から得られる全ての \mbox{$\rho$ 次行列式}の\mbox{共通約元}となるが、必ずしも最大公約元とはならない。しかし、この後者の性質も\mbox{保たれる}ように\mbox{特殊化}を\mbox{施すこと}が常にできる。実際、加群基底の存在が示すように、$\bar R^{(i)}(z,x)$ は $\Mmod^*_{i-1}$ から\mbox{得られる}\mbox{有限個}の $\rho$ 次\mbox{行列式}の最大公約元としても特徴づけられる。そのように特殊化された $\bar R^{(i)}(z,x)$ の分解は、(36) において $u_{\mu\nu}$ を $\bar u_{\mu\nu}$ で置き換えるだけで得られる。したがって、消去理論全体が保たれる。

\begin{center}
（1922 年 3 月 17 日受理。）
\end{center}

\endgroup

%
\section*{23. 代数的不変式と微分不変式}
\begin{center}
Jahresber. d. Deutschen Mathematiker-Vereinigung 32 (1923), S. 177--184
\end{center}

\begin{center}
{\Large\bfseries 代数的不変式と微分不変式。}\par
\vspace{0.45em}
1922年9月18日、ライプツィヒでの講演に基づく報告。\par
\vspace{1.0em}
ゲッティンゲンの \textsc{Emmy Noether} による。
\end{center}

\begingroup
\renewcommand{\thefootnote}{\arabic{footnote})}
\setcounter{footnote}{0}

代数的不変式および微分不変式の発展について報告するにあたり、私は両分野とも、\textsc{Hilbert} の言葉によれば素朴な時期と形式的な時期に続く批判的時期に話を限りたい。この批判的時期を特徴づける名は、代数的不変式については \textsc{Hilbert} 自身、微分不変式については \textsc{Riemann} である。内容の面からいえば、代数的不変式については、これらの問題をめぐって本質的に発達し、ここでその全き鋭さを発揮しえた代数学の算術的方法が、微分不変式については、形式的変分法の方法がそれを特徴づける。そこで私は、個々の対象を越える意義をもつ算術的方法について報告し、第二部では、まさに変分法の方法に即して達成される微分不変式の代数的不変式への従属だけを扱うことにする。

%
\textbf{1.} \emph{代数的不変式}の基礎には、周知のように、$x_1,\ldots,x_n$ における\emph{一般線型群}がある：
\srcnumdisplay{(1)}{x_i=\sum_k s_{ik}x'_k
\qquad\text{または略して}\qquad x=S(x').}
ここで $s_{ik}$ は不定元を表し、したがって $\D=|s_{ik}|$ は零でない。いま $f(a,x),g(b,x),\ldots$ を $x$ の形式で、それぞれ次数 $\alpha,\beta,\ldots$ をもち、不定元 $a,b,\ldots$ を係数とするものとする。このとき (1) によって $a,b,\ldots$ の\emph{誘導変換}が生じる。実際、
\[
\begin{aligned}
f(a,x)&=\sum a_{i_1\ldots i_n}x_1^{i_1}\cdots x_n^{i_n}
      =\sum a'_{j_1\ldots i_n}{x'_1}^{j_1}\cdots {x'_n}^{j_n}
      =f(a',x'),\\
g(b,x)&=\sum b_{r_1\ldots r_n}x_1^{r_1}\cdots x_n^{r_n}
      =\sum b'_{s_1\ldots s_n}{x'_1}^{s_1}\cdots {x'_n}^{s_n}
      =g(b',x')
\end{aligned}
\]
から、$a',b',\ldots$ は $a,b,\ldots$ の線型斉次関数となり、その係数は $s_{ik}$ に関してそれぞれ $\alpha,\beta,\ldots$ 次となる：
\srcnumdisplay{(2)}{a'=A_\alpha(a);\qquad b'=B_\beta(b)\ldots .}
ここで\emph{不変式}とは、変換 (1) および (2) に対する不変式をいう。すなわち、不定元 $a,b,\ldots,x,y,\ldots$（ここで $y=S(y')$ とも変換される）の整有理関数であって、変換を施すと置換行列式の冪を因子として自らを再現するものである：
\srcnumdisplay{(3)}{I(a',b',\ldots,x',y',\ldots)
=\D^\rho I(a,b,\ldots,x,y,\ldots).}
$I$ の係数は有理数、さらには有理整数であると仮定してよい。任意の数体 $P$ の係数をもつ不変式は、$P$ の係数を用いて、そのような特殊な不変式有限個によって線型的に表せるからである。同様に、$I$ が各系列について別々に斉次であると仮定しても制限にはならない。ここでも任意の不変式は、そのような特殊な不変式有限個の和として組み立てられるからである。

逆に、(1) と (2) は条件 (3) によって再び決定されることに注意しておく。\textsc{Ostrowski} と \textsc{Schur} が最近示したように、\footnote{A. \textsc{Ostrowski} と I. \textsc{Schur}, Über eine fundamentale Eigenschaft der Invarianten einer binären Form, Math. Zeitschr. 15 (1922), およびこれに続く \textsc{Ostrowski} の未刊行の研究。} 誘導変換は、各系列ごとに分離して働く線型変換の全体、すなわち、明示できる一定の例外を除いて、$x,y,\ldots$ を含まない任意の不変式 $I$ を許す変換の全体としても特徴づけられる。

二つの不変式の積は再び不変式であるが、和が不変式となるのは、重み、すなわち (3) の指数 $\rho$ が一致するとき、かつそのときに限られる。しかし行列式一の変換に制限すれば、任意の和も再び不変式となり、そのように制限された群に対するすべての不変式を尽くす。したがって整域、しかも\emph{斉次整域}が得られる。これは、ある多項式がその領域に属すれば、その斉次成分も一つずつ属するような領域であり、ここでは再び (3) に従う各系列ごとに斉次な不変式に帰着する。

ここで、不変式論、ひいては算術的代数学全体がその周囲で発展した根本問題が生じる：\emph{この整域は有限生成か、すなわち、任意の不変式を $I_1,\ldots,I_k$ によって整有理的に $I=G(I_1,\ldots,I_k)$ と表せるような有限の整基底 $I_1,\ldots,I_k$ をもつか。} この問題に対する \textsc{Hilbert} の解法は、整基底の問題をイデアル基底の問題へ還元することに基づく。二元の場合に対する \textsc{Gordan} の証明は、見通しのきかない記号計算のため拡張を許さず、\textsc{Mertens-Hilbert} の証明も二元形式を線型形式へ分解するため拡張を許さない。ここでは算術的な根本思想を強調しつつその論旨を手短に示し、次いで関連する三つの問題圏、\mbox{すなわち有限個}の手順による基底の実際の構成、部分群における有限性問題、および整数係数性を考慮した有限性問題へ進みたい。

%
\textbf{2.} （有限次）代数体における\emph{イデアルの定義}を思い起こそう。体の数からなる系 $\ma$ は、次を満たすときイデアルと呼ばれる。
\PXXIIIitem{1}{$\ma$ が、その体のすべての整数からなる領域 $\mo$ に属すること、}
\PXXIIIitem{2}{$\alpha$ と $\beta$ が属すれば、差 $\alpha-\beta$ も $\ma$ に属すること、}
\PXXIIIitem{3}{$\alpha$ が属すれば $\lambda\alpha$ も $\ma$ に属すること。ただし $\lambda$ は $\mo$ の任意の元とする。}
そして周知の定理によれば、すべてのイデアルは\emph{イデアル基底}をもち、$\ma=(\alpha_1,\ldots,\alpha_k)$ と書ける。すなわち、$\ma$ の元 $\alpha_1,\ldots,\alpha_k$ が有限個存在し、2. と 3. によって $\ma$ をそれらから導けるので、$\alpha=\lambda_1\alpha_1+\cdots+\lambda_k\alpha_k$ が $\ma$ の全要素を尽くす。\footnote{$k$ を二まで減らせることは、以下では問題にならない。}

このイデアルの定義は、$\mo$ が整域をなすということだけに基づく。したがってイデアルの概念もイデアル基底の概念も、$\mo$ を任意の整域に置き換えても文字どおり保たれる。

さて \textsc{Hilbert} の有限性証明は、次の三定理に分かれる。
\PXXIIIitem{1}{有理性領域を係数とする $n$ 個の不定元の全多項式の領域では、すべてのイデアルについてイデアル基底定理が成り立ち、したがって任意の系 $\mS$ についても成り立つ。}
\PXXIIIitem{2}{多項式からなる斉次整域 $\mJ$ が有限生成であるための必要十分条件は、$\mJ$ においてイデアル基底定理が成り立つことである。}
\PXXIIIitem{3}{不変式領域 $\mJ$ では、イデアル基底定理は次の強化形で成り立つ。全多項式の領域に関して作った任意のイデアル基底が同時に $\mJ$ のイデアル基底でもあり、したがって}
\[
 I=A_1I_1+\cdots+A_kI_k
 \qquad\text{とともに常に}\qquad
 I=i_1I_1+\cdots+i_kI_k
\]
が成り立つ。ここで $A$ は $a,b,\ldots,x,y,\ldots$ の多項式を、$i$ は不変式を表す。

1. の証明は、原理上、代数体の場合と同じ推論に基づく。どちらの場合も、イデアル基底の存在は、線型形式からなる加群についての一般定理の直接の帰結である。\footnote{これについては、代数関数の算術理論に関する私の比較報告、Jahresber. 28 (1920), S. 187 および Anm. 2), S. 188 を参照。} 1. から、多項式からなる任意の有限生成整域についてイデアル基底の存在が直ちに従い、斉次領域について逆も成り立つことは容易に分かる。特殊な不変式の性質が使われるのは 3. だけである。\textsc{Hilbert} では 3. は $\Omega$ 過程に関する周知の定理から従う。一方、近年 E. \textsc{Fischer} は、別の一般定理を基礎にすれば、一般線型群が各変換とともにその反傾変換、またはその共役複素変換を含むという性質だけから 3. を導けることを示した。\footnote{E. \textsc{Fischer}, Über die Endlichkeit der Invarianten. Göttinger Nachrichten 1915, およびライプツィヒ講演。}

%
\textbf{3.} \emph{基底の実際の構成}は、\emph{関数体}の理論を援用して有限性概念をさらに算術的に変形することに基づく。次が成り立つ。
\PXXIIIitem{4}{多項式からなる整域 $\mJ$ が有限生成であるための必要十分条件は、$\mJ$ の中に有限生成部分整域 $\mJ'$ があり、$\mJ$ が $\mJ'$ 上整であることである。このとき $\mJ'$ の任意の整基底に、これら整な量の基本系を加えれば $\mJ$ の基底となる。\footnote{\textsc{Hilbert} は（Über die vollen Invariantensysteme, Math. Annalen 42 (1893), § 1 および 2）、仮定した有限生成性からこれらの事実が従うことを示した。逆に、整依存性から有限生成性が従うことは、有理基底の存在から分かる（E. Noether, Körper und Systeme rationaler Funktionen, Math. Annalen 76 (1915), § 12）。}}
とくに不変式の場合のように、イデアル基底定理が強化形 3. で成り立つ斉次領域 $\mJ$ を考えるなら、そのような領域 $\mJ'$ は、常に存在する何らかの有限個の多項式 $I_1,\ldots,I_s$ を整基底として導かれるように選べる。しかも、それらが消えれば $\mJ$ の全多項式も消える。この事実は再びイデアル論の一般定理に基づく。
\PXXIIIitem{5}{多項式からなる任意のイデアル $\ma$ に対し、通常の有理整数 $r$ が存在して、$\ma$ のすべての零点で消える任意のイデアル $\mb$ について $\mb^r$ は $\ma$ で割り切れる。}

不変式の場合には、$I_1,\ldots,I_s$ の\emph{次数の上界}、したがってその個数 $s$ の上界も、基礎に置いた基本形式の次数だけに依存する形で決定できる。ここでも不変式の特殊な性質が使われる。すなわち、数値的に特殊化した基本形式の系が零でない不変式をもつための必要十分条件は、$\D=|s_{ik}$ が変換後の係数上整であることである。

この上界から K. \textsc{Hentzelt}\footnote{私はこの研究を近く遺稿から刊行する予定である。} は、\emph{整基底に属する不変式の次数の上界}を求めた。この上界も基本形式の次数だけに依存する。すなわち彼は、5. の\emph{指数} $r$ について、$\ma$ のイデアル基底に属する多項式の個数と最高次数だけに依存し、それらの係数には依存せず、しかも与えられた数から有限個の手順で計算できる\emph{上界}を与えたのである。原理的には、これによって提起した問題は完全に解決された。ただし、その上界はきわめて高い。例えば三変数二次形式では、判別式、すなわち三次の不変式だけですでに $x$ を含まない全不変式の基底となるのに、\textsc{Hilbert} と \textsc{Hentzelt} の評価では数百万にまで達する。

%
\textbf{4.} 一般線型群の\emph{部分群}の不変式に関する有限性問題では、まだ部分的成果に達したにすぎない。\mbox{その方法}はほとんどもっぱら、\textsc{Study} が直交群について採った手続きに従い、部分群の不変式を一般線型群の同時不変式へ還元することに基づく。\textsc{Weitzenböck} は運動不変式とアフィン不変式について、\textsc{Deruyts} は半不変式と剪断不変式についてこれを実行した。\footnote{文献については、\textsc{Weitzenböck} の不変式論に関する百科事典項目 III E 1 を参照。} この方法がどこまで及ぶかはまだ未決である。2. の末尾で述べた \textsc{Fischer} の指摘も、一連の部分群について有限性問題を解く。とくに直交群については、これにより最も簡単な証明が得られる。しかし半不変式の例は、\textsc{Fischer} が近年指摘したように、\footnote{ライプツィヒ講演および Math. Zeitschrift。} 部分群では、有限生成性にもかかわらず、イデアル基底に関する強化形 3. が成り立つ必要はないことを示す。

この有限性問題を体論の意味で捉えると、\textsc{Hilbert} の\emph{相対的に整な関数}の問題、すなわち、有理関数体の多項式全体からなる整域が常に有限生成であるかという問題へ至る。この方向では、有限群の不変式に対し、際立った整基底を直接与えられるという、初等的に証明できる事実がある。それはガロア分解式の係数系である。\footnote{E. Noether, Der Endlichkeitssatz der Invarianten endlicher Gruppen. Math. Ann. 77 (1916).}

%
\textbf{5.} 不変式の有限性定理を\emph{整数係数性}に関して強化できるかという問題も、一般にはまだ解決されていない。確かに \textsc{Hilbert} が示したように、イデアル基底定理は全整数係数多項式の領域についても成り立ち、イデアル基底と整基底との関係 2. もそのまま保たれる。\footnote{この関係から、とくに、私が Math. Ann. 83 (1921) で展開した一般イデアル論の分解定理が、すべての有限生成整域について成り立つことが従う。} しかし、有限性には必要でない強化形 3. の証明は通用しない。\emph{二元不変式}の場合には、\emph{整数係数での有限生成性}を証明できる。\footnote{E. Noether, Die Endlichkeit des Systems der ganzzahligen Invarianten binärer Formen. Göttinger Nachrichten 1919.} その証明は、二元の場合の \textsc{Mertens-Hilbert} の有限性証明を整数係数性に関して強化したものだが、整数係数イデアル基底の定理だけでなく、二元形式を線型形式へ分解することも用いるため、より多くの不定元へ移すことはできない。これに対し、同様の基礎の上で、先に述べた証明を強化すれば、\emph{整数係数での有限生成性}が\emph{有限群の不変式}について得られる。ただしここでも、基底を実際に与えるのではなく存在証明だけが得られる。\footnote{11) で引用した論文の § 4。} ここでも、関数体論と一般イデアル論のさらなる発展だけがおそらく目標へ導くであろう。

%
\textbf{6.} ここで\emph{微分不変式}へ移り、序論で提起した、それを代数的不変式論へ還元する問題に進む。基礎となる群は、周知のように $x_i=x_i(y)$ によって次のように定義される：
\srcnumdisplay{(4)}{d x_i=\sum_k \frac{\partial x_i}{\partial y_k}d y_k;\qquad
 d\delta x_i=\sum_k \frac{\partial x_i}{\partial y_k}d\delta y_k+
\sum_{k,l}\frac{\partial^2 x_i}{\partial y_k\partial y_l}d y_k\,\delta y_l;\ \ldots .}
(4) によって誘導される変換へ移るには、任意の関数 $f(x,d x)=g(y,d y)$ を基礎に置けばよい。(4) に対する $\delta f,\delta^2f,\ldots$ の不変性を要求すると、$x$ と $d x$ による $f$ の導関数についての変換が得られる。例えば $d x$ に関して斉次な形式に限れば、
\[
 f(x,d x)=\sum a_{i_1\ldots i_n}(x)d x_1^{i_1}\cdots d x_n^{i_n}
 =\sum a'_{j_1\ldots j_n}(y)d y_1^{j_1}\cdots d y_n^{j_n}
 =g(y,d y),
\]
$a'$ は $a$ に関して線型となり、$\frac{\partial a'}{\partial y}$ は $a$ と $\frac{\partial a}{\partial x}$ に関して線型となる、等々：
\srcnumdisplay{(5)}{a'=A_0(a),\qquad
\frac{\partial a'}{\partial y}=A_1\!\left(a,\frac{\partial a}{\partial x}\right),\qquad
\frac{\partial^2 a'}{\partial y^2}=A_2\!\left(a,\frac{\partial a}{\partial x},\frac{\partial^2 a}{\partial x^2}\right),\ldots}
そして問題は、変換 (4) と (5) に対する不変性である。このとき現れるすべての量
\[
a,\ \frac{\partial a}{\partial x},\ldots,\ d x,\ \delta x,\ d\delta x,\ldots;\ 
\frac{\partial x}{\partial y},\ \frac{\partial^2 x}{\partial y^2},\ldots
\]
は不定元とみなすべきであるから、各変換の行列式は確かに零でない。形式的に完成した理論を特殊な関数へ適用するときに初めて、ある領域で変換が一意に可逆なままであり、十分な数の微分商が存在するような関数へ制限する問題が生じる。これは、代数の場合、数値的特殊化に際して行列式 $|s_{ik}|$ が零でないと仮定しなければならないのとまったく同じである。

このように、現れるすべての引数を不定元とみなすとき、通常の代数的方法では直接扱えない特殊な線型群を考えることになる。この群では $d x$ の変換と $a'=A_0(a)$ の変換が (1) およびそこから誘導された (2) と一致する。そこで、一般に\emph{新しい引数を導入する}ことによって、変換 (4) と (5) を一般線型群 (1) と、それによって\emph{誘導される変換} (2) へ\emph{還元できる}か、という問題が生じる。

これは実際に可能である。高次微分 $d^2x,\ d\delta x\ldots$ は、$f(x,d x)$ に属する変分問題から生じるラグランジ\nolinebreak[4]ュの変分導関数と、さらにそれらを極化して得られる結合（\textsc{Weyl} と \textsc{Schouten} のいう接続）で置き換えられる。これらはすべて $d x$ と同じ変換則に従い、その変換は一般線型群 (1) によって与えられる。一方、$\frac{\partial a}{\partial x},\frac{\partial^2a}{\partial x^2},\ldots$（より一般には、斉次と仮定した $f$ の $x$ と $d x$ による導関数）は、「高次変分の正規形」から、上記の $d^2x,\ d\delta x\ldots$ の消去によって生じる一定の形式の係数と、それらの「共変微分」
\[
 [\Omega_i],\quad [\Omega_i^{(1)}],\quad [\Omega_i^{(2)}],\ldots\qquad (i=1,2,\ldots);
\]
で置き換えられる。そして $[\Omega_i^{(k)}]$ は (4) と (5) に対する不変式で $d x$ と $\delta x$ だけに依存するから、その係数は代数的に誘導される変換 (2) を満たす。\footnote{E. Noether, Invarianten beliebiger Differentialausdrücke. Göttinger Nachrichten. 1918.} この\emph{還元定理}により、微分不変式論は代数的不変式論に従属する。

$f(x,d x)$ が $d x$ に関して $\alpha$ 次の斉次形式であるとき、$[\Omega_i]$ の列は $[\Omega_\alpha]$ で終わる。二次形式では $[\Omega_2]$ は \textsc{Riemann} の曲率形式と一致し、ここで示した構成過程は、\textsc{Riemann} がラテン語の懸賞論文で与え、またそれとは独立に、\textsc{Riemann} の教授資格講演を受けて \textsc{Lipschitz} が展開した過程とまったく同じである。\footnote{7) で挙げた \textsc{Weitzenböck} の百科事典項目の第二部第28項にある、\textsc{Riemann} と \textsc{Lipschitz} に関する私の叙述も参照。} \textsc{Christoffel} と \textsc{Ricci} が二次微分形式について計算で示した還元定理の証明は、「\textsc{Riemann} 正規座標」の導入に基づく。この座標は一点から出る極値曲線を直線に変換する。これによって $x$ の任意の変換に正規座標の線型変換が対応するので、正規座標が一般線型群への移行をもたらす。

\textsc{Weyl} と \textsc{Schouten} に従い、群の定義に変分問題ではなく\emph{接続を基礎に置く}場合にも、このような還元定理があるかという問題が生じる。すなわち、ラグランジュの変分導関数に対応して、高次微分を消去する式 $\psi_i(d\delta x,d x,\delta x)$ を立て、二次微分形式との類比により、それ自体は不要な制限、すなわち $\psi_i$ が $d x$ に関して線型であるという制限を加える場合である（任意の斉次 $f(x,d x)$ では、極化したラグランジュの変分導関数は $d x$ に関して一次斉次にすぎない）。$\psi$ が $d x$ と同じ変換則に従うことを要求すると、その係数の変換が誘導され、それによって (5) と同様に群が定まる。最低階の不変式については、\textsc{Weitzenböck} が \textsc{Weyl} の場合に還元定理を計算で確認したが、一般的な構成法はまだない。\footnote{次を参照：\textsc{Weyl}, Raum, Zeit, Materie; \textsc{Schouten}, Math. Zeitschr. 13 (1922); \textsc{Weitzenböck}, Sitzungsber. d. Wiener Akademie, 1920 u. 1921.}

%
\begin{center}
（1922年10月26日受理。）
\end{center}
\endgroup

\clearpage
\setcounter{footnote}{0}

\section*{24. 消去理論と一般イデアル論}
\begin{center}
Math. Ann. 90 (1923), S. 229--261\\[0.9em]
{\Large\bfseries 消去理論と一般イデアル論。}\\[1.0em]
著者\\[0.45em]
ゲッティンゲンのエミー・ネーター。\\[0.6em]
\rule{4em}{0.4pt}
\end{center}

以下で問題となるのは、私が Hentzelt の叙述\footnote{Kurt Hentzelt, Zur Theorie der Polynomideale und Resultanten. E. Noether 編、Math. Ann. 88 (1922), pp. 53--79; 以下 H.-N. と引用する。}に与えた算術的な形、すなわち多項式領域におけるイデアル論と記述しうる形の消去理論を、一般イデアル論\footnote{E. Noether, Idealtheorie in Ringbereichen, Math. Ann. 83 (1921), pp. 24--66; 以下 Idealtheorie と引用する。}の中に位置づけることであり、同時に零点に関わる部分を新たに基礎づけることである。ここで参照できるよう、両理論の基本概念、および体論の基本概念\footnote{E. Steinitz, Algebraische Theorie der Körper, J. f. M. 137 (1910), pp. 167--309; 以下 Steinitz と引用する。}を § 1 で簡潔にまとめておく。

分類と新たな基礎づけは、四つの問題を中心に組織される。次元概念の算術的定式化、零点の理論、ノルムおよび単因子に関する分解定理と一般イデアル論の分解定理との関係、そしてノルムおよび単因子形式による素イデアルと準素イデアルの特徴づけである。

\emph{次元概念}（§ 4）の算術的定式化は、素イデアルの鎖によって与えられる。これは、曲面上には曲線があり、曲線上には点があるという事実の算術的対応物である。この算術的定式化は、消去理論における\emph{パラメータによる定義}と同一であることが証明されるが、わずかな修正を加えれば任意の環領域へ移すことができる。また他の諸問題の一部の移行と合わせて、そこではイデアルの構造を正確に見通す手段を与える。このことは別のところで示す予定である。

H.-N. では、\emph{零点}の問題は通常どおり、与えられたイデアルに対して直接に、すなわち逐次消去へ戻ることによって扱われるため、検算的性格を完全には避けられない。これと対照的なのが、一変数多項式で常に用いられる道筋である。すなわち、与えられた多項式を\emph{素関数（準素関数）の冪の積}として表し、各素関数に対して、その剰余体と\emph{同型な}零点体を構成する。素関数の次数はこの体の次数を与えるが、同じ零点をもつ準素関数の次数は、この準素関数を法とする剰余類環の次数を与える。各素関数についてガロア拡大体へ移れば、線型因子への分解が得られ、これにより最初に与えた多項式について、零点、線型因子分解、重複度の問題が解決される。これに完全に対応して、ここでは（§ 3）素イデアルを法とする剰余類系を、商の形成によって剰余体へ拡張し、それに同型な零点体を構成する。この零点体は、たとえば $x_{i+1},\ldots,x_n$ という不定元を添加して生成される部分体を含む。ノルムの次数はこの部分体に関する次数を与え、同時にガロア拡大体へ移ることによって、単因子形式、したがってノルムの線型因子への分解が得られる。対応する準素イデアルについては零点は同じであり、分解も同時に与えられる。なぜならノルムは素イデアルの単因子形式の冪となるからである（§ 2）。ノルムの次数は、不定元から得られる部分体に関して、準素イデアルを法とする剰余類環の次数を与える。

したがって任意のイデアルの零点の問題は、\emph{ノルムと単因子に関する分解定理を一般イデアル論の分解定理と結びつける問題}へ還元される（§ 5）。イデアルの零点が個々の随伴素イデアルの零点から構成されるという事実は、ノルムが\emph{最大準素因子}へ分解され、それらが個々の随伴素イデアルの単因子形式の冪に等しくなることに対応する。\footnote{消去理論と一般イデアル論のこの並行性は、Kronecker の消去理論では満たされない。H.-N., 注 4) を参照。} これによって一般の場合のノルムの線型因子分解が実現される。重複度は、H.-N. では基本イデアルとその成分によって解釈される。すなわち、ノルムの最大準素因子の次数は、$x_{i+1},\ldots,x_n$ を添加した後の、$(i-1)$ 段の基本イデアルの、$i$ 段の基本イデアルの一成分を法とする線型独立な剰余類の数に等しいとされる。ここでは、$(i-1)$ 段の基本イデアルが、$(n-i)$ 番目より大きい次元のすべての素イデアルによって一意に定まる分解の孤立成分と同一であることを示す。また $i$ 段の基本イデアルの一成分は、ノルム因子に対応する素イデアルと、より高次元のすべての素イデアルによって定まる分解の孤立成分と同一である。対応するノルム因子の次数は、$x_{i+1},\ldots,x_n$ を添加した後、この成分の剰余類のうち $(i-1)$ 段の基本イデアルの類だけからなる部分環の次数に等しくなる。

\emph{素イデアルおよび真の準素イデアルの特徴づけ}（§ 6）は、剰余体を考察することから直接に得られる。もとの係数領域が完全体ならば、素関数はノルムおよび単因子形式として素イデアルに対応し、真の準素関数は真の準素イデアルに対応し、逆も成り立つ。係数領域が不完全体である場合には、必要条件または十分条件しか与えられない。例が示すように、素イデアルのノルムが真に準素となることがあり、真の準素イデアルの単因子形式が素関数となることもある。より詳しく言えば、標数 $p$ では、素イデアルの\mbox{ノルム}はその単因子形式の $p^g$ 乗（$g\geqq0$）となる。他方、単因子形式が素関数である真の準素イデアルでは、\mbox{再び}例が示すように、ノルムは任意の冪になりうる。最後に\emph{絶対素イデアル}（§ 7）、すなわち係数領域を代数閉体へ拡張しても素イデアルであり続けるものを考察する。有限代数的数体からの係数をもつ絶対素イデアルについて、この特徴づけは、A. Ostrowski が絶対素関数について最初に述べた定理の移行を導く。すなわち、素イデアルである性質は、有限個の例外を除き、数体の任意の素イデアルを法として保たれる。

最後の問題と代数的数体におけるイデアル論との類似も指摘しておこう。そのようなイデアルの単因子形式としては、そのイデアルで割り切れる最小の有理整数を考えるべきである（注 10）。したがって素イデアルの場合、それは常に素数となる。逆に、ノルムが素数になれば直ちにイデアルは素イデアルとなる。証明もまったく並行して進む。したがって\emph{上の例}は、不完全体上では高次の素イデアルおよび分岐イデアルの類似物が現れるが、完全体上では一次の素イデアルだけが存在し、分岐イデアルは存在しないことを示している。

\subsection*{§ 1. 消去理論、一般イデアル論、および体論の基本概念}

\textbf{1. 領域。} $\bar{\mathfrak{J}}$ を、抽象的に定義された体 $P$ を係数とする $y_1,\ldots,y_n$ のすべての多項式からなる整域とする。$y$ を、不定元 $u_{\mu\nu}$ を係数とする次の変換に従わせる。\footnote{この変換は § 3 以降を見越して H.-N. よりやや一般的である。そこにあるすべての結果は、このためなおさら保たれる。}
\begin{align}
 y_1&=u_{11}x_1+\cdots+u_{1n}x_n,\quad \ldots,\quad
 y_n=u_{n1}x_1+\cdots+u_{nn}x_n, \tag{1}
\end{align}
逆変換は
\begin{align}
 x_1&=t_{11}y_1+\cdots+t_{n1}y_n,\quad \ldots,\quad
 x_n=t_{1n}y_1+\cdots+t_{nn}y_n. \tag{1'}
\end{align}
$u_{\mu\nu}$、また相互に有理的に表せるので同じことである $t_{\mu\nu}$ を $\mathsf{P}$ に添加し、係数領域 $\mathsf{P}(u)=\mathsf{P}(t)$ を得る。$\mathfrak{J}$ を、$P(u)$ を係数とする $x_1,\ldots,x_n$ のすべての多項式からなる整域とする。領域 $\mathfrak{J}$ と $\bar{\mathfrak{J}}$ が消去理論の基礎であり、$\frakS$ のイデアル $\fraka,\frakb,\ldots$ および $\bar{\frakS}$ のイデアル $\bar{\fraka},\bar{\frakb},\ldots$ を考える。任意の環におけるイデアルはここでは通常どおり、$\alpha$ と $\beta$ とともに $\alpha-\beta$ も含み、$\alpha$ とともに任意の環元 $\lambda$ に対する $\lambda\alpha$ も含むものとして定義する。$\frako$ は常に、すべての元からなる単位イデアルを表す。

\textbf{2. 変換されたイデアル。} $\frakS$ のイデアル $\frakm$ は、$\bar{\frakS}$ のイデアル $\barfrakm$ から、$u_{\mu\nu}$ または $t_{\mu\nu}$ を添加した後に (1) によって生じるとき、変換されたものと呼ばれる。したがって、$\bar f(y)$ または $\bar g(y)$ が $\barfrakm$ のすべての多項式を遍歴するとき、$\frakm$ は次のすべての線型結合からなる。
\[
 f(x)=\sum U_\lambda(u)\bar f_\lambda(y)=\sum U_\lambda(u)f_\lambda(x)
\]
または
\[
 g(x)=\sum T_\lambda(t)\bar g_\lambda(y)=\sum T_\lambda(t)g_\lambda(x);
\]
特に、すべての $f_\lambda(x)=\bar f_\lambda(y)$ が含まれ、同じことであるが、すべての $g_\lambda(x)=\bar g_\lambda(y)$ が含まれる。$f(x)$ および $g(x)$ は $\mathsf{P}(u)=\mathsf{P}(t)$ からの量を除いてしか定まらないので、$U_\lambda,T_\lambda$ は一般性を失わず $u$、それぞれ $t$ の冪積であると仮定できる。$U,T$ を任意の他の冪積に置き換えても、特に (1) の列または (1') の行を入れ換えても、対応して $u$ または $t$ に依存する係数の入れ換えを行う限り、$f(x),g(x)$ は再び $\frakm$ の多項式へ移る。したがって $\frakm$ は任意の多項式とともに、$x$ を入れ換えて得られるすべての多項式を含む。以下に現れるすべてのイデアルは、明示的に反対を述べない限り、変換されたものと仮定する。変換されていないイデアルは、(1) を
\[
 x_i=v_{i1}z_1+\cdots+v_{in}z_n,
\]
と合成することにより、$P(u,v)$ を係数とする $z$ の多項式領域における変換されたイデアルへ移る。一方、すでに変換されたイデアルに対しては、この合成は $u_{\mu\nu}$ を $u,v$ の双線型結合で置き換えることにすぎない。

\textbf{3. 記法。} $f^{(i)},g^{(i)},\ldots,a^{(i)},b^{(i)},\ldots$ は、以下つねに $x_1,\ldots,x_{i-1}$ を含まない $\mathfrak{J}$ の多項式を意味する。

\textbf{4. 基本イデアル。} 各イデアル $\frakm$ には、次の取り決めにより $n$ 個の基本イデアル $\frakg_0,\frakg_1,\ldots,\frakg_{n-1}$ が対応する。$(i-1)$ 段の基本イデアル $\frakg_{i-1}$ は、ある多項式 $b^{(i)}\ne0$（一般には $G(x)$ に依存してよい）が存在して $b^{(i)}G(x)\equiv0\pmod{\frakm}$ となるような多項式 $G(x)$ を、\emph{すべてかつそれだけ}含む。イデアル基底の存在から、すべての $G(x)$ に対して固定された少なくとも一つの $B^{(i)}$ が存在して、次が成り立つことも従う。
{\let\veqno\leqno\begin{equation}
 B^{(i)}\frakg_{i-1}\equiv0(\frakm),\qquad B^{(i)}\ne0\tag{2}
\end{equation}}
変換されたイデアルの基本イデアル自身も変換されたイデアルである（H.-N., 定理 VI）。$i$ 段の基本イデアル $\frakg_i$ は、$x_{i+1},\ldots,x_n$ を $P(u)$ に添加したときに $\frakm$ が移る先のイデアルとしても定義される。ただし、\mbox{そのとき}は $x_{i+1},\ldots,x_n$ に関して\emph{整な}多項式に制限してよく、これは一般性を失わない。$\frakg_0$ は\emph{常に}単位イデアル $\frako$ に等しい。

\textbf{5. イデアルの加群表示、単因子および個別ノルム。} 各多項式 $f(x)$ を、$x_1,\ldots,x_{i-1}$ の冪積 $\xi$ に関する線型形式とみなし、その係数を多項式 $a^{(i)}$ とみなすと、イデアル $\frakm$ は $a^{(i)}$ の領域に関する線型形式の加群 $\frakM_{i-1}$ へ移る。すなわち $\frakM_{i-1}$ は任意の二つの線型形式とともにその差を含み、線型形式 $l(\xi)$ とともに任意の $a^{(i)}l(\xi)$ も含む。そのような加群 $\frakM$ の基本加群とは、$b^{(i)}g(\xi)\equiv0\pmod{\frakM}$, $b^{(i)}\ne0$ となる線型形式 $g(\xi)$ 全体をいう。したがって基本イデアル $\frakg_{i-1}$ は、加群表示のもとで $\frakM_{i-1}$ の基本加群 $\frakG_{i-1}$ へ移る。$\frakM_{i-1}$ と $\frakG_{i-1}$ は無限個の不定元 $\xi$ に依存するが、各個の線型形式にはそのうち有限個しか現れない。$\frakM_{i-1}$ は、$x_{i+1},\ldots,x_n$ を $\mathsf{P}(u)$ に添加した後、\emph{単位と異なる単因子を有限個しかもたない}という特徴的性質をもつ。すなわち、この添加の後、$\frakG_{i-1}$ と $\frakM_{i-1}$ は基底表示を許す。ただし $\zeta$ は新しい不定元であり、$\xi$ と可逆な線型変換で結ばれ、その変換は $x_i$ に関して整であり、各 $\zeta_i=r_i(\xi),\ \xi_k=s_k(\zeta)$ にはそのつど有限個の不定元しか現れない。
{\csname tagsleft@true\endcsname\begin{align}
\frakG_{i-1}&=(\zeta_0,\zeta_1,\ldots,\zeta_\sigma,\zeta_{\sigma+1},\ldots,\zeta_{\sigma+\nu},\ldots), \notag\\
\frakM_{i-1}&=(E^{(i)}\zeta_0,E_1^{(i)}\zeta_1,\ldots,E_\sigma^{(i)}\zeta_\sigma,
\zeta_{\sigma+1},\ldots,\zeta_{\sigma+\nu},\ldots), \tag*{\raisebox{9.5pt}[0pt][0pt]{(3)}}
\end{align}}
ここで各 $E_\mu^{(i)}$ はその前のものを割る（H.-N. (24) および (28)）。

\emph{最高単因子} $E^{(i)}$ は、(2) に現れるすべての $B^{(i)}$ の最大公約因子（多項式の意味で）としても定義される。これは $x_{i+1},\ldots,x_n$ に関して整かつ原始的であると仮定でき、そのとき $x_i$ に関して正則である。すなわち、$x$ に関して次数 $\lambda$ ならば、$E^{(i)}$ は零でない係数をもつ項 $x_i^\lambda$ を含む。

(3) に現れるすべての単因子の積 $R^{(i)}$ を、$\frakM_{i-1}$ に関する $\frakG_{i-1}$ のノルム、またはイデアル $\frakm$ の\emph{個別ノルム}と呼ぶ。$x_{i+1},\ldots,x_n$ を添加すると $\frakm$ が $\frakg_i$ へ移ることを考慮すれば、記号で
\begin{equation*}
 R^{(i)}=E^{(i)}E_1^{(i)}\cdots E_\sigma^{(i)}
 =N(\frakG_{i-1}\mid\frakM_{i-1})=N(\frakg_{i-1}\mid\frakg_i).
\end{equation*}
(3) が示すように、$x_{i+1},\ldots,x_n$ を添加した後、$R^{(i)}$ は $\frakG_{i-1}$ から $\frakM_{i-1}$ への移行置換の行列式に等しくなる。また $R^{(i)}$ の次数は、$\mathsf{P}(u,x_{i+1},\ldots,x_n)$ 上線型独立な $\frakG_{i-1}$ の $\frakM_{i-1}$ を法とする剰余類の数、したがって $\frakg_{i-1}$ の $\frakg_i$ を法とする剰余類の数を与える。$R^{(i)}$ は $x_{i+1},\ldots,x_n$ に関して整かつ原始的であると仮定でき、そのとき $x_i$ に関して正則である。$R^{(i)}$ は、常に存在する階数 $\varrho$ の加群 $\frakM^*_{i-1}$ のすべての $\varrho$ 行列式の最大公約因子（多項式の意味で）として計算できる。ただしこの加群は、$\frakG^*_{i-1}$ を $\frakM^*_{i-1}$ の基本加群とするとき、剰余類系 $\frakG^*_{i-1}\mid\frakM^*_{i-1}$ が $\frakG_{i-1}\mid\frakM_{i-1}$ と同型となる性質をもつ（H.-N., 定理 VII および § 5）。

\textbf{6. イデアルの単因子形式とノルム（終結式形式）。} すべての最高単因子の積 $E_\frakm=E^{(1)}E^{(2)}\cdots E^{(n)}$ を\emph{単因子形式}と呼び、すべての個別ノルムの積 $R_\frakm=R^{(1)}R^{(2)}\cdots R^{(n)}$ を $\frakm$ の\emph{ノルム（終結式形式）}\footnote{H.-N. では終結式形式という名称だけを用いる。ノルムという名称は算術的性質に対応する。}と呼ぶ。単因子形式とノルムはイデアルで割り切れ、ノルムは単因子形式で割り切れ、後者のある冪はノルムで割り切れる。より一般には次も成り立つ。
{\csname tagsleft@true\endcsname\begin{align}
 E^{(i)}\frakg_{i-1}&\equiv0(\frakg_i),
& E^{(n)}\cdots E^{(i)}\frakg_{i-1}&\equiv0(\frakm)\quad\text{したがって}\notag\\
 R^{(i)}\frakg_{i-1}&\equiv0(\frakg_i),
& R^{(n)}\cdots R^{(i)}\frakg_{i-1}&\equiv0(\frakm),\quad\text{(H.-N., 定理 VIII),} \tag*{\raisebox{8.5pt}[0pt][0pt]{(4)}}
\end{align}}
さらに次が成り立つ。$\frakm$ が $\frakn$ で割り切れ、ノルム $R_\frakm$ と $R_\frakn$ が一致するならば、\emph{イデアル $\frakm$ と $\frakn$ も一致する}（H.-N., 定理 IX）。$\frakm$ の因子について、$R^{(1)},R^{(2)},\ldots$ の一致から基本イデアル $\frakg_1,\frakg_2,\ldots$ の一致が従い、また常に $\frakg_0=\frako$ であることを考慮すると、同様に次を得る。$\frakn$ が $\frakm$ の\emph{真の因子}なら、$\frakm$ の対応するものと異なる $\frakn$ の最初の個別ノルムは\emph{真の因子}となる。\footnote{これに対して $R_\frakn$ の後の因子は $R_\frakm$ の対応する因子の倍数となることがある。たとえば $\frakm$ が素イデアルで $\frakn$ が単位イデアルでない場合には常にそうである。} \emph{$\frakm$ が多項式 $G^{(i)}\ne0$ を含むなら}、定義により $\frakg_0,\frakg_1,\ldots,\frakg_{i-1}$ は単位イデアルに等しい。したがって 5 で述べた個別ノルムの性質、すなわち線型独立な剰余類の数を決定する性質により、$R^{(1)},\ldots,R^{(i-1)}$ は\emph{単位}に等しくなり、したがって $E^{(1)},\ldots,E^{(i-1)}$ も\emph{単位}に等しくなる。

\textbf{7. 個別ノルムと単因子の分解；基本イデアルの成分。} 素関数とは $\mathsf{P}(u)$ に関して既約な多項式をいう。準素関数とは、ここでは素関数の冪をいう。真の準素関数とは同時に素関数ではないもの、すなわち一次より高い冪が関わるものをいう。多項式の最大準素因子への分解とは、各因子が準素関数であるが任意の二因子の積はもはや準素でないような分解である。\footnote{素関数と準素関数の定義は、語「イデアル」を「多項式」に置き換えるだけで 11 と正確に類似して定式化することもできる。最大準素因子は 12 の分解における最大準素成分の類似物である。} 個別ノルム $R^{(i)}=N(\frakg_{i-1}\mid\frakg_i)$ の最大準素因子 $Q_\lambda^{(i)}$ への分解に対して、$\frakg_i$ の最小公倍数としての表示
\[
 \frakg_i=[\frakr_1,\frakr_2,\ldots,\frakr_s],
\]
が対応し、$Q_\lambda^{(i)}=N(\frakg_{i-1}\mid\frakr_\lambda)$ となる。ここで $\frakr_\lambda$ は $(i-1)$ 段の基本イデアルとして $\frakg_{i-1}$ をもち、$i$ 段の基本イデアルと一致する。この $\frakr_\lambda$ を\emph{基本イデアルの成分}と呼ぶ。$\frakr_\lambda$ の最高単因子は、$Q_\lambda^{(i)}$ と $E^{(i)}$ の最大公約因子（多項式の意味で）に等しくなり、したがって $E^{(i)}$ の最大準素因子に等しくなる。$\frakr_\lambda$ は、上に述べた基本イデアルに関する挙動と、個別ノルムまたは最高単因子がそれぞれ $R^{(i)}$ または $E^{(i)}$ の最大準素因子となるという要求によって一意に定まる（H.-N., 定理 X）。

\textbf{8. 次元概念。} イデアル $\frakm$ に対し、$R^{(i)}$ が単位と異なる最初の個別ノルムであるとき、また同値に 6 により $\frakg_i$ が単位イデアルと異なる最初の基本イデアルであるとき、最高次元 $n-i$ を割り当てる。単位イデアル $\frako$ には次元 $-1$ を割り当てる。単位と異なる個別ノルム $R^{(i)}$ が\emph{一つ}だけ存在し、したがって $\frakg_i=\frakg_{i+1}=\cdots=\frakm$ であるとき、最高次元を単に $\frakm$ の次元とも呼ぶ。$\frakm$ が多項式 $G^{(i)}\ne0$ を含むなら、6 の終わりの注意が示すように、その最高次元は高々 $n-i$ である。$\frakm$ の最高次元が $n-i$ であり、$\frakn$ が $\frakm$ の因子なら、$\frakn$ の最高次元も高々 $n-i$ である。

\textbf{9. ノルム（終結式形式）と零点。} イデアル $\frakm$ の零点とは、$\mathsf{P}(u;x_{i-1},\ldots,x_n)$ の適当な代数拡大体に属する $x_1,\ldots,x_i$ の値の系（$i=1,2,\ldots,n$）であって、$x_{i+1},\ldots,x_n$ を係数領域に添加した後、$\frakm$ のすべての多項式を零にするものをいう。この添加のもとで、ノルム（終結式形式）の各因子 $R^{(i)}$ には、$R^{(i)}$ の次数が示すだけの $\frakm$ の零点が対応する。また $R^{(i)}$ は、$u_{\mu\nu}$ とさらに別の不定元 $v$ の双線型結合 $w_{\mu\nu}$ で書くと、明示的な分解を許す。
\begin{equation}
 R^{(i)}=\prod\{z-(v_1\bar{x}_{1\nu}+\cdots-v_{i-1}\bar{x}_{i-1,\nu}-v_i\bar{x}_{i,\nu})\}^{\lambda_\nu},\notag
\end{equation}
ここで $\bar{x}_{1\nu},\ldots,\bar{x}_{i\nu}$ は互いに対応する一組の零点系を表す（H.-N., 定理 XIII。§ 3 で新たな基礎づけを与える）。したがって最高次元 $n-i$ のイデアルの零点の中には $(n-i)$ 個のパラメータに依存する零点が存在するが、それより多くのパラメータに依存するものはない。これにより 8 で与えた次元概念が通常の定式化と一致することが示される。

\textbf{10. 一般イデアル論の環領域。} 基礎となる領域は、\emph{有限鎖定理}が成り立つ可換環とする。すなわち、各 $\fraka_i$ が直前のものの真の因子（イデアルの意味での因子）であるようなイデアルの鎖 $\fraka,\fraka_1,\ldots,\fraka_\nu,\ldots$ は有限個の項の後に停止する。この要求は、すべてのイデアルがイデアル基底をもつという要求と同値である（Idealtheorie, 定理 I）。したがって特に多項式領域では満たされる。以下の 11, 12, 13 ではこの一般環領域を仮定する。

\textbf{11. 素イデアルと準素イデアル。} 積が $\frakp$ で割り切れることから少なくとも一つの因子が $\frakp$ で割り切れることが従うとき、イデアル $\frakp$ を素イデアルと呼ぶ。積が $\frakq$ で割り切れることから、一つの因子が割り切れるか、または各因子のある冪が割り切れることが従うとき、$\frakq$ を準素イデアル、または準素と呼ぶ。記号で書けば、
\[
 a\not\equiv0\pmod{\frakp},\quad b\not\equiv0\pmod{\frakp}\quad\hbox{ならば}\quad ab\not\equiv0\pmod{\frakp};
\]
\[
 a\not\equiv0\pmod{\frakq},\quad b^x\not\equiv0\pmod{\frakq}\ \hbox{すべての冪 }x\hbox{ について}\quad\hbox{ならば}\quad ab\not\equiv0\pmod{\frakq}.
\]
各準素イデアル $\frakq$ には、ただ一つの随伴素イデアル $\frakp$ が存在する。これは $\frakq$ の因子であり、そのある冪が $\frakq$ で割り切れる。すなわち $\frakq\equiv0(\frakp)$, $\frakp^\rho\equiv0(\frakq)$ である。ここで $\frakp$ は、そのある冪が $\frakq$ で割り切れるすべての環元からなる。定義が示すように、すべての素イデアルは同時に準素イデアルである。真の準素イデアルとは同時に素イデアルではないもの、したがって $\rho>1$ であるものをいう。

\textbf{12. 分解定理。} イデアルを最小公倍数として $\fraka=[\frakq_1,\ldots,\frakq_\alpha]$ と表す表示は、どの $\frakq$ も他のものの最小公倍数、すなわちその補集合の中に含まれないとき、最短であるという。また各 $\frakq$ が準素であるが任意の二つの $\frakq$ の最小公倍数は準素でないとき、最大準素成分による表示という。すべてのイデアルは\emph{有限個の最大準素成分}による最短表示をもつ。二つの異なるそのような表示では、\emph{成分数および随伴素イデアルは一致し、これらの素イデアルはすべて互いに異なる}。このように一意に定まる素イデアルを、そのイデアルの素イデアル、または随伴素イデアルと呼ぶ（Idealtheorie, 定理 IX）。

\textbf{13. 分解の孤立成分。} イデアル $\frakr=[\frakq_{i_1},\ldots,\frakq_{i_\tau}]$ は、$\frakq_{i_j}$ が $\fraka$ の最大準素成分による少なくとも一つの最短表示に現れ、かつそれらの随伴素イデアルのどれも $\fraka$ の他の素イデアルに含まれないという条件を満たすとき、$\fraka$ の分解の孤立成分と呼ばれる。分解の\emph{孤立成分}は\emph{その素イデアルによって一意に定まる}。特に孤立した最大準素成分は一意に定まる（Idealtheorie, 定理 XIII）。

\textbf{14. 標数、完全体と不完全体、第一種拡大。} 体 $\Omega$ は、単位から導かれ $\Omega$ に含まれる素体が有理数体型であるとき標数零であるといい、この素体が素数 $p$ を法とする剰余類系型であるとき標数 $p$ であるという。体 $\Omega$ は、$\Omega$ の係数をもつすべての素関数が適当な拡大体で互いに異なる線型因子に分解されるとき完全体と呼ばれ、そうでないとき不完全体と呼ばれる。標数零の体はすべて完全である。標数 $p$ の体が完全であるための必要十分条件は、各元とともにその $p$ 乗根も含むことである。不完全体における素関数は $x^{p^f}$ に関する次数 $r$ の整関数であり、適当な拡大体では $r$ 個の互いに異なる線型因子の $p^f$ 乗に分解される。$f$ を指数という。適当な拡大体で互いに異なる線型因子に分解され、したがって指数 $f$ が零であるような\mbox{素関数}を第一種という。すべての元が第一種、すなわち第一種の素関数の零点であるような拡大を第一種拡大という。第一種の元を添加して得られるすべての拡大は第一種であり、完全体上では第一種拡大しか存在しない（Steinitz, § 11 および § 13）。

\textbf{15. 有限拡大の縮約次数と指数。} $\mathfrak{L}$ が不完全体 $\Omega$ の有限拡大であるとき、$\mathfrak{L}$ は、$\Omega$ に関して第一種であり、$\frK$ における第一種の元すべてかつそれだけからなる次数 $r$ の部分体 $\frK_0$ を含む。$\frK$ の各元の $p^f$ 乗は $\mathfrak{L}_0$ に属し、$\mathfrak{L}$ にはちょうど次数 $rp^f$ の元が存在する。$r$ を縮約次数、$f$ を有限拡大の指数という（Steinitz, § 14, 1 および § 14, 5）。

\subsection*{§ 2. 素イデアルおよび準素イデアルの単因子形式とノルム。素イデアルの因子}

以後は § 1, 1 で定義した多項式領域を一貫して扱う。まず、素イデアルおよび準素イデアルについても変換されたイデアルに制限してよいことを示す。

\textbf{補題 I.} \emph{$\bar{\frakS}$ の素イデアル $\bar{\frakp}$ の変換されたイデアル $\frakp$ は素イデアルとなる。$\bar{\mathfrak{J}}$ の準素イデアル $\bar{\frakq}$ の変換された $\frakq$ は準素イデアルとなる。言い換えれば、素イデアルまたは準素イデアルであるという性質は、$u_{\mu\nu}$ を $P$ に添加しても保たれる。}

実際、$\fraka\not\equiv0\pmod{\frakp}$, $\frakb\not\equiv0\pmod{\frakp}$ と仮定する。ここで $\fraka$ と $\frakb$ は変換されたイデアルである必要はない。$\fraka$ と $\frakb$ からの多項式 $a(x)$, $b(x)$ について $a(x)\not\equiv0\pmod{\frakp}$, $b(x)\not\equiv0\pmod{\frakp}$ とする。(1) を (1') に従って反転すると、$T$ を一般性を失わず $T_{\mu\nu}$ の冪積と理解して、
\[
 a(x)=\sum T_j \bar a_j(y);\qquad b(x)=\sum T_j \bar b_j(y),
\]
を得る。そして少なくとも一つの $\bar a_j(y)$ と一つの $\bar b_j(y)$ は $\bar{\frakp}$ で割り切れない。たとえば、$T$ に関するある辞書式順序のもとでそれぞれの最高項を $\bar a_r(y)\not\equiv0\pmod{\bar{\frakp}}$ および $\bar b_s(y)\not\equiv0\pmod{\bar{\frakp}}$ とすれば、$\bar a_r(y)\bar b_s(y)\not\equiv0\pmod{\bar{\frakp}}$ でもある。したがって $a(x)b(x)\not\equiv0\pmod{\frakp}$、よって $\fraka\frakb\not\equiv0\pmod{\frakp}$ であり、$\frakp$ が素であることが分かる。この証明は、素イデアルを法とする剰余類系が零因子をもたない環をなすという事実に依拠しており、この性質は不定元の添加で保たれる。

同様に、$\fraka\not\equiv0\pmod{\frakq}$ かつすべての冪 $z$ について $\frakb^z\not\equiv0\pmod{\frakq}$ と仮定する。するとイデアル基底の存在から、$\fraka$ からの $a(x)$ と $\frakb$ からの $b(x)$ が存在して、$a(x)\not\equiv0\pmod{\frakq}$ かつすべての冪 $z$ について $b(x)^z\not\equiv0\pmod{\frakq}$ となることが従う。反対なら、$\frakb$ のある冪が $\frakq$ で割り切れるからである。

$a(x)b(x)=c(x)$ とおき、$a(x),b(x),c(x)$ の係数 $\bar a_j(y),\bar b_j(y),\bar c_j(y)$ からそれぞれ導かれるイデアルを $\fraka^*,\frakb^*,\frakc^*$ とする。するとすべての冪 $z$ について $\frakb^{*z}\not\equiv0\pmod{\bar{\frakq}}$ である。そうでなければ $b(x)^\varkappa$ が $\frakq$ で割り切れるからである。$\fraka^*\not\equiv0\pmod{\bar{\frakq}}$ なので、すべての冪 $\lambda$ について $\fraka^*\frakb^{*\lambda}\not\equiv0\pmod{\bar{\frakq}}$ でもなければならない。しかし Dedekind--Mertens の定理\footnote{H.-N., p. 63 および注 13) を参照。}により、$\fraka^*\frakb^{*h}=\frakb^{*h-1}\frakc^*$ となる指数 $h$ が存在する。したがって $\frakc^*\not\equiv0\pmod{\bar{\frakq}}$ である。これから $c(x)\not\equiv0\pmod{\frakq}$、従って $\fraka\frakb\not\equiv0\pmod{\frakq}$ が得られ、$\frakq$ が準素であることが示される。

イデアルを最小公倍数として表す場合にも、変換されたイデアルに制限してよい。

\textbf{補題 II.} \emph{$\bar{\frakS}$ における表示 $\barfrakm=[\bar{\frakc}_1,\ldots,\bar{\frakc}_\alpha]$ から、$\frakS$ における表示 $\frakm=[\frakc_1,\ldots,\frakc_\alpha]$ が従う。ここで $\frakc_\lambda$ は $\bar{\frakc}_\lambda$ の変換されたイデアルである。したがって特に、最大準素成分による任意の最短表示（§ 1, 12）では、それらを変換されたものと仮定してよい。}

実際、$\frakm$ は $[\frakc_1,\ldots,\frakc_\alpha]$ で割り切れる。なぜなら $\frakm$ からの任意の多項式 $f(x)$ は、必要なら適当な $u_{\mu\nu}$ の冪を掛けた後、$\sum U_i\bar f_i(y)$ の形をもち、各 $\bar f_i(y)$ は $\barfrakm$ からの多項式として仮定によりすべての $\bar{\frakc}_\lambda$ で割り切れるからである。逆に $f(x)$ がすべての $\frakc_\lambda$ で割り切れるなら、それは上の形をもち、したがって $\frakm$ で割り切れる。したがって $\bar{\mathfrak{J}}$ で最大準素成分への分解から出発すれば、$\frakS$ で $\frakm=[\frakq_1,\ldots,\frakq_\alpha]$ という分解を得る。ここで $\frakq_i$ は準素イデアルの変換されたイデアルとして補題 I により準素であり、また最大準素成分である。なぜなら $\bar{\frakS}$ の相異なる随伴素イデアル $\bar{\frakp}$ は $\mathfrak{J}$ の相異なる $\frakp$ に対応するからである。

いま、素と準素をまだ完全には分離しないまま、単因子形式とノルムによる素イデアルと準素イデアルの第一の特徴づけが必要となる。代数的数体におけるイデアル論との正確な類似で、\footnote{代数的数体のイデアルの最高単因子としては、そのイデアルで割り切れる最小の有理整数、したがって特に素イデアルで割り切れる素数、または準素イデアル（素イデアルの冪）で割り切れる素数冪を考えるべきである。実際、任意のイデアル $\fraka^*$ は、体基底の元 $\omega_1,\ldots,\omega_n$ に関する線型形式の加群 $A^*$ でもあり、$(\omega_1,\ldots,\omega_n)$ が $\mathfrak{A}^*$ の基本加群である。したがって $\fraka^*$ で割り切れる最小の有理整数はこの加群 $\mathfrak{A}^*$ の最高単因子となり、$\fraka^*$ のノルムは $\mathfrak{A}^*$ のすべての単因子の積となる。} 次が成り立つ。

\textbf{定理 I.} \emph{素イデアルの単因子形式は素関数に等しくなり、ノルムはこの素関数の冪となる。準素イデアルについては、単因子形式とノルムは準素関数、すなわち同じ素関数の冪となる。この素関数自体は\mbox{随伴素イデアル}の単因子形式である。素イデアルおよび準素イデアルについて、最高次元は単に次元と呼ぶものと一致する。この次元は準素イデアルとその随伴素イデアルで一致する。}

定理 I は、単因子形式とノルムが単位となる単位イデアルについて明らかに満たされる。いま\mbox{素イデアル} $\frakp$ の最高次元を $(n-i)$ とする。(4), § 1, 6 により
\[
 E^{(n)}\cdots E^{(i+1)}\frakg_i\equiv0(\frakp);\quad\hbox{しかし}\quad
 E^{(n)}\cdots E^{(i+1)}\not\equiv0(\frakp),
\]
である。そうでなければ、§ 1, 8 により $\frakp$ の最高次元は高々 $n-i-1$ になる。したがって $\frakg_i\equiv0(\frakp)$ であり、$\frakp$ はその $i$ 段の基本イデアルと同一で、従って $\frakg_{i+1},\frakg_{i+2},\ldots$ とも同一である。よって $R^{(i+1)}=N(\frakg_i\mid\frakg_{i+1}),R^{(i+2)},\ldots,R^{(n)}$ は単位となり、したがってそれらの因子である $E^{(i+1)},E^{(i+2)},\ldots$ も単位となる。ゆえに単因子形式は $E^{(i)}$ に等しくなり、これは素関数 $P^{(i)}$ に等しくなければならない。実際、§ 1, 5 と $\frakg_{i-1}$ が単位イデアルになることを考慮すれば、$E^{(i)}$ は $\frakp$ で割り切れるすべての $B^{(i)}$ の最大公約因子（多項式の意味で）として定義される。$E^{(i)}=C_1^{(i)}C_2^{(i)}\equiv0(\frakp)$ なら、$E^{(i)}$ は因子 $C_1^{(i)}$ または $C_2^{(i)}$ のどちらかに含まれなければならないからである。一方、$E^{(i)}$ のある冪は $R^{(i)}$ で割り切れるので、$R^{(i)}$ はこの素関数 $P^{(i)}$ の冪、場合によっては一次の冪となる。同時に $R^{(i)}$ は $\frakp$ のノルムとなる。単位と異なる個別ノルムは一つしか現れないので、最後に $\frakp$ の最高次元はその単なる次元と一致する。

同様に、準素イデアル $\frakq$ の最高次元を $n-i$ とすれば、上と同様に
\[
 E^{(n)}\cdots E^{(i+1)}\frakg_i\equiv0(\frakq);\quad
 (E^{(n)}\cdots E^{(i+1)})^\varkappa\not\equiv0(\frakq)
\]
がすべての冪 $\varkappa$ について成り立つ。従って上と同様に $\frakq=\frakg_i$ であり、その単因子形式は $E^{(i)}$、ノルムは $R^{(i)}$、最高次元は単なる次元である。この次元は随伴素イデアルの次元と一致する。実際、$\frakp^\rho\equiv0\pmod{\frakq}$, $\frakq\equiv0\pmod{\frakp}$ から、$\frakp$ の次元 $(n-i)$ は $\frakq$ の次元以下であり、$\frakq$ の次元は $\frakp^\rho$ の次元以下であることが従う。$P^{(i)}$ を $\frakp$ の単因子形式とすると、$(P^{(i)})^\rho\equiv0\pmod{\frakp^\rho}$ から後者の最高次元は高々 $n-i$ と分かる。したがって $n-i$ が $\frakp$ と $\frakq$ の次元である。さらに $(P^{(i)})^\rho\equiv0(\frakq)$ から、$E^{(i)}$ は $\frakq$ からのすべての $B^{(i)}$ の最大公約因子として $P^{(i)}$ の冪となり、その冪は一次の場合もある。同じことが $R^{(i)}$ にも成り立つ。これで定理 I のすべての部分が証明された。

次元概念による素イデアルの特徴づけは次で与えられる。

\textbf{定理 II.} \emph{素イデアルは同じ最高次元をもつ真の因子をもたない。逆に、この性質をもつ任意のイデアルは素である。}

証明は次に基づく。

\textbf{補題 III.} \emph{体 $\Omega$ を係数とする多項式の素イデアル $\frakp$ について、$\Omega$ 上線型独立な剰余類が有限個しかないなら、$\frakp$ は単位イデアルと異なる真の因子をもたない。言い換えれば、$\frakp$ を法とする剰余類系は体をなす。}

仮定は、任意の $(k+1)$ 個の剰余類の間に $\Omega$ からの係数、より正確には $\Omega$ のすべての元で表される剰余体 $(\Omega)$ からの係数による線型関係があるような有限数 $k$ が存在し、しかも $(\Omega)$ 上線型独立な $k$ 個の剰余類の系が少なくとも一つ存在する、ということである。そこから、すべての剰余類は選ばれた任意の $k$ 個の線型独立な系によって線型的に表されることが直ちに従う。いま $\fraka$ を $\frakp$ の真の因子とし、$a(z)\not\equiv0(\frakp)$ を $\fraka$ からの多項式とする。$c_1f_1(z)+\cdots+c_kf_k(z)\not\equiv0\pmod{\frakp}$ から
\[
 c_1a(z)f_1(z)+\cdots+c_ka(z)f_k(z)\not\equiv0\pmod{\frakp}.
\]
これは、$f_1,\ldots,f_k$ の剰余類とともに $af_1,\ldots,af_k$ の剰余類も $k$ 個の線型独立な系をなし、それを通じて特に単位類が線型的に表されることを意味する。したがって $\fraka$ は単位イデアルである。別の言い方をすれば、剰余類系は体をなす。なぜなら $a(z)\not\equiv0\pmod{\frakp}$ なら常に $1\equiv a(z)f(z)\pmod{\frakp}$ となる $f(z)$ が存在するからである。

定理 II の前半を証明するには、次を観察すればよい。次元 $(n-i)$ の任意の素イデアル、より一般に最高次元 $(n-i)$ の任意のイデアルは、$\mathsf{P}(u;x_{i+1},\ldots,x_n)$ 上線型独立な剰余類を有限個しかもたない。§ 1, 5 によりその数は $R^{(i)}$ の次数が示す数であり、ここでは $\frakg_{i-1}$ が単位イデアルになるからである。したがって $\frakp$ の任意の真の因子 $\fraka$ は、$x_{i+1},\ldots,x_n$ を $P(u)$ に添加すると単位イデアルとなる。別の言い方をすれば、$\fraka$ の $i$ 段の基本イデアルは単位イデアルとなり、これは $\fraka$ の最高次元が高々 $n-i-1$ であることを意味する。因子としての非変換イデアル $\fraka$ についても主張が成り立つようにするには、§ 1, 2 に従って新しい変換領域へ移ればよい。

逆に、$\frakp$ が同じ最高次元 $n-i$ の真の因子をもたないと仮定し、$\fraka\not\equiv0(\frakp)$ および $\frakb\not\equiv0(\frakp)$ とする。$\fraka$ と $\frakb$ は再び変換されたイデアルと仮定する。このとき仮定により、$(\fraka,\frakp)$ と $(\frakb,\frakp)$ はイデアルの意味での最大公約因子として $\frakp$ の真の因子となるので、
\pagebreak[4]
\[
 G^{(i+1)}\equiv0\pmod{(\fraka,\frakp)},\qquad
 H^{(i+1)}\equiv0\pmod{(\frakb,\frakp)},\qquad
 G^{(i+1)}\ne0,\qquad H^{(i+1)}\ne0.
\]
これから
\[
 G^{(i+1)}H^{(i+1)}\equiv0\pmod{(\fraka\frakb,\frakp)},\qquad
 G^{(i+1)}H^{(i+1)}\ne0,
\]
を得る。したがって必ず $\fraka\frakb\not\equiv0\pmod{\frakp}$ である。そうでなければ $\frakp$ の最高次元は $n-i$ より小さくなるからである。$\fraka,\frakb$ は変換されたイデアルであるから、$\bar{\frakp}$ は素イデアルであり、補題 I により $\frakp$ も素イデアルである。これで定理 II が証明された。

\subsection*{§ 3. 素イデアルおよび準素イデアルの零点}

零点の理論（§ 1, 9）の直接的な基礎づけへ進む最初の段階、まず素イデアルについての段階は、補題 III を拡張し任意の環の素イデアルに有効な次の補題である。

\textbf{補題 IV.} \emph{単位イデアルと異なる任意の素イデアルを法とする剰余類系は、商の形成によって体へ、すなわちその素イデアルの剰余体へ拡張できる。}

任意のイデアルを法とする剰余類系は環をなす。和、差、積が再び系に属し、剰余類へ移っても結合法則、交換法則、分配法則が保たれるからである。イデアルが特に素であれば、この環は零因子をもたない。素イデアルの定義（§ 1, 11）は、消えない剰余類の積も消えないことを述べているからである。素イデアルが単位イデアルと異なれば、この環は少なくとも零元と異なる元を含む。したがって商の形成、すなわち元の対の添加によって体へ拡張できる。\footnote{Steinitz, § 3。もとの整域に\emph{一つの非零元}が存在するという仮定で十分である。なぜなら、そのとき商の形成により単位の存在が保証され、もとの領域は体の部分領域となるからである。Steinitz は整域に単位が存在することを仮定している。} こうして剰余体が定義される。

まず、以下で一貫して用いる記法を定める。

\textbf{記法。} 剰余類は一般に、その任意の代表元を括弧に入れて $(x),\ldots,(a(x)),\ldots$ と表す。同様に、剰余類の体も括弧で表す。特に $(P)$, $(P(u))$, $(P(u,x_{i+1},\ldots,x_n))$ は、それぞれ $P$, $P(u)$, $P(u,x_{i+i},\ldots,x_n)$ の元によって表される剰余類すべてかつそれだけからなる剰余体を表す。

また次を思い出しておく。体 $\mathfrak{K}$ が基礎体 $\Omega$ に関して\emph{代数階数（超越次数）}$k$ をもつとは、$\mR$ が $\Omega$ 上代数的に独立な $k$ 個の量、すなわち超越量の系を少なくとも一つ含み、しかし $\mR$ の任意の $k+1$ 個の量は $\Omega$ 上代数的に従属であることをいう。$\mathfrak{K}$ が部分体 $\Omega(z_1,\ldots,z_s)$ の代数拡大体として表され、ここで $z$ が $\Omega$ 上代数的に独立、すなわち超越的であるなら、必ず $s=k$ である。\footnote{任意標数で有効なこのよく知られた事実の証明については Steinitz, § 22, 9 を参照。--- $u$ を添加しても代数階数は増えない。なぜなら、関与するのは $\Omega(u)$ からの係数をもつもとの元の有理結合だけだからである。また減ることもない。なぜなら、$\Omega(u)$ からの係数をもつ $\mR$ の元の間の関係は、$\Omega$ からの係数をもつ有限個の関係へ分解されるからである。} 同様に、$u$ が $\mR$ に含まれない不定元を表すとき、$\Omega(u)$ 上の $\mathfrak{K}(u)$ の代数階数は $k$ に等しい。

零点体の構成は、一不定元の素関数に対する通常の道筋\footnote{Steinitz, § 6。}と完全に類似して進む。

\textbf{定理 III.} \emph{次元 $n-i$ の素イデアル $\frakp$ の剰余体 $(\mR)$ には、同型的に $\frakp$ の零点体 $R$ を対応させることができる。この $R$ は代数階数 $n-i$、すなわち $n-i$ 個のパラメータに依存し、特に $\mathsf{P}(u,x_{i-1},\ldots,x_n)$ の有限拡大体とみなせる。$(\mathfrak{K})$ と $\mathfrak{K}$ の同型は、$(\mathsf{P})$ と $\mathsf{P}$、$(\mathsf{P}(u))$ と $\mathsf{P}(u)$ の同型を含み、$x_{i+1},\ldots,x_n$ をパラメータとして取るときには $(\mathsf{P}(u,x_{i+1},\ldots,x_n))$ と $\mathsf{P}(u,x_{i+1},\ldots,x_n)$ の同型も含む。逆に、代数階数 $n-i$、すなわち $n-i$ 個のパラメータに依存する任意の零点体は、剰余体 $(\mathfrak{K})$ と同型である。}

定理 III から、よく知られた定理が系として従う。すなわち、\emph{イデアル $\fraka$ が素イデアル $\frakp$ のすべての零点で消えるなら、$\fraka$ は $\frakp$ で割り切れる。}

証明のため、まず剰余体 $(\mathfrak{K})$ が $(P(u))$ に関して代数階数 $n-i$ をもつことに注意する。実際、$\frakp$ は次元 $n-i$ であるから $x_1,\ldots,x_i$ を含まない多項式を含まず、したがって類 $(x_{i+1}),\ldots,(x_n)$ は $(P(u))$ に関して代数的に独立である。一方、さらに任意の類はこれらに従属である。単因子形式 $P^{(i)}$ が $\frakp$ に属することから、§ 1, 2 により、$\lambda=1,2,\ldots,i$ についてそれぞれ $x_\lambda,x_{i+1},\ldots,x_n$ に依存する多項式が $\frakp$ に属することも従うからである。したがって $(\mathfrak{K})$ は $(\mathsf{P}(u,x_{i+1},\ldots,x_n))$ の有限拡大体となる。§ 1, 5 により、$\frakg_{i-1}$ が単位イデアルであり $\frakg_i$ が定理 I により $\frakp$ に等しいことを考慮すれば、この部分体上の次数はノルムの次数に等しい。

同型な零点体へ移るために、さらに $\mathsf{P}(u,x_{i+1},\ldots,x_n)$ と $(\mathsf{P}(u,x_{i+1},\ldots,x_n))$ が、各元にその元で表される剰余類を対応させることにより同型に関係していることを観察する。この対応のもとで、まず $x_{i+1},\ldots,x_n$ の多項式からなる整域と、$(x_{i+1}),\ldots,(x_n)$ の多項式からなる整域は同型に対応する。なぜなら $\frakp$ は $x_1,\ldots,x_i$ を含まない多項式を含まないので、$\mathsf{P}(u,x_{i+1},\ldots,x_n)$ の異なる多項式は異なる類に対応するからである。同型な整域からは商の形成により同型な体が生じる。$\mathsf{P}(u,x_{i+1},\ldots,x_n)$ と $(\mathsf{P}(u,x_{i+1},\ldots,x_n))$ の同型関係は、$\mathsf{P}$ と $(\mathsf{P})$、$\mathsf{P}(u)$ と $(\mathsf{P}(u))$ の関係を含む。この同型を $(\mathfrak{K})$ と零点体 $\mathfrak{K}$ の間のものへ拡張するため、類 $(x_1),\ldots,(x_i)$ に新しい元 $\bar{x}_1,\ldots,\bar{x}_i$ を同型的に対応させる。すなわち、各多項式 $(a(x))$ には多項式 $a(\bar{x}_1,\ldots,\bar{x}_i,x_{i+1},\ldots,x_n)$ が対応し、和と積は和と積に対応する。商の形成、ただし分母は $(\mathsf{P}(u,x_{i+1},\ldots,x_n))$ および $\mathsf{P}(u,x_{i+1},\ldots,x_n)$ から取るものに制限してよい、によって、剰余体 $(\mathfrak{K})$ には体 $\mathfrak{K}$ が対応する。これは $\mathsf{P}(u,x_{i+1},\ldots,x_n)$ の有限拡大で、その次数はノルムの次数であり、零点体と呼ぶことができる。この対応は、$f(x)\equiv0\pmod{\frakp}$ ならば $f(\bar{x}_1,\ldots,\bar{x}_i,x_{i+1},\ldots,x_n)$ が零になることを意味する。$(\mathsf{P}(u,x_{i+1},\ldots,x_n))$ の代わりに、$n-i$ 個の代数的に独立な量を添加して得られる $(\mathfrak{K})$ の\mbox{任意}の他の部分体を、議論全体に用いることもできる。

逆に、代数階数 $n-i$ の任意の零点体 $R$ が剰余体と同型であることは容易に示される。$\mathfrak{K}$ は $\mathsf{P}(u)$ からの係数を用いて $\bar{x}_1,\ldots,\bar{x}_n$ から有理的に導かれるので、これらの量の中には $P(u)$ に関して代数的に独立、したがって超越的な元が $n-i$ 個なければならない。特にこれは $\bar{x}_{i+1},\ldots,\bar{x}_n$ について成り立つ。なぜなら $P^{(i)}\equiv0\pmod{\frakp}$ なので、上と同様に $R$ は $\mathsf{P}(u,\bar{x}_{i+1},\ldots,\bar{x}_n)$ の代数拡大体であるからである。仮定により $\frakp$ の任意の多項式 $f(x)$ について $f(\bar{x})$ が零になるので、$(\mR)$ に含まれる $(x)$ の多項式領域の各類には、$\bar{x}$ の多項式である $R$ のただ一つの元が対応し、和と積は和と積に対応する。しかし異なる類は $\mathfrak{K}$ の異なる元に対応する。そうでなければ非零類、したがって適当に掛ければ $(\mathsf{P}(u,x_{i+1},\ldots,x_n))$ からの多項式で表される類が $\mathfrak{K}$ の零元に対応し、仮定に反して $\bar{x}_{i+1},\ldots,\bar{x}_n$ は代数的従属関係を満たすことになる。この対応は $\mathfrak{K}$ に含まれる $\bar{x}$ のすべての多項式を尽くす。商の形成、ここでも分母は $(\mathsf{P}(u,x_{i+1},\ldots,x_n))$ および $\mathsf{P}(u,\bar{x}_{i+1},\ldots,\bar{x}_n)$ から取るものに制限してよい、によって、これらの同型な整域から同型な体が生じる。

いまイデアル $\fraka$ が $\frakp$ のすべての零点、特に $n-i$ 個のパラメータに依存するすべての零点で消えるとする。$a(x)$ を $\fraka$ からの任意の多項式とすれば、$a(\bar{x})$ は消える。$\mathfrak{K}$ と $(\mathfrak{K})$ の同型により、類 $(a(x))$ は零類であり、したがって $a(x)$、従って $\fraka$ は $\frakp$ で割り切れる。これで定理 III とその系が証明された。

定理 III から § 1, 9 の分解へ、まず素イデアルについて移るため、また指数についてより精密な主張を行うためには、次の単因子形式に関する追加の検討が必要である。

\textbf{補題 V.} \emph{標数 $p$ において、素イデアルまたは準素イデアルの単因子形式 $E^{(i)}$、より一般に最高次元と単なる次元が一致するイデアル $\fraka$ の単因子形式が、$x_i$ に関して指数 $f$ をもち、したがって $x_i^{p^f}$ の整関数であるなら、それは $x_{i+1},\ldots,x_n$ に関しても、また $E^{(i)}$ に現れる不定元 $u_{\mu\nu}$ および $t_{\mu\nu}$ に関しても指数 $f$ をもつ。特に係数領域が完全体なら、素イデアルの単因子形式 $P^{(i)}$ は $x_i$ に関して常に指数零であり、従って第一種の素関数である。}

このためには次を観察する。$P$ が標数 $p$ の完全体であっても、$P(u,x_{i+1},\ldots,x_n)$ は不完全となるので、§ 1, 14 から直接に $P^{(i)}$ が第一種であるとは従わない。$E^{(i)}$ が $x_i$ に関して指数 $f$ をもち、別の不定元 $x_{i+\lambda}$ に関して指数 $f'<f$ をもつとする。このとき $u_{\mu\nu}$ または $t_{\mu\nu}$ を入れ換えることにより、§ 1, 2 によって、イデアルで割り切れる非零多項式 $G^{(i)}$ が存在し、これは $x_i^{p^{f'}}$ の整関数となり、単因子形式の定義により $E^{(i)}$ で割り切れる。しかし $G^{(i)}$ は $E^{(i)}$ と同じ次数をもつから、$\mathsf{P}(u)$ からの因子を除いて $E^{(i)}$ と一致しなければならない。したがって必ず $f'=f$ である。ゆえに、$t_{\mu r}$ に関して整かつ原始的と仮定できる $E^{(i)}$ は次の形をもつ。
\[
 E^{(i)}=\sum c_\lambda(t)\,x_i^{\lambda_i p^f}\cdots x_n^{\lambda_n p^f}
       =\sum c_\lambda(t)\,g_\lambda(y^{p^f},t^{p^f})
       =\sum T_\mu(t)\,\bar h_\mu(y^{p^f}),
\]
ここで $T_\mu$ は $t$ の冪積を表し、$\bar h$ は $\bar{\fraka}$ で割り切れる。したがって、$T_\mu$ における各 $t$ の指数を、その中に含まれる $p^f$ の最大倍数で置き換えれば、再び $\fraka$ に属する多項式を得る。上の表示が示すように、これは $c_\lambda(t)$ において各 $t$ の指数を、その中に含まれる $p^f$ の最大倍数で置き換えることに等しい。この置換によって得られる多項式は、単因子形式の定義により $E^{(i)}$ で割り切れ、仮定した原始性のため $t$ に関してもそうである。どの引数についても次数が増していないので、この多項式は $E^{(i)}$ と同一でなければならない。特に $\mathsf{P}$ が完全なら、$E^{(i)}$ は $x^{p^f},t^{p^f}$ の整関数として $p^f$ 乗である。したがって素イデアルの場合には必ず $f$ であり、単因子形式 $P^{(i)}$ は第一種の素関数となる。

分解は次のように与えられる。

\textbf{定理 IV.} \emph{素イデアルの単因子形式 $P^{(i)}$ は、零点体から導かれるガロア拡大体において、次の明示的分解を許す。}
\[
 P^{(i)}=\prod (x_i-t_{1i}\bar y_{1\nu}-\cdots-t_{ni}\bar y_{n\nu})^\delta
 =\prod (t_{1i}(y_1-\bar y_{1\nu})+\cdots+t_{ni}(y_n-\bar y_{n\nu}))^\delta,
\]
\emph{ここで $\bar y_{1\nu},\ldots,\bar y_{n\nu}$ はそれぞれ、もとの不定元 $y$ における互いに対応する一組の零点系を表し、$t_{1i},\ldots,t_{ni}$ には依存しない。異なる $\nu$ には異なる線型因子が対応する。また指数 $\delta$ は、完全体 $\mathsf{P}$ では一、不完全体では $p^f$（$f\geqq0$）である。新しい不定元 $v$ を添加すると、さらに、$u_{\mu\nu}$ の代わりに $u$ と $v$ の双線型結合を書いて得られる単因子形式の分解}
\[
 P^{(i)}(u,v)=\prod (z-v_1\bar x_{1\nu}-\cdots-v_i\bar x_{i\nu})^\delta,
\]
\emph{が得られ、ここで $\delta$ は上と同じ意味をもつ。素イデアルのノルム、または $\frakp$ を随伴素イデアルとする準素イデアルのノルムの、$P^{(i)}$ の冪としての分解も同様に与えられる。}

定理 IV の系として、\emph{二つの素イデアルが単因子形式 $P^{(i)}$ において一致するなら、それらは同一である。}

定理 IV の証明は、単因子形式 $P^{(i)}$ が拡大体 $(\mathsf{P}(u,x_i,x_{i+1},\ldots,x_n))$ の $(\mathsf{P}(u,x_{i+1},\ldots,x_n))$ 上の基本方程式と同一であることを示すことに帰着する。まず不定元 $u_{\mu\nu}$ および $t_{\mu\nu}$ への依存を見通すために、$\bar{\frakp}$ の\mbox{剰余体} $(\bar{\mathfrak{K}})$ へ移る。これは代数階数の定義で述べたことにより、$(\bar P)$ に関して代数階数 $n-i$ をもつ。ここで $(\bar{\mathsf{P}})$ は $\mathsf{P}$ の元で表される部分体を表す。$(\mathfrak{K})$ は、$\bar{\frakS}$ からの多項式で表せる類すべてかつそれだけから商形成によって得られる $(\mR)$ と同型な部分体に、不定元 $u$ または $t$ を添加して生じ、$(\bar{\mathsf{P}})$ と $(P)$ は\mbox{互いに対応する}。したがって $(P(u))$ に関する $(\mathfrak{K})$ の代数階数は、この部分体の $(\mathsf{P})$ に関する代数階数に等しく、同型により $(\bar{\mathfrak{K}})$ の $(\bar{\mathsf{P}})$ に関する代数階数に等しい。$(\bar{\mathfrak{K}})$ は、$(\bar{\mathsf{P}})$ からの係数を用いて類 $(y_1),\ldots,(y_n)$ から有理的に導かれるので、これらの類の中に $(n-i)$ 個の代数的に独立なものが存在し、$(\bar{\mathfrak{K}})$ はそれら $(n-i)$ 個の類を $(\bar{\mathsf{P}})$ に添加して生成される部分体の有限代数拡大体として生じる。\footnote{この注意は、定理 III とその系を $\bar{\frakp}$ の零点体 $\bar R$ について直接に述べ証明できたことを示している。変換されたイデアルへ移ることによって初めて、任意のイデアルの同じ次元の最大準素成分において、零点の中に同じパラメータが現れることが保証される。そしてそのことによって初めて、任意のイデアルのノルムおよび単因子形式との関連が生じる。}

$x_{i+1},\ldots,x_n$ に現れる不定元 $t_{\sigma\tau}$ だけを添加すると、再び剰余体が生じる。これは類 $(y_1),\ldots,(y_n)$ と $(x_{i+1}),\ldots,(x_n)$ を含み、$(\mathsf{P}(t_{\sigma\tau},x_{i+1},\ldots,x_n))$ の有限代数拡大体となる。これらは $(\mathfrak{K})$ の部分体と同型な体であるが、同一ではない。異なる不定元の場合、剰余類は同じ元で表されても同じ元を含むわけではないからである。したがって、$(P(t_{\sigma\tau},x_{i+1},\ldots,x_n))$ に対する類 $(y)$ の依存には、不定元 $t_{\sigma\tau}$ だけが入る。いま $t_{1i},\ldots,t_{ni}$ を添加して類 $(x_i)=(t_{1i}y_1+\cdots+t_{ni}y_n)$ を作り、$(x_{i\nu})=(t_{1i}y_{1\nu}+\cdots+t_{ni}y_{n\nu})$ を $(P(t_{\sigma\tau},t_{1i},\ldots,t_{ni},x_{i+1},\ldots,x_n))$ に関するガロア拡大体内にある $r$ 個の相異なる共役値とする。このとき、第一種拡大であるか否かに応じて、因子 $t-(x_{i\nu})$ の積、またはそれらの $p^f$ 乗の積は、$(P(t,x_{i+1},\ldots,x_n))$ に関する素関数 $(G)$ である。実際、次数 $r\cdot p^f$ の元が存在するので（§ 1, 15）、すべての元が $t_{1i},\ldots,t_{ni}$ の特殊化で生じる以上、$(x_i)$ は必ずそのような元でなければならず、$(x_i)$ は $(G)$ の零点である。$\frakp$ の同型な零点体およびそこから導かれるガロア拡大体へ移ると、$G$ は線型因子 $t-t_{1i}\bar y_{1\nu}-\cdots-t_{ni}\bar y_{n\nu}$、またはその $p^f$ 乗への分解を許す。いま $(G)$ は $t=(x_i)$ で消えるので、$G(x_i)$ は $\frakp$ で、したがって $P^{(i)}$ で割り切れる。また $G$ は $\mathsf{P}(u,x_{i+1},\ldots,x_n)$ に関する素関数であり、$x_{i+1},\ldots,x_n$ に関して原始的であるから、$P^{(i)}$ と同一にならなければならない。これで定理 IV の第一の分解が証明された。

第二の分解へ移るには、第一種拡大であるか否かという性質が不定元の添加で保たれることに注意する。したがって $v_1,\ldots,v_i$ とすべての $t_{\mu\nu}$ を添加した後、類 $(z)=(v_1x_1+\cdots+v_ix_i)$ と共役類 $(z_\nu)=(v_1x_{1\nu}+\cdots+v_ix_{i\nu})$ を作る。するとすべての $t-(z_\nu)$ の積、またはそれらの $p^f$ 乗の積は、再び $t=(z)$ で消える素関数 $(H)$ である。よって $H$ は上と同様に線型因子の積として表され、変換 (1) を $z=v_1x_1+\cdots+v_ix_i$ と合成して $\frakp$ から得られる素イデアルの単因子形式と同一である。ただしこの単因子形式は、相互の特殊化により、$P^{(i)}$ の $u_{\mu\nu}$ を $u$ と $v$ のある双線型結合で置き換えて得られる。\footnote{H.-N., § 7。} これらの分解により、ノルムの $P^{(i)}$ の冪としての分解も与えられる。\footnote{第一種拡大、特に完全体上では一次の冪である。不完全体上では、§ 6 の定理 XIII, XIV で示すように $p^g$ 乗となる。} 同様に、$\frakp$ を随伴素イデアルとする準素イデアルのノルムおよび単因子形式の分解も与えられる。

最後に、二つの素イデアルが単因子形式 $P^{(i)}$ において一致するなら、上の分解の帰結として零点において一致する。したがって互いに割り切れ、ゆえに同一である。

\subsection*{§ 4. 次元概念の算術的定式化}

変換されたイデアルの導入に依存しない次元概念の第一の定式化は、次によって与えられる。

\textbf{定理 V.} \emph{素イデアルの次元は、その剰余体の代数階数によって与えられる。イデアルの最高次元は、その随伴素イデアルの次元数の最大値に等しい。}

定理 V の第一部は、定理 IV への第一の注記で証明された。そこでは、$\bar{\frakp}$ の剰余体 $(\bar{\mathfrak{K}})$ の代数階数が $\frakp$ の剰余体 $(\mR)$ の代数階数に等しく、したがって § 1, 8 と一致して $\frakp$ の次元に等しいことが示されている。第二部の証明のため、$\frakm=[\frakq_1,\ldots,\frakq_\alpha]$ を最大準素成分による最短表示とする。まず、$\frakm$ の因子としてのどの $\frakq$ も、$\frakm$ より高い次元を持つことはできない。しかしすべての $\frakq$ の積は $\frakm$ で割り切れ、したがって個々の $\frakq$ のすべての単因子形式の積もまたそうである。そこで、$n-i$ が $\frakq$ たちの次元数の最大値であるならば、$\frakm$ は多項式 $G^{(i)}\ne0$ を含み、したがって $n-i$ より高い次元を持つことはできない。$\frakq$ の次元数は、定理 I により、それらの随伴素イデアルの次元数と一致する。従って、変換イデアルを導入せずに素イデアルの次元をその剰余体の代数階数によって定義するならば、定理 V の第二部が任意のイデアルの最高次元の定義を与える。

他方、素イデアルの鎖によって次元概念を把握するためには、ここでも変換イデアルの導入には依存せず、定理 II を補って次を示さなければならない。

\textbf{定理 VI.} \emph{単位イデアルと異なる任意の素イデアルは、必要ならば不定元を添加した後で、次元がちょうど一単位だけ減少した素イデアルを少なくとも一つ因子として持つ。}

実際、$\frakp$ が次元 $(n-i)>0$ を持ち、$v$ が新しい不定元を表すならば、
\[
 \frakc=(\frakp,x_i-v)
\]
は要求された性質を持つイデアルである。これは同型対応から従う。$\frakc=(\frakp^*,x_i-v)$ であり、ここで $\frakp^*$ は、$\frakp$ において $x_i$ を不定元 $v$ で置き換え、係数領域に $v$ を添加して得られる。剰余類も同様に、類 $(x_i)$ を $(v)$ で置き換えることによって得られる。しかし $x_i$ を $P(u)$ に添加しても $\frakp$ は自分自身に移る。というのは、
\[
 h(x_i)f(x)\equiv0\pmod{\frakp},\qquad h(x_i)\ne0
\]
からは必ず $f(x)\equiv0\pmod{\frakp}$ が得られるからである。そうでなければ $\frakp$ は $x_i$ だけに依存する多項式を含み、従って $x_n$ だけに依存する多項式も含むことになり、仮定に反して次元が高々零となる。従って割当て $v\sim x_i$ の下で、$\frakc$ を法とする零類は必ず $\frakp$ を法とする零類に対応し、もちろん逆も成り立つ。したがって $\frakp$ を法とする剰余類系と $\frakc$ を法とする剰余類系の間に同型がある。後者には零因子がないので、$\frakc$ は素イデアルである。この同型の下で、$(x_i)$ と $(\mathsf{P}(u,x_{i+1},\ldots,x_n))$ の間の代数従属性は、$(P(u,v))$ に関する $(x_{i+1}),\ldots,(x_n)$ の間の関係に対応する。一方、$\lambda=1,\ldots,i-1$ については、$(x_\lambda)$ と $(\mathsf{P}(u,x_{i+1},\ldots,x_n))$ の間の従属性はそのまま残り、従って代数階数をさらに下げることはない。代数階数はちょうど一単位だけ減少した。$\frakp$ が次元零であった場合、$\frakc$ は単位イデアルになる。したがって上の結果はそのまま有効である。従って元の領域 $\bar{\mathfrak{J}}$ では、
\[
 \bar{\fraka}=(\bar{\frakp},v+v_1y_1+\cdots+v_ny_n),
\]
ここで $v_\lambda=-t_{\lambda i}$ とおくと、これは要求された種類のイデアルである。特に $\mathsf{P}$ が無限個の元を含むならば、$v,v_1,\ldots,v_n$ は常に $\mathsf{P}$ の元に特殊化でき、その際、ノルムと単因子形式、したがって $\fraka$ の次元は、特殊化されたイデアル $\frakb$ のそれらに等しくなる。\footnote{H.-N., § 7 の最後の段落を参照。} ただし単因子形式は、素関数のままとは限らない。それにもかかわらず、定理 V により $\frakb$ は望ましい次元の随伴素イデアルを少なくとも一つ持つ。再び $\bar{\mathfrak{J}}$ に戻れば、したがって定理 VI はそこでは不定元を添加せずに成り立つ。

定理 II と定理 VI は、求める鎖による定式化をただちに与える。

\textbf{定理 VII.} \emph{素イデアル $\frakp$ が次元 $n-i$ を持つとは、必要ならば不定元を添加した後で、少なくとも一つの $1+(n-i)+1$ 個の素イデアルからなる鎖}
\[
 \frakp_0=\frakp,\quad \frakp_1,\ldots,\frakp_{n-i},\quad \frakp_{n-i+1},
\]
\emph{が存在し、その各項が直前の項の真の因子であり、これより多くの項を持つそのような鎖は存在しない、ということである。この定義は、$P$ が無限個の元を含むとき、元の領域においても、また不定元を導入せずにも成り立つ。}\footnote{この鎖による次元の定義は、代数的数体の場合、通常の素イデアルには次元零を、単位イデアルには次元 $-1$ を、零イデアルには次元一を与える。実際、後者も素イデアルの定義を満たす。体では非零量の積は常に零と異なるからである。この次元定義の下で、定理 II は代数的数体においても有効である。}

まず、鎖の最後の項は単位イデアルに等しくなければならないことが明らかである。単位イデアルは素イデアルの定義を満たし、そうでなければ $\frako$ を添加して鎖を延長できるからである。$\frako$ 自身については、定理 VII は § 1, 8 と一致して次元 $-1$ を与える。いま $\frakp$ を $\frako$ と異なり次元 $n-i$ のものとする。このとき鎖の項数は高々 $1+(n-i)+1$ である。定理 II により同じ次元の二つの素イデアルは現れえないからである。逆に、定理 VI により、必要ならば新しい不定元を添加した後で、この長さの鎖が少なくとも一つ存在する。鎖の各項で次元がちょうど一単位だけ減少するからである。定理 VII を次元の定義として採用するならば、この定義は明らかに元の領域でも全く同じ形で述べることができ、かつ定理 VI によって $\mathsf{P}$ が無限個の元を含むとき不定元の導入が不要になる。その場合にも、任意のイデアルの最高次元の定義は再び定理 V の第二部によって与えられる。

\subsection*{§ 5. 一般イデアル論における基本イデアルとその成分の分類。イデアルの零点}

問題は、基本イデアルとその成分 (§ 1, 7) を、一般イデアル論の意味での分解 (§ 1, 13) の孤立成分によって特徴づけることである。この特徴づけは、ノルムに対する分解定理と一般イデアル論の分解定理との平行性を示し、さらに素イデアルおよび準素イデアルについて § 3 で与えた零点の理論を任意のイデアルについて述べることを可能にする。基本イデアルについては次が成り立つ。

\textbf{定理 VIII.} \emph{第 $i$ 段の基本イデアル $\frakg_i$ は、次元 $n-1,n-2,\ldots,n-i$ の随伴素イデアル全体によって一意に決定される分解の孤立成分である。すべての随伴素イデアルが同じ次元を持つならば、イデアルもまたこの次元を持つ。すなわちノルムの因子 $R^{(i)}$ は一つだけ現れる。一般には、ノルムは、随伴素イデアルの間に現れる異なる次元数と同じ数だけ、単位と異なる因子 $R^{(i)},R^{(i+\sigma)},R^{(i+\tau)},\ldots$ を持つ。}

証明に先立って、

\textbf{補題 VI.} \emph{イデアルが最小公倍数として $\frakm=[\fraka_1,\ldots,\fraka_\alpha]$ と表示されるならば、その第 $i$ 段の基本イデアル $\frakg_i$ は、各 $\fraka$ の第 $i$ 段の基本イデアル $\frakh_i$ の最小公倍数である。}

というのは、
\[
 b^{(i+1)}f\equiv0\pmod{\frakm},\qquad b^{(i+1)}\ne0
\]
したがって
\[
 b^{(i+1)}f\equiv0\pmod{\fraka_\lambda},\qquad b^{(i+1)}\ne0;
\]
から、$\frakg_i$ は $[\frakh_{i1},\ldots,\frakh_{i\alpha}]$ で割り切れる。逆に、
\[
 c_\lambda^{(i+1)}f\equiv0(\fraka_\lambda),\qquad c^{(i+1)}\ne0
\]
したがって
\[
 c_1^{(i+1)}\cdots c_\alpha^{(i+1)}f\equiv0(\frakm),\qquad
 c_1^{(i+1)}\cdots c_\alpha^{(i+1)}\ne0,
\]
が得られ、従って逆向きの可除性も従い、補題が証明される。

いま $\frakm$ の随伴素イデアルを次元によって並べる（もちろん、ある次元が存在しないこともあり、そのときは $j_\lambda=0$）：
\[
 \frakp_{1,1},\ldots,\frakp_{1,j_1};\quad
 \frakp_{2,1},\ldots,\frakp_{2,j_2};\quad \ldots\quad
 \frakp_{n,1},\ldots,\frakp_{n,j_n},
\]
ここで一般に $\frakp_{\lambda,\kappa}$ は次元 $n-\lambda$ の素イデアルを表し、$\frakq_{\lambda,\kappa}$ は $\frakm$ の固定された一つの分解に現れる対応する準素イデアルを表す。このとき定義により、第 $i$ 段の基本イデアルは、$\lambda>i$ であるすべての $\frakq_{\lambda,\kappa}$ については単位イデアルに等しく、$\lambda\le i$ については定理 I により $\frakq_{\lambda,\kappa}$ に等しい。したがって補題 VI により、
{\let\veqno\leqno\begin{equation}
 \tag{5}
 \frakg_i=[\frakq_{1,1},\ldots,\frakq_{1,j_1};\ldots;\frakq_{i,1},\ldots,\frakq_{i,j_i}],
\end{equation}}
となり、この表示は $\frakg_i$ を分解の孤立成分として認識させる。なぜなら、$\frakg_i$ に属するどの素イデアルも、残りの素イデアルの一つに含まれることはできず、それらはすべてより低い次元を持つからである。

特に、次元 $n-i$ の随伴素イデアルだけが現れるならば、(5) により $\frakg_0,\ldots,\frakg_{i-1}$ は単位イデアルに等しく、$\frakg_i=\frakg_{i+1}=\cdots=\frakm$ である。従って $R_\frakm=R^{(i)}$ である。これに対して、異なる次元の素イデアルが現れる、たとえば $j_i,j_{i+\sigma},j_{i+\tau},\ldots$ が零でないならば、(5) により
\[
 \frakg_0\ne \frakg_i\ne \frakg_{i+\sigma}\ne \frakg_{i+\tau}\ne\cdots,
\]
一方、
\[
 \frakg_0=\frakg_1=\cdots=\frakg_{i-1}=\frako,
 \qquad
 \frakg_i=\frakg_{i+1}=\cdots=\frakg_{i+\sigma-1},\ldots;
\]
したがって因子 $R^{(i)},R^{(i+\sigma)},R^{(i+\tau)}$、そしてそれらだけが、単位と異なる。\footnote{補題 II を考慮すれば、定理 VIII は、基本イデアルが変換イデアルの最小公倍数として再び変換イデアルになることも含意する。定理 VIII は、この事実を用いない H.-N. の諸定理だけに基づくので、これは同時に内的理由を明らかにする新しい証明を与える。H.-N. ではこの定理は零点の理論においてだけ用いられるが、それはここで新たに展開される。}

定理 VIII によって、定理 I は次のように精密化される。

\textbf{定理 IX.} \emph{イデアルが準素であるための必要十分条件は、その単因子形式、したがってそのノルムが準素関数であることである。}

準素イデアルについてこの条件が満たされることは定理 I で示された。逆に、$\frakm$ の単因子形式が $Q^{(i)}=(P^{(i)})^\rho$ で与えられ、ここで $P^{(i)}$ が素関数を表すとする。定理 VIII により、$\frakm$ のすべての素イデアルは同じ次元 $n-i$ を持つ。$\frakm=[\frakq_1,\ldots,\frakq_\alpha]$ から、
\[
 Q^{(i)}\equiv0\pmod{\frakq_\lambda};
 \quad\hbox{したがって}\quad
 (P^{(i)})^\rho\equiv0\pmod{\frakp_\lambda},
 \quad\hbox{ゆえに}\quad P^{(i)}\equiv0\pmod{\frakp_\lambda}.
\]
すべての $\frakp_\lambda$ は次元 $n-i$ を持ち、$P^{(i)}$ は素関数であるから、$P^{(i)}$ は各 $\frakp_\lambda$ の単因子形式になる。従って定理 IV の系により、すべての $\frakp_\lambda$ は一致する。したがって最大準素成分を扱っているので、現れうる $\frakp$ は一つだけである。$\frakm$ は準素であり、$\rho=1$ の特別な素の場合は除外されない。定理 VIII と IX は、基本イデアルの成分の分類へと次のように導く。

\textbf{定理 X.} \emph{$R^{(i)}$ の最大準素因子 $Q_\lambda^{(i)}=(P_\lambda^{(i)})^{\rho_\lambda}$ に対応する基本イデアル $\frakg_i$ の成分 $\frakr_\lambda$ は、次元 $n-1,\ldots,n-i+1$ の随伴素イデアルと、単因子形式が $P_\lambda^{(i)}$ に等しくなる次元 $n-i$ のその素イデアルとによって一意に決定される分解の孤立成分によって与えられる。随伴素イデアルとノルムの最大準素因子とは一対一に対応する。}

§ 1, 7 により、$\frakr_\lambda$ はそれ自身の第 $i$ 段の基本イデアルと一致し、$(i-1)$ 段の基本イデアルとして $\frakg_{i-1}$ を持つ。従って定理 VIII、または (5) により、これらは最短表示
\[
 \frakr_\lambda=[\frakg_{i-1},\frakq_{\lambda,1},\ldots,\frakq_{\lambda,t_\lambda}],
\]
を持つ。ここで、すべての $\frakq$ は次元 $n-i$ を持ち、互いに異なる素イデアルに属する。$Q_\lambda^{(i)}$ の定義により、また $\frakr_\lambda$ がその第 $i$ 段の基本イデアルと一致することを考慮して、
\[
 Q_\lambda^{(i)}\frakg_{i-1}\equiv0\pmod{\frakq_{\lambda,\kappa}},
 \quad(\kappa=1,\ldots,t_\lambda),
\]
したがって
\[
\begin{aligned}
 (Q_\lambda^{(i)})^{\sigma_\lambda}&\equiv0\pmod{\frakq_{\lambda,\kappa}};
 \quad\hbox{したがって}\quad
 (P_\lambda^{(i)})^{\rho_\lambda\sigma_\lambda}\equiv0\pmod{\frakp_{\lambda,\kappa}},\\
 &\hbox{従って}\quad P_\lambda^{(i)}\equiv0\pmod{\frakp_{\lambda,\kappa}}.
\end{aligned}
\]
これから定理 IX の証明と同様に、$\frakr_\lambda$ の表示には次元 $n-i$ の $\frakq_\lambda$ と対応する $\frakp_\lambda$ が一つだけ現れうること、また $\frakp_\lambda$ の単因子形式が $P_\lambda^{(i)}$ で与えられることが従う。残るのは、$\frakp_\lambda$ が $\frakm$ の随伴素イデアルであり、$\frakq_\lambda$ が実際に $\frakm$ の最大準素成分によるある表示に現れることを証明することである。すると $\frakr_\lambda$ はまた、分解の孤立成分として認識される。すなわち、同じまたはより低い次元の素イデアルに吸収される素イデアルは存在しないので、$\frakp_\lambda$ と、より高い次元のすべての随伴素イデアルとによって決定される成分である。

この証明のためには、$\frakq_\lambda$ が $\frakg_i$ の最短表示に現れることを示せば十分である。というのも、そのような表示は定理 VIII により、常に $\frakm$ の最短表示へ補うことができるからである。したがって示すべきことは、
\[
 \frakg_i=[\frakr_1,\ldots,\frakr_\alpha]=[\frakg_{i-1},\frakq_1,\ldots,\frakq_\alpha]
\]
が最短表示になることである。例えば $\frakq_\lambda$ がその補集合によって吸収され、$\frakg_{i-1}\frakq_1\cdots\frakq_{\lambda-1}\frakq_{\lambda+1}\cdots\frakq_\alpha$ が $\frakq_\lambda$ で割り切れるとする。すると $\frakg_{i-1}\not\equiv0\pmod{\frakq_\lambda}$ であるから、$\mu\ne\lambda$ なるある $\frakq_\mu$ の冪が $\frakq_\lambda$ で割り切れることになる。従って、同じ次元を持つ随伴素イデアル $\frakp_\mu$ と $\frakp_\lambda$ は定理 II によって一致しなければならない。しかしこれは不可能である。なぜなら、それらの単因子形式 $P_\mu^{(i)}$ と $P_\lambda^{(i)}$ は定義により相異なるからである。さらに $\frakg_i$ の表示は、定理 VIII により次元 $n-i$ のすべての随伴素イデアルに対応するものであるから、そのような各素イデアルに対してノルムの最大準素因子が実際に対応することを示す。従って一対一対応が証明された。

定理 X は直ちに定理 IV の移行を与える。

\textbf{定理 XI.} \emph{ノルムの個々の最大準素因子 $Q_\lambda^{(i)}=(P_\lambda^{(i)})^{\rho_\lambda}$ は、積表示}
\[
 Q_\lambda^{(i)}=\prod \bigl(x_i-t_{1i}\bar y_{1\nu}-\cdots-t_{ni}\bar y_{n\nu}\bigr)^{\delta_\lambda\rho_\lambda}
 =\prod \bigl(t_{1i}(y_1-\bar y_{1\nu})+\cdots+t_{ni}(y_n-\bar y_{n\nu})\bigr)^{\delta_\lambda\rho_\lambda},
\]
\emph{を持つ。ここで $\bar y_{1\nu},\ldots,\bar y_{n\nu}$ は、元の不定元 $y$ における随伴素イデアル $\frakp_\lambda$ の互いに対応する一組の零点系を表し、$t_{1i},\ldots,t_{ni}$ には依存しない。また $\delta_\lambda$ は一または $p^{f_\lambda}$ の値を持つ。これにより、ノルムの線型因子への分解は完全に与えられる。$Q_\lambda^{(i)}$ の次数は、商形成によって得られた剰余類部分体 $(P(u,x_{i+1},\ldots,x_n))$ 上で、$\frakg_{i-1}$ からの剰余類によって表される剰余類の部分環の、孤立成分 $\frakr_\lambda$ を法とする次数を与える。準素イデアルの場合、これはしたがって全剰余類環である。新しい不定元 $v$ を添加すると、さらに分解}
\[
 Q_\lambda^{(i)}(u,v)=\prod (z-v_1\bar x_{1\nu}-\cdots-v_i\bar x_{i\nu})^{\delta_\lambda\rho_\lambda}.
\]

$P_\lambda^{(i)}$ は定理 X により $\frakp_\lambda$ の単因子形式として認識されているので、これは実際には定理 IV の分解の代入にすぎない。次数に関する注記は、§ 1, 5 および 7 で述べられたことの別の定式化にすぎない。\footnote{同様に、H.-N. の注 2 で述べられた重複度の定式化も、定理 X の直ちに出る帰結である。すなわち $x_{i+1},\ldots,x_n$ を添加した後、剰余類系 $\frakg_{i-1}\mid\frakr_\lambda$ は $\frakt_\lambda\mid\frakr_\lambda$ と同型である。ここで $\frakt_\lambda=[\frakg_{i-1},\frakr_1,\ldots,\frakr_{\lambda-1},\frakr_{\lambda+1},\ldots,\frakr_\alpha]$ とおく。このことは、因子 $Q_\mu^{(i)}$ の相対素性を考慮して、線型形式の加群へ戻ること（H.-N., 定理 V）から従う。直接には、$\frakt_\lambda\mid\frakr_\lambda$ が $\frakt_\lambda\mid\frakq_\lambda$ と同型であることが示される。ここで $\frakt_\lambda$ は、$x_{i+1},\ldots,x_n$ を添加した後の $\frakq_\lambda$ の補集合から生じる。}

\subsection*{§ 6. 単因子形式とノルムによる素イデアルおよび真の準素イデアルの特徴づけ}

定理 IX では準素イデアルが単因子形式とノルムによって特徴づけられたが、これは素の場合を準素の場合の特殊例と見る形での特徴づけにすぎなかった。今度は、素イデアルと真の準素イデアルを分ける必要がある。考察は § 3 の考察と平行に進む。必要条件または十分条件は容易に述べられる。

\textbf{定理 XII.} \emph{素イデアルであるための必要条件は、その単因子形式が素関数であることであり、十分条件は、そのノルムが素関数であることである。言い換えれば、素イデアルの単因子形式は常に素関数であり、真の準素イデアルのノルムは常に真の準素関数である。}

素イデアルの単因子形式が素関数であることは、すでに定理 I で示された。いまノルム $R_{\frakq}$ が素関数に等しいと仮定する。この仮定の下で、$\frakq$ の任意の真の因子 $\fraka$ がより低い最高次元を持つことを示す必要がある。そうすれば、定理 II により $\frakq$ は素イデアルとして認識される。実際、§ 1, 6 により、$R_\frakq$ の対応する因子と異なる $R_\fraka$ の最初の因子は真の因子である。$R_\frakq=P^{(i)}$ であるから、$R_\fraka$ の因子 $R^{(i)}$ は $P^{(i)}$ の真の因子、従って単位に等しくなければならない。よって $\fraka$ はより低い最高次元を持つ。

第二の簡単な証明は次の補題に基づき、次の定理もまたこの補題に基づく。

\textbf{補題 VII.} \emph{準素イデアル、より一般には次元 $n-i$ の随伴素イデアルだけを持つイデアルの単因子形式 $E^{(i)}$ を法として、多項式 $H^{(i)}$ によって表される剰余類の系は、このイデアルを法とする剰余類の部分系と同型である。単因子形式の次数は、商体 $(\mathsf{P}(u,x_{i+1},\ldots,x_n))$ に関するこの剰余類系の次数を与える。}

実際、同じ多項式 $H^{(i)}$ によって表される類を互いに対応させることによって対応を定義する。和と積は和と積に対応する。この対応は一対一でもある。というのは、イデアルを法として零類に属するすべての多項式 $H^{(i)}$ は、定義により $E^{(i)}$ で割り切れるからである。従って零類は $E^{(i)}$ を法とする零類に対応し、逆もまた成り立つ。

いま、あるイデアルのノルムが素関数 $P^{(i)}$、従って単因子形式と同一であるとする。このとき、単因子形式 $P^{(i)}$ を法とする剰余類系と $\frakq$ を法とする剰余類系の次数は、$(\mathsf{P}(u,x_{i+1},\ldots,x_n))$ に関して一致する。従って、$P^{(i)}$ を法とする剰余類系と同型な部分系は、$\frakq$ を法とする剰余類を尽くす。この系は零因子を持ちえないので、$\frakq$ は素イデアルである。

一般には、定理 XII を必要十分条件にまで鋭くすることはできない。これは後に例によって示される。しかし係数領域 $\mathsf{P}$ が完全体であるとき (§ 1, 14) には、これは可能である。

\textbf{定理 XIII.} \emph{係数領域 $\mathsf{P}$ が完全体であるならば、素イデアルのノルムと単因子形式は素関数となり、したがって互いに一致する。真の準素イデアルのノルムと単因子形式は真の準素関数となる。従って完全体上では、単因子形式が素関数であるという条件は、イデアルが素であるための必要十分条件である。この条件では、単因子形式の代わりにノルムを置いてもよい。}

定理 XII を考慮すれば、完全体上で、単因子形式が素関数であるとすぐにイデアルが素となり、そのときノルムがこの素関数に等しいことを示せば、定理 XIII は証明される。補題 V, § 3 では、完全体上で単因子形式がそもそも素関数であるならば、それは第一種の素関数になることが示された。従って、$(\mR)$ を、$(x_i)$ を $(\mathsf{P}(u,x_{i+1},\ldots,x_n))$ に添加して得られる剰余体、すなわち $P^{(i)}$ を法とする $(\mathsf{P}(u,x_i,x_{i+1},\ldots,x_n))$ とする。このとき $(x_i)$ は第一種の素関数 $(P^{(i)})$ の零点であり、従って例えば Lagrange の公式が示すように、$(y_1),\ldots,(y_n)$ は $(\mR)$ に含まれる。考えているイデアル $\frakq$ を法とする剰余類の部分系で、補題 VII により $(\mR)$ と同型なものは、分母として $(\mathsf{P}(u,x_{i+1},\ldots,x_n))$ からのものだけを含むが、したがって剰余類 $(y_1),\ldots,(y_n)$ を含み、よって $\frakq$ を法とする剰余類系を尽くす。これは $(\mR)$ と同型であるため零因子を持ちえないので、$\frakq$ は素イデアルである。$(P(u,x_{i+1},\ldots,x_n))$ に関する $\frakq$ のこの剰余体 $(\mR)$ の次数は、$\frakq$ のノルムの次数に等しい。他方、$(\mR)$ との同型により、それは単因子形式 $P^{(i)}$ の次数に等しい。従ってノルムと単因子形式は一致しなければならない。

不完全な係数領域に対しても、定理 XII はなお次の精密化を許す。

\textbf{定理 XIV.} \emph{係数領域 $\mathsf{P}$ が標数 $p$ の不完全体であるならば、素イデアルのノルムは単因子形式の $p^g$ 乗である。ここで $0\leqq g\leqq(i-1)f$ であり、$f$ は単因子形式の指数を表し、$n-i$ は素イデアルの次元を表す。}

定理 XIV は、$\frakp$ の剰余体 $(\mR)$ が、$(\mR)$ と同型な部分体 $(\mR^{\prime})=(P(u,x_i,x_{i+1},\ldots,x_n))$ に関して次数 $p^g$（$g\ge0$）を持つ、という主張に等しい。これは § 1, 15 で述べられた Steinitz の諸定理から直接従い、より正確には次から従う。

\textbf{補題 VIII.} \emph{$\mathfrak L$ を、不完全体 $\Omega$ の有限拡大で、縮約次数 $r$ と指数 $f$ を持つものとする。$\mathfrak L_0$ を $\mathfrak L$ に含まれる第一種の体とし、$\mathfrak M$ を $\mathfrak L_0$ と $\mathfrak L$ の間の中間体とする。このとき $\mathfrak L$ の任意の元の $\mathfrak M$ に関する次数は $p$ の冪である。ただし $p$ は $\Omega$ の標数を表す。}

なぜなら、$\mathfrak L$ の任意の元 $z$ の $p^{f^{\prime}}$ 乗（$f^{\prime}\le f$）は $\mathfrak L_0$ に属するからである。従って $z$ は、$\mathfrak L_0$ に係数を持つ素関数 $G(t)=(t-z)^{p^{f^{\prime}}}$ の零点である。$H(t)=(t-z)^\delta$ を、$z$ を零点に持つ $\mathfrak M$ 上の素関数とする。このとき $H(t)$ も $\mathfrak L_0(z)$ に係数を持つので、$H(t)$ の係数は $\mathfrak M$ と $\mathfrak L_0(z)$ の交わりの体に属し、従って $\mathfrak L_0$ と $\mathfrak L_0(z)$ の間の中間体に属する。よって $z$ の $\mathfrak M$ に関する次数はこの中間体に関する次数に等しく、従って $p^{f^{\prime}}$ の約数として $p$ の冪である。

定理 XIV の証明のためには、次のことだけを観察すればよい。すなわち、$r$ が縮約次数を、$f$ が $(\mathfrak{K})$ の指数を表すならば、体 $(\mR^{\prime})=(P(u,x_i,x_{i+1},\ldots,x_n))$ は $(P(u,x_{i+1},\ldots,x_n))$ に関して次数 $r\cdot p^f$ を持たなければならず、したがって $(\mR)$ に含まれる第一種の体 $(\mathfrak{K}_0)$ は $(\mR^{\prime})$ の部分体でなければならない。$(\mR)$ には次数 $r\cdot p^f$ の元が存在するので（§ 1, 15）、$(x_i)$ は必ずこの次数の元でなければならない。なぜなら $(\mR)$ のすべての元は $(x_i)$ における $t_{\mu\nu}$ の特殊化によって生じるからである。$(\mathfrak{K})$ の指数は単因子形式の指数と一致し、$(x_i)^{p^f}$ は次数 $r$ を持ち、$(\mathfrak{K}_0)$ の元であり、従って $(\mathfrak{K}_0)$ の原始元である。$(\mathfrak{K})$ は $(\mR^{\prime})$ に有限個の元、例えば $(x_1),\ldots,(x_{i-1})$ を添加して得られるので、補題 VIII を有限回繰り返して適用すれば、望む証明と同時に $g$ の評価が得られる。\footnote{A. Ostrowski, Zur arithmetischen Theorie der algebraischen Größen, Gött. Nachr. 1919, pp. 279--298, 特に p. 288 における平行した考察を参照。そこでは対応する定理が証明なしに述べられている。}

ここで、定理 XII が一般には定理 XIV を越えて鋭くできないことを示す例を付け加える。それらは、ノルムが単因子形式の平方に等しい素イデアル（$g>0$）と、単因子形式が素関数である真の準素イデアルに関するものである。後者の例はまた、ここでは定理 XIV の類似物は存在せず、ノルムが単因子形式の任意の冪になりうることを示す。

係数領域 $\mathsf{P}$ を、$2$ を法とする剰余体に二つの不定元 $\lambda,\mu$ を添加して得るものとし、従って $P$ を標数二の不完全体とする。イデアル
\[
 \bar{\frakp}=(y_1^2+\lambda,\,y_2^2+\mu),
\]
を考える。これは素イデアルでなければならない。剰余類系 $(\bar{\mR})$ が体 $\mathsf{P}(\sqrt\lambda,\sqrt\mu)$ と同型だからである。体 $(\bar{\mR})$ は $(\bar{\mathsf{P}})$ に関して次数四、縮約次数 $r=1$、指数 $f=2$ を持つ。従って単因子形式 $P^{(2)}$ は二次であり、ノルムは四次であるから $P^{(2)}$ の平方である。このことは § 1, 5 の加群表示 (3) から直接確認することもできる。得られるのは
\[
 P^{(2)}=(u_{11}u_{22}+u_{12}u_{21})^2x_2^2+(u_{11}^2\mu+u_{21}^2\lambda),
\]
また、適当に乗じた後には、
\[
 x_2^2+(t_{12}^2\lambda+t_{22}^2\mu).
\]

いま $\mathsf{P}$ を、$2$ を法とする剰余体に一つの不定元 $\lambda$ を添加して得るものとし、基礎としてイデアル
\[
 \bar{\frakq}=(y_1^2+\lambda,\,y_2^2+\lambda)=(y_1^2+\lambda,\,(y_1+y_2)^2).
\]
を取る。いま与えたものにおいて $\mu=\lambda$ と特殊化することにより、単因子形式として
\[
 P^{(2)}=x_2^2+\lambda(t_{12}^2+t_{22}^2),
\]
が得られ、これは素関数である。なぜなら $t_{12}=(0),\ t_{22}=(1)$ とすれば素関数 $x_2^2+\lambda$ が得られるからである。しかし $\bar{\frakq}$ は真の準素である。$(y_1+y_2)^2$ を含むが $(y_1+y_2)$ を含まないからである。$\bar{\frakq}$ を法とする、$P$ 上線型独立な剰余類の数が四であるため、ノルムは $P^{(2)}$ の平方、したがって $p$ 乗に等しくなければならない。

この最後の定理 XIV の類似が常に満たされるわけではないことは、次の例が示す。$P$ を、$3$ を法とする剰余体に不定元 $\lambda$ を添加して得るものとし、
\[
 \bar{\frakq}=(y_1^3+\lambda,\,(y_1-y_2)^2).
\]
とおく。このとき上と同様に $\bar{\frakq}$ は真の準素イデアルである。$y_2^3+\lambda$ も $\bar{\frakq}$ で割り切れることを考慮すると、単因子形式は
\[
 P^{(2)}=x_2^3+\lambda(t_{12}+t_{22})^3,
\]
という素関数になる。しかし、$\bar{\frakq}$ を法として六つの線型独立な剰余類があるので、ノルムは $P^{(2)}$ の平方に等しい。したがって単因子形式の $p^g$ 乗はここには現れない。

第一の例は、係数領域としての不完全体上では、代数的数体の場合と同じように、一次より高い次数の素イデアルが現れうることを示す。このとき $p^g$ を素イデアルの次数と呼ぶべきである。さらに続く例は分岐イデアルの類似物として理解される。すなわち、素数が素イデアルの高い冪、つまり準素イデアルで割り切れることがある。というのは、注 10 で示されたように、単因子形式は代数的数体のイデアルで割り切れる最小の有理整数に対応するからである。

\subsection*{§ 7. 絶対素イデアル}

$\mathsf{P}$ に係数を持つ素イデアルは、$P$ を拡大して得られる代数閉体 $A$ においても素イデアルであり続けるならば、\emph{絶対素イデアル}と呼ばれる。\footnote{代数閉体とは、よく知られているように、一つの不定元の任意の多項式が線型因子に分解する体である。任意の体に属する代数閉体の存在と本質的な一意性については Steinitz を見よ。} 対応して、$\mathsf{P}$ に係数を持つ素関数 $P$ は、$\mathsf{A}$ においても素関数であり続けるならば、\emph{絶対素関数}、すなわち絶対既約多項式と呼ばれる。絶対素イデアルの特徴づけは次によって与えられる。

\textbf{定理 XV.} \emph{絶対素イデアルの必要十分条件は、その単因子形式、これは $u$ に関して整かつ原始的と仮定してよいが、$x$ と $u$ の多項式として絶対素関数になることである。単因子形式はノルムと同一になり、したがって条件はノルムについて述べることもできる。}

証明のために、まず一般にノルムと単因子形式は係数領域の代数拡大の下で保存されることを注意する。実際、個々の因子 $R^{(i)}$ および $E^{(i)}$ は多項式の意味で最大公約数として定義されており、この性質は\mbox{係数領域}の代数拡大によって保存される。従って、$P$ から $A$ へ移るときにも $P^{(i)}$ は $\frakp$ の単因子形式であり続け、従って係数領域として完全体へ移るときにもそうである。なぜなら、代数閉体はすべて完全体であることが定義から直ちに分かるからである。定理 XIII により、$\frakp$ が絶対素イデアルであるための必要十分条件は、$P^{i)}$ が $\mathsf{A}(u)$ に関して素関数になることである。すなわち、$P^{(i)}$ を $u$ に関して整かつ原始的と仮定するとき、$x$ と $u$ の多項式として絶対素関数になることである。同時に、定理 XIII により、$A$ を\mbox{係数領域}とするとすぐに $P^{(i)}$ は $\frakp$ のノルムと同一になる。従って、いま注意したことにより、同じことは $P$ 上でも成り立つ。これで定理 XV が証明された。

（有限）代数的数体 $\mathfrak K$ に係数を持つ絶対素関数 $P$ については、いま次の定理\footnote{A. Ostrowski, loc. cit.（注 21）, 補題 p. 296。より簡単な証明は E. Noether, Ein algebraisches Kriterium für absolute Irreduzibilität, Math. Ann. 85 (1922), pp. 26--33, no. 7 にある。} がある。すなわち、$\mathfrak K$ の有限個を除くすべての素イデアルを法として、$P$ は素関数、より正確には絶対素関数のままである。この定理を絶対素イデアルへ移すことは、次に基づく。

\textbf{定理 XVI.} \emph{係数領域として体の代わりに（有限）代数的数体 $\mathfrak K$ のすべての代数的整数からなる環 $\mathfrak o^*$ を取るならば、多項式イデアルのノルムと単因子形式は、$\mathfrak K$ の有限個を除くすべての素イデアルを法として、ノルムおよび単因子形式として保存される。}

証明のためには、§ 1, 1 に対して、基礎となる多項式領域を修正する必要がある。$\bar{\mathfrak{J}}$ を、$\mathfrak o^*$、すなわち $\mathfrak K$ の代数的整数を係数とする $y_1,\ldots,y_n$ のすべての多項式からなるものとする。$\mathfrak{J}$ の係数領域を $\mathfrak o^*[u]$、すなわち $\mathfrak o^*$ に係数を持つ $u$ のすべての多項式とする。$\bar{\frakm}$ が $\bar{\mathfrak{J}}$ のイデアルであるならば、(1) によって $\frakS$ の変換イデアル $\frakm$ に移る。示すべきことは、ノルムを作る際に分母に現れる $\mathfrak o^*[u]$ の多項式は有限個だけである、ということである。そうすれば定理 XVI は直接従う。

まず、§ 1, 5 の加群表示 (3) は、個別ノルムの形成にあたって、$x_{i-1}$ の冪を $x_{i-1}$ に正則な多項式を法として簡約することに帰着することに注意する。
\[
 C^{(i-1)}=U_{i-1}(u)x_{i-1}^{k}+\hbox{低次項},
\]
\footnote{H.-N., § 4 を参照。} ここで $C^{(i-1)}$ のすべての係数、特に $U_{i-1}$ は $\mathfrak o^*[u]$ の元と仮定してよい。これは $x_1,\ldots,x_{i-1}$ の逐次簡約の問題であるから、$x$ に関して整な加群表示は、分母に冪積
\[
 U_1^{\lambda_1}\cdots U_{i-1}^{\lambda_{i-1}}
\]
を除けば、$\mathfrak o^*[u]$ においても整である。$U_1U_2\cdots U_n$ を $u$ の多項式として考えると、その係数によって $\mathfrak K$ のイデアル $\mathfrak r^*$ が定まり、したがってそれは $\mathfrak K$ の有限個の素イデアルでしか割り切れない。$\frakp^*$ をこれら有限個のものと異なるように選び、$\frakp^*$ を法とする剰余体を係数領域として用いるならば、$\mathfrak o^*$ の各数を $\frakp^*$ を法とするその剰余類で置き換えることによって、もとの加群表示は保存される。

さらに、ノルムと単因子形式は $\mathfrak o^*[u]$ の因子を除いてのみ定まるので、$\mathfrak o^*[u]$ に関して $\frakm$ で割り切れると仮定してもよい。この約束の下で、例えば
\[
 R^{(i)}=V_i(u)x_i^{l}+\hbox{低次項},
\]
と仮定する。このとき、$C^{(i-1)}$ によって定まる階数 $\rho$ の加群 $\mathfrak M_{i-1}^*$ のすべての $\rho$ 行小行列式の可除性において、分母には $U_1^{\lambda_1}\cdots U_{i-1}^{\lambda_{i-1}}$ のほかに $V_i$ の冪だけが現れる。$\frakp^*$ を、$V_1V_2\cdots V_n$ によって定まるイデアルに現れる $\mathfrak K$ の有限個の素イデアルともさらに異なるように選ぶ。すると、$\frakp^*$ を法とする剰余体を係数領域として取るとき、$R^{(i)}$ はこれらの小行列式の共通因子であり続け、また $\mathfrak M_{i-1}^*$ の階数は保存される。なぜなら $C^{(i)}$ は $\mathfrak M_{i-1}^*$ からの $\rho$ 行小行列式であったからである。

最後に、$R^{(i)}$ がそれぞれの場合に最大公約数であり続けるように $\frakp^*$ を選ばなければならない。$D^{(i)}$ が $\mathfrak M_{i-1}^*$ からのすべての $\rho$ 行小行列式を遍歴するとする。仮定により、
\[
 U_1^{\lambda_1}\cdots U_{i-1}^{\lambda_{i-1}}V_i^{\sigma}D^{(i)}=R^{(i)}T^{(i)}
\]
であり、ここで $T^{(i)}$ たちは多項式の意味で $x$ を含む共通因子を持たない。したがって、すべての $T^{(i)}$ から導かれるイデアルは多項式
\[
 G^{(i+1)}=W_i(u)x_{i+1}^{m}+\hbox{低次項},\qquad W_i\ne0.
\]
$G^{(i+1)}$ を $x_{i+1}$ に正則と仮定してよいことは、変換 (1) に、係数として新しい不定元を用いる $x_{i+1},\ldots,x_n$ の同種の変換を合成することから従う。そこで、$\frakp^*$ をさらに $W_1W_2\cdots W_n$ によって定まるイデアルに含まれる有限個の素イデアルとも異なるように選ぶならば、$\frakp^*$ を法とする剰余体を係数領域として\mbox{用いるとき}、$R_\frakm$ は実際にノルムとして保存される。全く同様に、$E_\frakm$ も単因子形式であり続けるように、$\frakp^*$ をさらに選ぶことができる。これで定理 XVI が証明された。

定理 XV と XVI から、絶対素関数についての定理の一般化として、またその定理を用いて、次を得る。

\textbf{定理 XVII.} \emph{$\frakp$ が（有限）代数的数体 $\mathfrak K$ の代数的整数を係数とする絶対素イデアルであるならば、$\mathfrak K$ の有限個を除くすべての素イデアルを法として、$\frakp$ は素イデアル、より正確には絶対素イデアルであり続ける。}

証明のため、$\frakp$ のノルム $R_\frakp$ を、定理 XVI と同じように、$\mathfrak o^*[u]$ に関しても $\frakp$ で割り切れるように選ぶ。すると
\[
 R_\frakp=T(u)P^{(i)}(x),
\]
ここで $P^{(i)}(x)$ は $u$ に関して整かつ原始的と仮定してよい。定理 XV により、$P^{(i)}$ は、$\mathfrak K$ に係数を持つ $u$ と $x$ の多項式と見れば、絶対素関数である。絶対素関数に関する定理により、これは $\mathfrak K$ の有限個を除くすべての素イデアルを法として、この性質を保つ。従って $\frakp^*$ をこれら有限個の素イデアルと異なり、さらに定理 XVI により $\mathfrak K$ の有限個の別の素イデアルとも異なるように選ぶ。すると、$P^{(i)}$ は、$\mathfrak o^*$ からの係数を $\frakp^*$ を法とする対応する剰余類で置き換えることで、$\frakp^*$ を法とする剰余体を係数領域 $P$ としたときの $\frakp$ のノルムとなり、このノルムは絶対素関数である。したがって $P$ を係数領域とすれば、定理 XV により $\frakp$ は絶対素イデアルであり続ける。これで定理 XVII が証明された。

\begin{center}
Göttingen, 1923 年 5 月 8 日。\\[1em]
（1923 年 3 月 9 日受理。）
\end{center}



\clearpage

\section*{25. 消去理論とイデアル論。}

\begin{center}
J. Ber. d. DMV 33 (1924), S. 116--120
\end{center}

\setcounter{footnote}{0}

\begin{center}
1923年9月25日、Marburg で行われた講演。\footnote{詳しい叙述と文献情報は \emph{Math. Ann.} 90 および近日刊行予定の論文に掲載される。}\\[0.5em]
Göttingen の \textsc{Emmy Noether} による。
\end{center}

以下では、一変数多項式の零点論と完全に平行して消去理論を構成できることを示したい。この構成は、\emph{Hentzelt} の消去理論を、\emph{Steinitz} が展開した体論の方法、および私自身が展開したイデアル論の方法と結びつけることに基づく。

まず、一変数多項式の場合について問題設定を概説する。係数領域としては、素関数などの概念がそれを基準として定まる一つの固定体 $P$ を、終始、基礎に置く。$P$ は Steinitz の意味で抽象的に定義されたものとしてよいが、特殊な場合には、消去理論で以前一般に仮定されていたように、有理数体または複素数全体からなる体と一致してもよい。一変数多項式の零点論は、次の四つの段階に分かれる。

1. \emph{分解定理}。これは、問題を素関数、すなわち $P$ 上既約な多項式に関する問題へ還元する。実際、
\[
 a(x)=p_1(x)^{\varrho_1}\cdots p_r(x)^{\varrho_r}=q_1(x)\cdots q_r(x)
\]
から、ここで $p(x)$ は互いに異なる素関数を、したがって $q(x)$ は「最大準素因子」を表すので、$a(x)$ は素関数 $p(x)$ のすべての零点で、かつそれらの零点でのみ消えることが従う。

2. \emph{素関数 $p(x)$ の零点体（零点が生成する体）の構成}。このような零点体 $\mathfrak K=P(\alpha)$（ここで $\alpha$ は $p(x)$ の零点）は、\emph{Steinitz} によれば、$p(x)$ の剰余体 $(\mathfrak K)$ と同型な $P$ の拡大体として常に構成できる。逆に、$P$ のある拡大体に含まれる任意の零点体は剰余体 $(\mathfrak K)$ と同型になる。したがって、零点体と剰余体は抽象的には同一である。

3. \emph{$p(x)$ によって定まるガロア拡大体における $p(x)$ の一次因子分解}。2. の構成を反復することにより、本質的に一意に定まるガロア拡大体 $P(\alpha_1,\ldots,\alpha_m)$ に到達する。この体は、$p(x)$ が一次因子に分解するためにちょうど十分なものである。

4. \emph{もとの多項式 $a(x)$ に対する重複度の意味}。1., 2., 3. により、任意の多項式 $a(x)$ について、零点の問題とそれに対応する一次因子分解は解決される。このとき重複度として、個々の準素因子 $q(x)$ の次数、すなわち $q(x)$ を法とする、$P$ 上線型独立な剰余類の個数を考えることができる。実際、係数領域 $P$ が代数閉で $p(x)$ が一次式になる場合、この個数はちょうど零点の重複度となる。

いま一般の消去理論へ移り、$P$ の係数をもつ $x_1,\ldots,x_n$ のすべての多項式からなる多項式環 $\mathfrak{J}$ を基礎に置く。このとき消去理論は、この多項式環の\emph{一つのイデアルの零点の理論}として定義できる。イデアル $\mathfrak a$ とは、通常どおり、$f$ と $g$ とともに $f-g$ も、また $f$ とともに $Af$ も $\mathfrak a$ に属するような多項式の系をいう。ここで $A$ は $\mathfrak{J}$ の任意の多項式である。$\mathfrak a$ の零点とは、$P$ の拡大体に属する値の系 $\alpha_1,\ldots,\alpha_n$ であって、イデアルに属するすべての多項式 $f$ に対して $f(\alpha)$ が消えるものをいう。$f_1,\ldots,f_t$ が、常に存在する $\mathfrak a$ のイデアル基底を表す、すなわち $\mathfrak a$ の任意の $f$ に対して $f=A_1f_1+\cdots+A_tf_t$ が成り立つならば、$\mathfrak a$ の零点は明らかに $f_1,\ldots,f_t$ の零点と同一である。逆に、有限個の多項式からなる任意の系は、それらが生成するイデアルの基底とみなせる。したがって、消去理論について上で与えた零点の定義は通常の定義と同一でありながら、基底を恣意的に特別扱いすることから生じる困難を避ける。一変数多項式の場合には主イデアルしか現れないため、この困難は生じない。基底多項式が、単元、すなわち $P$ の零でない元を因子として掛けることを除いて一意に定まるからである。それでも、この場合にもイデアルによる把握はより精密な理解を与える。

消去理論は、一変数の場合に対応して、次の四つの段階に分かれる。

I. \emph{分解定理}、すなわち一つのイデアルを最大準素成分によって表すこと。これは、問題を素イデアルに関する問題へ還元する。イデアルが、積を割る場合には必ず少なくとも一つの因子を割るならば、素イデアルと呼ぶ。また、積を割る場合に、少なくとも一つの因子または各因子のある冪を割るならば、準素イデアルと呼ぶ。各準素イデアル $\mathfrak q$ には、ただ一つの随伴素イデアル $\mathfrak p$ が存在し、$\mathfrak p$ は $\mathfrak q$ を割り、かつ $\mathfrak p$ のある冪は $\mathfrak q$ で割り切れる。最小公倍数としての表示
\[
 \mathfrak a=[\mathfrak q_1,\cdots,\mathfrak q_r];\qquad
 \mathfrak p_1^{\sigma_1}\cdots \mathfrak p_r^{\sigma_r}\equiv 0\,(\mathfrak a).
\]
が成り立つ。ここで $\mathfrak q$ は最大準素成分である。すなわち、各 $\mathfrak q$ は準素であるが、任意の二つの $\mathfrak q$ の最小公倍数は準素ではない。さらに、一般性を失うことなく、どの $\mathfrak q$ も他の $\mathfrak q$ たちの最小公倍数を割らない、したがってどの $\mathfrak q$ も単に省略できない、と仮定してよい。$\mathfrak p$ は $\mathfrak q$ の随伴素イデアルである。いま、$\mathfrak a$ によって一意に定まるのは $\mathfrak q$ ではなく、それらの随伴素イデアル $\mathfrak p$ であり、ここに一変数の場合との類比がある。そして各 $\mathfrak p$ は $\mathfrak a$ を割り、逆に $\mathfrak p$ たちの冪積は $\mathfrak a$ で割り切れるので、イデアルの零点はその随伴素イデアルの零点と一致する。

II. \emph{素イデアル $\mathfrak p$ に対する零点体の構成}。定義により、素イデアルを法とする剰余類は零因子をもたない環をなし、$\mathfrak p$ が単位イデアルと異なるならば、少なくとも零元とは異なる元を含む。したがって、この整域の商体をとることにより、この環を $\mathfrak p$ の剰余体 $(\mathfrak K)$ へ拡張できる。$(\mathfrak R)$ と同型な $P$ の拡大体 $\mathfrak K$ を常に構成でき、これは $\mathfrak p$ の零点体となる。すなわち $\mathfrak K=P(\alpha_1,\ldots,\alpha_n)$ であり、ここで $\alpha_1,\ldots,\alpha_n$ は $\mathfrak p$ の一つの零点を表す。この $\mathfrak R$ は $P$ 上の超越次数 $k$（$0\leq k<n$）をもち、したがって $P$ 上代数的に独立な $k$ 個の元を含む一方、任意の $(k+1)$ 個の元は $P$ 上代数的に従属する。逆に、$P$ のある拡大体に含まれる超越次数 $k$ の任意の零点体は剰余体 $(\mathfrak K)$ と同型になる。したがって、超越次数 $k$ の零点体と剰余体は抽象的には同一である。

III. \emph{$\mathfrak p$ によって定まるガロア拡大体における $\mathfrak p$ のノルムの一次因子分解}。不定元 $x_1,\ldots,x_n$ が、もとの不定元 $y_1,\ldots,y_n$ から、係数として不定元をもつ線型変換によって生成されると考える。このとき、超越次数 $k$ の零点体 $\mathfrak K$ は $P(x_{i+1},\ldots,x_n)$ の有限次代数拡大体とみなせる（$k=n-i$）。したがって、$P(x_{i+1},\ldots,x_n)$ に関するガロア拡大体 $\overline{\mathfrak K}$ へ移ることができ、その中で零点 $\alpha_i$ をもつ素関数 $P^{(i)}(x_i,x_{i+1},\ldots,x_n)$ は一次因子に分解する。この素関数 $P^{(i)}$ は、もとの不定元 $y$ へ戻るとき、基本方程式の役割を果たす。実際、$\alpha_i$ は基本形式 $u_1\beta_1+\cdots+u_n\beta_n$ に等しくなる。ここで $\beta$ は $y$ における零点系であり、パラメータ $x_{i+1},\ldots,x_n$ に代数的に依存する。したがって、一次因子分解はすべての零点にわたる積を与える。しかし、素関数 $P^{(i)}$ またはその冪は、$\mathfrak p$ のノルム $N(\mathfrak p)$ とみなすこともできる。そのためには、$\mathfrak p$ を $x_1,\ldots,x_{i-1}$ の冪積を変数とする線型形式からなる加群とみなし、$P$ の代わりに $P(x_{i+1},\ldots,x_n)$ を係数領域として基礎に置き、ノルムをすべての単因子の積として定義する。単位と異なる単因子は有限個しかない。係数領域 $P$ が完全体ならば、常に $N(\mathfrak p)=P^{(i)}$ となるが、不完全体では $P^{(i)}$ の冪がノルムとして現れることがある。

IV. \emph{もとのイデアルのノルムと重複度の意味}。対応して、イデアル $\mathfrak a$ を $x_1,\ldots,x_{i-1}$ の冪積を変数とする線型形式からなる加群として順次考える（$(i=1,2,\ldots n)$）。すると、$P(x_{i+1},\ldots,x_n)$ を係数領域として基礎に置いたとき、この加群が単位と異なる単因子をそれぞれ有限個しかもたないことが示される。これは、イデアル基底の存在から従う有限性条件である。各段階における有限個の単因子の積が、それぞれ個別ノルムを与え、ノルム $N(\mathfrak a)$ はすべての個別ノルムの積として定義される。$N(\mathfrak a)$ の最大準素因子 $Q^{(i)}(x_i,x_{i+1},\ldots,x_n)$ は、$\mathfrak a$ の随伴素イデアル $\mathfrak p$ と一対一に対応し、$Q^{(i)}$ は $P^{(i)}$ の冪となる。ここで $P^{(i)}$（または、場合によっては $P^{(i)}$ の冪）は、$\mathfrak p$ のノルム $N(\mathfrak p)$ を表す。したがって、分解
\[
 N(\mathfrak a)=N(\mathfrak p_1)^{\varrho_1}\cdots N(\mathfrak p_r)^{\varrho_r}=Q_1\cdots Q_r,
\]
が得られる。ここで、場合によっては $N(\mathfrak p)$ を最大単因子 $E(\mathfrak p)=P^{(i)}$ で置き換える。これにより、I, II, III に従って、$\mathfrak a$ のすべての零点に対するノルムの一次因子分解が完了する。$\mathfrak a$ の個々の素イデアル $\mathfrak p$ に対応する零点の重複度として、$Q^{(i)}$ の次数を考えることができる。これは、$P(x_{i+1},\ldots,x_n)$ 上線型独立な剰余類の個数と解釈できる。\newpage
\noindent より正確には、$\mathfrak p$ とより高次元の素イデアルとによって一意に定まる、分解におけるある孤立成分の剰余類が問題となる。すなわち、第 $i$ 段階の基本イデアルのうち $\mathfrak p$ によって定まる成分を法とする、第 $(i-1)$ 段階の基本イデアルの剰余類である。この第 $i$ 段階の成分は $\mathfrak p$ およびすべてのより高次元の素イデアルを随伴素イデアルにもつが、第 $(i-1)$ 段階の基本イデアルはより高次元の素イデアルだけを随伴素イデアルにもつ。

\emph{一変数の場合の完全な位置づけ}は、ここでも主イデアルへ移ることによって得られる。1. で与えた分解は、そのとき $a(x)$ を基底多項式とみなすか、ノルムとみなすかに応じて、二つの別々の分解に分かれる。すなわち、一つは I に従う、最大準素成分の最小公倍数としての主イデアルの表示であり、ここでは素イデアルの冪積としての表示に帰着する。もう一つは IV に従う、個々の素イデアルのノルムの冪積としてのノルムの分解である。2. による零点体の構成で実際に問題となっているのは、$p(x)$ によって定義された素イデアルの零点体であり、3. ではノルムが一次因子に分解される。4. による重複度の定義は IV で与えた定義に従属する。第零段階の基本イデアルは単位イデアルに等しく、一変数では第一段階の基本イデアルがすでにイデアル自身に等しいからである。

\emph{通常の消去理論の中での位置づけ}は、イデアルのノルムが終結式としての意味をもつことによって与えられる。この見方では、消去理論の古い問題は多項式環のイデアル論の一部、すなわち、イデアルの分解定理と並んで、それらのノルムの分解および二つの分解の関係を扱う部分となる。

\begin{center}
（1923年10月19日受理。）
\end{center}

\clearpage

\section*{26. 代数体におけるイデアル論の抽象的構成}

\begin{center}
J. Ber. d. DMV 33 (1924), S. 102
\end{center}

4. E. Noether, ゲッティンゲン：代数体におけるイデアル論の抽象的構成。

代数体におけるイデアル論と完全に同値であり、したがってそのようなイデアル論をもつすべての環を特徴づける抽象的性質が示された。それは、単位元をもち零因子をもたない環であって、二重鎖条件（Doppelkettensatz）が成り立つもの、すなわちイデアルからなる任意の約イデアル鎖（Teilerkette）が有限で止まり、また零イデアルと異なる任意のイデアルを法として任意の倍イデアル鎖（Vielfachenkette）も有限で止まるもの、かつその商体に関して代数的に整閉であるものである。最後の仮定を外すと、有限オーダー（order）におけるイデアル論についての定理が得られ、その一部は Dedekind（第3版）にすでに証明なしで見られる。

詳しい叙述は \emph{Mathematische Annalen} に現れる予定である。

\section*{27. イデアル論における Hilbert 数。}

\begin{center}
J. Ber. d. DMV 34 (1925), S. 101
\end{center}

--- \textbf{E. Noether}：イデアル論における Hilbert 数。Hilbert 数とは一般に、イデアルの剰余類を数え上げることに関係する数をいう。Hilbert 自身はその「特性関数」において、多項式領域の各イデアルに対し、それぞれある限界以上で線型独立な剰余類の正確な個数を与える関数を導入した。低い次数については、Macaulay がこの関数を生成関数によって補い、Ostrowski はその生成関数を明示的に求めることができた。しかし Macaulay はさらに別種の数も与えており、それらは抽象的に定式化できるため、一般のイデアル論にとって最も本質的なものと判明した。一般の分解定理によれば、準素イデアル $\mathfrak q$ について定義すれば十分である。ここで問題となるのは、付随する素イデアル $\mathfrak p$ の剰余体上で線型独立な剰余類の個数である。これらの数は部分系に分かれ、各部分系は $\mathfrak q,\mathfrak q/\mathfrak p,\mathfrak q/\mathfrak p^2,\ldots$ の既約成分の数によって与えられる。第 $i$ 部分系の要素数は、$\mathfrak q/\mathfrak p^i$ から $\mathfrak q/\mathfrak p^{i-1}$ へ至る組成列の項数によっても特徴づけられる。全体の組成列とその Hilbert 数は、準素イデアルを素イデアルの冪として表す表示――これは一般の場合、例えば代数的数体の有限オーダーでは、もはや成り立たない――に相当するものをなす。

\section*{28. 群指標とイデアル論}

\begin{center}
J. Ber. d. DMV 34 (1925), S. 144
\end{center}

\textbf{3. E. Noether, ゲッティンゲン：群指標とイデアル論。}

Frobenius の群指標の理論、すなわち有限群の表現論は、完全可約環（今日の半単純環）である群環のイデアル論として捉えられる。したがって、完全可約環の加法的分解における同型定理が一般に展開される。直和としての、直既約な単純イデアルへの表示は同型を除いて一意であり、同じ両側イデアルに属する直既約片側イデアルは互いに同型である。したがって同型なイデアルを一つの類にまとめると、両側直既約イデアルの数と同じだけの異なる直既約イデアル類が存在する。

完全可約な超複素量の系、特に群環へ特殊化すると、直既約イデアル類と既約表現類は互いに一意に対応する。さらに両側直既約イデアルと指標が対応し、両側イデアルの加法的直既約成分の数は直既約イデアルの階数に等しい。これにより Frobenius の理論は位置づけられる。

詳しい叙述は \emph{Math. Ann.} に現れる予定である。

\section*{29. 標数 $p$ の有限線型群の不変式に関する有限性定理}

\begin{center}
Nachr. v. d. Ges. d. Wiss. zu Göttingen 1926, pp. 28--35
\end{center}

1926年7月30日の会合で R. Courant により提出。

以下で私は、線型置換の有限群の不変式に関するよく知られた有限性定理\footnote{有限性の初等的証明については E. Noether, Der Endlichkeitssatz der Invarianten endlicher Gruppen. Math. Ann. 77 (1916), pp. 89--92 を参照。}が、係数領域として通常の数体の代わりに任意の体、特に標数 $p$ の体をとっても成り立つことを示す。\footnote{体論の概念については E. Steinitz, Algebraische Theorie der Körper. J. f. M. 137 (1910), pp. 167--309 を参照。標数は § 4、根体は § 12、商体は § 3。} ただし、群の位数が $p$ で割り切れると既知の有限性証明は失敗する。そのため、より深い算術的事実、すなわち次の、これまで標数零についてだけ知られていた基準を用いなければならない。

\textbf{有限性基準}\footnote{E. Noether, Algebraische und Differentialinvarianten. Jahresber. d. d. Math.-Ver. 32 (1923), pp. 177--184, 3番、定理4、およびそこに挙げた文献を参照。その基準の最後の文では、誤って「$\mathfrak R$ に関する加群基底」（整量の基本系）とあるが、「$\mathfrak R$ のある部分環に関する」とすべきである。}：体を係数にもつ多項式、より一般には有理関数からなる整域 $\mathfrak J$ が有限であるための必要十分条件は、$\mathfrak J$ の中に有限な部分環 $\mathfrak R$ が含まれ、$\mathfrak J$ が $\mathfrak R$ に代数的かつ整に依存することである。そのとき $\mathfrak R$ の任意の整基底は、$\mathfrak R$ のある部分環に関する $\mathfrak J$ の加群基底によって補われて、$\mathfrak S$ の整基底になる。

この基準は、それ自体独立の興味をもつが、有理基底の存在と、有限拡大の整量へ移行しても約元鎖定理が保たれるという事実に基づいている。この第二の定理は、直前の Artin と van der Waerden のノートで、任意の標数について、ただしもとの領域にある有限性仮定を置いた上で証明された。その仮定は有限群の場合には満たされる（根環は有限拡大を表す）。有理基底の存在については、任意の体を係数領域とする場合に有効な証明を与える。これにより、上の制限の下で有限性基準を証明できる。

有限群の不変式について有限性基準が満たされることは、対称関数の原理に基づく。標数零では Galois リゾルベントの係数がすでに整基底を与えるのに対し、一般の場合にはそれらは有限性基準で要求される有限な部分環を生成するだけである。この有限な部分環の存在から、同じ推論によって「相対」不変式および「モジュラー」不変式の有限性が従う。

ここで考える不変式には、Dickson とその学派によって扱われた「射影的不変式 mod $p$」が特殊な場合として含まれる。\footnote{Dickson, On Invariants and the theory of Numbers. The Madison Colloquium 1913 を参照。より新しい文献は、その後刊行された Transactions of the American Mathematical Society の諸巻にある。} ここでは、有限性定理はモジュラー不変式については知られていたが、一般の「形式的」不変式については知られていなかった。

\subsection*{§ 1. 有限性基準}

1. \textbf{有理基底の存在定理}。第一の定式化：体 $P$ を係数とする $n$ 個の不定元の有理関数の任意の系 $S$ は、有限な有理基底をもつ。すなわち、この系の有限個の関数が存在し、系の任意の関数はそれら有限個を用いて $\mathsf{P}$ の係数で有理的に表される。

第二の定式化：$\mathfrak K$ を $\mathsf{P}$ と $\mathsf{P}(x_1,\ldots,x_n)$ の中間体とし、ここで $x_1,\ldots,x_n$ は不定元を表すものとし、$\mathfrak K$ の超越次数を $t\leq n$ とする。$y_1,\ldots,y_t$ を $\mathfrak K$ の既約系\footnote{Steinitz に従い、代数的に独立な関数からなる系を既約という。ある系が超越次数 $t$ をもつとは、$t$ 個の元からなる既約系を少なくとも一つ含むが、それより多くの元からなる既約系を含まないことをいう。以下で用いる既約系に関する簡単な定理については Steinitz, § 22 を参照。}とする。このとき $\mathfrak K$ は $\mathsf{P}(y_1,\ldots,y_t)$ の有限代数拡大となる。

二つの定式化は実際には同値である。すべての中間体について有理基底の存在を証明すれば十分である。というのは、系 $S$ から導かれる体はそのような中間体 $\mathfrak K$ となり、その有理基底は直ちに $S$ の有理基底を与えるからである。実際 $\mathfrak K$ の任意の元は $S$ の有限個の関数の有理結合になる。他方、第二の定式化からは、$\mathfrak K$ の任意の既約系 $y_1,\ldots,y_t$ に、$\mathfrak K$ を $\mathsf{P}(y_1,\ldots,y_t)$ の有限拡大として特徴づける有限個の関数を加えれば、そのような有理基底が得られる。逆に、第一の定式化により $\mathfrak K$ の有理基底が与えられているなら、超越次数の定義により各基底元は $\mathsf{P}(y_1,\ldots,y_t)$ に代数的に従属しなければならない。したがって $\mathfrak K$ は $\mathsf{P}(y_1,\ldots,y_t)$ の有限代数拡大として特徴づけられる。

定理は第二の定式化で証明する。まず $\mathfrak K$ が $P(x_1,\ldots,x_n)$ と同じ超越次数 $n$ をもつとする。$y_1,\ldots,y_n$ を $\mathfrak K$ の既約系、したがって $\mathsf{P}(x_1,\ldots,x_n)$ の既約系とする。$y_1,\ldots,y_n$ の超越次数が $n$ であるため、各 $x_i$ は $\mathsf{P}(y_1,\ldots,y_n)$ 上代数的である。よって $\mathsf{P}(x_1,\ldots,x_n)$、したがって中間体としての $\mathfrak K$ は、$\mathsf{P}(y_1,\ldots,y_n)$ の有限代数拡大となる。

次に $t<n$ とする。$y_1,\ldots,y_t,x_{t+1},\ldots,x_n$ が $\mathsf{P}(x_1,\ldots,x_n)$ の既約系となるように $x$ の番号づけを選ぶ。すると $x_{t+1},\ldots,x_n$ は $\mathsf{P}(y_1,\ldots,y_t)$ に関して既約であり、代数的依存の推移律により、$\mathfrak K$ および $\mathsf{P}(y_1,\ldots,y_t)$ と $\mathfrak K$ の間の任意の中間体 $\mathfrak M$ に関しても既約である。これは、$\overline{\mathfrak K}=\mathfrak K(x_{t+1},\ldots,x_n)$ および $\overline{\mathfrak M}=\mathfrak M(x_{t+1},\ldots,x_n)$ が、それぞれ $\mathfrak K$ および $\mathfrak M$ の純超越拡大になることを意味する。したがって $\overline{\mathfrak K}$ のうち $\mathfrak K$ に属さない任意の元は $\mathfrak K$ に関して超越的であり、同様に $\overline{\mathfrak M}$ のうち $\mathfrak M$ に属さない任意の元は $\mathfrak M$ に関して超越的である。同じく $\overline{\mathsf{P}}$ を純超越拡大 $\mathsf{P}(x_{t+1},\ldots,x_n)$ とおく。すると $y_1,\ldots,y_t$ は $\overline{\mathsf{P}}$ に関して既約である。

ここで $\overline{\mathsf{P}}$ を係数領域として基礎に置くことにより、証明はすでに済んだ場合に帰着される。$\mathsf{P}<\mathfrak K< \mathsf{P}(x_1,\ldots,x_t;\,x_{t+1},\ldots,x_n)$\footnote{$\mathsf{P}<\mathfrak K$ は、$\mathsf{P}$ が $\mathfrak K$ に含まれる、すなわち $\mathfrak K$ の部分体である、という意味である。} から、$\overline{\mathsf{P}}<\overline{\mathfrak K}<\overline{\mathsf{P}}(x_1,\ldots,x_t)$ が従い、ここで $\overline{\mathfrak K}$ は $\overline{\mathsf{P}}$ に関して $\overline{\mathsf{P}}(x_1,\ldots,x_t)$ と同じ超越次数 $t$ をもつ。したがって $\overline{\mathfrak K}$ は、係数領域を $\overline{\mathsf{P}}$ として、有限な有理基底、例えば $y_1,\ldots,y_t; z_1,\ldots,z_s$ をもつ。ここで $\overline{\mathfrak K}$ は $\mathfrak K$ と $\overline{\mathsf{P}}$ の合成体であるから、$z$ は $\mathfrak K$ から選べる。示すべきことは、$\mathfrak K$ が $\mathfrak M=\mathsf{P}(y_1,\ldots,y_t;z_1,\ldots,z_s)$ に等しいことである。いま $\mathfrak M$ は $\mathsf{P}(y_1,\ldots,y_t)$ と $\mathfrak K$ の中間体である。したがって $\mathfrak K$ の任意の元は $\mathfrak M$ に関して代数的である。一方、$\overline{\mathfrak M}$ は $\overline{\mathfrak K}$ に等しく、したがって $\mathfrak K$ は $\overline{\mathfrak M}$ に含まれる。ゆえに $\mathfrak K$ は $\mathfrak M$ と $\overline{\mathfrak M}$ の中間体である。しかし $\overline{\mathfrak M}$ のうち $\mathfrak M$ に属さない任意の元は $\mathfrak M$ に関して超越的であるから、必然的に $\mathfrak K$ は $\mathfrak M$ と同一である。

系：$\mathsf{P}<\mathfrak K<\mathfrak L<\mathsf{P}(x_1,\ldots,x_n)$ かつ $\mathfrak L$ が $\mathfrak K$ に関して代数的なら、$\mathfrak L$ は $\mathfrak K$ の有限代数拡大である。実際、$\mathfrak L$ の有理基底の各元は $\mathfrak K$ に関して代数的であるから、$\mathfrak L$ は $\mathfrak K$ の有限拡大として特徴づけられる。

2. \textbf{有限性基準の証明}。条件が必要であることは明らかである。実際、$\mathfrak J$ が有限な整域なら、$\mathfrak J$ 自身を基準に現れる有限な部分環 $\mathfrak R$ とみなせる。

条件は、係数体が標数 $p$ の場合に根体 $P^{1/p}$ が係数体の有限拡大\footnote{p. 28, 注2を参照。}であるという仮定の下で十分である。標数零については制限仮定は必要ない。これを証明するには、$\mathfrak J$ が $\mathfrak R$ のある部分環に関して有限な加群基底をもつことを示せばよい。

$\mathfrak K$ を $\mathfrak R$ の商体\footnote{p. 28, 注2を参照。}、$\mathfrak L$ を $\mathfrak J$ の商体とする。これらの商体は、関係する環が零因子をもたないため存在する。そして $\mathsf{P}<\mathfrak K<\mathfrak L<\mathsf{P}(x_1,\ldots,x_n)$ が成り立つ。したがって 1. の系により、$\mathfrak L$ は $\mathfrak K$ の有限代数拡大体となる。さらに仮定により $\mathfrak J$ の各元は $\mathfrak R$ に関して整であり、したがって $\mathfrak J$ は、$\mathfrak L$ のうち $\mathfrak R$ 上整なすべての元からなる環 $\mathfrak S$ の部分環となる。しかし $\mathfrak R$ が $\mathfrak K$ の中で整閉であるとは何も仮定していないので、有限拡大で約元鎖定理が保たれるという事実を直接用いることはできない。むしろまず、$\mathfrak R$ の中のやはり有限な部分環 $\mathfrak T$ であって、この整閉性を満たすものへ移らなければならない。

まず係数領域 $\mathsf{P}$ が無限個の元を含むと仮定する。このとき $\mathfrak R$ から既約系 $g_1(x),\ldots,g_t(x)$ を、$\mathfrak R$、したがって $\mathfrak J$ が $\mathfrak T=\mathsf{P}[g_1(x),\ldots,g_t(x)]$ 上整となるように選べる。すなわち、仮定により存在する $\mathfrak R$ の整基底を $f_1(x),\ldots,f_k(x)$ とし、これが既約系をなさないなら、必要なら添字を付け替えたうえで、$f_k$ が満たす関係 $F(f_k;f_1,\ldots,f_{k-1})=0$ を取る。$f_1,\ldots,f_{k-1}$ を
\[
 f_1+t_1f_k,\ldots,f_{k-1}+t_{k-1}f_k;
\]
で置き換えると、$t_i$ を不定元として選ぶ場合、$f_k$、したがって $f_1,\ldots,f_{k-1}$ も、これら新しい関数の上で整に依存する。$\mathsf{P}$ は無限個の元をもつので、この整依存性が保たれるように $t_i$ を $\mathsf{P}$ の元へ特殊化できる。整依存の推移律を用い、有限回の手順の後、求める環 $\mathfrak T=\mathsf{P}[g_1(x),\ldots,g_t(x)]$\footnote{Hilbert, Über die vollen Invariantensysteme, § 1. Math. Ann. 42 (1893), pp. 313--373 を参照。} に到達する。すると $\mathfrak L$ は $\mathfrak T$ の商体 $\mathsf{P}(g_1(x),\ldots,g_t(x))$ の有限拡大体となり、$\mathfrak L$ のうち $\mathfrak R$ 上整なすべての元からなる環 $\mathfrak S'$ は、$\mathfrak L$ のうち $\mathfrak T$ 上整なすべての元からなる環と同一になる。

$\mathfrak T$ においては、すべてのイデアルの系について約元鎖定理が成り立ち、また $\mathfrak T$ はその商体の中で整閉である。いずれも $\mathfrak T$ が $t$ 個の不定元の多項式領域に同型だからである。さらに $\mathsf{P}^{1/p}$ が $\mathsf{P}$ に関して有限であるという仮定から、$\mathfrak T^{1/p}$ は $\mathfrak T$ に関して有限であることが従う。\footnote{序論で引用した Artin と van der Waerden のノートを参照。} したがって $\mathfrak S'$ においては、すべての $\mathfrak T$-加群の系について約元鎖定理が成り立つ。特に、$\mathfrak T$ を含む $\mathfrak S'$ の部分環としての $\mathfrak S$ は、有限な $\mathfrak T$-加群基底をもつ。$\mathfrak S$ に対するこの有限な $\mathfrak T$-加群基底の存在は、$\mathfrak L$ が $\mathsf{P}(g_1(x),\ldots,g_t(x))$ の第一種拡大であるとき、係数領域 $\mathsf{P}$ についていかなる制限仮定もなしに成り立ち続ける。したがって特に $P$ が標数零である場合には常にそうである。\footnote{例えば E. Noether, Abstrakter Aufbau der Idealtheorie in algebraischen Zahl- und Funktionenkörpern, § 3. Math. Ann. 96 (1926), pp. 26--61 を参照。}

$h_1(x),\ldots,h_s(x)$ を $\mathfrak S$ の $\mathfrak T$-加群基底とする。このとき $\mathfrak J$ の任意の $f(x)$ についての表示
\[
 f(x)=A_1(g)h_1(x)+\cdots+A_s(g)h_s(x)
\]
は、ここで $A$ は $\mathfrak T$ の元として $g$ の多項式を表すので、$g_1(x),\ldots,g_t(x),h_1(x),\ldots,h_s(x)$ が求める整基底をなすことを示している。

体 $\mathsf{P}$ が有限個の元しか含まない場合、ひとつの不定元 $u$、すなわち $\mathfrak J$ に関して超越的な元を添加して得られる体 $\mathsf{P}(u)$ は無限個の元を含む。したがって環 $\mathfrak J(u)$ は、$\mathsf{P}(u)$ を係数領域として有限な整基底をもつ。そして $\mathfrak J(u)$ は $\mathfrak J$ と $\mathsf{P}(u)$ の合併環であるから、この基底は $u$ の冪で分けることにより $\mathfrak J$ から選べる。ただしこのとき $\mathfrak S$ の基底元 $g_i$ はもはや既約系をなさないであろう。したがって $\mathfrak T$ の任意の多項式 $f(x)$ は表示 $C(u)f(x)=G(u,g_i,h_i)$ をもつ。ここで $C(u)$ は $\mathsf{P}[u]$ の零でない多項式を表し、$G$ は $u,g_i,h_i$ で整となる。適当な $u$ の冪の係数を比較すると、$g_i(x)$ と $h_i(x)$ が $\mathfrak J$ の整基底を与えることがわかる。これで\textbf{有限性基準は一般に証明された}。

\subsection*{§ 2. 不変式の有限性}

1. \textbf{標数 $p$ の有限線型群}とは、標数 $p$ の体 $P$ を係数とする、$h$ 個の線型置換 $A_1,\ldots,A_h$（行列式は零でない）からなる群 $\mathfrak H$ をいう。ここで $A_k$ は線型置換
\[
 x_1^{(k)}=a_{11}^{(k)}x_1+\cdots+a_{1n}^{(k)}x_n,
 \ldots,
 x_n^{(k)}=a_{n1}^{(k)}x_1+\cdots+a_{nn}^{(k)}x_n,
\]
または略して $(x^{(k)})=A_k(x)$ によって表される。したがって群 $\mathfrak H$ は、元 $x_1,\ldots,x_n$ をもつ行 $(x)$ を、元 $x_1^{(k)},\ldots,x_n^{(k)}$ をもつ行 $(x^{(k)})$ へ移す。$A_1,\ldots,A_h$ の中には恒等置換が含まれなければならないので、行 $(x)$ も行 $(x^{(k)})$ の中に含まれる。

群の\textbf{整有理（絶対）不変式}とは、$x_1,\ldots,x_n$ の整有理関数で、$\mathsf{P}$ の係数をもち\footnote{これは一般性の制限ではない。標数 $p$ の体の係数をもつ任意の不変式は、不変式を定義するために $\mathsf{P}$ と共通の包摂体をもたなければならないが、$\mathsf{P}$ の係数をもつ有限個の不変式の線型結合に分解できるからである。一般の場合には、係数を $\mathsf{P}$ に関して一次独立なものを用いて表せばよい。すると不変性の性質はこれら独立な係数に関して恒等的に満たされる。}、$A_1,\ldots,A_h$ を適用しても恒等的に変わらないものをいう。したがって、行 $(x^{(k)})$ の、$P$ の係数をもつすべての対称関数はこれらの不変式に属する。逆に、このような任意の不変式に対して
\[
 f(x)=f(x^{(1)})=\cdots=f(x^{(h)})
 \quad\text{または}\quad
 h\cdot f(x)=f(x^{(1)})+\cdots+f(x^{(h)})
\]
が成り立つ。ここから、群の位数 $h$ が標数で割り切れない場合、特に標数零の場合には、行 $(x^{(k)})$ の $\mathsf{P}$ の係数をもつ対称関数が不変式全体を尽くすことが従う。これら対称関数は有限な整基底をもつので、この場合には有限性の証明が完了する。同時に、整基底は Galois リゾルベントの係数の系 $G_{\alpha\alpha_1\ldots\alpha_n}(x)$ として得られる。\footnote{p. 28, 注1を参照。}
\[
 \Phi(z,u)=\prod_{k=1}^{h}\bigl(z-u_1x_1^{(k)}-\cdots-u_nx_n^{(k)}\bigr)
\]
\[
 =z^h+\sum G_{\alpha\alpha_1\ldots\alpha_n}(x)
 z^{\alpha}u_1^{\alpha_1}\ldots u_n^{\alpha_n}
 \quad(\alpha\ne h;\ \alpha+\alpha_1+\cdots+\alpha_n=h).
\]

2. $h$ が $p$ で割り切れる場合には、群のすべての不変式が行 $(x^{(k)})$ の対称関数である必要はない。もはや $h\cdot f(x)$ から $f(x)$ へ結論できないからである。しかし、有限個の不変式 $G_{\alpha\alpha_1\ldots\alpha_n}(x)$、すなわち Galois リゾルベントの係数から導かれる環 $\mathfrak R$ は、すべての不変式からなる系 $\mathfrak J$ の有限な部分環をなし、$\mathfrak J$ が $\mathfrak R$ に関して有限性基準の仮定を満たすことを示せばよい。

まず一般性を失うことなく\footnote{p. 33, 注を参照。}、すべての置換係数 $a_{\mu\nu}^{(k)}$ から導かれる $P$ の部分体 $\overline{\mathsf{P}}$ を係数領域として基礎に置いてよい。この体は有限個の元から導かれるので、素体に有限個の代数的または超越的な元を添加して得られる。素体は完全であるから、$\overline{\mathsf{P}}^{1/p}$ は $\overline{\mathsf{P}}$ に関して有限であることが従う。\footnote{p. 32, 注2を参照。}

さらに $\mathfrak J$ が $\mathfrak R$ に代数的かつ整に依存することを示さなければならない。1. により、$h$ 個の行 $x^{(k)}$ の中にはもとの $x_1,\ldots,x_n$ も現れる。したがって Galois リゾルベント $\Phi(z,u)$ は、$z=u_1x_1+\cdots+u_nx_n$ とおくと $u$ に関して恒等的に消える。各場合に、$u_i$ だけを一にし他のすべての $u$ を零に特殊化すると、関係
\[
 x_i^h+\sum x_i^\alpha G_{\alpha00\ldots\alpha_i\ldots00}(x)=0
 \quad (\alpha\ne h;\ \alpha+\alpha_i=h)
\]
が得られる。これは各 $x_i$ が $\mathfrak R$ に代数的かつ整に依存することを示す。すると同じことは $x$ の任意の多項式に対して成り立つ。特に $\mathfrak J$ は $\mathfrak R$ 上整である。これで有限性基準の仮定は満たされ、\textbf{不変式の有限性が証明された}。

3. 同じ推論により、\textbf{相対不変式およびモジュラー不変式の有限性}も得られることを注意しておく。多項式 $f(x)$ が群の\textbf{相対不変式}であるとは、すべての $f(x^{(k)})$ が $f(x)$ と $P$ の零でない因子だけ異なることをいう。したがって絶対不変式の系、特に $G_{\alpha\alpha_1\ldots\alpha_n}(x)$ から導かれる環 $\mathfrak R$ は、すべての相対不変式から導かれる環の有限な部分環をなす。そして上と同様に、$\mathfrak R$ に関して有限性基準の仮定が満たされる。

多項式 $f(x)$ が群の\textbf{モジュラー不変式}であるとは、$P$ からの任意の特殊な値の系 $x_1=\alpha_1,\ldots,x_n=\alpha_n$ に対して $f(x)$ が $f(x^{(k)})$ に等しくなることをいう。任意の絶対不変式は同時にモジュラー不変式である。しかし $P$ が有限個の元しか含まない場合、逆は成り立たないことがある。ここでも、$G_{\alpha\alpha_1\ldots\alpha_n}(x)$ から導かれる部分環 $\mathfrak R$ に関して有限性基準の仮定は満たされる。

\section*{30. 代数的数体および関数体におけるイデアル論の抽象的構成}
\noindent\emph{Math. Ann. 96 (1927), pp. 26--61.}

以下では、代数的数体のすべての整数のイデアル論と一致するイデアル論をもつすべての環、すなわちそのイデアルが素イデアルの冪の積として一意に表される環の抽象的特徴づけを与える。このような環のよく知られた例には、一変数の代数的関数体のすべての整数からなる環も含まれ、さらに一般に、Kronecker に従って最高次元のイデアルに制限する、すなわちある商環である関数領域へ移るなら、複数の不定元の場合も含まれる。

基礎に置く可換環について、このイデアル論と完全に同値な公理は次のものである。\footnote{イデアル論の基本概念の定義については、例えば E. Noether, \emph{Idealtheorie in Ringbereichen}, Math. Ann. 83 (1921), pp. 24--66（以下では \emph{イデアル論}として引用）を参照。なお、本論文、とくに §§1, 4, 5 でも、最も本質的な概念を簡潔に定式化している。完全を期して、本論文の直接の目的には必要でない材料もいくらか含めた。}

\begin{description}
\item[I. 約イデアル鎖条件。] 各イデアルが直前のものの真の約イデアルであるような任意のイデアル鎖は、有限段階の後に停止する。言い換えれば、イデアルの任意の約鎖には、以後すべてのイデアルが等しくなる添字がある。
\item[II. 任意の非零イデアルを法とする倍イデアル鎖条件。] 零イデアルと異なる固定イデアルの約イデアルだけからなり、各イデアルが直前のものの真の倍イデアルであるような任意のイデアル鎖は、有限段階の後に停止する。
\item[III.] 乗法の単位元の存在。
\item[IV.] 零因子をもたない環。
\item[V. 商体における整閉性。] 商体の元で環に関して整なものは、すべてその環に属する。
\end{description}

これらの公理から、望まれるものに至るまで、次第に制限されたイデアル論を段階的に展開する。逆に、これらすべての公理が実際にそこで満たされることを示す。

約イデアル鎖条件 I からは、私が \emph{イデアル論} §4 で示したように、イデアルが、相異なる素イデアルに属する有限個の準素イデアルの最小公倍数として表されることが従う。この証明をここで簡単に想起する（§6）。その理由は、イデアル商の概念を用いることで簡単化できるからであり、また一つの見落としを訂正したいからでもある。元の証明は一箇所で、仮定されていない乗法単位元の存在を用いていた。\footnote{P. Urysohn がこの見落としを私に指摘してくれた。私自身による証明の変更は、ここに与える E. Artin によるものよりやや煩雑であった。彼はすでに以前この欠落を自分で埋めており、また私とは独立に、イデアル商による簡単化にも気づいていた。「簡約表示」の導入を避けるさらなる簡単化は、W. Krull, \emph{Algebraische Theorie der Ringe III}, Math. Ann. 92 (1924), pp. 181--213, 定理 I の関連する問題から取った。}

倍イデアル鎖条件の仮定は（§7）、零イデアルと異なる任意のイデアルを法とする剰余類環において、素イデアルが、全元を含む単位イデアルを除いて、真の約イデアルをもたないという効果をもつ。従って、これらのイデアルの準素成分は、単位イデアルに属するものを除いて一意に定まる。やや弱い形ではこの結果はすでに Masazo Sono に見られる。\footnote{Masazo Sono, \emph{On the Reduction of Ideals}, Memoirs of the College of Science, Kyoto University, Series A, 7 (1924), pp. 191--204.} 彼の、剰余類環に組成列が存在するという仮定は、この環における「二重鎖条件」と同値であることを末尾（§10）で簡単に示す。二重鎖条件とは、零イデアルと異なるイデアルへの制限なしに、約イデアル鎖条件と倍イデアル鎖条件の双方が成り立つことをいう。

さらに公理 III、すなわち単位元の存在が満たされるなら、単位イデアルは随伴素イデアルとして現れない。従って準素成分は一意に定まり、互いに二つずつ互いに素であり、その最小公倍数は積に等しい。公理 I, II, III は代数的数体のすべての有限オーダーに対して成り立つ。すでに Dedekind（Dirichlet--Dedekind, \emph{Zahlentheorie}, 第3版, 第11補遺, §172）は、証明を完全には詳述しないまま、イデアルを互いに素な一種イデアル（準素成分）の積として一意に表すことを述べていた。

仮定 IV、すなわち零因子がないことからは、零イデアルでも単位イデアルでもないイデアルの異なる冪はすべて相異なること、および商体が存在することが従う。最後に、商体における整閉性の公理 V も満たされるなら、IV により一意に定まる準素イデアルは素イデアルの冪となり、通常の因数分解定理が得られる（§8）。この最後の証明は、Dedekind が数体の場合に（第3版, §172）すでに整閉性から引き出した帰結に基づく。すなわち、真に分数的な元は、分子の二乗さえ分母で割り切れないように、整数の商として表せるという帰結である。後の叙述では、この帰結は同じく整閉性の帰結として、はるかに見通しの悪い一般化 Gauss 定理、あるいはそれよりやや強い対応する加群定理によって置き換えられる。これらはいずれも基底元を前提として適用される定理であるが、ここではイデアルそのものを直接扱う。

逆方向（§9）では、公理 V 以外の公理はほとんど直接に従う。V は、真に分数的な元が、分子のいかなる冪も分母で割り切れないように常に表されるという証明から従い、これはもともと V から導かれた帰結を強めるものである。この証明が成功するのは、イデアル表示から通常の帰結を引き出した後に限られる。すなわち、固定イデアルを法とする主イデアルによる表示と、分数イデアルの理論、つまり体におけるイデアル商の理論である。素イデアルの冪による表示から整閉性が従うことは Dedekind にもすでに知られていた。第3版 §172, p. 522 下部の注記がそれを示している。

公理 I から V は、一般には異なる四つの分解、すなわち二つずつ互いに素な最小イデアルへの分解、互いに素なイデアルへの分解、最大準素成分への分解、既約イデアルへの分解が一致することを含意する。逆に、この一致と公理 I--IV から整閉性が従う。零因子をもつ環では、残りの公理を満たしても、V を商環における整閉性として定義した場合でさえ、必ずしもこれらの分解が一致するとは限らない。有限オーダーのあるイデアルを法とする剰余類環の例がこれを示す（§9）。

§1 では、環に関して整な量の理論を、整閉性から Dedekind が得た帰結まで概観する。§2 では、鎖条件が、その環を係数および乗数領域とする線型形式からなる加群へどのように移されるかを示す。\footnote{この形の加群定理は Prüfer（\emph{Neue Begründung der algebraischen Zahlentheorie}, §2, Math. Ann. 94 (1924), pp. 198--243）によるもので、基底定理を移す通常の方法より見通しがよい。} そこから（§3）、第一種の代数的拡大体の整元へ移ると、公理 I から IV は任意のオーダーについて保たれ、すべての整元からなるオーダーでは V も成り立つ、という結論を得る。この事実は、もちろん代数的数体では周知であるが、冒頭で述べた関数体を従属させることを可能にする。特に、複数の不定元の整多項式から導かれる関数領域が公理を満たすことの簡単な証明により、Kronecker のイデアル論は完全に包含される。\footnote{違いは判別式定理にだけ現れる。素数を法とする関数領域の剰余類環が不完全体になり、従って第二種拡大が起こり得るからである。}

後で展開されるイデアル論は、分類だけに用いられる §§2, 3 を前提としない。§4 と §5 では、公理 I から V に依存しない基礎を与える。これにより、理論のどの部分が有限性仮定に依存しないかが、元の基礎づけよりも鋭く浮き出る。

\subsection*{§1. 整量の理論}

可換環\footnote{よく知られているように、環は、集合上に等号関係が与えられ、さらに二つの演算、加法と乗法が与えられて通常の法則を満たすとき定義される。加法は結合的・可換的で一意に可逆であり、乗法は結合的で、環が可換であるときは可換的であり、加法に関して分配的である。基礎となる等号の定義が集合論的同一性でない場合には、すべての部分系で同じ等号の定義が保たれるように、その部分系（部分環、加群、イデアル）は任意の元とともにそれに等しいすべての元を含まなければならない。同様に、すべての拡大では新しい等号の定義が古いものを含むと仮定しなければならない。「有限個の元」「一意の対応」などのすべての概念は、この等号定義の意味で理解される。例えば「一つだけの元がある」とは、等号を除いて一つだけであることを意味する。なお、等しい元を一つの類にまとめ、これらの類を新しい元として導入することで、等号から集合論的同一性へ常に移れる。ここに述べた等号定義に関する注意は R. Hölzer による。} $T$ が与えられているとする。これは零因子をもたず、乗法の単位元 $e$ をもつ（公理 III, IV）。$T$ の中には、単位元を含む固定された部分環 $R$ が区別されている。加群、オーダー、整という概念は、常にこの固定環 $R$ に関するものとし、後には簡単のため記号から省く。\footnote{整量の理論は、零因子を許しても、それに代えて $R$ における約イデアル鎖条件を仮定すれば有効であり、§2 によりこの仮定は補題の証明も可能にする。この形の整量の理論は後で使わないので、注で指摘するにとどめる。ここで問題となる定義はすべて、零因子が存在しても有効である。その多くは、よく知られ容易に見られるように、さらに弱い仮定しか必要としない。保たれないのは Dedekind の「数は包絡をもつとき整である」という定式化である。零因子があると、二つの有限加群の商が有限であるとは限らないからである。したがってここではオーダーを加群商としてでなく直接定義する。} $R$ の元はラテン文字で、$T$ の一般の元はギリシア文字で表す。

1. $T$ の元からなる系 $M$ は、通常どおり、二つの元 $\alpha,\beta$ とともにその差を含み、$\alpha$ とともに $r\alpha$ も含むとき、$R$-加群、略して加群と呼ばれる。ただし $r$ は $R$ の任意の元である。$M$ が $N$ で割り切れるとは
\[
        M\equiv 0 \pmod N
\]
で表し、$M$ を倍加群、$N$ を約加群と呼ぶ。$R$-加群で、そのすべての元が $R$ に属するものを $R$ のイデアルと呼ぶ。

明らかに、$T$ の任意個の加群の共通部分、すなわち最小公倍数 $[\ldots,M_i,\ldots]$ は再び加群である。$T$ 自身も加群である。したがって $T$ の元からなる任意の系 $\Sigma$ から導かれる加群は、$\Sigma$ を含むすべての加群の共通部分として一意に存在する。任意の加群系の和、すなわち最大公約数 $(\ldots,M_i,\ldots)$ は、それらの和集合から導かれる加群である。二つの加群の積 $MN$ は、$\mu$ が $M$ のすべての元を、$\nu$ が $N$ のすべての元を走るときの元 $\mu\nu$ から導かれる加群として定義される。従って積は、環の元について成り立つため、結合律と可換律を満たす。さらに $T$ における分配律から、加群についても分配律
\[
        (\ldots,M_i,\ldots)N=(\ldots,M_iN,\ldots)
\]
が従う。まさに $T$ におけるこの法則によって、両辺はいずれもすべての元 $\mu_i\nu$ から導かれる加群だからである。

$R$ 自身が加群であるから、$R$ が乗数領域であるという条件は $RM\equiv0\pmod M$ と表せる。仮定により $R$ は単位元を含むので、$M\equiv0\pmod{RM}$ も成り立ち、従って $M=RM$ である。

加群 $M$ は、それを導く有限個の元からなる系 $\Sigma$ が少なくとも一つ存在するとき有限と呼ばれる。$\Sigma$ の元 $\mu_1,\ldots,\mu_n$ は
\[
        M=(\mu_1,\ldots,\mu_n)
\]
の加群基底と呼ばれる。二つ、従って有限個の有限加群の和および積は再び有限である。それぞれ基底元の和集合、あるいは基底元の積から導かれるからである。後者の主張は、$M$ のすべての元 $\mu$ が基底によって $r_1\mu_1+\cdots+r_n\mu_n$ と表されること、\footnote{$R$ が単位元を含まない場合、基底表示は $r_1\mu_1+\cdots+r_n\mu_n+n_1\mu_1+\cdots+n_n\mu_n$ の形になる。ここで整数 $n_i$ は環の元ではなく、反復加法または反復減法を表す記号である。}および $T$ における乗法の可換律から従う。

2. $T$ の中にある $R$ の拡大環、すなわち $R$ のすべての元を含む環を $R$-オーダー、略してオーダーと呼ぶ。$T$ における任意個のオーダーの共通部分は再びオーダーであり、$T$ もまたオーダーである。したがって $T$ の元からなる任意の系 $\Sigma$ から導かれるオーダーは一意に存在する。オーダーの和および積は加群の場合と同様に定義され、特に任意のオーダーについて $O^2=O$ である。

すべてのオーダーは同時に $R$-加群である。オーダーは有限加群であるとき有限と呼ばれる。和と積は加群の和と積を作ることで得られるので、1 により有限オーダーの和と積は再び有限オーダーである。さらに次の推移律が成り立つ。$A$ が有限 $R$-オーダーで、$B$ が有限 $A$-オーダーなら、$B$ は有限 $R$-オーダーでもある。実際、$B$ は同時に $R$-オーダーであり、従って $R$-加群である。$\alpha_1,\ldots,\alpha_m$ が $R$ に関する $A$ の加群基底で、$\beta_1,\ldots,\beta_n$ が $A$ に関する $B$ の加群基底なら、積 $\alpha_i\beta_j$ は明らかに $R$ に関する $B$ の加群基底をなす。

3. $\mT$ の元 $\alpha$ は、$\alpha$ から導かれるオーダー $\mR_\alpha$ が有限であるとき、$\mR$ に関して整、略して整と呼ばれる。同値に、加群鎖
\[
        \mA_\mu=(\alpha^0,\alpha,\ldots,\alpha^{\mu-1})
\]
が有限段階の後に停止し、
\[
        \mA_n=\mA_{n+1}=\cdots=\mA_{n+\nu}=\cdots
\]
となるときである。

この定義には、4. で用いる第二の定式化もある。$\mT$ の元 $\alpha$ は、零加群と異なる主加群 $\mC$、すなわちある $\gamma\ne0$ から導かれる $\mC$ が少なくとも一つ存在し、オーダー $\mR_\alpha$ と $\mC$ の積 $\mC\mR_\alpha$ が有限加群になるとき、整と呼ばれる。言い換えれば、加群鎖 $\mC\mA_\mu$ が有限段階の後に停止することである。\footnote{$\mR$ に零因子を許す場合には、正則な主加群の存在を要求しなければならない。すなわち $\gamma$ は零因子でない正則元であると仮定しなければならず、零因子のない環ではこれは単に $\gamma\ne0$ である。}

最後に、この定義は通常の定式化とも同一である。$\mT$ の元 $\alpha$ は、$\mR$ の元 $r_i$ を用いて少なくとも一つの方程式
\[
        \alpha^n+r_1\alpha^{n-1}+\cdots+r_n=0
\]
を満たすとき整と呼ばれる。

これらの定義は実際に同一である。$\mR_\alpha$ は $\alpha$ のすべての冪から導かれる加群と一致するから、$\mR_\alpha$ が有限であるとき、その冪の中から常に加群基底を選ぶことができる。そのような基底が、例えば $e=\alpha^0,\alpha,\ldots,\alpha^{n-1}$ で与えられるなら、これは $\mA_\mu$ の鎖が $\mA_n$ で停止することを意味する。逆に $\mA_n=\mA_{n+\nu}\ldots$ はこの基底を与える。さらに $\mA_\mu$ の鎖の停止は、上の方程式の存在とも同一である。$\mR_\alpha$ が有限であれば $\mC\mR_\alpha$ も有限であり、従って鎖 $\mC\mA_\mu$ は停止する。逆に $\mC\mR_\alpha$ の有限性、したがって $\mC\mA_\mu$ の停止から、$\mA_\mu$ も停止し、従って $\mR_\alpha$ が有限であることが従う。実際、$\mC$ は零加群と異なる主加群であると仮定されているので、
\[
        \mC\mA_n=\mC\mA_{n+\nu}
\]
は $\mA_n=\mA_{n+\nu}$ も与える。

整量の環性と整依存の推移律は、次の補題に基づく。

\medskip
\noindent\textbf{補題。} $\mS$ を $\mT$ の有限オーダーとし、$\alpha$ を $\mS$ の任意の元とする。このとき $\alpha$ から導かれるオーダー $\mR_\alpha$ は有限である。言い換えれば、有限オーダーのすべての元は整である。

証明は通常の議論による。$\sigma_1,\ldots,\sigma_n$ を $\mS$ の加群基底とする。$\mS$ は環であるから、$\alpha\sigma_i$ は $\mS$ に属する。従って
\[
        \alpha\sigma_i=r_{i1}\sigma_1+\cdots+r_{in}\sigma_n,
\]
したがって
\[
        \left|r_{ik}-\alpha e_{ik}\right|=0
\]
である。ただし $i\ne k$ なら $e_{ik}=0$、$e_{ii}=e$ である。これにより $\alpha$ が整であることが示される。行列式の消滅には、$\mT$、従って $\mS$ が零因子をもたないという仮定を用いる。\footnote{$\mR$ に約イデアル鎖条件を仮定しつつ零因子を許す場合、この補題は §2 の加群定理から直ちに従う。その定理によれば鎖条件は $\mS$ からなる加群へ移る。この証明も、数体については Dedekind（第4版, §173, II）に由来し、文献における鎖条件の最初の応用である。4. の下で与える帰結では、Dedekind は代数的数体におけるイデアルを法とする剰余類の数が有限であるという、より特殊な事実を用いている。}

$\mT$ におけるすべての整量からなる系 $\mS$ はオーダーをなす。$\mR$ の元は、単位元を基底とする有限 $\mR$-加群として $\mR$ が有限であるため整であり、したがって $\mS$ は $\mR$ を含む。また $\mS$ は環性ももつ。任意の二つの整量 $\alpha,\beta$ から導かれるオーダー
\[
        \mR_{\alpha,\beta}=\mR_\alpha\cdot\mR_\beta
\]
は有限である。ゆえに補題により、$\alpha\pm\beta$ と $\alpha\beta$ は有限オーダーの元として整である。

\medskip
\noindent\textbf{整依存の推移律。} $\mS$ が整量からなるオーダーであるなら、$\mS$ に関して整な $\mT$ のすべての元は $\mR$ に関しても整である。定義により $\alpha$ は方程式
\[
        \alpha^m+\sigma_1\alpha^{m-1}+\cdots+\sigma_m=0
\]
を満たし、従って $\sigma_1,\ldots,\sigma_m$ から導かれるオーダー
\[
        \overline{\mS}=\mR_{\sigma_1,\ldots,\sigma_m}
        =\mR_{\sigma_1}\cdot\mR_{\sigma_2}\cdots\mR_{\sigma_m}
\]
に関しても整である。このオーダーは有限である。2. で与えた推移律により、$\alpha$ から導かれる有限 $\overline{\mS}$-オーダー $\overline{\mS}_\alpha$ は従って有限 $\mR$-オーダーでもあり、$\alpha$ は有限オーダーの元として補題により整である。
4. 環 $R$ は、$R$ に関して整な $T$ のすべての元が $R$ に属するとき、$T$ において整閉であると呼ばれる。3. の整量定義の第二の定式化によれば、$T$ の元 $\alpha$ が $R$ に属するための必要十分条件は、零加群と異なる主加群 $C$ が少なくとも一つ存在して $CR_\alpha$ が有限加群となることである。\footnote{$R$ に零因子を許す場合、$C$ は再び正則主加群と仮定しなければならない。同様に、帰結 I の元 $c$ も正則と仮定しなければならない。}

整閉環の概念から Dedekind（第3版, §172）は、これをイデアル論に用いることを可能にする二つの重要な帰結を引き出した。

\medskip
\noindent\textbf{Dedekind の帰結 I。} $\mR$ が $\mT$ において整閉であり、さらに $\mR$ でイデアルの約イデアル鎖条件（公理 I）が成り立つとする。$\beta$ が $\mT$ の非整元で、$c\ne0$ が $\mR$ の元なら、指数 $\sigma\geqq0$ が存在し、元の列
\[
 c,\;c\beta,\ldots,c\beta^r,\ldots
\]
において $c,\ldots,c\beta^\sigma$ は整であり、残りはすべて非整である。

もしすべての元 $c\beta^r$ が整であれば、$\mR$ の整閉性によりそれらは $\mR$ に属する。$\mC$ を $c$ から導かれる加群、$\mB_r$ を $\beta^0,\ldots,\beta^{r-1}$ から導かれる加群とすると、加群鎖
\[
        \mC\mB_1,\mC\mB_2,\ldots,\mC\mB_r,\ldots
\]
はイデアル鎖 $\mc_1,\ldots,\mc_r,\ldots$ となる。ただし
\[
        \mc_r=(c,c\beta,\ldots,c\beta^{r-1}).
\]
これは仮定により停止するが、そのとき仮定に反して $\mR_\beta$ は有限となり、$\beta$ は整となる。従って、元 $c\beta^r$ のいくつかは非整である。さらに、一つの非整元があれば、それ以後のすべての元は非整である。実際、$c\beta^\tau$ が整、従って $\mR$ に属するなら、$\rho<\tau$ に対して元 $c\beta^\rho$ は
\[
(c\beta^\rho)^\tau-c^{\,\tau-\rho}(c\beta^\tau)^\rho=0
\]
を満たし、従って整である。ゆえに最初の非整元が $c\beta^{\sigma+1}$ であるような $\sigma\geqq0$ が存在し、これが求める指数である。
拡大環 $T$ を $R$ の商体、すなわちすべての元の対を添加して得られる体\footnote{Steinitz のよく知られた定式化、J. f. M. 137 において。$R$ に零因子を許す場合には、分母が $R$ の正則元である商対すべてを添加して、商体の代わりに商環を用いなければならない。この場合、Dedekind の帰結 II はもはや有効でない。$n$ は一般に正則元ではないからである。}に特殊化すると、次が得られる。

\medskip
\noindent\textbf{Dedekind の帰結 II。} $R$ がその商体 $K$ の中で整閉であり、$R$ でイデアルの約イデアル鎖条件が成り立つとする。$\beta$ が $K$ の非整元なら、$\beta$ は $R$ の元の商
\[
        \beta=m/n
\]
として表され、しかも $m^2/n$ も非整となるようにできる。

$\beta=b/c$ と置き、$c\ne0$ とする。このとき列 $c,c\beta=b,c\beta^2\ldots$ の最初の二つの元は確かに整であるから、帰結 I の指数を $\sigma$ とすれば $\sigma\geqq1$ である。そこで
\[
        m=c\beta^\sigma,\qquad n=c\beta^{\sigma-1}
\]
と置けば、$m/n=\beta$ であり、$m^2/n=c\beta^{\sigma+1}$ は非整である。

整閉性の推移律は次の形をとる。$R$ を $T$ の任意の環とし、$G$ を $R$ に関して整な $T$ のすべての量の系とする。このとき $G$ は、$G$ に関して整な $T$ のすべての量からもなる。実際、$G$ に関して整な任意の量は、3. の推移律により $R$ に関しても整であり、従って $G$ に属する。

\subsection*{§2. 有限加群領域における鎖条件}

鎖条件を移す加群定理は、ここでの応用では拡大環の加群についてだけ現れる。証明は一般の加群領域の場合でも同じであるから、そのような領域を出発点とする。

元 $\alpha,\beta,\ldots$ からなる系 $M$ は、環 $R$ に関して、$M$ の元へ一意に導く二つの演算、すなわち元の加法的結合と $R$ の元による乗法が与えられ、$M$ が加法について Abel 群であり、乗法が結合律を満たし、分配律が二つの形で成り立つとき、加群領域と呼ばれる。\footnote{\emph{イデアル論}, §9 を参照。}

加群領域については、§1, 1. で与えた加群定義のうち、乗法に関係しないものはすべて明らかに有効である。特に、$\mathsf{M}$ は、$\mathsf{M}$ の有限個の元 $\xi_1,\ldots,\xi_k$ が存在し、$\mathsf{M}$ が $\xi_1,\ldots,\xi_k$ から導かれる加群、従って $R$ が単位元をもつ場合にはすべての線型形式 $r_1\xi_1+\ldots+r_k\xi_k$ の系に等しいとき、有限加群領域と呼ばれる。そうでない場合には追加項 $n_i\xi_i$ が現れる。

\medskip
\noindent\textbf{加群定理。} $M$ を、単位元をもつ可換環 $R$ に関する有限加群領域とし、$R$ においてイデアルに対する約イデアル鎖条件、または倍イデアル鎖条件が成り立つとする。このとき $M$ では、加群に対する約加群鎖条件、または倍加群鎖条件が成り立つ。

証明は、$M$ の各加群に対して $R$ の有限個のイデアルからなる系を一意に対応させ、加群の真の約加群または真の倍加群に対して、そのたびに少なくとも一つの真の約イデアルまたは倍イデアルが対応するようにすることに基づく。

$\mathsf{M}$ の元が長さ $i$ をもつとは、それが少なくとも一つの仕方で $\xi_1,\ldots,\xi_i$ の線型形式として書けるが、$\xi_1,\ldots,\xi_{i-1}$ の線型形式としては表されないことをいう。$M$ の任意の加群 $\mA$ に、$\mR$ の $k$ 個のイデアル $\mathfrak a_1,\ldots,\mathfrak a_k$ を割り当てる。$\mA_i$ を $\mA$ のうち長さ $\le i$ のすべての元からなる加群とし、$\mathfrak a_i$ を $\mA_i$ における $\xi_i$ のすべての係数の系とする。まず次を得る。

$\mB$ が $\mA$ の真の約加群であり、$\mathfrak b_1\ldots,\mathfrak b_k$ が $\mB$ に対応するイデアルであるなら、$\mathfrak b$ の中には対応する $\mathfrak a$ の真の約イデアルが少なくとも一つある。$\mA\equiv0(\mB)$ から $\mA_\lambda\equiv0(\mB_\lambda)$、従って $\lambda=1,2,\ldots,k$ に対して $\mathfrak a_\lambda\equiv0(\mathfrak b_\lambda)$ が従う。仮定により $\mB$ には $\mA$ に含まれない元があり、その中には最小の長さ $i$ をもつものがある。このとき $\mathfrak b_i$ は $\mathfrak a_i$ の真の約イデアルである。実際、
\[
        \beta=b_1\xi_1+\cdots+b_i\xi_i
\]
が $\mB$ におけるそのような最小長の元なら、$b_i\equiv0(\mathfrak b_i)$ だが $b_i\not\equiv0(\mathfrak a_i)$ である。もし $b_i\equiv0(\mathfrak a_i)$ なら、元
\[
        \alpha=a_1\xi_1+\cdots+b_i\xi_i\equiv0(\mA)
\]
が存在し、従って $\beta-\alpha\equiv0(\mB)$、$\not\equiv0(\mA)$ で、$\beta$ より長さの短い元が存在することになる。

いま約鎖または倍鎖
\[
        \mA_1,\mA_2,\ldots,\mA_r,\ldots
\]
が与えられ、対応する鎖条件が $\mR$ で成り立つと仮定する。$\mA_r$ に割り当てられたイデアルを $\mathfrak a_{r1},\ldots,\mathfrak a_{rk}$ とすれば、各 $\lambda=1,\ldots,k$ について列
\[
        \mathfrak a_{1\lambda},\mathfrak a_{2\lambda},\ldots,
        \mathfrak a_{r\lambda},\ldots
\]
は約イデアル鎖または倍イデアル鎖である。したがって仮定により、すべての $\lambda$ について $\mathfrak a_{\mu\lambda}$ が以後のすべてのイデアルと等しくなる添字 $\mu$ が存在する。すると、いま証明したことにより、加群 $\mA_\mu$ は以後のすべての加群と等しい。こうして鎖条件は移される。

これから直ちに次が得られる。

\medskip
\noindent\textbf{加群定理の帰結。} $R$ において、零イデアルと異なる任意のイデアルを法として倍イデアル鎖条件が成り立つなら、$M$ では、割り当てられたイデアル $\mathfrak c_1,\ldots,\mathfrak c_k$ がすべて零イデアルと異なる任意の加群 $C$ を法として倍加群鎖条件が成り立つ。これらの仮定の下で、倍イデアル鎖 $\mathfrak a_{i1},\ldots,\mathfrak a_{i\nu},\ldots$ は有限段階の後に停止する。

\medskip
\noindent\textbf{補足。} $R$ が単位元をもたない環である場合にも、零イデアルと異なるイデアルを法とする約イデアル鎖条件および倍イデアル鎖条件は移される。すなわち、$M$ の代わりに、形
\[
        a_1\xi_1+\cdots+a_k\xi_k+n_1\xi_{k+1}+\cdots+n_k\xi_{2k}
\]
のすべての元からなる系 $M'$ を考える。この系は、追加の整倍 $n_i\xi_i$ を $M$ の元で置き換えると再び $M$ へ移る。この $M'$ に、上と対応して $2k$ 個のイデアルを割り当てることができる。ここで $\mathfrak a_i$ は $R$ のイデアルであり、$\mathfrak n_i$ は有理整数のイデアルである。$\mathfrak n_i$ については、約イデアル鎖条件と任意の非零イデアルを法とする倍イデアル鎖条件の双方が満たされるので、約イデアル鎖条件の証明は $M'$ に対して、従って $M$ に対しても有効である。しかし倍イデアル鎖条件については、$C$ あるいは $C'$ に割り当てられる $2k$ 個のイデアルがすべて零イデアルと異なることを要求しなければならない。

\subsection*{§3. 有限拡大体への移行}

数体と関数体の分類のためには、序で定式化した公理 I から V が、有限拡大体へ移る際にどのように移されるかを示す必要が残る。

1. \emph{$\mR$ において、公理 II、すなわち倍イデアル鎖条件を除くすべての公理 I から V が満たされているとする。$\mK$ を $\mR$ の商体とし、$\mL$ を $\mK$ 上の第一種有限拡大体とする。\footnote{Steinitz に従い、代数的拡大体は、そのすべての元が、適当な拡大体で相異なる線型因子に分解する素関数の零点であるとき、「第一種」と呼ばれる。} このとき、$\mR$ に関して整な $\mL$ のすべての元からなる系 $\mS$ では同じ公理が満たされ、$\mS$ の任意のオーダーでも整閉性を除くすべての公理がなお満たされる。}

§1, 3 により、$\mS$ は公理 III, IV、すなわち単位元の存在と零因子の不存在を満たす環をなす。§1, 4 の結論により、$\mS$ は $\mL$ において整閉である。同時に、$\mL$ は $\mS$ の商体になる。$\mL$ の任意の元は、$\mR$ の適当な元を掛けることで $\mS$ の元に移されるからである。したがって公理 V は満たされる。

I の証明はよく知られた議論に従う。第一種拡大として、$\mL$ は $\mK$ に一つの元 $\alpha$ を添加して得られ、この $\alpha$ は上により $\mR$ に関して整と仮定できる。$\alpha$ のほか、$\mK$ 上の $\mL$ のガロア体に含まれる共役量 $\alpha',\alpha'',\ldots$ も整である。したがってそれらの差の積も整であり、その差の積の二乗 $D$ も整である。$\alpha,\alpha',\ldots$ はすべて異なるので、$D$ は $\mK$ の非零量であり、$\mK$ における $\mR$ の整閉性により $\mR$ の元である。通常の議論により、$\mR$ は整閉であるから、$\mS$ は元 $\xi_i=\alpha^i/D$ によって生成される有限 $\mR$-加群領域 $\mathsf{M}$ に含まれる。このためには、$\mS$ の元を $\xi_i$ によって表した表示において、共役元へ移り、得られた線型方程式系を表示の係数について解けばよい。§2 の加群定理により、$\mathsf{M}$ のすべての $\mR$-加群について約鎖条件が成り立ち、特に $\mS$ のすべてのイデアルについて成り立つ。$\mS$ の任意のオーダーのイデアルについても同様に約鎖条件が成り立つ。ここでも $M$ の中の $\mR$-加群を扱っているからである。さらに任意のオーダーについて、公理 III と IV は満たされる。

2. \emph{$\mR$ においてすべての公理 I から V が満たされるなら、$\mS$ においても同じことが成り立つ。$\mS$ の任意のオーダーでも、整閉性を除いてすべての公理がなお成り立つ。}

1 を踏まえれば、証明すべきは公理 II だけである。§2 の加群定理の系により、次を示せば十分である。$\mS$ の非零イデアルを $M$ の中の $\mR$-加群 $\mC$ と見たとき、$\mC$ に対応する $\mR$ のイデアル $\mc_i$ はすべて零イデアルと異なる。いま $\mC$ は、$\mK$ 上で $\mathsf{M}$ と同じ有限線型階数をもつ。任意の元 $\gamma$ とともに、$\mC$ は $\gamma\alpha^i$ も含むからである。従って各 $\xi_i$ について、$c_i\xi_i$ が $\mC$ に属するような $\mR$ の元 $c_i\ne0$ が存在する。$\xi_i$ による表示の一意性により、ここで添字は体の次数より小さい値だけを走るものとするが、$c_i$ は $\mc_i$ に属し、従って $\mc_i$ は零イデアルではない。この議論は $\mathsf{M}$ と同じ階数のすべてのオーダーに対して有効である。階数の小さいオーダーについてはその商体へ移る。この商体は $\mK$ と $\mL$ の間の中間体としてやはり第一種であり、$\mK$ 上の次数はオーダーの線型階数と一致する。従って同じ議論が保たれる。最後に、$\mS$ と異なる $\mS$ のオーダーは決して整閉ではないことを注意しておく。どのオーダーも $\mR$ を含むので、整閉性はそのオーダーに $\mS$ を含むことを強制するからである。

数体と関数体を従属させるには、したがって基礎領域、すなわち有理整数、一変数多項式、および複数不定元多項式の関数領域について、公理 I から V を検証すれば十分である。

3. \emph{有理整数の環 $\mR$、または体係数の一変数多項式の環については、公理 I から V が満たされる。} ここで I と II は、非零数または一変数の非零多項式を法とすると、剰余類が有限個、または線型独立な剰余類が有限個しかないという事実からも、またすべてのイデアルが主イデアルであるという事実からも従う。このことから、イデアルの素イデアル冪の積としての一意分解が従い、従って非零イデアルは有限個のイデアルによってしか割り切られない。公理 III と IV は明らかに満たされ、後者は等号の定義の帰結である。整閉性 V は、$\mR$ の元が単元、すなわち 1 の約元を除いて既約元へ一意に分解されることから従う。$\mR$ に含まれない商体の任意の元は、$\mR$ の二つの互いに素な元の商として書ける。従って分子のいかなる冪も分母で割り切られない。このような元は、$\mR$ に関して整であることを特徴づける方程式を満たし得ない。

4. \emph{いま $\mR$ を、体係数または有理整数係数の複数不定元の全多項式の環とする。} このとき公理 III, IV, V は 3 と同様に満たされる。ここでも元の既約元への因数分解は単元を除いて一意だからである。約イデアル鎖条件 I は、イデアル基底の存在に関する Hilbert の定理の直接の帰結であるが、倍イデアル鎖条件 II はここでは成り立たない。

しかし、\emph{関数領域}へ移る、すなわち最高次元のイデアルに制限することによって、公理 II の有効性も得ることができる。その場合、I の有効性も Hilbert の定理を用いずに 3 と同様に証明できる。実際、$\mR$ に不定元 $u$ を添加し、$u$ の全多項式で係数を $\mR$ にもつ環 $\mR^*$ を考える。等号は係数ごとに定義される。通常どおり、$\mR^*$ の多項式は、その係数の最大公約数、ただしイデアル約元でなく $\mR$ の多項式としての約元、が $\mR$ の単元であるとき、$\mR$ に関して原始的と呼ぶ。このとき $\mR^*$ のすべての多項式は、$\mR$ の元と $\mR^*$ の原始多項式の積であり、$\mR^*$ の原始多項式の積は再び原始である。

\emph{従って、$\mR^*$ の商体の元のうち、分母が $\mR^*$ の原始多項式であるもの全体は一つの環、すなわち $\mR$ の関数領域 $\mF$ をなす。} この関数領域 $\mF$ では、単元はちょうど、分子と分母がともに $\mR^*$ の原始多項式である元である。従って $\mF$ の各元は、$\mR$ の元と $\mF$ の単元の積として一意に表される。ゆえに $\mR$ の元に対する因数分解定理は $\mF$ の元へ移る。従って $\mF$ では、公理 III と IV に加えて、公理 V も 3 と同様に満たされる。

さらに $\mF$ ではすべてのイデアルが主イデアルである。なぜなら、イデアルは $\mF$ の任意の元とともに、それに単元を掛けて得られる $\mR$ の元も含み、また $\mR$ の任意の二つの元とともにその最大公約数も含むからである。実際、$f(x)=t(x)\overline f(x)$ および $g(x)=t(x)\overline g(x)$ とすれば、イデアルは $t(x)(\overline f(x)+u\overline g(x))$ も含み、$\overline f$ と $\overline g$ が互いに素でその線型結合が従って単元であるなら、$t(x)$ も含む。したがって、イデアルの任意の元 $f(x)$ から始めて、もしイデアルの中にそれで割り切れない元 $g(x)$ があれば、最大公約数 $t(x)$ もイデアルに属し、$f(x)$ の真の約元である。この操作を有限回繰り返すとイデアルの基底元に到達し、そのイデアルが主イデアルであることが分かる。従って $\mF$ では、$\mR$ の元の因数分解に対応して、イデアルの素イデアル冪の積としての一意分解が成り立つ。3 と同様に、公理 I と II が満たされることが従う。ここで考慮すべき $\mF$ の商体の拡大体は、$\mS$ の元を添加して生成できるものだけである。\footnote{これらの拡大体において $\mF$ に関して整な量の系は、$\mR$ から $\mF$ へ移ったのと同じように、$\mS$ から対応する関数領域へ移ることによっても得られる。J. König, \emph{Algebraische Größen}, Leipzig 1903, pp. 468/69 を参照。} 1 と 2 を考慮すれば、数体および関数体について公理 I から V の有効性が確立される。

\subsection*{§4. 同型定理。直和}

以下では、環性または加群性だけを仮定し、それ以外の公理は前提としない。

1. $\mathsf{M}$ と $\overline{\mathsf{M}}$ が $\mR$ に関する加群領域（§2）であるとき、$M$ が $\overline M$ に準同型（より正確には加群準同型）である、$M\sim\overline M$、とは、$M$ の各元に $\overline M$ の一つかつただ一つの元が対応し、$\overline M$ が尽くされ、\footnote{$\overline M$ で有効な等号定義の意味において。注9を参照。} かつこの対応の下で差と同じ $\mR$ の元による乗法が対応することをいう。従って $\beta\sim\overline\beta$ および $\gamma\sim\overline\gamma$ からは常に $(\beta-\gamma)\sim(\overline\beta-\overline\gamma)$ および $r\beta\sim r\overline\beta$ が従う。

従って $M$ の中の $\mR$-加群 $\mB$ には、$\overline M$ の中の $\mR$-加群 $\overline{\mB}$ が準同型的に対応する。また、$\overline{\mathsf{M}}$ の $\mR$-加群 $\overline{\mC}$ の元に対応する $M$ の元全体 $\mC^*$ は、$\overline{\mC}$ によって一意に定まる $M$ の $\mR$-加群、すなわち $\overline{\mC}$ に随伴する加群 $\mC^*$ をなし、$\mC^*\sim\overline{\mC}$ である。特に $\mA$ が $\overline{\mathsf{M}}$ の零元に随伴する $M$ の加群であるなら、$\mC^*$ は $\mA$ の約加群であり、$\mA$ を法として合同な $\mathsf{M}$ の元の類に分かれる。それらの類は $\mC$ の元と一対一に対応する。$M$ の中の $\mB$ から $\overline M$ の中の $\overline{\mB}$ へ移り、さらに戻って $\mB^*$ へ行くと、$\mB^*=(\mB,\mA)$、すなわち $\mB$ と $\mA$ の最大公約数を得る。従って $\mB$ が $\mA$ の約加群であれば $\mB^*=\mB$ である。

$\mathsf{M}$ と $\overline{\mathsf{M}}$ の元の対応が可逆的一対一であるなら、これらの領域は同型（より正確には加群同型）であると呼ばれ、$\mathsf{M}\simeq\overline{\mathsf{M}}$ と書く。

$\mA$ が $M$ の任意の $\mR$-加群であるなら、$\mA$ を法とする合同を新しい等号関係とすることにより、$\mathsf{M}$ に準同型な加群領域 $\overline{\mathsf{M}}$、すなわち剰余類加群 $\mathsf{M}\mid\mA$ が生じる。$M$ の各元には、$\overline M$ においてそれに等しいすべての、かつそれらだけの元が随伴する。$\overline M$ において等号定義から同一性へ移るとは、$\overline M$ で等しいすべての元を一つの類、すなわち剰余類にまとめ、これらの剰余類を $\overline M$ の新しい元として取ることである。\footnote{これについては注9を参照。}

\emph{すべての準同型は剰余類加群への移行によって生じる。} 実際、$\mathsf{M}\sim\overline{\overline{\mathsf{M}}}$ であり、$\mA$ が $\overline{\overline{\mathsf{M}}}$ の零元に随伴する加群なら、上で示したように $\overline{\overline{\mathsf{M}}}$ は剰余類加群 $\mathsf{M}\mid\mA$ と同型である。

2. \textbf{第一同型定理。} \emph{$\overline{\mathsf{M}}$ を剰余類加群 $\mathsf{M}\mid\mA$ とし、$\mC$ を $\mA$ の約加群とする。このとき同型 $\overline{\mathsf{M}}\mid\overline{\mC}\simeq \mathsf{M}\mid\mC$ がある。} 実際、$\mA$ を法とする合同は同時に $\mC$ を法とする合同であり、$\mA$ を法として等しい元は $\mC$ を法としても等しいままである。したがって、剰余類加群 $\mathsf{M}\mid\mC$、すなわち $\mC$ を法とする等号は、まず $\mA$ を法として元を同一視し、$\overline{\mathsf{M}}$ へ移り、ついで $\overline M$ の中で $\mC$ を法として等しい元をまとめることによって形成できる。これは $\overline{\mathsf{M}}$ で $\overline{\mC}$ を法として同一視すること、すなわち $\overline{\mathsf{M}}\mid\overline{\mC}$ を作ることと同じである。

\textbf{第二同型定理。} \emph{$\mB$ と $\mA$ が $\mathsf{M}$ の加群であるなら、$(\mB,\mA)\mid\mA\simeq\mB\mid[\mB,\mA]$ である。} なぜなら 1 により、$(\mB,\mA)\mid\mA$ を $\overline{\mB}$ と置くと、$\mB$ は $\overline{\mB}$ に準同型となる。そして $\overline{\mB}$ の零元に対応するのは正確に $[\mB,\mA]$ の元なので、主張した同型は再び 1 から従う。

3. $\mR$ と $\overline{\mR}$ が（可換）環であるとき、$\mR$ が $\overline{\mR}$ に準同型（より正確には環準同型）である、$\mR\sim\overline{\mR}$、とは、$\mR$ の各元に $\overline{\mR}$ の一つかつただ一つの元が対応して $\overline{\mR}$ が尽くされ、この対応の下で差と積が対応することをいう。元の対応が可逆的一対一であるなら、これらの環は同型（より正確には環同型）と呼ばれ、$\mR\simeq\overline{\mR}$ と書く。

1 と 2 のすべての議論は、加群を $\mR$ のイデアルに置き換え、\footnote{非可換環では両側イデアルを用いなければならない。} 加群準同型および加群同型を環準同型および環同型に置き換えても有効である。特に、$\mc$ が $\ma$ の約イデアルなら、剰余類環 $\mc\mid\ma$ は、$\ma$ を法とする合同を新しい等号関係として導入することで得られる。ここでは、剰余類の積も定義されることに注意しなければならない。

従って $\mc$ は $\mc\mid\ma$ に準同型である。すべての準同型はこのようにして生じ、同型定理は 2 と同様に成り立つ。

\textbf{第一同型定理。} $\overline{\mR}$ が剰余類環 $\mR\mid\ma$ を表し、$\mc$ が $\ma$ の約イデアルであるなら、環同型の意味で $\overline{\mR}\mid\overline{\mc}\simeq\mR\mid\mc$ である。

\textbf{第二同型定理。} $\mb$ と $\ma$ が $\mR$ のイデアルであるなら、環同型 $(\mb,\ma)\mid\ma\simeq\mb\mid[\mb,\ma]$ が成り立つ。\footnote{群論でよく知られているこれらの同型定理は、加群についてはやや特殊な形で最初に Dedekind に現れる（第3版, p. 484 を参照）。上記のイデアルに対する形は Masazo Sono, \emph{On Congruences. I, II, III, IV}, Memoirs of the College of Science, Kyoto Imperial University, 2 (1917), 3 (1918, 1919), 2, p. 215 参照、に現れる。いずれの場合にも第二同型定理だけが明示的に述べられている。}

4. ここで、5 の直和に関する定理が従う互いに素なイデアルの計算法則\footnote{Dedekind（第4版, §178, III から VIII）に遡る形である。最近の叙述で一貫して現れる単位元を用いる計算は、はるかに煩雑である。}をまとめる。可換環 $\mR$ において、乗法単位元 $e$ の存在を仮定する。従って $\mR$ の任意のイデアルについて $\ma=\mo\ma$ であり、$\mo$ は単位イデアルを表す。二つのイデアル $\ma,\mb$ は、その最大公約数 $(\ma,\mb)$ が $\mo$ であるとき、互いに素と呼ばれる。

\emph{4$\alpha$. $(\ma,\mb)=\mo$ なら $(\ma,\mb\mc)=(\ma,\mc)$ である。} 実際、$(\ma,\mc)=(\ma,\mb)\cdot(\ma,\mc)=(\ma^2,\ma\mb,\ma\mc,\mb\mc)$ は $(\ma,\mb\mc)$ で割り切られ、逆も成り立つ。従って次が成り立つ。

\emph{4$\beta$. $\mc\mb\equiv0(\ma)$ かつ $(\ma,\mb)=\mo$ なら、$\mc\equiv0(\ma)$ である。} 実際、$\mc\mb\equiv0(\ma)$ は $(\ma,\mb\mc)=\ma$ と同値であり、$(\ma,\mb)=\mo$ であるから、$(\ma,\mc)=\ma$、すなわち $\mc\equiv0(\ma)$ も得られる。さらに次が成り立つ。

\emph{4$\gamma$. イデアル $\ma_1,\ldots,\ma_r$ の各々がイデアル $\mb_1,\ldots,\mb_s$ の各々と互いに素であるなら、$\ma_i$ の積は $\mb_i$ の積と互いに素である。} $(\ma,\mb)=\mo$ および $(\ma,\mc)=\mo$ から $(\ma,\mb\mc)=\mo$ が得られ、有限回繰り返せば主張が従う。

4$\alpha$ と 4$\gamma$ から次が従う。

\emph{4$\delta$. イデアル $\mb_1,\ldots,\mb_r$ が二つずつ互いに素、すなわち $i\ne k$ に対して $(\mb_i,\mb_k)=\mo$ であるなら、その最小公倍数は積に等しい。} $(\ma,\mb)=\mo$ とし、$\mv=[\ma,\mb]$ と置く。すると $\mv=\mo\mv=(\ma\mv,\mb\mv)\equiv0(\ma,\mb)$ である。$\ma\mb\equiv0(\mv)$ だから、$\mv=\ma\mb$ が従う。4$\gamma$ を用いて有限回繰り返せば主張が得られる。

\emph{4$\varepsilon$. イデアル $\mb_1,\ldots,\mb_r$ が二つずつ互いに素であり、$\ma_i=[\mb_1,\ldots,\mb_{i-1},\mb_{i+1},\ldots,\mb_r]=\mb_1\cdots\mb_{i-1}\mb_{i+1}\cdots\mb_r$ と置くなら、$(\ma_1,\ldots,\ma_{i-1},\ma_{i+1},\ldots,\ma_r)=\mb_i$ であり、従って $(\ma_1,\ldots,\ma_r)=\mo$ である。} 実際、分配律により $(\ma_1,\ma_2)=\mb_3\cdots\mb_r(\mb_2,\mb_1)=\mb_3\cdots\mb_r$ である。従って $(\ma_1,\ma_2,\ma_3)=\mb_4\cdots\mb_r(\mb_3,\mb_1\mb_2)=\mb_4\cdots\mb_r$ であり、適切な番号づけの下で有限回繰り返せば主張が証明される。

イデアル $\ma_i$ は表示 $\mm=[\mb_1,\ldots,\mb_r]$ における $\mb_i$ の補イデアルと呼ばれる。

5. 再び、可換環 $\mR$ が乗法単位元をもつとする。$\mR$ の各元 $c$ が、$a_i$ を $\ma_i$ の元として、ただ一通りに $c=a_1+a_2+\cdots+a_r$ の形で表されるとき、$\mR$ はイデアル $\ma_1,\ldots,\ma_r$ の直和であるといい、記号で $\mR=\ma_1+\ma_2+\ldots+\ma_r$ と書く。

\emph{零イデアルが二つずつ互いに素なイデアル $\mb_1,\ldots,\mb_r$ の最小公倍数であり、$\ma_i$ が $\mb_i$ の補イデアルであるなら、$\mR$ はイデアル $\ma_1,\ldots,\ma_r$ の直和である。} $\mo=(\ma_1,\ldots,\ma_r)$ であるから、$\mR$ の任意の元は少なくとも一つの $\ma_i$ による加法表示をもつ。$[\ma_i,\mb_i]=(0)$ であるから、この表示は一意である。$0=a_1+a_2+\ldots+a_r=a_i+b_i$ から $a_i=0$ が得られる。第二同型定理により、$\ma_i$ は剰余類環 $\mR\mid\mb_i$ と同型である。直交関係 $i\ne k$ に対して $\ma_i\ma_k=(0)$、および $\ma_i^2=\ma_i$ が成り立つ。後者は $\ma_i=\mo\ma_i$ だからである。

\emph{逆に、$\mR$ の直和表示 $\mR=\ma_1+\ma_2+\ldots+\ma_r$ があり、$\mb_i=(\ma_1,\ldots,\ma_{i-1},\ma_{i+1},\ldots,\ma_r)=\ma_1+\ldots+\ma_{i-1}+\ma_{i+1}+\ldots+\ma_r$ と置くなら、$\mb_i$ は二つずつ互いに素であり、その最小公倍数は零イデアルである。} 実際、$\mo=(\ma_i,\mb_i)$ から、$k\ne i$ について $\mo=(\mb_k,\mb_i)$ が従う。$\mb_k$ は $\ma_i$ の約イデアルだからである。さらに $c$ がすべての $\mb_i$ で割り切れるなら、
\[
 c=a_2^{(1)}+\cdots+a_r^{(1)}=a_1^{(2)}+a_3^{(2)}+\cdots+a_r^{(2)}=\cdots=a_1^{(r)}+\cdots+a_{r-1}^{(r)},
\]
である。ただし各場合に $a_i^{(\lambda)}\equiv0\pmod{\ma_i}$ である。表示の一意性により、各表示で一つの成分が零であるから、すべての $\lambda$ について $0=a_1^{(\lambda)},\ldots,0=a_r^{(\lambda)}$ が従い、従って $c=0$ である。\footnote{5 に述べた定理は、直和と直交共通部分との関係に関する一般定理の特殊な場合である。H. Prüfer, \emph{Theorie der Abelschen Gruppen I}, Math. Zeitschr. 20 (1924), pp. 165--187, §6 を参照。そこに与えられた議論は、加群とイデアルが特殊な場合として含まれるほど一般的に群概念を定式化する場合にも有効である。}

\subsection*{§5. 素イデアルと準素イデアル}

再び $\mR$ を可換環とするが、ここではさらなる公理は仮定せず、単位元の存在さえ仮定しない。素イデアルと準素イデアルについての基本事実をまとめる。\footnote{\emph{イデアル論}, §4 では公理 I、すなわち約イデアル鎖条件を仮定している。そのため、素イデアルおよび準素イデアルのどの性質がこの公理に依存しないかは、そこでは鋭く切り分けられていない。}

1. $\mR$ のイデアル $\mpideal$ は、剰余類環 $\mR\mid\mpideal$ が零因子をもたない環であるとき、\emph{弱素イデアル}と呼ばれる。すなわち、$\ma\not\equiv0\pmod{\mpideal}$ および $\mb\not\equiv0\pmod{\mpideal}$ から常に $\ma\mb\not\equiv0\pmod{\mpideal}$ が従う場合である。$\mR$ のイデアル $\mpideal^*$ は、剰余類環 $\mR\mid\mpideal^*$ が零因子イデアルをもたない環であるとき、\emph{強素イデアル}と呼ばれる。すなわち、$\ma\not\equiv0\pmod{\mpideal^*}$ および $\mb\not\equiv0\pmod{\mpideal^*}$ から常に $\ma\mb\not\equiv0\pmod{\mpideal^*}$ が従う場合である。

\emph{すべての強素イデアルは同時に弱素イデアルである}。これは $\ma,\mb$ を主イデアルに特殊化すれば分かる。逆に、\emph{すべての弱素イデアルも強素イデアルである}。従って単に素イデアルといってよい。実際、$\mpideal$ を弱素イデアルとし、$\ma\not\equiv0\pmod{\mpideal}$ および $\mb\not\equiv0\pmod{\mpideal}$ とする。$\ma$ と $\mb$ の中から $a\not\equiv0\pmod{\mpideal}$ および $b\not\equiv0\pmod{\mpideal}$ を満たす元 $a,b$ を選ぶ。すると $ab\not\equiv0\pmod{\mpideal}$ であり、従って $\ma\mb\not\equiv0\pmod{\mpideal}$ であるから、$\mpideal$ は強素イデアルでもある。

2. $\mR$ のイデアル $\mq$ は、剰余類環 $\mR\mid\mq$ において、すべての零因子のある冪が消えるとき、\emph{弱準素イデアル}と呼ばれる。同値に、$a\not\equiv0\pmod{\mq}$ かつすべての $\varkappa$ について $b^\varkappa\not\equiv0\pmod{\mq}$ なら常に $ab\not\equiv0\pmod{\mq}$ である。$\mR\mid\mq$ で零因子となる $\mR$ のすべての元の系 $\mpideal$ はイデアルをなし、しかも素イデアルである。それは $\mq$ の約イデアルであり、随伴素イデアルと呼ばれる。実際、$a$ とともに $ra$ も零因子であり、$a$ と $b$ とともに $a-b$ も零因子である。$a^\varkappa$ と $b^\lambda$ があれば、$(a-b)^{\varkappa+\lambda}$ は常に $\mq$ で割り切られるからである。さらに $a$ が零因子でないなら、定義により $a$ のいかなる冪も零因子ではない。$a\not\equiv0\pmod{\mpideal}$ と $b\not\equiv0\pmod{\mpideal}$ から $a^\varkappa b^\varkappa=(ab)^\varkappa\not\equiv0\pmod{\mq}$ が従い、従って $ab\not\equiv0\pmod{\mpideal}$ である。

$\mR$ のイデアル $\mq^*$ は、剰余類環 $\mR\mid\mq^*$ において、すべての零因子イデアルのある冪が消えるとき、\emph{強準素イデアル}と呼ばれる。すなわち、$\ma\not\equiv0\pmod{\mq^*}$ かつすべての $\varkappa$ について $\mb^\varkappa\not\equiv0\pmod{\mq^*}$ なら、常に $\ma\mb\not\equiv0\pmod{\mq^*}$ である。弱準素イデアルの場合と同様に、$\mR\mid\mq^*$ で零因子イデアルとなる $\mR$ のすべてのイデアルの最大公約数は、素イデアルであり、$\mq^*$ の約イデアルであり、随伴素イデアルをなすことが示される。逆に、同じ随伴素イデアル $\mpideal$ をもつすべての準素イデアルは、$\mpideal$ に属すると言われる。

\emph{すべての強準素イデアルは同時に弱準素イデアルである}。これは $\ma,\mb$ を主イデアルに特殊化すれば分かる。しかし一般には逆は偽である。\footnote{R. Hölzer が私に伝えた次の例がこれを示す。$\mR$ を体係数の可算個の不定元 $x_i$ の多項式領域とし、$\mq=(x_1^2,x_2^3,\ldots,x_\nu^{\nu+1},\ldots,x_i x_k\,[i\ne k]\ldots)$ とする。剰余類環における零因子はちょうど $\mpideal=(x_1,x_2,\ldots,x_\nu,\ldots)$ で割り切れる多項式であり、そのような各零因子のある冪は $\mq$ で割り切れる。従って $\mq$ は随伴素イデアル $\mpideal$ をもつ弱準素イデアルである。しかし $\mq$ は強準素イデアルではない。$\ma=(x_1,x_3,x_5,\ldots,x_{2\nu+1},\ldots)$ および $\mb=(x_2,x_4,\ldots,x_{2\nu},\ldots)$ と置くと、$\ma\mb\equiv0(\mq)$ であるが、$\ma$ または $\mb$ のどの冪も $\mq$ で割り切れない。} \emph{しかし $\mR$ に公理 I、すなわち約イデアル鎖条件を仮定すれば、すべての弱準素イデアルは強準素イデアルでもある。} 従って単に準素イデアルといってよい。$\mq$ を弱準素イデアルとし、$\ma\not\equiv0\pmod{\mq}$ かつすべての $\varkappa$ について $\mb^\varkappa\not\equiv0\pmod{\mq}$ とする。このとき、$a\not\equiv0\pmod{\mq}$ かつすべての $\varkappa$ について $b^\varkappa\not\equiv0\pmod{\mq}$ を満たす $\ma$ の元 $a$ と $\mb$ の元 $b$ が存在する。従って $ab\not\equiv0\pmod{\mq}$ であり、ゆえに $\ma\mb\not\equiv0\pmod{\mq}$ である。もし $\mb$ の各 $b$ について、ある指数 $\lambda$ が存在して $b^\lambda\equiv0\pmod{\mq}$ となるなら、有限個の基底元の指数を $\lambda_1,\ldots,\lambda_r$ として、$\mb^{\lambda_1+\cdots+\lambda_r}\equiv0\pmod{\mq}$ となるからである。\footnote{公理 I から有限イデアル基底を導くには、$\mR$ に固定された整列を用いなければならない。§6 を参照。} 同じ議論は、随伴素イデアルを $\mpideal$ とするとき、$\mpideal^\varrho\equiv0\pmod{\mq}$ となる最小指数 $\varrho$ が存在することも示す。この $\varrho$ は $\mq$ の指数と呼ばれる。

3. 以下では、約イデアル鎖条件を再び\emph{仮定しない}。有限個のイデアルの最小公倍数 $[\ma_1,\ldots,\ma_r]$ は、どの $\ma_i$ も残りのイデアルの最小公倍数に含まれない、すなわちどの $\ma_i$ も省けないとき、\emph{最短表示}と呼ばれる。

\emph{同じ素イデアル $\mpideal$ に属する有限個の弱準素イデアル、または強準素イデアルの最小公倍数は、再び $\mpideal$ を随伴素イデアルとする弱準素イデアルである。相異なる素イデアルに属する有限個の弱準素イデアル、または強準素イデアルによる最短表示は、弱準素イデアル、または強準素イデアルではない。} $\mf=[\mq_1,\ldots,\mq_r]$ とし、$\mpideal$ を弱準素イデアル $\mq_i$ の随伴素イデアルとする。このとき $\mpideal$ は、ある冪が $\mf$ で割り切れる元全体からも正確になる。$a\not\equiv0\pmod{\mf}$ かつすべての $\varkappa$ について $b^\varkappa\not\equiv0\pmod{\mf}$ なら、$b\not\equiv0\pmod{\mpideal}$ であり、また少なくとも一つの添字 $i$ について $a\not\equiv0\pmod{\mq_i}$ である。従って $ab\not\equiv0\pmod{\mq_i}$ であり、ゆえに $ab\not\equiv0\pmod{\mf}$ である。全体を通して元をイデアルに置き換えると、もはや $\mb\not\equiv0\pmod{\mpideal}$ を推論できない。

いま $\mm=[\mq_1,\ldots,\mq_r]$ を、$i\ne k$ なら $\mpideal_i\ne\mpideal_k$ である弱準素イデアルによる最短表示とし、$r\ge2$ とする。このとき少なくとも一つの随伴素イデアル、例えば $\mpideal_1$ は、他のどれにも含まれない。$\mpideal$ たちの任意の真の倍鎖は高々 $r$ 段で停止しなければならないからである。従って $\mpideal_i$ の元 $a_i$ で、$a_i^{\varrho_i}\equiv0\pmod{\mq_i}$ であり、かつ $a_2\not\equiv0\pmod{\mpideal_1},\ldots,a_r\not\equiv0\pmod{\mpideal_1}$ となるものが存在する。さらに $q_1\not\equiv0\pmod{\mm}$ を $\mq_1$ の元とすると、$q_1a_2^{\varrho_2}\cdots a_r^{\varrho_r}\equiv0\pmod{\mm}$ であるが、$a_2^{\varrho_2}\cdots a_r^{\varrho_r}$ のどの冪も $\mm$ で割り切れない。さもなければ $\mpideal_1$ による可除性が従うからである。従って $\mm$ は弱準素イデアルではない。全体を通して元をイデアルに置き換える、すなわち $q_1$ を $\mq_1$ に、$a_i$ を $\ma_i[\not\equiv0\pmod{\mpideal_1}$ だが $\ma_i^{\varrho_i}\equiv0\pmod{\mq_i}]$ に置き換えれば、強準素イデアルについての証明が得られる。\footnote{約イデアル鎖条件を不要にするこの私の元の証明の修正は、B. L. van der Waerden に負っている。}

4. \emph{加群商とイデアル商。} $\mT$ を $\mR$ の拡大環とし、$\mA,\mB,\mC,\ldots$ を $\mT$ の中の $\mR$-加群とする。商 $\mC=\mA:\mB$ は、条件 $\mC\mB\equiv0(\mA)$ と、もし $\mD\mB\equiv0(\mA)$ なら $\mD\equiv0(\mC)$ である、という条件を満たす加群として定義される。すなわち $\mC$ は、$\mC\mB\equiv0\pmod{\mA}$ が成り立つ最小の加群である。これは商を、$\mD\mB\equiv0\pmod{\mA}$ を満たすすべての加群 $\mD$ の最大公約数として一意に定める。そのような $\mD$ は存在する。零加群が確かにその一つだからである。

定義から次が従う。$\mB$ が $\overline{\mB}$ の倍加群であるなら、$\mA:\overline{\mB}$ は $\mA:\mB$ の倍加群である。$\mA$ が $\overline{\mA}$ の倍加群であるなら、$\mA:\mB$ は $\overline{\mA}:\mB$ の倍加群である。また
\[
        \mA:\mB\mC=(\mA:\mB):\mC=(\mA:\mC):\mB .
\]
拡大環 $\mT$ が $\mR$ と一致すれば、加群商はイデアル商 $\ma:\mb$ になる。従って上の規則もここで成り立つ。$\ma\mb\equiv0(\ma)$ であるから、$\ma\equiv0(\ma:\mb)$ でもある。

5. \emph{$\ma:\mb=\ma$ なら、$\mb$ は $\ma$ に対して素であると言う。} 従ってイデアル $\mb$ が $\ma$ に対して素であるための必要十分条件は、$\md\mb\equiv0(\ma)$ が常に $\md\equiv0(\ma)$ を含意することである。4 の計算法則から次が従う。$\mb$ と $\mc$ が $\ma$ に対して素なら、それらの積および最小公倍数も $\ma$ に対して素である。$\mb$ が $\ma$ に対して素で、$\ma$ が $\mb$ に対して素なら、$\ma$ と $\mb$ は互いに素と呼ばれる。§4, 4$\beta$ から、互いに素なイデアルは互いに素（相互素）でもあることが従う。

\setcounter{footnote}{25}

\subsection*{§6. 約イデアル鎖条件を仮定した場合のイデアル論}

基礎として、公理 I、すなわち約イデアル鎖条件を満たす可換環 $\mR$ を取る。以下ではさらに、$\mR$ の全元に一つの整列順序が固定されているものとする。これにより $\mR$ の全イデアルにも整列順序が定まる。実際、二つの仮定から、$\mR$ の各イデアルは有限個の元からなるイデアル基底をもつことが直ちに従う。$\mR$ の元の整列順序から、$\mR$ のすべての有限部分集合の準辞書式整列順序へ移り、
\footnote{これは次のように理解すべきである。有限部分集合 $\mathfrak W$ の元 $a_1,\ldots,a_n$ を、添字の順序が $\mR$ で指定された整列順序と一致するように書くとき、$n$ 個の元をもつ任意の部分集合は、$m>n$ 個の元をもつ部分集合に先行する。二つの部分集合が同じ個数の元をもつ場合は、整列されたリストが初めて異なる位置で、一方の該当元が他方の該当元に先行するとき、前者が後者に先行する。従って、有限部分集合の任意の系は第一元をもつ。$\mR$ の整列順序が与えられれば、それ以上の選択公理は不要である。とくに、代数的数体の場合のように $^{9)}$ が可算であれば、選択公理は関与しない。イデアル論における選択公理の役割は、詳述なしに \emph{イデアル論} 注 5 で述べておいた。}
各イデアルに対し、その可能な基底のうち、この有限部分集合の整列順序で第一のものを標準基底として割り当てるなら、有限部分集合の整列順序とともにイデアルの整列順序が得られる。

\textbf{定理 I.} \emph{約イデアル鎖条件の仮定の下で、$\mR$ の任意のイデアルは、有限個の既約イデアル、すなわち二つの真の約イデアルの最小公倍数として表せないイデアルの最小公倍数として表される。}

証明には次を示せば十分である。あるイデアル $\mm$ について定理 I が成り立たないなら、$\mm$ は、定理 I が同様に成り立たない真の約イデアルをもつ。このことから、約イデアル鎖条件に反して、有限段階で停止しない約イデアル鎖を構成できる。実際、$\mm$ は可約でなければならない。そうでなければ $\mm=[\mm]$ がすでに望む表示を与えるからである。$\mm=[\ma,\mb]$ とすると、定理 I が $\ma$ と $\mb$ の双方に同時に成り立つことはできない。もし成り立てば、$\mm$ についても対応する表示が従うからである。従って、定理 I が成り立たない $\mm$ の真の約イデアルが存在する。そのうち、イデアルの整列順序で第一のものを $\ma_1$ とする。同様に $\ma_1$ から真の約イデアル $\ma_2$ を構成し、一般に $\ma_i$ から真の約イデアル $\ma_{i+1}$ を構成すると、停止しない、しかもよく定まった約イデアル鎖が得られ、仮定に矛盾する。
\footnote{定理 I は、全加群の系に対して約加群鎖条件を仮定すれば、加群領域についてもまったく同様に成り立つ。この定理は純粋に集合論的な性格をもつ。同時に、証明においてイデアルの整列順序を用いる唯一の箇所でもある。抽象的には次の通りである。集合 $M$ において、部分集合の族 $\Sigma$ が区別され、整列されているとする。$\Sigma$-集合について鎖条件を仮定する。すなわち、$X_i$ が $X_{i-1}$ の真の上集合であるような鎖 $X_1,X_2,\ldots$ は必ず停止する。$\Sigma$-集合が、いずれも真の上集合である二つの $\Sigma$-集合の交わりであるとき可約といい、そうでないとき既約という。このとき任意の $\Sigma$-集合は、有限個の既約 $\Sigma$-集合の交わりである。}

約イデアル鎖条件を仮定しているので、§5, 2 により、単に準素イデアルと言うことができる。準素性と既約性の関係は次で与えられる。

\textbf{定理 II.} \emph{約イデアル鎖条件の仮定の下で、任意の既約イデアルは準素である。言い換えれば、準素でない任意のイデアルは可約である。}

剰余類環 $\mR/\mm$ へ移っても、約イデアル鎖条件は準同型により保たれる。$\mm$ の最小公倍数としての表示は、零イデアルの表示に対応し、また第一同型定理（§4, 3）により、$\mm$ の準素な約イデアルは $\mR\mid\mm$ の準素イデアルに対応し、逆も成り立つ。従って、剰余類環における零イデアルの分解を考えれば十分である。この環も同じ一般性をもつ環であるから、初めから分解すべきイデアルを $\mR$ の零イデアルと仮定してよい。

そこで、$\mR$ の零イデアルが準素でないと仮定する。少なくとも一対のイデアル $\ma,\mb$ が存在し、しかも $\mb$ は主イデアルと仮定してよく、次を満たす。
\[
a\ne(0),\qquad \mb^\varkappa\ne(0)\quad\text{任意の }x\text{ について},
\qquad\text{しかし}\qquad \ma\mb=(0).
\]
イデアル商の鎖
\[
\ma,\quad \ma:\mb,\quad \ldots,\quad \ma:\mb^\nu,\quad\ldots
\]
を作る。この鎖は停止する。例えば
\[
\mt=\ma:\mb^{m-1}=\ma:\mb^m=\cdots
\]
とする。従って $\mt=\mt:\mb$、すなわち $\mb$ は $\mt$ に対して素である。さらに $\mt$ は $\ma$ の約イデアルであって零イデアルと異なり、$\mb^\varkappa$ は任意の指数 $x$ について零と異なる。従って定理 II の証明には
\[
        (0)=[\mt,\mb^{m+1}]
\]
を示せば十分である。定義により
\[
        \mb^{m-1}\mt\equiv0\pmod{\ma},\qquad\text{従って}\qquad \mb^m\mt=(0),
\]
であるから、残るのは
\[
        \mc=[\mt,\mb^{m+1}]\equiv0\pmod{\mb^m\mt}
\]
を示すことだけである。これは、$\mb$ が主イデアルであり、$\mt$ に対して素であるという仮定から従う。$\mc$ の任意の元 $c$ は、$\mb^{m+1}$ で割り切れるので、$n$ を整数記号として
\[
        c=k b^{m+1}+n b^{m+1}=r b^m
\]
という表示をもつ。ただし $r$ は再び $\mR$ の元である。$c\equiv0\pmod{\mt}$ から $r b^m\equiv0\pmod{\mt}$ が従い、$\mt:\mb=\mt$ であるから $r\equiv0\pmod{\mt}$ が従う。よって $c\equiv0\pmod{\mt\mb^m}$ であり、定理 II は証明された。

\textbf{定理 III.} \emph{約イデアル鎖条件の仮定の下で、任意のイデアルは、相異なる素イデアルに属する有限個の最大準素成分の最小公倍数としての最短表示をもつ。}

このような表示を得るには、定理 I によって存在する既約イデアルによる表示、従って定理 II により準素イデアルによる表示を、可能な限り最短表示に置き換えればよい。同じ素イデアルに属する準素イデアルをまとめると、§5, 3 により主張された表示が得られる。これらの準素成分は最大と呼んでよい。なぜなら、それらのうち任意のものの最小公倍数はもはや準素ではないからである。
\footnote{対応する素イデアルと孤立成分は一意に定まる。\emph{イデアル論}、または W. Krull, ``Ein neuer Beweis für die Hauptsätze der allgemeinen Idealtheorie'', Math. Ann. 90 (1923), pp. 55--64 を参照。§7 では、そこに現れる単純な特殊場合について、一意性の直接証明を与える。そこに一意なものとして認識される準素イデアルは孤立成分である。}

\subsection*{§7. 二重鎖条件を仮定した場合のイデアル論}

これまで展開された理論に比べて得られる簡単化は、1 と 2 に掲げる補助結果に基づく。

1. \emph{零因子をもたない可換環で倍イデアル鎖条件が成り立つなら、その環はすでに体である。} 示すべきことは、$a\ne0$ に対して方程式 $ax=b$ が常に環の中に解をもつことである。解が高々一つであることは零因子がないことから従う。

$\ma$ を $a$ で生成される主イデアルとする。仮定により冪の列は停止する。$\ma^m$ が以後すべての冪に等しいとする。$\mb$ が $b$ で生成される主イデアルを表すなら、
\[
        \ma^m\mb=\ma^{m+1}\mb,
        \qquad \ma^{m+1}\mb=(a^{m+1}b).
\]
特に元 $a^m b$ は、ここでも $n$ を整数記号として、
\[
        a^m b=k a^{m+1}b+n a^{m+1}b=a^{m+1}c
\]
という表示をもつ。ただし $c$ は $\mR$ の元である。零因子がないので、これから $b=ac$ が従い、環が体であることが証明される。

2. 1 から直ちに次の補助結果が従う。\emph{倍イデアル鎖条件が可換環で成り立つなら、素イデアルは $\mo$ と異なる真の約イデアルをもたない。}

同様に、次の補足も得られる。\emph{公理 II だけ、すなわち零イデアルと異なる任意のイデアルを法とする倍イデアル鎖条件だけが満たされているなら、零イデアルと異なる任意の素イデアルは $\mo$ と異なる真の約イデアルをもたない。}

素イデアル $\mpideal$ を法とする剰余類環は、定義により零因子をもたない環である。倍イデアル鎖条件もそこで準同型により保たれるから、1 によりそれは体である。$\mpideal$ の約イデアルと剰余類環のイデアルとの一対一対応から、$\mpideal$ は $\mo$ 以外の真の約イデアルをもたない。
\footnote{$\mR$ に単位元が存在する必要はなく、従って $\mR$ の全元からなるイデアル $\mo$ に単位元が存在する必要もない。ただし 1 により、剰余類環には単位類が存在する。この種の最も単純な例は、補足の弱い場合では、すべての偶整数からなる系によって与えられる。}

\textbf{定理 IV.} \emph{可換環で二重鎖条件が成り立つなら、任意のイデアルの最短表示において、単位イデアル $\mo$ に属さない準素成分は一意に定まる。}

\emph{補足。} 公理 I と II の仮定の下で、定理 IV は零イデアルと異なる任意のイデアルについて成り立つ。

\[
        \mm=[\mq,\mq_1,\ldots,\mq_r]=[\mbar\mq,\mbar\mq_1,\ldots,\mbar\mq_{\mbar r}]
\]
を、最大準素成分による $\mm$ の最短表示とし、その随伴素イデアルを
\[
        \mo,\mpideal_1,\ldots,\mpideal_r,
        \qquad \mo,\mbar\mpideal_1,\ldots,\mbar\mpideal_{\mbar r}
\]
とする。$\mo$ に属する準素イデアルが現れないときは、成分 $\mq$、それぞれ $\mbar\mq$ を省くものとする。なぜなら
\[
        \mo^\varrho\mpideal_1^{\varrho_1}\cdots\mpideal_r^{\varrho_r}\equiv0(\mm)
\]
であるから、各 $\mbar\mpideal_i$ はある $\mpideal_i$ に含まれ、補助結果によりこれと同一である。同様に各 $\mpideal_k$ はある $\mbar\mpideal_k$ と一致するので、$\mo$ と異なる素イデアルは一致し、$\mpideal_\lambda=\mbar\mpideal_\lambda$ となる。さらに
\[
        \mq\mq_1\ldots\mq_r\equiv0(\mbar\mq_\lambda)
\]
から $\mq_\lambda\equiv0(\mbar\mq_\lambda)$ が従う。残りの成分は $\mpideal_\lambda$ で割り切れず、従って $\mbar\mq_\lambda$ に対して素だからである。同様に $\mbar\mq_\lambda\equiv0(\mq_\lambda)$ も従うので、$\mo$ に属さない準素成分は一意である。最後に $\mq$ が実際に現れるなら、最短表示であることから $\mbar\mq$ も実際に現れなければならず、随伴素イデアルの一意性も得られる。さらにこの証明は、$\mo$ に属する準素成分を含まない任意の表示が同時に最短であることも示している。

3. 二重鎖条件にさらに単位元の存在を加えると、定理 IV は次の定理 V へと強められる。

\textbf{定理 V.} \emph{二重鎖条件を満たす単位元をもつ可換環では、任意のイデアルが、有限個の二つずつ互いに素な準素イデアルの積として一意に表される。}

\emph{補足。} 公理 I, II, III の下で、定理 V は零イデアルと異なる任意のイデアルについて成り立つ。

単位元が存在するので $\mo^2=\mo$ である。従って $\mo$ に属する任意の準素イデアルは $\mo$ に等しく、$\mo$ と異なるイデアルの最短表示には現れ得ない。従って定理 IV により準素成分は一意に定まり、同時に任意の表示は $\mo$ を省くことで最短表示となる。2 の補助結果により、相異なる二つの素イデアルは常に互いに素である。§4, 4 の計算規則により、その準素成分についても同じことが成り立つ。従って最小公倍数は積になる。

同じ仮定から次の補助結果が従う。\emph{単位元をもつ可換環で二重鎖条件が成り立つとき、イデアル $\mm$ に対して素な素イデアルは、単位イデアルを除けば、$\mm$ に含まれていないものにちょうど限られる。素であることと互いに素であることは一致する。}

\emph{補足。} 公理 I, II, III の下では、イデアル $\mm$ は零イデアルと異なるものと仮定しなければならない。

$\mpideal$ が $\mm$ に含まれていなければ、それは $\mm$ に属するすべての素イデアルと異なり、2 の補助結果によりそれらのいずれでも割り切れない。従って各準素成分に対して素であり、従って $\mm$ に対して素である。同時に各成分と互いに素であり、従って $\mm$ とも互いに素である。

他方、$\mpideal\ne\mo$ が $\mm$ に含まれていれば、それは随伴素イデアルの一つと一致しなければならない。それが成分 $\mq=\mq_1$ に属するとする。このとき $\mq:\mpideal$ は $\mq$ の真の約イデアルである。なぜなら
\[
        \mpideal^{\varrho-1}\equiv0\pmod{\mq:\mpideal},\qquad
        \mpideal^{\varrho-1}\not\equiv0\pmod{\mq}
\]
だからである。同時に $\mq:\mpideal$ は、$\mpideal^{\varrho-1}$ の約イデアルとして、素イデアル $\mpideal\ne\mo$ だけで割り切られ、従って準素である。分解の一意性により、積
\[
        (\mq:\mpideal)\mq_2\cdots\mq_r
\]
は $\mm:\mpideal$ で割り切れるが、$\mm$ の真の約イデアルである。従って $\mpideal$ は $\mm$ に対して素でない。

従って、イデアル $\mt$ が $\mm$ に対して素であるのは、$\mt$ に属するどの素イデアルも $\mm$ に含まれないとき、かつそのときに限る。この場合 $\mt$ は $\mm$ とも互いに素である。逆に、互いに素なイデアルは常に相互に素である（§5, 5）ので、この仮定の下では二つの概念は一致する。

\subsection*{§8. 商体における整閉性の下でのイデアル論}

ここまで展開されたイデアル論は、最後の二つの公理を加えることで通常の理論へと精密化される。

1. \emph{補助結果。} $\mR$ を零因子をもたず単位元をもつ可換環とする。$\mR$ に約イデアル鎖条件を仮定する（公理 I, III, IV）と、$\ma=\ma\mb$ かつ $\ma\ne(0)$ から $\mb=\mo$ が従う。従って
\[
        \mc^\varrho\ne\mc^{\varrho+1}
\]
は、零イデアルおよび単位イデアルと異なるすべてのイデアルについて成り立つ。
\footnote{他方、同じ仮定の下でも、$\ma\mb=\ma\mc$ から $\mb=\mc$ を推論することはできない。例えば、$\mR$ を体上の $x,y$ の多項式領域とし、$\ma=(x,y)$、$\mb=(x^2,xy,y^2)$、$\mc=(x^2,y^2)$ とする。このとき $\ma\mb=\ma\mc=\ma^3$ であるが、$\mb\ne\mc$ である。}
証明は補助結果 §1, 3 と同じである。$\ma\mb$ を $\mb$ に関する有限加群と見なすことができるからである。I により存在する $\ma$ のイデアル基底を $\alpha_1,\ldots,\alpha_n$ とする。$\ma=\ma\mb$ から方程式系
\[
        \alpha_i=b_{i1}\alpha_1+\cdots+b_{in}\alpha_n,
        \qquad b_{ik}\equiv0(\mb),\quad i=1,\ldots,n
\]
が得られる。$\mR$ は零因子をもたないので、行列式 $|b_{ik}-e_{ik}|$ は消える。ただし $i\ne k$ では $e_{ik}=0$、$e_{ii}=e$ とする。
\[
        e^n+b_1e^{n-1}+\cdots+b_n=0
\]
で $b_i\equiv0(\mb)$ だから、$e\equiv0(\mb)$、従って $\mb=\mo$ が従う。

2. あとは、公理 I から IV に商体における整閉性（公理 V）を加えると、零イデアルと異なる任意の準素イデアルがその随伴素イデアルの冪であることを示すだけである。結果として得られる表示
\[
        \mm=\mpideal_1^{\varrho_1}\cdots\mpideal_r^{\varrho_r}
\]
の一意性は、零イデアルと異なる任意のイデアルについて、すでに定理 V と 1 の補助結果により証明されている。同時に、これらの素イデアルは単位イデアルと異なる真の約イデアルをもたないことも示されている。証明は、まず Dedekind の帰結 II（§1, 4）を用いて指数二の場合に行い、一般の場合は完全帰納法によって行う。

\emph{補助結果。} \emph{公理 I から V の下で、指数二の準素イデアルは $\mq=\mpideal^2$ 以外に存在しない。}

示すべきことは、
\[
 \mpideal^2\equiv0\pmod{\mq},
 \qquad \mq\equiv0\pmod{\mpideal},
 \qquad \mq\not\equiv0\pmod{\mpideal^2}
\]
から必ず $\mq=\mpideal$ が従うことを示す。そこで
\[
 c\equiv0\pmod{\mq},
 \qquad c\equiv0\pmod{\mpideal},
 \qquad c\not\equiv0\pmod{\mpideal^2}
\]
とする。補助結果 §7, 3 から、$\mo c:\mpideal$ は $\mo c$ の真の約イデアルになる。従って、ある元 $b$ が存在して
\[
        b\mpideal\equiv0(\mo c)\equiv0(\mq),
        \qquad b\not\equiv0(\mo c)
\]
を満たす。従って $\gamma=b/c$ は商体に属するが $\mR$ には属さず、整ではない。示すべきは $b\not\equiv0(\mpideal)$ である。そうすれば、$\mq$ は準素で $\mpideal$ に属するので、$\mpideal\equiv0(\mq)$ が従う。

Dedekind の帰結 II により、$\mR$ の元 $m,n$ が存在して
\[
        \gamma=b/c=m/n
\]
となり、さらに $m^2/n$ は整でない。すなわち
\[
        bn=mc;
        \qquad m\not\equiv0(\mo n);
        \qquad m^2\not\equiv0(\mo n)
\]
である。$b\mpideal$ に $m$ を掛けると
\[
        mb\mpideal\equiv0(\mo mc);
        \qquad\text{従って}\qquad \equiv0(\mo nb),
        \qquad\text{そこから}\qquad m\mpideal\equiv0(\mo n)
\]
を得る。最後の結論では、環が零因子をもたないことを用いた。ここから
\[
        m\not\equiv0(\mpideal)
\]
が従う。実際、$m^2\not\equiv0(\mo n)$ かつ $n\equiv0(\mpideal)$ であり、後者は $m\not\equiv0(\mo n)$ のため $\mo n:\mpideal$ が $\mo n$ の真の約イデアルになることから、§7, 3 の補助結果によって従う。四つの関係
\[
        m\not\equiv0(\mpideal);
        \quad n\equiv0(\mpideal);
        \quad c\equiv0(\mpideal);
        \quad {}\not\equiv0(\mpideal^2);
        \quad bn=mc
\]
は $b\not\equiv0(\mpideal)$ を与える。実際、$\mpideal^2$ は準素なので、まず $mc\not\equiv0(\mpideal^2)$ となり、次に $bn=mc$ から $b\not\equiv0(\mpideal)$ が従う。これで補助結果が証明された。

3. \emph{補助結果。} \emph{公理 I から V の下で、指数 $\varrho$ の準素イデアルは $\mq=\mpideal^\varrho$ 以外に存在しない。}

これは 2 の補助結果から、さらなる公理の使用なしに従う。
\footnote{補助結果 2 の妥当性を仮定した補助結果 3 の証明は、Masazo Sono, \emph{On Congruences II}, §§9, 10 にある。注 19 を参照。}
まず次を示す。
\[
        \text{もし }c\equiv0(\mpideal),\quad c\not\equiv0(\mpideal^2)
        \quad\text{ならば}\quad
        \mpideal^\sigma=(\mo c^\sigma,\mpideal^{\sigma+1})
\]
がすべての $\sigma$ について成り立つ。2 の結果により $\mpideal=(\mo c,\mpideal^2)$ である。そこで
\[
        \mpideal^{\sigma-1}=(\mo c^{\sigma-1},\mpideal^\sigma)
\]
と仮定し、$\mpideal$ を掛けて分配法則を用いると、
\[
        \mpideal^\sigma=(\mpideal c^{\sigma-1},\mpideal^{\sigma+1})
        =(\mo c^\sigma,\mpideal^2c^{\sigma-1},\mpideal^{\sigma+1})
        =(\mo c^\sigma,\mpideal^{\sigma+1})
\]
となる。証明すべき補助結果は次の第二の形に言い換えられる。
\[
        \mq\equiv0(\mpideal^\sigma);
        \quad \mq\not\equiv0(\mpideal^{\sigma+1});
        \quad\text{かつ}\quad
        \mpideal^{\sigma+\lambda}\equiv0(\mq),
        \quad \lambda\ge1
\]
ならば
\[
        \mpideal^{\sigma+\lambda-1}\equiv0(\mq)
\]
でもある。実際、$\mq\equiv0(\mpideal)$ だが $\not\equiv0(\mpideal^{\varrho+1})$ となる指数 $\sigma\geqq1$ は、各準素イデアルについて存在する。ここで $\mq\not\equiv0(\mpideal^{\varrho+1})$ は、$\mpideal^\varrho\equiv0(\mq)$ と、1 の補助結果による $\mpideal^\varrho\ne\mpideal^{\varrho+1}$ から従う。第二の形を有限回適用すれば $\mpideal^\sigma\equiv0(\mq)$、従って $\mq=\mpideal^\sigma$ が得られる。

第二の形の仮定と $\mpideal^\sigma$ の表示から、$\mq$ の各元 $q$ について
\[
        q\equiv bc^\sigma(\mpideal^{\sigma+1}),
        \qquad b\not\equiv0(\mpideal),
        \quad\text{または}\quad
        bc^\sigma\equiv0(\mq,\mpideal^{\sigma+1})
\]
である。$(\mq,\mpideal^{\sigma+1})$ は準素なので、ここから
\[
        c^\sigma\equiv0(\mq,\mpideal^{\sigma+1}),
        \quad\text{または}\quad
        c^{\sigma+\lambda-1}\equiv0(\mq,\mpideal^{\sigma+\lambda})
\]
が従い、従って $\mpideal^{\sigma+\lambda-1}\equiv0(\mq)$ である。

4. まとめると、次が得られる。

\textbf{定理 VI.} \emph{可換環で公理 I から V が成り立つなら、零イデアルおよび単位イデアルと異なる任意のイデアルは、零イデアルおよび単位イデアルと異なる有限個の素イデアルの冪の積として一意に表される。これらの素イデアルは、単位イデアルと異なる真の約イデアルをもたない。}

\subsection*{§9. 仮定された分解の帰結としての公理}

逆に、公理 I から V によって定理 VI から従う通常のイデアル分解の存在から、さらに公理を推論するためには、分解を次の形で仮定しなければならない。

\emph{仮定。} 可換環 $\mR$ において、零元または単位元によって生成されない任意のイデアルは、素イデアルの冪の積として一意に表される。これらの素イデアルは、単位元によって生成されない単純イデアル、すなわち $\mo$ と異なる真の約イデアルをもたないイデアルであり、逆にすべての単純イデアルは素イデアルである。一意性は、零元または単位元によって生成されない任意のイデアルについて
\[
        \ma^\varrho\ne\ma^{\varrho+1}
\]
が成り立つという鋭い形で成り立つものとする。

ここでは単位元の存在を仮定していない。単位元が存在しないなら、「単位元によって生成されない」という条件は制限ではない。

1. \emph{公理 III、単位元の存在の証明。} 公理 III が満たされていないと仮定する。仮定に反して、$\mo^2$ が素でない単純イデアルになることを示す。鋭い一意性仮定により $\mo\ne\mo^2$ である。従って $\mo\not\equiv0\pmod{\mo^2}$ だが $\mo^2\equiv0\pmod{\mo^2}$ であり、$\mo^2$ は素でない。しかしそれは単純である。$\mo$ と異なるどの素イデアルでも割り切れず、また仮定された積表示により、その約イデアルは $\mo$ と $\mo^2$ だけだからである。

2. \emph{公理 IV、零因子がないことの証明。} $a$ が零因子で、$a\ne0$, $b\ne0$ かつ $ab=0$ であるとする。すると $a$ と $b$ によって生成される主イデアルの積表示を掛け合わせて
\[
        (0)=\mpideal_1^{\varrho_1}\cdots\mpideal_k^{\varrho_k}
\]
を得る。ここで $\mpideal$ は零イデアルおよび単位イデアルと異なる。後者は、単位元の存在がすでに証明されているからである。表示は最短、すなわちどの因子を削除しても零イデアルの真の約イデアルを与えるものとして取る。零イデアルについては一意性を仮定していないので、これは自動的ではない。一つの素イデアルだけが現れるなら、$(0)=\mpideal^\varrho=\mpideal^{\varrho+1}$ となり、鋭い一意性要求に反する。複数の素イデアルが現れるなら、
\[
\begin{aligned}
\mpideal_1^{\varrho_1+1}
 &= (\mpideal_1^{\varrho_1+1},\mpideal_1^{\varrho_1}\ldots\mpideal_k^{\varrho_k}) \\
 &= \mpideal_1^{\varrho_1}(\mpideal_1,\mpideal_2^{\varrho_2}\ldots\mpideal_k^{\varrho_k}) \\
 &= \mpideal_1^{\varrho_1}.
\end{aligned}
\]
となる。なぜなら単位元の存在により、$\mpideal_1$ は他のすべての素イデアルと互いに素だからである。これも再び鋭い一意性条件に矛盾する。

3. \emph{公理 I と II、鎖条件の証明。} 仮定は、零イデアルと異なる任意のイデアルについて、可除性が積表示に反映されることを与える。従って、
\[
        \mm=\mpideal_1^{\varrho_1}\cdots\mpideal_r^{\varrho_r},
        \qquad
        \ma=\mbar\mpideal_1^{\bar\varrho_1}\cdots\mbar\mpideal_s^{\bar\varrho_s}
\]
で、$\mm\equiv0\pmod{\ma}$ かつ $\ma\ne\mo$ なら、各 $\mbar\mpideal_j$ はある $\mpideal_i$ と同一であり、$0\le\bar\varrho_i\le\varrho_i$ であることが従う。従って $\mm$ の約イデアルは、有限個の指数の組合せに対応して有限個しかない。よって、$\mm$ を法とする剰余類環において二重鎖条件が成り立つ。これにより公理 I と II が証明される。約イデアル鎖条件は $\mR$ 自身でも成り立つ。なぜなら零イデアルの真の約イデアルはすべて零イデアルと異なるからである。

公理 V、すなわち商体における整閉性の証明には、イデアル分解の標準的帰結を用いるので、まずそれらを導かなければならない。

4. \emph{零イデアルと異なる任意のイデアルを法とする剰余類環は主イデアル環である。} 剰余類環が零元だけからなる場合には主張は明らかなので、イデアルは単位イデアルと異なると仮定してよい。

まずイデアルが準素で、$\mq=\mpideal^\varrho$ とし、$c\equiv0(\mpideal)$ だが $\not\equiv0(\mpideal^2)$ であるとする。このとき 3 により $\mo c=\mpideal\ma$ で $(\ma,\mpideal)=\mo$ である。従って
\[
        \mpideal^\sigma=(\mo c^\sigma,\mpideal^\varrho)
        \qquad\text{すべての }\sigma\le\varrho\text{ について}
\]
であり、準素イデアルを法とする剰余類環の任意のイデアルは主イデアルである。一般に、
\[
        \mm=\mq_1\cdots\mq_r
\]
が表示であり、
\[
        \mR/\mm=\ma_1+\cdots+\ma_r
\]
が剰余類環の直和表示（§4, 5）であるなら、$\ma_i$ は $\mR/\mq_i$ と同型である。従って $\ma_i$ の任意のイデアルは主イデアルである。$\mc$ が剰余類環の任意のイデアルなら、$\ma_i\mc$ は主イデアル $\ma_i c_i$ であり、
\[
        \mc=\ma_1c_1+\cdots+\ma_r c_r=\mo c,
        \qquad c=c_1+\cdots+c_r.
\]
実際、$i\ne k$ で $\ma_i\ma_k=(0)$、かつ $c_i\equiv0\pmod{\ma_i}$ であるから、$\ma_i\mo c_i$ と $\ma_i\mc$ は $\ma_i$ の中で一致する。

剰余主イデアル環に関するこの定理は、次の形にも述べられる。$\mc$ が任意のイデアルなら、与えられたイデアル $\mb$ に対して素なイデアルを掛けることにより、$\mc$ を主イデアルに変えることができる。実際、$\mm=\mb\mc$ と置くと、$\mc$ は $\mm$ を法として主イデアルになる。すなわち
\[
        \mc=(\mo c,\mm)=(\mc\ma,\mc\mb)=\mc(\ma,\mb),
\]
従って $\mo c=\mc\ma$ かつ $(\ma,\mb)=\mo$ である。

5. \emph{分数イデアルの理論。} 分数イデアルとは、商体 $\mK$ 内の有限 $\mR$-加群を意味する。

\emph{$\mK$ 内の零でない有限 $\mR$-加群は、乗法に関してアーベル群をなす。}

二つの有限 $\mR$-加群の積は再び有限である。乗法は結合的かつ可換である。$\mK=\mo\mK$ であるから、単位イデアルが単位元である。残るのは、方程式
\[
        \mathfrak A\mX=\mo
\]
が有限 $\mR$-加群の系において常に解をもつことを示すだけである。

予備的注意。方程式 $\mathfrak A\mT=\mathfrak B$ がただ一つの解をもつなら、$\mT$ は加群商 $\mathfrak B:\mathfrak A$（§5, 4）である。なぜなら
\[
\mT\equiv0\pmod{\mathfrak B:\mathfrak A},\quad
\mathfrak B=\mathfrak A\mT\equiv0\pmod{\mathfrak A\cdot(\mathfrak B:\mathfrak A)}\equiv0\pmod{\mathfrak B}
\quad\text{または}\quad
\mathfrak A\cdot(\mathfrak B:\mathfrak A)=\mathfrak B.
\]

まず $\mathfrak A$ が主加群で、$\mathfrak A=(\alpha)$ であるなら、$\mX=(e/\alpha)$ は $\mathfrak A\mX=\mo$ の解であり、$\mX$ もまた有限 $\mR$-加群である。従って任意の有限 $\mR$-加群は、$\mR$ の二つのイデアルの加群商であり、このことが分数イデアルという名称を正当化する。実際、$\tau_1,\ldots,\tau_r$ が $\mT$ の加群基底であるなら、$\tau_i=t_i/a$ と置き、$\mt$ を $\mR$ において $t_i$ によって生成されるイデアルとする。このとき
\[
        \mo a\cdot\mT=\mt,
        \qquad \mT=(\mt:\mo a)
\]
である。ここで商は $\mK$ の中で取られる。

方程式 $\mathfrak A\mX=\mo$ の解は一般に次のように与えられる。
\[
        \mathfrak A c=\ma,
        \qquad \mathfrak A\mb c=\mo a,
        \qquad \mathfrak A\cdot(\mb c:\mo a)=\mo.
\]
ここで $\mX=(\mb c:\mo a)$ は、イデアルを主イデアルで割った商として、再び有限 $\mR$-加群である。群性が証明された。また任意の二つのイデアルの加群商 $\ma:\mb$ も、方程式 $\mb\mX=\ma$ の解として有限 $\mR$-加群であることが従う。
6. 群性からさらに次の帰結が従う。

6\,$\alpha$. \emph{消去が可能である。$\mathfrak T=\mathfrak A\mathfrak M:\mathfrak B\mathfrak M$ から $\mathfrak T=\mathfrak A:\mathfrak B$ が従い、逆も成り立つ。} 実際、$\mathfrak T\mathfrak B\mathfrak M=\mathfrak T\mathfrak A\mathfrak M$ から $\mathfrak T\mathfrak B=\mathfrak T\mathfrak A$ が従い、逆も同様である。

6\,$\beta$. \emph{任意の分数イデアルは、$\mathfrak C=\ma:\mb$ という表示をもつ。ただし $\ma$ と $\mb$ は互いに素であり、この条件により $\ma$ と $\mb$ は一意に定まる。} 存在は 5 と 6\,$\alpha$ から従う。$\mathfrak C=\ma:\mb=\mbar\ma:\mbar\mb$ から $\ma\mbar\mb=\mb\mbar\ma$ が従い、従って両方の組 $\ma,\mb$ および $\mbar\ma,\mbar\mb$ が互いに素であると仮定すれば一意性が得られる。

6\,$\gamma$. \emph{非整元、すなわち $\mR$ に属さない $\mK$ の元によって生成される任意の主加群は、分子のどの冪も分母で割り切れないような主イデアルの商として表示される。} 簡約表示を $\mo\eta=\mb:\mc$ とし、$(\mb,\mc)=\mo$ とする。4 に従って $\mo b=\mb\ma$ かつ $(\ma,\mc)=\mo$ となるように $\ma$ を選ぶ。すると
\[
        \mo\eta=\ma\mb:\ma\mc=\mo b:\ma\mc,
\]
であり、$\ma\mc=\mo b:\mo\eta$ なので、イデアル $\ma\mc$ もまた主イデアルである。従って
\[
        \mo\eta=\mo b:\mo c
\]
と書く。$(\ma\,\mc)=\mo$ かつ $(\mb,\mc)=\mo$ であるから、$b$ のどの冪も $c$ で割り切れない。

7. \emph{公理 V、商体における整閉性の証明。} 6\,$\gamma$ により、$\mK$ の任意の非整元は商表示
\[
        \eta=a/c
\]
をもち、$\mR$ において分子のどの冪も分母で割り切れないようにできる。このような元は、$\mR$ 上整であることを特徴づける方程式を満たすことができない。なぜなら
\[
        \eta^n+r_1\eta^{n-1}+\cdots+r_n=0
\]
からは $\mR$ において $a^n\equiv0\pmod{\mo c}$ が得られるからである。

8. \emph{帰結。} 可換環 $\mR$ において公理 I から IV が成り立ち、かつすべての準素イデアルが既約であるなら、公理 V、すなわち商体における整閉性も成り立つ。公理 I から III により、任意の素イデアル冪 $\mpideal^\varrho$ は準素である。特に $\mpideal^2$ は準素であり、仮定により既約である。ここから
\[
        \mpideal=(\mo c,\mpideal^2)
\]
を証明しなければならない。そうすれば、補助結果 §8, 3 によりすべての準素イデアルは素イデアル冪となり、他の仮定と合わせて 7 から整閉性が従う。

約イデアル鎖条件により
\[
        \mpideal=(\mo c_1,\ldots,\mo c_k,\mpideal^2)
\]
と書ける。ここで乗数としては $\mpideal$ を法とする剰余類だけが関係し、$c$ は $\mpideal$ を法とする剰余類体上線型独立であると仮定できる。$k>1$ なら、この線型独立性から
\[
        \mpideal^2=[\mq_1,\mq_2],
        \qquad
        \mq_1=(\mo c_1,\mpideal^2),
        \qquad
        \mq_2=(\mo c_2,\ldots,\mo c_k,\mpideal^2),
\]
という表示が得られ、$\mq_1$ と $\mq_2$ はいずれも $\mpideal^2$ の真の約イデアルである。従って $\mpideal^2$ は可約となり、仮定に反する。よって $\mpideal=(\mo c,\mpideal^2)$ であり、主張した帰結が証明された。

逆に、公理 I から V の帰結として、任意の準素イデアルは直ちに既約である。素イデアル冪 $\mpideal^\varrho$ は、$\mpideal^\sigma$ および $\mpideal^\tau$ という形の真の約イデアルの最小公倍数にはなり得ないからである。

9. 環に零因子を許すなら、公理 I から III と商環における整閉性から、任意の準素イデアルが既約であることは従わない。
\footnote{注 19 を参照。}
$\mT$ を、整閉性を除く公理 I から IV がすべて成り立つ環とする。このような環は存在する。例えば、$\mR$ で公理 I から V が成り立つとき、$\mR$ の商体の有限拡大体における整量全体の系とは異なる有限オーダーがそれである（§3, 2）。8 の帰結により、$\mT$ の中には可約な準素イデアル $\mq$ が存在する。$\mpideal^\sigma$ を $\mq$ の真の倍イデアルとし、$\mR$ を剰余類環 $\mT/\mpideal^\varrho$ とする。この環では、$\mq$ に対応するイデアルが非零の可約準素イデアルである。$\mR$ では公理 I から III が成り立ち、商環における整閉性も成り立つ。なぜなら、$\mT$ の元で $\mpideal$ で割り切れないものはすべて $\mpideal^\sigma$ と互いに素になるので、$\mR$ のすべての正則元は単元であり、$\mR$ はその商環と一致するからである。

\subsection*{§10. 二重鎖条件と組成列}

ここで、整列されていると仮定された任意の加群領域（§2）について、二重鎖条件、すなわち加群の任意の約鎖および任意の倍鎖が有限段階で停止することが、組成列の存在と同値であることを示す。
\footnote{領域の加群は加法に関してアーベル群であり、正確には乗数領域が $\mR$ の元からなるので一般化されたアーベル群である。それらはアーベル群であるから、ここでは組成列の概念は主列の概念と一致する。この段落の定理は、「加群」の下に任意の非可換、さらには一般化された群の正規約数全体を理解しても成り立つ。以前の記法とわずかに異なる記法は、群論を想起させるためのものである。}
加群領域を可換環に特殊化すれば、環の全イデアルの系について対応する事実が得られる。2 では全体を通じて、加群同型を環同型に置き換えなければならない。

加群領域 $\mG$ は、$\mG$ 自身と零加群 $\mE$ 以外に $\mR$-加群を含まないとき、\emph{単純}と呼ばれる。加群 $\mA$ は、$\mG/\mA$ が単純であるとき、\emph{$\mG$ において単純}と呼ばれる。

倍鎖
\[
        \mG,\mU_1,\ldots,\mU_r,\mE
\]
は、鎖中のすべての加群が相異なり、各加群が直前のものにおいて単純であるとき、$\mG$ の長さ $r$ の組成列と呼ばれる。これは、剰余類加群、すなわち商群 $\mU_{i-1}/\mU_i$ が単純な加群領域であり、また $\mG/\mU_1$ と $\mU_r/\mE$ も単純であることと同値である。剰余類加群 $\mG/\mF$ を作ることにより、$\mG$ から $\mF\ne\mG$ まで組成列が存在する場合は、絶対的な場合に帰着される。

1. \emph{$\mG$ に二重鎖条件を仮定すると、$\mG\ne\mE$ なら $\mG$ に組成列が存在する。} $\mG$ の仮定された整列順序と約加群鎖条件から、§6 と同様に、準辞書式順序により全加群の系の整列順序が得られる。この整列順序と約加群鎖条件により、$\mG$ において単純な加群が少なくとも一つ存在する。すなわち、$\mE_1$ を整列順序で $\mE$ の第一の真の約加群とし、一般に $\mE_i$ を $\mE_{i-1}$ の第一の真の約加群とする。よく定まった鎖
\[
        \mG,\mG_1,\ldots,\mG_i,\ldots
\]
は有限段階で停止し、従って必ず $\mG$ において単純な加群 $\mA$ に到達する。$\mU_1\ne\mE$ なら、同じ手続きにより $\mU_1$ において単純な加群 $\mU_2$ が得られる。一般に、$\mU_{i-1}\ne\mE$ なら、$\mU_{i-1}$ において単純な加群 $\mU_i$ が存在する。これにより、よく定まった倍鎖
\[
        \mG,\mU_1,\ldots,\mU_i,
\]
が得られ、これは倍加群鎖条件により停止する。従ってこれは $\mG$ の組成列である。

2. \emph{$\mG$ に組成列が存在するなら、$\mG$ で二重鎖条件が成り立つ。} 証明は完全帰納法によるもので、$\mG$ の整列順序を必要としない。帰納法の過程で、組成列の存在だけを仮定した Jordan--Hölder の定理の証明も同時に得られる。
\footnote{環の場合については Masazo Sono（注 21）, \emph{On Congruences I}, §§11--13 を参照。そこでは非可換群における組成列の類似物も扱われている。すなわち $\mR,\mU_1,\ldots,\mU_r,\mE$ という列で、$\mU_i$ はイデアルであり $\mU_{i-1}$ において単純であるが、必ずしも $\mR$ のイデアルである必要はない。現在の議論は、倍鎖を、各 $\mU_i$ が $\mU_{i-1}$ において正規であるような鎖と解し、約鎖を対応して定義するなら、非可換群に対しても有効である。Dedekind, ``Über die von drei Moduln erzeugte Dualgruppe'', Math. Ann. 53 (1900), pp. 371--403 も参照。そこでは、はるかに一般的な領域について、存在仮定だけの下で同様に Jordan--Hölder の定理が証明されている。第二同型定理の代わりにやや弱い対応関係が現れるため、定理の主張もやや弱い。証明の原理はここに与えるものと同一である。}

$\mG$ が単純で、従って $r=0$ であり、唯一の組成列が $\mG,\mE$ であるとき、次の主張が成り立つ。

$\alpha$) $\mG$ において単純な任意の加群を通って組成列を引くことができる。

$\beta$) Jordan--Hölder の定理：任意の組成列は同じ長さをもち、商群（剰余類加群）の系は順序を除いて一致する。対応する商群は同型である。

$\gamma$) $\mG$ と異なる任意の加群を通って組成列を引くことができる。

$\delta$) 倍加群鎖条件が成り立つ。

$\varepsilon$) 約加群鎖条件が成り立つ。

従って、$\alpha$)--$\varepsilon$) が、長さ $r<r+1$ の組成列をもつ任意の加群領域について証明されていると仮定する。$\mG$ に長さ $r+1$ の組成列
\[
        \mG,\mA,\mA_1,\ldots,\mA_r,\mE
\]
が存在するという仮定の下で、それらを証明する。

2$\alpha$. $\mG$ において $\mA$ 以外に単純な加群がなければ、証明すべきことはない。$\mB$ を $\mG$ において単純な加群で、$\mB\ne\mA$ とする。$\mA$ と $\mB$ はいずれも $\mG$ において単純で、かつ異なるので、必ず $(\mA,\mB)=\mG$ である。$\mM=[\mA,\mB]$ と置く。第二同型定理（§4, 2）により
\[
        \mG/\mB\cong \mA/\mM,
        \qquad
        \mG/\mA\cong \mB/\mM.
\]
従って $\mM$ は $\mA$ においても $\mB$ においても単純である。$\mA$ は長さ $r<r+1$ の組成列をもつので、帰納仮定 $\alpha)$ と $\beta)$ により、$\mA$ の中に $\mM$ を通る組成列、例えば
\[
        \mA,\mM,\mM_1,\ldots,\mM_{r-1},\mE
\]
が存在する。従って
\[
        \mG,\mB,\mM,\mM_1,\ldots,\mM_{r-1},\mE
\]
は $\mB$ を通る $\mG$ の組成列である。

2$\beta$. Jordan--Hölder の定理を証明するため、
\[
\begin{aligned}
&(1)\quad \mG,\mA,\mA_1,\ldots,\mA_r,\mE,\qquad\text{および}\qquad
(2)\quad \mG,\mB,\mB_1,\ldots,\mB_s,\mE
\end{aligned}
\]
を $\mG$ の二つの組成列とする。$\mA=\mB$ なら、結果は長さ $r<r+1$ に対する帰納仮定から従う。そうでなければ、2$\alpha$ で構成された列
\[
        (3)\quad \mG,\mB,\mM,\mM_1,\ldots,\mM_{r-1},\mE,
        \qquad
        (4)\quad \mG,\mA,\mM,\mM_1,\ldots,\mM_{r-1},\mE
\]
と比較する。帰納仮定により、(1) と (3) の商群の系は一致する。同様に、適切な商へ移ると (4) は $\mB$ を通る長さ $r$ の組成列であるから、(2) と (4) の商群の系も一致する。従って $s=r$ である。2$\alpha$ の同型から (3) と (4) の一致が得られ、従って (1) と (2) の一致が得られる。これで Jordan--Hölder の定理が証明された。

2$\gamma$. 予備的注意。$\mA,\mB,\mH$ を $\mG$ の加群とし、$\mB$ が $\mA$ において単純であると仮定する。このとき加群 $(\mB,\mH)$ と $(\mA,\mH)$ は一致するか、あるいは $(\mB,\mH)$ が $(\mA,\mH)$ において単純である。第二同型定理により
\[
        (\mA,\mB,\mH)/(\mB,\mH)
        \simeq
        \mA/[\mA,(\mB,\mH)]
\]
である。$\mB\equiv0\pmod{\mA}$ であるから、これは
\[
        (\mA,\mH)/(\mB,\mH)\simeq \mA/\mC,
        \qquad \mC=[\mA,(\mB,\mH)]
\]
である。ここで $\mC\equiv0(\mA)$ である。他方、$\mB\equiv0(\mC)$ でもある。なぜなら $\mB\equiv0(\mA)$ かつ $\mB\equiv0((\mB,\mH))$ だからである。従って $\mC$ は $\mA$ と $\mB$ の間にあり、仮定により $\mA$ または $\mB$ と同一である。これで予備的注意は証明された。

今、
\[
        \mG,\mU,\mU_1,\ldots,\mU_r,\mE
\]
を $\mG$ の組成列とし、$\mH$ を $\mG$ と異なる任意の $\mG$ の加群とする。$\mH$ を通る組成列があることを示すため、
\[
        \mG,(\mU,\mH),(\mU_1,\mH),\ldots,(\mU_r,\mH),\mH
\]
を作る。予備的注意により、この中の相異なる加群は $\mG$ から $\mH$ への組成列をなす。従って第一の真の項は $\mG$ において単純である（これは、それが $\mB$ と一致する場合の 2$\alpha$ の特殊場合である）。2$\alpha$ により、$\mG$ にはその項を通る長さ $r<r+1$ の組成列があり、従って帰納仮定により $\mH$ を通る組成列がある。$\mG$ を付け加えれば、$\mH$ を通る $\mG$ の組成列が得られる。この議論を繰り返すと、そのような列はとくに、いま $\mG$ から $\mH$ まで構成した列の延長として選べることが分かる。

2$\delta$. 倍加群鎖条件の証明。再び $\mG$ に長さ $r+1$ の組成列が存在すると仮定し、
\[
        \mG,\mB,\mB_1,\ldots,\mB_\nu,\ldots
\]
を倍鎖とする。各項はその直前の項の真の倍加群である。鎖が $\mG$ から始まると仮定しても一般性を失わない。$\mG$ は常に第一項として付け加えられるからである。2$\gamma$ により、$\mB$ を通る組成列がある。従って $\mB$ はより短い長さの組成列をもつ。帰納仮定により、鎖
\[
        \mB,\mB_1,\ldots,\mB_\nu,
\]
は有限段階で停止する。従って $\mG$ を付け加えた後も同じことが成り立つ。

2$\varepsilon$. 約加群鎖条件の証明。再び $\mG$ に長さ $r+1$ の組成列を仮定し、
\[
        \mG,\mC,\mC_1,\ldots,\mC_i,
\]
を約鎖とする。各項はその直前の項の真の約加群である。この鎖が $\mG$ から始まると仮定しても、やはり制限はない。2$\gamma$ により、$\mC$ を通る組成列がある。従って $\mG/\mC$ にはより短い長さの組成列がある。帰納仮定により、約鎖
\[
        \mG/\mC,\ \mC_1/\mC,\ldots,\mC_i/\mC,
\]
は有限段階で停止する。第一同型定理（§4, 2）が与える一対一対応により、$\mG$ から始まる元の鎖についても同じことが成り立ち、従って $\mE$ を付け加えた鎖についても同じことが成り立つ。これにより、組成列の存在と二重鎖条件という二つの仮定の同値性が証明された。

\emph{1925年8月13日受理。}

\clearpage
%
\begingroup
\renewcommand{\thefootnote}{\arabic{footnote})}

\section*{31. 代数的数体または関数体のオーダーに関する判別式定理}
\begin{center}
Journal f. d. reine u. angew. Math. 157 (1927), S. 82--104

Göttingen の \emph{Emmy Noether} による。\\[0.4em]
\rule{9em}{0.4pt}
\end{center}

\emph{Dedekind} の判別式定理は、周知のとおり次を述べる。ある素数が体の判別式を割るのは、その素数の生成する主イデアルが少なくとも一つの素イデアルの平方で割り切れるとき、かつそのときに限る。

以下では、この定理が（有限）代数的数体のすべての有限オーダー、すなわち\emph{代数的整数からなる環}であって有理整数環を含むものすべてについて、次の形で保たれることを示す。\emph{素数 $p$ があるオーダーの判別式を割るのは、そのオーダーにおいて成り立つ主イデアル $(p)$ の最大準素成分へのイデアル分解\footnote[1]{これについては E. Noether, \emph{Abstrakter Aufbau der Idealtheorie in algebraischen Zahl- und Funktionenkörpern}, Math. Ann. 96 (1926), S. 26--61 を参照されたい。すべての定義については、以下で「イデアル論」として引用するこの論文を参照する。以下で「オーダー」は常に「有限オーダー」を意味し、「環」は可換環を意味する。}の中に、少なくとも一つの真の準素イデアル、すなわち素イデアルではない準素イデアルが現れるとき、かつそのときに限る。}

主オーダー、すなわち体のすべての整数からなる系には、素イデアルの冪以外の準素イデアルは存在しないので、特殊化によって \emph{Dedekind} の定理が得られる。

定理の証明は、オーダーから $(p)$ を法とする剰余類環へ移ることにより、その環に含まれる一つの体に関して有限階数をもつ環の分解論へ帰着される。すなわち、同じ体を係数とし、単位元をもつ可換な超複素量の系をなす環の分解論へ帰着される。ここで係数体として現れるのは、とくに有理整数環の $p$ を法とする剰余体であり、$(p)$ の分解には剰余類環における零イデアルの分解が対応する。

そこで、このような有限階数の環を、その零イデアルが（有限個の）素イデアルの最小公倍数として表されるとき\emph{完全可約}と呼ぶならば、問題は完全可約性の判定条件を求めることになる。その一つの判定条件として、\emph{係数領域が完全体である場合}、もとの環が完全可約であることと、係数領域を代数閉体へ拡張した環が完全可約であることとが同値であることがまず得られる。そして、この拡張環については、判別式が零でないことが完全可約性の必要十分条件であるときわめて容易に認められる\footnote[1]{特別な場合には \emph{Hilbert} も同様に論じている。\emph{Zur Theorie der aus $n$ Haupteinheiten gebildeten komplexen Größen}, Gött. Nachr. 1896.}。さらに、\mbox{オーダー}の判別式は、任意の素数を法として剰余類環の判別式へ移るから、これで判別式定理が証明される\footnote[2]{非可換の場合をも含む関連研究は、ただし\emph{主オーダー}（極大領域）に限定し、有理整数係数を仮定するものとして、刊行されたばかりの \mbox{A. Speiser}, \emph{Allgemeine Zahlentheorie}, Naturf. Gesellschaft Zürich 71 (1926), S. 8--48 に見られる。そこに挙げられたアメリカの文献も参照されたい。}。

数体の代わりに関数体を考えるならば、オーダーは今度はすべての有理整関数からなる基礎領域、あるいは、複数の不定元の代数関数または整な代数関数の場合には、有理整係数をもつすべての整な汎関数からなる基礎領域を含まなければならない。この場合には、剰余類環において\emph{係数領域が不完全体である場合}も起こりうる。たとえば、整汎関数の領域で一つの素数を法とする剰余類環を考える場合がそうである。係数領域が不完全体である場合には、素イデアルを法とする剰余体が第一種または第二種の拡大体であるかに応じて\footnote[3]{\emph{Steinitz} のよく知られた定式化、J. f. M. 137 による。本文 §~2, 4 の注も参照されたい。}、第一種と第二種の素イデアルを区別しなければならず、それに対応して第一種と第二種の完全可約性を区別しなければならない。上に示した判定条件はいずれも第一種の完全可約性に関するものであり、完全体の場合には第一種の完全可約性だけが現れる。したがって、最も一般の場合の\emph{判別式定理}は次のように述べられる。\emph{基礎領域の素元 $p$ があるオーダーの判別式を割るのは、$(p)$ のイデアル分解の中に少なくとも一つの真の準素成分、または一つの第二種の素イデアルが現れるとき、かつそのときに限る。}

\endgroup
%
\begingroup
\renewcommand{\thefootnote}{\arabic{footnote})}

主オーダーについては、この事実は \emph{Kronecker} により、いくつかの例において初めて認められた。それに先立つ Festschrift には、数体の主オーダーについて成り立つ定式化がなお誤って記されており、しかも証明は遂行されていなかった。主オーダーについては、\emph{Ostrowski}\footnote[4]{A. Ostrowski, Zur arithmetischen Theorie der algebraischen Größen, Nachr. Gött. Ges. d. Wissensch. 1919.} が上の判別式定理の証明と、それに続くさらなる分岐諸定理を与え、同時に詳細な文献案内も掲げている。最後に、\emph{相対体のすべてのオーダーに対する判別式定理}が、基礎領域の各素イデアルに関する商環へ移ることによって直接の帰結として得られる。この移行の考えは、特殊な場合には知られていたが、\emph{W. Krull}\footnote[5]{同人の Danzig 講演を参照されたい。Jhrbr. d. D. Math. Ver. 34, S. 121.} によって原理的に強調された。より一般的な研究の枠内での商環の体系的理論は、\emph{H. Grell}\footnote[6]{\emph{Beziehungen zwischen den Idealen verschiedener Ringe}（Math. Annalen に掲載予定）。} が与えている。

ここで行う、すべてのオーダー、ならびに完全体または不完全体である係数領域の一様な扱いは、\mbox{イデアル論}の知識に加えて、初めから有限階数の環の理論、したがって有限拡大の一般概念を基礎に置き、\mbox{原始元}によって生成される拡大の概念を基礎に置かないことに基づいている。したがって、\mbox{オーダー}の\mbox{加群基底}、およびその基底元相互の関係は、それぞれ体の原始元の冪と定義方程式に代わるものであり\footnote[1]{この見方は、\emph{Dedekind} が高次複素数について示した見方と密接に通じる。Zur Theorie der aus $n$ Haupteinheiten gebildeten komplexen Größen, Gött. Nachr. 1885.}、あるいは、それぞれ基本形式の冪と基本方程式に代わるものである。

この加群基底は、基礎領域の各素元を法とする剰余類環の基底へ移る。これに対し、\mbox{原始元の冪}からなる基底では、主オーダーにおいてさえそうなるとは限らない。また、複数の不定元をもつ\mbox{整な代数関数}が関わると、基本形式の冪からなる基底についてもそうなるとは限らない\footnote[2]{\emph{Ostrowski}, 前掲論文、「第一種の不規則性」を参照されたい。}。言い換えれば、\mbox{基礎領域}を\mbox{係数}とする一変数多項式環を任意の定義方程式、または基本方程式で割った剰余類環は、$p$ を法として、\mbox{主オーダー}を $(p)$ で割った剰余類環と同型であるとは限らない。この同型が失われることに、通常の\mbox{表示法}における主要な困難がある。ここではこの補助多項式環が現れないので、これらの困難は自動的に回避される。同時に、剰余類環の加法的分解論は方程式の乗法的分解に相当するものとみなすことができ、ここでも方程式を互いに素な因子に分解することには、多項式剰余類環の加法的分解が対応する。

判別式定理は、オーダーの一般分岐論における最初の定理をなすにすぎないことに注意すべきである。主オーダーでは、判別式定理のほかに、より精密な指数決定を可能にする分岐イデアルの理論がある。この指数決定は方程式の準素因子の冪表示に基づき、$p$ から $q=p^e$ に至る組成列の長さの決定と解釈できる。そして一般分岐論でも、おそらくこの形で現れるであろう。

\endgroup
%
\begingroup
\renewcommand{\thefootnote}{\arabic{footnote})}

\begin{center}
\textbf{\S 1. 体の拡大環}
\end{center}

\textbf{1.} 任意の環について、いくつかの注意をあらかじめ述べる。$\mathfrak T$ を環、$S$ を $\mathfrak T$ の元からなる系とする。このとき「$\mathfrak T$ 内で $S$ から生成される環」とは、$\mathfrak T$ 内の $S$ を含むすべての環の共通部分をいう。これに応じて、「$\mathfrak T$ 内で $S$ から生成されるイデアル」も定義され、また $P$ が $\mathfrak T$ の部分環であるときには、「$\mathfrak T$ 内で $S$ から生成される $P$-加群」も定義される。\footnote[3]{これについては「イデアル論」\S~1 を参照されたい。}

$\mathfrak R$ が $\mathfrak T$ の部分環であり、$S$ が $\mathfrak T$ の元からなる系であるとする。このとき、$\mathfrak T$ 内で $\mathfrak R$ と $S$ から生成される環を、$\mathfrak R$ に $S$ を添加して得られる環と呼び、$\mathfrak R[S]$ と書く。この環は、$S$ の元に関する $\mathfrak R$ 係数多項式のすべてからなる。ただし、その等式関係および演算関係は、$\mathfrak T$ における等式関係および演算関係によって定められる。これらの等式関係と演算関係がすでに定まっているため、特別な場合には、多項式全体の一部分、たとえば $S$ の元に関する $\mathfrak R$ 係数の線型形式だけを作れば十分である。そのような特別な場合は、たとえば（2 を参照）、$\mathfrak R$ が単位元をもつ環であり、$S$ が同じ単位元をもつ体であるときに生じる。このとき、$c_i$ が $S$ に属し、$\gamma_i$ が $\mathfrak R$ に属するすべての線型結合 $\sum c_i\gamma_i$ が環 $\mathfrak R[S]$ をなす。

しかし、環 $\mathfrak R$ だけが与えられ、さらに $\mathfrak R$ に含まれない元からなる系 $S$ が、それらの等式関係または相互の演算関係とともに与えられている場合でも、$\mathfrak R$ に $S$ を添加して得られる拡大環 $\mathfrak R[S]$ を構成できる。実際、$\mathfrak R$ と $S$ を含む環は常に存在する。すなわち、通常の演算規則を定めたうえで、$S$ の元に関する $\mathfrak R$ 係数多項式のすべてから形式的に作られる環がそれである。等号の定義には、演算規則を別にすれば、$\mathfrak R$ と $S$ において成り立つ等式関係をすべて含めるという制約以外には、なお完全な自由がある。したがって、$\mathfrak R$ 内および $S$ 内での元の等しさと、定めた演算とに基づき、等しい元から同じ仕方で作られる多項式だけは、必ず等しいと置かなければならない。このように構成された、$\mathfrak R$ と $S$ を含む環 $\mathfrak T$ は、$\mathfrak T$ 内で $\mathfrak R$ と $S$ から生成される環と一致する。\footnote[1]{Steinitz による商体の構成を参照されたい。}

\endgroup
%
\begingroup
\renewcommand{\thefootnote}{\arabic{footnote})}

\textbf{2. 体の拡大環。} これ以後は、単位元をもち、体 $P$ を部分環として含む環 $\mathfrak R$、すなわち体の上環を仮定する。このとき、$\mathfrak R$ の単位元 $e$ は同時に $P$ の単位元となり、$\mathfrak R$ は $\mathsf{P}$-加群となる。

$\overline{\mathsf{P}}$ が $\mathsf{P}$ の拡大体であるとき、$\mathfrak R$ に $\overline{\mathsf{P}}$ を添加して得られる拡大環 $\overline{\mathfrak R}=\mathfrak R[\overline P]$ を、次のように定義する。

$\mathfrak R$ と $\overline P$ との集合としての共通部分 $[\mathfrak R,\overline P]$ は $\mathsf{P}$ であり、さらに、$\mathsf{P}$ に属さない $\mathfrak R$ の元と、$P$ に属さない $\overline{\mathsf{P}}$ の元との間には相互の演算関係が存在しないものと仮定する。この仮定は、$\overline{\mathsf{P}}$ を $\mathsf{P}$ の同値な拡大体で置き換えるか、または $\mathfrak R$ を同値な拡大環で置き換えることによって、常に満たすことができる。\footnote[2]{Steinitz に従い、$\mathsf{P}$ の二つの拡大体は、両者の間に $\mathsf{P}$ の各元を固定する同型写像が存在するとき「同値」と呼ぶ。$\mathsf{P}$ の同値な拡大環も、これに応じて定義する。新しい記号を導入すれば、同値な拡大は常に作ることができる。この同値な拡大を用いる自由は、以下で本質的である。}

\endgroup
%
\begingroup
\renewcommand{\thefootnote}{\arabic{footnote})}

形式的に作られる次の形の線型結合
\[
   \bar c_{i_1}\gamma_{i_1}+\cdots+\bar c_{i_\sigma}\gamma_{i_\sigma}
\]
をすべて $\overline{\mathfrak R}$ の元と定める。ここで、$\bar c_{i_\lambda}$ は $\overline{\mathsf{P}}$ のすべての元を遍歴し、$\gamma_{i_\lambda}$ は $\mathfrak R$ のすべての元を遍歴する。\footnote[1]{一般に、$\mathfrak R$ または $\overline{\mathfrak R}$ の元はギリシア文字で、$\mathsf{P}$ または $\overline{\mathsf{P}}$ の元はラテン文字で表すものとする。} この定義は一意でなければならない。すなわち、$\gamma$ を $\mathfrak R$ においてそれと等しい元で置き換え、$\bar c$ を $\overline{\mathsf{P}}$ においてそれと等しい元で置き換えたならば、得られる $\overline{\mathfrak R}$ の元も元のものと等しいと定める。

和を、次式によって一意に定める。ここで「一意に」とは、これから定める等号のもとで一つの項をそれと等しいものに置き換えても、結果は同じであるという意味である。
\[
   \sum \bar c_i\gamma_i+\sum \bar d_i\gamma_i
   =\sum(\bar c_i+\bar d_i)\gamma_i .
\]
（$\bar c$ と $\bar d$ の中には零が現れてよく、それによって形式上は同じ $\gamma_i$ が現れる。）

積を、同じ意味で次式によって一意に定める。
\[
   \sum \bar c_i\gamma_i\cdot \sum \bar d_k\gamma_k
   =\sum \bar c_i\bar d_k\gamma_i\gamma_k .
\]

\endgroup
%

等号を次のように定める。$\gamma_{i_1},\ldots,\gamma_{i_\sigma}$ が $\mathsf{P}$ に関して線型独立であるならば、二つの元
\[
   \bar c_{i_1}\gamma_{i_1}+\cdots+\bar c_{i_\sigma}\gamma_{i_\sigma}
   \quad\text{および}\quad
   \bar d_{i_1}\gamma_{i_1}+\cdots+\bar d_{i_\sigma}\gamma_{i_\sigma}
\]
が等しいのは、すべての $\lambda$ について $\bar c_{i_\lambda}$ と $\bar d_{i_\lambda}$ が $\overline{\mathsf{P}}$ において等しいとき、かつそのときに限る。言い換えれば、$\mathsf{P}$ に関して線型独立な $\mathfrak R$ の元は、$\overline{\mathsf{P}}$ に関しても線型独立であると定める。
%

和と積の定義によって、これで $\overline{\mathfrak R}$ のすべての元に対する等号の定義が与えられる。実際、それらの元を $\mathfrak R$ の線型独立な元によって表せることが従う。というのも、$\mathfrak R$ における関係
\[
   \mu=c_1\gamma_1+\cdots+c_s\gamma_s\quad\text{in }\mathfrak R
\]
から、$\bar b=\bar b e$ を用い、積の定義に従って両辺に $\bar b$ を掛けると、
\[
   \bar b\mu=\bar b e\cdot(c_1\gamma_1+\cdots+c_s\gamma_s)
      =\bar b c_1\gamma_1+\cdots+\bar b c_s\gamma_s,
\]
が得られ、したがって和の定義により、すべての元を $\mathfrak R$ の線型独立な元によって表せるからである。
%
\begingroup
\renewcommand{\thefootnote}{\arabic{footnote})}

したがって、すべての元が $\mathfrak R$ の線型独立な元によって表されているとき、等号の定義は $\overline{\mathsf{P}}$ に\mbox{おける}\mbox{係数}の等号となる。それゆえ、$\mathfrak R$ と $\overline{\mathfrak R}$ の双方、または $\overline{\mathsf{P}}$ と $\overline{\mathfrak R}$ の双方に属する元については、\mbox{それぞれ}そこで成り立つ等号と一致する。一方、$\overline{\mathfrak R}$ のその他の元については、共通部分に関する仮定と、\mbox{さらに}置いた条件とにより、$\mathfrak R$ および $\overline{\mathsf{P}}$ で成り立つ等号は、これらの元の定義で要求した一意性以上のことを何も述べない。\footnote[2]{この仮定がなければ、$\overline P$ に属するが $P$ には属さない元 $\bar c$ が少なくとも一つ存在して、$\bar c=\bar c e=\gamma$ が成り立つことになる。したがって、$\mathfrak R$ における等号の定義により、$e$ と $\gamma$ は $\mathsf{P}$ に関しては線型独立であるが、$\overline{\mathsf{P}}$ に関しては線型従属となる。$\overline{\mathsf{P}}$ の元と $\mathfrak R$ の元との間の関係のため、上の等号の定義は不可能になる。たとえば、基礎体を中間体で置き換えたときに有限拡大体の次数が下がることを考えればよい。\mbox{これに対して}、上の拡大は新しい\mbox{零因子}が現れうることによって特徴づけられる。たとえば、有理係数体 $\mathsf{P}$ 上の多項式環を $x^2+1$ で割った剰余類環は、\mbox{零因子}をもたない環であって体 $\mathsf{P}(\sqrt{-1})$ と同型であるが、\mbox{係数}を $\mathsf{P}(\sqrt{-1})$ に取る多項式環へ移ると、$(x+i)(x-i)=(x^2+1)$ である一方、\mbox{どちら}の因子も $x^2+1$ では割り切れないことに対応して、\mbox{零因子}が現れる。同様のことは、たとえば任意の有限数体についても成り立つ。これは、\mbox{高次複素数}について Dedekind が最初に表明した見解である。この点については \S~6, 4 の末尾を参照されたい。} 同様に、和と積の一意性を要求しても、それ以上の関係は生じない。等しい元は\mbox{係数ごとに}等しいと仮定でき、したがって形式的に同じ和と積の表示を与えるからである。

これで、環演算のすべての公理が満たされることを直ちに確かめられる。したがって、$\mathfrak R$ に $\overline P$ を添加して得られる拡大環 $\mathfrak R[\overline{\mathsf{P}}]$ は常に存在する。

\endgroup
%
\begingroup
\renewcommand{\thefootnote}{\arabic{footnote})}

\textbf{3. $\mathfrak R$ と $\overline{\mathfrak R}$ におけるイデアル。} $\mathfrak a$ が $\mathfrak R$ のイデアルであるとき、加群積 $\overline{\mathfrak R}\mathfrak a$ は $\overline{\mathfrak R}$ のイデアル $\overline{\mathfrak a}$ となり、これを $\mathfrak a$ の拡大イデアルという。同時に、$\overline{\mathfrak a}$ は $\overline{\mathfrak R}$ 内で $\mathfrak a$ から生成されるイデアルである。それは、和を毎回有限個の元について取った線型結合 $\sum \bar c_i a_i$ のすべてからなり、ここで $a_i$ は $\mathfrak a$ のすべての元を、$\bar c_i$ は $\overline{\mathsf{P}}$ のすべての元を遍歴する。

$\overline{\mathfrak b}$ が $\overline{\mathfrak R}$ のイデアルであるとき、共通部分 $[\overline{\mathfrak b},\mathfrak R]$ は $\mathfrak R$ のイデアル $\mathfrak b$ となり、これを $\overline{\mathfrak b}$ の縮約イデアルという。\footnote[1]{83頁の注6で引用した H. Grell の論文における呼称である。} $\mathfrak R$ のすべてのイデアル $\mathfrak a$ は、その拡大イデアル $\overline{\mathfrak a}$ の縮約イデアルである。実際、$[\mathfrak R,\overline{\mathfrak a}]$ の任意の元 $a$ は $\sum \bar c_i a_i$ の形をし、ここで $a_i$ は $\mathsf{P}$ または $\overline{\mathsf{P}}$ に関して線型独立であると仮定してよい。したがって、元 $a,a_i$ は $\overline{\mathfrak R}$ において線型従属となり、等号の定義により $\mathfrak R$ においても線型従属となる。それゆえ、$a$ は $\mathfrak a$ に属する。逆も成り立つので、$\mathfrak a=[\mathfrak R,\overline{\mathfrak a}]$ である。

\endgroup
%
\begingroup
\renewcommand{\thefootnote}{\arabic{footnote})}

$\overline{\mathfrak p}$ が $\overline{\mathfrak R}$ の素イデアルであるとき、$\mathfrak p=[\mathfrak R,\overline{\mathfrak p}]$ は $\mathfrak R$ の素イデアルとなる。実際、第二同型定理\footnote[2]{「イデアル論」\S~4, 3。同論文では、この定理は同じ環のイデアルについてだけ述べられている。しかし、同じ考察により、上の形の定理が $\overline{\mathfrak R}$ の任意のイデアル $\overline{\mathfrak a}$ についても成り立つことが分かる。実際、環準同型 $\overline{\mathfrak R}\sim\overline{\mathfrak R}/\overline{\mathfrak a}$ は、$\overline{\mathfrak R}$ の部分環 $\mathfrak R$ に対して環準同型 $\mathfrak R\sim(\mathfrak R,\overline{\mathfrak a})\,|\,\overline{\mathfrak a}$ を\mbox{誘導する}。この準同型のもとで零元に対応するのは、$[\mathfrak R,\overline{\mathfrak a}]$ のすべての元であり、またそれらだけである。したがって、環同型 $(\mathfrak R,\overline{\mathfrak a})\,|\,\overline{\mathfrak a}\simeq \mathfrak R\,|\,[\mathfrak R,\overline{\mathfrak a}]$ が\mbox{得られる}。このとき、右辺と左辺の類は $\mathfrak R$ の同じ元によって代表させることができる。証明の各段階で、各元にはそれが代表する類が\mbox{対応づけられる}\mbox{からである}。}により、剰余類環 $\mathfrak R\,|\,\mathfrak p$ は $\overline{\mathfrak R}\,|\,\overline{\mathfrak p}$ の一つの部分環と同型であり、したがって零因子をもたない環である。すなわち、次の環同型が成り立つ。
\[
   (\mathfrak R,\overline{\mathfrak p})/\overline{\mathfrak p}
     \simeq \mathfrak R/[\mathfrak R,\overline{\mathfrak p}],
\]
ここで $(\mathfrak R,\overline{\mathfrak p})$ は、加群 $\mathfrak R$ と $\overline{\mathfrak p}$ の和、すなわち加群の意味での両者の最大公約因子を表す。

\endgroup
%

\begin{center}
\textbf{\S 2. 有限階数の環。直和としての表示。}
\end{center}

\textbf{1. 有限階数の環。} $\mathfrak A$ が有限 $\mathsf{P}$-加群であるならば、$\mathfrak A$ は $P$ に関して有限階数、たとえば階数 $k$ をもつ。すなわち、$\mathfrak R$ には $P$ に関して線型独立な元が $k$ 個存在するが、それより多くは存在しない。したがって、$\mathfrak A$ の元からなる、$P$ に関して線型独立な $k$ 元系は、どれも $\mathfrak A$ の加群基底をなす。有限階数の $\mathsf{P}$-加群 $\mathfrak B$ が、同じ有限階数をもつ別の $P$-加群 $\mathfrak A$ で割り切れる、すなわち $\mathfrak B\equiv0\,(\mathfrak A)$ であるならば、両者は\mbox{同一である}。

\rightskip=0pt plus 4em\relax
これ以後、$\mathfrak R$ は $P$ に関して有限階数 $n$ をもつものと仮定する。このとき、$\mathfrak R$ 内の任意の \mbox{$\mathsf{P}$-加群}は\mbox{有限階数}をもつ。したがって、$\mathfrak R$ 内の $\mathsf{P}$-加群からなる任意の真の約鎖では、各加群の\mbox{階数}は\mbox{直前}の加群の階数より大きく、真の倍鎖では\mbox{直前}の加群の階数より小さくなければならない。\mbox{したがって}、\mbox{これら}の鎖はいずれも有限段で終わる。とくに、$\mathfrak R$ のすべてのイデアルからなる系に\mbox{ついて}「\mbox{二重鎖定理}」が\mbox{成り立つ}。
%
\begingroup
\renewcommand{\thefootnote}{\arabic{footnote})}

ここから、$\mathfrak R$ の素イデアル $\mathfrak p$ は、単位イデアルと異なる真の約イデアルをもたないことが\mbox{従う}\footnote[1]{「イデアル論」\S~7, 2。}。\mbox{したがって}、単位イデアルと異なる素イデアル $\mathfrak p$ を法とする剰余類環 $\mathfrak R\,|\,\mathfrak p$ は体となる。$\mathfrak q$ が $\mathfrak p$ に\mbox{属する}\mbox{準素イデアル}、\mbox{すなわち} $\mathfrak q\equiv0\,(\mathfrak p)$ かつ $\mathfrak p^{\varrho}\equiv0\,(\mathfrak q)$ であるならば、剰余類環 $\mathfrak R\,|\,\mathfrak q$ は\mbox{準素環となる}。\mbox{すなわち}、各零因子のある冪が零となり、この場合には各零因子の $\varrho$ 乗がすでに零となる環である。\mbox{さらに}、$\mathfrak R\,|\,\mathfrak q$ の元はすべて零因子または単元である。最後の主張は、$\mathfrak p$ で割り切れない任意のイデアルが $\mathfrak p$ と\mbox{互いに素}となり、\mbox{したがって} $\mathfrak q$ とも\mbox{互いに素}となることから直ちに従う。

\endgroup
%
\begingroup
\renewcommand{\thefootnote}{\arabic{footnote})}

二重鎖定理と単位元の存在からさらに、$\mathfrak R$ の任意のイデアル $\mathfrak a$ は、有限個の、どの二つも\mbox{互いに素}な\mbox{準素イデアル}の積として一意に表されることが従う\footnote[2]{「イデアル論」\S~7, 3、定理 V。}。これらが\mbox{互いに素}であるため、この積は\mbox{最小公倍数}に等しく、剰余類環 $\mathfrak R\,|\,\mathfrak a$ を準素環の直和として表すことに対応する。\footnote[3]{「イデアル論」\S~4, 5。本稿の現在のまとめ方で付け加えた内容は、文献の多くの箇所に散在している。} すなわち、$\mathfrak R\,|\,\mathfrak a$ はこれらの環の\mbox{最大公約因子}となり、$\mathfrak R\,|\,\mathfrak a$ の各元は、この直和において一意に和として表される。

\endgroup
%
\begingroup
\renewcommand{\thefootnote}{\arabic{footnote})}

\textbf{2. 有限階数の環を準素環の直和として表示すること。} $\mathfrak R$ の零イデアルの表示が
\[
   (0)=\mathfrak q_1\cdot\mathfrak q_2\cdots\mathfrak q_r
      =[\mathfrak q_1,\mathfrak q_2,\ldots,\mathfrak q_r]
\]
によって与えられているとし、これに対応する直和表示を
\[
   \mathfrak R=\mathfrak R_1+\mathfrak R_2+\cdots+\mathfrak R_r,
   \qquad
   \mathfrak R_i\mathfrak R_k=(0)\quad\hbox{$i\ne k$ のとき} i\ne k,
   \qquad
   \mathfrak R_i^2=\mathfrak R_i,
   \qquad
   \mathfrak R\mathfrak R_i=\mathfrak R_i
\]
とする。このとき $\mathfrak R_i$ は $\mathfrak R$ のイデアルとなり、実際、
\[
   \mathfrak R_i
   =\mathfrak q_1\cdots\mathfrak q_{i-1}\mathfrak q_{i+1}\cdots\mathfrak q_r
   =[\mathfrak q_1,\ldots,\mathfrak q_{i-1},\mathfrak q_{i+1},\ldots,\mathfrak q_r],
\]
\[
   \mathfrak q_i=\mathfrak R_1+\cdots+\mathfrak R_{i-1}+\mathfrak R_{i+1}+\cdots\mathfrak R_r
\]
となる。また、$\mathfrak R_i$ は剰余類環 $\mathfrak R\,|\,\mathfrak q_i$ と同型となる。すなわち、$\mathfrak R_i$ の各元 $\gamma_i$ には、$\mathfrak R\,|\,\mathfrak q_i$ において $\gamma_i$ \mbox{によって}代表される剰余類が対応する。\footnote[4]{「イデアル論」\S~4, 5。} $\mathfrak a$ を $\mathfrak R$ の任意のイデアルとすると、分配法則により
\[
   \mathfrak a=\mathfrak R\mathfrak a
   =\mathfrak R_1\mathfrak a+\cdots+\mathfrak R_r\mathfrak a
   =\mathfrak a_1+\cdots+\mathfrak a_r
\]
を得る。ここで $\mathfrak a_i=\mathfrak R_i\mathfrak a_i=\mathfrak R\mathfrak a_i$ であるから、成分 $\mathfrak a_i$ は $\mathfrak R_i$ のイデアルであると同時に $\mathfrak R$ の\mbox{イデアルでもある}。イデアルとして、$\mathfrak a_i$ は $\mathfrak R_i$ 自身と同様に有限 $\mathsf{P}$-加群である。和が直和であるため、$\mathfrak a_i$ の\mbox{階数}の和は $\mathfrak a$ の階数に等しく、$\mathfrak R_i$ の\mbox{階数}の和は $\mathfrak R$ の階数に等しい。

\endgroup
%

単位元の和表示が
\[
   e=\varepsilon_1+\cdots+\varepsilon_r,
   \qquad
   \varepsilon_i\varepsilon_k=0\quad\hbox{$i\ne k$ のとき} i\ne k,
   \qquad
   \varepsilon_i^2=\varepsilon_i,
   \qquad
   \varepsilon_i\equiv e\,(\mathfrak q_i)
\]
によって与えられているとき、$\mathfrak R$ の任意の元 $\gamma$ の表示は
\[
   \gamma=\gamma e=\gamma\varepsilon_1+\cdots+\gamma\varepsilon_r
      =\gamma_1+\cdots+\gamma_r
\]
によって一意に与えられる。ここで成分 $\gamma_i=\gamma_i\varepsilon_i$ は $\mathfrak R_i$ の元であり、同時に $\mathfrak R_i=\mathfrak R\varepsilon_i$ が従う。\mbox{「直交関係」から}
\[
   \gamma\delta=(\gamma_1\varepsilon_1+\cdots\gamma_r\varepsilon_r)
      (\delta_1\varepsilon_1+\cdots+\delta_r\varepsilon_r)
   =\gamma_1\delta_1\varepsilon_1+\cdots+\gamma_r\delta_r\varepsilon_r
\]
を得る。したがって $\gamma_i\delta_i$ は積の第 $i$ 成分となり、同様に $\gamma_i+\delta_i$ は和の第 $i$ 成分となる。これらの事実は、$\mathfrak R$ から $\mathfrak R_i$ への準同型からも従う。
%

環 $\mathfrak R_i$ は $\mathsf{P}$ と同型な部分体 $\mathsf{P}_i=\mathsf{P}\varepsilon_i$ を含み、$\mathsf{P}_i$ に関して有限階数をもつ。この階数は、$\mathfrak R_i$ が \mbox{$\mathsf{P}$-加群}としてもつ階数に等しい。実際、$P$ の各元 $c$ について $c\varepsilon_i\equiv ce\equiv c\,(\mathfrak q_i)$ となる。したがって、$c\ne0$ \mbox{ならば} $c\varepsilon_i\ne0$ であり、$P$ から $P_i$ への準同型は同型となる。さらに、$\mathfrak R_i$ の元の $P$ に関する任意の\mbox{線型従属関係}は、$\varepsilon_i$ を掛けることにより $\mathsf{P}_i$ に関する\mbox{線型従属関係}へ移る。逆に、$\mathsf{P}_i=\mathsf{P}\varepsilon_i$ であるため、$\mathsf{P}_i$ に関する\mbox{任意}の関係は同時に $\mathsf{P}$ に関する関係を与える。したがって、$\mathsf{P}$ に関する階数と $\mathsf{P}_i$ に関する階数は\mbox{一致する}。
%

\textbf{3. 拡大環の直和表示。} 
\[
   \mathfrak R=\mathfrak R_1+\cdots+\mathfrak R_r,
\]
であるとする。このとき、$\overline{\mathfrak R}$ には
\[
   \overline{\mathfrak R}=\overline{\mathfrak R}_1+\cdots+\overline{\mathfrak R}_r,
\]
という表示が成り立つ。ここで $\overline{\mathfrak R}_i$ は、$\overline{\mathfrak R}$ における $\mathfrak R_i$ の拡大イデアルを表し、同時に $\overline{\mathfrak R}$ 内で $\mathfrak R_i$ \mbox{から生成される環}を表す。実際、$(0)=\mathfrak q_1\cdots\mathfrak q_r$ に $\overline{\mathfrak R}$ を掛けると、積表示 $(\overline{0})=\overline{\mathfrak q}_1\cdots\overline{\mathfrak q}_r$ が得られ、\mbox{したがって} $\overline{\mathfrak q}_i$ は再びどの二つも互いに素となる。それゆえ $\overline{\mathfrak R}$ の直和表示が得られる。そして、この表示の第 $i$ \mbox{成分}は
\[
   \overline{\mathfrak q}_1\cdots\overline{\mathfrak q}_{i-1}\,
   \overline{\mathfrak q}_{i+1}\cdots\overline{\mathfrak q}_r
\]
で与えられるから、これは $\mathfrak R_i$ の拡大イデアルに等しい。しかし \S~1, 3 により、この拡大イデアルは、$\mathfrak R_i$ に $\overline{\mathsf{P}}_i=\overline{\mathsf{P}}\varepsilon_i$ を添加して得られる拡大環 $\mathfrak R_i[\overline{\mathsf{P}}_i]$ と同一である。したがって、拡大に際しては個々の成分をそれぞれ拡大すれば足りる。その際、$\overline{\mathfrak R}_i$ は一般にはもはや準素環ではない。
%
\begingroup
\renewcommand{\thefootnote}{\arabic{footnote})}

\textbf{4. 第一種および第二種の素イデアル。} 1 により、単位イデアルと異なる任意の素イデアル $\mathfrak p$ を法とする剰余類環 $\mathfrak R\,|\,\mathfrak p$ は体となる。この体は $\mathsf{P}$ と同型な部分体 $(\mathsf{P})=(\mathsf{P},\mathfrak p)\,|\,\mathfrak p\simeq \mathsf{P}\,|\,[\mathsf{P},\mathfrak p]\simeq \mathsf{P}$ を含む。というのも、$\mathfrak p$ は $\mathsf{P}$ の零でない元を含みえないからである。したがって、$(\mathsf{P})$ は、$\mathsf{P}$ の元を代表元にもつ剰余類全体にほかならない\footnote[1]{第二同型定理についての \S~1, 3 の注を参照。}。剰余類環 $\mathfrak R\,|\,\mathfrak p$ は $(\mathsf{P})$ に関して有限階数をもち、したがって有限代数拡大体となる。この拡大が第一種であるか第二種であるかに応じて\footnote[2]{Steinitz のよく知られた用語では、代数拡大は、その各元が、ある適当な拡大体においてすべて相異なる一次因子に分解する素関数の零点であるとき第一種と呼ばれ、そうでないとき第二種と呼ばれる。}、$\mathfrak p$ を第一種または第二種の素イデアルと呼ぶ。

\endgroup
%

\begin{center}
\textbf{\S~3. 完全可約環。}
\end{center}

1. 環 $\mathfrak R$ は、その零イデアルの積表示（\S~2, 2）における準素成分がすべて素イデアルであるとき、\emph{完全可約}であるという。また、\S~2, 1 および 2 により、この条件は $\mathfrak R$ が有限個の体の直和として表されることと同値である。完全可約環は、その素イデアルがすべて第一種であるとき（\S~2, 4）、第一種の完全可約環といい、そうでない場合には第二種の完全可約環という。

$P$ が完全体ならば、第一種の完全可約性しか起こらない。とりわけ、$P$ が代数閉体ならば常にそうである。以下では特に、$\overline P$ は $P$ の代数閉包を表すものとし、$\overline{\mathfrak R}=\mathfrak R[\overline P]$ とする。
%

2. \textbf{定理。} \emph{環 $\mathfrak R$ が第一種の完全可約環であることと、代数閉体 $\overline P$ による拡大環 $\overline{\mathfrak R}=\mathfrak R[\overline P]$ が完全可約（第一種）であることとは同値である。}

証明は三段階に分けて行う。
%

2a. $\overline{\mathfrak R}$ の完全可約性、より一般には任意の拡大環の完全可約性から、$\mathfrak R$ の完全可約性（第一種または第二種）が従う。実際、仮定により、$\overline{\mathfrak R}$ の零イデアルは素イデアルの最小公倍数として
\[
        (\overline{0})=[\overline{\mathfrak p}_1,\ldots,\overline{\mathfrak p}_s]
\]
と表される。共通部分をとることにより、$\mathfrak R$ では零イデアルの表示
\[
        (0)=[[\mathfrak R,\overline{\mathfrak p}_1],\ldots,[\mathfrak R,\overline{\mathfrak p}_s]]
\]
が得られ、イデアル $[\mathfrak R,\overline{\mathfrak p}_i]$ は \S~1, 3 により素イデアルである。
%

2b. $\mathfrak R$ の第一種の完全可約性から、$\overline{\mathfrak R}$ の第一種の完全可約性が従う。さらに一般に、$\mathfrak R$ と $\overline{\mathfrak R}$ の任意の中間環も第一種の完全可約環となる。

\S~2, 3 により、$\mathfrak R$ の直和表示には $\overline{\mathfrak R}$ の対応する直和表示が存在し、その各成分 $\overline{\mathfrak R}_i$ は拡大環 $\mathfrak R_i[\overline P_i]$ となる。ここで $\overline P_i$ は、成分 $\mathfrak R_i$ の $P$ と同型な部分体 $P_i$ の拡大体であり、$\overline P$ と同型である。したがって、個々の成分 $\mathfrak R_i$ だけを考えればよい。よって一般性を失わず、$\mathfrak R$ は体であり、しかも部分体 $P$ の第一種有限拡大体であると仮定できる。したがって、$\mathfrak R$ は $P$ に第一種の原始元を添加することによって生成される。言い換えると、$\mathfrak R$ は、一つの不定元 $x$ の多項式環 $P[x]$ における第一種の素関数 $f(x)$ による剰余類環 $\mathsf{P}[x]\,|\,f(x)$ と同型である。$\Omega$ を $\mathsf{P}$ の拡大体とすると、$\Omega[x]\,|\,f(x)$ は $\mathfrak R[\Omega]$ と同型となる。実際、$\mathsf{P}[x]$ から $\Omega[x]$ へ移るとき、剰余類環の階数は変わらず $f(x)$ の次数に等しいので、これはまさに \S~1, 2 で述べた拡大環の構成である。$\Omega[x]$ において $f(x)$ は、どの二つも互いに素な第一種の素関数に分解し、その各々が $\Omega[x]$ の素イデアルを生成する。したがって、$\Omega[x]\,|\,f(x)$ は第一種の完全可約環となり、同型であることから $\mathfrak R[\Omega]$ も第一種の完全可約環となる。これは任意の拡大体 $\Omega$、とくに代数閉体 $\overline P$ について成り立つので、2b が証明された。
%

2c. $\mathfrak R$ が第二種の完全可約環であるならば、$\overline{\mathfrak R}$ は完全可約ではない。

2b と同様に、$\mathfrak R$ を体、それも $P$ の第二種有限拡大体であると仮定すれば十分である。$\overline{\mathfrak R}$ の準素環の直和表示に、少なくとも一つの真の準素成分が現れることを示さなければならない。そのためには、ある冪が零となる $\overline{\mathfrak R}$ の零でない元が存在することを示せば十分である。実際、$\bar\beta\ne0$ ならば、少なくとも一つの成分 $\bar\beta_\lambda\ne0$ が零ではなく、$\bar\beta^\sigma=0$ ならば $\bar\beta_\lambda^\sigma=0$ である（\S~2, 2、$\overline{\mathfrak R}$ から $\overline{\mathfrak R}_\lambda$ への準同型）。したがって、$\overline{\mathfrak R}_\lambda$ は体ではない。
%

$\gamma$ を $\mathfrak R$ の指数 $f\geq 1$ の第二種の元とし、
\[
        F(\gamma)=\gamma^{kp^f}+c_1\gamma^{(k-1)p^f}+\cdots+c_k=0
\]
を、$\gamma$ が $\mathsf{P}$ に関して満たす最低次数の方程式とする。ここで $p$ は $\mathsf{P}$ の標数を表す。したがって、$\mathfrak R$ の元
\[
        e,\gamma,\gamma^2,\ldots,\gamma^{kp^f-1}
\]
は $P$ に関して線型独立であり、等号の定義により、$\overline{\mathfrak R}$ においても $\overline P$ に関して線型独立のままである。そこで
\[
        G(\gamma)=\gamma^{kp^{f-1}}+\sqrt[p]{c_1}\,\gamma^{(k-1)p^{f-1}}+\cdots+\sqrt[p]{c_k}
\]
とおく。このとき $G(\gamma)$ は $\overline{\mathfrak R}$ の零でない元となる。実際、$p>1$ かつ $f\geq1$ であるから、$kp^{f-1}$ は常に $kp^f$ より小さく、$G(\gamma)$ に現れる $\gamma$ の冪の間に従属関係はありえない。一方、$(G(\gamma))^p$ は $F(\gamma)$ に等しく、したがって零となる。これで 2c が証明された。

2a と 2c から、$\overline{\mathfrak R}$ の完全可約性（第一種）は $\mathfrak R$ の第一種の完全可約性を導く。2b から逆が従う。これで 2 の定理が証明された。
%
\begingroup
\renewcommand{\thefootnote}{\arabic{footnote})}

3. \textbf{系。} \emph{$\mathfrak R$ が第一種の完全可約環であるならば、$\overline{\mathfrak R}$ は階数一の環、したがって $\overline P$ と同型な体の直和となる。このような階数一の環への分解は、$\Omega$ を $P$ の有限拡大体とする環 $\mathfrak R[\Omega]$ の中にすでに存在する。}

実際、2b により $\overline{\mathfrak R}$ は体の直和となる。これらの体は、$\overline P$ と同型な部分体（\S~2, 2）の有限拡大体であるが、$\overline P$ が代数閉体であるため、その部分体と一致する。したがって
\[
        \overline{\mathfrak R}=\overline P\bar e_1+\cdots+\overline P\bar e_n,
        \qquad e=\bar e_1+\cdots+\bar e_n;
\]
となり、$e$ の成分 $\bar\varepsilon_\lambda$ は同時に $\overline{\mathfrak R}$ の（線型独立な）加群基底をなす。

\S~1, 2 に従って $\bar\varepsilon_\lambda$ を、$\mathfrak R$ の元 $\gamma$ の $\overline P$ 係数による線型結合として表すと、係数全体が $P$ の有限拡大体 $\Omega$ を定める。したがって $\bar\varepsilon_\lambda$ は $\mathfrak R[\Omega]$ に属し、加群基底の性質から
\[
        \mathfrak R[\Omega]=\Omega\bar e_1+\cdots+\Omega\bar e_n
\]
となる。ゆえに、$\mathfrak R[\Omega]$ の中ですでに階数一の環への分解が成り立ち、$\mathfrak R[\Omega]$ と $\overline{\mathfrak R}$ の任意の中間環における直和分解不能な成分は、$\Omega\bar\varepsilon_\lambda$ の拡大環として得られる。\footnote[1]{このような最小の拡大体 $\Omega$ は、$\mathfrak R$ の成分 $\mathfrak R_i=\mathfrak R e_i$ によって定まるガロア拡大体の合成体として特徴づけられる。実際、2b に従って $P[x],|,f_i(x)$ が各成分 $\mathfrak R_i$ と同型になるように $f_i(x)$ を選ぶ。$\Omega^{(i)}$ を、$\overline{\mathsf{P}}$ に含まれ、$f_i(x)$ を一次因子に分解するのにちょうど十分なガロア拡大体とし、$\Omega$ をすべての $\Omega^{(i)}$ の合成体とする。$f_i(x)$ の一次因子への分解には、$\Omega^{(i)}[x]\,|\,f_i(x)$ を階数一の体の直和として表す表示が対応し、しかも $\Omega^{(i)}$ はそのような表示が成り立つ最小の拡大体である。したがって、同型であることから $\mathfrak R_i[\Omega^{(i)}e_i]$ についても同じことが成り立ち、ひいては $\mathfrak R[\Omega]$ についても成り立つ。}

\endgroup
%
\begingroup
\renewcommand{\thefootnote}{\arabic{footnote})}

\begin{center}
\textbf{\S~4. 有限階数の環における行列表現、}\\
\textbf{トレースおよび判別式。}
\end{center}

この節と次節では、数体および関数体のオーダーも含めるため、基礎となる環に関する仮定をこれまでよりも少し一般的なものとする。本質的には既知の事柄を一般的かつ一様な形で定式化することが目的である。

\emph{仮定。} $\mathfrak R$ は零因子をもたない環 $\Sigma$ の拡大環であり、$\mathfrak R$ の単位元はすでに $\Sigma$ に属するものとする。$\mathfrak R$ の中の任意の $\Sigma$-加群、したがって特に $\mathfrak R$ 自身には、$\Sigma$ に関して線型独立な元からなる有限加群基底が少なくとも一つ存在するものとする。さらに、$\alpha_1,\ldots,\alpha_s$ と並んで $\beta_1,\ldots,\beta_s$ が同じ $\Sigma$-加群の線型独立な加群基底となるのは、その変換行列式が $\Sigma$ の単元であるとき、かつそのときに限る。

この仮定は、これまで考察してきた有限階数の環については明らかに満たされ、その場合には $\Sigma$ は体 $\mathsf{P}$ となる。また、数体および関数体のオーダーについても満たされ、その場合には $\Sigma$ はそれぞれ有理整数の環または汎関数領域となる。\footnote[2]{『イデアル論』\S~3 を参照。}

\endgroup
%
\begingroup
\renewcommand{\thefootnote}{\arabic{footnote})}

\textbf{1. $\mathfrak R$ に準同型な行列環。} $\alpha_1,\ldots,\alpha_s$ をイデアル $\mathfrak a$ の（線型独立な）$\Sigma$-加群基底とし、$\gamma$ を $\mathfrak R$ の任意の元とする。このとき $\gamma\alpha_i$ も $\mathfrak a$ に属するので、$\Sigma$ の係数をもつ次の等式が成り立つ。\footnote[3]{$\Sigma$ の元はラテン文字で表す。}
\[
\gamma\alpha_i=c_{1i}\alpha_1+\cdots+c_{si}\alpha_s,
\qquad (i=1,\ldots,s);
\]
あるいは、$(\tau_1,\ldots,\tau_s)$ が各場合にこれらの元からなる一行行列を表すものとすれば、行列の形では
\[
(\gamma\alpha_1,\ldots,\gamma\alpha_s)
=(\alpha_1,\ldots,\alpha_s)
\begin{pmatrix}
 c_{11}&\cdots&c_{1s}\\
 \vdots&&\vdots\\
 c_{s1}&\cdots&c_{ss}
\end{pmatrix}
=(\alpha_1,\ldots,\alpha_s)C
\]
となる。

\emph{$\gamma$ が $\mathfrak R$ のすべての元を遍歴するとき、基底 $a_i$ によって対応づけられる行列 $C$ 全体は、$\mathfrak R$ の準同型像である行列環 $R_{\mathfrak a}$ をなす。また $R_{\mathfrak a}$ は剰余類環 $\mathfrak R\,|\,((0):\mathfrak a)$ と同型である。ここで $(0):\mathfrak a$ は、零イデアルの $\mathfrak a$ によるイデアル商を表す。} 実際、差 $\beta-\gamma$ には明らかに $B-C$ が対応し、積についても同様である。というのも
\[
(\beta\gamma\alpha_1,\ldots,\beta\gamma\alpha_s)
=(\beta\alpha_1,\ldots,\beta\alpha_s)C
=(\alpha_1,\ldots,\alpha_s)BC\footnotemark[1]
\]
となるからである。
\footnotetext[1]{この表記は、$\mathfrak R$ が非可換環であり、$\Sigma$ の元がすべての元と可換であると仮定する場合にも、そのまま保たれるように定められている。}

とりわけ、$\Sigma$ の元 $c$ には対角行列 $cE$ が対応する。ここで $E$ は、単位元 $e$ に対応する行列
\[
\begin{pmatrix}
 e&&\\
 &\ddots&\\
 &&e
\end{pmatrix}
\]
を表す。さらに、零行列が対応するのは $\gamma\mathfrak a$ が零イデアルとなる元 $\gamma$ のすべてであり、かつそれらに限るので、上に述べた同型が得られる。

\emph{$\mathfrak a$ の基底 $\bar a$ は、$R_{\mathfrak a}$ と同型で、しかも同値な行列環 $R_{\bar a}=P^{-1}R_{\mathfrak a}P$ を生成する。} 実際、$\bar\alpha_1,\ldots,\bar\alpha_s$ を別の線型独立な $\Sigma$-加群基底とすれば、
\[
(\bar a_1,\ldots,\bar a_s)=(a_1,\ldots,a_s)P,
\]
であり、行列 $P$ の行列式は $\Sigma$ の単元である。従って
\[
(\gamma\bar\alpha_1,\ldots,\gamma\bar\alpha_s)
=(\alpha_1,\ldots,\alpha_s)CP
=(\bar a_1,\ldots,\bar a_s)P^{-1}CP,
\]
すなわち
\[
R_{\bar a}=P^{-1}R_{\mathfrak a}P.
\]

\endgroup
%
%
%
\begingroup
\renewcommand{\thefootnote}{\arabic{footnote})}

\textbf{2. イデアル類と表現類。} $\mathfrak R$ の二つのイデアル $\mathfrak a$ と $\mathfrak b$ は、$\mathfrak R$-加群とみなして同型であるとき、同じイデアル類に属する。すなわち、それらの元の間に、差と $\mathfrak R$ の同じ元による乗法とが対応するような一対一対応があり、任意の $\gamma\in\mathfrak R$ に対して $\alpha\simeq\beta$ から常に $\gamma\alpha\simeq\gamma\beta$ が従うときである。

1. の行列表現によって生成される二つの環 $R_{\mathfrak a}$ と $R_{\bar\alpha}$ は、同値であるとき、同じ表現類に属する。すなわち、
\[
R_{\bar a}=P^{-1}R_{\mathfrak a}P
\]
となるような（ユニモジュラー）行列 $P$ が存在するときである。\footnote[2]{特殊な場合には解決されている次の問題、すなわち、すべての $\mathfrak R$ の準同型像である行列環が、必要なら線型従属な基底も許すことによって、1. に従い $\mathfrak R$ のあるイデアルから生成されるかどうかという問題には、ここでは立ち入らない。定義上、表現類を構成するのは、1. に従い $\mathfrak R$ のイデアルから生成される行列環だけである。}

\endgroup
%
%
%

\emph{イデアル類と表現類は一対一に対応する。} 1. では、同一のイデアルの異なる加群基底から生成される環 $R_{\mathfrak a}$ と $R_{\bar\alpha}$ が、同じ表現類に属することを示した。しかしさらに、同じイデアル類に属する異なるイデアル $\mathfrak a$ と $\mathfrak b$ からも、同じ表現類に属する環が生成される。実際、$\mathfrak a$ の基底 $a_1,\ldots,a_s$ に同型によって対応する $\mathfrak b$ の基底を $\beta_1,\ldots,\beta_s$ とする。このとき $\gamma\alpha_i$ と $\gamma\beta_i$ が互いに対応し、同様に
\[
c_{1i}\alpha_1+\cdots+c_{si}\alpha_s
\quad\text{と}\quad
c_{1i}\beta_1+\cdots+c_{si}\beta_s
\]
も互いに対応する。したがって、環 $R_{\mathfrak a}$ と $R_{\beta}$ は同一になる。
%
%
%

逆に、$R_{\mathfrak a}$ と $R_{\overline{\beta}}$ が同じ表現類に属するとする。すなわち
\[
R_{\overline{\beta}}=P^{-1}R_{\mathfrak a}P
\]
であり、ここで $a_i$ は $\mathfrak a$ の基底を、$\overline{\beta}_i$ は $\mathfrak b$ の基底をなすものとする。このとき
\[
(\beta_1,\ldots,\beta_s)
=(\overline{\beta}_1,\ldots,\overline{\beta}_s)P^{-1}
\]
によっても $\mathfrak b$ の一つの基底が得られ、
\[
R_{\beta}=P R_{\overline{\beta}}P^{-1}=R_{\mathfrak a}
\]
となる。そこで、$\mathfrak a$ の各元 $\alpha=c_1\alpha_1+\cdots+c_s\alpha_s$ に $\mathfrak b$ の元 $\beta=c_1\beta_1+\cdots+c_s\beta_s$ を対応させる。この対応は一対一であり、和と差も互いに対応する。さらに $R_{\mathfrak a}=R_{\beta}$ から、$\gamma\alpha_i$ と $\gamma\beta_i$ も互いに対応することが従う。なぜなら、それらはそれぞれ $\alpha_i$ と $\beta_i$ の同じ線型形式によって表されるからである。したがって一般に $\gamma\alpha$ と $\gamma\beta$ が対応し、イデアル $\mathfrak a$ と $\mathfrak b$ は同じイデアル類に属する。

単位イデアルの類を \emph{主類} と呼ぶ。したがって、$\mathfrak R$ の任意の基底は、表現類として主類を与える。
%
%
%

\textbf{3. 類に関する元のトレース。}

$C$ を、$R_{\mathfrak a}$ において $\mathfrak R$ の元 $\gamma$ に対応する行列とする。$R_{\mathfrak a}$ から $R_{\bar\alpha}$ への移行により、行列式 $|C|$、さらに一般に、$t$ を不定元とするときの $|tE-C|$ は表現類の不変量であることが分かる。したがって 2 により、これは同時にイデアル類の不変量、すなわち、略していう \emph{類不変量} である。ゆえに $|tE-C|$ における $t$ の各係数は類不変量となる。特に、$(-t^{s-1})$ の係数を \emph{類に関する $\gamma$ のトレース} と呼ぶことにする。記号で書けば
\[
S_{(\mathfrak a)}(\gamma)=c_{11}+c_{22}+\cdots+c_{ss}
\]
である。$t$ を含まない係数 $\pm|C|$ を \emph{類に関する $\gamma$ のノルム} と呼ぶ。
%
%
%

\emph{固定した類 $(\mathfrak a)$ に関するトレース $S_{(\mathfrak a)}(\gamma)$ の全体は一つの $\Sigma$-加群をなし、$\Sigma$-加群とみなした $\mathfrak R$ はこの加群へ準同型に写る。}

実際、1 で証明した $\mathfrak R$ から $R_{\mathfrak a}$ への環準同型により
\[
(\beta-\gamma)\sim(B-C),
\qquad
c\beta\sim cEB=cB
\]
となる。したがってトレースについて
\[
S_{(\mathfrak a)}(\beta-\gamma)
=S_{(\mathfrak a)}(\beta)-S_{(\mathfrak a)}(\gamma),
\qquad
S_{(\mathfrak a)}(c\beta)=cS_{(\mathfrak a)}(\beta)
\]
が従い、加群性と準同型性が得られる。
%
%
%

\emph{注意：$\mathfrak R$ が零因子をもたない環ならば、元のトレースは類に依存せずに定まる。したがって単にトレースと呼ぶことができ、ノルムについても同様である。}

実際、零因子をもたない環の場合、零イデアルとは異なる各イデアルは環と同じ階数をもつ。したがって、イデアルの各基底は商体において同時に単位イデアルの基底となる。ゆえに、すべてのトレースとノルムは商体において同時に単位イデアルによって生成され、そのため $\mathfrak R$ において類に依存しない。
%
%
%

\textbf{4. 類に関するイデアルの判別式。}

イデアルと類の不変量として現れるのは、判別式そのものではなく、そこから $\Sigma$ において導かれる主イデアルである。
%
%
%

実際、$\varrho_1,\ldots,\varrho_r$ と $\bar\varrho_1,\ldots,\bar\varrho_r$ を、$\mathfrak R$ のイデアル $\mathfrak r$ の二組の（線型独立な）$\Sigma$-加群基底とする。\emph{行列式 $|S_{(\mathfrak a)}(\varrho_i\varrho_k)|$ と $|S_{(\mathfrak a)}(\bar\varrho_i\bar\varrho_k)|$ は、$\Sigma$ のある単元の平方を因子としてのみ異なる。}

実際、このことは二つの行列式 $|(\varrho_i\varrho_k)|$ と $|(\bar\varrho_i\bar\varrho_k)|$ について成り立つ。これらは、$\varrho$ と $\bar\varrho$ の間の変換行列式の平方を因子として異なるからである。3 で証明した加群準同型により、$S_{(\mathfrak a)}(\varrho_i\varrho_k)$ と $S_{(\mathfrak a)}(\bar\varrho_i\bar\varrho_k)$ の間には、$\varrho_i\varrho_k$ と $\bar\varrho_i\bar\varrho_k$ の間と同じ線型関係が成り立つ。したがって、合同（コグレディエント）変換により、行列式にも同じ因子が掛かる。
%
%
%

任意の行列式 $|S_{(\mathfrak a)}(\varrho_i\varrho_k)|$ を、$\Sigma$ の単元の平方倍を除いて定まる \emph{類 $(\mathfrak a)$ に関する $\mathfrak r$ の判別式} と呼ぶ。また、そこから $\Sigma$ において導かれる主イデアルを \emph{判別イデアル} と呼ぶ。
%
%
%

\emph{注意：$\mathfrak R$ が零因子をもたない環ならば、イデアルの判別式は、基礎として選んだ類に依存しない。}
%
%
%

というのも、3 により、この独立性は行列式をなす各要素について個別に成り立つからである。
%
%
%

\begin{center}
\textbf{\S 5. 類の直和。}
\end{center}
%
%
%

\textbf{1. イデアル類または表現類の直和。}

イデアル $\mathfrak a$ が直和 $\mathfrak a=\mathfrak b+\mathfrak c$ であるとする。同型により、この類に属する任意のイデアル $\overline{\mathfrak a}$ も直和となり、$\overline{\mathfrak a}=\overline{\mathfrak b}+\overline{\mathfrak c}$ と書ける。ここで、これらのイデアルを $\mathfrak R$-加群とみなすとき、$\overline{\mathfrak b}$ は $\mathfrak b$ に、$\overline{\mathfrak c}$ は $\mathfrak c$ にそれぞれ同型である。したがって、\emph{イデアル類の直和} およびその \emph{成分類} について語ることができる。
%
%
%

ある表現類に属する環 $R_{\mathfrak a}$ が、$R_{\mathfrak a}$ のイデアルでもある部分環の直和であるとする。同型により、このことはその表現類に属する任意の環 $R_{\bar\alpha}$ についても成り立ち、対応する成分はそれぞれ同値な環となる。したがって、\emph{表現類の直和} およびその \emph{成分類} について語ることができる。
%
%
%

\begingroup
\renewcommand{\thefootnote}{\arabic{footnote})}

\emph{イデアル類の直和には表現類の直和が対応する。この意味で、類の直和ということにする}\footnote[1]{表現類の直和がすべて、どの程度イデアル類の直和によって生成されうるかという問題には、ここでも立ち入らない。}。$\beta_1,\ldots,\beta_l$ を $\mathfrak b$ の基底、$\gamma_1,\ldots,\gamma_m$ を $\mathfrak c$ の基底とする。直和であることから、$\beta$ 族と $\gamma$ 族を合わせたものは $\mathfrak a$ の（線型独立な）基底 $a$ をなす。この基底によって生成される行列環を $R_{\mathfrak a}$ とすると、$\mathfrak R$ の任意の元 $\delta$ に対応する $R_{\mathfrak a}$ の行列は $(\begin{smallmatrix}D^{(\beta)}&0\\0&D^{(\gamma)}\end{smallmatrix})$ の形をとる。そこで、$R^{(\beta)}$ を $(\begin{smallmatrix}D^{(\beta)}&0\\0&0\end{smallmatrix})$ という形の行列すべてからなる系として定め、$R^{(\gamma)}$ も同様に定める。このとき示すべきことは、$R^{(\beta)}$ と $R^{(\gamma)}$ が $R_{\mathfrak a}$ の部分環をなし、$R_{\mathfrak a}$ がこれらの直和となること、また同時に、$R^{(\beta)}$ と $R^{(\gamma)}$ が、$\mathfrak b$ と $\mathfrak c$ から生成される環 $R_\beta$ と $R_\gamma$ にそれぞれ同型となることである。

\endgroup
%
%
%

実際、$R^{(\beta)}$ には準同型によって、$\delta\mathfrak c$ が零となる $\mathfrak R$ の元 $\delta$ 全体、すなわちイデアル商 $(0):\mathfrak c$ に属する元全体が対応する。したがって、この準同型により $R^{(\beta)}$ は $R_{\mathfrak a}$ のイデアルとなり、$R^{(\gamma)}$ についても同じことが成り立つ。$R_{\mathfrak a}$ はこれらのイデアルの和となる。しかも、二つの環の積 $R^{(\beta)}R^{(\gamma)}$ が零となるので、この和は直和である。最後に、行列 $D^{(\beta)}$ を、$D^{(\beta)}$ と零だけを含む $R^{(\beta)}$ の行列に対応させる写像は、$R^{(\beta)}$ と $R_\beta$ の間の同型を与える。この意味で、$\mathfrak a$ のイデアル類の成分類 $(\mathfrak b)$ および $(\mathfrak c)$ には、$R_{\mathfrak a}$ の成分類が対応する。
%
%
%

\textbf{2. 類の直和におけるトレースと判別式の挙動。} $\mathfrak a=\mathfrak b+\mathfrak c$ とし、トレース $S_{(\mathfrak a)}(\delta)$ を定める際に行列 $(\begin{smallmatrix}D^{(\beta)}&0\\0&D^{(\gamma)}\end{smallmatrix})$ を基礎に置くと、1. を考慮して次のことが読み取れる：
%
%
%

\emph{ある類に関して取った元のトレースは、その成分類に関して取ったトレースの和である。}
%
%
%

積 $\mathfrak b\mathfrak c$ が零となることから、さらに次のことが読み取れる：
%
%
%

\emph{$\beta$ が $\mathfrak b$ の元であるならば、$S_{(\mathfrak a)}(\beta)=S_{(\mathfrak b)}(\beta)$ となる。すなわち、類 $(\mathfrak a)$ に関して取った $\beta$ のトレースは、類 $(\mathfrak b)$ に関して取ったトレースに等しい。}
%
%
%

これより、\S 4, 4. を考慮すると、さらに次が従う：
%
%
%

\emph{類 $(\mathfrak a)$ に関して取った成分イデアル $\mathfrak b$ の判別イデアルは、成分類 $(\mathfrak b)$ に関して取った $\mathfrak b$ の判別イデアルに等しい。}
%
%
%

そして最後に、次が得られる：
%
%
%

\emph{類 $(\mathfrak a)$ に関して取った $\mathfrak a$ の判別イデアルは、成分類に関して取った各成分の判別イデアルの積に等しい。} 実際、判別式の構成にもトレースの構成にも、直和に対応する $\beta,\gamma$ の基底を採用すると、類 $(\mathfrak a)$ に関して取った $\mathfrak a$ の判別イデアルの一つの基底は、次で与えられる：
\[
\left|
\begin{array}{cc}
S_{(\mathfrak a)}(\beta_i\beta_k)&0\\
0&S_{(\mathfrak a)}(\gamma_\mu\gamma_\nu)
\end{array}
\right|
=
\left|
\begin{array}{cc}
S_{(\mathfrak b)}(\beta_i\beta_k)&0\\
0&S_{(\mathfrak c)}(\gamma_\mu\gamma_\nu)
\end{array}
\right|
\quad
\left(
\begin{matrix}
 i,k=1,\ldots,l,\\
\mu,\nu=1,\ldots,m
\end{matrix}
\right).
\]
%
%
%

\textbf{3. 拡大環への移行。} $\overline{\Sigma}$ を、$\Sigma$ の零因子をもたない拡大環であって、$\mathfrak R$ と $\overline{\Sigma}$ との共通部分が $\Sigma$ で与えられるものとする。また、$\overline{\Sigma}$ を添加して得られる拡大環 $\overline{\mathfrak R}=\mathfrak R[\overline{\Sigma}]$ は、\S 1, 2. と同様に、同じ階数をもつ環として定義されるものとする。
%
%
%

$\mathfrak a$ と $\mathfrak r$ を $\mathfrak R$ のイデアルとし、$\overline{\mathfrak a}$ と $\overline{\mathfrak r}$ を $\overline{\mathfrak R}$ におけるそれらの拡大イデアルとする。このとき、類 $(\overline{\mathfrak a})$ に関して取った $\overline{\mathfrak r}$ の判別イデアルもまた、類 $(\mathfrak a)$ に関して取った $\mathfrak r$ の判別イデアルを $\overline{\Sigma}$ に拡大したイデアルとなる。実際、$\mathfrak a$ または $\mathfrak r$ の任意の基底は、それぞれ $\overline{\mathfrak a}$ または $\overline{\mathfrak r}$ の $\overline{\Sigma}$-加群基底でもあるので、類 $(\overline{\mathfrak a})$ に関する $\overline{\mathfrak r}$ の判別式の構成には、$\mathfrak a$ と $\mathfrak r$ の基底を基礎として用いることができる。したがって特に、主類 $(\overline{\mathfrak R})$ に関して取った $\overline{\mathfrak R}$ の判別イデアルは、主類 $(\mathfrak R)$ に関して取った $\mathfrak R$ の判別イデアルの拡大イデアルとなる。このように定義された $\Sigma$ または $\overline{\Sigma}$ のイデアルを、簡単に \emph{環の判別イデアル} と呼ぶ。
%
%
%

\begin{center}
\textbf{\S 6. 第一種の完全可約性に対する判別式判定条件。}
\end{center}
%
%
%

今後、$\mathfrak R$ は再び体 $P$ の拡大環であると仮定し、$\overline P$ は $P$ の拡大体であるとする。第一種の完全可約性の必要十分条件は判別式が零でないことである、ということが以下で得られる。\S 5, 3. により、$\mathfrak R$ の判別イデアルは拡大環 $\overline{\mathfrak R}=\mathfrak R[\overline P]$ を用いて決定できる。さらに \S 5, 2. により、$\overline{\mathfrak R}$ を加法的に構成する準素環だけを考察すれば十分である。したがって、まず準素環の判別イデアルを考察する。
%
%
%

\textbf{1. 特殊な $P$-加群基底。} $\mathfrak a$ と $\mathfrak t$ が $\mathfrak R$ のイデアルであり、$\mathfrak a$ が $\mathfrak t$ で割り切られるなら、$\mathfrak a$ の任意の基底 $\alpha_1,\ldots,\alpha_s$ は、元 $\tau_{s+1},\ldots,\tau_t$ を添加することにより、$\mathfrak t$ の（線型独立な）加群基底へ補うことができる\footnote{この基底によって与えられる任意の元 $\gamma$ の行列表現は、そのとき $(\begin{smallmatrix}C_{\mathfrak a}&0\\ C_1&C_2\end{smallmatrix})$ の形となる。逆に、$\mathfrak t$ によって与えられるこのような行列表現から、$\mathfrak t$ の倍数 $\mathfrak a$ の存在が従う。完全可約な場合、直和分解不能な成分は零イデアルまたは単位イデアルと異なるイデアルをもたないので、これらの成分イデアルの表現類にはこのような表現は現れえない。さらに、\S 5, 1. で与えられた類の直和に関する関係を加味すれば、文献で通常用いられる完全可約性の定義に到達する。}。これは、体 $P$ に関して仮定した有限階数から直ちに従う。
%
%
%

ここで $\mathfrak R$ は\emph{準素環であると仮定する}。$\mathfrak p$ を零イデアルの随伴素イデアル、$\varrho$ を指数とする。すなわち、$\mathfrak p^{\varrho}=(0)$、$\mathfrak p^{\varrho-1}\ne(0)$ である。このとき、$\mathfrak p^{\varrho-1}$ の基底を $\mathfrak p^{\varrho-2}$ の基底へ、一般には $\mathfrak p^{\lambda}$ の基底を $\mathfrak p^{\lambda-1}$ の基底へと順次補っていくことにより、$\mathfrak R$ の特殊な $P$-加群基底が得られる。ただし、$\mathfrak R$ は $\mathfrak p^0$ と数える。このような基底を用いて、次を示すことができる：
%
%
%

\emph{$\mathfrak R$ が準素環ならば、随伴素イデアル $\mathfrak p$ で割り切れる任意の元について、主類（\S 4, 2.）に関して取ったトレースは零となる。} 実際、$\pi_{0,1},\ldots,\pi_{0,i_0};\ \pi_{1,1},\ldots,\pi_{1,i_1};\ \ldots;\ \pi_{\varrho-1,1},\ldots,\pi_{\varrho-1,i_{\varrho-1}}$ を、$\pi_{\lambda,j}\equiv0\;(\mathfrak p^{\lambda})$ かつ $\not\equiv0\;(\mathfrak p^{\lambda+1})$ を満たす、このような特殊な $P$-加群基底とする。$\pi\equiv0\;(\mathfrak p)$ から $\pi\pi_{\lambda,j}\equiv0\;(\mathfrak p^{\lambda+1})$ が従う。したがって、この基底によって $\pi$ に対応する行列は対角線上に零しかもたず、$\pi$ のトレースは零となる。
%
%
%

\textbf{2. 代数閉部分環 $\mathsf{P}$ をもつ準素環の判別イデアル。} $\mathfrak R$ の零イデアルが素イデアルであり、$P$ が代数閉ならば、$\mathfrak R$ は階数一となり、したがって $P$ と一致する（\S 3, 3.）。単位元 $e$ を加群基底に取ることができる。すると、$S_{(\mathfrak R)}(e^2)=S_{(\mathfrak R)}(e)=e$ であり、判別イデアルは単位イデアルに等しい。
%
%
%

さらに、代数閉部分環 $P$ をもつ真の準素環の判別イデアルは零イデアルに等しいことを示すことができる。$\mathfrak R$ について 1. で与えた特殊な加群基底を用いるか、あるいは単に $\mathfrak p$ の任意の基底を $\mathfrak R$ の基底へ補うと、剰余類環 $\mathfrak R\,|\,\mathfrak p$ が階数一の環となるため、$i_0=1$ となる。$\pi_{0,1}$ には、たとえば単位元 $e$ を取ることができる。すると、$e^2$ と異なるすべての基底積は $\mathfrak p$ で割り切られ、それらの主類に関して取ったトレースは零となる。したがって、判別式も少なくとも二行からなる行列式として零となる---$\mathfrak R$ は真の準素環と仮定されていたので、$e$ と異なる基底元を少なくとも一つもたなければならない---。この行列式では、$S_{(\mathfrak R)}(e^2)$ と異なるすべての行列要素が零となる。
%
%
%

\textbf{3. 定理。} \emph{部分体 $P$ に関して有限階数の環 $\mathfrak R$ が第一種の完全可約であるための必要十分条件は、その判別イデアルが $P$ における単位イデアルであることである---言い換えれば、単元の平方倍を除いて定まる判別式が零でないことである。}
%
%
%

第一種の完全可約性は、\S 3, 2. の定理により、拡大環 $\overline{\mathfrak R}=\mathfrak R[\overline P]$ が完全可約であること、すなわち、その直和分解不能な成分が $\overline P$ と同型な体となることと同じである。\S 5, 3. により、$\mathfrak R$ と $\overline{\mathfrak R}$ の判別イデアルは、同時に零イデアルとなるか、または同時に単位イデアルとなる。また、\S 5, 2. により、$\overline{\mathfrak R}$ の判別イデアルは、各成分自身に関して取った各成分の判別イデアルの積となる。ただし、これら個々の判別イデアルは、2. で考察したもの---そこでは各成分は、それ自体で考察される環として現れ、$\overline{\mathfrak R}$ のイデアルとしては現れない---から、それぞれ $\overline P$ と同型な体 $\overline P_i=\overline P\bar\varepsilon_i$ を $\overline P$ 自身で置き換えることによって得られる。したがって、$\overline{\mathfrak R}$ が完全可約であるとき、その判別イデアルは単位イデアルの積としてそれ自身単位イデアルとなる。これに対して、$\overline{\mathfrak R}$ が完全可約でないとき、少なくとも一つの成分について判別イデアルが零イデアルに等しくなるので、$\overline{\mathfrak R}$ の判別イデアルも零イデアルとなる。これで定理が証明された。
%
%
%

\textbf{4. 第一種の完全可約の場合における、共役元によるトレース、ノルムおよび判別式の表示。} \S 3, 3. により、第一種の完全可約の場合、$\mathfrak R[\Omega]=\Omega\bar\varepsilon_1+\cdots+\Omega\bar\varepsilon_n$ を階数一の環---すなわち $\Omega$ と同型な体---の直和表示とし、$\Omega$ をそのような表示が可能となる最小の拡大体とする。このとき、$\mathfrak R$ の $n$ 個の成分は $K_i\bar\varepsilon_i$ で与えられる。ここで各 $K_i$ は $\Omega$ の部分体であり、$\mathfrak R$ はそれぞれ $K_i$ に準同型に写る。以下、主類に関して取ったトレース、ノルムおよび判別式を、$K_i$ の元によって表示する。
%
%
%

$\bar\varepsilon_1,\ldots,\bar\varepsilon_n$ を加群基底として採ると、各元 $\gamma=c_1\bar\varepsilon_1+\cdots+c_n\bar\varepsilon_n$ には行列 $\left(\begin{smallmatrix}c_1&0\\\vdots&\vdots\\0&c_n\end{smallmatrix}\right)$ が対応する。このとき、$\gamma$ が $\mathfrak R$ の元ならば、$c_i$ は $K_i$ の元である。したがって、$S(\gamma)=c_1+\cdots+c_n,\ N(\gamma)=c_1\cdots c_n$ となる。ここで $S(\gamma)$ と $N(\gamma)$ は、それぞれ主類に関して取ったトレースとノルムを表す。
%
%
%

さらに、$\alpha_1,\ldots,\alpha_n$ が $\mathfrak R$ の元からなる $\mathfrak R$ の加群基底を表すなら、判別式 $|S(\alpha_i\alpha_k)|$ はすべての $K_\lambda$ の元によって表示される。$\alpha_i=\alpha_i^{(1)}\bar\varepsilon_1+\cdots+\alpha_i^{(n)}\bar\varepsilon_n$ とする。このとき、$\alpha_i\alpha_k=\alpha_i^{(1)}\alpha_k^{(1)}\bar\varepsilon_1+\cdots+\alpha_i^{(n)}\alpha_k^{(n)}\bar\varepsilon_n$ であり、$S(\alpha_i\alpha_k)=\alpha_i^{(1)}\alpha_k^{(1)}+\cdots+\alpha_i^{(n)}\alpha_k^{(n)}$ であるから、
\[
\left|S(\alpha_i\alpha_k)\right|
=
\left|\sum_{\lambda}\alpha_i^{(\lambda)}\alpha_k^{(\lambda)}\right|
=
\left|
\begin{matrix}
 \alpha_1^{(1)}&\cdots&\alpha_n^{(1)}\\
 \vdots&&\vdots\\
 \alpha_n^{(1)}&\cdots&\alpha_n^{(n)}
\end{matrix}
\right|^2.
\]
%
%
%

特に、$\mathfrak R$ 自身が体、すなわち部分体 $P$ の第一種拡大体であるなら、$\mathfrak R$ から $K_i$ への準同型は同型写像となる。このとき $\mathsf{P}$ の元は自分自身に対応する---なぜなら $\mathsf{P}$ の成分は $\mathsf{P}\varepsilon_i$ で与えられるからである。したがって $K_i$ は、すべて共通の包含体 $\Omega$ の中にあるので、$\mathsf{P}$ に関して互いに共役な体となる。しかも、上の零でない判別式 $|S(a_i a_k)|$ の行列式表示が示すように、$n$ 個の同型写像はすべて互いに異なる。
%
%
%

\emph{$\mathfrak R$ が体であるなら、$K_1,\ldots,K_n$ による直和表示は、ガロア拡大体の中にあり、$\mathfrak R$ と同型である、$P$ 上 $n$ 次の $n$ 個の共役拡大体を与える。トレース、ノルムおよび判別式の表示は、共役元による周知の表示に帰着する。} ここで改めて強調しておくが、$\mathfrak R$ は $K_i$ と同値な体であると仮定されているだけであって、ガロア拡大体 $\Omega$ の中にある体であるとは仮定されていない\footnote{これについては \S 1, 2. の最後の注を参照。}。
%
%
%

\begin{center}
\textbf{\S 7. 代数的数体または関数体のオーダーに対する判別式}\\
\textbf{定理。}
\end{center}
%
%
%

\textbf{1.} ここで考察する数体および関数体はすべて、整元の指定を伴う次の体型に従う---ただし、二つ以上の不定元をもつ代数関数の場合、および整代数関数の場合には、関数領域を拡大することになる\footnote[2]{\emph{Idealtheorie}, \S 3, 3. および 4.}：
%
%
%

$\mathfrak H$ を、零因子をもたず単位元をもつ\emph{主イデアル環}（有理整数の環、ある体を係数とする一つの不定元の多項式領域、汎関数領域）とし、$\mathfrak R'$ をその\emph{商体}とする。$\mathfrak L$ を $\mathfrak R'$ の\emph{第一種有限拡大}とし、$\mathfrak S$ を $\mathfrak L$ に属する\emph{$\mathfrak H$ に関して整なすべての元からなる環}\footnote[3]{環に関して整な元の定義については、\emph{Idealtheorie}, \S 1, 3. を参照。}とする。\emph{$\mathfrak S$ の部分環で $\mathfrak H$ を含むものをすべてオーダーと呼び、$\mathfrak S$ 自身を主オーダーと呼ぶ。}
%
%
%

$\mathfrak S$ 内の各 $\mathfrak H$-加群、特にオーダーのすべてのイデアルおよびオーダー自身は、$\mathfrak H$ に関する\emph{線型独立な加群基底}をもつ\footnote[4]{\emph{Idealtheorie}, \S 3, 1.}。$\mathfrak L$ が $\mathfrak R'$ に関して次数 $n$ であるならば、$\mathfrak{K}$ に関して階数 $n$ のオーダーに考察を限れば十分である。というのも、それより低い階数のオーダーはいずれも、商を作ることによって、$\mathfrak R'$ と $\mathfrak L$ の中間体で、まさにその階数と同じ次数をもつものを生成し、その中間体についてもすべての仮定が保たれるからである。
%
%
%

\emph{したがって、任意のオーダー $\mathfrak T$ は \S 4 および \S 5 のより一般的な仮定が当てはまる環であり、しかも零因子をもたない環である。それゆえ、オーダーの判別式 $D_{\mathfrak T}$ は、$\mathfrak H$ の単元の平方を因子として掛ける違いを除いて、一意に定義される。} この判別式はまた、$\mathfrak R'$ の単元を因子として掛ける違いを除いて、$\mathfrak R'$-加群とみなした体 $\mathfrak L$ の判別式と一致し、したがって\emph{零でない}。なぜなら、$\mathfrak L$ は第一種の拡大と仮定されているからである（\S 6, 3.）。同時に、この判別式は、\S 6, 4. で与えた共役基底元の行列式の平方による表示をもつ。
%
%
%

$\mathfrak T$ におけるイデアル論は、次の定理によって与えられる\footnote[5]{\emph{Idealtheorie}, \S 7, 3.}：\emph{$\mathfrak T$ において、零イデアルと異なる任意のイデアルは、有限個の、どの二つも互いに素な準素イデアルの積として一意に表される。しかも、それらの随伴素イデアルは、単位イデアルと異なる真の約イデアルをもたない。}
%
%
%

\textbf{2. オーダーから体に関して有限階数の環への移行。} $p$ を $\mathfrak H$ の素元とする。すなわち、$\mathfrak H$ において $p$ から生成される主イデアルは、$\mathfrak H$ の零イデアルおよび単位イデアルと異なる素イデアル $\mathfrak H p$ である。これに応じて、$\mathfrak T$ において $p$ から生成される主イデアルを $\mathfrak T p$ と記す\footnote[1]{今後、主イデアルにはこの表記を用い、それが考察される環を同時に示すことにする。}。\emph{剰余類環 $\mathfrak R=\mathfrak T/\mathfrak T p$ は、体 $\mathfrak H/\mathfrak H p$ と同型な部分体 $P$ に関して階数 $n$ の環となり、$\mathfrak T$ は $\mathfrak T/\mathfrak T p$ に準同型に写される。}
%
%
%

準同型 $\mathfrak T\sim\mathfrak T/\mathfrak T p$ は、任意の剰余類環について一般に成り立つ。$\mathfrak R$ が体 $\mathfrak H/\mathfrak H p$ と同型な部分体 $P$ を含むことは、第二同型定理\footnote[2]{\S 1 の最後の注を参照。}から従う。実際、部分環 $(\mathfrak H,\mathfrak T p)/\mathfrak T p$ は $\mathfrak H/[\mathfrak H,\mathfrak T p]=\mathfrak H/\mathfrak H p$ と同型となる。
%
%
%

最後に、$\mathfrak R$ は $P$ に関しても階数 $n$ の環となる。実際、たとえば $\alpha_1,\ldots,\alpha_n$ を $\mathfrak H$ に関する $\mathfrak T$ の線型独立な加群基底とし、$(\alpha_1),\ldots,(\alpha_n)$ を各元 $a_i$ が $\mathfrak T p$ を法として定める剰余類とする。このとき、$(a_i)$ の間の $P$ に関する任意の線型従属関係は、準同型によって合同式 $c_1a_1+\cdots+c_na_n\equiv0\;(\mathfrak T p)$ から生ずる。すなわち、$\gamma$ を $\mathfrak T$ の元として、関係式 $c_1\alpha_1+\cdots+c_n\alpha_n=p\gamma$ から生ずる。しかし、$\gamma$ が基底 $\alpha_i$ によって表示できることと、その表示の一意性から、各 $c_i$ は $p$ で割り切れる。したがって、$P$ の元 $(c_i)$ は零となる。これにより、$(a_i)$ が $P$ に関して線型独立であることが示された。
%
%
%

\emph{準同型によって、$\mathfrak T$ の判別式 $D_\mathfrak T$ は $\mathfrak R$ の判別式へ移り、$\mathfrak T p$ のイデアル分解は $\mathfrak R$ の零イデアルの分解へ移る。} 実際、先に証明した $(a_i)$ の線型独立性により、$\mathfrak R$ の判別式は、この基底を用いて $|S((\alpha_i)(\alpha_k))|$ として作ることができる。したがってそれは、準同型によって、基底 $a_i$ を用いて作られる $\mathfrak T$ の判別式 $|S(\alpha_i\alpha_k)|$ から生ずる。ここでトレースの定義の基礎に置かれているのは行列表現であって、$\mathfrak T$ には妥当するが $\mathfrak R$ には必ずしも妥当しない共役元による表示ではない。ゆえにトレース、したがって判別式は、環の元から環演算によって生ずる。
%
%
%

$\mathfrak T p=\mathfrak q_1\cdots\mathfrak q_r=[\mathfrak q_1,\ldots,\mathfrak q_r]$ が $\mathfrak T$ における $\mathfrak T p$ の表示であるとき、この表示は準同型によって $\mathfrak R$ における $(0)=[(\mathfrak q_1),\ldots,(\mathfrak q_r)]$ へ移る。このとき、第一同型定理\footnote[3]{\emph{イデアル論}, \S 4, 3 および \S 6（定理 II の後）を参照。}により、$\mathfrak R/(\mathfrak q_i)$ は $\mathfrak T/\mathfrak q_i$ と同型となる。したがって、$\mathfrak q_i$ と $(\mathfrak q_i)$ は同時に真の準素イデアルまたは素イデアルとなり、また準同型により $(\mathfrak q_i)$ は互いに素である。
%
%
%

\textbf{3. 判別式定理。} \emph{$\mathfrak H$ の素元 $p$ がオーダー $\mathfrak T$ の判別式 $D_\mathfrak T$ を割るのは、$\mathfrak T$ における主イデアル $\mathfrak T p$ の互いに素な準素成分への分解の中に、少なくとも一つの真の準素成分または少なくとも一つの第二種の素イデアルが現れるとき、かつそのときに限る\footnote[4]{第二種の素イデアルが現れる例は次のとおりである。$\mathfrak H$ を、有理整数係数をもつ $x$ のすべての多項式から導かれる関数領域とし、$\mathfrak R'$ をその商体、$\mathfrak L$ を $\sqrt{x}$ の添加によって生ずる拡大体とする。$\mathfrak H$ に関して加群基底 $1,\sqrt{x}$ をもつオーダー $\mathfrak T$ を考える（置換 $x=y^2$ が示すように、これは主オーダーである）。判別式は $\left|\begin{smallmatrix}1&\sqrt{x}\\1&-\sqrt{x}\end{smallmatrix}\right|^2=2^2x$ となる。主イデアル $\mathfrak T2$ は $\mathfrak T$ において素イデアルのままであるが、第二種である。実際、オーダー $\mathfrak T$ は、加群基底 $1,\sqrt{x}$ をもつため、剰余類環 $\mathfrak H[t]/(t^2-x)$ と同型である。ここで $t$ は新しい不定元を表す。したがって $\mathfrak T/\mathfrak T2=\mathfrak R$ は $\mathsf{P}[t]\,|\,t^2-x$ と同型となる。ここで $P$ は、法 $2$ の係数をもつ、$x$ および関数領域に現れる不定元 $u$ のすべての有理関数からなる体を表す。$t^2-x$ は $\mathsf{P}[t]$ の素関数、それも第二種の素関数であるから、$\mathfrak R$ は体、しかも $P$ の第二種の拡大体となる。したがって $\mathfrak T2$ は第二種の素イデアルとなる。他方、$\mathfrak T x=(\mathfrak T\sqrt{x})^2$ は真の準素イデアルとなる。$\mathfrak R$ が $\mathsf{P}$ に一つの元を添加して生成されることは、一般には成り立たない。一つの不定元 $x$ を $x$ と $y$ に置き換えれば、$1,\sqrt{x},\sqrt{y},\sqrt{xy}$ から導かれるオーダー---これも再び主オーダーとなる---において、主イデアル $\mathfrak T2$ が第二種の素イデアルであり、それに対する $\mathfrak R$ は $P$ に二つの元を添加して初めて生成されることが分かる（これらの例については Ostrowski, 前掲書を参照）。}。}
%
%
%

\emph{剰余体 $\mathfrak H/\mathfrak H p$ が完全体であるならば、$D_\mathfrak T$ を割るすべての素元 $p$（かつ、そのような素元だけ）が、少なくとも一つの真の準素成分をもつ。$\mathfrak H/\mathfrak H p$ が完全で、$\mathfrak T$ として主オーダー $\mathfrak S$ を選ぶならば、$p$ が判別式を割ることと、$p$ が $\mathfrak S$ のある素イデアルの平方で少なくとも割り切られることとは同値である。}
%
%
%

実際、$\mathfrak T$ における主イデアル $\mathfrak T p$ のイデアル分解と、$\mathfrak R=\mathfrak T/\mathfrak T p$ における零イデアルの分解との対応（2 を参照）により、$\mathfrak T p$ に真の準素成分または第二種の素イデアルが現れることは、$\mathfrak R$ が第一種の完全可約では\emph{ない}ことを意味する。しかし、\S 6, 3 の定理によれば、これは $\mathfrak R$ の判別式が零となるとき、かつそのときに限って成り立つ。そして 2 の準同型により、それは $D_\mathfrak T$ が $p$ で割り切られることを表す。
%
%
%

\begin{center}
\textbf{\S 8. 相対体のオーダーに対する判別式定理。}
\end{center}
%
%
%

主イデアル環 $\mathfrak H$ の代わりに、すでに数体または関数体の主オーダーを出発環として基礎に置くならば---より一般には、\emph{乗法環}\footnote[1]{この簡略な呼称「乗法環」は、W. Krull の小論 \,\emph{Über Multiplikationsringe}, \emph{Heidelberger Berichte}, 1925, 第5論文, 第3号から採った。ただし Krull は、ここに現れる零因子をもたない環を「正則乗法環」と呼んでいる。}、すなわち主オーダーのイデアル論が成り立つ環を基礎に置くならば---、判別式 $D_\mathfrak T$ の代わりに、$\mathfrak H$ に関する $\mathfrak T$ の判別イデアル\footnote[2]{文献では「相対数体」の場合、主オーダーの判別イデアル、すなわち「相対判別式」だけが現れる。Hilbert, \emph{Zahlbericht}, \S 14 および Hecke, \emph{Theorie der algebraischen Zahlen}, \S 38 を参照。}が現れる。以下の考察によって、判別式定理は完全に移される。
%
%
%

\textbf{1.} $\mathfrak H$ を\emph{乗法環}とする。すなわち、これは零因子をもたず単位元をもつ環であり、その中では零イデアルおよび単位イデアルと異なる任意のイデアルが素イデアルの冪積として一意に表され、しかもこれらの素イデアルは単位イデアルと異なる真の約イデアルをもたない\footnote[3]{乗法環の公理系については、\emph{Idealtheorie}, たとえば序論を参照。}。さらに $\mathfrak R',\mathfrak L,\mathfrak S,\mathfrak T$ を \S 7, 1 におけるとおりに定義する。このとき $\mathfrak T$ では、\S 7, 1 に記したイデアル論が成り立つ\footnote[4]{\emph{Idealtheorie}, \S 3, 2 および \S 7, 3。}。
%
%
%

特に、乗法環 $\mathfrak H$ が零イデアルおよび単位イデアルと異なる素イデアルをただ一つしかもたないならば、$\mathfrak H$ は主イデアル環となる。実際、$\mathfrak p$ をこの素イデアルとし、$p\equiv0\;(\mathfrak p)$ かつ $p\not\equiv0\;(\mathfrak p^2)$ とする。仮定により、零イデアルおよび単位イデアルと異なる任意のイデアルは $\mathfrak p$ の冪となるから、特に主イデアル $\mathfrak H p$ もそうである。したがって $p\not\equiv0\;(\mathfrak p^2)$ により、必然的に $\mathfrak H p=\mathfrak p$ となり、これによって $\mathfrak p$ の任意の冪も主イデアルとなる。零イデアルおよび単位イデアルも同様に主イデアルである。
%
%
%

\textbf{2. トレースと判別イデアル。} $\mathfrak T$ の各元 $\gamma$ のトレースは、\emph{$\mathfrak L$ を $\mathfrak R'$-加群とみなしたときの、$\mathfrak L$ における $\gamma$ のトレースとして定義される}。したがって、$\mathfrak L$ の元からなる、$\mathfrak R'$ に関して線型独立な任意の $n$ 元系を、トレースを定義する行列表現を作るための基底として選ぶことができる。$\alpha_1,\ldots,\alpha_n$ をそのような系とすると、$\mathfrak L$ は第一種の拡大であると仮定されているから、$|S(\alpha_i\alpha_k)|$ は $\mathfrak R'$ の非零元となり、しかも共役基底元の行列式の平方となる（\S 6, 4）。$\mathfrak H$ が $\mathfrak R'$ において整閉であること\footnote[1]{\emph{Idealtheorie}, \S 9, 7。}から、すべての $a_i$ が $\mathfrak S$ に属するならば、$|S(\alpha_i\alpha_k)|$ は $\mathfrak H$ の元となることが従う。
%
%
%

\emph{判別イデアル：} $\tau_1,\ldots,\tau_n$ に、$\mathfrak T$ の元をそれぞれ $n$ 個とるすべての系を走らせる。各系について行列式 $|S(\tau_i\tau_k)|$ を作ると、これは $\mathfrak H$ の元となり、しかも $\tau$ が線型独立ならば非零となる。\emph{これらの行列式全体によって $\mathfrak H$ 内で生成されるイデアルを、$\mathfrak H$ に関するオーダー $\mathfrak T$ の判別イデアルと呼ぶ。}したがって、このイデアルは零イデアルとは異なる。
%

さらに、$m\ge n$ として、$\beta_1,\ldots,\beta_m$ が $\mathfrak T$ の（常に存在する）\footnote[2]{\emph{Idealtheorie}, \S 9 および \S 3, 1。}$\mathfrak H$-加群基底を表すならば、各 $n$ 個の $\beta$ に対応する有限個の行列式が、判別イデアルのイデアル基底を与える。これは、$\mathfrak T$ から $\mathfrak T$ のトレース系への加群準同型性（\S 4, 3）から従う。実際、積 $\tau_i\tau_k$ から作られる行列式 $|(\tau_i\tau_k)|$ は、各 $n$ 個の $\beta$ に対応する有限個の行列式 $|(\beta_i\beta_k)|$ によって、$\mathfrak H$ の係数を用いて線型に表される。したがって、$|S(\tau_i\tau_k)|$ についても、$|S(\beta_i\beta_k)|$ に関して同じことが成り立つ。特に、オーダーが線型独立な加群基底 $\alpha_1,\ldots,\alpha_n$ をもつならば、行列式 $|S(\alpha_i\alpha_k)|$ を判別イデアルの基底として取ることができる。この定義は \S 7, 1 の定義と一致する。
%

\textbf{3. 商環への移行\footnote[3]{この番号で引用するすべての命題の証明については、原著83頁の注6で引用した H. Grell の論文を参照。}。} $\mathfrak a$ を $\mathfrak T$ の任意のイデアルとする。このとき商環 $\mathfrak T_\mathfrak a$ は、$\mathfrak T$ の商体 $\mathfrak L$ の元のうち、分母に $\mathfrak a$ と素な $\mathfrak T$ の元をもつもの全体として定義される\footnote[4]{$\mathfrak T$ の素イデアルは単位イデアルと異なる真の約イデアルをもたないので、ここでは「素」と「互いに素」とが---イデアルの意味で---同義となる。}。環 $\mathfrak T_\mathfrak a$ は、零イデアルおよび単位イデアルのほかには、$\mathfrak a$ に属する素イデアルの拡大イデアル以外の素イデアルをもたない。$\mathfrak T_\mathfrak a$ の零・単位イデアル以外のすべてのイデアルと、$\mathfrak T$ のイデアルのうち、それに属する素イデアルが $\mathfrak a$ に属する素イデアルの中に含まれるものとの間には、一対一の対応があり、対応する剰余類環は同型である。さらに、零イデアルと単位イデアルもそれぞれ一対一に対応する。
%

同じ考察は乗法環 $\mathfrak H$ にも当てはまる。特に $\mathfrak p$ を素イデアルとすると、$\mathfrak H_\mathfrak p$ は零イデアルおよび単位イデアルと異なる素イデアルをただ一つしかもたない。剰余類環の同型により、$\mathfrak H$ とともに $\mathfrak H_\mathfrak p$ も乗法環であり、したがって 1 により $\mathfrak H_\mathfrak p$ は主イデアル環となる。$\mathfrak p$ から導かれる素イデアルの基底は、$\mathfrak H_\mathfrak p$ の素元 $p$ となる。$\mathfrak H$ の任意のイデアル $\mathfrak c$ の拡大イデアルは、$\mathfrak p$ に属する $\mathfrak c$ の準素成分によって生成される。したがって、$\mathfrak p$ が $\mathfrak c$ を $\rho$ 乗まで割るか、またはまったく割らないかに応じて、この拡大イデアルは $(p)^\rho$ または単位イデアルに等しくなる。さらに次が従う。$\mathfrak H$ の二つのイデアル $\mathfrak b$ と $\mathfrak c$ の拡大が\emph{すべての}商環 $\mathfrak H_\mathfrak p$ で一致するならば、$\mathfrak b$ と $\mathfrak c$ は同一である。実際、両者は同じ素イデアル $\mathfrak p$ によって同じ冪まで割り切られる。
%

\textbf{4. 商環への移行における判別イデアル。} $\mathfrak p$ を $\mathfrak H$ の素イデアルとし、$\mathfrak P=\mathfrak T\mathfrak p$ で $\mathfrak T$ における拡大イデアルを表す。\emph{商環 $\mathfrak T_\mathfrak P$ は、主イデアル環 $\mathfrak H_\mathfrak p$ に関して線型独立な元からなる加群基底をもつ。同時に、$\mathfrak T_\mathfrak P$ は、$\mathfrak H_\mathfrak p$ の商体 $\mathfrak R'$ の拡大体 $\mathfrak L$ における有限 $\mathfrak H_\mathfrak p$-オーダーである\footnote[1]{証明については H. Grell を参照。}。したがって、$\mathfrak T_\mathfrak P$ では $\mathfrak H_\mathfrak p$ に関して判別式定理（\S 7, 3）が成り立つ。}
%

このとき、判別式 $D_{\mathfrak T_\mathfrak P}$ は、$\mathfrak H$ に関する $\mathfrak T$ の判別イデアルを $\mathfrak H_\mathfrak p$ に拡大して得られるイデアルの基底となる。実際、常に可能であるように、$\mathfrak T$ の元 $\alpha_1,\ldots,\alpha_n$ を $\mathfrak T_\mathfrak P$ の $\mathfrak H_\mathfrak p$-加群基底として取る。このとき $|S(\alpha_i \alpha_k)|$ は判別イデアルに属し、1 により、$\mathfrak H_\mathfrak p$ における任意の行列式 $|S(\tau_i\tau_k)|$ は $|S(\alpha_i\alpha_k)|$ で割り切られる\footnote[2]{ここから、$\mathfrak H$ が代数的数体の主オーダーである場合、Hilbert が定義した相対判別式との一致が得られる（これは Hecke の定義したものと一致することが知られている）。実際、$\omega_1,\ldots,\omega_m$ を相対体の主オーダーの $\mathfrak H$-加群基底とする---Hilbert は特に、有理整数全体の環に関する加群基底を選ぶ---。Hilbert に従い、各 $n$ 個の基底元とその共役元とからなるすべての行列式を作る。「相対判別式」は、これらの行列式を二つずつ（同一でも異なってもよい）掛けた積の全体から $\mathfrak H$ 内で生成されるイデアルである。したがって、相対判別式を $\mathfrak H_\mathfrak p$ に拡大すると、基底 $\alpha_1,\ldots,\alpha_n$ とその共役元との行列式の平方、すなわち $|S(\alpha_i\alpha_k)|$ となる。ゆえに、相対判別式と（相対主オーダーの）判別イデアルとの拡大は、すべての商環 $\mathfrak H_\mathfrak p$ で一致する。3 の末尾により、この二つのイデアルは同一となる。}。
%

\textbf{5. 一般判別式定理。} \emph{乗法環 $\mathfrak H$ の素イデアル $\mathfrak p$ が、$\mathfrak H$ に関するあるオーダーの判別イデアルを割るのは、$\mathfrak T\mathfrak p$ を $\mathfrak T$ の互いに素な準素成分へイデアル分解したとき、その中に少なくとも一つの真の準素成分または少なくとも一つの第二種の素イデアルが現れるとき、かつそのときに限る。}
%

3 と 4 から、$\mathfrak H$ の素イデアル $\mathfrak p$ が $\mathfrak T$ の判別イデアルを割るのは、$\mathfrak H_\mathfrak p$ の素元 $p$ が判別式 $D_{\mathfrak T_\mathfrak P}$ を割るとき、かつそのときに限る。これは、剰余類環 $\mathfrak T_\mathfrak P\,|\,\mathfrak T_\mathfrak P p$ が第一種の完全可約ではないことと同じである。3 により、この剰余類環は $\mathfrak T\,|\,\mathfrak T\mathfrak p$ と同型であるから、これによって一般判別式定理が証明された。
%

数体の相対判別式に特殊化すると、次を得る：

\emph{数体の主オーダーの素イデアル $\mathfrak p$ が、ある拡大体の相対判別式を割るのは、$\mathfrak p$ がその拡大体の主オーダーにおいて、ある素イデアルの平方で割り切られるとき、かつそのときに限る。}
%

Göttingen, 1926年3月30日。
%

\begin{center}
\rule{9em}{0.4pt}
\end{center}

\section*{32. R. Brauer と共同: 既約表現の極小分解体について}
\emph{Sitz. Ber. d. Preuß. Akad. d. Wiss. 1927, pp. 221--228.}

\begin{center}\emph{（Schur 氏により提出。）}\end{center}

基礎体 $\mathrm{P}$ 上で既約である有限群の表現が絶対既約成分に分解するような、$\mathrm{P}$ 上最小次数の数体の問題は、最初に I. Schur によって扱われた。\footnote{Arithmetische Untersuchungen über endliche Gruppen, Sitzungsber. d. Berl. Ak. d. Wiss. 1906, p. 164. Beiträge zur Theorie der Gruppen linearer Substitutionen, Transact. of the Am. Math. Soc. Ser. 2, Vol. XV, 1909, p. 159. 第二の論文では、結果が完全可約な超複素系へ移されている。一つの表現は、そのすべての変換された（相似な）表現とともに一つの類にまとめられる。} たとえば $\mathrm{P}$ がそのような成分の指標を含むと仮定する。このとき Schur は、これらの体の $\mathrm{P}$ に関する次数、すなわち指数が、同じ類に属する絶対既約成分の個数と一致すること、したがって同じ体は一つの絶対既約成分を切り離すという要求だけで既に定まること、さらに、そのような分解が起こる任意の体の $\mathrm{P}$ に関する次数は指数の倍数になることを示した。これらの体はすべて、$\mathrm{P}$ が指標を含まない場合にも、分解体と呼ぶことにする。Schur は一つの例によって、最小次数の分解体が同型である必要はないことを示した。$\mathrm{P}$ が対応する単純指標を含まない場合、分解体はそのような指標の体と、$\mathrm{P}$ に関するその共役体とを含まなければならない。異なる共役指標に属する絶対既約表現はたしかに共役であるが、明らかに異なる類に属する。一つの類の絶対既約表現の系を切り離すためには、この場合にも、対応する指標の体と、対応する分解体の一つで十分である。

以下では分解体の問題をさらに追究する。最小次数の分解体は、対応する非可換体（斜体）によって特徴づけることができる。一般の分解体は、これらの非可換体と不変的に結びついた両側単純環によって特徴づけられる。さらに、四元数体の例によって、極小分解体の次数は有界でないことを示す。ここで分解体を極小と呼ぶのは、そのいかなる真の部分体も分解体でない場合である。この非有界性は、本質的には、直後の H. Hasse の覚書で証明される数論的存在定理から従う。しかしこの例は、計算によって --- むしろ検証という形で --- 初等的にも扱われる。したがって、一般論や数論的存在定理を用いずに、より弱いが十分な存在定理を初等的に証明することで、次数の非有界性が示される。\footnote{この例は E. Noether によるものである。それは R. Brauer による例の修正に基づく。Brauer の例では、次数が指数を超える極小分解体の存在が一般に示されていた。この例の初等的取り扱いは R. Brauer による。分解体の特徴づけは、著者たちの研究に遡るが、その研究はそれぞれ別の目標をもつ。E. Noether にとっては、一般の有限性仮定のもとで、加群論とイデアル論を基礎として表現論を構成することが問題である。R. Brauer にとっては、因子系を用いて、与えられた可換な完全基礎体上の有限次数の非可換体を構成することが問題である。著者たちは、最小次数の分解体の特徴づけに独立に到達した。R. Brauer は完全基礎体について、E. Noether はこの制限なしに到達した。R. Brauer は、非可換埋め込みの可能性に関する E. Noether の問いに関連して、一般の分解体の特徴づけを見いだした。E. Noether はここでも、完全な基礎領域という制限を取り除くことができた。} また、この例で再び有限群、すなわち Schur の例の基礎にもなっている四元数群の分解体を扱っていることにも注意する。そこで与えられる表現は、同時に四元数体の（同型な）表現でもある。

\subsection*{1. 分解体の特徴づけ}

まず基本概念を簡単にまとめる。$\mS$ を体 $\mathrm{P}_0$ に関する超複素系とする。$\mS$ の表現 $\Gamma$ --- $\mathrm{P}$ における表現 --- とは、$\Gamma$ が $\mS$ の準同型像となり、かつ $P_0$ の元 $p_0$ に対して対角行列 $p_0E$ が対応するような、$P$ の元を成分にもつ行列系をいう。したがって $\mathrm{P}$ は $\mathrm{P}_0$ を含まなければならない。表現 $\Gamma$ は、そのすべての変換 $\mathit{P}^{-1}\Gamma \mathit{P}$ とともに一つの表現類をなす。このように定義された表現には、有限群 $\mathfrak G$ の表現も含まれる。対応する超複素系 $\mS$、すなわち「群環」は、すべての「群数」、つまり $\mathrm{P}_0$ の係数をもつ群元の線形結合全体から成る。ただしもとの群元を線形独立とみなし、群の乗法と算術法則によって乗法を定める。「可約」「既約」「完全可約」「絶対既約」な表現という概念は通常どおり定義される。この用語は表現類にも拡張する。系 $\mS$ の $\mathrm{P}$ におけるすべての表現類が完全可約であるとき、$\mS$ は「$P$ において完全可約」であるという。

いま $\mS$ が $P_0$ において完全可約であると仮定し、さらに $\mathrm{P}_0$ が完全体であると仮定する。すると $\mathrm{P}_0$ の任意の代数拡大 $\mathrm{P}$ についても同じことが成り立つ。表現類は $\mathrm{P}_0$ の任意の拡大 $\mathrm{P}$ のもとで完全可約のままである。\footnote{$\mathrm{P}_0$ が完全でない場合、完全可約性が保たれるのは、中心の成分 $\mathfrak{K}$ が $\mathrm{P}_0$ の「第一種の拡大」になるとき、かつそのときに限る。この仮定のもとでは、本文の以後の帰結はすべて有効である。} $\mS$ の中心は有限個の体の直和になる。このような体 $\mathfrak{K}$ に同型な $P_0$ の任意の拡大 $P$ は、ある（単純）指標の体である。$\mathrm{P}_0$ を $\mathfrak{K}$ に同型な $\mathrm{P}$ へ拡大すると、$\mS$ から、指標の絶対既約表現類が対応する両側単純環が切り離される。Wedderburn の定理により、この環は、$\mathrm{P}$ 上の行列環と非可換体 $\mA$ との直積である。$\mA$ は同型を除いて一意に定まり、その中心は $P$ であり、$P$ に関して次数 $m^2$ をもつ。ここで $m$ はその指標に属する Schur 指数を表す。$\mS$ の分解体 --- 指標に属する表現に関する分解体 --- は、同時に $P$ における表現に関する $\mA$ の分解体であり、逆も成り立つ。$\mS$ の共役表現は、$\mathrm{P}_0$ 上の対応する両側単純環から出発して得られる。この環に対応する非可換体 $\mA_0$ は $\mA$ に同型となり、$\mathfrak{K}$ を中心にもつ。こうして、分解体の問題は、対応する非可換体 $\mA$ とその分解の問題へ完全に還元される。\footnote{これは、$\mathrm{P}$ における $\mS$ の既約表現と可換な $\mathrm{P}$ 上の行列の系への Schur の還元に対応する。これらの行列の転置が、$P$ における $\mA$ の表現を与える。} 特に次の定理が得られる。

$\mA$ のすべての極大可換部分体は、$P$ に関して同じ次数 $m$ をもつ。ここで $m$ は Schur 指数である。これらの極大可換部分体はすべて最小次数の分解体であり、最小次数の任意の分解体はそれらの中に含まれる。任意の分解体の次数は $m$ の倍数である。環 $\mA$ はただ一つの絶対既約表現類をもち、また $\mathrm{P}$ に関してただ一つの既約表現類をもつ。

一般の分解体を特徴づけるため、$\mA_r$ により、$\mA$ と $P$ 上次数 $r^2$ の行列環（$\mathrm{P}$ の元を成分にもつ $r$ 行行列全体の系）との直積を表す。すると $\mA_r$ は両側単純であり、対応する非可換体は $\mA$ である。また、$\mathrm{P}$ 上超複素的な任意の両側単純環で、$\mA$ が対応するものは、この方法で得られる。$\mA_r$ は、$\mA$ の元を成分にもつ $r$ 行行列全体の系としても特徴づけられる。

$\mA$ の次数 $mr$ の分解体 --- そのような体が存在すれば --- は、$\mA_r$ の極大可換部分体に同型である。逆に、$\mA_r$ のすべての極大可換部分体は分解体であり、その次数はそれぞれ $ms$ で、$s$ は $r$ の約数である。正則の場合には、$\mA_r$ のすべての極大可換部分体の次数は $mr$ である。ここで正則の場合とは、基礎体 $\mathrm{P}$ が次の性質をもつことをいう。すなわち、$P$ 上有限な任意の可換体 $\Sigma$ に対して、$\Sigma$ 上の任意次数の可換拡大がなお存在することである。\footnote{したがって、たとえば $\mathrm{P}$ がすべての実数から成る体である場合には、正則の場合にはならない。}

これらの分解体のうち、$r>1$ に対して極小のものも現れるかどうかは、$\mA_r$ への埋め込みだけではまだ決まらない。次の例は、これが実際に無限に多くの $r$ について起こりうることを示す。

\subsection*{2. 極小分解体の次数の非有界性}

この非有界性の例として、四元数体を取り上げる。それを、有理数体 $\mathrm{P}$ 上の超複素系とみなし、単位元 $1$ のほかに基底元 $i,j,k$ をもち、よく知られた関係
\[
 i^2=j^2=k^2=-1;\qquad ij=k;\quad jk=i;\quad ki=j;\quad ji=-k;\quad kj=-i;\quad ik=-j.
\]
を満たすものとする。示すべきことは、代数的数体のうち、添加後に $r^2=r$ を満たす「冪等」元 $r$ が存在し、しかも $r$ が零元とも単位元とも異なるもの、そしてそれらだけが分解体である、ということである。これらの数体は、$(-1)$ が三平方和として表され、したがって二平方和としても表される体としても特徴づけられる。\footnote{$(-1)$ が三平方和として表されることから、常に二平方和としての表現が従うことは、恒等式
\[
(c^2+d^2)(a^2+b^2+c^2+d^2)=(ac+bd)^2+(ad-bc)^2+(c^2+d^2)^2
\]
から分かる。これは、たとえば四元数のノルムの積定理から従う。分解のためには、冪等元の存在を要求する代わりに、ノルム零の元の存在を要求しても十分であった。} 実際、
\[
        r=\alpha+\frac12\beta i+\frac12\gamma j+\frac12\delta k
\]
とおくと、
\[
        r^2=\left(\alpha^2-\frac14\beta^2-\frac14\gamma^2-\frac14\delta^2\right)
          +\alpha\beta i+\alpha\gamma j+\alpha\delta k
\]
であるから、$r^2=r$ は
\[
4\alpha^2-4\alpha-\beta^2-\gamma^2-\delta^2=0;\quad
2\alpha\beta-\beta=0;\quad 2\alpha\gamma-\gamma=0;\quad 2\alpha\delta-\delta=0
\]
と同値であり、したがって $\beta,\gamma,\delta$ が同時に零であってはならないことから、
\[
        (-1)=\beta^2+\gamma^2+\delta^2
\]
と同値である。

このような任意の冪等元が分解を与え、またそのような冪等元の存在が分解に必要であることは、次の事実から従う。完全可約系の任意の既約表現は、片側単純イデアル、より正確には対応するイデアル類によって生成され、またそのような各イデアル類は表現を生成する。これらの片側単純イデアルはすべて片側直和分解の中に現れ、この分解が、単位元の成分となる「原始的」冪等元を決定する。逆に、任意の冪等元は直和分解 $\mS=r\mS+(e-r)\mS$ を与える。ただし $e$ は単位元を表す。そのうち「原始的」なものが単純イデアルを切り離す。\footnote{原始的冪等元と既約表現との関係は、M. Herzberger の学位論文 Über Systeme hyperkomplexer Größen, Chapter 4, Berlin 1923 に既に見られる。単純イデアルが既約表現を生成することは、可換の場合には、たとえば E. Noether, Der Diskriminantensatz für Ordnungen, Crelle 157 (1926) により、非可換の場合には W. Krull, Theorie und Anwendung der verallgemeinerten Abelschen Gruppen, Heidelberger Berichte 1926 により示されている。Speiser の群論の末尾の注も参照されたい。$a_1,\cdots,a_m$ を分解体に関するイデアルの基底とし、$c$ を任意の環元として、$a_i c=\gamma^i_1 a_1+\cdots+\gamma_{im} a_m$ とすれば、$(\gamma_{ik})$ が $c$ に対応する行列である。すべての既約表現がイデアルによって生成されるという事実は、さらに深い。} 非可換体を、その中心に関する超複素系とみなすと、分解体に関する超複素系へ拡大したとき、$m$ 個の絶対単純な片側イデアルの直和になる。ここで $m$ は Schur 指数である。

四元数体の場合、$m$ は 2 である。したがって $0$ でも $1$ でもない任意の冪等元は、絶対単純イデアルへの分解を既に与え、それに絶対既約表現類への分解が対応する。これにより、先に述べた体が分解体として特徴づけられる。

さらに次が従う。$(-1)$ が二平方和として表される、有理数体上次数 $2^n$ のすべての巡回体 $\Omega_n$ は、四元数体の極小分解体である。実際、$(-1)$ が平方和として表されるため、分解体は実体ではありえない。一方、$\Omega_n$ の次数 $2^{n-1}$ の部分体 --- したがって $\Omega_n$ の任意の真の部分体 --- は必ず実である。なぜなら、それは $\Omega_n$ とすべての実代数的数の体との交わりだからである。すなわち、$\Omega_n$ はこの交わりの上の次数 2 のガロア体であり、したがってこの交わりは次数 $2^{n-1}$ の部分体と一致しなければならない。

Hasse の存在定理によれば\footnote{$\Omega_2$ の存在は、たとえば Weber (Kleines Lehrbuch der Algebra, § 93) に与えられている次数 4 の巡回体のパラメータ表示から既に読み取ることができる。}、このような体 $\Omega_n$ はすべての $n$ に対して存在する。したがって極小分解体の次数は有界でない。実際、指数の各冪がここでは極小分解体の次数として現れる。\footnote{その後 R. Brauer はさらに、四元数体には任意の偶数次数の極小分解体さえ存在すること、したがって分解体に対してそもそも可能なすべての次数で存在することを示すことができた。このような体 $\mathrm{P}(\vartheta)$ は、任意の $n\geqq2$ に対して方程式
\[
 f(\vartheta)=\vartheta^{2n}+a^2p^2\vartheta^2+b^2p^{2\beta}=0;
 \quad \text{したがって}\quad -1=\left(\frac{ap}{\vartheta^{n-1}}\right)^2+
       \left(\frac{bp^\beta}{\vartheta^n}\right)^2
\]
によって定義できる。ただし $a,b,p,\beta$ を適当に選ぶことで、1. $f(x)$ がすべての $n$ に対して既約であること、2. $\mathrm{P}(\vartheta)$ が $\mathrm{P}$ と異なる唯一の部分体として $\mathrm{P}(\vartheta^2)$ だけをもつこと、3. $\mathrm{P}(\vartheta^2)$ がその共役の中に実体をももち、したがって分解体ではありえないこと、を実現できる。$a,b,p,\beta$ の一つの可能な選択は、たとえば
\[
        f(x)=x^{2n}+(25\cdot9)4x^2+9\cdot64
\]
である。証明は、原始方程式を構成する Furtwängler の方法（Math. Ann. 85, p. 34, § 3）の修正に基づくが、細部の実行はかなり煩雑である。}

\subsection*{3. 四元数体の分解体の初等的取り扱い}

$\mathfrak{O}$ を、次の八つの行列から成る四元数群の表現とする。\footnote{Cf. Sylvester, Math. Papers Vol. III p. 647, Vol. IV p. 122.}
\[
\pm E=\pm\mat{1&0\\0&1},\qquad
\pm I=\pm\mat{i&0\\0&-i},\qquad
\pm J=\pm\mat{0&1\\-1&0},\qquad
\pm K=\pm\mat{0&i\\ i&0}.
\]
$\mathfrak{O}$ が生成する群環は、超複素系として見た四元数体の表現にほかならない。$\mathfrak{O}$ が体 $\mathrm{K}$ において表現可能であるための必要十分条件は、$-1$ が $\mathrm{K}$ において二平方和として表されることであることを、初等的に示す。

\emph{a)} $\mathfrak{O}^*=T^{-1}\mathfrak{O} T$ を、体 $\mathrm{K}$ 上有理的な表現とし、$I^*=T^{-1}IT$, $J^*=T^{-1}JT$ とおく。$J$ と $J^*$ は $\mathrm{K}$ 上有理的な二つの相似行列であるから、$\mathrm{K}$ 上有理的な変換 $R$ で、$J^*$ を $J$ へ移すものも存在する。
\[
        R^{-1}J^*R=J.
\]
最初から $\mathfrak{O}^*$ を同じく $K$-有理的な表現 $R^{-1}\mathfrak{O}^*R$ で置き換えれば、一般性を失わず $J^*=J$ と仮定できることが分かる。
\[
        I^*=\mat{a&b\\ c&d}\qquad (a,b,c,d\text{ は }K\text{ の数})
\]
とする。$IJ=-JI$ から $I^*J^*=-J^*I^*$ が従い、$J^*=J$ であるから
\[
        a=-d,
        \qquad b=c.
\]
を得る。また $I^2=-E$ から $I^{*2}=-E$ が従うので、
\[
        \mat{-1&0\\0&-1}=\mat{a&b\\ b&-a}^2
        =\mat{a^2+b^2&0\\0&a^2+b^2},
\]
すなわち $a^2+b^2=-1$ である。したがって $-1$ は確かに $\mathrm{K}$ において二平方和として表される。

\emph{b)} 逆に、$-1$ が体 $K$ で二平方和として表されると仮定する。たとえば $-1=a^2+b^2$ ならば、
\[
        I^*=\mat{a&b\\ b&-a},\qquad
        J^*=\mat{0&1\\-1&0},\qquad
        \mathit{K}^*=I^*J^*
\]
とおく。直接計算により、$\pm E$, $\pm I^*$, $\pm J^*$, $\pm \mathit{K}^*$ は、$K$ 上有理的で $\mathfrak{O}$ に相似な群をなすことが分かる。\footnote{分解体の特徴づけは、因子系を用いても容易に得られる。}

残るのは、無限に多くの $n$ に対して、次数 $2^n$ の巡回体で四元数群の極小分解体となるものが存在することを示すことである。

$p$ を、適当な $r>0$ に対して
\[
        (*)\qquad 2^r+1\equiv0\pmod p
\]
を満たす有理素数とし、$p-1$ を割る 2 の最大冪を $2^n$ とする。

まず、素数 $p$ が無限に多くの $n$ について現れることを示す。$k$ を任意の正の有理整数とし、$p$ を $2^{2^k}+1$ の素因数とする。このとき $p$ に要求された合同式は $r=2^k$ によって満たされる。さらに $2\pmod p$ の指数は $2^{k+1}$ であるから、$p-1$ を割る最大の 2 冪 $2^n$ について、$n$ は $k$ より大きい。したがって、少なくとも任意に大きい $n$ に対応する素数 $p$ が存在する。

$\varepsilon$ を原始 $p$ 乗根とする。いま $\mathrm{P}(\varepsilon)$ において $-1$ が二平方和として表されることを示す。
\[
        -1=a^2+b^2=(a+ib)(a-ib)\qquad (a,b\text{ は }P(\varepsilon)\text{ の数}).
\]
これは、$(4p)$ 乗根の体 $\mathrm{P}(\varepsilon,i)=\mathrm{P}(\varepsilon\cdot i)$ に、$\mathrm{P}(\varepsilon)$ に関する相対ノルムがちょうど $-1$ である数が存在する、ということと同値である。

数 $1+i\varepsilon^\nu$ $(\nu=1,2,\ldots,p-1)$ の相対ノルムは $(1+i\varepsilon^\nu)(1-i\varepsilon^\nu)=1+\varepsilon^{2\nu}$ である。したがって、すべての数 $1+\varepsilon^\nu$ $(\nu=1,2,\ldots,p-1)$ は、$\mathrm{P}(\varepsilon i)$ の数の、基礎体 $\mathrm{P}(\varepsilon)$ に関する相対ノルムである。したがって同じことが
\[
\Pi=\prod_{\nu=0}^{r-1}(1+\varepsilon^{2^\nu})
=(1+\varepsilon)(1+\varepsilon^2)(1+\varepsilon^4)\ldots(1+\varepsilon^{2^{r-1}})
=\sum_{\lambda=0}^{2^r-1}\varepsilon^\lambda
\]
についても成り立つ。$\Pi$ に項 $\varepsilon^{2^r}=\varepsilon^{-1}=\varepsilon^{p-1}$（$(*)$ による）を付け加えると、$\Pi+\varepsilon^{p-1}$ は $1+\varepsilon+\varepsilon^2+\ldots+\varepsilon^{p-1}=0$ という部分の和になる。したがって
\[
        \Pi=-\varepsilon^{-1}.
\]
さらに $\varepsilon$ も $\mathrm{P}(\varepsilon i)$ の数の相対ノルムであり、すなわち $\varepsilon^{\frac{p+1}{2}}$ のノルムである。ゆえに同様の主張は
\[
        \Pi\cdot\varepsilon=-1
\]
について成り立つ。したがって $-1$ は $\mathrm{P}(\varepsilon)$ において二平方和として表される。\footnote{$n8\pm3$ の形の素数に対しては、この議論は既に E. Landau, Über die Darstellung definiter Funktionen durch Quadrate, Math. Ann. vol. 62, pp. 272--285 に見られる。} $\mathrm{P}(\varepsilon)$ は四元数群の分解体である。

$\mathrm{K}$ を $\mathrm{P}(\varepsilon)$ に含まれる次数 $2^n$ の巡回体とする。$\mathrm{K}$ も分解体であると主張する。\footnote{次数 $2^n$ の巡回分解体を円分体の部分体として選ぶことは、Hasse の例を範としている。} $m$ を、$\mathrm{K}$ を基礎体としたときの四元数体の指数とする。このとき $m=1$ または $m=2$ である。しかし $K$ に関して奇数の相対次数 $(p-1)/2^n$ をもつ分解体 $\mathrm{P}(\varepsilon)$ が存在するので、第二の場合は不可能である。したがって $m=1$ である。これはまさに $\mathrm{K}$ が分解体であることを意味する。

こうして、無限に多くの $n$ について、次数 $2^n$ の巡回体で極小分解体となるものが存在する。

\section*{33. 超複素量と表現論の算術的把握}
\emph{Atti Congresso Bologna 2 (1928), pp. 71--73.}

\begingroup
\renewcommand{\thefootnote}{(\arabic{footnote})}

\begin{center}
\textsc{Emmy Noether}（Göttingen -- Germania）\\[-0.1em]
\rule{3em}{0.4pt}\\[0.8em]
\textsc{超複素量と表現論}\\
\textsc{算術的立場から}\footnote{詳しい叙述は \emph{Mathematische Zeitschrift} に掲載される予定であり、その続編では非可換ガロア理論も論じられる。分解体については R. Brauer--E. Noether, \emph{Berichte der Berliner Akademie} 1927 を、一般構造定理については本論集所収の G. Köthe の講演をも参照。}
\end{center}

ここでは、表現論、とりわけ群の表現論を、加群類およびイデアル類の理論として捉えられること、さらに後二者が、作用素をもつ群という一般化された群概念のもとに再び組み込まれることを示したい。\footnote{作用素をもつ群は W. Krull（\emph{Math. Z.} 23）および O. Schmidt（\emph{Math. Z.} 29）に遡る。与えられた表現の簡約問題をここで扱う方法も、可換な表現体に限れば、W. Krull, \emph{Theorie und Anwendung der verallgemeinerten Abelschen Gruppen}, \emph{Heidelberger Berichte} 1926 に由来する。}

あらかじめ指定された体 $P$ における\emph{群の表現}とは、周知のように次のものをいう。群 $\mathfrak G$ の元 $a,b$ に、$P$ の係数をもつ行列 $A,B$ を対応させ、積 $ab$ に積 $AB$ が対応するようにする。この行列系は群の準同型像となる。この表現と同時に、「群数」、すなわち線型結合
\[
        a\alpha+b\beta+\cdots+c\gamma
\]
の全体からなる「群環」の表現も与えられる。ここで $a,b,....,c$ はその都度有限個の群元を、$\alpha,\beta,....,\gamma$ はそれぞれ有限個の $P$ の元を走り、群元は線型独立とみなし、乗法は群の乗法と演算法則によって定義する。群環の表現は加法についても準同型である。こうして群の表現は、加法と乗法の双方に関する準同型を要求する一般の\emph{環の表現}問題に含まれる。

ここで本質的な問題は二つある。第一は、\emph{与えられた表現の簡約}と、その際に現れる既約成分の一意性である。第二は、与えられた環を、与えられた必ずしも可換でない体に表現する仕方の全体、あるいはそのうち既約表現だけの全体を問う問題である。

第一の問題ははるかに単純である。このことは、表現される環にも、一般には非可換である表現体にも、いかなる種類の特別な仮定も必要ないという事実から既に分かる。表現が固定次数の行列によって行われるという事実が、必要な有限性仮定をすべて与える。この問題は作用素をもつ群の理論によって完全に扱うことができるので、ここで簡単に述べる。

群 $\mathfrak{B}$ は、同時に記号系 $\varTheta,H$ --- 作用素 --- が与えられ、$\mathfrak{B}$ の各 $a$ に対して $\varTheta(a),H(a)$ が一意に定義されて $\mathfrak{B}$ の元を与え、さらに分配法則 $\varTheta(ab)=\varTheta(a)\varTheta(b)$ が満たされるとき、作用素領域をもつ群と呼ばれる。同じ作用素領域をもつ二つの群は、通常の意味で同型であり、さらに $a$ と $\bar a$ の対応から $\varTheta(a)$ と $\varTheta(\bar a)$、$H(a)$ と $H(\bar a)$ などの対応も従うとき、作用素同型であるという。そこで部分群としては、それ自身が作用素領域を許すものだけを認めることにする（したがって群が単純であるとは、単位元以外に許容される真の正規因子をもたないことをいう）。このとき合成列に関する Jordan--Hölder の定理は次の形でなお有効である。作用素領域をもつ群がそもそも合成列をもつならば、因子群の系は、作用素同型の意味で、順序を除いて一意に定まる。同じ仮定のもとで、直接分解不能因子への分解も、作用素同型の意味で一意である。

表現問題との関係は、作用素をもつ群の中に、とくにイデアルと加群が含まれるという事実によって与えられる。これらは加法に関してアーベル群とみなされ、作用素は環元による乗法で与えられる。しかし任意の表現は表現加群、すなわち環および表現体に関する加群によって生成される。より正確には、任意の加群類、すなわち互いに作用素同型な表現加群の任意の類は、同じ表現を生成する。したがって、表現の簡約の一意性の問題は、合成列定理および直積に関する定理によって答えられる。表現加群に適用すると、合成列は各行列に対し、次の形の簡約を与える。
\[
\begin{pmatrix}
B_1&0&0\\
&B_2&0\\
A_{ik}&....&B_r
\end{pmatrix},
\]
ここで因子群に対応する対角行列は「既約」表現、すなわちそれ自身はもはやこのような簡約を許さない表現を生成する。

これらの因子群の類は一意に定まるので、対角表現についても同じことが成り立つ（表現の同値の意味で一意）。同様に、直積に関する定理は、「分解不能」成分への「分解」の一意性を与える。
\[
\begin{pmatrix}
C_1&&&\\
&C_2&&\\
&&\ddots&\\
&&&C_s
\end{pmatrix},
\]
ここで各 $C_i$ は、さらに合成列定理に従って一意に簡約されうる。

第二の問題については、「二重鎖定理」を満たす超複素系、より一般には環の表現の場合を概説したい。結果は、すべての既約（単純）加群類はイデアル類である、というものである。したがって片側イデアルから得られる合成列は、その因子群を通じて、すべての既約表現を既に与える。より正確には、根基を法とする剰余類環における合成列で十分であり、後者は「完全可約」環になる。こうして問題は、完全可約環の構造研究に還元される。したがって、自己同型斜体における表現を出発点とするのが適切である。これは次のように理解される。こうした環の両側単純成分のすべての単純（右）イデアルは同じ類に属し、したがって同じ自己同型環をもつ。イデアルが単純であるため、この自己同型環は斜体になる。同時にこれは単純左イデアルの自己同型環でもある。単純環そのものは、自己同型斜体上のすべての行列から成る系に同型になる。そしてこの表現から、部分体に関するすべての表現、とくに超複素系の係数領域に関するすべての表現が、自己同型斜体の元そのものを対応する行列で置き換えることによって生じる。こうして、超複素系の場合に知られている Wedderburn の構造定理は、一般群論に由来する概念である自己同型斜体の概念によってここで再解釈されると同時に、超複素系の表現論の源泉として認識される。非可換体（斜体）のガロア理論という新しい問題が生じる。それは「分解体」の問題と密接に結びついている。同時に、分解不能イデアルの自己同型環を通じて、完全可約でない一般の環に対する構造定理への見通しも得られる。最も一般的な作用素をもつ群へ拡張された群論は、ここでも再び秩序づけの原理であることを示している。

\endgroup



\clearpage


\setcounter{footnote}{0}

\section*{34. 超複素量と表現論}
\emph{Math. Zs. 30 (1929), pp. 641--692.}

\subsection*{序論}

超複素系に関するもっとも重要な一般定理は、Molien（\emph{Math. Annalen} 第41巻および1897年の Dorpat 報告）に遡る。これと本質的に独立に、Frobenius はその少し後に、超複素系の理論とその表現の理論、とくに有限群の表現論を、一つの統一的な理論として展開した。基礎にあるのは Dedekind による群行列式の概念、より一般には任意の超複素系の行列式である。Frobenius は、この行列式の相異なる既約因子（係数領域はつねに全複素数体とする）に、相異なる既約表現類が対応し、この方法で既約表現類の全体が尽くされることを示す。そのような表現類は、その「指標系」によって完全に特徴づけられる。有限群の場合、この指標系は、群の共役元の類から作られる可換超複素系の行列式を分解することによって得られる。この行列式は、指標を係数にもつ一次因子に分解する。これは、Dedekind が最初に得た結果、すなわち有限アーベル群の群行列式は、そのアーベル群の相異なる指標を係数にもつ一次因子に分解するという結果（Frobenius との往復書簡）を直接に一般化するものである。\footnote{以下は、1927/28年冬学期の私の講義を B. L. van der Waerden が自由に展開したものである。われわれは共同で印刷用の本文を準備した。また、いくつかの批判的な指摘についても B. L. van der Waerden に負うている。}

しかし、概念的には単純で見通しのよいこれらの結果は、Frobenius においては労を要する計算によって得られている。その後の発展は、これらの結果をより簡単に導くこと、同時に係数領域として任意の体をとる場合へ拡張することを目指した。この後の発展は、超複素系と表現論とで、まったく別々の道を進んだ。超複素系は Wedderburn によって算術的な取り扱いを受けた。すなわち、任意の超複素系は、両側に直接分解不能な環の直和として一意的に表される。根基をもたない系では、これらの分解不能環は、指定された（必ずしも可換ではない）斜体の元を成分とするすべての $n$ 行行列の系と同型になる。その斜体は本質的に一意に定まる。\footnote{Dickson, \emph{Algebras and their Arithmetics}（ドイツ語版 \emph{Algebren und ihre Zahlentheorie}, Zürich 1927）の参考文献を参照。} 表現論は Burnside および I. Schur によって初等的な基礎づけを受けた。彼らは超複素系から独立に、与えられた表現から直接出発し、行列定理を用いて操作した。とくに Schur は、与えられた体上で既約な表現が絶対的に分解するような最小次数の数体の問題を扱った。これらの絶対既約成分は、有限個の共役な表現類に属する。求める数体の次数は、これらの類の数と、指数（各共役類に含まれる成分の数）との積に等しい。Schur の理論の主要な補助手段は、既約表現と可換なすべての行列からなる系である。\footnote{Speiser の群論、第10章の参考文献を参照。}

以下の純粋に算術的な基礎づけでは、超複素量と表現論とは再び一つの統一的全体として現れる。すなわち、ある有限性条件だけを満たす非可換環の一般理論の特殊な場合として現れる。より正確には、これはそれらの環に関する加群類とイデアル類の理論であり、その主要な結果は、既約加群類が対応するイデアル類によってすでに尽くされる、ということである。とくに根基をもたない環では、すべての加群類は既約なものに分解する（完全可約になる）。これは、正則な群行列式、あるいは系の行列式の既約因子が既約表現類の全体をすでに尽くし、また根基をもたない系について表現の完全可約性が成り立つ、という Frobenius の結果の算術的な対応物である。実際、系の行列式とその分解を考えることはノルムへの移行として解釈できるが、ここではイデアルの直和分解とその組成列を直接に扱う。表現論は、表現加群の概念によって、加群類の理論となる。

加群類とイデアル類の導入は、純粋に群論的な基礎をもっている。加群とイデアルは、加法に関してアーベル群と見なされ、さらに環の元による一定の乗法を許容するという条件で制限される。すなわち、それらは「作用素付き群」（§1）をなす。そのような群では、通常の同型のかわりに「作用素同型」（§2）が現れる。互いに作用素同型な群（固定された領域の中での同型）は一つの類にまとめられる。この類概念は、たとえば数体のイデアルの場合には通常のものと一致する。作用素付き群の理論は W. Krull と O. Schmidt に遡る（注6参照）。それは第I章で体系的に展開される。ここでもなお有効な定理、すなわち組成列、直積（それぞれ直和）、完全可約群、および先に述べた類分けの意味での一意性定理が、以下すべての基礎をなす。

そこからまず、任意の表現の既約対角成分についての一般的な一意性定理（第III章 §16）が得られる。これは、表現類と表現加群の類、したがって作用素付き群の類との一対一対応に基づいている。表現全体に関するさらなる問題には、表現されるべき環の構造、したがって特殊な場合には超複素系の構造を、より精密に調べることが必要である。第II章では、片側イデアルを前面に置く群論的見方に基づいて、Wedderburn の結果が新たに導かれ、さらに進められる。右イデアルについての「倍鎖定理」、すなわち同じこととしての「極小条件」（右イデアルの任意の集合の中に、その集合内で極小な元が少なくとも一つ存在すること）が、有限性条件として十分であることがわかる。\footnote{「二重鎖定理」を仮定すれば Wedderburn の証明方法は移すことができる。E. Artin, \emph{Zur Theorie der hyperkomplexen Zahlen}, Hamb. Abh. 1927 を参照。また A. Suschkewitsch, \emph{Über die endlichen Gruppen ohne das Gesetz der eindeutigen Umkehrbarkeit}, Math. Annalen 99 (1928), pp. 30--50 も参照。そこでは、一つの結合的で可逆とは限らない演算だけをもつ（有限）領域について、平行する構造定理が展開されている。Suschkewitsch における「核」の右群と左群への分裂は、単位元をもつ両側単純環を右イデアルと左イデアルの和として表すことに対応している。} その結果、極小条件を満たす根基なしの環は、単位元をもつ（右）完全可約環と一致することが得られる。そのような環では、両側分解不能環のすべての単純右イデアルは同じ類に属し、したがって作用素付き群として同じ自己同型環をもつ。イデアルが単純であるため、この自己同型環は斜体になる。この類の自己同型斜体こそが、Wedderburn の見いだした行列表現を媒介する。

この自己同型斜体における行列表現は、完全可約な場合の表現問題も解決する。ただし、自己同型斜体またはその部分体に関する表現を要求する場合、特に超複素系の係数領域に関する表現を要求する場合である。加群類をイデアル類に還元する定理の直接の帰結として、Wedderburn 表現およびそこから、自己同型斜体の元を自然な方法で部分体上の行列に置き換えて得られる表現以外に、既約表現は存在しない。

したがって完全可約な場合には、既約表現類の相異なる数はイデアル類の数と同じである。その数は、相異なる両側分解不能環、したがって単純環の数と一致する。さらにこの数は、可換体となる中心の分解不能成分の数と同一である。係数領域が代数閉である超複素系の場合、これらの体は中心の相異なる環準同型（既約表現）によって完全に定まり、数値因子を除けば、それらがちょうど指標である。

得られた結果は同時に、根基をもつ系の既約表現についての見通しも与える。すなわち、それらは根基で割った剰余類環の表現に還元でき、その剰余類環は根基をもたない系となる範囲においてである。

後に示すように、根基をもたない超複素系については、「分解体」の理論も自己同型斜体から従う。すなわち、表現が絶対既約表現に分解するような係数領域の可換拡大、同値に、絶対単純な右イデアルへの直和分解が生じるような可換拡大の理論である。分解体は自己同型斜体の分解体と一致し、最小次数のすべての分解体は自己同型斜体の極大可換部分体と同型である。上に述べた Schur の研究との関係は、表現と可換な行列の転置が、まさに自己同型斜体の表現を与えるという事実によって確立される。\footnote{R. Brauer--E. Noether, \emph{Über minimale Zerfällungskörper irreduzibler Darstellungen}, Sitz.-Ber. Preuß. Akad. Wiss. 1927, p. 221 を参照。} これらの定理は、非可換斜体のガロア理論の枠内で体系的に展開されるべきものである。

基礎にある概念、すなわち作用素準同型と自己同型環は、極小条件だけを満たす一般の根基付き環の構造も与える。これは別の側面から遂行される。

\section*{第I章　群論的基礎}

\subsection*{§ 1. 作用素付き群}

$\mG$ を（有限または無限の）群とし、その元を $a,b,\ldots$ とする。群 $\mG$ に対する作用素領域 $\Omega$ とは、新しい記号 $H,\Theta,\ldots$ の集合であって、$\mG$ の各 $a$ と $\Omega$ の各 $\Theta$ に対して、$\mG$ の中の一意に定まる $\Theta a$ が対応し、かつ分配法則
\[
        \Theta(ab)=\Theta a\cdot\Theta b
\]
が成り立つものをいう。したがって、各作用素は群からそれ自身への準同型を定め、そのとき群はそれ自身の上、またはある部分群の上に写される。単位元は単位元に移り、逆元は逆元に移る。

作用素領域は、異なる作用素が異なる準同型を定めるとき、\emph{絶対的}であるという。任意の作用素領域から、同じ準同型を生成する元を同一視することにより、絶対的なものが得られる。絶対的な作用素領域は、群からそれ自身へのすべての準同型からなる集合のある部分集合の一対一像である。

固定された作用素領域 $\Omega$ に関して、$\mG$ の\emph{許容部分群} $\mH$ とは、その作用素領域を許容する部分群、すなわち $\mH$ の任意の $a$ と $\Omega$ の任意の $\Theta$ について $\Theta a$ が $\mH$ に属するような部分群をいう。

\emph{例。} 1. 作用素を内部自己同型 $\Theta a=c^{-1}ac$ とする。許容部分群は正規因子（正規部分群）である。

2. 作用素をすべての自己同型とする。許容部分群は、すべての自己同型の下で自分自身へ移る「特性部分群」である。\footnote{ここで説明した概念は、W. Krull, \emph{Über verallgemeinerte endliche Abelsche Gruppen}, Math. Zeitschr. 23 (1925), pp. 161--196、および O. Schmidt, \emph{Über unendliche Gruppen mit endlicher Kette}, Math. Zeitschr. 29 (1928), pp. 34--41 に由来する。}

3. $\mG$ を環とする。すなわち、加法に関してアーベル群であり、さらに乗法が定義され、
\[
        r(a+b)=ra+rb,\qquad (a+b)r=ar+br,\qquad ab\cdot c=a\cdot bc
\]
を満たすものとする。各元 $r$ は同時に二つの作用素、すなわち $rx$ と $xr$ を定める。許容される部分群は「イデアル」$\mathfrak a$ である。すなわち、操作 $rx$ を許容する左イデアル $r\mathfrak a\subseteqq\mathfrak a$、操作 $xr$ を許容する右イデアル $\mathfrak a r\subseteqq\mathfrak a$、および両方の操作を許容する両側イデアルである。環 $\mG$ が斜体、すなわち $\mG-\{0\}$ が乗法に関して群であるとき、すべてのイデアルは自明 $(=\{0\}\text{ または }=\mG)$ となる。\footnote{したがって、環についても斜体についても、乗法の可換法則は仮定されない。}

4. \emph{環 $\mo$ に関する加群。} $\mo$ を環、$\mM$ を加法的に書かれたアーベル群とする。$\mo$ の元と $\mM$ の元との積 $r\cdot a$ が与えられ、その積は再び $\mM$ に属すると仮定する。次を要求する。
\[
\left.
\begin{aligned}
 r(a+b)&=ra+rb,\\
 rs\cdot a&=r\cdot sa,\\
 (r+s)a&=ra+sa,
\end{aligned}
\right\}\qquad r,s\in\mo,\quad a,b\in\mM.
\]
このとき $\mM$ は $\mo$-加群と呼ばれる。より正確には、$\mo$ からの乗数が左に書かれているので、左 $\mo$-加群である。環 $\mo$ は同時に作用素領域である。それが絶対的である、すなわち異なる環元が異なる操作を与えるなら、$\mM$ に対する\emph{絶対的な乗数領域}という。

上と同じく、任意の乗数領域から、同じ操作を生成する元を同一視することで絶対的なものへ移ることができる。絶対的な乗数領域は再び環である。加群 $\mM$ の許容部分群は\emph{部分加群}と呼ばれる。特に $\mo=\mM$ なら、再び左イデアルが得られる。\footnote{同様に、右加群については
\[
(a+b)r=ar+br,\qquad a\cdot rs=ar\cdot s,\qquad a(r+s)=ar+as
\]
を要求する。}

5. \emph{双加群。} $\mM$ が同時に左 $\mo$-加群かつ右 $\mo'$-加群であり、さらに $\mo$ の $a$、$\mM$ の $\alpha$、$\mo'$ の $a'$ についてつねに
\[
        a\cdot\alpha a'=a\alpha\cdot a'
\]
が成り立つなら、$\mM$ は双加群と呼ばれる。

6. \emph{アーベル群の自己同型環。}\footnote{A. Châtelet, \emph{Les Groupes Abéliens finis}, Paris, Gauthier-Villars, 1925, p. 99 を参照。} アーベル群（加法的に書く）の自己準同型については、次の式により加法と乗法を定義できる。
\[
        (\mathsf{H}+\Theta)a=\mathsf{H}a+\Theta a,\qquad (\mathsf{H}\Theta)a=\mathsf{H}(\Theta a).
\]
このように定義された操作は明らかに再び準同型を表し、したがってまた同じ系に属する。この系は環をなし、そのアーベル群は、4 に従って、この環に関する加群と見なすことができる。

以下では、「部分群」といえば常に許容部分群を意味する。二つの許容部分群の共通部分 $\mA\cap\mB$ は再び許容部分群である。積 $\mA\mB$ も、二つの部分群 $\mA,\mB$ の少なくとも一方が正規因子であるなら、再び許容部分群である。

\subsection*{§ 2. 同型定理}

群 $\mG$ から群 $\overline{\mG}$ の上への写像を考える。ただし $\mG$ と $\overline{\mG}$ は同じ作用素領域をもつものとする。この写像が、第一に通常の意味で準同型であり（$ab\rightarrow \bar a\bar b$）、第二に $a$ が $\bar a$ に写るとき $\Theta a$ が $\Theta\bar a$ に写るなら、それを作用素準同型という。\footnote{作用素準同型の概念は、群 $\mG,\overline{\mG}$ が別々の作用素領域をもつことを許すことで拡張できる。その場合、準同型は $\mG$ を $\overline{\mG}$ の上へ写すだけでなく、作用素領域どうしも写し、$a$ が $\bar a$ に、$\Theta$ が $\overline\Theta$ に写るなら、$\overline\Theta\bar a$ が $\Theta a$ の像となるようにする。この一般形では、作用素準同型の概念は環準同型の概念も含む。§8 参照。} 記法は $\mG\sim\overline{\mG}$ とする。対応が一対一なら、それを作用素同型という。記法は $\mG\simeq\overline{\mG}$ である。通常の群の場合と同様に、準同型定理が証明される。

\emph{$\overline{\mG}$ が $\mG$ の作用素準同型像であるなら、$\overline{\mG}$ は因子群 $\mG/\mN$ と作用素同型である。ここで $\mN$ は許容正規因子であり、$\overline{\mG}$ の単位元に対応する $\mG$ のすべての元からなる。逆に、任意の許容正規因子 $\mN$ は、$\mG$ の作用素領域を許容する因子群 $\mG/\mN$ を定め、それは $\mG$ の作用素準同型像である。}

群は、自分自身と単位群以外の正規因子をもたないとき\emph{単純}であるという。単純群の任意の準同型は、同型であるか、さもなければすべての元を単位元に対応させる。

\emph{第一同型定理。} \emph{$\overline{\mG}$ を $\mG$ の準同型像、$\overline{\mA}$ を $\overline{\mG}$ の正規因子、$\mA$ を $\overline{\mA}$ の元に対応する $\mG$ の元全体とする。このとき $\mA$ は再び正規因子であり、}
\[
        \overline{\mG}/\overline{\mA}\simeq \mG/\mA. \tag{1}
\]
\emph{証明。} $\mG\sim\overline{\mG}$、$\overline{\mG}\sim\overline{\mG}/\overline{\mA}$ であるから、$\mG\sim\overline{\mG}/\overline{\mA}$ である。したがって $\overline{\mG}/\overline{\mA}$ は $\mG$ のある因子群と同型である。対応する正規因子は、$\overline{\mG}/\overline{\mA}$ の単位元に対応する元全体であり、それがまさに $\mA$ である。

\emph{補足。} 準同型定理に従って $\overline{\mG}=\mG/\mN$ とおくと、$\mA$ は確かに $\mN$ を含む。$\mA$ から $\overline{\mA}$ を回復できる。すなわち $\overline{\mA}=\mA/\mN$ である。対応 $\mA\leftrightarrow\overline{\mA}$ は一対一である。式 (1) はまた
\[
        (\mG/\mN)/(\mA/\mN)\simeq \mG/\mA
\]
と書ける。

\emph{第二同型定理。} \emph{$\mA$ を部分群、$\mB$ を $\mG$ の正規因子とする。このとき $\mA\cap\mB$ は $\mA$ の正規因子であり、}
\[
        \mA\mB/\mB\simeq \mA/\mA\cap\mB.
\]
\emph{証明。} 準同型 $\mG\sim\mG/\mB$ の下で、とくに $\mA$ の元 $a$ には剰余類 $a\mB$ が対応し、それら全体が群 $\mA\mB/\mB$ をなす。したがって $\mA\sim\mA\mB/\mB$ であり、主張は準同型定理から従う。

群 $\mG$ からそれ自身の上、またはある部分群の上への作用素準同型写像を、$\mG$ の自己準同型という。特に $\mG$ が（作用素付きの）アーベル群で加法的に書かれ、準同型の和と積が上（§1、例6）のように定義されるなら、群からそれ自身への作用素準同型は環をなし、これを自己同型環という。

\emph{単純なアーベル群（作用素付き）の自己同型環は斜体である。}

\emph{証明。} 任意の準同型は、群をそれ自身の上へ写すか、零群の上へ写すかのどちらかである（他の許容部分群がないからである）。すべてを零に写さない準同型は、単純群について上で述べたことにより同型であり、群からそれ自身への同型写像は群をなす。したがって零作用素を除けば環は群となり、ゆえに環そのものは斜体である。

特に $\mG$ が左 $\mo$-加群であり、作用素準同型を右作用素として書くなら、$\mG$ は\emph{左に $\mo$、右に自己同型環をもつ双加群}となる。というのは、新しい作用素 $\Gamma$ を準同型として選んだことは、式
\[
        (a+b)\Gamma=a\Gamma+b\Gamma,
        \qquad ra\cdot\Gamma=r\cdot a\Gamma
\]
によって表されるからである。これらは、先にすでに解釈した残りの式とともに、双加群を特徴づける。

逆に、双加群が与えられていれば、同じ式は、右の乗数領域の各元が左作用素に関する作用素準同型を誘導し、またその逆も成り立つことを表す。

\subsection*{§ 3. 組成列}

$\mG$ の中に許容部分群の有限列
\[
        \mG\supset \mA_1\supset \mA_2\supset \ldots\supset \mA_r=\mE\tag{2}
\]
があり（$\mE$ は単位元 $e$ だけからなる群である）、各 $\mA_i$ が直前の項の正規因子で、かつ列の連続する二項の間に同じ性質をもつ別の部分群を挿入することができないなら、この列を組成列という。

因子群 $\mG/\mA_1,\ \mA_1/\mA_2,\ldots,\ \mA_{r-1}/\mE$ を組成因子という。これらは単純群、すなわち自分自身と単位群以外の許容正規因子をもたない群である。

作用素の中に、とくにすべての内部自己同型を含めるなら、列 (2) はすべて正規因子からなり、これを主列という。すべての自己同型を含めるなら、それは特性列と呼ばれる。§1 の後の例では、これは $\mo$ のイデアルの組成列、あるいは加群または双加群の組成列を意味することになる。

\emph{Jordan--Hölder の定理。} \emph{群 $\mG$ について二つの異なる組成列}
\[
        \mG\supset \mA\supset \mA_2\supset \ldots\supset \mA_r=\mE,
        \qquad
        \mG\supset \mB\supset \mB_2\supset \ldots\supset \mB_s=\mE,
\]
\emph{が存在するなら、それらは同じ長さ $r=s$ をもち、組成因子}
\[
        \mG/\mA,\ \mA/\mA_2,\ldots,\mA_{r-1}/\mE
\]
\emph{は、ある順序で、}
\[
        \mG/\mB,\ \mB/\mB_2,\ldots,\mB_{s-1}/\mE
\]
\emph{と同型である。}
証明は、例えば E. Noether, \emph{Abstrakter Aufbau} usw., Math. Annalen 96 (1926), §10, p. 57 を参照。同時にそこでは次も証明されている。

\emph{$\mG$ に組成列が存在するなら、$\mG$ の任意の正規因子 $\mH$ を通る組成列を引くことができる。}

組成列の存在は有限群について明らかであり、より一般には極大条件と極小条件が成り立つ群について明らかである。

\emph{極大条件}とは、部分群の任意の集合の中に極大部分群、すなわちその集合の他の部分群にこれ以上含まれないものが存在する、ということである。同値に、部分群の任意の昇鎖 $\mA_1\subset\mA_2\subset\mA_3\ldots$ は有限個の項の後で停止する。

\emph{極小条件}とは、部分群の任意の集合の中に極小部分群、すなわちその集合の他の部分群をこれ以上含まないものが存在する、ということである。同値に、任意の降鎖 $\mA_1\supset\mA_2\supset\mA_3\ldots$ は有限個の項の後で停止する。

\emph{部分群として正規因子だけを認める場合}（例えばアーベル群の場合）、\emph{逆に組成列の存在から極大条件と極小条件が従う。}

\emph{証明。} 任意の正規因子は組成列をもち、その長さをその正規因子の長さという。$\mA\subset\mB$ なら、$\mA$ を通って $\mB$ の組成列を引くことができるので、$\mA$ の長さは $\mB$ の長さより小さい。したがって、最短の長さをもつ部分群は同時に極小であり、最長の長さをもつ部分群は同時に極大である。

\subsection*{§ 4. 直積と交わり}

群 $\mG$ は、次の条件を満たすとき、二つの因子の直積 $\mG=\mA\times\mB$ と呼ばれる。
\[
\begin{array}{ll}
1.& \mA,\mB\text{ は }\mG\text{ の正規因子である},\\
2.& \mA\mB=\mG,\\
3.& \mA\cap\mB=\mE.
\end{array}
\]
\emph{この定義は明らかに次と同値である。} $\mG$ の任意の元 $g$ は $g=ab$（$a$ は $\mA$ に、$b$ は $\mB$ に属する）として一意的に表され、さらに $\mA$ の元は $\mB$ の元と可換である：$ab=ba$。

群 $\mG$ が $n$ 個の因子の直積 $\mG=\mA_1\times\ldots\times\mA_n$ であるとは、$\mB_i=\mA_1\ldots\mA_{i-1}\mA_{i+1}\ldots\mA_n$ とおいたとき、各 $i$ について $\mG=\mA_i\times\mB_i$ が直積であることをいう。

\emph{この定義もまた次と同値である。} $\mG$ の任意の $g$ は
\[
        g=a_1a_2\ldots a_n,\qquad a_i\in\mA_i
\]
として一意的に表され、$\mA_i$ の元は $\mA_k$ の元と可換である。

次の事実はしばしば用いられる。
\begin{enumerate}
\item $\mA_1\times\cdots\times\mA_n=\mR$ かつ $\mR\times\mH=\mG$ なら、$\mA_1\times\cdots\times\mA_n\times\mH=\mG$ である。
\item $\mG=\mA_1\times\ldots\times\mA_n=\mC_1\times\ldots\times\mC_n$ かつ $\mC_i\subseteqq\mA_i$ なら、$\mC_i=\mA_i$ である。（これは $\mA_i$ の元を $\mC_j$ によって表すことから従う。）
\item $\mG=\mA\times\mB$ かつ $\mR\supseteq\mA$ なら、$\mR=\mA\times(\mR\cap\mB)$ である。（これは $\mR$ の元を $ab$ の形に表すことから従う。このとき第二因子は $\mR$ にも $\mB$ にも属する。）
\item $\mG=\mA\times\mB$ なら、$\mA\simeq\mG/\mB$ かつ $\mB\simeq\mG/\mA$ である（第二同型定理）。さらに次が従う。
\item $\mG=\mA\times\mB$ で、$\mA$ と $\mB$ がそれぞれ長さ $m$ と $n$ の組成列をもつなら、$\mG$ は長さ $m+n$ の組成列をもつ。
\item $\mA\simeq\overline{\mA}$ かつ $\mB\simeq\overline{\mB}$ なら、$\mA\times\mB\simeq\overline{\mA}\times\overline{\mB}$ である。
\end{enumerate}
群は、$\ne\mE$ と異なる因子の直積として表せないとき、直接分解不能であるという。

明らかに、\emph{極小条件を満たす任意の群は、有限個の直接分解不能群の直積である}。\footnote{W. Krull が可換の場合に、O. Schmidt が一般の場合に示したように（注6参照）、極大条件と極小条件の仮定の下で、この表示は同型を除いて一意である。直接分解不能性の概念を変えずにすべての内部自己同型を作用素領域に加えることができるので、正規因子に対する有限性条件、または主列の存在を要求すれば十分である。}

対象となる群を加法的に書くなら、積および直積の概念は和および直和に移る。記法として、群の和には $(\mA_1,\mA_2,\ldots)$、直和には $\mA_1+\mA_2+\ldots$ を用いる。

「直接交わり」の概念は後では用いないが、直接交わりと直積との関係についての次の定理は、次章以降のイデアル論を和ではなく交わりで定式化することもでき、それによってよく知られた可換理論との関係が回復されることを示している。

群 $\mD$ は、$\mG$ に関して二つの群 $\mR$ と $\mS$ の直接交わりであるとは、
\[
\begin{array}{ll}
1.& \mR,\mS\text{ は }\mG\text{ の正規因子である},\\
2.& \mR\cap\mS=\mD,\\
3.& \mR\mS=\mG
\end{array}
\]
を満たすときをいう。同様に、$\mD$ が $n$ 個の群 $\mS_1,\ldots,\mS_n$ の直接交わりであるとは、$\mR_i=\mS_1\cap\ldots\cap\mS_{i-1}\cap\mS_{i+1}\cap\ldots\cap\mS_n$ とおいたとき、各 $i$ について交わり $\mD=\mR_i\cap\mS_i$ が直接であることをいう。

定義から、$\mD$ は $\mG$ の正規因子でなければならないことが従う。$\mD=\mS_1\cap\ldots\cap\mS_n$ が直接であり、準同型 $\mG\sim\mG/\mD$ の下で群 $\mG,\mS,\mD$ が $\overline{\mG},\overline{\mS},\overline{\mE}$ に移るなら、$\overline{\mE}=\overline{\mS}_1\cap\ldots\cap\overline{\mS}_n$ は直接になる。逆に、そのような任意の表示 $\overline{\mE}=\overline{\mS}_1\cap\ldots\cap\overline{\mS}_n$ から、直接な表示 $\mD=\mS_1\cap\ldots\cap\mS_n$ に戻ることができる。この対応によって、任意の $\mD$ についての直接交わりに関する定理は、特殊な場合 $\mD=\mE$ に還元される。

$\mD=\mE$ かつ因子が二つの場合、直積と直接交わりの条件は一致する。$n$ 個の因子の場合には、直積と交わりとの間に次の一対一対応がある。

\emph{I. $\mG=\mA_1\times\ldots\times\mA_n=\mA_i\times\mB_i$ なら、$\mE=\mB_1\cap\ldots\cap\mB_n=\mB_i\cap\mC_i$ は直接であり、$\mC_i=\mA_i$ である。}

\emph{II. $\mE=\mS_1\cap\ldots\cap\mS_n=\mR_i\cap\mS_i$ が直接なら、$\mG=\mR_1\times\ldots\times\mR_n=\mR_i\times\mT_i$ であり、$\mT_i=\mS_i$ である。}

\emph{I の証明。} $\mB=\mB_1\cap\ldots\cap\mB_n=\mB_i\cap\mC_i$ とする。このとき $\mC_i\supseteqq\mA_i$ である。$\mC_i\subseteqq\mA_i$ を示す。たとえば $c$ が $\mC_1$ に属するなら、$c$ は $\mB_2,\ldots,\mB_n$ に属するので、$\mA$ に関する成分表示から
\[
        c=a_1a_2\ldots a_n=a_1ea_3\ldots a_n=a_1a_2e\ldots a_n=\ldots=a_1\ldots a_{n-1}e
\]
が得られる。すべての $\mA_i$ の積は直積なので、これらの表示は一致する。第一の位置を除く各位置には一度ずつ $e$ が現れるから、$c=a_1$、すなわち $\mC_1\subseteqq\mA_1$ であり、ゆえに $\mC_1=\mA_1$、同様に $\mC_i=\mA_i$ である。したがって $\mB=\mB_i\cap\mC_i=\mB_i\cap\mA_i=\mE$ であり、さらに $\mB_i\mC_i=\mB_i\mA_i=\mG$ であるから、この交わりは直接である。

\emph{II の証明。} $\mH=\mR_1\ldots\mR_n=\mR_i\mT_i$ とする。このとき $\mT_i\subseteqq\mS_i$ である。$\mS_i\subseteqq\mT_i$ を示す。たとえば $s$ が $\mS_1$ に属するとする。$\mG=\mS_i\times\mR_i$ によって
\[
        s=s_1e=s_2r_2=\cdots=s_nr_n
\]
を得る。$t_1=er_2\ldots r_n=t_ir_i$ を作り、$\mR_i$ と $\mS_i$ の元ごとの可換性を考慮すると、
\[
        st_1^{-1}=s_it_i^{-1}
\]
が得られる。したがってこれはすべての $\mS_i$ の元であり、ゆえに $e$ に等しい。よって $\mS_1\subseteqq\mT_1$、すなわち $\mT_1=\mS_1$ であり、同様に $\mT_i=\mS_i$ である。したがって $\mH=\mR_i\mT_i=\mR_i\mS_i=\mG$ かつ $\mR_i\cap\mT_i=\mR_i\cap\mS_i=\mE$ であり、したがって $\mG$ は $\mR_i$ たちの直積である。


% 組版用マクロ
\providecommand{\ma}{\mathfrak{a}}
\providecommand{\mb}{\mathfrak{b}}
\providecommand{\mc}{\mathfrak{c}}
\providecommand{\md}{\mathfrak{d}}
\providecommand{\mf}{\mathfrak{f}}
\providecommand{\mh}{\mathfrak{h}}
\providecommand{\ml}{\mathfrak{l}}
\providecommand{\mn}{\mathfrak{n}}
\providecommand{\mr}{\mathfrak{r}}
\providecommand{\mt}{\mathfrak{t}}
\providecommand{\mx}{\mathfrak{x}}
\providecommand{\mz}{\mathfrak{z}}
\providecommand{\mZ}{\mathfrak{Z}}
\providecommand{\mL}{\mathfrak{L}}
\providecommand{\ann}{\operatorname{Ann}}
\providecommand{\tuple}[1]{(#1)}

\subsection*{§ 5. 完全可約群}

群は、有限個の単純群の直積であるとき、\emph{完全可約}であるという：
\[
        \mG=\mA_1\times\ldots\times\mA_n.
\]
この場合、群
\[
        \mA_1\times\ldots\times\mA_n
        \supset \mA_1\times\ldots\times\mA_{n-1}
        \supset\ldots\supset \mA_1\supset \mE
\]
は組成列をなす。その長さは $n$ であり、組成因子は
\[
        \sim \mA_n,\ \mA_{n-1},\ldots,\mA_1
\]
である。

\emph{$\mG$ が完全可約なら、任意の正規因子は直因子であり、他方の因子は、$\mG$ の与えられた積分解に現れる単純群の積として選ぶことができる。正規因子 $\mH$ 自身も完全可約である。}

\emph{証明。} つぎを得る：
\begin{equation}
        \mG=\mH\mA_1\mA_2\ldots\mA_n.\tag{3}
\end{equation}
$\mH_i=\mH\mA_1\mA_2\ldots\mA_i$ とおく。このとき $\mH_{i+1}=\mH_i\mA_{i+1}$ である。あるいは $\mA_{i+1}\subseteqq\mH_i$ であり、そのとき $\mH_{i+1}=\mH_i$ だから、因子 $\mA_{i+1}$ を (3) から省く。あるいは $\mH_i\cap\mA_{i+1}$ が $\mA_{i+1}$ の真の正規因子であり、したがって $=\mE$ であって、このとき $\mH_i\times\mA_{i+1}$ は直積である。余分な因子を省いたのち、(3) は
\[
        \mG=\mH\times\mA_{i_1}\times\ldots\times\mA_{i_r}
\]
となるので、$\mH$ は直因子である。さらに
\[
        \mH\simeq \mG/\mA_{i_1}\times\ldots\times\mA_{i_r}
        \simeq \mA_{i_1}\times\ldots\times\mA_{j_\mu},
\]
ただし $\mA_{j_1},\ldots,\mA_{j_\mu}$ は残りの $\mA_i$ である。ゆえに $\mH$ は完全可約である。

\emph{系。} \emph{完全可約群を直接分解不能因子に分解する任意の分解は、単純因子への分解であり、したがって組成列を与える。これから因子が同型を除いて一意に定まることが従う。}

\subsection*{§ 6. 体に関する加群。超複素系}

作用素付き完全可約群のほとんど自明な例は、必ずしも可換とは限らない体、すなわち斜体に関する有限基底をもつ加群によって与えられる。

$\mG$ を右 $\mathsf{K}$-加群とし、$K$ の単位元 $e$ は同時に恒等作用素であるとする：$\mG$ の $a$ について $ae=a$。ある $a$ から作られる任意の部分加群 $aK$ は単純であり、すなわちその任意の非零元から作られる加群に等しい。したがって、$\mH$ がいくつかの元から作られる部分加群で、$a$ が $\mH$ に属さない元であるなら、$\mH\cap a\mathsf{K}=\mE$ であり、$\mH+a\mathsf{K}$ は直和である。そこで任意の基底元 $a_1$ から始め、新しい基底元 $a_2,a_3,\ldots$ を順に付け加えることができ、$\mG$ は直和
\[
        \mG=a_1K+a_2K+\cdots+a_nK
\]
として得られる。\emph{したがって $\mG$ は完全可約である。}

数 $n$、すなわち組成列の長さを、$K$ に関する $\mG$ の\emph{階数}という。$\mG$ の元の表示の一意性により（また $e$ が恒等作用素であると仮定したことにより）、基底 $a_1,\ldots,a_n$ は線型独立である。

任意の部分加群 $\mH$ は直和因子である：
\[
        \mG=\mH+\mR=(h_1\mathsf{K}+\ldots+h_s\mathsf{K})+(r_1\mathsf{K}+\ldots+r_{n-s}\mathsf{K}),
\]
すなわち、\emph{$\mH$ の任意の線型独立な基底は、$\mG$ の線型独立な基底へ拡張できる。} しかも $r_i$ は、もとの基底元 $a_i$ の中から選ぶこともできる。

$(a_1,\ldots,a_n)$ と $(c_1,\ldots,c_n)$ を線型独立な基底とする。このとき
\[
        c_k=\sum a_i\pi_{ik},
\]
すなわち行列で書けば、
\[
        (c_1,\ldots,c_n)=(a_1,\ldots,a_n)P,
        \qquad
        P=\begin{pmatrix}
        \pi_{11}&\ldots&\pi_{1n}\\
        \vdots&&\vdots\\
        \pi_{n1}&\ldots&\pi_{nn}
        \end{pmatrix}.
\]
逆に、
\[
        (a_1,\ldots,a_n)=(c_1,\ldots,c_n)Q.
\]
したがって
\[
        (a_1,\ldots,a_n)=(a_1,\ldots,a_n)PQ,
        \qquad
        (c_1,\ldots,c_n)=(c_1,\ldots,c_n)QP,
\]
であり、$a_i$ と $c_i$ が線型独立であるから、
\[
        PQ=QP=E.
\]

同時に環でもある特殊な $\mathsf{K}$-加群が「超複素系」である。

環 $\mo$ が同時に可換体 $\mathsf{K}$ に関する右加群であるとき、それが次を満たせば、$\mathsf{K}$ に関する\emph{超複素系}（文献では「$\mathsf{K}$ 上の代数」とも呼ばれる）という：
\begin{enumerate}
\item 階数が有限である（線型独立な基底を $u_1,\ldots,u_n$ とする）。
\item $ab\cdot\varkappa=a\cdot b\varkappa=a\varkappa\cdot b$。これは、$\mathsf{K}$ が $\mo$ と可換に結合されている、と言い表される。
\item $K$ の単位元 $e$ は同時に恒等作用素である：$\mo$ の $a$ について $ae=a$。
\end{enumerate}

\[
        \Bigl(\sum_i u_i\alpha_i\Bigr)\Bigl(\sum_k u_k\beta_k\Bigr)
        =\sum_{i,k}(u_i u_k)(\alpha_i\beta_k)
\]
であるから、この系は、$\mathsf{K}$ に加えて\emph{乗法表}が与えられると一意に定まる。すなわち各積 $u_i u_k$ が $u_j$ によってどのように表されるか、
\[
        u_i u_k=\sum_j u_j\gamma^j_{ik}
\]
が知られればよい。$\gamma^j_{ik}$ は、結合法則から従うよく知られた関係式を満たさなければならない。

$\mo$ が、しばしば仮定するように、単位元 $e$ をもつなら（すべての $a$ について $ea=ae=a$）、$K$ の元 $x$ は $e\varkappa$ と同一視でき、$\mathsf{K}$ は $\mo$ の部分体とみなせる。そのとき超複素系は、\emph{中心に含まれる体に関して有限階数をもつ環}とも記述できる。

超複素系の例は、有限群の\emph{群環}である。群の元を基底元 $u_i$ とし、群の乗法を乗法として用いる。体 $\mathsf{K}$ は任意である。

\footnotetext{条件 2 は、$K$ の元による乗法が、$\mo$ を $\mo$-左加群としても $\mo$-右加群としても見たときの $\mo$ の準同型である、ということである。3 によって、これらの準同型はさらに同型である。}

\subsection*{§ 7. 行列}

環 $\mo$ の元 $\pi_{ik}$ を成分にもつ正方行列 $(\pi_{ik})$ は、通常の行列積 $\sum \pi_{ij}\varrho_{jk}=\sigma_{ik}$ と加法 $\pi_{ik}+\varrho_{ik}=\sigma_{ik}$ によって環をなす。

体 $\mathsf{K}$ の元を成分とする行列で、右逆元と左逆元の両方が存在するものを\emph{正則}という。

§6 で見たように、\emph{$\mathsf{K}$-加群の二つの線型独立な基底の間の変換行列は正則である。}

\emph{正則性には右逆元で十分である。} 実際、線型独立な基底 $(a_1,\ldots,a_n)$ をもつ右 $K$-加群（不定元 $a_1,\ldots a_n$ における線型形式の加群）を作り、
\[
        (c_1,\ldots,c_n)=(a_1,\ldots,a_n)P
\]
とおく。このとき
\[
        (c_1,\ldots,c_n)P^{-1}=(a_1,\ldots,a_n)PP^{-1}=(a_1,\ldots,a_n),
\]
であるから、$(c_1,\ldots,c_n)$ は再び基底であり、したがって再び線型独立であり、行列 $P$ は正則である。さらに右逆元は左逆元に等しい。

同様に、左逆元で十分である。

\emph{$P$ が左逆元をもつなら、$P$ は左零因子ではない。}
\emph{証明。} $PA=0$ から $A=P^{-1}PA=0$ が従う。

したがって $P$ は右逆元もただ一つしかもたず、それは §6 によって左逆元と一致する。

通常どおり $P_\varkappa$ が行列 $(\pi_{ik}\varkappa)$ を表し、$C_{ik}$ が $ik$ の位置に単位元をもち他のすべての位置に零をもつ行列を表すなら、
\[
        (\pi_{ik})=\sum C_{ik}\pi_{ik};
\]
したがって $\mathsf{K}$ 上の行列は階数 $n^2$ の $K$-加群をなす。$C_{ik}$ は関係式
\[
        C_{ij}C_{jk}=C_{ik},
        \qquad
        C_{ij}C_{\bar j k}=0,\quad j\ne\bar j
\]
を満たす。したがって $\mathsf{K}$ 上の行列は有限 $\mathsf{K}$-加群をなし、$\mathsf{K}$ が可換なら超複素系をなす。

\section*{第II章　非可換イデアル論}

\subsection*{§ 8. 環に対する準同型定理}

$\ma$ が $\mo$ の両側イデアルなら、剰余類領域 $\mo/\ma$ は、容易にわかるように $\mo$-加群であるだけでなく環でもある。$\mo$ の任意の準同型像は再び環であり、両側イデアル $\ma$ による剰余類環 $\mo/\ma$ と同型である。証明は群の場合と同じである。

とくに、$\mo$ が体であれば、$(0)$ と $\mo$ 以外のイデアルは存在しない。したがってこの場合の準同型は、同型であるか、すべての元を零へ写すかのいずれかである。

とくに、環 $\mo$ が右 $\mathsf{K}$-加群であり、$\mathsf{K}$ は $\mo$ と元ごとに可換な環であると仮定する：
\[
        ab\cdot x=a\cdot bx=ax\cdot b
\]
（たとえば、可換体 $\mathsf{K}$ に関する超複素系の場合）。このとき、許容される右、左、または両側イデアルとしては、同時に $\mathsf{K}$-加群であるものだけを考え、環準同型に加えて、$\mathsf{K}$ による乗法に関する作用素準同型も要求する。許容イデアルは $\mo$ と $\mathsf{K}$ に関する双加群になる。（超複素系では、単に「イデアル」といえば、基礎にある係数体 $\mathsf{K}$ に関して許容されるものだけを常に意味する。）この拡張された意味でも、上の準同型定理は有効である。

環準同型の一例は、乗数領域 $\mo$ から絶対乗数領域への移行である（§1、例4）。したがって絶対乗数領域は剰余類環 $\mo/\md$ と同型であり、ここで $\md$ は $\mM$ を消す $\mo$ の元からなる両側イデアルである。

準同型定理が成り立つことから、§2 と同様に、環同型および作用素同型について第一および第二同型定理が従う。

$\mo$ の部分集合のすべての積 $\mA\mB$ は、加群積として理解する。すなわち $\mA\mB$ は、$a$ が $\mA$ に、$b$ が $\mB$ に属するすべての和 $\sum ab$ の全体であり、$a\mB$ は $b$ が $\mB$ に属するすべての和 $\sum ab$ の全体である。$\mB$ が、ある環を作用素領域とする右加群なら、積 $\mA\mB$、それぞれ $a\mB$ もまた右加群である。とくに、右イデアル、係数領域 $\mathsf{K}$ に関して許容されるイデアルなどである。$(\mA,\mB)$ が加群 $\mA$ と $\mB$ の和を表すなら、計算法則は
\[
        \mA\mB\cdot\mC=\mA\cdot\mB\mC,
\]
\[
        \mA(\mB,\mC)=(\mA\mB,\mA\mC),
        \qquad
        (\mA,\mB)\mC=(\mA\mC,\mB\mC)
\]
である。直和は $\mA+\mB$ と書く（§4）。

以下では、右イデアルについて極大条件または極小条件（§3）を満たす環をしばしば考える。とくに超複素系はこれに属する。いま述べた約束により、イデアルとしては $\mathsf{K}$-加群だけが許容され、その階数が有界に保たれるからである（§6）。

\subsection*{§ 9. 冪等元。右イデアルへの直和分解}

$\mo$ を単位元をもつ環とする。

$\mr$ が右イデアルなら、$\mr\mo\subseteqq\mr$ であり、単位元の存在によって実は $\mr\mo=\mr$ である。左イデアルについては対応して $\mo\ml=\ml$ である。

元 $c$ は $c^2=c$ のとき冪等元と呼ばれる。

\emph{$\mo=\mr_1+\ldots+\mr_n$ が右イデアルへの分解であり、かつ}
\[
        e=e_1+\ldots+e_n,
        \qquad e_i\in\mr_i,
\]
\emph{なら、}
\[
        e_i^2=e_i,
        \qquad e_i e_k=0\quad(i\ne k)
        \quad\hbox{（「直交関係」）},
\]
\[
        \mr_i=e_i\mo
\]
\emph{である。}

\emph{証明。} $r$ を $\mr_1$ の元とする。このとき
\[
        r=er=e_1r+\ldots+e_nr,
        \qquad e_ir\in\mr_i,
\]
であるが、一方で
\[
        r=r+0+\ldots+0
\]
でもある。したがって
\[
        e_1r=r,
        \qquad e_ir=0\quad(i\ne1).
\]
第一式から $\mr_1\subseteqq e_1\mo$ が従う。他方 $e_1\mo\subseteqq\mr_1$ であるから、$e_1\mo=\mr_1$。$r=e_1$ と特殊化すると、
\[
        e_1^2=e_1,
        \qquad e_i e_1=0\quad(i\ne1)
\]
を得る。添字 $1$ を $k$ で置き換えても同じである。

\emph{逆に、$e=\sum e_i$、$e_i e_k=0$ $(i\ne k)$、$e_i^2=e_i$ であり、$\mr_i=e_i\mo$ とおくなら、$\mo=\mr_1+\ldots+\mr_n$ である。}

\emph{証明。} $\mo$ の任意の元 $r$ は
\[
        r=er=e_1r+\ldots+e_nr
\]
であるから、
\[
        \mo=(e_1\mo,\ldots,e_n\mo).
\]
しかし表示は一意である。なぜなら
\[
        0=e_1a_1+\ldots+e_na_n
\]
から、$e_i$ を掛けると
\[
        e_i^2a_i=e_i a_i=0
\]
が従うからである。

右側表示
\[
        \mo=\mr_1+\ldots+\mr_n=e_1\mo+\ldots+e_n\mo
\]
から、それに応じて左側表示
\[
        \mo=\ml_1+\ldots+\ml_n=\mo e_1+\ldots+\mo e_n
\]
が従う。

$\mr_i$ が直接分解不能なら、$\ml_i$ もそうである。たとえば $\ml_1$ の分解があれば、
\[
        \ml_1=\mo e_1=\mo e_1'+\mo e_1'',
        \qquad
        \mo=\mo e_1'+\mo e_1''+\mo e_2+\ldots+\mo e_n
\]
を意味する。このことから右分解
\[
        \mo=e_1'\mo+e_1''\mo+e_2\mo+\ldots+e_n\mo
        =\mr_1'+\mr_1''+\mr_2+\ldots+\mr_n,
\]
\[
        \mr_1\cong \mo/(\mr_2+\ldots+\mr_n)=\mr_1'+\mr_2'
\]
が生じ、$\mr_1$ は分解可能になってしまう。

\footnotetext{もう一つの逆もある。$e_i e_k=0$、$e_i^2=e_i$、$c=\sum e_i$ なら、$c$ は環 $e_1\mo+\cdots+e_n\mo$ の左単位元である。和が直和であることは、上と全く同じ方法でわかる。}

\subsection*{§ 10. 両側イデアルへの分解}

再び $\mo$ を単位元をもつ環とし、
\begin{equation}
        \mo=\ma_1+\ldots+\ma_n\tag{1}
\end{equation}
を $\mo$ の両側イデアルへの分解とする。このとき直交関係はイデアルについても成り立つ：
\begin{equation}
        \ma_i\ma_k=0\quad(i\ne k),
        \qquad
        \ma_i^2=\ma_i.                                             \tag{2}
\end{equation}

\emph{証明。} $\mb_1=\ma_2+\ldots+\ma_n$、$\mo=\ma_1+\mb_1$ とする。このとき $\ma_1\mb_1\subseteqq \ma_1\cap\mb_1=0$ であるから $\ma_1\mb_1=0$、したがってなおさら
\[
        \ma_1\ma_i=0\quad(i\ne1);
\]
また
\[
        \ma_1=\ma_1\mo=\ma_1(\ma_1+\mb_1)=(\ma_1^2,\ma_1\mb_1)=\ma_1^2.
\]

逆に、\emph{(1) と (2) から、加群 $\ma_i$ は両側イデアルであることが従う。}
\emph{証明。}
\[
        \mo\ma_i=(\ma_1+\cdots+\ma_n)\ma_i=\ma_i^2=\ma_i,
        \qquad
        \ma_i\mo=\ma_i.
\]

\emph{環 $\ma_i$ における右イデアルは、$\mo$ における右イデアルである。}
\emph{証明。}
\[
        \mr\mo=\mr(\ma_i+\mb_i)=(\mr\ma_i,\mr\mb_i)
 \subseteq(\mr,\ma_i\mb_i)=(\mr,0)=\mr.
\]

\emph{$\mr$ が $\mo$ の右イデアルなら、$\mr$ は $\ma_i$ に含まれる右イデアル $\mr\ma_i$ の直和である。}
\emph{証明。} $\mr=\mr\mo=(\mr\ma_1,\ldots,\mr\ma_n)$ であり、$\mr\ma_i\mo=\mr\ma_i$ だから、$\mr\ma_i$ は $\mr$ に含まれる右イデアルである。$\mr\ma_i\leq\ma_i$ なので、和は直和である：
\[
        \mr=\mr\ma_1+\ldots+\mr\ma_n.
\]

\emph{とくに、$\mr$ が直接分解不能なら、$\mr$ は一つの成分しかもつことができない。したがって $\mr\subseteqq$ はある $\ma_i$ に含まれなければならない。}

これらの定理によって、個々の $\ma_i$ のイデアル論がわかれば、$\mo$ におけるイデアル論は制御される。したがって以下では、ほとんどの場合、両側分解不能な環に制限する。

\emph{異なる $\ma_i$ に含まれる二つの右イデアルは、決して作用素同型ではありえない。}
\emph{証明。} $\ma_1$ に含まれるイデアルは、他のすべての $\ma_k$ によって消されるが、$\ma_1$ によっては消されない。これに対して、$\ma_2$ に含まれるイデアルは $\ma_1$ によって消される。

したがって、環 $\mo$ を両側分解不能イデアル $\ma_1+\ldots+\ma_n$ に分解し、各 $\ma_i$ をさらに片側の直接分解不能な右イデアル $\mr_i'+\mr_i''+\ldots$ に分解できるなら、作用素同型な $\mr$ は同じ $\ma_i$ に属するものの中で探さなければならない。しかし、同じ $\ma_i$ の中にもなお異なる $\mr$ の類、すなわち作用素同型でない類が存在することがある。次の例がそれを示す。

$\mathsf{K}$ を有理数体とし、
\[
        \mo=e_1\mathsf{K}+e_2\mathsf{K}+u\mathsf{K}
\]
を超複素系とする。$e_i$ と $u$ は $\mathsf{K}$ の数と可換であり、乗法表は
\[
\begin{array}{c|ccc}
     & e_1&e_2&u\\ \hline
 e_1 & e_1&0&u\\
 e_2 & 0&e_2&0\\
 u   & 0&u&0
\end{array}
\]
である。単位元は $e=e_1+e_2$ である。$e_i$ は直交関係を満たす。したがって右分解は
\[
        \mo=e_1\mo+e_2\mo=(e_1,u)+(e_2),
\]
左分解は
\[
        \mo=\mo e_1+\mo e_2=(e_1)+(e_2,u)
\]
である。$e_2\mo$ と $\mo e_1$ は見て明らかに分解不能なので、$\mo e_2$ と $e_1\mo$ も分解不能でなければならない。

両側分解は不可能である。というのは、その場合、一方の成分は $e_1\mo$ を含み、他方は $e_2\mo$ を含まなければならないが、前者は $u$ を含み、後者も $ue_2=u$ を含むことになるからである。

しかし $e_1\mo$ と $e_2\mo$ は作用素同型ではない。$\mathsf{K}$ に関する階数が異なるからである。

\emph{$\mo=\ma_1+\ldots+\ma_n$ とし、$\ma_i$ を両側分解不能とする。このとき $\ma_i$ は一意に定まる。}

\emph{証明。} さらに $\mo=\mc_1+\ldots+\mc_n$ と仮定する。このとき
\[
        \ma_i=\ma_i\mo=(\ma_i\mc_1,\ldots,\ma_i\mc_n).
\]
この後の和は、$\ma_i\mc_k\subseteqq\mc_k$ かつ $\subseteqq\ma_i$ であるため直和である。しかし $\ma_i$ は分解不能なので、$\ma_i\mc_k$ は一つ、$\ma_i\mc_{j_i}$ を除いてすべて零でなければならない。したがって
\[
        \ma_i=\ma_i\mc_{j_i}.
\]
同様に、
\[
        \mc_{j_i}=\mo\mc_{j_i}=\ma_1\mc_{j_i}+\cdots+\ma_n\mc_{j_i},
\]
であり、$\ma_i\mc_{j_i}\ne0$ であるから、残りの項はすべて $=0$ である。ゆえに
\[
        \mc_{j_i}=\ma_i\mc_{j_i}=\ma_i,
\]
q.e.d.

\subsection*{§ 11. 中心}

環 $\mo$ の中心 $\mZ$ とは、すべての環元と可換なすべての $z$ の全体である（すべての $a$ について $za=az$）。$\mZ$ は可換環である。

$\mo$ の任意の片側イデアル $\ma$ は、$\mZ$ の中に「収縮イデアル」$\ma\cap\mZ$ をもつ。$\ma$ と $\mZ$ はどちらも $\mZ$-加群であるから、$\ma\cap\mZ$ もそうである。

$\mZ$ の任意のイデアル $\mA$ は拡張イデアルをもつ。すなわち $\mA$ によって $\mo$ の中に生成されるイデアル $\ma=(\mA,\mo\mA)$ である。これは両側である。なぜなら
\[
        \ma\mo=(\mo\mA,\mA)\mo=(\mo\mA\mo,\mA\mo)
        = (\mo^2\mA,\mo\mA)\subseteqq \mo\mA\subseteqq\ma
\]
だからである。$\mo$ が単位元をもつと仮定すれば、$\ma=\mo\mA$ と書ける。この場合、次の定理が成り立つ。

\emph{$\mo$ の任意の両側分解は、中心の分解と一対一に対応する：}
\begin{equation}
        \mo=\ma_1+\ldots+\ma_n,
        \qquad \mA_i=\ma_i\cap\mZ                         \tag{1}
\end{equation}
\emph{から}
\begin{equation}
        \mZ=\mA_1+\ldots+\mA_n,
        \qquad \mo\mA_i=\ma_i,                              \tag{2}
\end{equation}
\emph{が従い、また逆も成り立つ。}

\emph{証明。} $z$ を中心元とし、
\begin{equation}
        z=z_1+\ldots+z_n\tag{3}
\end{equation}
を (1) に従うその分解とする。このとき $z_i$ は中心元である。実際、
\[
        za=z_1a+\ldots+z_na=az=az_1+\ldots+az_n.
\]
和が直和であり、かつ $\ma_i$ が両側であることから、
\[
        z_i a=az_i
\]
が従う。したがって $z_i$ は $\ma_i\cap\mZ=\mA_i$ に属し、(3) は主張された和分解を与える。とくに $e=\sum e_i$ であるから、$e_i$ は中心元である。最後に
\[
        z=ze_1+\ldots+ze_n,
\]
したがって
\[
        \mA_i=\mZ e_i,
        \qquad
        \mo\mA_i=\mo\mZ e_i=\mo e_i=\ma_i.
\]

逆に、$\mZ=\mA_1+\ldots+\mA_n$ が中心の分解であるなら、§9 により
\[
        \mA_i=\mZ e_i,
        \qquad e_i^2=e_i,
        \qquad e_i e_k=0\quad(i\ne k)
\]
である。したがって、直交関係 \hbox{§9} により、
\[
\begin{aligned}
        \mo&=\mo e=\mo e_1+\ldots+\mo e_n
        =\mo\mZ e_1+\ldots+\mo\mZ e_n\\
        &=\mo\mA_1+\ldots+\mo\mA_n
        =\ma_1+\ldots+\ma_n .
\end{aligned}
\]
いま $\ma_i\cap\mZ=\mA_1'$ とすると、主張の前半から
\[
        \mZ=\mA_1'+\ldots+\mA_n'\quad\hbox{が直和である}
\]
ことがわかる。しかし $\mA_i'\supseteqq\mA_i$ であるから、§4, 2 により $\mA_i'=\mA_i$ である。

\emph{系。} $\ma_i$ と $\mA_i$ は、同時に両側直接分解可能または分解不能である。

\subsection*{§ 12. 冪零イデアル}

イデアル $\mc$ は $\mc^\varrho=0$ となるとき\emph{冪零}であるという。

\emph{例。} §10 のイデアル $(u)$。

\emph{$\mo$ に冪零右イデアル $\mr$ が存在するなら、冪零左イデアル、実際には両側イデアルも存在する。}

\emph{証明。} $\mr^\varrho=0$ とする。このとき $\mo\mr$、あるいは $\mo$ が単位元を含まず $\mo\mr$ が零になりうるなら $(\mo\mr,\mr)$ は左イデアルであり、
\[
        (\mo\mr)^\varrho=\mo\mr\mo\mr\ldots\mo\mr\subseteqq\mo\mr\mr\ldots\mr=\mo\mr^\varrho=0.
\]

\emph{二つの冪零右イデアルの和は、再び冪零右イデアルである。}

\emph{証明。} $\mc^\varrho=0$、$\md^\sigma=0$ とする。
\[
        (\mc,\md)^{\varrho+\sigma-1}=(\ldots,\md\mc\ldots\md\ldots\mc\ldots,\ldots)
\]
の各項には、少なくとも $\varrho$ 個の因子 $\mc$、または少なくとも $\sigma$ 個の因子 $\md$ のどちらかが含まれる。前者の場合、
\[
        \md\mc\ldots\md\ldots\mc\ldots\subseteqq\md\mc\mc\mc\ldots\mc=\md\mc^\varrho=0;
\]
後者の場合には $\md$ の因子について同じ議論を行う。したがって
\[
        (\mc,\md)^{\varrho+\sigma-1}=0.
\]

いま $\mo$ で右イデアルについて極大条件が満たされるなら、極大な冪零右イデアルが存在する。それは他のすべての冪零右イデアルを含む。そうでなければ、そのようなイデアルとの和を作ることができるからである。それを $\mc$ と呼ぶ。$\mo\mc$ もまた冪零右イデアルであるから、$\mo\mc\subseteqq\mc$ であり、したがって $\mc$ は左イデアルである。任意の冪零左イデアル $\md$ も同様に $\mc$ に含まれる。なぜなら $(\md,\md\mo)$ は冪零右イデアルだからである。したがって $\md\subseteqq\mc$。\emph{こうして、他のすべての右および左の冪零イデアルを含む極大両側冪零イデアル $\mc$、すなわち「根基」が存在する。}

§10 の例では、$(u)$ が極大冪零イデアルである。

根基が零イデアルであるとき、「根基なしの環」（「半単純環」、「Dedekind 系」）という。根基を法とする剰余類環は常に根基なしの環である。

\emph{$\mZ$ が $\mo$ の中心であり、$\mc$ が $\mo$ の根基であるなら、$\mC=\mc\cap\mZ$ は $\mZ$ の根基である。}

\emph{証明。} $\mC$ が冪零であることは明らかである。もし $\mZ$ の中により大きい冪零イデアル $\overline{\mC}$ が存在するなら、それを $\mc$ に完全には含まれない冪零イデアル $\overline{\mc}$ に拡張できる。なぜなら $((\overline{\mC}\mo)^\varrho=\overline{\mC}^{\varrho}\mo^\varrho=0)$ だからである。

\emph{帰結。} $\mo$ が根基なしの環なら、$\mZ$ もそうである。しかし逆は正しくない。§10 の例が示している。$\mo/\mc$ の中心は $\mZ/\mC$ と同型な環を含むが、これは真の部分環でありうる（§10 の例）。

\subsection*{§ 13. 完全可約環}

§5 に従って、環は有限個の単純右イデアルの直和であるとき、右完全可約であるという。

\emph{単位元をもつ右完全可約環は、$\ne(0)$ と異なる冪零イデアルをもたず、したがって根基をもたない。}

\emph{証明。} 任意の右イデアル $\mr$ は直和因子である（§5）：
\[
        \mo=\mt+\mr=e_1\mo+e_2\mo.
\]
$e_2^2=e_2$ なので、$e_2^\varrho=e_2$ でもある。もし $\mr$ が冪零なら、$e_2^\varrho=0$ であり、したがって $e_2=0$、ゆえに $\mr=0$ である。q.e.d.

次に二つの逆を示す。第一に、\emph{右イデアルについて極小条件を満たす根基なしの環は、右イデアルに関して完全可約である。}

\emph{証明。} $\mt$ を $0$ でない極小右イデアルとする。$\mt$ が冪等元を含むことを示す。

$\mt^2\subseteqq\mt$ かつ $\mt^2\ne0$ であるから、$\mt^2$ は右イデアルなので $\mt^2=\mt$ である。したがって、$a\mt\ne0$ となる $a$ が $\mt$ の中に存在する。このとき $a\mt$ は右イデアルだから、必ず $a\mt=\mt$ である。

$a$ によって消される $\mt$ のすべての $b$（$ab=0$）の全体は右イデアルであり、$\mt$ とは等しくないので $0$ である。したがって $ab=0$ なら $b=0$ である。

$\mt=a\mt$ であるから、$a$ は $a=ac$、$c\ne0$ と表されなければならない。すると
\[
        ac=ac^2,
        \qquad a(c-c^2)=0,
        \qquad c-c^2=0,
        \qquad c^2=c
\]
が従う。つぎに $\mt$ が直和因子であることを示す。集合 $c\mo$ は $\mt$ に含まれる非零右イデアルである。実際、$c^2=c$ は $c\mo$ に属する。したがってこれは $\mt$ である。$\mo$ の任意の元 $r$ は
\[
        r=cr+(r-cr) \qquad \hbox{（片側 Peirce 分解）}
\]
と表される。元 $cr$ はイデアル $\mt$ をなす。元 $r-cr$ は別の右イデアル $\mr$ をなし、$c$ によって消される：$c\mr=0$。したがって
\[
        \mo=(\mt,\mr).
\]
$\mt$ の元は $c$ によって消されないので、この和は直和である：
\begin{equation}
        \mo=\mt+\mr.                                             \tag{1}
\end{equation}

極小条件により $\mo$ を直接分解不能イデアルの直和
\[
        \mo=\mr_1+\ldots+\mr_n
\]
として表すと、$\mr_i$ は単純、すなわち極小でなければならない。というのは、たとえば $\mr_1$ が単純でなく、$\mt$ が $\mr_1$ の中の極小イデアルなら、(1) は
\[
        \mr_1=\mt+\mr\cap\mr_1 \qquad \hbox{§4, 3}
\]
を与え、$\mr_1$ が分解可能になってしまうからである。したがって $\mo$ は完全可約である。

第二に、\emph{極小条件を満たす根基なしの環は単位元をもつ。}

まず左単位元の存在を示すには、次を証明すれば十分である。\emph{右イデアル $\ma=c_1\mo\ne\mo$ が左単位元 $c_1$ をもつなら、より大きい長さをもつ右イデアル $c\mo$ が存在し、それも左単位元 $c$ をもつ。} すでに完全可約性、したがってすべての長さの有界性が示されているので、有限回で $\mo$ 自身の左単位元に達するからである。

したがって $\ma=c_1\mo$、$c_1^2=c_1$ とする。式
\[
        r=c_1r+(r-c_1r)
\]
により、上と同様に Peirce 分解 $\mo=c_1\mo+\mr$ が得られる。上と同様に $\mr$ の冪等元を $c_2$ とする。すなわち $c_2^2=c_2$ であり、$c_1$ は $\mr$ のすべての元を消すので $c_1c_2=0$ である。$e_1=c_1$、$e_2=c_2-c_2c_1$ とおく。このとき
\[
        e_1^2=e_1,
        \quad e_1e_2=0,
        \quad e_2e_1=0,
        \quad e_2c_2=c_2,
        \quad e_2\ne0,
        \quad e_2^2=e_2.
\]
上の注で述べたことにより、$c=e_1+e_2$ は環
\[
        c\mo=e_1\mo+e_2\mo
\]
の左単位元であり、$e_2^2=e_2\ne0$ であるから、この環は右イデアルとして $\ma$ より大きい長さをもつ。

構成された左単位元 $e$ が右単位元でもあることを示すため、左イデアルへの Peirce 分解
\[
        \mo=\mo e+\ml
\]
を行う。$\ml e=0$、$\ml=e\ml$ であるから、$\ml^2=\ml e\ml=0$ である。仮定により冪零左イデアルは存在しえないので、$\ml=0$。したがって $e$ は右単位元でもあり、結局、単位元である。

\emph{定理。右完全可約性と単位元の存在から両側完全可約性が従う：}
\[
        \mo=\md_1+\cdots+\md_s,
        \qquad \md_i \hbox{ は両側単純である。}
\]

前の証明と同じく、任意の極小両側イデアル $\ma$ が直和因子であることを示せば十分である。右イデアルとして、$\ma$ は直和因子である：
\[
        \mo=\ma+\mr=e_1\mo+e_2\mo.
\]
対応する左分解は
\[
        \mo=\mc+\ml=\mo e_1+\mo e_2
\]
である。すると
\[
        \mr\mc\subseteqq \mr\ma\subseteqq \mr\cap\ma=0,
        \qquad
        \ml\ma=\mo e_2e_1\mo=0,
\]
\[
        (\ma\ml)^2=\ma\ml\ma\ml=0,
        \qquad \hbox{したがって } \ma\ml=0;
\]
\[
        \mr=\mr\mo=(\mr\mc,\mr\ml)=\mr\ml,
        \qquad
        \ml=\mo\ml=(\ma\ml,\mr\ml)=\mr\ml.
\]
ゆえに $\mr=\ml$ であり、$\mr$ は両側である。したがって $\ma$ は両側直和因子である。ここから先の証明は、片側完全可約性の場合と同じである。

$\md_i$ は環としても両側単純である。$\md_i$ のすべてのイデアルは $\mo$ のイデアルだからである。これからそれらの構造を調べる。

\subsection*{§ 14. 単位元をもつ両側単純完全可約環}

$\mo$ をそのような環とする：
\[
        \mo=\mr_1+\ldots+\mr_n=e_1\mo+\ldots+e_n\mo,
        \qquad \mr_i\hbox{ は単純である。}
\]
このとき
\[
        \mo=\ml_1+\ldots+\ml_n=\mo e_1+\ldots+\mo e_n,
        \qquad \ml_i\hbox{ は直接分解不能である }\hbox{§9}.
\]
さらに $\ml_i\mr_i$ は非零両側イデアルである（$e_i^2$ が $\ml_i\mr_i$ に属するから）。したがってそれは $=\mo$ である。

$\mA_{ik}=\mr_i\ml_k$ とおく。このとき
\[
        \mo=(\mr_1,\ldots,\mr_n)\cdot(\ml_1,\ldots,\ml_n)
        = (\ldots,\mA_{ij},\ldots)
\]
は直和である。実際、$0=\sum_{i,j}a_{ij}$ であれば、$a_{ij}$ は $\mr_i$ に属し、$\sum\mr_i$ は直和であるから、
\[
        0=\sum_j a_{ij},
\]
さらに $a_{ij}$ は $\ml_j$ に属するので、
\[
        0=a_{ij}
\]
である。

さらに $\mA_{ij}\mr_j=\mr_i\ml_j\mr_j=\mr_i\mo=\mr_i$ であるから、$\mA_{ij}\ne0$。$a_{ij}\ne0$ を $\mA_{ij}$ の元とする。このとき
\[
        a_{ij}\mr_j\subseteqq\mr_i,
        \qquad
        a_{ij}\mr_j\ne0 \hbox{ である。なぜなら } a_{ij}e_j=a_{ij},
\]
ゆえに
\[
        a_{ij}\mr_j=\mr_i.
\]
$\mr_j$ の各 $r_j$ について
\[
        a_{ij}r_j=r_i
\]
とおくと、写像 $r_j\to r_i$ は $\mr_j$ から $\mr_i$ への作用素準同型である。零へ写される元は $\mr_j$ に含まれるイデアルをなし、したがって零イデアルである。ゆえにこの写像は\emph{同型}である。したがって、

\emph{$\mo$ の一つの表示に現れる任意の二つの単純右イデアル $\mr_i$ と $\mr_j$ は作用素同型であり、その同型は $\mA_{ij}$ の元によって媒介される。} 異なる任意の二つの単純右イデアル $\mr,\mr'$ は、少なくとも一つの $\mo$ の表示にともに現れるので、そのような任意の二つのイデアルは同型である。

\emph{$\mr_j$ から $\mr_i$ へのすべての準同型は、$\mA_{ij}$ の元によって媒介される。}

\emph{証明。} $e_j\to a_{ij}$ なら、$e_j^2\to a_{ij}e_j$ であるから、$a_{ij}=a_{ij}e_j$ は $\mr_i\ml_j=\mA_{ij}$ に属する。さらに
\[
        \mr_j=e_j\mr_j\to a_{ij}\mr_j=\mr_i.
\]

\emph{とくに、$\mr_i$ からそれ自身への準同型は $\mA_{ii}$ によって媒介される。}

$\mA_{ii}$ の二元の積は準同型の積に、和は和に対応する（定義は §1, 6 参照）。異なる $a_{ii}$ は異なる自己同型を与える。なぜなら $e_i$ は $a_{ii}e_i=a_{ii}$ に送られるからである。したがって環 $\mA_{ii}$ は $\mr_i$ の自己同型環と同型である。

しかし単純イデアルの自己同型環は斜体である（§2）。したがって $\mA_{ii}$ は斜体である。すべての $\mr_i$ は作用素同型であるから、それらの自己同型環は環同型である。すなわち $\mA_{ii}$ は $\mo$ によって環同型を除いて一意に定まる。現れうる異なる $\mA_{ii}$ は、単純右イデアルの抽象的に定義された自己同型環の、異なる具体的実現にすぎない。

\emph{主定理。単位元をもつ任意の両側単純完全可約環 $\mo$ は、斜体 $K$ 上の次数 $n$ の行列環と同型である。（斜体 $\mathsf{K}$ は、$\mo$ の右イデアルの自己同型斜体と同型である。）}

\emph{証明。行列単位 $c_{ik}$ の構成。} $\Gamma_{11}$ を $\mr_1$ の恒等自己同型とし、$\Gamma_{i1}\mr_1=\mr_i$ である任意の同型を $\Gamma_{i1}$ とし、最後に
\[
        \Gamma_{ik}=\Gamma_{i1}\Gamma_{k1}^{-1}
\]
とおく。このとき一般に
\[
        \Gamma_{ik}\Gamma_{kl}=\Gamma_{il}
\]
である。$\Gamma_{ik}$ が $\mA_{ik}$ の元 $c_{ik}$ によって媒介されるなら、さらに
\[
        c_{ik}=e_i c_{ik}=c_{ik}e_k,
        \qquad
        c_{ik}c_{kl}e_l=c_{il}e_l,
\]
したがって
\[
        c_{ik}c_{kl}=c_{il},
        \qquad
        c_{ij}c_{kl}=c_{ij}e_j e_k c_{kl}=0\quad(j\ne k).
\]
ゆえに $c_{ik}$ は行列単位である（§7）。

\footnotetext{最近 B. L. van der Waerden は、抽象的な自己同型斜体を直接用いる、より簡単な証明を見いだした。この証明は彼の代数学の本（Grundlehren d. Math. Wiss., Berlin: Julius Springer）に掲載予定である。[1929年6月11日。]}

\emph{$\mathsf{K}$ の構成。} 対応
\[
        a_{ii}=c_{i1}a_{11}c_{1i}
\]
は、$\mA_{11}$ の各 $a_{11}$ に、$\mA_{ii}$ の「共役元」$a_{ii}$ を対応させる。

この対応は環同型である。なぜなら、$a$ と $c$ に対応する同型 $\Delta,\Gamma$ は
\[
        \Delta_{ii}=\Gamma_{i1}\Delta_{11}\Gamma_{i1}^{-1}
\]
によって結ばれており、この関係は環同型を表すものとしてよく知られているからである。いま元
\[
        \alpha=a_{11}+\ldots+a_{nn}=a_{11}+c_{21}a_{11}c_{12}+\ldots+c_{n1}a_{11}c_{1n}
\]
を作る。各 $a_{11}$ はちょうど一つの $\alpha$ を定める（和が直和であるため一意である）。$\alpha$ の全体は斜体 $K\sim\mA_{11}$ である。実際、
\[
        \alpha=a_{11}+\ldots+a_{nn},
        \qquad
        \beta=b_{11}+\ldots+b_{nn}
\]
から
\[
        \alpha+\beta=(a_{11}+b_{11})+\ldots+(a_{nn}+b_{nn}),
        \qquad
        \alpha\cdot\beta=a_{11}b_{11}+\ldots+a_{nn}b_{nn}
\]
を得るので、対応 $a_{11}\to\alpha$ は同型である。

さらに
\[
        c_{ik}\alpha=c_{ik}a_{kk}=c_{ik}c_{k1}a_{11}c_{1k}=c_{i1}a_{11}c_{1k},
\]
\[
        \alpha c_{ik}=a_{ii}c_{ik}=c_{i1}a_{11}c_{1i}c_{ik}=c_{i1}a_{11}c_{1k},
\]
であるから、任意の $\alpha$ はすべての $c_{ik}$ と可換である。

最後に、同じ式から
\[
        c_{ik}\mathsf{K}=c_{ik}\mA_{kk}=c_{ik}\mr_k\ml_k=\mr_i\ml_k=\mA_{ik},
\]
したがって
\[
        \mo=\sum c_{ik}\mathsf{K}
\]
が得られる。

\emph{対応。} $\mo$ の各元を $\sum c_{ik}\alpha_{ik}$ の形に書き、その元に行列 $(\alpha_{ik})$ を対応させると、この対応は行列乗法の定義から直ちに同型であることが従う。これで主定理が証明された。

\emph{逆。斜体 $A$ 上の次数 $n$ の行列環は両側単純かつ完全可約である。$A$ は右イデアルの自己同型環と同型であり、$Ae$ は前の定理で構成された $K$ である。}

\emph{証明。} 行列単位（§7）を $c_{ik}$ とする。
\[
        \mr_i=\sum_k c_{ik}A
\]
とおく。このとき明らかに $\mo=\mr_1+\ldots+\mr_n$ である。$\mr_i$ は単純右イデアルである。任意の非零元 $\sum c_{ik}\alpha_k\ne0$ が $\mr_i$ 全体を生成するからである。実際、$\alpha_j\ne0$ なら
\[
        \Bigl(\sum c_{ik}\alpha_k\Bigr)\cdot(c_{jl}\alpha_j^{-1})=c_{il},
\]
であり、$\sum c_{il}\beta_l$ は $\mr_i$ のすべての元を走る。

$\mr_i$ は作用素同型である：$\mr_i=c_{ik}\mr_k$。

したがって $\mo$ は両側分解不能である。ゆえに、§13 の定理とすでに証明した完全可約性により、また単位元が存在するので、$\mo$ は両側単純である。このことは、任意の $\sum\sum c_{ik}\alpha_{ik}\ne0$ が両側イデアル $\mo$ 全体を生成することからもわかる。

さらに
\[
        \mr_i=c_{ii}\mo,
        \qquad
        \ml_i=\mo c_{ii},
\]
\[
        e=\sum c_{ii},
        \qquad c_{ii}=e_i,
\]
\[
        \mA_{ii}=\mr_i\cap\ml_i=c_{ii}\Lambda\simeq\Lambda.
\]
最後に、$a_{11}=c_{11}\lambda$ とおくと、
\[
        \alpha=a_{11}+c_{21}a_{11}c_{12}+\ldots
        =c_{11}\lambda+c_{22}\lambda+\ldots=e\lambda,
\]
が得られ、主張はすべて証明される。

\emph{$\mo$ の中心は $\mathsf{K}$ の中心であり、したがって体である。}

\emph{証明。} $\mZ(\mathsf{K})$ はすべての $x$ と可換なすべての $\zeta$ からなる。しかしこれらはすべての $c_{ik}$ とも可換であり、したがってすべての和 $\sum c_{ik}a_{ik}$ と可換である。ゆえに
\[
        \mZ(\mathsf{K})\subseteqq\mZ(\mo).
\]
逆に、$z$ を $\mZ(\mo)$ の元とする：
\[
        z=\sum_\mu\sum_\nu c_{\mu\nu}\gamma_{\mu\nu}.
\]
このとき
\[
        zc_{ik}=c_{ik}z,
        \qquad
        \sum_\mu c_{\mu k}\gamma_{\mu i}=\sum_\nu c_{i\nu}\gamma_{k\nu},
\]
したがって
\[
        \gamma_{\mu i}=0\ (\mu\ne i),
        \qquad
        c_{ik}\gamma_{ii}=c_{ik}\gamma_{kk},
        \qquad
        \gamma_{ii}=\gamma_{kk}=\gamma;
\]
\[
        z=\sum c_{ii}\gamma=e\gamma=\gamma.
\]
したがって $z$ は $\mathsf{K}$ に属し、すべての $\varkappa$ と可換であるから、$z$ は $\mZ(\mathsf{K})$ に属する。

行列環はもちろん左完全可約でもあり、単位元をもつ任意の右完全可約環は行列環の直和であるから、次が従う：

\emph{単位元をもつ任意の右完全可約環は左完全可約でもあり、逆も成り立つ。}

また、\emph{単位元をもつ完全可約環の中心は、両側単純行列環に対応する可換体の直和である。}

次の定理が残る\srcfnmark{14)}：

\emph{任意の二つの異なる分解 $\mo=\sum c_{ik}\mathsf{K}$、$\mo=\sum c'_{ik}\mathsf{K}'$ は、内部自己同型 $\alpha'\to x^{-1}\alpha'x$、$c'_{ik}= x^{-1}c_{ik}x$ によって互いに移される。}
\srcfntext{14)}{E. Artin, 前掲書, p. 258 を参照。}

\emph{証明。} 
\[
        \mo=\mr_1+\ldots+\mr_n=\mr_1'+\cdots+\mr_n'
\]
とし、$a,b$ を、イデアル $\mr_1,\mr_1'$ の二つの逆向きの同型を媒介する元とする：
\[
        \mr_1=a\mr_1',
        \qquad
        \mr_1'=b\mr_1,
        \qquad
        ab=c_{11},
        \qquad
        ba=c'_{11}.
\]
次をおく：
\[
        x=\sum_i c_{i1}ac'_{1i},
        \qquad
        y=\sum_i c'_{i1}bc_{1i}.
\]
すると
\[
        xy=e,
        \qquad
        y=x^{-1},
        \qquad
        x^{-1}c_{ik}x=c'_{ik},
        \qquad
        x^{-1}\alpha x=\alpha'.
\]


\clearpage


\subsection*{続き：第III章 §§15--19 および第IV章 §§20--26}

\subsection*{第III章．加群論と表現論}

\subsection*{§ 15. 表現と表現加群}

$\mo$ を環、$\mathsf{K}$ を単位元をもつ環とする。後の応用では $\mathsf{K}$ は常に体である。ただし非可換の文脈では、これは非可換体、すなわち斜体として読むべきである。

$\mo$ の $K$ における $n$ 次の表現とは、準同型
\[
        \mo \sim \mathfrak D,
\]
であり、ここで $\mathfrak D$ は $K$ 上の $n$ 次行列からなる環である。

$\mo$ の $\mathsf{K}$ に関する\emph{表現加群}とは、$\mo$ 上の左加群かつ $\mathsf{K}$ 上の右加群である双加群 $\mM$、
\[
        \mo\mM\subseteqq\mM,
        \qquad
        \mM \mathsf{K}\subseteqq\mM,
\]
であって、さらに有限個の一生成元 $\mathsf{K}$-加群の直和
\[
        \mM=x_1\mathsf{K}+\ldots+x_n\mathsf{K}
\]
であり、$\mathsf{K}$ の単位元が恒等作用素であるものをいう。

\emph{すべての表現加群は表現を与える。} $c\in\mo$ とし、
\[
        cx_k=\sum x_i\gamma_{ik},
\]
または行列形式で
\[
        c(x_1,\ldots,x_n)=(x_1,\ldots,x_n)C,
        \qquad C=(\gamma_{ik})
\]
と書く。すると行列 $C$ は $c$ の表現をなす。実際、
\[
        (b+c)x_k=\sum x_i(\beta_{ik}+\gamma_{ik}),
\]
かつ
\[
\begin{aligned}
        bcx_k
        &=b\sum_jx_j\gamma_{jk}
          =\sum b x_j\gamma_{jk}
          =\sum_j\sum_i x_i\beta_{ij}\gamma_{jk}       \\
        &=\sum_i x_i\biggl(\sum_j\beta_{ij}\gamma_{jk}\biggr),
\end{aligned}
\]
すなわち
\[
        bc(x_1,\ldots,x_n)=(x_1,\ldots,x_n)BC
\]
である。

逆に、$\mo$ の任意の表現 $\mathfrak D$ は一つの表現加群に属し、しかもその加群のある特定の基底に属する。すなわち、$\mM$ を $x_1,\ldots,x_n$ の形式的線形形式全体
\[
        y=\sum x_i\alpha_i,
        \qquad \alpha_i\in K
\]
と解する。すると $\mM$ は右 $\mathsf{K}$-加群である。さらに、元 $c$ に行列 $C=(\gamma_{ik})$ が対応するとき、
\[
        cx_k=\sum_i x_i\gamma_{ik},
        \qquad
        c\biggl(\sum_k x_k\alpha_k\biggr)
        =\sum_i x_i\sum_k\gamma_{ik}\alpha_k
\]
と定める。準同型関係
\[
        c+d\sim C+D,
        \qquad
        cd\sim CD
\]
から
\[
        (c+d)x_i=cx_i+dx_i,
        \qquad
        cdx_i=c(dx_i)
\]
が従い、同じことは和 $\sum x_i\alpha_i$ についても成り立つ。残る双加群の法則
\[
        c(y+z)=cy+cz,
        \qquad
        c(y\alpha)=(cy)\alpha
\]
は自明である。したがって $\mM$ は表現加群であり、構成によりちょうど表現 $\mo\sim\mathfrak D$ に属する。

$(y_1,\ldots,y_n)$ を同じ表現加群の別の基底とする：
\[
        (y_1,\ldots,y_n)=(x_1,\ldots,x_n)\cdot P,
        \qquad
        (x_1,\ldots,x_n)=(y_1,\ldots,y_n)\cdot P^{-1}.
\]
$x$-基底に関して $b\sim B$ であれば、$y$-基底に関しては
\[
        b(y_1,\ldots,y_n)
        =b(x_1,\ldots,x_n)P
        =(x_1,\ldots,x_n)BP
        =(y_1,\ldots,y_n)P^{-1}BP
\]
である。表現 $b\sim B$ と $b\sim P^{-1}BP$ を\emph{同値}と呼び、同じ表現類に数える。任意の正則行列 $P$ は $\mM$ の一つの基底を別の基底に移すので、つぎが示された。

\emph{すべての表現加群は、一意に定まる一つの表現類を与える。}\textsuperscript{15a)}

{\footnotesize\noindent
15) 可換体に関する表現加群は、W. Krull, \emph{Theorie und Anwendung der verallgemeinerten Abelschen Gruppen}, Sitzungsberichte der Heidelberger Akademie, 1926, 第1論文、に初めて現れる。\par
15a) 校正時追記（1929年4月14日）。B. L. van der Waerden が私に知らせてくれたように、「線形変換」と「行列」という概念を分ければ、特別な基底の選択に依存しない不変的な結びつきを得ることができる。線形変換とは二つの線形形式加群の準同型であり、行列とは、特定の基底の選択に対するこの準同型の表現である。I. すべての表現加群は、乗法子領域 $\mo$ に、加群自身の中への線形変換からなる一意な準同型系を割り当てる。実際、§2 の終わりにより、双加群 $\mM$ の左乗法子は、右作用素（ここでは $K$ による乗法）に関する $\mM$ から自身への作用素準同型を生成し、この対応は準同型である。II. 線形形式加群の自身の中への線形変換の系で、$\mo$ に準同型であるものはすべて、表現加群を与える。\par}

二つの表現加群が作用素同型であれば、それらが同じ表現類に対応することは明らかである。逆に、二つの表現加群 $(x_1,\ldots,x_n)$ と $(\bar x_1,\ldots,\bar x_n)$ が同じ表現を生成するなら、対応
\[
        x_i\to \bar x_i,
        \qquad
        \sum_i x_i\alpha_i\to \sum_i \bar x_i\alpha_i
\]
は同型を与える。乗法規則は行列 $C$ によって完全に決定されるからである。

特殊な場合として、$\mM$ の二つの基底が同じ表現を生成するなら、それらは $\mM$ から自身への作用素同型によって互いに移される。

同値な言い方をすれば、すべての $C$ について一つの固定された $P$ により $C=PCP^{-1}$ であるなら、対応
\[
        y_i\mapsto x_i,
        \qquad
        y_i=\sum_k x_k\pi_{ik}
\]
は $\mM$ から自身への同型である。逆に、双加群 $\mM$ の自身への任意の同型 $y_i\mapsto x_i$ は、すべての $C$ と可換な行列 $P$ によって媒介される。

表現の概念は、可換な乗法子領域 $\mathsf{P}$ とまだ可換的に結びついている環 (§8) にも拡張できる。その場合、表現環 $\mathsf{K}$ としては、その中心に $\mathsf{P}$ を含むものだけを取り、表現には単に環準同型 $\mo\sim\mathfrak D$ であることだけではなく、作用素準同型であることを要求する。すなわち $a\sim A$ から
\[
        a\rho\sim A\rho\qquad(\rho\in P)
\]
が従うべきである。表現加群に対してこの要求は、追加の規則
\[
        am\cdot\rho=a(m\rho)=a\rho\cdot m
\]
を意味する。このとき加群と環は $P$ と\emph{可換的に結びついている}という。

\subsection*{§ 16. 可約表現}

$\mU$ が表現加群 $\mM$ の部分加群であり、$\mU$ の基底 $z_1,\ldots,z_t$ に $y_1,\ldots,y_r$ を付け加えて $\mM$ の基底を選べるとする。すなわち
\[
        \mM=y_1K+\cdots+y_rK+z_1K+\cdots+z_tK,        \tag{2a}
\]
である。このとき表現は
\[
        C=\begin{pmatrix}R&0\\ S&T\end{pmatrix}
\]
の形をもつ。ここで行列 $T$ だけで、$\mU$ によって媒介される $t$ 次の表現をなし、行列 $R$ は $\mM/\mU$ によって媒介される $r$ 次の表現をなす。

実際、
\[
        (cy_1,\ldots,cy_r,cz_1,\ldots,cz_t)
        =(y_1,\ldots,y_r,z_1,\ldots,z_t)\cdot C.
\]
$cz_i$ は $\mU$ の元として $z_i$ だけで表されるので、$C$ の右上ブロックは零である。$C$ の他の部分を上に示した順に $R,S,T$ と名づけるなら、とくに
\[
        (cz_1,\ldots,cz_t)=(z_1,\ldots,z_t)\cdot T
\]
である。したがって $T$ は $\mU$ によって媒介される表現をなす。さらに
\[
        (cy_1,\ldots,cy_r)\equiv (y_1,\ldots,y_r)R \pmod{\mU},
\]
一方で $y_i$ は $\mU$ を法として線形独立な基底をなす。よって $R$ は $\mM/\mU$ によって生成される表現をなす。

逆に、
\[
        C=\begin{pmatrix}R&0\\ S&T\end{pmatrix}
\]
という「可約」表現が与えられ、$R$ と $T$ が正方行列であるなら、対応する表現加群において、各 $c$ と最後の $t$ 個の基底元 $z_1,\ldots,z_t$ との積はそれらの元だけで表される。したがって $\mU=(z_1,\ldots,z_t)$ は部分加群である。

系。 
\[
        \mM\supset \mU_1\supset \mU_2\supset\cdots\supset \mU_e=0
\]
を $\mM$ の組成列とし、各 $\mU_i$ は直前のものの $K$-加群としての直和因子であると仮定する。このとき基底
\[
        z_{11},\ldots,z_{1r_1},\quad
        z_{21},\ldots,z_{2r_2},\quad\ldots\quad
        z_{e1},\ldots,z_{er_e}
\]
を、$i>\nu$ の $z_{ik}$ が $\mU_\nu$ の基底をなすように選ぶことができる。$\mM$ によって媒介される表現はそのとき
\[
 C=\begin{pmatrix}
 R_{11}&0&\cdots&0\\
 R_{12}&R_{22}&\cdots&0\\
 \vdots&\vdots&\ddots&\vdots\\
 R_{1e}&R_{2e}&\cdots&R_{ee}
 \end{pmatrix},                                      \tag{2}
\]
の形をもち、$R_{\nu\nu}$ は組成因子 $\mU_{\nu-1}/\mU_\nu$ によって媒介される表現をなす。

個々の表現 $R_{\nu\nu}$ は既約である。なぜなら組成因子 $\mU_{\nu-1}/\mU_\nu$ が単純だからである。

$K$ が体であると仮定し、したがって任意の部分加群が直和因子であるなら、逆に、形 (2) で可能な限り簡約された任意の表現は組成列を与える。Jordan--Hölder の定理により、組成因子は作用素同型を除いて一意に定まる。したがって $R_{\nu\nu}$ は同値な表現および順序を除いて一意に定まる。\textsuperscript{16)}

{\footnotesize\noindent 16) 言い換えれば、$\mU$ が $K$-加群として直和因子であるということである。$K$ が体なら、この仮定は常に満たされる。\par}

$\mM=\mA+\mB$ が二つの表現加群 $\mA=(y_1,\ldots,y_r)$ と $\mB=(z_1,\ldots,z_s)$ の直和であるなら、基底 $(y_1,\ldots,y_r,z_1,\ldots,z_s)$ によって媒介される表現は
\[
        C=\begin{pmatrix}R&0\\0&S\end{pmatrix}
\]
の形をもつ。ここで $R$ は $\mA$ によって媒介される元 $c$ の表現、$S$ は $\mB$ によって媒介される表現を表す。逆も同様である。まず $\mM$ を直接分解不能な成分の直和として書き、それらの中に上のような組成列を引けば、適切な基底に対して表現は、外側の大きな箱が直接分解不能な成分に対応し、対角の小ブロックが組成因子に対応するブロック形をもつ。

\[
\begin{array}{@{}r@{\quad}c@{}}
(3)&
C=\left(
\begin{array}{ccc}
\boxed{\begin{matrix}
R_{11}&0\\
\vdots&\ddots\\
R_{ik}&R_{rr}
\end{matrix}} & & 0\\[1.4em]
& \boxed{\begin{matrix}
S_{11}&0\\
\vdots&\ddots\\
S_{ik}&S_{ss}
\end{matrix}} & \\[1.4em]
0 & & \boxed{\begin{matrix}
T_{11}&0\\
\vdots&\ddots\\
T_{ik}&T_{tt}
\end{matrix}}
\end{array}
\right).
\end{array}
\]

ふたたび $\mathsf{K}$ が体であるなら、このように可能な限り簡約された任意の表現は、逆に、直接分解不能な和因子とその内部の組成因子による加群の表現を与える。Krull と Otto Schmidt の定理により、直接分解不能な表現成分の類は順序を除いて一意に定まる。

加群 $\mM$ が完全可約であるなら、ブロック形のすべての箱は一つの既約表現だけからなり、この表現を完全可約と呼ぶ。

\subsection*{§ 17. 単位元をもつ環に対する表現加群の直和分解}

$\mo$ を単位元をもつ環とし、$\mM$ を $\mo$-加群とする。このとき
\[
        \mM=\mM_e+\mM_0
\]
であり、単位元は $\mM_e$ 上で恒等作用素、$\mM_0$ 上で零作用素である。$\mM$ がさらに右 $K$-加群であるなら、$\mM_e$ と $\mM_0$ も右 $\mathsf{K}$-加群である。

証明。任意の $m\in\mM$ は
\[
        m=em+(m-em)
\]
と書ける。$em$ 全体を $\mM_e$、$m-em$ 全体を $\mM_0$ とする。これらは明らかに加法群であり、$\mM$ が右 $K$-加群であれば右 $K$-加群でもある。さらに
\[
        r(em)=erm
\]
なので $\mM_e$ は $\mo$-加群でもあり、
\[
        r(m-em)=rm-rem=rm-rm
\]
なので $\mM_0$ も $\mo$-加群である。$em$ については $e$ が恒等作用素であり、$(m-em)$ については零作用素である。したがって $\mM_e$ と $\mM_0$ の両方に属するものは零だけで、和 $\mM_e+\mM_0$ は直和である。

表現加群については、$\mM_0$ は各元に零を割り当てる自明表現を与える。和因子 $\mM_0$ を分離すれば、$\mM_e$ によって媒介される表現が残り、そこで単位元は単位行列に対応する。したがって以後はそのような表現に制限してよいし、そうする。

\[
        \mo=\ma_1+\ldots+\ma_s
\]
を両側イデアルの直和とし、
\[
        e=e_1+\ldots+e_s
\]
を対応する $e$ の分解とする。また $\mM$ を、$e$ が恒等作用素である $\mo$-加群とする。このとき
\[
        \mM=\ma_1\mM+\ldots+\ma_s\mM
\]
は直和である。

実際、明らかに $\mM=\mo\mM=(\ma_1\mM,\ldots,\ma_s\mM)$ である。$\mb_i=\ma_1+\cdots+\widehat{\ma_i}+\cdots+\ma_s$ とおく。すると $\mb_i$ は対応して両側イデアルであり、
\[
        \mM=(\ma_i\mM,\mb_i\mM)
\]
で、この和は直和である。なぜなら $e_i$ は $\ma_i\mM$ 上で恒等作用素、$\mb_i\mM$ 上で零作用素だからである。

この定理に基づき、通常は両側分解不能な環の表現に制限する。

\subsection*{§ 18. 完全可約環の加群論と表現論}

$\mo$ を完全可約かつ両側単純とし、$\mM$ を $\mo$ の単位元が恒等作用素である有限 $\mo$-加群とする。このとき $\mM$ は完全可約であり、その単純構成要素は単純左イデアル $\ml_i$ と作用素同型である。

証明。
\[
        \mo=\ml_1+\ldots+\ml_n,
        \qquad
        \mM=(\mo m_1,\ldots,\mo m_k)
\]
とする。したがって
\[
        \mM=(\ldots,\ml_i m_k,\ldots).                         \tag{4}
\]
零でない $\ml_i m_k$ は $\ml_i$ と同型である。写像 $a\mapsto am_k$ が作用素同型だからである。したがって $\ml_i m_k$ は単純である。(4) から、すでに先行する和に含まれている $\ml_i m_k$ を省けば、和は直和になる。

したがって、さらに $\mM$ が直接分解不能であるなら、$\mM$ は単純であり、ある $\ml_i$ と同型である。

この単純な $\mM$ の自己同型斜体を $\Delta$、$\ml_i$ のそれを $\Lambda$ とする。$\mM$ と $\ml_i$ の作用素同型により、環同型 $\Delta\simeq\Lambda$ が得られる。§1 により、$\mM$ と $\ml$ はそれぞれ $\Delta$ と $\Lambda$ を右側の領域とする双加群とみなすことができ、したがって表現加群でもある。それらによって媒介される表現は、環同型 $\Delta\simeq\Lambda$ の下で互いに移る。ゆえにこの表現類を研究するには、$\ml_1$ によって、その自己同型斜体において媒介される表現に制限してよい。

具体的に、$\ml_1$ について行列単位からなる基底
\[
        \ml_1=(c_{11},\ldots,c_{n1})
\]
を基礎に置き、自己同型斜体 $\Lambda$ の実現として、以前（§14）$\mathsf{K}$ と記した $\mo$ の部分斜体を取る（その際、先の考察では右イデアルを左イデアルで置き換え、それに応じて自己同型を右に書かなければならない）。
\[
        a=\sum c_{ik}\alpha_{ik}
\]
と表せば、
\[
        a\longmapsto(\alpha_{ik})
\]
が $\ml_1$ によって媒介される $K$ における表現である。実際、
\[
\begin{aligned}
        a\cdot c_{k1}
        &=\left(\sum\sum c_{ij}\alpha_{ij}\right)\cdot c_{k1}\\
        &=\sum c_{ik}\alpha_{ik}c_{k1}\\
        &=\sum c_{i1}\alpha_{ik},
\end{aligned}
\]
すなわち
\[
        a(c_{11},\ldots,c_{n1})
        =(c_{11},\ldots,c_{n1})(\alpha_{ik})
\]
である。

\[
\left\{
\begin{aligned}
&\text{$\Gamma$ を $K$ の部分斜体で、$K$ が $\Gamma$ 上有限階数 $t$ をもつものとする：}\\
&\qquad K=z_1\Gamma+\cdots+z_t\Gamma.\\
&\text{このとき $K$ は $\Gamma$ に関する自分自身の表現加群となる。$K$ の元 $\beta$ は、}\\
&\text{次のように得られる $\Gamma$ 上の行列 $B$ によって表される：}\\
&\qquad \beta z_j=\sum z_i\beta_{ij};\quad B=(\beta_{ij}).
\end{aligned}
\right.
\]
いま $\ml_1$ を部分斜体 $\Gamma$ に関する表現加群とみなすと、次を得る。

\emph{$\ml_1$ によって媒介される $\mo$ の $\Gamma$ における表現は、上に与えた $\mo$ の $\mathsf{K}$ における表現
\[
        a\to(\alpha_{ik})
\]
において、各元 $\alpha_{ik}$ を、$K$ の $\Gamma$ における表現でそれに対応する行列 $A_{ik}$ に置き換えることによって得られる。}

証明。
\[
        \ml_1=\sum c_{i1}\mathsf{K}=\sum\sum c_{i1}\varkappa_j\Gamma.
\]
基底 $(\ldots,c_{i1}\varkappa_j,\ldots)$ によって、
\[
c_{ik}\to
\begin{array}{c@{\quad}c}
 & \begin{array}{ccc} \boxed{\scriptstyle 1} & \boxed{\scriptstyle k} & \boxed{\scriptstyle n} \end{array}\\[-.2em]
\begin{array}{c}\boxed{\scriptstyle 1}\\ \boxed{\scriptstyle i}\\ \boxed{\scriptstyle n}\end{array}
&
\begin{pmatrix}
0&\ldots&0&\ldots&0\\
\vdots&&\vdots&&\vdots\\
0&\ldots&E_t&\ldots&0\\
\vdots&&\vdots&&\vdots\\
0&\ldots&0&\ldots&0
\end{pmatrix}
\end{array},
\]
となる。ここで $E_t$ は単位行列、$0$ は $t$ 次零行列であり、
\[
        \alpha\longmapsto
        \begin{pmatrix}
        A& &0\\
         &\ddots&\\
        0& &A
        \end{pmatrix},
\]
したがって
\[
        a=\sum c_{ik}\alpha_{ik}\to
        \begin{pmatrix}
        A_{11}&\ldots&A_{1n}\\
        \vdots&&\vdots\\
        A_{n1}&\ldots&A_{nn}
        \end{pmatrix}.
\]
任意の $\mM$ への移行は、表現について再び環同型を意味する。$\mo$ が両側単純な環の直和であるなら、これらの環は順次同型的に表現され、残りは零によって表現される。

\textbf{注意。} ここで研究した、左イデアルの自己同型斜体およびその有限次数部分斜体における表現は、同型を除けば、単純 $\mo$-加群（または左イデアル）によって媒介され得る表現のすべてである。実際、先の注意（§2）によれば、$\mo$-加群 $\mM$ を任意の斜体 $\mathsf{K}$ に関する双加群とみなすとき、$\mathsf{K}$ の元は加群準同型を生成する。したがって斜体 $\mathsf{K}$ は、この $\mo$-加群の自己同型斜体の部分斜体と準同型に対応していなければならない。単位元が恒等作用素であるから、この準同型は同型である。また自己同型斜体全体は、その部分斜体上有限次数でなければならない。さもなければ表現加群も無限階数となり、これは不可能だからである。

最後に、ここで調べた既約表現はまだすべてではなく、単純 $\mo$-加群によって媒介されるものだけであることを注意しておく。一般にはほかにも表現加群がある（例えば、斜体 $\Omega$ を、部分斜体 $\Sigma$ を左側領域、自身を右側領域とする双加群とみなしたもの）。それらは $\mo$-加群としては可約であるが、双加群としては既約であり、そのためやはり既約表現を与える。

\subsection*{§ 19. 加群および表現加群における単純組成因子}

$\mo$ を左イデアルについて極大条件と極小条件を満たす環、$\mc$ をその極大冪零イデアルとする。このとき次が成り立つ。

\emph{任意の単純 $\mo$-加群 $\mM$ は、$\mo$ によって零化されるか、または根基をもたない環 $\mo/\mc$ の単純左イデアル $\ml$ と同型である。後者の場合、この加群は $\mo/\mc$ の両側イデアルのうち $\ml$ を含まないものすべてによって零化され、絶対乗数環（§1）は $\ml$ を含む両側単純環と同型である。}

\textbf{証明。} $\mc^{\varrho}=0$ とする。$\mc\mM=0$ でなければならない。さもなければ $\mc\mM=\mM$ であり、したがって
\[
        \mM=\mc\mM=\mc^2\mM=\ldots=0
\]
となるからである。ゆえに $\mM$ は $\mo/\mc$-加群ともみなせる（$\mc$ を法とする一つの剰余類のすべての元は、$\mM$ の元に掛けたとき同じ結果を与える）。根基をもたない環として $\mo/\mc$ は単位元をもち、したがって $\mM$ は、$\mo/\mc$ によって零化される加群と、単位元が恒等作用素となる加群との直和である（§17）。$\mM$ は単純なので、この二つの和因子の一方だけが現れ得る。単位元が恒等作用素である場合、§18 と同じく単純左イデアルとの同型が従う。実際、$m\ne0$ を $\mM$ の元とすれば $(\mo/\mc)m\ne0$ であり、したがって $\mo/\mc$ の少なくとも一つの単純左イデアル $\ml$ について $\ml m\ne0$ となり、ゆえに $\ml m$ は $\mM$ を尽くさなければならない。そのほかの主張はこのような左イデアルについて明らかであり、それと同型なすべての加群へただちに移る。

\textbf{帰結。} \emph{$\mM$ が組成列をもつ $\mo$-加群なら、その単純組成因子は $\mo$ によって零化されるか、または $\mo/\mc$ の単純左イデアルと同型である。}

$\mathsf{P}$ が $\mo$ と可換的に結びついた環である場合も、すべての結果は保たれる。その場合には、加群の代わりに、右側の $\mathsf{P}$ と左側の $\mo$ に関する双加群であって、条件
\[
        am\cdot\varrho=a\cdot m\varrho=a\varrho\cdot m\qquad(\S~15)
\]
を満たすものを考えればよい。さらに $\mathsf{P}$ の単位元も恒等作用素であるとする。部分加群としては、常に $\mathsf{P}$ による乗法を許すものだけを考える（イデアルの場合と同様である）。われわれの証明で構成された部分加群 $\mc\mM,\mc^2\mM,\ldots$ はこの条件を満たす。

とくに $\mathsf{P}$ が体で、$\mM$ が有限階数の $\mathsf{P}$-加群なら、$\mM$ は同時に表現加群である。$\mM$ によって媒介される表現は環準同型であるだけでなく、$\mathsf{P}$ に関して作用素準同型でもある（§15）。この性質をもつ $\mathsf{P}$ におけるすべての表現は、このような加群によって与えられる。既約表現は単純加群に属し、逆も成り立つ。したがってこの節の定理は、ただちに $\mo$ の $\mathsf{P}$ における表現についての定理として解釈できる。\emph{すべての既約表現（零表現を除く）は $\mo/\mc$ の単純左イデアルによって媒介される。したがって、任意の表現を §16 に従って組成列により簡約すると、その主対角線上には（零を除き）$\mo/\mc$ の単純左イデアルによって媒介される表現と同値な表現だけが現れる。}

\textbf{注意。} $\mo$ と $\mathsf{P}$ について以上の仮定の下では、$\mo$ に有限性条件を仮定する必要はない。すべての表現は同時に絶対乗数環の同型表現であり、この絶対乗数環は有限次数行列環の同型逆像として、それ自身が有限階数の $\mathsf{P}$-加群となる。また定義により許されるイデアルとして $\mathsf{P}$-加群だけを考えるので、この絶対乗数環について極大条件と極小条件が満たされる。

$\mathsf{P}$ を可換体と仮定すれば、この絶対乗数環は $\mathsf{P}$ に関する超複素系となる。主対角線上に零だけでないものが現れるなら、この仮定は常に満たされる。実際、そのとき §18 の注意により、$\mathsf{P}$ は $\ml$ の自己同型斜体の部分斜体と同型であり、$\mo/\mc$ の部分斜体 $K$ による実現（§14）では、可換的結合のために $K$ の中心に含まれなければならない。
\subsection*{第IV章．群と超複素系の表現}

\subsection*{§ 20. 超複素系の組み入れ}

ここから、可換体 $P$ に関する超複素系 $\mo$ を考える。約束により、$\mo$ のイデアルとして数えるのは $\mathsf{P}$-加群だけである。したがって左イデアルおよび右イデアルについて極大条件と極小条件が満たされ、第II章の環論全体が適用される。

以下で超複素系 $\mo$ の表現というときは、常に体 $\mathsf{P}$ における表現、より正確には、$a\sim A$ から
\[
        a\rho\sim A\rho\qquad(\rho\in P)
\]
が従うものを意味する。すなわち $P$ に関する加群準同型性である。表現加群ではこれは
\[
        a\rho\cdot m=a(m\rho)
\]
と読まれる。したがって §19 の結果が適用される。すべての既約表現は
\[
        \mo_0=\mo/\mc
\]
の単純左イデアルによって媒介される。ただし $\mc$ は根基である。したがって、分解
\[
        \mo_0=\ma_1+\cdots+\ma_s
\]
に現れる両側単純和因子 $\ma$ の数だけ、互いに非同値な既約表現がある。

そのような左イデアルによって定まる表現を明示的に決定するには、まず $\mo$ の元から $\mo_0$ における剰余類へ移る。$P$ の元による乗法は $\mo_0$ のすべてのイデアルについて同型である。したがって $P$ は、各単純左イデアルの自己同型斜体 $K_\nu$ の部分体 $e^{(\nu)}\mathsf{P}$ に同型である。$\mathsf{P}$ 上で $\ml_\nu$ によって媒介される表現は、まずこの部分体 $e^{(\nu)}\mathsf{P}$ 上の表現を考え、ついで同型 $e^{(\nu)}\mathsf{P}\simeq \mathsf{P}$ によってそれを $\mathsf{P}$ に移すことで得られる。体 $\mathsf{K}_\nu$ は $\mathsf{P}$ 上有限階数をもつので、これは $\ml_\nu$ の自己同型斜体の有限次数部分体における表現である。

\[
        c=c_1+\ldots+c_s
\]
を $\mo_0$ の元の両側成分への分解とする。$\ml_\nu$ によって媒介される表現では、$c_\nu$ の行列だけを考える。ほかの $c_i$ は $\ml_\nu$ を零化し、零行列として表現されるからである。$c_\nu$ を行列単位で
\[
        c_\nu=\sum c_{ik}^{(\nu)}\alpha_{ik}^{(\nu)}
\]
と書くなら、各 $\alpha_{ik}^{(\nu)}$ は、$P$ 上の $\mathsf{K}_\nu$ の表現においてそれを表す行列で置き換えられるべきである。

$\mathsf{K}_\nu$ は $\mathsf{P}$ 上有限階数をもつので、$\mathsf{K}_\nu$ の任意の元 $x$ は、$e^{(\nu)}P$ に係数をもつ方程式 $f(x)=0$ を満たす。なぜならその冪が線形従属になるからである。$P$ が代数閉であれば、この方程式は一次因子に分解する。$K_\nu$ は体であるから、$\varkappa$ はすでにその一次因子の一つの根であり、したがって $e^{(\nu)}\mathsf{P}$ に属する。ゆえに $\mathsf{K}_\nu=e^{(\nu)}\mathsf{P}$ である。この場合には行列 $(\alpha_{ik}^{(\nu)})$ そのものが既約表現である。

これを §19 の終わりと合わせると Burnside の定理を得る。

\emph{$\mathsf{P}$ が代数閉であり、環 $\mo$ と可換的に結びついているとする。このとき、$P$ 上の任意の $n$ 次既約表現は、ちょうど $n^2$ 個の線形独立な行列を含む。}

ここから直ちに通常の形が従う。$P$ 上の $n$ 次行列の系で乗法について閉じているものは、すべての行列の成分 $x_{ik}$ が満たす $P$ 係数の非自明な斉次線形関係
\[
        \sum \alpha_{ik}x_{ik}=0
\]
が存在しないとき、かつそのときに限り既約である。実際、すべての線形結合を付け加えると、環かつ $P$-加群が得られ、この環をそれ自身の表現とみなすことができる。

絶対乗法子領域は、単位元をもつ両側単純完全可約環に同型であり、したがってその左イデアルの自己同型斜体上の完全行列環に同型である。$P$ の代数閉性により、その自己同型斜体は $P$ に等しくなる。既約表現は単純左イデアルによって生成される。このイデアルの階数を $m$、表現の次数を $n$ とすると、$m=n$ であり、行列環にはちょうど $n^2$ 個の線形独立な元がある。

同様に一般化 Burnside 定理が従う。同じ仮定の下で、$\mathfrak D$ が次数
\[
        n_1,n_2,\ldots,n_s
\]
の互いに非同値な成分に分解された完全可約表現であるなら、$\mathfrak D$ は $P$ に関して階数
\[
        n_1^2+n_2^2+\cdots+n_s^2
\]
をもつ。

$\mo$ が任意の体上の根基をもたない超複素系であるなら、$\mo$ は完全可約で単位元をもつ。§§17, 18 から、$\mo$ の任意の表現は完全可約である。

逆に、$\mathsf{P}$ と可換的に結びついた環 $\mo$ の任意の完全可約表現は、絶対乗法子領域として根基をもたない超複素系をもつ。絶対乗法子領域は、$P$ 上の行列からなる環かつ $P$-加群に同型であり、したがって超複素系である。その根基の元は任意の既約構成成分で零として表現され、したがって全表現でも零として表現される。ゆえに根基は零である。

超複素系 $\mo$ の\emph{正則表現}とは、単位イデアル $\mo$ 自身によって媒介される表現である。群環では、とくに群元を基底として選ぶ。$\mo/\mc$ のすべての左イデアルは $\mo$ の組成因子として現れるので、§16 のように組成列によって正則表現を簡約すれば、すべての既約表現はすでに正則表現の対角上に現れる。

表現は、基底元 $u_1,\ldots,u_h$ を表す行列が分かれば知られる。系が群環であり、したがって $u_i$ が群元であるなら、群の任意の準同型的な行列表現は群環の表現でもある。したがって群を表現する問題は、超複素系を表現する問題の特殊な場合である。

\subsection*{§ 21. 基礎体の拡大。中心の表現}

\[
        \mo=a_1P+\cdots+a_hP
\]
を超複素系とし、$\Omega$ を $\mathsf{P}$ の拡大体とする。基底元に対する古い乗法規則を用いて
\[
        \mo\Omega=a_1\Omega+\ldots+a_h\Omega
\]
を作ることができる。\textsuperscript{18a)} このとき $\mo\Omega$ は $\Omega$ に関してふたたび超複素的である。

$\mo$ の任意の表現は $\mo\Omega$ の表現を与える。すべての元の表現は、基底元 $a_i$ の表現が分かれば知られるからである。\textsuperscript{19)}

{\footnotesize\noindent 18a) より正確には、古い乗法規則をもつ新しい記号 $a_i$ を導入する。すると $a_1\Omega+\cdots+a_h\Omega$ は $\mo$ に同型な部分環を含むので、その後で新しい記号を古い $a_i$ と同一視してよい。\par
19) もちろん、ここでも $P$-加群準同型である表現だけを考える。すなわち $a\sim A$ から $\rho\in P$ について $a\rho\sim A\rho$ が従うべきであり、$\Omega$ についても対応して同様である。\par}

イデアルまたは表現加群、および対応する表現は、代数閉体へ移った後も既約にとどまるなら、\emph{絶対既約}と呼ばれる。

$\mo\Omega$ が根基をもたないなら、$\mo$ も根基をもたない。$\mo$ の冪零イデアル $\mc$ は $\mo\Omega$ の拡張イデアル $\mc\Omega$ を与えるからである。逆は成り立たないことが下で分かる。

$\mo\Omega$ の中心は、拡張された中心 $\mZ\Omega$ に等しい。$\mo$ が完全可約であるなら、$\mZ$ は体の直和である：
\[
        \mZ=\mZ_1+\cdots+\mZ_s.
\]
$\Omega$ へ移っても $\mo\Omega$ が完全可約にとどまるなら（すなわち冪零イデアルが加わらないなら）、新しい中心は
\[
        \mZ\Omega=\mZ_1\Omega+\cdots+\mZ_s\Omega
\]
と再び体の直和に分解する。$\Omega$ が代数閉であれば、これらの体は $\Omega$ 上階数 $1$ をもち、$\Omega$ に同型である。

中心 $\mZ$ の表現については、次の定理が成り立つ。

代数閉体 $\Omega$ における可換超複素系 $\mZ$ の任意の既約表現は次数 $1$ である。同値に、既約表現は準同型 $\mZ\to\Omega$ と同一である。

証明。$\mZ$ の任意の既約表現は $\mZ\Omega$ の表現を一つ与え、したがって $\mZ\Omega/\mc$ の表現も与える。ただし $\mc$ は $\mZ\Omega$ の根基である。商 $\mZ\Omega/\mc$ は根基をもたない可換系であり、§14 の終わりにより体の直和である。$\Omega$ は代数閉なので、これらは次数 $1$ であり、すべての既約表現を生成する（§20）。

同時に、次数 $1$ の同値表現は必ず等しくなるので、相異なる準同型 $\mZ\to\Omega$ の数は $\mZ\Omega/\mc$ の階数に等しい。

特別な場合として、$\mZ$ が $\mathsf{P}$ 上の体であれば、次の定理が得られる。$\mZ$ が完全可約、すなわち根基をもたないことは、$\mZ$ が $P$ の第一種拡大であることと同値である。体では準同型は同型である。その数は $\mZ\Omega/\mc$ の階数であり、$\mc=0$ のときちょうど体の次数に等しい。

\[
        \mZ=\mZ_1+\cdots+\mZ_s
\]
が完全可約であるなら、各表現において個々の体 $\mZ_i$ も準同型的に表現される。既約表現では、それらの体の一つが同型的に写され、他は零で写される。

根基をもたない系について、これらの事実を超複素系に適用すると、まず次を得る。

$\mo\Omega$、したがって $\mo$ も根基をもたない系であるなら、$\mo$ の任意の絶対既約表現において、中心元 $z$ は対角行列 $E\xi$ によって表現され、$z\mapsto \xi$ は $\Omega$ における $\mZ$ の一次表現である。こうして得られる $\mo$ の絶対既約表現類と $\mZ$ のそれとの対応は一対一である。したがって、これらの表現類の数は中心の階数に等しい。\textsuperscript{20)}

{\footnotesize\noindent 20) 対応する定理は、代数閉体 $\Omega$ における表現だけでなく、個々の自己同型斜体における表現についても成り立つ。これは別の所で展開される予定である。\par}

実際、両側分解不能イデアルへの分解
\[
        \mo\Omega=\ma_1+\cdots+\ma_s
\]
は、中心の体への分解
\[
        \mZ\Omega=\mZ_1+\cdots+\mZ_s
\]
に一対一に対応する。ここで $\mZ_i$ は $\ma_i$ の中心である。$\mo$ の既約表現では、一つの $\ma_i$ が同型的に写され、他は零で写される。表現類は $\ma_i$ によって一意に決まる。成分 $\ma_i$ は完全行列環によって表現され、そのような完全行列環の中心は対角行列 $E\xi$ からなる。対応 $z\mapsto \xi$ は $\mZ$ の準同型であり、一つの $\mZ_\nu$ は同型的に写され、他は零で写される。これで主張が証明される。

根基のない $\mo$ が根基のない $\mo_\Omega$ を与えるのはいつかという問題については、$\mZ$ が $P$ の第二種拡大である成分 $\mZ_i$ をもてば、$\mo\Omega$ は確かに根基をもつことが従う。\textsuperscript{21)} この場合には $\mZ_i\Omega$ がすでに根基をもつからである。

{\footnotesize\noindent 21) $\mZ$ が第一種の成分だけをもつなら、$\mo\Omega$ は根基をもたない。このことも後で証明される。標数零については、I. Schur がすでに同値な定理、すなわち $P$ 上の既約表現は係数領域の任意の拡大の下で完全可約にとどまる、という定理を証明している。I. Schur, Beiträge zur Theorie der Gruppen linearer Substitutionen, \emph{Trans. Amer. Math. Soc.} 15 (1909), p. 159 を参照。\par}

\subsection*{§ 22. アーベル群への応用}

\[
        G=G_1\times G_2\times\cdots\times G_r
\]
を、位数 $h_i$ の巡回群 $G_i$ に分解された有限アーベル群とする。その元は
\[
        a=a_1^{\alpha_1}\cdots a_r^{\alpha_r}
\]
である。$P$ を、標数が群の位数
\[
        h=h_1h_2\cdots h_r
\]
を割らない体とする。群環 $\mZ$ はすべての和
\[
        \sum_{\alpha_1,\ldots,\alpha_r} A_{\alpha_1\ldots\alpha_r}
        a_1^{\alpha_1}\cdots a_r^{\alpha_r}
        \qquad (A\in P)
\]
からなり、$x_i\mapsto a_i$ によって多項式環 $P[x_1,\ldots,x_r]$ の準同型像である。したがって
\[
        \mZ\simeq P[x_1,\ldots,x_r]/(x_1^{h_1}-1,\ldots,x_r^{h_r}-1).
\]
拡大体 $\Omega$ では、多項式 $x_i^{h_i}-1$ は相異なる因子 $x_i-\xi_i$ に分解する。したがって定義イデアルは、相異なる極大イデアル
\[
        (x_1-\xi_1^{(\alpha_1)},\ldots,x_r-\xi_r^{(\alpha_r)})
\]
の積、または交わりである。ただしこれは各零点 $(\xi_1^{(\alpha_1)},\ldots,\xi_r^{(\alpha_r)})$ に一つずつ対応する。この交わりに対して直和分解
\[
        \mZ\Omega=\mZ_1+\cdots+\mZ_h
\]
が対応し、各成分は $\Omega$ に同型であって、表現
\[
        a_i\mapsto \xi_i^{(\alpha_i)},
        \qquad\hbox{またはより一般に}\qquad
        a\mapsto \chi^{(\alpha)}(a)
\]
を与える。これらの表現が指標である。その数は、一般理論が要求する通り、$h_1\cdots h_r=h$ である。

\subsection*{§ 23. 超複素系の行列式}

\[
        \mo=a_1P+\cdots+a_hP
\]
を超複素系とする。$\mathsf{P}$ に $h$ 個の不定元 $x_1,\ldots,x_h$ を付け加え、
\[
        \mo^*=a_1\mathsf{P}(x)+\ldots+a_h\mathsf{P}(x)
\]
を作る。ただし基底元 $a_i$ には古い乗法規則を用いる。$x$ は $a_i$ と可換とし、これによって $\mo^*$ における計算規則が定まる。

$\mo^*$ の中には $\mo$ の「一般元」
\[
        w=a_1x_1+\ldots+a_hx_h
\]
がある。表現において $a_i\to A_i$ なら、$w$ には
\[
        W=A_1x_1+\ldots+A_hx_h
\]
が対応する。行列 $W$ をその表現に属する系行列という。$a_i$ が群をなし、$\mo$ が群環であるなら、それは群行列である。正則表現では正則系行列が得られる。

$W$ の成分は $x$ の線形形式である。したがって系行列式 $|W|$ は、次数 $n$ の表現に対して次数 $n$ である。とくに正則系行列式は次数 $h$ をもつ。

系行列式は同値な表現へ移っても変わらない。なぜなら
\[
        |PWP^{-1}|=|P|\,|W|\,|P^{-1}|=|W|
\]
だからである。基底 $(a_1,\ldots,a_h)$ から新しい基底 $(b_1,\ldots,b_h)$ へ移り、$w=\sum a_ix_i$ から $w=\sum b_iy_i$ へ移るとき、$W$ の新しい元、したがって新しい行列式も、古いものから変数の正則置換
\[
        x_i=\sum \varrho_{ik}y_k
\]
によって得られる。

表現加群の組成列が与えられているなら、適切な基底の選択に対して行列 $W$ は下三角ブロック形
\[
        W=\begin{pmatrix}
        W_1&0&\cdots&0\\
        *&W_2&\cdots&0\\
        \vdots&\vdots&\ddots&\vdots\\
        *&*&\cdots&W_r
\end{pmatrix}
\]
をもつ。従って
\[
        |W|=|W_1|\cdot|W_2|\cdots|W_r|
\]
である。任意の表現の行列式 $|W_i|$ の中には、$\mo$ の正則表現、あるいは根基を $\mc$ とする $\mo/\mc$ の正則表現に現れる因子以外は現れない。

$P$ が代数閉であるなら、既約表現に属する行列式 $|W_i|$ は $x$ の素関数であり、非同値な表現は異なる素因子を与える。

証明。$\mo$ のすべての既約表現は $\mo/\mc$ の表現でもあるから、根基をもたない環 $\mo/\mc$ に制限してよい。その中で行列単位 $c_{ik}$ を基底に取ると、一般元は
\[
        w=\sum_{\nu,i,k} c_{ik}^{(\nu)}x_{ik}^{(\nu)}
\]
である。既約表現の行列は
\[
        W_\nu=(x_{ik}^{(\nu)})
\]
である。関数 $|W_\nu|=|x_{ik}^{(\nu)}|$ は既約であることが知られており、明らかに互いに異なる。

$|W_\nu|$ を計算するには、$\mo/\mc$ の正則系行列を因数分解すればよい。各素因子は、対応する既約表現の次数だけ現れる。対応するイデアルが組成列の中にその回数だけ現れるからである。

$\mo$ の正則系行列から出発することもできるが、その場合、各既約因子はより多く、すなわち対応する単純イデアルが左イデアルの中で組成因子として現れる回数だけ得られる。この正則表現では $\mo$ は左イデアルとして見られていた。$\mo$ を右イデアルとして見るなら、第二の正則系行列、Frobenius の「反対形行列」が現れる。これは同じ既約因子、すなわちすべての既約表現の系行列式を含むが、指数は異なり得る。

可換系の系行列式は一次因子に分解する。すべての既約表現が次数 $1$ だからである。これらの一次因子は既約表現そのものであり、アーベル群の場合には指標を与える。この事実は、Dedekind が非アーベル群の群行列式を研究する際の出発点であった。

\subsection*{§ 24. トレースと指標}

超複素系 $\mo$ の表現 $\mo\sim\mathfrak D$ において、元 $a$ に行列 $A$ が対応するとき、
\[
        \operatorname{Sp}_D(a)=\operatorname{Sp} A
\]
とおく。トレースは線形関数である：
\[
        \mathit{Sp}(c+d)=\mathit{Sp}(c)+\mathit{Sp}(d),
        \qquad
        \mathit{Sp}(c\alpha)=\alpha\mathit{Sp}(c).
\]
同値な表現は同じトレースをもつ。可約表現において、トレースは組成因子によって媒介される表現におけるトレースの和である。実際、
\[
        \mathit{Sp}\begin{pmatrix}
        A_{11}&0&\cdots&0\\
        *&A_{22}&\cdots&0\\
        \vdots&\vdots&\ddots&\vdots\\
        *&*&\cdots&A_{ss}
        \end{pmatrix}
        =\operatorname{Sp}(A_{11})+\operatorname{Sp}(A_{22})+\cdots+\operatorname{Sp}(A_{ss}).
\]

\emph{主トレース}とは正則表現におけるトレースである。\emph{被約トレース}とは、相異なる既約表現におけるトレースの和である。

$c$ が極大冪零イデアル $\mc$ の元であるなら、任意の表現について $\mathit{Sp}(c)=0$ である。

証明。$\mo/\mc$ の単純左イデアルによる表現だけを考えれば十分である。任意の他の表現では、これらの組成因子だけが現れ、トレースは加法的に構成されるからである。しかし $\mathfrak{c}$ は $\mo/\mc$ のすべての元を零化する。したがって各 $c\in\mc$ には零行列が割り当てられ、$\mathit{Sp}(c)=0$ となる。

$\mo$ が両側単純で、$\mathsf{P}$ が代数閉であり、したがって行列環
\[
        \mo=\sum c_{ik}P
\]
であるなら、$\mo$ の既約表現において、元
\[
        a=\sum c_{ik}\alpha_{ik}
\]
には行列 $(\alpha_{ik})$ が割り当てられる。したがって
\[
        \operatorname{Sp}_{\rm red}(a)=\sum_i\alpha_{ii},
        \qquad
        \operatorname{Sp}_{\rm pr}(a)=n\cdot\operatorname{Sp}_{\rm red}(a)=n\cdot\sum_i\alpha_{ii}.
\]
とくに
\[
        \operatorname{Sp}_{\rm red}(c_{ik})=0\quad(i\ne k),\qquad
        \operatorname{Sp}_{\rm red}(c_{ii})=e,
\]
一方で
\[
        \operatorname{Sp}_{\rm pr}(c_{ik})=0\quad(i\ne k),\qquad
        \operatorname{Sp}_{\rm pr}(c_{ii})=n\cdot e.
\]
$P$ から代数閉体 $\Omega$ へ移っても主トレースは変わらない。同じ基底から計算できるからである。これは被約トレースについては成り立たない。

根基をもたない系 $\mo$ の元 $a$ の、絶対既約表現におけるトレースを指標と呼び、$\chi(a)$ と書く。表現を指定する場合には $\chi^{(\nu)}(a)$ と書く。

次数 $n_\nu$ の既約表現において、中心元は §21 により対角行列 $E\theta^{(\nu)}$ として表される。ここで $z\mapsto\theta^{(\nu)}(z)$ は中心の次数 $1$ の表現、または中心から $\Omega$ への準同型である。したがってトレースの値は $n_\nu\zeta^{(\nu)}$ である。ゆえに中心の準同型 $\Theta$ は指標 $\chi$ と
\[
        \chi^{(\nu)}(z)=n_\nu\theta^{(\nu)}(z).              \tag{1}
\]
によって結びついている。可換の場合には $n_\nu=1$ であり、指標自身が準同型を与える。

以下では常にそう仮定するように、$\Omega$ の標数が零であれば、(1) を $n$ で割ることができる：
\[
        \theta^{(\nu)}(z)=\frac{\chi^{(\nu)}(z)}{n_\nu}.
\]
$\theta^{(\nu)}(z)$ の準同型性は
\[
        \frac{\chi^{(\nu)}(z+z')}{n_\nu}
        =\frac{\chi^{(\nu)}(z)}{n_\nu}
         +\frac{\chi^{(\nu)}(z')}{n_\nu},
\]
および
\[
        \frac{\chi^{(\nu)}(zz')}{n_\nu}
        =\frac{\chi^{(\nu)}(z)}{n_\nu}\,
         \frac{\chi^{(\nu)}(z')}{n_\nu}
\]
で表される。

表現類は、行列のトレースだけですでに一意に定まる。すべての行列のトレースを知るには、もちろん、与えられた表現における基底元のトレースを知れば十分である。

証明。表現加群 $R$ は、各既約表現加群 $\mM_\nu$ がそこに直和因子として何回現れるかを知れば分かる。この数を $p_\nu$ とすると、表現における $e^{(\nu)}$ のトレースは $p_\nu n_\nu$ である。ただし $n_\nu$ は既約表現の次数である。したがって
\[
        \mu_\nu=\frac{\operatorname{Sp}(e^{(\nu)})}{n_\nu}.
\]

\subsection*{§ 25. 判別式}

\[
        \mo=a_1P+\cdots+a_hP
\]
を超複素系とする。\emph{判別行列}とは、成分が主トレース $\operatorname{Sp}(a_i a_j)$ である行列である。\emph{被約判別行列}とは、代数閉体の中で作られた、被約トレースに関する対応する行列である。この行列の行列式を、それぞれ判別式、被約判別式という。

判別式は基礎体を拡大しても同じである。別の基底に移ると、変換行列式の平方が掛かる。

証明。
\[
        (b_1,\ldots,b_h)=(a_1,\ldots,a_h)P
\]
とする。すると
\[
        \bigl(\operatorname{Sp}(a_i a_j)\bigr)=\bigl(\operatorname{Sp}(a_i b_j)\bigr)P.       \tag{1}
\]
同様に、$P^t$ が転置行列を表すなら
\[
        \bigl(\operatorname{Sp}(a_i b_j)\bigr)=P^t\bigl(\operatorname{Sp}(b_i b_j)\bigr),      \tag{1a}
\]
であり、(1) と (1a) から
\[
        \bigl(\operatorname{Sp}(a_i a_j)\bigr)=P^t\bigl(\operatorname{Sp}(b_i b_j)\bigr)P.
\]
行列式に移れば主張を得る。

したがって行列式は $P$ からの平方因子を除いてしか定まらないが、それが零であるか非零であるかは不変な性質である。(1) からはまた、非零因子を除けば、判別式を二つの異なる基底から、すなわち $|\operatorname{Sp}(a_i b_j)|$ として計算できることも従う。

$\mo$ が冪零イデアル $\mc$ をもてば、判別式は零になる。

証明。$\mo$ の基底元を
\[
        c_1,\ldots,c_t,d_1,\ldots,d_s
\]
の形に選び、$c_1,\ldots,c_r$ が $\mc$ の基底をなすようにする。すると判別行列は、行列式に必要な上ブロックの意味で
\[
        \begin{pmatrix}
        \operatorname{Sp}(c_i c_j)&\operatorname{Sp}(c_i d_j)\\
        \operatorname{Sp}(d_i c_j)&\operatorname{Sp}(d_i d_j)
        \end{pmatrix}
        =
        \begin{pmatrix}
        0&0\\
        0&\operatorname{Sp}(d_i d_k)
        \end{pmatrix}
\]
の形をもつ。したがって行列式は零である。被約判別式についても同じである。

その結果、代数閉体へ移った後に冪零イデアルが現れる場合にも、判別式は零になる。

\[
        \mo=\ma_1+\ldots+\ma_s
\]
を両側直和とし、$M,M_1,\ldots,M_s$ を環 $\mo,\ma_1,\ldots,\ma_s$ の判別行列とする。このとき $M$ はブロック対角であり、
\[
        M=\begin{pmatrix}
        M_1&0&\cdots&0\\
        0&M_2&\cdots&0\\
        \vdots&\vdots&\ddots&\vdots\\
        0&0&\cdots&M_s
        \end{pmatrix},
\]
したがって
\[
        |M|=|M_1|\cdot\,|M_2|\cdots |M_s|.
\]
実際、$\mo$ の基底を $\ma_1,\ldots,\ma_s$ の基底から作って選ぶ。$a\in\ma_i$ なら
\[
        \mathit{Sp}_{\mo}(\ma)=\mathit{Sp}_{\ma_i}(a),
\]
である。ただしどちらも主トレースを意味し、前者は環 $\mo$ におけるもの、後者は環 $\ma_i$ におけるものである。さらに、因子が異なる両側成分に属すれば $a_i a a_j=0$ である。したがって混合トレースブロックは零である。同じ議論が被約判別式にも適用される。

行列環 $\sum c_{ik}\mathsf{P}$ の判別式を作るには、二つの異なる基底を選ぶ。すなわち $c_{ik}$ を一度は第一添字で、もう一度は第二添字で順序づける。原典ページに現れる積表は次の通りである：
\[
\begin{array}{c|ccc}
& \overbrace{c_{11}\;\ldots\;c_{1n}}^{\mathfrak r_1}
& \overbrace{c_{21}\;\ldots\;c_{2n}}^{\mathfrak r_2}
& \ldots \\ \hline
\left.\begin{array}{c}c_{11}\\[-1mm]\vdots\\[-1mm]c_{n1}\end{array}\right\}t_1
& (c_{ik}) & 0 & 0 \\[1.3em]
\left.\begin{array}{c}c_{12}\\[-1mm]\vdots\\[-1mm]c_{n2}\end{array}\right\}t_2
& 0 & (c_{ik}) & 0 \\[1.3em]
\vdots & 0 & 0 & (c_{ik})
\end{array}
\]
これから
\[
        D=\left|\begin{array}{cccc}
        \mathit{Sp}(c_{11})& & &0\\
        &\mathit{Sp}(c_{22})& & \\
        & &\ddots& \\
        0& & &\mathit{Sp}(c_{nn})
        \end{array}\right|^{n}
        =\begin{cases}
        e, & \text{被約トレースの場合},\\
        n^{n^2\cdot}e, & \text{主トレースの場合}.
        \end{cases}
\]
すなわち
\[
        D_{\rm red}=e
        \quad\hbox{被約トレースについて},
        \qquad
        D=n^{n^2}e
        \quad\hbox{主トレースについて}.
\]

したがって、代数閉拡大体上で $\mo$ が次数 $n_1,\ldots,n_s$ の行列環に分解するなら、その判別式は
\[
        D=n_1^{n_1^2}n_2^{n_2^2}\ldots n_s^{n_s^2}\cdot e,
\]
またその被約判別式は
\[
        D_{\rm red}=e.
\]
したがって、根基をもたない系について被約判別式は常に非零である。主判別式は、どの $n_i$ も体の標数で割り切れないとき、かつそのときに限り非零である。ゆえに次を得る。

\emph{被約判別式が非零であることは、体 $\mathsf{P}$ の代数閉包へ移った後に根基をもたない系であるための必要十分条件である。}

\emph{標数零では、判別式が非零であることは、体 $\mathsf{P}$ の代数閉包へ移った後に根基が現れないための必要十分条件である。}\textsuperscript{22)}

{\footnotesize\noindent 22) 可換系については、E. Noether, \emph{Diskriminantensatz für Ordnungen ...}, J. reine angew. Math. 157 (1927), pp. 82--104, §§4--6 を参照。そこの証明方法は同じであるが、細部ではより複雑である。基底から別の基底への移行（同論文 §4, 4）は、ここでこの節の冒頭に与えたものに置き換えるべきである。そこに現れる行列式は零になるからである。\par}

\subsection*{§ 26. 群環の位置づけ}

$a_1,\ldots,a_h$ を有限群の元とし、標数が $h$ を割らない体上で群環を作る。このとき判別式 $D\ne0$ であり、したがって群環は根基をもたない環である。

証明。まず、以後 $\mathit{Sp}$ を主トレースとすることにして、
\[
        \operatorname{Sp}(e)=he,
        \qquad
        \operatorname{Sp}(a_i)=0\quad (a_i\ne e)
\]
を示す。正則表現では
\[
        e\mapsto\begin{pmatrix}e&0&\cdots&0\\0&e&\cdots&0\\\vdots&\vdots&\ddots&\vdots\\0&0&\cdots&e\end{pmatrix},
        \qquad \operatorname{Sp}(e)=he
\]
である。$a_i\ne e$ に対して、積 $a_i a_1,\ldots,a_i a_h$ は $a_1,\ldots,a_h$ の固定点のない置換である。したがって、これらの積を $a_1,\ldots,a_h$ によって表す行列は主対角上に零だけをもつので、$\mathit{Sp}(a_i)=0$ である。

いま基底 $a_1,\ldots,a_h$ と $a_1^{-1},\ldots,a_h^{-1}$ を取る。行列 $\bigl(\operatorname{Sp}(a_i a_j^{-1})\bigr)$ は対角行列で、対角成分は $he$ である。したがって
\[
        D=h^h e\ne0.
\]

群元 $a$ の類とは、群環の中で作られる和
\[
        K_a=\sum_s s^{-1}as,                              \tag{1}
\]
をいう。ただし和は相異なる共役元 $s^{-1}as$ だけにわたる。類 $K_a$ は中心元である。すべての群元と可換だからである。これらは中心を生成する。実際、群環の元 $\sum_i \alpha_i a_i$ がすべての $s$ と可換であるなら
\[
        \sum_i\alpha_i a_i=s^{-1}\biggl(\sum_i\alpha_i a_i\biggr)s
        =\sum_i\alpha_i(s^{-1}a_i s),
\]
であるから、係数は各共役類上で一定である。

したがって中心の階数は共役な群元の類の数に等しい。ゆえに絶対既約表現の数もこの類数に等しい。

行列 $A$ と $S^{-1}AS$ は同じトレースをもつ。したがって群元 $a$ と $s^{-1}as$ は同じトレースをもつ。(1) の両辺でトレースを取ると
\[
        \chi(K_a)=h_a\chi(a),
\]
を得る。ここで $h_a$ は類 $K_a$ に含まれる元の数である。言葉で言えば、群元の指標は、その類の指標を類に含まれる元の数で割ったものに等しい。

\begin{center}
（1928年8月12日受理。）
\end{center}



\clearpage
\clearpage
\setcounter{footnote}{12}
\section*{35. 整係数関数からなる極大領域について}
\begin{center}
E. Noether\\
\textit{Rec. Soc. Math. Moscou \textbf{36} (1929), pp. 65--72}
\end{center}

以下の考察は、相対的に整な関数に関する Hilbert の問題\footnote{D. Hilbert, \emph{Mathematische Probleme} (Gott. Nachr. 1900, pp. 253--297), 第14問題。}に密接に関わるものであり、E. Hecke の折に触れた質問に由来する。Hecke は、自分が補題として必要とした一つの特殊な場合を、直接計算によって処理することができたのである。\footnote{E. Hecke, ``Über die Konstruktion relativ-Abelscher Zahlkörper durch Modulfunktionen von zwei Variabeln'' (Math. Ann. 74 (1913), pp. 465--510), \S\,2。Hecke では、有限な整域 $\mathfrak{I}$ が、二変数の冪級数の三つの商によって生成され、それらの間には一つの代数関係が成り立つ。}

問題設定は次の通りである。与えられた領域 $\mathfrak{Z}$ の内部で、$z_1,\ldots,z_r$ の整係数関数、すなわち多項式、冪級数、またはそれらの商からなる\emph{極大領域}とは、整数分母を付け加えてももはや拡大できない領域のことである。すなわち、$g(x)$ が $\mathfrak{Z}$ に属し、$a$ が非零の有理整数であるとき、$a\cdot g(x)$ の所属から常に $g(x)$ の所属が従うものをいう。$\mathfrak{Z}$ 内の任意の整域 $\mathfrak{I}$ は、一意的に極大領域 $M(\mathfrak{I})$ を生成する。これは、ある $a\ne0$ について $a\cdot g(x)$ が $\mathfrak{I}$ に属するような、$\mathfrak{Z}$ のすべての関数 $g(x)$ から成る。

問われるのは、もとの整域 $\mathfrak{I}$ が有限であった場合、すなわち $\mathfrak{Z}$ の有限個の関数 $f_1(x),\ldots,f_t(x)$ に関する整係数多項式全体から成っていた場合、この極大な $M(\mathfrak{I})$ の分母 $a$ について何が言えるかである。このとき $\mathfrak{Z}$ としては、$z_1,\ldots,z_r$ のすべての整係数多項式または冪級数、あるいは $f(x)$ に現れる分母とその冪以外の分母を含まない商を取るものとする。

私は、冪級数またはその商の場合には、$f(x)$ 相互の依存関係が代数的なものだけであるという付加仮定のもとで、$a$ を常に、そこに現れる相異なる素数が有限個に限られるように選べることを示す。ただし、一般にはそれらの素数は任意の高い冪で現れ得る。

後者は直ちに明らかである。$a\cdot g(x)$ は $\mathfrak{I}$ に属するが、どの指数 $\varrho$ に対しても $a^\varrho g(x)$ があらかじめ制限された形で $\mathfrak{I}$ に入らないならば、すべての冪 $a^\varrho$ が現れる。例えば $\mathfrak{I}=[2x]$ のとき、$x=(2x)/2$、$x^2=(2x)^2/2^2$ である。

有界性の問題は次のようにしか定式化できない。一般には無限であるような $M(\mathfrak{I})$ の生成系を、対応する分母における素数の指数が有界に留まるように与えられるか、という問題である。関与する素数は有限個であるから、この問題は既知の議論\footnote{Hilbert 前掲、または E. Noether, \emph{Die Endlichkeit des Systems der ganzzahligen Invarianten binärer Formen} (Gött. Nachr. 1919, pp. 138--156), \S\,1, IV のより詳しい説明を参照。}によって $M(\mathfrak{I})$ の有限性の問題と同値であり、したがって整係数の形での、相対的に整な関数に関する Hilbert の問題そのものである。この問題は、ここで用いる方法では決定できない。\footnote{したがって例えば、基本形式系の整（射影的）不変式の場合には、それらが通常の意味での完全な不変式系によって、しかも分母に有限個の相異なる素数だけを含んで表されることが示される。これは通常の $\Omega$ 過程の方法からは従わない結果である。しかし整な有限性そのものについては何も述べていない。それはこれまで二元基本形式の場合にだけ知られている（3) で引用した注を参照）。}

分母 $a$ に現れる相異なる素数 $p$ が有限個だけであることの証明は、各そのような $p$ に対し、$\mathfrak{I}$ を生成する $f(x)$ の間に少なくとも一つの法 $p$ の関係が対応する、という事実に基づく。この関係は整数係数をもち、$x$ において恒等的に満たされるものではない。別の言い方をすれば、$\mathfrak{o}$ を不定元 $u_1,\ldots,u_t$ に関する整係数多項式領域とすると、対応 $u_i\mapsto f_i(x)$ によって $\mathfrak{I}$ は $\mathfrak{o}$ の準同型像となる。この準同型は $\mathfrak{o}$ の素イデアル $\mathfrak{a}$ によって媒介される、すなわち $\mathfrak{I}$ は剰余類環 $\mathfrak{o}/\mathfrak{a}$ と同型になる。これに対応して、$\mathfrak{I}_p$ は $\mathfrak{o}_p$ の準同型像となる。添字は法 $p$ の剰余類への移行を意味し、これは $\mathfrak{I}$ の分母に現れる有限個の素数を除いて可能である。この準同型は再び $\mathfrak{o}_p$ の素イデアル $\mathfrak{c}_p$ によって媒介され、これは $\mathfrak{a}_p$ を含む。すると、$p$ が $\mathfrak{M}(\mathfrak{I})$ の分母に現れることと、$\mathfrak{c}_p$ が $\mathfrak{a}_p$ の真の上位イデアル（イデアルの意味での真の約数）になることとは同値である。これは、$\mathfrak{I}_p$ の超越次数が $\mathfrak{I}$ のそれに比べて下がるか、あるいは $\mathfrak{a}_p$ が素イデアル性を失う場合にだけ起こり得る。\footnote{系の「超越次数」とは、周知のように、その系が代数的に依存する代数的独立な関数の個数、しかも不変な個数のことである。代数的独立な関数から成る系を既約系と呼ぶ。素イデアルの超越次数（次元）とは、その剰余類体の超越次数をいう。素イデアルは同じ超越次数をもつ真の素約数をもたない。例えば B. L. v. d. Waerden, ``Zur Nullstellentheorie der Polynomideale'', Math. Ann. 96 (1926), pp. 183--208 を参照。} これが有限個の $p$ についてしか起こらないことは、法 $p$ においても代数的依存関係が関数行列式の消滅をもたらすこと（\S\,1）、さらに $\mathfrak{a}$ が絶対素イデアルであり、したがってその素イデアル性を失うのは有限個の $p$ を法とする場合だけであることから従う。\footnote{E. Noether, ``Eliminationstheorie und allgemeine Idealtheorie'', Math. Ann. 90 (1923), pp. 229--261, \S\,7。この定理は消去理論により\textsuperscript{*)}、絶対既約多項式がその性質を失うのは有限個の素数を法とする場合に限られるというより簡単な定理から従う（A. Ostrowski, ``Zur arithmetischen Theorie der algebraischen Größen'', Gött. Nachr. 1919, pp. 279--298, p. 296。より簡単な証明は E. Noether, ``Ein algebraisches Kriterium für absolute Irreduzibilität'', Math. Ann. 85 (1922), pp. 26--33, \S\,7）。$\mathfrak{I}$ がその超越次数よりちょうど一つ多い関数を含む整基底を持つ場合、例えば Hecke が扱った場合（注2）のように、$\mathfrak{a}$ は主イデアルになるので、消去理論を避けてこのより簡単な定理を用いるだけで足りる。素数の代わりに、どこでも（有限な）代数的数体の素イデアルを置くことができる。この定理は $\mathfrak{a}$ が零イデアルの場合にも明らかに成り立つ。}
\srcfntext{*)}{すなわち、$\mathfrak{a}$ は標数零の係数をもつので、本論文の補助定理 V とは独立である。この補助定理の証明は正しくない（そこで示された代入では、すべてが恒等的に消える可能性がある）。正しい証明は B. L. v. d. Waerden, \emph{Eine Verallgemeinerung des Bézoutschen Theorems}, Math. Ann. 99 (1928), pp. 497--541 に初めて与えられている。同論文の注 37 を参照。}

さらに、本論文の p. 257 には、容易に分かる別の訂正が必要である（定理 XIV も \S\,7 では用いられていない）。そこでは、$(\mathfrak R)$ のすべての元は、$(x_i)$ における $t_{\mu,\nu}$ を $(\mathfrak R')$ の元へ特殊化することで生じる、と述べられている。実際には、それらは、より一般的な線型形式
\[
\sum s_{\lambda_1\ldots\lambda_n}(y_1)^{\lambda_1}\cdots(y_n)^{\lambda_n}
\]
を特殊化することで生じるが、その指数は $(x_i)$ の指数と同じである。

なお、すべての考察は、有理整数の代わりに（有限な）代数的数体の整数とその素イデアルを用いる場合にも、わずかな修正で保たれる（\S\,4）。

\subsection*{\S\,1. 特性 $p$ の係数領域における関数行列式}

\noindent\textbf{1.}
$f_1(x_1,\ldots,x_r),\ldots,f_t(x_1,\ldots,x_r)$ を、有理関数、あるいはより一般に、形式的に定義された冪級数またはその商とし、その係数は特性 $p$ の体 $P$ に属するとする。導関数 $\partial f_i/\partial x_k$ も形式的に定義され、通常のすべての計算法則が保たれるものとする。$\varOmega$ は $\mathrm{P}$ の代数閉拡大体を意味する。

そのとき、特性零の場合と同様に次が成り立つ。

\noindent\textbf{定理.}
\emph{$f(x)$ の間に $\varOmega$ 係数の代数的依存関係が存在するなら、関数行列 $(\partial f_i/\partial x_k)$（$i=1,\ldots,t$、$k=1,\ldots,n$）の階数は $t$ より小さい。}\footnote{元の形では、係数領域 $P$ を完全体と仮定していた。$P$ から $\Omega$ へ移ることによって、任意の体、完全でない体にまで拡張できることは B. L. v. d. Waerden に負う。特性 $p$ では、$P$ が完全であっても逆は成り立たないことは、例 $f=x^p$, $\partial f/\partial x=0$ が示している。}

\noindent\emph{証明.}
多項式領域 $\varOmega[u_1,\ldots,u_t]$ において、$G(f)=0$ が $x$ に関して恒等的に成り立つようなすべての多項式 $G(u)$ の素イデアルを $\mathfrak{m}$ とする。仮定により $\mathfrak{m}$ は零イデアルではない。$\mathfrak{m}$ の中で最小次数の多項式 $G(u)\ne0$（したがって定数ではない）を取る。特に $G(u)$ は分解不能である。すると通常通り
\[
\begin{aligned}
\left(\frac{\partial G}{\partial u_1}\right)_{f}\frac{\partial f_1}{\partial x_1}
 +\cdots+
\left(\frac{\partial G}{\partial u_t}\right)_{f}\frac{\partial f_t}{\partial x_1}&=0,\\[-1mm]
&\vdots\\[-1mm]
\left(\frac{\partial G}{\partial u_1}\right)_{f}\frac{\partial f_1}{\partial x_n}
 +\cdots+
\left(\frac{\partial G}{\partial u_t}\right)_{f}\frac{\partial f_t}{\partial x_n}&=0
\end{aligned}
\]
が得られる。したがって、関数行列の階数が $t$ より小さいか、またはすべての $\left(\partial G/\partial u_i\right)_{f}$ が消えるかのいずれかである。しかし $G(u)$ は最小次数の多項式として選ばれていたから、後者ならすべての $\partial G/\partial u_i$ が消えることになる。

よって $G(u)=H(u_1^p,\ldots,u_t^p)$ となる。さらに係数領域 $\varOmega$ は代数閉、したがって完全であるから、$G(u)=H(u_i^p)=(T(u))^p$ となり、$G(u)$ の選び方に反する。証明はもちろん特性零の場合にも通用し、その場合には最後の推論が不要になる。

\medskip\noindent\textbf{2. 系.}
$F_1(x_1,\ldots,x_n),\ldots,F_t(x_1,\ldots,x_n)$ を、関数行列がちょうど階数 $t$ をもつ整係数関数（多項式、冪級数またはその商）とする。素数 $p$ は、$F_i(x)$ の分母にも、関数行列のすべての $t$ 行小行列式にも入らないように選ぶ。

すると $F_1(x),\ldots,F_t(x)$ は法 $p$ でも代数的独立のままである。

\noindent\emph{証明.}
$\bar P$ を法 $p$ の剰余類体とし、$f_1(x),\ldots,f_t(x)$ を、$F_1(x),\ldots,F_t(x)$ の整係数をその法 $p$ の剰余類で置き換えて得られる $\bar P$ 上の関数とする。$p$ は $F_i(x)$ の分母に入っていなかったので $f_i(x)$ は定義される。また、関数行列 $(\partial f_i/\partial x_k)$ も、$(\partial F_i/\partial x_k)$ の整係数を法 $p$ の剰余類で置き換えて得られる。$(\partial F_i/\partial x_k)$ には $F_i(x)$ にある以外の分母は現れない。$p$ の選び方により、$(\partial f_i/\partial x_k)$ は階数 $t$ を保つ。したがって $f_i(x)$ は $\varOmega$ で代数的独立であり、まして $\bar P$ でも代数的独立である。

\subsection*{\S\,2. 定理の定式化}

\noindent\textbf{1.}
導入部と同様、$\mathfrak{I}$ は $x_1,\ldots,x_n$ の整係数関数、すなわち有理整数係数の多項式または冪級数、あるいはそれらの商からなる整域（零因子をもたない環）を表すものとする。冪級数の場合には、形式的に定義されたものでも、ある領域で収束するものでもよい。用いられるのは、現れる関数が整域をなすという事実だけである。

$\mathfrak{Q}$ を $\mathfrak{I}$ の商体とし、$\mathfrak{Z}$ を $\mathfrak{I}$ と $\mathfrak{Q}$ の間にある整域とする。導入部と同様に、$M$ が $\mathfrak{Z}$ において\emph{極大}であるとは、$a\ne0$ が整数で $g(x)$ が $\mathfrak{Z}$ に属するとき、$a\cdot g(x)$ の所属から常に $g(x)$ の所属が従うことをいう。$\mathfrak{I}$ に対しては、$\mathfrak{Z}$ 内で $\mathfrak{I}$ によって生成される極大領域 $M(\mathfrak{I})$ が一意的に存在する。これは $\mathfrak{Z}$ 内の $\mathfrak{I}$ を含むすべての極大領域の共通部分として得られる。なぜなら $\mathfrak{Z}$ 自身がそのような領域であり、任意個のそのような領域の共通部分も極大で、かつ $\mathfrak{I}$ を含むからである。$M(\mathfrak{I})$ は、ある $a\ne0$ について $a\cdot g(x)$ が $\mathfrak{I}$ に属するような $\mathfrak{Z}$ のすべての関数 $g(x)$ から成る。

実際、これらの $g(x)$ は、$\mathfrak{I}$ を含む $\mathfrak{Z}$ 内の極大領域 $N$ をなす。というのも、$b\cdot h(x)$ が $N$ に属し、$b\ne0$、$h(x)\in\mathfrak{Z}$ なら、ある $a\ne0$ に対して $ab\cdot h(x)$ が $\mathfrak{I}$ に属し、$ab\ne0$ であるから、$h(x)$ は $N$ に属する。しかも $\mathfrak{I}$ を含む任意の極大領域 $M$ は $N$ を含む。なぜなら $a\cdot g(x)$ が $\mathfrak{I}$ に属すれば、それは $M$ にも属し、したがって $g(x)$ も $M$ に属するからである。

\medskip\noindent\textbf{2.}
いま $\mathfrak{I}$ を、単位元をもつ有限な整域と仮定し、$f_1(x),\ldots,f_s(x)$ をその整基底とする。すなわち $\mathfrak{I}$ は $f_i(x)$ のすべての整係数多項式から成る。さらに $\mathfrak{Z}$ は、$\mathfrak{Q}$ の整係数多項式または冪級数、あるいは有限個の $f_i(x)$ に現れる分母以外の分母を含まない商から成るものとする。ただし当然、それらの分母は任意の冪で現れてよい。

また、少なくとも一つの $f_i(x)$ --一般にはすべて-- に $x$ が実際に現れると仮定する。そうでなければ、$\mathfrak{I}$ と $\mathfrak{Z}$ は単に一つの分数 $1/c$ に関する整係数多項式全体から成り、証明すべきことはない。

多項式または有理関数の場合には、商体 $\mathfrak{Q}$ は有理数体 $K$ に関して超越次数 $t$（$1\le t\le r$）をもつ。例えば $f_1(x),\ldots,f_t(x)$ が既約系\footnote{「既約系」とは、周知のように、代数的独立な関数からなる系をいう。}をなすなら、$\mathfrak{Q}$ は純超越体 $\mathrm{K}(f_1(x),\ldots,f_t(x))$ の有限代数拡大である。同時に、$f_1(x),\ldots,f_t(x)$ の関数行列は正確に階数 $t$ をもつ。\footnote{この関数行列式定理の純代数的証明については --- ここで扱っているのは標数零である --- B. L. v. d. Waerden, \emph{Algebraische Theorie der Differentiation}, Nieuw Archief v. Wiskunde (2) 15 (1926), pp. 111--120 を参照。}

冪級数またはその商の場合にもこの $\mathfrak{Q}$ の構造を保つには、付加仮定が必要である。すなわち、$t$ が $f_i(x)$ の関数行列の階数であり、さらに $f_1(x),\ldots,f_t(x)$ の関数行列の階数でもあるなら（したがって \S\,1.1 によりこれらは既約系をなす）、$\mathfrak{Q}$ は再び純超越体 $\mathbf{K}(f_1(x),\ldots,f_t(x))$ の有限代数拡大であると仮定する。言い換えれば、$f_{t+1}(x),\ldots,f_r(x)$ は $\mathbf{K}(f_1(x),\ldots,f_t(x))$ に関して代数的でなければならない。\footnote{この付加仮定は、Hecke が扱った場合（注2）のように解析関数を扱い、$f_1(x),\ldots,f_t(x)$ が解析的に独立で、$f_{t+1}(x),\ldots,f_s(x)$ がこれらに代数的に依存する場合には満たされる。} このときも $t\le r$ である。なぜなら、特性零ではすべての $\partial f_i/\partial x_k$ が同時に消えることはないからである。

同時に、この条件は、関数行列が正確に階数 $t$ をもつような整基底の任意の $t$ 個の関数系について従う。というのも、それらは既約系をなしており、残りの関数はそれらに代数的に依存するからである。

\medskip\noindent\textbf{3.}
いま $\mathfrak{o}$ を不定元 $u_1,\ldots,u_s$ に関する整係数多項式領域とする。すると対応 $u_i\mapsto f_i(x)$ が直ちに示すように、$\mathfrak{I}$ は $\mathfrak{o}$ の環準同型像である。したがって $\mathfrak{I}$ は $\mathfrak{o}/\mathfrak{a}$ と環同型である。ここで $\mathfrak{a}$ は素イデアルである。なぜなら $\mathfrak{I}$ は零因子をもたない環だからである。$\mathfrak{a}$ は、$G(f)$ が $x$ において恒等的に消えるようなすべての多項式 $G(u)$ から成る。

導入部ですでに述べたように、問題は次の定理である。

\medskip\noindent\textbf{定理.}
\emph{2 の付加仮定のもとで、$\mathfrak{o}_p$, $\mathfrak{a}_p$, $\mathfrak{I}_p$ を、$\mathfrak{o}$, $\mathfrak{a}$, $\mathfrak{I}$ から整係数を法 $p$ の剰余類で置き換えて得られる領域とする。すると、有限個を超えない例外素数を除いて --この中には、$\mathfrak{I}_p$ が定義されるために、$f(x)$ の分母に現れるものが含まれる--、$\mathfrak{o}_p$ から $\mathfrak{I}_p$ への準同型は $\mathfrak{a}_p$ によって媒介される。したがって $\mathfrak{I}_p$ は $\mathfrak{o}_p/\mathfrak{a}_p$ と、さらに $\mathfrak{o}/(\mathfrak{a},p\mathfrak{o})$ と環同型になる。}

この定理から、$M(\mathfrak{I})$ の分母 $a$ に現れる素数が有限個だけであるという事実は、次の系によって得られる。

\medskip\noindent\textbf{系.}
\emph{$\mathfrak{I}_p$ が $\mathfrak{o}_p/\mathfrak{a}_p$ と同型になるように $p$ を選ぶなら、$p\cdot g(x)$ が $\mathfrak{I}$ に属し、$g(x)$ が $\mathfrak{Z}$ に属するとき、常に $g(x)$ は $\mathfrak{I}$ に属する。}

別の言い方をすれば、$\mathfrak{I}$ は、有限個の例外素数以外のすべての素数に関して、$\mathfrak{Z}$ 内で極大である。

$\mathfrak{o}_p/\mathfrak{a}_p$ は $\mathfrak{o}/(\mathfrak{a},p\mathfrak{o})$ と同型である。なぜなら、定義により $\mathfrak{o}_p$ は剰余類環 $\mathfrak{o}/p\mathfrak{o}$ であり、対応して $\mathfrak{a}_p$ は $(\mathfrak{a},p\mathfrak{o})/p\mathfrak{o}$ であるからである。したがって、第一同型定理により、$\mathfrak{o}_p/\mathfrak{a}_p$ は、仮定により $\mathfrak{I}_p$ とともに、$\mathfrak{o}/(\mathfrak{a},p\mathfrak{o})$ と同型になる。

これは次を意味する。$\mathfrak{I}$ において $H(f_i(x))\equiv0\pmod p$ なら、$\mathfrak{o}$ において $H(u)\equiv0\pmod{(\mathfrak{a},p\mathfrak{o})}$ である。

これは次のように言い換えられる。$H(f(x))=p\cdot g(x)$（$H(f(x))$ は $\mathfrak{I}$ に、$g(x)$ は $\mathfrak{Z}$ に属する）なら、$H(u_i)-pF(u_i)\equiv0\ (\mathfrak{a})$ が従い、したがって $H(f_i(x))=pF(f_i(x))$ であり、$g(x)=F(f_i(x))$ となる。ゆえに $g(x)$ は $\mathfrak{I}$ に属する。

有限回これを繰り返すと、$a$ がどの例外素数にも割られないなら、$a\cdot g(x)$ が $\mathfrak{I}$ に属し、$g(x)$ が $\mathfrak{Z}$ に属することから、常に $g(x)$ が $\mathfrak{I}$ に属することが従う。これで系は証明された。

\subsection*{\S\,3. 定理の証明}

\noindent\textbf{1. 補題.}
\emph{$\mathfrak{I}$ と $\mathfrak{o}$ の間の準同型を媒介する素イデアル $\mathfrak{a}$（\S\,2.3）は絶対素イデアルである。}

$\mathfrak{o}^*$ を、$u_1,\ldots,u_r$ に関する多項式領域であって、任意の代数数 --分数を含む-- を係数として許すものとする。また $\mathfrak{a}^*$ を、$\mathfrak{a}$ によって $\mathfrak{o}^*$ 内に生成されるイデアルとする。すなわち $\mathfrak{a}^*$ は、$\alpha_i$ が任意の代数数、$G_i(u)$ が $\mathfrak{a}$ の多項式であるような線型結合 $\alpha_1G_1(u)+\cdots+\alpha_sG_s(u)$ 全体から成る。証明すべきことは、$a^*$ が素イデアルであることである。

これは、$a^*$ が $\mathfrak{o}^*$ から $\mathfrak{I}^*$ への準同型を媒介することを示せばよい。ここで $\mathfrak{I}^*$ とは、$f_i(x)$ の多項式で係数として任意の代数数を許すもの全体からなる整域をいう。すなわち、$H(u)$ が $\mathfrak{o}^*$ に属し、$H(f_i(x))$ が消えるなら、$H(u)$ は $\mathfrak{a}^*$ に属することを示せばよい。

$H(u)$ の有限個の係数が生成する有限数体について、$\alpha_1,\ldots,\alpha_s$ を線型独立な基底とする。すると
\[
\alpha\cdot H(u)=\alpha_1H_1(u)+\cdots+\alpha_sH_s(u)
\]
となる。ただし $H_1(u),\ldots,H_s(u)$ は $\mathfrak{o}$ に属し、$\alpha\ne0$ は有理整数である。$H(f)=0$ から
$\omega_1 H_1(f) + \cdots + \omega_s H_s(f) = 0$ が従い、したがって $H_1(f)=0,\ldots,H_s(f)=0$ となる。よって $H_1(u),\ldots,H_s(u)$ は $\mathfrak{a}$ に属し、したがって $H(u)$ は $\mathfrak{a}^*$ に属する。

\medskip\noindent\textbf{2.}
定理の証明は、いまや次の形で行える。$p$ を次のように選んだとする。
\begin{enumerate}
\item[I.] $\mathfrak{I}_p$ が定義されている。
\item[II.] $\mathfrak{a}_p$ が素イデアル、しかも絶対素イデアルのままである。
\item[III.] $\mathfrak{I}_p$ の超越次数が $\mathfrak{I}$ のそれに比べて下がっていない。
\end{enumerate}
このとき $\mathfrak{o}_p$ から $\mathfrak{I}_p$ への準同型は $\mathfrak{a}_p$ によって媒介される。これらの条件は、有限個を超えない例外素数を除き、すべての素数について満たされる。

条件 I, II, III に関して例外素数が有限個だけであることは、前述から直ちに従う。I は明らかである。例外素数は $f_i(x)$ の分母に入る有限個だけである。II は、補題により $\mathfrak{a}$ が絶対素イデアルであり、この性質を有限個を超えない素数を除くすべての素数を法として保つことから従う（注6参照）。最後に III は \S\,1.2 の系で示されている。なぜなら、$f_1(x),\ldots,f_t(x)$ の関数行列は \S\,2.2 により正確に階数 $t$ だからである。

いま $p$ をこれらの例外素数とは異なるように選ぶ。$\mathfrak{o}_p$ から $\mathfrak{I}_p$ への準同型は、あるイデアル $\mathfrak{c}_p$ によって媒介される。このイデアルは素イデアルである。なぜなら $\mathfrak{I}_p$ も零因子をもたない環だからである。またその超越次数は $\mathfrak{I}_p$ の超越次数によって与えられ、したがって少なくとも $t$ である。実際、$\mathfrak{I}_p$ は $\mathfrak{o}_p/\mathfrak{c}_p$ と同型であるから、$\mathfrak{c}_p$ の剰余類体は $\mathfrak{I}_p$ の商体と同型になる。

さらに $\mathfrak{a}_p$ も素イデアルで、$\mathfrak{c}_p$ の倍イデアル（部分イデアル）であるから、$\mathfrak{a}_p$ と $\mathfrak{c}_p$ は、その超越次数が一致するとき、かつそのときに限り一致する。\footnote{例えば B. L. v. d. Waerden, ``Zur Nullstellentheorie der Polynomideale'', Math. Ann. 96 (1926), pp. 183--208, \S\,3 を参照。} $\mathfrak{c}_p$ の倍イデアルとしての $\mathfrak{a}_p$ の超越次数も少なくとも $t$ であり、しかも $u_1,\ldots,u_t$ の $\mathfrak{a}_p$ を法とする類は既約系をなす。この超越次数がちょうど $t$ であることを示せば、すべてが証明される。

ここで、冪級数またはその商の場合には、$\mathfrak{I}$ の商体 $\mathfrak{Q}$ の構造に関する付加仮定（\S\,2.2）を用いる。それによれば、素イデアル $\mathfrak{a}$ はいずれの場合にも超越次数 $t$ であった。しかも番号付けは、$u_1,\ldots,u_t$ の $\mathfrak{a}$ を法とする類が既約系をなし、$u_{t+1},\ldots,u_s$ の類がそれらに代数的に依存するように選ばれていた。したがって $\mathfrak{a}$ には、整係数かつ原始的な多項式
\[
G_1(u),\ldots,G_{s-t}(u)\tag{1}
\]
が含まれており、これらは、原始性により消えない多項式
\[
\bar G_1(u),\ldots,\bar G_{s-t}(u)\tag{2}
\]
として $\mathfrak{a}_p$ に移る。しかし各 $\bar G_j$ には $u_{t+j}$ が実際に現れなければならない。そうでなければ、$p$ の選び方に反して、$u_1,\ldots,u_t$ の $\mathfrak{a}_p$ を法とする類の間に代数的依存関係が存在したことになるからである。こうして $\mathfrak{a}_p$ が超越次数 $t$ をもつことが示された。したがって $\mathfrak{a}_p$ は $\mathfrak{c}_p$ と同一であり、その準同型を媒介する。これで証明は完了する。

\subsection*{\S\,4. 全代数的係数への拡張}

導入部の終わりですでに述べたように、すべての考察は、有理整数の代わりに（有限な）代数的数体 $\mathbf{K}$ の整数を、また素数 $p$ の代わりにその数体の素イデアル $\mathfrak{p}$ を基礎に置いても、わずかな修正で保たれる。このとき \S\,1 は完全にそのまま残り、\S\,2 における定理の定式化へ導く考察も同様である。ただし証明（\S\,3.2）および \S\,2.3 の系にはいくつかの補足が必要になる。

\medskip\noindent\textbf{1.}
\S\,3.2 では、さらに有限個の素イデアル、すなわち多項式 $G_1(u),\ldots,G_{r-t}(u)$（\S\,3.2, (1)）を割るものを例外素イデアルとして加えなければならない。というのも、いまやこれらを原始的と仮定することはできないからである。これによって $\bar G_1(u),\ldots,\bar G_{r-t}(u)$（\S\,3.2, (2)）が消えないことが保証される。これは条件 III の強化である。条件 II は保たれる。絶対既約性に関する定理は変わらないからである。もちろん I も保たれる。

\medskip\noindent\textbf{2.}
最終結果を表す \S\,2.3 の系は、そこでの最終形に対応して次のように定式化される。

\noindent\textbf{系.}
\emph{$K$ の整数 $\alpha$ が有限個の例外素イデアルのいずれによっても割られないなら、$\alpha\cdot g(x)$ が $\mathfrak{I}$ に属し、$g(x)$ が $\mathfrak{Z}$ に属することから、常に $g(x)$ が $\mathfrak{I}$ に属することが従う。}

$K$ の整数 $\lambda_1,\ldots,\lambda_s$ を、それぞれ一つの例外素イデアルで割られるが、その平方でも他の例外素イデアルでも割られないように選べば、分母 $\alpha$ としてはこれらの $\lambda$ の冪積で十分である。

\noindent\emph{証明.}
$(\alpha) = \mathfrak{p}_1^{\varrho_1} \ldots \mathfrak{p}_g^{\varrho_g}$ を $\alpha$ の素イデアル分解とし、$\mathfrak{p}_1,\ldots,\mathfrak{p}_g$ は例外素イデアルではないとする。したがって各 $\mathfrak{p}_i$ について $\mathfrak{I}_p$ は $\mathfrak{o}_p/\mathfrak{a}_p$ と同型である。$\mathbf{K}$ のすべての整数の環から、$(\alpha)$ によって定められる商環 $\mathfrak R$ へ移る。これは、$(\alpha)$ と素、すなわち $\mathfrak{p}_1,\ldots,\mathfrak{p}_g$ のいずれによっても割られない整数を、かつそれらだけを分母に加えることによって得られる。$\mathfrak R$ では $\mathfrak{p}_1,\ldots,\mathfrak{p}_g$ は素イデアルのままであるが主イデアルになり、他のすべての素イデアルは単位イデアルになる。また $(\alpha)$ の分解は同じまま残る。\footnote{例えば H. Grell, ``Zur Theorie der Ordnungen in algebraischen Zahl- und Funktionenkörpern'', Math. Ann. 97 (1927), pp. 524--558, \S\,2.1 を参照。} したがって、係数領域を $\mathfrak R$ に移しても、系とその証明は \S\,2.3 の形そのままで保たれる。

つまり、$\alpha\cdot g(x)$ が $\mathfrak{I}$ に属するなら、$g(x)=F(f_i(x))$ となり、ここで多項式 $F(u)$ の係数は $\mathfrak R$ に属する。これらの係数の共通分母を $\delta$ とすれば、$\delta$ は $\alpha$ と素である。したがって次の形にも言い換えられる。$\alpha\cdot g(x)$ が $\mathfrak{I}$ に属すれば、$\alpha$ と素な $\delta$ について $\delta\cdot g(x)$ が $\mathfrak{I}$ に属する。ゆえに $(\delta,\alpha)$ は単位イデアルであり、その結果、$\alpha\cdot g(x)$ が $\mathfrak{I}$ に属することから $g(x)$ が $\mathfrak{I}$ に属することが従う。これで上の系の第一主張が証明された。

第二主張を証明するため、$\mathfrak R^*$ を有限個の例外素イデアルによって $\mathrm{K}$ 内に定められる商環とし、$\lambda_1,\ldots,\lambda_s$ を $\mathfrak R^*$ におけるこれらの素イデアルの基底元とする。これらは整数として取ってよい。$\mathfrak R^*$ では、$\alpha$ は $\mathfrak R^*$ の単元を除いて $\lambda$ の冪積になる。整数へ戻ると
\[
\gamma\alpha=\beta\lambda_1^{\varrho_1}\cdots\lambda_s^{\varrho_s}
\]
が得られる。ただし $\gamma$ と $\beta$ は、どの例外素イデアルによっても割られない。$\alpha\cdot g(x)$ が $\mathfrak{I}$ に属するなら、$\beta\lambda_1^{\varrho_1}\cdots\lambda_s^{\varrho_s}g(x)$ が $\mathfrak{I}$ に属し、さらに系の第一主張により $\lambda_1^{\sigma_1} \ldots \lambda_s^{\sigma_s} g(x)$ も $\mathfrak{I}$ に属する。
\begin{center}
（1931年11月11日受理）
\end{center}

\bigskip
\noindent\rule{\textwidth}{0.4pt}
\medskip
\noindent\textbf{英語要旨}（原文のロシア語要旨からの翻訳）:

\smallskip
\emph{任意の整域 $\mathfrak{I}$（Integritätsbereich、零因子をもたない環）は、同じ種類のある極大領域 $M(\mathfrak{I})$ を定める。それは、ある $a\ne0$ について $a\cdot g(x)$ が $\mathfrak{I}$ に含まれるようなすべての関数 $g(x)$ から成る。}

\emph{本論文では、この極大領域 $M(\mathfrak{I})$ の「分母」$a$ の問題を、領域 $\mathfrak{I}$ が有限であったという仮定のもとで考察する。本論文で証明されるのは、これらの $a$ を、そこに因子として現れる素数が全体として有限個だけになるように常に選べる、ということである。}

\emph{同時に、$g(x)$ が冪級数またはそのような級数の商である場合には、領域 $\mathfrak{I}$ を定義する関数 $f_1(x),\ldots,f_r(x)$ が相互に代数関係だけで結ばれ得るものと仮定する。}

\emph{証明は、問題を一般イデアル論のいくつかの問題へ帰着することに基づく。この帰着は、分母 $a$ に素数 $p$ が現れることを、関数 $f_i$ の間の法 $p$ の関係と見なすことによって行われる。}

\begin{center}
(Rec. Math. XXXVI:1; 1929).
\end{center}



\clearpage
\setcounter{footnote}{0}
\clearpage
\section*{36. イデアル微分と共役差積}

\begin{center}
\emph{J. Ber. d. DMV 39 (1930), S. 17}
\end{center}

2. E. Noether，ゲッティンゲン：イデアル微分と共役差積。

代数的数体の共役差積を，定義イデアルの\emph{微分商}として捉える方法を示す。この定義が通常の定義と一致することは，それ自体興味ある構造定理から従う。これらの構造定理は，類似の意味でラグランジュの補間公式の一般化をなす。

詳しい叙述は \emph{Mathematische Annalen} に掲載される予定である。

\section*{37. 高次分岐をもたない体における正規基底}

\begin{center}
\emph{Journal f. d. reine u. angew. Math. 167 (1932), pp. 147--152}

\textsc{Emmy Noether}（Göttingen）
\end{center}

ガロア数体 \(K/k\) の正規基底とは、\(K\) の主整環 \(\frakO\) が \(k\) の主整環 \(\frako\) に関してもつ、共役元からなる基底をいう。そのような基底が存在するための必要条件は、よく知られているように、\footnote{A. Speiser, Gruppendeterminante und Körperdiskriminante, § 6. Math. Ann. 77 (1916), pp. 546--562 を見よ。}\(K/k\) に高次分岐体が現れないことである。これは一つの素点に限って考えてもなお成り立つ。すなわち、\(k\) の \(\frakp\)-進拡大 \(k_\frakp\) と、それに対応する \(K\) の \(K_\frakp\) に移っても成り立つ。以下ではその逆が成り立つことを示す。

\emph{\(\frakp\) が \(K/k\) の次数を割らないすべての素点 \(\frakp\) について、正規基底が存在する。したがって特に、\(K/k\) が高次分岐をもたなければ、各分岐素点で正規基底が存在し、そこで判別式は群行列式の平方となる。}

この最後の主張には、E. Artin の予想に反して、彼の導手へ移るには\footnote{E. Artin, Über die gruppentheoretische Struktur der Diskriminante. J. f. Math. 164 (1931), pp. 1--11.}さらに考察が必要であることを付け加えなければならない。群行列式を既約指標に従って分解すると一般化された根数が得られ、これらを適切にまとめると、指標によって拡大された基礎体における分解が得られる。最も単純な場合には、根数のイデアル論的分解によって、これらの新しい因子を Artin の導手と同一視できる。\footnote{[1931年10月6日] 一般の場合にもなおイデアル論的考察が必要になることは、M. Deuring から教えられた次の例、しかも任意に変形できる例、によって示される。これは、判別式が群行列式の平方として表せる非極大整環でありながら、共役指標に異なるイデアル因子が対応する例である。\(k\) を五次単位根の体、\(K=k(\sqrt[5]{2})\) とする。整環 \([1,2\sqrt[5]{2},\sqrt[5]{2^2},\sqrt[5]{2^3},\sqrt[5]{2^4}]\) を考える。場所は \((2)\) であり、この注の考察によりそこでは正規基底をもつ。指標による分解は、対応する群行列式に対して、ちょうど上の基底元を因子として与える。したがって「導手」は、\(1\) を除けば \(2\sqrt[5]{2}\cdot\sqrt[5]{2^4}\), \(\sqrt[5]{2^2}\cdot\sqrt[5]{2^3}\), \(\sqrt[5]{2^3}\cdot\sqrt[5]{2^2}\), \(\sqrt[5]{2^4}\cdot(2\sqrt[5]{2})\)、すなわち \(4,2,2,4\) となり、共役指標に対して異なる。} さらに、正規基底が存在しない一般の場合にも、「ガロア加群による組成列」によって一般化された根数への分裂は常に可能であり、そこから上と同様に拡大された基礎体での分解が得られることを述べておきたい。ここで再びイデアル論的問題が始まる。

上の仮定のもとで正規基底が存在することの証明は、\(\frakO/\frako\) をガロア加群、すなわちガロア群の置換を作用域とする \(\frako\)-加群、とみなすと、それが \(\frako\) で拡張された整数群環の片側イデアルに作用素同型となる、という事実に基づく。この群環は通常の分岐素点ではすべて極大整環を与える。\(\frakp\)-進半単純系の極大整環におけるイデアルは主イデアルであるから、\footnote{H. Hasse, Über \(\frakp\)-adische Schiefkörper und ihre Bedeutung für die Arithmetik hyperkomplexer Zahlsysteme, 定理 46, 47. Math. Ann. 104 (1931), pp. 495--534 を見よ。可換系では、主イデアル性は \(k\) における \(\frakp\) に関する商環へ移るだけですでに成り立つ。したがって、アーベル群およびアーベル体を扱う場合には、素点を定義するためにこの弱い拡大に制限してよい。}ここでは群または群環の置換により一つの元から生じる。これをガロア加群へ戻して読み替えると、ちょうど正規基底の存在が得られる。高次分岐素点では、整数群環は非極大整環となり、\(\frakO/\frako\) に対応するイデアルは非主イデアルとなる。

これらの考察は、定数体の標数が \(K/k\) の次数を割らない限り、一変数関数体の場合にも明らかに保たれる。

\subsection*{§1. \(\frakp\)-進的に拡張された整数群環}

\(\frako\) および \(\frako_\frakp\) は、それぞれ代数的数体 \(k\) とその \(\frakp\)-進拡大 \(k_\frakp\) の主整環を表すものとする。ここで \(\frakp\) は \(k\) の素イデアルである。また、位数 \(n\) の群 \(\Gg\) に対し、\((\Gg)\) は有理整数係数の群環、\([\Gg]\) は整数群環を表す。\((\Gg)_k\) および \((\Gg)_{k_\frakp}\) は、それぞれ \(k\) および \(k_\frakp\) 係数の群環を意味し、同様に \([\Gg]_\frako\) および \([\Gg]_{\frako_\frakp}\) は、整数群環を \(\frako\) および \(\frako_\frakp\) 係数へ拡張したものを意味する。

したがって、\((\Gg)_k\) および \((\Gg)_{k_\frakp}\) は根基をもたない系（半単純系）であり、\([\Gg]_\frako\) および \([\Gg]_{\frako_\frakp}\) はそれぞれ \((\Gg)_k\) および \((\Gg)_{k_\frakp}\) の整環である。

\paragraph{定理 1.}
\emph{素イデアル \(\frakp\) が \(\Gg\) の位数 \(n\) を割らなければ、\([\Gg]_{\frako_\frakp}\) は \((\Gg)_{k_\frakp}\) の極大整環となる。\(\frakp\) が \(n\) を割るならば、\([\Gg]_{\frako_\frakp}\) は極大でない。}

整数群環の判別式、したがって \([\Gg]_\frako\) および \([\Gg]_{\frako_\frakp}\) の判別式も、位数 \(n\) のある冪に等しい。\footnote{例えば E. Noether, Hyperkomplexe Größen und Darstellungstheorie, § 26. Math. Ztschr. 30 (1929), pp. 641--692 を見よ。} したがって \(\frakp\) が \(n\) を割らなければ、\([\Gg]_{\frako_\frakp}\) の判別式は単元である。これは、\(\frako_\frakp\) を固定したまま \([\Gg]_{\frako_\frakp}\) をより大きな整環へ拡張できないことを意味する。なぜなら、そのような拡張には、判別式から平方の非単元因子を分離することが対応するからである。しかも \(\frako_\frakp\) はすでに主整環、すなわち \(k_\frakp\) において極大、と仮定されていた。これで、以下で唯一用いられる定理 1 の第一部が証明された。

これに対して \(\frakp\) が \(n\) を割るならば、単位表現に対応する直和
\[
        \mathfrak C=E^{(1)}[\Gg]_{\frako_\frakp}+(1-E^{(1)})[\Gg]_{\frako_\frakp},
        \qquad
        E^{(1)}=\frac1n\sum_{S\in\Gg} S,
\]
は \([\Gg]_{\frako_\frakp}\) の真の拡大である。実際、\(E^{(1)}=\frac1n\sum S\) であるから、\(\mathfrak C\) は真の拡大となる。\(\mathfrak C\) は従属生成元 \(E^{(1)}\) および \((1-E^{(1)})S\)（\(S\in\Gg\), \(E^{(1)}S=E^{(1)}\)）をもつ整環である。

\paragraph{定理 2.}
\emph{\(\frakp\) が \(n\) を割らなければ、\([\Gg]_{\frako_\frakp}\) に関する任意の整イデアルまたは分数イデアルは主イデアルである。}

なぜなら、定理 1 により \([\Gg]_{\frako_\frakp}\) は \(\frakp\)-進半単純系の極大整環であり、そのイデアルは主イデアルだからである。\footnote{Hasse, 前掲。そこでは主イデアル性は単純系についてだけ証明されているが、既知の議論により半単純系にも従う。極大整環の成分は再び極大であり、したがって考えているイデアルの成分は、たとえば基底 \(a_i\) をもつ主イデアルである。このとき \(a=\sum a_i\) はもとのイデアルの基底である。アーベル群については注 3 も参照。}

\subsection*{§2.ガロア加群、作用素同型、正規基底}

\(K/k\) をガロア拡大とし、\(\Gg\) をその群とする。\(z^S\) は、\(\Gg\) の置換 \(S\) によって \(z\) から得られる \(K\) の元を表す。\(\Gg\) の置換は、\(K/k\) を \(k\)-加群、すなわち \(k\) を作用域とする加法アーベル群、と見たものの上に作用素自己同型を生じる。したがって \(K/k\) は、次の規則によって \((\Gg)_k\) に関する加群とみなせる。
\[
        z\Bigl(\sum_i S_i c_i\Bigr)=\sum_i z^{S_i}c_i,
        \qquad z\in K,\; S_i\in\Gg,\; c_i\in k.
\]

\paragraph{定義.}
\emph{\(K/k\) から得られる任意の \((\Gg)_k\)-加群を有理ガロア加群という。同様に、\(K/k\) から得られる \([\Gg]_\frako\)-加群を整ガロア加群という。特に \(\frakO/\frako\) は整ガロア加群である。}

\paragraph{定理 3.}
\emph{有理ガロア加群として、\(K/k\) は \((\Gg)_k\) に作用素同型である。}

これは、\(K/k\) が常に体としての正規基底、すなわち共役元 \(z^S\) からなる \(k\)-基底をもつことから従う。\footnote{\(a_1,\ldots,a_n\) を \(K/k\) の基底とし、\(u_i\) を不定元とする。このとき \(S\) が \(\Gg\) の元を動くと、\(\sum u_i a_i^S\) は線型独立である。したがって、線型独立性が保たれるように \(u_i\) を \(k\) の元に特殊化できる。} 対応
\[
        S\to z^S,
        \qquad
        \sum S_i c_i\to\sum z^{S_i}c_i
        \quad(c_i\in k),
        \quad\hbox{したがって } ST\to z^{ST}
\]
は作用素準同型を与え、\(k\) 上の階数が等しいので同型となる。

\paragraph{定理 4.}
\emph{整ガロア加群として、\(\frakO/\frako\) は \((\Gg)_k\) から得られる最高階数 \(n\) の \([\Gg]_\frako\)-加群に作用素同型である。同様に、\(\frakO_\frakp/\frako_\frakp\) は \((\Gg)_{k_\frakp}\) から得られる階数 \(n\) の \([\Gg]_{\frako_\frakp}\)-加群に作用素同型である。}

定理 3 で与えられた \((\Gg)_k\)-同型 \(K/k\to(\Gg)_k\) は、\([\Gg]_\frako\)-加群 \(\frakO/\frako\) に \((\Gg)_k\) の中の階数 \(n\) の \([\Gg]_\frako\)-加群を対応させる。さらに、定理 3 の対応における係数 \(c_i\) に \(k_\frakp\) のすべての元を走らせることにより、この \((\Gg)_k\)-同型は \(K_\frakp/k_\frakp\) から \((\Gg)_{k_\frakp}\) への \((\Gg)_{k_\frakp}\)-同型に拡張される。この同型は \(\frakO_\frakp/\frako_\frakp\) の \([\Gg]_{\frako_\frakp}\)-同型を誘導する。このとき \(K_\frakp/k_\frakp\) は \(k_\frakp\) 上の超複素系と見なすべきである。これは \(K/k\) から係数拡大によって生じる。一般には零因子をもつ系、すなわち同型な体の和となり、したがって半単純系となる。

\paragraph{定理 5.}
\emph{\(K/k\) の次数 \(n\) を割らない任意の素点 \(\frakp\) において、\(\frakO_\frakp/\frako_\frakp\) は正規基底をもつ。}\footnote{アーベル体の場合には、注 3 および 5 に従って、商環によっても素点を定義できる。}

定理 1 により、そこで \([\Gg]_{\frako_\frakp}\) は極大整環である。したがって \((\Gg)_{k_\frakp}\) から得られる階数 \(n\) の \([\Gg]_{\frako_\frakp}\)-加群は整イデアルまたは分数イデアルであり、定理 2 により主イデアルである。\(\frakO_\frakp/\frako_\frakp\) に対応するイデアルが \(W[\Gg]_{\frako_\frakp}\) であるなら、それは \(\frako_\frakp\)-基底 \(WS_i\) をもつ。作用素同型は、\(W\) に対応する \(\frakO_\frakp/\frako_\frakp\) の元を \(w\) とするとき、\(w^{S_i}\) が \(\frakO_\frakp/\frako_\frakp\) の \(\frako_\frakp\)-基底であることを述べている。したがって \(n\) 個の共役元 \(w^{S_i}\) は \(\frakO_\frakp/\frako_\frakp\) の正規基底をなす。

\(\frakp\) が \(n\) を割る場合、\(\frakO_\frakp/\frako_\frakp\) に対応する加群は主加群でありえない。なぜなら、その場合には正規基底は存在しえないからである。

\paragraph{定理 3 への補足。}
定理 3 で証明された \((\Gg)_k\) と \(K/k\) の作用素同型から、\footnote{証明は明らかに、\(K\) が \(k\) 上第一種のガロア拡大であり、かつ \(k\) が無限個の元をもつことだけを仮定している。この後者の制限は不要である。M. Deuring は、有限個の元をもつ体にも有効な定理 3 の証明を考え、それから逆に体の正規基底の存在が従う。同時に、有限体および無限体に対して一様に働く原始元の存在証明も得られる。} Klein の形式問題に関する Speiser の結果が従う。\footnote{前掲。ガロア加群の概念はそこから抽象されたものである。\(k\) の標数が \(n\) を割る場合には、ガロア加群による組成列とその既約因子を扱うことになる。} それは次のように定式化される。

\emph{\(K/k\) のガロア群 \(\Gg\) の \(k\) 上既約な表現は、\(K/k\) から得られる既約有理ガロア加群によって生成される。絶対既約表現は、\(K_Z/Z\) から得られるガロア加群によって生成される。}

ここで \(Z\) は、群環の対応する成分の \(k\) を含む分解体を意味する。\(K_Z/Z\) は \(Z\) 上の超複素系として見なされ、\(K/k\) から係数拡大によって生じる。特に、\(Z\) を \(k\) に関して \(K\) と素に選べるなら、\(K_Z\) は体のままであり、\(K_Z/Z\) は Speiser が扱った場合、すなわち付随的拡大となる。

主張そのものは、作用素同型により、\((\Gg)_k\) または \((\Gg)_Z\) に関する対応する主張から従う。そこでは既約イデアルが表現を生成する。

\(v_1,\ldots,v_t\) をそのようなガロア加群の基底とし、\(S\mapsto\bar S\) を対応する表現とすると、
\[
        \begin{pmatrix}\vdots\\ v_i^S\\ \vdots\end{pmatrix}
        =\bar S\begin{pmatrix}\vdots\\ v_i\\ \vdots\end{pmatrix}
\]
が成り立つ。これが Klein と Speiser の定式化である。

\subsection*{§3. 高次分岐をもたない体における群行列式としての判別式}

判別式の分解には、よく知られているように、各分岐素点での分解を考えれば十分である。高次分岐をもたない体では、正規基底をもつ素点だけを扱うことになる。

\paragraph{定理 6.}
\emph{高次分岐をもたないガロア体 \(K/k\) では、各分岐素点において判別式は、ガロア群の群行列式に正規基底を不定元の代わりに代入したものの平方に等しい。既約表現に応じて群行列式は一般化された根数に分解し、判別式は指標によって拡大された基礎体のイデアルに分解する。このとき複素共役指標には同じイデアルが対応する。}

\(w^S\) を正規基底とすると、判別式 \(\Delta\) は
\[
        D=\bigl|w^{ST^{-1}}\bigr|,
        \qquad S,T\in\Gg
\]
の平方に等しい。\(D\) は、不定元 \(x_S\) の代わりに \(w^S\) を入れた群行列式である。

根数への分解は、Speiser 前掲の議論から得られる。
\[
        D=D_1^{f_1}\cdots D_t^{f_t},
        \qquad
        M=f_1M_1+\cdots+f_tM_t
\]
を、絶対既約表現に対応する群行列式および群行列の分解とする。\(M_\lambda\) に対応する表現を \(S\mapsto\bar S\) とし、\(\bar S\) が拡大された基礎体に属するとすると、
\[
        M_\lambda=\sum_S w^S\bar S,
        \qquad
        M_\lambda^{T^{-1}}=\sum_S w^{ST^{-1}}\bar S
          =\sum_R w^R\bar R\bar T=M_\lambda\bar T,
\]
したがって行列式へ移ると
\[
        D_\lambda^{T^{-1}}=D_\lambda|\bar T|=D_\lambda\varepsilon_T.
\]
ゆえに \(D_\lambda\) は一般化された根数である。\(\varepsilon_T\) は \(\bar T\) の行列式として、\(\Gg\) の交換子群による商群の既約表現、すなわち単位根になる。循環群の場合、\(D_\lambda\) は単に Lagrange の根数 \(\sum w^S\chi_\lambda(S)\) である。

一般に \(D_\lambda\) は、\(K_\frakp/k_\frakp\) から、係数体 \(k_\frakp\) を対応する指標で拡大して得られる超複素系に属し、しかもその系の整元である。実際、不定元 \(x_S\) で作った行列式は、指標体の代数的整数を係数にもつ。一方、\(w^S\) は \(K_\frakp\) の整元である。

群行列式の分解から判別式の分解へ移るために、各表現をその随伴表現とまとめる。\(\lambda,\bar\lambda\) を随伴表現とすると、同時に
\[
        D_\lambda^{T^{-1}}=D_\lambda\varepsilon_T,
        \qquad
        D_{\bar\lambda}^{T^{-1}}=D_{\bar\lambda}\varepsilon_T^{-1}
\]
が成り立つ。また、不定元 \(x_S\) で作られた \(\lambda\) と \(\bar\lambda\) に属する行列式は複素共役であり、一般に共役指標には、不定元 \(x_S\) のもとで共役な行列式が対応する。

そこで \(\Delta_\lambda=D_\lambda D_{\bar\lambda}\) とおくと、\(\lambda\) 番目の指標の実部分体だけを付加すればよく、
\[
        \Delta_\lambda^T=\Delta_\lambda
        \qquad\hbox{すべての }T\hbox{ について}\Gg
\]
が成り立つ。したがって\footnote{これは体の超複素係数拡大を扱っているので、通常のガロア理論の結論は修正されなければならない。\(P\) を考えている \(k_\frakp\) の拡大体とし、
\[
        (K_\frakp)_{\mathsf{P}}=(K_\frakp)_{\mathsf{P}} e^{S_1}+\cdots+(K_\frakp)_{\mathsf{P}} e^{S_r}
\]
を体への直和分解とする。このとき \((K_\frakp)_P e^{S_i}/Pe^{S_i}\) はガロア拡大であり（その群は \(\frakp\) の一つの素因子の分解群の部分群である）、\(e^{S_i}\) は剰余類群に関して共役である。\(\Delta_\lambda^T=\Delta_\lambda\) から、まず \(\Delta_\lambda\) の第 \(i\) 成分は \(Pe^{S_i}\) に属し、\(P\) の元 \(\gamma_i\) を用いて \(\gamma_i e^{S_i}\) と書ける。さらに \(e^{S_i}\) の共役性からすべての \(\gamma_i\) は等しく、実際 \(\Delta_\lambda=\gamma\sum e^{S_i}=\gamma\) は \(P\) に属する。}: \(\Delta_\lambda\) は、\(\lambda\) 番目の指標の実部分体で拡大された基礎体に属する。

\emph{判別式の分解}
\[
        \Delta=\Delta_1^{f_1}\cdots\Delta_t^{f_t}
\]
\emph{は、したがって、指標の実部分体によって拡大された基礎体における分解であり、複素共役指標には同じ因子が対応する。その他の共役指標については、さらなることはそのままでは結論できない。}

これで定理 6 はすべての部分で証明された。\(\Delta_\lambda\) が生成するイデアルが実際にはすでに基礎体 \(k\) に属し、共役指標に対して等しいことが示せれば、合成指標に対応する \(\Delta_\lambda\) の冪積を割り当てることにより、それらは Artin の導手と一致する。というのは、関数等式も、また正規基底から直接従う事実（Speiser 前掲参照）、すなわち部分群の単位表現には対応する部分体の判別式が対応する、という事実も成り立つからである。各素点での一致から、完全なイデアルの一致が従う。有理既約指標に限れば、述べた事実により一致が得られる。絶対既約指標については、最も単純な場合に直接確認しておく。

\paragraph{定理 7.}
\emph{\(k\) を有理数体とし、\(K/k\) を判別式に素な素数次数 \(l\) の循環拡大、したがって高次分岐をもたない拡大とする。このとき、単位表現でない表現に対する \(\Delta_\lambda\) から生成されるイデアルは互いに等しく、\(K/k\) の導手と一致する。}

これは根数の既知の分解から従う。\footnote{Hilbert, Zahlbericht, § 108。定理 5 により正規基底の存在は確立している。素点は注 7 に従って商環で簡単に定義できるので、Zahlbericht で用いられた、絶対循環体が円分体であるという事実は、もはや援用する必要がない。}\(K/k\) の分岐素点 \(\frakp\) において、\(k_\varepsilon\) を \(l\) 次単位根の体、\(K_\varepsilon\) を対応する体とする（\(k_\varepsilon\) は \(K\) と素であり、注 7 によりここでは \(\frakp\)-進拡大は不要なので、\(K_\varepsilon\) は体のままである）。このとき
\[
        (p)=\frakp_1\cdots\frakp_{l-1},
        \qquad
        \frakp_i=\mathfrak P_i^{\,l}
\]
であり、\(\frakp\) は \(k_\varepsilon\) に、\(\mathfrak P\) は \(K_\varepsilon\) にある。さらに根数について
\[
        (\Omega_\lambda)=\mathfrak P_1^{r_1}\cdots\mathfrak P_{l-1}^{r_{l-1}},
        \qquad
        (\Omega_{\bar\lambda})=
        \mathfrak P_1^{l-r_1}\cdots\mathfrak P_{l-1}^{l-r_{l-1}},
\]
が成り立つ。ここで \(r_1,\ldots,r_{l-1}\) は \(1,\ldots,l-1\) の置換である。しかしこの場合、根数 \(\Omega_\lambda\) は群行列式の因子 \(D_\lambda\) と同一である。したがって
\[
        (\Delta_\lambda)=(D_\lambda D_{\bar\lambda})
        =(\Omega_\lambda\Omega_{\bar\lambda})
        =\mathfrak P_1^{l}\cdots\mathfrak P_{l-1}^{l}
        =\frakp_1\cdots\frakp_{l-1}=(p)
\]
となり、\(K/k\) の導手と一致する。

\begin{center}
1931年8月24日受理。
\end{center}


\clearpage

\section*{38. R. Brauer および H. Hasse との共著：代数論における主定理の証明}

\begin{center}
\emph{Journal f. d. reine u. angew. Math. 167 (1932), pp. 399--404}

\begin{center}
\textbf{代数論における主定理の証明}\\
R. Brauer（Königsberg）、H. Hasse（Marburg）、E. Noether（Göttingen）\footnote{本ノートの執筆は H. Hasse が担当した。}
\end{center}
\end{center}

われわれ三人の共同の努力により，ついに次の定理を証明することができた。この定理は，代数的数体上の代数の構造論にとって，またそれを越えて，基本的な重要性をもつ。

\paragraph{主定理.}
代数的数体上の任意の正規斜体は巡回的である。言い換えれば，Dickson 型である。

この結果は本質的に \(p\)-進的方法に負う成果であり，その方法の創始者 Kurt Hensel の七十歳の誕生日に献呈できることは，われわれにとって特別の喜びである。

証明は三つの還元から成り，それぞれはわれわれの一人によって与えられた。\footnote{三つの還元は成立した順に述べる。この順序は体系的な順序とは逆である。}

\paragraph{1.}
第一の還元は，H. Hasse が，代数的数体上の巡回代数について最近展開した理論に基づいて与えたものである。\footnote{H. Hasse, \emph{Theorie der zyklischen Algebren über einem algebraischen Zahlkörper}, Gött. Nachr. 1931. 証明の詳しい叙述は，E. Noether が講義で展開した分解体および接合積の理論をも含めて，Trans. Amer. Math. Soc. に \emph{Theory of cyclic algebras over an algebraic number field} として現れる予定である。以下ではこれを H と引用する。H, 1--6 は，言語が異なる点を除き，前記のノートに相当する。}

\paragraph{還元 1.}
主定理は，次を示せば証明される：
\[
\text{I. }\Omega\text{ 上で至る所分解する任意の代数は }\sim\Omega\text{ である。}
\]
ここで，以下「\(\Omega\) 上の代数 \(A\)」とは常に「代数的数体 \(\Omega\) 上の正規単純代数 \(A\)」，すなわち中心 \(\Omega\) をもつ単純超複素系を意味する。さらに \(A\sim\Omega\) は，\(A\) が \(\Omega\) 上の全行列代数であること，すなわち相似の意味で，\(\Omega\) 自身を \(\Omega\) 上の斜体と見ることで定まる特殊な類に属することを表す。また，\(\Omega\) 上で至る所分解する代数とは，\(\Omega\) の各素点 \(\mathfrak p\) について，係数体 \(\Omega\) を対応する \(\mathfrak p\)-進体 \(\Omega_{\mathfrak p}\) へ拡大して得られる \(\mathfrak p\)-進拡大が
\[
A_{\mathfrak p}\sim \Omega_{\mathfrak p}
\]
を満たす代数である。

\paragraph{証明.}
\(D\) を \(\Omega\) 上の斜体とする。\(Z/\Omega\) が巡回体で，各素点 \(\mathfrak p\) において \(Z\) の \(\mathfrak p\)-次数が \(D\) の \(\mathfrak p\)-指数の倍数であるとする。このとき，\(Z\) の対応する素因子 \(\mathfrak P\) に対し，\(\mathfrak P\)-進体 \(Z_{\mathfrak P}\) は \(D_{\mathfrak p}\) の分解体である [H, 18.1]。したがって \(D\) の係数体を \(Z\) に拡大して得られる \(D_Z\) は \(Z\) 上で至る所分解する。I がすべての \(\Omega\) について既知であれば，\(D_Z\sim Z\) となる。これは \(D\) が巡回分解体 \(Z\) をもつ，すなわち巡回的に表せることを意味する [H, 5]。すると \(D\) はさらに，次数が \(D\) 自身の次数と一致する巡回分解体をもつ。従って \(D\) は巡回的である [H, 6, 定理 6]。ゆえに I のもとで主定理が従う。

\paragraph{2.}
第二の還元への決定的な契機は R. Brauer によって与えられた。彼は H. Hasse に宛てた書簡で，斜体の指数の正確な値という関連問題が，Sylow の群定理により，可解な分解体をもつ場合へ還元されることを示した。この問題も主定理によって解決される。その考えは，巡回性の問題についても対応するさらなる還元に用いることができた。\footnote{その間に A. A. Albert も関連する結果に到達したことを知らせてきた。すでに印刷中であったノート \emph{On the construction of Riemann matrices, II}, Ann. of Math. 32 (1931) には，a) \(\Omega\) 上の次数 \(p^\alpha\) の任意の斜体が \(p\) と素な次数の分解体をもつこと，b) 後の 定理 1，すなわちブラウア類の位数（exponent）\(=\)指数（index），c) 還元 2 の基本思想に加えて還元 3 の基本思想が，もちろん還元 1 との関係なしに，含まれている。従ってそこでは，次数 \(p^\alpha\) の斜体 \(D/\Omega\) に対し，\(\Omega\) 上 \(p\) と素な次数の拡大 \(\Omega'\) が存在して \(D_{\Omega'}\) が巡回的となるという結果が得られる。これら三つの結果は，その後のわれわれの主定理の証明により当然包含されるが，A. A. Albert にも主定理の証明への独立した寄与があることを示している。さらに Albert は，われわれの証明を知った後，中心的な 定理 I が，印刷中の彼の論文 Bull. Amer. Math. Soc. 37 (1931) の定理 13, 10, 9 から数行で従うことを指摘した。それらの証明は本質的に，われわれの還元 2, 3 と同じ推論に基づいている。}

\paragraph{還元 2.}
定理 I は次を示せば証明される：
\[
\text{II. 可解な分解体をもち，至る所分解する }\Omega\text{ 上の任意の代数は }\sim\Omega\text{ である。}
\]

\paragraph{証明.}
\(A\) を \(\Omega\) 上で至る所分解する代数とし，\(K\) を \(A\) のガロア分解体とする。任意の有理素数 \(p\) と，\(K/\Omega\) の \(p\) に関する Sylow 体 \(X\)，すなわちガロア群の \(p\)-Sylow 部分群の不変体をとる。すると \(A_X\) も至る所分解し，しかも可解な分解体 \(K\) をもつ。II がすべての \(\Omega\) について成り立てば，\(A_X\sim X\) である。これは \(A\) が \(p\) と素な次数の分解体 \(X\) をもつことを意味する。従って \(A\) の指数はあらゆる \(p\) と素であり，指数は \(1\) である。よって \(A\sim\Omega\)。II から I が従う。

\paragraph{3.}
第三の還元は、H. Hasse が――最初の二つの還元を得たのち――最終的な証明へ導くものと認識したものであり、H. Hasse の知らせを契機として E. Noether が与えた。その知らせでは、可換分解体という特別な場合の主定理が、一般的な還元 1 と、それに属する因子系の公式による還元とを用いて証明されていた。E. Noether はその核心として、まさに第三の還元を取り出した。なお R. Brauer も、E. Noether とは独立に、この第三の還元をすでに考えていた。

\paragraph{還元 3.}
定理 II は次を示せば証明される：
\[
\text{III. 素数次数の巡回分解体をもち、至る所分解する }\Omega\text{ 上の任意の代数は }\sim\Omega\text{ である。}
\]

\paragraph{証明.}
\(A\) を、\(\Omega\) 上で可解な分解体 \(K\) をもち、至る所分解する代数とする。
\[
K=\Lambda_0\supset\Lambda_1\supset\cdots\supset\Lambda_r=\Omega
\]
を、\(K\) と \(\Omega\) の間の体の鎖で、各 \(\Lambda_i/\Lambda_{i+1}\) が素数次数の巡回拡大であるものとする。まず \(A_{\Lambda_1}\) は \(\Lambda_1\) 上で同じく至る所分解し、素数次数の巡回分解体 \(K=\Lambda_0\) をもつ。III がすべての \(\Omega\) について既知であれば、\(A_{\Lambda_1}\sim\Lambda_1\) である。これは \(\Lambda_1\) が \(A\) の分解体であることを意味する。そこで \(\Lambda_0\) の代わりに \(\Lambda_1\) から出発して同じ議論を行えば \(A_{\Lambda_2}\sim\Lambda_2\) を得る。以下同様にして、最後には \(A_{\Lambda_r}\sim\Lambda_r\)、すなわち \(A\sim\Omega\) を得る。従って III を仮定すれば II が証明される。

\paragraph{4.}
ところが III は H. Hasse の結果 [H, 24] によって既に成り立っている。したがって、還元 3, 2, 1 を逆にたどれば主定理の正しさが従う。

E. Noether はさらに次の点を指摘している。III の正しさ――より一般には、任意次数の巡回分解体に対する正しさ――は Hasse のノルム定理\footnote{H, 3.11 に引用され、H. Hasse, \emph{Beweis eines Satzes und Widerlegung einer Vermutung über das allgemeine Normenrestsymbol}, Gött. Nachr. 1931 で証明されている。}に帰着する。しかし還元 3 に必要なのは素数次数の場合、すなわちもとの Hilbert--Furtwängler のノルム定理\footnote{H. Hasse, \emph{Bericht über neuere Untersuchungen und Probleme in der Theorie der algebraischen Zahlkörper II}, Jahresber. der D. M.-V., Erg.-Bd. 6 (1930), § 8、およびそこで引用される文献を参照。}だけである。従って還元 III は Hasse のノルム定理の新しい簡明な証明を与える。その証明は次のようになる。

\(Z\) を \(\Omega\) 上の巡回体、\(\alpha\) を \(\Omega\) の数とし、\(\Omega\) の各素点 \(\mathfrak p\) に対してノルム剰余記号が
\[
\left(\frac{\alpha,Z}{\mathfrak p}\right)=1
\]
であるとする。このとき巡回代数
\[
A=(\alpha,Z)
\]
\footnote{H, 1 の記法である。\(\alpha\) と結び付く \(Z\) の自己同型 \(S\) の指定は、ここでは本質的でないため省略した。}は至る所分解する [H, 17.7]。従って還元 3 により――Hilbert--Furtwängler のノルム定理だけを用いて――\(A\sim\Omega\) となる。ゆえに \(\alpha\) は \(Z\) の元のノルムである [H, 15.4]。

さらにノルム定理との関連で、定理 I は、高次の場合（非可換の場合をも含む）へのその正しい一般化とみなすべきである。一方、字義どおりの一般化は、H. Hasse が示したように、一般には正しくない。
\begin{center}
\textbf{帰結（H. Hasse）}
\end{center}

\paragraph{5.}
以上によって、H. Hasse が代数的数体上の巡回代数について展開した理論は、正規単純代数、ことに正規斜体の一般的な構造論および不変量論としての意義を得た。とくに、まず次の定理が一般に証明された。

\paragraph{定理 1.}
代数的数体上の正規単純代数のブラウア類の位数（exponent）は、その指数（index）に等しい。

さらに 定理 I から次の注目すべき事実が従う。

\paragraph{定理 2.}
\(\Omega\) 上の真の正規斜体の基本イデアルは，\(\Omega\) の少なくとも一つの素点で割り切れ，実際には少なくとも二つで割り切れる。

ここで基本イデアルは，例えば reduced different の reduced norm として定義される。\footnote{有理四元数代数の場合，これは H. Brandt の導入した「基本数」の概念に帰着する。H. Brandt, \emph{Idealtheorie in Quaternionenalgebren}, Math. Ann. 99 (1928) を参照。ここでは数ではなくイデアルであるから「基本イデアル」と言う必要がある。これは Dedekind の different と判別式に対する「基本イデアル」「基本数」という用語とは混同してはならない。この用語は相対体ではここに述べた理由で不適当である。reduced different の定義については次の注を参照。E. Noether はイデアルの reduced norm も「素点ごと」に定義する。すなわち，各素点ではすべてのイデアルが主イデアルであることに対応して，主イデアルの基底数の reduced numerical norm を各素点について取る。} また無限素点は，その点で代数が実数体または複素数体ではなく四元数代数へ還元される場合，かつその場合に限って，基本イデアルに含められる。

\paragraph{証明.}
H. Hasse の結果\footnote{H. Hasse, \emph{Über \(p\)-adische Schiefkörper und ihre Bedeutung für die Arithmetik hyperkomplexer Zahlsysteme}, Math. Ann. 104 (1931), 定理 42 および 定理 59.}によれば，\(\Omega\) の素点 \(\mathfrak p\) が基本イデアルに現れるのは，\(\mathfrak p\)-指数が \(1\) と異なる場合，かつその場合に限る。すべての \(\mathfrak p\)-指数が \(1\) であれば，代数は至る所分解し，定理 I により真の斜体ではない。また，ただ一つの \(\mathfrak p\)-指数だけが \(1\) と異なることもない。なぜなら，その位数が \(\mathfrak p\)-指数であるノルム剰余記号 [H, 17.7] は積公式，すなわち相互法則 [H, 3.8] を満たすからである。

\paragraph{6.}
定理 I はさらに，\(\Omega\) 上の代数 \(A\) に属する分解体 \(K\) の算術的特徴付けを与える。H. Hasse が巡回的な \(A,K\) の特別な場合に得た結果 [H, 6, 定理 2] の正確な一般化として，次を得る。

\paragraph{定理 3.}
\(\Omega\) 上の代数 \(A\) に対し，代数的数体 \(K/\Omega\) が分解体であるための必要十分条件は，\(\Omega\) の各素点 \(\mathfrak p\) と，その上にある \(K\) の各素因子 \(\mathfrak P_i\) について，\(K\) の \(\mathfrak P_i\)-次数 \(n_{\mathfrak P_i}\) が \(A\) の \(\mathfrak p\)-指数 \(m_{\mathfrak p}\) の倍数であることである。

ここで \(\mathfrak P_i\)-次数とは，\(\mathfrak P_i\)-進体 \(K_{\mathfrak P_i}\) の \(\mathfrak p\)-進体 \(\Omega_{\mathfrak p}\) 上の次数である。
\[
\mathfrak p=\prod_i \mathfrak P_i^{\,e_{\mathfrak P_i}},
        \qquad \mathsf{N}_{\mathsf{K}\Omega}(\mathfrak P_i)=\mathfrak p^{f_{\mathfrak P_i}}
\]
ならば \(n_{\mathfrak P_i}=e_{\mathfrak P_i}f_{\mathfrak P_i}\) である [H, 4]。また \(A\) の \(\mathfrak p\)-指数とは，\(\mathfrak p\)-進拡大 \(A_{\mathfrak p}\) の指数 \(m_{\mathfrak p}\) である [H, 6]。

\paragraph{証明.}
\(K\) が \(A\) の分解体であること，すなわち \(A_K\sim K\) は，定理 I により，各 \(\mathfrak P_i\) について \(A_{K_{\mathfrak P_i}}\sim K_{\mathfrak P_i}\) であることと同値である。

有限素点 \(\mathfrak p\) について、
\[
        A_{\mathfrak p}\sim(\pi,W_{\mathfrak p})
\]
を \(A_{\mathfrak p}\) の算術的に特別な巡回表示とする [H, 16; Hasse, loc. cit., 定理 38 も参照]。ここで \(W_{\mathfrak p}\) は \(\Omega_{\mathfrak p}\) 上次数 \(m_{\mathfrak p}\) の不分岐体、\(\pi\in\Omega_{\mathfrak p}\) はちょうど \(\mathfrak p^1\) で割り切れる数である。さらに \(K_{\mathfrak P_i}W_{\mathfrak p}\) を合成体とし、
\[
        \Delta_{\mathfrak P_i}=K_{\mathfrak P_i}\cap W_{\mathfrak p}
\]
を \(W_{\mathfrak p}\) と \(K_{\mathfrak P_i}\) の交わりとする。\(K_{\mathfrak P_i}\) に含まれる最大の不分岐部分体は \(\Omega_{\mathfrak p}\) 上次数 \(f_{\mathfrak P_i}\) であり、また \(\Omega_{\mathfrak p}\) 上では各次数について不分岐体が一意であるから、\(\Delta_{\mathfrak P_i}\) は
\[
        \text{\(\Omega_{\mathfrak p}\) 上次数 }d_{\mathfrak P_i}=(m_{\mathfrak p},f_{\mathfrak P_i})
\]
をもつ。従って \(K_{\mathfrak P_i}W_{\mathfrak p}\) は \(K_{\mathfrak P_i}\) 上次数 \(m_{\mathfrak p}/d_{\mathfrak P_i}\) の不分岐拡大である。

いま
\[
        A_{K_{\mathfrak P_i}}=(A_{\mathfrak p})_{K_{\mathfrak P_i}}
        \sim(\pi,K_{\mathfrak P_i}W_{\mathfrak p})
\]
である [H, 15.5]。従って \(A_{K_{\mathfrak P_i}}\sim K_{\mathfrak P_i}\) であることは、\(\pi\) が \(K_{\mathfrak P_i}W_{\mathfrak p}\) から \(K_{\mathfrak P_i}\) へのノルムであることと同値である。拡大 \(K_{\mathfrak P_i}W_{\mathfrak p}/K_{\mathfrak P_i}\) は不分岐であるから、これは \(\pi\) の \(\mathfrak P_i\) における位数 \(e_{\mathfrak P_i}\) が次数 \(m_{\mathfrak p}/d_{\mathfrak P_i}\) の倍数であることと同値である。また
\[
\left(\frac{m_{\mathfrak p}}{d_{\mathfrak P_i}},
      \frac{f_{\mathfrak P_i}}{d_{\mathfrak P_i}}\right)=1
\]
であるから、これはさらに
\[
\frac{f_{\mathfrak P_i}}{d_{\mathfrak P_i}}e_{\mathfrak P_i}
 =\frac{n_{\mathfrak P_i}}{d_{\mathfrak P_i}}
\]
が \(m_{\mathfrak p}/d_{\mathfrak P_i}\) の倍数であること、すなわち \(n_{\mathfrak P_i}\) が \(m_{\mathfrak p}\) の倍数であることと同値である。これが主張である。

\(\mathfrak p\) が無限素点で、自明な場合 \(m_{\mathfrak p}=1\) でないときには、\(m_{\mathfrak p}=2\) であり、\(A_{\mathfrak p}\) は実数体 \(\Omega_{\mathfrak p}\) 上の四元数代数である。このとき \(A_{K_{\mathfrak P_i}}=(A_{\mathfrak p})_{K_{\mathfrak P_i}}\sim K_{\mathfrak P_i}\) であることは、\(K_{\mathfrak P_i}\) が複素数体であること、すなわち
\[
        n_{\mathfrak P_i}=f_{\mathfrak P_i}=2=m_{\mathfrak p}
\]
であることと同値である（ここでは \(e_{\mathfrak P_i}\) は現れない）。従ってここでも同じ主張を得る（\(n_{\mathfrak P_i}\) は \(m_{\mathfrak p}\) と同様に 1 または 2 しか値をとらない）。
\paragraph{7.}
固定した代数 \(A\) のすべての分解体 \(K\) を問う 定理 3 の問題を逆にし，固定したガロア代数体 \(K/\Omega\) について，\(K\) によって分解されるすべての \(\Omega\)-代数 \(A\) を考えると，重要な結果に到達する。それらの代数は，\(\Omega\) 上のすべての代数，より正確には相似類全体からなる R. Brauer 群 \(\mathfrak A\) の，\(K\) によって定まる部分群 \(\mathfrak K\) をなす。\footnote{R. Brauer, \emph{Über Systeme hyperkomplexer Größen}, Jahresber. d. D. M.-V. 38 (1929), pp. 47/48; H, 13.1 も参照。} \(K/\Omega\) がガロアであると仮定すると，相対 Abel 数体の理論である類体論の主定理を，一般の相対ガロア数体へ一般化するものと見ることのできる定理が得られる。

\paragraph{定理 4.}
\emph{分解定理.} \(K\) の相対判別式に現れない \(\Omega\) の素イデアル \(\mathfrak p\) について，\(K\) におけるその素因子の相対次数 \(f\) は，\(K\) に対応する群 \(\mathfrak K\) のすべての代数 \(A\) に対して
\[
        A_{\mathfrak p}^{\,f}\sim\Omega_{\mathfrak p}
\]
となる最初の指数に等しい。

\paragraph{証明.}
\(f\) は \(\mathfrak p\) のすべての素因子について一致する \(K\) の \(\mathfrak p\)-次数であり，一方 \(A\) の \(\mathfrak p\)-指数は \(A_{\mathfrak p}\) の指数（index）であり、定理 1 によりその冪指数（exponent）でもある。従って 定理 3 により，\(f\) は少なくともその最初の指数の倍数である。残るのは，\(\mathfrak K\) 内に正確に \(\mathfrak p\)-指数 \(f\) をもつ代数 \(A\) が存在することを示すことである。

Frobenius の密度定理により，\(K\) の相対判別式に現れない別の素イデアル \(\mathfrak p'\) が存在し，その \(K\) における素因子も相対次数 \(f\) をもつ。\(Z/\Omega\) を，\(\mathfrak p\)-次数および \(\mathfrak p'\)-次数がいずれも \(f\) である巡回体とする。さらに、\(\alpha\in\Omega\) を、ノルム剰余記号
\[
\left(\frac{\alpha,Z}{\mathfrak p}\right)\quad\text{および}\quad
\left(\frac{\alpha,Z}{\mathfrak p'}\right)
\]
が位数 \(f\) の互いに逆の値をとり、\(Z\) の導手の他のすべての因子 \(\mathfrak q\) について
\[
\left(\frac{\alpha,Z}{\mathfrak q}\right)=1
\]
となるように選ぶ。算術級数に関する一般化された定理により，\(\alpha\) は，場合により \(\mathfrak p,\mathfrak p'\) とそれらの \(\mathfrak q\) を除けば，ただ一つの素イデアル \(\mathfrak r\) をちょうど一乗で含むようにも選べる。この \(\mathfrak r\) では、ノルム剰余記号の積定理、すなわち相互法則により、同じく
\[
\left(\frac{\alpha,Z}{\mathfrak r}\right)=1
\]
である [H, 3.8--10]。代数
\[
        (\alpha,Z)=A
\]
は \(\mathfrak p\)-指数および \(\mathfrak p'\)-指数 \(f\) をもち，他のすべての素点では指数 \(1\) をもつ [H, 17.7]。定理 3 と \(\mathfrak p\) についての仮定および \(\mathfrak p'\) の選択により，\(K\) は \(A\) の分解体である。従って \(\mathfrak K\) 内に \(\mathfrak p\)-指数 \(f\) の代数が実際に存在する。

\paragraph{定理 5.}
\emph{一意性および順序定理.} \(K\leqq K'\) ならば \(\mathfrak K\leqq\mathfrak K'\) であり，逆も成り立つ。従ってガロア体 \(K\) への代数群 \(\mathfrak K\) の対応は，一対一で順序を保つ。

\paragraph{証明.}
\(K\leqq K'\) から \(\mathfrak K\leqq\mathfrak K'\) は明らかである。\(K\) で分解する任意の代数は，なおさら \(K'\) で分解する。逆に \(\mathfrak K\leqq\mathfrak K'\) とする。分解定理により，相対判別式の非因子 \(\mathfrak p\) ごとに対応する相対次数の割り切り関係が得られる。特に \(f'_{\mathfrak p}=1\) であるほとんどすべての \(\mathfrak p\) について \(f_{\mathfrak p}=1\) となる。通常の解析的議論\footnote{例えば H. Hasse, 注 6 の文献, § 25, III を参照。}から \(K\leqq K'\) が従う。

\paragraph{8.}
最後に，主定理は I. Schur\footnote{I. Schur, \emph{Arithmetische Untersuchungen über endliche Gruppen linearer Substitutionen}, Berl. Akad.-Ber. 1906.} が扱った，有限群の絶対既約表現がどの数体で可能かという問題にも本質的な進展を与える。

\paragraph{定理 6.}
有限群 \(\mathfrak G\) の絶対既約表現はすべて円分体で可能である。たとえば \(n=|\mathfrak G|\) とし、\(h\) を十分大きく取れば、\(n^h\) 乗根の体で常に可能である。

\paragraph{証明.}
\(\mathfrak G\) の有理群環 \(G\) へ移ると、\(\mathfrak G\) の絶対既約表現 \(\Gamma_i\) は、半単純代数 \(G\) の単純成分 \(G_i\) の絶対既約表現となる。それらの中心は対応する指標の体 \(\Omega_i\) である。\footnote{a) R. Brauer and E. Noether, \emph{Über minimale Zerfällungskörper irreduzibler Darstellungen}, Berl. Akad.-Ber. 1927, § 1; b) R. Brauer, \emph{Über Systeme hyperkomplexer Zahlen}, Math. Z. 30 (1929), 定理 3; c) E. Noether, \emph{Hyperkomplexe Größen und Darstellungstheorie}, Math. Z. 30 (1929), §§ 21, 24, 26 を参照。} \(\mathfrak G\) は有限であるから、これらは円分体であり、少なくとも \(n\) 乗根の体の部分体である。

定理 3 により、指標体 \(\Omega_i\) 上の巡回体 \(Z_i\) は、各素点 \(\mathfrak p\) における \(\mathfrak p\)-次数 \(n_{i\mathfrak p}\) が \(G_i\) の \(\mathfrak p\)-指数 \(m_{i\mathfrak p}\) の倍数であれば、\(G_i\) の分解体である。\(m_{i\mathfrak p}\ne1\) となるのは \(G_i\) の基本イデアルの素因子、したがって \(G_i\) の絶対判別式の素因子に限られる。その絶対判別式は \(n^n\) を割る。なぜなら \(n^n\) は群環 \(G\) の非極大オーダーの判別式だからである。\footnote{E. Noether, 注 12c の文献, § 26 を参照。} 従って問題となるのは \(n\) の素因子 \(\mathfrak p\) だけである。これらについて \(n_{i\mathfrak p}\) を \(m_{i\mathfrak p}\) の倍数にするには、\(\mathfrak p\) が \(Z_i\) でそれぞれ \(m_{i\mathfrak p}\) で割り切れる位数だけ分岐するよう指定すればよい。これは \(h\) を十分大きく取った \(n^h\) 乗根の体によって達成される。
引用された Schur の論文の最後の文で，Schur は既知のすべての場合ではすでに \(n\) 乗根の体で十分であると述べている。これが常に正しいか，またここで展開した方法がこの問題を決定するのに十分かは，今後の研究に残される。

\begin{center}
1931年11月11日受理。
\end{center}

\clearpage

\section*{39. 可換代数および数論との関係における超複素系}

\begin{center}
\emph{Verhandl. Intern. Math.-Kongreß Zürich 1 (1932), pp. 189--194}

\textsc{Emmy Noether（Göttingen）}
\end{center}

\paragraph{1.}
超複素系，すなわち代数の理論は，近年大きく発展した。しかし，この理論が可換的な問題に対してもつ意味が明らかになったのはごく最近である。ここでは，非可換なものが可換なものに対してもつ意味を報告し，Gauss に遡る二つの古典的問題，主種定理とそれに密接に関係するノルム定理を通して具体的に述べたい。これらの問題は時代とともに何度も形を変えてきた。Gauss では二次形式論の締めくくりとして現れ，その後，類体論による相対巡回および Abel 数体の特徴付けに本質的な役割を果たし，最後には自己同型および代数の分解に関する定理として述べられる。この最後の定式化は，同時に，それらの定理を任意の相対ガロア数体へ移すことを可能にする。

この概略を後で具体化しながら，非可換なものを可換なものへ応用する原理も説明したい。代数の理論によって，二次形式または巡回体に関する既知の事実に対し，代数の構造的性質だけに依存する不変で単純な定式化を求めるのである。そのような不変な定式化が一度証明されれば，上記の例でそうであるように，それらの事実は自動的に任意のガロア体へ移される。

\paragraph{2.}
詳細に入る前に，諸方法とさらに得られる結果を概観する。まず注意すべきは，一般のガロア体に対する定式化を得ることが主要な困難であり，超複素的方法なしには出発点すらなかったことである。上の例における非可換への移行は，接合積とその乗法定数，すなわち「因子系」（3 を参照）によって，体と群を同時に考えることから得られる。こうして基礎体上の正規単純代数が得られ，任意のそのような代数は本質的にこのように作れる。このような接合積は Dickson によって最初に考えられた。\footnote{彼の著書 \emph{Algebren und ihre Zahlentheorie}, Zürich 1927, § 34 を参照。} 一方，因子系の理論は Speiser, Schur, R. Brauer によって\footnote{例えば R. Brauer, \emph{Untersuchungen über die arithmetischen Eigenschaften von Gruppen linearer Substitutionen}, および同論文注 2 の文献, Math. Z. 28 (1928) を参照。}全く別の出発点，すなわち絶対既約表現の問題から発展した。両者の融合によって初めて，可換的問題に取り組むのに十分単純で広い構成が得られた。\footnote{私はこの構成を 1929/30 年冬学期の講義で初めて展開した。これは H. Hasse, \emph{Theory of cyclic algebras}, Trans. Amer. Math. Soc. 34 (1932), 第2章に再現されている。この講演で扱う領域全体については，M. Deuring の超複素数と数論的応用に関する報告が \emph{Ergebnisse der Mathematik} に現れる予定で，そこに概説される。}

同時に，既知の事実についても新しく見通しのよい証明が得られる。H. Hasse が与え，まもなく Math. Ann. に現れる予定の，巡回体に対する相互法則の超複素的証明を挙げたい。\footnote{H. Hasse, \emph{Die Struktur der R. Brauerschen Algebrenklassengruppe über einem algebraischen Zahlkörper}, 特にノルム剰余記号と相互法則, Math. Ann. 107 (1932/33).} これは接合積の理論に基づく，Hasse のノルム剰余記号の不変な定式化によるものである。さらに，同じ基礎の上で，Chevalley が最近与えた局所類体論の超複素的基礎付けを挙げる。ただしそこでは因子系に関する新しい代数的定理を展開する必要があった。\footnote{C. Chevalley, \emph{Sur la théorie du symbole de restes normiques}, J. f. Math. 169 に掲載予定。} しかし制限も述べておかなければならない。接合積の方法だけでは，どうやらガロア数体の完全な理論は得られない。これは，Artin の未公表の新しい結果から従う。その結果は，上の原理の意味で Hasse の証明に接続するが，完全な同型定理の代わりに個数の等式だけを与える。

完全な同型，ただし作用素同型を与える方法は，すでに代数の中にある。これは A. Speiser の着想\footnote{A. Speiser, \emph{Gruppendeterminante und Körperdiskriminante}, Math. Ann. 77 (1916).}の発展であり，ガロア体を「ガロア加群」，すなわちガロア群の置換を作用素として許す基礎体上の加群と見る立場である。体と群環（群代数）の間には作用素同型があり，基礎体上の線型形式が対応し，体におけるガロア群の置換は群環における乗法に対応する。この私が述べた定理\footnote{E. Noether, \emph{Normalbasis bei Körpern ohne höhere Verzweigung}, Satz 3, J. f. Math. 167 (1932); 証明には欠陥がある。}は M. Deuring により証明され，\footnote{M. Deuring, \emph{Galoissche Theorie und Darstellungstheorie}, Math. Ann. 107 (1932).}彼はそれに基づいてガロア理論の証明を与えた。さらに Deuring による構造定理は，Artin の \(L\)-級数に関する形式的事実と平行し，Artin 導手への構造的接近を与える。一般群指標で作られる Artin の \(L\)-級数および導手\footnote{E. Artin, \emph{Über eine neue Art von \(L\)-Reihen}, Math. Sem. Hamburg 3 (1924); \emph{Zur Theorie der \(L\)-Reihen mit allgemeinen Gruppencharakteren}, 同 8 (1931); \emph{Die gruppentheoretische Struktur der Diskriminanten algebraischer Zahlkörper}, J. f. Math. 164 (1931).}は，Speiser と並んで，数論と表現論の最初の結合，Abel 体を越える最初の前進をなす。これらは全体の発展に強い刺激を与え，とくにガロア加群の理論はそれに導かれた。

\paragraph{3.}
ここから，冒頭に掲げたノルム定理と主種定理を詳しく追う。まず接合積の定義である。\(K/k\) を次数 \(n\) のガロア体，\(\Gg\) をその群とする。接合積とは，\(K\) と \(\Gg\) を同時に代数 \(A\) の中へ埋め込み，\(K\) の自己同型が内的になるようにすることである。\(u_{S_1},\ldots,u_{S_n}\) を \(n\) 個の群要素に対応する記号とする。まず \(A\) を \(K\) 上ランク \(n\) の線型形式加群として
\[
(1)\qquad A=u_{S_1}K+\cdots+u_{S_n}K
\]
と置く。すなわち \(A\) は任意の \(a_i\in K\) に対する \(u_{S_1}a_1+\cdots+u_{S_n}a_n\) からなる。\(u_S\)，より一般には \(u_SK^\ast\)\footnote{\(K^\ast\) は \(K\) から零を除いて得られる。この記法は一般に用いられる。}によって内的自己同型を要求することで，\(A\) は環となり，従って \(k\) 上ランク \(n^2\) の代数となる。この条件は，各 \(z\in K\) について
\[
(2)\qquad u_S^{-1}zu_S=z^S,\qquad \text{または}\qquad zu_S=u_Sz^S
\]
と表される。\footnote{\(z^S\) は通常通り，置換 \(S\) によって \(z\) から得られる元を意味する。} また
\[
(3)\qquad u_Su_T=u_{ST}a_{S,T},\quad a_{S,T}\in K^\ast
\]
であり，結合法則は因子系条件
\[
\begin{gathered}
(4)\qquad a_{ST,R}\,a_{S,T}^{\,R}=a_{S,TR}\,a_{T,R},\\
\text{（結合法則 }[u_Su_T]u_R=u_S[u_Tu_R]\text{ から）}
\end{gathered}
\]
を与える。この \(A\) が、因子系 \(a_{S,T}\) に関する \(K\) と \(\Gg\) の接合積である。\(A\) は \(k\) 上の正規単純代数、したがって随伴斜体 \(D\) 上の次数 \(r\) の行列環 \(D_r\) であり、\(K\) は極大可換部分体、従って分解体となる。係数領域 \(k\) を \(K\) と同型な体で拡大すると，分解代数，すなわち中心上の全行列環が得られる。逆に，与えられた斜体 \(D\) に対しては，常にこのような接合積として作れる行列環 \(D_r\) が存在する。

\(u_S\) から \(v_S=u_Sc_S\), \(c_S\in K^\ast\) へ移ると，同じ自己同型を与え，随伴する因子系
\[
(5)\qquad \bar a_{S,T}=a_{S,T}\frac{c_S^{\,T}c_T}{c_{ST}}
\]
が得られる。随伴因子系を一つの類 \((a)\) にまとめ，同様に \(A\) と相似なすべての代数，すなわち \(D_r\), \(r=1,2,\ldots\)，を一つの類 \(\mathfrak A\) にまとめる。類 \(\mathfrak A\) と \((a)\) は一対一に対応する。固定された分解体 \(K\) をもつ類は，直積に関して Abel 群をなし，因子系の類の成分ごとの積と同型である。単位元は分解代数の類，あるいは同じこととして，すべての変換量
\[
        \frac{c_S^{\,T}c_T}{c_{ST}}
\]
の系である。これは R. Brauer が見出した代数類群である。

\paragraph{4.}
次に巡回分解体に特殊化して，ノルム概念との関係に進み，冒頭の原理に従って一般化されたノルム定理の定式化に到達する。\(Z/k\) が巡回的で，\(S\) がその群の生成置換であるとき，\(S\) の冪に \(u\) の冪を対応させることができる：
\[
(1')\qquad A=Z+uZ+\cdots+u^{n-1}Z,
\]
\[
(2')\qquad zu=uz^S,\qquad
(3')\quad u^n=a,\qquad
(4')\quad a\in k^\ast,
\]
また \(v=uc\) とおけば
\[
(5')\qquad \bar a=a\cdot N(c).
\]
従ってここでは各因子系は基礎体の一つの元 \(a\) だけからなり，\(A=(a,Z)\) と書く。因子系の単位類は \(Z^\ast\) からのノルムで与えられ，代数類群は
\[
        k^\ast/N(Z^\ast)
\]
と同型になる。巡回代数 \((a,Z)\) が分解するのは，\(a\) が \(Z\) の元のノルムである場合，かつその場合に限る。

このノルムと分解の関係が「ノルム定理」の定式化，すなわち分解代数に関する定理を与える。代数が各点で分解すれば，大域的にも分解する。ここで「点」とは，数論で通常そうするように，基礎体 \(k\) を素イデアル \(\mathfrak p\) に関する \(p\)-進拡大 \(k_{\mathfrak p}\) に置き換えることにより定義される。有限個の無限点では，\(k\) とその共役を実数体で拡大することに対応する。

実際，巡回体についてのノルム定理はここに含まれる。巡回代数 \((a,Z)\) では，この主張は次と同値である：\(a\) が各有限または無限点で \(p\)-進ノルムであれば，\(a\) は \(Z\) の数のノルムである。あるいは \(p\)-進への移行なしに言えば，\(a\) が \(k\) の各素イデアル \(\mathfrak p\) に関してノルム剰余であり，さらに一定の符号条件を満たせば，\(a\) は \(Z\) の数のノルムである。この後者の形こそ，既知の解析的手段を用いて類体論で証明されるノルム定理である。分解代数についての一般定理の証明は，この巡回的特別場合から純代数的・算術的考察によって得られる。\footnote{R. Brauer, H. Hasse, E. Noether, \emph{Beweis eines Hauptsatzes in der Theorie der Algebren}, J. f. Math. 167 (1932).} Hasse はそこから最初の重要な帰結を引き出した：代数的数体上の任意の正規単純代数は巡回的である。この長く予想されていた事実の証明を探す過程で，一般的な定式化が生じた。

\paragraph{5.}
分解代数に関する定理からの第二の帰結は，やはり純代数的・算術的推論による，冒頭に掲げた主種定理である。\footnote{証明は Math. Ann. に現れる予定である。} その不変な定式化は，接合積を定義する関係 (2) から (5) が純粋に乗法的であることに基づく。従って \(K^\ast\) を Abel 群 \(\mathfrak J\) に置き換えても意味をもつ。ただし \(\mathfrak J\) の自己同型群が \(\Gg\) と同型な部分群を含むことだけを仮定する。(1) の代わりには，群論の意味での \(\Gg\) の \(\mathfrak J\) による拡大が現れる。\(\mathfrak J\) として \(K\) のすべてのイデアルの群を取ると，因子系はイデアルの系となる。\(\mathfrak J\) の類分けは因子系の類分けを誘導し，これは一般に元のイデアル類分けより細かい。なぜなら、\(u_S\) に絶対イデアル類を掛けることの一意性から言えるのは、変換量 \(c_S^{\,T}c_T/c_{ST}\)、すなわち元の因子系の単位類が、イデアル因子系の単位類に属することだけだからである。しかし既知の場合への特殊化が示すように、実際にはやや粗い類分けで十分である。そこで次のように定義する。

因子系の単位類には，\(K\) のすべての有限および無限分岐点で分解する代数を生成する元 \(a_{S,T}\) を入れる。このようにして得られる \(\Gg\) の拡大を \(\Gg^\ast\) と書き，\((c)\) を \(c\) の絶対イデアル類とする。すると主種定理の不変な定式化は次である：

この類分けのもとで置換
\[
        v_S=u_S(c_S)
\]
が \(\Gg^\ast\) の自己同型を生じるならば，これらすべての \((c_S)\) が主種をなし，その自己同型は内的であって，あるイデアル類 \((\mathfrak b)\) によって生成される。

既知の特殊場合は，次の同値でやや明示的な形から従う。イデアル類 \((c_S)\) から作られる変換量
\[
        \frac{(c_S^{\,T})(c_T)}{(c_{ST})}
\]
が因子系の単位イデアル類に属するならば，あるイデアル類 \((\mathfrak b)\) が存在して，すべての \(S\in\Gg\) について
\[
        (c_S)=\frac{(\mathfrak b)}{(\mathfrak b)^S}
\]
となる。仮定はまさに自己同型性を表し，その自己同型が内的であることは
\[
        v_S=(\mathfrak b)^{-1}u_S(\mathfrak b)
            =u_S(\mathfrak b)^{\,1-S}=u_S(c_S)
\]
によって表される。巡回の場合に特殊化し，正規化を考慮すると，\(N([c])\) が因子系の単位イデアル類に属するならば，\((c)\) は記号的な \((1-S)\) 乗となる：
\[
        (c)=(\mathfrak b)^{\,1-S}.
\]

\paragraph{6.}
ここから巡回体および二次形式についての既知の形へ進むため，まずこの定理は完全なイデアル類について述べられているが，通常のように分岐点と素なイデアルに制限しても内容は同じであることを注意する。

これにより，ここで定義された二次体の主種は Gauss の主種へ移る。イデアル類には二次形式が対応し，類のノルムにはその形式によって表される数が対応する。単位類が \(K\) の分岐点で分解する代数を生成するということは，その表される数が分岐点で二次剰余になることを意味する。対応する形式は主形式の全指標をもち，従って Gauss の主種を形成する。\((1-S)\) 乗の記号的冪が重複に移ることはよく知られている。

巡回体では，主張は次の形になる：主種は，そのノルムが有限および無限の分岐点でノルム剰余となるすべてのイデアル類から成る。これは「導手に関するノルム剰余」と同値であり，従って通常の定理が特殊化として現れる。

なお，一般の任意のガロア体の場合にも，分岐点だけから構成される導手を導入できる。そのようにすれば，因子系を正規化した後，これら各点の ray は単位類の元だけを含む。

ここで，概観 2 で述べた Artin 導手との関係が問題になる。それらは同じ素イデアルから構成されるからである。そして同時に，第二の超複素的方法であるガロア加群の理論との関係も問題になる。この二つの方法がどこまで到達するかは，今後に委ねられる。

%
\clearpage
\section*{40. 非可換代数}

\begin{center}
\emph{Math. Zs. 37 (1933), 514--541}
\end{center}

可換代数の主要定理は、よく知られているように、ガロア理論に含まれており、その前段には添加体と分解体の理論が置かれる。すなわち、与えられた多項式が一次因子を一つ分離するため、または完全に一次因子へ分解するために十分な体の理論である。

ここではそれに対応する代数の部分を、非可換の場合、とくに超複素の場合に展開する。原理的には非可換的方法、すなわち非可換体（斜体）における表現を用いる。最後に、上に述べた可換代数の定理が、可換体における表現によって完全に平行な仕方で証明されることを示す。

用いられる表現論は、短い自己同型論（§1）を前置したうえで、表現加群に基づく理論をさらに発展させたものである。とくに、直接表現のほかに相反表現を考える。両者は互いに還元できるが、ここでは相反表現は相反表現加群によって生成される。この利点は、その相反表現加群が、したがって双加群が、拡大環上の片側加群としても見られる点にある（§2）。こうして斜体における表現は拡大環におけるイデアル論に還元される。既約表現類は拡大環の既約イデアル類に対応し、これは超複素系そのものを可換な係数体上で表現したときに成り立つ事実と正確に同じである（§§3, 4）。

この立場から、非可換体上の行列環についての構造定理が得られる。すなわち、任意の部分環はこの体による表現を、同値にいえば相反同型な体による相反表現を与える、という観察（§5）による。この相反同型な体は、可換の場合における最小添加体および分解体の第一の類似物である。これはすべての一次の相反表現を媒介し、また中心上の階数が有限の場合には階数 \(1\) の直和因子への完全分解も媒介する。これにより斜体のガロア理論（§6）が得られ、さらに、相反同型な除法代数、すなわち中心上有限階数の体が、R. Brauer の代数類群において逆類を生成するという事実（§7）も表される。

より一般に単純系に有効なガロア理論の第二の証明は、可換化する部分環に関する定理のほとんど直接の帰結であり、この定理自体も前の自己同型に関する考察とほぼ直接につながっている。ここまでは議論は純粋に非可換である。しかし可換な分解体の問題も、§5 で準備した観察を用い、分解体を除法代数を通じて表現することにより非可換的に扱われる（§7）。最後の部分では可換の場合への移行（§8）と任意の系に対する分解体の理論（§9）を与える。

R. Brauer は分解体の理論を、可換の場合のこの表現論と非有理的な因子系の上に築いた。この可換的基礎のために、彼、そして後に Albert は、中心が完全体であることを仮定しなければならなかったが、この制限は非可換的基礎では不要である。同じ制限のもとで、やはり可換の場合の表現論を用いて、R. Brauer と K. Shoda は理論をさらに発展させた。

最後に、斜体については、1928 年夏の私の講義に従った van der Waerden の短い叙述があることを述べておく。van der Waerden の叙述はその講義に対していくつかの簡略化を導入した。その一部は二度目の講義で採用し、一部はここで初めて採用した。可換化する部分環に関する定理（§5）と、それに基づく非可換ガロア理論の第二の証明（§6, 1 および 2）は、後になって加えられたものである。第一の講義におけるある議論、すなわち差分イデアルの交わりを作る方法は、より複雑ではあるが、無限階数の系および整な系に移行できるという利点を持つ。可換な整系については、これがイデアル微分と共役差積（Different）との関係へ導く。

分解体の理論の主要な応用は、接合積とその因子系の理論にあり、それはさらに数論への応用の基礎となる。ここではこれ以上追わない。

\subsection*{§1. 自己同型、加群、双加群について}

表現加群に基づく表現論は、作用素を持つ、または持たないアーベル群の自己同型環の理論に依拠する。そこで暗に用いられる議論、すなわち写像と計算法則との関係を、反復を避けるためにいくつかの単純な命題として定式化しておく。

\paragraph{1. 乗法的写像。結合法則。}
まず \(\frG\) を作用素を持たない群とし、\(\frA\) をその絶対自己同型領域、すなわち \(\frG\) から自身へのすべての準同型の系とする。\(\frA\) は乗法に関して閉じている。なぜなら積 \(\sigma\tau\) は
\[
        g(\sigma\tau)=(g\sigma)\tau,
        \qquad g\in\frG,\quad \sigma,\tau\in\frA,                       \tag{1}
\]
によって定義され、明らかに結合法則を満たすからである。

記号 \(O,H,\ldots\) からなる集合 \(\frB\) が与えられ、積 \(gO,gH,\ldots\) が \(\frG\) の元として定義され、かつ \(\frG\) から自身への準同型を生成するとき、群 \(\frG\) は作用素を持つ群となる：
\[
        (gh)O=gO\cdot hO.
\]
作用素領域 \(\frB\) が乗法的に閉じているなら、自己同型領域への誘導写像が乗法的準同型であることは、結合関係
\[
        g(OH)=(gO)H,
        \qquad g\in\frG,
        \ O,H\in\frB                                                   \tag{1a}
\]
が成り立つことと同値である。対応して、左作用素については相反準同型が得られる。

\paragraph{2. 作用素準同型写像。}
\(\frG\) が作用素を持つ群であるとき、作用素準同型性は
\[
        (gO)\sigma=(g\sigma)O,                                          \tag{2}
\]
によって、左作用素については
\[
        (Og)\sigma=O(g\sigma)                                           \tag{2*}
\]
によって定義される。二つの作用素領域 \(\frB\) と \(\frC\) について、\(\frC\) の元が \(\frB\)-自己同型を生成し、同時に \(\frB\) の元が \(\frC\)-自己同型を生成するのは、それらの領域が \(\frG\) と可換的に結びついているとき、すなわち走る結合法則
\[
        (gO)H=(gH)O,                                                    \tag{2a}
\]
または
\[
        (OH)g=O(Hg)                                                     \tag{2a*}
\]
が満たされるとき、かつそのときに限る。

\paragraph{3. 環上の加群と双加群。}
環 \(\frR\) 上の右加群 \(\frM\) とは、\(\frR\) の元を右作用素として持ち、結合関係と分配関係
\[
        (g+h)\rho=g\rho+h\rho,
        \qquad
        g(\rho+\sigma)=g\rho+g\sigma                                  \tag{3a}
\]
を満たす加法的アーベル群である。左加群についても対応する定義が成り立つ。

双加群には二種類ある。二つの環 \(\frR\) と \(\frS\) 上の右加群は、個別の加群法則に加えて
\[
        (m\rho)\sigma=(m\sigma)\rho
\]
を満たすとき双加群と呼ばれる。\(\frR\)-左、\(\frS\)-右加群は、走る結合法則
\[
        \rho(m\sigma)=(\rho m)\sigma
\]
を加えるとき双加群である。

加法的に書かれたアーベル群 \(\frM\) の作用素領域 \(\frR\) が環であるとき、\(\frR\) から絶対自己同型環への写像が環準同型であることは、\(\frM\) が右 \(\frR\)-加群であることとちょうど同値である。左加群の場合には相反的な環準同型になる。\(\frM\) が \(\frR,\frS\) 上の右加群であるとき、\(\frM\) が双加群であることは、\(\frR\) が \(\frS\)-自己同型環の部分環へ準同型に写り、\(\frS\) が \(\frR\)-自己同型環の部分環へ準同型に写ることとちょうど同値である。\(\frR\)-左、\(\frS\)-右双加群では、\(\frR\) の写像は相反であり、\(\frS\) の写像は直接である。

とくに次が成り立つ。

1) \(\frR\) が単位元を持つ環なら、\(\frR\) は左 \(\frR\)-加群としての自己同型環に直接同型であり、右 \(\frR\)-加群としての自己同型環に相反同型である。

2) \(\frR\) が一般には非可換な体 \(A\) 上の全行列環で、
\[
        \frR=\sum c_{ik}\mathsf{A}=\sum \mathsf{A} c_{ik},
\]
ならば、\(A\) は単純右イデアルの自己同型体に相反同型であり、単純左イデアルの自己同型体に直接同型である。

\paragraph{4. 右双加群から積環上の片側加群への移行。}
右双加群の意義は本質的にこの移行に基づく。

\paragraph{定理。}
\(\frM\) を、互いに元ごとに可換な二つの部分環 \(\frR\) と \(\frS\) を含む環 \(\frT\) 上の右加群とする。このとき \(\frM\) は \(\frR\) と \(\frS\) 上の双加群とも見なせる。逆に、\(\frR\) と \(\frS\) 上の双加群が与えられ、\(\frR\) と \(\frS\) が元ごとに可換となる積環 \(\frT\) が存在するなら、\(\frM\) は \(\frT\) 上の右加群とも見なせ、その \(\frT\)-作用は与えられた \(\frR\)-作用と \(\frS\)-作用を延長する。

絶対自己同型環の中の積環は、
\[
        \sum \rho_i\sigma_i+\rho+\sigma
\]
の形のすべての元からなる（\(\frR,\frS\) が単位元を持つなら、余分な項 \(\rho,\sigma\) は消える）。準同型
\[
        \frR\to\overline{\frR},\qquad
        \frS\to\overline{\frS}
\]
は
\[
        \sum \rho_i\sigma_i+\rho+\sigma
        \to
        \sum \overline\rho_i\,\overline\sigma_i+\overline\rho+\overline\sigma
\]
によって \(\frT\to\overline{\frT}\) の準同型へ延長される。したがって
\[
        m\Bigl(\sum \rho_i\sigma_i+\rho+\sigma\Bigr)
        =\sum (m\rho_i)\sigma_i+m\rho+m\sigma
\]
はよく定義され、与えられた \(\frR,\frS\)-加群構造を含む \(\frT\)-加群を定める。

\subsection*{§2. 相反表現と直接表現}

\paragraph{1. 線形形式の加群。}
\(\frS\) を単位元を持つ環とする。右 \(\frS\)-加群 \(\frM\) が \(\frS\) における線形形式の加群と呼ばれるのは、\(\frM\) が \(n\) 個の一項 \(\frS\)-加群の直和
\[
        \frM=m_1\frS+\cdots+m_n\frS
\]
であり、各 \(m_i\frS\) が \(\frS\) と作用素同型、したがって \(m_i s=0\) なら常に \(s=0\) となるときである。

\paragraph{定理 1.}
\(\frM\) が \(\frS\) における線形形式の加群なら、\(\frM\) の \(\frS\)-自己同型環 \(\frU\) は、\(\frS\) 上の \(n\) 行全行列環 \(\overline{\frU}\) に、単なる準同型ではなく相反同型である。

\(\frM\) の任意の \(\frS\)-自己同型 \(\alpha\) は
\[
        m_i\longmapsto \overline m_i=m_i\alpha
\]
によって完全に記述される。したがって
\[
        \sum m_i s_i\to
        \sum \overline m_i s_i
        =\sum (m_i\alpha)s_i
        =\sum (m_i s_i)\alpha.
\]
さらに
\[
        (m_1\alpha,\ldots,m_n\alpha)=(m_1,\ldots,m_n)A
\]
と書けば、対応 \(\alpha\mapsto A\) は一対一であり、
\[
        \alpha+\beta\mapsto A+B,
        \qquad
        \alpha\beta\mapsto BA.
\]
よってこの同型は相反である。従って次を得る。

\paragraph{定理 1'.}
\(\frS\) における次数 \(n\) の線形形式の加群の自己同型環は、\(\frS\) 上の \(n\) 行全行列環によって、同型で忠実な相反表現を許す。

\paragraph{2. 表現と表現加群。}
\(\frS\) における線形形式の加群 \(\frM\) は、\(\frR\) と \(\frS\) 上の双加群であり、同時に \(\frR\) と \(\frS\) の双方について右加群であるとき、\(\frR\) の \(\frS\) における相反表現加群と呼ばれる。\(\frM\) が双加群であり、同時に左 \(\frR\)-加群かつ右 \(\frS\)-加群であるとき、直接表現加群と呼ばれる。

\paragraph{定理 2.}
任意の相反表現加群、また直接表現加群は、それぞれ \(\frR\) の \(\frS\) における同値な相反表現、また直接表現の一つの類を生成し、そのようなすべての表現はこの仕方で得られる。

相反表現加群に対して、\(\frR\) は \(\frM\) の \(\frS\)-自己同型環の部分環へ直接準同型に写る。そこで定理 1' が相反表現を与える。逆に、相反表現を \(\frS\)-自己同型環への相反同型と合成すると、\(\frR\) からの直接準同型が得られ、したがって表現加群が得られる。直接の場合は準同型性を反転して得られる。\(\frR\) の直接表現は、\(\frR\) の相反環の相反表現であり、その逆も同様である。

\paragraph{3. 相反表現加群から拡大環上の加群への移行。}
表現加群に対して移行定理は次の形をとる。

\(\frM\) が、互いに元ごとに可換な二つの部分環 \(\frR\) と \(\frS\) を含む環 \(\frT\) 上の右加群であり、かつ \(\frS\) 上の線形形式の加群であるなら、\(\frM\) は \(\frR\) の \(\frS\) における相反表現加群とも見なせる。逆に、相反表現加群は、積環 \(\frT\) が存在する限り \(\frT\)-加群と見なせる。

\(\frM\) が \(\frR\) の \(\frS\) における相反表現加群であり、かつ \(\frT\) を積環とする \(\frT\)-加群でもあるなら、\(\frM\) と加群 \(\frN\) との \(\frR,\frS\)-同型は \(\frT\)-同型と同値である。従って、同型な \(\frT\)-加群の各類は \(\frR\) の \(\frS\) における相反表現類に対応し、逆も同様である。

\paragraph{可換化する行列に関する定理。}
\(\frR\mapsto\frR^*\) を \(\frR\) の \(\frS\) における次数 \(n\) の相反表現とし、\(\frB^*\) を \(\frS\) 上の \(n\) 行行列のうち \(\frR^*\) と元ごとに可換なすべてのものの環とする。このとき \(\frB^*\) は、それによって生成される表現加群の \(\frR,\frS\)-自己同型、すなわち同時に \(\frR\)-自己同型でもある \(\frS\)-自己同型の相反同型表現を与える。積環 \(\frT\) が存在すれば、それらは \(\frT\)-自己同型である。

\(\frR,\frS\)-自己同型は、同じ表現を生じる対応 \(m_i\mapsto m_i'\) にほかならない。元 \(m_i'\) は \(\frM\) の \(\frS\)-基底をなす必要はなく、したがって \(\frB^*\) の行列も可逆である必要はない。この拡張された意味では
\[
        (m_1',m_2',\ldots,m_n')a=(m_1',m_2',\ldots,m_n')A
\]
はなお正しい。ただし \(A\) はもはや \(a\) によって一意には決まらない。

対角行列 \(E\cdot s\) は \(\frS\) の直接同型な表現、すなわち \(\frS\) の恒等表現を与える。したがって \(\frR\) と \(\frS\) によって定義される \(\frT\) から \(\frS\) 上の行列への対応は、環 \(\frT\) の表現ではなく、\(\frR\) の相反表現を、\(\frS\) に関して作用素準同型なものへ延長したものである：
\[
        \sum r_i s_i\longmapsto \sum R_i s_i.
\]
とくに、\(\frR\) の相反表現は、\(\frR\) と \(\frS\) の交わり \([\frR,\frS]\) に関して作用素準同型であり、この交わりは \(\frR\) の中心に含まれる。

\subsection*{§3. 斜体に関する加群}

以下では、一般には非可換な斜体における表現だけを扱う。そのため、斜体上の線形形式の加群についての簡単な結果を集めておく。

\paragraph{1. 全加群の与えられた基底に関する部分加群の正規基底。}
\(A\) を斜体とし、
\[
        \frN=x_1A+\cdots+x_nA
\]
を階数 \(n\) の線形形式の加群とする。また
\[
        \frL=z_1A+\cdots+z_\ell A
\]
を階数 \(\ell\le n\) の部分加群とする。\(z_i\) が \(x_i\) に関する正規基底と呼ばれるのは、それらが
\[
        z_i=x_i-(x_{l+1}\alpha_{i,l+1}+\ldots+x_n\alpha_{i,n}),
        \qquad (i=1,\ldots,l)
\]
の形をしているときである。任意の加群 \(\frL\) は、\(x_i\) の適切な番号づけの後に、このような正規基底を持つ。実際、
\[
        \frN=\frL+x_{l+1}\mathsf{A}+\ldots+x_n\mathsf{A}
\]
であり、したがって
\[
        x_i\equiv x_{l+1}\alpha_{i,l+1}+\ldots+x_n\alpha_{i,n}
        \pmod{\frL},\qquad i=1,\ldots,\ell.
\]
従って対応する \(z_i\) は \(\frL\) に属し、\(\frL\) を尽くす。

\paragraph{2. 拡大加群。}
再び \(A\) を斜体とし、\(P\) を部分体とする。階数 \(n\) の \(A\)-加群 \(N\) が \(P\)-加群 \(M\) の拡大加群と呼ばれ、\(N=M_A\) と書かれるのは、\(M\) が \(N\) の同じ階数を持つ部分加群であるときである。もし
\[
        M=y_1P+\cdots+y_nP,
\]
なら
\[
        M_A=y_1A+\cdots+y_nA.
\]
逆に、任意の \(P\)-加群 \(M\) に対して、拡大加群は加群同型を除いて一意に存在する。

\paragraph{系 a).}
\(z_1,\ldots,z_s\) が \(M\) の中で \(P\) 上線形独立なら、それらは \(M_A\) の中で \(A\) 上も線形独立である。従って \(M\) の中の任意の \(P\)-加群
\[
        T=z_1P+\cdots+z_sP
\]
は拡大加群
\[
        T_A=z_1A+\cdots+z_sA\subset M_A
\]
を生成する。

\paragraph{系 b).}
\(M\) の中の任意の \(P\)-加群 \(T\) は収縮加群、すなわちその拡大加群と \(M\) との交わりである：
\[
        T=T_A\cap M .
\]

\paragraph{3. 不変加群に関する定理。}
\(N=M_A\) を \(P\)-加群 \(M\) の拡大加群とし、\(P\) を \(A\) の環自己同型群 \(\frG\) に対する完全な不変体とする。群 \(\frG\) は \(N\) 上の作用素領域として
\[
        Gm=m,
        \qquad m\in M,
\]
および
\[
        G\Bigl(\sum m_i\alpha_i\Bigr)=\sum m_i\,G(\alpha_i)
\]
により定義される。補題：\(M\) は \(\frG\) によって固定される \(N\) の元の全体である。

\paragraph{定理。}
\(M\) の部分加群 \(L\) の拡大加群 \(L_A\)、かつそれらだけが、\(\frG\) を作用素領域とする許容部分加群である。

一方向は明らかである。逆に、\(L\) を \(\frG\) に関して許容とし、\(z_1,\ldots,z_\ell\) を \(M\) の \(P\)-基底 \(x_1,\ldots,x_n\) に関する \(L\) の正規基底とする。任意の \(G\in\frG\) に対して
\[
        G(z_i)=x_i-(x_{\ell+1}G(\alpha_{i,\ell+1})+\cdots+x_nG(\alpha_{i,n}))
\]
は再び \(L\) に属する。\(x_1,\ldots,x_\ell\) の係数を比較すると \(G(z_i)=z_i\) を得る。従って補題により \(z_i\) は \(M\) に属し、\(L\) は \(M\) の中の \(P\)-加群
\[
        T=z_1P+\cdots+z_\ell P
\]
の拡大であり、\(T=L\cap M\) である。

\subsection*{§4. 超複素系とその表現類}

\paragraph{1. 拡大定理。}
いま §2 の表現定理と §3 の加群定理を組み合わせる。任意の環 \(\frR\) の代わりに、可換体 \(P\) 上の単位元を持つ超複素系 \(S,T,\ldots\) を考える：
\[
        S=x_1\mathsf{P}+\ldots+x_n\mathsf{P}.
\]
任意の表現環 \(\frS\) の代わりに、特に断らない限り中心 \(P\) を持つ斜体 \(A,B,\ldots\) だけを考える。この場合、\(S\) と \(A\) が元ごとに可換で、交わりが \(P\) である積環が常に存在する。これは直接積 \(S\times A\) であり、同時に \(S\) の拡大加群 \(S_A\) でもある：
\[
        S_{\mathsf{A}}=x_1\mathsf{A}+\ldots+x_n\mathsf{A}.
\]
従ってこれを拡大環とも呼ぶ。

この特殊化のもとで、不変加群についての定理は次のようになる。

\paragraph{拡大定理。}
\(S_A\) の任意の両側イデアルは \(S\) のあるイデアルの拡大イデアル、すなわち \(S\) との交わりの拡大である。

実際、選んだ群 \(\frG\) を \(A\) の内部自己同型群とする。すると \(A\) の中心である \(P\) は完全な不変体であり、\(A\) と \(S\) の元ごとの可換性は、\(\frG\) が \(S_A\) へ拡張されて \(S\) が完全な不変領域となることを意味する。\(S_A\) の任意の両側イデアル \(\mathfrak a\) は、\(\alpha^{-1}\mathfrak a\alpha\subseteq \mathfrak a\) により許容部分群である。よってそれは交わり加群 \(\mathfrak a\cap S\) の拡大であり、この交わりは \(S\) のイデアルである。

補足：\(S\) が可換なら、\(S\) は \(S_A\) の中心である。より一般に、\(S\) の中心は \(S_A\) の中心でもある。\(A_r,B_n,\ldots\) は常に、\(A,B,\ldots\) 上の指定された次数の行列環を意味する：
\[
        A_r=\sum A c_{ik}=\sum c_{ik}A.
\]

本質的に用いるのは単純系に対する拡大定理である。

\(S\) が単純なら、\(S_A\) も単純である：
\[
        S_A=\sum Bc_{ik}=\sum c_{ik}B=B_t,
\]
ここで \(B\) は随伴する斜体であり、その中心は \(P\) を含む。ここで \(A\)、したがって \(B\) は \(P\) 上無限階数でもよい。仮定するのは \(S\) が超複素であることだけである。もし \(S\) と \(A\) が共通の中心 \(P\) を持つなら、\(P\) は \(S_A\) の中心でもある。より一般に、\(S\) が単純な超複素系なら、直接積 \(S\times A_r\) は単純である。

\paragraph{2. 表現類。}
超複素系の相反表現または直接表現とは、通常どおり、零表現ではなく、係数体 \(P\) に関して作用素準同型な表現をいう。

\paragraph{表現類に関する定理。}
\(S\) を単位元と係数体 \(P\) を持つ超複素系とし、\(A\) を中心 \(P\) を持つ斜体とする。このとき、\(S\) の \(A\) における相異なる既約相反表現類の数は、\(S_A\) をその根基で割った商環における単純右イデアル類の数に等しい。\(S\) が単純系なら、\(A\) において既約相反表現類を一つ、既約直接表現類を一つだけ持つ。

移行定理の系により、単純な相反表現類と \(S_A\) 上の単純加群類は一対一に対応する。\(S_A\) は斜体 \(A\) 上有限階数を持つので、片側イデアルについて最大条件と最小条件を満たす。\(S\) が単純なら \(S_A\) も単純である。従って単純な片側イデアル類は一つだけであり、既約相反表現類も一つだけである。相反系は直接表現についての対応する主張を与える。

\paragraph{3. 階数関係。}
\(S\) が単純で、
\[
        S_A=\sum Bc_{ik}
\]
が \(A\) 上および \(B\) 上の双方で有限階数を持つなら、
\[
        n=rt,
        \qquad n=(S:P),
        \qquad t^2=(S_A:B),                                             \tag{1}
\]
が成り立つ。ここで \(r\) は既約相反表現の次数である。

\paragraph{4. \(A\) が \(P\) 上有限階数を持つ場合の階数関係の精密化。}
この場合には三つの関係
\[
        (S:P)=n=rt,                                                     \tag{1}
\]
\[
        (B:P)\,t=(A:P)\,r,                                               \tag{2}
\]
\[
        (S:P)(B:P)=(A:P)\,r^2                                           \tag{3}
\]
がある。ここで (2) は \(S_A\) の単純右イデアルの \(P\) 上の階数を、\(B\) 上および \(A\) 上の階数で表したものである。(3) は (1) を用いて (2) に \(r\) を掛けることで従う。

\subsection*{§5. 表現定理の行列環への応用}

§4 で考えた \(S\) の \(A\) における相反表現は、\(S\) から \(A_r\) の部分環への直接準同型としても見られる。ここで \(\overline A\) は \(A\) に相反同型な斜体を表す。以下の段落の原理はその逆である。すなわち行列環 \(A_r\) の研究を、相反同型な斜体 \(A\) における表現論へ還元する。

\paragraph{1. \(A_r\) への可約および既約な埋め込み。}
\(A\) と \(\overline A\) を、その中心 \(P\) 上有限または無限階数を持つ相反同型な斜体とする。\(P\) 上の単純系 \(S\) が \(A_r\) に既約に、または可約に埋め込めるというのは、\(S\) が \(A\) において次数 \(r\) の既約、または可約な相反表現を持つときである。\(S\) が \(A_r\) に既約に埋め込めるなら、すべての行列環 \(A_{rs}\), \(s>1\), に可約に埋め込める。この同型は常に \(P\) の恒等写像を延長する。§4, 3 により、\(r\) は \(S\) の \(P\) 上の階数 \(n\) を割る。

このようにして、\(P\) を含み \(P\) 上有限階数を持つすべての単純部分環が、種々の \(A_r\) の中で得られる。実際、そのような任意の部分環は \(\overline A_r\) の部分環に相反同型であり、したがって \(A\) において可約または既約な相反表現を持つ。

\paragraph{2. 内部自己同型に関する定理。}
\(S^{(1)}\) と \(S^{(2)}\) を、\(P\) を含み \(P\) 上有限階数を持つ \(A_f\) の二つの単純部分環とする。もし \(S^{(1)}\) と \(S^{(2)}\) が \(P\) 上同型なら、この同型は \(A_f\) の内部自己同型によって誘導される。

実際、\(S^{(1)}\) と \(S^{(2)}\) は単純系 \(S\) の \(A_f\) への一般には可約な埋め込みを与え、従って \(A\) における次数 \(f\) の直接表現を与える。それらは同じ可約直接表現類に属し、変換行列が自己同型を誘導する。\(A\) 自身が \(P\) 上有限階数、従って超複素である場合には、これはよく知られた定理、すなわち中心 \(P\) を含み \(P\) 上同型な \(A_f\) の二つの単純部分系は \(A_f\) の内部自己同型によって互いに移される、という定理になる。とくに \(A_f\) の任意の自己同型は内部的である。

\paragraph{3. 元ごとに可換な部分環に関する定理。}
\(S\) を \(A_f\) に含まれ \(P\) を含む単純系とし、\(R\) を \(A_f\) の元のうち \(S\) と元ごとに可換なものすべての全体とする。このとき \(R\) も行列環である：
\[
        R=B_s.
\]
\(S_A\) の随伴斜体 \(B\) と \(R\) の随伴斜体 \(B\) は相反同型である。\(R\) と \(S\) の交わりは \(S\) の中心である。環 \(R\) が斜体であることと、\(S\) が \(A_f\) に既約に埋め込まれていることとは同値である。

§2, 3 により、環 \(R\) は、\(S\) の \(A\) における相反表現加群 \(\frM\) の自己同型環に直接同型である。これは、\(\frM\) を \(S_A\)-加群と見たときの自己同型環である。\(\frM\) が \(s\) 個の単純 \(S_A\)-加群へ分解し、§4, 1 と同じく
\[
        S_A=\sum Bc_{ik}
\]
と書けば、\(R\) は \(\sum Bc_{ik}\) に同型であり、二つの斜体は相反同型である。したがって \(R=B\) が斜体であるのは、\(s=1\)、すなわち \(S\) が既約に埋め込まれるとき、かつそのときに限る。定義により、\(R\cap S\) は \(S\) の中心である。

\paragraph{4. 超複素行列環に対する精密な交換定理。}
この場合、§4, 4 の階数関係は次の精密化を与える。

\(A\) が中心 \(P\) 上有限階数を持つとき、\(P\) を含む \(A_f\) の単純部分系は \(S,\overline S\) の対に分かれ、\(\overline S\) は \(S\) と元ごとに可換な元全体からなり、その逆も同様である。\(S_A\) と \(\overline S\) の随伴斜体は相反同型であり、\(S\) と \(\overline S_A\) の随伴斜体も同様である。\(S\) と \(\overline S\) の交わりは共通中心である。一方の部分環が斜体であることと、他方が既約に埋め込まれることとは同値である。\(S\) と \(\overline S\) の \(P\) 上の階数の積は \(A_f\) の階数である。

\(f=rs=r's'\) で、\(S\) が \(A_r\) に、\(\overline S\) が \(A_{r'}\) に既約に埋め込めるなら、\(S\) と \(\overline S\) の随伴斜体は、\(g=ss'\) として、\(A_g\) の中の可換化する部分環の一対に同型となる。

積に関する主張は階数関係から直接従う。\(S\) が既約に埋め込まれているなら、\(\overline S\) は斜体であり、\(S_A=B_t\) における \(B\) に相反同型である。§4, 4 の階数関係 (3) が主張を与える。\(S\) が可約に埋め込まれていて、\(\overline S=B_s\) なら、\(f=rs\) であり、(3) に \(s^2\) を掛けると主張を得る。これにより \(\overline S\) が \(S\) と元ごとに可換な元全体であることも示され、残りの主張は可換化する部分環に関する定理から従う。

\subsection*{§6. 単純系のガロア理論}

\paragraph{1. 中心上有限階数の斜体のガロア理論。}
\(A\) のガロア群 \(\Gcal\) は、中心 \(P\) の恒等写像を延長するすべての自己同型の群、したがって §5, 2 によりすべての内部自己同型の群として定義される。従って
\[
        \Gcal \simeq A^*/P^*,
\]
ここで \(A^*\) と \(P^*\) は \(A\) と \(P\) の非零元の乗法群を表す。したがって \(\Gcal\) の部分群 \(\mathfrak H\) は、\(P^*\) を含む \(A^*\) の部分群 \(H^*\) と
\[
        \mathfrak H\sim H^*/P^*
\]
によって一対一に対応する。\(\Gcal\) の部分群 \(\mathfrak H\) は、零を添えると \(H^*\) が斜体 \(H\) になるとき閉じていると呼ばれる。\(C\) が群 \(\mathfrak H\) を許すことは、\(C\) が \(H\) と元ごとに可換であることと同値である。従って §5, 4 の交換定理は次の形をとる。

\paragraph{非可換斜体のガロア理論の主定理。}
\(P\) と \(A\) の間の斜体 \(C\) と \(\Gcal\) の閉部分群 \(\mathfrak H\) とは、\(C\) が \(\mathfrak H\) の完全な不変斜体であり、\(\mathfrak H\) が \(C\) の完全な不変群であるように、一対一に対応する。

\paragraph{2. 単純系のガロア理論。}
\(A_f\) が中心 \(P\) を持つ単純系なら、ガロア群 \(\Gcal\) はふたたび内部自己同型群であり、従って
\[
        A_f^*/P^*
\]
に同型である。ここで \(A_f^*\) は \(A_f\) の正則元の乗法群を表す。\(\Gcal\) の部分群 \(\mathfrak H\sim H^*/P^*\) は、\(H^*\) によって生成される部分環 \(H\) が単純系で、かつ \(H^*\) が \(H\) の正則元全体であるとき、単純閉と呼ばれる。交換定理からガロア理論への移行を与えるのは次の補題である。

\paragraph{補題。}
\(P\) を含む \(A_f\) の任意の単純系は、その正則元の群 \(H^*\) によって生成される。これは \(P\) が無限個の元を持つ場合には明らかである。不定元で作った「一般元」は正則であり、適当な特殊化により \(P\) 上の正則な基底元が得られる。一般の場合にも補題は Shoda によって成り立つ。\footnote{注 3) で引用された Shoda の仕事。そこでは不定元の導入が避けられている。}

補題により、\(S\) が単純系 \(\mathfrak S\) と可換化することは、\(S\) が \(\mathfrak S\) の正則元によって誘導される自己同型を許すこととふたたび同値である。従って §5, 4 の交換定理はここでも次の形をとる。

\paragraph{単純系のガロア理論の主定理。}
\(P\) を含む \(A_f\) の単純系 \(S\) と \(\Gcal\) の単純閉部分群 \(\mathfrak H\) とは、\(S\) が \(\mathfrak H\) の完全な不変領域であり、\(\mathfrak H\) が \(S\) の完全な不変群であるように、一対一に対応する。

\paragraph{3. 拡大原理による証明。}
斜体の場合について、主定理の第二の証明を与える。それは同型、すなわち一次表現を延長する原理に基づいており、階数の数え上げを使わないので有限性仮定が少ない。とくに、無限階数の斜体 \(S\) への移行がどの方向で可能かを示す。この証明は既約相反表現類が一つだけであるという事実も用いないため、可換の場合 §8, 2 にもまったく同じように適用される。まずいくつかの補題を述べる。

\paragraph{仮定。}
\(P\) 上の単純系 \(S\) は \(A\) における一次の相反表現を許すものとし、従って斜体である。ここで \(A\) は中心 \(P\) 上有限階数でも無限階数でもよい。§4, 3 のように、\(S_A\) が \(n\) 個の単純で作用素同型な右イデアルへ分解されるものとし、
\[
        S_{\mathsf{A}}=\mathfrak r_1+\ldots+\mathfrak r_n
            =e_1S_{\mathsf{A}}+\ldots+e_nS_{\mathsf{A}}
            =e_1\mathsf{A}+\ldots+e_n\mathsf{A} ;
\]
\(e_i\) によって生成される相反表現は
\[
        e_i s=e_i\sigma
\]
により定義される。

\paragraph{補題 1.}
\(e_1,\ldots,e_n\) によって生成される \(n\) 個の相反表現は相異なる。したがって同一の表現類に属する相異なる表現である。従って \(S\) はその階数と少なくとも同数の相異なる表現を持つ。

証明のため、§2 の終わりと同じく、\(S\) からの表現を、\(A\) 上で作用素準同型な \(S_A\) の対応へ延長する。ここで対応は各 \(w\in S_A\) に対して
\[
        e_iw=e_i\omega
\]
（\(\omega\in A\)）によって定義される。もし \(e_i\) と \(e_j\), \(i\ne j\), が同じ表現を生成したなら、すべての \(w\in S_A\) について \(e_iw=e_jw\) でなければならない。これは \(e_ie_i=e_i\) かつ \(e_je_i=0\) に反する。

\paragraph{補題 2.}
\(T\) を \(P\) と \(S\) の間の斜体とし、\(s\) を \(S\) の \(T\) 上の階数とする。このとき \(T\) から \(A\) への任意の相反同型は、少なくとも \(s\) 個の異なる延長を許す。

その相反同型、すなわち \(T\) の \(A\) における相反表現が、分解
\[
        T_{\mathsf{A}}=E_1T_{\mathsf{A}}+\ldots+E_hT_{\mathsf{A}}
            =E_1A+\cdots+E_hA,
        \qquad E_it=E_i\tau
\]
によって媒介されるとする。ここで \(E_i\) は冪等元である。これは制限ではない。基底元 \(E_i\alpha\) によって生成される表現は \(\alpha^{-1}E_i\alpha\) によって、従って冪等元によっても生成されるからである。\(S\) による右乗法は
\[
        S_A=E_1S_A+\cdots+E_hS_A
\]
を与える。各 \(E_iS_A\) は \(A\) 上階数 \(s\) を持つ。実際、\(S\) に \(T\)-基底を代入し、\(S\) と \(A\) の可換性を用いれば、\(E_iS_A\) の階数は高々 \(s\) である：
\[
        S=Tw_1+\cdots+Tw_s,
\]
\[
        E_iS_{\mathsf{A}}=E_iT_{\mathsf{A}}w_1+\ldots+E_iT_{\mathsf{A}}w_s
              =E_iw_1\mathsf{A}+\ldots+E_iw_s\mathsf{A} .
\]
全体の和 \(S_A\) は階数 \(n=hs\) を持つので、各 \(E_iS_A\) はちょうど階数 \(s\) を持つ。従って
\[
        E_iS_A=r_{i1}+\cdots+r_{is}
              =e_{i1}A+\cdots+e_{is}A,
\]
であり、補題 1 により \(e_{i1},\ldots,e_{is}\) によって生成される \(s\) 個の表現はすべて相異なる。それらはすべて \(E_i\) によって生成される \(T\) の表現の延長である。なぜなら \(E_it=E_i\tau\) は
\[
        (e_{i1}+\cdots+e_{is})t=(e_{i1}+\cdots+e_{is})\tau
\]
を含意し、直和分解により \(e_{ij}t=e_{ij}\tau\) が従うからである。

\paragraph{定義。}
\(T\) によって誘導される、\(S\) から \(A\) への相反同型 \(\mathfrak I\) の分割 \(\mathfrak I_T\) を次のように定義する。\(\mathfrak I\) の二つの同型を、それらが \(T\) 上で同じ同型を誘導するとき、かつそのときに限って同値とみなす。

\paragraph{主定理。}
\(\mathfrak I_T\) が \(T\) によって誘導される分割なら、\(T\) はこの分割に関して極大部分斜体である。

実際、\(Z\) が \(T\) と \(S\) の間の中間斜体で、\(T\) 上の次数が \(l\) なら、補題 2 により \(\mathfrak I_T\) の各類は \(\mathfrak I_Z\) の少なくとも \(l\) 個の類に分かれる。この形から通常の形への移行は、集合 \(\mathfrak I\) の元を、固定したその一つの逆で掛けることによって得られる。集合は自己同型群となり、\(\mathfrak I_T\) の一つの類は \(T\) の不変群となる。閉部分群が完全な不変群であるという主張もこれに含まれる。\(\mathfrak H\sim H^*/P^*\) については、中間斜体 \(T\) の代わりに斜体 \(H\) を置けばよい。

\subsection*{§7. 類似代数の類。分解体}

これ以後は、中心 \(P\) 上有限階数を持つ斜体 \(A,\overline A\) と行列環 \(A_f\) だけを扱う。従ってそれらは \(P\) 上の単純正規代数であり、\(P\) 上の代数、または単に代数と呼ぶ。\(A,\overline A\) は除法代数である。同じ随伴 \(A\) を持つすべての \(A_f\) は、類似代数 \(A,A_f\) の一つの類へまとめられる。同型な \(A\) は同一視する。

\paragraph{1. 代数類の群。}
\(A,B\) が \(P\) 上の代数なら、それらの積、すなわち \(P\) 上の直接積についても同じことが成り立つ。§3 の終わりにより
\[
        A_r\times B_s=C_t,
\]
ここで \(C\) は中心 \(P\) を持つ除法代数で、同型を除き一意に定まる。\(A_f,B_g\) に \(P\) 上の行列環を掛けることにより、この乗法が類について一意に定まることがわかる。したがって類は乗法的に閉じた可換な系をなす。さらに R. Brauer の定理が成り立つ。類似代数の類は直接積に関してアーベル群をなす。この系は恒等類を持ち、それは \(P\) 上の行列環、すなわち \(P\) に類似な代数からなる。また各類 \((A)\) は逆類 \((\overline A)=(A)^{-1}\) を持つ。これは \(A\) に相反同型な \(\overline A\) に類似な代数からなる。このことは、§5, 3 の交換定理において \(S=A\) と特殊化することから従う。実際 \(A\) は自身に既約に埋め込めるので、\(B=P\) および \(A_A=P_f\) が得られる。\footnote{代数類の群に関するより精密な定理は因子系の理論を用いる。ここでは扱わない。ここまで、可換体に関する考察は一切現れていないことに注意せよ。例えば、階数 \((A:P)\) が平方であるという事実は用いられても証明されてもいない。}

\paragraph{2. 代数類の分解体。}
分解体の理論は、§5 の初めで述べた原理、すなわち可換拡大体を \(A\) または \(\overline A\) の中で表現するという原理に、ふたたび完全に基づく。まず次の補題を必要とする。

\paragraph{補題。}
\(A\) を中心 \(P\) を持つ斜体とし、\(Z\) を \(P\) の有限可換拡大体とする。このとき \(A_Z\) は中心 \(Z\) を持つ単純系である。実際 §4, 1 によりこれは \(A_Z\) に相反同型な系 \(Z_{\overline A}\) について成り立つ。従って、\(A\) に類似な \(P\) 上の代数類 \((A)\) に対し、体 \(Z\) は \(A_Z\) に類似な \(Z\) 上の単純代数の拡大類 \((A)_Z\) を与える。

拡大類の随伴除法代数 \(D\) は、§5, 3 の交換定理からただちに従う。

\paragraph{拡大類の随伴除法代数に関する定理。}
\((A)_Z\) を拡大類とし、\(Z\) を \(A_f\) における \(Z\) の既約埋め込み、すなわち表現とする。このとき
\[
        (A)_Z=(D),
\]
ここで \(D\) は \(A_f\) の元のうち \(Z\) と元ごとに可換なものすべての集合である。§5, 3 の単純系 \(S\) を \(Z\) に置き換え、\(Z_A\) と \(A_Z\) が相反同型であることを観察すればよい。\(Z\) が可約に埋め込まれるなら、\(Z\) と元ごとに可換な元全体は \(D\) 上の行列環である。

\(P\) の可換拡大体 \(Z\) は、拡大類 \((A)_Z\) が完全に分解する、すなわち \(Z\) 上の恒等類、つまり \(Z\) 上の全行列環の類になるとき、類 \((A)\) の分解体と呼ばれる。

\paragraph{分解体の特徴づけに関する定理。}
\(P\) の有限可換拡大体 \(Z\) が類 \((A)\) の分解体であることは、その \(A_f\) への既約埋め込みが \(A_f\) の極大可換部分体を与えることと同値である。

実際、\(Z\) が \((A)\) の分解体であるという性質は \((D)=(Z)\) と同値である。従って随伴除法代数の定理により、既約埋め込み \(Z\) は極大可換部分体でなければならない。この条件が満たされるなら、任意の \(a\in D\) を \(Z\) に添加すると可換拡大体を生じるので、必ず \(D=Z\) である。

\paragraph{注意。}
1. 同時に、極大可換部分体 \(Z\) は、既約に埋め込まれているとき、極大可換部分環を生成することが従う。なぜなら \(Z\) と元ごとに可換な元は斜体をなすからである。

2. この定理は分解体の存在も与える。随伴除法代数 \(A\) は確かに極大可換部分体を持つからである。

\paragraph{3. 階数関係、指数、無限可換拡大。}
§5, 4 の階数の主張はまず、中心上の単純代数の階数が平方であるという既知の事実を与える。\(Z\) が \(A\) の極大可換部分体なら、\(A\) へのすべての埋め込みは既約であるから、
\[
        (Z:P)(Z:P)=(A:P),
\]
従って
\[
        (A:P)=m^2,
        \qquad (Z:P)=m,
\]
ここで \(m\) は Schur の指数であり、\(A\) 自身の中に埋め込めるすべての分解体の次数でもある。一般の分解体の次数 \(n\) については
\[
        n=mr
\]
を得る。これは例えば \((A_r:P)=m^2r^2\) から、または §4, 3 の階数関係から従う。指数 \(m\) はまた、\(A_Z\) の単純成分の数でもあり、これは \(Z\) 上の階数 \(m^2\) の行列環になるからである。より一般に \(A_\Lambda\sim D\) であり、\(L\) が \(\Lambda\) の \(A_r\) への既約埋め込みで、\(d^2\) が中心 \(L\) 上の除法代数 \(D\) の階数 \((D:L)\) であるなら、\(l\cdot d=m\cdot r\) となる。

実際、§5, 4 と 2 の随伴除法代数の定理により
\[
        (A_r:P)=(L:P)(D:P),
\]
したがって
\[
        m^2r^2=l^2d^2 .
\]

分解体の存在から、\(P\) の無限な可換、代数的または超越的拡大体 \(\Omega\) について次が従う。拡大類 \((A)_\Omega\) は中心 \(\Omega\) を持つ単純代数の類である。とくに \(\Omega\) が代数的閉体なら、\(A_\Omega\) は次数 \(m\) の行列環である。従って指数は「絶対成分数」に等しい。実際、\(Z\) が \(A\) の分解体で、\(\Omega'\) が \(\Omega\) と \(Z\) の合成体なら、\(A_{\Omega'}\) は \(\Omega'\) 上の行列環であり、従って根基を持たず、中心は \(\Omega'\) である。したがって \(A_\Omega\) も根基を持たず、その中心は \(\Omega\) である。\(A_\Omega\) の中心に属し \(\Omega\) に属さない元があれば、それは \(A_{\Omega'}\) の中心にも属してしまう。従って \(A_\Omega\) は \(\Omega\) 上の単純代数であり、\(\Omega\) が代数的閉体なら確かに \(\Omega\) 上の行列環である。\footnote{一般には「極小」分解体、すなわち真の部分体が分解体でないような分解体も、あらゆる可能な次数で存在する。注 2) で引用された Brauer--Noether を参照。}

\paragraph{4. 分離的分解体の存在。}
任意の類 \((A)\) は分離的分解体を持ち、しかも \(A\) 自身の中に埋め込めるものを持つ。\footnote{G. Köthe, Über Schiefkörper mit Unterkörpern zweiter Art über dem Zentrum, Journ. f. Math. 166 (1932), pp. 182--184 を参照。このずっと簡単な証明は M. Zorn の観察に由来する。}

基礎体 \(P\) の標数を \(p\) とし、指数を
\[
        m=s p^f
\]
と書く。ただし \(s\) は \(p\) と素である。\(Z\) が \(A\) の次数 \(m\) の分解体で、\(Z_1\) がその中に含まれる第一種拡大、すなわち
\[
        (Z:Z_1)=p^f
\]
を満たすものなら、\(A_{Z_1}\sim D\) の指数は \(p^f\) である。\(D_{Z_2}\) の指数が \(p^f\) の真の約数になるような、\(Z_1\) の分離的拡大体 \(Z_2\) の存在を示す。\(\Omega\) を代数的閉体とすると、3 により \(D_\Omega\) は根基を持たない。従って \(D\) の被約判別式は消えない。よって被約跡が非零である元 \(d\in D\) が少なくとも一つ存在する。この \(d\) はすでに基礎体 \(Z_1\) に属することはできない。なぜなら \(Z_1\) の任意の元 \(\alpha\) の跡は \(p^f\alpha\) であり、従って消えるからである。ゆえに \(Z_2=Z_1(d)\) は真の、しかも分離的な拡大体であり、\(D_{Z_2}\) の指数は \(p^f\) の真の約数である。これを有限回繰り返すと、次数 \(m\) が \(A\) の指数に一致する \((A)\) の分離的分解体が得られる。

\subsection*{§8. 可換の場合における分解体、部分分解体、ガロア理論}

中心上単純な系の分解体から任意の単純系の分解体へ移るためには、なお中心の分解理論を展開し、それに §6, 3 と平行なガロア理論を加えなければならない。この段落に現れるすべての体と系は可換である。

\paragraph{1. 可換単純系の分解体と部分分解体。}
可換系 \(Z\) を \(P\) 上単純、従って体とし、\(\Omega\) を \(P\) の代数的閉拡大体とする。\(Z_\Omega\) が根基を持たない系であり続けることは、\(Z\) が \(P\) 上分離的、すなわち第一種拡大であることと同値であることが知られている。このときに限り、\(P\) 上の階数と同じだけの一次同型写像を \(\Omega\) へ持ち、従って同じだけの相異なる表現を持つ。この場合それらは表現類と一致する。

\(Z_\Omega\) に根基が現れる可能性、従って絶対単純成分への直和分解が必ずしも起こらない可能性は、次の定義へ導く。\(P\) の拡大体 \(\Lambda\) は、\(Z_\Lambda\) が一次の組成因子へ分解する、言い換えればすべての絶対既約表現がすでに \(\Lambda\) にあるとき、分解体と呼ばれる。\(P\) の拡大体 \(T\) は、\(Z_T\) が少なくとも一つの一次組成因子を分離する、言い換えれば \(Z\) の少なくとも一つの絶対既約表現が \(T\) に存在するとき、部分分解体と呼ばれる。

分解体と部分分解体は次の定理によって特徴づけられる。

\paragraph{定理。}
体 \(T\) が部分分解体であることは、\(T\) が \(Z\) に同型な部分体 \(Z'\) を含むことと同値である。体 \(\Lambda\) が分解体であることは、\(Z\) に属するガロア体に同型な部分体 \(\Gamma\) を含むことと同値である。

これは、絶対既約表現による分解体と部分分解体の定義から直接従う。とくに非可換の場合との類比で次の系を得る。体 \(Z'\) が極小部分分解体、すなわち真の部分体が部分分解体でないものであることは、\(Z'\) が \(Z\) に同型であることと同値である。極小部分分解体が同時に分解体であることは、\(Z\) が正規、従って \(P\) 上ガロア体であることと同値である。

これが、単純代数をその中心上正規と呼ぶ理由である。\(P\) 上の固定された代数的閉体 \(\Omega\) の中では、極小分解体はただ一つ、すなわち \(Z\) に属するガロア体だけである。これは、非可換の場合に一般に無限個の非同型体があることと対照的である。その理由は、可換の場合の直和分解の一意性にあり、非可換の場合にはそれが作用素同型を除いてしか一意でないことにある。同じことを別の形でいえば、可換の場合には相異なる絶対既約表現が有限個しかないのに対し、非可換の場合には同じ類に属していても相異なるものが無限個あるからである。

\paragraph{2. 分離的体 \(Z/P\) の場合のガロア理論の超複素的証明。}
§6, 3 と正確に同じ類比で、任意の \(P\) 上分離的な \(Z\) について同型の理論が成り立つ。\(Z\) がガロアなら、通常の自己同型群へ移ることができる。

\paragraph{仮定。}
\(P\) 上の単純系 \(Z\) は有限分離拡大体であり、\(\Omega\) は \(P\) 上有限または無限の分解体、\(Z'=Z^{(1)}\) は \(Z\) に同型な部分体、\(\Gamma\) は対応するガロア体を表す。1 により直和分解
\[
        Z_\Omega=r^{(1)}+\cdots+r^{(n)}
        =e^{(1)}Z_\Omega+\ldots+e^{(n)}Z_\Omega
        =e^{(1)}\Omega+\ldots+e^{(n)}\Omega 
\]
がある。\(e^{(i)}\) によって生成される表現 \(Z\to Z^{(i)}\), ただし \(Z^{(i)}\leqq\Gamma\), は
\[
        e^{(i)}z=e^{(i)}z^{(i)}
\]
により定義される。

\paragraph{補題 1.}
\(e^{(1)},\ldots,e^{(n)}\) によって生成される \(n\) 個の表現はすべて相異なる。証明は §6, 3 と同じであり、また 1 からも従う。それらは相異なる表現類に属する。

\paragraph{補題 2.}
\(T\) が \(P\) と \(Z\) の間の中間体で、\(s\) が \(Z\) の \(T\) 上の階数なら、\(T\) から \(\Omega\) への任意の同型は少なくとも、従って正確に、\(s\) 個の延長を許す。

証明は §6, 3 と同じである。「少なくとも」が「正確に」へ強められるのは、\(Z_\Omega\) の分解の一意性により、\(Z\) が \(\Omega\) の中に少なくともではなく正確に \(n\) 個の同型を持つからである。

§6, 3 と同じく、\(Z\) の同型からなる有限集合 \(\mathfrak I\) の、\(T\) によって誘導される分割 \(\mathfrak I_T\) を定義する。すなわち、\(\mathfrak I\) のうち \(T\) 上で同じ同型を誘導するものを、ちょうど同値とみなす。すると次を証明できる。

\paragraph{主定理。}
\(\mathfrak I_T\) が \(T\) によって誘導される分割なら、\(T\) はこの分割に関する極大部分体である。

ここから、\(Z\) がガロアである場合には、§6, 3 とまったく同じように、同型の逆と合成することによって自己同型群と通常の定式化へ移る。部分群についての対応する主張は、冪等元の考察から従う。この考察はそれ自体にも意義がある。

\paragraph{3. \(Z_\Omega\) の冪等元。}
\(S\) を、\(Z\) を個々の \(Z^{(i)}\) へ送る \(n\) 個の置換とする。つまり、\(Z\to Z_1\) と \(Z\to Z^{(i)}\) の同型の合成であり、前者は \(Z\to Z_1\) の相反である。体 \(Z\) はガロアである必要はない。\(a\in Z_Z\) に対し、置換 \(S\) は \(Z\) 上恒等的であると約束することにより作用素として定義される。したがって各 \(a\) について、\(Z_{Z^{(i)}}\) に属する共役 \(a^S\) が定義される。この約束のもとで次の定理が成り立つ。

\paragraph{定理。}
\(Z_\Omega\) の \(n\) 個の冪等元 \(e^{(i)}\) は共役である：
\[
        e^{(i)}=e^S,
        \qquad e=e^{(1)};
\]
それらは \(n\) 個の共役な拡大系 \(Z_{Z^{(i)}}\) に属する。

実際、\(S\) を作用させると、定義関係
\[
        ez=ez^{(1)},
        \qquad e^2=e
\]
は
\[
        e^S z=e^S z^{(i)},
        \qquad (e^S)^2=e^S
\]
へ移る。分解の一意性により冪等元自身も一意に定まるので、\(e^S=e^{(i)}\) である。さらに \(Z\) は部分分解体でもあり、表現 \(ez=ez^{(1)}\) はすでに媒介されているから、\(e\) はすでに \(Z_Z\) に属し、従って \(e^S\) は \(Z_{Z^{(i)}}\) に属する。

\(Z\) がガロアなら、ガロア理論の主定理のうち部分群に関する部分はただちに従う。\(\mathfrak H\) がガロア群 \(\Gcal\) の部分群なら、直和の成分の一意性により、元
\[
        \sum e^H
\]
は \(\mathfrak H\) のすべての置換を許し、それ以外は許さない。群 \(\Gcal\) は \(Z\) に対して定義されたので、これは \(Z\) のガロア理論であり、同型により \(Z'\) の理論も決定する。

\paragraph{4. 冪等元と補基底との関係。}
ここでも \(Z\) がガロアであるとは仮定しない。

\paragraph{定理。}
\(a_1,\ldots,a_n\) を \(P\) 上の \(Z\) の基底とし、冪等元を
\[
        e=a_1\beta_1+\cdots+a_n\beta_n
\]
と書く。すると
\[
        e^S=a_1\beta_1^S+\cdots+a_n\beta_n^S .
\]
このとき
\[
        \beta_1^S,\ldots,\beta_n^S
\]
は、\(e^S\) によって媒介される写像のもとで、\(a_1,\ldots,a_n\) に相補的な基底 \(b_1,\ldots,b_n\) の像である。

\(e^S\) によって媒介される写像では \(e^Sx=e^Sx^S\)、とくに \(e^Sa_i=e^Sa_i^S\) であり、関係
\[
        e^Se^S=e^S,
        \qquad e^Se^R=0\quad(S\ne R)
\]
は割当て \(e^S\mapsto1\), \(e^R\mapsto0\) を与える。従って
\[
        1=a_1^S\beta_1^S+\cdots+a_n^S\beta_n^S,
        \qquad
        0=a_1^R\beta_1^S+\cdots+a_n^R\beta_n^S\quad(S\ne R).
\]
これを行列で書くと、相補基底の定義式
\[
\begin{pmatrix}
 a_1^1 & \cdots & a_n^1\\
 \vdots & \ddots & \vdots\\
 a_1^n & \cdots & a_n^n
\end{pmatrix}
\begin{pmatrix}
 \beta_1^S\\ \vdots\\ \beta_n^S
\end{pmatrix}
=
\begin{pmatrix}
0\\ \vdots\\ 1\\ \vdots\\ 0
\end{pmatrix}
\]
である。跡による定義への移行は、
\[
        a_i=\sum_S a_i^S e^S,
        \qquad b_i=\sum_S b_i^S e^S
\]
について、行を入れ替えた後の行列関係が
\[
        \delta_{ij}=\Sp(a_i b_j)
\]
になることから従う。ここで \(a_i\) と \(b_i\) はすべての置換 \(S\) を許すので \(Z\) に属する。\footnote{表現論 §25 を参照。}\footnote{相補基底と冪等元との関係は Dedekind, Zur Theorie der aus \(n\) Haupteinheiten gebildeten komplexen Größen に初めて現れる。全集第 II 巻、pp. 4--5 を参照。}

相補基底の反変性は、ここでは冪等元 \(e^S\) が \(Z\) の不変元であり、選ばれた基底に依存しないことから従う。したがって、
\[
        a_1\beta_1^S+\cdots+a_n\beta_n^S
\]
はすべて、\(Z/P\) の別の基底
\[
        a_1',\ldots,a_n'
\]
へ移っても自分自身に変換される。

\subsection*{§9. 任意の系の分解体と部分分解体}

\paragraph{1. 定義と単純な場合への還元。}
§8 で与えた定義に、一般の場合には次の定義が対応する。\(S\) を \(P\) 上超複素とする。\(P\) の可換拡大体 \(\Lambda\) は、片側イデアルによる組成列において \(S_\Lambda\) が絶対単純な組成因子へ分解するとき、\(S\) の分解体と呼ばれる。言い換えれば、\(S\) の \(\Lambda\) におけるすべての既約表現がすでに絶対既約であるときである。\(P\) の拡大体 \(T\) は、\(S_T\) が少なくとも一つの絶対単純組成因子を分離するとき、同値にいえば、\(T\) において既約な少なくとも一つの表現がすでに絶対既約であるとき、部分分解体と呼ばれる。

これらの定義は、可換の場合について §8 で、単純代数について §7 で与えられた定義を明らかに含む。後者については、中心の任意の可換拡大に対して \(A_\Lambda\) が完全可約であり、従って組成因子と直和分解の成分が一致するからである。中心上の全行列環、そしてそれだけが、絶対単純成分を与え、従って絶対既約表現も与える。さらに \(A_\Lambda\) は両側単純であり、作用素同型な成分へ分解する。従って一つの絶対単純成分を分離することはすでに完全分解を与え、部分分解体と分解体は一致する。

定義からただちに次の事実が従う。\(S\) の分解体は、\(P\) において既約な各表現に対して一つずつ分解体を取った合成体である。\(S\) の部分分解体は、\(P\) において既約な表現の任意の部分分解体である。単純系、すなわち根基で割った剰余類環の成分と、\(P\) において既約な表現との間には一対一対応があるので、単純系を考えれば十分である。その場合の結果は §7 と §8 を組み合わせて得られる。

\paragraph{2. 単純系の分解体と部分分解体。}
まず、\(P\) 上分離的な中心 \(Z\) を持つ単純系 \(A\) の場合を考える。§8, 1 および §8, 3 から次が得られる。\(\Omega\) を \(Z\) の分解体とすると、\(Z_\Omega\) の直和分解に対応する \(A_\Omega\) の両側分解は、\(A\) にそれぞれ同型な \(n\) 個の共役成分
\[
        e^S A_\Omega
\]
を与え、それらはそれぞれ中心
\[
        e^S Z_\Omega=e^S Z
\]
上の単純代数となる。置換 \(S\) は、\(A\) 上恒等的であると約束することにより、\(A_\Omega\) に対する作用素として定義される。

実際、§8, 3 により冪等元は共役であるから、\(A\) に対する置換の定義により成分 \(e^S A_\Omega\) も共役である。\(A\) は両側単純なので、\(e^S A_\Omega\) は環として \(A\) に同型である。さらに \(e^SA_\Omega=e^S A_{Z^S}\) である。なぜなら \(e^S Z_\Omega=e^S Z\) だからである。従って中心は係数領域と一致し、成分は正規代数である。

§7, 2 の結果からさらに次が従う。\(A\) が \(P\) 上の単純系で、中心 \(Z\) が \(P\) 上分離的なら、任意の有限または無限拡大体 \(\Omega\) に対して \(A_\Omega\) も根基を持たず、従って完全可約である。実際、§7, 2 により \(e^S A_\Omega\) は任意の係数拡大の後も単純であり、根基を持たない。従って直和についても同じである。この直和は \(A_\Omega\) であり、また \(\Omega\) が \(Z\) の分解体でないときには \(A_\Omega\) から係数拡大によって得られる。

§7, 2 から次の特徴づけも従う。

\paragraph{部分分解体と分解体の特徴づけ。}
\(A\) が分離的な中心 \(Z\) を持つ除法代数なら、部分分解体の既約埋め込みは、\(Z\) を含む \(A_f\) の極大可換部分体のすべて、かつそれだけによって与えられる。分解体は、部分分解体と中心の分解体、特に中心に属するガロア体との合成体のすべて、かつそれだけである。極小部分分解体は、中心がガロアでなくても分解体であり得る。

部分分解体は \(Z\) に同型な体を含まなければならない。部分分解体の埋め込みに関する主張は、前の事実と §7, 2 から従う。さらに分解体は \(Z\) の分解体 \(\Omega\) を含まなければならない。しかし \(\Omega\) 上の任意の部分分解体は同時に分解体である。なぜなら成分 \(e^S A_\Omega\) は単純であり、\(\Omega\) に同型な中心を持つからである。\(Z\) がガロアなら、極小部分分解体と分解体は一致する。しかし非ガロアの場合にも、可換の場合とは違って、同時に分解体である極小部分分解体が存在し得る。すなわち、そのような極小部分分解体が \(Z\) に属するガロア体を含む場合である。例：中心が \(P(\sqrt2)\) に同型な四元数除法代数。ただし \(P\) は有理数体とする。このとき \(P(\sqrt2)\) に属するガロア体は極小部分分解体かつ分解体である。

最後に、中心が非分離的な場合には、\(Z_\Omega\) も、従って \(A_\Omega\) も根基を持つ系となる。\(\mathfrak C\) を \(Z_\Omega\) の根基、\(\mathfrak C A_\Omega\) を \(A_\Omega\) における拡大イデアルとする。すると \(Z_\Omega/\mathfrak C\) は、§8, 2 と正確に同じ類比で、共役な成分への直和分解を持つ。これにより上と同じく \(A_\Omega/\mathfrak C A_\Omega\) の直和分解が得られ、その成分 \(e^SA_\Omega\) は \(A\) に同型で、中心は \(e^SZ^\Omega\) である。従って部分分解体と分解体に関する定理はなお成り立つ。さらに階数の数え上げから、\(\mathfrak C A_\Omega\) の組成列に対応する組成因子は、単なる準同型ではなく、\(A_\Omega/\mathfrak C A_\Omega\) と作用素同型であることがわかる。

既約表現について述べると、次のまとめになる。\(P\) 上既約な超複素系の表現が与えられているとき、その表現が絶対的に完全可約であることは、対応する中心が \(P\) 上分離的であることと同値である。いずれの場合にも、その表現の分解体と部分分解体は、対応する単純系への埋め込みによって上のように特徴づけられる。とくに、部分分解体の最低次数は、中心の次数と中心上の対応する代数類の指数との積であり、任意の部分分解体の次数はその倍数である。表現は、組成列の意味で、中心に含まれる最大分離体の次数だけの、相異なるが共役な絶対既約表現類へ分解する。各類は、指数と、この分離拡大上の中心の次数との積で与えられる回数だけ現れる。

完全基礎体の場合、すなわち分離拡大体だけが現れる場合には、これは I. Schur の既知の結果を返す。ただし、R. Brauer にすでに現れる埋め込みによる特徴づけが付け加わっている。

\begin{center}
(1932年6月8日受理。)
\end{center}
%
\providecommand{\Sp}{\operatorname{Sp}}
\providecommand{\Tr}{\operatorname{Tr}}
\providecommand{\Norm}{\operatorname{N}}
\providecommand{\Mat}{\operatorname{Mat}}
\providecommand{\Gal}{\operatorname{Gal}}
\providecommand{\Cl}{\operatorname{Cl}}
\providecommand{\Gg}{\mathfrak G}
\providecommand{\Hh}{\mathfrak H}
\providecommand{\Ii}{\mathfrak I}
\providecommand{\Jj}{\mathfrak J}
\providecommand{\Aa}{\mathfrak A}
\providecommand{\Bb}{\mathfrak B}
\providecommand{\Cc}{\mathfrak C}
\providecommand{\Oo}{\mathfrak O}
\providecommand{\oo}{\mathfrak o}
\providecommand{\pp}{\mathfrak p}
\providecommand{\PP}{\mathfrak P}
\providecommand{\LL}{\mathfrak L}
\providecommand{\RR}{\mathfrak R}
\providecommand{\ZZ}{\mathfrak Z}
\providecommand{\ee}{\mathfrak e}

\clearpage

\section*{41. 相対ガロア数体に対する主種定理}

\begin{center}
\emph{Math. Ann. 108 (1933), pp. 411--419}
\end{center}

近年になって、非可換的方法、とりわけ代数の理論が、相対巡回数体および相対 Abel 数体について知られていた諸定理を、任意の相対ガロア体に対して定式化し証明する可能性を与えることが明らかになった。\footnote{これについては私のチューリヒ講演 \emph{Hyperkomplexe Systeme in ihren Beziehungen zur kommutativen Algebra und Zahlentheorie}, Verhandlungen des Internationalen Mathematiker-Kongresses Zürich 1932, Bd. I を参照。分裂代数に関する定理については、さらに H. Hasse, \emph{Die Struktur der R. Brauerschen Algebrenklassengruppe}, Math. Ann. 107 (1932) を参照。} とりわけノルム定理を思い出しておく。この定理は、一般には分裂代数に関する定理として述べられる。以下で私は、これに対応して主種定理もまた、内部自己同型、あるいは交叉表現についての定理として、超複素的に定式化できることを示す。証明は、先に述べた分裂代数に関する定理から、純粋に代数的かつ算術的な考察によって得られる。一方で、前者の定理の中には、巡回の場合、しかもすでに素数次の場合に、ノルム定理が超越的な核として含まれている。

定式化を理解しやすくするため、以下では用いない最小の場合の主種定理を先に述べておく。\footnote{ここで「最小の場合」と言うのは、任意の体における主種定理の代数的類似物を問題にしているからであり、「小さい場合」ではすでに \(p\)-進基礎体への移行の意味が入ってくる。} これは初等的な代数定理であり、\footnote{この定理は A. Speiser, \emph{Zahlentheoretische Sätze aus der Gruppentheorie}, Math. Z. 5 (1919), pp. 1--6 に見られる。そこですでに交叉表現に関する定理として述べられている。} 巡回の特殊な場合には Hilbert の \emph{Zahlbericht} のよく知られた定理 90、すなわち \(\Norm(a)\) が一に等しくなるのはそのとき、かつそのときに限り \(a=b^{1-S}\) となる、という定理に移る。そしてこの定理もすでに、内部自己同型または交叉表現に関する定理として述べ、代数の内部自己同型および接合積表示に関する一般定理によって証明する。

本来の主種定理では、後者の構成の代わりに、イデアル類群とガロア群との接合積が現れる。ただし、因子系に対してより細かい誘導されたイデアル類別を伴う。これによって、類体論における射線類別の類似物が捉えられる。ここでは自己同型および表現についての一般定理がまだ欠けているので、主種定理はこの方向への第一歩と見なせる。したがって、先に述べたように、分裂代数に関する定理による還元を与える。このことにより、定理はイデアル類ではなくイデアルに対する対応する定理へ還元される。ここでは、代数の極大オーダー相互の関係に関する Brandt の定理\footnote{極大オーダーに関する Brandt の理論は、H. Hasse, \emph{Über \(p\)-adische Schiefkörper und ihre Bedeutung für die Arithmetik hyperkomplexer Zahlsysteme}, Math. Ann. 104 (1931) において述べられ、また新たに基礎づけられている。}を、自己同型定理の同等物と見ることができる。接合積の極大オーダーに関する考察\footnote{C. Chevalley、H. Hasse、E. Noether の各一篇のノートが予告された。後二者は Herbrand 記念論文集に収められる予定であった。}とあわせて、それらは、個々の素点へ戻るためより複雑ではあるが、最小の場合の証明と本質的に平行な証明を与える。そこで私は、Speiser の証明のここではさらに簡単な類似物である、E. Artin によって示されたほとんど自明な証明を採る。

\subsection*{§1. 最小の場合の主種定理}

この節では、\(k\) は任意の可換体、必ずしも数体でなくてよいものを表し、\(K/k\) は \(n\) 次の分離ガロア拡大、\(\Gg\) は対応するガロア群を表す。

まず接合積についての既知の事実をまとめる。\footnote{この理論は、私の講義 (1929/30) に続いて、H. Hasse, \emph{Theory of cyclic algebras}, Trans. Amer. Math. Soc. 34 (1932), 第2章に述べられている。さらに超複素数についての M. Deuring も参照。} 接合積とは、\(K\) と \(\Gg\) とを同時に一つの代数 \(A\) に埋め込み、\(K\) の自己同型が内部的になるようにすることである。\(u_{S_1},\ldots,u_{S_n}\) を \(n\) 個の群要素に対応する記号とすると、まず \(A\) を \(K\) 上階数 \(n\) の線形形式の加群として
\[
        A=u_{S_1}K+\cdots+u_{S_n}K .                         \tag{1}
\]
とおく。\(u_S\)、より一般には \(u_S K^*\) によって生成される内部自己同型を要求することにより、\(A\) は環、したがって \(k\) 上階数 \(n^2\) の代数となる。この要求は、各 \(z\in K\) に対して
\[
        u_S^{-1}zu_S=z^S,\qquad\text{または}\qquad
        z u_S=u_S z^S,                                      \tag{2}
\]
および
\[
        u_Su_T=u_{ST}a_{S,T},\qquad a_{S,T}\in K^*,           \tag{3}
\]
さらに結合法則
\[
        a_{S,T}^{\,R}a_{ST,R}=a_{S,TR}a_{T,R}                 \tag{4}
\]
によって表される。任意の二つの元の積を
\[
 \Bigl(\sum_S u_S b_S\Bigr)\Bigl(\sum_T u_T c_T\Bigr)
   =\sum_{S,T}u_{ST}\,a_{S,T}\,b_S^{\,T}c_T
\]
で定めると、因子系 \(a_{S,T}\) に関する \(K\) と \(\Gg\) との接合積が得られる。これを
\[
        A=(a_{S,T},K,\Gg)
\]
と書く。\(A\) は \(k\) 上の単純正規代数、したがって対応する斜体 \(D\) 上の \(r\) 次の行列環 \(D_r\) であり、また \(K\) は極大可換部分体、すなわち分解体であることが示される。逆に、そのような正規斜体 \(D\) ごとに、このように接合積として生成できる行列環 \(D_r\) が存在する。

\(u_S\) から \(v_S=u_Sc_S\)（\(c_S\in K^*\)）へ移ると、随伴する因子系
\[
        \bar a_{S,T}
        =a_{S,T}\frac{c_S^{\,T}c_T}{c_{ST}}                  \tag{5}
\]
が生じる。とくに (5) によって、因子系
\[
        \frac{c_S^{\,T}c_T}{c_{ST}}
\]
は一に随伴する。これらの量を変換量という。随伴する因子系を一つの類 \((a)\) にまとめる。類 \((a)\) は Abel 群の直積形成に一対一に対応し、直積形成に関して Abel 群をなし、因子系の類の成分ごとの積と同型である。単位元は分裂代数類、すなわちすべての変換量からなる系である。これは R. Brauer の代数類群である。

巡回代数については次の特殊な場合が得られる。\(Z/k\) が巡回的で \(S\) が生成置換であれば、\(S\) の冪に \(u\) の冪が対応し、
\[
        A=Z+uZ+\cdots+u^{n-1}Z,                              \tag{1'}
\]
\[
        zu=uz^S,\qquad u^n=\alpha,                            \tag{2'--3'}
\]
で、
\[
        \alpha\in k^*,\qquad
        \bar\alpha=\alpha\,\Norm(c)\quad\text{ただし }v=uc     \tag{4'--5'}
\]
となる。したがって各因子系はここでは基礎体の一つの元だけからなり、\(A=(\alpha,Z,S)\) と記す。単位類は \(Z^*\) からのノルムによって与えられ、代数類の群は \(k^*/\Norm(Z^*)\) と同型である。

最小の場合の主種定理の定式化は、関係 (2) から (5) が純粋に乗法的であり、したがって同時に、\(n\) 個の複合体 \(u_SK^*\) の元からなる \(\Gg\) の拡大群 \(\Gg^*\) を定義する、という事実を用いる。
\[
        \Gg^*=\{\,u_{S_1}K^*,\ldots,u_{S_n}K^*\,\}.            \tag{6}
\]
\(\Gg^*\) は \(K^*\) を Abel 正規部分群としてもち、\(\Gg^*/K^*\simeq\Gg\) である。\(\Gg^*\) は \(K\) を自分自身に変換する正則元 \(g^*\) の全体、すなわち \(g^{*-1}Kg^*=K\) を満たす元の全体である。したがって \(K\) の環自己同型は、すべて、かつそれらだけが、\(\Gg^*\) の元によって生成される。実際、\(v=uw\) が \(K\) の内部自己同型を生成するならば、\(w\) は \(K\) のすべての元と可換であり、\(K\) は極大可換部分環であるから、\(w\in K^*\) である。

\paragraph{最小の場合の主種定理。}
第一の形：\(K^*\) の恒等自己同型の延長である \(\Gg^*\) の任意の群自己同型は内部的であり、\(K^*\) の一元によって生成される。言い換えれば、変換量
\[
        \frac{c_S^{\,T}c_T}{c_{ST}}
\]
がすべて値一をもつならば、\(c_S\) は記号的な \((1-S)\) 乗、すなわち
\[
        c_S=b^{1-S}=\frac{b}{b^S}\qquad(S\in\Gg)
\]
である。
第三の形：群 \(\Gg\) は、因子系一に属する \(K^*\) における一次の交叉表現類をただ一つだけもつ。

写像 \(u_S\mapsto C_S\) は、
\[
        C_S^{\,T}C_T=C_{ST}a_{S,T}
\]
を満たすとき、因子系 \(a_{S,T}\) に関して交叉している表現と呼ばれる。二つの表現は、
\[
        C_S=B^{-S}D_SB
\]
であるとき同じ類に属する。

私は第一の形を証明し、その後で第二および第三の形との一致を示す。証明は、指定された種類の任意の群自己同型が \(A\) の環自己同型に延長でき、この自己同型がよく知られているように内部自己同型である、という事実に基づく。この自己同型における \(u_S\) の像を \(v_S\) とする。このとき \(\sum v_SK\) は、同じ因子系に関する \(\Gg\) と \(K\) との接合積を再び成す。なぜなら、仮定によりすべての関係 (2) から (4) が保たれるからである。対応
\[
        \sum u_Sb_S\longmapsto \sum v_Sb_S
\]
は、\(K\) の元を固定する環自己同型を与える。ゆえにそれは \(K^*\) の一元によって生成される。

第二の形への移行は、(2) により \(v_S\) が \(v_S=u_Sc_S\) の形でなければならないことに基づく。因子系も保存されるので、(5) により、対応する変換量は一に等しくならなければならない。逆に (2) から (5) は、置換 \(v_S=u_Sc_S\) が指定された種類の自己同型を生成することを示す。したがって第一の形から
\[
        v_S=b^{-1}u_Sb=u_S\,\frac{b}{b^S},
        \qquad c_S=b^{1-S}
\]
を得る。巡回の場合には、仮定はとくに \(\Norm(c)=1\) を意味し、\(c=b^{1-S}\) がよく知られた定理である。第三の形は第二の形と直ちに同値である。仮定は \(u_S\mapsto c_S\) が因子系一に関する一次の交叉表現をなすことをいい、\(c_S=b^{1-S}\) は単位表現との同値性を意味するからである。\par ただし、この第三の形は、次のはるかに一般的な定理の一つの特殊な場合にすぎないことに注意されたい。所与の因子系 \(a_{S,T}\) に対して、既約交叉表現類はただ一つしか存在し、その次数は \(m\) である。ここで \(m\) は対応する代数の Schur 指数、すなわち対応する斜体の次数である。

\subsection*{§2. 主種定理}

\paragraph{類別についての予備的注意。}
巡回相対体に関する類体論の主種定理では、よく知られているように、二つの異なる類別が現れる。すなわち、上体における絶対イデアル類と、下体における射線類である。一般の場合には、上体のイデアルの類別が、因子系に対してより細かいイデアル類別を誘導する、ということがこれに対応する。

まず \(\Ii\) が \(K\) のすべてのイデアルからなる群を表すものとする。\(\Gg\) は \(\Ii\) に自己同型を誘導するので、\(\Gg\) による拡大は関係 (2) から (5) によって定義され、\(n\) 個の複合体 \(\{\,\cdots u_S\Ii\cdots\,\}\) から成る。\(\Ii\) において同値関係を課して絶対イデアル類の群へ移るときにも、\(n\) 個の複合体 \(u_S\bar\Ii\) は定義されている。なぜなら、\(\Gg\) は類群にも自己同型を誘導し、主類は自分自身に変換されるからである。したがって、\(\Ii\) から \(\bar\Ii\) へ導く同値関係は、\(\{\,\cdots u_S\Ii\cdots\,\}\) から \(\{\,\cdots u_S\bar\Ii\cdots\,\}\) へ移される。とくに、\(u_S\) は、\((c_S)\) が主イデアル全体を走るとき、すべての積 \(u_S(c_S)\) と同値である。

いま、積関係 (3) がこの同値性の意味で一意的であることを要求すると、これは、主イデアルから来るすべての変換量が一意的に同値になることを意味する。言い換えれば次のとおりである。

\(\bar\Ii\) が \(K\) の絶対イデアル類群を表し、\(H\) がとくに主類を表すとき、\(n\) 個の複合体 \(\{\,\cdots u_S\bar\Ii\cdots\,\}\) の系において、\(u_SH\) に対する積関係 (3) が一意的に定義されるのは、因子系の誘導されたイデアル類別において、主類が主イデアルから来るすべての変換量を含むとき、かつそのときに限る。

\paragraph{定義。}
因子系の誘導されたイデアル類別は、主類が次の性質をもつすべての主イデアル \((a_{S,T})\) から成ることによって定められる。すなわち、少なくとも一つの基礎元系 \(a_{S,T}\) に対して、代数
\[
        (a_{S,T},K,\Gg)
\]
が \(K\) のすべての分岐素点で分裂することである。

これらの主イデアルは、因子系の結合に関して群をなす。この群は変換量
\[
        \frac{(c_S)^{\,T}(c_T)}{(c_{ST})}
\]
を含む。なぜなら、因子系 \(c_S^{\,T}c_T/c_{ST}\) はすべての素点で分裂する代数を生成するからである。

\paragraph{主種定理の定式化。}
この定理は、最小の場合の定理に対応して、三つの同値な形で述べられる。

第一の形：因子系の誘導されたイデアル類別を基礎にして、置換
\[
        v_S=u_S c_S
\]
によって \(n\) 個の複合体 \(\{\,\cdots u_S\bar\Ii\cdots\,\}\) の自己同型が生じる、すなわちこの類別の意味で \(v_S\) に対して同じ関係 (3) が成り立つならば、この自己同型は内部的であり、イデアル類 \(b\) によって生成される。

第二の形：イデアル類 \(c_S\) から作られる変換量
\[
        \frac{c_S^{\,T}c_T}{c_{ST}}
\]
がすべて、因子系の誘導された類別の主類に属するならば、イデアル類からなるベクトル \(\{\,\cdots c_S\cdots\,\}\) は主種を成し、類 \(c_S\) は記号的な \((1-S)\) 乗である。すなわち、あるイデアル類 \(b\) が存在して
\[
        c_S=b^{1-S}=\frac{b}{b^S}\qquad(S\in\Gg)
\]
となる。

第三の形：因子系に対する誘導されたイデアル類別を基礎にするとき、群 \(\Gg\) は、因子系の単位類に属するイデアル類群内の一次の交叉表現類をただ一つだけもつ。

これらの形が同値であることは §1 と同様に従う。ここでは自己同型がすでに \(v_S=u_Sc_S\) によって生成されると仮定されているので、第一の形から第二の形への移行はさらに簡単である。第三の形については、イデアル類群における表現は乗法的にのみ定義されることを注意すればよい。一次表現に関する主張はこれによって制限されない。

\paragraph{主種定理の証明。}
第二の形の証明は二つの補題に基づく。

\paragraph{補題 1.}
イデアル \(c_S\) から作られる変換量
\[
        \frac{c_S^{\,T}c_T}{c_{ST}}
\]
が単位イデアルを与えるならば、\(c_S\) は記号的な \((1-S)\) 乗である。
\[
        c_S=b^{1-S}\qquad(S\in\Gg).
\]
これは、そこで自己同型が \(v_S=u_Sc_S\) によって生成されると仮定するならば、第一の形と同値である。証明：\(b=\sum c_S\) が要求を満たす。ここで和は加群和、すなわち最大公約数を意味する。実際、
\[
        b=\sum_S c_S,\qquad
        b^T=\sum_S c_S^{\,T},\qquad
        b^T c_T=\sum_S c_S^{\,T}c_T=\sum_S c_{ST}=b,
\]
したがって \(c_T=b^{1-T}\) である。

\paragraph{補題 2.}
代数
\[
        A=(a_{S,T},K,\Gg)
\]
が \(K\) のすべての有限および無限の分岐素点で分裂し、さらに主イデアル \((a_{S,T})\) がイデアルの変換量、すなわち
\[
        (a_{S,T})=\frac{c_S^{\,T}c_T}{c_{ST}}
\]
であるとする。このとき \(A\) はすべての素点で、したがって絶対的に分裂する。すると \((a_{S,T})\) は主イデアルの変換量、
\[
        (a_{S,T})=\frac{(d_S)^{\,T}(d_T)}{(d_{ST})}
\]
となる。逆に、すべての素点で分裂する任意の代数はこの条件を満たす。

証明は \(p\)-進拡大のもとでの \(A\) の既知の振る舞いに基づく。\(k_p\) を \(k\) の \(p\)-進拡大、\(K_\PP\) を \(K\) の \(\PP\)-進拡大とする。ただし \(\PP\) は \(p\) の上にある素イデアルである。\(A_p\) を、\(k\) から \(k_p\) への移行によって得られる \(A\) の拡大とする。さらに \(\mathfrak Z\) を \(\PP\) に対する分解群、\(a_{\mathfrak Z}\) を因子系の対応する部分とする。このとき
\[
        A_p\sim (a_{\mathfrak Z},K_\PP/k_p,\mathfrak Z)
\]
が成り立つ。とくに \(p\) が非分岐、したがって \(\mathfrak Z\) が巡回であるならば、これは非分岐慣性体 \(K_\PP/k_p\) による \(A_p\) の標準的な表示である。

示すべきことは、補題 2 の仮定のもとで、\(A_{\mathfrak p}\) がこの非分岐の場合にも分裂することである。無限素点 \(\mathfrak p\) では \(K_{\mathfrak P}=k_{\mathfrak p}\) なので、\(A_{\mathfrak p}\) は分裂する。有限の \(\mathfrak p\) については、まず因子系 \(a_{\mathfrak Z}\) が \(K_{\mathfrak P}\) の単数から成る因子系に随伴することが従う。実際、\(K_{\mathfrak P}\) ではすべてのイデアル \(\mathfrak c_S\) が主イデアル \((c_S)\) となるので、\((a_{S,T})\) に関する仮定は
\[
a_{S,T}=e_{S,T}\frac{c_S^Tc_T}{c_{ST}}
\]
と書けることを意味する。ここで \(e_{S,T}\) は単数であり、従って \(a_{\mathfrak Z}\) は単数から成る因子系 \(e_{S,T}\) に随伴する。

さらに、正規化された巡回表示 (1') から (5') へ移ってもこの性質は保たれる。すなわち、\(R\) を \(\mathfrak Z\) の生成元とすると、\(A_{\mathfrak p}=(\alpha,K_{\mathfrak P}/k_{\mathfrak p},R)\) において \(\alpha\) は単数になる。\(e_{S,T}\) が因子系であることは、関係
\[
u_{R^i}u_{R^j}=u_{R^{i+j}}e_{R^i,R^j}
\]
によって表される。\(u_R=u\) とおけば
\[
u^i=u_{R^i}e_{R,R^{i-1}}e_{R,R^{i-2}}\ldots e_{R,R}
\]
を得る。従って、\(u_E=e_{E,E}\) であり、\(f\) が \(\mathfrak Z\) の位数なら、とくに
\[
u^f=\alpha=e_{E,E}e_{R,R^{f-1}}\ldots e_{R,R}
\]
である。非分岐拡大 \(K_{\mathfrak P}/k_{\mathfrak p}\) ではすべての単数がノルムなので、\(A_{\mathfrak p}\) は分裂する。従って \(A\) はすべての素点で分裂し、分裂代数に関する定理により絶対的に分裂する。

言い換えれば、仮定から \(a_{S,T}\) が元から来る変換量、すなわち
\[
        a_{S,T}=\frac{d_S^{\,T}d_T}{d_{ST}},
        \qquad d_S\in K^*
\]
であることが従う。弱めて言えば、\((a_{S,T})\) は主イデアルの変換量になる。逆に \(A\) がすべての素点で分裂するなら、主イデアル \((a_{S,T})\) も各素点においてイデアルの変換量であり、\(A\) が分裂することは、分裂代数に関する定理の仮定の別の形そのものになる。

二つの補題から主種定理の証明は直ちに従う。\(c_S\) を類 \(c_S\) に属する一つのイデアルとする。仮定は、誘導された類別の定義により、\(c_S\) の変換量が補題 2 の仮定を満たす主イデアル \((a_{S,T})\) になることをいう。したがって、
\[
        \frac{c_S^{\,T}c_T}{c_{ST}}
\]
が \((d_S)\) によって単位になるような主イデアル \((d_S)\) が存在する。同じ類 \(c_S\) に属するイデアル \(c_S/(d_S)\) は、したがって補題 1 の仮定を満たす。ゆえに、あるイデアル類 \(b\) が存在して
\[
        c_S=b^{1-S}
\]
となる。これは、類 \(c_S\) が類 \(b\) の記号的な \((1-S)\) 乗であることを意味する。

巡回の場合には、正規化を基礎にしたとき、誘導された類別の主類が、\(\alpha\) がすべての分岐素点でノルム剰余となるような主イデアル \((\alpha)\) の群へ移るならば、主種定理はよく知られた定理に移る。

\begin{center}
(1932年10月27日受理。)
\end{center}

\section*{42. 分裂する接合積とその極大オーダー}

\begin{center}
\emph{Actualités scientifiques et industrielles 148 (1934), pp. 5--15}
\end{center}

分裂する接合積のもっとも簡単な場合、\footnote{接合積の理論は、私の講義に従って、H. Hasse, \emph{Theory of cyclic algebras over an algebraic number field}, Trans. Amer. Math. Soc. 34 (1932) の第 II 節に再録されている。ここでは最初の基礎概念だけを用いる。Brandt の極大オーダー論は、H. Hasse, \emph{Über $p$-adische Schiefkörper und ihre Bedeutung für die Arithmetik hyperkomplexer Zahlsysteme}, Math. Ann. 104 (1931)（以下 H と略記）に、さらに広い結果の枠内で詳しく叙述され、新たに基礎づけられている。} すなわち分裂代数を与え、したがって一に随伴する因子系をもつものの場合には、事情をかなり明示的に追跡することができる。これは、非ガロアの分解体、すなわち極大可換部分体の場合にもなお成り立つ。ここに属する単純な代数的事実を §1 で展開する。§2 では、基礎体を代数的数体とする場合に、基礎にある分解体 \(k\) の加群および相補加群によって、すべての極大オーダーとそのイデアルの明示的表示を与える。\footnote{これらの表示は私には長く知られていた。続く考察への契機は H. Hasse の一つの定理の知らせであった。}

さらに領域への分類を与える。一つの領域には、\(k\) との同じ交わりをもつすべての極大オーダーを入れる。これにより領域は \(k\) のオーダーと一対一に対応し、主オーダーには主領域が対応する。一つの領域の極大オーダーは、対応するオーダーの加群から形成されるイデアルによって相互に変換される。とくに主領域では、これは \(k\) のイデアルの拡大である。§3 では任意の単純代数へ移る。領域への分類を精密化し、交わりだけでなく分岐素点における極大オーダーの \(p\)-進成分も一致しなければならないとすることにより、分裂する場合に得られた諸定理を移すことができる。とくにその場合にも、主領域の極大オーダーは \(k\) のイデアルの拡大によって相互に変換される。

主領域に関する定理は Chevalley と Hasse に見られる。\footnote{C. Chevalley, \emph{Sur certains idéaux d'une algèbre simple}, Abh. Math. Sem. Hamburg, 1934; H. Hasse, \emph{Über gewisse Ideale in einer einfachen Algebra}, 本論文に先行する論文。} Chevalley はまた、分岐素点に関するこの精密化された分類を行わなければ、一般には定理が成り立たなくなることを示している。

\subsection*{§I. 分裂接合積の行列単位と分解体の相補基底}

この節では \(\Omega\) を任意の基礎体とする。\(k/\Omega\) は \(n\) 次のガロア分離拡大、\(\Gg\) は元 \(S,T,\ldots\) をもつガロア群とし、\(u_S,u_T,\ldots\) は対応する作用素とする。因子系一をもつ接合積、したがって分裂代数
\[
        K=\Gg\ast k=u_{S_1}k+\cdots+u_{S_n}k,\qquad
        u_Su_T=u_{ST},
\]
を考える。ただし
\[
        z u_S=u_S z^S \qquad(z\in k)
\]
である。ここで
\[
        E_\Gg=\sum_{S\in\Gg}u_S
\]
とおくと、因子系一から関係
\[
        E_\Gg u_S=u_S E_\Gg=E_\Gg,                            \tag{1}
\]
ここで \(S\) は \(\Gg\) のすべての元を走る。すなわち \(E_\Gg\) によって \(\Gg\) の単位表現が生成される。\footnote{\(k\) の標数が \(n\) を割らない場合には、\(\frac1nE_\Gg\) が単位表現の冪等元となる。しかし標数が \(n\) を割る場合にも \(K\) の標準形は成り立ち、唯一の仮定は \(k/\Omega\) が分離的であることである。}
\[
        zE_\Gg=\sum_Su_Sz^S,                                  \tag{2a}
\]
および
\[
        E_\Gg z E_\Gg
        =E_\Gg\sum_Su_Sz^S
        =E_\Gg\,\Sp(z),                                      \tag{2b}
\]
が得られる。ここで \(\Sp(z)\) は \(k/\Omega\) の元としての \(z\) のトレースを意味する。

\paragraph{定理。}
\(a_1,\ldots,a_n\) が \(k/\Omega\) の基底で、\(\bar a_1,\ldots,\bar a_n\) がそれに相補的な基底を表すならば、
\[
        c_{ik}=\bar a_iE_\Gg a_k
\]
によって \(K\) の行列単位系が与えられる。\(K\) は加群積として
\[
        K=kE_\Gg k
          =k\Bigl(\sum_Su_S\Bigr)k
          =\sum_{i,k}(\bar a_iE_\Gg a_k)\,\Omega
\]
の形に表示できる。実際、(2b) により
\[
        c_{ij}c_{lk}
        =\bar a_iE_\Gg a_j\bar a_lE_\Gg a_k
        =\bar a_iE_\Gg a_k\,\Sp(a_j\bar a_l)
        =c_{ik}\,\Sp(a_j\bar a_l)
\]
となる。これは、相補基底の定義式
\[
        \Sp(a_j\bar a_l)=\delta_{jl}
\]
によって、まさに行列単位の性質である。

\paragraph{注意 1.}
\(Z\) が \(\Omega\) の任意の可換拡大体であるなら、\(a_i\) が \(k_Z/Z\) の基底を表し、\(K_Z\) が体の直和になる場合にも、すべては保たれる。そのとき群の置換は、\(Z\) 上で恒等を誘導することによって定義される。

\paragraph{注意 2.}
\(K=kE_\Gg k\) であることは、右辺が単純系 \(K\) の零でないイデアルであることからも直ちに従う。

いま \(\Omega\) の標数が零であるとする。分裂代数の表示 \(K=kE_\Gg k\) は、\(k\) が \(K/\Omega\) の非ガロア分解体である場合にも有効なままである。これはガロア分解体へ移ることによって示される。\(\bar k\) を \(k\) に属するガロア体、\(\bar\Gg\) をその群、\(\bar u_S\) を対応する作用素、\(\bar K=\bar kE_{\bar\Gg}\bar k\) を分裂接合積とする。部分体 \(k\) は位数 \(h\) の部分群 \(\Hh\) に属するとする。次のようにおく。
\[
        \bar\Gg=\Hh S_1+\cdots+\Hh S_n
              =S_1^{-1}\Hh+\cdots+S_n^{-1}\Hh,
        \qquad
        E_\Hh=\frac1h\sum_{H\in\Hh}\bar u_H,
\]
および
\[
        E_\Gg=\frac1h E_{\bar\Gg}
              =\sum_i u_{S_i},\qquad
        u_{S_i}=E_\Hh\bar u_{S_i}.
\]
すると、とくに冪等元 \(E_\Hh\) は \(\Hh\) の単位表現を生成し、\(E_\Gg\) も \(E_{\bar\Gg}\) と同様に \(\bar\Gg\) の単位表現を生成する。このように定義された \(E_\Gg\) と \(u_{S_i}\) に対して、関係 (1) は一方の側について、関係 (2) は完全に成り立つ。ここで \(z\) は \(k\) のすべての元を走り、従って \(z^S\) は共役体 \(k^S\) のすべての元を走る。\(k^S\) は \(\bar k\) の部分体として \(\bar K\) の部分体でもある。まず
\[
        E_\Gg=E_\Gg E_\Hh=E_\Hh E_\Gg=E_\Hh E_\Gg E_\Hh,       \tag{3}
\]
\[
        zE_\Hh=E_\Hh z,                                      \tag{4}
\]
従って
\[
        z u_{S_i}=u_{S_i}z^{S_i}                               \tag{5}
\]
である。(3) は、\(E_{\bar\Gg}\) については \(E_\Hh\) の定義と (1) から成り立ち、従って \(E_\Gg\) についても成り立つ。一方 (4) は、\(z\) が \(k\) の元であるため \(\Hh\) の置換を許すこと、すなわち \(z\bar u_H=\bar u_Hz\) をいう。(4) の右から \(\bar u_{S_i}\) を掛ければ (5) を得る。さらに (1) は一方の側について
\[
        E_\Gg u_{S_i}=E_\Gg,
        \qquad
        E_\Gg u_{S_i}=E_\Gg E_\Hh\bar u_{S_i}
        =E_\Gg\bar u_{S_i}=E_\Gg                              \tag{1'}
\]
として成り立つ。(2a) は (5) を総和して得られ、そこから (1') により (2b) が従う。

\(K=kE_\Gg k\) の証明は今や 1 とまったく同じである。実際、\(\bar a_iE_\Gg a_k\) が行列単位であることの証明に用いたのは (2b) だけであった。\(\Omega\) 上の全行列環として、\(K\) は任意の \(n\) 次体を部分体として含み、とくに \(k\) を含む。

§3 のために、\(K\) が分裂しない場合にも \(K\) から \(\bar K\) への移行が可能であることを付け加える。\footnote{ここから非ガロアの場合の接合積表示の類似物が得られる。この点は、おそらくなお学位論文で詳述されるであろう。} \(k\) のほかに \(\bar k\) も分解体であるから、\(K\) は \(\bar\Gg\) と \(\bar k\) との接合積 \(\bar K\) に相似である。\(k\) は分解体であるため、\(k\) に対応する部分群 \(\Hh\) には因子系一が対応し、従って冪等元 \(E_\Hh\) は上と同じように定義される。逆にここから、\(K\) は \(E_\Hh\bar K E_\Hh\) と同型であることが従う。この代数は左イデアル \(\bar K E_\Hh\) の自己同型環と同型であり、従って \(\bar K\) と相似である。その階数は \(K\) の階数と一致する。実際、\(\bar K E_\Hh\) は \(\bar k\) 上階数 \(n\) であり（\(\bar u_{S_i}^{-1}E_\Hh\) は \(\bar k\) 上一次独立である）、\(\bar K\) は階数 \(nh\) である。すなわち \(\bar K\) は \(\bar K E_\Hh\) と同型な \(h\) 個のイデアルに分解し、自己同型環に対しては次数 \(h\) の全行列代数を掛けることになる。

\(K\)、従って \(\bar K\) が分裂する場合には、この方法で 2 の結果を再び得る。実際、
\[
\begin{aligned}
        E_\Hh\bar K E_\Hh
        &=E_\Hh\bar kE_{\bar\Gg}\bar kE_\Hh\\
        &=(E_\Hh\bar kE_\Hh)E_{\bar\Gg}(E_\Hh\bar kE_\Hh)\\
        &=\Sp(\bar k/k)E_\Hh E_{\bar\Gg}E_\Hh\Sp(\bar k/k)\\
        &=kE_\Gg k.
\end{aligned}
\]
\subsection*{§II. 分裂接合積の極大オーダーとその領域への分類}

以後、\(\Omega\) は有限代数的数体\footnote{本節の考察は、\(\Omega\) が一不定元の関数体である場合にも保たれる。}、\(K\) は \(\Omega\) 上次数 \(n\) の分裂代数、\(k\) はガロアまたは非ガロアの極大可換部分体で、\(K=kE_\Gg k\) とする。\(a,\bar a\) を \(k\) における有限な相補的 \(\oo\)-加群の対とする。ここで \(\oo\) は \(\Omega\) の主オーダー、\(o\) は \(k\) の主オーダーである。\(k\) に関する \(K\) の極大オーダーの構造は、次の定理によって記述される。

\paragraph{定理 1.}
加群積
\[
        \mathfrak O_a=\bar a E_\Gg a
\]
は \(K\) の極大オーダーである。とくに
\[
        \mathfrak O=\mathfrak O_o=\bar oE_\Gg o
\]
は極大オーダーである。

\paragraph{定理 2.}
\(\mathfrak O\) から \(\mathfrak O_a\) への変換は、\(\mathfrak O\) に関する左イデアル
\[
        \mathfrak L=\bar oE_\Gg a
\]
およびその逆右イデアル
\[
        \mathfrak R=\bar aE_\Gg o
\]
によって行われる。
\paragraph{注意。}
積 \(\mathfrak L\mathfrak R\) の中に二つの相補加群 \(a,\bar a\) が現れることは、\(E_\Gg\) によって引き起こされるトレース形成のために、ちょうど逆関係へ導く。

\paragraph{定理 3.}
\(\mathfrak O\) に関する左および右イデアルは、定理 2 に示されたものによって尽くされる。したがって極大オーダー \(\mathfrak O_a\) は \(K\) のすべての極大オーダーを尽くす。任意の二つの極大オーダー \(\mathfrak O_a\) と \(\mathfrak O_b\) は
\[
        \mathfrak L_{ab}=\bar aE_\Gg b,
        \qquad
        \mathfrak R_{ba}=\bar bE_\Gg a
\]
によって相互に変換される。これらの \(\mathfrak L_{ab}\) と \(\mathfrak R_{ba}\) は \(\mathfrak O_a\) に関する左および右イデアルを尽くす。

\paragraph{定理 4.}
\(\mathfrak O_a\) が \(K\) のすべての極大オーダーを走るとき、交わり
\[
        [\mathfrak O_a,k]
\]
は \(k\) のすべてのオーダーを走る。この交わりは各場合に \(a\) のオーダーに等しい。\(k\) の同じオーダーに属するすべての極大オーダーを一つの領域にまとめるならば、一つの領域内の変換は、\(k\) の対応するオーダーの加群 \(a\) と \(b\) によって、\(\mathfrak O_a\) に関する左イデアル \(\bar aE_\Gg b\) および逆右イデアルによって行われる。

\paragraph{定理 4a.}
とくに、主領域、すなわち \(k\) の主オーダー \(o\) を含むすべての極大オーダーからなる領域の極大オーダーの変換は、
\[
        \bar aE_\Gg b
\]
というイデアルによって行われる。ここで \(a\) と \(b\) は \(k\) のイデアルである。言い換えれば、\(\mathfrak L\) が主領域の極大オーダーの左イデアルであり、\(\mathfrak L\) の右オーダーに \(o\) が含まれているならば、\(\mathfrak L\) は \(o\)-イデアルの拡大である。

証明は個々の素点への移行に基づく。\(p\) が \(\Omega\) のすべての素イデアルを走るとして、各場合に \(p\) に関する商環へ移れば十分である。\(\mathcal C\) を \(K\) の任意の \(\oo\)-加群とするとき、その成分 \(\mathcal C_p\) として拡大加群 \(\mathcal C\oo_p\) を表すことにする。すると次が成り立つ。

\paragraph{補題。}
\(\mathcal C\) はすべての拡大加群 \(\mathcal C_p\) の交わりである。実際、\(w\) が \(\mathcal C_p\) に属するならば、\(p\) で割れない \(\sigma\in\oo\) があって \(\sigma w\in\mathcal C\) となる。\(w\) がすべての \(\mathcal C_p\) の交わりに属し、\(g\) を \(\sigma w\in\mathcal C\) を満たすすべての \(\sigma\in\oo\) からなるイデアルとすれば、いま示したように \(g\) はどの \(p\) でも割れない。従って \(g=\oo\) である。\footnote{素点を \(p\)-進拡大によって定義するならば、対応する補題は「すべての \(p\)-進成分 \(\mathcal C^{(p)}\) と \(K\) との交わりは \(\mathcal C\) に等しい」となる。これは \([\mathcal C^{(p)},K]=\mathcal C_p\) から従う。この等式は、\(\mathcal C_p\) の \(\oo_p\)-基底であって同時に \(\mathcal C^{(p)}\) の基底となるものを用いれば分かる。H, Satz 66 では、この補題は（最大階数の）イデアルについてのみ述べられている。} この補題は加群が最大階数をもつことを仮定しない。
\paragraph{系。}
任意の加群はその成分によって決定される。\(\mathcal D\) が \(\mathcal C\) の真の上加群であるならば、少なくとも一つの \(\mathcal D_p\) は \(\mathcal C_p\) の真の上加群である。とくに \(\mathcal M\) が \(K\) のオーダーであり、すべての成分 \(\mathcal M_p\) が \(\oo_p\) 上の極大オーダーであるならば、\(\mathcal M\) 自身も極大オーダーである。

\paragraph{定理 1 の証明。}
\(\mathfrak O_a\) はオーダーである。実際、\(a\) と \(\bar a\) が有限 \(\oo\)-加群なので \(\mathfrak O_a\) も有限 \(\oo\)-加群であり、しかも最大階数をもつ。さらに
\[
\begin{aligned}
        \mathfrak O_a^2
        &=\bar aE_\Gg a\,\bar aE_\Gg a\\
        &=\bar aE_\Gg a\,\Sp(\bar a a)
        \subseteq \bar aE_\Gg a=\mathfrak O_a.
\end{aligned}
\]
ここで \(\Sp(\mathfrak c)\) は一般に、\(\mathfrak c\) の全要素のトレースからなる \(\oo\) のイデアルを意味する。極大性を示すには、上の系により、各成分の極大性を示せば十分である。\(\oo_p\) は主イデアル環であるから、\(a_p\) と \(\bar a_p\) は互いに相補的な基底 \(a_1,\ldots,a_n\) と \(\bar a_1,\ldots,\bar a_n\) をもつ。§1 により、積 \(\bar a_iE_\Gg a_k\) は行列単位をなす。従って
\[
        \mathfrak O_{a,p}=\sum_{i,k}(\bar a_iE_\Gg a_k)\oo_p
\]
は極大である。\(\mathfrak O_a\) が極大であるから、\(\Sp(\bar a a)=\oo\) も従う。
\paragraph{定理 2 の証明。}
トレース規則から次を得る。
\[
\begin{aligned}
 \mathfrak O\mathfrak L&=\bar oE_\Gg o\,\bar oE_\Gg a=\bar oE_\Gg a=\mathfrak L,\\
 \mathfrak R\mathfrak O&=\bar aE_\Gg o\,\bar oE_\Gg o=\bar aE_\Gg o=\mathfrak R,\\
 \mathfrak L\mathfrak R&=\bar oE_\Gg a\,\bar aE_\Gg o=\bar oE_\Gg o=\mathfrak O,\\
 \mathfrak R\mathfrak L&=\bar aE_\Gg o\,\bar oE_\Gg a=\bar aE_\Gg a=\mathfrak O_a.
\end{aligned}
\]
\paragraph{定理 3 の証明。}
\(\mathfrak L\) を \(\mathfrak O\) に関する左イデアルとする。このとき、とくに
\[
        \mathfrak L=\bar oE_\Gg\mathfrak L
        =(\bar oE_\Gg l_1,\ldots,\bar oE_\Gg l_t)
\]
である。ただし \(l_1,\ldots,l_t\) は、たとえば \(\mathfrak L\) の \(\oo\)-基底である。\(l_\mu=\sum_i u_{S_i}z_{\mu,i}\)、\(z_{\mu,i}\in k\) とおくと、\footnote{\(k\) が非ガロアならば、\(z\) を \(\bar zE_\Hh\) で、従って \(a\) を \(\bar aE_\Hh\) で置き換えなければならない。このとき \(a\in k\) に対して
\[
E_\Gg l=E_\Gg\bar aE_\Hh=E_\Gg a
\]
を得る。実際、(3), (4), (5) により
\[
kE_\Gg k=kE_\Gg kE_\Hh=\sum_i u_{S_i}\{k^{S_i},k\}E_\Hh
\]
である。}
\[
        E_\Gg l_\mu=E_\Gg\sum_i z_{\mu,i}=E_\Gg a_\mu
\]
となり、従って \(\mathfrak L=\bar oE_\Gg a\) である。正則な \(\mathfrak L\) では \(a\) は最大階数 \(n\) をもたなければならず、そのとき \(\bar aE_\Gg o\) が逆右イデアルである。合成により、\(\bar aE_\Gg b\) は \(\mathfrak O_a\) に関する左イデアルを尽くし、\(\bar bE_\Gg a\) はそれらに逆な右イデアルを尽くす。
\paragraph{定理 4 の証明。}
交わり \(t=[k,\mathfrak O_a]\) はオーダーである。\(\mathfrak O_a\) の部分集合として、\(t\) は右乗数領域であり、この条件を満たす \(o\) の中の最大のオーダーである。したがって \(t\) は確かに \(a\) のオーダー \(o_a\) を含む。\(t=o_a\) を示すには、成分ごとの等式を示せば十分である。\(\tau\in t_p\) ならば、各 \(c_{ik}=\bar a_iE_\Gg a_k\) について \(\tau c_{ik}\) は \(\oo_p\) の係数をもつ \(c_{ij}\) の線形結合である。他方、\(a_k\tau\) は \(\oo_p\) の係数をもつ \(a_j\) の線形結合になる。\(c_{ij}\) による表示の一意性のため、これらの係数は \(\oo_p\) に属さなければならない。したがって \(\tau\) は \(a_p\) のオーダーに属する。これで \(t\) は \(a\) のオーダーに等しい。\(a\) がとくに \(k\) のすべてのオーダーを走ることから、各オーダーが交わりとして実際に現れることも従う。変換についての主張は、いまや定理 3 から直ちに従う。

\paragraph{定理 4a の証明。}
この定理は定理 4 の主領域への特殊化にすぎない。まず出発する極大オーダーを \(\mathfrak O=\bar oE_\Gg o\) とする。\(o\) が \(\mathfrak L=\bar oE_\Gg a\) の右オーダーにも含まれるならば、\(\bar aE_\Gg a\) も主領域に属する。従って \(a\) は \(k\) のイデアルであり、\(\mathfrak L=\mathfrak Oa\) となる。\(\mathfrak O_a\) から出発すれば、対応して \(\mathfrak L=\bar aE_\Gg b\) となり、\(b\)、従って \(a^{-1}b\) は \(k\) のイデアルである。ゆえに \(\mathfrak L=\bar aE_\Gg aa^{-1}b=\mathfrak O_a a^{-1}b\) となる。
\subsection*{§III. 任意の接合積の極大オーダー}

分裂接合積について導かれた定理は、領域への分類を精密化すれば、一般の場合にも成り立つ。主領域についての事実はそのとき正確に成り立ち、その他の事実は個々の素点についての主張として定式化しなければならない。ここで素点は \(p\)-進拡大として理解される。

\paragraph{定義。}
\(K\) を \(\Omega\) 上の単純正規代数とし、\(k\) を極大可換部分体とする。\(k\) に関する一つの領域には、\(K\) の極大オーダーのうち、\(k\) との交わりが同じで、さらに \(K\) の分岐素点において同じ \(p\)-進成分をもつものをすべて入れる。\(k\) との交わりが主オーダーであれば、それは主領域である。

主領域については次が成り立つ。

\paragraph{定理 1.}
一つの主領域の極大オーダーは、\(k\) のイデアルの拡大イデアルによる変換によって相互に移る。言い換えれば、主領域の極大オーダーの左イデアルで、共役差積に素であり、その右オーダーに \(o\) を含むものは、\(k\) のイデアルの拡大である。

\paragraph{定理 2.}
\(K\) が素数次であれば、主領域はただ一つである。この主領域は、\(k\)-イデアルの拡大イデアルによって自分自身へ移る。固定された極大オーダー \(\mathfrak O\) の対応するイデアルは、\(\mathfrak O\) の両側イデアルによって拡大された \(k\) のイデアルである。

\paragraph{証明。}
一つの主領域の二つの極大オーダーは有限個の素点でのみ異なる。定義により、これらは \(K\) が分裂する素点である。しかしそこでは \(K\) は §§1--2 で考えた形になる。これは \(k\) がガロア的であれば明らかである。因子系が一に随伴するならば、それは単位系で置き換えられる。§1 により、これは非ガロアの場合にも成り立つ。したがって、各そのような素点で、したがって全体として\footnote{H, §8、または先行する Hasse の論文（注4）を参照。}、拡大イデアルが問題となる。

\(K\) が素数次であれば、\(K\) は分裂するか斜体である。後者の場合、それはすべての分岐素点で斜体のままであり、したがってただ一つの主領域をもつので、拡大イデアルによる変換が再び成り立つ。もっとも一般的な変換は、拡大イデアルを極大オーダーを自分自身に変換するイデアル、すなわち両側イデアルと合成することによって得られる。これらは、よく知られているように、分岐素点においてのみ中心イデアル、したがって拡大イデアルと異なる。

さらに次が成り立つ。

\paragraph{定理 3.}
任意の領域の極大オーダーは、各素点で §2 の定理 3 に示したように対応するオーダーの加群から組み立てられる、共役差積に素なイデアルによって相互に移る。固定された一つのイデアルに対して、これらの加群の成分は有限個の素点でのみオーダーと異なる。

これは §2 から直ちに従う。ただし、個々の素点で作用素 \(u_S\) を対応する同値なものに置き換えなければならないことに注意すればよい。

\(k\) の任意のオーダーに対して、そのオーダーを交わりとしてもつ極大オーダーがここでも存在するかどうかは、分岐素点でより精密な調査を要する。主領域は常に存在する。\footnote{別証明は Hasse および Chevalley（注4、注3）にある。}これは、因子系を整と仮定すると \(o\) を含むオーダーへ導く、接合積表示、場合によっては一般化された接合積表示から読み取れる。

さらに、算術的な領域分類によって代数的分解体をどの程度まで特徴づけることが可能か、という問題がある。言い換えれば、異なる分解体は常に異なる領域分類を生じるのか、すでに異なる主領域を生じるのか。あるいは、すべての分解体の主領域が極大オーダー全体を尽くすのか。

単一の極大オーダーによってそのような特徴づけが確実には不可能であることは、もっとも単純な例が示している。例えば、四元数体の極大オーダー
\[
        \left[\,1,i,j,\frac{1+i+j+k}{2}\,\right]
\]
は、主オーダー \([1,i]\) のほかに、三次単位根の体の主オーダー
\[
        \left[\,1,\frac{1+i+j+k}{2}\,\right]
\]
も含んでいる。

\begin{center}
Göttingen, 1932年8月。
\end{center}
%

\clearpage

\section*{43. イデアル微分と共役差積}

\begin{center}
\emph{Journal f. d. reine u. angew. Math. 188 (1950), S. 1--21}\\[0.75ex]
\emph{エミー・ノイエター} \(\dagger\)\srcfnmark{1)}.
\end{center}

\srcfntext{1)}{本稿は E. ノイエターの遺稿で、1927/28年冬に執筆され、H. グレルの厚意により編者へ提供された。これは1929年プラハの D. M. V. 大会で行われた同題講演の詳細な叙述である。\emph{Jahresbericht D. M. V.} 39 (1930), p. 17（斜体頁）を参照。§6, 3 以後は明らかになお精密な改訂が予定されていた。しかし専門家なら、小幅な整文のほかは原形を保ったこの概略的叙述を難なく理解できるであろう。}

\subsection*{序論。}

代数的数体の分岐理論の主定理は、周知のように、共役差積（分岐イデアル）が、素数 $p$ の分解に $\varrho$ 乗で現れる素イデアルの少なくとも $(\varrho-1)$ 乗で割り切れることをいう。より正確には、$\varrho$ が $p$ で割り切れないときにはちょうど $(\varrho-1)$ 乗で割り切れ、$\varrho$ が $p$ で割り切れるときにはそれより高い冪で割り切れる。

この定理は、多項式 $f(x)$ の微分商 $f'(x)$ が、$f(x)$ に $\varrho$ 乗で現れる一次因子の少なくとも $(\varrho-1)$ 乗で割り切れるという事実に対応する。すなわち、$\varrho$ が係数領域の標数で割り切れないときにはちょうど $(\varrho-1)$ 乗で割り切れ、$\varrho$ がその標数で割り切れるときにはそれより高い冪で割り切れる。

以下で、この対応が単なる形式的な類似にとどまらないことを示す。\emph{共役差積は、数体 $\mathsf{K}$ の定義イデアルを、$x_1,\ldots,x_n$ に関する対応する整係数多項式環の中で微分した微分商とみなすことができる。この微分商は点 $x=\omega$、すなわち $x_1=\omega_1,\ldots,x_n=\omega_n$ で取る。ここで $\omega_1,\ldots,\omega_n$ は、$\mathsf{K}$ の整数全体の系、すなわち主オーダー $\mathfrak o$ の加群基底を表す。} このとき定義イデアル $\mathfrak M$ は $\omega$ の間の関係全体に対応し、したがって $f(\omega)$ が零となるすべての整係数多項式 $f(x)$ からなる。一方、微分商 $\mathfrak M'[x\to\omega]$ は、点 $x=\omega$ で取った差分商として定義される。

実際、$f(\omega)=0$ であるから、イデアル $\mathfrak M$ をすべての差 $f(x)-f(\omega)$ からなるものとみなし、さらに差分イデアル $\mathfrak P=(x_1-\omega_1,\ldots,x_n-\omega_n)$ と差分商 $\mathfrak A=\mathfrak M:\mathfrak P$ を作る。ここで商は通常のイデアル商を意味し、$K$ の整数、すなわち $\mathfrak o$ の元を係数とする $x$ の多項式環の中で取られる。このとき共役差積 $\mathfrak D$ は
\[
\mathfrak D=\mathfrak A[x\to\omega].
\]
によって定義される。差分イデアルと差分商をそもそも定義できるようにするためには、係数領域を拡大する必要がある。したがって、これは一つの不定元 $x$ の多項式の微分商
\[
f'(\xi)=(f(x)-f(\xi)):(x-\xi),
\]
を点 $x=\xi$ で取るという構成の直接の一般化である。この場合にも、差を定義するために、$\xi$ を添加して多項式環を拡大しなければならない。実際、$\omega$ の基底が一つの元の冪からなるならば、イデアル微分商は多項式微分商に帰着する。

ここで与えた共役差積の定義は、このように既知の概念へ直接接続する。しかし、イデアル $\mathfrak M$ が選んだ基底に依存するため、不変でない中間対象を導入している。これと同値な \emph{完全に不変な定義} は、整係数多項式環の代わりに、主オーダー $\mathfrak o$ と同型な環 $\mathfrak O$、すなわち $\mathfrak M$ を法とする剰余類環と同型な環を導入し、その係数領域を $\mathfrak o$ によって拡大して $\mathfrak O_{\mathfrak o}$ とすることにより得られる。\srcfn{2)}{係数領域の拡大、すなわち直積形成は、§ 1 では非可換環についても扱われており、本論文の目的に必要な範囲より一般的である。本論文の目的には、§ 2 を可換環かつ有限基底の場合に特殊化すれば十分である。} ここでは、差分イデアル $\mathfrak P$ は、$\mathfrak O$ と $\mathfrak o$ の同型によって互いに対応する $x$ と $\xi$ のすべての差 $(x-\xi)$ から導かれる。差分商 $\mathfrak A$ は零イデアルの $\mathfrak P$ による商に等しく、$\mathfrak A$ の各 $x$ を同型に対応する $\xi$ で置き換えることによって共役差積が得られる。

以下では、この不変な定義を出発点とする。こうして、係数拡大の可能性に関する一定の仮定だけのもとで、任意の可換環の共役差積を部分環に関して直接定義する。この仮定は、たとえば加群基底が存在すれば、必要に応じて適当な商環へ移ることにより常に満たされる。これによって、とりわけ数体の任意のオーダーの共役差積、たとえば素数冪を法とするオーダーの剰余類環の共役差積、さらに相対共役差積、および複数の不定元の代数関数の場合の共役差積も定義される。共役差積の微分的定義が通常の定義と一致することは、直積形成によって数体に付随するガロア拡大環について、直和分解に基づく精密な構造解析を行うことから従う。\srcfn{3)}{ここで行うのは、私の判別式に関する論文の基礎をなす考察の発展であるが、その論文の知識は前提としない。E. Noether, Der Diskriminantensatz für die Ordnungen eines algebraischen Zahl- und Funktionenkörpers, J. Math. 157 (1927), 82--104 を参照。} 一つの不定元の多項式を法とする剰余類環という特殊な場合、この解析はラグランジュ補間公式の解析に帰着する。より詳しく述べれば、次のようになる。

数体 $\mathsf{K}$ と同型で $\mathfrak O$ を含む環 $\mathfrak K$ を導入し、その係数領域である有理数体を $K$ のガロア閉包 $F$ によって拡大する。この拡大された系 $\mathfrak K_F$ は、$K$ の $n$ 個の同型、すなわち各共役体への移行に対応する $n$ 個の単純成分の直和となり、これらの成分は差分商 $\mathfrak A$ とその各共役の拡大である。同時に、$\mathfrak K_\Gamma$ の零イデアルは、差分イデアルとその各共役から導かれる $n$ 個のイデアルの共通部分となる。単純成分は $\Gamma e^{(i)}$ の形をもち、ここで $e^{(i)}$ は単位元の成分を表す。$e^{(1)}$ は $x=\xi$ とすると単位元へ移るので、差分商は $\mathfrak d e^{(1)}$ に等しい。ここで $\mathfrak d$ は $\mathfrak o$ の共役差積を表す。共役差積 $\mathfrak d$ は、$e^{(1)}$ を掛けると $\mathfrak O_{\mathfrak o}$、すなわちオーダーとそれ自身との直積に属するような体の元すべてからなる。各共役についても同様である。

ここから補加群による定義を得るには、$e^{(i)}$ が、$\mathfrak O$ または $\mathfrak K$ の、$\omega$ に対応する基底 $x_1,\ldots,x_n$ に関する線形形式となることに注意すればよい。各 $e^{(i)}$ における $x$ の係数は、まさに $\mathfrak o$（およびその各共役）の補加群の一つの基底を与える。差分商 $\mathfrak A$ の係数が $\mathfrak o$ に属するという事実から、共役差積が $\mathfrak o$ のその補加群による商に等しいことが直ちに従い、これによってデデキントの定義が得られる。以上の考察は主オーダーに限られず、任意のオーダーの共役差積を、そのオーダーの補加群による商として特徴づける。

主オーダーの場合、微分的定義と基本方程式による定義との一致は、基底が一つの元、ここでは基本形式の冪からなるならばイデアル微分商が多項式微分商に帰着するという、先に述べた事実から従う。これによって、後者の定義とデデキントの定義との概念的な関係も明らかになる。

主オーダーの場合、基本方程式を基礎に置けば、冒頭で述べた主定理が直接従うことはよく知られている。指数を正確に決定するには、十分大きな $t$ に対して、主オーダーの $p^t$ を法とする剰余類環の共役差積 $\mathfrak d[p^t]$ を考えればよい。$t=n$ で十分であり、二次体の場合にはこれが必要でもある。ここで本質的なのは、共役差積 $\mathfrak d$ を mod $p^t$ で考えると、それが正確に $\mathfrak d[p^t]$ へ移ることである。さらに、環の直和には共役差積の直和が対応するので、$\mathfrak p^{t\varrho}$ を法とする剰余類環に制限してよい。この剰余類環は、mod $p^t$ の係数をもつ一つの不定元の多項式を法とする剰余類環と同型になる。そして、別の方法でこのような多項式に到達した Ö. Ore が示したように、\srcfn{4)}{刊行時の注記：原稿ではこの引用が補われていない。本論文の概略だけが記された § 7 の文献指示を参照。} この多項式を微分すれば、補助数を用いて、起こりうるすべての指数を正確に見渡すことができる。

最後に、この構造解析から、$\mathsf{K}$ の上体の共役差積は、その相対共役差積と $\mathsf{K}$ の共役差積との積に等しいという既知の定理も得られる。これは、直和分解における単位元の成分、したがって補加群についても、同じ乗法性が成り立つことの帰結である。

以上の事実はすべて、私が判別式に関する論文で相対判別式について行ったのとまったく同様に、適当な商環へ移ることによって相対共役差積にも成り立つ。さらに、関数体へ移れば、係数体の標数が零である代数関数体についてもすべてそのまま成り立つ。一方、標数 $p$ の場合および整代数関数の場合には、構造的性質は保たれるものの（本質的には関数体へ移ることさえ不要である）、非分離拡大体が現れるために分岐理論が変わる。この点はここではこれ以上扱わない。最も本質的な成果として挙げたいのは、上に概説した完全に不変な構造定理（§ 5 および § 6）であり、これは複数の不定元の代数関数にも成り立つ。補加群への移行は、すでに係数の恒等式へ分解することを意味し、既知の結果を位置づけるためにのみ行ったものである。
% 原典：論文43、原刊頁 4--6 （雑誌の頁 693--695）
\subsection*{§ 1. 環の係数領域の拡大または直積形成。定義イデアル\srcfnmark{5)}}
\srcfntext{5)}{本節について批判的な助言を寄せられた B. L. van der Waerden に感謝する。}

本節では、乗法の交換法則を仮定しない一般の環を扱う。一方、簡単のために単位元の存在を仮定する。この仮定は、明示的に反対のことを記さない限り、本論文全体を通じて置かれる。

\paragraph{1. 直積の定義。}
環 $\mathfrak O\times\mathfrak o$ または $\mathfrak O_{\mathfrak o}$ が、$\mathfrak h$ に関する環 $\mathfrak O$ と $\mathfrak o$ の \emph{直積}、あるいは \emph{$\mathfrak O$ の係数領域 $\mathfrak h$ を $\mathfrak o$ へ拡大して得られた環}\srcfn{6)}{以下で主に用いる第二の記法は係数領域の拡大を想起させるためのものであり、第一の記法は直積の通常の表記である。} と呼ばれるのは、次の条件が満たされるときである。

\begin{enumerate}
\item[(1)] $\mathfrak O_{\mathfrak o}$ は $\mathfrak O$ と $\mathfrak o$ を部分環として含み、これら二環の積、すなわち $\mathfrak O_{\mathfrak o}$ の中で $\mathfrak O$ と $\mathfrak o$ を含むすべての環の共通部分に等しい。

\item[(2)] $\mathfrak O$ と $\mathfrak o$ の共通部分は $\mathfrak h$ であり、$\mathfrak h$ は $\mathfrak O_{\mathfrak o}$ の単位元を含む。

\item[(3)] $\mathfrak O$ と $\mathfrak o$ は元ごとに可換である。したがって、$\mathfrak O_{\mathfrak o}$ はすべての双線形結合
\[
x_{i_1}y_{i_1}+\cdots+x_{i_r}y_{i_r},
\]
からなり、ここで $x$ と $y$ はそれぞれ $\mathfrak O$ と $\mathfrak o$ の有限個の元からなる任意の系を走る。また $\mathfrak h$ は $\mathfrak O$ と $\mathfrak o$ の中心に、したがって $\mathfrak O_{\mathfrak o}$ の中心にも含まれる。

\item[(4)] $\mathfrak O_{\mathfrak o}$ の二つの元が等しいための必要十分条件は、$\mathfrak O$ において成り立つ関係と $\mathfrak o$ において成り立つ関係とによって、少なくとも一通り形式的に等しくできることである。より正確には、
\[
x_1y_1+\cdots+x_ny_n=0
\quad\hbox{において }\mathfrak O_{\mathfrak o},
\]
が成り立つならば、$h$ は $\mathfrak h$、$\gamma$ は $\mathfrak O$、$\delta$ は $\mathfrak o$ の元として、形式的に零となる結合 $(\gamma h)\delta-\gamma(h\delta)$ の和を加えることにより、この式を
\[
(z_1\alpha_1+\cdots+z_r\alpha_r)+(t_1\beta_1+\cdots+t_s\beta_s)
\]
の形へ変形できる。ここで $z_1=0,\ldots,z_r=0$ および $\beta_1=0,\ldots,\beta_s=0$ であり、$\alpha,\beta$ は $\mathfrak o$、$z,t$ は $\mathfrak O$ の元である。したがって、$z=0$ は $\mathfrak O$ における $x,\gamma,h$ の間の関係を、$\beta=0$ は $\mathfrak o$ における $y,\delta,h$ の間の関係を表す。\srcfn{7)}{単位元の存在を仮定しない場合には、(3) と (4) にさらに線形項が加わる。}
\end{enumerate}

$\mathfrak O_{\mathfrak o}$ の二つの元が等しいための必要十分条件は、その差が上述の仕方で消えることである。

したがって、直積の要請により、等号関係と演算関係（加法および乗法）は各因子における関係から一意に定まる。これより次が従う。$\overline{\mathfrak O},\overline{\mathfrak o}$ およびその共通部分 $\overline{\mathfrak h}$ が、それぞれ $\mathfrak O,\mathfrak o$ および $\mathfrak h$ と同型であり、直積 $\mathfrak O\times\mathfrak o$ が存在するならば、$\overline{\mathfrak O}\times\overline{\mathfrak o}$ も存在し、$\mathfrak O\times\mathfrak o$ と同型になる。

任意の環 $\mathfrak S$ が $\overline{\mathfrak O}$ と $\overline{\mathfrak o}$ の積に等しく、しかも $\overline{\mathfrak O}$ と $\overline{\mathfrak o}$ が元ごとに可換であるならば、$\mathfrak S$ は $\mathfrak O\times\mathfrak o$ の準同型像となる。このとき、$\overline{\mathfrak O}$ から $\mathfrak O$ への写像、および $\overline{\mathfrak o}$ から $\mathfrak o$ への写像は同型のままである。実際、$\mathfrak O\times\mathfrak o$ における等号関係は $\mathfrak S$ における等号関係に含まれ、$\overline{\mathfrak O}$ および $\overline{\mathfrak o}$ のそれぞれに限れば、$\mathfrak O\times\mathfrak o$ の中での $\mathfrak O$ および $\mathfrak o$ の等号関係と一致する。\srcfn{8)}{§ 1, 4 も参照。}

\paragraph{2. 部分環に関する直積の存在の必要十分条件。}
$\mathfrak O$ と $\mathfrak o$ を、その共通部分が両者の中心に含まれる環 $\mathfrak h$ であるような環とする。このとき、次の例が示すように、$\mathfrak h$ に関する直積は必ずしも存在しない。$\mathfrak h$ を有理整数環、$\mathfrak O=\mathfrak h[x]$ を関係 $1-2x=0$ をもつ環、$\mathfrak o=\mathfrak h[y]$ を関係 $1-2y=0$ をもつ環とし、$y$ は新しい記号であるとする。すると $\mathfrak O$ と $\mathfrak o$ は $\mathfrak h$ に関して同値な環であり、集合論的な共通部分は $\mathfrak h$ であるが、記号 $x$ と $y$ の間には何の演算関係も与えられていない。ところが、$\mathfrak O$ と $\mathfrak o$ を含み、$R$ の中では $x$ と $y$ の間にも演算関係が存在し、しかも両者の共通部分が $\mathfrak h$ となるような環 $R$ は存在しない。実際、$R$ では
\[
0=(1-2x)y-x(1-2y)=y-x,
\]
が従う。この推論には $x$ と $y$ の可換性さえ用いていない。\srcfn{9)}{固定された環の中にある可換環の商環、ここでは $\frac12$ の添加によって得られる環は一意だからである。これは、同値ではあるが相異なる二つのガロア体拡大を一つの共通上体へ埋め込めないことに対応する。ただし、ここでは直積形成によって、零因子をもつ環へ埋め込むことはできる。}

\emph{直積が存在するための必要十分条件} は次のとおりである。

\emph{$\mathfrak O$ と $\mathfrak o$ を含む環 $R$ であって、その中で $\mathfrak O$ と $\mathfrak o$ の共通部分が $\mathfrak h$ であり、かつ $\mathfrak O$ と $\mathfrak o$ が元ごとに可換となるものが、少なくとも一つ存在しなければならない。}

実際、直積が存在すれば、それ自身がそのような環 $R$ である。逆に直積を構成する際には、$\mathfrak h$ に属さない $\mathfrak O$ の元と $\mathfrak o$ の元との間に演算関係が存在しないと仮定してよい。以下では常にこの仮定を置くが、これは $\mathfrak h$ の同値な拡大へ移り、別の記号を導入することによっていつでも実現できる。\srcfn{9a)}{§ 2, 1 および脚注 12) を参照。} さらに、$\mathfrak h$ は $\mathfrak O$ と $\mathfrak o$ の単位元 $e$ を含むものとする（ただし注 7) も参照）。いま $\mathfrak T$ を、すべての双線形結合
\[
x_{i_1}y_{i_1}+\cdots+x_{i_r}y_{i_r}
 =y_{i_1}x_{i_1}+\cdots+y_{i_r}x_{i_r},
\]
の集合として定義する。ここで $x$ と $y$ はそれぞれ $\mathfrak O$ と $\mathfrak o$ の有限個の元からなる任意の系を走り、したがって $\mathfrak T$ の中では $x$ と $y$ は可換とされる。このような二つの結合が等しいとは、1, (4) で明確にした意味において、$\mathfrak O$ と $\mathfrak o$ における関係によって形式的に等しくなることと定める。さらに、
\[
\alpha_{i_1}x_{i_1}+\cdots+\alpha_{i_r}x_{i_r}
\quad\hbox{および}\quad
\beta_{i_1}x_{i_1}+\cdots+\beta_{i_r}x_{i_r},
\]
について、ここで $\alpha,\beta$ には零が現れてもよいものとし、その和と積をそれぞれ
\[
(\alpha_{i_1}+\beta_{i_1})x_{i_1}+\cdots+(\alpha_{i_r}+\beta_{i_r})x_{i_r}
\quad\hbox{および}\quad
\sum\alpha_{i_\lambda}\beta_{i_\mu}x_{i_\lambda}x_{i_\mu}.
\]
によって定義する。この定義は、等号の意味に関して一意である。実際、任意の零関係（1, (4) の意味での零関係）に $\mathfrak o$ または $\mathfrak O$ の元を右または左から掛けても再び同種の零関係となり、それらの零関係の和と差についても同じことが成り立つ。系 $\mathfrak T$ は $\mathfrak O$ と $\mathfrak o$ を含む環となり、等号、加法および乗法について直積の条件を満たす。さらに $\mathfrak T$ は、一般の埋め込みに関する仮定から、共通部分についての条件も満たす。実際、$\mathfrak S$ を $R$ の中での $\mathfrak O$ と $\mathfrak o$ の積とすれば、1 により $\mathfrak S$ は $\mathfrak T$ の準同型像であり、その写像は $\mathfrak O$ と $\mathfrak o$ のそれぞれに限れば同型である。したがって、$x$ が $\mathfrak O$ に、$y$ が $\mathfrak o$ に属するとき、$\mathfrak T$ における関係 $x=y$ から $\mathfrak S$ における同じ関係が従う。仮定により、$\mathfrak S$ では両者とも $\mathfrak h$ に属するので、$\mathfrak O$ と $\mathfrak o$ のそれぞれに関する同型性から、$\mathfrak T$ でも同じことが成り立つ。ゆえに、この埋め込み条件は必要十分である。

\paragraph{3. 部分環に関する環の定義イデアル。}
$\mathfrak O$ の元 $\ldots,z,\ldots$ からなる系 $S$ が、$\mathfrak h$ に関する $\mathfrak O$ の \emph{生成系} であるとは、$\mathfrak O$ が $\mathfrak h$ と $S$ によって $\mathfrak O$ の中で生成される環、すなわち $\mathfrak O$ の中で $\mathfrak h$ と $S$ を含むすべての環の共通部分に等しいことをいう。記号では $\mathfrak O=\mathfrak h[S]$ と書く。再び $\mathfrak h$ が $\mathfrak O$ の中心の部分環であると仮定する。このとき $\mathfrak O$ は非可換多項式環 $\mathfrak h[..Z..]$ の準同型像となる。ここで $Z$ は不定元を表し、その不定元系は $S$ と同じ濃度をもつ。非可換多項式環とは、すべての有限和 $\sum h_iP_i=\sum P_i h_i$ からなる系をいう。ここで $P_i$ は、指数 $0$ に対応する単位元も含めた冪積
\[
P_i=Z_{i_1}^{\varrho_1}Z_{i_2}^{\varrho_2}\cdots Z_{i_n}^{\varrho_n}
\]
であり（もちろん同じ添字が繰り返されてもよい）、等号は異なる冪積に関する係数の等しさとして定義される。この準同型から、$\mathfrak O$ は剰余類環 $\mathfrak h[..Z..]/\mathfrak M$ と同型になる。ここで両側イデアル $\mathfrak M$ は、$\mathfrak h[\ldots Z\ldots]$ の元であって、$\mathfrak O$ において $F(z)$ が零となるようなすべての多項式 $F(Z)$ からなる。ただし $Z$ と $z$ はこの準同型によって対応する。\emph{イデアル $\mathfrak M$ を $\mathfrak h$ に関する定義イデアル、より詳しくは $z$ の生成系に対応する定義イデアルという。} この定義から、同型に対応する部分環 $\mathfrak h,\overline{\mathfrak h}$ と生成系 $S$ および $\overline S$ に関する同型な環は、同型な定義イデアルをもつことが分かる。とくに $\mathfrak h$ が $\overline{\mathfrak h}$ と一致する場合には、それらの定義イデアルを同一とみなしてよい。

$\mathfrak h$ は $\mathfrak O$ の中心に含まれると仮定したので、生成系 $S$ が一つの元 $z$ だけからなるならば、$\mathfrak O$ は必ず可換である。さらに $\mathfrak M$ が主イデアルであるならば、$\mathfrak M$ の基底多項式 $G(Z)$ に対応する方程式 $G(z)=0$ を、$\mathfrak h$ に関する $\mathfrak O$ の定義方程式という。

\paragraph{4. 直積の定義イデアル。}
3 と同じく、$\mathfrak O$ は $\mathfrak h[..Z..]/\mathfrak M$ と同型であり、$\mathfrak M$ は定義イデアルであるとする。$\mathfrak h[..Z..]$ と並んで、多項式環 $\mathfrak o[..Z..]$、すなわちすべての有限和
\[
\sum \gamma_iP_i=\sum P_i\gamma_i
\qquad\hbox{ただし }\gamma_i\hbox{ は }\mathfrak o\hbox{ に属する},
\]
からなる系を考え、等号を係数の等しさとして定義する。$\mathfrak o[Z]$ は $\mathfrak h[Z]$ を含み、$\mathfrak h$ に関する $\mathfrak h[..Z..]$ と $\mathfrak o$ の直積である。言い換えれば、$\mathfrak h[..Z..]$ の係数領域を拡大して得られる。\srcfn{10)}{実際、等号関係と演算関係についての条件は満たされ、$\gamma=h+h_1P_1+\cdots+h_rP_r$ からは、$\mathfrak o[..Z..]$ における等号の定義により $\gamma=h$ が従う。§ 2, 2 も参照。ただし、$\mathfrak o[..Z..]$ が直積であるという事実は、本項では用いない。} $\mathfrak M_{\mathfrak o}$ を、$\mathfrak o[..Z..]$ に関する $\mathfrak M$ の拡大イデアル、すなわち $\mathfrak o[..Z..]$ の中で $\mathfrak M$ を含むすべての両側イデアルの共通部分とする。これは $\mathfrak M$ の元の、$\mathfrak o$ の係数によるすべての線形結合からなる系であり、$\mathfrak o$ と $Z$ が可換であるため、すでに両側イデアルをなす。\emph{$\mathfrak h$ に関する直積 $\mathfrak O_{\mathfrak o}$ が存在するならば、$\mathfrak O_{\mathfrak o}$ は $\mathfrak o[..Z..]/\mathfrak M_{\mathfrak o}$ と同型である。} 実際、$\mathfrak O_{\mathfrak o}$ は $\mathfrak o[..Z..]$ の準同型像であり、したがって $z$ の生成系に対応する、$\mathfrak o$ に関する $\mathfrak O_{\mathfrak o}$ の定義イデアルを法とする剰余類環と同型である。しかし等号の定義 1, (4) により、$\mathfrak O_{\mathfrak o}$ において $z$ の間に成り立つ関係は、$\mathfrak O$ において成り立つ関係の、$\mathfrak o$ の係数をもつ線形結合として書ける。したがって $\mathfrak M_{\mathfrak o}$ が定義イデアルとなる。

\paragraph{備考。}
同じ考察から、直積とそれに属する定義イデアルについて、対称的ではあるがより煩雑な表示も得られる。$\mathfrak o$ が $\mathfrak h[..Y..]/\mathfrak N$ と同型であり、$Y$ は $Z$ と異なる記号を表すものとして、多項式環 $\mathfrak h[..Y..Z..]$ を定義する。このとき $\mathfrak O_{\mathfrak o}$ は
\[
\mathfrak h[..Y..Z..]/(\mathfrak M,\mathfrak N,..YZ-ZY..),
\]
と同型である。ここで右辺の定義イデアルは $\mathfrak M,\mathfrak N$ とすべての結合 $YZ-ZY$ によって生成される。後者は $\mathfrak o$ と $\mathfrak O$ の元どうしの可換性を表し、前二者は $\mathfrak O_{\mathfrak o}$ における関係が $\mathfrak O$ と $\mathfrak o$ のそれぞれにおける関係から従うことを表す。
% 原典：論文43、原刊頁 7--10 （雑誌の頁 696--699）
\subsection*{§ 2. 独立な加群基底をもつ環。直積の構成。}

\paragraph{1. 直積の構成。}
$\mathfrak O$ の元 $\ldots,t,\ldots$ からなる系 $T$ が $\mathfrak O$ の $\mathfrak h$-加群基底であるとは、$\mathfrak O$ が $T$ によって $\mathfrak O$ の中で生成される $\mathfrak h$-加群に等しいことをいう。このとき、$\mathfrak h$ に単位元が存在するので、$\mathfrak O$ はすべての線形結合 $h_{i_1}t_{i_1}+\cdots+h_{i_r}t_{i_r}$ からなる。加群基底が独立であるとは、
\[
h_{i_1}t_{i_1}+\cdots+h_{i_r}t_{i_r}=0
\quad\hbox{常に}\quad
h_{i_1}=0,\ldots,h_{i_r}=0
\]
が従うことをいう。\emph{$\mathfrak O$ が単位元を含む独立な加群基底をもつならば、$\mathfrak h$ の任意の拡大環 $\mathfrak o$ に対して直積 $\mathfrak O_{\mathfrak o}$ が存在する。ただし、$\mathfrak h$ は中心の部分環であり、集合論的な共通部分 $[\mathfrak O,\mathfrak o]$ は $\mathfrak h$ であるとする。}\srcfn{11)}{もちろん、これと対称に、$\mathfrak o$ に基底の仮定を置くこともできる。§ 2, 2 により、一般性を失わずに、$\mathfrak O$ と $\mathfrak o$ のうち $\mathfrak h$ に属さない元の間には演算関係が存在しないと仮定する。E. Noether, Der Diskriminantensatz für die Ordnungen eines algebraischen Zahl- oder Funktionenkörpers, dies. Journ. 157 (1927), 82--104, §§ 1, 2 も参照。} § 1, 2 により、対応する埋め込み環 $\mathfrak R$ を構成すれば十分であり、ここではそれがすでに直積であることも分かる。$\mathfrak R$ を、すべての線形結合
\[
\gamma_{i_1}t_{i_1}+\cdots+\gamma_{i_r}t_{i_r}
=t_{i_1}\gamma_{i_1}+\cdots+t_{i_r}\gamma_{i_r},
\]
の全体として定義する。ここで $t_{i_r}$ は独立な基底 $T$ の有限個の元からなる任意の系を走り、それに応じて $\gamma$ は $\mathfrak o$ を走る。等号は $t_{i_\nu}$ に関する係数の等しさとして定義し、和と積は § 1, 2 と同様に通常の計算法則に従って定義する。基底の性質により、任意の二つの $t_{i_\nu}$ の積は有限個の基底元の、$\mathfrak h$ の係数をもつ線形結合であるから、$\mathfrak R$ は環となる。$\mathfrak h$ が $\mathfrak o$ の部分環なので、$\mathfrak R$ は $\mathfrak O$ を含み、$T$ が単位元を含むので、$\mathfrak o$ も含む。さらに、$\gamma_{i_\nu}$ は $h_{i_\nu}$ とも $t_{i_\nu}$ とも元ごとに可換であるから、$\mathfrak R$ の中で $\mathfrak o$ と $\mathfrak O$ は元ごとに可換となる。最後に、$\mathfrak R$ における $\mathfrak O$ と $\mathfrak o$ の共通部分は $\mathfrak h$ である。実際、関係
\[
\gamma=\gamma e=h_1t_1+\cdots+h_rt_r
\]
から、等号の定義と、$e$ が仮定により基底の元であることによって、一つの $t$ が $e$ に等しく、他の項は消えることが従い、したがって $\gamma=h$ となる。ゆえに直積は存在し、$\mathfrak O$ と $\mathfrak o$ の積である $\mathfrak R$ に等しい。この環は可換性と等号についての条件を満たす。基底の元が独立であるため、後者の条件は係数の等しさに帰着する。この $\mathfrak O$ の基底 $T$ は、同時に $\mathfrak O_{\mathfrak o}$ の $\mathfrak o$-基底となる。

\paragraph{2. 拡大加群と縮約加群。}
$\mathfrak B$ が $\mathfrak O$ の $\mathfrak h$-加群であるとき、その拡大 $\mathfrak B_{\mathfrak o}$ とは、$\mathfrak O_{\mathfrak o}$ の $\mathfrak o$-加群であって、$\mathfrak O_{\mathfrak o}$ において $\mathfrak B$ を含むすべての $\mathfrak o$-加群の共通部分であるものをいう。$\mathfrak C$ が $\mathfrak O_{\mathfrak o}$ の $\mathfrak o$-加群であるとき、その $\mathfrak O$ における縮約加群 $\mathfrak C_{\mathfrak h}$ は、共通部分 $[\mathfrak C,\mathfrak O]$ として定義される。$\mathfrak B$ が同時に $\mathfrak O$ のイデアルであれば、$\mathfrak B_{\mathfrak o}$ も $\mathfrak O_{\mathfrak o}$ のイデアルとなる。また $\mathfrak C$ が $\mathfrak O_{\mathfrak o}$ のイデアルであれば、$[\mathfrak C,\mathfrak O]$ も $\mathfrak O$ のイデアルとなる。

\paragraph{定理。}
\emph{$\mathfrak O$ の加群 $\mathfrak B$ が独立な $\mathfrak h$-加群基底をもち、それを $\mathfrak O$ の独立な $\mathfrak h$-加群基底へ延長できるならば、$\mathfrak B$ は $\mathfrak O_{\mathfrak o}$ におけるその拡大加群 $\mathfrak B_{\mathfrak o}$ の縮約である。}

\paragraph{証明。}
$T$ を $\mathfrak B$ の独立な $\mathfrak h$-加群基底とする。仮定により、$\mathfrak O$ の元 $\overline t$ からなる系 $\overline T$ を付け加えて、$\mathfrak O$ の独立な $\mathfrak h$-加群基底へ延長できる。このとき $T$ は同時に $\mathfrak B_{\mathfrak o}$ の独立な $\mathfrak o$-基底となり、$T$ と $\overline T$ を合わせたものは $\mathfrak O_{\mathfrak o}$ の同様の独立な基底となる。

$\mathfrak B$ は共通部分 $[\mathfrak B_{\mathfrak o},\mathfrak O]$ に含まれる。逆に、この縮約加群の元 $c$ を取る。$\mathfrak B_{\mathfrak o}$ の元として、$c$ は $T$ の元による表示をもち、それは同時に $\mathfrak O_{\mathfrak o}$ の $\mathfrak o$-基底 $T,\overline T$ による表示でもある。すなわち、
\[
c=\gamma_1t_1+\cdots+\gamma_rt_r
\qquad\hbox{ただし }\gamma\hbox{ は }\mathfrak o\hbox{ に属する}.
\]
一方、$\mathfrak O$ の元として、$c$ は $\mathfrak O$ の $\mathfrak h$-基底 $T,\overline T$ による表示
\[
c=h_{i_1}t_{i_1}+\cdots+h_{i_r}t_{i_r}+h_{j_1}\overline t_{j_1}+\cdots+h_{j_s}\overline t_{j_s},
\]
をもつ。$\mathfrak O$ は $\mathfrak O_{\mathfrak o}$ の部分環であり、$\mathfrak h$ は $\mathfrak o$ の部分環なので、これは同時に $\mathfrak O_{\mathfrak o}$ の $\mathfrak o$-基底 $T,\overline T$ による表示でもある。したがって二つの表示の係数は一致しなければならず、$\overline t$ の係数はすべて零である。ゆえに
\[
c=h_{i_1}t_{i_1}+\cdots+h_{i_r}t_{i_r}
\]
となり、$c$ が $\mathfrak B$ の元であることが分かる。したがって、$\mathfrak B$ は縮約加群である。

\paragraph{補足。}
この証明は、定理に述べた仮定よりわずかに弱い仮定しか用いておらず、残りの部分も同時に証明している。したがって、次の強められた形が成り立つ。

\emph{$\mathfrak O$ に独立な $\mathfrak h$-加群基底 $T,\overline T$ が存在し、$T$ が $\mathfrak B$ に含まれ、しかも $T$ が同時に $\mathfrak B_{\mathfrak o}$ の（独立な）$\mathfrak o$-加群基底であるならば、$\mathfrak B$ はその拡大 $\mathfrak B_{\mathfrak o}$ の縮約加群である。このとき $T$ は $\mathfrak B$ の独立な $\mathfrak h$-加群基底でもあることが分かる。}

\paragraph{3. 直積の存在に関する十分条件。}
（たとえば相対共役差積の場合のように）係数領域 $\mathfrak h$ を拡大して初めて独立な加群基底が存在することがある。このとき、直積の存在は次の十分条件から導ける。

\emph{$\mathfrak h$ の拡大環 $\mathfrak f$ が存在し、$\mathfrak h$ に関する直積 $\mathfrak O_{\mathfrak f}$ と $\mathfrak o_{\mathfrak f}$ が存在し、さらに $\mathfrak f$ に関する $\mathfrak O_{\mathfrak f}\times\mathfrak o_{\mathfrak f}$ が存在するならば、$\mathfrak h$ に関する $\mathfrak O\times\mathfrak o$ が存在する。}

$\mathfrak R=\mathfrak O_{\mathfrak f}\times\mathfrak o_{\mathfrak f}$ が § 1, 2 の意味で埋め込み環となることを示せばよい。実際、$\mathfrak O$ と $\mathfrak o$ はそれぞれ $\mathfrak O_{\mathfrak f}$ と $\mathfrak o_{\mathfrak f}$ の部分環として、元ごとに可換となる。さらに
\[
[\mathfrak O,\mathfrak o]\subseteqq[\mathfrak O_{\mathfrak f},\mathfrak o_{\mathfrak f}]=\mathfrak f;
\]
が成り立つ。同様に $[\mathfrak O,\mathfrak o]\subseteqq\mathfrak O$ であるから、$[\mathfrak O,\mathfrak o]\subseteqq[\mathfrak O,\mathfrak f]=\mathfrak h$、したがって $[\mathfrak O,\mathfrak o]=\mathfrak h$ である。これにより、$\mathfrak R$ が埋め込み環であることが示された。

\paragraph{備考。}
この証明で共通部分に関する仮定を用いるのは $\mathfrak O_{\mathfrak f}$ についてだけであり、$\mathfrak o_{\mathfrak f}$ については用いない。したがって後者の環は、単に $\mathfrak o$ と $\mathfrak f$ の積であってもよい。
\subsection*{§ 3. 部分環に関する環の共役差積。定義イデアルの微分商。}

以後、現れる環はすべて可換であると仮定する。\srcfn{12)}{この仮定がなければ、右共役差積と左共役差積を定義し、それらが互いにどの程度一致するか、また半単純系のオーダーの場合に既知の構成とどの程度一致するかを調べなければならない（§ 4 の考察の一般化として）。} さらに、環 $\mathfrak O$ と $\mathfrak o$ は部分環 $\mathfrak h$ の同型かつ同値な拡大であること（すなわち、$\mathfrak O$ から $\mathfrak o$ への写像が $\mathfrak h$ 上の恒等写像の延長であること）、および $\mathfrak h$ に関する直積 $\mathfrak O_{\mathfrak o}$ が存在することを常に仮定する。

\paragraph{1. 共役差積の定義。}
$\mathfrak h$ に関する直積 $\mathfrak O_{\mathfrak o}$ が存在すると仮定する。$x$ を $\mathfrak O$ のすべての元にわたらせ、$\xi$ はそのつど同型によって $\mathfrak o$ に対応する元を表すものとする。$\mathfrak P$ を、すべての差 $x-\xi$ によって生成される $\mathfrak O_{\mathfrak o}$ の両側 \emph{差分イデアル} とする：
\[
\mathfrak P=\{..(x-\xi)..\}\srcfnmark{13)}
\]
\srcfntext{13)}{この記法は、集合によって生成されるイデアルに対して常に用いる。}
\emph{差分商} $\mathfrak A$ は、零イデアルを $\mathfrak P$ で割ったイデアル商として定義される。すなわち、$\mathfrak O_{\mathfrak o}$ の元 $a$ で $a\mathfrak P$ が零となるものすべてからなる：
\[
\mathfrak A=(0):\mathfrak P.
\]
$\mathfrak h$ に関する $\mathfrak o$ の \emph{共役差積} $\mathfrak d$ は $\mathfrak A[x\to\xi]$ と定義される。すなわち、$\mathfrak A$ から、$\mathfrak A$ の各元 $a=\alpha_{i_1}x_{i_1}+\cdots+\alpha_{i_r}x_{i_r}$ に現れる $x$ を、同型によって対応する $\xi$ で置き換えることによって得られる。同様に、$\mathfrak h$ に関する $\mathfrak O$ の共役差積 $\mathfrak D$ は $\mathfrak A[\xi\to x]$ と定義される。$\mathfrak A$ は $x$ と $\xi$ に関して対称に定義されているから、$\xi$ と $x$ を交換して得られる $\mathfrak d$ と $\mathfrak D$ は、$\mathfrak o$ から $\mathfrak O$ への同型によって互いに対応する。また、$\mathfrak A$ は $\mathfrak O_{\mathfrak o}$ のイデアルであり、したがってとくに $\mathfrak o$-加群かつ $\mathfrak O$-加群であるから、$\mathfrak d$ と $\mathfrak D$ はそれぞれ $\mathfrak o$ と $\mathfrak O$ のイデアルとなる。

\paragraph{2. 定義イデアルの微分商。}
(§ 1, 3) の意味で、$\mathfrak M$ を生成系 $z$ に対応する、$\mathfrak h$ に関する $\mathfrak O$ の定義イデアルとする。すると (§ 1, 3)、$\mathfrak M$ は同型な生成系 $\xi$ に対応する、$\mathfrak h$ に関する $\mathfrak o$ の定義イデアルでもある。また (§ 1, 4)、$\mathfrak O_{\mathfrak o}$ は $\mathfrak o[..Z..]/\mathfrak M_{\mathfrak o}$ と同型である。ここで、多項式環 $\mathfrak o[..Z..]$ における差分イデアルと差分商を、それぞれ次のように定義する：
\[
\mathfrak Q=\{..(Z-\xi)..\}
\]
および
\[
\mathfrak C=\mathfrak M_{\mathfrak o}:\mathfrak Q
=\{..(F(Z)-F(\xi))..\}:\{..(Z-\xi)..\}.
\]
\emph{微分商} $\mathfrak M'$ を $\mathfrak C[\xi\to Z]$ と定義する。この定義は、$\mathfrak o$ の元を生成元 $\xi$ で表す \emph{表示} に依存せず、一意であることに注意しなければならない。実際、$\mathfrak C$ はイデアル商なので $\mathfrak M_{\mathfrak o}$ を含む。$\mathfrak o$ の一つの元の二つの表示は多項式 $F(\xi)$ だけ異なる。したがって、これらに $\xi$ を $Z$ で置き換える操作を施すと、二つの表示は、$\mathfrak M$ に、ゆえにすでに $\mathfrak C$ に属する多項式 $F(Z)$ だけ異なる。こうして $\mathfrak M'$ は、$\mathfrak M$ を含む $\mathfrak h[..Z..]$ のイデアルとなる。\emph{微分商 $\mathfrak M'$ は、$Z\to\xi$ によって $\mathfrak o$ の共役差積 $\mathfrak d$ に移る。} なぜなら、準同型写像のもとでは、イデアルだけでなくその商も互いに対応するからである。実際、$Z$ と $z$ を対応させる $\mathfrak o[\ldots Z\ldots]$ から $\mathfrak O_{\mathfrak o}$ への準同型により、$\mathfrak Q$ と $\mathfrak M_{\mathfrak o}$ がそれぞれ $\mathfrak B$ と零イデアルに移るので、$\mathfrak A$ は $\mathfrak C$ の準同型像となる。\srcfn{13a)}{刊行時の注：商が対応することを結論するには、$\mathfrak M_{\mathfrak o}$ が $(0)$ の逆像全体であることが重要である。} したがって、この準同型のもとで、$\mathfrak C$ における置換 $\xi\to Z$ は $\mathfrak A$ における置換 $\xi\to z$ に対応する。ゆえに、$\mathfrak O$ の共役差積 $\mathfrak D$ は、$Z$ と $z$ を対応させる写像による $\mathfrak M'$ の準同型像である。したがって
\[
\mathfrak M'[Z\to z]=\mathfrak D
\qquad\hbox{および}\qquad
\mathfrak M'[Z\to\xi]=\mathfrak D[z\to\xi]=\mathfrak d.
\]

\paragraph{3. 定義方程式がある場合の共役差積。}
$\mathfrak O$ が定義方程式 $G(z)=0$ をもつとする（§ 1, 3 の末尾）。さらに、最高次項の係数が単位元であることを \emph{仮定} する。すなわち
\[
G(Z)=Z^n+h_1Z^{n-1}+\cdots+h_n
\]
であり、係数 $h_i$ は $\mathfrak h$ に属するものとする。このとき、多項式の微分商との関係は次の定理で与えられる。

\paragraph{定理。}
\emph{$\mathfrak h$ に関する $\mathfrak o$ および $\mathfrak O$ の共役差積は定義される（すなわち、$\mathfrak h$ に関する直積 $\mathfrak O_{\mathfrak o}$ が存在する）。それらは、それぞれ $G'(\xi)$ および $G'(z)$ を生成元とする主イデアルである。ここで $G'(Z)$ は多項式の微分商を表す。イデアルの微分商 $\mathfrak M'$ は $(G'(Z),G(Z))$ に等しい。}

\paragraph{証明。}
仮定からまず、$e,z,\ldots,z^{n-1}$ は $\mathfrak O$ の独立な加群基底をなす。実際、$G(Z)$ の最高次係数が単位元なので、これらの冪は $\mathfrak h$-加群基底をなす。この基底は独立でもある。なぜなら、$\mathfrak M$ は次数が $n$ より低い多項式を含みえないからである。実際、この最高次係数に関する仮定により、多項式 $K(Z)G(Z)$ の次数は両因子の次数の和に等しい。したがって加群基底が存在するので、$\mathfrak h$ に関する直積 $\mathfrak O_{\mathfrak o}$ が存在し、共役差積も定義される。残りは、$\mathfrak o[\ldots Z\ldots]$ における差分商 $\mathfrak C$ の考察から従う（2 を参照）。この差分商は、次の多項式を生成元とする主イデアルである：
\[
\begin{aligned}
C(Z,\xi)={}&(Z^{n-1}+Z^{n-2}\xi+\cdots+\xi^{n-1})\\
&+h_1(Z^{n-2}+\cdots+\xi^{n-2})+\cdots+h_{n-1}.
\end{aligned}
\]
実際、$C$ は $G(Z)-G(\xi)$ を $Z-\xi$ で形式的に割って得られるから、$\mathfrak Q$ の生成元が $Z-\xi$ であることにより $\mathfrak C$ に属する。任意の多項式は、$C$ を法として、$Z$ に関する次数が $n-1$ に達しない多項式 $D$ に簡約できる。したがって $D(Z-\xi)$ の次数は $n$ に達せず、$\mathfrak M$ に属するためには零でなければならない。しかし $Z-\xi$ の $Z$ の係数は単位元なので、$D(Z-\xi)=0$ から $D=0$ が従う。したがって $C$ は $\mathfrak C$ の生成元である。ここで、$C$ とその倍数に置換 $\xi\to Z$ を施し、さらに $G(Z)=C(Z-\xi)$ とその倍数を考えると、
\[
\mathfrak M'=(G'(Z),G(Z)),
\]
が得られ、したがって
\[
\mathfrak d=(G'(\xi))
\qquad\hbox{および}\qquad
\mathfrak D=(G'(z)).
\]
% 原典：論文43、原刊頁 10--13 （雑誌の頁 699--702）
\subsection*{\S{} 4. 直和の共役差積。}

$\mathfrak O$ について、次の仮定を置く：

\begin{enumerate}
\item[(1)] $\mathfrak O$ はイデアルの直和であり、
\[
\mathfrak O=\mathfrak R_1+\cdots+\mathfrak R_r
=\mathfrak R_i+\mathfrak S_i
=\mathfrak Oe_1+\cdots+\mathfrak Oe_r,
\]
ここで $e_i$ は単位元の成分を表す。

\item[(2)] 各 $\mathfrak R_i$ は独立な $\mathfrak h_i$-加群基底 $T_i$ をもつ。ここで $\mathfrak h_i$ は $\mathfrak R_i$ における $\mathfrak h$ の成分、すなわち $\mathfrak h_i=\mathfrak h e_i$ である。この基底 $T_i$ は $\mathfrak R_i$ の単位元 $e_i$ を元として含む。

\item[(3)] 各共通部分 $[\mathfrak h,\mathfrak S_i]$ は零元に等しい（$\mathfrak S_i$ は定数を含まない）。
\end{enumerate}

同型性により、$\mathfrak o$ についても同じ仮定が成り立つ。

\paragraph{定理。}
\emph{以上の仮定のもとで、$\mathfrak h$ に関する $\mathfrak O$ の共役差積は、$\mathfrak h_i$ に関する各 $\mathfrak R_i$ の共役差積の直和となる。}

証明は、以下の一連の事実を示すことによって行う。

\paragraph{I.}
まず、$\mathfrak O$ 自身も単位元 $e$ を含む独立な $\mathfrak h$-加群基底をもち、したがって $\mathfrak h$ に関する $\mathfrak O$ および $\mathfrak o$ の共役差積が定義されることが従う。この基底は、$\mathfrak R_i$ の基底 $T_i$ すべての合併 $T$ からなる。実際、
\[
x=xe_1+\cdots+xe_r
\]
および
\[
xe_i=(h_1e_i)t_{1i}e_i+\cdots+(h_\lambda e_i)t_{\lambda i}e_i
=h_1(t_{1i}e_i)+\cdots+h_\lambda(t_{\lambda i}e_i)
\]
から、$T$ が確かに $\mathfrak O$ の $\mathfrak h$-基底であることが従う。この基底は独立でもある。実際、$T$ の元の間の関係は、直和性により各 $T_i$ の元の間の関係を与え、したがって各係数について $h e_i=0$ を与える。しかし仮定 (3) により、これは $h=0$ を意味する。なぜなら、
\[
h\equiv h e_i\ (\mathfrak S_i)
\]
なので、$h e_i=0$ ならば $h$ は $[\mathfrak h,\mathfrak S_i]$ に属するからである。さらに $T$ は、$\mathfrak R_i$ の単位元 $e_i$ をすべて元として含み、その和は $e$ である。したがって $T$ に含まれる一つの $e_i$ を $e$ で置き換えると、すべての条件を満たす $\mathfrak h$-基底 $\overline T$ が得られる。これにより $\mathfrak O_{\mathfrak o}$ が存在し、$\mathfrak h$ に関する $\mathfrak O$ の共役差積が定義される。

\paragraph{II.}
共通部分に関する条件 (3) から、さらに強い条件
\[
[\mathfrak o,\mathfrak S_{i(\mathfrak o)}]=0,
\]
が従う。ここで $\mathfrak S_{i(\mathfrak o)}$ は、$\mathfrak S_i$ から $\mathfrak O_{\mathfrak o}$ において生成されるイデアルを表す。いま
\[
\alpha=\alpha e=\beta_1t_1+\cdots+\beta_\lambda t_\lambda,
\]
とする。ここで $t_1,\ldots,t_\lambda$ は、$T$ から $T_i$ を除いて得られる基底の元であり、$\alpha$ と $\beta$ は $\mathfrak o$ の元である。$e_i$ を $e$ で置き換えて $\overline T$ を作ると、$e,t_1,\ldots,t_\lambda$ は $\mathfrak O$ の $\mathfrak h$-基底 $\overline T$ の互いに異なる元となる。一方、$\mathfrak O_{\mathfrak o}$ における等式は、このような任意の基底について係数の等式に帰着するから、関係
\[
\alpha=\alpha e=\beta_1t_1+\cdots+\beta_\lambda t_\lambda,
\]
から $\alpha=0$ が従う。これで \emph{強められた共通部分性} が証明された。とくに、上と同様に、
\[
\alpha\text{ は }\mathfrak o\text{ の元}\quad\text{かつ}\quad\alpha\ne0
\]
ならば $\alpha e_i\ne0$ である。

\paragraph{III.}
この共通部分条件から直ちに、$\mathfrak o$ は $\mathfrak O$ の分解に対応する \emph{成分} $\mathfrak o e_i$ と \emph{環同型} であることが従う。この同型は
\[
\mathfrak h\simeq\mathfrak h_i=\mathfrak h e_i.
\]
という同型を含む。実際、$e_i^2=e_i$ であり、しかも $e_i$ は $\mathfrak o$ のすべての元と可換なので、$\mathfrak o e_i$ は $\mathfrak o$ の環同型像である。この対応は一対一でもある。なぜなら、$\alpha\ne0$ ならば $\alpha e_i\ne0$ だからである。このとき部分環 $\mathfrak h$ には成分 $\mathfrak h_i=\mathfrak h e_i$ が対応する。

\paragraph{IV.}
$\mathfrak o$ に対する対応する分解
\[
\mathfrak o=\mathfrak r_1+\cdots+\mathfrak r_r
=\mathfrak r_i+\mathfrak s_i
=\mathfrak o\varepsilon_1+\cdots+\mathfrak o\varepsilon_r.
\]
も用いると、\emph{さらに次の同型} が得られる：
\[
\mathfrak o\simeq\mathfrak o e_i,
\quad \mathfrak O\simeq\mathfrak O\varepsilon_i,
\quad \mathfrak r_\lambda\simeq\mathfrak r_\lambda e_i,
\quad \mathfrak R_\lambda\simeq\mathfrak R_\lambda\varepsilon_i;
\quad \mathfrak h\simeq\mathfrak h e_i\simeq\mathfrak h\varepsilon_i\simeq\mathfrak h e_i\varepsilon_i.
\]
$\mathfrak O\simeq\mathfrak O\varepsilon_i$ という関係は、上の考察で $\mathfrak O$ の代わりに $\mathfrak o$ から出発したとき、$\mathfrak o\simeq\mathfrak o e_i$ に正確に対応して得られる。同様に、続く二つの関係も互いに対応する。その最初のものは、$\mathfrak o\simeq\mathfrak o e_i$ のもとで部分環 $\mathfrak r_\lambda$ に $\mathfrak r_\lambda e_i$ が対応することから従う。また同じように、
\[
\mathfrak h\simeq\mathfrak h e_i\simeq\mathfrak h\varepsilon_i.
\]
が得られる。さらに $\mathfrak h\varepsilon_i$ は $\mathfrak o$（および $\mathfrak r_i$）の部分環なので、最後に次の同型も得られる：
\[
\mathfrak h\varepsilon_i\simeq\mathfrak h\varepsilon_i e_i.
\]

\paragraph{V.}
$\mathfrak O$ の直和表示から、次が従う。\emph{$\mathfrak C=(0):\mathfrak B$ が $\mathfrak O$ における零イデアルの $\mathfrak B$ によるイデアル商であるならば、各成分 $\mathfrak C_i$ は、$\mathfrak R_i$ の零イデアルの $\mathfrak B_i$ によるイデアル商に等しい。} 実際、
\[
\mathfrak B\mathfrak C=(\mathfrak B e_1)(\mathfrak C e_1)+\cdots+(\mathfrak B e_r)(\mathfrak C e_r)=0
\]
から、\emph{直和性} により各成分について
\[
(\mathfrak B e_i)(\mathfrak C e_i)=0.
\]
が従う。逆に、$\mathfrak R_i$ の元 $r e_i$ が条件 $(\mathfrak B e_i)(r e_i)=0$ を満たすならば、$i\ne k$ のとき $(\mathfrak B e_k)(r e_i)=0$ も成り立つ。

したがって $r e_i$ は $\mathfrak C$ の元であると同時に $\mathfrak R_i$ の元でもあり、成分 $\mathfrak C_i=\mathfrak C e_i$ に属する。これにより $\mathfrak C_i=\mathfrak C e_i$ は、零イデアルの $\mathfrak B_i$ によるイデアル商であることが分かる。また、
\[
r e_i\mathfrak B_i\equiv0\quad(\mathfrak S_i)
\]
から常に $r e_i\mathfrak B_i=0$ が従うので、これは $\mathfrak R_i$ の零イデアルの $\mathfrak B_i$ による商でもある。\srcfn{14)}{これは、準同型写像のもとではイデアルだけでなくその商も対応するという、2 で用いた定理から直接にも従う。} さらに、$\mathfrak O$ だけでなく
\[
\mathfrak O_{\mathfrak o}=\mathfrak O_{\mathfrak o}e_1+\cdots+\mathfrak O_{\mathfrak o}e_r
\]
も直和なので、\emph{この命題は $\mathfrak O_{\mathfrak o}$ のイデアルについても成り立つ。}

\paragraph{VI. 共役差積の成分分解。}
定義により（§ 3, 1 を参照）、$\mathfrak D=\mathfrak A[\xi\to x]$ である。したがって
\[
\mathfrak D=\mathfrak D e_1+\cdots+\mathfrak D e_r
=\mathfrak A[\xi\to x]e_1+\cdots+\mathfrak A[\xi\to x]e_r.
\]
一方、$e_i$ は $\mathfrak O$ に属するので置換のもとで不変であり、また $\varepsilon_i$ は同型によって $e_i$ に移り、$e_i^n=e_i$ である。したがって
\[
\mathfrak A[\xi\to x]e_i=\mathfrak A e_i[\xi\to x]
=\mathfrak A e_i\varepsilon_i[\xi\to x].
\]
V により、$\mathfrak A e_i\varepsilon_i$ は商 $(0):\mathfrak P e_i\varepsilon_i$ に等しく、さらに $\mathfrak R_i\mathfrak r_i$ の零イデアルの $\mathfrak P e_i\varepsilon_i$ による商にも等しい。実際、$\mathfrak O$ と $\mathfrak o$ の直和分解を掛け合わせると、
\[
\mathfrak O_{\mathfrak o}=\sum\mathfrak O_{\mathfrak o}e_i\varepsilon_k
=\sum\mathfrak R_i\mathfrak r_k,
\]
という分解が得られる。この和は、$\sum e_i\varepsilon_k=e$ となり、単位元の成分について直交関係が成り立つので、直和である。

\paragraph{VII.}
\emph{$\mathfrak D e_i=\mathfrak A e_i\varepsilon_i[\xi\to x]=(0):\mathfrak P e_i\varepsilon_i$ が、$\mathfrak h_i$ に関する $\mathfrak R_i$ の共役差積であることを示す。} この証明は、次の事実に基づく。$\mathfrak h e_i\varepsilon_i$ に関して、$\mathfrak R_i\mathfrak r_i$ は $\mathfrak R_i\varepsilon_i\times\mathfrak r_i e_i$ に等しい。その差分イデアルは $\mathfrak P e_i\varepsilon_i$ であり、その差分商は $\mathfrak A e_i\varepsilon_i$ である。一方、$\mathfrak D e_i\varepsilon_i$ は $\mathfrak R_i\varepsilon_i$ の共役差積となり、同様に $\mathfrak d\varepsilon_i e_i$ は $\mathfrak r_i e_i$ の共役差積となる。さらに、同型
\[
\mathfrak R_i\varepsilon_i\simeq\mathfrak R_i
\qquad\text{および}\qquad
\mathfrak r_i\simeq\mathfrak r_i e_i
\]
により、$\mathfrak D e_i$ は $\mathfrak h e_i$ に関する $\mathfrak R_i$ の共役差積となり、$\mathfrak d\varepsilon_i$ は $\mathfrak h\varepsilon_i$ に関する $\mathfrak r_i$ の共役差積となる。

\paragraph{VIII.}
\emph{$\mathfrak h e_i\varepsilon_i$ に関して、$\mathfrak R_i\mathfrak r_i$ は直積 $\mathfrak R_i\varepsilon_i\times\mathfrak r_i e_i$ に等しい。} 実際、$\mathfrak R_i$ は $\mathfrak h_i$-加群基底 $T_i$ をもつので、同型 (IV) $\mathfrak R_i\simeq\mathfrak R_i\varepsilon_i$ により、$T_i\varepsilon_i$ は $\mathfrak R_i\varepsilon_i$ の $\mathfrak h e_i\varepsilon_i$-加群基底となる。また IV により、$\mathfrak r_i e_i$ は $\mathfrak R_i\varepsilon_i$ と同型であり（実際、$\mathfrak o\simeq\mathfrak O$ により $\mathfrak r_i\simeq\mathfrak R_i$ である）、$\mathfrak h e_i\varepsilon_i$ の拡大環である。したがって直積は、$T_i\varepsilon_i$ の元の線形形式で係数を $\mathfrak r_i e_i$ から取るものすべてからなり、等号は係数の等号として定義される。ところが、$T_i=T_i e_i$、$\mathfrak r_i=\mathfrak r_i\varepsilon_i$ であり、$T_i$ は $\mathfrak O$ の加群基底の一部をなすので、直積 $\mathfrak O\times\mathfrak o$ の中でこれはちょうど $\mathfrak R_i\mathfrak r_i$ となる。

\paragraph{IX.}
\emph{$\mathfrak R_i\varepsilon_i\times\mathfrak r_i e_i$ の差分イデアルは $\mathfrak P e_i\varepsilon_i$ で与えられ、その差分商は $\mathfrak A e_i\varepsilon_i$ で与えられる。} これを見るため、$\mathfrak P$ を次の形に書く：
\[
\mathfrak P=\{\ldots,xe_1-\xi\varepsilon_1,\ldots,xe_r-\xi\varepsilon_r,\ldots\}.
\]
すると
\[
\mathfrak P e_i\varepsilon_i
=\{\ldots,xe_i\varepsilon_i-\xi\varepsilon_i e_i,\ldots\},
\]
となり、これは実際に $\mathfrak R_i\varepsilon_i\times\mathfrak r_i e_i$ の差分イデアルに等しい。また (VI. により)、$\mathfrak A e\varepsilon_i$ はこの環の零イデアルの $\mathfrak P e_i\varepsilon_i$ による商に等しいので、差分商であることが分かる。

\paragraph{X.}
\emph{$\mathfrak R_i\varepsilon_i$ の共役差積は $\mathfrak D e_i\varepsilon_i$ に等しい。} $\mathfrak R_i\varepsilon_i$ の共役差積は、$\mathfrak A e_i\varepsilon_i$ において $\mathfrak r_i e_i$ の各元を同型によって対応する元で置き換えることにより得られる。すなわち、各 $(\xi\varepsilon_i)e_i$ を $(xe_i)\varepsilon_i$ で置き換える。$\xi\varepsilon_i$ 自身も $\mathfrak o$ の元であり、それは $xe_i$ に移るので、これは $\mathfrak A e_i\varepsilon_i$ に置換 $\xi\to x$ を施し --- それによって $(\xi\varepsilon_i)e_i$ は $(xe_i)e_i=xe_i$ に移る ---、その後 $\varepsilon_i$ を掛けることに等しい。ところで (VI. により) $\mathfrak A e_i\varepsilon_i[\xi\to x]=\mathfrak D e_i$ であったから、$\mathfrak h e_i\varepsilon_i$ に関する $\mathfrak R_i\varepsilon_i$ の共役差積は $\mathfrak D e_i\varepsilon_i$ に等しい。これと対称的に、$\mathfrak r_i e_i$ の共役差積は $\mathfrak d e_i\varepsilon_i$ に等しい。

\paragraph{XI.}
\emph{$\mathfrak R_i$ の共役差積は $\mathfrak D e_i$ に等しい。} 実際、同型によって、$\mathfrak R_i$ の共役差積は $\varepsilon_i$ を掛けることで $\mathfrak R_i\varepsilon_i$ の共役差積に移る。後者は $\mathfrak d e_i\varepsilon_i$ に等しいから、前者は $\mathfrak D e_i$ に等しい。したがって
\[
\mathfrak D=\mathfrak D e_i+\cdots+\mathfrak D e_r,
\]
\emph{かつ $\mathfrak D e_i$ は $\mathfrak R_i$ の共役差積である。} これと対称的に、
\[
\mathfrak d=\mathfrak d\varepsilon_1+\cdots+\mathfrak d\varepsilon_r,
\]
\emph{かつ $\mathfrak d\varepsilon_i$ は $\mathfrak r_i$ の共役差積である。} 以上で定理が証明された。
% 原典：論文43、原刊頁 13--17 （雑誌の頁 702--706）
\subsection*{\S{} 5. 体のガロア拡大環。構造定理。}

\paragraph{1. 基礎となる環。}
$\mathsf{K}$ を基礎体 $P$ 上 $n$ 次の体とし、$K/P$ は分離拡大であると仮定する。$\Gamma$ をこれに属するガロア体とする。これは同型を除いて一意に定まり、$K=K^{(1)}$ とおけば、$n$ 個の共役体 $K^{(i)}$ を含み、それらの合成体となる。$\mathfrak K$ は $\mathsf{K}$、したがってすべての $\mathsf{K}^{(i)}$ と同型で、$\mathsf{P}$ に関して同値な $\mathsf{P}$ の拡大であり、共通部分 $[\mathfrak K,\Gamma]$ が $\mathsf{P}$ となるものとする（すなわち、$\Gamma$ と $\mathfrak K$ をともに含む共通の上体は考えない）。\srcfn{15)}{ここでは当然、一般性を失うことなく、$\mathsf{P}$ に属さない $\mathsf{K}$ の元と $\mathfrak K$ の元との間には何の結合関係も課されていないと仮定する（§1, 2、§2, 1 および注11参照）。} $\mathfrak K$ は単位元を含む有限独立な $\mathsf{P}$-加群基底をもつから、§2, 1 により直積 $\mathfrak K_\Gamma$、$\mathfrak K_{\mathsf{K}}$、$\mathfrak K_{\mathsf{K}^{(i)}}$ が存在し、$\mathfrak K_\Gamma$ はその $n$ 個の部分環 $\mathfrak K_{\mathsf{K}^{(i)}}$ の積となる。$\mathfrak K_\Gamma$ を $\mathfrak K$ の\emph{ガロア拡大環}と呼ぶ。また $\mathfrak K$ の任意の独立な $P$-基底が $\mathfrak K_K$ の独立な $K$-基底となり、共役についても同様であるから、$\mathfrak K_\Gamma$ は同時に、$K$ に関する $\mathfrak K_{\mathsf{K}}$ と $\Gamma$ の直積でもある。\srcfn{16)}{すべての共役について一様に成り立つ定理は、通常 $\mathfrak K_{\mathsf{K}}$ についてだけ、すなわち添字を省いて述べる。} これらの環のイデアルの間には次の関係がある。

\paragraph{定理。}
\emph{$\mathfrak K$ の任意のイデアルは、$\mathfrak K_{\mathsf{K}}$ または $\mathfrak K_\Gamma$ におけるその拡大イデアルの縮約イデアルである。また $\mathfrak K_{\mathsf{K}}$ の任意のイデアルは、$\mathfrak K_\Gamma$ におけるその拡大イデアルの縮約イデアルである。}

\paragraph{証明。}
$\mathsf{P}$ と $\mathsf{K}$ は体であるから、$\mathfrak K$ の任意の $P$-部分加群、および $\mathfrak K_K$ の任意の $K$-部分加群は、それぞれ有限独立な $\mathsf{P}$-加群基底、$\mathsf{K}$-加群基底をもち、それを $\mathfrak K$、$\mathfrak K_{\mathsf{K}}$ の同種の基底へ延長できる。したがって定理は §2, 2 の定理から従う。

$\mathfrak P_{\mathsf{K}^{(i)}}^{(i)}$ で $\mathfrak K_{\mathsf{K}^{(i)}}$ の差分イデアルを表す。すなわち
\[
\mathfrak P_{\mathsf{K}^{(i)}}^{(i)}=\{\ldots,x-\xi^{(i)},\ldots\},
\]
とくに
\[
\mathfrak P_{\mathsf{K}}=\{\ldots,x-\xi,\ldots\}.
\]
ここで $\xi$ または $\xi^{(i)}$ は、同型によって $x$ に対応する $\mathsf{K}$ または $\mathsf{K}^{(i)}$ の元である。したがって、この $n$ 個の差分イデアルは互いに共役である。

$\mathfrak A_{\mathsf{K}}$ および $\mathfrak A_{\mathsf{K}^{(i)}}$ で、$\mathfrak K_K$ における差分商 $(0):\mathfrak P_{\mathsf{K}}$ およびその共役を表す。$\mathfrak K_\Gamma$ における拡大イデアルを $\mathfrak P_\Gamma$、$\mathfrak A_\Gamma$ と書く。この項の拡大・縮約イデアルに関する定理から、とくに次が従う。

\emph{差分イデアル $\mathfrak P_{\mathsf{K}}$ と差分商 $\mathfrak A_{\mathsf{K}}$ は、それぞれ拡大イデアル $\mathfrak P_\Gamma$ と $\mathfrak A_\Gamma$ の縮約イデアルである。}

\paragraph{2.ガロア拡大環の直和分解。}
差分イデアルおよび差分商とガロア拡大環の構造との関係は、次の定理によって与えられる。

\paragraph{構造定理。}
\emph{ガロア拡大環 $\mathfrak K_\Gamma$ の零イデアルは、$n$ 個の共役な差分イデアルの拡大の共通部分であり、$\mathfrak K_\Gamma$ 自身は、$n$ 個の共役な差分商の拡大の直和である。$\mathfrak K_{\mathsf{K}^{(i)}}$ は対応する差分イデアルと差分商との直和であり、その零イデアルは両者の共通部分である。}

\paragraph{証明。}
各 $\mathfrak P_\Gamma^{(i)}$ による剰余類環は $\Gamma$ 上階数一である。実際、すべての元は $\mathfrak P_\Gamma^{(i)}$ を法として $\Gamma$ のある元と合同である一方、$\mathfrak P_\Gamma^{(i)}$ は $\Gamma$ の非零元を含み得ない（$x\to\xi^{(i)}$ とすれば $\mathfrak P_\Gamma^{(i)}$ は零になる）。分離性により $\mathfrak P_\Gamma^{(i)}$ はすべて相異なるので、これらは二つずつ互いに素である（任意の二つの和は $\mathfrak K_\Gamma$ に等しい）。したがってイデアル
\[
\mathfrak P_\Gamma^{(1)},
[\mathfrak P_\Gamma^{(1)},\mathfrak P_\Gamma^{(2)}],\ldots,
[\mathfrak P_\Gamma^{(1)},\ldots,\mathfrak P_\Gamma^{(n)}]
\]
の階数は各段階で少なくとも一ずつ減少する。すべての共通部分は零イデアルとなるが、共役イデアルを一つでも欠く共通部分は、互いに素であるため零イデアルとは異なる（したがって各段階で階数はちょうど一ずつ減少した）。既知の定理により、ここから $\mathfrak K_\Gamma$ の次の一意な直和表示が従う。\srcfn{17)}{例えば E. Noether, \emph{Abstrakter Aufbau der Idealtheorie in algebraischen Zahl- und Funktionenkörpern}, Math. Ann. 96 (1927), 26--61, §§4--5 を参照。}
\begin{align}
\mathfrak K_\Gamma
&=\mathfrak C_\Gamma^{(1)}+\cdots+\mathfrak C_\Gamma^{(n)}
=\Gamma e^{(1)}+\cdots+\Gamma e^{(n)}
=\mathfrak C_\Gamma^{(i)}+\mathfrak B_\Gamma^{(i)}, \tag{1}\\
e&=e^{(1)}+\cdots+e^{(n)},\qquad
\mathfrak C_\Gamma^{(i)}=
[\mathfrak B_\Gamma^{(1)},\ldots,\mathfrak B_\Gamma^{(i-1)},
\mathfrak B_\Gamma^{(i+1)},\ldots,\mathfrak B_\Gamma^{(n)}].\notag
\end{align}
$\mathfrak C_\Gamma$ が $\mathfrak A_\Gamma$ に等しいことを示そう。そのためには $\mathfrak K_{\mathsf{K}}$ における対応する一致を証明すればよい。ここでは
\begin{equation}
\mathfrak K_K=\mathfrak C_K+\mathfrak B_K
=\mathfrak K_K e^{(1)}+\mathfrak K_K(e-e^{(1)}). \tag{2}
\end{equation}
という分解が成り立つ。実際 $e^{(1)}$ は $\mathfrak K_{\mathsf{K}}$ に属する。なぜなら $\mathfrak C_\Gamma^{(1)}$ は $\mathfrak P_\Gamma^{(2)},\ldots,\mathfrak P_\Gamma^{(n)}$ とのすべての交換を許し、その単位元成分 $e^{(1)}$ についても同じことが成り立つからである。そこで $e^{(1)}$ を $\mathfrak K_K$ の独立な $K$-基底で表す。この基底は同時に $\mathfrak K_\Gamma$ の独立な $\Gamma$-基底でもあるから、この表示は可能であり、各係数は $\mathsf{K}^{(2)},\ldots,\mathsf{K}^{(n)}$ とのすべての交換を許すので $\mathsf{K}$ に属する。$(e^{(1)})^2=e^{(1)}$ であるため、和
\[
\mathfrak K_{\mathsf{K}}=\mathfrak K_{\mathsf{K}} e^{(1)}+\mathfrak K_{\mathsf{K}}(e-e^{(1)})
\]
は直和である。その $\mathfrak K_\Gamma$ への拡大も直和であり、(1) により $\mathfrak K_\Gamma=\mathfrak C_\Gamma^{(1)}+\mathfrak P_\Gamma^{(1)}$ を与える。ゆえに §2, 2 の末尾から
\[
\mathfrak K_{\mathsf{K}} e^{(1)}=[\mathfrak C_\Gamma,\mathfrak K_{\mathsf{K}}]
=[\Gamma e^{(1)},\mathfrak K_{\mathsf{K}}]=\mathsf{K}e^{(1)}=\mathfrak C_{\mathsf{K}};
\]
\[
\mathfrak K_{\mathsf{K}}(e-e^{(1)})=[\mathfrak P_\Gamma,\mathfrak K_{\mathsf{K}}]=\mathfrak P_{\mathsf{K}},
\]
となり、(2) が証明された。

(2) から $\mathfrak C_{\mathsf{K}}=\mathfrak A_{\mathsf{K}}$、したがって $\mathfrak C_\Gamma=\mathfrak A_\Gamma$ が従う。実際 $e^{(1)}(e-e^{(1)})=0$ なので $\mathfrak C_{\mathsf{K}}\mathfrak P_{\mathsf{K}}=0$ である。逆に $b\mathfrak P_{\mathsf{K}}=0$ なら $b(e-e^{(1)})=0$ であり、従って $b=be=be^{(1)}$ となる。ゆえに
\[
\mathfrak C_{\mathsf{K}}=(0):\mathfrak P_{\mathsf{K}}=\mathfrak A_{\mathsf{K}}
\]
また拡大によって $\mathfrak C_\Gamma=\mathfrak A_\Gamma$ である。したがって
\begin{align}
\mathfrak K_\Gamma
&=\mathfrak A_\Gamma^{(1)}+\cdots+\mathfrak A_\Gamma^{(n)}
=\Gamma e^{(1)}+\cdots+\Gamma e^{(n)};
&0&=[\mathfrak B_\Gamma^{(1)},\ldots,\mathfrak B_\Gamma^{(n)}]; \tag{3}\\
\mathfrak K_{K^{(i)}}
&=\mathfrak A_{K^{(i)}}^{(i)}+\mathfrak B_{K^{(i)}}^{(i)}
=K^{(i)}e^{(i)}+\mathfrak B_{K^{(i)}}^{(i)};
&0&=[\mathfrak A_{K^{(i)}}^{(i)},\mathfrak B_{K^{(i)}}^{(i)}].\notag
\end{align}
これで定理のすべての部分が証明された。なお、ここで現れる構成 $\mathfrak K_{\mathsf{K}}$、$\mathfrak K_\Gamma$、$\mathfrak P_{\mathsf{K}}$、$\mathfrak A_{\mathsf{K}}$ はすべて基底によらず定義されているので、これは\emph{完全に不変な構造定理}である。\srcfn{18)}{ここで可能な、定義方程式と対応する多項式の一次因子への分解による定式化は不変ではない。}

\paragraph{3. 共役元による $\mathfrak K$ の成分表示。}
共役な差分イデアルと差分商には、$\mathfrak K$ の元の共役成分による表示が対応する。

\paragraph{定理。}
\emph{$x\in\mathfrak K$ の $\mathfrak K$-成分表示が
\[
x=\xi^{(1)}e^{(1)}+\cdots+\xi^{(n)}e^{(n)}
\]
であるなら、その $n$ 個の成分は互いに共役である。とくに $\xi^{(1)},\ldots,\xi^{(n)}$ は、$\mathfrak K$ から $\mathsf{K}^{(1)},\ldots,\mathsf{K}^{(n)}$ への同型によって $x$ に対応する $n$ 個の共役な値である。}

$e^{(1)},\ldots,e^{(n)}$ が互いに共役であることは、2 の関係式 (3) から分かる。実際、$\mathfrak P_{\mathsf{K}}$ を $\mathfrak P_{\mathsf{K}^{(i)}}^{(i)}$ に移す $\Gamma$ の自己同型は、同時に $e^{(1)}$ を $e^{(i)}$ に移し、後者は $\mathfrak K_{\mathsf{K}^{(i)}}$ の元である。同じ関係式と $\mathfrak B_{K^{(i)}}^{(i)}e^{(i)}=0$ から $xe^{(i)}=\xi^{(i)}e^{(i)}$、従って
\[
(x-\xi^{(i)})e^{(i)}=0
\]
が得られる（あるいは、体としての $\mathfrak K$ がその成分と同型であり、その成分が $\mathsf{K}^{(i)}e^{(i)}$ で、$x$ と $\xi^{(i)}e^{(i)}$ が同型に対応することからも従う）。

この定理から次が従う。

\paragraph{系。}
\emph{$\mathfrak B$ を $\mathfrak K$ の絶対加群（すなわち、任意の二元とともにその差も含む加群）とし、$\mathfrak B^{(1)},\ldots,\mathfrak B^{(n)}$ を $\mathfrak K_\Gamma$ におけるその成分とする。このとき $\mathfrak B$ は共通部分 $[\mathfrak K,\mathfrak B^{(1)}+\cdots+\mathfrak B^{(n)}]$ に等しい。}

$\mathfrak B$ がこの共通部分に含まれることは明らかである。逆に $\mathfrak B^{(1)}+\cdots+\mathfrak B^{(n)}$ の元が $\mathfrak K$ に属するなら、定理によりそれは共役な成分の和でなければならず、従って $\mathfrak B$ の元となる。

\paragraph{4. $\mathfrak K$ の補基底と単位元成分との関係。}
不変な構造から基底表示へ移ると、この構造は基底を補基底の対にまとめる。

\paragraph{定理。}
\emph{$\mathfrak K$ の任意の $P$-基底 $t_1,\ldots,t_n$ には補基底 $T_1,\ldots,T_n$ が対応し、その補基底は再びもとの $t$-基底である。補基底 $T_1,\ldots,T_n$ に同型に対応する $K^{(i)}$ の基底 $\mathsf{A}_1^{(i)},\ldots,\mathsf{A}_n^{(i)}$ は、$t$ による単位元成分の基底表示
\[
e^{(i)}=\mathsf{A}_1^{(i)}t_1+\cdots+\mathsf{A}_n^{(i)}t_n.
\]
の係数として定義される。$\alpha_1^{(i)},\ldots,\alpha_n^{(i)}$ を $t_i$ に対応する基底とすれば、$\alpha$ の行列と $\mathsf{A}$ の行列は互いに逆行列である。二つの基底 $(s)$ と $(t)$ が変換
\[
(s)=(t)P
\]
で結ばれるとき、補基底の間には反変関係
\[
(S)=P^{-1}(T).
\]
が成り立つ。}

\paragraph{証明。}
$\mathfrak K$ の $P$-基底 $t_1,\ldots,t_n$ を固定する。これは $\mathfrak K_\Gamma$ の $\Gamma$-基底でもある。単位元の $n$ 個の成分の基底表示を $e^{(i)}=A_1^{(i)}t_1+\cdots+A_n^{(i)}t_n$ とする。2 と 3 により $A_1^{(i)},\ldots,A_n^{(i)}$ は $K^{(i)}$ に属し、異なる $e^{(i)}$ には共役な係数系 $A$ が対応する。$\alpha_1^{(i)},\ldots,\alpha_n^{(i)}$ を、$K^{(i)}$ において $t_1,\ldots,t_n$ と同型に対応する基底とする。このときさらに
\begin{equation}
e^{(i)}[t\to\alpha^{(i)}]=e;
\qquad e^{(i)}[t\to\alpha^{(j)}]=0\quad(i\ne j). \tag{1}
\end{equation}
が成り立つ。実際 $\mathfrak B_{K^{(i)}}^{(i)}$ は $t_1-\alpha_1^{(i)},\ldots,t_n-\alpha_n^{(i)}$ を含み、これらは $\mathfrak B_{K^{(i)}}^{(i)}$ の従属な $P$-加群基底をなすので、代入 $t=\alpha^{(i)}$ によって零になる。従って
\[
e\equiv e^{(i)}\;(\mathfrak P_{\mathsf{K}^{(i)}}^{(i)})
\]
から (1) の第一の関係が従う。第二の関係は、$i\ne j$ のとき $e^{(i)}$ がすべての $\mathfrak P_{\mathsf{K}^{(j)}}^{(j)}$ に属することから従う。係数について解けば、(1) は行列の形で
\begin{equation}
\begin{pmatrix}
\alpha_1^{(1)}&\cdots&\alpha_n^{(1)}\\
\vdots&&\vdots\\
\alpha_1^{(n)}&\cdots&\alpha_n^{(n)}
\end{pmatrix}
\begin{pmatrix}
A_1^{(1)}&\cdots&A_1^{(n)}\\
\vdots&&\vdots\\
A_n^{(1)}&\cdots&A_n^{(n)}
\end{pmatrix}=E_n. \tag{2}
\end{equation}
となる。ここで $E_n$ は単位行列であり、これによって $\mathsf{A}$ は $\alpha$ から一意に定まる。(2) は両行列の行列式が零でないことを示す。また $\mathsf{A}_1^{(i)},\ldots,\mathsf{A}_n^{(i)}$ が $\mathsf{P}$ 上線型独立であることも示す。もし線型従属なら、共役な係数系についても同じ関係が成り立ち、$\mathsf{A}$-行列の階数が低下して、$E_n$ の階数が $n$ であることに反するからである。従って各 $\mathsf{A}_1^{(i)},\ldots,\mathsf{A}_n^{(i)}$ は $\mathsf{K}^{(i)}$ の基底である。これらの基底は互いに共役なので、すべて $\mathfrak K$ の同じ基底 $T_1,\ldots,T_n$ と同型に対応し、3 により
\[
T_\lambda=\mathsf{A}_\lambda^{(1)}e^{(1)}+\cdots+\mathsf{A}_\lambda^{(n)}e^{(n)}
\]
で定まる。この $T_1,\ldots,T_n$ を $t_1,\ldots,t_n$ の補基底と呼び、同様に $\mathsf{A}_1^{(i)},\ldots,\mathsf{A}_n^{(i)}$ を $\alpha_1^{(i)},\ldots,\alpha_n^{(i)}$ の補基底と呼ぶ。(2) により $\mathsf{A}$ はすでに一意に定まり、また転置行列へ移れば同じ関係が $T$ の基底、すなわち $\mathsf{A}$ から $\alpha$ を定めるので、もとの $t$-基底は $T$-基底の補基底である。このように、単位元の成分による対応はすべての基底を補基底の対に分ける。

$s,S$ をもう一つの補基底の対とし、
\[
e^{(i)}=\mathsf{B}_1^{(i)}s_1+\cdots+\mathsf{B}_n^{(i)}s_n
\qquad\text{かつ}\qquad (s)=(t)P,
\]
とする。
\[
\mathsf{A}_1t_1+\cdots+\mathsf{A}_nt_n=\mathsf{B}_1s_1+\cdots+\mathsf{B}_ns_n
\]
に $(s)=(t)P$ を代入し、$t$ に関する恒等式として比較すれば、$B,A$ は $s,t$ と反変的に変換される。従って同型により $S,T$ についても同じ関係が成り立つ。これで定理はすべて証明された。

\paragraph{注意。}
$c=C_1t_1+\cdots+C_nt_n$ を $P$-基底による $c\in\mathfrak K$ の表示、$c=c_1T_1+\cdots+c_nT_n$ を補基底による同じ元の表示とすれば
\[
C_\lambda=\operatorname{Sp}(cT_\lambda),
\qquad c_\lambda=\operatorname{Sp}(ct_\lambda),
\]
となる。ここで $\operatorname{Sp}(x)$ は通常どおり共役の和、すなわち任意の $x\in\mathfrak K$ に対して
\[
\operatorname{Sp}(x)=\xi^{(1)}+\xi^{(2)}+\cdots+\xi^{(n)}
\]
を意味する。実際
\[
c=\gamma^{(1)}e^{(1)}+\cdots+\gamma^{(n)}e^{(n)}
=\sum_\lambda\gamma^{(1)}A_\lambda^{(1)}t_\lambda+
\cdots+\sum_\lambda\gamma^{(n)}A_\lambda^{(n)}t_\lambda,
\]
であるから、係数を比較すれば $C_\lambda=\operatorname{Sp}(cT_\lambda)$ を得る。第二の関係は $t$ と $T$ を交換して同様に従う。\srcfn{19)}{この項の考察は、不変な構造定理が非不変な基底定理より本質的に簡単であることを示す。後者は既知の事柄との関係を示すためだけに必要であった。実際、ガロア拡大環の分解と単位元成分の導入は、補基底（次節の補加群）が果たすことをすべて果たす。次節で共役差積と補加群との関係を与えるのも、分類上の位置を示すためにすぎない。5 におけるラグランジュ補間公式の考察についても同様である。}

\paragraph{5. 定義方程式を基礎とする場合。ラグランジュ補間公式。}
［刊行時注：原稿ではこの項は空欄のままである。］
% 原典：論文43、原刊頁 17--21 （雑誌の頁 706--710）
\subsection*{\S{} 6. オーダーの共役差積。}

ここで、体とそのガロア拡大環だけに関する構造定理を、$\mathsf{P}$ の特定の部分環と、それらに関して $\mathfrak K$ および $\mathsf{K}$ で整な元、すなわちオーダーを指定することにより、整性の意味で強化する。§ 5 の記号はすべて引き続き用いる。\srcfn{20)}{環上整および整閉の概念については、脚注 17 に引用した論文の § 1 を参照。}

\paragraph{1. 基礎となる環。}
$\mathfrak h$ を、その商体が $\mathsf{P}$ となるような $\mathsf{P}$ の部分環とし、さらに $\mathfrak h$ は $\mathsf{P}$ において \emph{整閉} であるとする（すなわち、$\mathsf{P}$ の元で $\mathfrak h$ 上整なものはすべて $\mathfrak h$ に属する）。また $\mathfrak O$ を $\mathfrak K$ の $\mathfrak h$-オーダー、すなわち $\mathfrak h$ を含み、商体が $\mathfrak K$ であり、有限な（必ずしも独立でない）$\mathfrak h$-加群基底をもつ $\mathfrak K$ の部分環とする。このとき $\mathfrak O$ のすべての元は $\mathfrak h$ 上整であり、その $\mathfrak h$-加群基底は同時に $\mathfrak K$ の（一般には従属な）$P$-加群基底となる。$\mathfrak o$ を $K$ において $\mathfrak O$ と同型に対応するオーダー、$\mathfrak o^{(i)}$ を $K^{(i)}$ におけるその共役とする。\emph{$n$ 個の共役オーダーが $\Gamma$ の中で生成する環を $\mathfrak g$ とする。}このとき $\mathfrak g$ もオーダーである。実際、$\mathfrak g$ は $\mathfrak h$ を含み、各 $\mathfrak o^{(i)}$ の $\mathfrak h$-基底の元をすべて掛け合わせて得られる積は、$\mathfrak g$ の $\mathfrak h$-基底であると同時に $\Gamma$ の $P$-基底となる。

\emph{直積 $\mathfrak O_{\mathfrak o}$ と $\mathfrak O_{\mathfrak g}$ は存在する。}これは § 2, 3 の十分条件から従う（そこに現れる環 $\mathfrak f$ を $\mathsf{P}$ で置き換える）。まず $\mathfrak K_\Gamma$ の中に、$\mathfrak h$ に関する直積 $\mathfrak O_{\mathsf{P}}=\mathfrak K$ が存在する。実際、$\mathfrak h$ は交わり $[\mathfrak O,\mathsf{P}]$ に含まれる。逆に、この交わりの元は $\mathfrak O$ の元として整であり、かつ $\mathsf{P}$ の元であるから、$\mathfrak h$ の整閉性により $\mathfrak h$ に属する。$\mathfrak O$ の任意の $n$ 個の線形独立な元の $P$-係数線形結合は、直積として $\mathfrak K$ を与える。同様に
\[
\mathfrak o_{\mathsf{P}}=\mathsf{K}
\qquad\text{および}\qquad
\mathfrak g_{\mathsf{P}}=\Gamma;
\]
$\mathfrak h$ に関する直積 $\mathfrak K_{\mathsf{K}}=\mathfrak O_{\mathsf{P}}\times\mathfrak o_{\mathsf{P}}$ から、§ 2, 3 の判定法により $\mathfrak h$ に関する $\mathfrak O_{\mathfrak o}$ の存在が従う。同様に、$\mathfrak K_\Gamma=\mathfrak O_{\mathsf{P}}\times\mathfrak g_{\mathsf{P}}$ から $\mathfrak O_{\mathfrak g}$ の存在が従う。$\mathfrak O_{\mathfrak o}$ とともに共役環 $\mathfrak O_{\mathfrak o^{(i)}}$ も存在する。$\mathfrak O_{\mathfrak g}$ を \emph{オーダー $\mathfrak O$ のガロア拡大環} と呼ぶ。

したがって、$\mathfrak h$ に関する $\mathfrak O$ および $\mathfrak o$ の共役差積と、それらの共役な共役差積とが定義される。ここで
\[
\mathfrak P=\{\ldots,y-\eta,\ldots\}
\]
を $\mathfrak O_{\mathfrak o}$ の差分イデアルとし、
\[
\mathfrak P^{(i)}=\{\ldots,y-\eta^{(i)},\ldots\}
\]
を $\mathfrak O_{\mathfrak o^{(i)}}$ の共役な差分イデアルとする。また
\[
\mathfrak A^{(i)}=(0):\mathfrak P^{(i)}
\quad\text{--- 商は $\mathfrak O_{\mathfrak o^{(i)}}$ において取る ---}
\]
を共役な \emph{差分商} とする。特に
\[
\mathfrak A=(0):\mathfrak P\quad\text{において }\mathfrak O_{\mathfrak o}.
\]
さらに
\[
\mathfrak D=\mathfrak A[\eta\to y]
\]
を $\mathfrak O$ の共役差積とし、
\[
\mathfrak d^{(i)}=\mathfrak A[y\to\eta^{(i)}]
\]
を $n$ 個の共役オーダー $\mathfrak o^{(i)}$ の共役差積とする。特に
\[
\mathfrak d=\mathfrak A[y\to\eta]
\]
は $\mathfrak o$ の共役差積である。\srcfn{21)}{$\mathfrak h$ が $P$ において整閉であるという仮定は、たとえば $\mathfrak h$ が数体の主オーダーである場合、特に $\mathfrak h$ が有理整数全体の環である場合に満たされる。したがって、相対体の場合も含め、任意のオーダーの共役差積と相対共役差積が定義される。また $\mathfrak h$ を $n$ 個の不定元の多項式領域とすることもでき、この場合には $n$ 不定元の代数関数体における共役差積が得られる。}

\paragraph{2. 構造定理。}
体のガロア拡大環に対する構造定理（§ 5, 2）に対応して、次が成り立つ。

\paragraph{オーダーのガロア拡大環に対する構造定理。}
\emph{オーダーと体に関する差分イデアルおよび差分商は、次の関係にある。$\mathfrak P_{\mathsf{K}}$ は $\mathfrak O_{\mathfrak o}$ の差分イデアル $\mathfrak P$ の拡大であり、同様に $\mathfrak A_K$ は $\mathfrak O_{\mathfrak o}$ の差分商 $\mathfrak A$ の拡大である。また $\mathfrak A$ は $\mathfrak A_{\mathsf{K}}$ の縮約である。差分商、共役差積、および単位元の成分の間には次の関係がある。}
\[
\mathfrak A=\mathfrak d e^{(1)},
\qquad
\mathfrak A^{(i)}=\mathfrak d^{(i)}e^{(i)},
\]
\emph{したがって}
\[
\mathfrak d=[\mathfrak o,\mathfrak A^{(1)}+\cdots+\mathfrak A^{(n)}].
\]
\emph{$\mathfrak o$ の共役差積 $\mathfrak d$ は、$\varkappa e^{(1)}$ が $\mathfrak O_{\mathfrak o}$ に属するような $K$ の元 $x$ 全体からなり、零イデアルとは異なる。}

$\mathfrak P_{\mathsf{K}}$ が $\mathfrak P$ の拡大であることは、$\mathfrak K$ が $\mathfrak O_{\mathsf{P}}$ に等しく、また $\mathsf{K}$ が $\mathfrak o_{\mathsf{P}}$ に等しいため、$\mathfrak B_K$ の各 $x-\xi$ が $\mathfrak B$ の元 $y-\eta$ の $K$-係数線形結合として表されることから従う。さらに、$\mathfrak A_{\mathsf{K}}=(0):\mathfrak P_{\mathsf{K}}$ は $\mathfrak K_K$ における商 $(0):\mathfrak P$ としても定義される。したがって $\mathfrak O_{\mathfrak o}$ における $\mathfrak A=(0):\mathfrak P$ は、$\mathfrak A_K$ に属する $\mathfrak O_{\mathfrak o}$ の元全体からなり、$\mathfrak A=[\mathfrak O_{\mathfrak o},\mathfrak A_{\mathsf{K}}]$、すなわち $\mathfrak A_{\mathsf{K}}$ の縮約である。$\mathfrak A_{\mathsf{K}}=\mathsf{K}e^{(1)}$ であるから、$\mathfrak A=\mathfrak z e^{(1)}$ と書ける。ここで $\mathfrak A$ は $\mathfrak o$-加群でもあるので、$\mathfrak z$ は $K$ の $\mathfrak o$-加群である。一方、共役差積の定義から
\[
\mathfrak d=\mathfrak A[x\to\xi]
=(\mathfrak z e^{(1)})[x\to\xi]
=\mathfrak z\bigl(e^{(1)}[x\to\xi]\bigr)=\mathfrak z e=\mathfrak z
\]
となる（§ 5, 4, (1) 参照）。ゆえに $\mathfrak z=\mathfrak d$ である。差分商 $\mathfrak A$ は、$\mathfrak o$ の共役差積 $\mathfrak d$ の $K$ における成分であり、共役な場合にも同様である。§ 5, 3 の系により、ここから上の $\mathfrak d$ に関する関係が従う。また
\[
\mathfrak A=\mathfrak d e^{(1)}=[\mathsf{K}e^{(1)},\mathfrak O_{\mathfrak o}]
\]
は、$\mathfrak d$ が、$\varkappa e^{(1)}$ が $\mathfrak O_{\mathfrak o}$ に属するような $K$ の元 $x$ 全体からなることを示す。また $\mathfrak d$ は零イデアルではない。実際、$e^{(1)}$ は $\mathfrak K_{\mathsf{K}}=\mathfrak O_{\mathsf{K}}$ の元として、$\mathfrak O$ の有限個の元の $\mathsf{K}$-係数線形結合であり、したがって $\mathfrak O_{\mathfrak o}$ の元を $\mathfrak o$（さらには $\mathfrak h$）の元で割った商である。この分母は $\mathfrak d$ の零でない元を与える。したがって $\mathfrak d$ は $P$ 上階数 $n$、すなわち $\mathfrak d_{\mathsf{P}}=\mathsf{K}$ であり、$\mathfrak A_{\mathsf{K}}$ は $\mathfrak A$ の拡大となる。

\begin{center}
\textbf{ここから先は概略である。}
\end{center}

\paragraph{3. 独立基底の場合の補加群との関係。}
$t_1,\ldots,t_n$ が $\mathfrak O$ の独立な $\mathfrak h$-基底であるなら、$\mathfrak K$ の基底 $(s)=(t)P$ も、$P$ の成分が $\mathfrak h$ に属し $P$ が単模である限り、独立な $\mathfrak h$-基底である。このとき $(S)=P^{-1}(T)$ から、$T$ と $S$ は同一の $\mathfrak h$-加群の基底となることが分かる。この加群を $\mathfrak O$ の補加群と呼ぶ。

2 の構造定理の最後の部分は、$\mathfrak d$ が $K$ における商
\[
\mathfrak o:\mathfrak c
\]
に $\mathsf{K}$ において等しいことを述べる。ここで $\mathfrak e$ は補加群である。実際、
\[
e^{(1)}=\mathsf{A}_1t_1+\cdots+\mathsf{A}_nt_n,
\]
とし、$\mathsf{A}_\lambda$ が $\mathfrak e$ の基底であるとする。このとき $\mathfrak d$ は、$\delta e^{(1)}$ が $\mathfrak O_{\mathfrak o}$ に属する元 $\delta$ 全体からなる。これは、すべての $i$ について $\delta \mathsf{A}_\lambda$ が $\mathfrak o$ に属することと同値である。したがって $\mathfrak d=\mathfrak o:\mathfrak e$ である。

\paragraph{特殊な場合。}
$e,z,\ldots,z^{n-1}$ を加群基底とする正則オーダーの場合には、
\[
\mathfrak d=\{f'(z)\}.
\]

\paragraph{備考。}
§ 5, 4 の備考から、$\mathfrak O$ の補加群 $\mathfrak E$ に関する別のよく知られた定義が得られる。\emph{$\mathfrak C$ は、任意の $t$、すなわち $\mathfrak O$ の任意の $o$ に対してトレース $\operatorname{Sp}(ct)$、すなわち $\operatorname{Sp}(co)$ が $\mathfrak h$ に属するような $\mathfrak K$ の元 $c$ 全体からなる。}$\mathfrak o$ と $K$ に関する $\mathfrak c$ についても同様である。

実際、$\mathfrak E$ の元 $c$ を、$\mathfrak K$ の $t$ に対する補基底を用いて $c=c_1T_1+\cdots+c_nT_n$ と表示し、各 $c_i$ は $\mathfrak h$ に属するとする。このとき $c_i=\operatorname{Sp}(ct_i)$ であるから $c_i$ は $\mathfrak h$ に属する。逆に、すべての $\operatorname{Sp}(ct_\lambda)$ が $\mathfrak h$ に属すれば $c_\lambda$ は $\mathfrak h$ に属し、したがって $c$ は $\mathfrak E$ に属する。

\paragraph{補足。}
\emph{オーダーに独立な $\mathfrak h$-加群基底が存在するなら、差分商 $\mathfrak A^{(i)}$ は $\mathfrak P^{(i)}$ 以外のすべての差分イデアルの交わりである。より正確には}
\[
\mathfrak A^{(i)}=
[\mathfrak O_{\mathfrak o^{(i)}},
\mathfrak P_{\mathfrak g}^{(1)},\ldots,
\mathfrak P_{\mathfrak g}^{(i-1)},
\mathfrak P_{\mathfrak g}^{(i+1)},\ldots,
\mathfrak P_{\mathfrak g}^{(n)}].
\]
実際、§ 2, 2 の強化された形により、$\mathfrak P$ は $\mathfrak P_{\mathsf{K}}$ の縮約イデアルでもある。$e,t_2,\ldots,t_n$ を $\mathfrak O$、したがって $\mathfrak O_{\mathfrak o}$ の加群基底とすれば、
\[
e,t_2-\alpha_2,\ldots,t_n-\alpha_n;
\]
一方、$t_2-\alpha_2,\ldots,t_n-\alpha_n$ は同時に $\mathfrak B_K$ の $K$-加群基底である。したがって § 2, 2（補足）の仮定が満たされる。ところで
\[
\mathfrak A_\Gamma=[\mathfrak P_\Gamma^{(2)},\ldots,\mathfrak P_\Gamma^{(n)}];
\qquad
\mathfrak A_\mathfrak g=[\mathfrak A_\Gamma,\mathfrak O_\mathfrak g]
=[\ldots,[\mathfrak P_\Gamma^{(\lambda)},\mathfrak O_\mathfrak g],\ldots]
=[\mathfrak P_\mathfrak g^{(2)},\ldots,\mathfrak P_\mathfrak g^{(n)}].
\]
しかし $\mathfrak A=[\mathfrak A_\mathfrak g,\mathfrak O_{\mathfrak o}]$ であるから、
\[
\mathfrak A=[\mathfrak O_{\mathfrak o},\mathfrak P_\mathfrak g^{(2)},\ldots,\mathfrak P_\mathfrak g^{(n)}].
\]

\paragraph{4. 基本方程式が存在する場合の、その微分商との関係。}
$n$ 個の不定元 $u_1,\ldots,u_n$ を添加してオーダー $\mathfrak O$ を拡大し、直積 $\mathfrak O_{\mathfrak U}$ を作る。ここで $\mathfrak U$ は多項式領域 $\mathfrak h[u_1,\ldots,u_n]$ を表し、$\mathfrak O_{\mathfrak U}=\mathfrak O[u_1,\ldots,u_n]$ である。このとき $\mathfrak O$ の任意のイデアル $\mathfrak C$ の拡大 $\mathfrak C_{\mathfrak U}$ は、$\mathfrak C$ の元を係数とする $u$ の多項式全体からなる。逆に $\mathfrak C$ は、$\mathfrak C_{\mathfrak U}$ のすべての係数からなる系として特徴づけられる。

いま基本方程式が存在するとする。すなわち、$\mathfrak h[u]$ の一つまたは複数の原始多項式 $H(u)$ による商環へ移るとき、$e,U,\ldots,U^{n-1}$ が $\mathfrak h$-加群基底を成すとする。そのときイデアル $\mathfrak C_{\mathfrak U}$ は、$\mathfrak C_{\mathfrak U}$ に単元（$H(u)$ およびその正負の冪）を乗じて生じる諸系へ移る。したがって $\mathfrak C$ を回収するには、まずこれらの単元を乗じて $\mathfrak C_{\mathfrak U}$（$u$ について整）に達し、そこから $\mathfrak C$ に達することができる。特に $\mathfrak O$ の共役差積 $\mathfrak D$ は $\mathfrak D_{\mathfrak U}$ へ移り、$\mathfrak D_{\mathfrak U}$ は $\mathfrak O_{\mathfrak U}$ の共役差積である。しかし $\mathfrak O_{\mathfrak U}$ には、原始多項式 $H(u)$ による商を取ることで、定義方程式 $G(U)=0$ が存在する。したがって $G'(U)$ は共役差積の基底多項式となり、すなわち共役差積は $G'(U)$ の係数から導かれる。（$(G'(U))$ が実際に $\mathfrak D_{\mathfrak U}$ だけを含み、それ以上を含まないことは、なお厳密に示す必要がある。これはおそらく $K$ における $\mathfrak o$ の整閉性のもとでのみ成り立つ。$H(u)$ の約元を掛けることによって、$\mathfrak O_{\mathfrak U}$ に属する別の元が $\mathfrak D_{\mathfrak U}$ として生じる可能性があるからである。整閉性（五公理）のもとでは、係数イデアル（内容）が乗法的であるという定理から従う。これにより、基本方程式を $p$ を法として考察する分岐理論が可能になる。）

\subsection*{\S{} 7. 主オーダーの $p^t$ を法とする分岐理論。}

数体を基礎に置き、$\mathfrak h$ をその整数環とする。あるいは、主オーダーを素イデアルで割った商環に関する相対体を考えてもよい。この場合 $p$ は、その商環の素イデアルの基底元を表す。不定元 $u_1,\ldots,u_n$ を添加したとする。このとき主オーダーも $p^t$ を法とする剰余類環も、$\mathfrak h$-加群基底 $e,U,\ldots,U^{n-1}$ をもつ。したがって、いずれの場合にも共役差積は $G'(u)$ により与えられる。\emph{主オーダーの共役差積は、$p^t$ を法として、$p^t$ を法とする剰余類環の共役差積へ移る。}（これは任意のオーダーについて成り立つか。）いま
\[
p=\mathfrak p_1^{\varrho_1}\cdots\mathfrak p_r^{\varrho_r},
\]
とする。このとき $p^t$ を法とする剰余類環 $\mathfrak R$ は直和であり、その各成分は $\mathfrak p_i^{t\varrho_i}$ を法とする剰余類環と同型である。すなわち
\[
\mathfrak R=\mathfrak R e_1+\cdots+\mathfrak R e_r;
\]
ここで $i\ne k$ なら $e_k$ は $\mathfrak p_i^{t\varrho_i}$ で割り切れる。したがって
\[
\mathfrak d[p^t]=\mathfrak d[p^t]e_1+\cdots+\mathfrak d[p^t]e_r
=\mathfrak d_1+\cdots+\mathfrak d_r
\]
かつ $t$ が十分大きいなら、$\mathfrak p_i$ は $\mathfrak d_i$ において $\mathfrak d[p^t]$、したがって $\mathfrak d$ におけるのと同じ冪で現れる。一方 § 4 により $\mathfrak d_i$ は $\mathfrak R_i$ の共役差積である。ゆえに $\mathfrak p^{\varrho t}$ を法とする剰余類環の共役差積だけを考えればよい。$\xi$ を、$\xi$ が整数であり（Zahlbericht, 定理 29 および 30）、
\[
\varphi(\xi)\equiv0(\mathfrak p),
\quad \not\equiv0(\mathfrak p^2),
\]
となるように選ぶ。ここで $\varphi(x)$ は次数 $f$ の $p$ を法とする素関数である。$t\geqq2$ のとき、この $\xi$ を $\xi^*=\xi e_1$ で置き換える。ここで $e_1$ は、$p^t$ を法とする剰余類環 $\mathfrak R$ の分解に現れる最初の冪等元である。このとき $\varphi(\xi^*)\equiv0(\mathfrak p)$、$\not\equiv0(\mathfrak p^2)$ に加えて、$\mathfrak p_1\equiv(p,\varphi(\xi^*))$ さえ成り立つ。$p^t$ を法とする整数に関する基底は、$\mu=0,1,\ldots,f-1$、$\nu=0,1,\ldots,t-1$ として $\xi^\mu\varphi(\xi)^\nu$ からなる。独立な $\mathfrak h$-基底（$p^t$ を法とする剰余類）は
\begin{equation}
\xi^{\mu_0}\varphi(\xi)^{\nu_0}
\quad\text{ただし}\quad
\mu_0=0,1,\ldots,f-1;
\quad \nu_0=0,1,\ldots,\varrho-1. \tag{1}
\end{equation}
で与えられる。実際、$(\varphi(\xi),\mathfrak p^{\varrho t})=\mathfrak p$ であり、剰余類環において $(p)=\mathfrak p^\varrho$ である。したがって主イデアルから元へ移れば、$\varphi(\xi)=pM(\xi)$ と書ける。ここで $M(\xi)$ は剰余類環の単元、すなわち $M(\xi)\not\equiv0(\mathfrak p)$ である。ゆえに $p^t=0$ を用いて、まず $M(\xi)$ において、ついで剰余類環において、$\xi^{\varrho t}$ を順次より低い冪で表すことができる。したがって（1）は実際に基底である。この基底に関する定義方程式は
\[
F(x)=\varphi(x)^\varrho+pM(x),
\]
である。ここで $M(x)$ は $p$ を法として $\varphi(x)$ で割り切れない（Ö. Ore, Math. Ann. 96, p. 348）。すなわち
\[
M(\xi)\not\equiv0(\mathfrak p).
\]
また、この基底は独立である。実際、
\[
a_0(\xi)+a_1(\xi)\varphi(\xi)+\cdots+a_{\varrho-1}(\xi)\varphi(\xi)^{\varrho-1}
\equiv0(\mathfrak p^{t\varrho})
\]
からまず $\equiv0(\mathfrak p)$、したがって $a_0(\xi)\equiv0(\mathfrak p)$、ゆえに
\[
a_0(x)\equiv0(p);
\]
$\mathfrak p$ を割り去ることを繰り返すと、順次
\[
a_0(x)\equiv0(p),\ldots,a_{\varrho-1}(x)\equiv0(p)
\]
を得る。さらに $p$ を割り去ると
\[
a_0(\alpha)\equiv0(p^2),\ldots,a_{\varrho-1}(\alpha)\equiv0(p^2)
\]
となり、以下同様にして最後には
\[
a_0(\alpha)\equiv0(p^t),\ldots,a_{\varrho-1}(\alpha)\equiv0(p^t).
\]
したがって、剰余類環の共役差積は多項式微分商 $F'(\xi)$ によって与えられ、問題はすべて Ore, p. 345, § 3 に帰着する。Ore の定理 10（Dedekind--Hensel 参照）から、各 $\mathfrak p$ について $p^t=p^n$ と取れば十分であることが従う。さらに Ore により補助数が得られる。

商環へ移ることにより、以上はすべて相対共役差積についても同時に成り立つ。したがって Ore, Math. Ann. 97, p. 594, § 4 は Ore, Math. Ann. 96 から直ちに従う。付け加えるべきなのは、上体の共役差積が下体の共役差積と相対共役差積との積として表されるという主張だけである（本論文の序論参照）。

\begin{center}
1949年10月25日受理。
\end{center}

%
\input{t0}
\input{t1}
\input{t2}
\input{t3}
\input{t4}
\input{t5}
\input{t6}
\input{t7}

\section*{論文40・41の補注}
\addcontentsline{toc}{section}{論文40・41の補注}
\noindent 以下は論文40・41の原著に付された注記である。

\subsection*{論文40}
\begin{enumerate}
\item E. Noether, 「超複素量と表現論」, Math. Z. 30 (1929), 641--692；van der Waerden, \emph{Moderne Algebra II} 参照。
\item Brauer--Noether「既約表現の最小分解体について」, Sitzungsberichte Preuss. Akad. Wiss. 1927, 221--228。p. 222 の注は優先関係の報告である。主要定理は独立に、ほぼ同時に得られた。R. Brauer, Math. Z. 30 (1929), 79--107 も参照。A. A. Albert は後に独立に再発見した。
\item R. Brauer, J. reine angew. Math. 166 (1932), 241--252；K. Shoda, Math. Ann. 107 (1932), 252--258。J. Levitzki もこれらの結果を展開していた。
\item 第一講義は G. Köthe が、第二講義は M. Deuring が書き起こした。
\item この方法は G. Köthe, Math. Ann. 105 (1931), 15--39, とくに §5 に本質的に再録されている。注10も参照。この共通部分法は現在、両側イデアルが不変加群に属するという、van der Waerden が初めて指摘したほとんど自明な事実（§4, 1）で置き換えられている。これによりすべてをより簡単な直和分解で基礎づけられるが、この方法は無限階数の場合および整の場合には成り立たない。
\item プラハ講演, Jahresbericht DMV 39 (1930), p. 17（斜めに付された頁番号）を参照。
\item 接合積の理論は第二講義で展開され、H. Hasse, Trans. Amer. Math. Soc. 34 (1932), 171--214 に小変更を伴って再録された。講義により近い形の記述は M. Deuring の報告として予告された。
\item 非可換群の場合は一般化された環である；H. Fitting, Math. Ann. 107, 514--542 参照。
\item ここで「積環」とは、共通の上環をもつ二つの環から生成される環を意味するだけであり、その積は直積でなくてもよい。指定された共通部分に対して積環が常に存在するとは限らないことを、次の例が示す。$\mathfrak o$ を整数環とし、$\mathfrak R=\mathfrak o[x]$、$\mathfrak S=\mathfrak o[y]$、$2x-1=0$、$2y-1=0$ と置く。$x$ と $y$ の間に関係を課さなければ、$\mathfrak R$ と $\mathfrak S$ の共通部分は $\mathfrak o$ である。しかし両者が共通の上環に入るなら、$(2x-1)y-x(2y-1)=0$、従って $x=y$ となる。ゆえに共通部分が $\mathfrak o$ である積環は存在しない。
\item 当該段落の定理は、$P$ 上の $M$ の階数が無限でも、単純な整列法により成立する；G. Köthe, Göttinger Nachrichten 1931, 195--207 参照。
\item 「単純系」は通常どおり両側単純を意味する。
\item $A$ が $S$ の付随自己同型斜体である場合について、表現論の注20参照。
\item 逆同型の体または系は、対応するギリシア文字とラテン文字で表す。
\item 以下の同型は、断らない限り、$P$ の恒等写像の延長である。
\item $A$ が $P$ 上無限階数のとき、これはもはや成り立たない；Köthe 参照。
\item R. Brauer--E. Noether, 前掲。
\end{enumerate}
\subsection*{論文41}
\begin{enumerate}
\item $K^*$ は $K$ から零を除いたものを表す。
\item 関係式 (2) から (5) は、群の拡大における定義関係式にほかならない。このような拡大は、乗法群としてのみ見た $K^*$ の完全自己同型群に、$\mathfrak G$ の準同型像となる部分群が存在すれば成立する。接合積で特殊なのは、この準同型が同型となり、さらに問題の部分群として、乗法的かつ加法的な自己同型の全体、すなわち $K$ のガロア群を取ることで、群拡大とともに環拡大も得られる点である。
\item ここで与えた証明は、とくに $k$ が有限個の元しかもたない場合にも有効である。この場合 Hilbert の証明は成り立たないが、Speiser の証明はなお有効である。
\item I. Schur, Math. Z. 5 (1919), 7--10。交叉表現加群による証明については Deuring 参照。
\item 一般に、$\mathfrak I$ のうち $\mathfrak G$ を許す部分群を主類として取れば、$\{\ldots u_S\overline{\mathfrak I}\ldots\}$ は定義される。このとき、その部分群のイデアルから生じる変換量が、主イデアルの変換量に代わって現れる。
\item $\bar c_S$ は $c_S$ のイデアル類を表す。
\item 導入の末尾を参照。
\item より正確には $a_{S,T}=d_S^Td_T/d_{ST}$ が成り立つ。ただし主種定理の証明では、主イデアルについての弱められた主張だけを用いる。
\item Hasse（注6）, 14 を参照。より簡単な証明は次のとおりである。$Z$ を $\mathfrak P$ に対する分解体とすると、$k_{\mathfrak p}$ はそれと同型な $Z$ を含む。このとき $A_Z\sim(a_{\mathfrak Z},K/Z,\mathfrak Z)$ である。実際、$A_Z$ は $A$ のうち $Z$ のすべての元と可換な元の全体に相似である（van der Waerden, Bd. II, p. 210、または Math. Zeitschr. 掲載予定の非可換代数に関する著者の論文を参照）。$A_Z$ から $A_{\mathfrak p}$ への拡大により上の結果が得られる。共役な $\mathfrak P$ のどれを選ぶかは問題でない。$\mathfrak G=\sum\mathfrak ZT$ のとき、得られる代数は同型で、$u_T$ による変換で互いに移り合う。
\item 注4に引用された H. Hasse の仕事 §3 を参照。
\end{enumerate}

\end{document}

%% END NOETHER SOURCE
%% BEGIN NOETHER SOURCE "ja/t0.tex"
\begingroup
\IfFontExistsTF{FreeSerif}{\fontspec{FreeSerif}}{\fontspec{Latin Modern Roman}}
\let\rOneOneOnemathbf\mathbf
\renewcommand{\mathbf}[1]{\ensuremath{\rOneOneOnemathbf{#1}}}
\let\rOneOneOnemathfrak\mathfrak
\renewcommand{\mathfrak}[1]{\ensuremath{\rOneOneOnemathfrak{#1}}}
\let\rOneOneOnemathsf\mathsf
\renewcommand{\mathsf}[1]{\ensuremath{\rOneOneOnemathsf{#1}}}
\let\rOneOneOnemathbb\mathbb
\renewcommand{\mathbb}[1]{\ensuremath{\rOneOneOnemathbb{#1}}}
\let\rOneOneOnemathcal\mathcal
\renewcommand{\mathcal}[1]{\ensuremath{\rOneOneOnemathcal{#1}}}
\providecommand{\mathscr}[1]{\mathcal{#1}}
\let\rOneOneOnemathscr\mathscr
\renewcommand{\mathscr}[1]{\ensuremath{\rOneOneOnemathscr{#1}}}
\providecommand{\texorpdfstring}[2]{#1}
\providecommand{\tightlist}{\setlength{\itemsep}{0pt}\setlength{\parskip}{0pt}}
\clearpage
\begin{center}
\vspace*{3.5em}
{\LARGE\bfseries 超複素量の代数\par}
\vspace{3.5em}
{\large E. Noether 教授講義\par}
{\large 1929/30 年冬学期\par}
\vspace{1.5em}
{\large M. Deuring 教授編\par}
\vspace{4em}
{\Large\bfseries 目次\par}
\end{center}
\vspace{1em}
\begingroup
\small
\setlength{\parindent}{0pt}
\newcommand{\tocline}[2]{\noindent #1\dotfill #2\par}
\newcommand{\tocsec}[3]{\noindent\hspace*{1.2em}\S\,#1.\quad #2\dotfill #3\par}
\tocline{\textbf{序論}}{3}
\medskip
\tocline{\textbf{第 I 章：表現}}{3}
\tocsec{1}{直接表現の定義 --- 相反表現の定義}{3}
\tocsec{2}{表現類}{4}
\tocsec{3}{表現加群}{4}
\tocsec{4}{表現加群と表現との関係}{5}
\medskip
\tocline{\textbf{第 II 章：可換体のガロア理論}}{7}
\tocsec{5}{超複素系の係数拡大}{7}
\tocsec{6}{可換系の既約表現}{8}
\tocsec{7}{一つの体の同型写像}{8}
\tocsec{8}{分解体}{9}
\tocsec{9}{$\mathsf{Z}_1$ から各 $\mathsf{Z}_i$ への同型写像}{9}
\tocsec{10}{ガロア群}{9}
\tocsec{11}{ガロア理論の基本定理}{9}
\tocsec{12}{成分 $e_i$ の形式的意味}{11}
\tocsec{13}{一般延長定理}{12}
\medskip
\tocline{\textbf{第 III 章：アーベル群}}{13}
\tocsec{14}{群環}{13}
\tocsec{15}{アーベル群の群環}{14}
\tocsec{16}{指標関係式}{15}
\tocsec{17}{アーベル群のガロア理論}{15}
\medskip
\tocline{\textbf{第 IV 章：両側単純環}}{18}
\tocsec{18}{補題}{18}
\tocsec{19}{係数体の非可換拡大体における両側単純超複素系の表現}{19}
\tocsec{20}{非可換体}{22}
\tocsec{21}{非可換体のガロア理論}{26}
\clearpage
\tocline{\textbf{第 V 章：因子系}}{29}
\tocsec{22}{中心を指定した体の群}{29}
\tocsec{23}{因子系}{30}
\tocsec{24}{因子系の乗法}{33}
\tocsec{25}{ガロア的な極大可換部分体をもつ $\mathfrak K_r$ の標準表現}{39}
\tocsec{26}{交叉表現の乗法}{43}
\medskip
\tocline{\textbf{第 VI 章：接合積の理論}}{44}
\tocsec{27}{$\mathfrak K_r$ の接合積としての表現}{44}
\tocsec{28}{因子系の積定理}{47}
\tocsec{29}{最小の場合の主種定理}{50}
\tocsec{30}{巡回分解体への特殊化}{50}
\tocsec{31}{巡回特殊例の応用}{52}
\endgroup
\clearpage

\section*{序論}\label{einleitung}

本講義は環の一般表現論を扱う。この理論は、一方では有限群の表現論から、他方では超複素数系の代数的・算術的理論から発展したものである。このように表現論を一般化して捉える立場には、従来の立場（Frobenius、Schur）に比べて次の利点がある。

群そのものではなく群環、さらに一般には超複素系または任意の環を考察するため、環とイデアルの理論を利用でき、多くの見通しの悪い計算から表現論を解放できる。一般の環論・イデアル論を用いることは、理論を簡明にするだけではない。たとえば群とその部分群との関係という問題を前進させ、さらに新たな応用領域を開く。一般表現論によって、ガロア理論に新しくきわめて簡明な基礎づけを与えることができ、また、おそらくそれ以上に重要なこととして、非可換体のガロア理論を構成できる。

群論・環論・イデアル論の基本概念、および表現論の一般定理は、Noether の論文 ``Hyperkomplexe Größen und Darstellungstheorie'', \emph{Mathematische Zeitschrift} 30, 641 にある。以下 M. Z. と略記するこの論文は、本講義録を補うものとみなすべきである。M. Z. と同じ内容は、Noether の講義「超複素量と群指標」（1927/28 年冬学期）の講義録にも見いだされる。

ここでは、形式的な点でのみ新しい相反表現加群および相反表現を導入するために、表現と表現加群の概念を簡潔に説明する。

\section*{第 I 章　表現}\label{kapitel-i.-darstellungen}

$\mathfrak{o}$ を表現される環、$\mathsf T$ を表現が行われる別の環とする。

\subsection*{§ 1. 直接表現の定義}\label{definition-der-direkten-darstellung}

$\mathfrak{o}$ の $\mathsf T$ における $n$ 次の\emph{直接表現}とは、$\mathsf T$ の元を成分とする $n$ 次行列環 $\mathsf T_n$ の部分環 $\mathfrak D$ であって、$\mathfrak{o}$ がそこへ環準同型に写されるものをいう。記号では

\[
\mathfrak{o}\to\mathfrak D\subseteqq\mathsf T_n.
\]

\subsection*{相反表現の定義}\label{definition-der-reziproken-darstellung}

$\mathfrak{o}$ の $\mathsf T^*$ における $n$ 次の\emph{相反表現}とは、$\mathsf T^*$ の元を成分とする $n$ 次行列環 $\mathsf T_n^*$ の部分環 $\mathfrak D^*$ であって、$\mathfrak{o}$ がそこへ相反環準同型に写されるものをいう。相反環準同型とは次の意味である。$\mathfrak{o}$ の元 $c,d$ に $\mathfrak D^*$ の元 $c^*,d^*$ が対応するとき、$\mathfrak{o}$ の元 $cd$ には $\mathfrak D^*$ の元 $d^*\cdot c^*$ が、$\mathfrak{o}$ の元 $c+d$ には $c^*+d^*$ が対応する。記号では

\[
\mathfrak{o}\nrightarrow\mathfrak D^*\subseteqq T_n^*.
\]

相反準同型の例を一つ挙げよう。可換体上の $n$ 行行列にその転置行列を対応させる写像は、相反環同型である。

一般に、任意の環 $\mathsf T$ に対して、それと相反同型な環 $\mathsf T^*$ を次のように構成できる。

$\mathsf T$ の各元 $c$ に記号 $c^*$ を対応させ、これらの記号の加法を $c^*+d^*=(c+d)^*$、乗法を

\[
c^*\cdot d^*=(dc)^*
\]

によって定める。記号 $c^*$ 全体は、$\mathsf T$ と相反同型な環 $\mathsf T^*$ をなす。

\subsection*{§ 2. 表現類}\label{darstellungsklassen}

$\mathfrak{o}$ の $\mathsf T$ における直接表現 $\mathfrak D$ は、正則行列 $P$ で各元を変換することによって同型な表現 $\mathfrak D'$ に移せる。すなわち、$\mathfrak{o}$ の元 $c$ に $C\in\mathfrak D$ を対応させる代わりに、$P^{-1}CP\in P^{-1}\mathfrak D P$ を対応させる。

このような変換によって互いに移り合う二つの表現を\emph{同値}という。

一つの表現と同値な表現全体を一つの\emph{表現類}（d. D. K.）という。

相反表現の同値性も同様に定義する。一方が他方から正則行列による変換で得られるとき、二つの相反表現は同値であるという。一つの固定した相反表現と同値なもの全体を一つの\emph{相反表現類}という。

\subsection*{\texorpdfstring{$\S$ 3. 表現加群}{\textbackslash S 3. Darstellungsmoduln}}\label{s-3.-darstellungsmoduln}

\emph{直接表現加群}（d. D. M.）の定義。

$\mathfrak{o}$ の $\mathsf T$ に関する直接表現加群とは、$\mathfrak{o},\mathsf T$-双加群 $\mathfrak M$、すなわち加法的に書かれたアーベル群であって、$\mathfrak{o}$ に関して左加群、$\mathsf T$ に関して右加群であり、混合結合則
$c\in\mathfrak{o}$、$m\in\mathfrak M$、$\tau\in\mathsf T$ に対して $(cm)\tau=c(m\tau)$
を満たすもののうち、有限個の一元 $\mathsf T$-加群の直和として書け、かつ $x_i\tau=0$ ならば $\tau=0$ となるものをいう：

\[
\mathfrak M=x_1\cdot\mathsf T+x_2\cdot\mathsf T+\cdots+x_n\cdot\mathsf T.
\]

\emph{相反表現加群}（r. D. M.）の定義。

$\mathfrak{o}$ の $\mathsf T^*$ に関する相反表現加群とは、$\mathfrak{o}$-左加群かつ $\mathsf T^*$-左加群である $\mathfrak M^*$ であって、可換則（$c\in\mathfrak{o}$、$m\in\mathfrak M^*$、$\tau\in\mathsf T^*$ ならば $c(\tau^*m^*)=\tau^*(cm)$）を満たし、有限個の単純 $\mathsf T^*$-加群の直和として書け、かつ $\tau^*\cdot x_i^*=0$ ならば $\tau^*=0$ となるものをいう：

\[
\mathfrak M^*=\mathsf T^*\cdot x_1+\cdots+\mathsf T^*\cdot x_n.
\]

\subsection*{§ 4. 表現加群と表現との関係}\label{der-zusammenhang-zwischen-darstellungsmoduln-und-darstellungen}

\textbf{1.} \emph{各表現加群は一意的に一つの表現類を与える（d. D. M. は d. D. K. を、r. D. M. は r. D. K. を与える）。}

その仕方は次のとおりである。

d. D. M. $\mathfrak M$ に $\mathfrak{o}$ の元 $c$ を掛けることは、$\mathfrak M$ からそれ自身への準同型を意味する。結合則 $(cm)\tau=c(m\tau)$ により、この準同型は $\mathsf T$ の元を作用素として許す。したがってこれは $\mathfrak M$ からそれ自身への\emph{線型変換}である。すなわち、対応 $m\mapsto cm=m'$ は次の性質をもつ：

\[
\text{$y_i\mapsto y_i'(=cy_i)$ ならば、同時に }
\sum y_i\tau_i\mapsto \sum y_i'\tau_i\left(=c\sum y_i\tau_i\right).
\]

二つの線型変換

\[
\mathfrak T_1:\ m\mapsto m'(=c_1m)
\quad\text{および}\quad
\mathfrak T_2:\ m\mapsto m''(=c_2m)
\]

の和 $\mathfrak T_1+\mathfrak T_2$ を変換 $m\mapsto m'+m''(=(c_1+c_2)m)$ として、その積 $\mathfrak T_1\cdot\mathfrak T_2$ を合成変換 $m\mapsto(m'')'=c_1c_2m$ として定める。このとき、$\mathfrak{o}$ の元が生成する $\mathfrak M$ の相異なる線型変換全体は、$\mathfrak{o}$ の環準同型像 $\overline{\mathfrak D}$ をなす。

次に $\mathfrak M$ の任意の基底 $x_1,\ldots,x_n$ を取る。$c$ が生成する $\mathfrak M$ の線型変換は、基底元 $x_i$ がどのように変換されるか、すなわち方程式系

\[
x'_j=cx_j=\sum x_i\gamma_{ij}
\]

の行列 $(\gamma_{ij})=C$ がわかれば、明らかに一意に定まる。他の元 $x=\sum x_j\tau_j$ に対しては単に $x'=\sum x'_j\tau_j$ だからである。これらの行列 $C$ は相異なる線型変換と一意に対応し（互いに他方を定める）、容易にわかるように $\overline{\mathfrak D}$ と同型な環をなす。したがってこれは $\mathfrak{o}$ の環準同型像、すなわち $\mathfrak{o}$ の $n$ 次表現である。

$\overline{\mathfrak D}$ とこの表現との同型は、次の式に表れている。$c\in\mathfrak{o}$ が生成する線型変換に行列 $C$ が対応して

\[
c(x_1,\ldots,x_n)=(x_1,\ldots,x_n)C
\]

が成り立ち、同様に $d\in\mathfrak{o}$ と $D\in\mathsf T_n$ が対応して

\[
d(x_1,\ldots,x_n)=(x_1,\ldots,x_n)D
\]

が成り立つならば、明らかに

\[
\begin{aligned}
(c+d)(x_1,\ldots,x_n)
&=(cx_1+dx_1,\ldots,cx_n+dx_n)\\
&=(x_1,\ldots,x_n)(C+D)
\end{aligned}
\]

および

\[
\begin{aligned}
cd(x_1,\ldots,x_n)
&=c(dx_1,\ldots,dx_n)=c((x_1,\ldots,x_n)D)\\
&=(cx_1,\ldots,cx_n)D=(x_1,\ldots,x_n)CD
\end{aligned}
\]

を得る。このように、$\mathfrak{o}$ が生成する $\mathfrak M$ の自己線型変換を、$\mathfrak M$ の特定の基底 $x_\nu$ によって実現すると、$\mathfrak{o}$ の $\mathsf T$ における $n$ 次表現が得られる。

$(y_1,\ldots,y_n)$ を $\mathfrak M$ の別の基底とする。それは正則行列を掛けることによって $x_1,\ldots,x_n$ から得られる：

\[(y_1,\ldots,y_n)=(x_1,\ldots,x_n)P.\]

$\mathfrak M$ のある線型変換が基底 $x_\nu$ に関して行列 $C$ で実現されるならば、基底 $y_\nu$ に関しては $P^{-1}CP$ で実現される。実際、

\[c(x_1,\ldots,x_n)=(x_1,\ldots,x_n)C\]

ならば

\[
\begin{aligned}
c(y_1,\ldots,y_n)
&=c(x_1,\ldots,x_n)P=((x,\ldots,x_n)C)P\\
&=(x_1,\ldots,x_n)CP=((y_p,\ldots,y_n)P^{-1})CP\\
&=(y_1,\ldots,y_n)P^{-1}CP
\end{aligned}
\]

が成り立つ。ここから直ちに次が従う。

\emph{$\mathfrak M$ のすべての基底を用いて、$\mathfrak{o}$ が $\mathfrak M$ 上に誘導する線型変換を実現すれば、$\mathfrak M$ は一つの表現類に属するすべての表現を与える。}

\begin{enumerate}
\def\labelenumi{\arabic{enumi}.}
\setcounter{enumi}{1}
\tightlist
\item
  \emph{各表現類は、1. で述べた仕方で一つの表現加群によって生成される。}（再び直接表現に限って論じる。）
\end{enumerate}

\emph{証明。} 類の中から一つの表現を固定し、それが準同型対応 $c\mapsto C$ で与えられているとする。表現の次数を $n$ とし、$n$ 個の不定元 $x_1,\ldots,x_n$ によって $\mathsf T$-右加群

\[\mathfrak{M} = x_1 \mathsf{T} + \dots + x_n \mathsf{T}\]

を作る（通常の計算法則 $\sum x_i\tau_i+\sum x_i\overline{\tau}_i=\sum x_i(\tau_i+\overline{\tau}_i)$、$\sum x_i\tau_i=\sum x_i\overline{\tau}_i$ が成り立つのは、すべての $i$ について $\tau_i=\overline{\tau}_i$ のとき、かつそのときに限ること、ならびに $(\sum x_i\tau_i)\varrho=\sum x_i\tau_i\varrho$ を用いる）。そして、$(x_1',\ldots,x_n')=(x_1,\ldots,x_n)C$ であり、ここで $C$ は $\mathfrak{o}$ の元 $c$ に表現によって対応する行列であるとき、$c\sum x_i\tau_i=\sum x_i'\tau_i$ と定めて、これを $\mathfrak{o}$-左加群にする。

この規定により $\mathfrak M$ は $\mathfrak{o},\mathsf T$-双加群になる。第一に、準同型関係

\[
c+d\mapsto C+D \quad \text{および} \quad cd\mapsto CD
\]

から

\[
(c+d)x_i=cx_i+dx_i
\]

および

\[
cd\cdot x_i=c\cdot dx_i,\quad\text{さらに}\quad
c\left(\sum x_i\tau_i+\sum x_i\overline{\tau}_i\right)
=c\sum x_i\tau_i+c\sum x_i\overline{\tau}_i
\]

が従う。これらは左加群の性質である。第二に、明らかに

\[c(\sum x_i \, \tau_i \, \varrho) = (c \, \sum x_i \, \tau_i) \cdot \varrho\]

が成り立つ。これは混合結合則である。

したがって $\mathfrak M$ は表現加群であり、その線型変換の実現に基底 $x_1,\ldots,x_n$ を用いれば、使用した表現が 1. で述べた仕方によって生成される。他の基底は、その類に属する残りの表現を与える。

相反表現および相反表現加群についてもすべて同様である。その結果、加群

\[\mathfrak{M}^* = \mathsf{T}^* x_1^* + \dots + \mathsf{T}^* x_n^*\]

は、$\mathfrak{o}$ の元 $c$ に、等式

\[
c\begin{pmatrix}x_1\\ \vdots\\ x_n\end{pmatrix}
=\begin{pmatrix}cx_1\\ \vdots\\ cx_n\end{pmatrix}
=C\begin{pmatrix}x_1\\ \vdots\\ x_n\end{pmatrix}
\]

に現れる行列 $C$ を対応させることによって、$\mathfrak{o}$ の $\mathsf T^*$ における相反表現を生成する。異なる基底は、ここでも同じ一つの類に属する異なる表現を与える。

%% END NOETHER SOURCE
%% BEGIN NOETHER SOURCE "ja/t1.tex"
\section*{第 II 章　可換体のガロア理論}\label{kapitel-ii.-galoissche-theorie-kommutativer-koerper}

ここまでに展開した表現論だけでガロア理論を基礎づけることができる。この基礎づけには、特別な生成元系や不定元の多項式に関する定理へ立ち戻る必要がないという、通常の基礎づけに対する利点がある。\srcfn{1}{以下については M. Z. §§ 15--20 を参照。M. Z. § 21 は、後のガロア理論に関する節の最初の諸定理を含む。超複素系の定義については M. Z. § 6 を参照。}

\subsection*{§ 5. 超複素系の係数拡大}\label{koeffizientenerweiterung-hyperkomplexer-systeme}

ほとんど自明なので証明を省く、次の一般定理を定式化する。$\mathfrak{o}=c_1\mathsf P+\cdots+c_n\mathsf P$ を可換体 $\mathsf P$ 上の超複素系とする。$\Omega$ を、その中心が $\mathsf P$ を含む $\mathsf P$ の拡大環とする。\emph{拡大系} $\mathfrak{o}_{\Omega}$ を、すべての形式和 $\sum c_i\omega_i$（$\omega_i\in\Omega$）の集合として定義する。まず規定

\[
\begin{array}{l}
\sum c_i\omega_i=\sum c_i\omega_i' \quad\text{であるのは、すべての $i$ について $\omega_i=\omega_i'$ のとき、かつそのときに限る};\\
\sum c_i\omega_i+\sum c_i\omega_i'=\sum c_i(\omega_i+\omega_i');\\
\omega\sum c_i\omega_i=\sum c_i\omega\omega_i;\quad
\sum c_i\omega_i\cdot\omega=\sum c_i\omega_i\omega;\\
\sum c_i\omega_i=\sum \omega_i c_i
\end{array}
\]

によってこれを $\Omega$ と可換的に結合した $\Omega$-加群とし、さらに規定

\[
\sum c_i\omega_i\cdot \sum c_i\omega_i
=\sum_i\sum_j c_ic_j\omega_i\omega_j
\]

によって $\mathfrak{o}$ の拡大環とする。ここで $c_ic_j$ は $\mathfrak{o}$ の元としてすでに定義されている。$\mathfrak{o}_{\Omega}$ は $\Omega$ 上階数 $n$ であり、

\[
\mathfrak{o}_{\Omega}=c_1\Omega+\cdots+c_n\Omega
\]

と書ける。

\emph{$\mathfrak{o}_{\Omega}$ は基底 $c_1,\ldots,c_n$ に依存せず定義される。$\Omega$ が可換体ならば、$\mathfrak{o}_{\Omega}$ は $\Omega$ 上の超複素系である。}

\subsection*{§ 6. 可換系の既約表現}\label{die-irreduziblen-darstellungen-kommutativer-systeme}

$\mathfrak Z=z_1\mathsf P+\cdots+z_n\mathsf P$ を $\mathsf P$ 上の可換系とする。$\mathsf P$ の代数閉拡大体 $\Omega$ における $\mathfrak Z$ の既約表現をすべて求めよう。$\mathfrak Z$ の $\Omega$ における各表現から $\mathfrak Z_\Omega=z_1\Omega+\cdots+z_n\Omega$ の表現が得られる。実際、すでに $\mathsf P$ に属する基底元の表現がわかれば表現全体を定めるのに十分である。逆に、$\mathfrak Z_\Omega$ の各表現は $\mathfrak Z$ の表現を含む。したがって、$\mathfrak Z_\Omega$ の表現を求めてもよい。M. Z. § 20 によれば、それらは $\mathfrak Z_\Omega$ の正則表現、すなわち $\mathfrak Z_\Omega$ 自身を表現加群として用いるときに生じる表現に含まれる。さらに、$\mathfrak C$ を $\mathfrak Z_\Omega$ の根基とすれば、表現加群を $\mathfrak Z_\Omega/\mathfrak C$ に限ってもよい（M. Z. § 19 参照）。

ところが $\mathfrak Z_\Omega/\mathfrak C$ は完全可約系、したがって $\Omega$ 上有限次数の体の直和である。$\Omega$ は代数閉であるから、これらの直和因子は $\Omega$ と同型である：

\[
\mathfrak Z_\Omega/\mathfrak C=e_1\Omega+\cdots+e_t\Omega.
\]

$e_\nu$ は単位元の成分である。したがって次が成り立つ。

\emph{$\mathfrak Z_\Omega$ は $\Omega$ においてちょうど $t$ 個の相異なる既約表現類をもち、それらはすべて一次である。ここで $t$ は $\mathfrak Z_\Omega/\mathfrak C$ の階数であり、したがって高々 $n$ である。}

これらの表現類はいずれも一つの表現しか含まない。可換体における一次表現を変換しても表現は変わらないからである：$\pi^{-1}\cdot\alpha\cdot\pi=\alpha$。

したがって、$\mathfrak Z_\Omega$ から $\Omega$ の部分環への相異なる準同型写像 $\Theta_1,\ldots,\Theta_t$ がちょうど $t$ 個存在する。

\subsection*{§ 7. 一つの体の同型写像}\label{die-isomorphismen-eines-koerpers}

ここで $\mathfrak Z$ を体とする。$\mathfrak Z_\Omega$ から $\Omega$ の部分環への各準同型 $\Theta_i$ は、$\mathfrak Z$ を $\Omega$ の部分体 $Z_i$ へ同型に写さなければならない。実際、$Z_i=0$ であるか、または $Z_i$ は $\mathfrak Z$ の同型像である。前者はここでは起こらない。$\mathfrak Z$ の単位元には $\Omega$ の単位元が対応するからである。ゆえに次を得る。

\emph{$\mathfrak Z$ が $\mathsf P$ の $n$ 次有限体拡大ならば、$\Omega$ の中に $\mathfrak Z$ と同型な相異なる体 $Z_i$ が、$\mathfrak Z_\Omega/\mathfrak C$ の階数 $t$ と同じ個数だけ存在する。したがって $Z_i$ の個数は高々 $n$ である。}

$Z_i$ の個数が $n$ に等しいとき $\mathfrak Z$ は $\mathsf P$ 上\emph{第一種}であるといい、その個数が $n$ より小さいとき\emph{第二種}であるという。

\emph{$\mathfrak Z$ が $\mathsf P$ 上第一種であるのは、$\mathfrak Z_\Omega$ が完全可約であるとき、かつそのときに限る。}

ここに、このガロア理論の基礎づけが通常のものより優れている点が見える。通常のガロア理論では、一つの $Z_i$ から他の $Z_i$ への同型写像を考えるが、ここではこの非対称性が避けられる。また、一つの体がその次数を超える数の共役体をもたないことを示すため、通常は原始元を用いる。ここでは、この事実は $\mathfrak Z_\Omega$ のすべての既約表現が正則表現に含まれることからただちに従う。

\subsection*{§ 8. 分解体}\label{der-zerfaellungskoerper}

$Z_i$ は $\mathfrak Z$ の $\Omega$ における既約表現である。したがって、$\mathsf P$ と $\Omega$ の中間体 $\Gamma$、$\mathsf P\subseteq\Gamma\subseteq\Omega$ がすべての $Z_i$ を含むならば、すでに $\mathfrak Z_\Gamma/\mathfrak C$ は $\Gamma$ 上一次の $t$ 個の体に分解しなければならない。逆に、$\mathfrak Z_\Gamma/\mathfrak C$ が $t$ 個の体へ分解するためには、それらは $\mathfrak Z_\Gamma$ の一次既約表現となるので、$Z_i$ が $\Gamma$ に含まれなければならない。$\mathfrak Z_\Gamma/\mathfrak C$ が $t$ 個の体の直和となる性質をもつ体 $\Gamma$ を $\mathfrak Z$ の分解体と呼べば、次を得る。

\emph{$\mathfrak Z$ と同型な体 $Z_i$ のガロア体、すなわちそれらの合成体は、$\mathfrak Z$ の最小分解体である。}

\subsection*{§ 9. $Z_1$ から各 $Z_i$ への同型写像}\label{die-isomorphismen-von-z1-auf-die-zi}

$\Theta_i$ は、$\mathfrak Z$ に制限すれば、$\mathfrak Z$ から $Z_i$ への同型写像である。したがって $S_i=\Theta_i\Theta_1^{-1}$ は $Z_1$ から $Z_i$ への同型写像である。$\Theta_i$ は互いに異なるから、$S_i$ も互いに異なる。$S_i$ は、通常のガロア理論で考える、一つの体 $Z_1$ からその共役体への同型写像である。他の $\Theta_i$ が存在しないので、他の同型写像も存在しない。

\subsection*{§ 10. ガロア群}\label{die-galoissche-gruppe}

$\mathfrak Z$ がガロア体、すなわち $Z_i$ が元の集合として一致するならば、$S_i$ は $\mathsf P$ を固定する $Z$ の自己同型すべてからなる群をなす。実際、各積 $S_iS_j$ は $Z$ の自己同型であり、$S_i$ 以外にはないから、ある $S_k$ に等しい。したがって $S_i$ は、$\mathsf P$ を元ごとに固定する $Z$ の自己同型すべてからなる群、すなわち\emph{$\mathsf P$ に関する $Z$ のガロア群}をなす。

\subsection*{§ 11. ガロア理論の基本定理}\label{der-hauptsatz-der-galoisschen-theorie}

$\mathfrak Z$ を $\mathsf P$ の第一種ガロア拡大とする。ガロア理論の基本定理を証明する。

\emph{$\mathsf P$ と $\mathfrak Z$ の中間体 $\Gamma$ と、$\mathsf P$ に関する $Z$ のガロア群 $\mathfrak G$ の部分群 $\mathfrak H$ との間には、次の性質をもつ一対一対応がある。体 $\Gamma$ に対応する群 $\mathfrak H$ は $\Gamma$ を元ごとに固定する $Z$ の自己同型すべてからなり、$\Gamma$ は $\mathfrak H$ に属するすべての自己同型によって固定される $Z$ の元すべてからなる。}

簡潔には次のように述べられる。

\begin{enumerate}
\item \emph{$\mathfrak H$ が $\Gamma$ の不変群ならば、$\Gamma$ は $\mathfrak H$ の不変体である。}
\item \emph{$\Gamma$ が $\mathfrak H$ の不変体ならば、$\mathfrak H$ は $\Gamma$ の不変群である。}
\end{enumerate}

\paragraph{1. は次の補題へ帰着される。}\label{laesst-sich-zurueckfuehren-auf-den-hilfssatz}

\emph{$\mathsf P \subseteq \Sigma \subseteq \mathsf T \subseteq \mathsf Z$ であり、$\mathsf T$ が $\Sigma$ 上次数 $t$ であるとする。このとき、$\Sigma$ から一つの（必ず $\mathsf Z$ に含まれる）共役体への各同型写像は、$\mathsf T$ から共役体への相異なる同型写像へ $t$ 通りに延長できる。}

$\mathfrak H$ が $\Sigma$ の不変群であり、$\mathsf T$ が $\mathfrak H$ の不変体であると仮定すれば、次のように論じられる。$\mathfrak H$ の元は、$\Sigma$ の恒等同型を $\mathsf Z$ 全体へ延長したものである。これらの延長を $\mathsf T$ にだけ適用すると、補題により、それらは $t$ 個の類に分かれる。同じ類の同型写像は $\mathsf T$ 上で一致し、異なる類の同型写像は $\mathsf T$ 上で異なる。しかし $\mathsf T$ は $\mathfrak H$ の不変体であるから、$t=1$、$\mathsf T=\Sigma$ でなければならない。証明された。

補題を証明するため、$\mathsf P$ と $\mathsf Z$ の中間体 $\mathfrak S,\mathfrak T$ であって、$\Theta_1$ によりそれぞれ $\Sigma,\mathsf T$ に対応するものへ戻る。示すべきことは次である。

\emph{$\mathfrak S$ から一つの $\Sigma_i$ への各同型写像 $H_i$ は、$\mathfrak T$ から $\mathsf Z$ の部分体 $\mathsf T_i$ への同型写像へ、ちょうど $t$ 通りに延長できる。}

$\mathfrak Z$ が $\mathsf P$ 上第一種であるという仮定により、$\mathfrak Z_\Omega/\mathfrak C$ は $\mathfrak Z_\Omega$ と一致する。したがって $\mathfrak S_\Omega$ は、その次数 $s$ と同じ個数の体に分解する：
\[
\mathfrak S_\Omega = E_1\Omega+\cdots+E_s\Omega .
\]
$E_\nu\Omega$ は表現加群として、$\mathfrak S$ から $s$ 個の部分体 $\Sigma_1,\ldots,\Sigma_s$ への $s$ 個の同型写像を生成する。$\mathfrak T$ の同型写像を求めるには $\mathfrak T_\Omega$ を考えればよい。$E_\nu\Omega$ は $\mathfrak S_\Omega$ のイデアルであるから、$E_\nu\Omega=\mathfrak S_\Omega E_\nu$ である。また明らかに
\[
\mathfrak T_\Omega=\mathfrak T\cdot\Omega=\mathfrak T\cdot\mathfrak S_\Omega .
\]
したがって
\[
\mathfrak T_\Omega
=E_1\mathfrak S_\Omega\cdot\mathfrak T+\cdots+E_s\mathfrak S_\Omega\cdot\mathfrak T
=E_1\mathfrak T_\Omega+\cdots+E_s\mathfrak T_\Omega .
\]
これは $\mathfrak T_\Omega$ のイデアルへの分解である。$E_\nu\cdot\mathfrak T_\Omega$ は $t$ 個の単純イデアルへ分解する。実際、次数の関係
\[
[\mathfrak T_\Omega\cdot E_\nu:\mathfrak S_\Omega E_\nu]
\leqq [\mathfrak T_\Omega:\mathfrak S_\Omega]=[\mathfrak T:\mathfrak S]=t
\]
が成り立つ。$\mathfrak S_\Omega E_\nu=\Omega E_\nu\simeq\Omega$ だから、
\[
[\mathfrak T_\Omega\cdot E_\nu:\Omega]
=[\mathfrak T_\Omega E_\nu:\mathfrak S_\Omega E_\nu]\leqq t.
\]
そして、すべての次数 $[\mathfrak T_\Omega E_\nu:\Omega]$ の和は $s\cdot t$ でなければならないから、
\[
[\mathfrak T_\Omega\cdot E_\nu:\Omega]=t,\qquad
E_\nu\mathfrak T_\Omega=e_\nu^{(1)}\cdot\Omega+\cdots+e_\nu^{(t)}\Omega .
\]
各 $e_\nu^{(i)}\Omega$ は $\mathfrak T_\Omega$ の既約表現を与える。これらの表現を $\mathfrak S_\Omega$ に制限すると、$E_\nu\cdot\Omega$ が生成する表現と一致することを\emph{さらに示せば}証明は終わる。しかしこれは明らかである。なぜなら
\[
s\cdot E_\nu=E_\nu\sigma_\nu\quad
(s\in\mathfrak S_\Omega\text{ は $\sigma_\nu$ で表現される})
\]
から
\[
\sum s\cdot e_\nu^{(i)}=\sum e_\nu^{(i)}\sigma_\nu
\]
が従い、和の表示の一意性により
\[
s\cdot e_\nu^{(i)}=e_\nu^{(i)}\sigma_\nu
\]
となるからである。したがって新しい表現も元 $s$ に元 $\sigma_\nu$ を対応させる。

\paragraph{2. $\mathsf T$ を $\mathfrak H$ の不変体とする。}
$\mathfrak H$ が $\mathsf T$ の不変群であることを示したい。そのためには、$\mathfrak H$ に属さない各自己同型によって固定されない $\mathsf T$ の元を一つ示せば\emph{十分}である。

$\mathsf Z$ のガロア群 $\mathfrak G$ の元 $S_i$ は、$\mathsf P$ を固定する $\mathsf Z$ の自己同型である。規定
\[
S_i z=z,\quad \text{$z$ が $\mathfrak Z$ に属するとき}
\]
によって、これを $\mathfrak Z_{\mathsf Z}$ の自己同型にする。$S_i$ を
\[
\mathfrak Z_{\mathsf Z}=e_1\mathsf Z+\cdots+e_n\mathsf Z
\]
へ適用すると、一意に定まる直和因子 $e_\nu\mathsf Z$ を置換しなければならない。特に
\[
S_i e_\nu=e_{\nu'},\qquad\text{かつ}\qquad S_i e_\nu\ne S_k e_\nu
\quad\text{$i\ne k$ のとき}
\]
である。$S_1,\ldots,S_h$ を $\mathfrak H$ の元とする。任意の $e_\nu$、たとえば $e'$ を一つ選び、$S_i$（$i\leqq h$）によってそこから生じる $e_\nu$ をすべて考える：
\[
e_1=S_1e',\ldots,e_h=S_he'.
\]
$i\leqq h,j\leqq h$ ならば
\[
S_i e_j=S_iS_j e'=S_ke'=e_k \quad \text{ただし}\quad k\leqq h
\]
であるから、$\mathfrak H$ の各 $S_i$ は $e_1,\ldots,e_h$ を互いに置換する。したがって
\[
E_1=e_1+\cdots+e_h
\]
は、すべての $S_i\in\mathfrak H$ に対して不変な $\mathfrak Z$ の元である。$\mathfrak Z$ の一つの基底によって $E_1$ を表すと、その係数は $\mathfrak H$ に対して不変な $\mathsf Z$ の元でなければならず、したがって $\mathsf T$ に属する。

$E_1$ は $\mathfrak H$ に属さないどの $S_i$ に対しても不変ではない。実際、そのような $S_i$ に対して
\[
S_i e_1=e_i+\cdots
\]
は $e_1,\ldots,e_h$ の中に現れない成分 $e_i$ を含むが、$E_1$ の和表示は一意だからである。

\subsection*{§ 12. 成分 $e_i$ の形式的意味}\label{die-formale-bedeutung-der-komponenten-ei}

$\mathfrak Z$ の一つの固定した基底 $z_1,\ldots,z_n$ を取る：
\[
\mathfrak Z=z_1\mathsf P+\cdots+z_n\mathsf P.
\]
（ここでは $\mathfrak Z$ が第一種であることだけを仮定する。）同型写像 $\Theta_i$ は、関係式
\[
z_\nu e_i=e_i\zeta_\nu^{(i)}
\]
によって、系 $z_1,\ldots,z_n$ を $\mathsf Z_i$ の基底 $\zeta_1^{(i)},\ldots,\zeta_n^{(i)}$ へ写す。

逆に、単位元の成分 $e_j$ は $\mathsf Z$ の係数を用いて $z_\nu$ によって表されなければならない：
\[
e_j=\sum_\nu \varrho_\nu^{(j)}z_\nu .
\]
したがって
\[
e_j=e_j^2=\sum \varrho_\nu^{(j)}z_\nu e_j
=\sum \varrho_\nu^{(j)}\cdot\zeta_\nu^{(j)}\cdot e_j
\]
および
\[
0=e_je_k=\sum \varrho_\nu^{(j)}\cdot z_\nu\cdot e_k
=\sum \varrho_\nu^{(j)}\cdot\zeta_\nu^{(k)}\cdot e_k
\]
を得る。ゆえに
\[
\sum_\nu \varrho_\nu^{(j)}\cdot\zeta_\nu^{(j)}=1;
\qquad
\sum_\nu \varrho_\nu^{(j)}\cdot\zeta_\nu^{(k)}=0,\quad
\text{$j\ne k$ のとき}.
\]
これは、行列
\[
\begin{pmatrix}
\zeta_1^{(1)} & \cdots & \zeta_n^{(1)}\\
\cdots & \cdots & \cdots\\
\cdots & \cdots & \cdots\\
\cdots & \cdots & \cdots\\
\zeta_1^{(n)} & \cdots & \zeta_n^{(n)}
\end{pmatrix}
\quad\text{と}\quad
\begin{pmatrix}
\varrho_1^{(1)} & \cdots & \varrho_1^{(n)}\\
\cdots & \cdots & \cdots\\
\cdots & \cdots & \cdots\\
\cdots & \cdots & \cdots\\
\varrho_n^{(1)} & \cdots & \varrho_n^{(n)}
\end{pmatrix}
\]
が互いに逆であることを意味する。別の言い方をすれば、$\varrho_1^{(i)},\ldots,\varrho_n^{(i)}$ は $\zeta_1^{(i)},\ldots,\zeta_n^{(i)}$ の相補基底である。この基底は数体の差積の理論で役割を果たす。

\subsection*{§ 13. 一般延長定理}\label{der-allgemeine-fortsetzungssatz}

10 頁で証明した $H_i$ の延長に関する定理は一般化できる。

$\mathfrak Z=z_1\mathsf P+\cdots+z_n\mathsf P$ を可換な完全可約超複素系とする。$\mathfrak S$ と $\mathfrak T$ を、この状況では必然的にやはり完全可約となる $\mathfrak Z$ の部分系とし、
\[
\mathsf P\subseteqq\mathfrak S\subseteqq\mathfrak T\subseteqq\mathfrak Z 
\]
とする。M. Z. § 13 により、三つの系はそれぞれ単位元をもつ。ただし、これら三つの単位元は必ずしも一致しない（$\mathfrak S$ と $\mathfrak T$ が $\mathfrak Z$ 自身とは異なる $\mathfrak Z$ のイデアルならば、確実に一致しない）。$E$ を $\mathfrak S$ の単位元とする。系 $\mathfrak T\cdot E$ は有限階数の $\mathfrak S$-加群である：
\[
E\mathfrak T=a_1\mathfrak S+\cdots+a_t\mathfrak S
\]
（一般には $\mathfrak T$ 自身はそうではない）。

\emph{証明。} $\mathfrak S$ は体の直和である：$\mathfrak S=\mathfrak K_1+\cdots+\mathfrak K_r$。これに対応して $\mathfrak T E$ も直和に分解する。$\mathfrak T E$ の各直和因子は体 $\mathfrak K_i$ のいずれかを含まなければならない。そうでなければ $E$ によって零化されてしまうが、$E$ は $\mathfrak T E$ の単位元だからである。したがって $\mathfrak T E$ は $\mathfrak K_i$ の拡大環の直和、すなわち
\[
\mathfrak T E=a_{11}\mathfrak K_1+\cdots+a_{1p_1}\mathfrak K_1+\cdots+a_{rp_r}\mathfrak K_r
\]
という形の直和である。したがって $E\mathfrak T=a_{11}\mathfrak S+\cdots+a_{rp_r}\mathfrak S$ となる。証明された。

\begin{center}
一般延長定理。
\end{center}

\emph{準同型 $H_i$ によって与えられる $\mathfrak S$ の $\Omega$ における各既約表現は、$\mathfrak T$ の表現へちょうど $t$ 通りに拡張できる。}

証明は 10--11 頁の証明と同様である。すなわち
\[
\mathfrak S_\Omega=E_1\Omega+\cdots+E_s\Omega,\qquad
E_\nu\Omega=\mathfrak S_\Omega E_\nu,
\]
\[
\mathfrak T_\Omega\cdot E=\mathfrak T\cdot E\Omega
=\mathfrak T\mathfrak S\Omega=\mathfrak T\mathfrak S_\Omega,
\]
\[
\mathfrak T_\Omega\cdot E
=\mathfrak T(E_1\mathfrak S_\Omega+\cdots+E_s\mathfrak S_\Omega)
=E_1\mathfrak T_\Omega+\cdots+E_s\mathfrak T_\Omega,
\]
\[
[\mathfrak T_\Omega E_\nu:\Omega]
=[\mathfrak T_\Omega E_\nu:\Omega E_\nu]
=[\mathfrak T_\Omega E_\nu:E_\nu\mathfrak S_\Omega]
\leqq[\mathfrak T_\Omega\cdot E:\mathfrak S_\Omega]
=[\mathfrak T\cdot E:\mathfrak S]=t,
\]
\[
\sum[\mathfrak T_\Omega\cdot E_\nu:\Omega]=st
\]
であり、したがって
\[
[\mathfrak T_\Omega E_\nu:\Omega]=t.
\]
$\mathfrak T_\Omega E_\nu$ はそれゆえ $t$ 個のイデアルへ分解する。その各々は $\mathfrak T_\Omega E$ の、したがって $\mathfrak T_\Omega$ の直和因子なので、$\mathfrak T_\Omega$ の一つの表現を与え、10--11 頁と同様に、それが $H_\nu$ により与えられる表現の延長であることがわかる。$\mathfrak T_\Omega$ の残りの既約表現、すなわち $\mathfrak T_\Omega\cdot E_\nu$ に含まれない $\mathfrak T_\Omega$ の直和因子から生成される表現は、$H_\nu$ が生成する $\mathfrak S_\Omega$ の表現の延長ではありえない。これらのイデアルは $\mathfrak T_\Omega\cdot E_\nu$ によって、したがってなおさら $\mathfrak S_\Omega$ によって零化されるからである。


%% END NOETHER SOURCE
%% BEGIN NOETHER SOURCE "ja/t2.tex"
\section*{第 III 章　アーベル群}\label{kapitel-iii.-abelsche-gruppen}

\subsection*{§ 14. 群環}\label{der-gruppenring}

$\mathfrak G$ を元 $a_1,\ldots,a_h$ からなる有限群、$\mathsf P$ を可換体とする。群の乗法を乗法として超複素系
\[
\mathfrak o[\mathfrak G]=a_1\mathsf P+\cdots+a_h\mathsf P
\]
を作り、これを\emph{$\mathsf P$ における $\mathfrak G$ の群環}という（M. Z. § 6、Speiser, \emph{Theorie d. Gruppen end. Ord.} § 48）。

$\mathfrak o[\mathfrak G]$ の $\mathsf P$ における表現、または $\mathsf P$ の代数閉拡大体 $\Omega$ における表現を求める。すでに述べたように、既約表現は $\mathfrak o[\mathfrak G]/\mathfrak C$ の単純左イデアルによって生成される。したがって、ガロア理論の場合と同じく、$\mathfrak o[\mathfrak G]$ がいつ根基をもつかが問題となる。これについて次の一般定理が成り立つ。

\begin{center}
\emph{$\mathfrak o[\mathfrak G]$ が完全可約であるのは、$\mathfrak G$ の位数 $h$ が $\mathsf P$ の標数 $p$ で割り切れないとき、かつそのときに限る。}
\end{center}

まず次を証明する：
\[
\text{$h\equiv0(p)$ ならば $\mathfrak C\ne 0$ である。}
\]
\[
\mathfrak a=\mathsf P(a_1+\cdots+a_h)
\]
は $\mathfrak o$ の両側イデアルである。実際、
\[
\mathsf P\cdot(a_1+\cdots+a_h)a
=\mathsf P\cdot(a_1a+\cdots+a_ha)
=\mathsf P(a_1+\cdots+a_h)=\mathfrak a
\]
および
\[
a\cdot\mathsf P(a_1+\cdots+a_h)
=\mathsf P(aa_1+\cdots+aa_h)
=\mathsf P(a_1+\cdots+a_h)=\mathfrak a
\]
が、$a_i$ の群の性質により成り立つ。$\mathfrak a$ はもちろん零ではないが、冪零である：
\[
\left(\sum_i a_i\right)^2=\sum_i\sum_j a_i a_j
=h\sum_j a_j=0.
\]

逆、すなわち
\[
\text{$h\not\equiv0(p)$ ならば $\mathfrak C=0$ である}
\]
ことは、ここではアーベル群についてだけ証明する（M. Z. § 26 では定理のこの部分が一般に証明されている）。

\subsection*{§ 15. アーベル群の群環}\label{die-gruppenringe-abelscher-gruppen}

$\mathfrak A=\{a_1,\ldots,a_h\}$ をアーベル群とする。このとき $\mathfrak o$ は可換なので、可換超複素系に関する一般定理を適用できる。

\emph{$h$ が $\mathsf P$ の標数 $p$ で割り切れないならば、$\mathfrak o[\mathfrak A]$ は $\Omega$ においてちょうど $h$ 個の相異なる既約表現をもつ（$\mathfrak o$ は可換なので、それらは一次である）。}

ここから $\mathfrak o[\mathfrak A]/\mathfrak C$ が $\mathfrak o$ と同じ階数をもち、したがって $\mathfrak C=0$ であることが従う。これは先ほど証明を保留した定理の第二の部分を、アーベル群について与える。

\emph{証明。} $\mathfrak A$ は、位数 $h_i$ の巡回群 $\mathfrak Z_i$ の直積である：
\[
\mathfrak A=\mathfrak Z_1\times\cdots\times\mathfrak Z_t,\qquad
h=h_1h_2\ldots h_t .
\]
$z_i$ を $\mathfrak Z_i$ の生成元とすると、一次既約表現においてそれに対応する元は $h_i$ 乗根 $\varepsilon_\nu^{(h_i)}$ でなければならない。逆に、各 $z_i$ に一つの $h_i$ 乗根 $\varepsilon_\nu^{(h_i)}$ を対応させれば、$\mathfrak o$ の $\Omega$ における既約表現が得られる。いま $h\not\equiv0(p)$、したがってなおさら $h_i\not\equiv0(p)$ であるから、$\mathfrak o$ にはちょうど $h_1h_2\cdots h_t=h$ 個の相異なる表現がある。証明された。

（$h=0(p)$、したがって少なくとも一つの $h_i$ について $h_i=0(p)$ ならば、相異なる $\varepsilon_\nu^{(h_i)}$ は高々 $h_i/p$ 個であり、したがって表現は $h$ 個より少ないことも容易にわかる。ゆえに $\mathfrak o_\Omega$ は零でない根基を含まなければならない。ただし、これだけから $\mathfrak o$ の根基について直ちに何かが従うわけではない。）

$\mathfrak o_\Omega$ から $\Omega$ の部分環への $h$ 個の相異なる準同型写像を $\Theta_1,\ldots,\Theta_h$ と書く。各 $\Theta_i$ には、表現加群として $\mathfrak o_\Omega$ の直和因子 $e_i\Omega$ が属する。$\Theta_i$ を $\mathfrak A$ にだけ適用して、$\mathfrak A$ の表現と呼ぶこともある。$\mathfrak A$ から単位根の群へのこれらの写像を、\emph{$\mathfrak A$ の $h$ 個の指標}という。特別な指標として、$\mathfrak A$ の各元に $1$ を対応させる\emph{主指標}がある。これに属する $\Theta_i$ を $\Theta_1$ とする。

主指標による表現は、アーベル群に限らず、どの群についても可能である。この表現はすでに $\mathsf P$ にあるから、$\mathfrak o$ が完全可約である場合には、$\mathfrak o$ はこの表現を与える単純直和因子 $E\cdot\mathfrak o$ を分離しなければならない。$E$ は容易に計算できる。$E=\sum x_i a_i$ でなければならない。第一に
\[
aE=E\cdot 1\quad\text{すなわち}\quad \sum x_i a a_i=\sum x_i a_i
\]
である。$a=a_j a_i^{-1}$ とすれば
\[
x_1a_ja_i^{-1}a_1+\cdots+x_i a_j+\cdots+x_na_ja_i^{-1}a_n
=x_1a_1+\cdots+x_na_n
\]
となり、したがって
\[
x_i=x_j
\]
が従う。

第二に
\[
E=E^2=x^2\left(\sum a_i\right)^2=x^2h\sum a_i=xhE.
\]
$h=0(p)$ の場合には $E=0$ が従う。すなわち $\mathfrak o$ は主指標の表現加群を含まず、したがって完全可約ではありえない。これで 13 頁の定理の一方向が改めて証明された。

一方、$h\not\equiv0(p)$、したがって $\mathfrak o$ が完全可約ならば、直ちに
\[
x=\frac{1}{h},\qquad E=\frac{1}{h}(a_1+\cdots+a_h)
\]
が従う。

\subsection*{§ 16. 指標関係式}\label{die-charakterenrelationen}

$E=(1/h)\cdot\sum a_i$ が主指標を与えることから、指標関係式が直ちに得られる。すなわち
\[a \cdot e_i = e_i \cdot \Theta_i a\]
でなければならず、したがって
\[e_i \cdot \Theta_i e_j = \begin{cases} e_i & i = j \\ 0 & i \neq j \end{cases}\]
または
\[\Theta_i e_j = \begin{cases} 1 & i = j \\ 0 & i \neq j \end{cases}\]
である。特に
\[\Theta_i e_1 = \begin{cases} 1 & i = 1 \\ 0 & i \neq 1 \end{cases}\]
である。$e_1=(1/h)\sum a_i$ だから、
\[
\sum_\nu \Theta_i a_\nu=\begin{cases} h & i=1,\\ 0 & i\ne 1. \end{cases}
\]
これが周知の指標関係式である。

\emph{一つの指標をすべての群元にわたって足した和は、主指標を除けば零であり、主指標については $h$ である。}

\subsection*{§ 17. アーベル群のガロア理論}\label{die-galoissche-theorie-abelscher-gruppen}

アーベル群の群環について、可換体のガロア理論で証明したものと類似の定理を述べることができる。ガロア理論では $\Theta_i\Theta_1^{-1}$ が群 $\mathfrak G$ をなした。基本定理は、ある不変性に基づいて $\mathfrak G$ の部分群と $Z$ の部分体との間に一対一対応があることを表していた。

$\mathfrak{o}[\mathfrak{A}]$ の準同型 $\Theta_i$ は一つの群 $\mathsf A$ にまとめられ、その部分群は $\mathfrak A$ の部分群と類似の関係をもつ。

\emph{積 $\Theta_i \Theta_k=\Theta$ を}
\[
\Theta a = \Theta_i a \cdot \Theta_k a
\]
\emph{によって定めると、$\Theta_i$ 全体は $\mathfrak A$ と同型な群 $\mathsf A$ をなす。}

\emph{証明。} $\Theta=\Theta_i\cdot\Theta_k$ は、
\[\begin{aligned}
\Theta a \cdot \Theta b &= \Theta_i a \cdot \Theta_k a \cdot \Theta_i b \cdot \Theta_k b \\
&= \Theta_i a \cdot \Theta_i b \cdot \Theta_k a \cdot \Theta_k b \\
&= \Theta_i a b \cdot \Theta_k a b = \Theta a b
\end{aligned}\]
により $\mathfrak A$ から $\mathsf A$ への準同型である。$\mathfrak A$ の準同型には $\Theta_i$ 以外のものがないから、$\Theta$ はある $\Theta_j$ に等しい。したがって $\Theta_i$ は実際に $h$ 元からなる群をなす。

$\varepsilon_{h_\nu}$ を原始 $h_\nu$ 乗根とする。$\Theta_i$ は $\mathfrak A$ の巡回因子 $\mathfrak Z_\nu$ の生成元 $z_\nu$ を $\varepsilon_{h_\nu}$ のある冪へ写す：
\[
\Theta_i z_\nu=\varepsilon_{h_\nu}^{\varrho_\nu^{(i)}}.
\]
$\Theta_i$ に $\mathfrak A$ の元 $\prod_\nu z_\nu^{\varrho_\nu^{(i)}}$ を対応させる。これは $\mathsf A$ から $\mathfrak A$ への一対一写像である。実際、$\Theta_i$ は $\varrho_\nu^{(i)}$ によって定まるので、異なる $\Theta_i$ には異なる $\prod_\nu z_\nu^{\varrho_\nu^{(i)}}$ が属する。この写像は同型である。一方では
\[
\Theta_i\Theta_j z_\nu=\Theta_i z_\nu\cdot \Theta_j z_\nu
=\varepsilon_{h_\nu}^{\varrho_\nu^{(i)}+\varrho_\nu^{(j)}}
\]
であり、他方では
\[
\prod_\nu z_\nu^{\varrho_\nu^{(i)}}\cdot
\prod_\nu z_\nu^{\varrho_\nu^{(j)}}=
\prod_\nu z_\nu^{\varrho_\nu^{(i)}+\varrho_\nu^{(j)}}.
\quad\text{証明された。}
\]

ここで次を定義する。

\emph{$\mathsf A$ の部分群 $\mathsf B$ の不変領域：}

$\mathsf B$ の各 $\Theta$ によって $1$ へ写される $\mathfrak A$ の元からなる部分群 $\mathfrak B$。

\emph{$\mathfrak A$ の部分群 $\mathfrak B$ の不変群：}

$\mathfrak B$ の各 $a$ を $1$ へ写す $\mathsf A$ の元からなる部分群 $\mathsf B$。

そして基本定理を次のように述べる。

\emph{$\mathfrak A$ の部分群 $\mathfrak B$ と $\mathsf A$ の部分群 $\mathsf B$ との間には一対一対応があり、$\mathfrak B$ と $\mathsf B$ が対応するとき、$\mathsf B$ は $\mathfrak B$ の不変群、$\mathfrak B$ は $\mathsf B$ の不変領域である。}

したがって、次を示さなければならない。
\begin{enumerate}
\item \emph{$\mathsf B$ が $\mathfrak B$ の不変群ならば、$\mathfrak B$ は $\mathsf B$ の不変領域である。}
\item \emph{$\mathfrak B$ が $\mathsf B$ の不変領域ならば、$\mathsf B$ は $\mathfrak B$ の不変群である。}
\end{enumerate}

\paragraph{1.}
これは 10 頁と同様に一般延長定理から従う。$\mathsf B$ を $\mathfrak B$ の不変群、$\mathfrak B'$ を $\mathsf B$ の不変領域とする。$\mathfrak B'$ は $\mathfrak B$ を含む。$\mathsf B$ の元は、$\mathfrak B$ を $1$ へ写す一つの準同型を $\mathfrak A$ 全体へ延長したものである。これらの延長を $\mathfrak B'$ にだけ適用すると、一般延長定理により、それらは $t$ 個の類に分かれる。同じ類の準同型は $\mathfrak B'$ 上で一致し、異なる類の準同型は $\mathfrak B'$ 上で異なる。ここで $t$ は $\mathfrak o[\mathfrak B']$ の $\mathfrak o[\mathfrak B]$ 上の次数、すなわち $\mathfrak B'$ における $\mathfrak B$ の指数である。$\mathfrak o[\mathfrak B]$ と $\mathfrak o[\mathfrak B']$ は同じ単位元、すなわち $\mathfrak A$ 全体の単位元をもつからである。ところが $\mathfrak B'$ は $\mathsf B$ の不変領域なので、$t=1$、したがって $\mathfrak B'=\mathfrak B$ でなければならない。証明された。

\textbf{2.} を 1. へ帰着させる。

$\mathfrak A$ の元 $a$ を $\mathsf A$ の準同型とみなす。すなわち $a\Theta$ を単位根 $\Theta a$ に等しいと置き、簡潔に $a\Theta=\Theta a$ と書く。これによって $a$ は実際に $\mathsf A$ の準同型となる。なぜなら
\[
a\Theta\cdot a\Theta'=\Theta a\cdot\Theta' a=\Theta\Theta' a=a\Theta\Theta'
\]
だからである。

二つの準同型 $a,b$ の積は通常の群の積である：
\[
\begin{array}{cccccc}
ab\Theta&=&a\Theta\cdot b\Theta&=&\Theta a\cdot\Theta b=\Theta ab&=ab\Theta\\
&&\text{準同型の積}&&\text{群の積}&
\end{array}
\]

群元として異なる $a,b$ は、準同型としても異なる。すべての $\Theta$ について $a\Theta=b\Theta$ ならば、すべての $\Theta$ について $\Theta a=\Theta b$、すなわち $\Theta ab^{-1}=1$ となる。したがって $\Theta$ は、$ab^{-1}$ が生成する $\mathfrak A$ の巡回部分群 $\mathfrak Z$ の恒等準同型を $\mathfrak A$ 全体へ延長したものである。延長定理により、$\mathfrak A$ における $\mathfrak Z$ の指数は $h$、すなわち $\mathfrak Z=e$ であり、$a=b$ となる。

新しい見方で述べた次の命題
\[
\begin{array}{ll}
\text{a)}& \mathfrak B\ \text{は }\mathsf B\text{ の不変群である。}\\
\text{b)}& \mathsf B\ \text{は }\mathfrak B\text{ の不変領域である。}
\end{array}
\]
を、もとの見方へ翻訳する。

a) について、$\mathfrak B$ は、すべての $\Theta\in\mathsf B$ に対して $a\Theta=1$ となる $a\in\mathfrak A$ すべてからなる。すなわち、すべての $\Theta\in\mathsf B$ に対して $\Theta\cdot a=1$ となる $a\in\mathfrak A$ すべてからなる。したがって a) の翻訳は
\[
\text{a')} \quad \mathfrak B\ \text{は }\mathsf B\text{ の不変領域である。}
\]
である。

同様に b) の翻訳は
\[
\text{b')} \quad \mathsf B\ \text{は }\mathfrak B\text{ の不変群である。}
\]
である。

したがって、1. ですでに証明した命題

$\mathfrak B$ が $\mathsf B$ の不変群ならば、$\mathsf B$ は $\mathfrak B$ の不変領域である

を翻訳したものは、まさに証明すべき命題

$\mathfrak B$ が $\mathsf B$ の不変領域ならば、$\mathsf B$ は $\mathfrak B$ の不変群である

となる。

15 頁の指標関係式を $\mathfrak A$ の代わりに $\mathsf A$ について書けば、
\[
\sum_{i=1}^{h} \Theta_i a_\nu =
\begin{cases}
h & \nu=1,\\
0 & \nu\ne 1
\end{cases}
\]
を得る。

すべての指標を一つの群元について足した和は、単位元を除けば零であり、単位元については $h$ である。


%% END NOETHER SOURCE
%% BEGIN NOETHER SOURCE "ja/t3.tex"
\section*{第 IV 章　両側単純環}\label{kapitel-iv.-zweiseitig-einfache-ringe}

\subsection*{§ 18. 補題}\label{ein-hilfssatz}

$\mathsf K$ を体、$\mathfrak G$ を $\mathsf K$ の自己同型（それ自身への同型写像）からなる群とする。$\mathfrak G$ に属する各自己同型によって不変な $\mathsf K$ の元全体は $\mathsf K$ の部分体 $\mathsf P$ をなし、これを $\mathfrak G$ の\emph{不変体}という。

$\tau\ne0$ を $\mathsf K$ の任意の元とすると、$\mathsf K$ の元 $\alpha$ に $\tau^{-1}\alpha\tau$ を対応させることによって $\mathsf K$ の自己同型が得られる。群論の対応する用語にならって、このような自己同型を\emph{内部自己同型}という。一つの体 $\mathsf K$ の内部自己同型すべてからなる群の不変体 $\mathsf P$ は、$\mathsf K$ の中心である。

いま、体 $\mathsf K$、その自己同型群 $\mathfrak G$、およびその不変体 $\mathsf P$ が与えられているとする。

\[
\mathfrak M=x_1\mathsf P+\cdots+x_n\mathsf P
\]

を階数 $n$ の $\mathsf P$-加群とする。$\mathfrak G$ の元 $\gamma$ の $\mathfrak M$ の元への作用を

\[
\gamma\cdot x_i=x_i,\qquad\text{したがって}\qquad
\gamma\sum \varrho_i x_i=\sum \varrho_i x_i
\]

によって定める。すると $\mathfrak G$ は拡大加群

\[
\mathfrak M_{\mathsf K}=x_1\mathsf K+\cdots+x_n\mathsf K
\]

の自己同型群となり、$\mathfrak M_{\mathsf K}$ における $\mathfrak G$ の不変領域は $\mathfrak M$ となる。この仮定のもとで、次の定理が成り立つ。

\begingroup\itshape
\[
\mathfrak C=z_1\mathsf K+\cdots+z_r\mathsf K
\]

が $\mathfrak M_{\mathsf K}$ の部分加群であり、各 $\gamma\in\mathfrak G$ によってそれ自身へ移される（各元が固定される必要はない）ならば、$\mathfrak C$ は $\mathfrak M$ のある部分加群 $\mathfrak c$ を $\mathsf K$ によって拡大した加群である：

\[
\mathfrak C=\mathfrak c_{\mathsf K}.
\]
\endgroup

\emph{証明。} 1. 予備的注意として、$\mathfrak c=\mathfrak C\cap\mathfrak M$ でなければならない。一般に
\[
\mathfrak c=\mathfrak c_{\mathsf K}\cap\mathfrak M
\]
だからである。

$\mathfrak c\subseteq\mathfrak c_{\mathsf K}\cap\mathfrak M$ は明らかである。逆の包含を示すため、$\mathfrak c$ の基底 $y_1,\ldots,y_r$ を取り、
\[
\mathfrak c=y_1\mathsf P+\cdots+y_r\mathsf P
\]
とする。このとき
\[
\mathfrak c_{\mathsf K}=y_1\mathsf K+\cdots+y_r\mathsf K.
\]
$\mathfrak c_{\mathsf K}\cap\mathfrak M$ の各元 $c$ は、$\mathfrak c_{\mathsf K}$ の元としての基底表示
\[
c=\sum_{i=1}^{r}y_i\varkappa_i
\]
と、$\mathfrak M$ の元としての基底表示
\[
c=\sum_{i=1}^{r}y_i\varrho_i+\sum_{j=r+1}^{n}y_j\varrho_j
\]
をもつ。ここで $y_{r+1},\ldots,y_n$ は、$y_1,\ldots,y_r$ を $\mathfrak M$ の基底に補う $\mathfrak M$ の元である（M. Z. § 6）。両方の表示を $\mathfrak M_{\mathsf K}$ における基底表示とみなせるので、それらは一致し、$\chi_i=\varrho_i$（$i=1,\ldots,r$）、$\varrho_j=0$（$j=r+1,\ldots,n$）である。したがって
\[
c=\sum_{i=1}^{r}y_i\varrho_i,
\qquad\text{証明された。}
\]

\textbf{2.} 定理を証明するため、元 $z_i$ が $\mathfrak M$ に属するような $\mathfrak C$ の基底 $z_1,\ldots,z_r$ を構成する。このとき
\[
\mathfrak c=z_1\mathsf P+\cdots+z_r\mathsf P
\]
と置けば、実際に
\[
\mathfrak C=\mathfrak c_{\mathsf K}
\]
となる。

$\mathfrak C$ の任意の基底 $y_1,\ldots,y_r$ を取り、$\mathfrak M$ の一つの基底 $x_1,\ldots,x_n$ の適当な元 $x_{r+1},\ldots,x_n$ を付け加えて、$\mathfrak M_{\mathsf K}$ の基底にする。この基底は同時に $\mathfrak M_{\mathsf K}$ の基底でもある。この選択は M. Z. § 5 により可能である：
\[
\mathfrak M_{\mathsf K}=\mathfrak C+x_{r+1}\mathsf K+\cdots+x_n\mathsf K.
\]
ここから、$\mathfrak M$ の基底の最初の $r$ 元 $x_1,\ldots,x_r$ について
\[
x_i=z_i-\sum_{j=r+1}^{n}x_j\varkappa_j^{(i)},\qquad
z_i\in\mathfrak C,\quad \varkappa_j^{(i)}\in\mathsf K
\]
が得られる。$z_i$ は $x_{r+1},\ldots,x_n$ とともに明らかに $\mathfrak M_{\mathsf K}$ の基底をなすので線型独立であり、$z_1,\ldots,z_r$ は $\mathfrak C$ の基底である。この基底が求める形であること、すなわち $z_i\in\mathfrak M$ を示す。

$\gamma$ を $\mathfrak G$ の元とする。一方では
\[
\begin{array}{rcl}
(1)\qquad \gamma z_i
&=&\gamma x_i+\displaystyle\sum_{j=r+1}^{n}\gamma x_j\cdot \gamma\varkappa_j^{(i)}
=x_i+\displaystyle\sum_{j=r+1}^{n}x_j\gamma\varkappa_j^{(i)}.
\end{array}
\]
他方では、$\mathfrak C$ がすべての $\gamma\in\mathfrak G$ を許すから、
\[
\begin{array}{rcl}
(2)\qquad \gamma z_i&=&\displaystyle\sum_{\mu=1}^{r}z_\mu\sigma_\mu^{(i)},
\qquad \sigma_\mu^{(i)}\in\mathsf K.
\end{array}
\]
(2) はさらに
\[
\begin{array}{rcl}
(3)\qquad \gamma z_i
&=&\displaystyle\sum_{\mu=1}^{r}x_\mu\sigma_\mu^{(i)}
+\displaystyle\sum_{j=r+1}^{n}x_j\sum_{\mu=1}^{r}\varkappa_j^{(\mu)}\sigma_\mu^{(i)}
\end{array}
\]
と書ける。(1) と (3) の係数を比較すると
\[
\sigma_j^{(i)}=
\begin{cases}
1 & i=j,\\
0 & i\ne j
\end{cases}
\]
を得る。したがって
\[
\gamma\cdot z_i=z_i.
\]
ゆえに $z_i$ は $\mathfrak G$ の不変領域、すなわち $\mathfrak M$ に属する。証明された。

\subsection*{§ 19. 係数体の非可換拡大体における両側単純超複素系の表現}\label{darstellungen-zweiseitig-einfacher-hyperkomplexer-systeme-in-nichtkommutativen-erweiterungskoerpern-ihrer-koeffizientenkoerper}

\[
\mathfrak S=x_1\mathsf P+\cdots+x_n\mathsf P
\]

を両側単純、したがって完全可約な超複素系とし、$\mathsf K$ を中心 $\mathsf P$ の体とする。

\textbf{定理 1.} \emph{$\mathfrak S_{\mathsf K}$ は両側単純である（したがって完全可約でもある）。}

\emph{証明。} 1. $\mathfrak S_{\mathsf K}$ の各両側イデアル $\mathfrak a$ は、$\mathsf K$ の内部自己同型群 $\mathfrak G$ に関して不変である。実際、零でない $\mathsf K$ の元 $\chi$ に対して $\chi^{-1}\mathfrak a\chi=\bar{\mathfrak a}$ と置けば、$\mathfrak a$ のイデアル性により $\bar{\mathfrak a}\subseteq\mathfrak a$ である。しかし $\bar{\mathfrak a}$ は $\mathsf K$-加群として $\mathfrak a$ と同じ $\mathsf K$-階数をもつので、$\bar{\mathfrak a}=\mathfrak a$ でなければならない。

\textbf{2.} 直前の補題により、$\mathfrak S_{\mathsf K}$ の両側イデアル $\mathfrak a$ は $\mathfrak S$ の部分加群 $\mathfrak A$ の拡大イデアル、$\mathfrak a=\mathfrak A_{\mathsf K}$ でなければならない。$\mathfrak A=\mathfrak S\cap\mathfrak a$ であるから、
\[
\mathfrak S\mathfrak A\subseteqq\mathfrak S\mathfrak S\cap\mathfrak S\mathfrak a\subseteqq\mathfrak S\cap\mathfrak a=\mathfrak A
\]
および
\[
\mathfrak A\mathfrak S\subseteqq\mathfrak S\mathfrak S\cap\mathfrak a\mathfrak S\subseteqq\mathfrak S\cap\mathfrak a=\mathfrak A
\]
となり、$\mathfrak A$ は $\mathfrak S$ の両側イデアルである。ところが $\mathfrak S$ の両側イデアルは $0$ と $\mathfrak S$ だけなので、$\mathfrak S_{\mathsf K}$ の両側イデアルも $0$ と $\mathfrak S_{\mathsf K}$ だけである。証明された。

次に $\mathfrak S$ の $\mathsf K$ における既約表現を求める。形式上の利点があるので、相反表現を構成する。直接表現と相反表現は一対一に対応するため、どちらを取っても結局は同じである。

$\mathfrak M$ を $\mathfrak S$ の $\mathsf K$ における相反表現加群、すなわち有限階数の $\mathfrak S$-左加群かつ $\mathsf K$-左加群
\[
\mathfrak M=\mathsf K\cdot y_1+\cdots+\mathsf K\cdot y_r
\]
で、交換則
\[
s(\varkappa m)=\varkappa (sm)
\]
を満たすものとする。

$\mathfrak M$ が $\mathfrak S_{\mathsf K}$-左加群ともみなせることを示す。まず $\mathfrak S_{\mathsf K}$ の元と $\mathfrak M$ の元との積を定めなければならない。$s\in\mathfrak S$、$x\in\mathsf K$、$m\in\mathfrak M$ に対して
\[
s\varkappa \cdot m=s\cdot(\varkappa \cdot m)
\]
と置き、一般の $\mathfrak S_{\mathsf K}$ の元 $\sum s_ix_i$ には
\[
\sum s_i\varkappa _i\cdot m=\sum s_i\cdot(\varkappa _i\cdot m)
\]
と定める。

この定義は一意であり、$\mathfrak S_{\mathsf K}$ の元を積 $s_ix_i$ の和として表す仕方に依存しない。すなわち
\[
\begin{array}{rcl}
(1)\qquad \displaystyle\sum_{i=1}^{p}s_i\varkappa _i&=&0
\end{array}
\]
から
\[
\sum_{i=1}^{p}s_i(\varkappa _i m)=0
\]
が従うことを示せばよい。$s_1,\ldots,s_q$ が線型独立で、$s_{q+1},\ldots,s_p$ が
\[
s_j=\sum_{i=1}^{q}s_i\varrho_i^{(j)},\qquad j=q+1,\ldots,p
\]
と表されるとする。表示は一意で、しかもすでに $\mathfrak S$ の中で可能だから、係数は $\mathsf P$ に属する。(1) から
\[
\varkappa _i+\sum_{j=q+1}^{p}\varrho_i^{(j)}\varkappa _j=0,\qquad i=1,\ldots,q
\]
を得る。したがって
\[
\begin{aligned}
\sum_{i=1}^{p}s_i(\varkappa _i m)
&=s_1(\varkappa _1m)+\cdots+s_q(\varkappa _qm)
  +\left(\sum_i s_i\varrho_i^{(q+1)}\right)(\varkappa _{q+1}m)\\
&\quad+\cdots+\left(\sum_i s_i\varrho_i^{(p)}\right)(\varkappa _pm)\\
&=s_1(\varkappa _1m)+\cdots+s_q(\varkappa _qm)\\
&\quad+s_1\varrho_1^{(q+1)}(\varkappa _{q+1}m)+\cdots+s_q\varrho_q^{(q+1)}(\varkappa _{q+1}m)
  \qquad\text{加群の性質により}\\
&\quad+\cdots\\
&\quad\cdots\\
&\quad+s_1\varrho_1^{(p)}(\varkappa _pm)+\cdots+s_q\varrho_q^{(p)}(\varkappa _pm)\\
&=s_1(\varkappa _1m)+\cdots+s_q(\varkappa _qm)\\
&\quad+s_1(\varrho_1^{(q+1)}\varkappa _{q+1}m)+\cdots+s_q(\varrho_q^{(q+1)}\varkappa _{q+1}m)\\
&\quad+\cdots\\
&\quad\cdots\\
&\quad+s_1(\varrho_1^{(p)}\varkappa _pm)+\cdots+s_q(\varrho_q^{(p)}\varkappa _pm)
  \qquad\begin{gathered}\text{結合則 }s'(sm)=s's\cdot m\text{ および}\\ \varkappa'(\varkappa m)=\varkappa'\varkappa\cdot m\text{ により}\end{gathered}\\
&=s_1\left(\varkappa _1m+\sum\varrho_1^{(q+1)}\varkappa _{q+1}\cdot m\right)+\cdots\\
&=s_1\left(\left(\varkappa _1+\sum\varrho_1^{q+1}\varkappa _{q+1}\right)m\right)+\cdots\\
&=0.
\end{aligned}
\]

$\mathfrak M$ の $\mathfrak S_{\mathsf K}$-加群としての分配則 $(S+S')m=Sm+S'm$、$S(m+m')=Sm+Sm'$ は明らかである。注意を要するのは結合則 $S(S'm)=SS'\cdot m$ である。分配則が成り立つので、
\[
s\varkappa \cdot(s'\varkappa ' m)=s\varkappa (s'\varkappa '\cdot m)
\]
の場合だけを考えればよい。一方、$sx\cdot m$ の定義により
\[
s\varkappa \cdot(s'\varkappa '\cdot m)=s\cdot(\varkappa \cdot(s'\cdot(\varkappa '\cdot m)))
\]
である。他方、
\[
s\varkappa  s'\varkappa '\cdot m=ss'\varkappa \varkappa '\cdot m=ss'\cdot(\varkappa \varkappa '\cdot m)
=s\cdot(s'\cdot(\varkappa \cdot(\varkappa '\cdot m)))
\]
は、二つの交換関係
\[
\varkappa \cdot s'=s'\cdot \varkappa
\]
および
\[
s'\cdot(\varkappa \cdot m)=\varkappa \cdot(s'\cdot m)
\]
により $s\cdot(x\cdot(s'\cdot(x'\cdot m)))$ に等しい。

したがって $\mathfrak M$ は実際に有限 $\mathfrak S_{\mathsf K}$-左加群である。逆に、各有限 $\mathfrak S_{\mathsf K}$-左加群は $\mathfrak S$ の $\mathsf K$ における相反表現加群である。実際、明らかに有限 $\mathsf K$-加群であり、
\[
s(\varkappa \cdot m)=s\varkappa \cdot m=\varkappa s\cdot m=\varkappa (s\cdot m)
\]
により交換条件も満たす（§ 24 のより簡単な証明も参照）。

M. Z. § 18 により、各有限 $\mathfrak S_{\mathsf K}$-加群は、$\mathfrak S_{\mathsf K}$ の単純左イデアルと $\mathfrak S_{\mathsf K}$-作用素同型な単純加群の和である。逆に $\mathfrak S_{\mathsf K}$ の各左イデアルは $\mathfrak S$ の $\mathsf K$ における表現加群であり、$\mathfrak S_{\mathsf K}$ の単純左イデアルはすべて作用素同型なので、次が成り立つ。

\textbf{定理 2.} \emph{$\mathfrak S$ の $\mathsf K$ における既約相反表現類はただ一つであり、$\mathfrak S_{\mathsf K}$ の一つの単純左イデアルによって生成される。}

この既約表現の次数を $r$ とすると、任意の表現の次数は $r$ の倍数である。もちろん既約表現を並べて結合すれば、$r$ の任意の倍数 $rs$ に対して $rs$ 次表現を構成できる。これは $s$ 個の単純加群の直和である表現加群を作ることに対応する。

$\mathfrak S$ 自身を表現加群として生成される $\mathfrak S$ の $\mathsf K$ における表現（主表現、M. Z. § 20）はすでに知られており、これは $\mathsf P$ においてすでに存在する。その次数は $n$ なので、$n$ は $r$ で割り切れる：
\[
n=rt.
\]
商 $t$ は容易に定まる。$r$ は一つの左イデアルの $\mathsf K$ 上の階数であり、$n$ は $\mathfrak S_{\mathsf K}$ 全体の階数だから、$t$ は $\mathfrak S_{\mathsf K}$ の左イデアルの個数である。したがって $t^2$ は、その単純イデアルの自己同型体上の $\mathfrak S_{\mathsf K}$ の階数である。

以下、$\mathfrak S$ が可換体である場合に特殊化する。

\textbf{定理 3.} \emph{$\mathfrak S$ が可換体ならば、$\mathfrak S$ は $\mathfrak S_{\mathsf K}$ の中心である。}

\emph{証明。} $\mathfrak S$ の各元は $\mathfrak S$ の各元および $\mathsf K$ の各元と、したがって $\mathfrak S_{\mathsf K}$ の各元と可換なので、$\mathfrak S$ は $\mathfrak S_{\mathsf K}$ の中心に属する。逆に、$\mathfrak S_{\mathsf K}$ の元がそのすべての元と可換であるならば、少なくとも $\mathsf K$ の各元と可換でなければならない。したがってその係数は $\mathsf K$ の中心 $\mathsf P$ に属し、その元は $\mathfrak S$ に属する。証明された。

可換環では相反同型と直接同型に差がないので、各相反表現は直接表現ともみなせる。したがって $\mathfrak S$ の $\mathsf K$ における既約表現は、直接表現加群によっても与えられなければならない。この直接表現加群は $\mathfrak S_{\mathsf K}$ の単純右イデアルである。実際、$\mathfrak r$ は階数 $r$ の有限 $\mathsf K$-右加群であり、$\mathfrak r\mathfrak S=\mathfrak S\mathfrak r$ により $\mathfrak S$-左加群でもあり（もちろん混合結合則を満たす）、階数 $r$ の直接表現加群である。既約表現類は一つだけなので、それは与えられた表現を生成しなければならない。

\subsection*{§ 20. 非可換体}\label{nichtkommutative-koerper}

ここで
\[
\mathfrak K=y_1\mathsf P+\cdots+y_m\mathsf P
\]
を、中心が $\mathsf P$ である非可換体とする。$\mathfrak K$ について、可換な有限代数体拡大の理論における諸定理と類似する一連の定理を展開する。

$Z=\xi_1\mathsf P+\cdots+\xi_n\mathsf P$ を $\mathsf P$ 上の超複素系とすると、二つの拡大系 $\mathfrak K_Z$ と $Z_{\mathfrak K}$ は一致する。実際、
\[
\mathfrak K_{\mathsf{Z}}=y_1\xi_1\mathsf P+\cdots+y_m\xi_n\mathsf P
\]
および
\[
\mathsf{Z}_{\mathfrak K}=\xi_1y_1\mathsf P+\cdots+\xi_ny_m\mathsf P
\]
と書け、$y_i\xi_j=\xi_jy_i$ により両者は一致する。

以下では $Z$ を常に可換体とするので、§ 19 の定理を適用できる。§ 19 の定理 1 により $\mathfrak K_Z$ は常に両側単純である。

\textbf{定理 1.} \emph{$\Omega$ を $\mathsf P$ の代数閉拡大体とする。このとき $\mathfrak K_\Omega$ の中心は $\Omega$ に等しい。}

\emph{証明。} § 19 の定理 3 を用いる。$\Omega$ が $\mathfrak K_\Omega$ の中心に含まれることは明らかである。$\alpha$ を $\mathfrak K_\Omega$ の中心の任意の元とすると、その係数は $\mathsf P$ のある有限代数拡大体 $Z$（係数が生成する体）に属する。$\alpha$ はもちろん $\mathfrak K_Z$ の中心にも属する。ところがその中心は $Z$ なので、$\alpha\in Z$、したがって $\alpha\in\Omega$ である。証明された。

全く同様に次を証明する。

\textbf{定理 2.} \emph{$\mathfrak K_\Omega$ は両側単純である。}

\emph{証明。} $\mathfrak a=z_1\Omega+\cdots+z_r\Omega$ を両側イデアルとする。$z_1,\ldots,z_r$ の係数は $\mathsf P$ のある有限拡大体 $Z$ に属する。$\mathfrak K_Z$ の中で $\mathfrak a'=z_1Z+\cdots+z_rZ$ は両側イデアルである。$\mathfrak K_Z$ は両側単純なので、$\mathfrak a'$ は零または $\mathfrak K_Z$ に等しい。したがって $\mathfrak a$ は零または $\mathfrak K_\Omega$ に等しく、$\mathfrak K_\Omega$ は両側単純である。

この定理により、$\mathfrak K_\Omega$ はその中心 $\Omega$ のある拡大体上の行列環である。この体は $\mathfrak K_\Omega$ の部分環として $\Omega$ 上有限次数でなければならず、したがって $\Omega$ 自身に等しい。

\textbf{定理 3.} \emph{$\mathfrak K_\Omega$ は $\Omega$ 上の行列環である：}
\[
\mathfrak K_\Omega=\sum\Omega c_{ik}.
\]
\emph{したがって、中心 $\mathsf P$ 上の $\mathfrak K$ の次数は平方数 $m=t^2$ である。}

$\mathfrak K_\Omega$ が分解する単純右（または左）イデアルの個数 $t$ を、$\mathfrak K$ の絶対成分数と呼ぶ。

$\mathfrak K_\Omega$ が $t$ 個の単純イデアルへ分解するという $\Omega$ の性質は、$\mathsf P$ のある有限拡大 $Z$ によってすでに得られる。たとえば $c_{ik}$ の係数を $\mathsf P$ に添加して得られる $Z$ がそうである。次にこれらの分解体を調べる。

\textbf{定義。} $\mathsf P$ の有限代数拡大体 $Z$ を $\mathfrak K$ の\emph{分解体}というのは、$\mathfrak K_Z$ が分解する単純右イデアルの個数がすでに絶対成分数 $t$ に等しいときである。\srcfn{1}{9 頁の分解体の概念と比較せよ。}

\textbf{定理 4.} \emph{$Z$ が $\mathfrak K$ の分解体ならば、$Z$ は $\mathfrak K_Z$ の右イデアルの自己同型体である。逆に、$\mathfrak K_Z$ の右イデアルの自己同型体が $Z$ ならば、$Z$ は分解体である。}

\emph{証明。} 1. $Z$ を分解体、$\mathfrak T$ を $\mathfrak K_Z$ の右イデアルの自己同型体とする。すなわち
\[
\mathfrak K_{\mathsf{Z}}=\sum\mathfrak T c_{ik}.
\]
ここから
\[
\mathfrak K_\Omega=\sum\mathfrak T_\Omega c_{ik}
\]
が従う。したがって $\Omega$ 上の $\mathfrak K_\Omega$ の階数は $t^2$ と $\mathfrak T_\Omega$ の階数との積である。他方、その階数は $t^2$ に等しいから、$\mathfrak T_\Omega$、したがって $\mathfrak T$ の階数は $1$ であり、$\mathfrak T=Z$ である。

\textbf{2.} $Z$ が $\mathfrak K_Z$ の単純イデアルの自己同型体ならば
\[
\mathfrak K_{\mathsf{Z}}=\sum\mathsf Zc_{ik},
\]
したがって
\[
\mathfrak K_\Omega=\sum\Omega c_{ik}.
\]
ゆえに $c_{ik}$ の個数は $t^2$ でなければならない。証明された。

\textbf{定理 5.} \emph{$\mathfrak K$ の分解体 $Z$ の $\mathfrak K$ における既約表現 $\mathfrak Z$ は、行列環 $\mathfrak K_r$ の極大可換部分体である（22 頁の定理 2 により既約表現類は一つだけであり、$r$ は $Z$ の既約表現の次数である）。}

\emph{逆に、$\mathsf P$ の可換拡大体 $Z$ の既約表現 $\mathfrak Z$ が $\mathfrak K_r$ の極大可換部分体ならば、$Z$ は分解体である。}

\emph{証明。} 1. $\mathsf P$ 上の分解体 $Z$ の次数を $n$ とすると、$Z$ の既約表現 $\mathfrak Z$ の次数 $r$ は $n/t$ である（22 頁）。$\mathfrak Z\subseteq\mathfrak Z^*\subseteq\mathfrak K_r$ とし、$\mathfrak Z^*$ を $\mathsf P$ 上次数 $n^*\geq n$ の極大可換体とする。$\mathfrak Z^*=\mathfrak Z$ を示せばよいが、その代わりに $n^*=n$ を示せば十分である。$Z$ の拡大体で $\mathfrak Z^*$ と同型な $Z^*$ の既約表現の次数を $r^*$ とすると、再び $n^*=r^*t$ である。他方、$\mathfrak Z^*$ の中には $Z^*$ の $r$ 次表現があるので、$r$ は $r^*$ で割り切れ、$r^*\leq r$ である。したがって
\[
n\leq n^*=r^*t\leq rt=n \qquad\text{より}\qquad n^*=n.
\]

2. $\mathsf P$ の可換拡大体 $Z$ の $\mathfrak K$ における $r$ 次既約表現 $\mathfrak Z$ が、$\mathfrak K_r$ の極大可換部分体であるとする。$\mathfrak K_Z$ の単純イデアルの自己同型体を $\mathfrak T$ とし、
\[
\mathfrak K_{\mathsf{Z}}=\sum_{1}^{t'}\mathfrak T c_{ik}
\]
とする。$\mathfrak T$ は $Z$ を含む。$\mathfrak T=Z$ を示さなければならない。$\mathfrak K_r$ の単純右イデアル $\mathfrak r$ は、$Z$ の $\mathfrak K$ における表現加群であるだけでなく、$\mathfrak T$ の表現加群でもある。実際、
\[
\mathfrak r=\mathfrak T c_{i1}+\cdots+\mathfrak T c_{ik}
\]
であるから、$\mathfrak T\mathfrak r=\mathfrak r$ であり、$\mathfrak r$ は $\mathfrak T$-左加群かつ $\mathfrak K$-右加群である。

したがって $\mathfrak K_r$ には $\mathfrak T$ と同型で $\mathfrak Z$ を含む部分体 $\overline{\mathfrak T}$ がある。$\mathfrak T$ が $Z$ と異なり、したがって $\overline{\mathfrak T}$ が $\mathfrak Z$ と異なるならば、$\mathfrak Z$ に属さない $\overline{\mathfrak T}$ の元 $\alpha$ を $\mathfrak Z$ に添加すると、$\mathfrak K_r$ 内で $\mathfrak Z$ の真の可換拡大体が得られる。これは $\mathfrak Z$ の極大性に反する。証明された。

定理 5 から直ちに次が従う。

\textbf{定理 6.} \emph{$\mathfrak K$ の極大可換部分体は $\mathsf P$ 上次数 $t$ であり、分解体である。\srcfn{1}{各分解体が $\mathfrak K$ の極大可換部分体を含むという推測は誤りである。その性質をもたない分解体もありうる。そのとき対応する $r$ は必ず $1$ より大きい。R. Brauer and E. Noether, ``Über minimale Zerfällungskörper irreduzibler Darstellungen'', Sitzungsberichte d. preuss. Akademie der Wissenschaften 1927 XXXII, 221 を参照。}}

次に $\mathfrak K_r$ の部分環の同型に関する定理を導く。非可換体 $\mathfrak K$ のガロア群の構造と、それに関連する諸事実の精密な理解はこの定理に基づく。

\textbf{定理 7.} \emph{$\mathfrak S_1,\mathfrak S_2$ を、$\mathsf P$ を含む、$\mathfrak K$ 上の行列環 $\mathfrak K_r$ の両側単純部分環とし、両者の間に $\mathsf P$ を元ごとに固定する同型 $\mathfrak S_1\cong\mathfrak S_2$ があるとする。このとき、その同型は $\mathfrak K_r$ の正則元 $\tau$ による変換である。}

\emph{$\sigma_1\in\mathfrak S_1$ と $\sigma_2\in\mathfrak S_2$ が対応するならば、}
\[
\tau^{-1}\sigma_1\tau=\sigma_2.
\]

\emph{証明。} $\mathfrak S_1$ と $\mathfrak S_2$ の双方に相反同型な環 $\Sigma$ を構成し、その中の $\mathsf P$ と同型な部分体を $\mathsf P$ と同一視する。すると $\mathfrak S_1,\mathfrak S_2$ は $\Sigma$ の $\mathfrak K$ における同じ次数 $r$ の二つの相反表現である。22 頁の定理 2 により既約相反表現類は一つだけであり、したがって $\Sigma$ の $\mathfrak K$ における所定の次数の表現類も一つだけである（既約類の倍数）。ゆえに $\mathfrak S_1$ と $\mathfrak S_2$ は同じ類に属し、変換によって互いに移り合う。証明された。

定理 7 から直ちに次が従う。

\textbf{定理 8.} \emph{$\mathsf P$ を元ごとに固定する $\mathfrak K$ の自己同型群、すなわち中心に関する $\mathfrak K$ のガロア群 $\mathfrak G$ は、$\mathfrak K$ の内部自己同型群である。}

定理 7 で $r=1$、$\mathfrak S_1=\mathfrak S_2=\mathfrak K$ と置けばよい。

次の二定理は、ここまでに展開した一般理論が個別の問題に有用であることを示す。

\textbf{定理 9.} \emph{中心が実数体 $R$ である有限次数の非可換体は、四元数体だけである。}

\emph{証明。} $\mathfrak K$ を中心 $R$ の非可換体（$R$ 上有限階数）とする。定理 6 により、$R$ 上の $\mathfrak K$ の次数は、$\mathfrak K$ の極大可換部分体の次数 $t$ の平方である。この部分体は複素数体と同型な体しかありえないから、$t=2$ であり、$\mathfrak K$ は $R$ 上次数 $4$ である。$\mathfrak Z_1$ を $\mathfrak K$ の極大可換部分体とし、$i^2=-1$ として $\mathfrak Z_1=R(i)$ と置ける。$\mathfrak Z_2=R(-i)$ も $\mathfrak Z_1$ と同型であり、集合としては $\mathfrak Z_1$ と一致する体で、その同型は $i\mapsto-i$ で与えられる。したがって定理 7 により、
\[
j^{*-1}ij^*=-1
\]
を満たす $j^*\in\mathfrak K$ が存在する。

$j^*$ は $\mathfrak Z_1$ に属しえないので、$\mathfrak Z_1$ 上の $\mathfrak Z_1(j^*)$ の階数は少なくとも $2$ である。したがって $\mathfrak Z_1(j^*)$ は $R$ 上階数 $4$ で、$\mathfrak K$ に等しく、
\[
\mathfrak K=\mathsf R+\mathsf Ri+\mathsf Rj^*+\mathsf Rij^*
\]
と書ける。

\[
\begin{aligned}
j^{*-2}1j^{*2}&=1,\\
j^{*-2}ij^{*2}&=i,\\
j^{*-2}jj^{*2}&=j^*,\\
j^{*-2}ij^*j^{*2}&=ij^*
\end{aligned}
\]
であるから、$j^{*2}$ は $\mathfrak K$ のすべての元と可換で、中心 $R$ に属し、実数 $r$ である。$r$ は正ではありえない。正ならば、可換体 $R(j^*)$ の多項式 $x^2-r$ が $+\sqrt r,-\sqrt r,+j^*,-j^*$ という四つの根をもつことになるからである。したがって $j^{*2}=-m^2$、$m\in R$ である。$j=j^*m^{-1}$ と置けば $j^2=-1$ となり、$1,i,j,k=ij$ は四元数単位である。証明された。

この証明は、$R$ の代わりに任意の代数的実閉体について明らかに成り立つ。

\textbf{定理 10（Wedderburn の定理）。} \emph{有限個の元からなる体はすべて可換である。}

\emph{証明。} $\mathfrak K$ を有限体、$\mathsf P$ をその中心とし、$\mathsf P$ 上の $\mathfrak K$ の次数を $t^2$ とする。$\mathfrak K$ の極大可換部分体はすべて $\mathsf P$ 上次数 $t$ なので、同じ個数の元をもつ。したがってガロア体の理論の周知の定理（Moore の定理）により互いに同型である。さらに各ガロア体はその素体のガロア拡大なので、同型は $\mathsf P$ の各元を固定するように選べる。定理 7 により、$\mathfrak K$ の極大可換部分体はすべて、一つの $\mathfrak Z$ から変換によって得られる。一方、$\mathfrak K$ はその極大可換部分体すべての合併である（$\mathsf P$ に\emph{一つ}の元を添加すれば可換体が生じる）。ここで $\mathfrak K$ の乗法群 $\mathfrak K^*$（零を除いた $\mathfrak K$）だけを考えると、$\mathfrak K^*$ は部分群 $\mathfrak Z^*$ と共役な部分群すべての合併である。しかし $t>1$ ならば矛盾が生じる。$\mathfrak K^*$ が $M$ 個、$\mathfrak Z^*$ が $N$ 個の元をもち、$L$ を $\mathfrak K^*$ における $\mathfrak Z^*$ の指数とすると、$M=NL$ であり、相異なる共役部分群の個数は高々 $L$ である。この個数が $L$ に等しいとしても、それらの合併が群全体を尽くすには、どの二つも共通の元をもってはならない。しかしそれらはすべて単位元を共有する。

\subsection*{§ 21. 非可換体のガロア理論}\label{die-galoissche-theorie-der-nichtkommutativen-koerper}

$G$ を中心 $\mathsf P$ に関する $\mathfrak K$ のガロア群とする。定理 8 により $G$ は $\overline G=\mathfrak K^*/\mathsf P^*$ と同型であり、この同型を固定する。$G$ の部分群 $H$ を\emph{閉}と呼ぶのは、同型 $G\simeq\mathfrak K^*/\mathsf P^*$ によって対応する $\mathfrak K^*$ の部分群 $\mathfrak H^*$ に\emph{零}を加えたものが体となるときである。

\textbf{定理 11　ガロア理論の基本定理。} \emph{$G$ の閉部分群 $H$ と $\mathsf P$ と $\mathfrak K$ の中間体 $\mathfrak S$ との間には一対一対応がある。$H$ と $\mathfrak S$ が対応するとき、$H$ は $\mathfrak S$ の不変群であり、$\mathfrak S$ は $H$ の不変体である。}

\emph{証明。} 定理を三つの部分命題に分ける。
\begin{enumerate}
\item 体 $\mathfrak S$ の不変群 $H$ は閉である。
\item $\mathfrak S$ の不変群が $H$ ならば、$\mathfrak S$ は $H$ の不変体である。
\item $H$ の不変体が $\mathfrak S$ ならば、$H$ は $\mathfrak S$ の不変群である。
\end{enumerate}

\begin{enumerate}
\item 同型 $G\simeq\mathfrak K^*/\mathsf P^*$ により、体 $\mathfrak S$ の不変群 $H$ には、$\mathfrak S$ のすべての元と可換な零でない $\mathfrak K$ の元全体 $\mathfrak H^*$ が対応する。これらに $0$ を加えれば明らかに体となる。したがって $H$ は閉である。

\item 第三の部分命題を第二へ帰着させる。体 $\mathfrak S$ には $G$ の部分群 $S$、すなわち $\mathfrak S^*/\mathsf P^*$ に対応する部分群が属する。同様に群 $H$ には $\mathfrak K$ の部分体 $\mathfrak H$ が属し、$H\simeq\mathfrak H^*/\mathsf P^*$ である。$\mathfrak S$ が $H$ の不変体ならば、$\mathfrak S$ は $\mathfrak H$ の各元と可換な $\mathfrak K$ の元すべてからなる。したがって、$S$ は $\mathfrak H$ の不変群であるともいえる。同様に、$H$ が $\mathfrak S$ の不変群ならば、$\mathfrak H$ は $S$ の不変体である。ゆえに第三の部分命題は次のようにも述べられる。

\item[$(3')$.] $\mathfrak H$ の不変群が $S$ ならば、$\mathfrak H$ は $S$ の不変体である。これはまさに、なお証明すべき第二の部分命題である。

\item[3.] この証明は若干の小さな相違を除いて可換の場合と同様に進む。
\end{enumerate}

$\mathfrak K$ と相反同型な体 $K$ を作り、その中心を $\mathsf P$ と同一視する。$\mathfrak K$ の自己同型の代わりに、$\mathfrak S$ から $K$ への相反同型を考えられる。

可換の場合に用いた延長定理に類似する補題を再び証明する。$\mathfrak S,\mathfrak T$ を $\mathsf P$ と $\mathfrak K$ の中間体とし、
\[
\mathsf P \subseteqq \mathfrak S \subseteqq \mathfrak T \subseteqq \mathfrak K
\]
とする。$\mathfrak S$ から $K$ の部分体への各相反同型（$\mathfrak S$ の $K$ における一次相反表現）は、$\mathfrak S$ 上の $\mathfrak T$ の次数 $s$ と少なくとも同じ数の仕方で、$\mathfrak T$ の相反同型へ延長できると主張する。

体 $\mathsf P,\mathfrak S,\mathfrak T,\mathfrak K$ を $K$ によって拡大する：
\[
\mathsf{K}\subseteqq \mathfrak S_{\mathsf{K}}\subseteqq \mathfrak T_{\mathsf{K}}\subseteqq \mathfrak K_{\mathsf{K}}.
\]
20 頁の定理 1 により、各 $\mathfrak S_K$ は両側単純かつ完全可約であり、互いに作用素同型な単純左イデアルは、$\mathfrak S$ の $K$ におけるただ一つの既約表現類の表現加群である。この表現類は一次である。$K$ の構成により、$K$ の中には $\mathfrak S$ と相反同型な体が存在するからである。したがって単純左イデアルはすべて階数一であり、$\mathfrak S_K$ はその $\mathsf P$ 上の次数と同じ個数の左イデアルの和となる：
\[
\mathfrak S_{\mathsf{K}}=\mathsf K E_1+\cdots+\mathsf K E_p,\qquad \mathsf K E_\nu=\mathfrak S_{\mathsf{K}} E_\nu.
\]
すると
\[
\begin{aligned}
\mathfrak T_{\mathsf{K}}&=\mathfrak T\cdot \mathsf{K}=\mathfrak T\cdot \mathfrak S_{\mathsf{K}},\\
\mathfrak T_{\mathsf{K}}&=\mathfrak T\mathfrak S_{\mathsf{K}} E_1+\cdots+\mathfrak T\mathfrak S_{\mathsf{K}} E_p
              =\mathfrak T_{\mathsf{K}} E_1+\cdots+\mathfrak T_{\mathsf{K}} E_p
\end{aligned}
\]
である。各イデアル $\mathfrak T_K\cdot E_\nu$ は $s$ 個の単純イデアルへ分解する。実際、
\[
[\mathfrak T_{\mathsf{K}} E_\nu:\mathsf{K}]=[\mathfrak T_{\mathsf{K}} E_\nu:\mathsf K E_\nu]=[\mathfrak T_{\mathsf{K}} E_\nu:\mathfrak S_{\mathsf{K}} E_\nu]\leqq[\mathfrak T_{\mathsf{K}}:\mathfrak S_{\mathsf{K}}]=[\mathfrak T:\mathfrak S]=s
\]
であり、
\[
\sum [\mathfrak T_{\mathsf{K}} E_\nu:\mathsf{K}]=sp
\]
だから、
\[
[\mathfrak T_{\mathsf{K}} E_\nu:\mathsf{K}]=s,
\qquad
\mathfrak T_{\mathsf{K}} E_\nu=\mathsf K e_1^{(\nu)}+\cdots+\mathsf K e_s^{(\nu)}.
\]

$K e_i^{(\nu)}$ が与える $\mathfrak T$ の表現は、$K E_\nu$ が与える $\mathfrak S$ の表現の延長である。実際、$\alpha\in\mathfrak S$ ならば
\[
\alpha E_\nu=\sigma E_\nu,\qquad
\sum \alpha e_i^{(\nu)}=\sigma\sum e_i^{(\nu)},\qquad
\alpha e_i^{(\nu)}=\sigma e_i^{(\nu)}.
\]

残るのは、$K e_i^{(\nu)}$ が与える表現がすべて相異なることの証明だけである。可換の場合には、それらが異なる表現類に属したので明らかであった。ここでは逆にすべて同じ表現類に属するため、表現そのものを調べなければならない。$K e_i^{(\nu)}$ が与える $\mathfrak T_K$ 全体の $K$-準同型写像を考える。それらは $\mathfrak T$ の表現がわかれば完全に定まる。写像全体を構成するには、基底元がどのように表現されるかを知れば十分だからである。したがって、$K e_i^{(\nu)}$ が与える $\mathfrak T_K$ の写像がすべて異なることを示せばよい。$e_i^{(\nu)}$ を一度は表現加群 $K e_i^{(\nu)}$ により、別には $K e_j^{(\nu)}$ により表現すると、
\[
e_i^{(\nu)2}=e_i^{(\nu)}
\]
により前者では $1$ が、
\[
e_i^{(\nu)}e_j^{(\nu)}=0
\]
により後者では $0$ が、$e_i^{(\nu)}$ を表す元となる。したがって二つの表現は実際に異なる（全く同じ議論は可換の場合にも使えるが、そこでは不要である）。

ここからは 10 頁および 16 頁と全く同じように論じる。$H$ を $\mathfrak S$ の不変群、$\mathfrak T$ を $H$ の不変体とする。$H$ の元は、$\mathfrak S$ の恒等自己同型を $\mathfrak K$ 全体へ延長したものすべてである。直前の補題により、その中には $\mathfrak T$ 上で相異なるものが少なくとも $s$ 個あるから、$s=1$、$\mathfrak T=\mathfrak S$ でなければならない。証明された。


%% END NOETHER SOURCE
%% BEGIN NOETHER SOURCE "ja/t4.tex"
\section*{第 V 章　因子系}\label{kapitel-v.-faktorensysteme}

\subsection*{\S 22. 与えられた中心をもつ体の群}\label{die-gruppe-der-koerper-mit-gegebenem-zentrum}

以下の考察では、固定した可換体 \(\mathsf P\) を基礎とする。自己同型体 \(\mathfrak K\) が \(\mathsf P\) 上有限階数で、しかも \(\mathsf P\) 自身を中心とする体であるような行列環 \(\mathfrak K_r\) 全体は、「直積形成」に関して乗法的に閉じた系をなす。このような二つの環 \(\mathfrak K_r\) と \(\mathfrak L_s\) の直積とは、単に \(\mathfrak K_r\) の基礎体を \(\mathfrak L_s\) へ拡大して得られる環 \(\mathfrak K_r\cdot\mathfrak L_s\) のことである：
\[
\mathfrak K_r\times\mathfrak L_s=\mathfrak K_r\cdot\mathfrak L_s.
\]
22 頁の注意により、この直積形成は可換である：
\[
\mathfrak K_r\times\mathfrak L_s=\mathfrak L_s\times\mathfrak K_r.
\]
すべての \(\mathfrak K_r\) の系が乗法的に閉じているという主張は、したがって \(\mathfrak K_r\times\mathfrak L_s=\mathfrak K_r\mathfrak L_s\) が常に再び行列環であり、その自己同型体 \(\mathfrak M\) の中心が \(\mathsf P\) である、ということにほかならない。\(\mathsf P\) に関する有限階数性は明らかである。20 頁の定理 1 により、\(\mathfrak K_r\times\mathfrak L_s\) は両側単純かつ完全可約であり、したがって \(\mathsf P\) の拡大体 \(\mathfrak M\) 上の行列環である。\(\mathfrak M\) の中心が \(\mathsf P\) に等しいことは、第一に \(\mathsf P\) が確かに中心に含まれ、第二にその中心が \(\mathfrak K\)（または \(\mathfrak L\)）の中心 \(\mathsf P\) に含まれなければならないことから直ちに従う。積 \(\mathfrak K_r\times\mathfrak L_s\) は明らかに結合的である。\(\mathfrak K_r\) は乗法的に閉じた系をなすが、群ではない。実際、\(\mathfrak K_r\times\mathfrak L_s=\mathfrak M_m\) ならば必ず \(m\geq rs\) であり、したがって \(r>m\) のとき、方程式 \(\mathfrak K_r\times\mathfrak L_s=\mathfrak M_m\) を満たす \(\mathfrak L_s\) は見いだせない。

ここで、同じ \(\mathfrak K\) をもつすべての \(\mathfrak K_r\) を一つの類 \(\{\mathfrak K\}\) にまとめる。すると次が成り立つ。\(\mathfrak K_r\) と \(\mathfrak K_{r'}\) が同じ類 \(\{\mathfrak K\}\) に属し、\(\mathfrak L_s\) と \(\mathfrak L_{s'}\) が同じ類 \(\{\mathfrak L\}\) に属するならば、\(\mathfrak K_r\times\mathfrak L_s\) と \(\mathfrak K_{r'}\times\mathfrak L_{s'}\) も同じ類に属する。\(\mathfrak K_r\times\mathfrak L_s\) が \(\mathfrak K\times\mathfrak L\) と同じ類に属すること、すなわち \(\mathfrak K\times\mathfrak L=\mathfrak M_m\) ならば \(\mathfrak K_r\times\mathfrak L_s\) も \(\mathfrak M\) 上の行列環であることを示せば十分である。これは容易にわかる。実際、
\[
\mathfrak K_r\times\mathfrak L
=\left(\sum_1^r\mathfrak K c_{ik}\right)\mathfrak L
=\sum_1^r\mathfrak K_{\mathfrak L}c_{ik}
=\sum_1^r\mathfrak M_m c_{ik}.
\]
これは \(\mathfrak M_m\) 上の次数 \(r\) の行列環であり、もちろん \(\mathfrak M\) 上の次数 \(mr\) の行列環である。さらに全く同様に、
\[
\mathfrak K_r\times\mathfrak L_s
=\left(\sum_1^s\mathfrak L d_{ik}\right)_{\mathfrak K_r}
=\sum_1^s\mathfrak L_{\mathfrak K_r}d_{ik}
=\sum_1^s\mathfrak M_{mr}d_{ik}
=\mathfrak M_{mrs}
\]
を得て、主張は証明された。\(\mathfrak L\times\mathfrak K=\mathfrak M_m\) から \(\mathfrak K_r\times\mathfrak L_s=\mathfrak M_{mrs}\) が従うことを記憶しておこう。

以上から、二つの類 \(\{\mathfrak K\}\) と \(\{\mathfrak L\}\) の積を、\(\{\mathfrak K\}\) の代表元と \(\{\mathfrak L\}\) の代表元との積の類として定義すれば、この定義は代表元の選び方に依存しない。

したがって写像 \(\mathfrak K_r\mapsto\{\mathfrak K\}\) は、すべての \(\mathfrak K_r\) の系からすべての類 \(\{\mathfrak K\}\) の集合への準同型であり、類 \(\{\mathfrak K\}\) の集合 \(\mathscr K\) は乗法的に閉じた系である。\(\mathscr K\) が群であることを証明しよう。単位元が類 \(\{\mathsf P\}\) であることは、\(\mathfrak K_r\times\mathsf P=\mathfrak K_r\) から直ちにわかる。

そこで、各類 \(\{\mathfrak K\}\) に対して逆類 \(\{\mathfrak K\}^{-1}\) が存在することだけを示せばよい。\(\{\overline{\mathfrak K}\}\) が \(\mathfrak K\) と相反同型な体を表すとき、
\[
\{\mathfrak K\}^{-1}=\{\overline{\mathfrak K}\}
\]
と置かなければならないことがわかる。言い換えれば、\(\mathfrak K\times\overline{\mathfrak K}\) が \(\mathsf P\) 上の行列環でなければならない。あるいは、\(\mathfrak K\) が \(\mathsf P\) 上階数 \(t^2\)、したがって \(\mathfrak K\times\overline{\mathfrak K}\) が階数 \(t^4\) をもつならば、\(\mathfrak K\times\overline{\mathfrak K}\) は \(t^2\) 個の単純左イデアルに分解しなければならない。これは次のことから従う。\(\mathfrak K\times\overline{\mathfrak K}\) の単純左イデアルは、\(\overline{\mathfrak K}\) から \(\mathfrak K\) への既約な相反表現（これは次数 1 である）の表現加群として、\(\mathfrak K\) 上階数 1、したがって \(\mathsf P\) 上階数 \(t^2\) をもつ。

\(\mathsf P\) と \(\mathfrak K\) の間に、中心がやはり \(\mathsf P\) である体が存在しないとき、\(\mathfrak K\) を\emph{単純体}と呼ぶことにする。

\textbf{定理.} \emph{任意の体 \(\mathfrak K\) は単純部分体 \(\mathfrak K^{(i)}\) の直積である：}
\[
\mathfrak K=\mathfrak K^{(1)}\times\cdots\times\mathfrak K^{(p)}.
\]

\emph{証明.} \(\mathfrak K^{(1)}\) が中心 \(\mathsf P\) をもつ \(\mathfrak K\) の部分体ならば、中心 \(\mathsf P\) をもつ \(\mathfrak K\) の部分体 \(\mathfrak S\) で
\[
\mathfrak K^{(1)}\times\mathfrak S=\mathfrak K
\]
を満たすものが存在することを示せば十分である。いずれにせよ、
\[
\{\mathfrak K^{(1)}\}\times\{\mathfrak S\}=\{\mathfrak K\}
\]
という性質をもつ類 \(\{\mathfrak S\}\) は存在する。実際、
\[
\{\mathfrak S\}=\{\overline{\mathfrak K^{(1)}}\}\times\{\mathfrak K\}.
\qquad
\text{ここで}\qquad
\overline{\mathfrak K^{(1)}}\times\mathfrak K=\mathfrak S_\lambda
\]
とする。すると
\[
\mathfrak K^{(1)}\times\mathfrak S_\lambda
=\mathfrak K^{(1)}\times\overline{\mathfrak K^{(1)}}\times\mathfrak K
=\mathsf P_{t^2}\times\mathfrak K
=\mathfrak K_{t^2},
\qquad
t^2=\mathfrak K^{(1)}\text{ の階数}.
\]
ところが \(\mathfrak K^{(1)}\times\mathfrak S_\lambda\) は少なくとも \(\lambda\) に等しい次数の行列環であるから、\(\lambda\leq t^2\) である。\(\overline{\mathfrak K^{(1)}}\) は \(\mathfrak K\) において次数 1 の相反表現を許すので、\(\mathfrak S_\lambda=\overline{\mathfrak K^{(1)}}\times\mathfrak K\) は \(\mathfrak K\) 上階数 1 の単純イデアルに分解しなければならない。したがって \(\lambda\geq t^2\)、ゆえに \(\lambda=t^2\) であり、
\[
\mathfrak K_{t^2}=\mathfrak K^{(1)}\times\mathfrak S_{t^2}.
\]
それゆえ \(\mathfrak K^{(1)}\times\mathfrak S\) 自身が体でなければならず、
\[
\mathfrak K^{(1)}\times\mathfrak S_{t^2}
=\bigl(\mathfrak K^{(1)}\times\mathfrak S\bigr)_{t^2}
=\mathfrak K_{t^2}.
\]
\(\mathfrak K^{(1)}\times\mathfrak S\) は \(\mathfrak K_{t^2}\) のある分解における単純左イデアルの自己同型体であり、\(\mathfrak K\) は \(\mathfrak K_{t^2}\) の別の分解における単純左イデアルの自己同型体である。しかし \(\mathfrak K_{t^2}\) の任意の二つの異なる左イデアルは一つの分解に組み入れられるから、\(\mathfrak K\) と \(\mathfrak K^{(1)}\times\mathfrak S\) は同型であり、この同型は \(\mathsf P\) の各元を固定するように選べる。したがって、25 頁の定理 7 により、この同型は \(\mathfrak K_{t^2}\) の内部自己同型から生じる。
\[
\mathfrak K
=\lambda^{-1}\bigl(\mathfrak K^{(1)}\times\mathfrak S\bigr)\lambda
=\lambda^{-1}\mathfrak K^{(1)}\lambda\times\lambda^{-1}\mathfrak S\lambda,
\qquad
\lambda\in\mathfrak K_{t^2}.
\]
\(\lambda^{-1}\mathfrak K^{(1)}\lambda\) は \(\mathfrak K^{(1)}\) と同型な \(\mathfrak K\) の部分体であるから、25 頁の定理 7 により、適当な \(\mu\in\mathfrak K\) を用いて
\[
\mu^{-1}\lambda^{-1}\mathfrak K^{(1)}\lambda\mu=\mathfrak K^{(1)}
\]
となる。したがって
\[
\mu^{-1}\mathfrak K\mu
=\mathfrak K
=\mathfrak K^{(1)}\times\mu^{-1}\lambda^{-1}\mathfrak S\lambda\mu,
\]
となり、主張は証明された。

\subsection*{\S 23. 因子系}\label{faktorensysteme}

以下の考察では、基礎体 \(\mathsf P\) が完全であると仮定する。その代わりに、中心 \(\mathsf P\) をもつ非可換体 \(\mathfrak K\) のうち、各 \(\mathfrak K_r\) が \(\mathsf P\) 上第一種の極大可換部分体を含むものだけを考える、というより一般的な仮定を置いてもよい。（その後 Köthe が示したように、これはすべての \(\mathfrak K\) について成り立つ。）

\(Z\) を、中心 \(\mathsf P\) をもつ非可換体 \(\mathfrak K\) の分解体とする。\(Z\) から \(\mathfrak K\) への既約表現 \(\mathfrak Z\) の次数を \(r\) とすると、これは \(\mathfrak K_r\) の極大可換部分体である。\(\Gamma\) を \(Z\) のガロア体とする。すなわち
\[
\mathfrak Z_\Gamma=e_1\Gamma+\cdots+e_n\Gamma
\qquad\text{（9 頁）}
\]
である。\(n\) は \(Z\) の \(\mathsf P\) 上の次数である。すると、\(\Gamma\) によって拡大した環 \(\mathfrak K_r\) にも分解
\[
\mathfrak K_{r\Gamma}=\mathfrak K_{r\Gamma}e_1+\cdots+\mathfrak K_{r\Gamma}e_n
\]
すなわち
\[
\mathfrak K_{r\Gamma}=\mathfrak l_1+\cdots+\mathfrak l_n,\qquad
\mathfrak l_i=\mathfrak K_{r\Gamma}e_i
\]
が得られる。\(\mathfrak l_i\) は明らかに左イデアルである。しかも単純左イデアルである。実際、\(Z\) は分解体であるから、\(\mathfrak K_{r\Gamma}\) は \(\Gamma\) 上次数 \(r\cdot t=n\) の行列環である（23 頁の定理 4、および 22 頁の \(n=rt\)）。このように得た \(\mathfrak K_{r\Gamma}\) の単純左イデアルへの分解は、\(\mathfrak K\) と \(Z\) の代数的性質に特に適合している。実際、18 頁と同様に、\(\Gamma/\mathsf P\) のガロア群 \(\mathfrak G\) の置換 \(S\) をすべての \(\mathfrak K_{r\Gamma}\) の元へ作用させる際、各 \(\mathfrak K_r\) の元 \(z\) に対して
\[
S(z)=z
\]
と定めれば、イデアル \(\mathfrak l_i\) は共役である。すなわち、ある \(\mathfrak l_j\) に \(S\) を作用させると別のイデアル
\[
S(\mathfrak l_j)=\mathfrak l_i
\]
が得られる。これは、単位元の各既約成分 \(e_j\) が別の \(e_i\) へ移らなければならないからである（11 頁も参照）。

表現加群 \(e_i\Gamma\) によって、\(\mathfrak Z\) は \(\Gamma\) の部分体 \(Z_i\) へ写される。\(\mathfrak Z\) の表現 \(Z_i\) はすでに \(Z_i\) に含まれるので、\(\mathfrak Z_{Z_i}\) は表現加群 \(e_iZ_i\) を分離しなければならず、したがって \(e_i\) はすでに \(\mathfrak Z_{Z_i}\) に属する。\(\{Z_i,Z_k\}\) を \(Z_i\) と \(Z_k\) の合成体とする。\(\mathfrak K_{r\{Z_i,Z_k\}}\) は二つの左イデアル
\[
\bar{\mathfrak l}_i=\mathfrak K_{i\{\mathsf{Z}_i,\mathsf{Z}_k\}}e_i,\qquad
\bar{\mathfrak l}_k=\mathfrak K_{r\{\mathsf{Z}_i,\mathsf{Z}_k\}}e_k
\]
を分離し、これらは \(\mathfrak l_i=\mathfrak K_{r\Gamma}\bar{\mathfrak l}_i\) であるため単純である。したがって
\[
\bar{\mathfrak l}_i=\sum_{\lambda=1}^m\{\mathsf{Z}_i,\mathsf{Z}_k\}c_{\lambda i}^{(ik)},\qquad
\bar{\mathfrak l}_k=\sum_{\lambda}\{\mathsf{Z}_i,\mathsf{Z}_k\}c_{\lambda k}^{(ik)}.
\]
ここで \(c_{\lambda\rho}^{(ik)}\) は \(\mathfrak K_{r\{Z_i,Z_k\}}\) の行列単位系である。

\(\mathfrak K_{r\{Z_i,Z_k\}}\) を作用素領域として、\(\bar{\mathfrak l}_i\) から \(\bar{\mathfrak l}_k\) への各同型 \(a_i\mapsto a_k\) は、\(e_i\mapsto p_{ik}\) となる元 \(p_{ik}\) によって特徴づけられる。実際、そのとき \(a_i=a_ie_i\mapsto a_ip_{ik}\)、したがって \(\bar{\mathfrak l}_i\cdot p_{ik}=\bar{\mathfrak l}_k\) である。このような \(p_{ik}\) はすべて明らかに
\[
p_{ik}=c_{ik}^{(ik)}\gamma_{ik},\qquad
\gamma_{ik}\ne0\quad\text{かつ}\quad \gamma_{ik}\in\{Z_i,Z_k\}
\]
である。\(a_i\mapsto a_ip_{ik}\) によって、\(\mathfrak l_i\) から \(\mathfrak l_k\) への作用素同型も与えられる。ここで、このような \(p_{ik}\) の特定の系を、共役性条件を課して選ぶ。

\(S\) の添字 \(i\) への作用を
\[
S(i)=k,\qquad\text{ただし}\qquad S(\mathsf{Z}_i)=\mathsf{Z}_k
\]
によって定める。各添字対 \((\nu,\mu)\) には、\(S(\nu,\mu)=(1,k)\) という形の共役な対が少なくとも一つ存在する。共役な添字対の各類から固定した対 \((1,k)\) を一つ選び、この \((1,k)\) に対して
\[
p_{1k}=c_{1k}^{(1k)}\gamma_{1k},\qquad
\gamma_{1k}\ne0\quad\text{かつ}\quad \gamma_{1k}\in\{Z_1,Z_k\}
\]
を、それ以外は任意に定める。そして \(S(1,k)=(\nu,\mu)\) のとき
\[
p_{\nu\mu}=S(p_{1k})
\]
と置く。この \(p_{\nu\mu}\) は実際
\[
p_{\nu\mu}=c_{\nu\mu}^{(\nu\mu)}\gamma_{\nu\mu},\qquad
\gamma_{\nu\mu}\ne0\quad\text{かつ}\quad \gamma_{\nu\mu}\in\{Z_\nu,Z_\mu\}
\]
という形をもつ。なぜなら、\(\bar{\mathfrak l}_1\) から \(\bar{\mathfrak l}_k\) への同型
\[
a_1\mapsto a_1p_{1k},\qquad
\text{特に}\qquad e_1\mapsto p_{1k}
\]
に \(S\) を作用させると、\(\bar{\mathfrak l}_\nu\) から \(\bar{\mathfrak l}_\mu\) への同型
\[
a_\nu\mapsto a_\nu p_{\nu\mu},\qquad
\text{特に}\qquad e_\nu\mapsto p_{\nu\mu}
\]
が得られるからである。

\(c_{\lambda\rho}\) を \(\mathfrak K_{r\Gamma}\) の行列単位系とすると、\(p_{\nu\mu}\) は
\[
p_{\nu\mu}=c_{\nu\mu}\bar\gamma_{\nu\mu}
\]
とも書けることは明らかである。したがって関係式
\[
c_{ij}c_{1k}=0,\quad\text{もし }\quad j\ne1,
\qquad
c_{ij}c_{jk}=c_{ik}
\]
から、関係式
\[
\begin{aligned}
p_{ij}p_{1k}&=0,\quad &&\text{もし }j\ne1,\\
p_{ij}p_{jk}&=\alpha_{ik}^{(j)}p_{ik},
&&\alpha_{ik}^{(j)}\in\Gamma
\quad\text{（より正確には }\alpha_{ik}^{(j)}\in\{Z_i,Z_j,Z_k\}\text{）}
\end{aligned}
\]
が従う。\(n^3\) 個の量 \(\alpha\) は、「擬行列単位」系 \(p_{ik}\) に属する、\emph{\(Z\) を基礎とした \(\mathfrak K\) の因子系}をなす。

\(\alpha\) は共役である。\(S(i,k,j)=(\nu,\mu,\tau)\) ならば
\[
S\bigl(\alpha_{ik}^{(j)}\bigr)=\alpha_{\nu\mu}^{\tau}.
\]

擬行列単位の任意の系 \(p_{ik}^*\) は、最初の \(p_{ik}\) に \(\{Z_i,Z_k\}\) の任意の \(\delta_{ik}\ne0\) を掛けることによって、ある系 \(p_{ik}\) から得られる。すなわち
\[
p_{1k}^*=p_{1k}\delta_{1k}
\]
とし、さらに \(S(1,k)=(\nu,\mu)\) のとき
\[
p_{\nu\mu}^*=S(p_{1k}^*)
\]
と置く。したがって
\[
p_{\nu\mu}^*=p_{\nu\mu}\delta_{\nu\mu},
\qquad \delta_{\nu\mu}\in\{Z_\nu,Z_\mu\}.
\]
\(\delta_{\nu\mu}\) は共役であり、\(S(i,k)=(\nu,\mu)\) ならば \(S(\delta_{ik})=\delta_{\nu\mu}\) である。\(\alpha^*\) を \(p_{ik}^*\) に属する因子系とすれば、
\[
\alpha_{ik}^{(i)*}=\frac{\delta_{ij}\delta_{jk}}{\delta_{ik}}\alpha_{ik}^{(j)}.
\]

因子系 \(\alpha^*\) と \(\alpha\) を\emph{同伴する}という（\(p_{ik}^*\) と \(p_{ik}\) についても同様である）。\(Z\) を基礎とした \(\mathfrak K\) のすべての因子系は、同伴する因子系の一つの類 \(\{\alpha\}\) をなす。

類 \(\{\alpha\}\) は、その定義の基礎とする \(Z\) から \(\mathfrak K_r\) への表現 \(\mathfrak Z\) の選び方に依存しない。実際、ある \(\mathfrak Z\) から別の \(\mathfrak Z'\) への移行は、\(\mathfrak K_r\) の元、したがって \(\mathfrak G\)-不変な元 \(\lambda\) による変換で行えるので、\(\mathfrak Z\) に属する系 \(p_{ik}\) から、\(\mathfrak Z'\) に属する系
\[
p'_{ik}=\lambda^{-1}p_{ik}\lambda
\]
が得られる。そして \(p_{ik}\) と \(p'_{ik}\) の因子系は同一である。

因子系とは \(\mathfrak K_{n\Gamma}\) の乗法表にほかならないが、\(\mathfrak K\) の特別な代数的性質に適合した乗法表である。

\subsection*{\S 24. 因子系の乗法}\label{multiplikation-von-faktorensystemen}

非可換体 \(\mathfrak K\) の分解体という代わりに、\(\{\mathfrak K\}\) の分解体ともいう。この言い方には意味がある。というのも、\(\mathfrak K_Z\) とともに各環 \(\mathfrak K_{rZ}\) も絶対既約成分に分解し、逆もまた成り立つからである。

中心 \(\mathsf P\) をもつ任意の二体 \(\mathfrak K,\mathfrak S\) は共通の分解体をもつ。例えば、\(\mathfrak K\) の分解体と \(\mathfrak S\) の分解体との合成体がそれである。

\textbf{定理 1.} \emph{ある体 \(Z\) を分解体にもつ \(\{\mathfrak K\}\) は、すべての \(\{\mathfrak K\}\) の群 \(\mathscr K\) の部分群 \(\mathscr K_Z\) をなす。}

\emph{証明.} \(\mathfrak K\) と \(\mathfrak S\) が \(Z\) を分解体にもつならば、\(\mathfrak K_Z\) と \(\mathfrak S_Z\)、したがって \((\mathfrak K\times\mathfrak S)_Z\) も絶対既約成分に分解する。すなわち \(Z\) は \(\{\mathfrak K\times\mathfrak S\}\) の分解体である（23 頁の定理 4）。

相反同型な体はすべての分解体を共有する。

補助的考察を行う。
\[
\mathfrak K_r=\mathfrak L_1+\cdots+\mathfrak L_r,\qquad
\mathsf P_m=\mathfrak z_1+\cdots+\mathfrak z_m
\]
を、\(\mathfrak K_r\) と \(\mathsf P_m\) の単純左イデアルへの分解とする。このとき
\[
\mathfrak K_{rm}=\mathfrak K_r\times\mathsf P_m
=\mathfrak L_1\mathfrak z_1+\cdots+\mathfrak L_i\mathfrak z_j+\cdots+\mathfrak L_r\mathfrak z_m
\]
は、\(\mathfrak K_{rm}\) を零でない \(rm\) 個、したがって単純な左イデアルへ分解したものである。それゆえ、\(\mathfrak K_{rm}\) の任意の \(r\)-成分左イデアル \(A\) は、\(r\)-成分左イデアル
\[
\mathfrak L_1\mathfrak z_1+\cdots+\mathfrak L_r\mathfrak z_1=\mathfrak K_r\mathfrak z_1
\]
と \(\mathfrak K_r\) の内部自己同型によって結ばれている：
\[
\mathsf{A}=\tau^{-1}\mathfrak K_r\tau\cdot\tau^{-1}\mathfrak z_1\tau
=\mathfrak K'_r\cdot\mathfrak z'_1.
\]

\(\mathfrak K_{r\Gamma}=\mathfrak l_1+\cdots+\mathfrak l_n\) を、\(\mathfrak K_r\) 内のある \(\mathfrak Z\) に属する、\(\mathfrak K_{r\Gamma}\) の共役な左イデアルへの分解とし、\(\mathfrak l_i=\mathfrak K_{r\Gamma}e_i\) とする。このとき
\[
S(\tau^{-1}\mathfrak l_i\tau\cdot\mathfrak z'_1)
=S(\tau^{-1}\mathfrak l_i\tau)\,\mathfrak z'_1
=\tau^{-1}(S\mathfrak l_i)\tau\cdot\mathfrak z'_1
=\tau^{-1}\mathfrak l_k\tau\cdot\mathfrak z'_1
\]
であるから、
\[
\tag{1}
\mathsf{A}_\Gamma=\mathfrak l'_1\mathfrak z'_1+\cdots+\mathfrak l'_n\mathfrak z'_1
\]
は \(A_\Gamma\) の共役な左イデアルへの分解である。\(\tau^{-1}e_i\tau=e'_i\) と置き、\(E'\) を \(\mathfrak z'_1\) の単位元とすれば、
\[
\mathfrak l'_i\mathfrak z'_1=\mathfrak K_{rm\Gamma}\cdot e'_iE'
\]
である。

共役な左イデアルへの任意の別の分解
\[
\tag{2}
\mathsf{A}_\Gamma=\mathfrak q_1+\cdots+\mathfrak q_n,\qquad
S\mathfrak q_i=\mathfrak q_k\quad\text{もし}\quad
S\mathfrak l'_i\mathfrak z'_1=\mathfrak l'_k\mathfrak z'_1
\]
は、ある \(\mathfrak K_{rm}\) の元 \(\lambda\) による変換によって (1) から得られる：
\[
\mathfrak q_i=\lambda^{-1}\mathfrak l'_i\mathfrak z'_1\lambda,
\qquad
\text{したがって、\(e_i^*\) を \(\mathfrak q_i\) の単位元とすれば}\quad
e'_i=\lambda^{-1}e'_iE'\lambda.
\]

その証明は、次の補題から直ちに従う。

\textbf{補題.} \emph{\(A\) を \(\mathfrak K_n\) の左イデアルとし、}
\[
\mathsf{A}=\mathfrak l_1+\cdots+\mathfrak l_n
=\mathfrak K_{n\Gamma}e_1+\cdots+\mathfrak K_{n\Gamma}e_n
\]
\emph{を、仮定により存在する共役な左イデアルへの分解とする。この際、一般性を失うことなく \(\mathfrak l_i\) を共役と仮定してよい。}

\emph{\(\mathfrak z\) が \(A\) と作用素同型であり、同じ仮定が満たされているとする：}
\[
\mathfrak z=\mathfrak q_1+\cdots+\mathfrak q_n
=\mathfrak K_{n\Gamma}e_1^*+\cdots+\mathfrak K_{n\Gamma}e_n^*
\]
\emph{（したがって \(\mathfrak q_i\) と \(e_i^*\) は共役である）。このとき \(\mathfrak K_n\) の正則元 \(\lambda\) で、\(\mathfrak q_i=\lambda^{-1}\mathfrak l_i\lambda\) を満たすものが存在する。}

\emph{証明. 1. \(e_i\) を共役と仮定してよいことは次のように従う。} \(e_1\) を \(\mathfrak l_1\) の任意の生成冪等元、すなわち \(\mathfrak l_1=\mathfrak K_{n\Gamma}e_1\) とする。\(\Gamma\) の自己同型を作用させれば \(\mathfrak l_i=\mathfrak K_{n\Gamma}e_i\) が従い、ここで \(e_i\) は \(e_1\) と共役で、やはり生成冪等元である。

2. \emph{任意の \(\mathfrak l\) は任意の \(\mathfrak q\) と同型である。} 実際、\(\mathfrak l\) は共役であるためすべて同じ（絶対）長をもち、\(\mathfrak z\) と \(A\) は同型と仮定され、しかも \(\mathfrak l\) と \(\mathfrak q\) の個数は一致するからである。

3. 対応
\[
\tag{3}
e_1\mapsto s_1=e_1s_1;\qquad
e_1^*\mapsto t_1=e_1^*t_1;\qquad
s_1t_1=e_1;\qquad
t_1s_1=e_1^*
\]
によって、\(\mathfrak l_1\) と \(\mathfrak q_1\) の間の同型とその逆を固定する。すべての自己同型を作用させることにより、(3) は
\[
e_i\mapsto s_i=e_is_i;\qquad
e_i^*\mapsto t_i=e_i^*t_i;\qquad
s_it_i=e_i;\qquad
t_is_i=e_i^*
\]
へ移る。こうして同型 \(\mathfrak q_i=\mathfrak l_is_i\) と \(\mathfrak l_i=\mathfrak q_it_i\) が定められる。

4. ここで
\[
s=s_1+\cdots+s_n;\qquad t=t_1+\cdots+t_n
\]
と置く。すると
\[
\mathfrak q_i=\mathfrak l_is;\qquad
\mathfrak l_i=\mathfrak q_it;\qquad
st=e_1+\cdots+e_n=\xi;\qquad
ts=e_1^*+\cdots+e_u^*=\xi^*
\]
を得る。ここで \(\xi\) と \(\xi^*\) は \(\mathfrak K_n\) に属し、
\[
\mathsf{A}\xi=\mathsf{A},\qquad \mathfrak z\xi^*=\mathfrak z\quad
\quad\text{（}\mathfrak l_i\xi=\mathfrak l_i\mathfrak l_i=\mathfrak l_i\text{ であるため）},
\]
\[
\xi^2=\xi,\qquad \xi^{*2}=\xi^*
\]
となる。いま \(e=\overline E+E=\overline E^*+E^*\) と置き、したがって
\[
\mathfrak K_n=\overline{\mathsf A}+\mathsf{A}
=\mathfrak K_n\overline\xi+\mathfrak K_n\xi
=\overline{\mathfrak z}+\mathfrak z
=\mathfrak K_n\overline\xi^*+\mathfrak K_n\xi^*
\]
とすれば、さらに \(\overline A\) と \(\overline{\mathfrak z}\) は同型になる。そこで
\[
\overline E\mapsto\overline s=\overline E\,\overline s
\qquad\text{および}\qquad
\overline E^*\mapsto\overline t=\overline E^*\overline t
\]
を相互に逆な同型とし、
\[
\lambda=s+\overline s,\qquad \mu=t+\overline t
\]
と置けば、
\[
\lambda\mu=\mu\lambda=e,\qquad
\mathfrak l_i\lambda=\mathfrak q_i,\qquad
\lambda^{-1}\mathfrak q_i=\mathfrak q_i
\quad
(\lambda^{-1}\mathfrak q_i\subseteqq\mathfrak q_i
=\lambda\cdot\lambda^{-1}\mathfrak q_i\leqq\lambda^{-1}\mathfrak q_i)
\]
となり、したがって
\[
\lambda^{-1}\mathfrak l_i\lambda=\mathfrak q_i,\qquad \text{証明終。}
\]

分解 (1) に対して、再び元
\[
r_{ik}=c_{ik}^{\prime(ik)}\gamma_{ik},\qquad \gamma_{ik}c\{\mathsf{Z}_i\mathsf{Z}_k\}
\]
を定めることができる。これらは \(e'_iE'\mapsto r_{ik}\) によってイデアル \(\mathfrak l'_i\mathfrak z'_1\) と \(\mathfrak l'_k\mathfrak z'_1\) の間の作用素同型を与え、また \(S(i,k)=(\nu,\mu)\) のとき共役性条件 \(S(r_{ik})=r_{\nu\mu}\) を満たす。これらは
\[
r_{ij}r_{jk}=\alpha_{ik}^{(j)}r_{ik}
\]
によって \(A_\Gamma\) の因子系 \(\alpha\) を定める。ここで重要なのは、\(A_\Gamma\) の因子系を定めるために任意の分解 (2) も同様に使えることである。実際、(1) に属する系 \(r_{ik}\) を \(r'_{ik}=\lambda^{-1}r_{ik}\lambda\) で置き換えると、\(e_i^*\mapsto r'_{ik}\) は \(\mathfrak q_i\) から \(\mathfrak q_k\) への作用素同型を与える。\(\lambda\) は \(\mathfrak q\)-不変なので、\(r'_{ik}\) は共役性条件を満たし、しかも \(r'_{ik}\) は \(r_{ik}\) と同じ因子系をもつ。

\(\mathfrak K\) の擬行列単位系 \(p_{ik}\) を取れば、
\[
r_{ik}=\tau^{-1}p_{ik}\tau\cdot E'
\]
は同じ因子系をもつ \(A_\Gamma\) の \(r_{ik}\)-系である。したがって次の補助定理を得る。

\textbf{定理.} \emph{\(A_\Gamma\) の共役な左イデアルへの任意の分解}
\[
\mathsf{A}_\Gamma=\mathfrak q_1+\cdots+\mathfrak q_n
\]
\emph{を用いて作られる \(A_\Gamma\) の因子系は、\(\mathfrak K_r\) の因子系と同じである。}

次の定理は、\(\alpha_{ik}^{(j)}\) と \(\beta_{ik}^{(j)}\) がそれぞれ、\(Z\) を基礎とする \(\{\mathfrak K\}\) と \(\{\mathfrak S\}\) の因子系ならば、\(\varepsilon_{ik}^{(j)}=\alpha_{ik}^{(j)}\cdot\beta_{ik}^{(j)}\) が、\(Z\) を基礎とする \(\{\mathfrak K\times\mathfrak S\}\) の因子系であることを示す。したがって、同伴する因子系の二類 \(\{\alpha\}\)、\(\{\beta\}\) の\emph{積 \(\{\alpha\}\times\{\beta\}\)} を、類 \(\{\alpha_{ik}^{(j)}\beta_{ik}^{(j)}\}\) として定義することには意味がある。

\textbf{定理 2.} \emph{\(\{\mathfrak K\}\) が、\(Z\) を基礎として同伴する因子系の類 \(\{\alpha\}\) をもつとする。分解体 \(Z\) をもつすべての \(\{\mathfrak K\}\) の類 \(\{\alpha\}\) は群をなし、\(\{\mathfrak K\}\mapsto\{\alpha\}\) によって \(\mathscr K_Z\) と同型に対応する。}

\textbf{証明のために}次の二点を示さなければならない。

1. \(\{\mathfrak K\}\) が因子系 \(\alpha_{ik}^{(j)}\) を、\(\{\mathfrak S\}\) が \(\beta_{ik}^{(j)}\) をもつならば、\(\varepsilon_{ik}^{(j)}=\alpha_{ik}^{(j)}\beta_{ik}^{(j)}\) は \(\{\mathfrak K\times\mathfrak S\}\) の因子系である。

2. \(\{\mathfrak K\}\) が因子系 \(\alpha_{ik}^{(j)}=1\) をもつならば、\(\mathfrak K=\mathsf P\) である。

1. \(Z\) から \(\mathfrak K\)、\(\mathfrak S\) への既約表現の次数をそれぞれ \(r,s\) とし、
\[
\mathfrak K_{r\Gamma}=\mathfrak l_1+\cdots+\mathfrak l_n,\qquad
\mathfrak S_{s\Gamma}=\mathfrak m_1+\cdots+\mathfrak m_n
\]
を共役なイデアルへの分解とする。ここで \(\mathfrak l_i\) と \(\mathfrak m_i\) は互いに対応するイデアルとする（それぞれ \(\mathfrak Z_r\) および \(\mathfrak Z'_r\) の対応する成分から導かれる）。すると
\[
(\mathfrak K_r\times\mathfrak S_s)_\Gamma
=\mathfrak l_1\mathfrak m_1+\cdots+\mathfrak l_i\mathfrak m_j+\cdots+\mathfrak l_n\mathfrak m_n
\]
は、零でない \(n^2\) 個、したがって単純なイデアルへの分解である。それゆえ
\[
\mathsf{A}_1=\mathfrak l_1\mathfrak m_1+\mathfrak l_2\mathfrak m_2+\cdots+\mathfrak l_n\mathfrak m_n
\]
は絶対長 \(n\) のイデアルである。\(\mathfrak m_i\) は \(\mathfrak l_i\) と同じように変換されるので、\(A_1\) はすべての \(S\) に対して不変である。したがって
\[
\mathsf{A}_1=\mathsf{A}_\Gamma
\]
となる \(\mathfrak K_r\times\mathfrak S_s\) のイデアル \(A\) が存在する（18 頁の補題）。

\(\mathfrak K\times\mathfrak S\) の単純イデアルの自己同型体 \(\mathfrak M\) における \(\mathfrak Z\) の既約表現の次数を \(u\) とすると、\(A\) は \(u\) 個の単純イデアルに分解する（絶対長が \(n\) だからである）。したがって補助定理により、\(A\) の因子系は \(\mathfrak M\) の因子系でもある。

ところで、定義に用いる \(A_\Gamma\) の共役なイデアルへの分解は
\[
\tag{3}
\mathsf{A}_\Gamma=\mathfrak l_1\mathfrak m_1+\mathfrak l_2\mathfrak m_2+\cdots+\mathfrak l_n\mathfrak m_n
\]
である。\(\alpha_{ik}^{(j)}\)、\(\beta_{ik}^{(j)}\) がそれぞれ \(\mathfrak K\)、\(\mathfrak S\) の擬行列単位 \(p_{ik}\)、\(q_{ik}\) に属するならば、\(r_{ik}=p_{ik}q_{ik}\) は明らかに、分解 (3) に対する \(A_\Gamma\) の擬行列単位系である。これに属する因子系 \(\varepsilon_{ik}^{(j)}=\alpha_{ik}^{(j)}\beta_{ik}^{(j)}\) は、したがって \(\mathfrak M\) の因子系でもある。

2. \(\{\mathfrak K\}\) が因子系 \(\alpha_{ik}^{(j)}=1\) をもち、対応する擬行列単位を \(p_{ik}\)、共役なイデアルへの分解を
\[
\mathfrak K_{r\Gamma}=\mathfrak l_1+\cdots+\mathfrak l_n
\]
とする。\(\mathfrak K=\mathsf P\)、または \(\mathfrak K_r=\mathsf P_r\)、または \(\mathfrak K_r\) に \(n\)-成分左イデアル \(\mathfrak l\) が存在することを示さなければならない。なぜなら、単純な \(\mathfrak K_r\)-左イデアルは階数 \(t^2r\) をもつので、整数 \(h\) によって \(n=t^2rh\) となるか、または \(t=1\) でなければならないからである。

\(\mathfrak K_{r\Gamma}\)-イデアル
\[
\mathfrak L=\mathfrak l_1(p_{11}+\cdots+p_{1n})
\]
は零でなく（\(e_1\sum_i p_{1i}\ne0\)）、単純であり、したがって階数 \(n\) である。\(\Omega\) は \(\mathfrak G\)-不変である。実際、\(\alpha_{ik}^{(j)}=1\) であるから、
\[
S\mathfrak L
=\mathfrak L\mathfrak l_1\sum_\nu S p_{1\nu}
=\mathfrak l_i\sum_\mu p_{i\mu}
=\mathfrak l_i\sum_\mu p_{i1}p_{1\mu}.
\]
したがって
\[
S\mathfrak L=\mathfrak l_ip_{i1}\sum_\mu p_{1\mu}
=\mathfrak l_1\sum_\mu p_{1\mu}
=\mathfrak L.
\]
ゆえに \(\Omega=\mathfrak l_\Gamma\) であり、\emph{ここで \(\mathfrak l\) は階数 \(n\) の \(\mathfrak K_r\)-イデアルである。}

\textbf{定理 3.} \emph{群 \(\mathscr K\) の各元 \(\{\mathfrak K\}\) は有限位数 \(\lambda\) をもち、それは \(\mathfrak K\) の絶対成分数 \(t\) の約数である。}

\emph{第一証明（Brauer）.} \(\mathfrak K_{r\Gamma}\) の共役な左イデアルへの二つの分解
\[
\mathfrak K_{r\Gamma}=\mathfrak l_1+\cdots+\mathfrak l_n
=\mathfrak z_1+\cdots+\mathfrak z_n,\qquad
\mathfrak l_i=\mathfrak K_{r\Gamma}e_i,\quad
\mathfrak z_i=\mathfrak K_{r\Gamma}e_i^*
\]
を取る。\(a_i\mapsto a_ir_{ik}\)、特に \(e_i\mapsto r_{ik}\) によって、\(\mathfrak l_i\) から \(\mathfrak z_k\) への作用素同型が与えられるとする。その逆写像で \(e_k^*\) は \(\mathfrak l_i\) の元 \(s_{ki}\) へ移り、\(s_{ki}\cdot r_{ik}=e_k^*\) となる。\(r_{ik}\) を共役に選び、
\[
S(r_{ik})=r_{\nu\mu}\quad\text{もし}\quad S(i,k)=(\nu,\mu)
\]
とすれば、\(s_{ik}\) も共役である。実際、
\[
S(s_{ki})S(r_{ik})=S(s_{ki}r_{ik})=S(e_k^*)=e_\mu^*
=s_{\mu\nu}r_{\nu\mu}=s_{\mu\nu}S(r_{ik}),
\]
したがって \(S(s_{ki})=s_{\mu\nu}\) である。\(e_i\mapsto e_ir_{ij}s_{jk}\) によって \(\mathfrak l_i\) から \(\mathfrak l_k\) への作用素同型が与えられるから、
\[
r_{ij}s_{jk}=p_{ik}\xi_{ik}^{(j)},\qquad
\xi_{ik}^{(j)}\ne0\in\Gamma
\]
でなければならない。

\(r\)、\(s\)、\(p\) が共役であるから、\(\xi\) も共役である。\(S(i,j,k)=(\nu,\tau,\mu)\) ならば
\[
S(\xi_{ik}^{(j)})=\xi_{\nu\mu}^{\tau}.
\]

\[
r_{i\lambda}s_{\lambda j}r_{j\lambda}\varepsilon_{\lambda k}
=r_{i\lambda}e_\lambda^*s_{\lambda k}
=r_{i\lambda}s_{\lambda k}
\]
から
\[
p_{ij}\xi_{ij}^{(\lambda)}\cdot p_{jk}\xi_{ik}^{(\lambda)}
=p_{ik}\xi_{ik}^{(\lambda)}
\]
が従う。あるいは \(p_{ij}p_{jk}=p_{ik}\alpha_{ik}^{(j)}\) であるから、
\[
\alpha_{ik}^{(j)}
=\frac{\xi_{ik}^{(\lambda)}}{\xi_{ij}^{(\lambda)}\xi_{jk}^{(\lambda)}}.
\]
\(\delta_{\nu\mu}=\prod_\lambda\xi_{\nu\mu}^{(\lambda)}\) と置けば、\(\delta_{\nu\mu}\) は \(\nu\) と \(\mu\) を固定するすべての \(S\) に対して不変であるから \(\{Z_\nu,Z_\mu\}\) に属し、さらに \(\xi_{\nu\mu}^{(j)}\) が共役であるため \(\delta_{\nu\mu}\) も共役である。いま
\[
\alpha_{ik}^{(j)n}
=\frac{\prod_\lambda\xi_{ik}^{(\lambda)}}
       {\prod_\lambda\xi_{ij}^{(\lambda)}\prod_\lambda\xi_{jk}^{(\lambda)}}
=\frac{\delta_{ik}}{\delta_{ij}\delta_{jk}}
\]
であるから、同伴する因子系の公式（32 頁）により、\(\alpha_{ik}^{(j)n}\) は系 \(\beta_{ik}^{(j)}=1\) と同伴する。したがって \(\{\mathfrak K\}\) の位数は \(\{\alpha\}\) の位数に等しいので、\(\{\mathfrak K\}\) の位数は \(n\) の約数である。特に \(\mathfrak Z\) を \(\mathfrak K\) の極大可換部分体として選べば \(n=t\) であるから、\(\lambda\) は \(t\) の約数である。

\emph{第二証明（Schur）.} この証明は、\(p_{ik}\) を \(\Gamma\) 上の正則な \(n\) 次行列 \(P_{ik}\) で表現することに基づく。この行列は \(p_{ik}\) と同じ共役性条件を満たす。\(P_{ik}\) は \(\{Z_i,Z_k\}\) に属し、\(p_{ik}\) に \(P_{ik}\) が対応するならば、\(S(p_{ik})\) には \(S(P_{ik})\) が対応する。この表現は準同型である。すなわち \(p_{ik}\) に \(P_{ik}\)、\(p_{kj}\) に \(P_{kj}\) が対応するならば、\(p_{ik}p_{kj}\) に \(P_{ik}P_{kj}\) が対応し、したがって
\[
P_{ik}P_{kj}=\alpha_{ij}^{(k)}P_{ij}
\]
となる（\(k_1\ne k_2\) に対して \(P_{ik_1}P_{k_2j}=0\) となることはありえない）。この \(p_{ik}\) の表現が得られれば、次のように結論する。行列式を取れば
\[
|P_{ik}|\cdot |P_{kj}|=\alpha_{ij}^{(k)n}\cdot |P_{ij}|.
\]
行列 \(P_{ik}\) は正則なので \(|P_{\nu\mu}|\ne0\) であり、
\[
\alpha_{ij}^{(k)n}=\frac{\delta_{ik}\delta_{kj}}{\delta_{ik}},
\]
ここで \(|P_{\nu\mu}|=\delta_{\nu\mu}\) と置いた。その後は Brauer の証明と同じである。

\emph{行列 \(P_{ik}\) による \(p_{ik}\) の表現：}

分解
\[
\mathfrak K_{r\Gamma}=\mathfrak l_1+\cdots+\mathfrak l_n
\]
のイデアル \(\mathfrak l_i\) は、\(\Gamma\) における \(\mathfrak K_r\) の相反表現加群である（32 頁の定理 2）。\(\mathfrak Z\) に関する \(\mathfrak K_r\) の固定した基底 \(k_1,\ldots,k_n\) を基礎とすれば、
\[
\mathfrak K_{r\Gamma}=k_1\mathfrak Z_\Gamma+\cdots+k_n\mathfrak Z_\Gamma,
\]
したがって
\[
\mathfrak l_i=\mathfrak K_{r\Gamma}e_i
=k_1\mathfrak Z_\Gamma e_i+\cdots+k_n\mathfrak Z_\Gamma e_i
=k_1e_i\Gamma+\cdots+k_ne_i\Gamma.
\]
これから
\[
\mathfrak l_k=\mathfrak l_i p_{ik}
=k_1p_{ik}\Gamma+\cdots+k_np_{ik}\Gamma
\]
が従う。\(\mathfrak l_k\) の二基底 \((k_1e_i,\ldots,k_ne_i)\) と \((k_1p_{ik},\ldots,k_np_{ik})\) は、正則行列 \(P_{ik}\) を掛けることによって互いに移される：
\[
P_{ik}\cdot
\begin{pmatrix}
k_1e_i\\ \vdots\\ k_ne_i
\end{pmatrix}
=
\begin{pmatrix}
k_1p_{ik}\\ \vdots\\ k_np_{ik}
\end{pmatrix},
\qquad
\text{略記すれば}\quad
P_{ik}(k_\varrho e_i)=(k_\varrho p_{ik}).
\]
これから \(P_{ik}\) の共役性が従う：
\[
S(P_{ik})(k_\varrho S(e_k))=(k_\varrho S(p_{ik}))
\quad\text{または}\quad
S(P_{ik})(k_\varrho e_\mu)=(k_\varrho p_{\nu\mu})
=P_{\nu\mu}(k_\varrho e_\mu).
\]
すなわち
\[
S(P_{ik})=P_{\nu\mu}\quad\text{もし}\quad S(i,k)=(\nu,\mu).
\]
これからさらに、\(P_{ik}\) が \(\{Z_i,Z_k\}\)、すなわち \(i\) と \(k\) を固定する \(S\) の不変体に属することが従う。準同性も容易にわかる：
\[
\begin{split}
P_{ik}P_{kj}(k_\varrho e_j)
&=P_{ik}(k_\varrho p_{kj})
=(k_\varrho p_{ik}p_{kj})\\
&=\alpha_{ij}^{(k)}(k_\varrho p_{ij})
=\alpha_{ij}P_{ij}(k_\varrho e_j),
\quad\text{したがって}\\
&\cdot\, P_{ik}P_{kj}=\alpha_{ij}^{(k)}P_{ij},
\quad\text{証明終。}
\end{split}
\]

定理 3 は次のように強められる。

\textbf{定理 4.} \emph{\(p\) が \(\mathfrak K\) の絶対成分数 \(t\) の素因数ならば、\(p\) は \(\{\mathfrak K\}\) の位数 \(\lambda\) も割り切る。（Brauer）}

\emph{証明.} \(p^\sigma\) を、分解体 \(Z\) の次数 \(n\) を割る \(p\) の最高冪とし、\(\mathfrak H\) を \(\mathfrak G\) の位数 \(p^\sigma\) の部分群（これは \emph{Sylow の定理}、Speiser 第 1 版、\S 19、42 頁、定理 43 により存在する）、\(T\) を \(\mathfrak H\) の不変体とする。

このとき \([\Gamma:T]=p^\sigma\) であり、\([T:\mathsf P]=n/p^\sigma\) は \(p\) で割れず、まして \(t\) では割れない。したがって \(T\) は分解体ではなく、\(\mathfrak K_T=\mathfrak S_m\) の自己同型体 \(\mathfrak S\) は \(\mathsf P\) と異なる。しかし \(\mathfrak S\) は分解体 \(\Gamma\) をもち、その \(\mathfrak S\) の中心 \(T\) 上の次数は \(p^\sigma\) であるから、\(\mathfrak S\) の絶対成分数は \(p\) の冪 \(p^\beta>1\) である。したがって定理 3 により、\(\{\mathfrak S\}\) の位数は \(p\) の冪 \(p^\alpha>1\) である。
\[
\{\mathfrak K\}^{\lambda}=\{\mathsf P\},
\qquad
\{\mathfrak K_{\mathsf{T}}\}^{\lambda}=\{\mathsf{T}\}
\]
が従うので、\(\lambda\equiv0\pmod {p^\alpha}\) である。証明終。（\(\lambda=t\) である必要がないことについて、Brauer は例を与えた\srcfn{1}{Brauer: Untersuchungen über die arithmetischen Eigenschaften von Gruppen linearer Substitutionen II. Math. Zeitschr. 31, S. 733 \S 5, (1930).}。）

\textbf{定理 5.} \emph{\(t\) を素因数分解して \(t=\prod_i p_i^{\varrho_i}\)、\(\varrho_i\ge1\)、\(i\ne j\) ならば \(p_i\ne p_j\) とする。このとき \(\mathfrak K=\mathfrak K^{(1)}\times\mathfrak K^{(2)}\times\cdots\) であり、\(\mathfrak K^{(i)}\) は絶対成分数 \(p_i^{\varrho_i}\) をもつ体である。}

\emph{証明.} 定理 3 と 4 により、\(\{\mathfrak K\}\) の位数 \(\lambda\) は
\[
\lambda=\prod_i p_i^{\alpha_i},\qquad 1\leqq\alpha_i\leqq\varrho_i.
\]
\(\{\mathfrak K\}\) が生成する \(\mathscr K\) の巡回部分群 \(\mathfrak Z\) は、位数 \(p_i^{\alpha_i}\) の部分群 \(\mathfrak Z_i\) の直積である。したがって
\[
\{\mathfrak K\}=\prod_i\{\mathfrak K^{(i)}\}.
\]
これから、\(\mathfrak K\) の各分解体 \(Z\) は \(\mathfrak K^{(i)}\) の分解体でもあることがわかる。したがって \(\mathfrak K^{(i)}\) の絶対成分数 \(p_i^{\sigma_i}\) は \(t\) を割り、\(\sigma_i\le\varrho_i\) である。\(\prod_i\mathfrak K^{(i)}=\mathfrak K_m\) とすれば、
\[
\prod p_i^{\varrho_i}=\prod p_i^{\varrho_i}\cdot m,
\qquad
\sigma_i\geqq\varrho_i,
\]
したがって
\[
\sigma_i=\varrho_i,\qquad m=1,\qquad
\mathfrak K=\prod \mathfrak K^{(j)},
\quad\text{証明終。}
\]

\subsection*{\S 25. ガロア的極大可換部分体をもつ \(\mathfrak K_r\) の正規表現\protect\footnotemark}
\label{normaldarstellung-von-r_r-mit-galoisschem-maximalem-kommutativem-teilkoerper-1}
\footnotetext{より簡単な根拠は \S 27 以降を参照。}

\(\mathfrak K_r\) の極大可換部分体 \(\mathfrak Z\) がガロア的であるとする。\(\mathfrak Z/\mathsf P\) の \(n\) 個の自己同型 \(H_i\) は、25 頁の定理 7 により \(\mathfrak K_r\) の内部自己同型である。すなわち、元 \(u_i\) が存在して
\[
\mathsf{H}_i z=u_i^{-1}zu_i\quad\text{すべての }z\text{ に対して}\mathfrak Z
\]
となる。各 \(u_i\) は、\(\mathfrak Z\) の零でない因子 \(y\) を除いて定まる。実際、すべての \(z\in\mathfrak Z\) に対して \(u^{-1}zu=u'^{-1}zu'\) ならば、すべての \(z\in\mathfrak Z\) に対して \((u^{-1}u')^{-1}z(u^{-1}u')=z\) が従う。したがって \(u^{-1}u'\) を \(\mathfrak Z\) に添加すると \(\mathfrak K_r\) の可換部分体が得られ、それは \(\mathfrak Z\) の極大性により \(\mathfrak Z\) と一致する。ゆえに \(u^{-1}u'\in\mathfrak Z\) である。

次の定理を証明する：
\[
\mathfrak K_r=u_1\mathfrak Z+\cdots+u_n\mathfrak Z.
\]
そのためには、明らかに \(u_i\) が \(\mathfrak Z\) に関して線形独立であることだけが必要である。証明は \(u_i\) を直接与えることによって行う。いま
\[
\mathfrak Z_{\mathsf{Z}}=e_1\mathsf{Z}+\cdots+e_n\mathsf{Z},\qquad
\mathfrak K_{r\mathsf{Z}}=\mathfrak l_1+\cdots+\mathfrak l_n,\quad
\mathfrak l_i=\mathfrak K_{r\mathsf{Z}}e_i
\]
とし、\(p_{ik}\) をこの分解に対する擬行列単位系とする。8 頁では、表現加群 \(e_iZ\) が与える \(\mathfrak Z\) から \(Z\) への同型を \(\Theta_i\) と表したので、
\[
S_i=\Theta_i\Theta_1^{-1}
\]
は \(Z\) の \(n\) 個の自己同型である（ただし \(Z\) 自身に作用させる限りにおいてのみであり、ここで常に基礎としている \(S_i\) の \(\mathfrak K_{rZ}\) 全体への作用についてではない！）。簡単のため、\(S_i=\Theta_i\Theta_1^{-1}\) における \(\Theta_i\) の代わりに \(\Theta_{S_i}\) と書き、さらに \(p_{ik}\) の添字 \(i,k\) も対応する自己同型 \(S\) で置き換える：
\[
p_{ik}=p_{S_iS_k}.
\]

\(\mathfrak Z\) の \(n\) 個の自己同型は（\(E\) は恒等置換）
\[
\mathsf{H}=\Theta_E^{-1}\Theta
\]
である。ここで \(H\) の番号付けを
\[
\mathsf{H}_S=\Theta_E^{-1}\Theta_{S^{-1}}
\]
となるように選ぶ。

このもとで
\[
u_T=\sum_S p_{S,ST}=\sum_S S(p_{E,T})
\]
と置くことができる。実際、

1. \(u_T^{-1}\) は存在する：
\[
u_T^{-1}=\sum_S p_{ST,S}
\frac{1}{\alpha_{S,S}^{ST}\overline{\gamma}_{S,ST}\overline{\gamma}_{ST,S}}.
\]
（\(\bar\gamma\) の意味については 32 頁を参照。）

2. \(u_T\) は \(\mathfrak G\)-不変であるから、\(\mathfrak K_r\) の元である：
\[
R(u_T)=\sum_S RS(p_{E,T})=\sum_{S'} S'(p_{E,T})=u_T.
\]

3. \(u_T\) は \(H_T\) を生成する。すなわち \(z\in\mathfrak Z\) ならば \(u_T^{-1}zu_T=H_T(z)\)、または
\[
z u_T=u_T \mathsf{H}_T(z)
\]
である。

実際、
\[
z=\sum_S\Theta_S(z)e_S,\qquad
u_T=\sum_S S(e_Ep_{E,T})=\sum_S e_S S(p_{E,T})
\]
であるから、
\[
zu_T=\sum_{S,S'}\Theta_S(z)e_Se_{S'}S(p_{E,T})
=\sum_S\Theta_S(z)S(p_{E,T}).
\]
他方、
\[
u_T=\sum_S S(p_{E,T}e_T)=\sum_S S(p_{E,T})e_{ST}
\]
であり、\(\mathfrak Z\) の元に作用させるときには次の等式が成り立つ：
\[
\Theta_S \mathsf{H}_T=\Theta_S\Theta_E^{-1}\Theta_{T^{-1}}
=S\Theta_{T^{-1}}=ST^{-1}\Theta_E=\Theta_{ST^{-1}}.
\]
したがって
\[
\mathsf{H}_T(z)=\sum_S e_S\Theta_S \mathsf{H}_T(z)=\sum_S e_S\Theta_{ST^{-1}}(z),
\]
すなわち
\[
\begin{aligned}
u_T\mathsf{H}_T(z)
&=\sum_{S,S'} S(p_{E,T})e_{ST}\cdot e_{S'}\Theta_{S'T^{-1}}(z)\\
&=\sum_S S(p_{E,T})\Theta_S(z).
\end{aligned}
\]
ゆえに
\[
z\cdot u_T=u_T\cdot \mathsf{H}_T(z).
\]

4. \(u_S\) は \(\mathfrak Z\) に関して線形独立である。

\[
\sum_T z_Tu_T=0
\]
ならば、\(z=\sum_S\Theta_S(z)e_S\) により
\[
\sum_{T,S,L}\Theta_S(z_T)e_Sp_{L,LT}
=\sum_{S,T}\Theta_S(z_T)p_{S,ST},
\quad\text{したがって}\quad
\Theta_S(z_T)=0,
\]
ゆえに \(z_T=0\) である。

以下では \(H_S(z)\) の代わりに \(z^S\) と書き、これを「記号的冪」と呼ぶ。\(u_S\) については次の形の関係式が成り立たなければならない：
\[
u_Ru_T=u_{RT}a_{R,T},\qquad a_{R,T}\ne0\quad\text{かつ}\mathfrak Z.
\]
なぜなら
\[
\begin{aligned}
u_Ru_T
&=\sum_{S,S'} S(p_{E,R})S'(p_{E,T})
=\sum_S p_{S,SR}p_{SR,SRT}\\
&=\sum_S p_{S,SRT}\alpha_{S,SRT}^{(SR)}
=\sum_S p_{S,SRT}\alpha_{S,SRT}^{(SR)}e_{SRT}
\end{aligned}
\]
であり、また
\[
a_{R,\,T}=\sum_S e_{SRT}\,\Theta_{SRT}\left(a_{R,\,T}\right),
\]
したがって
\[
u_{RT}a_{R,T}=\sum_S p_{S,SRT}\Theta_{SRT}(a_{R,T})
\]
であるから、
\[
\Theta_{SRT}(a_{R,\,T})=\alpha_{S,SRT}^{(SR)},
\quad
a_{R,\,T}=\sum_S e_{SRT}\alpha_{S,SRT}^{(SR)}
=\sum_L e_L\alpha_{LT^{-1}R^{-1},L}^{(LT^{-1})}
\]
が従う。

\(a_{R,T}\) の系を \(\mathfrak K_r\) の\emph{小因子系}と呼ぶことにする。大因子系 \(\alpha_{BC}^{(A)}\) と小因子系 \(a_{R,T}\) は互いに一意に定め合う。二つの大因子系の積 \(\varepsilon_{BC}^{(A)}=\alpha_{BC}^{(A)}\beta_{BC}^{(A)}\) に属するのは、\(\alpha_{BC}^{(A)}\)、\(\beta_{BC}^{(A)}\) に対応する小因子系 \(a_{R,T}\)、\(b_{R,T}\) の積 \(f_{R,T}=a_{R,T}b_{R,T}\) である。これは上の公式から直ちに従う。大単位因子系 \(\alpha_{BC}^{(A)}=1\) には小単位因子系 \(a_{R,T}=1\) が属し、逆もまた成り立つ。

\(u_S\) を \(u_S'=u_Sr_S\) で置き換えることができる。ここで \(r_S\ne0\) は \(\mathfrak Z\) から任意に選ぶ。\(u_S'\) に属する小因子系を \(a_{R,T}'\) とし、
\[
\text{(1)}\qquad u_R'u_T'=u_{RT}'a_{R,T}'
\]
とする。\(a_{R,T}'\) は \(a_{R,T}\) と\emph{同伴する}という。いま
\[
u_R r_R u_T r_T=u_{RT}r_{RT}a'_{R,T}
=u_Ru_T r_R^T r_T=u_{RT}a_{R,T}r_R^T r_T
\]
であるから、
\[
\text{(2)}\qquad a'_{R,T}=a_{R,T}\frac{r_R^T r_T}{r_{RT}}
\]
が従う。

\(u_S'=u_Sr_S\) に属する擬行列単位系は
\[
p_{S,ST}'=p_{S,ST}r_T
\]
である。したがって小因子系 \(a'_{R,T}\) は、\(a_{R,T}\) に属する大因子系 \(\alpha_{BC}^{(A)}\) と同伴する大因子系 \(\alpha_{BC}^{(A)\prime}\) に属する。同伴する小因子系は同伴する大因子系に対応し、逆もまた成り立つ。同伴する大因子系の類 \(\{\alpha\}\) から、それに属する同伴小因子系の類 \(\{a\}\) への対応は群同型である。

ここで \(u_S\) を \(\mathfrak Z\) 上の \(n\) 次行列で表すことに進む。この表現は、これまで考えた表現とは本質的に異なる。

\(\mathfrak K_r\) の \(\mathfrak Z\)-基底
\[
\text{(I)}\qquad k_{S_1},\ldots,k_{S_n}
\]
に右から \(u_S\) を掛けると、
\[
\text{(II)}\qquad k_{S_1}u_S, \ldots, k_{S_1}u_S
\]
が得られる。(I) を (II) へ移す（正則）行列 \(A_S\)、すなわち
\[
\text{(III)}\qquad
(k_{S_1}u_S,\ldots,k_{S_n}u_S)
=(k_{S_1},\ldots,k_{S_n})A_S
\]
を、\(u_S\) を表す行列と呼ぶ。この表現は、\(u_S\) と同様に \(\mathfrak Z/\mathsf P\) の自己同型を生成するすべての \(\mathfrak K_r\)-元、すなわち \(r_S\ne0\) を \(\mathfrak Z\) から任意に取ったすべての元 \(u_Sr_S\) に及ぶものとする。これらの元はすべて一つの群 \(\mathfrak G^*\) をなし、その正規部分群として \(\mathfrak Z\) の乗法群 \(\mathfrak Z^*\) を含む。商群 \(\mathfrak G^*/\mathfrak Z^*\) は、\(\mathfrak G^*\) の各剰余類の各元がそれに対応する \(\mathfrak Z\) の自己同型を生成する、という仕方で \(\mathfrak Z/\mathsf P\) のガロア群 \(\mathfrak G\) と同型である。この群 \(\mathfrak G^*\) は、いま述べた仕方で
\[
u_S\mapsto A_S
\]
と表現される。しかしこの表現は準同型でも相反準同型でもない。\(u_S\mapsto A_S\)、\(u_T\mapsto A_T\) から \(u_Su_T\mapsto A_SA_T\) が従うのではなく、
\[
\begin{aligned}
(k_{S_1},\ldots,k_{S_n})u_Su_T
&=(k_{S_1},\ldots,k_{S_n})A_Su_T
=(k_{S_1},\ldots,k_{S_n})u_TA_S^T\\
&=(k_{S_1},\ldots,k_{S_n})A_TA_S^T:
\quad u_Su_T\mapsto A_TA_S^T
\end{aligned}
\]
となるからである。

この \(\mathfrak G^*\) の表現を\emph{交叉表現}と呼ぶ。\(h_{S_1},\ldots,h_{S_n}\) が \(\mathfrak Z\) に関する \(\mathfrak K_r\) の別の基底で、
\[
\text{(IV)}\qquad
(h_{S_1},\ldots,h_{S_n})=(k_{S_1},\ldots,k_{S_n})M
\]
ならば、(III) と (IV) から
\[
(h_{S_1},\ldots,h_{S_n})u_S
=(h_{S_1},\ldots,h_{S_n})M^{-1}A_SM^S
\]
を得る。すなわち、\(u_S\) が \(k\)-基底に関して \(A_S\) で表されるならば、\(k\)-基底を \(h\)-基底へ変換する行列を \(M\) とすると、\(h\)-基底に関しては \(M^{-1}A_SM^S\) で表される。

交叉表現を定義する公式 (III) は、交叉表現と環 \(\mathfrak K_r\) との関係が、通常の相反表現（ただし左からでなく右から書かれる）とその表現加群との関係に似ていることを示す。違いが生じるのは、相反表現加群における交換関係
\[
(m^*\tau^*)c=(m^*c)\tau^*
\]
（5 頁参照）の代わりに、ここでは
\[
(kz)u_S=(ku_S)z^S,\qquad
k\in\mathfrak K_r,\quad z\in\mathfrak Z
\]
が現れるからである。

\(\mathfrak K_r\) の右イデアル \(\mathfrak r\) の基底によっても、\(\mathfrak G^*\) の交叉表現を作ることができる。\(\mathfrak K_r\) の各右イデアルは、この仕方で同値な表現の一類を生成する。ここで表現 \(u_S\mapsto A_S\)、\(u_S\mapsto B_S\) が同値であるとは、\(B_S=M^{-1}A_SM^S\) となる行列 \(M\) が存在することをいう。

ここで、\(\mathfrak K_r\) の右イデアルが生成する \(\mathfrak G^*\) から \(\mathfrak Z\) への表現だけに限らず、一般に \(\mathfrak G^*\) から \(\mathfrak Z\) への次数 \(m\) の交叉表現を次のように定義できる。
\begin{enumerate}
\item 各 \(u_S\) に \(m\) 次の \(\mathfrak Z\)-行列 \(A_S\) が対応する。
\item \(u_S\mapsto A_S\)、\(u_T\mapsto A_T\) から \(u_Su_T\mapsto A_TA_S^T\) が従う。
\end{enumerate}

階数 \(m\) の \(\mathfrak Z\)-右加群
\[
\mathfrak M=x_1\mathfrak Z+\cdots+x_m\mathfrak Z
\]
を取り、
\[
(x_1u_S,\ldots,x_mu_S)=(x_1,\ldots,x_m)A_S
\]
によって、交換関係
\[
(mz)u_S=(mu_S)z^S
\]
をもつ \(\mathfrak G^*\)-右加群とする。これは 5 頁で、与えられた表現から表現加群を構成したのと全く同様である。さらに 20 頁と同じように、
\[
m\cdot zu_S=mz\cdot u_S=mu_S\cdot z^S,
\]
一般に
\[
m\cdot\sum z_Su_S=\sum mz_S\cdot u_S
=\sum mu_S\cdot z_S^S
\]
と定めることによって、\(\mathfrak M\) が有限 \(\mathfrak K_r\)-右加群になることを示す。この定義には特定の基底を用いたので、20 頁の計算に相当するものは全く必要ない。

しかし M.Z. \S 18 により、各有限単純 \(\mathfrak K_r\)-加群 \(\mathfrak M\) は \(\mathfrak K_r\) の単純右イデアルと作用素同型である。したがって 22 頁の定理 2 と同様に、\(\mathfrak G^*\) から \(\mathfrak Z\) への既約交叉表現の類は\emph{一つ}しかなく、それはまさに単純右イデアルに属する類である。

この既約表現の次数は \(t\) である。ここで \(t\) は \(\mathfrak K\) の指数（絶対指数）である。これは単純な階数計算から得られる。実際、その次数は単純右イデアルの \(\mathfrak Z\)-階数 \((\mathfrak r:\mathfrak Z)\) に等しい。22 頁により
\[
(\mathfrak Z:\mathsf P)=n=rt;\qquad
(\mathfrak K_r:\mathsf P)\equiv n^2;\qquad
(\mathfrak r:\mathsf P)\equiv n\cdot t
\]
である。最後の等式は \(\mathfrak K_r=\mathfrak r_1+\cdots+\mathfrak r_r\) による。
\[
(\mathfrak r:\mathsf P)
=(\mathfrak r:\mathfrak Z)\cdot(\mathfrak Z:\mathsf P)
\]
すなわち \(n\cdot t=(\mathfrak r:\mathfrak Z)\cdot n\) から、\((\mathfrak r:\mathfrak Z)=t\) が従う。

\subsection*{\S 26. 交叉表現の乗法}\label{multiplikation-der-verschraenkten-darstellungen}

これまでは交叉表現を常に \(\mathfrak G^*\) の表現とみなした。しかし、それを \(\mathfrak Z/\mathsf P\) のガロア群 \(\mathfrak G\) の表現と考えることもできる。すなわち、\(\mathfrak G\) の各元 \(S\) に、変換 \(u_S^{-1}zu_S=z^S\) によって自己同型 \(S\) を生成する \(\mathfrak K_r\) の元 \(u_S\) を表す行列 \(A_S\) を対応させる。

このような表現 \(S\mapsto A_S\) を特徴づけるには\srcfn{1}{このような表現（超複素系とは独立に）
\[
S\mapsto A_S;\quad T\mapsto A_T;\quad
ST\mapsto A_TA_s^T=A_{ST}\cdot a_{S,T}^{-1}
\]
（またはその転置）は、Speiser（Math. Zschr. 5）が因子系 \(a_{S,T}\) を定義する出発点であった。したがって、この因子系は、ここで別の基礎から導入した因子系と同一であることがわかる。}、それをどの \(\mathfrak K_r\) で取ったか、また \(\mathfrak Z\) 上の \(\mathfrak K\) のどの \(u_S\) とどの \(k_{S_1},\ldots,k_{S_n}\) を用いたかを述べなければならない。基底の変更は同値な表現を与え、\(u\) の変更は同伴する因子系 \(a_{S,T}\) を与える。したがって本質的なのは、表現が取られる \(\mathfrak K_r\)、すなわち同伴する因子系の類 \(\{a_{S,T}\}\) である。

ここで、このような \(\mathfrak G\) の表現が二つ与えられたとする：
\[
\begin{array}{lll}
S\mapsto u_S\mapsto A_S
&\text{対応する小因子系}&a_{R,T},\\
S\mapsto v_S\mapsto B_S
&\text{対応する小因子系}&b_{R,T}.
\end{array}
\]

二つの行列 \(A=(a_{ik})\)、\(B=(b_{\mu\nu})\) の Kronecker 積 \(A\times B\) とは、行列 \(A\times B=(a_{ik}b_{\mu\nu})\) をいう。

\noindent\textbf{主張.} \emph{\(S\mapsto A_S\times B_S\) も \(\mathfrak G\) の交叉表現である。さらに \(A_S\) が \(\mathfrak K_r\) から、\(B_S\) が \(\mathfrak S_s\) から生じるならば、この表現は積 \(\mathfrak K_r\times\mathfrak S_s\) に属する。}

\noindent\emph{証明.} 表現 \(S\mapsto A_S\times B_S\) は交叉表現である。実際、
\[
S\mapsto A_S\times B_S,\qquad
T\mapsto A_T\times B_T
\]
であり、
\[
ST\mapsto A_{ST}\times B_{ST}
=(A_TA_S^T\times B_TB_S^T)a_{S,T}^{-1}b_{S,T}^{-1}.
\]
したがって証明すべきことは
\[
(A_T\times B_T)(A_S\times B_S)^T a_{S,T}^{-1}b_{S,T}^{-1}
=(A_TA_S^T\times B_TB_S^T)a_{S,T}^{-1}b_{S,T}^{-1}
\]
である。これは次の考察から従う。

行列 \(A\) と \(B\) を行列単位 \(c_{ik}\) と \(d_{\mu\nu}\) で表す（これらは \(\mathfrak K_r\) や \(\mathfrak S_s\) の行列単位とは無関係である）：
\[
A=\sum c_{ik}\alpha_{ik};\qquad
B=\sum d_{\mu\nu}\beta_{\mu\nu}.
\]
すると
\[
A\times B=B\times A
=\sum c_{ik}d_{\mu\nu}\cdot\alpha_{ik}\cdot\beta_{\mu\nu},
\]
すなわち新しい行列単位 \(c_{ik}d_{\mu\nu}=d_{\mu\nu}c_{ik}\) による行列に等しい。したがって \(A\)-行列と \(B\)-行列の乗法は、二つの行列環の直積
\[
\sum c_{ik}\mathfrak Z\times\sum d_{\mu\nu}\mathfrak Z
\]
における元の乗法に一対一に対応する。ゆえに各 \(A\)-行列は各 \(B\)-行列と交換可能である（ただし \(A\)-行列同士または \(B\)-行列同士が交換可能なのではない）。そして、この各 \(A\)-行列と各 \(B\)-行列との交換可能性から、\(a_{S,T}^{-1}b_{S,T}^{-1}\) を掛ければ、まさに証明すべき関係式が得られる。

%% END NOETHER SOURCE
%% BEGIN NOETHER SOURCE "ja/t5.tex"
\section*{第 VI 章　接合積の理論}\label{kapitel-vi.-theorie-der-verschraenkten-produkte}

\subsection*{\S 27. \(\mathfrak K_r\) の接合積としての表示}

\S 24 で因子系から展開した接合積の理論を、ここで改めて独立に基礎づける。これらの接合積の意義は、ガロア体とその群が \(\mathfrak K_r\) の中で\emph{同時に}記述されることにある。

\noindent\textbf{定理 1.} \emph{\(\mathfrak Z\) が \(\mathfrak K_r\) のガロア的分解体、したがって極大可換部分体であり、\(\mathfrak G\) がそのガロア群ならば、\(\mathfrak K_r\) は接合積 \(\mathfrak Z\cdot\mathfrak G\) に等しい。逆に、各接合積 \(\mathfrak Z\cdot\mathfrak G\) はこの種の \(\mathfrak K_r\) を生成する。}

\noindent\textbf{証明と定義. 定義 1.} \(\mathfrak G=\{E,S,\ldots,T\}\) を \(\mathfrak Z\) のガロア群とする。\(\mathfrak K_r\) が接合積に等しいとは、
\[
\mathfrak K_r=\mathfrak Zu_E+\cdots+\mathfrak Zu_T
\]
であり、\(zu_S=u_Sz^S\) および \(u_Su_T=a_{S,T}u_{ST}\) が成り立つことをいう（\(\mathfrak Z\) の元 \(a_{S,T}\) は小因子系である）。

\noindent\textbf{定義 2.} 生成群 \(\mathfrak G^*\) を、変換によって \(\mathfrak Z\) 全体をそれ自身へ移す \(\mathfrak K_r\) の正則元 \(g\) の全体として定義する。\(g^{-1}\mathfrak Zg=\mathfrak Z\)、したがってある \(S\in\mathfrak G\) に対して \(g^{-1}zg=z^S\) である。\(\mathfrak Z\) の各自己同型はある \(g\) によって生成されるので、\(\mathfrak G\) は \(\mathfrak G^*\) の群準同型像、\(\mathfrak G^*\to\mathfrak G\) となり、さらに
\[
\mathfrak G\cong\mathfrak G^*/\mathfrak T^*
\quad\text{ただし}\quad
\mathfrak T^*\supseteqq\mathfrak Z^*.
\]
（\(\mathfrak T^*=\mathfrak Z^*\) であることは後で示す。\(\mathfrak Z\) が極大可換部分環でもあることはまだ証明されていない。これは独立にも証明できる。）

\noindent\textbf{1.} \emph{接合積としての表示の証明.} \(E,\ldots,T\) という自己同型を生成する \(\mathfrak K_r\) の元を \(u_E,u_S,\ldots,u_T\) とする。すなわち \(zu_S=u_Sz^S\) であり、特に 2 により \(u_E=e\) である。和
\[
\mathfrak M=\{\mathfrak Zu_E,\ldots,\mathfrak Zu_T\}
\]
が直和 \(\mathfrak Zu_E+\cdots+\mathfrak Zu_T\) に等しいことを示せばよい。そうすれば階数が一致するため、それは \(\mathfrak K_r\) に等しい。これは次のように従う（Schwarz）。

\(\mathfrak M\) は \(\mathfrak Z\) の左加群かつ右加群、したがって \(\mathfrak Z\) から \(\mathfrak Z\) への表現加群であり、しかも \(\mathfrak M\) は \(n\) 個の異なる表現を生成する。したがって、\(\mathfrak Z\) は第一種可換体であるから、\(\mathfrak M\) の階数は少なくとも \(n\) である。それ以外は証明は全く同じように進む。

\noindent\textbf{2.} \(\mathfrak G\cong\mathfrak G^*/\mathfrak Z^*\)、すなわち \(\mathfrak T^*=\mathfrak Z^*\) である。変換によって \(\mathfrak Z\) 上の恒等写像を誘導するのは、\(\mathfrak Z\) の元だけである。

実際、各 \(z\in\mathfrak Z\) に対して \(tz=zt\) であるとし、
\[
t=z_0+z_1u_S+\cdots+z_{n-1}u_T
\]
と置く。すると
\[
zt=zz_0+\cdots+zz_{n-1}u_T
\]
であり、また
\[
tz=z_0z+\cdots+z_{n-1}u_Tz
=zz_0+z^{S^{-1}}z_1u_S+\cdots+z^{T^{-1}}z_{n-1}u_T.
\]
したがって直和性により
\[
(z-z^{S^{-1}})z_1u_S=0,\ldots,
(z-z^{T^{-1}})z_{n-1}u_T=0,
\]
さらに \(u_S,\ldots\) が正則なので、
\[
(z-z^{S^{-1}})z_1=0,\ldots,
(z-z^{T^{-1}})z_{n-1}=0,
\qquad\text{すべての }z\text{ に対して}
\]
を得る（あるいは階数が \(n\) に等しいことによる）。しかし \(S,\ldots\) など（いずれも \(E\) と異なる）について、定義により \(z-z^{S^{-1}}\ne0\) となる \(z\) が少なくとも一つあり、また \(\bar z-\bar z^{T^{-1}}\ne0\) となる \(\bar z\) がある。ゆえに \(z_i=0\)、\(i=1,\ldots,n-1\) であり、したがって \(t=z_0\in\mathfrak Z^*\) である。証明終。

\noindent\textbf{3.} \(u\) は乗法則 \(u_Su_T=a_{S,T}u_{ST}\) を満たし、\(a_{S,T}\) は接合結合則
\[
a_{R,S}\cdot a_{RS,T}=a_{S,T}^{R^{-1}}\cdot a_{R,ST}
\]
を満たす（\(a_{S,T}\) は小因子系である）。

実際、2 により
\[
\mathfrak Z^*u_S\cdot\mathfrak Z^*u_T=\mathfrak Z^*u_{ST},
\]
したがって
\[
u_Su_T=a_{S,T}u_{ST}
\]
である。ここで \(a_{S,T}\in\mathfrak Z^*\)、したがって \(a_{S,T}\in\mathfrak Z\) かつ \(a_{S,T}\ne0\) である。さらに \(u_R\cdot u_Su_T=u_Ru_S\cdot u_T\) により
\[
\begin{aligned}
u_Ra_{S,T}u_{ST}
&=a_{S,T}^{R^{-1}}\cdot u_Ru_{ST}
=a_{S,T}^{R^{-1}}a_{R,ST}u_{RST}\\
&=a_{R,S}u_{RS}\cdot u_T
=a_{R,S}a_{RS,T}u_{RST},
\end{aligned}
\]
となり、これが上の乗法則である。

\noindent\textbf{4. 逆方向：} 接合積 \(=\mathfrak K_r\)。直和
\[
\mathfrak M=\mathfrak Zu_E+\cdots+\mathfrak Zu_T
\]
を、関係式 \(zu_S=u_Sz^S\)、\(u_Su_T=a_{S,T}u_{ST}\) によって環とする（ここで \(a_{S,T}\) は接合結合関係を満たし、したがって \(u_Ru_Su_T\) は一意に定義される）：
\[
\begin{array}{ll}
1.&z_1\cdot z_2u_S=z_1z_2\cdot u_S;\\[2pt]
2.&zu_S\cdot u_T=z\cdot u_Su_T.
\end{array}
\]
これによって \(\mathfrak M\) は結合環となる。

\noindent\textbf{逆定理 2.} \emph{接合積 \(\mathfrak Z\cdot\mathfrak G\) は両側単純であり、中心は \(\mathsf P\) である（R. Brauer の着想による）。} 証明は次の補題に基づく。

\noindent\textbf{補題 1.} \emph{\(\mathfrak Z\) の各元と交換する元はすべて \(\mathfrak Z\) に属する。}

\noindent\emph{証明.} \(w=\sum_S c_{\lambda_S}u_S\)（\(c_{\lambda_S}\in\mathfrak Z\)）とし、すべての \(z\in\mathfrak Z\) に対して \(wz=zw\) と仮定する。すると
\[
\sum_S zc_{\lambda_S}u_S
=\sum_S c_{\lambda_S}u_Sz
=\sum_S c_{\lambda_S}z^{S^{-1}}u_S,
\]
したがって、すべての \(z\) と \(S\) に対して \(zc_{\lambda_S}=c_{\lambda_S}z^{S^{-1}}\)、または \(c_{\lambda_S}(z-z^{S^{-1}})=0\) である（\(u\) の線形独立性！）。各 \(S\ne E\) について \(z-z^{S^{-1}}\ne0\) となる \(z\) が少なくとも\emph{一つ}あるので、\(c_{\lambda_S}=0\) が従う。したがって \(w\in\mathfrak Z\) である。

\noindent\textbf{補題 2.} \emph{\(\mathfrak Z\) を左右の作用素としたとき、\(\mathfrak Zu_S=u_S\mathfrak Z\) と双同型である各 \(\mathfrak Z\)-双加群 \(\mathfrak A\) は、\(\mathfrak Zu_S\) と同一である。}

\noindent\emph{2 の証明.} \(\mathfrak A\) は左加群として \(\mathfrak Zu_S\) と同型である。したがって \(u_S\mapsto a\in\mathfrak A\) ならば \(\mathfrak Zu_S\mapsto\mathfrak Za\)、ゆえに \(\mathfrak A=\mathfrak Za\) である。双同型性から (1) \(za=az^S\) が従い（\(u_Sz\mapsto az\) による）、これから \(\mathfrak A=\mathfrak Zu_S\) を得る。実際、\(a=u_Su_S^{-1}a=u_Sw\) と置けば、\(za=zu_Sw=u_Sz^Sw\) である一方、(1) により \(za=u_Swz^S\) である。したがって（\(u_S\) を掛けて）すべての \(z\in\mathfrak Z\) に対して \(z^Sw=wz^S\) となる。\(z^S\) も \(\mathfrak Z\) 全体を動くから、\(w\) は \(\mathfrak Z\) の各元と交換し、1 により \(w\in\mathfrak Z\) である。

\noindent\textbf{補題 3.} \emph{\(\mathfrak Z\cdot\mathfrak G\) の各 \(\mathfrak Z\)-双加群は \(\mathfrak Zu_{S_1}+\cdots+\mathfrak Zu_{S_k}\) という形である。ここで \(u_{S_1}\) は \(u_S,\ldots,u_T\) のうち相異なる任意の \(k\) 個を表す。}

\noindent\emph{証明.} 各直和因子 \(\mathfrak Zu_S\) は階数 1 なので、接合積 \(\mathfrak Z\cdot\mathfrak G\) は \(\mathfrak Z\)-双加群として完全可約である。さらに双加群として、すべての直和因子は異なる類に属し、したがって互いに双同型ではない。ゆえに各 \(\mathfrak Z\)-双加群 \(\mathfrak M\) は、完全可約性によって直和因子であり、それ自身完全可約である。\(\mathfrak M=\sum\mathfrak A_{S_i}\) で、\(\mathfrak A_{S_i}\) は \(\mathfrak Zu_{S_i}\) と同型である。したがって 2 により、それは \(\mathfrak Zu_{S_i}\) と同一である（すべての \(\mathfrak A_{S_i}\) は異なる類に属する）。

\noindent\textbf{補題 4.} \emph{\(\mathfrak Z\cdot\mathfrak G\) のうち \(\mathfrak Z\) を含む各環、特に \(\mathfrak Z\cdot\mathfrak G\) 自身は両側単純である。さらに、\(\mathfrak Z\) を含む部分環と \(\mathfrak G\) の部分群とは一対一に対応する。}

\noindent\emph{証明.} \(\mathfrak Z\) を含む各環 \(\mathfrak o\) は \(\mathfrak Z\)-双加群であるから、3 により \(\sum\mathfrak Z\cdot u_{S_i}\) という形である。\(\mathfrak o\) は環なので、二つの \(u_{S_i}\) の積も \(\mathfrak o\) に属し、したがって対応する \(\mathfrak Z\)-加群も属する。個数は有限なので、\(u_{S_i}\) は因子系を別にすれば群をなす。したがって \(\mathfrak o=\mathfrak Z\cdot\mathfrak H\) であり、\(\mathfrak H\) は \(\mathfrak G\) の部分群である（より正確には、\(\mathfrak Z\cdot u_{S_i}\) が乗法的な系をなす）。いま \(\mathfrak a\ne0\) を \(\mathfrak o\) の任意の両側イデアルとすると、それは \(\mathfrak o\)-双加群であり、したがって \(\mathfrak Z\)-双加群かつ \(\mathfrak o\) の部分加群である。それゆえ、ある \(\mathfrak Zu_{S_i}\) とともに群 \(\mathfrak H\) 全体を含み、したがって \(\mathfrak o=\mathfrak Z\cdot\mathfrak H\) となる。

\noindent\textbf{補題 5.} \emph{\(\mathsf P\) は \(\mathfrak Z\cdot\mathfrak G\) の中心である。}

\noindent\emph{証明.} \(w\) が \(\mathfrak Z\cdot\mathfrak G\) の各元と交換するならば、特に \(\mathfrak Z\) と交換するので、1 により \(w\in\mathfrak Z\) である。さらに \(wu_S=u_Sw^S\) であるから、交換可能ならば \(w=w^S\) となる。したがって \(w\in\mathsf P\) である。

\noindent\emph{注意 1.} この証明は、\(\mathfrak Z\) が \(\mathsf P\) 上の無限代数的ガロア体である場合にも定理が成り立つことを示す。実際、完全可約性に関する定理は保たれる（Krull, Zahlringe, Math. Ann. 99）。この場合、\(\mathfrak Z\cdot\mathfrak G\) はすべての有限和 \(\sum_{S_\lambda}u_{S_\lambda}\cdot c_{S_\lambda}\) の系である。

ただし 4 では、\(\mathfrak o=\mathfrak Z\cdot\mathfrak H\) を \(\mathfrak o\) について仮定しなければならない。これは全積 \(\mathfrak Z\cdot\mathfrak G\) については成り立つが、\(\mathfrak o\) が \(\mathfrak Z\) 上無限ならば、\(u_{S_i}u_{S_k}\) が属することから、その系が群であることをもはや結論できないからである。

\noindent\emph{注意 2.} ここで直接に基礎づける場合には、\S\S 25/26 で展開した交叉行列表現の理論を接続すべきである。以下、\S 28 の末尾では、既約交叉表現の次数が \(t\)、すなわち \(\mathfrak K\) の絶対成分数（指数）であることを用いる。この表現論によって、ここで乗法定数として導入した因子系が文献で通常用いられる概念と同一になることを、もう一度指摘しておく。

\subsection*{\S 28. 因子系の積定理}\label{produktsatz-fuer-faktorensysteme}

いま
\[
\mathfrak K_r=\mathfrak Z\cdot\mathfrak G
=\mathfrak Zu_E+\cdots+\mathfrak Zu_T,
\qquad u_Su_T=u_{ST}a_{S,T},
\]
\[
\overline{\mathfrak K}_r
=\overline{\mathfrak Z}\,\overline{\mathfrak G}
=\overline{\mathfrak Z}\bar u_E+\cdots+\overline{\mathfrak Z}\bar u_T,
\qquad
\bar u_S\bar u_T=\bar u_{ST}\bar b_{S,T}
\]
とする（\(\mathfrak Z\) と \(\overline{\mathfrak Z}\) は同型なガロア体であり、\(\mathfrak G\)、\(\overline{\mathfrak G}\) はその群である）。すると
\[
\mathfrak K_r\times\overline{\mathfrak K}_r
=\widetilde{\mathfrak Z}\,\widetilde{\mathfrak G}\times\mathsf P_n
=\mathfrak A_f\times\mathsf P_n
=\mathfrak A_j
\]
であり、\(\tilde u_S\tilde u_T=\tilde u_{ST}\tilde c_{S,T}\) となる（\(\widetilde{\mathfrak Z}\) は \(\mathfrak Z\) および \(\overline{\mathfrak Z}\) と同型である）。これは階数計算と、\(\mathfrak Z\) も分解体であることから従う。

\noindent\textbf{主張.} \emph{\(\tilde c_{S,T}=\tilde a_{S,T}\tilde b_{S,T}\) である。ここで \(\tilde a,\tilde b\) は、それぞれ \(a\) および \(\bar b\) と同型な \(\widetilde{\mathfrak Z}\) の元である。}

証明は、絶対長 \(n\) をもつ \(\mathfrak A_j\) の任意の左イデアルの自己同型環が \(\mathfrak A_f=\widetilde{\mathfrak Z}\widetilde{\mathfrak G}\) と同型であること、また適当な \(\mathfrak A_j\) の左イデアルについてこの自己同型環を直接に定められることに基づく。より一般には、\(n\) の代わりに \(s\) として、便宜上 \(\mathfrak K_r\) について述べる次が成り立つ。

\noindent\textbf{1. 定理 1.} \emph{\(\mathfrak K_{rs}=\mathfrak K_r\times\mathsf P_s\) の絶対長 \(n\) の左イデアル \(\mathfrak l\) の自己同型環は、\(\mathfrak K_r\) と同型である。}

証明は次の補題から従う。

\noindent\textbf{補題 1.} \emph{\(\mathfrak K_{rs}\) の絶対長 \(n\) の左イデアルはすべて作用素同型である。したがってそれらの自己同型環は環同型である。特に \(\mathfrak b\) を \(\mathsf P_s\) の単純イデアルとすれば、\(\mathfrak l=\mathfrak K_r\times\mathfrak b\) は絶対長 \(n\) をもつ。}

\noindent\emph{証明.} 左イデアルが絶対長 \(n\) をもつ、すなわち \(n\) 個の絶対既約単純イデアルに分解することから、\(\mathfrak K_{rs}\) における長さも同じであること、すなわち同数の \(\mathfrak K_{rs}\) の単純イデアル直和因子が現れることが従う。これからそれらの作用素同型性、したがって自己同型環の環同型性が得られる。さらに \(\mathfrak K_r=\mathfrak Z\cdot\mathfrak G\) の絶対長は \(n\)、すなわち階数は \(n^2\) である。ゆえに \(\mathfrak K_r\times\mathsf P_s\) の絶対長は \(n\cdot s\)、したがって \(\mathfrak K_r\times\mathfrak b\) の絶対長は \(n\) である。

\noindent\textbf{補題 2.} \emph{\(\mathfrak o\) を単位元をもつ環、\(\mathfrak l=\mathfrak o e_1\) を直和因子である左イデアル、したがって \(e_1\) を冪等元とすると、\(\mathfrak l\) の自己同型環は \(e_1\mathfrak o e_1\) と同型である。}

\noindent\emph{証明.} 元 \(e_1re_1\) を掛けることによって \(\mathfrak l\) からそれ自身への作用素準同型が生じる。異なる元 \(e_1re_1\) は異なる準同型を生じる。なぜなら、そのとき \(e_1\) に異なる元が対応するからである。また各準同型はこのように生成される。すなわち \(e_1\mapsto re_1\) であり、\(e_1e_1=e_1\) により \(re_1=e_1re_1\) である。しかも \(e_1\mathfrak o e_1\) における和と積は、対応する準同型の和と積に対応する。

\noindent\textbf{補題 3.} \emph{行列単位 \(c_{ik}\) を用いて \(\mathsf P_s=\sum_1^s c_{ik}\mathsf P\)、\(\mathfrak b=c_{11}\mathsf P+\cdots+c_{s1}\mathsf P\) と置けば、\(\mathfrak K_r\times\mathfrak b\) の自己同型環として \(\mathfrak K_rc_{11}\) が得られ、これは \(\mathfrak K_r\) と同型である。}

\noindent\emph{証明.} \(c_{11}\) は \(\mathfrak K_r\) の各元と交換し、\(\mathfrak K_r\) のいかなる元も零化しない。なぜなら \(c_{11}\) は \(\mathfrak K_r\times\mathsf P_s\) の生成冪等元であり、
\[
c_{11}\cdot\mathfrak K_r\times\mathsf P_s\cdot c_{11}
=\mathfrak K_r\times c_{11}\mathsf Pc_{11}
=\mathfrak K_rc_{11}
\]
だからである。

これらの補題から定理 1 が直ちに従う。

\noindent\textbf{2. \(\mathfrak K_r\times\overline{\mathfrak K}_r\) の決定と、絶対長 \(n\) の左イデアルへの分解.} 定義により
\[
\mathfrak A_f=\mathfrak K_r\times\overline{\mathfrak K}_r
=\mathfrak Z\cdot\mathfrak G\times\overline{\mathfrak Z}\cdot\overline{\mathfrak G},
\]
ここで無標のものと上線付きのものとは各元ごとに交換可能である。したがって
\[
\mathfrak A_j
=(\mathfrak G\times\overline{\mathfrak G})\cdot
(\mathfrak Z\times\overline{\mathfrak Z})
=\mathfrak G\times\overline{\mathfrak G}\cdot
(\mathfrak Ze_1+\cdots+\mathfrak Ze_n)
\]
であり、\(\mathfrak Ze_i=\overline{\mathfrak Z}e_i\) である（ガロア体とそれ自身との積）。すべての直和因子は \(\mathfrak Z\) に関して同じ階数をもち、したがって同じ絶対長をもつ。一方、\(\mathfrak A_f\) の絶対長は \(n^2\) であるから、各直和因子の絶対長は \(n\) である。以下では
\[
\mathfrak l=(\mathfrak G\times\overline{\mathfrak G})\mathfrak Ze_1
\]
を基礎とし、\(ze_1=e_1\bar z\) によって \(\mathfrak Z\) と \(\overline{\mathfrak Z}\) の間の特定の同型を固定する。

\noindent\textbf{3. \(e_1\mathfrak A_fe_1\) の交叉表現、したがって積定理の証明.} 1 と 2 により、積 \(\mathfrak K_r\times\overline{\mathfrak K}_r\) の因子系は、\(e_1\mathfrak A_fe_1=\mathfrak A\) の交叉表現の因子系と同型である。

\noindent\textbf{定理 2.} \emph{\(\mathfrak A=\widetilde{\mathfrak Z}\cdot\widetilde{\mathfrak G}\) であり、ここで \(\widetilde{\mathfrak Z}=\mathfrak Ze_1=\overline{\mathfrak Z}e_1\)、\(\widetilde{\mathfrak G}=\{\ldots,e_1u_S\bar u_Se_1,\ldots\}\) である。しかも}
\[
e_1u_S\bar u_Se_1\cdot e_1u_T\bar u_Te_1
=e_1u_{ST}\bar u_{ST}e_1\cdot a_{S,T}e_1\bar b_{S,T}e_1.
\]
\emph{したがって \(\tilde a_{S,T}=a_{S,T}e_1\)、\(\tilde b_{S,T}=\bar b_{S,T}e_1\) と置けば積定理が証明される（\(\tilde a\) と \(\tilde b\) は同じ体 \(\widetilde{\mathfrak Z}\) に属する）。}

\noindent\emph{証明.} \(e_1\) は環の単位元となる。したがって \(\tilde z=ze_1\) と共役な元は \(\tilde z^Se_1\) で与えられる。元 \(u_S\bar u_Se_1\) がこれらの置換を誘導することを示せばよく、それによって交叉表現が証明される。

\noindent\textbf{補題 1.} \(e_1u_S\bar u_S=u_S\bar u_Se_1\)、したがってこれは \(e_1u_S\bar u_Se_1\) に等しい。

\noindent\emph{証明.} \(u_S^{-1}e_1u_S=e^S\)\srcfn{1}{この規約により \(e_1=e^E\) となる。}、または \(e_1u_S=u_Se^S\) と置く。このとき \(e^S\) は共役な \(e_1,e_2,\ldots,e_n\) 全体を動く。すると (1)\srcfn{2}{\(\sum a_i\bar A_i\bar u_R=\sum a_i\bar u_R\bar A_i^R=\bar u_R\sum(a_i^R)^{R^{-1}}\bar A_i^{-R}\) による。}
\[
e_1\bar u_R=\bar u_Re^{R^{-1}}
\]
が成り立つ。実際、\(e_1=\sum a_i\bar A_i=\sum a_i^R\bar A_i^R\) であり、ここで \(a_i\) は任意の \(\mathfrak Z\)-基底を、\(\bar A_i\) はそれに相補的な \(\overline{\mathfrak Z}\)-基底を動くので、\(a_i^R\) と \(\bar A_i^R\) についても同じことが成り立つ。したがって \(e_1\bar u_R=\bar u_Re^{R^{-1}}\) であり、(1) が証明された。(1) から
\[
e_1u_S\bar u_S=u_Se^S\bar u_S
=u_S\bar u_S\cdot e^{SS^{-1}}=u_S\bar u_Se_1
\]
が従い、補題 1 が証明された。

\noindent\textbf{補題 2.} \(ze_1\cdot e_1u_S\bar u_Se_1=e_1u_S\bar u_Se_1\cdot z^Se_1\)。

\noindent\emph{証明.} 実際、
\[
ze_1\cdot e_1u_S\bar u_Se_1
=zu_S\bar u_Se_1\quad(\text{補題 1})
=u_Sz^S\bar u_Se_1
=u_S\bar u_Se_1\cdot z^Se_1
\equiv b=e_1b
\]
である（\(z\) は \(\bar u\) および \(e_1\) と交換可能である）。\(u'=u_S\bar u_Se_1\) とする。補題 2 は、\(u'\) が \(\mathfrak Ze_1\) の置換を指示することを述べており、したがって定理に示した交叉表現が成り立つ。

\noindent\textbf{補題 3.} \emph{因子系は \(a_{S,T}e_1\cdot\bar b_{S,T}e_1\) に等しくなる。}

\noindent\emph{証明.}
\[
\begin{aligned}
e_1u_S\bar u_Se_1\cdot e_1u_T\bar u_Te_1
&=u_S\bar u_S\cdot u_T\bar u_Te_1
=u_{ST}a_{S,T}\cdot\bar u_{ST}\bar b_{S,T}e_1\\
&=u_{ST}\bar u_{ST}\cdot a_{S,T}\bar b_{S,T}e_1
=e_1u_{ST}\bar u_{ST}e_1\cdot a_{S,T}e_1\cdot\bar b_{S,T}e_1.
\end{aligned}
\]
これで 2 の定理、したがって因子系の積定理が証明された。

\noindent\textbf{4. 同伴する因子系の類についての積定理：}
\[
\{a_{S,T}\}\cdot\{b_{S,T}\}=\{a_{S,T}\cdot b_{S,T}\}.
\]
これは 3 から直ちに従う。実際、新しい生成元 \(v\) と \(\bar v\) への移行は、\(v\bar ve_1\) への移行をも意味する。なぜなら \(u=vc\) および \(\bar u=\bar v\bar d\) ならば
\[
u\bar ue_1=vc\cdot\bar v\bar de_1
=v\bar v\cdot c\bar de_1
=v\bar ve_1\cdot c\bar de_1
\]
だからである。

\noindent\textbf{5.} \(\{a\}\) が \(\{\mathfrak K\}\) の因子系であり、\(t\) が \(\{\mathfrak K\}\) の指数ならば、\(\{a\}^t\) は単位類に等しい。これは、\(\mathfrak G\) の既約交叉表現の次数が \(t\) であることから従い、\(p_{ik}\) の表現に続く Schur の証明（\S 24、37 頁）と全く同じである。

\noindent\textbf{6.} \(\mathfrak Z\) に属する \(\{\mathfrak K\}\)、すなわち \(\mathfrak Z\) を分解体にもつものの群 \(\mathscr K_{\mathfrak Z}\) と、同伴する因子系の類 \(\{a_{S,T}\}\) の群との同型。\S 27 の構成により、\(\{\mathfrak K\}\) と \(\{a_{S,T}\}\) の間には一対一の対応がある。4 により、この対応は準同型なので、同型である。特に単位類が互いに対応することが従う：

\noindent\emph{接合積が \(\mathsf P\) に等しくなる、すなわち中心上の行列環となるのは、因子系が単位系と同伴するとき、かつそのときに限る。}

\subsection*{\S 29. 最小における主種定理}\label{hauptgeschlechtssatz-im-minimalen}

\(\mathfrak K_r=\mathfrak Z\cdot\mathfrak G\) を接合積とし、\(\mathfrak G^*\) を、内部自己同型によって \(\mathfrak Z\) 全体をそれ自身へ移す元すべての群とする（剰余類で書けば \(\mathfrak Z^*u_E,\mathfrak Z^*u_S,\ldots\)）。

\noindent\textbf{定義.} 「\emph{最小における主種}」\(\mathscr H\) は、\(\mathfrak Z^*\) 上の恒等写像の延長である \(\mathfrak G^*\) の群自己同型すべてからなる。この名称の理由は、巡回体への特殊化から明らかになる。\S 30 を参照。

\noindent\textbf{定理.} \emph{主種の各自己同型は \(u_S\mapsto u_S\cdot b^Sb^{-1}=u_Sb^{S-1}\) という形である。ここで \(b\in\mathfrak Z\) である。逆に、このような対応はすべて主種の自己同型を生成する。}

\noindent\emph{証明.} 群自己同型において \(u_S\mapsto v_S\)、また各 \(z\in\mathfrak Z\) に対して \(z\mapsto z\) とし、積が対応するとする。この自己同型は、対応
\[
\sum u_Sc_{\lambda_S}\mapsto\sum v_Sc_{\lambda_S},
\]
ここで \(c_{\lambda_S}\in\mathfrak Z\)、によって \(\mathfrak K_r\) の環自己同型へ拡張できる。しかし、このような自己同型はすべて内部自己同型であることが知られている。

\(b\) を生成元とする。すなわち \(v_S=bu_Sb^{-1}\)、\(z=bzb^{-1}\) である。\(b\) はすべての \(z\) と交換するので \(\mathfrak Z\) に属し、したがって
\[
v_S=bu_Sb^{-1}=u_Sb^Sb^{-1}
\]
となる。

逆に、\(b\in\mathfrak Z\) として \(v_S=u_Sb^Sb^{-1}\) と置けば \(v_S=bu_Sb^{-1}\) となる。したがってこれは、\(\mathfrak Z\) の各元をそれ自身へ移し、定義により \(\mathfrak G^*\) 全体もそれ自身へ移す \(\mathfrak K_r\) の自己同型、すなわち主種の自己同型である。

\subsection*{\S 30. 巡回分解体への特殊化}\label{spezialisierung-auf-zyklische-zerfaellungskoerper}

特に巡回体 \(\mathfrak Z\) を用いて接合積を作ると、既知のノルム定理の解釈が得られる。まず次が成り立つ。

\noindent\textbf{定理 1.} \emph{\(\mathsf P\) 上の固定した巡回分解体 \(\mathfrak Z\) をもつ \(\{\mathfrak K\}\) の群 \(\mathscr K_{\mathfrak Z}\) は、群 \(\mathsf P^*/N\mathfrak Z^*\) と同型である。ここで \(N\mathfrak Z^*\) は、\(\mathfrak Z\) のすべての元 \(z\ne0\) のノルム \(Nz\) がなす群である。}

\noindent\emph{証明.} 証明は交叉表現の正規化と、それに関するいくつかの補題に基づく。

\noindent\textbf{定義.} 巡回交叉表現が
\[
\mathfrak K_r=\mathfrak Z+\mathfrak Zu+\cdots+\mathfrak Zu^{n-1}
\]
という形で、\(zu=uz^S\)、\(u^n=a\) であるとき、その表現を\emph{正規化されている}という。ここで \(S\) は \(\mathfrak G\) の生成元である（すなわち、より一般に \(b_r\in\mathfrak Z\) として \(u_{S^r}=b_ru^r\) と置く代わりに、\(r=0,\ldots,n-1\) に対して \(u_{S^r}=u^r\) と置く）。

\noindent\textbf{補題 1.} \emph{巡回交叉表現は常に正規化されていると仮定できる。その際 \(a\in\mathsf P\) であり、\(\mathsf P\) の各 \(a\) がこのような表現を与える。}

\noindent\emph{証明.} 第一の部分は明らかである。実際、\(u^{-1}zu=z^S\) から \(u^{-r}zu^r=z^{S^r}\) が従い、したがって \(u^r\) が置換 \(z\mapsto z^{S^r}\) を生成する。さらに \(u^n\) は恒等写像 \(z\mapsto z^{S^n}\) を誘導するから、\(u^n=a\) で \(a\in\mathfrak Z\) となる。しかし \(u^n\) はすべての \(u^r\) と交換するので、\(a\) についても同じことが成り立ち、したがって \(a^{S^r}=a\)、ゆえに \(a\in\mathsf P\) である。

したがって \(a\) が唯一の本質的な乗法定数である。それゆえ結合条件は内容を失い、\(\mathsf P\) の各 \(a\) が結合的な系、したがって一つの \(\mathfrak K_r\) を与える。

\noindent\textbf{補題 2.} \emph{同伴する因子系の全類 \(\{a_{S,T}\}\) の代わりに、その類に含まれる正規化された因子系、すなわち正規化された表現から生じる因子系だけを考えれば十分である。この類の部分集合を正規化類と呼ぶ。全類 \(\{a_{S,T}\}\) の群は正規化類の群と同型である。}

\noindent\emph{証明.} 補題 1 により空の正規化類はないので、対応は一対一である。結合は任意の代表元によって定義されるため、従来のものと同じである。

\noindent\textbf{補題 3.} \emph{正規化因子系の群は乗法群 \(\mathsf P^*\) と同型であり、正規化類の群は \(\mathsf P^*/N\mathfrak Z^*\) と同型である。}

\noindent\emph{証明.} 補題 1 により、正規化因子系はすべて \([1,1,\ldots,a,a,\ldots]\) という形である。実際、
\[
u^\nu\cdot u^\mu=u^{\nu+\mu}\quad(\nu+\mu<n),
\qquad
u^\nu\cdot u^\mu=au^\varrho
\quad(\nu+\mu=n+\varrho, 0\leq\varrho\leq n-2).
\]
対応
\[
[1,\ldots,a,a,\ldots]\mapsto a,
\qquad
[1,1,\ldots,b,b,\ldots]\mapsto b,
\]
\[
[1,\ldots,ab,\ldots]\mapsto ab
\]
は \(\mathsf P^*\) への群同型である。補題 1 により、\(a\) は \(\mathsf P^*\) のすべての元を動くからである。二つの同伴する正規化因子系は、ある変換 \(v=uc\)（\(c\in\mathfrak Z\)）によって互いに移る。これから
\[
v=uc,\qquad v^2=u^2\cdot c^S\cdot c,
\]
\[
v^r=u^r c^{S^{r-1}}c^{S^{r-2}}\ldots c
v^n=u^nN(c)
\]
を得る。したがって \(u^n=a\) から \(v^n=aN(c)\) が従う。ゆえに、一つの正規化類は、\(c\) が \(\mathfrak Z\) の零でないすべての元を動くときの因子系 \([1,\ldots,aN(c),aN(c),\ldots]\) からなる。これで補題の第二の部分も証明された。

補題 2 と 3、およびそこでの \(\{\mathfrak K\}\) と類 \(\{a_{S,T}\}\) との同型を考慮すれば、上の定理 1 が従う。

\noindent\textbf{定理 2.} \emph{最小における主種の群は、\(N(c)=1\) となる \(\mathfrak Z\) の元 \(c\) の群と同型である。したがって \(N(c)=1\) ならば、\(c=b^S\cdot b^{-1}=b^{S-1}\)（記号的な \(S-1\) 乗）となる \(\mathfrak Z\) の元 \(b\ne0\) が存在する。}

\noindent\emph{証明.} 補題 3 の証明で、\(N(c)=1\) であることは、変換 \(v=uc\) が主種の自己同型を表すことと同値であるとわかった。異なる \(c\) は異なる自己同型に対応し、逆もまた成り立ち、積には積が対応する。

\noindent\textbf{注意.} 定理 2 は主種という名称を正当化する。意味に即して、商群
\[
\mathfrak Z^*/\mathscr H
\quad(\mathscr H=\text{主種})
\cong\mathfrak Z^*/\mathfrak Z^{*S-1}
\]
を最小における種の群と呼ぶことができる。これは \(N(\mathfrak Z^*)\) と同型であり、したがって定理 1 と合わせて
\[
\begin{aligned}
\text{最小における種群}
&=\mathfrak Z^*/\mathscr H
\cong\mathfrak Z^*/\mathfrak Z^{*S-1}
\cong N(\mathfrak Z^*)\\
&=\mathfrak Z\text{ に属する }\{\mathfrak K\}\text{ の群}
=\mathscr K_{\mathfrak Z}
\cong\mathsf P^*/N(\mathfrak Z^*).
\end{aligned}
\]

\subsection*{\S 31. 巡回特殊場合の応用}\label{anwendungen-des-zyklischen-spezialfalles}

巡回交叉表現は、\S 20 で証明した二つの定理、すなわち有限体はすべて可換であること、および実数体の真の非可換体拡大は四元数体以外に存在しないことに、新しい証明を与える。いわば、最小における類体論の証明と呼べるものである。

\noindent\textbf{1. 有限個の元からなる各体は可換である（\S 20、定理 10）。}

証明は、可換有限体（Galois 体）についての次の二つの周知の事実に基づく。

\noindent\textbf{1.} \emph{Galois 体は、その任意の部分体上巡回的である。}

\noindent\textbf{2.} \emph{\(\mathfrak Z\) が Galois 体、したがって \(\mathsf P\) 上巡回的ならば、\(\mathsf P\) の各元はある \(\mathfrak Z\) の元のノルムである。\(\mathsf P^*=N(\mathfrak Z^*)\)。}

いま \(\mathfrak K\) を有限とし、\(\mathsf P\) をその中心、\(\mathfrak Z\) を巡回分解体とする。すると \(\mathfrak Z\) に属する \(\{\mathfrak K\}\) の群 \(\mathscr K_{\mathfrak Z}\) は（\S 30 の末尾）、1 により \(\mathsf P^*/N\mathfrak Z^*\) と同型であり、したがって 2 により単位 \(\{\mathsf P\}\) に等しい。これは、中心 \(\mathsf P\) をもつ \(\mathsf P\) と異なる体が存在しないこと、すなわち \(\mathfrak K=\mathsf P\) を意味する。

\noindent\textbf{2. 実数体 \(\mathsf P\) 上の真の非可換体拡大は四元数体だけである（\S 20、定理 9）。}

\noindent\emph{証明.} \(\mathsf P\) は中心でなければならない。なぜなら \(\mathsf P\) の可換拡大は複素数体 \(\mathfrak Z\) しかなく、\(\mathfrak Z\) は代数的閉体なので、それを中心とする非可換体拡大は存在しないからである。同時に、分解体として考えられるのも \(\mathfrak Z\) だけである。そして \(\mathfrak Z\) は \(\mathsf P\) 上巡回的である。さらに各正数はノルムなので、群 \(\mathsf P^*/N\mathfrak Z^*\) の次数は 2 である。したがって \(\{\mathsf P\}\) のほかには類 \(\{\mathfrak K\}\) が一つしかなく、それが四元数体である。さらに同伴する類の正規化因子系として \(-1\) を取れば、同時に正規形も得られる。実際、\(\mathfrak K=\mathfrak Z+\mathfrak Zu\)、\(u^2=-1\) となる。したがって \(i^2=-1\)、\(u^2=-1\) として \(\mathfrak K=\mathsf P+\mathsf Pi+(\mathsf P+\mathsf Pi)u\) である。

\noindent\textbf{1 への注意.} 周知の事実 2 の証明も、\S 30 の末尾に基づいて次のように行える。\(\mathscr H=\mathfrak Z^{*S-1}\cong\mathfrak Z^*/\mathsf P^*\) となる。なぜなら、\(\mathsf P^*\) の元だけが、かつそのすべてが主種の恒等自己同型を生成し、あるいは \(S-1\) を作用させると単位元へ移るからである。したがって有限な \(\mathfrak Z\) では、\(\mathscr H\) の元の個数は \(\mathfrak Z^*\) における \(\mathsf P^*\) の指数に等しい。それゆえ \(\mathfrak Z^*/\mathscr H\) の元の個数は \(\mathsf P^*\) の元の個数に等しく、したがって \(N(\mathfrak Z^*)=\mathfrak Z^*/\mathscr H\) についても同じである。ゆえに \(\mathsf P^*=N(\mathfrak Z^*)\) である。（種の数＝両義類の数を参照。ただし、この数論的な言い方には常に「最小における」を付け加えなければならない。）

\noindent\textbf{3.} \emph{\(\mathsf P\) が \(p\)-進基礎体であり、\(\mathfrak Z\) が \(\mathsf P\) 上次数 \(n\) の巡回体ならば、\(\mathfrak Z\) に属する \(\{\mathfrak K\}\) の群 \(\mathscr K_{\mathfrak Z}\) は \(\mathfrak Z\) のガロア群と同型である。}

これは \S 30 の末尾により、
\[
\mathsf P^*/N\mathfrak Z^*\cong\mathfrak G
\]
と述べる「小域の類体論」から直ちに従う。

\noindent\textbf{3 への注意.} \(\mathscr K_{\mathfrak Z}\) の元の群位数、すなわち \(\{\mathfrak K\}\) の指数がその標数に等しくなることが従う。実際、\(\mathscr K_{\mathfrak Z}\) の生成元の標数は \(n\) であり、したがって指数も、\(n\) の約数かつ標数の倍数として \(n\) になる。ゆえに群 \(\mathscr K_{\mathfrak Z}\) が含むのは、指数が \(n\) の約数で、\(\mathfrak Z\) が分解体となる類だけである。さらに H. Hasse は Math. Ann. 104 において、\(\mathscr K_{\mathfrak Z}\) が、指数が \(n\) の約数であるすべての類 \(\{\mathfrak K\}\) からなることを示した。したがって、\(p\)-進体 \(\mathsf P\) 上の次数 \(n\) の各巡回分解体 \(\mathfrak Z\) は、体の類 \(\{\mathfrak K\}\) の同じ群 \(\mathscr K_{\mathfrak Z}\) を与える。

次の二事実、

\noindent\textbf{1.} \(\mathsf P\) 上の、指数が \(n\) の約数であるすべての \(\{\mathfrak K\}\) の群は次数 \(n\) の巡回群である。

\noindent\textbf{2.} \(\mathsf P\) 上の次数 \(n\) の各巡回分解体は、\(\{\mathfrak K\}\) の全群を生成する。

は、小域の類体論の逆定理の本質的内容とみなすことができる。

%% END NOETHER SOURCE
%% BEGIN NOETHER SOURCE "ja/t6.tex"
\clearpage
\subsection*{代数関数に関する Noether の基本定理のための必要十分な重複度条件：H. Kapferer}\label{notwendige-und-hinreichende-multiplizitaetsbedingungen-zum-noetherschen-fundamentalsatz-der-algebraischen-funktionen}

\begin{center}
（E. Noether との共同補遺を付す）

Math. Ann. 97, (1927), S. 559--567
\end{center}

以下では、Noether の基本定理を必要十分な重複度条件という形に強める\srcfn{1)}{詳しい叙述については、A. Loewy（Freiburg i. Br.）の企画による Heidelberger Akademie-Berichte 1927 の特別号「Beiträge zur Algebra」に掲載された、同名の私の論文を参照されたい。この特別号には複数の著者の論文が収められている。}。ここおよび以下ではイデアル論の概念や定理を用いないが、この定理のイデアル論的理解\srcfn{2)}{例えば B. L. van der Waerden, ``Zur Nullstellentheorie der Polynomideale''. Math. Annalen 96 (1926), S. 203 を参照。}に従って、Noether の基本定理に含まれる一連の事実を、次の二つの独立した命題\srcfn{3)}{引用した論文の \S 1 では、二つの命題を個別かつ直接に、すなわちイデアル論的概念を用いずに証明した。}にまとめる。

\textbf{定理 Ia.} \(x,y\) の互いに素な二多項式 \(\varphi,\psi\)（係数は任意の複素数）の各共通零点 \(x=\alpha, y=\beta\) に対して、最小の正整数 \(\varrho\) が存在し、加群 \((\varphi,\psi,p^\varrho)\) が表すすべての多項式の集合は、\(\sigma>\varrho\) を任意の整数とした加群 \((\varphi,\psi,p^\sigma)\) が表す集合とちょうど同じである。ここで \(p=(x-\alpha,y-\beta)\) である\srcfn{4)}{文字 \(p\) と \(q\) を選んだのは、それらが表す内容が、近年の文献で通常用いられる素イデアル \(p\) および準素イデアル \(q\) の概念と一致することを想起させるためである。}。

\textbf{定理 Ib.} この概念的に一意に定まる多項式集合を \(q\) と表し、また \(s\) 個の異なる零点の第 \(i\) 番目 \(x=\alpha_i, y=\beta_i\) に関する場合には \(q_i\) と表す。このとき、\(x,y\) の所与の多項式 \(K\) が可除性 \(K\equiv0(\varphi,\psi)\) をもつための必要十分条件は、次の \(s\) 個の合同式が同時に満たされることである：
\[
K\equiv0(q_1),\quad K\equiv0(q_2),\quad\ldots,\quad K\equiv0(q_s).
\]

歴史的には、この主題に関する Max Noether の基礎論文の序論\srcfn{5)}{Max Noether, ``Über einen Satz aus der Theorie der algebraischen Funktionen'', Math. Annalen 6 (1872).}を参照すれば、基本定理は次の基本問題を解こうとする努力から生まれた。\emph{\(x,y\) の各所与の多項式 \(K\) について、それが \((\varphi,\psi)\) で割り切れるか否かを判定できなければならない。}すなわち、さらに二つの \(x,y\) の多項式 \(\lambda,\mu\) が存在して、等式
\[
K=\lambda\cdot\varphi+\mu\cdot\psi
\]
が \(x,y\) の恒等式となるか否かを判定することである。

この問題は Noether の定理によって解かれてはいない。むしろ同定理は、提出された問題をより簡単な問題へ移す、基本的ではあるが一つの\emph{還元}\srcfn{6)}{この還元は、初等整数論から誰もが知る次の還元に比べられる。自然数 \(A\) が自然数 \(B\) で割り切れるという性質は、\(A\) が \(p_1^{a_1}\)、\(p_2^{a_2}\)、\(\ldots\)、\(p_s^{a_s}\) のそれぞれで割り切れることと同値である。ここで \(p_1,p_2,\ldots,p_s\) は \(B=p_1^{a_1}p_2^{a_2}\cdots p_s^{a_s}\) の相異なる素因数である。素数冪に対応するものとして、まさに「準素イデアル」――本文では \(q\) と表す――の概念を理解すべきである。［E. Noether, ``Idealtheorie in Ringbereichen'', Math. Annalen 83 (1921), Einleitung 参照。］}とみなすべきである。そのより簡単な問題とは、各合同式 \(K\equiv0(q)\) が成り立つための必要十分条件を問うことである。では、いつ \(K\equiv0(q)\) なのか。この後者の問いが答えられて初めて、M. Noether の当初の問題も解決される。これが本論文の目的であり、実際に達成される。すなわち、\(K\equiv0(q)\) のために必要十分な、有限個の、実際にも容易に検査できる\emph{重複度条件}を与える。\(K\equiv0(q)\) の十分条件ではあるが必要条件ではない条件は既知で、文献ではしばしば次の形で適用・引用される。

合同式 \(K\equiv0(\varphi,\psi)\) が成り立つためには、\(\varphi=0\)、\(\psi=0\) の各交点が、多項式 \(K\) において「十分」高い重複度で現れることが\emph{十分である}。

この定理は定理 I に特殊場合として含まれる。実際、\(K\) における \(\varrho\) 重点は \(K\equiv0(p^\varrho)\) と同値であり、なおさら \(K\equiv0(\varphi,\psi,p^\varrho)\)、すなわち \(K\equiv0(q)\) である。しかし逆は成り立たない。

予告した問題の解法は、簡潔には次のように特徴づけられる。各 \(q\) には \(t\) 個の\emph{正}整数の系 \((g_1),(g_2),\ldots,(g_t)\) を対応させることができ、加群 \((K,q)\) には同数の\emph{正または負}の整数 \((K_1),(K_2),\ldots,(K_t)\) を対応させることができる。そして次の\emph{主定理}が成り立つ。

\textbf{定理 II.} \emph{\(K\equiv0(q)\) のための必要十分条件は、\(t\) 個の条件}
\[
(K_1)\geq(g_1),\quad (K_2)\geq(g_2),\quad\ldots,\quad (K_t)\geq(g_t)
\]
\emph{が同時に満たされることである。このとき}
\[
(g_1)+(g_2)+\ldots+(g_t)=\mu,
\]
\emph{ここで \(\mu\) は \(q\) に対応し、終結式によって定義される交点重複度である。}

証明のため、一般性を失うことなく、\((\varphi,\psi)\) の \(q\) に対応する零点を \(x=0, y=0\) と仮定する。さらに一般に、\(A(x,y)\) が \(x,y\) の任意の多項式\srcfn{7)}{\(A(x,y)\) として有理分数関数 \(C(x,y):D(x,y)\) も許せば、\((A)=(C)-(D)\) となる。この後者の場合、\((A)\) は負の値を取ることもある。}であるとき、記号 \((A)\) は \(A(0,y)\) に因子として \(y\) が何回現れるかを示す数を表すものとする。ただし \(A(0,y)\) が恒等的に消える場合は \((A)=\infty\) と置く。\(q\) に対応する数 \((g_1),(g_2),\ldots,(g_t)\) は、手続の方法論的基礎をなす次の方程式系によって定義される：
\begin{equation}
\tag{1}
\begin{aligned}
f_1\cdot g_1(0,y):y^{(g_1)}-g_1\cdot f_1(0,y):y^{(g_1)}&=x\cdot h_1,\\
f_2\cdot g_2(0,y):y^{(g_2)}-g_2\cdot f_2(0,y):y^{(g_2)}&=x\cdot h_2,\\
&\vdots \qquad\qquad\vdots\qquad\qquad\vdots\\
f_r\cdot g_r(0,y):y^{(g_r)}-g_r\cdot f_r(0,y):y^{(g_r)}&=x\cdot h_r.
\end{aligned}
\end{equation}

これは所与の多項式対 \(\varphi,\psi\) だけから構成される。すなわち多項式 \(f_1,g_1\) は、順序を除けば \(\varphi,\psi\) と同一と定義する。その順序は \((f_1)\geq(g_1)\) となるようにする。この規約を置くとき、そしてそのときに限り、\(h_1\) は整式となる。

多項式 \(f_2,g_2\) は、付帯条件 \((f_2)\geq(g_2)\) のもとで \(h_1,g_1\) と同一と定義する。一般に多項式対 \(f_r,g_r\) は、付帯条件 \((f_r)\geq(g_r)\) のもとで \(h_{r-1},g_{r-1}\) と同一と定義する。

方程式系 (1) は、説明した仕方で無限に続けることができる。しかし主定理のために作る必要があるのは、この系の最初の \(t\) 本の方程式だけであり、\(t\) は次の事実によって定まる。

\textbf{補題 I.} 自然数 \(t\) が存在し［しかも \(t\leqq\mu\)。ここで \(\mu\) は \(\varphi=0,\ \psi=0\) における点 \(x=0, y=0\) の交点重複度である］、(1) により定義される最初の \(t\) 個の数、すなわち \((g_1),(g_2),\ldots,(g_t)\) はすべて\emph{零より大きい}一方、\((g_{t+1})\) およびその後のすべては\emph{零に等しい}。

実際、終結式の積定理を考慮すると、点 \(x=0, y=0\) における \(f_r=0, g_r=0\) の交点重複度は
\[
\mu-(g_1)-(g_2)-\ldots-(g_r)
\]
に等しい。このことから、\emph{遅くとも \(\mu\) 段階で停止する}ことが従う。同時に
\begin{equation}
\tag{2}
\mu=(g_1)+(g_2)+\ldots+(g_t)
\end{equation}
を得る。この公式は、該当する零点が既知ならば、交点重複度を定めるための一般的で有理的、かつ実際にも容易に実行できる方法を与える点で注目に値する。

こうして数 \((g_i)\) がわかれば、数 \((K_i)\) は (1) に対応する次の方程式系によって定まる：
\begin{equation}
\tag{3}
\begin{aligned}
K_1\cdot g_1(0,y):y^{(g_1)} - g_1\cdot K_1(0,y):y^{(g_1)} &= x\cdot K_1,\\
K_2\cdot g_2(0,y):y^{(g_2)} - g_2\cdot K_2(0,y):y^{(g_2)} &= x\cdot K_2,\\
&\vdots \qquad\qquad\vdots\qquad\qquad\vdots\\
K_r\cdot g_r(0,y):y^{(g_r)}-g_r\cdot K_r(0,y):y^{(g_r)}&=x\cdot K_{r+1}.
\end{aligned}
\end{equation}

\(K_1\) は \(K\) と同一と定義する。\(K_2\) が整式となるのは、明らかに \((K_1)\geq(g_1)\) のとき、かつそのときに限る。後者が満たされるならば、さらに \(K_3\) が整式となるのは \((K_2)\geq(g_2)\) のとき、かつそのときに限る。一般に、\(K_r\) がすでに整式であるならば、\(K_{r+1}\) が整式となるのは \((K_r)\geq(g_r)\) のとき、かつそのときに限る。

主定理（定理 II）の証明は、さらに次の補題に基づく。

\textbf{補題 II.} \((f_i,g_i,p^e)=q^{(i)}\) と置けば、無限整数列 \(1,2,\ldots\) の各 \(i\) に対して次の定理が成り立つ。

合同式
\[
K_i\equiv0(\mathfrak{q}^{(i)})
\]
のための\emph{必要十分条件は、次の二条件が同時に満たされることである：}
\[
(K_i)\geqq(g_i)\quad \text{かつ}\quad K_{i+1}\equiv0(\mathfrak{q}^{(i+1)}).
\]

補題の証明\srcfn{8)}{「補遺」の注意 I における証明も参照。}は、方程式系 (1) と (3) を組み合わせることによって得られる。補題 I により \(g_{t+1}\) は原点でもはや消えず、したがって加群 \(q^{(t+1)}\) は\emph{すべて}の多項式からなり（単位イデアルであり）、ゆえに \(K_{t+1}\equiv0(q^{(t+1)})\) は自明となる。これを考慮すれば、補題 II を \(t\) 回適用することで主定理が証明される。

\begin{center}
\rule{4cm}{0.4pt}
\end{center}

証明は前掲論文に譲り、主定理に二つの補足を付け加える。

\textbf{補足 I.} 主定理の数 \(t\) は、
\[
x^v\equiv0(\mathfrak{q})
\]
となる最小の正数 \(\nu\) と一致する\srcfn{9)}{ここでは、Frl. Grete Hermann［``Die Frage der endlich vielen Schritte in der Theorie der Polynomideale''（K. Hentzelt の遺稿定理を利用）、Math. Annalen 95 (1926)］による評価とは対照的に、\emph{正確な}指数が得られる。そこでは、問題設定ははるかに一般的であるものの、上界だけが与えられている。}。

\textbf{補足 II.} 方程式系 (1) と (3)、およびそれらで \(x\) と \(y\) を交換して得る系は、適切に組み合わせれば、\(q\) に「属する」指数 \(\varrho\)、すなわち定理 Ia の最小指数 \(\varrho\) を数値的に正確に定める有理的方法を与える。

E. Bertini\srcfn{10)}{E. Bertini, Math. Annalen 34 (1889).} は \(\varrho\) の一般公式を与えた。しかし、この公式が必ずしも \(\varrho\) の最小値を与えないことは、引用した論文で例 \(\varphi=x^3+y^4;\ \psi=x^2+y^3\) により示される。最後に、Bertini の公式そのものも、われわれの主定理によって新しい仕方で基礎づけられることを述べておく。

\begin{center}
\rule{4cm}{0.4pt}
\end{center}

主定理の興味深い\emph{特殊化}は、該当する \(q=(\varphi,\psi,p^e)\) の零点が、二曲線 \(\varphi=0,\ \psi=0\) の少なくとも一方において\emph{単純点}であるという制限的仮定から得られる。例えばそれが \(\psi\) の単純点ならば、主定理の数 \(t\) は\emph{ちょうど} \(\mu\) に等しく、主定理の \(t\) 個の不等式を\emph{一つ}の重複度条件で置き換えられることが示せる。これは次の一般定理を導く。

\textbf{定理 III.} \emph{\(\varphi=0,\ \psi=0\) の\emph{すべての}交点が \(\psi\) において\emph{単純点}であるならば、\(K\) が \(\psi\) と互いに素であると仮定して、合同式 \(K\equiv0(\varphi,\psi)\) が成り立つための\emph{必要十分条件}は、\(K\) がこれらすべての点で消え、しかも曲線対 \(K=0,\ \psi=0\) における各交点重複度が、曲線対 \(\varphi=0,\ \psi=0\) におけるものと少なくとも同じ大きさであることである。}

この定理 III は、さらに特別な理由から注目に値する。すなわち、そこで要求される重複度条件が満たされるか否かは、(1) に基づく公式 (2) とは独立に、有限回の\emph{微分}によって判定できる。このことは、私が 1923 年に立てた次の「定理」\srcfn{11)}{H. Kapferer, ``Über die Multiplizität der Schnittpunkte von zwei algebraischen Kurven''. Jahresber. der D. M.-V. 32 (1923).}を引用すれば直ちに明らかである。

\(A,B\) を \(x,y\) の任意の二つの互いに素な多項式とし、\(P\) を共通点であって\emph{\(B\) の単純点}であるものとする。この点の交点としての重複度が数 \(\mu\) に等しいための必要十分条件は、同じ点の曲線対 \(A'=0, B=0\) における重複度が\emph{ちょうど} \(\mu-1\) であることである。ここで
\[
A'=\frac{\partial A}{\partial x}\cdot\frac{\partial B}{\partial y}
-\frac{\partial A}{\partial y}\cdot\frac{\partial B}{\partial x}.
\]

\begin{center}
\rule{4cm}{0.4pt}
\end{center}

最後に、主定理（定理 II）はわずかな変更によって三変数 \(x,y,z\) の斉次多項式にも移せることを述べておく。この移行は、上の特殊場合には特に簡潔に行われる。すなわち次が示せる。

\textbf{定理 III への補足.} \(\varphi,\psi,K\) を\emph{三}変数の\emph{斉次}多項式と解するときも、定理 III は文言まで含めて全く変更なく成り立つ\srcfn{12)}{三変数の斉次合同式 \(K(x,y,z)\equiv0(\varphi(x,y,z),\psi(x,y,z))\) が成り立つ条件を問う問題は、強く特殊化された仮定のもとではあるが、A. Ostrowski, ``Über die Maxwellsche Erzeugung der Kugelfunktionen'', Jahresb. der D. M.-V. 33 (1925) の研究の中心問題でもある。この論文では、\(\varphi\) がすべて一次因子に分解すること、\(\varphi\) と \(K\) が同じ次数であること、\(\psi=x^2+y^2+z^2\) であることを仮定する。したがって、ここでは \(\psi\) が全く特異点をもたないので、Ostrowski の論文の主定理はわれわれの定理 III の特殊場合とみなせる。}。

\subsection*{E. Noether との共同補遺\srcfn{13)}{イデアル論的解釈は E. Noether により、基礎となる (5) と (6) の証明は H. Kapferer による。}}\label{zusatz-gemeinsam-mit-e.-noether.}

方程式系 (1) は同時に、\emph{\(q\) を法とする完全剰余系を決定する}有理的方法を与える。完全剰余系とは、係数領域 \(\mathsf P\) に関して、\(q\) を法として\emph{線形独立}な\emph{有限個}の剰余類の完全代表系をいう。これによって主定理の証明の各段階が非常に明瞭になる。同時に、副結果として、\emph{終結式重複度 \(\mu\) がイデアル \(q\) の長さと一致する}ことが得られる。

すなわち、それは剰余類環の階数、または \(\mathsf P\) に関して線形独立な剰余類の個数に等しい\srcfn{14)}{一般性を失うことなく、体 \(\mathsf P\) は代数的閉体と仮定する。この範囲では、微分条件を除く先行するすべての定理が実際に成り立つ。\(q\) の長さは、\(p\) から \(q\) へ至るイデアルの組成列の長さとしても定義される。E. Noether, Jahresber. d. d. Math.-Ver. 34 (1925), S. 101（斜め頁付け）を参照。より詳しくは H. Grell, Math. Annalen 97 (1927), S. 529 を参照。二つの重複度数が一致することについて文献にある証明（\(n\) 変数の場合）は非常に複雑であり、最も簡単なのは未刊の Hentzelt の遺稿証明である。}。

完全剰余系は、(1) と補題 II の記号を用いれば次の定理によって特徴づけられる。

\textbf{定理 IV.} \emph{\(q\) を法とする完全剰余系は、次の多項式によって与えられる：}
\[
h_t,\ h_ty,\ldots,h_ty^{(g_t)-1};\ \ldots;\
h_i,\ h_iy,\ldots,h_iy^{(g_i)-1};\ \ldots;\
h_1,\ h_1y,\ldots,h_1y^{(g_1)-1}.
\]
\emph{ここで各 \(h_i,h_iy,\ldots,h_iy^{(g_i)-1}\) は、\(q^{(i)}\) を法とした \(q^{(i+1)}\) の完全剰余系を与える。}

定義から直ちに \(q^{(i+1)}=(q^{(i)},h_i)\)、したがって \(q^{(i)}\equiv0(q^{(i+1)})\) が従う。したがって定理 IV は、次を示せば証明される：
\begin{equation}
\tag{4}
h_i x\equiv0(q^{(i)});\qquad
h_i y^{(g_i)}\equiv0(q^{(i)})
\end{equation}
および
\begin{equation}
\tag{5}
 h_iF(y)\equiv0(q^{(i)})
 \text{ から }F(y)\equiv0(y^{(g_i)})\text{ が従う。}
\end{equation}

実際、関係式 (4) は、\(h_i,h_iy,\ldots,h_iy^{(g_i)-1}\) の線形結合が、\(q^{(i)}\) を法とする \(q^{(i+1)}\) の剰余類を尽くすことを述べ、(5) は線形独立性を表す。\(q^{(t+1)}\) は単位イデアルに等しくなる――これはいまや \(q\) の有限長から従う――ので、これから定理 IV の証明が得られ、同時に \(q\) の長さが \((g_1)+\cdots+(g_t)\) に等しく、したがって (2) により終結式重複度と一致することが得られる。

関係式 (4) の第一式は (1) から読み取れる。第二式は次のように得られる。

\(h_ig_i\equiv0(q^{(i)})\) である。したがって (4) の第一関係式により、\(h_ig_i(0,y)\equiv0(q^{(i)})\)、または
\[
h_iy^{(g_i)}\cdot(g_i(0,y):y^{(g_i)})\equiv0(q^{(i)})
\]
でもある。\(g_i(0,y):y^{(g_i)}\not\equiv0(y)\) であるから、ここから \(h_iy^{(g_i)}\equiv0(q^{(i)})\) が従う。

(5) の証明に先立って、(1) から直ちに読み取れる次の\emph{注意}を述べる。\(f_i\) と \(g_i\) が共通因子としてもちうるのは、\(y\) で割れない \(y\) の多項式 \(d_i(y)\) だけである：
\[
f_i=d_i(y)\bar f_i,\qquad g_i=d_i(y)\bar g_i,
\]
そして \(\bar f_i\) と \(\bar g_i\) は互いに素である。\(d_i(y)\not\equiv0(p)\) であるから、\(q^{(i)}\) は \((\bar f_i,\bar g_i,p^e)\) にも等しい。さらに、\(q\) の約数として、これは各 \(\sigma\geq\varrho\) について \((\bar f_i,\bar g_i,p^\sigma)\) に等しい。したがって \(q^{(i)}\) は、\((\bar f_i,\bar g_i)\) の原点に属する準素成分である。

\(\gamma_1,\delta_1;\ldots;\gamma_r,\delta_r\) を、\((\bar f_i,\bar g_i)\) の \(x=0, y=0\) と異なる零点とし、\(\tau_1,\ldots,\tau_r\) を \((\bar f_i,\bar g_i)\) の対応する準素成分の指数とする。さらに積
\[
P=((x-\gamma_1)+u(y-\delta_1))^{\tau_1}\cdots
((x-\gamma_r)+u(y-\delta_r))^{\tau_r}
\]
において、不定元 \(u\) を \(\mathsf P\) のある元へ特殊化し、\(P\not\equiv0(p)\) が保たれるようにする。\(G(x,y)\) を、特殊化した \(P\) と \(d_i(y)\) との積に等しいと置けば、\(G(x,y)\not\equiv0(p)\) かつ \(q^{(i)}\cdot G(x,y)\equiv0(f_i,g_i)\) となる。

仮定 \(h_iF(y)\equiv0(q^{(i)})\) から、\(h_iF(y)G(x,y)\equiv0(f_i,g_i)\) が従う。これを (1) と合わせて恒等式の形に書けば、
\[
h_iF(y)G=A f_i+B g_i;\qquad
h_ix=(g_i(0,y):y^{(g_i)})\cdot f_i
-(f_i(0,y):y^{(g_i)})g_i.
\]
これから
\[
\bar f_i(xA-F(y)H)\equiv0(\bar g_i);\qquad
H=G(x,y)\cdot g_i(0,y):y^{(g_i)}\not\equiv0(\mathfrak{p})
\]
を得る。

\(\bar f_i\) と \(\bar g_i\) が互いに素であることから、さらに
\[
xA-F(y)H\equiv0(\bar g_i)
\]
が従う。したがって
\[
F(y)H\equiv0(\bar g_i,x)
\equiv0(\bar g_i,x,p^\sigma)=(y^{(g_i)},x),
\]
ゆえに \(F(y)\equiv0(y^{(g_i)})\) である。これで (5)、したがって定理 IV が証明された。

\textbf{注意 I.} 関係式 (4)、(5) および補題 II は、イデアル商の間で互いに相反する関係式\srcfn{15)}{商 \(\mathfrak a:\mathfrak b\) とは、\(\mathfrak b c\equiv0(\mathfrak a)\) となるような多項式 \(c(x,y)\) 全体をいうことは周知である。}として書ける。ここで第二関係式の証明は、同時に補題 II の証明を与える：
\begin{equation}
\tag{6}
\mathfrak{q}^{(i)}:\mathfrak{q}^{(i+1)}=(\mathfrak{q}^{(i)},x)=(y^{(g_i)},x);\qquad
\mathfrak{q}^{(i)}:(\mathfrak{q}^{(i)},x)=\mathfrak{q}^{(i+1)}.\srcfnmark{16)}
\end{equation}
\srcfntext{16)}{互いに相反するこの二関係式のうち一方は、既約イデアルの一般理論により他方を含意する。Macaulay, Modular Systems, Nr. 72 und 73. Cambridge Tracts 19 (1916) を参照。}

第一関係式は (4) と (5) をまとめたものである。第二関係式については、\(x\cdot q^{(i+1)}\equiv0(q^{(i)})\) は (4) で証明済みであり、その逆は (5) と同様に得られることに注意する。\(xM\equiv0(q^{(i)})\) から、上と同じように
\[
xMG\equiv0(f_i,g_i);\quad\text{したがって}\quad
xMG\cdot g_i(0,y):y^{(g_i)}\equiv0(xh_i,g_i)
\]
が従う。

定義により \(g_i\not\equiv0(x)\)――\((f_i)\geq(g_i)\)――なので、ここから
\[
M\cdot(G\cdot g_i(0,y):y^{(g_i)})
\equiv0(h_i,g_i)\equiv0(q^{(i+1)}),
\quad\text{したがって}\quad M\equiv0(\mathfrak{q}^{(i+1)})
\]
を得る。

補題 II の条件が必要であることは、いまや次のように従う。\(K_i\equiv0(q_i)\) から \(K_i(0,y)\equiv0(y^{(g_i)},x)\) が従い、したがって \((K_i)\geq(g_i)\) である。ゆえに (3) により
\[
xK_{i+1}\equiv0(q^{(i)});\quad
\text{したがって――(6) の第二関係式により――}\quad
K_{i+1}\equiv0(\mathfrak{q}^{(i+1)}).
\]

条件が十分であることは直ちに読み取れる。補題 II を有限回適用すれば主定理が得られる。同様に、(6) の第二関係式を反復すれば補足 I、すなわち \(q:x^t=q^{(t+1)}\)、したがって単位イデアルに等しいことが得られる。しかも \(q:x^{t-1}=q^{(t)}\) であるため、\(t\) はこの性質をもつ最低の冪である。

\textbf{注意 II.} \(q\) で割り切れない多項式 \(K\) について、さらに次がいえる。\(K\equiv0(q^{(i+1)})\)、\(K\not\equiv0(q^{(i)})\) ならば、\((K,q^{(i)})=(h_iy^{(K)},q^{(i)})\) となる。実際、定理 IV により、\(c(y)\not\equiv0(y)\) として \(K\equiv h_iy^\lambda c(y)\pmod {q^{(i)}}\) であり、したがって
\[
(K,\mathfrak{q}^{(i)})=(h_iy^\lambda,\mathfrak{q}^{(i)}).
\]

このとき \(y^{(g_i)-\lambda}K\equiv0(q^{(i)})\) であり、(5) により、これはその性質をもつ最低の冪である。一方、補題 II により \((g_i)-(K)\) がその最低の冪であるから、\(\lambda=(K)\) となる。

\begin{center}
（1926 年 10 月 24 日受理。）
\end{center}

%% END NOETHER SOURCE
%% BEGIN NOETHER SOURCE "ja/t7.tex"
\clearpage
\clearpage
\begin{center}
{\Large\bfseries 文献目録}
\end{center}
\label{bibliographie}

\begin{enumerate}
\setlength{\itemsep}{0.35em}
\item 三変数双二次形式の形式系の構成について。\\
Sitz. Ber. d. Physikal.-mediz. Sozietät in Erlangen 39 (1907), S. 176--179
\item 三変数双二次形式の形式系の構成について。\\
J. Reine Angew. Math. 134 (1908), S. 23--90 および二つの表
\item \(n\) 変数形式の不変式論について。\\
J. Ber. d. DMV 19 (1910), S. 101--104
\item \(n\) 変数形式の不変式論について。\\
J. Reine Angew. Math. 139 (1911), S. 118--154
\item 有理関数体。\\
J. Ber. d. DMV 22 (1913), S. 316--319
\item 体と有理関数系。\\
Math. Ann. 76 (1915), S. 161--196
\item 有限群の不変式の有限性定理。\\
Math. Ann. 77 (1916), S. 89--92
\item 任意個の基本形式からなる系の不変式の整有理表示について。\\
Math. Ann. 77 (1916), S. 93--102.（第 16 番も参照）
\item 整超越数からなる最も一般的な領域。\\
Math. Ann. 77 (1916), S. 103--128.（第 16 番も参照）
\item 同型写像の関数方程式。\\
Math. Ann. 77 (1916), S. 536--545
\item 指定された群をもつ方程式。\\
Math. Ann. 78 (1918), S. 221--229.（第 16 番も参照）
\item 任意の微分式の不変式。\\
Nachr. v. d. Ges. d. Wiss. zu Göttingen 1918, S. 37--44
\item 不変変分問題。\\
Nachr. v. d. Ges. d. Wiss. zu Göttingen 1918, S. 235--257
\item 一変数代数関数の算術理論と、他の諸理論および代数体論との関係。\\
J. Ber. d. DMV 28 (1919), S. 182--203
\item 二変数形式の整係数不変式系の有限性。\\
Nachr. v. d. Ges. d. Wiss. zu Göttingen 1919, S. 138--156
\item 形式論における級数展開について。\\
Math. Ann. 81 (1920), S. 25--30
\item W. Schmeidler と共著：非可換領域における加群、特に微分式および差分式からなるもの。\\
Math. Z. 8 (1920), S. 1--35
\item 戦没した K. Hentzelt の消去理論に関する論文について。\\
J. Ber. d. DMV 30 (1921), S. 101
\item 環領域におけるイデアル論。\\
Math. Ann. 83 (1921), S. 24--66
\item 絶対既約性の代数的判定基準。\\
Math. Ann. 85 (1922), S. 26--33
\item 形式変分法と微分不変式。\\
Encyklopädie d. math. Wiss. III, 3 (1922), S. 68--71（R. Weitzenböck, Differentialinvarianten 所収）
\item K. Hentzelt \(\dagger\) の論文の編集：多項式イデアルと終結式の理論について。\\
Math. Ann. 88 (1923), S. 53--79
\item 代数不変式と微分不変式。\\
J. Ber. d. DMV 32 (1923), S. 177--184
\item 消去理論と一般イデアル論。\\
Math. Ann. 90 (1923), S. 229--261
\item 消去理論とイデアル論。\\
J. Ber. d. DMV 33 (1924), S. 116--120
\item 代数体におけるイデアル論の抽象的構成。\\
J. Ber. d. DMV 33 (1924), S. 102
\item イデアル論における Hilbert 数。\\
J. Ber. d. DMV 34 (1925), S. 101
\item 群指標とイデアル論。\\
J. Ber. d. DMV 34 (1925), S. 144
\item 標数 \(p\) の有限線形群の不変式の有限性定理。\\
Nachr. v. d. Ges. d. Wiss. zu Göttingen 1926, S. 28--35
\item 代数的な数体および関数体におけるイデアル論の抽象的構成。\\
Math. Ann. 96 (1927), S. 26--61
\item 代数的数体または関数体の整環に対する判別式定理。\\
J. Reine Angew. Math. 157 (1927), S. 82--104
\item R. Brauer と共著：既約表現の最小分解体について。\\
Sitz. Ber. d. Preuß. Akad. d. Wiss. 1927, S. 221--228
\item 算術的観点における超複素量と表現論。\\
Atti Congresso Bologna 2 (1928), S. 71--73
\item[34.] 超複素量と表現論。\\
Math. Z. 30 (1929), S. 641--692
\item[35.] 整関数からなる極大領域について。\\
Rec. Soc. Math. Moscou 36 (1929), S. 65--72
\item[36.] イデアル微分と共役差積。\\
J. Ber. d. DMV 39 (1930), S. 17
\item[37.] 高次分岐をもたない体の正規基底。\\
J. Reine Angew. Math. 167 (1932), S. 147--152
\item[38.] R. Brauer、H. Hasse と共著：代数の理論における主定理の証明。\\
J. Reine Angew. Math. 167 (1932), S. 399--404
\item[39.] 可換代数および数論との関係における超複素系。\\
Verhandl. Intern. Math.-Kongreß Zürich 1 (1932), S. 189--194
\item[40.] 非可換代数。\\
Math. Z. 37 (1933), S. 514--541
\item[41.] 相対ガロア数体に対する主種定理。\\
Math. Ann. 108 (1933), S. 411--419
\item[42.] 分解する接合積とその極大整環。\\
Actualités scientifiques et industrielles 148 (1934)（15 頁）
\item[43.] イデアル微分と共役差積。\\
J. Reine Angew. Math. 188 (1950), S. 1--21
\end{enumerate}

\noindent 代数関数に関する Noether の基本定理のための必要十分な重複度条件：H. Kapferer\\
（E. Noether との共同補遺を付す）。\\
Math. Ann. 97 (1927), S. 559--567

\clearpage
\begin{center}
{\Large\bfseries Emmy Noether の短報および書評一覧}
\end{center}
\label{liste-der-kurzmitteilungen-und-buchbesprechungen-von-emmy-noether}

\begin{center}
\emph{Deutsche Mathematiker-Vereinigung 年報掲載の短報}
\end{center}

\begingroup
\setlength{\parindent}{0pt}
\setlength{\parskip}{0.15em}
任意の微分式の不変式について、27, Teil II, (1918), S. 28\par
不変変分問題、27, Teil II, (1918), S. 47\par
二変数整係数不変式の有限性、28, Teil II, (1919), S. 29\par
整係数多項式と冪級数について、28, Teil II, (1919), S. 29--30\par
加群論の分解定理、29, Teil II, (1920), S. 54\par
単因子と一般イデアル論、30, Teil II, (1921), S. 32\par
Lipschitz の公式に関する指摘、30, Teil II, (1921), S. 48\par
一般イデアル論における Hilbert の定理の類似、32, Teil II, (1923), S. 20--21\par
公理論の問題に対する群論的注意、32, Teil II, (1923), S. 21\par
任意の体における Galois 理論、33, Teil II, (1925), S. 120\par
イデアル論における Hilbert 数、34, Teil II, (1926), S. 101\par
イデアルの微分商と分岐理論、38, Teil II, (1929), S. 81\par
\endgroup

\begin{center}
\emph{Deutsche Mathematiker-Vereinigung 年報掲載の書評}
\end{center}

\begingroup
\setlength{\parindent}{0pt}
\setlength{\parskip}{0.35em}
E. Study『ベクトル解析に基づく線形変換の不変式論入門』第 1 部（Die Wissenschaft, Bd. 71）、Braunschweig 1923, Vieweg。\\
36, Teil II, (1927), S. 71\par
Leopold Kronecker 著作集、第 4 巻、K. Hensel 編、VI + 509 頁、Leipzig und Berlin 1929, B. G. Teubner。\\
40, Teil II, (1931), S. 11--12\par
Leopold Kronecker 著作集、第 5 巻、K. Hensel 編、527 頁、Leipzig und Berlin 1930, B. G. Teubner。\\
41, Teil II, (1932), S. 17\par
Leopold Kronecker 著作集、第 III.2 巻、K. Hensel 編、215 頁、Leipzig und Berlin 1931, B. G. Teubner。\\
42, Teil II, (1933), S. 38--39\par
D. E. Rutherford『Modular Invariants』Cambridge Tracts No. 27、Cambridge 1932、VII + 84 頁。\\
43, Teil II, (1934), S. 94--95\par
\endgroup

\clearpage
\begin{center}
\emph{Emmy Noether の協力のもとに成立した書籍}
\end{center}
\label{buecher-entstanden-unter-mitwirkung-von-emmy-noether}

\begingroup
\setlength{\parindent}{0pt}
\setlength{\parskip}{0.5em}
R. Dedekind『数学著作集 I--III』R. Fricke、Emmy Noether、Øystein Ore 編、Vieweg, Braunschweig 1930--1932\\
（第 I 巻：1930、第 II 巻：1931、第 III 巻：1932）\par
G. Cantor『R. Dedekind との往復書簡』Emmy Noether、J. Cavaillès 編、60 頁、Actualités scient. et ind. 581, Paris 1937\par
\endgroup

\endgroup

%% END NOETHER SOURCE
